% Generated by matins/tools/build_latex.py from data/corpus.json — do not edit.
\documentclass[10pt,twoside,openright]{book}
\usepackage{matins}
\begin{document}
\BodySize
\frontmatter
\thispagestyle{empty}
\vspace*{0.25\textheight}
{\centering
{\Huge\rubr{Lectiones ad Matutinum}}\par\vspace{0.4em}
{\large\itshape ex Breviario Romano anni MCMLIV}\par\vspace{2em}
{\LARGE The Lessons of Matins}\par\vspace{0.4em}
{\large\itshape from the Roman Breviary of 1954, in Latin and English}\par\vspace{3em}
{\Large\rubr{Pars Aestiva}}\par\vspace{0.3em}{\large\itshape Summer}\par\vspace{2em}
{\small }\par}
\clearpage
\thispagestyle{empty}
{\footnotesize\setlength{\parindent}{0pt}\setlength{\parskip}{0.35em}\raggedright
\par\addvspace{0.6em}\textsc{Divinum Officium}\par
The texts, and the program that arranges them by the rubrics, are the work of the Divinum Officium Project (divinumofficium.com), created by the late Laszlo Kiss and continued since by the Project's team of volunteers. Its code and data are freely available at github.com/DivinumOfficium/divinum-officium under the MIT License.\par
\par\addvspace{0.6em}\textsc{Latin}\par
The Latin text derives from the Ratisbon 1888 edition of the Breviarium Romanum, scanned by the University of St Michael's College, Toronto, amended from the 1906 edition, and corrected by Laszlo Kiss and the Project team against later printings: Pustet (Regensburg, 1943), Benziger Brothers (New York, 1945), Burns Oates \& Washbourne (London, 1946), La Presse Catholique Panaméricaine (Montréal, 1943), and Desclée and Mame (1962).\par
The accented text of the Sixto-Clementine Vulgate was supplied by Frère Romain-Marie of the Abbaye Saint-Joseph de Clairval, Flavigny-sur-Ozerain.\par
\par\addvspace{0.6em}\textsc{English}\par
Holy Scripture is given in the Douay-Rheims translation.\par
The other lessons and the responsories are from the translation of the Roman Breviary by John, Marquess of Bute (edition of 1908), made available by the University of St Michael's College, Toronto.\par
\par\addvspace{0.6em}\textsc{This edition}\par
The lessons were gathered by running the Divinum Officium program, unchanged, for every day from 1950 to 2100 under the rubrics of 1954, and tracing each lesson it read to the office it belongs to. Errors of arrangement in this edition are not the Project's.\par
\par\addvspace{0.6em}\textsc{Typeface}\par
Set in EB Garamond, designed by Georg Duffner and Octavio Pardo, under the SIL Open Font License 1.1.\par
}
\mainmatter
\PartTitle{Proprium de Tempore}{Proper of Time}

\addcontentsline{toc}{chapter}{Proprium de Tempore \textit{— Proper of Time}}

\SeasonTitle{Tempus post Pentecosten}{Time after Pentecost}

\addcontentsline{toc}{section}{Tempus post Pentecosten \textit{— Time after Pentecost}}

\Office{Dominica Sanctissimæ Trinitatis}{Trinity Sunday}{Duplex I classis}
\addcontentsline{toc}{subsection}{Dominica Sanctissimæ Trinitatis \textit{— Trinity Sunday}}
\begin{paracol}{2}
\LA
\LessonHead{Lectio I}
\switchcolumn
\EN
\LessonHead{Lesson I}
\switchcolumn*

\LA
\LessonTitle{De Isaía Prophéta}
\switchcolumn
\EN
\LessonTitle{Lesson from the book of Isaias}
\switchcolumn*

\LA
\Cite{Isa 6:1-4}
\switchcolumn
\EN
\Cite{Isa 6:1-4}
\switchcolumn*

\LA
\noindent In anno, quo mórtuus est rex Ozías, vidi Dóminum sedéntem super sólium excélsum et elevátum: et ea, quæ sub ipso erant, replébant templum. Séraphim stabant super illud: sex alæ uni, et sex alæ álteri, duábus velábant fáciem ejus, et duábus velábant pedes ejus, et duábus volábant. Et clamábant alter ad álterum, et dicébant: Sanctus, sanctus, sanctus, Dóminus exercítuum, plena est omnis terra glória ejus. Et commóta sunt superliminária cárdinum a voce clamántis, et domus repléta est fumo.\par
\switchcolumn
\EN
\noindent In the year that king Ozias died, I saw the Lord sitting upon a throne high and elevated: and his train filled the temple. Upon it stood the seraphims: the one had six wings, and the other had six wings: with two they covered his face, and with two they covered his feet, and with two they hew. And they cried one to another, and said: Holy, holy, holy, the Lord God of hosts, all the earth is full of his glory. And the lintels of the doors were moved at the voice of him that cried, and the house was filled with smoke.\par
\switchcolumn*

\LA
\begin{Resp}\Rsign Vidi Dóminum sedéntem super sólium excélsum et elevátum, et plena erat omnis terra majestáte ejus: \Star{} Et ea, quæ sub ipso erant, replébant templum.\end{Resp}
\begin{Resp}\Vsign Séraphim stabant super illud: sex alæ uni, et sex alæ álteri.\end{Resp}
\begin{Resp}\Rsign Et ea, quæ sub ipso erant, replébant templum.\end{Resp}
\switchcolumn
\EN
\begin{Resp}\Rsign I saw the Lord sitting upon a throne, high and lifted up, and the whole earth was full of His glory; \Star{} And His train filled the temple.\end{Resp}
\begin{Resp}\Vsign Above it stood the Seraphim each one had six wings.\end{Resp}
\begin{Resp}\Rsign And His train filled the temple.\end{Resp}
\switchcolumn*

\LA
\LessonHead{Lectio II}
\switchcolumn
\EN
\LessonHead{Lesson II}
\switchcolumn*

\LA
\Cite{Isa 6:5-8}
\switchcolumn
\EN
\Cite{Isa 6:5-8}
\switchcolumn*

\LA
\noindent Et dixi: Væ mihi, quia tácui, quia vir pollútus lábiis ego sum, et in médio pópuli pollúta lábia habéntis ego hábito, et Regem Dóminum exercítuum vidi óculis meis. Et volávit ad me unus de Séraphim, et in manu ejus cálculus, quem fórcipe túlerat de altári. Et tétigit os meum, et dixit: Ecce tétigit hoc lábia tua, et auferétur iníquitas tua, et peccátum tuum mundábitur. Et audívi vocem Dómini dicéntis: Quem mittam? et quis ibit nobis? Et dixi: Ecce ego, mitte me.\par
\switchcolumn
\EN
\noindent And I said: Woe is me, because I have held my peace; because I am a man of unclean lips, and I dwell in the midst of a people that hath unclean lips, and I have seen with my eyes the King the Lord of hosts. And one of the seraphims flew to me, and in his hand was a live coal, which he had taken with the tongs off the altar. And he touched my mouth, and said: Behold this hath touched thy lips, and thy iniquities shall be taken away, and thy sin shall be cleansed. And I heard the voice of the Lord, saying: Whom shall I send? and who shall go for us? And I said: Lo, here am I, send me.\par
\switchcolumn*

\LA
\begin{Resp}\Rsign Benedíctus Dóminus Deus Israël, qui facit mirabília magna solus: \Star{} Et benedíctum nomen majestátis ejus in ætérnum.\end{Resp}
\begin{Resp}\Vsign Replébitur majestáte ejus omnis terra: fiat, fiat.\end{Resp}
\begin{Resp}\Rsign Et benedíctum nomen majestátis ejus in ætérnum.\end{Resp}
\switchcolumn
\EN
\begin{Resp}\Rsign Blessed be the Lord God of hosts, Who only doeth wondrous things. \Star{} And blessed be His glorious Name for ever.\end{Resp}
\begin{Resp}\Vsign And let the whole earth be filled with His glory. Amen. Amen.\end{Resp}
\begin{Resp}\Rsign And blessed be His glorious Name for ever.\end{Resp}
\switchcolumn*

\LA
\LessonHead{Lectio III}
\switchcolumn
\EN
\LessonHead{Lesson III}
\switchcolumn*

\LA
\Cite{Isa 6:9-12}
\switchcolumn
\EN
\Cite{Isa 6:9-12}
\switchcolumn*

\LA
\noindent Et dixit: Vade, et dices pópulo huic: Audíte audiéntes, et nolíte intellégere: et vidéte visiónem, et nolíte cognóscere. Excǽca cor pópuli hujus, et aures ejus ággrava, et óculos ejus claude, ne forte vídeat óculis suis, et áuribus suis áudiat, et corde suo intéllegat, et convertátur, et sanem eum. Et dixi: Usquequo, Dómine? Et dixit: Donec desoléntur civitátes absque habitatóre, et domus sine hómine, et terra relinquétur desérta, et longe fáciet Dóminus hómines, et multiplicábitur quæ derelícta fúerat in médio terræ.\par
\switchcolumn
\EN
\noindent And he said: Go, and thou shalt say to this people: Hearing, hear, and understand not: and see the vision, and know it not. Blind the heart of this people, and make their ears heavy, and shut their eyes: lest they see with their eyes, and hear with their ears, and understand with their heart, and be converted and I heal them. And I said: How long, O Lord? And he said: Until the cities be wasted without inhabitant, and the houses without man, and the land shall be left desolate. And the Lord shall remove men far away, and she shall be multiplied that was left in the midst of the earth.\par
\switchcolumn*

\LA
\begin{Resp}\Rsign Benedícat nos Deus, Deus noster, benedícat nos Deus: \Star{} Et métuant eum omnes fines terræ.\end{Resp}
\begin{Resp}\Vsign Deus misereátur nostri, et benedícat nos Deus.\end{Resp}
\begin{Resp}\Rsign Et métuant eum omnes fines terræ.\end{Resp}
\begin{Resp}\Vsign Glória Patri, et Fílio, \Star{} et Spirítui Sancto.\end{Resp}
\begin{Resp}\Rsign Et métuant eum omnes fines terræ.\end{Resp}
\switchcolumn
\EN
\begin{Resp}\Rsign Let God, even our own God, bless us; let God bless us. \Star{} And let all the ends of the earth fear Him.\end{Resp}
\begin{Resp}\Vsign God be merciful unto us, and bless us.\end{Resp}
\begin{Resp}\Rsign And let all the ends of the earth fear Him.\end{Resp}
\begin{Resp}\Vsign Glory be to the Father, and to the Son, \Star{} and to the Holy Ghost.\end{Resp}
\begin{Resp}\Rsign And let all the ends of the earth fear Him.\end{Resp}
\switchcolumn*

\LA
\LessonHead{Lectio IV}
\switchcolumn
\EN
\LessonHead{Lesson IV}
\switchcolumn*

\LA
\LessonTitle{Ex libro sancti Fulgéntii Epíscopi, de fide ad Petrum}
\switchcolumn
\EN
\LessonTitle{From the Book on the Faith, addressed to Peter by St. Fulgentius, Bishop [of Ruspa]}
\switchcolumn*

\LA
\Source{Inter Opera Augustini, tom. 3}
\switchcolumn
\EN
\Source{Found in the Works of Augustine, tom. 3}
\switchcolumn*

\LA
\noindent Fides, quam sancti Patriárchæ atque Prophétæ ante incarnatiónem Fílii Dei divínitus accepérunt, quam étiam sancti Apóstoli ab ipso Dómino in carne pósito audiérunt, et Spíritus Sancti magistério instrúcti non solum sermóne prædicavérunt, verum étiam ad instructiónem salubérrimam posterórum scriptis suis índitam reliquérunt; unum Deum prǽdicat Trinitátem, id est, Patrem, et Fílium, et Spíritum Sanctum. Sed Trínitas vera non esset, si una eadémque persóna dicerétur Pater et Fílius et Spíritus Sanctus.\par
\switchcolumn
\EN
\noindent The Faith which the holy Patriarchs and Prophets received from God before His Son was made Flesh, the Faith which the holy Apostles heard from the Lord Himself when Present in the Flesh, the Faith which the same Apostles learnt by the teaching of the Holy Ghost not only to preach by word of mouth, but also to leave behind them in their writings for the healthful instruction of all that should come after, that Faith teacheth that the Trinity, that is to say, the Father, the Son, and the Holy Ghost, is but One God. But we could not truly call the Father, the Son, and the Holy Ghost a Trinity, if One and the Selfsame Person were named Father, Son, and Holy Ghost.\par
\switchcolumn*

\LA
\begin{Resp}\Rsign Quis Deus magnus sicut Deus noster? \Star{} Tu es Deus, qui facis mirabília.\end{Resp}
\begin{Resp}\Vsign Notam fecísti in pópulis virtútem tuam: redemísti in brácchio tuo pópulum tuum.\end{Resp}
\begin{Resp}\Rsign Tu es Deus, qui facis mirabília.\end{Resp}
\switchcolumn
\EN
\begin{Resp}\Rsign Who is so great a God as our God? \Star{} Thou art the God that doest wonders.\end{Resp}
\begin{Resp}\Vsign Thou hast declared thy strength among the people Thou hast with thine arm redeemed thy people.\end{Resp}
\begin{Resp}\Rsign Thou art the God that doest wonders.\end{Resp}
\switchcolumn*

\LA
\LessonHead{Lectio V}
\switchcolumn
\EN
\LessonHead{Lesson V}
\switchcolumn*

\LA
\noindent Si enim, sicut est Patris, et Fílii, et Spíritus Sancti una substántia, sic esset una persóna; nihil omníno esset in quo veráciter Trínitas dicerétur. Rursus quidem Trínitas esset vera, sed unus Deus Trínitas ipsa non esset, si quemádmodum Pater, et Fílius, et Spíritus Sanctus personárum sunt ab ínvicem proprietáte distíncti, sic fuíssent naturárum quoque diversitáte discréti. Sed quia in illo uno vero Deo Trinitáte, non solum quod unus Deus est, sed étiam quod Trínitas est, naturáliter verum est; proptérea ipse verus Deus in persónis Trínitas est, et in una natúra unus est.\par
\switchcolumn
\EN
\noindent Nor if as the Being of the Father, the Son, and the Holy Ghost is One Being, so were there but One Person, then were it untrue to say that God is a Trinity. On the other hand, if, as the Persons of the Father, the Son, and the Holy Ghost are distinguished One from Another by that which is proper to Each, so were They diverse by difference of nature, then were it untrue to say that God is One. But since concerning the nature of the One True God, Who is a Trinity, it is the Truth to say that God is One, and the Truth to say that God is a Trinity, therefore the True God is a Trinity in Persons, and an Unity in nature.\par
\switchcolumn*

\LA
\begin{Resp}\Rsign Tibi laus, tibi glória, tibi gratiárum áctio in sǽcula sempitérna, \Star{} O beáta Trínitas.\end{Resp}
\begin{Resp}\Vsign Et benedíctum nomen glóriæ tuæ sanctum: et laudábile et superexaltátum in sǽcula.\end{Resp}
\begin{Resp}\Rsign O beáta Trínitas.\end{Resp}
\switchcolumn
\EN
\begin{Resp}\Rsign To thee be praise, to thee be glory, to thee be thanksgiving for ever and ever, \Star{} O Blessed Trinity.\end{Resp}
\begin{Resp}\Vsign And blessed is thy glorious and Holy Name, and to be praised and exalted above all for ever.\end{Resp}
\begin{Resp}\Rsign O Blessed Trinity.\end{Resp}
\switchcolumn*

\LA
\LessonHead{Lectio VI}
\switchcolumn
\EN
\LessonHead{Lesson VI}
\switchcolumn*

\LA
\noindent Per hanc unitátem naturálem totus Pater in Fílio et Spíritu Sancto est, totus Fílius in Patre et Spíritu Sancto est, totus quoque Spíritus Sanctus in Patre et Fílio. Nullus horum extra quémlibet ipsórum est: quia nemo álium aut præcédit æternitáte, aut excédit magnitúdine, aut súperat potestáte: quia nec Fílio nec Spíritu Sancto, quantum ad natúræ divínæ unitátem pértinet, aut antérior aut major Pater est; nec Fílii ætérnitas atque imménsitas, velut antérior aut major, Spíritus Sancti immensitátem æternitatémque aut præcédere aut excédere naturáliter potest.\par
\switchcolumn
\EN
\noindent Through this Oneness of nature All That the Father is is in the Son and the Holy Ghost, All That the Son is is in the Father and the Holy Ghost, and All That the Holy Ghost is is in the Father and the Son. Of the Father, the Son, and the Holy Ghost, None is without Other, None is before Other, None is Greater than Other, None is Mightier than Other. The Father, as touching the One Divine Nature, is neither before nor greater than the Son and the Holy Ghost neither is it possible that the Eternity and Infinity of the Son, whether as before or greater, should be before or greater than the Eternity and Infinity of the Spirit.\par
\switchcolumn*

\LA
\begin{Resp}\Rsign Magnus Dóminus, et laudábilis nimis: \Star{} Et sapiéntiæ ejus non est númerus.\end{Resp}
\begin{Resp}\Vsign Magnus Dóminus, et magna virtus ejus: et sapiéntiæ ejus non est finis.\end{Resp}
\begin{Resp}\Rsign Et sapiéntiæ ejus non est númerus.\end{Resp}
\begin{Resp}\Vsign Glória Patri, et Fílio, \Star{} et Spirítui Sancto.\end{Resp}
\begin{Resp}\Rsign Et sapiéntiæ ejus non est númerus.\end{Resp}
\switchcolumn
\EN
\begin{Resp}\Rsign Great is the Lord, and greatly to be praised; \Star{} And His Wisdom is unsearchable.\end{Resp}
\begin{Resp}\Vsign Great is our Lord, and of great power, and His understanding is infinite.\end{Resp}
\begin{Resp}\Rsign And His Wisdom is unsearchable.\end{Resp}
\begin{Resp}\Vsign Glory be to the Father, and to the Son, \Star{} and to the Holy Ghost.\end{Resp}
\begin{Resp}\Rsign And His Wisdom is unsearchable.\end{Resp}
\switchcolumn*

\LA
\LessonHead{Lectio VII}
\switchcolumn
\EN
\LessonHead{Lesson VII}
\switchcolumn*

\LA
\LessonTitle{Léctio sancti Evangélii secúndum Matthǽum}
\switchcolumn
\EN
\LessonTitle{From the Holy Gospel according to St. Matthew}
\switchcolumn*

\LA
\Cite{Matt 28:18-20}
\switchcolumn
\EN
\Cite{Matt 28:18-20}
\switchcolumn*

\LA
\noindent In illo témpore: Dixit Jesus discípulis suis: Data est mihi omnis potéstas in cælo et in terra. Eúntes ergo docéte omnes gentes, baptizántes eos in nómine Patris, et Fílii, et Spíritus Sancti. Et réliqua.\par
\switchcolumn
\EN
\noindent At that time, Jesus said unto His disciples: All power is given unto Me in heaven and in earth. Go ye, therefore, and teach all nations, baptizing them in the Name of the Father, and of the Son, and of the Holy Ghost. And so on.\par
\switchcolumn*

\LA
\LessonTitle{Homilía sancti Gregórii Nazianzéni}
\switchcolumn
\EN
\LessonTitle{Homily by St. Gregory of Nazianzus, Patriarch of Constantinople}
\switchcolumn*

\LA
\Source{In Tractatu de fide, post init.}
\switchcolumn
\EN
\Source{Treatise on the Faith}
\switchcolumn*

\LA
\noindent Quis catholicórum ignórat Patrem vere esse Patrem, Fílium vere esse Fílium, et Spíritum Sanctum vere esse Spíritum Sanctum? sicut ipse Dóminus ad Apóstolos suos dicit: Eúntes baptizáte omnes gentes in nómine Patris, et Fílii, et Spíritus Sancti. Hæc est perfécta Trínitas in unitáte consístens, quam scílicet uníus substántiæ profitémur. Non enim nos secúndum córporum conditiónem, divisiónem in Deo fácimus; sed secúndum divínæ natúræ poténtiam, quæ in matéria non est, et nóminum persónas vere constáre crédimus, et unitátem divinitátis esse testámur.\par
\switchcolumn
\EN
\noindent There is no Catholic but knoweth that the Father is a Very Father, the Son a Very Son, and the Holy Ghost a Very Holy Ghost, even as the Lord Himself saith unto His Apostles: “Go ye and baptize all nations in the Name of the Father, and of the Son, and of the Holy Ghost.” This is that Perfect Trinity Who is but One being, and of Whom therefore we testify that His Substance is one. For we make no division in God, as divisions are made in bodies, but we testify, that, according to the power of the Divine Nature, Which standeth not in matter, the Persons named have a real existence, and that God is One.\par
\switchcolumn*

\LA
\begin{Resp}\Rsign Benedicámus Patrem et Fílium cum Sancto Spíritu: \Star{} Laudémus et superexaltémus eum in sǽcula.\end{Resp}
\begin{Resp}\Vsign Benedíctus es, Dómine, in firmaménto cæli: et laudábilis et gloriósus in sǽcula.\end{Resp}
\begin{Resp}\Rsign Laudémus et superexaltémus eum in sǽcula.\end{Resp}
\switchcolumn
\EN
\begin{Resp}\Rsign Bless we the Father, and the Son, and the Holy Ghost. \Star{} Let us praise and exalt Him above all for ever.\end{Resp}
\begin{Resp}\Vsign Blessed art Thou, O Lord, in the firmament of heaven, and above all to be praised and glorified for ever.\end{Resp}
\begin{Resp}\Rsign Let us praise and exalt Him above all for ever.\end{Resp}
\switchcolumn*

\LA
\LessonHead{Lectio VIII}
\switchcolumn
\EN
\LessonHead{Lesson VIII}
\switchcolumn*

\LA
\noindent Nec extensiónem partis alicújus ex parte, ut quidam putavérunt, Dei Fílium dícimus: nec verbum sine re, velut sonum vocis, accípimus: sed tria nómina et tres persónas uníus esse esséntiæ, uníus majestátis atque poténtiæ crédimus. Et ídeo unum Deum confitémur: quia únitas majestátis, plúrium vocábulo deos próhibet appellári. Dénique Patrem et Fílium cathólice nominámus; duos autem Deos dícere, nec póssumus, nec debémus. Non quod Fílius Dei Deus non sit, immo verus Deus de Deo vero; sed quia non aliúnde, quam de ipso uno Patre, Dei Fílium nóvimus, perínde unum Deum dícimus. Hoc enim Prophétæ, hoc Apóstoli tradidérunt: hoc ipse Dóminus dócuit, cum dicit: Ego et Pater unum sumus. Unum ad unitátem divinitátis, ut dixi, refert; Sumus autem, persónis assígnat.\par
\switchcolumn
\EN
\noindent We do not say, as some have dreamt, that the Begetting of the Son of God is an outgrowing from one part to another part neither do we say that He is the Word in the sense of a mere sound uttered by a voice, but we do believe that these three Names and the Persons meant by them are all of only One Being, One Majesty, and One Power. And therefore we testify that God is one, because this One-ness of His Majesty forbiddeth that we should use the Plural form of speech and say, “Gods.” It is Catholic language to say, “Father and Son,” but we cannot and must not say that the Father and the Son are two gods. And that, not because the Son of God is not by Himself God; yea, He is Very God of Very God but because we know that the Son of God is not from elsewhere, but from the One Father Himself, and therefore we say that God is One. This is the doctrine which Prophets and Apostles have delivered; this is the doctrine which the Lord Himself taught when He said: “I and the Father are One,” (John x. 30,) that is, He meant, as touching the one Divine Being, but as touching Persons, We are distinct.\par
\switchcolumn*

\LA
\begin{Resp}\Rsign Duo Séraphim clamábant alter ad álterum: \Star{} Sanctus, sanctus, sanctus Dóminus Deus Sábaoth: * Plena est omnis terra glória ejus.\end{Resp}
\begin{Resp}\Vsign Tres sunt qui testimónium dant in cælo: Pater, Verbum, et Spíritus Sanctus: et hi tres unum sunt.\end{Resp}
\begin{Resp}\Rsign Sanctus, sanctus, sanctus Dóminus Deus Sábaoth.\end{Resp}
\begin{Resp}\Vsign Glória Patri, et Fílio, \Star{} et Spirítui Sancto.\end{Resp}
\begin{Resp}\Rsign Plena est omnis terra glória ejus.\end{Resp}
\switchcolumn
\EN
\begin{Resp}\Rsign One Seraph cried unto another: \Star{} Holy, Holy, Holy is the Lord God of hosts; the whole earth is full of His glory.\end{Resp}
\begin{Resp}\Vsign There are Three That bear record in heaven: the Father, the Word, and the Holy Ghost, and these Three are One.\end{Resp}
\begin{Resp}\Rsign Holy, Holy, Holy is the Lord God of hosts.\end{Resp}
\begin{Resp}\Vsign Glory be to the Father, and to the Son, \Star{} and to the Holy Ghost.\end{Resp}
\begin{Resp}\Rsign The whole earth is full of His glory.\end{Resp}
\switchcolumn*

\LA
\LessonHead{Lectio IX}
\switchcolumn
\EN
\LessonHead{Lesson IX}
\switchcolumn*

\LA
\LessonTitle{Commemoratio Dominicæ}
\switchcolumn
\EN
\LessonTitle{Commemoration for the First Sunday after Pentecost}
\switchcolumn*

\LA
\Cite{Luc 6:36-42}
\switchcolumn
\EN
\Cite{Luke 6:36-42}
\switchcolumn*

\LA
\noindent Léctio sancti Evangélii secúndum Lucam\par
\switchcolumn
\EN
\noindent From the Holy Gospel according to Luke\par
\switchcolumn*

\LA
In illo témpore: Dixit Jesus discípulis suis: Estóte misericórdes, sicut et Pater vester miséricors est. Et réliqua.\par
\switchcolumn
\EN
At that time: Jesus said unto His disciples Be ye merciful, as your Father also is merciful. And so on.\par
\switchcolumn*

\LA
\LessonTitle{Homilía sancti Augustíni Epíscopi}
\switchcolumn
\EN
\LessonTitle{Homily by St. Augustine, Bishop of Hippo}
\switchcolumn*

\LA
\Source{Serm. 15. in Evang. Matthǽi de Verbis Dómini, post inítium}
\switchcolumn
\EN
\Source{15th Sermon on Matthew, Words of the Lord}
\switchcolumn*

\LA
\noindent Duo sunt ópera misericórdiæ, quæ nos líberant, quæ bréviter ipse Dóminus pósuit in Evangélio: Dimíttite, et dimittétur vobis: date, et dábitur vobis. Dimíttite, et dimittétur vobis, ad ignoscéndum pértinet: Date, et dábitur vobis, ad præstándum benefícium pértinet. Quod ait de ignoscéndo, et tu vis tibi ignósci quod peccas, et habes álium, cui tu possis ignóscere. Rursus, quod pértinet ad tribuéndum benefícium, petit te mendícus, et tu es Dei mendícus. Omnes enim quando orámus, mendíci Dei sumus: ante jánuam magni Patrisfamílias stamus, immo et prostérnimur, súpplices ingemíscimus, áliquid voléntes accípere; et ipsum áliquid ipse Deus est. Quid a te petit mendícus? Panem. Et tu quid petis a Deo, nisi Christum, qui dicit: Ego sum panis vivus, qui de cælo descéndi? Ignósci vobis vultis? ignóscite: Remíttite, et remittétur vobis. Accípere vultis? date, et dábitur vobis.\par
\switchcolumn
\EN
\noindent There are two works of mercy which free us, and which the Lord Himself hath briefly named in the Gospel “Forgive, and ye shall be forgiven give, and it shall be given unto you.” “Forgive, and ye shall be forgiven”, such is the promise of pardon. “Give, and it shall be given unto you”, such is the promise of favour. As touching forgiveness thou hast trespasses which thou wouldest fain have forgiven, and them which have trespassed against thee, whom thou canst forgive. As touching favour, there are beggars that beg from thee, and thou art a beggar to God. When we pray, we are all beggars to God, standing at the door of the Great Householder, yea, falling down on our knees, and beseeching Him to give us somewhat and that somewhat is God Himself. What doth a beggar ask of thee Bread. And what dost thou ask of God but that Christ Who saith: “I am the Living Bread Which came down from heaven?” If ye will be forgiven, forgive ye. If ye will be pardoned, pardon ye. If ye will receive, “give, and it shall be given unto you.”\par
\switchcolumn*

\LA
\TeDeum{Te Deum}
\switchcolumn
\EN
\TeDeum{Te Deum}
\switchcolumn*

\end{paracol}

\Office{Feria Secunda infra Hebdomadam I post Octavam Pentecostes}{Monday of the First Week after Pentecost}{Feria}
\addcontentsline{toc}{subsection}{Feria Secunda infra Hebdomadam I post Octavam Pentecostes \textit{— Monday of the First Week after Pentecost}}
\begin{paracol}{2}
\LA
\LessonHead{Lectio I}
\switchcolumn
\EN
\LessonHead{Lesson I}
\switchcolumn*

\LA
\LessonTitle{Incipit liber primus Regum}
\switchcolumn
\EN
\LessonTitle{Lesson from the first book of Samuel}
\switchcolumn*

\LA
\Cite{1 Reg 1:1-3}
\switchcolumn
\EN
\Cite{1 Sam 1:1-3}
\switchcolumn*

\LA
\noindent Fuit vir unus de Ramáthaim Sophim, de monte Ephraim, et nomen ejus Elcana, fílius Iéroham, fílii Eliu, fílii Thohu, fílii Suph, Ephrathǽus: et hábuit duas uxóres, nomen uni Anna, et nomen secúndæ Phenénna. Fuerúntque Phenénnæ fílii: Annæ autem non erant líberi. Et ascendébat vir ille de civitáte sua statútis diébus, ut adoráret et sacrificáret Dómino exercítuum in Silo. Erant autem ibi duo fílii Heli, Ophni et Phínees, sacerdótes Dómini.\par
\switchcolumn
\EN
\noindent There was a man of Ramathaimsophim, of mount Ephraim, and his name was Elcana, the son of Jeroham, the son of Eliu, the son of Thohu, the son of Suph, an Ephraimite: And he had two wives, the name of one was Anna, and the name of the other Phenenna. Phenenna had children: but Anna had no children. And this man went up out of his city upon the appointed days, to adore and to offer sacrifice to the Lord of hosts in Silo. And the two sons of Heli, Ophni and Phinees, were there priests of the Lord.\par
\switchcolumn*

\LA
\begin{Resp}\Rsign Præparáte corda vestra Dómino, et servíte illi soli: \Star{} Et liberábit vos de mánibus inimicórum vestrórum.\end{Resp}
\begin{Resp}\Vsign Convertímini ad eum in toto corde vestro, et auférte deos aliénos de médio vestri.\end{Resp}
\begin{Resp}\Rsign Et liberábit vos de mánibus inimicórum vestrórum.\end{Resp}
\switchcolumn
\EN
\begin{Resp}\Rsign Prepare your hearts unto the Lord, and serve Him Only \Star{} And He will deliver you out of the hand of your enemies.\end{Resp}
\begin{Resp}\Vsign Return unto Him with all your hearts, and put away the strange gods from among you.\end{Resp}
\begin{Resp}\Rsign And He will deliver you out of the hand of your enemies.\end{Resp}
\switchcolumn*

\LA
\LessonHead{Lectio II}
\switchcolumn
\EN
\LessonHead{Lesson II}
\switchcolumn*

\LA
\Cite{1 Reg 1:4-8}
\switchcolumn
\EN
\Cite{1 Sam 1:4-8}
\switchcolumn*

\LA
\noindent Venit ergo dies, et immolávit Elcana, dedítque Phenénnæ uxóri suæ, et cunctis fíliis ejus et filiábus, partes: Annæ autem dedit partem unam tristis, quia Annam diligébat. Dóminus autem conclúserat vulvam ejus. Affligébat quoque eam ǽmula ejus, et veheménter angébat, in tantum ut exprobráret quod Dóminus conclusísset vulvam ejus: sicque faciébat per síngulos annos: cum redeúnte témpore ascénderent ad templum Dómini, et sic provocábat eam: porro illa flebat, et non capiébat cibum. Dixit ergo ei Elcana vir suus: Anna, cur fles? et quare non cómedis? et quam ob rem afflígitur cor tuum? numquid non ego mélior tibi sum, quam decem fílii?\par
\switchcolumn
\EN
\noindent Now the day came, and Elcana offered sacrifice, and gave to Phenenna his wife, and to all her sons and daughters, portions: But to Anna he gave one portion with sorrow, because he loved Anna. And the Lord had shut up her womb. Her rival also afflicted her, and troubled her exceedingly, insomuch that she upbraided her, that the Lord had shut up her womb: And thus she did every year, when the time returned that they went up to the temple of the Lord: and thus she provoked her: but Anna wept, and did not eat. Then Elcana her husband said to her: Anna, why weepest thou? and why dost thou not eat? And why dost thou afflict thy heart? Am not I better to thee than ten children?\par
\switchcolumn*

\LA
\begin{Resp}\Rsign Deus ómnium exaudítor est: ipse misit Angelum suum, et tulit me de óvibus patris mei: \Star{} Et unxit me unctióne misericórdiæ suæ.\end{Resp}
\begin{Resp}\Vsign Dóminus, qui erípuit me de ore leónis, et de manu béstiæ liberávit me.\end{Resp}
\begin{Resp}\Rsign Et unxit me unctióne misericórdiæ suæ.\end{Resp}
\switchcolumn
\EN
\begin{Resp}\Rsign God, Which heareth all, even He sent His Angel, and took me from keeping my father's sheep, and \Star{} Anointed me with the oil of His mercy.\end{Resp}
\begin{Resp}\Vsign The Lord That delivered me out of the mouth of the lion, and out of the paw of the bear\end{Resp}
\begin{Resp}\Rsign And anointed me with the oil of His mercy.\end{Resp}
\switchcolumn*

\LA
\LessonHead{Lectio III}
\switchcolumn
\EN
\LessonHead{Lesson III}
\switchcolumn*

\LA
\Cite{1 Reg 1:9-11}
\switchcolumn
\EN
\Cite{1 Sam 1:9-11}
\switchcolumn*

\LA
\noindent Surréxit autem Anna postquam coméderat et bíberat in Silo. Et Heli sacerdóte sedénte super sellam ante postes templi Dómini, cum esset Anna amáro ánimo, orávit ad Dóminum, flens lárgiter, et votum vovit, dicens: Dómine exercítuum, si respíciens víderis afflictiónem fámulæ tuæ, et recordátus mei fúeris, nec oblítus ancíllæ tuæ, dederísque servæ tuæ sexum virílem: dabo eum Dómino ómnibus diébus vitæ ejus, et novácula non ascéndet super caput ejus.\par
\switchcolumn
\EN
\noindent So Anna arose after she had eaten and drunk in Silo: And Heli the priest sitting upon a stool, before the door of the temple of the Lord: As Anna had her heart full of grief, she prayed to the Lord, shedding many tears, And she made a vow, saying: O Lord, of hosts, if thou wilt look down on the affliction of thy servant, and wilt be mindful of me, and not forget thy handmaid, and wilt give to thy servant a man child: I will give him to the Lord all the days of his life, and no razor shall come upon his head.\par
\switchcolumn*

\LA
\begin{Resp}\Rsign Dóminus, qui erípuit me de ore leónis, et de manu béstiæ liberávit me, \Star{} Ipse me erípiet de mánibus inimicórum meórum.\end{Resp}
\begin{Resp}\Vsign Misit Deus misericórdiam suam et veritátem suam: ánimam meam erípuit de médio catulórum leónum.\end{Resp}
\begin{Resp}\Rsign Ipse me erípiet de mánibus inimicórum meórum.\end{Resp}
\begin{Resp}\Vsign Glória Patri, et Fílio, \Star{} et Spirítui Sancto.\end{Resp}
\begin{Resp}\Rsign Ipse me erípiet de mánibus inimicórum meórum.\end{Resp}
\switchcolumn
\EN
\begin{Resp}\Rsign The Lord That delivered me out of the mouth of the lion, and out of the paw of the bear \Star{} He will deliver me out of the hand of mine enemies.\end{Resp}
\begin{Resp}\Vsign God hath sent forth His mercy and His truth, and delivered my soul from among the lion's whelps.\end{Resp}
\begin{Resp}\Rsign He will deliver me out of the hand of mine enemies.\end{Resp}
\begin{Resp}\Vsign Glory be to the Father, and to the Son, \Star{} and to the Holy Ghost.\end{Resp}
\begin{Resp}\Rsign He will deliver me out of the hand of mine enemies.\end{Resp}
\switchcolumn*

\end{paracol}

\Office{Feria Tertia infra Hebdomadam I post Octavam Pentecostes}{Tuesday of the First Week after Pentecost}{Feria}
\addcontentsline{toc}{subsection}{Feria Tertia infra Hebdomadam I post Octavam Pentecostes \textit{— Tuesday of the First Week after Pentecost}}
\begin{paracol}{2}
\LA
\LessonHead{Lectio I}
\switchcolumn
\EN
\LessonHead{Lesson I}
\switchcolumn*

\LA
\LessonTitle{De libro primo Regum}
\switchcolumn
\EN
\LessonTitle{Lesson from the first book of Samuel}
\switchcolumn*

\LA
\Cite{1 Reg 1:12-18}
\switchcolumn
\EN
\Cite{1 Sam 1:12-18}
\switchcolumn*

\LA
\noindent Factum est autem, cum illa multiplicáret preces coram Dómino, ut Heli observáret os ejus. Porro Anna loquebátur in corde suo, tantúmque lábia illíus movebántur, et vox pénitus non audiebátur. Æstimávit ergo eam Heli temuléntam, dixítque ei: Usquequo ébria eris? dígere paulísper vinum, quo mades. Respóndens Anna: Nequáquam, inquit, dómine mi: nam múlier infélix nimis ego sum: vinúmque et omne quod inebriáre potest, non bibi, sed effúdi ánimam meam in conspéctu Dómini. Ne réputes ancíllam tuam quasi unam de filiábus Bélial: quia ex multitúdine dolóris et mæróris mei locúta sum usque in præsens. Tunc Heli ait ei: Vade in pace: et Deus Israël det tibi petitiónem tuam quam rogásti eum. Et illa dixit: Utinam invéniat ancílla tua grátiam in óculis tuis.\par
\switchcolumn
\EN
\noindent And it came to pass, as she multiplied prayers before the Lord, that Heli observed her mouth. Now Anna spoke in her heart, and only her lips moved, but her voice was not heard at all. Heli therefore thought her to be drunk, And said to her: How long wilt thou, be drunk? digest a little the wine, of which thou hast taken too much. Anna answering, said: Not so, my lord: for I am an exceeding unhappy woman, and have drunk neither wine nor any strong drink, but I have poured out my soul before the Lord. Count not thy handmaid for one of the daughters of Belial: for out of the abundance of my sorrow and grief have I spoken till now. Then Heli said to her: Go in peace: and the God of Israel grant thee thy petition, which thou hast asked of him. And she said: Would to God thy handmaid may find grace in thy eyes.\par
\switchcolumn*

\LA
\begin{Resp}\Rsign Percússit Saul mille, et David decem míllia: \Star{} Quia manus Dómini erat cum illo: percússit Philisthǽum, et ábstulit oppróbrium ex Israël.\end{Resp}
\begin{Resp}\Vsign Nonne iste est David, de quo canébant in choro, dicéntes: Saul percússit mille, et David decem míllia?\end{Resp}
\begin{Resp}\Rsign Quia manus Dómini erat cum illo: percússit Philisthǽum, et ábstulit oppróbrium ex Israël.\end{Resp}
\switchcolumn
\EN
\begin{Resp}\Rsign Saul hath slain his thousands, and David his ten thousands. \Star{} Because the hand of the Lord was with him, he smote the Philistine, and took away the reproach from Israel.\end{Resp}
\begin{Resp}\Vsign Is not this David? Did they not sing one to another of him in dances, saying: Saul hath slain his thousands, and David his ten thousands?\end{Resp}
\begin{Resp}\Rsign Because the hand of the Lord was with him, he smote the Philistine, and took away the reproach from Israel.\end{Resp}
\switchcolumn*

\LA
\LessonHead{Lectio II}
\switchcolumn
\EN
\LessonHead{Lesson II}
\switchcolumn*

\LA
\Cite{1 Reg 1:18-22}
\switchcolumn
\EN
\Cite{1 Sam 1:18-22}
\switchcolumn*

\LA
\noindent Et ábiit múlier in viam suam, et comédit, vultúsque illíus non sunt ámplius in divérsa mutáti. Et surrexérunt mane, et adoravérunt coram Dómino: reversíque sunt, et venérunt in domum suam Rámatha. Cognóvit autem Elcana Annam uxórem suam: et recordátus est ejus Dóminus. Et factum est post círculum diérum, concépit Anna, et péperit fílium: vocavítque nomen ejus Sámuel, eo quod a Dómino postulásset eum. Ascéndit autem vir ejus Elcana, et omnis domus ejus, ut immoláret Dómino hóstiam solémnem, et votum suum. Et Anna non ascéndit: dixit enim viro suo: Non vadam donec ablactétur infans, et ducam eum, ut appáreat ante conspéctum Dómini, et máneat ibi júgiter.\par
\switchcolumn
\EN
\noindent So the woman went on her way, and ate, and her countenance was no more changed. And they rose in the morning, and worshipped before the Lord: and they returned, and came into their house at Ramatha. And Elcana knew Anna his wife: and the Lord remembered her. And it came to pass when the time was come about, Anna conceived and bore a son, and called his name Samuel: because she had asked him of the Lord. And Elcana her husband went up, and all his house, to offer to the Lord the solemn sacrifice, and his vow. But Anna went not up: for she said to her husband: I will not go till the child be weaned, and till I may carry him, that he may appear before the Lord, and may abide always there.\par
\switchcolumn*

\LA
\begin{Resp}\Rsign Montes Gélboë, nec ros nec plúvia véniant super vos, \Star{} Ubi cecidérunt fortes Israël.\end{Resp}
\begin{Resp}\Vsign Omnes montes, qui estis in circúitu ejus, vísitet Dóminus: a Gélboë autem tránseat.\end{Resp}
\begin{Resp}\Rsign Ubi cecidérunt fortes Israël.\end{Resp}
\switchcolumn
\EN
\begin{Resp}\Rsign Ye mountains of Gilboa, let there be no dew, neither let there be rain upon you \Star{} For there are the mighty of Israel fallen.\end{Resp}
\begin{Resp}\Vsign All ye mountains that stand round about, the Lord look upon you but let Him pass by Gilboa\end{Resp}
\begin{Resp}\Rsign For there are the mighty of Israel fallen.\end{Resp}
\switchcolumn*

\LA
\LessonHead{Lectio III}
\switchcolumn
\EN
\LessonHead{Lesson III}
\switchcolumn*

\LA
\Cite{1 Reg 1:23-28}
\switchcolumn
\EN
\Cite{1 Sam 1:23-28}
\switchcolumn*

\LA
\noindent Et ait ei Elcana vir suus: Fac quod bonum tibi vidétur, et mane donec abláctes eum: precórque ut ímpleat Dóminus verbum suum. Mansit ergo múlier, et lactávit fílium suum, donec amovéret eum a lacte. Et addúxit eum secum, postquam ablactáverat, in vítulis tribus, et tribus módiis farínæ, et ámphora vini, et addúxit eum ad domum Dómini in Silo. Puer autem erat adhuc infántulus: et immolavérunt vítulum, et obtulérunt púerum Heli. Et ait Anna: Obsecro mi dómine, vivit ánima tua, dómine: ego sum illa múlier, quæ steti coram te hic orans Dóminum. Pro púero isto orávi, et dedit mihi Dóminus petitiónem meam quam postulávi eum. Idcírco et ego commodávi eum Dómino cunctis diébus quibus fúerit commodátus Dómino. Et adoravérunt ibi Dóminum.\par
\switchcolumn
\EN
\noindent And Elcana her husband said to her: Do what seemeth good to thee, and stay till thou wean him: and I pray that the Lord may fulfill his word. So the woman stayed at home, and gave her son suck, till she weaned him. And after she had weaned him, she carried him with her, with three calves, and three bushels of flour, and a bottle of wine, and she brought him to the house of the Lord in Silo. Now the child was as yet very young: And they immolated a calf, and offered the child to Heli. And Anna said: I beseech thee, my lord, as thy soul liveth, my lord: I am that woman who stood before thee here praying to the Lord. For this child did I pray, and the Lord hath granted me my petition, which I asked of him. Therefore I also have lent him to the Lord all the days of his life, he shall be lent to the Lord. And they adored the Lord there.\par
\switchcolumn*

\LA
\begin{Resp}\Rsign Ego te tuli de domo patris tui, dicit Dóminus, et pósui te páscere gregem pópuli mei: \Star{} Et fui tecum in ómnibus ubicúmque ambulásti, firmans regnum tuum in ætérnum.\end{Resp}
\begin{Resp}\Vsign Fecíque tibi nomen grande, juxta nomen magnórum, qui sunt in terra: et réquiem dedi tibi ab ómnibus inimícis tuis.\end{Resp}
\begin{Resp}\Rsign Et fui tecum in ómnibus ubicúmque ambulásti, firmans regnum tuum in ætérnum.\end{Resp}
\begin{Resp}\Vsign Glória Patri, et Fílio, \Star{} et Spirítui Sancto.\end{Resp}
\begin{Resp}\Rsign Et fui tecum in ómnibus ubicúmque ambulásti, firmans regnum tuum in ætérnum.\end{Resp}
\switchcolumn
\EN
\begin{Resp}\Rsign Thus saith the Lord: I took thee out of thy father's house, and appointed thee to be ruler over My people, over Israel. \Star{} And I was with thee whithersoever thou wentest, to establish thy kingdom for ever.\end{Resp}
\begin{Resp}\Vsign And I have made thee a great name, like unto the name of the great men that are in the earth, and have caused thee to rest from all thine enemies.\end{Resp}
\begin{Resp}\Rsign And I was with thee whithersoever thou wentest, to establish thy kingdom for ever.\end{Resp}
\begin{Resp}\Vsign Glory be to the Father, and to the Son, \Star{} and to the Holy Ghost.\end{Resp}
\begin{Resp}\Rsign And I was with thee whithersoever thou wentest, to establish thy kingdom for ever.\end{Resp}
\switchcolumn*

\end{paracol}

\Office{Feria Quarta infra Hebdomadam I post Octavam Pentecostes}{Wednesday of the First Week after Pentecost}{Feria}
\addcontentsline{toc}{subsection}{Feria Quarta infra Hebdomadam I post Octavam Pentecostes \textit{— Wednesday of the First Week after Pentecost}}
\begin{paracol}{2}
\LA
\LessonHead{Lectio I}
\switchcolumn
\EN
\LessonHead{Lesson I}
\switchcolumn*

\LA
\LessonTitle{De libro primo Regum}
\switchcolumn
\EN
\LessonTitle{Lesson from the first book of Samuel}
\switchcolumn*

\LA
\Cite{1 Reg 2:12-14}
\switchcolumn
\EN
\Cite{1 Sam 2:12-14}
\switchcolumn*

\LA
\noindent Porro fílii Heli, fílii Bélial, nesciéntes Dóminum, neque offícium sacerdótum ad pópulum: sed quicúmque immolásset víctimam, veniébat puer sacerdótis, dum coqueréntur carnes, et habébat fuscínulam tridéntem in manu sua, et mittébat eam in lebétem, vel in caldáriam, aut in ollam, sive in cácabum: et omne quod levábat fuscínula, tollébat sacérdos sibi: sic faciébant univérso Israéli veniéntium in Silo.\par
\switchcolumn
\EN
\noindent Now the sons of Heli were children of Belial, not knowing the Lord, nor the office of the priests to the people: but whosoever had offered a sacrifice, the servant of the priest came, while the flesh was in boiling, with a flesh-hook of three teeth in his hand, and thrust it into the kettle, or into the cauldron, or into the pot, or into the pan: and all that the flesh-hook brought up, the priest took to himself. Thus did they to all Israel that came to Silo.\par
\switchcolumn*

\LA
\begin{Resp}\Rsign Peccávi super númerum arénæ maris, et multiplicáta sunt peccáta mea: et non sum dignus vidére altitúdinem cæli præ multitúdine iniquitátis meæ: quóniam irritávi iram tuam, \Star{} Et malum coram te feci.\end{Resp}
\begin{Resp}\Vsign Quóniam iniquitátem meam ego cognósco: et delíctum meum contra me est semper, quia tibi soli peccávi.\end{Resp}
\begin{Resp}\Rsign Et malum coram te feci.\end{Resp}
\switchcolumn
\EN
\begin{Resp}\Rsign My sins are many, yea, they are more in number than the sands of the sea; I am not worthy to look up toward heaven because of the multitude of my iniquities; for I have provoked thee to anger; \Star{} And done evil in thy sight.\end{Resp}
\begin{Resp}\Vsign For I acknowledge my transgression, and my sin is ever before me, for against thee only have I sinned.\end{Resp}
\begin{Resp}\Rsign And done evil in thy sight.\end{Resp}
\switchcolumn*

\LA
\LessonHead{Lectio II}
\switchcolumn
\EN
\LessonHead{Lesson II}
\switchcolumn*

\LA
\Cite{1 Reg 2:15-17}
\switchcolumn
\EN
\Cite{1 Sam 2:15-17}
\switchcolumn*

\LA
\noindent Etiam ántequam adolérent ádipem, veniébat puer sacerdótis, et dicébat immolánti: Da mihi carnem, ut coquam sacerdóti: non enim accípiam a te carnem coctam, sed crudam. Dicebátque illi ímmolans: Incendátur primum juxta morem hódie adeps, et tolle tibi quantumcúmque desíderat ánima tua. Qui respóndens ajébat ei: Nequáquam: nunc enim dabis, alióquin tollam vi. Erat ergo peccátum puerórum grande nimis coram Dómino: quia retrahébant hómines a sacrifício Dómini.\par
\switchcolumn
\EN
\noindent Also before they burnt the fat, the servant of the priest came, and said to the man that sacrificed: Give me flesh to boil for the priest: for I will not take of thee sodden flesh, but raw. And he that sacrificed said to him: Let the fat first be burnt today according to the custom, and then take as much as thy soul desireth. But he answered and said to him: Not so: but thou shalt give it me now, or else I will take it by force. Wherefore the sin of the young men was exceeding great before the Lord: because they withdrew men from the sacrifice of the Lord.\par
\switchcolumn*

\LA
\begin{Resp}\Rsign Exaudísti, Dómine, oratiónem servi tui, ut ædificárem templum nómini tuo: \Star{} Bénedic et sanctífica domum istam in sempitérnum, Deus Israël.\end{Resp}
\begin{Resp}\Vsign Dómine, qui custódis pactum cum servis tuis, qui ámbulant coram te in toto corde suo.\end{Resp}
\begin{Resp}\Rsign Bénedic et sanctífica domum istam in sempitérnum, Deus Israël.\end{Resp}
\switchcolumn
\EN
\begin{Resp}\Rsign O Lord, Thou hast hearkened unto the prayer of thy servant, that I might build a temple unto thy Name, \Star{} O God of Israel, bless Thou, and hallow this house for ever.\end{Resp}
\begin{Resp}\Vsign O Lord, Who keepest covenant with thy servants that walk before thee in all their heart.\end{Resp}
\begin{Resp}\Rsign O God of Israel, bless Thou, and hallow this house for ever.\end{Resp}
\switchcolumn*

\LA
\LessonHead{Lectio III}
\switchcolumn
\EN
\LessonHead{Lesson III}
\switchcolumn*

\LA
\Cite{1 Reg 2:18-21}
\switchcolumn
\EN
\Cite{1 Sam 2:18-21}
\switchcolumn*

\LA
\noindent Sámuel autem ministrábat ante fáciem Dómini, puer accínctus ephod líneo. Et túnicam parvam faciébat ei mater sua, quam afferébat statútis diébus, ascéndens cum viro suo, ut immoláret hóstiam solémnem. Et benedíxit Heli Elcanæ et uxóri ejus: dixítque ei: Reddat tibi Dóminus semen de mulíere hac, pro fœ́nore quod commodásti Dómino. Et abiérunt in locum suum. Visitávit ergo Dóminus Annam, et concépit, et péperit tres fílios, et duas fílias: et magnificátus est puer Sámuel apud Dóminum.\par
\switchcolumn
\EN
\noindent But Samuel ministered before the face of the Lord: being a child girded with a linen ephod. And his mother made him a little coat, which she brought to him on the appointed days, when she went up with her husband, to offer the solemn sacrifice. And Heli blessed Elcana and his wife: and he said to him: The Lord give thee seed of this woman, for the loan thou hast lent to the Lord. And they went to their own home. And the Lord visited Anna, and she conceived, and bore three sons and two daughters: and the child Samuel became great before the Lord.\par
\switchcolumn*

\LA
\begin{Resp}\Rsign Audi, Dómine, hymnum et oratiónem, quam servus tuus orat coram te hódie, ut sint óculi tui apérti, et aures tuæ inténtæ, \Star{} Super domum istam die ac nocte.\end{Resp}
\begin{Resp}\Vsign Réspice, Dómine, de sanctuário tuo, et de excélso cælórum habitáculo.\end{Resp}
\begin{Resp}\Rsign Super domum istam die ac nocte.\end{Resp}
\begin{Resp}\Vsign Glória Patri, et Fílio, \Star{} et Spirítui Sancto.\end{Resp}
\begin{Resp}\Rsign Super domum istam die ac nocte.\end{Resp}
\switchcolumn
\EN
\begin{Resp}\Rsign Hearken, O Lord, unto the cry and to the prayer which thy servant prayeth before thee today, that thine eyes may be open and thine ears attend; \Star{} Toward this house day and night.\end{Resp}
\begin{Resp}\Vsign Look down from thine high and holy place, O Lord, even from heaven thy dwelling.\end{Resp}
\begin{Resp}\Rsign Toward this house, day and night.\end{Resp}
\begin{Resp}\Vsign Glory be to the Father, and to the Son, \Star{} and to the Holy Ghost.\end{Resp}
\begin{Resp}\Rsign Toward this house, day and night.\end{Resp}
\switchcolumn*

\end{paracol}

\Office{Festum Sanctissimi Corporis Christi}{Corpus Christi}{Duplex I classis cum Octava privilegiata II ordinis}
\addcontentsline{toc}{subsection}{Festum Sanctissimi Corporis Christi \textit{— Corpus Christi}}
\begin{paracol}{2}
\LA
\LessonHead{Lectio I}
\switchcolumn
\EN
\LessonHead{Lesson I}
\switchcolumn*

\LA
\LessonTitle{De Epístola prima beáti Pauli Apóstoli ad Corínthios}
\switchcolumn
\EN
\LessonTitle{Lesson from the first letter of St. Paul the Apostle to the Corinthians}
\switchcolumn*

\LA
\Cite{1 Cor 11:20-22}
\switchcolumn
\EN
\Cite{1 Cor 11:20-22}
\switchcolumn*

\LA
\noindent Conveniéntibus ergo vobis in unum, jam non est Domínicam cenam manducáre. Unusquísque enim suam cenam præsúmit ad manducándum. Et álius quidem ésurit, álius autem ébrius est. Numquid domos non habétis ad manducándum et bibéndum? aut Ecclésiam Dei contémnitis, et confúnditis eos, qui non habent? Quid dicam vobis? Laudo vos? In hoc non laudo.\par
\switchcolumn
\EN
\noindent When you come therefore together into one place, it is not now to eat the Lord's supper. For every one taketh before his own supper to eat. And one indeed is hungry and another is drunk. What, have you not houses to eat and to drink in? Or despise ye the church of God; and put them to shame that have not? What shall I say to you? Do I praise you? In this I praise you not.\par
\switchcolumn*

\LA
\begin{Resp}\Rsign Immolábit hædum multitúdo filiórum Israël ad vésperam Paschæ: \Star{} Et edent carnes et ázymos panes.\end{Resp}
\begin{Resp}\Vsign Pascha nostrum immolátus est Christus: ítaque epulémur in ázymis sinceritátis et veritátis.\end{Resp}
\begin{Resp}\Rsign Et edent carnes et ázymos panes.\end{Resp}
\switchcolumn
\EN
\begin{Resp}\Rsign The whole assembly of the children of Israel shall kill the lamb toward the evening of the Passover. \Star{} And they shall eat the flesh, and unleavened bread.\end{Resp}
\begin{Resp}\Vsign Even Christ our Passover is sacrificed for us therefore let us keep the feast with the unleavened bread of sincerity and truth.\end{Resp}
\begin{Resp}\Rsign And they shall eat the flesh, and unleavened bread.\end{Resp}
\switchcolumn*

\LA
\LessonHead{Lectio II}
\switchcolumn
\EN
\LessonHead{Lesson II}
\switchcolumn*

\LA
\Cite{1 Cor 11:23-26}
\switchcolumn
\EN
\Cite{1 Cor 11:23-26}
\switchcolumn*

\LA
\noindent Ego enim accépi a Dómino quod et trádidi vobis, quóniam Dóminus Jesus, in qua nocte tradebátur, accépit panem, et grátias agens fregit, et dixit: Accípite, et manducáte: hoc est corpus meum, quod pro vobis tradétur: hoc fácite in meam commemoratiónem. Simíliter et cálicem, postquam cœnávit, dicens: Hic calix novum testaméntum est in meo sánguine: hoc fácite, quotiescúmque bibétis, in meam commemoratiónem. Quotiescúmque enim manducábitis panem hunc, et cálicem bibétis, mortem Dómini annuntiábitis donec véniat.\par
\switchcolumn
\EN
\noindent For I have received of the Lord that which also I delivered unto you, that the Lord Jesus, the same night in which he was betrayed, took bread. And giving thanks, broke, and said: Take ye, and eat: this is my body, which shall be delivered for you: this do for the commemoration of me. In like manner also the chalice, after he had supped, saying: This chalice is the new testament in my blood: this do ye, as often as you shall drink, for the commemoration of me. For as often as you shall eat this bread, and drink the chalice, you shall show the death of the Lord, until he come.\par
\switchcolumn*

\LA
\begin{Resp}\Rsign Comedétis carnes, et saturabímini pánibus: \Star{} Iste est panis, quem dedit vobis Dóminus ad vescéndum.\end{Resp}
\begin{Resp}\Vsign Non Móyses dedit vobis panem de cælo, sed Pater meus dat vobis panem de cælo verum.\end{Resp}
\begin{Resp}\Rsign Iste est panis, quem dedit vobis Dóminus ad vescéndum.\end{Resp}
\switchcolumn
\EN
\begin{Resp}\Rsign Ye shall eat flesh, and shall be filled with bread. \Star{} This is the bread which the Lord hath given you to eat.\end{Resp}
\begin{Resp}\Vsign Moses gave you not that Bread from heaven, but My Father giveth you the true Bread from heaven.\end{Resp}
\begin{Resp}\Rsign This is the bread which the Lord hath given you to eat.\end{Resp}
\switchcolumn*

\LA
\LessonHead{Lectio III}
\switchcolumn
\EN
\LessonHead{Lesson III}
\switchcolumn*

\LA
\Cite{1 Cor 11:27-32}
\switchcolumn
\EN
\Cite{1 Cor 11:27-32}
\switchcolumn*

\LA
\noindent Itaque quicúmque manducáverit panem hunc, vel bíberit cálicem Dómini indígne, reus erit córporis et sánguinis Dómini. Probet autem seípsum homo: et sic de pane illo edat, et de cálice bibat. Qui enim mandúcat et bibit indígne, judícium sibi mandúcat et bibit, non dijúdicans corpus Dómini. Ideo inter vos multi infírmi et imbecílles, et dórmiunt multi. Quod, si nosmetípsos dijudicarémus, non útique judicarémur. Dum judicámur autem, a Dómino corrípimur, ut non cum hoc mundo damnémur.\par
\switchcolumn
\EN
\noindent Therefore whosoever shall eat this bread, or drink the chalice of the Lord unworthily, shall be guilty of the body and of the blood of the Lord. But let a man prove himself: and so let him eat of that bread, and drink of the chalice. For he that eateth and drinketh unworthily, eateth and drinketh judgment to himself, not discerning the body of the Lord. Therefore are there many infirm and weak among you, and many sleep. But if we would judge ourselves, we should not be judged. But whilst we are judged, we are chastised by the Lord, that we be not condemned with this world.\par
\switchcolumn*

\LA
\begin{Resp}\Rsign Respéxit Elías ad caput suum subcinerícium panem: qui surgens comédit et bibit: \Star{} Et ambulávit in fortitúdine cibi illíus usque ad montem Dei.\end{Resp}
\begin{Resp}\Vsign Si quis manducáverit ex hoc pane, vivet in ætérnum.\end{Resp}
\begin{Resp}\Rsign Et ambulávit in fortitúdine cibi illíus usque ad montem Dei.\end{Resp}
\begin{Resp}\Vsign Glória Patri, et Fílio, \Star{} et Spirítui Sancto.\end{Resp}
\begin{Resp}\Rsign Et ambulávit in fortitúdine cibi illíus usque ad montem Dei.\end{Resp}
\switchcolumn
\EN
\begin{Resp}\Rsign Elijah looked, and, behold, there was a cake baken on the coals at his head, and he arose, and did eat and drink; \Star{} And went in the strength of that meat [forty days and forty nights] unto the mount of God.\end{Resp}
\begin{Resp}\Vsign If any man eat of this Bread, he shall live for ever.\end{Resp}
\begin{Resp}\Rsign And went in the strength of that meat [forty days and forty nights] unto the mount of God.\end{Resp}
\begin{Resp}\Vsign Glory be to the Father, and to the Son, \Star{} and to the Holy Ghost.\end{Resp}
\begin{Resp}\Rsign And went in the strength of that meat [forty days and forty nights] unto the mount of God.\end{Resp}
\switchcolumn*

\LA
\LessonHead{Lectio IV}
\switchcolumn
\EN
\LessonHead{Lesson IV}
\switchcolumn*

\LA
\LessonTitle{Sermo sancti Thomæ Aquinátis}
\switchcolumn
\EN
\LessonTitle{From the Sermons of St. Thomas Aquinas}
\switchcolumn*

\LA
\Source{Lectio iv in opusculo 57}
\switchcolumn
\EN
\Source{17th or 57th of his Opuscula, or Lesser Works}
\switchcolumn*

\LA
\noindent Imménsa divínæ largitátis benefícia, exhíbita pópulo christiáno, inæstimábilem ei cónferunt dignitátem. Neque enim est, aut fuit aliquándo tam grandis nátio, quæ hábeat deos appropinquántes sibi, sicut adest nobis Deus noster. Unigénitus síquidem Dei Fílius, suæ divinitátis volens nos esse partícipes, natúram nostram assúmpsit, ut hómines deos fáceret factus homo. Et hoc ínsuper, quod de nostro assúmpsit, totum nobis cóntulit ad salútem. Corpus namque suum pro nostra reconcilatióne in ara crucis hóstiam óbtulit Deo Patri, sánguinem suum fudit in prétium simul et lavácrum; ut redémpti a miserábili servitúte, a peccátis ómnibus mundarémur. Ut autem tanti benefícii jugis in nobis manéret memória, corpus suum in cibum, et sánguinem suum in potum, sub spécie panis et vini suméndum fidélibus derelíquit.\par
\switchcolumn
\EN
\noindent The immeasurable benefits, which the goodness of God hath bestowed on Christian people, have conferred on them also a dignity beyond all price. “For what nation is there so great, who hath gods so nigh unto them, as the Lord, our God, is unto us?” (Deut. iv. 7.) The Only-begotten Son of God, being pleased to make us “partakers of the Divine nature,” (2 Pet. i. 4,) took our nature upon Him, being Himself made Man that He might make men gods. And all, as much of ours as He took, He applied to our salvation. On the Altar of the Cross He offered up His Body to God the Father as a sacrifice for our reconciliation He shed His Blood as the price whereby He redeemeth us from wretchedness and bondage, and the washing whereby He cleanseth us from all sin. And for a noble and abiding memorial of that so great work of His goodness, He hath left unto His faithful ones the Same His very Body for Meat, and the Same His very Blood for Drink, to be fed upon under the appearance of bread and wine.\par
\switchcolumn*

\LA
\begin{Resp}\Rsign Cenántibus illis, accépit Jesus panem, et benedíxit, ac fregit, dedítque discípulis suis, et ait: \Star{} Accípite et comédite: hoc est corpus meum.\end{Resp}
\begin{Resp}\Vsign Dixérunt viri tabernáculi mei: Quis det de cárnibus ejus, ut saturémur?\end{Resp}
\begin{Resp}\Rsign Accípite et comédite: hoc est corpus meum.\end{Resp}
\switchcolumn
\EN
\begin{Resp}\Rsign As they were eating, Jesus took bread, and blest it, and brake it, and gave it to the disciples, and said: \Star{} Take, eat this is My Body.\end{Resp}
\begin{Resp}\Vsign The men of my tabernacle said: O that we had of his flesh, we cannot be satisfied.\end{Resp}
\begin{Resp}\Rsign Take, eat this is My Body.\end{Resp}
\switchcolumn*

\LA
\LessonHead{Lectio V}
\switchcolumn
\EN
\LessonHead{Lesson V}
\switchcolumn*

\LA
\noindent O pretiósum et admirándum convívium, salutíferum et omni suavitáte replétum! Quid enim hoc convívio pretiósius esse potest? in quo non carnes vitulórum et hircórum, ut olim in lege, sed nobis Christus suméndus propónitur verus Deus. Quid hoc Sacraménto mirabílius? In ipso namque panis et vinum in Christi corpus et sánguinem substantiáliter convertúntur; ideóque Christus, Deus et homo perféctus, sub módici panis et vini spécie continétur. Manducátur ítaque a fidélibus, sed mínime lacerátur; quinímmo, divíso Sacraménto, sub quálibet divisiónis partícula ínteger persevérat. Accidéntia autem sine subjécto in eódem subsístunt, ut fides locum hábeat, dum visíbile invisibíliter súmitur aliéna spécie occultátum; et sensus a deceptióne reddántur immúnes, qui de accidéntibus júdicant sibi notis.\par
\switchcolumn
\EN
\noindent How precious a thing then, how marvelous, how health-giving, how furnished with all dainties, is the Supper [of the Lord!] Than His Supper can anything be more precious? Therein there is put before us for meat, not, as of old time, the flesh of bulls and of goats, but Christ Himself, our very God. Than this Sacrament can anything be more marvelous? Therein it cometh to pass that bread and wine are bread and wine no more, but in the stead thereof there is the Body and there is the Blood of Christ; that is to say, Christ Himself, Perfect God and Perfect Man, Christ Himself is there, under the appearance of a little bread and wine. His faithful ones eat Him, but He is not mangled; nay, when [the veil which shroudeth Him in] this Sacrament is broken, in each broken piece thereof remaineth whole Christ Himself, Perfect God and Perfect Man. All that the senses can reach in this Sacrament, [look, taste, feel, smell, and the like, all these] abide of bread and wine, but the Thing is not bread and wine. And thus room is left for faith; Christ Who hath a Form That can be seen, is here taken and received not only unseen, but seeming to be bread and wine, and the senses, which judge by the wonted look, are warranted against error.\par
\switchcolumn*

\LA
\begin{Resp}\Rsign Accépit Jesus cálicem, postquam cenávit, dicens: Hic calix novum testaméntum est in meo sánguine: \Star{} Hoc fácite in meam commemoratiónem.\end{Resp}
\begin{Resp}\Vsign Memória memor ero, et tabéscet in me ánima mea.\end{Resp}
\begin{Resp}\Rsign Hoc fácite in meam commemoratiónem.\end{Resp}
\switchcolumn
\EN
\begin{Resp}\Rsign Jesus took the cup, after supper, saying: This cup is the New Testament in My Blood. \Star{} This do in remembrance of Me.\end{Resp}
\begin{Resp}\Vsign My soul hath them still in remembrance, and is humbled in me.\end{Resp}
\begin{Resp}\Rsign This do in remembrance of Me.\end{Resp}
\switchcolumn*

\LA
\LessonHead{Lectio VI}
\switchcolumn
\EN
\LessonHead{Lesson VI}
\switchcolumn*

\LA
\noindent Nullum étiam sacraméntum est isto salúbrius, quo purgántur peccáta, virtútes augéntur, et mens ómnium spirituálium charísmatum abundántia impinguátur. Offértur in Ecclésia pro vivis et mórtuis, ut ómnibus prosit, quod est pro salúte ómnium institútum. Suavitátem dénique hujus Sacraménti nullus exprímere súfficit, per quod spirituális dulcédo in suo fonte gustátur; et recólitur memória illíus, quam in sua passióne Christus monstrávit, excellentíssimæ caritátis. Unde, ut árctius hujus caritátis imménsitas fidélium córdibus infigerétur, in última cena, quando Pascha cum discípulis celebráto, transitúrus erat de hoc mundo ad Patrem, hoc Sacraméntum instítuit, tamquam passiónis suæ memoriále perénne, figurárum véterum impletívum, miraculórum ab ipso factórum máximum; et de sua contristátis abséntia solátium singuláre relíquit.\par
\switchcolumn
\EN
\noindent Than this Sacrament can anything be more health-giving? Thereby are sins purged away, strength renewed, and the soul fed upon the fatness of spiritual gifts. This Supper is offered up in the Church both for the quick and dead it was ordained to the health of all, all get the good of it. Than this Sacrament can anything be more furnished with dainties The glorious sweetness thereof is of a truth such that no man can fully tell it. Therein ghostly comfort is sucked from its very well-head. Therein a memorial is made of that exceeding great love which Christ showed in time of His sufferings. It was in order that the boundless goodness of that His great love might be driven home into the hearts of His faithful ones, that when He had celebrated the Passover with His disciples, and the last Supper was ended, “the Lord Jesus, knowing that His hour was come that He should depart out of this world unto the Father, having loved His own which were in the world, He loved them unto the end,” (John xiii. 1,) and instituted this Sacrament, this Sacrament, the everlasting “forth-showing of His death until He come again,” (1 Cor. xi. 26,) this Sacrament, the embodied fulfilment of all the ancient types and figures, this Sacrament, the greatest miracle which He ever wrought, and the one mighty joy of them that now have sorrow, till He shall come again, and their heart shall rejoice, and their joy no man take from them. (John xvi. 22.)\par
\switchcolumn*

\LA
\begin{Resp}\Rsign Ego sum panis vitæ; patres vestri manducavérunt manna in desérto, et mórtui sunt: \Star{} Hic est panis de cælo descéndens, ut, si quis ex ipso mandúcet, non moriátur.\end{Resp}
\begin{Resp}\Vsign Ego sum panis vivus, qui de cælo descéndi: si quis manducáverit ex hoc pane, vivet in ætérnum.\end{Resp}
\begin{Resp}\Rsign Hic est panis de cælo descéndens, ut, si quis ex ipso mandúcet, non moriátur.\end{Resp}
\begin{Resp}\Vsign Glória Patri, et Fílio, \Star{} et Spirítui Sancto.\end{Resp}
\begin{Resp}\Rsign Hic est panis de cælo descéndens, ut, si quis ex ipso mandúcet, non moriátur.\end{Resp}
\switchcolumn
\EN
\begin{Resp}\Rsign I am that Bread of life. Your fathers did eat manna in the wilderness, and are dead. \Star{} This is the Bread Which cometh down from heaven, that a man may eat thereof, and not die.\end{Resp}
\begin{Resp}\Vsign I am the living Bread Which came down from heaven if any man eat of this Bread, he shall live for ever.\end{Resp}
\begin{Resp}\Rsign This is the Bread Which cometh down from heaven, that a man may eat thereof, and not die.\end{Resp}
\begin{Resp}\Vsign Glory be to the Father, and to the Son, \Star{} and to the Holy Ghost.\end{Resp}
\begin{Resp}\Rsign This is the Bread Which cometh down from heaven, that a man may eat thereof, and not die.\end{Resp}
\switchcolumn*

\LA
\LessonHead{Lectio VII}
\switchcolumn
\EN
\LessonHead{Lesson VII}
\switchcolumn*

\LA
\LessonTitle{Léctio sancti Evangélii secúndum Joánnem}
\switchcolumn
\EN
\LessonTitle{From the Holy Gospel according to John}
\switchcolumn*

\LA
\Cite{Joannes 6:56-59}
\switchcolumn
\EN
\Cite{John 6:56-59}
\switchcolumn*

\LA
\noindent In illo témpore: Dixit Jesus turbis Judæórum: Caro mea vere est cibus, et sanguis meus vere est potus. Et réliqua.\par
\switchcolumn
\EN
\noindent At that time: Jesus said unto the multitudes of the Jews: My Flesh is meat indeed, and My Blood is drink indeed. And so on.\par
\switchcolumn*

\LA
\LessonTitle{Homilía sancti Augustíni Epíscopi}
\switchcolumn
\EN
\LessonTitle{Homily by St. Augustine, Bishop (of Hippo)}
\switchcolumn*

\LA
\Source{Tractatus 26 in Joannem, sub finem}
\switchcolumn
\EN
\Source{26th Tract on John}
\switchcolumn*

\LA
\noindent Cum cibo et potu id áppetant hómines, ut neque esúriant, neque sítiant: hoc veráciter non præstat, nisi iste cibus et potus, qui eos, a quibus súmitur, immortáles et incorruptíbiles facit; id est, socíetas ipsa Sanctórum, ubi pax erit et únitas plena atque perfécta. Proptérea quippe, sicut étiam ante nos hoc intellexérunt hómines Dei, Dóminus noster Jesus Christus corpus et sánguinem suum in eis rebus commendávit, quæ ad unum áliquid redigúntur ex multis. Namque áliud in unum ex multis granis confícitur: áliud in unum ex multis ácinis cónfluit. Dénique jam expónit, quómodo id fiat, quod lóquitur; et quid sit manducáre corpus ejus, et sánguinem bíbere.\par
\switchcolumn
\EN
\noindent By use of meat and drink men would fain that “they shall hunger no more, neither thirst any more,” (Apoc. vii. 16,) and yet there is but one Meat and one Drink, Which doth work in them that feed thereon that “this corruptible must put on incorruption, and this mortal put on immortality,” (1 Cor. xv. 53,) namely communion with that general assembly and Church of God's holy children, who are “kept in perfect peace,” (Isa. xxvi. 3,) and are “all one,” (John xvii. 11,) fully and utterly. And therefore it is, as men of God before our time have taken it, that our Lord Jesus Christ hath set before us His Body and His Blood in the likeness of things which, from being many, are reduced into one. In one loaf are many grains of corn, and one cup of wine the juice of many grapes. And now He giveth us to know how that which He spake cometh to pass, and how indeed “this Man can give us His Flesh to eat,” and His Blood to drink.\par
\switchcolumn*

\LA
\begin{Resp}\Rsign Qui mandúcat meam carnem et bibit meum sánguinem, \Star{} In me manet, et ego in eo.\end{Resp}
\begin{Resp}\Vsign Non est ália nátio tam grandis, quæ hábeat deos appropinquántes sibi, sicut Deus noster adest nobis.\end{Resp}
\begin{Resp}\Rsign In me manet, et ego in eo.\end{Resp}
\switchcolumn
\EN
\begin{Resp}\Rsign He that eateth My Flesh and drinketh My Blood, \Star{} Dwelleth in Me, and I in him.\end{Resp}
\begin{Resp}\Vsign What nation is there so great, who hath gods so nigh unto them, as the Lord our God is to us\end{Resp}
\begin{Resp}\Rsign Dwelleth in Me, and I in him.\end{Resp}
\switchcolumn*

\LA
\LessonHead{Lectio VIII}
\switchcolumn
\EN
\LessonHead{Lesson VIII}
\switchcolumn*

\LA

\switchcolumn
\EN
\LessonTitle{“He that eateth My Flesh, and drinketh My Blood, dwelleth in Me, and I in him.” To dwell in Christ, therefore, and to have Him dwelling in us, is to “eat of that Bread and drink of that Cup,” (1 Cor. xi. 28,) and he which dwelleth not in Christ, and in whom Christ dwelleth not, without all doubt doth not spiritually eat His Flesh nor drink His Blood, although he do carnally and visibly press the Sacrament with his teeth but, contrariwise, he “eateth and drinketh damnation to himself,” because he dareth to draw nigh filthy to that secret and holy thing of Christ, whereunto none draweth nigh worthily, save he which is pure, even he which is of them concerning whom it is said: “Blessed are the pure in heart, for they shall see God.” (Matth. v. 8.)}
\switchcolumn*

\LA
\noindent Qui mandúcat carnem meam et bibit meum sánguinem, in me manet, et ego in illo. Hoc est ergo manducáre illam escam, et illum bíbere potum, in Christo manére, et illum manéntem in se habére. Ac per hoc, qui non manet in Christo, et in quo non manet Christus, proculdúbio nec mandúcat spiritáliter carnem ejus, nec bibit ejus sánguinem, licet carnáliter et visibíliter premat déntibus Sacraméntum córporis et sánguinis Christi: sed magis tantæ rei sacraméntum ad judícium sibi mandúcat et bibit, quia immúndus præsúmpsit ad Christi accédere Sacraménta, quæ áliquis non digne sumit, nisi qui mundus est; de quibus dícitur: Beáti mundo corde, quóniam ipsi Deum vidébunt.\par
\switchcolumn
\EN

\switchcolumn*

\LA
\begin{Resp}\Rsign Misit me vivens Pater, et ego vivo propter Patrem: \Star{} Et qui mandúcat me, vivet propter me.\end{Resp}
\begin{Resp}\Vsign Cibávit illum Dóminus pane vitæ et intelléctus.\end{Resp}
\begin{Resp}\Rsign Et qui mandúcat me, vivet propter me.\end{Resp}
\begin{Resp}\Vsign Glória Patri, et Fílio, \Star{} et Spirítui Sancto.\end{Resp}
\begin{Resp}\Rsign Et qui mandúcat me, vivet propter me.\end{Resp}
\switchcolumn
\EN
\begin{Resp}\Rsign As the living Father hath sent Me, and I live by the Father, \Star{} So he that eateth Me, even he shall live by Me.\end{Resp}
\begin{Resp}\Vsign With the bread of life and understanding hath the Lord fed him.\end{Resp}
\begin{Resp}\Rsign So he that eateth Me, even he shall live by Me.\end{Resp}
\begin{Resp}\Vsign Glory be to the Father, and to the Son, \Star{} and to the Holy Ghost.\end{Resp}
\begin{Resp}\Rsign So he that eateth Me, even he shall live by Me.\end{Resp}
\switchcolumn*

\LA
\LessonHead{Lectio IX}
\switchcolumn
\EN
\LessonHead{Lesson IX}
\switchcolumn*

\LA

\switchcolumn
\EN
\LessonTitle{“As the living Father hath sent Me, and I live by the Father, so he that eateth Me, even he shall live by Me.” This is as though He said: The Father hath sent Me into the world (John x. 36,) and I have emptied Myself [and taken upon Me the form of a servant, and being found in fashion as a man], (Phil. ii. 7, 8.) I have My life from the Father, as One That is greater than I. (John xiv. 28.) He that eateth Me, even he, by thereby taking part in Me, shall live by Me. It is as having humbled Myself (Phil. ii. 8) that I live by the Father, but he that eateth Me, him will I raise up, (John vi. 55,) and so he shall live by Me. It is said “I live by the Father” that is to say, He is of the Father, not the Father of Him, and yet not so, but that the Father and the Son are co-equal together. Also it is said “So he that eateth Me, even he shall live by Me,” whereby He showeth the gracious work towards His people of Him Who is the “one Mediator between God and man,” (1 Tim. ii. 5,) and not that He Which is eaten and he which eateth Him are co-equal together.}
\switchcolumn*

\LA
\noindent Sicut, inquit, misit me vivens Pater, et ego vivo propter Patrem: et qui mandúcat me, et ipse vivet propter me. Ac si díceret: Ut ego vivam propter Patrem, id est, ad illum tamquam ad majórem réferam vitam meam, exinanítio mea fecit, in qua me misit: ut autem quisquam vivat propter me, participátio facit, qua mandúcat me. Ego ítaque humiliátus vivo propter Patrem: ille eréctus vivit propter me. Si autem ita dictus est, Vivo propter Patrem, quia ipse de illo, non ille de ipso est; sine detriménto æqualitátis dictum est. Nec tamen dicéndo, Et qui mandúcat me, et ipse vivet propter me; eándem suam et nostram æqualitátem significávit, sed grátiam mediatóris osténdit.\par
\switchcolumn
\EN

\switchcolumn*

\LA
\TeDeum{Te Deum}
\switchcolumn
\EN
\TeDeum{Te Deum}
\switchcolumn*

\end{paracol}

\Office{Feria Sexta infra octavam Corporis Christi}{Friday within the Octave of Corpus Christi}{Semiduplex II classis}
\addcontentsline{toc}{subsection}{Feria Sexta infra octavam Corporis Christi \textit{— Friday within the Octave of Corpus Christi}}
\begin{paracol}{2}
\LA
\LessonHead{Lectio I}
\switchcolumn
\EN
\LessonHead{Lesson I}
\switchcolumn*

\LA
\LessonTitle{De libro primo Regum}
\switchcolumn
\EN
\LessonTitle{Lesson from the first book of Samuel}
\switchcolumn*

\LA
\Cite{1 Reg 2:27-29}
\switchcolumn
\EN
\Cite{1 Sam 2:27-29}
\switchcolumn*

\LA
\noindent Venit autem vir Dei ad Heli, et ait ad eum: Hæc dicit Dóminus: Numquid non apérte revelátus sum dómui patris tui, cum essent in Ægýpto in domo Pharaónis? Et elégi eum ex ómnibus tríbubus Israël mihi in sacerdótem, ut ascénderet ad altáre meum, et adoléret mihi incénsum, et portáret ephod coram me: et dedi dómui patris tui ómnia de sacrifíciis filiórum Israël. Quare calce abiecístis víctimam meam, et múnera mea quæ præcépi ut offerréntur in templo: et magis honorásti fílios tuos quam me, ut comederétis primítias omnis sacrifícii Israël pópuli mei?\par
\switchcolumn
\EN
\noindent And there came a man of God to Heli, and said to him: Thus saith the Lord: Did I not plainly appear to thy father's house, when they were in Egypt in the house of Pharao? And I chose him out of all the tribes of Israel to be my priest, to go up to my altar, and burn incense to me, and to wear the ephod before me: and I gave to thy father's house of all the sacrifices of the children of Israel. Why have you kicked away my victims, and my gifts which I commanded to be offered in the temple: and thou hast rather honoured thy sons than me, to eat the firstfruits of every sacrifice of my people Israel?\par
\switchcolumn*

\LA
\begin{Resp}\Rsign Immolábit hædum multitúdo filiórum Israël ad vésperam Paschæ: \Star{} Et edent carnes et ázymos panes.\end{Resp}
\begin{Resp}\Vsign Pascha nostrum immolátus est Christus: ítaque epulémur in ázymis sinceritátis et veritátis.\end{Resp}
\begin{Resp}\Rsign Et edent carnes et ázymos panes.\end{Resp}
\switchcolumn
\EN
\begin{Resp}\Rsign The whole assembly of the children of Israel shall kill the lamb toward the evening of the Passover. \Star{} And they shall eat the flesh, and unleavened bread.\end{Resp}
\begin{Resp}\Vsign Even Christ our Passover is sacrificed for us therefore let us keep the feast with the unleavened bread of sincerity and truth.\end{Resp}
\begin{Resp}\Rsign And they shall eat the flesh, and unleavened bread.\end{Resp}
\switchcolumn*

\LA
\LessonHead{Lectio II}
\switchcolumn
\EN
\LessonHead{Lesson II}
\switchcolumn*

\LA
\Cite{1 Reg 2:30-33}
\switchcolumn
\EN
\Cite{1 Sam 2:30-33}
\switchcolumn*

\LA
\noindent Proptérea ait Dóminus Deus Israël: Loquens locútus sum, ut domus tua, et domus patris tui, ministráret in conspéctu meo usque in sempitérnum. Nunc autem dicit Dóminus: Absit hoc a me: sed quicúmque glorificáverit me, glorificábo eum: qui autem contémnunt me, erunt ignóbiles. Ecce dies véniunt, et præcídam brácchium tuum, et brácchium domus patris tui, ut non sit senex in domo tua. Et vidébis ǽmulum tuum in templo, in univérsis prósperis Israël: et non erit senex in domo tua ómnibus diébus. Verúmtamen non áuferam pénitus virum ex te ab altári meo: sed ut defíciant óculi tui, et tabéscat ánima tua: et pars magna domus tuæ moriétur cum ad virílem ætátem vénerit.\par
\switchcolumn
\EN
\noindent Wherefore thus saith the Lord the God of Israel: I said indeed that thy house, and the house of thy father should minister in my sight, for ever. But now saith the Lord: Far be this from me: but whosoever shall glorify me, him will I glorify: but they that despise me, shall be despised. Behold the days come: and I will cut off thy arm, and the arm of thy father's house, that there shall not be an old man in thy house. And thou shalt see thy rival in the temple, in all the prosperity of Israel, and there shall not be an old man in thy house for ever. However I will not altogether take away a man of thee from my altar: but that thy eyes may faint and thy soul be spent: and a great part of thy house shall die when they come to man's estate\par
\switchcolumn*

\LA
\begin{Resp}\Rsign Comedétis carnes, et saturabímini pánibus: \Star{} Iste est panis, quem dedit vobis Dóminus ad vescéndum.\end{Resp}
\begin{Resp}\Vsign Non Móyses dedit vobis panem de cælo, sed Pater meus dat vobis panem de cælo verum.\end{Resp}
\begin{Resp}\Rsign Iste est panis, quem dedit vobis Dóminus ad vescéndum.\end{Resp}
\switchcolumn
\EN
\begin{Resp}\Rsign Ye shall eat flesh, and shall be filled with bread. \Star{} This is the bread which the Lord hath given you to eat.\end{Resp}
\begin{Resp}\Vsign Moses gave you not that Bread from heaven, but My Father giveth you the true Bread from heaven.\end{Resp}
\begin{Resp}\Rsign This is the bread which the Lord hath given you to eat.\end{Resp}
\switchcolumn*

\LA
\LessonHead{Lectio III}
\switchcolumn
\EN
\LessonHead{Lesson III}
\switchcolumn*

\LA
\Cite{1 Reg 2:34-36}
\switchcolumn
\EN
\Cite{1 Sam 2:34-36}
\switchcolumn*

\LA
\noindent Hoc autem erit tibi signum, quod ventúrum est duóbus fíliis tuis, Ophni et Phínees: in die uno moriéntur ambo. Et suscitábo mihi sacerdótem fidélem, qui juxta cor meum et ánimam meam fáciet: et ædificábo ei domum fidélem, et ambulábit coram christo meo cunctis diébus. Futúrum est autem, ut quicúmque remánserit in domo tua, véniat ut orétur pro eo, et ófferat nummum argénteum, et tortam panis, dicátque: Dimítte me, óbsecro, ad unam partem sacerdotálem, ut cómedam buccéllam panis.\par
\switchcolumn
\EN
\noindent And this shall be a sign to thee, that shall come upon thy two sons, Ophni and Phinees: In one day they shall both of them die. And I will raise me up a faithful priest, who shall do according to my heart, and my soul, and I will build him a faithful house, and he shall walk all days before my anointed. And it shall come to pass, that whosoever shall remain in thy house, shall come that he may be prayed for, and shall offer a piece of silver, and a roll of bread, and shall say: Put me, I beseech thee, to somewhat of the priestly office, that I may eat a morsel of bread.\par
\switchcolumn*

\LA
\begin{Resp}\Rsign Respéxit Elías ad caput suum subcinerícium panem: qui surgens comédit et bibit: \Star{} Et ambulávit in fortitúdine cibi illíus usque ad montem Dei.\end{Resp}
\begin{Resp}\Vsign Si quis manducáverit ex hoc pane, vivet in ætérnum.\end{Resp}
\begin{Resp}\Rsign Et ambulávit in fortitúdine cibi illíus usque ad montem Dei.\end{Resp}
\begin{Resp}\Vsign Glória Patri, et Fílio, \Star{} et Spirítui Sancto.\end{Resp}
\begin{Resp}\Rsign Et ambulávit in fortitúdine cibi illíus usque ad montem Dei.\end{Resp}
\switchcolumn
\EN
\begin{Resp}\Rsign Elijah looked, and, behold, there was a cake baken on the coals at his head, and he arose, and did eat and drink; \Star{} And went in the strength of that meat [forty days and forty nights] unto the mount of God.\end{Resp}
\begin{Resp}\Vsign If any man eat of this Bread, he shall live for ever.\end{Resp}
\begin{Resp}\Rsign And went in the strength of that meat [forty days and forty nights] unto the mount of God.\end{Resp}
\begin{Resp}\Vsign Glory be to the Father, and to the Son, \Star{} and to the Holy Ghost.\end{Resp}
\begin{Resp}\Rsign And went in the strength of that meat [forty days and forty nights] unto the mount of God.\end{Resp}
\switchcolumn*

\LA
\LessonHead{Lectio IV}
\switchcolumn
\EN
\LessonHead{Lesson IV}
\switchcolumn*

\LA
\LessonTitle{De Sermóne sancti Thomæ Aquinátis}
\switchcolumn
\EN
\LessonTitle{From the Sermons of St. Thomas Aquinas.}
\switchcolumn*

\LA
\Source{In eodem Opusculo 57.}
\switchcolumn
\EN
\Source{Lesser works, 57}
\switchcolumn*

\LA
\noindent Cónvenit ítaque devotióni fidélium, solémniter recólere institutiónem tam salutíferi tamque mirábilis Sacraménti: ut ineffábilem modum divínæ præséntiæ in Sacraménto visíbili venerémur; et laudétur Dei poténtia, quæ in Sacraménto eódem tot mirabília operátur; nec non et de tam salúbri tamque suávi benefício exsolvántur Deo gratiárum débitæ actiónes. Verum etsi in die Cœnæ, quando Sacraméntum prædíctum nóscitur institútum solémnia de institutióne ipsíus speciális méntio habeátur; totum tamen resíduum ejúsdem diéi offícium ad Christi passiónem pértinet, circa cujus veneratiónem Ecclésia illo témpore occupátur.\par
\switchcolumn
\EN
\noindent It serveth well therefore to the edifying of the faithful to make memorial of the institution of this so health-giving and so wonderful a Sacrament, that we may worship the unspeakable way by the which the Divine Presence in-dwelleth in this Sacrament, Which we see, and may praise the power of God whereby in this Sacrament are wrought so many wonders, yea, and also give God some of those thanks which we owe unto Him for this so health-giving gift of His loving-kindness. It is true that on the Day of the Supper of the Lord, on which day we know it to have been that He ordained this Sacrament, at the solemn celebration of the Mass, the memory of the instituting thereof is more particularly mentioned, but all the rest of the Services on that day deal chiefly with Christ Suffering, to worshipping of Whom the Church doth at that season give all her mind.\par
\switchcolumn*

\LA
\begin{Resp}\Rsign Cenántibus illis, accépit Jesus panem, et benedíxit, ac fregit, dedítque discípulis suis, et ait: \Star{} Accípite et comédite: hoc est corpus meum.\end{Resp}
\begin{Resp}\Vsign Dixérunt viri tabernáculi mei: Quis det de cárnibus ejus, ut saturémur?\end{Resp}
\begin{Resp}\Rsign Accípite et comédite: hoc est corpus meum.\end{Resp}
\switchcolumn
\EN
\begin{Resp}\Rsign As they were eating, Jesus took bread, and blest it, and brake it, and gave it to the disciples, and said: \Star{} Take, eat this is My Body.\end{Resp}
\begin{Resp}\Vsign The men of my tabernacle said: O that we had of his flesh, we cannot be satisfied.\end{Resp}
\begin{Resp}\Rsign Take, eat this is My Body.\end{Resp}
\switchcolumn*

\LA
\LessonHead{Lectio V}
\switchcolumn
\EN
\LessonHead{Lesson V}
\switchcolumn*

\LA
\noindent Ut autem íntegro celebritátis offício institutiónem tanti Sacraménti recóleret plebs fidélium, Románus Póntifex Urbánus quartus, hujus Sacraménti devotióne afféctus, pie státuit præfátæ institutiónis memóriam prima quinta féria post Octávam Pentecóstes a cunctis fidélibus celebrári; ut qui per totum anni círculum hoc Sacraménto útimur ad salútem, ejus institutiónem illo témpore speciáliter recolámus, quo Spíritus Sanctus corda discipulórum edócuit ad plene cognoscénda hujus mystéria Sacraménti. Nam et in eódem témpore cœpit hoc Sacraméntum a fidélibus frequentári.\par
\switchcolumn
\EN
\noindent But, that the congregation of Christ's faithful people might celebrate with an whole Festal Office all to itself the institution of this so great Sacrament, Urban IV., Bishop of Rome, being touched with love toward this said Sacrament, hath made a godly ordinance that the memory of the said institution should be celebrated by all the faithful upon the Fifth Day of the week, next after the Eighth Day of Pentecost. From one end of the year to the other we use this Sacrament to our souls' health, and we more particularly celebrate the institution thereof at that season wherein the Holy Ghost taught the hearts of the disciples to acknowledge the mysteries thereof, for then it was, as we read, that “they continued steadfastly in [the Apostles' doctrine and fellowship, and in] breaking of Bread, [and in prayers.”] (Acts ii. 42.)\par
\switchcolumn*

\LA
\begin{Resp}\Rsign Accépit Jesus cálicem, postquam cenávit, dicens: Hic calix novum testaméntum est in meo sánguine: \Star{} Hoc fácite in meam commemoratiónem.\end{Resp}
\begin{Resp}\Vsign Memória memor ero, et tabéscet in me ánima mea.\end{Resp}
\begin{Resp}\Rsign Hoc fácite in meam commemoratiónem.\end{Resp}
\switchcolumn
\EN
\begin{Resp}\Rsign Jesus took the cup, after supper, saying: This cup is the New Testament in My Blood. \Star{} This do in remembrance of Me.\end{Resp}
\begin{Resp}\Vsign My soul hath them still in remembrance, and is humbled in me.\end{Resp}
\begin{Resp}\Rsign This do in remembrance of Me.\end{Resp}
\switchcolumn*

\LA
\LessonHead{Lectio VI}
\switchcolumn
\EN
\LessonHead{Lesson VI}
\switchcolumn*

\LA
\noindent Ut autem prædícta quinta féria, et per Octávas sequéntes, ejus salutáris institutiónis honorificéntius agátur memória, et solémnitas de hoc celébrior habeátur: loco distributiónum materiálium, quæ in ecclésiis cathedrálibus largiúntur exsisténtibus Horis Canónicis, noctúrnis paritérque diúrnis, præfátus Románus Póntifex eis, qui hujúsmodi Horis in hac solemnitáte personáliter in ecclésiis interfúerint, stipéndia spirituália apostólica largitióne concéssit; quátenus per hæc fidéles ad tanti festi celebritátem avídius et copiósius convenírent.\par
\switchcolumn
\EN
\noindent And, moreover, to the end that on the aforesaid Thursday, and the seven days next following, the memory of this health-giving Institution might be the more honourably celebrated, and the Feast thereby be held in more excellent worship, the abovenamed Bishop of Rome, after the manner of the doles which in Cathedral Churches are given to such as come to the singing or saying of the Canonical Hours by night and day, hath out of his Apostolic bounty granted spiritual rewards to all such as in their own persons are present in the Church at the diverse Canonical Hours during all this Festival, thereby to stir up the faithful to come to the keeping of this great Feast in greater eagerness and numbers.\par
\switchcolumn*

\LA
\begin{Resp}\Rsign Ego sum panis vitæ; patres vestri manducavérunt manna in desérto, et mórtui sunt: \Star{} Hic est panis de cælo descéndens, ut, si quis ex ipso mandúcet, non moriátur.\end{Resp}
\begin{Resp}\Vsign Ego sum panis vivus, qui de cælo descéndi: si quis manducáverit ex hoc pane, vivet in ætérnum.\end{Resp}
\begin{Resp}\Rsign Hic est panis de cælo descéndens, ut, si quis ex ipso mandúcet, non moriátur.\end{Resp}
\begin{Resp}\Vsign Glória Patri, et Fílio, \Star{} et Spirítui Sancto.\end{Resp}
\begin{Resp}\Rsign Hic est panis de cælo descéndens, ut, si quis ex ipso mandúcet, non moriátur.\end{Resp}
\switchcolumn
\EN
\begin{Resp}\Rsign I am that Bread of life. Your fathers did eat manna in the wilderness, and are dead. \Star{} This is the Bread Which cometh down from heaven, that a man may eat thereof, and not die.\end{Resp}
\begin{Resp}\Vsign I am the living Bread Which came down from heaven if any man eat of this Bread, he shall live for ever.\end{Resp}
\begin{Resp}\Rsign This is the Bread Which cometh down from heaven, that a man may eat thereof, and not die.\end{Resp}
\begin{Resp}\Vsign Glory be to the Father, and to the Son, \Star{} and to the Holy Ghost.\end{Resp}
\begin{Resp}\Rsign This is the Bread Which cometh down from heaven, that a man may eat thereof, and not die.\end{Resp}
\switchcolumn*

\LA
\LessonHead{Lectio VII}
\switchcolumn
\EN
\LessonHead{Lesson VII}
\switchcolumn*

\LA
\LessonTitle{Léctio sancti Evangélii secúndum Joánnem}
\switchcolumn
\EN
\LessonTitle{From the Holy Gospel according to John}
\switchcolumn*

\LA
\Cite{Joannes 6:56-59}
\switchcolumn
\EN
\Cite{John 6:56-59}
\switchcolumn*

\LA
\noindent In illo témpore: Dixit Jesus turbis Judæórum: Caro mea vere est cibus, et sanguis meus vere est potus. Et réliqua.\par
\switchcolumn
\EN
\noindent At that time, Jesus said unto the multitudes of the Jews: My Flesh is meat indeed and My Blood is drink indeed. And so on.\par
\switchcolumn*

\LA
\LessonTitle{De Homilía sancti Augustíni Epíscopi}
\switchcolumn
\EN
\LessonTitle{Homily by St. Augustine, Bishop of Hippo}
\switchcolumn*

\LA
\Source{Tractatus 27 in Joánnem}
\switchcolumn
\EN
\Source{27th Tract on John.}
\switchcolumn*

\LA
\noindent Verba Dómini ex Evangélio, quæ sermónem prístinum consequúntur, audívimus. Hinc sermo debétur áuribus et méntibus vestris, et hodiérna die non importúnus est: est enim de córpore Dómini, quod dicébat se dare ad manducándum propter ætérnam vitam. Expósuit autem modum attributiónis hujus et doni sui, quómodo daret carnem suam manducáre, dicens: Qui mandúcat carnem meam et bibit sánguinem meum, in me manet, et ego in illo. Signum quia manducávit et bibit, hoc est: si manet, et manétur; si hábitat, et inhabitátur; si hæret, ut non deserátur.\par
\switchcolumn
\EN
\noindent We have heard from the Gospel the words of the Lord which follow [those that formed the subject of my earlier discourse. To your ears and understandings we owe a discourse on these also, and to-day it becometh very well, for it is upon that Body of the Lord, Which He professeth Himself that He “will give for the life of the world,” “that a man may eat thereof and not die.” He hath made manifest how He giveth, and What is His Gift, where He saith: “He that eateth My Flesh and drinketh My Blood, dwelleth in Me and I in him.” The sign to show whether a man hath or hath not eaten that Flesh and drunk that Blood, is whether or not he dwelleth in Christ and Christ in him, whether or not he is a guest of Christ and Christ of his, whether or not he so cleaveth unto Christ, that Christ be not parted from him.\par
\switchcolumn*

\LA
\begin{Resp}\Rsign Qui mandúcat meam carnem et bibit meum sánguinem, \Star{} In me manet, et ego in eo.\end{Resp}
\begin{Resp}\Vsign Non est ália nátio tam grandis, quæ hábeat deos appropinquántes sibi, sicut Deus noster adest nobis.\end{Resp}
\begin{Resp}\Rsign In me manet, et ego in eo.\end{Resp}
\switchcolumn
\EN
\begin{Resp}\Rsign He that eateth My Flesh and drinketh My Blood, \Star{} Dwelleth in Me, and I in him.\end{Resp}
\begin{Resp}\Vsign What nation is there so great, who hath gods so nigh unto them, as the Lord our God is to us\end{Resp}
\begin{Resp}\Rsign Dwelleth in Me, and I in him.\end{Resp}
\switchcolumn*

\LA
\LessonHead{Lectio VIII}
\switchcolumn
\EN
\LessonHead{Lesson VIII}
\switchcolumn*

\LA
\noindent Hoc ergo nos dócuit et admónuit mýsticis verbis, ut simus in ejus córpore sub ipso cápite in membris ejus, edéntes carnem ejus, non relinquéntes unitátem ejus. Sed qui áderant, plures non intelligéndo scandalizáti sunt; non enim cogitábant, hæc audiéndo, nisi carnem, quod ipsi erant. Apóstolus autem dicit, et verum dicit: Sápere secúndum carnem, mors est. Carnem suam dat nobis Dóminus manducáre: et sápere secúndum carnem, mors est. Cum de carne sua dicat quia ibi est vita ætérna; ergo nec carnem debémus sápere secúndum carnem, sicut in his verbis: Multi ítaque audiéntes, non ex inimícis, sed ex discípulis ejus, dixérunt: Durus est hic sermo, et quis potest eum audíre?\par
\switchcolumn
\EN
\noindent This hath He taught, and warned us, by words of deep meaning, to be in His Body as members whose Head is He, eating His Flesh, and cleaving alway to His Oneness.” Many of His disciples when they had heard this.. went back, and walked no more with Him “for they understood not by Flesh any flesh other than such as they themselves were made of. The Apostle saith, (and very true it is) “To be carnally minded is death.” (Rom. viii. 6.) The Lord giveth us His Flesh to eat, and to understand it carnally is death. Where He saith: “Whoso eateth My Flesh is life” we. must not eat His Flesh carnally, whom it is written “Many of His disciples” (not His enemies) “when they heard this, said This is an hard saying who can hear it?”\par
\switchcolumn*

\LA
\begin{Resp}\Rsign Misit me vivens Pater, et ego vivo propter Patrem: \Star{} Et qui mandúcat me, vivet propter me.\end{Resp}
\begin{Resp}\Vsign Cibávit illum Dóminus pane vitæ et intelléctus.\end{Resp}
\begin{Resp}\Rsign Et qui mandúcat me, vivet propter me.\end{Resp}
\begin{Resp}\Vsign Glória Patri, et Fílio, \Star{} et Spirítui Sancto.\end{Resp}
\begin{Resp}\Rsign Et qui mandúcat me, vivet propter me.\end{Resp}
\switchcolumn
\EN
\begin{Resp}\Rsign As the living Father hath sent Me, and I live by the Father, \Star{} So he that eateth Me, even he shall live by Me.\end{Resp}
\begin{Resp}\Vsign With the bread of life and understanding hath the Lord fed him.\end{Resp}
\begin{Resp}\Rsign So he that eateth Me, even he shall live by Me.\end{Resp}
\begin{Resp}\Vsign Glory be to the Father, and to the Son, \Star{} and to the Holy Ghost.\end{Resp}
\begin{Resp}\Rsign So he that eateth Me, even he shall live by Me.\end{Resp}
\switchcolumn*

\LA
\LessonHead{Lectio IX}
\switchcolumn
\EN
\LessonHead{Lesson IX}
\switchcolumn*

\LA
\noindent Si discípuli durum habuérunt istum sermónem, quid inimíci? Et tamen sic oportébat ut dicerétur, quod non ab ómnibus intelligerétur. Secrétum Dei inténtos debet fácere, non advérsos: isti autem cito defecérunt, tália loquénte Dómino Jesu Christo. Non credidérunt áliquid magnum dicéntem, et verbis illis áliquam grátiam cooperiéntem: sed prout voluérunt, ita intellexérunt, et more hóminum: quia póterat Jesus, aut hoc disponébat Jesus, carnem, qua indútum erat Verbum, véluti concísam, distribúere credéntibus in se. Durus est, ínquiunt, hic sermo: quis potest eum audíre?\par
\switchcolumn
\EN
\noindent If His disciples took His words for an hard saying, how did His enemies take them? And, nevertheless, thus it behoved to speak them if all men were not to understand them. A Divine mystery ought to make us thoughtful, not to repel us and yet, when the Lord Jesus Christ spoke thus in mystery, many of His disciples went back and walked no more with Him. They believed not that He was speaking of some great thing, and darkly announcing in these hath eternal understand this as did they of “Many of His words a bounty. They understood but as they pleased, even after the manner of men, that Jesus was able, or that Jesus meant, to give that Flesh wherewith the Word is clothed on, as it were, in slices, to them that believe on Him. And they said “This is an hard saying who can hear it?”\par
\switchcolumn*

\LA
\TeDeum{Te Deum}
\switchcolumn
\EN
\TeDeum{Te Deum}
\switchcolumn*

\end{paracol}

\Office{Sabbato infra octavam Corporis Christi}{Saturday within the Octave of Corpus Christi}{Semiduplex II classis}
\addcontentsline{toc}{subsection}{Sabbato infra octavam Corporis Christi \textit{— Saturday within the Octave of Corpus Christi}}
\begin{paracol}{2}
\LA
\LessonHead{Lectio I}
\switchcolumn
\EN
\LessonHead{Lesson I}
\switchcolumn*

\LA
\LessonTitle{De libro primo Regum}
\switchcolumn
\EN
\LessonTitle{Lesson from the first book of Samuel}
\switchcolumn*

\LA
\Cite{1 Reg 3:1-7}
\switchcolumn
\EN
\Cite{1 Sam 3:1-7}
\switchcolumn*

\LA
\noindent Puer autem Sámuel ministrábat Dómino coram Heli, et sermo Dómini erat pretiósus in diébus illis: non erat vísio manifésta. Factus est ergo in quadam die, Heli jacébat in loco suo, et óculi ejus caligáverant, nec póterant vidére. Lucérna Dei ántequam exstinguerétur, Sámuel dormiébat in templo Dómini, ubi erat arca Dei. Et vocávit Dóminus Sámuel, qui respóndens ait: Ecce ego. Et cucúrrit ad Heli et dixit: Non vocávi: revértere et dormi. Et ábiit et dormívit. Et adjécit Dóminus rursum vocáre Samuélem. Consurgénsque Sámuel ábiit ad Heli, et dixit: Ecce ego, quia vocásti me. Qui respóndit: Non vocávi te, fili mi: revértere et dormi. Porro Sámuel necdum sciébat Dóminum, neque revelátus fúerat ei sermo Dómini.\par
\switchcolumn
\EN
\noindent Now the child Samuel ministered to the Lord before Heli, and the word of the Lord was precious in those days, there was no manifest vision. And it came to pass one day when Heli lay in his place, and his eyes were grown dim, that he could not see: Before the lamp of God went out, Samuel slept in the temple of the Lord, where the ark of God was. And the Lord called Samuel. And he answered: Here am I. And he ran to Heli and said: Here am I: for thou didst call me. He said: I did not call: go back and sleep. And he went and slept. And the Lord called Samuel again. And Samuel arose and went to Heli, and said: Here am I: for thou calledst me. He answered: I did not call thee, my son: return and sleep. Now Samuel did not yet know the Lord, neither had the word of the Lord been revealed to him.\par
\switchcolumn*

\LA
\begin{Resp}\Rsign Immolábit hædum multitúdo filiórum Israël ad vésperam Paschæ: \Star{} Et edent carnes et ázymos panes.\end{Resp}
\begin{Resp}\Vsign Pascha nostrum immolátus est Christus: ítaque epulémur in ázymis sinceritátis et veritátis.\end{Resp}
\begin{Resp}\Rsign Et edent carnes et ázymos panes.\end{Resp}
\switchcolumn
\EN
\begin{Resp}\Rsign The whole assembly of the children of Israel shall kill the lamb toward the evening of the Passover. \Star{} And they shall eat the flesh, and unleavened bread.\end{Resp}
\begin{Resp}\Vsign Even Christ our Passover is sacrificed for us therefore let us keep the feast with the unleavened bread of sincerity and truth.\end{Resp}
\begin{Resp}\Rsign And they shall eat the flesh, and unleavened bread.\end{Resp}
\switchcolumn*

\LA
\LessonHead{Lectio II}
\switchcolumn
\EN
\LessonHead{Lesson II}
\switchcolumn*

\LA
\Cite{1 Reg 3:8-12}
\switchcolumn
\EN
\Cite{1 Sam 3:8-12}
\switchcolumn*

\LA
\noindent Et adjécit Dóminus, et vocávit adhuc Samuélem tértio. Qui consúrgens ábiit ad Heli, et ait: Ecce ego, quia vocásti me. Intelléxit ergo Heli quia Dóminus vocáret púerum: et ait ad Samuélem: Vade, et dormi: et si deínceps vocáverit te, dices: Lóquere, Dómine, quia audit servus tuus. Abiit ergo Sámuel, et dormívit in loco suo. Et venit Dóminus, et stetit: et vocávit, sicut vocáverat secúndo: Sámuel, Sámuel. Et ait Sámuel: Lóquere, Dómine, quia audit servus tuus. Et dixit Dóminus ad Samuélem: Ecce ego fácio verbum in Israël, quod quicúmque audíerit, tínnient ambæ aures ejus. In die illa suscitábo advérsum Heli ómnia quæ locútus sum super domum ejus: incípiam, et complébo.\par
\switchcolumn
\EN
\noindent And the Lord called Samuel again the third time. And he arose up and went to Heli. And said: Here am I: for thou didst call me. Then Heli understood that the Lord called the child, and he said to Samuel: Go, and sleep: and if he shall call thee any more, thou shalt say: Speak, Lord, for thy servant heareth. So Samuel went and slept in his place. And the Lord came and stood: and he called, as he had called the other times: Samuel, Samuel. And Samuel said: Speak, Lord, for thy servant heareth. And the Lord said to Samuel: Behold I do a thing in Israel: and whosoever shall hear it, both his ears shall tingle. In that day I will raise up against Heli all the things I have spoken concerning his house: I will begin, and I will make an end.\par
\switchcolumn*

\LA
\begin{Resp}\Rsign Comedétis carnes, et saturabímini pánibus: \Star{} Iste est panis, quem dedit vobis Dóminus ad vescéndum.\end{Resp}
\begin{Resp}\Vsign Non Móyses dedit vobis panem de cælo, sed Pater meus dat vobis panem de cælo verum.\end{Resp}
\begin{Resp}\Rsign Iste est panis, quem dedit vobis Dóminus ad vescéndum.\end{Resp}
\switchcolumn
\EN
\begin{Resp}\Rsign Ye shall eat flesh, and shall be filled with bread. \Star{} This is the bread which the Lord hath given you to eat.\end{Resp}
\begin{Resp}\Vsign Moses gave you not that Bread from heaven, but My Father giveth you the true Bread from heaven.\end{Resp}
\begin{Resp}\Rsign This is the bread which the Lord hath given you to eat.\end{Resp}
\switchcolumn*

\LA
\LessonHead{Lectio III}
\switchcolumn
\EN
\LessonHead{Lesson III}
\switchcolumn*

\LA
\Cite{1 Reg 3:15-20}
\switchcolumn
\EN
\Cite{1 Sam 3:15-20}
\switchcolumn*

\LA
\noindent Dormívit autem Sámuel usque mane, aperuítque óstia domus Dómini. Et Sámuel timébat indicáre visiónem Heli. Vocávit ergo Heli Samuélem, et dixit: Sámuel fili mi? Qui respóndens ait: Præsto sum. Et interrogávit eum: Quis est sermo, quem locútus est Dóminus ad te? oro te ne celáveris me: hæc fáciat tibi Deus, et hæc addat, si abscónderis a me sermónem ex ómnibus verbis quæ dicta sunt tibi. Indicávit ítaque ei Sámuel univérsos sermónes, et non abscóndit ab eo. Et ille respóndit: Dóminus est: quod bonum est in óculis suis fáciat. Crevit autem Sámuel, et Dóminus erat cum eo, et non cécidit ex ómnibus verbis ejus in terram. Et cognóvit univérsus Israël, a Dan usque Bersabée, quod fidélis Sámuel prophéta esset Dómini.\par
\switchcolumn
\EN
\noindent And Samuel slept till morning, and opened the doors of the house of the Lord. And Samuel feared to tell the vision to Heli. Then Heli called Samuel, and said: Samuel, my son. And he answered: Here am I. And he asked him: What is the word that the Lord hath spoken to thee? I beseech thee hide it not from me. May God do so and so to thee, and add so and so, if thou hide from me one word of all that were said to thee. So Samuel told him all the words, and did not hide them from him. And he answered: It is the Lord: let him do what is good in his sight. And Samuel grew, and the Lord was with him, and not one of his words fell to the ground. And all Israel from Dan to Bersabee, knew that Samuel was a faithful prophet of the Lord.\par
\switchcolumn*

\LA
\begin{Resp}\Rsign Respéxit Elías ad caput suum subcinerícium panem: qui surgens comédit et bibit: \Star{} Et ambulávit in fortitúdine cibi illíus usque ad montem Dei.\end{Resp}
\begin{Resp}\Vsign Si quis manducáverit ex hoc pane, vivet in ætérnum.\end{Resp}
\begin{Resp}\Rsign Et ambulávit in fortitúdine cibi illíus usque ad montem Dei.\end{Resp}
\begin{Resp}\Vsign Glória Patri, et Fílio, \Star{} et Spirítui Sancto.\end{Resp}
\begin{Resp}\Rsign Et ambulávit in fortitúdine cibi illíus usque ad montem Dei.\end{Resp}
\switchcolumn
\EN
\begin{Resp}\Rsign Elijah looked, and, behold, there was a cake baken on the coals at his head, and he arose, and did eat and drink; \Star{} And went in the strength of that meat [forty days and forty nights] unto the mount of God.\end{Resp}
\begin{Resp}\Vsign If any man eat of this Bread, he shall live for ever.\end{Resp}
\begin{Resp}\Rsign And went in the strength of that meat [forty days and forty nights] unto the mount of God.\end{Resp}
\begin{Resp}\Vsign Glory be to the Father, and to the Son, \Star{} and to the Holy Ghost.\end{Resp}
\begin{Resp}\Rsign And went in the strength of that meat [forty days and forty nights] unto the mount of God.\end{Resp}
\switchcolumn*

\LA
\LessonHead{Lectio IV}
\switchcolumn
\EN
\LessonHead{Lesson IV}
\switchcolumn*

\LA
\LessonTitle{Sermo sancti Joánnis Chrysóstomi}
\switchcolumn
\EN
\LessonTitle{From the Sermons of St. John Chrysostom, Patriarch of Constantinople.}
\switchcolumn*

\LA
\Source{Hom. 61 ad pópulum Antiochenum}
\switchcolumn
\EN
\Source{61st Homily to the people of Antioch.}
\switchcolumn*

\LA
\noindent Necessárium est, dilectíssimi, mysteriórum díscere miráculum, quodnam sit, et quare sit datum, et quæ ejus rei utílitas. Unum corpus effícimur: membra, inquit, ex carne ejus et ex óssibus ejus. Sequámur autem initiáti, quæ dicúntur. Ut ítaque non tantum per caritátem hoc fiámus, verum étiam ipsa re, in illam misceámur carnem: hoc namque per escam effícitur, quam largítus est nobis volens osténdere desidérium quod erga nos habet. Proptérea semetípsum nobis immíscuit, et corpus suum in nos contemperávit, ut unum quid simus tamquam corpus cápiti coaptátum; ardénter enim amántium hoc est.\par
\switchcolumn
\EN
\noindent Dearly beloved brethren, it behoveth us to learn the miracle of the Mysteries what the Gift is, and why It was given, and what is the use thereof. “We, being many, are one body,” saith [the Apostle Paul, (1 Cor. x. 17,) and again] “We are members of His Body, of His Flesh, and of His Bones.” (Eph. v. 30.) Only the initiated will now understand what I say. That this union may take place, not by love only, but verily and indeed, we ought to mingle our own with His Flesh. And this is done by eating that Food Which He hath given unto us, being fain to manifest that exceeding great love which He beareth to us-ward. To this end He hath mingled Himself with us, and infused His Body into our bodies, that we may be one together, like as the limbs of a man and his head are all of one body. Such union do they long for that love much.\par
\switchcolumn*

\LA
\begin{Resp}\Rsign Cenántibus illis, accépit Jesus panem, et benedíxit, ac fregit, dedítque discípulis suis, et ait: \Star{} Accípite et comédite: hoc est corpus meum.\end{Resp}
\begin{Resp}\Vsign Dixérunt viri tabernáculi mei: Quis det de cárnibus ejus, ut saturémur?\end{Resp}
\begin{Resp}\Rsign Accípite et comédite: hoc est corpus meum.\end{Resp}
\switchcolumn
\EN
\begin{Resp}\Rsign As they were eating, Jesus took bread, and blest it, and brake it, and gave it to the disciples, and said: \Star{} Take, eat this is My Body.\end{Resp}
\begin{Resp}\Vsign The men of my tabernacle said: O that we had of his flesh, we cannot be satisfied.\end{Resp}
\begin{Resp}\Rsign Take, eat this is My Body.\end{Resp}
\switchcolumn*

\LA
\LessonHead{Lectio V}
\switchcolumn
\EN
\LessonHead{Lesson V}
\switchcolumn*

\LA
\noindent Tamquam leónes ígitur ignem spirántes ab illa mensa recedámus, facti diábolo terríbiles, et caput nostrum mente revolvéntes, et caritátem quam erga nos osténdit. Nam paréntes quidem áliis sæpe fílios tradunt aléndos: ego autem, inquit, non ita, sed cárnibus meis alo, et meípsum vobis appóno, vos omnes generósos esse volens, et spem bonam de futúris vobis præbens: quippe qui vobis hic meípsum trádidi, multo magis id in futúro fáciam. Vólui frater vester fíeri, carni propter vos et sánguini communicávi: vobis vicíssim ipsam carnem et sánguinem, per quæ cognátus vester factus sum, trado.\par
\switchcolumn
\EN
\noindent When we come back from that Table we ought to be like so many lions breathing fire, dreadful to the devil. Our thoughts ought to be concentrated on our Great Head and the love which He showeth us. Many fathers and mothers there are who give their children to others to nurse, but I, saith the Lord to His children, I am not so, but I feed you with Mine Own Flesh, and join Myself to you, fain that ye all should be sons of noble blood now, and giving you a noble hope of that which ye shall be hereafter. I was content to become your Brother, I for your sakes have taken unto Me Flesh, and Blood, and that Flesh and Blood wherein I am become your Brother, the Same give I in turn unto you.\par
\switchcolumn*

\LA
\begin{Resp}\Rsign Accépit Jesus cálicem, postquam cenávit, dicens: Hic calix novum testaméntum est in meo sánguine: \Star{} Hoc fácite in meam commemoratiónem.\end{Resp}
\begin{Resp}\Vsign Memória memor ero, et tabéscet in me ánima mea.\end{Resp}
\begin{Resp}\Rsign Hoc fácite in meam commemoratiónem.\end{Resp}
\switchcolumn
\EN
\begin{Resp}\Rsign Jesus took the cup, after supper, saying: This cup is the New Testament in My Blood. \Star{} This do in remembrance of Me.\end{Resp}
\begin{Resp}\Vsign My soul hath them still in remembrance, and is humbled in me.\end{Resp}
\begin{Resp}\Rsign This do in remembrance of Me.\end{Resp}
\switchcolumn*

\LA
\LessonHead{Lectio VI}
\switchcolumn
\EN
\LessonHead{Lesson VI}
\switchcolumn*

\LA
\noindent Attendámus ítaque nobis ipsis, dilectíssimi, tálibus fruéntes bonis: et, cum áliquid turpe dícere voluérimus, vel nos ab ira córripi vidérimus, vel álio quópiam hujúsmodi vítio, considerémus quibus facti sumus digni; talíque cogitátio nobis irrationabílium mótuum sit corréctio. Quotquot ígitur hujus partícipes córporis effícimur, quotquot sánguinem degustámus; cogitémus quod illum sursum sedéntem, qui ab Angelis adorátur incorruptíbili vicínus virtúti, hunc degustámus. Hei mihi, quot ad salútem nobis viæ! Nos corpus suum effécit; nobis suum communicávit corpus: et horum nos nihil a malis avértit.\par
\switchcolumn
\EN
\noindent Let us then, dearly beloved brethren, take good heed to ourselves, as unto the holders of so great mercies, and when any foul word springeth to our lips, or we feel anger taking possession of us, or the sting of any other sinful passion, let us call to mind of What we have been counted worthy, and let that remembrance still the unruly motion. As often as we take that Body, as often as we taste that Blood, let us think how that we feed on Him Who is sitting on high, adored of Angels, at the right hand of the Eternal Power. Ah me, how many a way is open to us whereby we may be saved He hath made us His He hath given His Body to us and we still are not turned away from evil\par
\switchcolumn*

\LA
\begin{Resp}\Rsign Ego sum panis vitæ; patres vestri manducavérunt manna in desérto, et mórtui sunt: \Star{} Hic est panis de cælo descéndens, ut, si quis ex ipso mandúcet, non moriátur.\end{Resp}
\begin{Resp}\Vsign Ego sum panis vivus, qui de cælo descéndi: si quis manducáverit ex hoc pane, vivet in ætérnum.\end{Resp}
\begin{Resp}\Rsign Hic est panis de cælo descéndens, ut, si quis ex ipso mandúcet, non moriátur.\end{Resp}
\begin{Resp}\Vsign Glória Patri, et Fílio, \Star{} et Spirítui Sancto.\end{Resp}
\begin{Resp}\Rsign Hic est panis de cælo descéndens, ut, si quis ex ipso mandúcet, non moriátur.\end{Resp}
\switchcolumn
\EN
\begin{Resp}\Rsign I am that Bread of life. Your fathers did eat manna in the wilderness, and are dead. \Star{} This is the Bread Which cometh down from heaven, that a man may eat thereof, and not die.\end{Resp}
\begin{Resp}\Vsign I am the living Bread Which came down from heaven if any man eat of this Bread, he shall live for ever.\end{Resp}
\begin{Resp}\Rsign This is the Bread Which cometh down from heaven, that a man may eat thereof, and not die.\end{Resp}
\begin{Resp}\Vsign Glory be to the Father, and to the Son, \Star{} and to the Holy Ghost.\end{Resp}
\begin{Resp}\Rsign This is the Bread Which cometh down from heaven, that a man may eat thereof, and not die.\end{Resp}
\switchcolumn*

\LA
\LessonHead{Lectio VII}
\switchcolumn
\EN
\LessonHead{Lesson VII}
\switchcolumn*

\LA
\LessonTitle{Léctio sancti Evangélii secúndum Joánnem}
\switchcolumn
\EN
\LessonTitle{From the Holy Gospel according to John}
\switchcolumn*

\LA
\Cite{Joannes 6:56-59}
\switchcolumn
\EN
\Cite{John 6:56-59}
\switchcolumn*

\LA
\noindent In illo témpore: Dixit Jesus turbis Judæórum: Caro mea vere est cibus, et sanguis meus vere est potus. Et réliqua.\par
\switchcolumn
\EN
\noindent At that time, Jesus said unto the multitudes of the Jews: My Flesh is meat indeed and My Blood is drink indeed. And so on.\par
\switchcolumn*

\LA
\LessonTitle{De Homilía sancti Augustíni Epíscopi}
\switchcolumn
\EN
\LessonTitle{Homily by St. Augustine, Bishop of Hippo.}
\switchcolumn*

\LA
\Source{Tract. 27 in Joánn. ante medium.}
\switchcolumn
\EN
\Source{27th Tract on John.}
\switchcolumn*

\LA
\noindent Díximus, fratres, hoc Dóminum commendásse in manducatióne carnis suæ et potatióne sánguinis sui, ut in illo maneámus, et ipse in nobis. Manémus autem in illo, cum sumus membra ejus; manet autem ipse in nobis, cum sumus templum ejus. Ut autem simus membra ejus, únitas nos compáginat: ut compáginet únitas, quæ facit, nisi cáritas? Et cáritas Dei unde? Apóstolum intérroga. Cáritas, inquit, Dei diffúsa est in córdibus nostris per Spíritum Sanctum, qui datus est nobis.\par
\switchcolumn
\EN
\noindent I have said, my brethren, that what the Lord hath set before us, in eating of His Flesh and drinking of His Blood, is that we should dwell in Him, and He in us. We dwell in Him when we are His members, and He dwelleth in us when we are His temple. But the bond whereby we are made His members is oneness and what is the cause of oneness but love And love of God, whence is it? Ask the Apostle. “The love of God,” saith he, “is shed abroad in our hearts by the Holy Ghost, Which is given unto us.” (Rom. v. 5.)\par
\switchcolumn*

\LA
\begin{Resp}\Rsign Qui mandúcat meam carnem et bibit meum sánguinem, \Star{} In me manet, et ego in eo.\end{Resp}
\begin{Resp}\Vsign Non est ália nátio tam grandis, quæ hábeat deos appropinquántes sibi, sicut Deus noster adest nobis.\end{Resp}
\begin{Resp}\Rsign In me manet, et ego in eo.\end{Resp}
\switchcolumn
\EN
\begin{Resp}\Rsign He that eateth My Flesh and drinketh My Blood, \Star{} Dwelleth in Me, and I in him.\end{Resp}
\begin{Resp}\Vsign What nation is there so great, who hath gods so nigh unto them, as the Lord our God is to us\end{Resp}
\begin{Resp}\Rsign Dwelleth in Me, and I in him.\end{Resp}
\switchcolumn*

\LA
\LessonHead{Lectio VIII}
\switchcolumn
\EN
\LessonHead{Lesson VIII}
\switchcolumn*

\LA
\noindent Ergo spíritus est qui vivíficat; spíritus enim facit viva membra: nec viva membra spíritus facit, nisi quæ in córpore, quod végetat ipse spíritus, invénerit. Nam spíritus, qui est in te, o homo, quo constas ut homo sis, numquid vivíficat membrum quod separátum invénerit a carne tua? Spíritum tuum dico ánimam tuam. Anima tua non vivíficat, nisi membra quæ sunt in carne tua: unum si tollas, jam non vivificátur ex ánima tua, quia unitáti córporis tui non copulátur.\par
\switchcolumn
\EN
\noindent So “it is the spirit that quickeneth.” It is the spirit that maketh lively the limbs, nor is the quickening power of the spirit shed through any limbs but such as remain in union with the body whose the spirit is. The spirit that thou hast in thee, O man, and whereby thou art a man, doth that spirit shed life through any limb cut off from thy flesh By “spirit,” I mean soul. The soul quickeneth no limb but such as remain attached to the body. Cut one off, and the soul quickeneth it no more, for it is separate from the oneness of thy body.\par
\switchcolumn*

\LA
\begin{Resp}\Rsign Misit me vivens Pater, et ego vivo propter Patrem: \Star{} Et qui mandúcat me, vivet propter me.\end{Resp}
\begin{Resp}\Vsign Cibávit illum Dóminus pane vitæ et intelléctus.\end{Resp}
\begin{Resp}\Rsign Et qui mandúcat me, vivet propter me.\end{Resp}
\begin{Resp}\Vsign Glória Patri, et Fílio, \Star{} et Spirítui Sancto.\end{Resp}
\begin{Resp}\Rsign Et qui mandúcat me, vivet propter me.\end{Resp}
\switchcolumn
\EN
\begin{Resp}\Rsign As the living Father hath sent Me, and I live by the Father, \Star{} So he that eateth Me, even he shall live by Me.\end{Resp}
\begin{Resp}\Vsign With the bread of life and understanding hath the Lord fed him.\end{Resp}
\begin{Resp}\Rsign So he that eateth Me, even he shall live by Me.\end{Resp}
\begin{Resp}\Vsign Glory be to the Father, and to the Son, \Star{} and to the Holy Ghost.\end{Resp}
\begin{Resp}\Rsign So he that eateth Me, even he shall live by Me.\end{Resp}
\switchcolumn*

\LA
\LessonHead{Lectio IX}
\switchcolumn
\EN
\LessonHead{Lesson IX}
\switchcolumn*

\LA
\noindent Hæc dicúntur, ut amémus unitátem et timeámus separatiónem. Nihil enim sic debet formidáre Christiánus quam separári a córpore Christi. Si enim separátur a córpore Christi, non est membrum ejus, non vegetátur Spíritu ejus. Quisquis autem, inquit Apóstolus, Spíritum Christi non habet, hic non est ejus. Spíritus ergo est, qui vivíficat, caro autem non prodest quidquam. Verba, quæ ego locútus sum vobis, spíritus et vita sunt. Quid est, Spíritus et vita sunt? Spiritáliter intelligénda sunt. Intellexísti spiritáliter? Spíritus et vita sunt. Intellexísti carnáliter? Etiam sic illa spíritus et vita sunt, sed tibi non sunt.\par
\switchcolumn
\EN
\noindent These things I say, that we may love oneness and dread division. In sooth, there is nothing which a Christian ought so much to dread, as to be cut off from the Body of Christ. If he be cut off from the Body of Christ, he is no longer a member of Christ, and the Spirit of Christ no longer quickeneth him. “Now, if any man,” saith the Apostle, “have not the Spirit of Christ, he is none of His.” (Rom. viii. 9.) “It is the Spirit that quickeneth the flesh profiteth nothing the words that I speak unto you, they are spirit and they are life.” “Spirit and life”, what meaneth this? It is to be taken spiritually. Hast thou taken it spiritually? Then the words the Lord spake, unto thee they are spirit and they are life. Hast thou taken it carnally? Then the words of the Lord are still indeed spirit and life but not for thee.\par
\switchcolumn*

\LA
\TeDeum{Te Deum}
\switchcolumn
\EN
\TeDeum{Te Deum}
\switchcolumn*

\end{paracol}

\Office{Dominica II Post Pentecosten infra Octavam Corporis Christi}{Dominica II. Post Pentecosten infra Octavam Corporis Christi}{Semiduplex I classis}
\addcontentsline{toc}{subsection}{Dominica II Post Pentecosten infra Octavam Corporis Christi \textit{— Dominica II. Post Pentecosten infra Octavam Corporis Christi}}
\begin{paracol}{2}
\LA
\LessonHead{Lectio I}
\switchcolumn
\EN
\LessonHead{Lesson I}
\switchcolumn*

\LA
\LessonTitle{De libro primo Regum}
\switchcolumn
\EN
\LessonTitle{Lesson from the first book of Samuel}
\switchcolumn*

\LA
\Cite{1 Reg 4:1-3}
\switchcolumn
\EN
\Cite{1 Sam 4:1-3}
\switchcolumn*

\LA
\noindent Et factum est in diébus illis, convenérunt Philísthiim in pugnam: et egréssus est Israël óbviam Philísthiim in prǽlium, et castrametátus est juxta lápidem Adiutórii. Porro Philísthiim venérunt in Aphec, et instruxérunt áciem contra Israël. Inito autem certámine, terga vertit Israël Philisthǽis: et cæsa sunt in illo certámine passim per agros, quasi quátuor míllia virórum. Et revérsus est pópulus ad castra: dixerúntque majóres natu de Israël: Quare percússit nos Dóminus hódie coram Philísthiim? afferámus ad nos de Silo arcam fœ́deris Dómini, et véniat in médium nostri, ut salvet nos de manu inimicórum nostrórum.\par
\switchcolumn
\EN
\noindent And it came to pass in those days, that the Philistines gathered themselves together to fight: and Israel went out to war against the Philistines, and camped by the Stone of help. And the Philistines came to Aphec, And put their army in array against Israel. And when they had joined battle, Israel turned their backs to the Philistines, and there was slain in that fight here and there in the fields about four thousand men. And the people returned to the camp: and the ancients of Israel said: Why hath the Lord defeated us today before the Philistines? Let us fetch unto us the ark of the covenant of the Lord from Silo, and let it come in the midst of us, that it may save us from the hand of our enemies.\par
\switchcolumn*

\LA
\begin{Resp}\Rsign Immolábit hædum multitúdo filiórum Israël ad vésperam Paschæ: \Star{} Et edent carnes et ázymos panes.\end{Resp}
\begin{Resp}\Vsign Pascha nostrum immolátus est Christus: ítaque epulémur in ázymis sinceritátis et veritátis.\end{Resp}
\begin{Resp}\Rsign Et edent carnes et ázymos panes.\end{Resp}
\switchcolumn
\EN
\begin{Resp}\Rsign The whole assembly of the children of Israel shall kill the lamb toward the evening of the Passover. \Star{} And they shall eat the flesh, and unleavened bread.\end{Resp}
\begin{Resp}\Vsign Even Christ our Passover is sacrificed for us therefore let us keep the feast with the unleavened bread of sincerity and truth.\end{Resp}
\begin{Resp}\Rsign And they shall eat the flesh, and unleavened bread.\end{Resp}
\switchcolumn*

\LA
\LessonHead{Lectio II}
\switchcolumn
\EN
\LessonHead{Lesson II}
\switchcolumn*

\LA
\Cite{1 Reg 4:4-6}
\switchcolumn
\EN
\Cite{1 Sam 5:4-11}
\switchcolumn*

\LA
\noindent Misit ergo pópulus in Silo, et tulérunt inde arcam fœ́deris Dómini exercítuum sedéntis super chérubim: erántque duo fílii Heli cum arca fœ́deris Dei, Ophni et Phínees. Cumque venísset arca fœ́deris Dómini in castra, vociferátus est omnis Israël clamóre grandi, et persónuit terra. Et audiérunt Philísthiim vocem clamóris, dixerúntque: Quænam est hæc vox clamóris magni in castris Hebræórum? Et cognovérunt quod arca Dómini venísset in castra.\par
\switchcolumn
\EN
\noindent So the people sent to Silo, and they brought from thence the ark of the covenant of the Lord of hosts sitting upon the cherubims: and the two sons of Heli, Ophni and Phinees, were with the ark of the covenant of God. And when the ark of the covenant of the Lord was come into the camp, all Israel shouted with a great shout, and the earth rang again. And the Philistines heard the noise of the shout, and they said: What is this noise of a great shout in the camp of the Hebrews? And they understood that the ark of the Lord was come into the camp.\par
\switchcolumn*

\LA
\begin{Resp}\Rsign Comedétis carnes, et saturabímini pánibus: \Star{} Iste est panis, quem dedit vobis Dóminus ad vescéndum.\end{Resp}
\begin{Resp}\Vsign Non Móyses dedit vobis panem de cælo, sed Pater meus dat vobis panem de cælo verum.\end{Resp}
\begin{Resp}\Rsign Iste est panis, quem dedit vobis Dóminus ad vescéndum.\end{Resp}
\switchcolumn
\EN
\begin{Resp}\Rsign Ye shall eat flesh, and shall be filled with bread. \Star{} This is the bread which the Lord hath given you to eat.\end{Resp}
\begin{Resp}\Vsign Moses gave you not that Bread from heaven, but My Father giveth you the true Bread from heaven.\end{Resp}
\begin{Resp}\Rsign This is the bread which the Lord hath given you to eat.\end{Resp}
\switchcolumn*

\LA
\LessonHead{Lectio III}
\switchcolumn
\EN
\LessonHead{Lesson III}
\switchcolumn*

\LA
\Cite{1 Reg 4:7-11}
\switchcolumn
\EN
\Cite{1 Sam 4:7-11}
\switchcolumn*

\LA
\noindent Timuerúntque Philísthiim, dicéntes: Venit Deus in castra. Et ingemuérunt, dicéntes: Væ nobis: non enim fuit tanta exultátio heri et nudiustértius: væ nobis. Quis nos salvábit de manu deórum sublímium istórum? hi sunt dii, qui percussérunt Ægýptum omni plaga in desérto. Confortámini, et estóte viri, Philísthiim, ne serviátis Hebrǽis, sicut et illi serviérunt vobis: confortámini, et belláte. Pugnavérunt ergo Philísthiim, et cæsus est Israël, et fugit unusquísque in tabernáculum suum: et facta est plaga magna nimis, et cecidérunt de Israël trigínta míllia péditum. Et arca Dei capta est: duo quoque fílii Heli mórtui sunt, Ophni et Phínees.\par
\switchcolumn
\EN
\noindent And the Philistines were afraid, saying: God is come into the camp. And sighing, they said: Woe to us: for there was no such great joy yesterday and the day before: Woe to us. Who shall deliver us from the hand of these high gods? these are the gods that struck Egypt with all the plagues in the desert. Take courage and behave like men, ye Philistines: lest you come to be servants to the Hebrews, as they have served you: take courage and fight. So the Philistines fought, and Israel was overthrown, and every man fled to his own dwelling: and there was an exceeding great slaughter; for there fell of Israel thirty thousand footmen. And the ark of God was taken: and the two sons of Heli, Ophni and Phinees, were slain.\par
\switchcolumn*

\LA
\begin{Resp}\Rsign Respéxit Elías ad caput suum subcinerícium panem: qui surgens comédit et bibit: \Star{} Et ambulávit in fortitúdine cibi illíus usque ad montem Dei.\end{Resp}
\begin{Resp}\Vsign Si quis manducáverit ex hoc pane, vivet in ætérnum.\end{Resp}
\begin{Resp}\Rsign Et ambulávit in fortitúdine cibi illíus usque ad montem Dei.\end{Resp}
\begin{Resp}\Vsign Glória Patri, et Fílio, \Star{} et Spirítui Sancto.\end{Resp}
\begin{Resp}\Rsign Et ambulávit in fortitúdine cibi illíus usque ad montem Dei.\end{Resp}
\switchcolumn
\EN
\begin{Resp}\Rsign Elijah looked, and, behold, there was a cake baken on the coals at his head, and he arose, and did eat and drink; \Star{} And went in the strength of that meat [forty days and forty nights] unto the mount of God.\end{Resp}
\begin{Resp}\Vsign If any man eat of this Bread, he shall live for ever.\end{Resp}
\begin{Resp}\Rsign And went in the strength of that meat [forty days and forty nights] unto the mount of God.\end{Resp}
\begin{Resp}\Vsign Glory be to the Father, and to the Son, \Star{} and to the Holy Ghost.\end{Resp}
\begin{Resp}\Rsign And went in the strength of that meat [forty days and forty nights] unto the mount of God.\end{Resp}
\switchcolumn*

\LA
\LessonHead{Lectio IV}
\switchcolumn
\EN
\LessonHead{Lesson IV}
\switchcolumn*

\LA
\LessonTitle{Sermo sancti Joánnis Chrysostomi}
\switchcolumn
\EN
\LessonTitle{From the Sermons of St. John Chrysostom Patriarch of Constantinople.}
\switchcolumn*

\LA
\Source{Ex Hom. 60 ad pop. Antioch}
\switchcolumn
\EN
\Source{60th Homily to the people of Antioch.}
\switchcolumn*

\LA
\noindent Quóniam Verbum dicit: Hoc est corpus meum; et assentiámur et credámus et intellectuálibus ipsum óculis intueámur. Nihil enim sensíbile nobis Christus trádidit; sed sensibílibus quidem rebus, at ómnia intelligibília. Itidem et in baptísmate: per rem nempe sensíbilem, aquam, donum confértur; intelligíbile vero quod perfícitur, generátio et renovátio. Si enim incorpóreus esses, nuda et incorpórea tibi dedísset ipse dona; sed quóniam ánima córpori consérta est, in sensibílibus intelligibília tibi præbet. Quot nunc dicunt: Vellem ipsíus formam aspícere, figúram, vestiménta, calceaménta? Ecce eum vides, ipsum tangis, ipsum mandúcas. Et tu quidem vestiménta cupis vidére; ipse vero tibi concédit non tantum vidére, verum et manducáre, et tángere, et intra te súmere.\par
\switchcolumn
\EN
\noindent His Word saith: “This is My Body.” This we confess, and believe, and, with spiritual eyes, do see. Christ hath not left unto us Himself in such form as that we can see, hear, touch, smell, or taste Him and yet hath He left Himself unto us in things which we can see, hear, touch, smell, and taste, and which all men may understand. Thus also is it in baptism by mean of water, which men perceive outwardly, is given unto them a gift which they can grasp only inwardly, that is, a new birth. If we had no bodies, then would these things be given us without any outward and visible signs, but since we are here made up of souls and bodies, there are given unto our souls gifts which they can grasp, in outward signs which our bodies may perceive. How many there be which say I would that I could see His comely presence, His Face, His garments, even His shoes Behold, thou dost see and touch Him, yea, thou dost feed upon Him. And wouldest thou behold His raiment Lo, He hath given unto thee not only to behold it, but to feed upon it, and handle it, and take it into thyself.\par
\switchcolumn*

\LA
\begin{Resp}\Rsign Cenántibus illis, accépit Jesus panem, et benedíxit, ac fregit, dedítque discípulis suis, et ait: \Star{} Accípite et comédite: hoc est corpus meum.\end{Resp}
\begin{Resp}\Vsign Dixérunt viri tabernáculi mei: Quis det de cárnibus ejus, ut saturémur?\end{Resp}
\begin{Resp}\Rsign Accípite et comédite: hoc est corpus meum.\end{Resp}
\switchcolumn
\EN
\begin{Resp}\Rsign As they were eating, Jesus took bread, and blest it, and brake it, and gave it to the disciples, and said: \Star{} Take, eat this is My Body.\end{Resp}
\begin{Resp}\Vsign The men of my tabernacle said: O that we had of his flesh, we cannot be satisfied.\end{Resp}
\begin{Resp}\Rsign Take, eat this is My Body.\end{Resp}
\switchcolumn*

\LA
\LessonHead{Lectio V}
\switchcolumn
\EN
\LessonHead{Lesson V}
\switchcolumn*

\LA
\noindent Igitur accédat nemo cum náusea, nemo resolútus; omnes accénsi, omnes fervéntes et excitáti. Nam si Judǽi stantes, et calceaménta in pédibus habéntes, et báculos mánibus gestántes, agnum cum festinatióne comedébant; te multo magis opórtet esse solértem. Nam illi quidem in Palæstínam erant profectúri, et proptérea viatórum figúram habébant: tu vero debes in cælum migráre. Quaprópter in ómnibus opórtet te vigiláre; nec enim parva pœna propónitur indígne suméntibus. Cógita quantum advérsus proditórem indignáris, et contra eos qui illum crucifixérunt: ítaque consídera, ne tu quoque sis reus córporis et sánguinis Christi. Illi sanctíssimum corpus occidérunt, tu vero pollúta súscipis ánima, post tot benefícia. Neque enim illi satis fuit, hóminem fíeri, cólaphis cædi, et crucifígi; verum et semetípsum nobis commíscet; et non fide tantum, verum et ipsa re, nos suum éfficit corpus.\par
\switchcolumn
\EN
\noindent At this table of the Lord let none dare to draw near with squeamishness or carelessness. Let all be fiery, all hot, all roused. To the Jews it was commanded touching the Paschal lamb, (Exod. xii. 11): “And thus shall ye eat it with your loins girded, your shoes on your feet, and your staff in your hand and ye shall eat it in haste it is the Lord's Passover.” But thou needest to be more watchful than they. They were just about to travel from Egypt to Palestine, and therefore they bore the guise of travellers but the journey that lieth before thee is from earth to heaven. And therefore it behoveth thee in all things to be on thy guard, for the punishment of him that eateth or drinketh unworthily is no light one (1 Cor. xi. 27.) Bethink thee how thou art indignant against him which betrayed, and them that crucified the Lord and look to it well that thou also be not “Guilty of the Body and Blood of the Lord.” As for them, they slew His Most Holy Body but thou, after all that He hath done for thee, dost thrust Him into thy polluted soul. For His love, it was not enough to be made Man, to be buffeted, and to be crucified He hath also mingled Himself with us, by making us His Body, and that not by faith only, but verily and indeed.\par
\switchcolumn*

\LA
\begin{Resp}\Rsign Accépit Jesus cálicem, postquam cenávit, dicens: Hic calix novum testaméntum est in meo sánguine: \Star{} Hoc fácite in meam commemoratiónem.\end{Resp}
\begin{Resp}\Vsign Memória memor ero, et tabéscet in me ánima mea.\end{Resp}
\begin{Resp}\Rsign Hoc fácite in meam commemoratiónem.\end{Resp}
\switchcolumn
\EN
\begin{Resp}\Rsign Jesus took the cup, after supper, saying: This cup is the New Testament in My Blood. \Star{} This do in remembrance of Me.\end{Resp}
\begin{Resp}\Vsign My soul hath them still in remembrance, and is humbled in me.\end{Resp}
\begin{Resp}\Rsign This do in remembrance of Me.\end{Resp}
\switchcolumn*

\LA
\LessonHead{Lectio VI}
\switchcolumn
\EN
\LessonHead{Lesson VI}
\switchcolumn*

\LA
\noindent Quo non opórtet ígitur esse puriórem, tali fruéntem sacrifício? quo solári rádio non splendidiórem manum, carnem hanc dividéntem? os quod igni spiritáli replétur, linguam quæ treméndo nimis sánguine rubéscit? Cógita quali sis insignítus honóre, quali mensa fruáris. Quod Angeli vidéntes horréscunt, neque líbere audent intuéri propter emicántem inde splendórem; hoc nos páscimur, huic nos unímur, et facti sumus unum Christi corpus, et una caro. Quis loquétur poténtias Dómini, audítas fáciet omnes laudes ejus? Quis pastor oves próprio pascit cruóre? Et quid dico, pastor? Matres multæ sunt, quæ post partus dolóres, fílios áliis tradunt nutrícibus. Hoc autem ipse non est passus; sed ipse nos próprio sánguine pascit, et per ómnia nos sibi coagméntat.\par
\switchcolumn
\EN
\noindent Anything be purer than that man ought to be, who eateth of this great Sacrifice Can sun-beam be clearer than that hand ought to be which breaketh this Flesh? that mouth, which is filled with that spiritual fire? that tongue, which is reddened by that Blood, awful exceedingly? That whereon the Angels quail to look, neither dare to gaze steadfastly upon It, because of the blinding glory that shineth therefrom, upon This we feed, with This we become one, and are made one body of Christ, and one flesh. “Who can utter the mighty acts of the Lord; who can show forth all His praise?” (Ps. cv. 2.) Where is the shepherd which feedeth his flock with his own blood Nay, why should I say, shepherd Many mothers there be, who after all the pains of travail, give their own little ones to strangers to nurse. But so would not He, but feedeth us with His Own Blood, and maketh us to grow up in His Own substance.\par
\switchcolumn*

\LA
\begin{Resp}\Rsign Ego sum panis vitæ; patres vestri manducavérunt manna in desérto, et mórtui sunt: \Star{} Hic est panis de cælo descéndens, ut, si quis ex ipso mandúcet, non moriátur.\end{Resp}
\begin{Resp}\Vsign Ego sum panis vivus, qui de cælo descéndi: si quis manducáverit ex hoc pane, vivet in ætérnum.\end{Resp}
\begin{Resp}\Rsign Hic est panis de cælo descéndens, ut, si quis ex ipso mandúcet, non moriátur.\end{Resp}
\begin{Resp}\Vsign Glória Patri, et Fílio, \Star{} et Spirítui Sancto.\end{Resp}
\begin{Resp}\Rsign Hic est panis de cælo descéndens, ut, si quis ex ipso mandúcet, non moriátur.\end{Resp}
\switchcolumn
\EN
\begin{Resp}\Rsign I am that Bread of life. Your fathers did eat manna in the wilderness, and are dead. \Star{} This is the Bread Which cometh down from heaven, that a man may eat thereof, and not die.\end{Resp}
\begin{Resp}\Vsign I am the living Bread Which came down from heaven if any man eat of this Bread, he shall live for ever.\end{Resp}
\begin{Resp}\Rsign This is the Bread Which cometh down from heaven, that a man may eat thereof, and not die.\end{Resp}
\begin{Resp}\Vsign Glory be to the Father, and to the Son, \Star{} and to the Holy Ghost.\end{Resp}
\begin{Resp}\Rsign This is the Bread Which cometh down from heaven, that a man may eat thereof, and not die.\end{Resp}
\switchcolumn*

\LA
\LessonHead{Lectio VII}
\switchcolumn
\EN
\LessonHead{Lesson VII}
\switchcolumn*

\LA
\LessonTitle{Léctio sancti Evangélii secúndum Lucam}
\switchcolumn
\EN
\LessonTitle{From the Holy Gospel according to Luke}
\switchcolumn*

\LA
\Cite{Luc 14:16-24}
\switchcolumn
\EN
\Cite{Luke 14:16-24}
\switchcolumn*

\LA
\noindent In illo témpore: Dixit Jesus pharisǽis parábolam hanc: Homo quidam fecit cenam magnam, et vocávit multos. Et réliqua.\par
\switchcolumn
\EN
\noindent At that time: Jesus spake unto the Pharisees this parable: A certain man made a great supper, and bade many. And so on.\par
\switchcolumn*

\LA
\LessonTitle{Homilía sancti Gregórii Papæ}
\switchcolumn
\EN
\LessonTitle{Homily by Pope St. Gregory the Great}
\switchcolumn*

\LA
\Source{Homilia 36 in Evangelia}
\switchcolumn
\EN
\Source{36th upon the Gospels}
\switchcolumn*

\LA
\noindent Hoc distáre, fratres caríssimi, inter delícias córporis et cordis solet: quod corporáles delíciæ, cum non habéntur, grave in se desidérium accéndunt; cum vero ávide edúntur, comedéntem prótinus in fastídium per satietátem vertunt. At contra, spiritáles delíciæ, cum non habéntur, in fastídio sunt; cum vero habéntur, in desidério: tantóque a comedénte ámplius esuriúntur, quanto et ab esuriénte ámplius comedúntur. In illis appetítus placet, experiéntia dísplicet; in istis appetítus saturitátem, satúritas fastídium génerat: in istis autem appetítus saturitátem, satúritas appetítum parit.\par
\switchcolumn
\EN
\noindent Dearly beloved brethren, between the dainties of the body and the dainties of the mind there is this difference, that the dainties of the body, when we lack them, raise up a great hunger after them, and when we devour them, straightway our fulness worketh in us niceness. But about the dainties of the mind we are nice while as yet we lack them, and when we fill ourselves with them, then are we an-hungered after them, and the more, being an-hungered, we feed thereon, the more are we an-hungered thereafter. In the bodily dainties, the hunger is keener than the fullness, but in the spiritual the fulness is keener than the hunger. In the bodily, hunger gendereth fullness, and fulness niceness in the spiritual, hunger indeed gendereth fullness, but fullness gendereth hunger.\par
\switchcolumn*

\LA
\begin{Resp}\Rsign Qui mandúcat meam carnem et bibit meum sánguinem, \Star{} In me manet, et ego in eo.\end{Resp}
\begin{Resp}\Vsign Non est ália nátio tam grandis, quæ hábeat deos appropinquántes sibi, sicut Deus noster adest nobis.\end{Resp}
\begin{Resp}\Rsign In me manet, et ego in eo.\end{Resp}
\switchcolumn
\EN
\begin{Resp}\Rsign He that eateth My Flesh and drinketh My Blood, \Star{} Dwelleth in Me, and I in him.\end{Resp}
\begin{Resp}\Vsign What nation is there so great, who hath gods so nigh unto them, as the Lord our God is to us\end{Resp}
\begin{Resp}\Rsign Dwelleth in Me, and I in him.\end{Resp}
\switchcolumn*

\LA
\LessonHead{Lectio VIII}
\switchcolumn
\EN
\LessonHead{Lesson VIII}
\switchcolumn*

\LA
\noindent Augent enim spiritáles delíciæ desidérium in mente, dum sátiant: quia quanto magis eárum sapor percípitur, eo ámplius cognóscitur quod avídius amétur; et idcírco non hábitæ amári non possunt, quia eárum sapor ignorátur. Quis enim amáre váleat quod ignórat? Proínde Psalmísta nos ádmonet, dicens: Gustáte et vidéte, quóniam suávis est Dóminus. Ac si apérte dicat: Suavitátem ejus non cognóscitis, si hanc mínime gustátis; sed cibum vitæ ex paláto cordis tángite, ut probántes ejus dulcédinem, amáre valeátis. Has autem homo delícias tunc amísit, cum in paradíso peccávit; extra éxiit, cum os a cibo ætérnæ dulcédinis clausit.\par
\switchcolumn
\EN
\noindent Spiritual dainties, in the very eating, do stir up the keenness of hunger in the mind which they fill, for, the more we taste their sweetness, the better we know how well they deserve to be loved and, if we taste them not, we cannot love them, for we know not how sweet they be. And who can love that whereof he knoweth nothing? Hence saith the Psalmist “O taste and see that the Lord is good,” (Ps. xxxiii. 9,) that is, as it were, “If ye taste not, ye shall not see His goodness but let your heart once taste the bread of life, and then indeed, having tasted and proved His sweetness, ye shall be able to love Him.” But these were the dainties which man lost when he sinned in Eden, and when he had shut his own mouth against the sweet bread whereof if any man eat he shall live for ever, he forsook paradise.\par
\switchcolumn*

\LA
\begin{Resp}\Rsign Homo quidam fecit cœnam magnam, et misit servum suum hora cœnæ dícere invitátis ut venírent, \Star{} Quia paráta sunt ómnia.\end{Resp}
\begin{Resp}\Vsign Veníte, comédite panem meum, et bíbite vinum quod míscui vobis.\end{Resp}
\begin{Resp}\Rsign Quia paráta sunt ómnia.\end{Resp}
\begin{Resp}\Vsign Glória Patri, et Fílio, \Star{} et Spirítui Sancto.\end{Resp}
\begin{Resp}\Rsign Quia paráta sunt ómnia.\end{Resp}
\switchcolumn
\EN
\begin{Resp}\Rsign A certain man made a great supper, and sent his servant at suppertime to say to them that were bidden: Come \Star{} For all things are now ready.\end{Resp}
\begin{Resp}\Vsign Come, eat of my bread, and drink of the wine which I have mingled for you.\end{Resp}
\begin{Resp}\Rsign For all things are now ready.\end{Resp}
\begin{Resp}\Vsign Glory be to the Father, and to the Son, \Star{} and to the Holy Ghost.\end{Resp}
\begin{Resp}\Rsign For all things are now ready\end{Resp}
\switchcolumn*

\LA
\LessonHead{Lectio IX}
\switchcolumn
\EN
\LessonHead{Lesson IX}
\switchcolumn*

\LA
\noindent Unde nos quoque, nati in hujus peregrinatiónis ærúmna, huc fastidiósi jam vénimus, nescímus quid desideráre debeámus. Tantóque se ámplius fastídii nostri morbus exággerat, quanto se magis ab esu illíus dulcédinis ánimus elóngat; et eo jam intérnas delícias non áppetit, quo eas comédere, diu longéque desuévit. Fastídio ergo nostro tabéscimus, et longa inédiæ peste fatigámur. Et quia gustáre intus nólumus parátam dulcédinem, amámus foris míseri famem nostram.\par
\switchcolumn
\EN
\noindent And we that, from the first man, are born under the afflictions of this pilgrimage, are come into the world smitten with niceness we know not what we ought to want, and the disease of our niceness groweth the worse, as our soul draweth itself the more away from that bread of sweetness. We are no longer an-hungered after inward dainties, since we have lost the use of feeding on them. And so in our niceness we starve, and the sickness of long famishing maketh prey of our health. We will not eat of that inward sweetness which is made ready for us, and being enamoured only of things outward we sink into the wretchedness of loving starvation.\par
\switchcolumn*

\LA
\TeDeum{Te Deum}
\switchcolumn
\EN
\TeDeum{Te Deum}
\switchcolumn*

\end{paracol}

\Office{Feria Secunda infra Octavam Corporis Christi}{Monday within the Octave of Corpus Christi}{Semiduplex II classis}
\addcontentsline{toc}{subsection}{Feria Secunda infra Octavam Corporis Christi \textit{— Monday within the Octave of Corpus Christi}}
\begin{paracol}{2}
\LA
\LessonHead{Lectio I}
\switchcolumn
\EN
\LessonHead{Lesson I}
\switchcolumn*

\LA
\LessonTitle{De libro primo Regum}
\switchcolumn
\EN
\LessonTitle{Lesson from the first book of Samuel}
\switchcolumn*

\LA
\Cite{1 Reg 5:1-5}
\switchcolumn
\EN
\Cite{1 Sam 5:1-5}
\switchcolumn*

\LA
\noindent Philísthiim autem tulérunt arcam Dei et asportavérunt eam a Lápide adjutórii in Azótum. Tulerúntque Philísthiim arcam Dei et intulérunt eam in templum Dagon et statuérunt eam juxta Dagon. Cumque surrexíssent dilúculo Azótii áltera die, ecce Dagon jacébat pronus in terra ante arcam Dómini; et tulérunt Dagon et restituérunt eum in locum suum. Rursúmque mane die áltera consurgéntes invenérunt Dagon jacéntem super fáciem suam in terra coram arca Dómini; caput autem Dagon et duæ palmæ mánuum ejus abscíssæ erant super limen; porro Dagon solus truncus remánserat in loco suo.\par
\switchcolumn
\EN
\noindent And the Philistines took the ark of God, and carried it from the Stone of help into Azotus. And the Philistines took the ark of God, and brought it into the temple of Dagon, and set it by Dagon. And when the Azotians arose early the next day, behold Dagon lay upon his face on the ground before the ark of the Lord: and they took Dagon, and set him again in his place. And the next day again, when they rose in the morning, they found Dagon lying upon his face on the earth before the ark of the Lord: and the head of Dagon, and both the palms of his hands were cut off upon the threshold: And only the stump of Dagon remained in its place.\par
\switchcolumn*

\LA
\begin{Resp}\Rsign Immolábit hædum multitúdo filiórum Israël ad vésperam Paschæ: \Star{} Et edent carnes et ázymos panes.\end{Resp}
\begin{Resp}\Vsign Pascha nostrum immolátus est Christus: ítaque epulémur in ázymis sinceritátis et veritátis.\end{Resp}
\begin{Resp}\Rsign Et edent carnes et ázymos panes.\end{Resp}
\switchcolumn
\EN
\begin{Resp}\Rsign The whole assembly of the children of Israel shall kill the lamb toward the evening of the Passover. \Star{} And they shall eat the flesh, and unleavened bread.\end{Resp}
\begin{Resp}\Vsign Even Christ our Passover is sacrificed for us therefore let us keep the feast with the unleavened bread of sincerity and truth.\end{Resp}
\begin{Resp}\Rsign And they shall eat the flesh, and unleavened bread.\end{Resp}
\switchcolumn*

\LA
\LessonHead{Lectio II}
\switchcolumn
\EN
\LessonHead{Lesson II}
\switchcolumn*

\LA
\Cite{1 Reg 5:6-8}
\switchcolumn
\EN
\Cite{1 Sam 5:6-8}
\switchcolumn*

\LA
\noindent Aggraváta est autem manus Dómini super Azótios, et demolítus est eos. Et ebulliérunt villæ et agri in médio regiónis illíus, et nati sunt mures, et facta est confúsio mortis magnæ in civitáte. Vidéntes autem viri Azótii hujuscémodi plagam dixérunt: Non máneat arca Dei Israël apud nos, quóniam dura est manus ejus super nos et super Dagon deum nostrum. Et mitténtes congregavérunt omnes sátrapas Philisthinórum ad se et dixérunt: Quid faciémus de arca Dei Israël? Responderúntque Gethǽi: Circumducátur arca Dei Israël.\par
\switchcolumn
\EN
\noindent And the hand of the Lord was heavy upon the Azotians, and he destroyed them, and afflicted Azotus and the coasts thereof with emerods. And in the villages and fields in the midst of that country, there came forth a multitude of mice, and there was the confusion of a great mortality in the city. And the men of Azotus seeing this kind of plague, said: The ark of the God of Israel shall not stay with us: for his hand is heavy upon us, and upon Dagon our god. And sending, they gathered together all the lords of the Philistines to them, and said: What shall we do with the ark of the God of Israel? And the Gethrites answered: Let the ark of the God of Israel be carried about.\par
\switchcolumn*

\LA
\begin{Resp}\Rsign Comedétis carnes, et saturabímini pánibus: \Star{} Iste est panis, quem dedit vobis Dóminus ad vescéndum.\end{Resp}
\begin{Resp}\Vsign Non Móyses dedit vobis panem de cælo, sed Pater meus dat vobis panem de cælo verum.\end{Resp}
\begin{Resp}\Rsign Iste est panis, quem dedit vobis Dóminus ad vescéndum.\end{Resp}
\switchcolumn
\EN
\begin{Resp}\Rsign Ye shall eat flesh, and shall be filled with bread. \Star{} This is the bread which the Lord hath given you to eat.\end{Resp}
\begin{Resp}\Vsign Moses gave you not that Bread from heaven, but My Father giveth you the true Bread from heaven.\end{Resp}
\begin{Resp}\Rsign This is the bread which the Lord hath given you to eat.\end{Resp}
\switchcolumn*

\LA
\LessonHead{Lectio III}
\switchcolumn
\EN
\LessonHead{Lesson III}
\switchcolumn*

\LA
\Cite{1 Reg 5:8-12}
\switchcolumn
\EN
\Cite{1 Sam 5:8-12}
\switchcolumn*

\LA
\noindent Et circumduxérunt arcam Dei Israël. Illis autem circumducéntibus eam, fiébat manus Dómini per síngulas civitátes interfectiónis magnæ nimis; et percutiébat viros uniuscujúsque urbis a parvo usque ad majórem. Misérunt ergo arcam Dei in Accaron. Cumque venísset arca Dei in Accaron, exclamavérunt Accaronítæ dicéntes: Adduxérunt ad nos arcam Dei Israël, ut interfíciat nos et pópulum nostrum. Misérunt ítaque et congregavérunt omnes sátrapas Philisthinórum, qui dixérunt: Dimíttite arcam Dei Israël, et revertátur in locum suum et non interfíciat nos cum pópulo nostro. Fiébat enim pavor mortis in síngulis úrbibus et gravíssima valde manus Dei.\par
\switchcolumn
\EN
\noindent And they carried the ark of the God of Israel about. And while they were carrying it about, the hand of the Lord came upon every city with an exceeding great slaughter: and he smote the men of every city, both small and great, and they had emerods in their secret parts. And the Gethrites consulted together, and made themselves seats of skins. Therefore they sent the ark of God into Accaron. And when the ark of God was come into Accaron, the Accaronites cried out, saying: They have brought the ark of the God of Israel to us, to kill us and our people. They sent therefore and gathered together all the lords of the Philistines: and they said: Send away the ark of the God of Israel, and let it return into its own place, and not kill us and our people. For there was the fear of death in every city, and the hand of God was exceeding heavy.\par
\switchcolumn*

\LA
\begin{Resp}\Rsign Respéxit Elías ad caput suum subcinerícium panem: qui surgens comédit et bibit: \Star{} Et ambulávit in fortitúdine cibi illíus usque ad montem Dei.\end{Resp}
\begin{Resp}\Vsign Si quis manducáverit ex hoc pane, vivet in ætérnum.\end{Resp}
\begin{Resp}\Rsign Et ambulávit in fortitúdine cibi illíus usque ad montem Dei.\end{Resp}
\begin{Resp}\Vsign Glória Patri, et Fílio, \Star{} et Spirítui Sancto.\end{Resp}
\begin{Resp}\Rsign Et ambulávit in fortitúdine cibi illíus usque ad montem Dei.\end{Resp}
\switchcolumn
\EN
\begin{Resp}\Rsign Elijah looked, and, behold, there was a cake baken on the coals at his head, and he arose, and did eat and drink; \Star{} And went in the strength of that meat [forty days and forty nights] unto the mount of God.\end{Resp}
\begin{Resp}\Vsign If any man eat of this Bread, he shall live for ever.\end{Resp}
\begin{Resp}\Rsign And went in the strength of that meat [forty days and forty nights] unto the mount of God.\end{Resp}
\begin{Resp}\Vsign Glory be to the Father, and to the Son, \Star{} and to the Holy Ghost.\end{Resp}
\begin{Resp}\Rsign And went in the strength of that meat [forty days and forty nights] unto the mount of God.\end{Resp}
\switchcolumn*

\LA
\LessonHead{Lectio IV}
\switchcolumn
\EN
\LessonHead{Lesson IV}
\switchcolumn*

\LA
\LessonTitle{De Sermóne sancti Joánnis Chrysostomi}
\switchcolumn
\EN
\LessonTitle{From the Sermons of St. John Chrysostom, Patriarch of Constantinople}
\switchcolumn*

\LA
\Source{Homilia 60 ad populum Antiochenum}
\switchcolumn
\EN
\Source{Continuation of the 60th Homily}
\switchcolumn*

\LA
\noindent Unicuíque fidélium Christus semetípsum per mystéria commíscet, et quos génuit, per semetípsum enútrit, nec álteri tradit; per hoc tibi rursum persuádens, quod carnem tuam assúmpsit. Ne torpeámus ígitur tanta digni caritáte et honóre putáti. Nonne vidétis quanta promptitúdine párvuli papíllas cápiunt, et quanto ímpetu lábia ubéribus infígunt? Accedámus cum tanta nos quoque alacritáte ad hanc mensam et ad úbera póculi spiritális: quinímmo cum longe majóri trahámus, tamquam infántes lacténtes, spíritus grátiam: et unus sit nobis dolor hac esca privári. Non sunt humánæ virtútis ópera, hæc quæ proponúntur: qui tunc ipsa fecit in illa cœna, idem ea nunc quoque facit. Nos ministrórum tenémus locum; qui vero sanctíficat ea et immútat, ipse est. Nullus ítaque Judas assístat, nullus avárus; nam tales mensa non súscipit. Si quis est discípulus, adsit; ait enim: Cum discípulis meis fácio Pascha. Hæc est illa mensa, et minus nihil habet. Non enim illam quidem Christus, hanc autem homo pérficit; verum et hanc ipse quoque.\par
\switchcolumn
\EN
\noindent In this mysterious Sacrament Christ doth mingle Himself with all and each of His faithful ones. They are His children, and He nurseth them Himself, and giveth them not over unto another, herein again assuring us that the Flesh He hath taken unto Himself is ours. We then, who have been deemed meet to be treated with such love and such honour, let us be wakeful See ye not how eagerly the sucklings seize on the breasts, how readily they fix their mouths on the paps Let us, with like eagerness, draw nigh to that Table, and suck at that spiritual Cup. Yea, let us prize that gracious Food as the suckling doth its mother's breast, and hold it the great woe of life to be cut off from that Banquet. Here there are set before us no works of man's power He That worked at that Last Supper, the Same worketh the same here still. As for us Priests, we hold the place of His ministers, but He Which halloweth and changeth is He. Hither let there draw nigh no Judas, nor covetous one this is no Table for him. But he which is Christ's disciple, let him come for the Lord saith “I will keep the Passover with My disciples,” (Matth. xxvi. 18.) This is that Passover Table, and it is all Christ's what is wrought there is not some of it Christ's work, and some of it man's work, but it is all His work and not another's.\par
\switchcolumn*

\LA
\begin{Resp}\Rsign Cenántibus illis, accépit Jesus panem, et benedíxit, ac fregit, dedítque discípulis suis, et ait: \Star{} Accípite et comédite: hoc est corpus meum.\end{Resp}
\begin{Resp}\Vsign Dixérunt viri tabernáculi mei: Quis det de cárnibus ejus, ut saturémur?\end{Resp}
\begin{Resp}\Rsign Accípite et comédite: hoc est corpus meum.\end{Resp}
\switchcolumn
\EN
\begin{Resp}\Rsign As they were eating, Jesus took bread, and blest it, and brake it, and gave it to the disciples, and said: \Star{} Take, eat this is My Body.\end{Resp}
\begin{Resp}\Vsign The men of my tabernacle said: O that we had of his flesh, we cannot be satisfied.\end{Resp}
\begin{Resp}\Rsign Take, eat this is My Body.\end{Resp}
\switchcolumn*

\LA
\LessonHead{Lectio V}
\switchcolumn
\EN
\LessonHead{Lesson V}
\switchcolumn*

\LA
\noindent Inhumánus accédat nemo, nemo crudélis et immiséricors, nemo prorsus immúndus. Hæc ad communicántes dico, et ad vos, ministrántes. Nam et ad vos sermónem convértere necessárium est, ut multo cum stúdio hæc dona distribuátis. Non parva vobis ímminet últio, si quemquam, illíus culpæ cónscii, hujus mensæ partícipem esse concedátis: sanguis ejus de mánibus vestris exquirétur. Sive quis dux milítiæ sit, sive præféctus, sive princeps diadémate coronátus, indígne autem accédat, próhibe: majórem illo potestátem habes. Proptérea vos Deus hoc insignívit honóre, ut tália discernátis. Hoc vestra dígnitas est, hoc secúritas, hoc omnis coróna; non ut albam et splendéntem túnicam circumeátis indúti. Verum et tu, láice, cum sacerdótem víderis offeréntem, ne ut sacerdótem esse putes hoc faciéntem, sed Christi manum invisibíliter exténsam.\par
\switchcolumn
\EN
\noindent Wither let there draw nigh none brutal, none cruel, none merciless in good sooth, none unclean. I speak to all that take that Holy Communion, and to you also, O ye that do administer the same To you now I turn my speech, to warn you with how great care that Gift is to be given. No slight vengeance is that which awaiteth you if ye admit for a partaker at the Lord's Table the sinner whose guiltiness ye know. At your hands will his blood be required. If a man be a General, a Governor, a crowned Monarch, yet if he come there unworthily, forbid him thou hast greater power than he. To this end hath God exalted you to the honour ye hold, that ye may judge in such matters. This office is your dignity, this is your strength, this is all your crown, this, and not the going about in white robes and glittering vestments. And thou, O layman when thou seest the Priest making the oblation, think not that He Which is then the real Worker is such a Priest as thou seest, but know of a surety that it is Christ's Hand Which is stretched out, albeit unseen by thee.\par
\switchcolumn*

\LA
\begin{Resp}\Rsign Accépit Jesus cálicem, postquam cenávit, dicens: Hic calix novum testaméntum est in meo sánguine: \Star{} Hoc fácite in meam commemoratiónem.\end{Resp}
\begin{Resp}\Vsign Memória memor ero, et tabéscet in me ánima mea.\end{Resp}
\begin{Resp}\Rsign Hoc fácite in meam commemoratiónem.\end{Resp}
\switchcolumn
\EN
\begin{Resp}\Rsign Jesus took the cup, after supper, saying: This cup is the New Testament in My Blood. \Star{} This do in remembrance of Me.\end{Resp}
\begin{Resp}\Vsign My soul hath them still in remembrance, and is humbled in me.\end{Resp}
\begin{Resp}\Rsign This do in remembrance of Me.\end{Resp}
\switchcolumn*

\LA
\LessonHead{Lectio VI}
\switchcolumn
\EN
\LessonHead{Lesson VI}
\switchcolumn*

\LA
\noindent Audiámus ígitur, et sacerdótes et súbditi, quali esca facti sumus digni: audiámus et horreámus. Sanctis cárnibus suis nos dedit impléri, semetípsum appósuit immolátum. Quænam ígitur erit nobis excusátio, cum tálibus pasti tália peccémus, cum lupi fiámus Agnum comedéntes, cum tamquam oves pasti more leónum diripiámus? Hoc enim mystérium non a rapína tantum, verum et ab omni vel ténui inimicítia purum esse pénitus jubet; est enim pacis mystérium. Judǽis quidem annuátim propriórum monuménta beneficiórum solemnitátes Deus alligávit; tibi vero síngulis diébus per hæc mystéria. Nullus ítaque Judas hanc mensam petat, nullus Simon. Hi namque duo propter avarítiam periérunt: hoc ígitur bárathrum fugiámus.\par
\switchcolumn
\EN
\noindent Let us hear, all of us, both Priests and laymen, let us hear What Food it is whereof we are made worthy let us hear, I say, and let us quake. The Lord satisfieth us with His Own holy Flesh, setting Himself slain before us. What excuse therefore shall we have, if, being so fed as we are, we sin as we do If, eating of the Lamb, we are still wolves If, pastured as the sheep of the flock, we raven like lions This mysterious Sacrament forbiddeth unto us not outrage only, but any the least enmity it is the Mystery of peace. Upon the Jews God laid it to make year by year by solemn festivals a yearly commemoration of His mercies unto them, but upon thee to do this in remembrance of His love to thee, day by day. To this Table then let there draw nigh no Judas Iscariot, no Simon Magus. These men fell through covetousness let us fly that bottomless pit.\par
\switchcolumn*

\LA
\begin{Resp}\Rsign Ego sum panis vitæ; patres vestri manducavérunt manna in desérto, et mórtui sunt: \Star{} Hic est panis de cælo descéndens, ut, si quis ex ipso mandúcet, non moriátur.\end{Resp}
\begin{Resp}\Vsign Ego sum panis vivus, qui de cælo descéndi: si quis manducáverit ex hoc pane, vivet in ætérnum.\end{Resp}
\begin{Resp}\Rsign Hic est panis de cælo descéndens, ut, si quis ex ipso mandúcet, non moriátur.\end{Resp}
\begin{Resp}\Vsign Glória Patri, et Fílio, \Star{} et Spirítui Sancto.\end{Resp}
\begin{Resp}\Rsign Hic est panis de cælo descéndens, ut, si quis ex ipso mandúcet, non moriátur.\end{Resp}
\switchcolumn
\EN
\begin{Resp}\Rsign I am that Bread of life. Your fathers did eat manna in the wilderness, and are dead. \Star{} This is the Bread Which cometh down from heaven, that a man may eat thereof, and not die.\end{Resp}
\begin{Resp}\Vsign I am the living Bread Which came down from heaven if any man eat of this Bread, he shall live for ever.\end{Resp}
\begin{Resp}\Rsign This is the Bread Which cometh down from heaven, that a man may eat thereof, and not die.\end{Resp}
\begin{Resp}\Vsign Glory be to the Father, and to the Son, \Star{} and to the Holy Ghost.\end{Resp}
\begin{Resp}\Rsign This is the Bread Which cometh down from heaven, that a man may eat thereof, and not die.\end{Resp}
\switchcolumn*

\LA
\LessonHead{Lectio VII}
\switchcolumn
\EN
\LessonHead{Lesson VII}
\switchcolumn*

\LA
\LessonTitle{Léctio sancti Evangélii secúndum Joánnem}
\switchcolumn
\EN
\LessonTitle{From the Holy Gospel according to John}
\switchcolumn*

\LA
\Cite{Joannes 6:56-59}
\switchcolumn
\EN
\Cite{John 6:56-59}
\switchcolumn*

\LA
\noindent In illo témpore: Dixit Jesus turbis Judæórum: Caro mea vere est cibus, et sanguis meus vere est potus. Et réliqua.\par
\switchcolumn
\EN
\noindent At that time, Jesus said unto the multitudes of the Jews: My Flesh is meat indeed, and My Blood is drink indeed. And so on.\par
\switchcolumn*

\LA
\LessonTitle{De Homilía sancti Augustíni Epíscopi}
\switchcolumn
\EN
\LessonTitle{Homily by St. Augustine, Bishop of Hippo.}
\switchcolumn*

\LA
\Source{Tract. 26. in Joánnem, post medium}
\switchcolumn
\EN
\Source{26th Tract on John.}
\switchcolumn*

\LA
\noindent Hic est panis, qui de cælo descéndit. Hunc panem significávit manna, hunc panem significávit altáre Dei. Sacraménta illa fuérunt: in signis divérsa sunt; sed in re, quæ significátur, pária sunt. Apóstolum audi: Nolo enim vos, inquit, ignoráre, fratres, quia patres nostri omnes sub nube fuérunt, et omnes mare transiérunt, et omnes per Móysen baptizáti sunt in nube et in mari, et omnes eámdem escam spiritálem manducavérunt. Spiritálem útique eámdem; nam corporálem álteram; quia illi manna, nos áliud: spiritálem vero, quam nos, sed patres nostri, non patres illórum, quibus nos símiles sumus, non quibus illi símiles fuérunt. Et adjúngit: Et omnes eúmdem potum spiritálem bibérunt. Aliud illi, áliud nos, sed spécie visíbili quidem, tamen hoc idem significánte virtúte spiritáli. Quómodo enim eúmdem potum? Bibébant, inquit, de spiritáli, sequénte petra: petra autem erat Christus. Inde panis, inde potus. Petra Christus in signo, verus Christus in verbo et in carne. Et quómodo bibérunt? Percússa est petra de virga bis: gémina percússio, duo ligna crucis signíficat.\par
\switchcolumn
\EN
\noindent “This is the bread which cometh down from heaven,” (v. 50.) By “this bread” the Lord here signifieth both the manna, and That Which we receive at the Altar of God. Both these are, as it were, Sacramental signs, differing indeed somewhat in their outward and visible part, but pointing to the Same Thing signified. Hear what the Apostle saith: “Moreover, brethren, I would not that ye should be ignorant how that all our fathers were under the cloud, and all passed through the sea, and were all baptized unto Moses in the cloud and in the sea, and did all eat the same spiritual meat.” (1 Cor. x. 1-3.) This meat was the same spiritually but not really they ate manna we eat Something else. Spiritually they ate What we eat but our fathers not their fathers; unto whom we are like not unto whom they are like. And it is added “And did all drink the same Spiritual drink.” They drank one thing, and we drank Another, the difference being in the outer show, the sameness in that the Same Thing is pointed to by both. And what was that Same Drink? “They drank of the spiritual Rock that followed them, and that Rock was Christ.” Him did bread and rock alike signify. The Rock was a figure, but by the Word and in the Flesh there is the very Christ Himself. And how came they to drink of that rock “Moses lift up his hand, and with his rod he smote the rock twice, and the water came out abundantly.” (Num. xx. 11.) These two strokes of the rod upon the rock are a figure of the two beams whereof the Cross was made.\par
\switchcolumn*

\LA
\begin{Resp}\Rsign Qui mandúcat meam carnem et bibit meum sánguinem, \Star{} In me manet, et ego in eo.\end{Resp}
\begin{Resp}\Vsign Non est ália nátio tam grandis, quæ hábeat deos appropinquántes sibi, sicut Deus noster adest nobis.\end{Resp}
\begin{Resp}\Rsign In me manet, et ego in eo.\end{Resp}
\switchcolumn
\EN
\begin{Resp}\Rsign He that eateth My Flesh and drinketh My Blood, \Star{} Dwelleth in Me, and I in him.\end{Resp}
\begin{Resp}\Vsign What nation is there so great, who hath gods so nigh unto them, as the Lord our God is to us\end{Resp}
\begin{Resp}\Rsign Dwelleth in Me, and I in him.\end{Resp}
\switchcolumn*

\LA
\LessonHead{Lectio VIII}
\switchcolumn
\EN
\LessonHead{Lesson VIII}
\switchcolumn*

\LA
\noindent Norunt fidéles corpus Christi, si corpus Christi non négligant esse. Fiant corpus Christi, si volunt vívere de Spíritu Christi. De Spíritu Christi non vivit, nisi corpus Christi. Intellígite, fratres mei, quid díxerim. Homo es, et spíritum habes, et corpus habes. Spíritum dico, quæ ánima vocátur, qua constas quod homo es; constas enim ex ánima et córpore. Habes enim spíritum invisíbilem, corpus visíbile. Dic mihi, quid ex quo vivat? Spíritus tuus vivit ex córpore tuo, an corpus tuum ex spíritu tuo? Respóndet omnis, qui vivit: qui autem hoc non potest respondére, néscio si vivit. Quid respóndet omnis, qui vivit? Corpus útique meum vivit de spíritu meo. Vis ergo et tu vívere de Spíritu Christi? In córpore esto Christi.\par
\switchcolumn
\EN
\noindent Christ's faithful ones discern the Lord's Body while they remain watchful members of His Body. They remain members of His Body as long as they will to live according to His Spirit. The Spirit of Christ giveth life to nothing but the body of Christ. Now, my brethren, understand what I am going to say. Thou art a man, and hast a body and a spirit. By spirit I mean the soul, which causeth thee to be a man at all. Thou art a man, made up of soul and body. thy spirit is unseen, thy body seen. Tell me, which of them is it which giveth animation to the other Doth thy spirit derive animation from thy body, or thy body from thy spirit Every one who liveth will answer for if any one cannot answer this, I know not if he be alive. What will whosoever hath life answer: “Verily, it is my spirit which doth animate my body.” Wilt thou then live by the Spirit of Christ? Be of the Body of Christ.\par
\switchcolumn*

\LA
\begin{Resp}\Rsign Misit me vivens Pater, et ego vivo propter Patrem: \Star{} Et qui mandúcat me, vivet propter me.\end{Resp}
\begin{Resp}\Vsign Cibávit illum Dóminus pane vitæ et intelléctus.\end{Resp}
\begin{Resp}\Rsign Et qui mandúcat me, vivet propter me.\end{Resp}
\begin{Resp}\Vsign Glória Patri, et Fílio, \Star{} et Spirítui Sancto.\end{Resp}
\begin{Resp}\Rsign Et qui mandúcat me, vivet propter me.\end{Resp}
\switchcolumn
\EN
\begin{Resp}\Rsign As the living Father hath sent Me, and I live by the Father, \Star{} So he that eateth Me, even he shall live by Me.\end{Resp}
\begin{Resp}\Vsign With the bread of life and understanding hath the Lord fed him.\end{Resp}
\begin{Resp}\Rsign So he that eateth Me, even he shall live by Me.\end{Resp}
\begin{Resp}\Vsign Glory be to the Father, and to the Son, \Star{} and to the Holy Ghost.\end{Resp}
\begin{Resp}\Rsign So he that eateth Me, even he shall live by Me.\end{Resp}
\switchcolumn*

\LA
\LessonHead{Lectio IX}
\switchcolumn
\EN
\LessonHead{Lesson IX}
\switchcolumn*

\LA
\noindent Numquid enim corpus meum vivit de spíritu tuo? Meum vivit de spíritu meo, et tuum de spíritu tuo. Non potest vívere corpus Christi, nisi de Spíritu Christi. Inde est quod expónens nobis Apóstolus Paulus hunc panem: Unus panis, inquit, unum corpus multi sumus. O sacraméntum pietátis, o signum unitátis, o vínculum caritátis! Qui vult vívere, habet ubi vivat, habet unde vivat. Accédat, credat, incorporétur, ut vivificétur. Non abhórreat a compáge membrórum, non sit putre membrum quod resecári mereátur, non sit distórtum de quo erubescátur. Sit pulchrum, sit aptum, sit sanum: hǽreat córpori, vivat Deo de Deo. Nunc labóret in terra, ut póstea regnet in cælo.\par
\switchcolumn
\EN
\noindent Is it not my spirit which doth animate my body My spirit doth animate my body, and thy spirit doth animate thy body. The Body of Christ liveth not save by the Spirit of Christ. Hence it is that the Apostle Paul saith, touching this Bread: “We, being many, are one bread, and one body, for we are all partakers of that one Bread.” (1 Cor. x. 17.) O what a Sacrament of love O what a seal of union O what a bond of charity He that willeth to live hath here where to live, and whence to live. Let him come near, let him believe, let him enter into that Body, that he may be quickened. Let him not sever himself from the fit joining-together of all the members let him not be as a mortifying limb, that must needs be cut off, nor a mis-shapen limb, a cause to blush. Let him be goodly, and useful, and healthy. Let him cleave unto the body let him live by God to God let him labour now on earth, that he may reign hereafter in heaven.\par
\switchcolumn*

\LA
\TeDeum{Te Deum}
\switchcolumn
\EN
\TeDeum{Te Deum}
\switchcolumn*

\end{paracol}

\Office{Feria Tertia infra Octavam Corporis Christi}{Tuesday within the Octave of Corpus Christi}{Semiduplex II classis}
\addcontentsline{toc}{subsection}{Feria Tertia infra Octavam Corporis Christi \textit{— Tuesday within the Octave of Corpus Christi}}
\begin{paracol}{2}
\LA
\LessonHead{Lectio I}
\switchcolumn
\EN
\LessonHead{Lesson I}
\switchcolumn*

\LA
\LessonTitle{De libro primo Regum}
\switchcolumn
\EN
\LessonTitle{Lesson from the first book of Samuel}
\switchcolumn*

\LA
\Cite{1 Reg 6:1-3}
\switchcolumn
\EN
\Cite{1 Sam 6:1-3}
\switchcolumn*

\LA
\noindent Fuit ergo arca Dómini in regióne Philisthinórum septem ménsibus, et vocavérunt Philísthiim sacerdótes et divínos, dicéntes: Quid faciémus de arca Dómini? Indicáte nobis quómodo remittámus eam in locum suum. Qui dixérunt: Si remíttitis arcam Dei Israël, nolíte dimíttere eam vácuam, sed quod debétis réddite ei pro peccáto: et tunc curabímini et sciétis quare non recédat manus ejus a vobis.\par
\switchcolumn
\EN
\noindent Now the ark of God was in the land of the Philistines seven months. And the Philistines called for the priests and the diviners, saying: What shall we do with the ark of the Lord? tell us how we are to send it back to its place? And they said: If you send back the ark of the God of Israel, send it not away empty, but render unto him what you owe for sin, and then you shall be healed: and you shall know why his hand departeth not from you\par
\switchcolumn*

\LA
\begin{Resp}\Rsign Immolábit hædum multitúdo filiórum Israël ad vésperam Paschæ: \Star{} Et edent carnes et ázymos panes.\end{Resp}
\begin{Resp}\Vsign Pascha nostrum immolátus est Christus: ítaque epulémur in ázymis sinceritátis et veritátis.\end{Resp}
\begin{Resp}\Rsign Et edent carnes et ázymos panes.\end{Resp}
\switchcolumn
\EN
\begin{Resp}\Rsign The whole assembly of the children of Israel shall kill the lamb toward the evening of the Passover. \Star{} And they shall eat the flesh, and unleavened bread.\end{Resp}
\begin{Resp}\Vsign Even Christ our Passover is sacrificed for us therefore let us keep the feast with the unleavened bread of sincerity and truth.\end{Resp}
\begin{Resp}\Rsign And they shall eat the flesh, and unleavened bread.\end{Resp}
\switchcolumn*

\LA
\LessonHead{Lectio II}
\switchcolumn
\EN
\LessonHead{Lesson II}
\switchcolumn*

\LA
\Cite{1 Reg 6:6-10}
\switchcolumn
\EN
\Cite{1 Sam 6:6-10}
\switchcolumn*

\LA
\noindent Quare aggravátis corda vestra, sicut aggravávit Ægýptus et phárao cor suum? Nonne, postquam percússus est, tunc dimísit eos, et abiérunt? Nunc ergo arrípite et fácite plaustrum novum unum, et duas vaccas fœtas, quibus non est impósitum jugum, júngite in plaustro et reclúdite vítulos eárum domi. Tolletísque arcam Dómini et ponétis in plaustro et vasa áurea, quæ exsolvétis ei pro delícto, ponétis in capséllam ad latus ejus, et dimíttite eam ut vadat. Et aspiciétis, et, si quidem per viam fínium suórum ascénderit contra Béthsames ipse fecit nobis hoc malum grande; sin autem mínime, sciémus quia nequáquam manus ejus tétigit nos, sed casu áccidit. Fecérunt ergo illi hoc modo.\par
\switchcolumn
\EN
\noindent Why do you harden your hearts, as Egypt and Pharao hardened their hearts? did not he, after he was struck, then let them go, and they departed? Now therefore take and make a new cart: and two kine that have calved, on which there hath come no yoke, tie to the cart, and shut up their calves at home. And you shall take the ark of the Lord, and lay it on the cart, and the vessels of gold, which you have paid him for sin, you shall put into a little box, at the side thereof: and send it away that it may go. And you shall look: and if it go up by the way of his own coasts towards Bethsames, then he hath done us this great evil: but if not, we shall know that it is not his hand hath touched us, but it hath happened by chance. They did therefore in this manner.\par
\switchcolumn*

\LA
\begin{Resp}\Rsign Comedétis carnes, et saturabímini pánibus: \Star{} Iste est panis, quem dedit vobis Dóminus ad vescéndum.\end{Resp}
\begin{Resp}\Vsign Non Móyses dedit vobis panem de cælo, sed Pater meus dat vobis panem de cælo verum.\end{Resp}
\begin{Resp}\Rsign Iste est panis, quem dedit vobis Dóminus ad vescéndum.\end{Resp}
\switchcolumn
\EN
\begin{Resp}\Rsign Ye shall eat flesh, and shall be filled with bread. \Star{} This is the bread which the Lord hath given you to eat.\end{Resp}
\begin{Resp}\Vsign Moses gave you not that Bread from heaven, but My Father giveth you the true Bread from heaven.\end{Resp}
\begin{Resp}\Rsign This is the bread which the Lord hath given you to eat.\end{Resp}
\switchcolumn*

\LA
\LessonHead{Lectio III}
\switchcolumn
\EN
\LessonHead{Lesson III}
\switchcolumn*

\LA
\Cite{1 Reg 6:12-15}
\switchcolumn
\EN
\Cite{1 Sam 6:12-15}
\switchcolumn*

\LA
\noindent Ibant autem in diréctum vaccæ per viam, quæ ducit Béthsames, et itínere uno gradiebántur pergéntes et mugiéntes et non declinábant neque ad déxteram neque ad sinístram; sed et sátrapæ Philísthiim sequebántur usque ad términos Béthsames. Porro Bethsamítæ metébant tríticum in valle, et elevántes óculos suos vidérunt arcam et gavísi sunt cum vidíssent. Et plaustrum venit in agrum Jósue Bethsamítæ et stetit ibi. Erat autem ibi lapis magnus, et concidérunt ligna plaustri, vaccásque imposuérunt super ea holocáustum Dómino. Levítæ autem deposuérunt arcam Dei.\par
\switchcolumn
\EN
\noindent And the kine took the straight way that leadeth to Bethsames, and they went along the way, lowing as they went: and turned not aside neither to the right hand nor to the left: and the lords of the Philistines followed them as far as the borders of Bethsames. Now the Bethsamites were reaping wheat in the valley: and lifting up their eyes they saw the ark, and rejoiced to see it. And the cart came into the field of Josue a Bethsamite, and stood there. And there was a great stone, and they cut in pieces the wood of the cart, and laid the kine upon it a holocaust to the Lord. And the Levites took down the ark of God.\par
\switchcolumn*

\LA
\begin{Resp}\Rsign Respéxit Elías ad caput suum subcinerícium panem: qui surgens comédit et bibit: \Star{} Et ambulávit in fortitúdine cibi illíus usque ad montem Dei.\end{Resp}
\begin{Resp}\Vsign Si quis manducáverit ex hoc pane, vivet in ætérnum.\end{Resp}
\begin{Resp}\Rsign Et ambulávit in fortitúdine cibi illíus usque ad montem Dei.\end{Resp}
\begin{Resp}\Vsign Glória Patri, et Fílio, \Star{} et Spirítui Sancto.\end{Resp}
\begin{Resp}\Rsign Et ambulávit in fortitúdine cibi illíus usque ad montem Dei.\end{Resp}
\switchcolumn
\EN
\begin{Resp}\Rsign Elijah looked, and, behold, there was a cake baken on the coals at his head, and he arose, and did eat and drink; \Star{} And went in the strength of that meat [forty days and forty nights] unto the mount of God.\end{Resp}
\begin{Resp}\Vsign If any man eat of this Bread, he shall live for ever.\end{Resp}
\begin{Resp}\Rsign And went in the strength of that meat [forty days and forty nights] unto the mount of God.\end{Resp}
\begin{Resp}\Vsign Glory be to the Father, and to the Son, \Star{} and to the Holy Ghost.\end{Resp}
\begin{Resp}\Rsign And went in the strength of that meat [forty days and forty nights] unto the mount of God.\end{Resp}
\switchcolumn*

\LA
\LessonHead{Lectio IV}
\switchcolumn
\EN
\LessonHead{Lesson IV}
\switchcolumn*

\LA
\LessonTitle{Ex Epístola sancti Cypriáni Epíscopi et Mártyris ad Cæcílium}
\switchcolumn
\EN
\LessonTitle{From the Letter written to Caecilius by the Holy Martyr Cyprian, Bishop of Carthage.}
\switchcolumn*

\LA
\Source{Liber 2. Epistola 3. sub init.}
\switchcolumn
\EN
\Source{Bk. ii. Ep. 3}
\switchcolumn*

\LA
\noindent In sacerdóte Melchísedech sacrifícii Domínici sacraméntum præfigurátum vidémus, secúndum quod Scriptúra divína testátur, et dicit: Et Melchísedech rex Salem prótulit panem et vinum. Fuit autem sacérdos Dei summi, et benedíxit Abraham. Quod autem Melchísedech typum Christi portáret, declárat in Psalmis Spíritus Sanctus, ex persóna Patris ad Fílium dicens: Ante lucíferum génui te: Tu es sacérdos in ætérnum secúndum órdinem Melchísedech. Qui ordo útique hic est, de sacrifício illo véniens et inde descéndens, quod Melchísedech sacérdos Dei summi fuit, quod panem et vinum óbtulit, quod Abraham benedíxit.\par
\switchcolumn
\EN
\noindent In the deed of the Priest Melchisedek we see a type of the Sacrament of the Lord's Sacrifice. For thus it is written in the writings of God “And Melchisedek King of Salem brought forth bread and wine for he was the Priest of the Most High God and he blessed Abraham.” (Gen. xiv. 18.) That Melchisedek was a type of Christ, the Holy Ghost Himself doth testify in the Psalms, where the First Person of the Holy Trinity, even the Father, is set before us as saying unto the Second Person, that is, the Son “Before the daystar have I begotten thee. Thou art a Priest for ever, after the order of Melchisedek.” (cix. 3, 4.) And verily that sameness of order cometh of this sacrifice, and procedeth from this, that Melchisedek was the Priest of the Most High God that he offered bread and wine and that he blessed Abraham.\par
\switchcolumn*

\LA
\begin{Resp}\Rsign Cenántibus illis, accépit Jesus panem, et benedíxit, ac fregit, dedítque discípulis suis, et ait: \Star{} Accípite et comédite: hoc est corpus meum.\end{Resp}
\begin{Resp}\Vsign Dixérunt viri tabernáculi mei: Quis det de cárnibus ejus, ut saturémur?\end{Resp}
\begin{Resp}\Rsign Accípite et comédite: hoc est corpus meum.\end{Resp}
\switchcolumn
\EN
\begin{Resp}\Rsign As they were eating, Jesus took bread, and blest it, and brake it, and gave it to the disciples, and said: \Star{} Take, eat this is My Body.\end{Resp}
\begin{Resp}\Vsign The men of my tabernacle said: O that we had of his flesh, we cannot be satisfied.\end{Resp}
\begin{Resp}\Rsign Take, eat this is My Body.\end{Resp}
\switchcolumn*

\LA
\LessonHead{Lectio V}
\switchcolumn
\EN
\LessonHead{Lesson V}
\switchcolumn*

\LA
\noindent Nam quis magis sacérdos Dei summi, quam Dóminus noster Jesus Christus? qui sacrifícium Deo Patri óbtulit; et óbtulit hoc idem, quod Melchísedech obtúlerat, id est, panem et vinum, suum scílicet corpus et sánguinem. Et circa Abraham benedíctio illa præcédens, ad nostrum pópulum pertinébat. Nam si Abraham Deo crédidit, et deputátum est ei ad justítiam; útique quisquis Deo credit, et fide vivit, justus invenítur, et jam pridem in Abraham fidéli benedíctus et justificátus osténditur, sicut beátus Apóstolus Paulus probat, dicens: Crédidit Abraham Deo, et deputátum est ei ad justítiam. Cognóscitis ergo quia qui ex fide sunt, hi sunt fílii Abrahæ. Próvidens autem Scriptúra, quia ex fide justíficat gentes Deus, prænuntiávit Abrahæ quia benedicéntur in illo omnes gentes.\par
\switchcolumn
\EN
\noindent What Priest of the Most High God is there, more so than is our Lord Jesus Christ He Who hath made an offering unto God the Father, and the same offering that Melchisedek made, bread and wine, that is to say, His Own Flesh and His Own Blood. And, as touching Abraham, that ancient blessing was spoken likewise by fore-knowledge upon us. For if Abraham believed God and it was accounted to him for righteousness, verily, whosoever believeth God and liveth by faith, the same is found righteous, and is shown unto us that he is already blessed in faithful Abraham, and justified as the Apostle Paul proveth, where he saith: “Abraham believed God and it was accounted unto him for righteousness. Know ye therefore that they which are of faith, the same are the children of Abraham. And the Scripture, foreseeing that God would justify the heathen through faith, preached before the Gospel unto Abraham, saying In thee shall all nations be blessed.” (Gal. iii. 6-8.)\par
\switchcolumn*

\LA
\begin{Resp}\Rsign Accépit Jesus cálicem, postquam cenávit, dicens: Hic calix novum testaméntum est in meo sánguine: \Star{} Hoc fácite in meam commemoratiónem.\end{Resp}
\begin{Resp}\Vsign Memória memor ero, et tabéscet in me ánima mea.\end{Resp}
\begin{Resp}\Rsign Hoc fácite in meam commemoratiónem.\end{Resp}
\switchcolumn
\EN
\begin{Resp}\Rsign Jesus took the cup, after supper, saying: This cup is the New Testament in My Blood. \Star{} This do in remembrance of Me.\end{Resp}
\begin{Resp}\Vsign My soul hath them still in remembrance, and is humbled in me.\end{Resp}
\begin{Resp}\Rsign This do in remembrance of Me.\end{Resp}
\switchcolumn*

\LA
\LessonHead{Lectio VI}
\switchcolumn
\EN
\LessonHead{Lesson VI}
\switchcolumn*

\LA
\noindent Ut ergo in Génesi per Melchísedech sacerdótem posset rite celebrári, præcédit ante imágo sacrifícii, in pane et vino scílicet constitúta. Quam rem perfíciens et adímplens Dóminus, panem et cálicem mixtum vino óbtulit; et, qui est plenitúdo, veritátem præfigurátæ imáginis adimplévit. Sed et per Salomónem Spíritus Sanctus typum Domínici sacrifícii ante præmónstrat, immolátæ hóstiæ et panis et vini, sed et altáris et Apostolórum fáciens mentiónem: Sapiéntia, inquit, ædificávit sibi domum, et súbdidit colúmnas septem; mactávit suas hóstias, míscuit in cratére vinum suum, et parávit mensam suam. Et misit servos suos cónvocans cum excélsa prædicatióne ad cratérem, dicens: Qui est insípiens, declínet ad me. Et egéntibus sensu dixit: Veníte, édite de meis pánibus, et bíbite vinum quod míscui vobis.\par
\switchcolumn
\EN
\noindent In Genesis, therefore, in order that the Priest Melchisedek might in due order pronounce the blessing upon Abraham, there was first offered a typical sacrifice, consisting of bread and wine. This was the offering which our Lord Jesus Christ completed and fulfilled, when He offered up bread and a cup of wine mingled with water. This fulfilment by Him Who came to fulfill (Matth. v. 17,) utterly satisfied the truth of the image which had gone before. The Holy Ghost doth by Solomon also clearly foreshadow, as it were in a parable, the Lord's Sacrifice, pointing to the victim slain, and the bread and the wine, and the Altar likewise, and the Apostles as it is written “Wisdom hath builded her house, she hath hewn out her seven pillars she hath killed her beasts, she hath mingled her wine, she hath also furnished her table. She hath sent forth her servants, she crieth upon the highest places of the city, saying Whoso is simple, let him turn in hither unto me. As for them that want understanding, she saith to them. Come, eat of my bread, and drink of the wine which I have mingled for you.” (Prov. ix. 1-5.)\par
\switchcolumn*

\LA
\begin{Resp}\Rsign Ego sum panis vitæ; patres vestri manducavérunt manna in desérto, et mórtui sunt: \Star{} Hic est panis de cælo descéndens, ut, si quis ex ipso mandúcet, non moriátur.\end{Resp}
\begin{Resp}\Vsign Ego sum panis vivus, qui de cælo descéndi: si quis manducáverit ex hoc pane, vivet in ætérnum.\end{Resp}
\begin{Resp}\Rsign Hic est panis de cælo descéndens, ut, si quis ex ipso mandúcet, non moriátur.\end{Resp}
\begin{Resp}\Vsign Glória Patri, et Fílio, \Star{} et Spirítui Sancto.\end{Resp}
\begin{Resp}\Rsign Hic est panis de cælo descéndens, ut, si quis ex ipso mandúcet, non moriátur.\end{Resp}
\switchcolumn
\EN
\begin{Resp}\Rsign I am that Bread of life. Your fathers did eat manna in the wilderness, and are dead. \Star{} This is the Bread Which cometh down from heaven, that a man may eat thereof, and not die.\end{Resp}
\begin{Resp}\Vsign I am the living Bread Which came down from heaven if any man eat of this Bread, he shall live for ever.\end{Resp}
\begin{Resp}\Rsign This is the Bread Which cometh down from heaven, that a man may eat thereof, and not die.\end{Resp}
\begin{Resp}\Vsign Glory be to the Father, and to the Son, \Star{} and to the Holy Ghost.\end{Resp}
\begin{Resp}\Rsign This is the Bread Which cometh down from heaven, that a man may eat thereof, and not die.\end{Resp}
\switchcolumn*

\LA
\LessonHead{Lectio VII}
\switchcolumn
\EN
\LessonHead{Lesson VII}
\switchcolumn*

\LA
\LessonTitle{Léctio sancti Evangélii secúndum Joánnem}
\switchcolumn
\EN
\LessonTitle{From the Holy Gospel according to John}
\switchcolumn*

\LA
\Cite{Joannes 6:56-59}
\switchcolumn
\EN
\Cite{John 6:56-59}
\switchcolumn*

\LA
\noindent In illo témpore: Dixit Jesus turbis Judæórum: Caro mea vere est cibus, et sanguis meus vere est potus. Et réliqua.\par
\switchcolumn
\EN
\noindent At that time: Jesus said unto the multitudes of the Jews: “My Flesh is meat indeed, and My Blood is drink indeed.” And so on.\par
\switchcolumn*

\LA
\LessonTitle{De Homilía sancti Augustíni Epíscopi}
\switchcolumn
\EN
\LessonTitle{Homily by St. Augustine, Bishop of Hippo.}
\switchcolumn*

\LA
\Source{Tractatus 26 in Joánnem, circa medium}
\switchcolumn
\EN
\Source{26/A Tract on John.}
\switchcolumn*

\LA
\noindent Non sicut manducavérunt patres vestri manna, et mórtui sunt. Quare manducavérunt, et mórtui sunt? Quia quod vidébant, credébant; quod non vidébant , non intelligébant. Ideo patres vestri, quia símiles estis illórum. Nam quantum pértinet, fratres mei, ad mortem istam visíbilem et corporálem, numquid nos non mórimur, qui manducámus panem de cælo descendéntem? Sic sunt mórtui et illi, quemádmodum et nos sumus moritúri; quantum áttinet, ut dixi, ad mortem hujus córporis visíbilem atque carnálem.\par
\switchcolumn
\EN
\noindent “Not as your fathers did eat manna, and are dead.” Wherefore did they eat and die Because they believed only that which they saw and that which they saw not, they understood not. Therefore were they your fathers, because ye are like unto them. Doth this death, my brethren, mean that death which is outward and bodily And do not we also die, who eat of that Bread Which cometh down from heaven That death died they, and so shall we also, as far, as I have said, as is meant that death which is outward and bodily.\par
\switchcolumn*

\LA
\begin{Resp}\Rsign Qui mandúcat meam carnem et bibit meum sánguinem, \Star{} In me manet, et ego in eo.\end{Resp}
\begin{Resp}\Vsign Non est ália nátio tam grandis, quæ hábeat deos appropinquántes sibi, sicut Deus noster adest nobis.\end{Resp}
\begin{Resp}\Rsign In me manet, et ego in eo.\end{Resp}
\switchcolumn
\EN
\begin{Resp}\Rsign He that eateth My Flesh and drinketh My Blood, \Star{} Dwelleth in Me, and I in him.\end{Resp}
\begin{Resp}\Vsign What nation is there so great, who hath gods so nigh unto them, as the Lord our God is to us\end{Resp}
\begin{Resp}\Rsign Dwelleth in Me, and I in him.\end{Resp}
\switchcolumn*

\LA
\LessonHead{Lectio VIII}
\switchcolumn
\EN
\LessonHead{Lesson VIII}
\switchcolumn*

\LA
\noindent Quantum autem pértinet ad illam mortem, de qua terret Dóminus, qua mórtui sunt patres istórum; manducávit manna et Móyses, manducávit manna et Aaron, manducávit manna et Phínees, manducavérunt ibi multi, qui Dómino placuérunt, et mórtui non sunt. Quare? Quia visíbilem cibum spiritáliter intellexérunt, spiritáliter esuriérunt, spiritáliter gustavérunt, ut spiritáliter satiaréntur. Nam et nos hódie accípimus visíbilem cibum; sed áliud est sacraméntum, áliud virtus sacraménti.\par
\switchcolumn
\EN
\noindent But the death whereof the Lord doth sound the alarm, the death that their fathers died, is another death than that which is outward and bodily. Moses ate manna, Aaron ate manna, Phinehas ate manna, many ate manna in whom the Lord was well pleased and these are not dead. Wherefore Because they understood spiritually that outward bread, spiritually hungered thereafter, spiritually tasted thereof, and spiritually were satisfied therewith. So also do we this day feed on a visible food, but the Sacrament is one thing, and the grace of the Sacrament is another.\par
\switchcolumn*

\LA
\begin{Resp}\Rsign Misit me vivens Pater, et ego vivo propter Patrem: \Star{} Et qui mandúcat me, vivet propter me.\end{Resp}
\begin{Resp}\Vsign Cibávit illum Dóminus pane vitæ et intelléctus.\end{Resp}
\begin{Resp}\Rsign Et qui mandúcat me, vivet propter me.\end{Resp}
\begin{Resp}\Vsign Glória Patri, et Fílio, \Star{} et Spirítui Sancto.\end{Resp}
\begin{Resp}\Rsign Et qui mandúcat me, vivet propter me.\end{Resp}
\switchcolumn
\EN
\begin{Resp}\Rsign As the living Father hath sent Me, and I live by the Father, \Star{} So he that eateth Me, even he shall live by Me.\end{Resp}
\begin{Resp}\Vsign With the bread of life and understanding hath the Lord fed him.\end{Resp}
\begin{Resp}\Rsign So he that eateth Me, even he shall live by Me.\end{Resp}
\begin{Resp}\Vsign Glory be to the Father, and to the Son, \Star{} and to the Holy Ghost.\end{Resp}
\begin{Resp}\Rsign So he that eateth Me, even he shall live by Me.\end{Resp}
\switchcolumn*

\LA
\LessonHead{Lectio IX}
\switchcolumn
\EN
\LessonHead{Lesson IX}
\switchcolumn*

\LA
\noindent Quam multi de altári accípiunt, et moriúntur, et accipiéndo moriúntur! Unde dicit Apóstolus: Judícium sibi mandúcat et bibit. Nonne buccélla Domínica venénum fuit Judæ? Et tamen accépit. Et cum accépit, in eum inimícus intrávit; non quia malum accépit, sed quia bonum male malus accépit. Vidéte ergo, fratres, panem cæléstem spiritáliter manducáte, innocéntiam ad altáre apportáte. Peccáta, etsi sunt quotidiána, vel non sint mortífera; ántequam ad altáre accedátis, atténdite quod dicátis: Dimítte nobis débita nostra, sicut et nos dimíttimus debitóribus nostris. Si dimíttis, dimittétur tibi: secúrus accéde, panis est, non venénum.\par
\switchcolumn
\EN
\noindent That eateth and drinketh unworthily eateth and drinketh damnation to himself.” \textasciitilde{}(1 Cor. xi. 29.) Is it not written “When Jesus had dipped the sop, he gave it to Judas Iscariot, the son of Simon, and after the sop Satan entered into him” (John xiii. 26, 27.) And yet he took it. And when he had eaten it, the enemy entered in and possessed him. Not because what he ate was evil, but because he, being evil, dared to eat that which was good. Look to it well, then, brethren, that ye take spiritually the Bread Which cometh down from heaven. Bring innocency with you to the Altar. Though your sins be daily, let them not be deadly. Before ye draw near to the Altar, think well what it is that ye say “Forgive us our trespasses, as we forgive them that trespass against us.” “For, if ye forgive men their trespasses, your heavenly Father will also forgive you” (Matth. vi. 14) and ye may draw near boldly, for unto you It is Bread, and not poison.\par
\switchcolumn*

\LA
\TeDeum{Te Deum}
\switchcolumn
\EN
\TeDeum{Te Deum}
\switchcolumn*

\end{paracol}

\Office{Feria Quarta infra Octavam Corporis Christi}{Wednesday within the Octave of Corpus Christi}{Semiduplex II classis}
\addcontentsline{toc}{subsection}{Feria Quarta infra Octavam Corporis Christi \textit{— Wednesday within the Octave of Corpus Christi}}
\begin{paracol}{2}
\LA
\LessonHead{Lectio I}
\switchcolumn
\EN
\LessonHead{Lesson I}
\switchcolumn*

\LA
\LessonTitle{De libro primo Regum}
\switchcolumn
\EN
\LessonTitle{Lesson from the first book of Samuel}
\switchcolumn*

\LA
\Cite{1 Reg 6:19-21; 7:1}
\switchcolumn
\EN
\Cite{1 Sam 6:19-21; 7:1}
\switchcolumn*

\LA
\noindent Percússit autem de viris Bethsamítibus, eo quod vidíssent arcam Dómini; et percússit de pópulo septuagínta viros et quinquagínta míllia plebis. Luxítque pópulus, eo quod Dóminus percussísset plebem plaga magna. Et dixérunt viri Bethsamítæ: Quis póterit stare in conspéctu Dómini Dei sancti hujus? et ad quem ascéndet a nobis? Miserúntque núntios ad habitatóres Cariathíarim, dicéntes: Reduxérunt Philísthiim arcam Dómini: descéndite, et redúcite eam ad vos. Venérunt ergo viri Cariathíarim, et reduxérunt arcam Dómini, et intulérunt eam in domum Abínadab in Gábaa: Eleázarum autem fílium ejus sanctificavérunt, ut custodíret arcam Dómini.\par
\switchcolumn
\EN
\noindent But he slew of the men of Bethsames, because they had seen the ark of the Lord: and he slew of the people seventy men, and fifty thousand of the common people. And the people lamented, because the Lord had smitten the people with a great slaughter. And the men of Bethsames said: Who shall be able to stand before the Lord this holy God? and to whom shall he go up from us? And they sent messengers to the inhabitants of Cariathiarim, saying: The Philistines have brought back the ark of the Lord, come ye down and fetch it up to you. So they came as they were bidden, the men of Cariathiarim, and brought back the ark with them, housing it with a certain Abinadab in Gabaa; and they set apart his son Eleazar to keep watch over the Lord's ark.\par
\switchcolumn*

\LA
\begin{Resp}\Rsign Immolábit hædum multitúdo filiórum Israël ad vésperam Paschæ: \Star{} Et edent carnes et ázymos panes.\end{Resp}
\begin{Resp}\Vsign Pascha nostrum immolátus est Christus: ítaque epulémur in ázymis sinceritátis et veritátis.\end{Resp}
\begin{Resp}\Rsign Et edent carnes et ázymos panes.\end{Resp}
\switchcolumn
\EN
\begin{Resp}\Rsign The whole assembly of the children of Israel shall kill the lamb toward the evening of the Passover. \Star{} And they shall eat the flesh, and unleavened bread.\end{Resp}
\begin{Resp}\Vsign Even Christ our Passover is sacrificed for us therefore let us keep the feast with the unleavened bread of sincerity and truth.\end{Resp}
\begin{Resp}\Rsign And they shall eat the flesh, and unleavened bread.\end{Resp}
\switchcolumn*

\LA
\LessonHead{Lectio II}
\switchcolumn
\EN
\LessonHead{Lesson II}
\switchcolumn*

\LA
\Cite{1 Reg 7:2-4}
\switchcolumn
\EN
\Cite{1 Sam 7:2-4}
\switchcolumn*

\LA
\noindent Et factum est, ex qua die mansit arca Dómini in Cariathíarim, multiplicáti sunt dies (erat quippe jam annus vigésimus), et requiévit omnis domus Israël post Dóminum. Ait autem Sámuel ad univérsam domum Israël dicens: Si in toto corde vestro revertímini ad Dóminum, auférte deos aliénos de médio vestri Báalim et Astaroth et præparáte corda vestra Dómino et servíte ei soli, et éruet vos de manu Philísthiim. Abstulérunt ergo fílii Israël Báalim et Astaroth, et serviérunt Dómino soli.\par
\switchcolumn
\EN
\noindent And it came to pass, that from the day the ark of the Lord abode in Cariathiarim days were multiplied, (for it was now the twentieth year,) and all the house of Israel rested following the Lord. And Samuel spoke to all the house of Israel, saying: If you turn to the Lord with all your heart, put away the strange gods from among you, Baalim and Astaroth: and prepare your hearts unto the Lord, and serve him only, and he will deliver you out of the hand of the Philistines. Then the children of Israel put away Baalim and Astaroth, and served the Lord only.\par
\switchcolumn*

\LA
\begin{Resp}\Rsign Comedétis carnes, et saturabímini pánibus: \Star{} Iste est panis, quem dedit vobis Dóminus ad vescéndum.\end{Resp}
\begin{Resp}\Vsign Non Móyses dedit vobis panem de cælo, sed Pater meus dat vobis panem de cælo verum.\end{Resp}
\begin{Resp}\Rsign Iste est panis, quem dedit vobis Dóminus ad vescéndum.\end{Resp}
\switchcolumn
\EN
\begin{Resp}\Rsign Ye shall eat flesh, and shall be filled with bread. \Star{} This is the bread which the Lord hath given you to eat.\end{Resp}
\begin{Resp}\Vsign Moses gave you not that Bread from heaven, but My Father giveth you the true Bread from heaven.\end{Resp}
\begin{Resp}\Rsign This is the bread which the Lord hath given you to eat.\end{Resp}
\switchcolumn*

\LA
\LessonHead{Lectio III}
\switchcolumn
\EN
\LessonHead{Lesson III}
\switchcolumn*

\LA
\Cite{1 Reg 7:5-8}
\switchcolumn
\EN
\Cite{1 Sam 7:5-8}
\switchcolumn*

\LA
\noindent Dixit autem Sámuel: Congregáte univérsum Israël in Masphath, ut orem pro vobis Dóminum. Et convenérunt in Masphath hauserúntque aquam et effudérunt in conspéctu Dómini et jejunavérunt in die illa atque dixérunt ibi: Peccávimus Dómino. Judicavítque Sámuel fílios Israël in Masphath. Et audiérunt Philísthiim quod congregáti essent fílii Israël in Masphath, et ascendérunt sátrapæ Philisthinórum ad Israël. Quod cum audíssent fílii Israël, timuérunt a fácie Philisthinórum. Dixerúntque ad Samuélem: Ne cesses pro nobis clamáre ad Dóminum Deum nostrum, ut salvet nos de manu Philisthinórum.\par
\switchcolumn
\EN
\noindent And Samuel said: Gather all Israel to Masphath, that I may pray to the Lord for you. And they gathered together to Masphath: and they drew water, and poured it out before the Lord, and they fasted on that day, and they said there: We have sinned against the Lord. And Samuel judged the children of Israel in Masphath. And the Philistines heard that the children of Israel were gathered together to Masphath, and the lords of the Philistines went up against Israel. And when the children of Israel heard this, they were afraid of the Philistines. And they said to Samuel: Cease not to cry to the Lord our God for us, that he may save us out of the hand of the Philistines.\par
\switchcolumn*

\LA
\begin{Resp}\Rsign Respéxit Elías ad caput suum subcinerícium panem: qui surgens comédit et bibit: \Star{} Et ambulávit in fortitúdine cibi illíus usque ad montem Dei.\end{Resp}
\begin{Resp}\Vsign Si quis manducáverit ex hoc pane, vivet in ætérnum.\end{Resp}
\begin{Resp}\Rsign Et ambulávit in fortitúdine cibi illíus usque ad montem Dei.\end{Resp}
\begin{Resp}\Vsign Glória Patri, et Fílio, \Star{} et Spirítui Sancto.\end{Resp}
\begin{Resp}\Rsign Et ambulávit in fortitúdine cibi illíus usque ad montem Dei.\end{Resp}
\switchcolumn
\EN
\begin{Resp}\Rsign Elijah looked, and, behold, there was a cake baken on the coals at his head, and he arose, and did eat and drink; \Star{} And went in the strength of that meat [forty days and forty nights] unto the mount of God.\end{Resp}
\begin{Resp}\Vsign If any man eat of this Bread, he shall live for ever.\end{Resp}
\begin{Resp}\Rsign And went in the strength of that meat [forty days and forty nights] unto the mount of God.\end{Resp}
\begin{Resp}\Vsign Glory be to the Father, and to the Son, \Star{} and to the Holy Ghost.\end{Resp}
\begin{Resp}\Rsign And went in the strength of that meat [forty days and forty nights] unto the mount of God.\end{Resp}
\switchcolumn*

\LA
\LessonHead{Lectio IV}
\switchcolumn
\EN
\LessonHead{Lesson IV}
\switchcolumn*

\LA
\LessonTitle{Ex libro sancti Ambrósii Epíscopi de Sacraméntis}
\switchcolumn
\EN
\LessonTitle{From the Book upon the Sacraments written by St. Ambrose, Bishop of Milan}
\switchcolumn*

\LA
\Source{Lib. 4, Cap. 4}
\switchcolumn
\EN
\Source{Bk. iv. ch. 4}
\switchcolumn*

\LA
\noindent Auctor sacramentórum quis est, nisi Dóminus Jesus? De cælo ista sacraménta venérunt. Consílium enim omne de cælo est. Vere autem magnum est et divínum miráculum, quod pópulo pluit Deus manna de cælo, et non laborábat pópulus, et manducábat. Tu forte dicis: Meus panis est usitátus. Sed panis iste, panis est ante verba sacramentórum: ubi accésserit consecrátio, de pane fit caro Christi. Hoc ígitur astruámus. Quómodo potest, qui panis est, corpus esse Christi? Consecratióne. Consecrátio ígitur quibus verbis est, et cujus sermónibus? Dómini Jesu. Nam réliqua ómnia quæ dicúntur, laudem Deo déferunt, orátio præmíttitur pro pópulo, pro régibus, pro céteris: ubi venítur ut conficiátur venerábile Sacraméntum, jam non suis sermónibus sacérdos, sed útitur sermónibus Christi.\par
\switchcolumn
\EN
\noindent Who invented the Sacraments but the Lord Jesus? The Sacraments came down from heaven, for all counsel is from heaven. Nevertheless, it was a great and wonderful work of God when He rained down manna upon His people, and the people laboured not, and yet were fed. Perchance, thou sayest Here, it is my bread which is used. But that bread is bread only till the Sacramental words are spoken at the Consecration, instead of bread, there cometh to be the Body of Christ. This therefore let us establish. How cometh it that that which was bread becometh the Body of Christ Through the Consecration. And in what words and in Whose language doth the Consecration take place In those of the Lord Jesus. All the other things which are said (in the Liturgy), the ascription of praise to God (in the Preface), the prayer for the people, for kings, and for others which formeth the first part [of the Canon, these are put in the mouth of the Priest.] But when that point is reached when this worshipful Sacrament is to be consecrated, then the Priest useth no more his own words, but Christ's.\par
\switchcolumn*

\LA
\begin{Resp}\Rsign Cenántibus illis, accépit Jesus panem, et benedíxit, ac fregit, dedítque discípulis suis, et ait: \Star{} Accípite et comédite: hoc est corpus meum.\end{Resp}
\begin{Resp}\Vsign Dixérunt viri tabernáculi mei: Quis det de cárnibus ejus, ut saturémur?\end{Resp}
\begin{Resp}\Rsign Accípite et comédite: hoc est corpus meum.\end{Resp}
\switchcolumn
\EN
\begin{Resp}\Rsign As they were eating, Jesus took bread, and blest it, and brake it, and gave it to the disciples, and said: \Star{} Take, eat this is My Body.\end{Resp}
\begin{Resp}\Vsign The men of my tabernacle said: O that we had of his flesh, we cannot be satisfied.\end{Resp}
\begin{Resp}\Rsign Take, eat this is My Body.\end{Resp}
\switchcolumn*

\LA
\LessonHead{Lectio V}
\switchcolumn
\EN
\LessonHead{Lesson V}
\switchcolumn*

\LA
\noindent Ergo sermo Christi hoc cónficit Sacraméntum. Quis sermo Christi? Nempe is, quo facta sunt ómnia. Jussit Dóminus, et factum est cælum: jussit Dóminus, et facta est terra: jussit Dóminus, et facta sunt mária: jussit Dóminus, et omnis creatúra generáta est. Vides ergo quam operatórius sit sermo Christi. Si ergo tanta vis est in sermóne Dómini Jesu, ut incíperent esse quæ non erant; quanto magis operatórius est, ut quæ erant, in áliud commuténtur? Cælum non erat, terra non erat. Sed audi dicéntem: Ipse dixit, et facta sunt: ipse mandávit, et creáta sunt. Ergo tibi ut respóndeam, non erat corpus Christi ante consecratiónem; sed post consecratiónem dico tibi quod jam corpus est Christi. Ipse dixit, et factum est: ipse mandávit, et creátum est.\par
\switchcolumn
\EN
\noindent It is the word of Christ, therefore, Which doth the needful work in this Sacrament. And what is the word of Christ? It is the word of Him at Whose bidding all things were made. The Lord commanded, and the heavens were created the Lord commanded, and the earth was formed the Lord commanded, and the seas were made the Lord commanded, and all creatures sprang into being. Thou seest, then, how mightily working a word is the word of Christ. If then the word of Christ hath such power that it can make that to be which hath never been, wherein doth it appear greater that it maketh one thing to be changed into Another There was once no heaven there was once no sea there was once no earth. But hear him who saith: “He spake, and it was done He commanded, and it stood fast.” (Ps. xxxii. 9.) If, then, I am to answer thee, I tell thee, that before the Consecration it is not the Body of Christ, but after the Consecration it is the Body of Christ, for Himself “hath spoken, and it is done He hath commanded, and it standeth fast.”\par
\switchcolumn*

\LA
\begin{Resp}\Rsign Accépit Jesus cálicem, postquam cenávit, dicens: Hic calix novum testaméntum est in meo sánguine: \Star{} Hoc fácite in meam commemoratiónem.\end{Resp}
\begin{Resp}\Vsign Memória memor ero, et tabéscet in me ánima mea.\end{Resp}
\begin{Resp}\Rsign Hoc fácite in meam commemoratiónem.\end{Resp}
\switchcolumn
\EN
\begin{Resp}\Rsign Jesus took the cup, after supper, saying: This cup is the New Testament in My Blood. \Star{} This do in remembrance of Me.\end{Resp}
\begin{Resp}\Vsign My soul hath them still in remembrance, and is humbled in me.\end{Resp}
\begin{Resp}\Rsign This do in remembrance of Me.\end{Resp}
\switchcolumn*

\LA
\LessonHead{Lectio VI}
\switchcolumn
\EN
\LessonHead{Lesson VI}
\switchcolumn*

\LA
\noindent Jam redi mecum ad propositiónem meam. Magnum quidem et venerábile, quod manna Judǽis pluit e cælo. Sed intéllige quid est ámplius, manna de cælo, an corpus Christi? Corpus Christi útique, qui auctor est cæli. Deínde, manna qui manducávit, mórtuus est: qui manducáverit hoc corpus, fiet ei remíssio peccatórum, et non moriétur in ætérnum. Ergo non otióse, cum áccipis, tu dicis, Amen; jam in spíritu cónfitens quod accípias corpus Christi. Dicit tibi sacérdos, Corpus Christi; et tu dicis, Amen, hoc est, Verum. Quod confitétur lingua, téneat afféctus.\par
\switchcolumn
\EN
\noindent And now I come back to my text. It is indeed a great and worshipful fact that manna was rained down upon the Jews but, think thou, which was the more great and worshipful, the manna from heaven or the Body of Christ ”the Body of that Same Christ by Whom the heavens were made And, again the fathers “did eat manna, and are dead he that eateth of this Bread,” (John vi. 58,) It is unto him “the remission of sins,” (Matth. xxvi. 28,) and “he shall never die.” (John xi. 26.) Therefore it is not idly that, when thou art receiving, thou sayest “Amen” testifying in thine heart that That Which thou art taking is the Body of Christ. The Priest saith unto thee: “The Body of Christ” ”and thou answerest: “Amen” that is to say, “It is true.” What then thy tongue confesseth, let thine heart hold to.\par
\switchcolumn*

\LA
\begin{Resp}\Rsign Ego sum panis vitæ; patres vestri manducavérunt manna in desérto, et mórtui sunt: \Star{} Hic est panis de cælo descéndens, ut, si quis ex ipso mandúcet, non moriátur.\end{Resp}
\begin{Resp}\Vsign Ego sum panis vivus, qui de cælo descéndi: si quis manducáverit ex hoc pane, vivet in ætérnum.\end{Resp}
\begin{Resp}\Rsign Hic est panis de cælo descéndens, ut, si quis ex ipso mandúcet, non moriátur.\end{Resp}
\begin{Resp}\Vsign Glória Patri, et Fílio, \Star{} et Spirítui Sancto.\end{Resp}
\begin{Resp}\Rsign Hic est panis de cælo descéndens, ut, si quis ex ipso mandúcet, non moriátur.\end{Resp}
\switchcolumn
\EN
\begin{Resp}\Rsign I am that Bread of life. Your fathers did eat manna in the wilderness, and are dead. \Star{} This is the Bread Which cometh down from heaven, that a man may eat thereof, and not die.\end{Resp}
\begin{Resp}\Vsign I am the living Bread Which came down from heaven if any man eat of this Bread, he shall live for ever.\end{Resp}
\begin{Resp}\Rsign This is the Bread Which cometh down from heaven, that a man may eat thereof, and not die.\end{Resp}
\begin{Resp}\Vsign Glory be to the Father, and to the Son, \Star{} and to the Holy Ghost.\end{Resp}
\begin{Resp}\Rsign This is the Bread Which cometh down from heaven, that a man may eat thereof, and not die.\end{Resp}
\switchcolumn*

\LA
\LessonHead{Lectio VII}
\switchcolumn
\EN
\LessonHead{Lesson VII}
\switchcolumn*

\LA
\LessonTitle{Léctio sancti Evangélii secúndum Joánnem}
\switchcolumn
\EN
\LessonTitle{From the Holy Gospel according to John}
\switchcolumn*

\LA
\Cite{Joannes 6:5-59}
\switchcolumn
\EN

\switchcolumn*

\LA

\switchcolumn
\EN
\Source{John 6: 55-59)}
\switchcolumn*

\LA
\noindent In illo témpore: Dixit Jesus turbis Judæórum: Caro mea vere est cibus, et sanguis meus vere est potus. Et réliqua.\par
\switchcolumn
\EN
\noindent At that time, Jesus said unto the multitudes of the Jews: My Flesh is meat indeed, and My Blood is drink indeed. And so on.\par
\switchcolumn*

\LA
\LessonTitle{Homilía sancti Hilárii Epíscopi}
\switchcolumn
\EN
\LessonTitle{Homily by St. Hilary, Bishop of Poitiers.}
\switchcolumn*

\LA
\Source{Liber 8 de Trinitáte, ante medium}
\switchcolumn
\EN
\Source{Bk. viii. on the Trinity)}
\switchcolumn*

\LA
\noindent Non est humáno aut sǽculi sensu in Dei rebus loquéndum. Quæ scripta sunt, legámus, et quæ legérimus, intelligámus; et tunc perféctæ fídei offício fungémur. De naturáli enim in nobis Christi veritáte quæ dícimus, nisi ab eo díscimus, stulte atque ímpie dícimus. Ipse enim ait: Caro mea vere est esca, et sanguis meus vere est potus. Qui edit carnem meam et bibit sánguinem meum, in me manet, et ego in eo. De veritáte carnis et sánguinis non relíctus est ambigéndi locus.\par
\switchcolumn
\EN
\noindent When we speak concerning the things of God, we must not speak after the manner of men, nor after the manner of the world. Let us read those things which are written, and understand those things which we read and then let us act as having a perfect faith. We shall speak but folly and godlessness if we speak concerning the natural truth of Christ in use and have not learnt at Christ's School how we should speak. He Himself saith “My Flesh is meat indeed, and My Blood is drink indeed. He that eateth My Flesh and drinketh My Blood, dwelleth in Me, and I in him.” There is here no room left for doubt as to What is His Flesh and what is His Blood.\par
\switchcolumn*

\LA
\begin{Resp}\Rsign Qui mandúcat meam carnem et bibit meum sánguinem, \Star{} In me manet, et ego in eo.\end{Resp}
\begin{Resp}\Vsign Non est ália nátio tam grandis, quæ hábeat deos appropinquántes sibi, sicut Deus noster adest nobis.\end{Resp}
\begin{Resp}\Rsign In me manet, et ego in eo.\end{Resp}
\switchcolumn
\EN
\begin{Resp}\Rsign He that eateth My Flesh and drinketh My Blood, \Star{} Dwelleth in Me, and I in him.\end{Resp}
\begin{Resp}\Vsign What nation is there so great, who hath gods so nigh unto them, as the Lord our God is to us\end{Resp}
\begin{Resp}\Rsign Dwelleth in Me, and I in him.\end{Resp}
\switchcolumn*

\LA
\LessonHead{Lectio VIII}
\switchcolumn
\EN
\LessonHead{Lesson VIII}
\switchcolumn*

\LA
\noindent Nunc enim, et ipsíus Dómini professióne, et fide nostra, vere caro est, et vere sanguis est. Et hæc accépta atque hausta id effíciunt, ut et nos in Christo, et Christus in nobis sit. An ne hoc véritas non est? Contíngat plane his verum non esse, qui Christum Jesum verum esse Deum negant. Est ergo in nobis ipse per carnem, et sumus in eo, dum secum hoc quod nos sumus, in Deo est. Quod autem in eo per sacraméntum communicátæ carnis et sánguinis simus, ipse testátur, dicens: Et hic mundus jam me non videt, vos autem me vidébitis: quóniam ego vivo, et vos vivétis: quóniam ego in Patre meo, et vos in me, et ego in vobis.\par
\switchcolumn
\EN
\noindent Now we know by the declaration of the Lord Himself and by (the teaching of) our Faith, the reality of His Flesh and Blood. And when we eat the One and drink the Other, They work effectually in us to make us dwell in Him and He in us. Is “not this a reality Surely it befalleth not them to find it true, who deny that Christ Jesus is Very God. He is in us by means of His Flesh, and we are in Him when that which we are is with Him in God. That we dwell in Him through that Sacrament wherein His Flesh and Blood are given unto us, He Himself doth testify, where He saith “Yet a little while, and the world seeth Me no more but ye see Me because I live ye shall live also. (At that day ye shall know that) I am in My Father, and ye in Me, and I in you.” (John xiv. 19, 20.)\par
\switchcolumn*

\LA
\begin{Resp}\Rsign Misit me vivens Pater, et ego vivo propter Patrem: \Star{} Et qui mandúcat me, vivet propter me.\end{Resp}
\begin{Resp}\Vsign Cibávit illum Dóminus pane vitæ et intelléctus.\end{Resp}
\begin{Resp}\Rsign Et qui mandúcat me, vivet propter me.\end{Resp}
\begin{Resp}\Vsign Glória Patri, et Fílio, \Star{} et Spirítui Sancto.\end{Resp}
\begin{Resp}\Rsign Et qui mandúcat me, vivet propter me.\end{Resp}
\switchcolumn
\EN
\begin{Resp}\Rsign As the living Father hath sent Me, and I live by the Father, \Star{} So he that eateth Me, even he shall live by Me.\end{Resp}
\begin{Resp}\Vsign With the bread of life and understanding hath the Lord fed him.\end{Resp}
\begin{Resp}\Rsign So he that eateth Me, even he shall live by Me.\end{Resp}
\begin{Resp}\Vsign Glory be to the Father, and to the Son, \Star{} and to the Holy Ghost.\end{Resp}
\begin{Resp}\Rsign So he that eateth Me, even he shall live by Me.\end{Resp}
\switchcolumn*

\LA
\LessonHead{Lectio IX}
\switchcolumn
\EN
\LessonHead{Lesson IX}
\switchcolumn*

\LA
\noindent Quod autem in nobis naturális hæc únitas sit, ipse ita testátus est: Qui edit carnem meam et bibit sánguinem meum, in me manet, et ego in eo. Non enim quis in eo erit, nisi in quo ipse fúerit; ejus tantum in se assúmptam habens carnem, qui suam súmpserit. Perféctæ autem hujus unitátis sacraméntum supérius jam docúerat, dicens: Sicut me misit vivens Pater, et ego vivo per Patrem; et qui mandúcat meam carnem, et ipse vivet per me. Vivit ergo per Patrem: et quo modo per Patrem vivit, eódem modo nos per carnem ejus vivémus.\par
\switchcolumn
\EN
\noindent But that this union in us is a real one, He testifieth thus: “He that eateth My Flesh and drinketh My Blood, dwelleth in Me, and I in him.” For no one dwelleth in Him in whom He doth not dwell, since he which receiveth (the Body of Christ) hath but received that Flesh of (the same nature as) his own, which Christ hath taken into Himself. The mystery of this perfect union He had taught before, when He said: “As the living Father hath sent Me, and I live by the Father, so, he that eateth Me, even he shall live by Me.” He therefore liveth by the Father, and, as He liveth by the Father, so shall we live by Him.\par
\switchcolumn*

\LA
\TeDeum{Te Deum}
\switchcolumn
\EN
\TeDeum{Te Deum}
\switchcolumn*

\end{paracol}

\Office{In Octava Sanctissimi Corporis Christi}{Octava Sanctissimi Corporis Christi}{Duplex majus}
\addcontentsline{toc}{subsection}{In Octava Sanctissimi Corporis Christi \textit{— Octava Sanctissimi Corporis Christi}}
\begin{paracol}{2}
\LA
\LessonHead{Lectio I}
\switchcolumn
\EN
\LessonHead{Lesson I}
\switchcolumn*

\LA
\LessonTitle{De libro primo Regum}
\switchcolumn
\EN
\LessonTitle{Lesson from the first book of Samuel}
\switchcolumn*

\LA
\Cite{1 Reg 8:4-6}
\switchcolumn
\EN
\Cite{1 Sam 8:4-6}
\switchcolumn*

\LA
\noindent Congregáti ergo univérsi majóres natu Israël, venérunt ad Samuélem in Rámatha, dixerúntque ei: Ecce tu senuísti, et fílii tui non ámbulant in viis tuis: constítue nobis regem, ut júdicet nos, sicut et univérsæ habent natiónes. Displícuit sermo in óculis Samuélis, eo quod dixíssent: Da nobis regem, ut júdicet nos. Et orávit Sámuel ad Dóminum.\par
\switchcolumn
\EN
\noindent Then all the ancients of Israel being assembled, came to Samuel to Ramatha. And they said to him: Behold thou art old, and thy sons walk not in thy ways: make us a king, to judge us, as all nations have. And the word was displeasing in the eyes of Samuel, that they should say: Give us a king, to judge us. And Samuel prayed to the Lord.\par
\switchcolumn*

\LA
\begin{Resp}\Rsign Immolábit hædum multitúdo filiórum Israël ad vésperam Paschæ: \Star{} Et edent carnes et ázymos panes.\end{Resp}
\begin{Resp}\Vsign Pascha nostrum immolátus est Christus: ítaque epulémur in ázymis sinceritátis et veritátis.\end{Resp}
\begin{Resp}\Rsign Et edent carnes et ázymos panes.\end{Resp}
\switchcolumn
\EN
\begin{Resp}\Rsign The whole assembly of the children of Israel shall kill the lamb toward the evening of the Passover. \Star{} And they shall eat the flesh, and unleavened bread.\end{Resp}
\begin{Resp}\Vsign Even Christ our Passover is sacrificed for us therefore let us keep the feast with the unleavened bread of sincerity and truth.\end{Resp}
\begin{Resp}\Rsign And they shall eat the flesh, and unleavened bread.\end{Resp}
\switchcolumn*

\LA
\LessonHead{Lectio II}
\switchcolumn
\EN
\LessonHead{Lesson II}
\switchcolumn*

\LA
\Cite{1 Reg 8:7-9}
\switchcolumn
\EN
\Cite{1 Sam 8:7-9}
\switchcolumn*

\LA
\noindent Dixit autem Dóminus ad Samuélem: Audi vocem pópuli in ómnibus quæ loquúntur tibi: non enim te abjecérunt sed me, ne regnem super eos. Juxta ómnia ópera sua, quæ fecérunt a die qua edúxi eos de Ægýpto usque ad diem hanc, sicut dereliquérunt me et serviérunt diis aliénis, sic fáciunt étiam tibi. Nunc ergo vocem eórum audi; verúmtamen contestáre eos et prædic eis jus regis, qui regnatúrus est super eos.\par
\switchcolumn
\EN
\noindent And the Lord said to Samuel: Hearken to the voice of the people in all that they say to thee. For they have not rejected thee, but me, that I should not reign over them. According to all their works, they have done from the day that I brought them out of Egypt until this day: as they have forsaken me, and served strange gods, so do they also unto thee. Now therefore hearken to their voice: but yet testify to them, and foretell them the right of the king, that shall reign over them\par
\switchcolumn*

\LA
\begin{Resp}\Rsign Comedétis carnes, et saturabímini pánibus: \Star{} Iste est panis, quem dedit vobis Dóminus ad vescéndum.\end{Resp}
\begin{Resp}\Vsign Non Móyses dedit vobis panem de cælo, sed Pater meus dat vobis panem de cælo verum.\end{Resp}
\begin{Resp}\Rsign Iste est panis, quem dedit vobis Dóminus ad vescéndum.\end{Resp}
\switchcolumn
\EN
\begin{Resp}\Rsign Ye shall eat flesh, and shall be filled with bread. \Star{} This is the bread which the Lord hath given you to eat.\end{Resp}
\begin{Resp}\Vsign Moses gave you not that Bread from heaven, but My Father giveth you the true Bread from heaven.\end{Resp}
\begin{Resp}\Rsign This is the bread which the Lord hath given you to eat.\end{Resp}
\switchcolumn*

\LA
\LessonHead{Lectio III}
\switchcolumn
\EN
\LessonHead{Lesson III}
\switchcolumn*

\LA
\Cite{1 Reg 8:10-14}
\switchcolumn
\EN
\Cite{1 Sam 8:10-14}
\switchcolumn*

\LA
\noindent Dixit ítaque Sámuel ómnia verba Dómini ad pópulum, qui petíerat a se regem, et ait: Hoc erit jus regis, qui imperatúrus est vobis. Fílios vestros tollet et ponet in cúrribus suis, faciétque sibi équites et præcursóres quadrigárum suárum; et constítuet sibi tribúnos et centuriónes et aratóres agrórum suórum et messóres ségetum et fabros armórum et cúrruum suórum; fílias quoque vestras fáciet sibi unguentárias, et focárias, et paníficas: agros quoque vestros, et víneas, et olivéta, óptima tollet, et dabit servis suis.\par
\switchcolumn
\EN
\noindent Then Samuel told all the words of the Lord to the people that had desired a king of him, And said: This will be the right of the king, that shall reign over you: He will take your sons, and put them in his chariots, and will make them his horsemen, and his running footmen to run before his chariots, And he will appoint of them to be his tribunes, and centurions, and to plough his fields, and to reap his corn, and to make him arms and chariots. Your daughters also he will take to make him ointments, and to be his cooks, and bakers. And he will take your fields, and your vineyards, and your best oliveyards, and give them to his servants.\par
\switchcolumn*

\LA
\begin{Resp}\Rsign Respéxit Elías ad caput suum subcinerícium panem: qui surgens comédit et bibit: \Star{} Et ambulávit in fortitúdine cibi illíus usque ad montem Dei.\end{Resp}
\begin{Resp}\Vsign Si quis manducáverit ex hoc pane, vivet in ætérnum.\end{Resp}
\begin{Resp}\Rsign Et ambulávit in fortitúdine cibi illíus usque ad montem Dei.\end{Resp}
\begin{Resp}\Vsign Glória Patri, et Fílio, \Star{} et Spirítui Sancto.\end{Resp}
\begin{Resp}\Rsign Et ambulávit in fortitúdine cibi illíus usque ad montem Dei.\end{Resp}
\switchcolumn
\EN
\begin{Resp}\Rsign Elijah looked, and, behold, there was a cake baken on the coals at his head, and he arose, and did eat and drink; \Star{} And went in the strength of that meat [forty days and forty nights] unto the mount of God.\end{Resp}
\begin{Resp}\Vsign If any man eat of this Bread, he shall live for ever.\end{Resp}
\begin{Resp}\Rsign And went in the strength of that meat [forty days and forty nights] unto the mount of God.\end{Resp}
\begin{Resp}\Vsign Glory be to the Father, and to the Son, \Star{} and to the Holy Ghost.\end{Resp}
\begin{Resp}\Rsign And went in the strength of that meat [forty days and forty nights] unto the mount of God.\end{Resp}
\switchcolumn*

\LA
\LessonHead{Lectio IV}
\switchcolumn
\EN
\LessonHead{Lesson IV}
\switchcolumn*

\LA
\LessonTitle{Sermo sancti Cyrílli Epíscopi Jerosolymitáni}
\switchcolumn
\EN
\LessonTitle{From the Sermons of the Blessed Patriarch of Jerusalem Cyril.}
\switchcolumn*

\LA
\Source{Catechesis mystagog. 4}
\switchcolumn
\EN
\Source{Catechetical Lectures, 40}
\switchcolumn*

\LA
\noindent Ipsa beáti Pauli doctrína abúnde suffícere vidétur, ut certam vobis de divínis mystériis fidem fáciat, quibus digni rédditi, concorpórei, ut ita dicam, et consanguínei Christi facti estis. Ipse enim modo clamábat, quod in nocta qua tradebátur Dóminus noster Jesus Christus, accípiens panem, et grátias agens fregit, et dedit discípulis suis, dicens: Accípite, et manducáte: hoc est corpus meum. Et accípiens cálicem, et grátias agens, dixit: Accípite, et bíbite: hic est sanguis meus. Cum ígitur ipse de pane pronuntiáverit ac díxerit: Hoc est corpus meum; quis audébit deínceps ambígere? Et cum idem ipse tam asseveránter díxerit: Hic est sanguis meus; quis umquam dubitáverit, ut dicat non esse ejus sánguinem?\par
\switchcolumn
\EN
\noindent The teaching of the blessed Paul seemeth of itself enough instruction for you concerning those Divine Mysteries, whereof, if ye be made worthy, ye become therein, so to speak, of one Body and of one Blood with Christ. Paul saith that our Lord Jesus Christ, “the same night in which He was betrayed, took bread and, when He had given thanks, He brake it, and gave it unto His disciples, saying Take, eat this is My Body. After the same manner also He took the cup,” and gave thanks, “and said:” Take this, and drink it this is My Blood. Since therefore it is He Who hath definitely stated and said, touching that Bread “This is My Body” -who will dare any longer to doubt that It is so And since it is He again that hath absolutely affirmed and said, touching that cup “This is My Blood” who is he that will doubt any longer, or say that It is not His Blood\par
\switchcolumn*

\LA
\begin{Resp}\Rsign Cenántibus illis, accépit Jesus panem, et benedíxit, ac fregit, dedítque discípulis suis, et ait: \Star{} Accípite et comédite: hoc est corpus meum.\end{Resp}
\begin{Resp}\Vsign Dixérunt viri tabernáculi mei: Quis det de cárnibus ejus, ut saturémur?\end{Resp}
\begin{Resp}\Rsign Accípite et comédite: hoc est corpus meum.\end{Resp}
\switchcolumn
\EN
\begin{Resp}\Rsign As they were eating, Jesus took bread, and blest it, and brake it, and gave it to the disciples, and said: \Star{} Take, eat this is My Body.\end{Resp}
\begin{Resp}\Vsign The men of my tabernacle said: O that we had of his flesh, we cannot be satisfied.\end{Resp}
\begin{Resp}\Rsign Take, eat this is My Body.\end{Resp}
\switchcolumn*

\LA
\LessonHead{Lectio V}
\switchcolumn
\EN
\LessonHead{Lesson V}
\switchcolumn*

\LA
\noindent Aquam olim in vinum convértit in Cana Galilǽæ, quod habet quamdam cum sánguine propinquitátem: et eum parum dignum existimábimus, cui credámus quod vinum in sánguinem transmutárit? Ad eas núptias, quibus córpora copulántur, vocátus, præter opiniónem ómnium hoc fecit miráculum: et non multo magis sic eum corpus et sánguinem suum fruénda nobis donásse persuásum fírmiter habébimus, ut ea cum omni certitúdine tamquam corpus ipsíus et sánguinem sumámus? Nam in spécie panis dat nobis corpus, et in spécie vini dat nobis sánguinem: ut, cum súmpseris, gustes corpus et sánguinem Christi, factus ejúsdem córporis et sánguinis párticeps. Sic enim effícimur Christíferi, hoc est, Christum in corpóribus nostris feréntes, cum corpus ejus et sánguinem in membra nostra recípimus: sic secúndum beátum Petrum, divínæ natúræ consórtes réddimur.\par
\switchcolumn
\EN
\noindent At the beginning of His ministry, at Cana in Galilee, the Lord turned water into wine, a thing which hath some qualities in common with blood and shall we deem Him less worthy that we should believe Him, when He turneth wine into Blood When He was bidden to that marriage wherein twain were made one flesh, He did the beginning of His miracles to the amazement of all men; and shall we less surely hold that He hath given us His Body and Blood to be our meat and drink, or take them with weaker faith that they are indeed His Body and His Blood? Under the appearance of bread He giveth unto us His Body, and, under the appearance of wine, His Blood and when thou shalt come to receive, it is on the Body and Blood of Christ that thou wilt feed, being made a partaker of His Body and of His Blood. Thus indeed it is that we become Christbearers, namely, by carrying about Christ in our bodies, when we receive His Body and Blood into our own frames. Thus, as the blessed Peter hath it, we are “partakers of the Divine nature.” (2 Pet. i. 4.)\par
\switchcolumn*

\LA
\begin{Resp}\Rsign Accépit Jesus cálicem, postquam cenávit, dicens: Hic calix novum testaméntum est in meo sánguine: \Star{} Hoc fácite in meam commemoratiónem.\end{Resp}
\begin{Resp}\Vsign Memória memor ero, et tabéscet in me ánima mea.\end{Resp}
\begin{Resp}\Rsign Hoc fácite in meam commemoratiónem.\end{Resp}
\switchcolumn
\EN
\begin{Resp}\Rsign Jesus took the cup, after supper, saying: This cup is the New Testament in My Blood. \Star{} This do in remembrance of Me.\end{Resp}
\begin{Resp}\Vsign My soul hath them still in remembrance, and is humbled in me.\end{Resp}
\begin{Resp}\Rsign This do in remembrance of Me.\end{Resp}
\switchcolumn*

\LA
\LessonHead{Lectio VI}
\switchcolumn
\EN
\LessonHead{Lesson VI}
\switchcolumn*

\LA
\noindent Olim cum Judǽis Christus dísserens, Nisi manducavéritis, inquit, carnem meam, et bibéritis meum sánguinem, non habébitis vitam in vobis. Cum autem illi, quæ dicta fúerant, non spiritáliter accepíssent, offénsi abiérunt retro; putábant enim quod eos ad manducándas carnes hortarétur. Erant et in véteri testaménto panes propositiónis; verum illi cum fúerint véteris testaménti, finem jam accepérunt. In novo vero testaménto panis est cæléstis et calix salutáris, qui et ánimam et corpus sanctíficant. Quam ob rem non sic hæc atténdas velim, tamquam sint nudus et simplex panis, nudum et simplex vinum; corpus enim sunt et sanguis Christi. Nam étiam si sensus illud tibi renúntiat, fides tamen te confírmet. Ne júdices rem ex gustu: sed te citra ullam dubitatiónem fides certum reddat, quod sis dignus factus, qui córporis et sánguinis Christi párticeps fíeres.\par
\switchcolumn
\EN
\noindent Christ once said, in conversing with the Jews “Except ye eat the Flesh of the Son of Man, and drink His Blood, ye have no life in you.” (John vi. 53.) But they took not spiritually that which He said, and “from that time many of His disciples went back, and walked no more with Him.” (66.) They thought that He had bidden them to eat flesh. The Old Testament also had Shewbread, but this Old Testament bread was now to have an end. The bread of the New Testament is “the Bread Which cometh down from heaven” (50), the cup of the New Testament, the Cup of Salvation, that Bread and that Cup Which hallow both souls and bodies. Wherefore I will have thee to understand that the Bread and Wine whereunto thou art to come, are not mere common bread or mere common wine for they are the Body and the Blood of Christ. Even if thy senses do indeed deny this fact, yet let faith make thee right sure of it. Judge not the Thing by the taste thereof, but let faith assure thee beyond all doubt thou art partaking of the Body and Blood of Christ.\par
\switchcolumn*

\LA
\begin{Resp}\Rsign Ego sum panis vitæ; patres vestri manducavérunt manna in desérto, et mórtui sunt: \Star{} Hic est panis de cælo descéndens, ut, si quis ex ipso mandúcet, non moriátur.\end{Resp}
\begin{Resp}\Vsign Ego sum panis vivus, qui de cælo descéndi: si quis manducáverit ex hoc pane, vivet in ætérnum.\end{Resp}
\begin{Resp}\Rsign Hic est panis de cælo descéndens, ut, si quis ex ipso mandúcet, non moriátur.\end{Resp}
\begin{Resp}\Vsign Glória Patri, et Fílio, \Star{} et Spirítui Sancto.\end{Resp}
\begin{Resp}\Rsign Hic est panis de cælo descéndens, ut, si quis ex ipso mandúcet, non moriátur.\end{Resp}
\switchcolumn
\EN
\begin{Resp}\Rsign I am that Bread of life. Your fathers did eat manna in the wilderness, and are dead. \Star{} This is the Bread Which cometh down from heaven, that a man may eat thereof, and not die.\end{Resp}
\begin{Resp}\Vsign I am the living Bread Which came down from heaven if any man eat of this Bread, he shall live for ever.\end{Resp}
\begin{Resp}\Rsign This is the Bread Which cometh down from heaven, that a man may eat thereof, and not die.\end{Resp}
\begin{Resp}\Vsign Glory be to the Father, and to the Son, \Star{} and to the Holy Ghost.\end{Resp}
\begin{Resp}\Rsign This is the Bread Which cometh down from heaven, that a man may eat thereof, and not die.\end{Resp}
\switchcolumn*

\LA
\LessonHead{Lectio VII}
\switchcolumn
\EN
\LessonHead{Lesson VII}
\switchcolumn*

\LA
\LessonTitle{Léctio sancti Evangélii secúndum Joánnem}
\switchcolumn
\EN
\LessonTitle{From the Holy Gospel according to John}
\switchcolumn*

\LA
\Cite{Joannes 6:56-59}
\switchcolumn
\EN
\Cite{John 6:56-59}
\switchcolumn*

\LA
\noindent In illo témpore: Dixit Jesus turbis Judæórum: Caro mea vere est cibus, et sanguis meus vere est potus. Et réliqua.\par
\switchcolumn
\EN
\noindent At that time Jesus said unto the multitudes of the Jews My Flesh is meat indeed, and My Blood is drink indeed. And so on.\par
\switchcolumn*

\LA
\LessonTitle{Homilía sancti Cyrílli Epíscopi Alexandríni}
\switchcolumn
\EN
\LessonTitle{Homily by St. Cyril, Pope of Alexandria.}
\switchcolumn*

\LA
\Source{Liber 4 in Joánnem, cap. 17}
\switchcolumn
\EN
\Source{Book iv. on John, ch. 17}
\switchcolumn*

\LA
\noindent Qui mandúcat, inquit, carnem meam et bibit sánguinem meum, in me manet, et ego in illo. Sicut enim si quis liquefáctæ ceræ áliam ceram infúderit, álteram cum áltera per totum commísceat necésse est; sic qui carnem et sánguinem Dómini récipit, cum ipso ita conjúngitur, ut Christus in ipso, et ipse in Christo inveniátur. Símili quodam modo apud Matthǽum compéries. Símile est, inquit, regnum cælórum ferménto, quod accéptum abscóndit múlier in farínæ satis tribus. Sicut parum, ut Paulus ait, ferménti totam massam ferméntat, sic párvula benedíctio totum hóminem in seípsam áttrahit, et sua grátia replet; et hoc modo in nobis Christus manet, et nos in Christo.\par
\switchcolumn
\EN
\noindent “He that eateth My Flesh and drinketh My Blood,” saith the Lord, “dwelleth in Me, and I in him.” If a man take two pieces of wax and melt them, and pour the one into the other, they necessarily mingle so also, he that receiveth the Body and Blood of the Lord doth become so joined with the Lord that he is to be found in Christ and Christ in him. Another comparison thou wilt find in Matthew. The Lord there saith “The kingdom of heaven is like unto leaven which a woman took, and hid in three measures of meal, [till the whole was leavened,]” (xiii. 33,) because, as Paul saith, “a little leaven leaveneth the whole lump.” (Gal. v. 9.) So also doth a little of this Blessing draw the whole man unto Itself, and fill him with Its grace and thus doth Christ dwell in us, and we in Christ.\par
\switchcolumn*

\LA
\begin{Resp}\Rsign Qui mandúcat meam carnem et bibit meum sánguinem, \Star{} In me manet, et ego in eo.\end{Resp}
\begin{Resp}\Vsign Non est ália nátio tam grandis, quæ hábeat deos appropinquántes sibi, sicut Deus noster adest nobis.\end{Resp}
\begin{Resp}\Rsign In me manet, et ego in eo.\end{Resp}
\switchcolumn
\EN
\begin{Resp}\Rsign He that eateth My Flesh and drinketh My Blood, \Star{} Dwelleth in Me, and I in him.\end{Resp}
\begin{Resp}\Vsign What nation is there so great, who hath gods so nigh unto them, as the Lord our God is to us\end{Resp}
\begin{Resp}\Rsign Dwelleth in Me, and I in him.\end{Resp}
\switchcolumn*

\LA
\LessonHead{Lectio VIII}
\switchcolumn
\EN
\LessonHead{Lesson VIII}
\switchcolumn*

\LA
\noindent Nos vero, si vitam ætérnam cónsequi vólumus, si largitórem immortalitátis habére in nobis desiderámus, ad recipiéndam benedictiónem libénter concurrámus; caveamúsque ne loco láquei damnósam religiónem diábolus nobis præténdat. Recte (inquit) dicis; scriptum tamen esse non ignorámus, judícium sibi comédere atque bíbere illum, qui de pane cómedit et de cálice bibit indígne. Ego ígitur probo meípsum, et indígnum invénio. Quando ígitur, quicúmque tu es qui ista dicis, dignus eris? quando Christo teípsum ófferes? Nam si peccándo indígnus es, et peccáre non désinis, (quis enim delícta intélligit? secúndum Psalmístam) expers omníno eris vivíficæ hujus sanctificatiónis.\par
\switchcolumn
\EN
\noindent As for ourselves, if we would win life everlasting; if we would that the Giver of immortality should dwell in us, let us run freely to receive this Blessing, and let us beware that the devil succeed not in laying a stumbling-block in our way, in the shape of a mistaken reverence. Thou rightly sayest, and we know well, how that it is written “Whosoever shall eat this Bread and drink this Cup of the Lord unworthily eateth and drinketh damnation to himself.” (Cor. xi. 27, 29.) I therefore examine myself and find myself unworthy. And I ask thee, who citest these words to me, who shall ever be found worthy When wilt thou be such an one as may be worthy to be offered to Christ If by sin thou art unworthy, and thou ceasest not to sin, (for, as the Psalmist hath it, “Who can understand his errors?” (Ps. xviii. 13) then shalt thou for ever lack this means of life and sanctification.\par
\switchcolumn*

\LA
\begin{Resp}\Rsign Misit me vivens Pater, et ego vivo propter Patrem: \Star{} Et qui mandúcat me, vivet propter me.\end{Resp}
\begin{Resp}\Vsign Cibávit illum Dóminus pane vitæ et intelléctus.\end{Resp}
\begin{Resp}\Rsign Et qui mandúcat me, vivet propter me.\end{Resp}
\begin{Resp}\Vsign Glória Patri, et Fílio, \Star{} et Spirítui Sancto.\end{Resp}
\begin{Resp}\Rsign Et qui mandúcat me, vivet propter me.\end{Resp}
\switchcolumn
\EN
\begin{Resp}\Rsign As the living Father hath sent Me, and I live by the Father, \Star{} So he that eateth Me, even he shall live by Me.\end{Resp}
\begin{Resp}\Vsign With the bread of life and understanding hath the Lord fed him.\end{Resp}
\begin{Resp}\Rsign So he that eateth Me, even he shall live by Me.\end{Resp}
\begin{Resp}\Vsign Glory be to the Father, and to the Son, \Star{} and to the Holy Ghost.\end{Resp}
\begin{Resp}\Rsign So he that eateth Me, even he shall live by Me.\end{Resp}
\switchcolumn*

\LA
\LessonHead{Lectio IX}
\switchcolumn
\EN
\LessonHead{Lesson IX}
\switchcolumn*

\LA
\noindent Quare pias, quæso, cogitatiónes suscípias, studióse sanctéque vivas, et benedictiónem partícipes: quæ (mihi crede) non mortem solum, verum étiam morbos omnes depéllit. Sedat enim, cum in nobis máneat Christus, sæviéntem membrórum nostrórum legem: pietátem corróborat, perturbatiónes ánimi exstínguit, ægrótos curat, collísos redíntegrat: et sicut pastor bonus, qui ánimam suam pro óvibus pósuit, ab omni nos érigit casu.\par
\switchcolumn
\EN
\noindent Therefore, I counsel thee to betake thee to godly thoughts, and to live carefully and holily, and so to receive that Blessing a Blessing which, believe me, doth banish, not death only, but all diseases likewise. For when Christ dwelleth in us, He stilleth the law of death in our members, which warreth against the law of our mind, (Rom. vii. 23,) He giveth strength to godliness, He turneth to calm the turbulent surging of our mind, He cureth them which are sick, He raiseth up them which are fallen, and, like the Good Shepherd, Which giveth His life for the sheep, He prevaileth that the sheep perish not.\par
\switchcolumn*

\LA
\TeDeum{Te Deum}
\switchcolumn
\EN
\TeDeum{Te Deum}
\switchcolumn*

\end{paracol}

\Office{Sacratissimi Cordis Domini Nostri Jesu Christi}{The Most Sacred Heart of Our Lord Jesus Christ}{Duplex I classis cum Octava privilegiata III ordinis}
\addcontentsline{toc}{subsection}{Sacratissimi Cordis Domini Nostri Jesu Christi \textit{— The Most Sacred Heart of Our Lord Jesus Christ}}
\begin{paracol}{2}
\LA
\LessonHead{Lectio I}
\switchcolumn
\EN
\LessonHead{Lesson I}
\switchcolumn*

\LA
\LessonTitle{De Jeremía Prophéta}
\switchcolumn
\EN
\LessonTitle{Lesson from the book of Jeremias}
\switchcolumn*

\LA
\Cite{Jer 24:5-7}
\switchcolumn
\EN
\Cite{Jer 24:5-7}
\switchcolumn*

\LA
\noindent Hæc dicit Dóminus, Deus Israël: Cognóscam transmigratiónem Juda, quam emísi de loco isto in terram Chaldæórum, in bonum. Et ponam óculos meos super eos ad placándum, et redúcam eos in terram hanc; et ædificábo eos, et non déstruam; et plantábo eos et non evéllam. Et dabo eis cor ut sciant me, quia ego sum Dóminus; et erunt mihi in pópulum, et ego ero eis in Deum, quia reverténtur ad me in toto corde suo.\par
\switchcolumn
\EN
\noindent Thus saith the Lord the God of Israel: Like these good figs, so will I regard the captives of Juda, whom I have sent forth out of this place into the land of the Chaldeans, for their good. And I will set my eyes upon them to be pacified, and I will bring them again into this land: and I will build them up, and not pull them down: and I will plant them, and not pluck them up. And I will give them a heart to know me, that I am the Lord: and they shall be my people, and I will be their God: because they shall return to me with their whole heart.\par
\switchcolumn*

\LA
\begin{Resp}\Rsign Fériam eis pactum sempitérnum et non désinam eis benefácere et timórem meum dabo in corde eórum \Star{} Ut non recédant a me.\end{Resp}
\begin{Resp}\Vsign Et lætábor super eis cum bene eis fécero in toto Corde meo.\end{Resp}
\begin{Resp}\Rsign Ut non recédant a me.\end{Resp}
\switchcolumn
\EN
\begin{Resp}\Rsign I will make an everlasting Covenant with them, and I will not cease from doing them good, and I will put my fear in their hearts, \Star{} So that they shall not depart from me.\end{Resp}
\begin{Resp}\Vsign Yea, I will rejoice over them, to do them good with my whole Heart.\end{Resp}
\begin{Resp}\Rsign So that they shall not depart from me.\end{Resp}
\switchcolumn*

\LA
\LessonHead{Lectio II}
\switchcolumn
\EN
\LessonHead{Lesson II}
\switchcolumn*

\LA
\Cite{Jer 30:18-19; 30:21-24}
\switchcolumn
\EN
\Cite{Jer 30:18-19; 30:21-24}
\switchcolumn*

\LA
\noindent Hæc dicit Dóminus: Ecce ego convértam conversiónem tabernaculórum Jacob, et tectis ejus miserébor, et ædificábitur cívitas in excélso suo, et templum juxta órdinem suum fundábitur, et egrediétur de eis laus, voxque ludéntium. Et erit dux ejus ex eo, et princeps de médio ejus producétur; et, applicábo eum et accédet ad me. Quis enim iste est qui ápplicet cor suum ut appropínquet mihi? ait Dóminus. Et éritis mihi in pópulum, et ego ero vobis in Deum. Ecce turbo Dómini, furor egrédiens, procélla ruens; in cápite impiórum conquiéscet. Non avértet iram indignatiónis Dóminus, donec fáciat et cómpleat cogitatiónem Cordis sui: in novíssimo diérum intellegétis ea.\par
\switchcolumn
\EN
\noindent Thus saith the Lord: Behold I will bring back the captivity of the pavilions of Jacob, and will have pity on his houses, and the city shall be built in her high place, and the temple shall be founded according to the order thereof. And out of them shall come forth praise, and the voice of them that play: And their leader shall be of themselves: and their prince shall come forth from the midst of them: and I will bring him near, and he shall come to me: for who is this that setteth his heart to approach to me, saith the Lord? And you shall be my people: and I will be your God. Behold the whirlwind of the Lord, his fury going forth, a violent storm, it shall rest upon the head of the wicked. The Lord will not turn away the wrath of his indignation, till he have executed and performed the thought of his heart: in the latter days you shall understand these things.\par
\switchcolumn*

\LA
\begin{Resp}\Rsign Si inimícus meus maledixísset mihi, sustinuíssem útique \Star{} Tu vero homo unánimis qui simul mecum dulces capiébas cibos.\end{Resp}
\begin{Resp}\Vsign Et si is qui me óderat super me magna locútus fuísset, abscondíssem me fórsitan ab eo.\end{Resp}
\begin{Resp}\Rsign Tu vero homo unánimis qui simul mecum dulces capiébas cibos.\end{Resp}
\switchcolumn
\EN
\begin{Resp}\Rsign It was not an open enemy that has done me this dishonour, for then I could have borne it, \Star{} But it was even thou, mine own familiar friend, who did also eat of my Bread.\end{Resp}
\begin{Resp}\Vsign Neither was it mine adversary that did magnify himself against me, for then peradventure I would have hid myself from him.\end{Resp}
\begin{Resp}\Rsign But it was even thou, mine own familiar friend, who did also eat of my Bread.\end{Resp}
\switchcolumn*

\LA
\LessonHead{Lectio III}
\switchcolumn
\EN
\LessonHead{Lesson III}
\switchcolumn*

\LA
\Cite{Jer 31:1-3; 31:31-33}
\switchcolumn
\EN
\Cite{Jer 31:1-3; 31:31-33}
\switchcolumn*

\LA
\noindent In témpore illo, dicit Dóminus, ero Deus univérsis cognatiónibus Israël, et ipsi erunt mihi in pópulum. Hæc dicit Dóminus: Invénit grátiam in desérto pópulus qui remánserat a gládio; vadet ad réquiem suam Israël. Longe Dóminus appáruit mihi. Et in caritáte perpétua diléxi te: ídeo attráxi te, míserans. Ecce dies vénient, dicit Dóminus: et fériam dómui Israël et dómui Juda fœdus novum: non secúndum pactum, quod pépigi cum pátribus eórum in die, qua apprehéndi manum eórum, ut edúcerem eos de Terra Ægýpti: pactum quod írritum fecérunt, et ego dominátus sum eórum, dicit Dóminus. Sed hoc erit pactum, quod fériam cum domo Israël: post dies illos dicit Dóminus: Dabo legem meam in viscéribus eórum, et in corde eórum scribam eam: et ero eis in Deum, et ipsi erunt mihi in pópulum.\par
\switchcolumn
\EN
\noindent At that time, saith the Lord, I will be the God of all the families of Israel, and they shall be my people. Thus saith the Lord: The people that were left and escaped from the sword, found grace in the desert: Israel shall go to his rest. The Lord hath appeared from afar to me. Yea I have loved thee with an everlasting love, therefore have I drawn thee, taking pity on thee. Behold the days shall come, saith the Lord, and I will make a new covenant with the house of Israel, and with the house of Juda: Not according to the covenant which I made with their fathers, in the day that I took them by the hand to bring them out of the land of Egypt: the covenant which they made void, and I had dominion over them, saith the Lord. But this shall be the covenant that I will make with the house of Israel, after those days, saith the Lord: I will give my law in their bowels, and I will write it in their heart: and I will be their God, and they shall be my people.\par
\switchcolumn*

\LA
\begin{Resp}\Rsign Cum essémus mórtui peccátis, convivificávit nos Deus in Christo \Star{} Propter nímiam caritátem suam qua diléxit nos.\end{Resp}
\begin{Resp}\Vsign Ut osténderet in sǽculis superveniéntibus abundántes divítias grátiæ suæ.\end{Resp}
\begin{Resp}\Rsign Propter nímiam caritátem suam qua diléxit nos.\end{Resp}
\begin{Resp}\Vsign Glória Patri, et Fílio, \Star{} et Spirítui Sancto.\end{Resp}
\begin{Resp}\Rsign Propter nímiam caritátem suam qua diléxit nos.\end{Resp}
\switchcolumn
\EN
\begin{Resp}\Rsign And we, being dead in our sins, hath God quickened together with Christ, \Star{} For his great love wherewith he hath loved us.\end{Resp}
\begin{Resp}\Vsign That in the ages to come he might shew the exceeding riches of his grace.\end{Resp}
\begin{Resp}\Rsign For his great love wherewith he hath loved us.\end{Resp}
\begin{Resp}\Vsign Glory be to the Father, and to the Son, \Star{} and to the Holy Ghost.\end{Resp}
\begin{Resp}\Rsign For his great love wherewith he hath loved us.\end{Resp}
\switchcolumn*

\LA
\LessonHead{Lectio IV}
\switchcolumn
\EN
\LessonHead{Lesson IV}
\switchcolumn*

\LA
\noindent Inter mira sacræ doctrínæ pietatísque increménta, quibus divínæ Sapiéntiæ consília clárius in dies Ecclésiæ manifestántur, vix áliud magis conspícuum est quam triumphális progréssio cultus sacratíssimi Cordis Jesu. Sǽpius quidem, priórum decúrsu témporum, Patres, Doctóres, Sancti, Redemptóris nostri amórem celebrárunt: vulnus in látere Christi apértum ómnium gratiárum arcánum dixérunt fontem. At inde a médio ævo, cum tenerióre quadam erga sanctíssimam Salvatóris Humanitátem religióne fidéles áffici cœpti sunt, ánimæ contemplatívæ per plagam illam ad ipsum Cor, amóre hóminum vulnerátum, penetráre fere solébant. Atque ex eo témpore hæc contemplátio sanctíssimis quibúsque ita familiáris evásit, ut neque régio neque ordo religiósus sit, in quibus non insígnia, hac ætáte, ejus reperiántur testimónia. Próximis demum sǽculis, eóque potíssimum témpore quo hærétici, sub falsæ pietátis título, a sanctíssima Eucharístia Christiános detérrere conabántur, cultus sacratíssimo Cordi públice exhibéri cœptus est, ópera imprímis sancti Joánnis Eudes, qui auctor litúrgici cultus sacrórum Córdium Jesu et Maríæ haud immérito nuncupátur.\par
\switchcolumn
\EN
\noindent Among the wonderful developments of sacred teaching and piety, by which the plans of the divine Wisdom are daily made clear to the Church, hardly any is more manifest than the triumphant progress made by the devotion of the most Sacred Heart of Jesus. Very often indeed, during the course of past ages, Fathers, Doctors, and Saints have celebrated our Redeemer's love: and they have said, that the wound opened in the side of Christ was the hidden fountain of all graces. Moreover, from the Middle Ages onward, when the faithful began to show a more tender piety towards the most sacred Humanity of the Saviour, contemplative souls became accustomed to penetrate through that wound almost to the very Heart itself, wounded for the love of men. And from that time, this form of contemplation became so familiar to all persons of saintly life, that there was no country or religious order in which, during this period, witnesses to it were not to be found. Finally, during recent centuries, and most especially at that period when heretics, in the name of a false piety, strove to discourage Christians from receiving the most Holy Eucharist, the veneration of the most Sacred Heart began to be openly practised, principally through the exertions of St. John Eudes, who is by no means unworthily called the founder of the liturgical worship of the Sacred Hearts of Jesus and Mary.\par
\switchcolumn*

\LA
\begin{Resp}\Rsign Prope est Dóminus ómnibus invocántibus eum, \Star{} Omnibus invocántibus eum in veritáte.\end{Resp}
\begin{Resp}\Vsign Miserátor et miséricors Dóminus, pátiens et multum miséricors.\end{Resp}
\begin{Resp}\Rsign Omnibus invocántibus eum in veritáte.\end{Resp}
\switchcolumn
\EN
\begin{Resp}\Rsign The Lord is nigh unto all them that call upon him, \Star{} Yea, unto all such as call upon him faithfully.\end{Resp}
\begin{Resp}\Vsign The Lord is full of compassion and mercy, long-suffering and of great goodness.\end{Resp}
\begin{Resp}\Rsign Yea, unto all such as call upon him faithfully.\end{Resp}
\switchcolumn*

\LA
\LessonHead{Lectio V}
\switchcolumn
\EN
\LessonHead{Lesson V}
\switchcolumn*

\LA
\noindent Verum, ad cultum sacratíssimi Cordis Jesu plene perfectéque constituéndum, eundémque per totum orbem propagándum, Deus ipse sibi instruméntum elégit humíllimam ex órdine Visitatiónis vírginem, sanctam Margarítam Maríam Alacóque, cui, a prima quidem ætáte jam in Eucharístiæ Sacraméntum amóre flagránti, Christus Dóminus sæpenúmero appárens, divíni Cordis sui et divítias et optáta significáre dignátus est. Quarum apparitiónum celebérrima illa est, qua ei ante Eucharístiam oránti Jesus conspiciéndum se dedit, sacratíssimum Cor osténdit et conquéstus quod, pro imménsa sua caritáte, nihil nisi ingratórum hóminum contumélias recíperet, ipsi præcépit ut novum festum, féria sexta post Octávam Córporis Christi, instituéndum curáret, quo Cor suum honóre débito colerétur, atque injúriæ sibi in Sacraménto amóris a peccatóribus illátæ dignis expiaréntur obséquiis. Quot autem quantásque Dei fámula in Christi mandátis exsequéndis expérta sit difficultátes, nemo est qui ignóret; sed ab ipso Dómino confirmáta, atque a religiósis ánimæ suæ moderatóribus, qui incredíbili quodam ardóre ad hunc cultum promovéndum laborárunt, strénue adjúta, múnere sibi cǽlitus commísso fidéliter fungi ad mortem usque non déstitit.\par
\switchcolumn
\EN
\noindent But in order to establish fully and entirely the worship of the most Sacred Heart of Jesus, and to spread the same throughout the whole world, God himself chose as his instrument a most humble virgin from the order of the Visitation, St. Margaret Mary Alacoque, who even in her earliest years already had a burning love for the Sacrament of the Eucharist, and to whom Christ the Lord had very many times appeared, and was pleased to make known the riches and the desires of his divine Heart. The most famous of these apparitions was that in which Jesus revealed himself to her in prayer before the blessed Sacrament, shewed her his most Sacred Heart, and, complaining that in return for his unbounded love, he met with nothing but outrages and ingratitude from mankind, he ordered her to concern herself with the establishment of a new feast, on the Friday after the Octave of Corpus Christi, on which his Heart should be venerated with due honour, and that the insults offered him by sinners in the Sacrament of love should be expiated by worthy satisfaction. But there is no one who knoweth not how many and how great were the obstacles which the handmaid of God experienced, in carrying out the commands of Christ; but, endowed with strength by the Lord himself, and actively aided by her pious spiritual directors, who exerted themselves with an almost unbelievable zeal, up to the time of her death she never ceased faithfully to carry out the duty entrusted to her by heaven.\par
\switchcolumn*

\LA
\begin{Resp}\Rsign Confíteor tibi, Pater, Dómine cæli et terræ, quia abscondísti hæc a sapiéntibus et prudéntibus \Star{} Et revelásti ea párvulis.\end{Resp}
\begin{Resp}\Vsign Ita, Pater, quóniam sic fuit plácitum ante te.\end{Resp}
\begin{Resp}\Rsign Et revelásti ea párvulis.\end{Resp}
\switchcolumn
\EN
\begin{Resp}\Rsign I thank thee, O Father, Lord of heaven and earth, because thou hast hid these things from the wise and prudent; \Star{} Yea, thou hast revealed them unto babes.\end{Resp}
\begin{Resp}\Vsign Even so, Father, for so it seemed good in thy sight.\end{Resp}
\begin{Resp}\Rsign Yea, thou hast revealed them unto babes.\end{Resp}
\switchcolumn*

\LA
\LessonHead{Lectio VI}
\switchcolumn
\EN
\LessonHead{Lesson VI}
\switchcolumn*

\LA
\noindent Anno tandem millésimo septingentésimo sexagésimo quinto, Clemens décimus tértius Póntifex Máximus offícium et missam in honórem sacratíssimi Cordis Jesu approbávit; Pius vero nonus festum ad univérsam Ecclésiam exténdit. Exínde, cultus sacratíssimi Cordis, quasi flumen exúndans, prolútis impediméntis ómnibus, per totum se orbem effúdit, et, novo illucescénte sǽculo, jubilǽo indícto, Leo décimus tértius humánum genus univérsum sacratíssimo Cordi devótum vóluit. Quæ consecrátio, in ómnibus quidem cathólici orbis ecclésiis, solémni ritu perácta, ingens áttulit devotiónis hujus increméntum, et ad eam non solum pópulos, verum étiam singuláres famílias addúxit, quæ Divíno Cordi innumerábiles se dévovent, regióque ejus império subíciunt. Dénique, Pius undécimus Póntifex Máximus, quo plénius festi solémnitas pópuli christiáni devotióni tam late paténti respondéret, sacratíssimi Cordis Jesu festum ad ritum dúplicem primæ classis cum octáva evéxit; ac prætérea, ut violáta jura Christi summi Regis ac Dómini amantíssimi resarciréntur, populorúmque peccáta defleréntur, eódem festo die piaculárem precatiónem in ómnibus christiáni orbis templis quotánnis recitándam mandávit.\par
\switchcolumn
\EN
\noindent At length, in the year 1765, the Supreme Pontiff Clement XIII approved the Mass and Office in honour of the most Sacred Heart of Jesus; and Pius IX extended the feast to the universal Church. From then on the worship of the most Sacred Heart, like an overflowing river, washing away all obstacles, hath poured itself forth over all the earth, and, at the dawn of the new century, Leo XIII, having proclaimed a jubilee, decided to dedicate the whole human race to the most Sacred Heart. This consecration was actually carried out with solemn rites in all the churches of the Catholic world, and brought about a great increase of this devotion, leading not only nations but even private families to it, who in countless numbers dedicated themselves to the Divine Heart, and submitted themselves to its royal sway. Lastly, the Sovereign Pontiff Pius XI, in order that, by its solemnity, the feast might answer more fully to the greatly widespread devotion of the Christian people, raised the feast of the most Sacred Heart of Jesus to the rite of a double of the first class, with an octave; and moreover, that the violated rights of Christ, the supreme King and most loving Lord, might be repaired, and that the sins of the nations might be bewailed, he ordered that annually, on that same feast-day, there should be recited an expiatory form of prayer in all the churches of the Christian world.\par
\switchcolumn*

\LA
\begin{Resp}\Rsign Omnes gentes quascúmque fecísti, vénient \Star{} Et adorábunt coram te, Dómine.\end{Resp}
\begin{Resp}\Vsign Et glorificábunt nomen tuum, quóniam magnus es tu, et fáciens mirabília.\end{Resp}
\begin{Resp}\Rsign Et adorábunt coram te, Dómine.\end{Resp}
\begin{Resp}\Vsign Glória Patri, et Fílio, \Star{} et Spirítui Sancto.\end{Resp}
\begin{Resp}\Rsign Et adorábunt coram te, Dómine.\end{Resp}
\switchcolumn
\EN
\begin{Resp}\Rsign All nations whom thou hast made shall come, \Star{} And they shall worship thee, O Lord.\end{Resp}
\begin{Resp}\Vsign Yea, they shall glorify thy Name, for thou art great, and doest wondrous things.\end{Resp}
\begin{Resp}\Rsign And they shall worship thee, O Lord.\end{Resp}
\begin{Resp}\Vsign Glory be to the Father, and to the Son, \Star{} and to the Holy Ghost.\end{Resp}
\begin{Resp}\Rsign And they shall worship thee, O Lord.\end{Resp}
\switchcolumn*

\LA
\LessonHead{Lectio VII}
\switchcolumn
\EN
\LessonHead{Lesson VII}
\switchcolumn*

\LA
\LessonTitle{Léctio sancti Evangélii secúndum Joánnem}
\switchcolumn
\EN
\LessonTitle{From the Holy Gospel according to John}
\switchcolumn*

\LA
\Cite{Joannes 19:31-37}
\switchcolumn
\EN
\Cite{John 19:31-37}
\switchcolumn*

\LA
\noindent In illo témpore: Judǽi, quóniam parascéve erat, ut non remanérent in cruce córpora sábbato (erat enim magnus dies ille sábbati) rogavérunt Pilátum, ut frangeréntur eórum crura et tolleréntur. Et réliqua.\par
\switchcolumn
\EN
\noindent At that time: The Jews, because it was the Preparation, that the bodies should not remain upon the cross on the Sabbath Day, for that Sabbath Day was an high day, besought Pilate that their legs might broken, and that they might be taken away. And so on.\par
\switchcolumn*

\LA
\LessonTitle{Homilía sancti Bonaventúræ Epíscopi}
\switchcolumn
\EN
\LessonTitle{A Homily by St. Bonaventure the Bishop}
\switchcolumn*

\LA
\Source{Liber de ligno vitæ, num. 30}
\switchcolumn
\EN
\Source{Book of the Tree of Life, num. 30}
\switchcolumn*

\LA
\noindent Ut de látere Christi dormiéntis in cruce formarétur Ecclésia, et Scriptúra implerétur quæ dicit: Vidébunt in quem transfixérunt, divína est ordinatióne indúltum ut unus mílitum láncea latus illud sacrum aperiéndo perfóderet, quátenus sánguine cum aqua manánte, prétium effunderétur nostræ salútis quod a fonte scílicet Cordis arcáno profúsum, vim daret sacraméntis Ecclésiæ ad vitam grátiæ conferéndam, essétque jam in Christo vivéntibus póculum fontis vivi, saliéntis in vitam ætérnam. Surge ígitur, ánima amíca Christi, vigiláre non cesses, ibi os appóne, ut háurias aquas de fóntibus salvatóris.\par
\switchcolumn
\EN
\noindent In order that the Church might be taken out of the side of Christ, in his deep sleep on the Cross, and that the Scripture might be fulfilled which saith: They shall look on him whom they pierced: it was divinely ordained that one of the soldiers should pierce his sacred side with a spear, and open it. Then forthwith there came flowing out blood and water, which was the price of our salvation, pouring forth from its mountain-source, in sooth, from the secret places of his Heart, to give power to the Sacraments of the Church, to bestow the life of grace, and to be as a saving drink of living waters, flowing up to life eternal for those who were already quickened in Christ. Arise, then, O soul beloved of Christ. Cease not thy vigilance, place there thy lips, and drink the waters from the fount of salvation.\par
\switchcolumn*

\LA
\begin{Resp}\Rsign Ego si exaltátus fúero a terra \Star{} Omnia traham ad meípsum.\end{Resp}
\begin{Resp}\Vsign Hoc autem dicébat signíficans qua morte esset moritúrus.\end{Resp}
\begin{Resp}\Rsign Omnia traham ad meípsum.\end{Resp}
\switchcolumn
\EN
\begin{Resp}\Rsign If I be lifted up, \Star{} I will draw all men unto me.\end{Resp}
\begin{Resp}\Vsign This he said, signifying what death he should die.\end{Resp}
\begin{Resp}\Rsign I will draw all men unto me.\end{Resp}
\switchcolumn*

\LA
\LessonHead{Lectio VIII}
\switchcolumn
\EN
\LessonHead{Lesson VIII}
\switchcolumn*

\LA
\Source{De vite mystica, cap. 3}
\switchcolumn
\EN
\Source{Of the mystical vine, Chap. 3}
\switchcolumn*

\LA
\noindent Quia semel vénimus ad Cor Dómini Jesu dulcíssimi, et bonum est nos hic esse, non fácile evellámur ab eo. O quam bonum et jucúndum habitáre in Corde hoc. Bonus thesáurus, pretiósa margaríta Cor tuum, óptime Jesu, quam, fosso agro córporis tui, invenímus. Quis hanc margarítam abíciat? Quin pótius, dabo omnes margarítas, cogitatiónes et affectiónes meas commutábo et comparábo illam mihi, jactans omnem cogitátum meum in Cor boni Jesu, et sine fallácia illud me enútriet. Hoc ígitur tuo et meo Corde, dulcíssime Jesu, invénto, orábo te Deum meum: admítte in sacrárium exauditiónis preces meas: immo me totum trahe in Cor tuum.\par
\switchcolumn
\EN
\noindent Because we are now come to the sweet Heart of Jesus, and because it is good for us to be here, let us not too soon turn away therefrom. O how good and joyful a thing it is to dwell in this Heart. What a good treasure, what a precious pearl, is thy Heart, O most excellent Jesu, which we have found hidden in the pit which hath been dug in this field, namely, in thy body. Who would cast away such a pearl? Nay, rather, for this same I would give all my pearls. I will sell all my thoughts and affections, and buy the same for myself, turning all my thoughts to the Heart of the good Jesus, and without fail it will support me. Therefore, o most sweet Jesu, finding this Heart that is thine and mine, I will pray to thee, my God: admit my prayers into the shrine of hearkening: and draw me even more altogether into thy Heart.\par
\switchcolumn*

\LA
\begin{Resp}\Rsign Simus ergo imitatóres Dei \Star{} Et ambulémus in dilectióne.\end{Resp}
\begin{Resp}\Vsign Sicut et Christus diléxit nos et trádidit semetípsum pro nobis.\end{Resp}
\begin{Resp}\Rsign Et ambulémus in dilectióne.\end{Resp}
\begin{Resp}\Vsign Glória Patri, et Fílio, \Star{} et Spirítui Sancto.\end{Resp}
\begin{Resp}\Rsign Et ambulémus in dilectióne.\end{Resp}
\switchcolumn
\EN
\begin{Resp}\Rsign Be ye therefore followers of God; \Star{} And walk ye therefore in love.\end{Resp}
\begin{Resp}\Vsign For Christ also hath loved us and hath given himself for us.\end{Resp}
\begin{Resp}\Rsign And walk ye therefore in love.\end{Resp}
\begin{Resp}\Vsign Glory be to the Father, and to the Son, \Star{} and to the Holy Ghost.\end{Resp}
\begin{Resp}\Rsign And walk ye therefore in love.\end{Resp}
\switchcolumn*

\LA
\LessonHead{Lectio IX}
\switchcolumn
\EN
\LessonHead{Lesson IX}
\switchcolumn*

\LA
\noindent Ad hoc enim perforátum est latus tuum, ut nobis páteat intróitus. Ad hoc vulnerátum est Cor tuum, ut in illo ab exterióribus turbatiónibus absolúti habitáre possímus. Nihilóminus et proptérea vulnerátum est, ut per vulnus visíbile, vulnus amóris invisíbile videámus. Quómodo hic ardor mélius posset osténdi, nisi quod non solum corpus, verum étiam ipsum Cor láncea vulnerári permísit? Carnále ergo vulnus, vulnus spirituále osténdit. Quis illud Cor tam vulnerátum non díligat? quis tam amántem non rédamet? quis tam castum non amplectátur? Nos ígitur adhuc in carne manéntes, quantum póssumus, amántem redamémus, amplectámur vulnerátum nostrum, cujus ímpii agrícolæ fodérunt manus et pedes, latus et Cor; oremúsque ut cor nostrum, adhuc durum et impœ́nitens, amóris sui vínculo constríngere et jáculo vulneráre dignétur.\par
\switchcolumn
\EN
\noindent For to this end was thy side pierced, that an entry might be open unto us. To this end was thy Heart wounded, that in it we might be able to dwell secure from alarms from without. And it was wounded none the less on this account that, because of the visible wound, we may perceive the wound of love which is invisible. How could this fire of love better shine forth than for him to permit that not only his body, but that even his Heart, should be wounded with the spear? Who would not love that Heart so wounded? Who would not, in return, love one who is so loving? Who would not embrace one so chaste? Wherefore let us who are in the flesh love in return, as much as we can, him who so loveth, embrace our wounded one, whose hands and feet, side and Heart, have been pierced by wicked husbandmen; and let us pray that he may deign to bind our hearts, still hard and impenitent, with the chain of his love, and wound them with the dart thereof.\par
\switchcolumn*

\LA
\TeDeum{Te Deum}
\switchcolumn
\EN
\TeDeum{Te Deum}
\switchcolumn*

\end{paracol}

\Office{Die II infra Octavam SSmi Cordis Jesu}{Second Day within the Octave of the Sacred Heart}{Semiduplex}
\addcontentsline{toc}{subsection}{Die II infra Octavam SSmi Cordis Jesu \textit{— Second Day within the Octave of the Sacred Heart}}
\begin{paracol}{2}
\LA
\LessonHead{Lectio I}
\switchcolumn
\EN
\LessonHead{Lesson I}
\switchcolumn*

\LA
\LessonTitle{De libro primo Regum}
\switchcolumn
\EN
\LessonTitle{Lesson from the first book of Samuel}
\switchcolumn*

\LA
\Cite{1 Reg 9:1-4}
\switchcolumn
\EN
\Cite{1 Sam 9:1-4}
\switchcolumn*

\LA
\noindent Et erat vir de Bénjamin nómine Cis, fílius Abiel, fílii Seror, fílii Béchorath, fílii Aphia, fílii viri Jémini, fortis róbore. Et erat ei fílius vocábulo Saul eléctus et bonus, et non erat vir de fíliis Israël mélior illo: ab húmero et sursum eminébat super omnem pópulum. Períerant autem ásinæ Cis patris Saul, et dixit Cis ad Saul fílium suum: Tolle tecum unum de púeris et consúrgens vade et quære ásinas. Qui, cum transíssent per montem Ephraim et per terram Salísa et non inveníssent, transiérunt étiam per terram Salim, et non erant, sed et per terram Jémini, et mínime reperérunt.\par
\switchcolumn
\EN
\noindent Now there was a man of Benjamin whose name was Cis, the son of Abiel, the son of Seror, the son of Bechorath, the son of Aphia, the son of a man of Jemini, valiant and strong. And he had a son whose name was Saul, a choice and goodly man, and there was not among the children of Israel a goodlier person than he: from his shoulders and upward he appeared above all the people. And the asses of Cis, Sauls father, were lost: and Cis said to his son Saul: Take one of the servants with thee, and arise, go, and seek the asses. And when they had passed through mount Ephraim, And through the land of Salisa, and had not found them, they passed also through the land of Salim, and they were not there: and through the land of Jemini, and found them not.\par
\switchcolumn*

\LA
\begin{Resp}\Rsign Peccávi super númerum arénæ maris, et multiplicáta sunt peccáta mea: et non sum dignus vidére altitúdinem cæli præ multitúdine iniquitátis meæ: quóniam irritávi iram tuam, \Star{} Et malum coram te feci.\end{Resp}
\begin{Resp}\Vsign Quóniam iniquitátem meam ego cognósco: et delíctum meum contra me est semper, quia tibi soli peccávi.\end{Resp}
\begin{Resp}\Rsign Et malum coram te feci.\end{Resp}
\switchcolumn
\EN
\begin{Resp}\Rsign My sins are many, yea, they are more in number than the sands of the sea; I am not worthy to look up toward heaven because of the multitude of my iniquities; for I have provoked thee to anger; \Star{} And done evil in thy sight.\end{Resp}
\begin{Resp}\Vsign For I acknowledge my transgression, and my sin is ever before me, for against thee only have I sinned.\end{Resp}
\begin{Resp}\Rsign And done evil in thy sight.\end{Resp}
\switchcolumn*

\LA
\LessonHead{Lectio II}
\switchcolumn
\EN
\LessonHead{Lesson II}
\switchcolumn*

\LA
\Cite{1 Reg 9:5-8}
\switchcolumn
\EN
\Cite{1 Sam 9:5-8}
\switchcolumn*

\LA
\noindent Cum autem veníssent in terram Suph, dixit Saul ad púerum, qui erat cum eo: Veni et revertámur, ne forte dimíserit pater meus ásinas et sollícitus sit pro nobis. Qui ait ei: Ecce vir Dei est in civitáte hac, vir nóbilis: omne quod lóquitur, sine ambiguitáte venit; nunc ergo eámus illuc, si forte índicet nobis de via nostra, propter quam vénimus. Dixítque Saul ad púerum suum: Ecce íbimus: quid ferémus ad virum Dei? panis defécit in sitárciis nostris, et spórtulam non habémus, ut demus hómini Dei, nec quidquam áliud. Rursum puer respóndit Sauli et ait: Ecce invénta est in manu mea quarta pars statéris argénti: demus hómini Dei, ut índicet nobis viam nostram.\par
\switchcolumn
\EN
\noindent And when they were come to the land of Suph, Saul said to the servant that was with him: Come, let us return, lest perhaps my father forget the asses, and be concerned for us. And he said to him: Behold there is a man of God in this city, a famous man: all that he saith, cometh certainly to pass. Now therefore let us go thither, perhaps he may tell us of our way, for which we are come. And Saul said to his servant: Behold we will go: but what shall we carry to the man of God? the bread is spent in our bags: and we have no present to make to the man of God, nor any thing at all. The servant answered Saul again, and said: Behold there is found in my hand the fourth part of a sicle of silver, let us give it to the man of God, that he may tell us our way.\par
\switchcolumn*

\LA
\begin{Resp}\Rsign Exaudísti, Dómine, oratiónem servi tui, ut ædificárem templum nómini tuo: \Star{} Bénedic et sanctífica domum istam in sempitérnum, Deus Israël.\end{Resp}
\begin{Resp}\Vsign Dómine, qui custódis pactum cum servis tuis, qui ámbulant coram te in toto corde suo.\end{Resp}
\begin{Resp}\Rsign Bénedic et sanctífica domum istam in sempitérnum, Deus Israël.\end{Resp}
\switchcolumn
\EN
\begin{Resp}\Rsign O Lord, Thou hast hearkened unto the prayer of thy servant, that I might build a temple unto thy Name, \Star{} O God of Israel, bless Thou, and hallow this house for ever.\end{Resp}
\begin{Resp}\Vsign O Lord, Who keepest covenant with thy servants that walk before thee in all their heart.\end{Resp}
\begin{Resp}\Rsign O God of Israel, bless Thou, and hallow this house for ever.\end{Resp}
\switchcolumn*

\LA
\LessonHead{Lectio III}
\switchcolumn
\EN
\LessonHead{Lesson III}
\switchcolumn*

\LA
\Cite{1 Reg 9:14-17}
\switchcolumn
\EN
\Cite{1 Sam 9:14-17}
\switchcolumn*

\LA
\noindent Et ascendérunt in civitátem. Cumque illi ambulárent in médio urbis, appáruit Sámuel egrédiens óbviam eis, ut ascénderet in excélsum. Dóminus autem reveláverat aurículam Samuélis, ante unam diem quam veníret Saul, dicens: Hac ipsa hora, quæ nunc est, cras mittam virum ad te de terra Bénjamin, et unges eum ducem super pópulum meum Israël, et salvábit pópulum meum de manu Philisthinórum, quia respéxi pópulum meum; venit enim clamor eórum ad me. Cumque aspexísset Sámuel Saulem Dóminus dixit ei: Ecce vir, quem díxeram tibi: iste dominábitur pópulo meo.\par
\switchcolumn
\EN
\noindent And they went up into the city. And when they were walking in the midst of the city, behold Samuel was coming out over against them, to go up to the high place. Now the Lord had revealed to the ear of Samuel the day before Saul came, saying: Tomorrow about this same hour I will send thee a man of the land of Benjamin, and thou shalt anoint him to be ruler over my people Israel: and he shall save my people out of the hand of the Philistines: for I have looked down upon my people, because their cry is come to me. And when Samuel saw Saul, the Lord said to him: Behold the man, of whom I spoke to thee, this man shall reign over my people.\par
\switchcolumn*

\LA
\begin{Resp}\Rsign Audi, Dómine, hymnum et oratiónem, quam servus tuus orat coram te hódie, ut sint óculi tui apérti, et aures tuæ inténtæ, \Star{} Super domum istam die ac nocte.\end{Resp}
\begin{Resp}\Vsign Réspice, Dómine, de sanctuário tuo, et de excélso cælórum habitáculo.\end{Resp}
\begin{Resp}\Rsign Super domum istam die ac nocte.\end{Resp}
\begin{Resp}\Vsign Glória Patri, et Fílio, \Star{} et Spirítui Sancto.\end{Resp}
\begin{Resp}\Rsign Super domum istam die ac nocte.\end{Resp}
\switchcolumn
\EN
\begin{Resp}\Rsign Hearken, O Lord, unto the cry and to the prayer which thy servant prayeth before thee today, that thine eyes may be open and thine ears attend; \Star{} Toward this house day and night.\end{Resp}
\begin{Resp}\Vsign Look down from thine high and holy place, O Lord, even from heaven thy dwelling.\end{Resp}
\begin{Resp}\Rsign Toward this house, day and night.\end{Resp}
\begin{Resp}\Vsign Glory be to the Father, and to the Son, \Star{} and to the Holy Ghost.\end{Resp}
\begin{Resp}\Rsign Toward this house, day and night.\end{Resp}
\switchcolumn*

\LA
\LessonHead{Lectio IV}
\switchcolumn
\EN
\LessonHead{Lesson IV}
\switchcolumn*

\LA
\LessonTitle{Ex lítteris Encýclicis Pii Papæ undécimi}
\switchcolumn
\EN

\switchcolumn*

\LA
\Source{Encýclica Miserentíssimus Redémptor}
\switchcolumn
\EN
\Source{From the Encyclical Letter of Pope Pius XI}
\switchcolumn*

\LA
\noindent Inter cétera infinítæ Redemptóris nostri benignitátis documénta, illud potíssimum elúcet, quod, defervescénte christifidélium caritáte, ipsa Dei cáritas ad honorándum peculiári cultu propósita est, ejúsque bonitátis divítiæ late patefáctæ sunt per eam religiónis formam qua sacratíssimum Cor Jesu cólitur, « in quo sunt omnes thesáuri sapiéntiæ et sciéntiæ abscónditi ». Nam, ut quondam humánæ genti a Noética arca exeúnti, amíci fœ́deris signum illucéscere Deus vóluit « arcum apparéntem in núbibus », sic turbulentíssimis recentióris ævi tempóribus, cum vaférrima ómnium sérperet hǽresis illa janseniána, amóri in Deum pietatíque inimíca, quæ Deum non tam diligéndum ut patrem quam extimescéndum ut implacábilem júdicem prædicábat, benigníssimus Jesus Cor suum sacratíssimum, quasi pacis et caritátis vexíllum elátum géntibus osténdit, haud dúbiam porténdens in certámine victóriam.\par
\switchcolumn
\EN
\noindent Among all other proofs of the infinite kindness of our Redeemer, this one is especially conspicuous, that, as the love of the Christian believers grew cold, he, Divine Love itself, was proposed to be honoured by a special devotion, and that the rich treasures of his goodness were thrown wide open by means of that form of worship with which we honour the most Sacred Heart of Jesus, “in whom are hid all the treasures of wisdom and knowledge.” For, as formerly God wished to give light to the human race as they came out of Noah's ark by the signal of a treaty of friendship, “a bow appearing in the clouds,” so, in those most troublous times of a more recent age, when that most subtle of heresies, Jansenism, was everywhere creeping in, and enemy of the love of God and of piety, preaching that God was not so much to be loved as a father, as to be feared as an unrelenting judge, the most kind Jesus manifested unto the nations his most Sacred Heart, borne on high like unto a banner of peace and love, an augury of certain victory in battle.\par
\switchcolumn*

\LA
\begin{Resp}\Rsign Prope est Dóminus ómnibus invocántibus eum, \Star{} Omnibus invocántibus eum in veritáte.\end{Resp}
\begin{Resp}\Vsign Miserátor et miséricors Dóminus, pátiens et multum miséricors.\end{Resp}
\begin{Resp}\Rsign Omnibus invocántibus eum in veritáte.\end{Resp}
\switchcolumn
\EN
\begin{Resp}\Rsign The Lord is nigh unto all them that call upon him, \Star{} Yea, unto all such as call upon him faithfully.\end{Resp}
\begin{Resp}\Vsign The Lord is full of compassion and mercy, long-suffering and of great goodness.\end{Resp}
\begin{Resp}\Rsign Yea, unto all such as call upon him faithfully.\end{Resp}
\switchcolumn*

\LA
\LessonHead{Lectio V}
\switchcolumn
\EN
\LessonHead{Lesson V}
\switchcolumn*

\LA
\noindent Síquidem appósite felícis recordatiónis decéssor Noster Leo décimus tértius in Lítteris Encyclicis « Annum Sacrum », tantam cultus sacratíssimi Cordis Jesu opportunitátem admirátus, edícere non dubitávit: « Cum Ecclésia per próxima origínibus témpora cæsáreo jugo premerétur, conspécta sublíme adolescénti imperatóri crux, amplíssimæ victóriæ, quæ mox est consecúta, auspex simul atque efféctrix. En álterum hódie oblátum óculis auspicatíssimum divinissimúmque signum: vidélicet Cor Jesu Sacratíssimum, superimpósita cruce, splendidíssimo candóre inter flammas elúcens. In eo omnes collocándæ spes; ex eo hóminum peténda atque exspectánda salus ».\par
\switchcolumn
\EN
\noindent Because Our predecessor, Leo XIII, of happy memory, desiring to obtain the advantages of such a great devotion to the most Sacred Heart of Jesus, in his Encyclical Letter Annum Sacrum most fittingly did not hesitate to proclaim: “When the Church, in the early period of her history, was oppressed by the yoke of the Caesars, a cross appeared in the heavens to a youthful emperor, which was at the same time both the sign and the cause of that most complete victory, which was soon to follow. Behold this day another most auspicious and most holy sign presented to our eyes: that is to say, the most Sacred Heart of Jesus, with a Cross set upon it, shining forth among flames of a most brilliant radiance. In this, all our hopes are to be placed; from this, the salvation of mankind is to be asked for and to be awaited.”\par
\switchcolumn*

\LA
\begin{Resp}\Rsign Confíteor tibi, Pater, Dómine cæli et terræ, quia abscondísti hæc a sapiéntibus et prudéntibus \Star{} Et revelásti ea párvulis.\end{Resp}
\begin{Resp}\Vsign Ita, Pater, quóniam sic fuit plácitum ante te.\end{Resp}
\begin{Resp}\Rsign Et revelásti ea párvulis.\end{Resp}
\switchcolumn
\EN
\begin{Resp}\Rsign I thank thee, O Father, Lord of heaven and earth, because thou hast hid these things from the wise and prudent; \Star{} Yea, thou hast revealed them unto babes.\end{Resp}
\begin{Resp}\Vsign Even so, Father, for so it seemed good in thy sight.\end{Resp}
\begin{Resp}\Rsign Yea, thou hast revealed them unto babes.\end{Resp}
\switchcolumn*

\LA
\LessonHead{Lectio VI}
\switchcolumn
\EN
\LessonHead{Lesson VI}
\switchcolumn*

\LA
\noindent Ac jure id quidem; in illo enim auspicatíssimo signo atque in ea, quæ exínde conséquitur, pietátis forma nonne totíus religiónis summa atque ádeo perfectióris vitæ norma continétur, quippe quæ et ad Christum Dóminum pénitus cognoscéndum mentes condúcat expedítius et ad eúndem veheméntius diligéndum pressiúsque imitándum ánimos infléctat efficácius? Nemo ígitur mirétur, hanc probatíssimam religiónis formam decessóres Nostros continénter et a calumniatórum criminatiónibus vindicásse et summis láudibus extulísse et veheménti provexísse stúdio, prout témporum rerúmque ratiónes postulárent. Dei autem adspiránte númine factum est ut pia christifidélium erga sacratíssimum Cor Jesu volúntas majóra in dies increménta cáperet.\par
\switchcolumn
\EN
\noindent And it is indeed justly so; for in this most auspicious sign and in that which doth follow from it, is there not contained the highest model of piety of the whole of religion, and therefore the rule of the more perfect life, inasmuch as it leadeth our minds the more easily to a deeper knowledge of Christ the Lord, and to a more vehement love of him, and moveth our souls more effectually to a more exact imitation of him? Therefore, no one will be surprised, that Our predecessors have contínuously vindicated this most approved form of devotion from the accusations of objéctors, that they have extolled it with the highest praises, and have promoted it with the most ardent zeal, according as considerations of the period and of affairs in general have demanded. And it hath come to pass by the providence of God, that the devout affection of Christ's faithful people towards the most Sacred Heart of Jesus obtaineth daily a great increase.\par
\switchcolumn*

\LA
\begin{Resp}\Rsign Omnes gentes quascúmque fecísti, vénient \Star{} Et adorábunt coram te, Dómine.\end{Resp}
\begin{Resp}\Vsign Et glorificábunt nomen tuum, quóniam magnus es tu, et fáciens mirabília.\end{Resp}
\begin{Resp}\Rsign Et adorábunt coram te, Dómine.\end{Resp}
\begin{Resp}\Vsign Glória Patri, et Fílio, \Star{} et Spirítui Sancto.\end{Resp}
\begin{Resp}\Rsign Et adorábunt coram te, Dómine.\end{Resp}
\switchcolumn
\EN
\begin{Resp}\Rsign All nations whom thou hast made shall come, \Star{} And they shall worship thee, O Lord.\end{Resp}
\begin{Resp}\Vsign Yea, they shall glorify thy Name, for thou art great, and doest wondrous things.\end{Resp}
\begin{Resp}\Rsign And they shall worship thee, O Lord.\end{Resp}
\begin{Resp}\Vsign Glory be to the Father, and to the Son, \Star{} and to the Holy Ghost.\end{Resp}
\begin{Resp}\Rsign And they shall worship thee, O Lord.\end{Resp}
\switchcolumn*

\LA
\LessonHead{Lectio VII}
\switchcolumn
\EN
\LessonHead{Lesson VII}
\switchcolumn*

\LA
\LessonTitle{Léctio sancti Evangélii secúndum Joánnem}
\switchcolumn
\EN
\LessonTitle{From the Holy Gospel according to John}
\switchcolumn*

\LA
\Cite{Joannes 19:31-37}
\switchcolumn
\EN
\Cite{John 19:31-37}
\switchcolumn*

\LA
\noindent In illo témpore: Judǽi, quóniam parascéve erat, ut non remanérent in cruce córpora sábbato (erat enim magnus dies ille sábbati) rogavérunt Pilátum, ut frangeréntur eórum crura et tolleréntur. Et réliqua.\par
\switchcolumn
\EN
\noindent At that time: The Jews, because it was the Preparation, that the bodies should not remain upon the cross on the Sabbath Day, for that Sabbath Day was an high day, besought Pilate that their legs might broken, and that they might be taken away. And so on/\par
\switchcolumn*

\LA
\LessonTitle{Homilía sancti Joánnis Chrysóstomi}
\switchcolumn
\EN
\LessonTitle{Homily by St. John Chrysostom}
\switchcolumn*

\LA
\Source{Homilia 85, alias 84 in Joánnem, num. 3}
\switchcolumn
\EN
\Source{Homily 85, alias 84 in Joannem, num. 3}
\switchcolumn*

\LA
\noindent Viden, quam fortis sit véritas? Per Judæórum stúdia implétur prophetía. Alia quoque prædíctio finem áccipit. Veniéntes enim mílites aliórum fregérunt crura, Christi non item. Attamen hi, in grátiam Judæórum, ejus latus láncea pérforant, et mórtuo córpori contuméliam ínferunt. O sceléstum et exsecrándum fácinus ! Sed ne turbéris, ne dejiciáris, dilécte. Nam quæ mala illi voluntáte faciébant, veritátem propugnábant; prophetía namque erat: Vidébunt in quem transfixérunt. Neque hoc tantum; sed étiam iis qui infidéles futúri erant, hoc fácinus demonstratióni fuit, ut Thomæ et ipsi simílibus. Ad hæc étiam mystérium ineffábile consummabátur. Exívit enim sanguis et aqua. Non sine causa vel casu hi fontes manárunt, sed quia ex hoc utróque Ecclésia constitúta est.\par
\switchcolumn
\EN
\noindent See ye not how mighty is the Truth? Through the zeal of the Jews the prophecy is fulfilled. And more than one prophecy was fulfilled. For when the soldiers came and brake the legs of the others, they brake not the legs of Christ. But yet these soldiers, to please the Jews, did pierce his side with a lance, and treat his body with contumely. O wicked and accursed crime! But be not troubled, beloved, or cast down. They indeed did it in ill-will, but they unwittingly contended for the truth, as verily the prophecy foretold: They shall look on him whom they have pierced. And more than this, the evil deed served as a demonstration even afterwards to those who were without faith, such as Thomas and others like him. This ineffable mystery was also consummated to another end: Forthwith came there out blood and water. Nor causelessly or by mere chance did these fountains flow, but because the Church was founded with Water and Blood.\par
\switchcolumn*

\LA
\begin{Resp}\Rsign Ego si exaltátus fúero a terra \Star{} Omnia traham ad meípsum.\end{Resp}
\begin{Resp}\Vsign Hoc autem dicébat signíficans qua morte esset moritúrus.\end{Resp}
\begin{Resp}\Rsign Omnia traham ad meípsum.\end{Resp}
\switchcolumn
\EN
\begin{Resp}\Rsign If I be lifted up, \Star{} I will draw all men unto me.\end{Resp}
\begin{Resp}\Vsign This he said, signifying what death he should die.\end{Resp}
\begin{Resp}\Rsign I will draw all men unto me.\end{Resp}
\switchcolumn*

\LA
\LessonHead{Lectio VIII}
\switchcolumn
\EN
\LessonHead{Lesson VIII}
\switchcolumn*

\LA
\noindent Hoc sciunt initiáti, qui per aquam regenerántur, ac per sánguinem et carnem nutriúntur. Hinc inítium mystéria sumunt, ut, cum ad treméndum póculum accésseris, sic vénias ac si ex hoc látere potatúrus esses. Et qui vidit testimónium perhíbuit, et verum est testimónium ejus. Hoc est: Non ab áliis audívi, sed ipse præsens vidi, et verum est testimónium. Mérito sane. Contuméliam narrat, non magnum quid et mirábile, ut possis contra suspicári; verum ille hæreticórum ora cómprimens et futúra prænúntians mystéria, atque conténtum in ipsis thesáurum consíderans, minutátim recénset et quæ gesta sunt. Impléta est autem prophetía illa: Os ejus non commínuent. Etiámsi enim hoc de Judæórum agno dictum sit, tamen propter veritátem figúra præcéssit, et in hoc magis complétum est. Ideo prophétam in médium addúcit.\par
\switchcolumn
\EN
\noindent This is well-known to those who have been initiated, namely, to all who have been regenerated by the Water, and nourished with the Flesh and Blood, so that when thou dost approach to the awesome cup, thou shouldst come as if thou wert about to drink from this very side of Christ. And he that saw it bare record, and his record is true: as though to say: Not from others have I heard it, but I myself was present, and saw it, and therefore my record of it is true. Truly indeed doth he thus speak. For he speaketh to us as of an insult, and not as of something great and wonderful, else thou mightest doubt his testimony; but he, (thus shutting the mouth of heretics, and foretelling future mysteries, and mindful of the treasure to be contained in them,) doth enúmerate one by one the events as they took place. These things were done that the Scriptures should be fulfilled: A bone of him shall not be broken. For even though this was written concerning the lamb which the Jews used for their Passover, nevertheless this lamb was a figure which came first to shew forth the reality yet to come, wherein the prophecy was to be perfectly fulfilled; and that is why the Evangelist quoteth the passage as a prophecy.\par
\switchcolumn*

\LA
\begin{Resp}\Rsign Simus ergo imitatóres Dei \Star{} Et ambulémus in dilectióne.\end{Resp}
\begin{Resp}\Vsign Sicut et Christus diléxit nos et trádidit semetípsum pro nobis.\end{Resp}
\begin{Resp}\Rsign Et ambulémus in dilectióne.\end{Resp}
\begin{Resp}\Vsign Glória Patri, et Fílio, \Star{} et Spirítui Sancto.\end{Resp}
\begin{Resp}\Rsign Et ambulémus in dilectióne.\end{Resp}
\switchcolumn
\EN
\begin{Resp}\Rsign Be ye therefore followers of God; \Star{} And walk ye therefore in love.\end{Resp}
\begin{Resp}\Vsign For Christ also hath loved us and hath given himself for us.\end{Resp}
\begin{Resp}\Rsign And walk ye therefore in love.\end{Resp}
\begin{Resp}\Vsign Glory be to the Father, and to the Son, \Star{} and to the Holy Ghost.\end{Resp}
\begin{Resp}\Rsign And walk ye therefore in love.\end{Resp}
\switchcolumn*

\LA
\LessonHead{Lectio IX}
\switchcolumn
\EN
\LessonHead{Lesson IX}
\switchcolumn*

\LA
\noindent Cum sese testem áfferens ubíque non viderétur fide dignum habéri, addúcit Móysen, ut ínnuat hoc non casu factum esse, sed jam olim scripto fuísse prænuntiátum. Hoc illud dictum est: Os ejus non comminuétur. Rursúmque ex seípso Prophétæ fidem facit. Hæc dixi, inquit, ut discátis magnam esse affinitátem inter figúram et veritátem. Viden, quantam curam adhíbeat, ut credátur illud quod turpe et ignominiósum vidétur? Nam corpus a mílite contumélia áffici, longe pejus erat quam crucifígi. Attamen et hæc dixi, inquit, et cum magna diligéntia dixi, ut credátis. Nemo ítaque fidem neget, neque præ pudóre nostris nóceat. Nam quæ máxime contumeliósa vidéntur, hæc sunt bonórum nostrórum honestíssima.\par
\switchcolumn
\EN
\noindent Since the testimony he himself beareth might not everywhere be held worthy of belief, he citeth Moses, to intimate that this thing was not done by chance, but had already long ago been foretold in writing. By Moses it was said: A bone of him shall not be broken. And again he resteth his faith on the same Prophet: These things I have said, saith he, that ye may learn how great is the resemblance between the figure and the reality. Seest thou what great care he taketh, that what appeareth as disgraceful and ignominious may be believed. For that the body should be treated with contempt by the soldier, was far worse than its crucifixion. But nevertheless, saith he, I have both said these things, and have said them most emphatically, that ye may believe. Let no one, therefore, deny credence to these things, nor in shame tamper with our beliefs. For those things which seem to be the most dishonouring, are in fact our greatest pride.\par
\switchcolumn*

\LA
\TeDeum{Te Deum}
\switchcolumn
\EN
\TeDeum{Te Deum}
\switchcolumn*

\end{paracol}

\Office{Dominica III Post Pentecosten infra Octavam SSmi Cordis D.N.J.C}{Dominica III Post Pentecosten, infra Octavam SSmi Cordis D.N.J.C}{Semiduplex}
\addcontentsline{toc}{subsection}{Dominica III Post Pentecosten infra Octavam SSmi Cordis D.N.J.C \textit{— Dominica III Post Pentecosten, infra Octavam SSmi Cordis D.N.J.C}}
\begin{paracol}{2}
\LA
\LessonHead{Lectio I}
\switchcolumn
\EN
\LessonHead{Lesson I}
\switchcolumn*

\LA
\LessonTitle{De libro primo Regum}
\switchcolumn
\EN
\LessonTitle{Lesson from the first book of Samuel}
\switchcolumn*

\LA
\Cite{1 Reg 9:18-21}
\switchcolumn
\EN
\Cite{1 Sam 9:18-21}
\switchcolumn*

\LA
\noindent Accéssit autem Saul ad Samuélem in médio portæ, et ait: Indica, oro, mihi, ubi est domus vidéntis. Et respóndit Sámuel Sauli, dicens: Ego sum videns: ascénde ante me in excélsum, ut comedátis mecum hódie, et dimíttam te mane: et ómnia quæ sunt in corde tuo indicábo tibi. Et de ásinis quas nudiustértius perdidísti, ne sollícitus sis, quia invéntæ sunt. Et cujus erunt óptima quæque Israël? nonne tibi et omni dómui patris tui? Respóndens autem Saul, ait: Numquid non fílius Jémini ego sum de mínima tribu Israël, et cognátio mea novíssima inter omnes famílias de tribu Bénjamin? quare ergo locútus est mihi sermónem istum?\par
\switchcolumn
\EN
\noindent And Saul came to Samuel in the midst of the gate and said: Tell me, I pray thee, where is the house of the seer? And Samuel answered Saul, saying: I am the seer, go up before me to the high place, that you may eat with me today, and I will let thee go in the morning: and tell thee all that is in thy heart. And as for the asses, which were lost three days ago, be not solicitous, because they are found. And for whom shall be all the best things of Israel? Shall they not be for thee and for all thy father's house? And Saul answering, said: Am not I a son of Jemini of the least tribe of Israel, and my kindred the last among all the families of the tribe of Benjamin? Why then hast thou spoken this word to me?\par
\switchcolumn*

\LA
\begin{Resp}\Rsign Fériam eis pactum sempitérnum et non désinam eis benefácere et timórem meum dabo in corde eórum \Star{} Ut non recédant a me.\end{Resp}
\begin{Resp}\Vsign Et lætábor super eis cum bene eis fécero in toto Corde meo.\end{Resp}
\begin{Resp}\Rsign Ut non recédant a me.\end{Resp}
\switchcolumn
\EN
\begin{Resp}\Rsign I will make an everlasting Covenant with them, and I will not cease from doing them good, and I will put my fear in their hearts, \Star{} So that they shall not depart from me.\end{Resp}
\begin{Resp}\Vsign Yea, I will rejoice over them, to do them good with my whole Heart.\end{Resp}
\begin{Resp}\Rsign So that they shall not depart from me.\end{Resp}
\switchcolumn*

\LA
\LessonHead{Lectio II}
\switchcolumn
\EN
\LessonHead{Lesson II}
\switchcolumn*

\LA
\Cite{1 Reg 9:22-25}
\switchcolumn
\EN
\Cite{1 Sam 9:22-25}
\switchcolumn*

\LA
\noindent Assúmens ítaque Sámuel Saulem et púerum ejus, introdúxit eos in triclínium, et dedit eis locum in cápite eórum qui fúerant invitáti: erant enim quasi trigínta viri. Dixítque Sámuel coco: Da partem quam dedi tibi, et præcépi ut repóneres seórsum apud te. Levávit autem cocus armum, et pósuit ante Saul. Dixítque Sámuel: Ecce quod remánsit: pone ante te, et cómede, quia de indústria servátum est tibi quando pópulum vocávi. Et comédit Saul cum Samuéle in die illa. Et descendérunt de excélso in óppidum, et locútus est cum Saule in solário: stravítque Saul in solário, et dormívit.\par
\switchcolumn
\EN
\noindent Then Samuel taking Saul and his servant, brought them into the parlour, and gave them a place at the head of them that were invited. For there were about thirty men. And Samuel said to the cook: Bring the portion, which I gave thee, and commanded thee to set it apart by thee. And the cook took up the shoulder, and set it before Saul. And Samuel said: Behold what is left, set it before thee, and eat: because it was kept of purpose for thee, when I invited the people. And Saul ate with Samuel that day. And they went down from the high place into the town, and he spoke with Saul upon the top of the house: and he prepared a bed for Saul on the top of the house, and he slept.\par
\switchcolumn*

\LA
\begin{Resp}\Rsign Si inimícus meus maledixísset mihi, sustinuíssem útique \Star{} Tu vero homo unánimis qui simul mecum dulces capiébas cibos.\end{Resp}
\begin{Resp}\Vsign Et si is qui me óderat super me magna locútus fuísset, abscondíssem me fórsitan ab eo.\end{Resp}
\begin{Resp}\Rsign Tu vero homo unánimis qui simul mecum dulces capiébas cibos.\end{Resp}
\switchcolumn
\EN
\begin{Resp}\Rsign It was not an open enemy that has done me this dishonour, for then I could have borne it, \Star{} But it was even thou, mine own familiar friend, who did also eat of my Bread.\end{Resp}
\begin{Resp}\Vsign Neither was it mine adversary that did magnify himself against me, for then peradventure I would have hid myself from him.\end{Resp}
\begin{Resp}\Rsign But it was even thou, mine own familiar friend, who did also eat of my Bread.\end{Resp}
\switchcolumn*

\LA
\LessonHead{Lectio III}
\switchcolumn
\EN
\LessonHead{Lesson III}
\switchcolumn*

\LA
\Cite{1 Reg 9:26-27; 10:1}
\switchcolumn
\EN
\Cite{1 Sam 9:26-27;10:1}
\switchcolumn*

\LA
\noindent Cumque mane surrexíssent, et jam elucésceret, vocávit Sámuel Saulem in solário, dicens: Surge, et dimíttam te. Et surréxit Saul: egressíque sunt ambo, ipse vidélicet, et Sámuel. Cumque descénderent in extréma parte civitátis, Sámuel dixit ad Saul: Dic púero ut antecédat nos et tránseat: tu autem subsíste paulísper, ut índicem tibi verbum Dómini. Tulit autem Sámuel lentículam ólei, et effúdit super caput ejus: et deosculátus est eum, et ait: Ecce unxit te Dóminus super hereditátem suam in príncipem, et liberábis pópulum suum de mánibus inimicórum ejus qui in circúitu ejus sunt.\par
\switchcolumn
\EN
\noindent And when they were risen in the morning, and it began now to be light, Samuel called Saul on the top of the house, saying: Arise, that I may let thee go. And Saul arose: and they went out both of them, to wit, he and Samuel. And as they were going down in the end of the city, Samuel said to Saul: Speak to the servant to go before us, and pass on: but stand thou still a while, that I may tell thee the word of the Lord. And Samuel took a little vial of oil and poured it upon his head, and kissed him, and said: Behold, the Lord hath anointed thee to be prince over his inheritance, and thou shalt deliver his people out of the hands of their enemies, that are round about them. And this shall be a sign unto thee, that God hath anointed thee to be prince.\par
\switchcolumn*

\LA
\begin{Resp}\Rsign Cum essémus mórtui peccátis, convivificávit nos Deus in Christo \Star{} Propter nímiam caritátem suam qua diléxit nos.\end{Resp}
\begin{Resp}\Vsign Ut osténderet in sǽculis superveniéntibus abundántes divítias grátiæ suæ.\end{Resp}
\begin{Resp}\Rsign Propter nímiam caritátem suam qua diléxit nos.\end{Resp}
\begin{Resp}\Vsign Glória Patri, et Fílio, \Star{} et Spirítui Sancto.\end{Resp}
\begin{Resp}\Rsign Propter nímiam caritátem suam qua diléxit nos.\end{Resp}
\switchcolumn
\EN
\begin{Resp}\Rsign And we, being dead in our sins, hath God quickened together with Christ, \Star{} For his great love wherewith he hath loved us.\end{Resp}
\begin{Resp}\Vsign That in the ages to come he might shew the exceeding riches of his grace.\end{Resp}
\begin{Resp}\Rsign For his great love wherewith he hath loved us.\end{Resp}
\begin{Resp}\Vsign Glory be to the Father, and to the Son, \Star{} and to the Holy Ghost.\end{Resp}
\begin{Resp}\Rsign For his great love wherewith he hath loved us.\end{Resp}
\switchcolumn*

\LA
\LessonHead{Lectio IV}
\switchcolumn
\EN
\LessonHead{Lesson IV}
\switchcolumn*

\LA
\Source{Ex lítteris Encyclicis Pii Papæ undécimi}
\switchcolumn
\EN
\Source{From the Encyclical Letter of Pope Pius XI}
\switchcolumn*

\LA
\noindent At certe inter cétera illa, quæ próprie ad sacratíssimi Cordis cultum pértinent, pia éminet ac memoránda est consecrátio, qua, nos nostráque ómnia ætérnæ Núminis caritáti accépta referéntes, divíno Jesu Cordi devovémus. Verum, áliud accédat opórtet, honéstæ satisfactiónis, ínquimus, seu reparatiónis, quam dicunt, offícium sacratíssimo Cordi Jesu præstándum. Nam, si illud est in consecratióne primum ac præcípuum ut amóri Creatóris creatúræ amor rependátur, álterum sponte hinc séquitur, ut eídem increáto Amóri, si quando aut oblivióne negléctus, aut offénsa violátus sit, illátæ quoquo modo injúriæ compensári débeant: quod quidem débitum reparatiónem vulgáto nómine vocámus.\par
\switchcolumn
\EN
\noindent Among the various devotions paid to the Sacred Heart, the one foremost in importance and interest is assuredly the Act of Consecration, whereby we give to the divine Heart of Jesus both ourselves and all that is ours; in recognition of the truth that all we have cometh unto us out of the infinite charity of the eternal Deity. But it is expedient that any attempt of ours at self-consecration be accompanied with the purpose of making expiation (otherwise called reparation) to the most Sacred Heart of Jesus. In consecration the predominant intention may be said to be the purpose to repay (as it were) the love of the Creator by the love of us his creatures. But since Love Uncreate is passed over by human forgetfulness, and dishonoured by the sins of mankind, we should endeavour to repair such outrages; and the performance of this duty is ordinarily known as reparation.\par
\switchcolumn*

\LA
\begin{Resp}\Rsign Prope est Dóminus ómnibus invocántibus eum, \Star{} Omnibus invocántibus eum in veritáte.\end{Resp}
\begin{Resp}\Vsign Miserátor et miséricors Dóminus, pátiens et multum miséricors.\end{Resp}
\begin{Resp}\Rsign Omnibus invocántibus eum in veritáte.\end{Resp}
\switchcolumn
\EN
\begin{Resp}\Rsign The Lord is nigh unto all them that call upon him, \Star{} Yea, unto all such as call upon him faithfully.\end{Resp}
\begin{Resp}\Vsign The Lord is full of compassion and mercy, long-suffering and of great goodness.\end{Resp}
\begin{Resp}\Rsign Yea, unto all such as call upon him faithfully.\end{Resp}
\switchcolumn*

\LA
\LessonHead{Lectio V}
\switchcolumn
\EN
\LessonHead{Lesson V}
\switchcolumn*

\LA
\noindent Quodsi ad utrámque rem iísdem prorsus ratiónibus impéllimur, reparándi tamen expiandíque offício ob validiórem quemdam justítiæ et amóris títulum tenémur: justítiæ quidem, ut irrogáta Deo nostris flagítiis expiétur offénsa et violátus ordo pœniténtia redintegrétur; amóris vero, ut Christo patiénti ac « saturáto oppróbriis » compatiámur eíque nonníhil solácii pro tenuitáte nostra afferámus. Peccatóres enim cum simus omnes, multísque oneráti culpis, non eo solo cultu Deus noster nobis est honorándus, quo vel ejus summam Majestátem débitis obséquiis adorémus, vel ejus suprémum domínium precándo agnoscámus, vel ejus infinítam largitátem gratiárum actiónibus laudémus; sed prætérea Deo justo víndici satisfaciámus opórtet « pro innumerabílibus peccátis et offensiónibus et negligéntiis » nostris. Consecratióni ígitur, qua Deo devovémur et sancti Deo vocámur, ea sanctitáte ac firmitáte quæ, ut docet Angélicus, consecratiónis est própria, addénda est expiátio, qua pénitus peccáta exstinguántur, ne forte indignitátem nostram impudéntem revérberet summæ justítiæ sánctitas, munúsque nostrum pótius árceat invísum quam gratum suscípiat.\par
\switchcolumn
\EN
\noindent If we are, for the aforesaid reasons, to undertake both of those practices, we must recognize that we are impelled to the duty of reparation by the most powerful motives of justice and love: of justice, in order to expiate the injury done to God by our sins, and to re-establish through penance the divine order which was violated by them; of love, in order to suffer together with Christ, (who patiently endured all possible dishonour,) so that we may offer him some solace in return for his sufferings. For it is our duty to do more than honour God by the worship of adoration, whereby we adore his infinite Majesty, or by means of prayer, when we recognize his supreme dominion over us, or by acts of thanksgiving, when we praise his infinite generosity towards us. Because we are sinners, burdened with many offences, we must also make satisfaction to the offended justice of God, because of the numberless sins, offences and negligences we have committed. Wherefore, we must add to the act of consecration, by which we offer ourselves to God, and become thereby, as it were, sacred unto God by reason of the holiness which naturally floweth from an act of consecration, as the Angelic Doctor teacheth. We must add the act of reparation, by means of which all our faults are blotted out, lest perchance the sanctity of Infinite Justice spurn our arrogant unworthiness, and look upon the gift of ourselves as something to be rejected rather than accepted.\par
\switchcolumn*

\LA
\begin{Resp}\Rsign Confíteor tibi, Pater, Dómine cæli et terræ, quia abscondísti hæc a sapiéntibus et prudéntibus \Star{} Et revelásti ea párvulis.\end{Resp}
\begin{Resp}\Vsign Ita, Pater, quóniam sic fuit plácitum ante te.\end{Resp}
\begin{Resp}\Rsign Et revelásti ea párvulis.\end{Resp}
\switchcolumn
\EN
\begin{Resp}\Rsign I thank thee, O Father, Lord of heaven and earth, because thou hast hid these things from the wise and prudent; \Star{} Yea, thou hast revealed them unto babes.\end{Resp}
\begin{Resp}\Vsign Even so, Father, for so it seemed good in thy sight.\end{Resp}
\begin{Resp}\Rsign Yea, thou hast revealed them unto babes.\end{Resp}
\switchcolumn*

\LA
\LessonHead{Lectio VI}
\switchcolumn
\EN
\LessonHead{Lesson VI}
\switchcolumn*

\LA
\noindent Hoc autem expiatiónis offícium humáno géneri univérso incúmbit, quippe quod, ut christiána docémur fide, post Adæ miserándum casum, hereditária labe inféctum, concupiscéntiis obnóxium et misérrime depravátum, in perníciem detrudéndum fuísset sempitérnam. Id quidem supérbi hac nostra ætáte sapiéntes, véterem Pelágii errórem secúti, inficiántur, natívam quandam virtútem humánæ natúræ jactántes quæ suápte vi ad altióra usque progrediátur; sed falsa hæc humánæ supérbiæ comménta réjicit Apóstolus, illud nos ádmonens: « natúra erámus fílii iræ ». Et sane jam ab inítio commúnis illíus expiatiónis débitum quasi agnovére hómines et Deo sacrifíciis vel públicis placándo, naturáli quodam sensu ducti, óperam dare cœpérunt.\par
\switchcolumn
\EN
\noindent All men are under obligation to make reparation; for our souls are disfigured as the Christian faith teacheth by original sin as a result of the pitiable fall of Adam. We are also subject to passions, whereby we are corrupted in a truly sad state, and have thus made ourselves worthy of everlasting condemnation. It is true that the proud philosophers of this world deny the aforesaid verities, and in their place do raise up again the ancient heresy of Pelagius: which taught that in human nature there is a certain innate goodness wherewith, by our own powers, we are raised up to ever higher levels of perfection; but such false theories, born of human pride, have been condemned by the Apostle in his saying that all men are by nature the children of wrath. As a matter of fact, from the very beginning of the creation of the world, mankind recognized, in one way or another, the obligation of making reparation, impelled thereto, as by a natural instinct, in an endeavour to placate God by offering public sacrifices unto him.\par
\switchcolumn*

\LA
\begin{Resp}\Rsign Omnes gentes quascúmque fecísti, vénient \Star{} Et adorábunt coram te, Dómine.\end{Resp}
\begin{Resp}\Vsign Et glorificábunt nomen tuum, quóniam magnus es tu, et fáciens mirabília.\end{Resp}
\begin{Resp}\Rsign Et adorábunt coram te, Dómine.\end{Resp}
\begin{Resp}\Vsign Glória Patri, et Fílio, \Star{} et Spirítui Sancto.\end{Resp}
\begin{Resp}\Rsign Et adorábunt coram te, Dómine.\end{Resp}
\switchcolumn
\EN
\begin{Resp}\Rsign All nations whom thou hast made shall come, \Star{} And they shall worship thee, O Lord.\end{Resp}
\begin{Resp}\Vsign Yea, they shall glorify thy Name, for thou art great, and doest wondrous things.\end{Resp}
\begin{Resp}\Rsign And they shall worship thee, O Lord.\end{Resp}
\begin{Resp}\Vsign Glory be to the Father, and to the Son, \Star{} and to the Holy Ghost.\end{Resp}
\begin{Resp}\Rsign And they shall worship thee, O Lord.\end{Resp}
\switchcolumn*

\LA
\LessonHead{Lectio VII}
\switchcolumn
\EN
\LessonHead{Lesson VII}
\switchcolumn*

\LA
\LessonTitle{Léctio sancti Evangélii secúndum Lucam}
\switchcolumn
\EN
\LessonTitle{From the Holy Gospel according to Luke}
\switchcolumn*

\LA
\Cite{Luc 15:1-10}
\switchcolumn
\EN
\Cite{Luke 15:1-5}
\switchcolumn*

\LA
\noindent In illo témpore: Erant appropinquántes ad Jesum publicáni et peccatóres, ut audírent illum. Et réliqua.\par
\switchcolumn
\EN
\noindent At that time: the publicans and sinners drew near unto him to hear him. And so on.\par
\switchcolumn*

\LA
\LessonTitle{Homilía sancti Gregórii Papæ}
\switchcolumn
\EN
\LessonTitle{Homily by Pope St. Gregory (the Great)}
\switchcolumn*

\LA
\Source{Homilia 34 in Evangelium, n. 2-3, post initium}
\switchcolumn
\EN
\Source{34th on the Gospels}
\switchcolumn*

\LA
\noindent Audístis in lectióne evangélica, fratres mei, quia peccatóres et publicáni accessérunt ad Redemptórem nostrum; et non solum ad colloquéndum, sed étiam ad convescéndum recépti sunt. Quod vidéntes pharisǽi dedignáti sunt. Ex qua re collígite quia vera justítia compassiónem habet, falsa justítia dedignatiónem. Quamvis et justi sóleant recte peccatóribus dedignári: sed áliud est quod ágitur typho supérbiæ, áliud quod zelo disciplínæ.\par
\switchcolumn
\EN
\noindent Ye have heard, my brethren, from the Gospel which hath but now been read, how that the publicans and sinners drew near unto our Redeemer, and how that He received them, not only to converse, but also to eat with Him. And when the Pharisees and Scribes saw it, they murmured. From this learn ye, that true righteousness is merciful, and false righteousness is contemptuous, albeit that the righteous also oft-times feel moved with just indignation at sinners. But it is one thing to feel thus indignant through pride, and another to feel so through love of law.\par
\switchcolumn*

\LA
\begin{Resp}\Rsign Ego si exaltátus fúero a terra \Star{} Omnia traham ad meípsum.\end{Resp}
\begin{Resp}\Vsign Hoc autem dicébat signíficans qua morte esset moritúrus.\end{Resp}
\begin{Resp}\Rsign Omnia traham ad meípsum.\end{Resp}
\switchcolumn
\EN
\begin{Resp}\Rsign If I be lifted up, \Star{} I will draw all men unto me.\end{Resp}
\begin{Resp}\Vsign This he said, signifying what death he should die.\end{Resp}
\begin{Resp}\Rsign I will draw all men unto me.\end{Resp}
\switchcolumn*

\LA
\LessonHead{Lectio VIII}
\switchcolumn
\EN
\LessonHead{Lesson VIII}
\switchcolumn*

\LA
\noindent Dedignántur étenim, sed non dedignántes: despérant, sed non desperántes: persecutiónem cómmovent, sed amántes: quia etsi foris increpatiónes per disciplínam exággerant, intus tamen dulcédinem per caritátem servant. Præpónunt sibi in ánimo ipsos plerúmque quos córrigunt: melióres exístimant eos quoque quos júdicant. Quod vidélicet agéntes, et per disciplínam súbditos, et per humilitátem custódiunt semetípsos.\par
\switchcolumn
\EN
\noindent The righteous indeed look down upon sinners, and yet, as not despising them; they abandon them, and yet, as not without hope; they fight against them, and yet, as loving them all the while; for if they be behoven to chasten them grievously as touching the outer man, yet is it through charity which offereth sweetness to their inner man. In their hearts they prefer before themselves them whom they are correcting; they hold as better than themselves them whom they judge. And thus doing, they watch by carefulness over them, which are committed unto their charge, and, by lowly-mindedness, over themselves.\par
\switchcolumn*

\LA
\begin{Resp}\Rsign Simus ergo imitatóres Dei \Star{} Et ambulémus in dilectióne.\end{Resp}
\begin{Resp}\Vsign Sicut et Christus diléxit nos et trádidit semetípsum pro nobis.\end{Resp}
\begin{Resp}\Rsign Et ambulémus in dilectióne.\end{Resp}
\begin{Resp}\Vsign Glória Patri, et Fílio, \Star{} et Spirítui Sancto.\end{Resp}
\begin{Resp}\Rsign Et ambulémus in dilectióne.\end{Resp}
\switchcolumn
\EN
\begin{Resp}\Rsign Be ye therefore followers of God; \Star{} And walk ye therefore in love.\end{Resp}
\begin{Resp}\Vsign For Christ also hath loved us and hath given himself for us.\end{Resp}
\begin{Resp}\Rsign And walk ye therefore in love.\end{Resp}
\begin{Resp}\Vsign Glory be to the Father, and to the Son, \Star{} and to the Holy Ghost.\end{Resp}
\begin{Resp}\Rsign And walk ye therefore in love.\end{Resp}
\switchcolumn*

\LA
\LessonHead{Lectio IX}
\switchcolumn
\EN
\LessonHead{Lesson IX}
\switchcolumn*

\LA
\noindent At contra, hi qui de falsa justítia superbíre solent, céteros quosque despíciunt, nulla infirmántibus misericórdia condescéndunt: quo se peccatóres esse non credunt, eo detérius peccatóres fiunt. De quorum profécto número pharisǽi exstíterant, qui diiudicántes Dóminum quod peccatóres suscíperet, arénti corde ipsum fontem misericórdiæ reprehendébant. Sed quia ægri erant, ita ut ægros se esse nescírent; quátenus quod erant agnóscerent, cæléstis eos médicus blandis foméntis curat, benígnum paradígma óbjicit, et in eórum corde vúlneris tumórem premit.\par
\switchcolumn
\EN
\noindent On the other hand, they whose exaltation cometh of a false righteousness, look down upon their neighbour, but are softened by no mercy toward his misery, and are all the more sinful, because they perceive not that they themselves are sinners. Of such were those Pharisees who judged the Lord because He received sinners, and, in the dryness of their own heart, rebuked the very Fountain of mercy. They were sick of so desperate a sickness that they knew not of themselves that they were sick; but, that they might know that they were so, the Heavenly Physician applied to them His tender ointments, and, by means of a gracious parable, lanced the boil of their pride of heart.\par
\switchcolumn*

\LA
\TeDeum{Te Deum}
\switchcolumn
\EN
\TeDeum{Te Deum}
\switchcolumn*

\end{paracol}

\Office{Die IV infra Octavam SSmi Cordis Jesu}{Fourth Day within the Octave of the Sacred Heart}{Semiduplex}
\addcontentsline{toc}{subsection}{Die IV infra Octavam SSmi Cordis Jesu \textit{— Fourth Day within the Octave of the Sacred Heart}}
\begin{paracol}{2}
\LA
\LessonHead{Lectio I}
\switchcolumn
\EN
\LessonHead{Lesson I}
\switchcolumn*

\LA
\LessonTitle{De libro primo Regum}
\switchcolumn
\EN
\LessonTitle{Lesson from the first book of Samuel}
\switchcolumn*

\LA
\Cite{1 Reg 10:17-19}
\switchcolumn
\EN
\Cite{1 Sam 10:17-19}
\switchcolumn*

\LA
\noindent Et convocávit Sámuel pópulum ad Dóminum in Maspha: et ait ad fílios Israël: Hæc dicit Dóminus Deus Israël: Ego edúxi Israël de Ægýpto, et érui vos de manu Ægyptiórum et de manu ómnium regum, qui affligébant vos. Vos autem hódie projecístis Deum vestrum, qui solus salvávit vos de univérsis malis et tribulatiónibus vestris, et dixístis: Nequáquam, sed regem constítue super nos. Nunc ergo state coram Dómino per tribus vestras et per famílias.\par
\switchcolumn
\EN
\noindent And Samuel called together the people to the Lord in Maspha: And he said to the children of Israel: Thus saith the Lord the God of Israel: I brought up Israel out of Egypt, and delivered you from the hand of the Egyptians, and from the hand of all the kings who afflicted you. But you this day have rejected your God, who only hath saved you out of all your evils and your tribulations: and you have said: Nay: but set a king over us. Now therefore stand before the Lord by your tribes, and by your families\par
\switchcolumn*

\LA
\begin{Resp}\Rsign Fériam eis pactum sempitérnum et non désinam eis benefácere et timórem meum dabo in corde eórum \Star{} Ut non recédant a me.\end{Resp}
\begin{Resp}\Vsign Et lætábor super eis cum bene eis fécero in toto Corde meo.\end{Resp}
\begin{Resp}\Rsign Ut non recédant a me.\end{Resp}
\switchcolumn
\EN
\begin{Resp}\Rsign I will make an everlasting Covenant with them, and I will not cease from doing them good, and I will put my fear in their hearts, \Star{} So that they shall not depart from me.\end{Resp}
\begin{Resp}\Vsign Yea, I will rejoice over them, to do them good with my whole Heart.\end{Resp}
\begin{Resp}\Rsign So that they shall not depart from me.\end{Resp}
\switchcolumn*

\LA
\LessonHead{Lectio II}
\switchcolumn
\EN
\LessonHead{Lesson II}
\switchcolumn*

\LA
\Cite{1 Reg 10:20-24}
\switchcolumn
\EN
\Cite{1 Sam 10:20-24}
\switchcolumn*

\LA
\noindent Et applícuit Sámuel omnes tribus Israël, et cécidit sors tribus Bénjamin; et applícuit tribum Bénjamin et cognatiónes ejus, et cécidit cognátio Metri, et pervénit usque ad Saul fílium Cis. Quæsiérunt ergo eum, et non est invéntus. Et consuluérunt post hæc Dóminum utrúmnam ventúrus esset illuc. Respondítque Dóminus: Ecce abscónditus est domi. Cucurrérunt ítaque et tulérunt eum inde; stetítque in médio pópuli, et áltior fuit univérso pópulo ab húmero et sursum. Et ait Sámuel ad omnem pópulum: Certe vidétis quem elégit Dóminus, quóniam non sit símilis illi in omni pópulo. Et clamávit omnis pópulus et ait: Vivat rex.\par
\switchcolumn
\EN
\noindent And Samuel brought to him all the tribes of Israel, and the lot fell on the tribe of Benjamin. And he brought the tribe of Benjamin and the kindreds thereof, and the lot fell upon the kindred of Metri, and it came to Saul the son of Cis. They sought him therefore and he was not found. And after this they consulted the Lord whether he would come thither. And the Lord answered: Behold he is hidden at home. And they ran and fetched him thence: and he stood in the midst of the people, and he was higher than any of the people from the shoulders and upward. And Samuel said to all the people: Surely you see him whom the Lord hath chosen, that there is none like him among all the people. And all the people cried and said: God save the king.\par
\switchcolumn*

\LA
\begin{Resp}\Rsign Si inimícus meus maledixísset mihi, sustinuíssem útique \Star{} Tu vero homo unánimis qui simul mecum dulces capiébas cibos.\end{Resp}
\begin{Resp}\Vsign Et si is qui me óderat super me magna locútus fuísset, abscondíssem me fórsitan ab eo.\end{Resp}
\begin{Resp}\Rsign Tu vero homo unánimis qui simul mecum dulces capiébas cibos.\end{Resp}
\switchcolumn
\EN
\begin{Resp}\Rsign It was not an open enemy that has done me this dishonour, for then I could have borne it, \Star{} But it was even thou, mine own familiar friend, who did also eat of my Bread.\end{Resp}
\begin{Resp}\Vsign Neither was it mine adversary that did magnify himself against me, for then peradventure I would have hid myself from him.\end{Resp}
\begin{Resp}\Rsign But it was even thou, mine own familiar friend, who did also eat of my Bread.\end{Resp}
\switchcolumn*

\LA
\LessonHead{Lectio III}
\switchcolumn
\EN
\LessonHead{Lesson III}
\switchcolumn*

\LA
\Cite{1 Reg 10:25-27}
\switchcolumn
\EN
\Cite{1 Sam 10:25-27}
\switchcolumn*

\LA
\noindent Locútus est autem Sámuel ad pópulum legem regni, et scripsit in libro, et repósuit coram Dómino; et dimísit Sámuel omnem pópulum, síngulos in domum suam. Sed et Saul ábiit in domum suam in Gábaa; et ábiit cum eo pars exércitus, quorum tetígerat Deus corda. Fílii vero Bélial dixérunt: Num salváre nos póterit iste? Et despexérunt eum, et non attulérunt ei múnera; ille vero dissimulábat se audíre.\par
\switchcolumn
\EN
\noindent And Samuel told the people the law of the kingdom, and wrote it in a book, and laid it up before the Lord: and Samuel sent away all the people, every one to his own house. Saul also departed to his own house in Gabaa: and there went with him a part of the army, whose hearts God had touched. But the children of Belial said: Shall this fellow be able to save us? And they despised him, and brought him no presents, but he dissembled as though he heard not.\par
\switchcolumn*

\LA
\begin{Resp}\Rsign Cum essémus mórtui peccátis, convivificávit nos Deus in Christo \Star{} Propter nímiam caritátem suam qua diléxit nos.\end{Resp}
\begin{Resp}\Vsign Ut osténderet in sǽculis superveniéntibus abundántes divítias grátiæ suæ.\end{Resp}
\begin{Resp}\Rsign Propter nímiam caritátem suam qua diléxit nos.\end{Resp}
\begin{Resp}\Vsign Glória Patri, et Fílio, \Star{} et Spirítui Sancto.\end{Resp}
\begin{Resp}\Rsign Propter nímiam caritátem suam qua diléxit nos.\end{Resp}
\switchcolumn
\EN
\begin{Resp}\Rsign And we, being dead in our sins, hath God quickened together with Christ, \Star{} For his great love wherewith he hath loved us.\end{Resp}
\begin{Resp}\Vsign That in the ages to come he might shew the exceeding riches of his grace.\end{Resp}
\begin{Resp}\Rsign For his great love wherewith he hath loved us.\end{Resp}
\begin{Resp}\Vsign Glory be to the Father, and to the Son, \Star{} and to the Holy Ghost.\end{Resp}
\begin{Resp}\Rsign For his great love wherewith he hath loved us.\end{Resp}
\switchcolumn*

\LA
\LessonHead{Lectio IV}
\switchcolumn
\EN
\LessonHead{Lesson IV}
\switchcolumn*

\LA
\Source{Ex lítteris Encyclicis Pii Papæ undécimi}
\switchcolumn
\EN
\Source{From the Encyclical Letter of Pope Pius XI}
\switchcolumn*

\LA
\noindent At nulla creáta vis hóminum sceléribus expiándis erat satis, nisi humánam natúram Dei Fílius reparándam assumpsísset. Quod quidem ipse hóminum Salvátor sacri Psaltis ore nuntiávit: Hóstiam et oblatiónem noluísti, corpus autem aptásti mihi: holocautómata pro peccáto non tibi placuérunt: tunc dixi: Ecce vénio. Et reápse vere languóres nostros ipse tulit et dolóres nostros ipse portávit; vulnerátus est propter iniquitátes nostras et peccáta nostra ipse pértulit in córpore suo super lignum; delens quod advérsus nos erat chirógraphum decréti, quod erat contrárium nobis, et ipsum tulit de médio affígens illud cruci, ut peccátis mórtui justítiæ vivámus. Quamquam vero copiósa Christi redémptio abúnde nobis ómnia delícta donávit; ob miram tamen illam divínæ Sapiéntiæ dispensatiónem, qua in carne nostra adimplénda sunt quæ desunt passiónum Christi pro córpore ejus quod est Ecclésia, étiam láudibus et satisfactiónibus, quas Christus in nómine peccatórum Deo persólvit, nostras quoque laudes et satisfactiónes adjícere póssumus, immo étiam debémus.\par
\switchcolumn
\EN
\noindent No effort on our part could have availed for the expiation of human sin, if the Son of God had not assumed human nature in order to redeem men from their sins. This truth the Saviour of mankind hath made known in the words of the Psalmist: Sacrifice and offering thou wouldest not, but a body hast thou prepared me; in burnt offerings and sacrifices for sin thou hast had no pleasure; then said I, Lo, I come. Surely he hath borne our griefs, and carried our sorrows: he was wounded for our transgressions, he was bruised for our iniquities: who his own self bare our sins in his body on the Tree; and blotting out the handwriting of ordinances that was against us, which same was contrary to us, he took it out of the way, nailing it to his Cross: that we, being dead unto sins, should live unto righteousness. But although this plenteous redemption of Christ was more than sufficient to satisfy for all our offenses, nevertheless, (due to the wondrous dispensation of the divine Wisdom, whereby in our own flesh we may fill up that which is behind of the afflictions of Christ for his body's sake, which is the Church,) we can, and in fact, we should, add our own acts of praise and satisfaction to those acts of praise and satisfaction which Christ, in the name of sinners, presented to God.\par
\switchcolumn*

\LA
\begin{Resp}\Rsign Prope est Dóminus ómnibus invocántibus eum, \Star{} Omnibus invocántibus eum in veritáte.\end{Resp}
\begin{Resp}\Vsign Miserátor et miséricors Dóminus, pátiens et multum miséricors.\end{Resp}
\begin{Resp}\Rsign Omnibus invocántibus eum in veritáte.\end{Resp}
\switchcolumn
\EN
\begin{Resp}\Rsign The Lord is nigh unto all them that call upon him, \Star{} Yea, unto all such as call upon him faithfully.\end{Resp}
\begin{Resp}\Vsign The Lord is full of compassion and mercy, long-suffering and of great goodness.\end{Resp}
\begin{Resp}\Rsign Yea, unto all such as call upon him faithfully.\end{Resp}
\switchcolumn*

\LA
\LessonHead{Lectio V}
\switchcolumn
\EN
\LessonHead{Lesson V}
\switchcolumn*

\LA
\noindent At semper meminérimus opórtet, totam expiatiónis virtútem ab uno Christi cruénto sacrifício pendére, quod sine témporis intermissióne in nostris altáribus incruénto modo renovátur; síquidem una eadémque est Hóstia, idem nunc ófferens sacerdótum ministério, qui seípsum tunc in cruce óbtulit, sola offeréndi ratióne divérsa; quamóbrem cum hoc augustíssimo Eucharístico sacrifício et ministrórum et aliórum fidélium immolátio conjúngi debet, ut ipsi quoque hóstias vivéntes, sanctas, Deo placéntes sese exhibeant. Quin immo, sanctus Cypriánus affirmáre non dúbitat sacrifícium domínicum legítima sanctificatióne non celebrári, nisi oblátio et sacrifícium nostrum respónderit passióni. Quaprópter nos monet Apóstolus, ut mortificatiónem Jesu in córpore nostro circumferéntes, atque cum Christo consepúlti et complantáti similitúdini mortis ejus, non modo carnem nostram crucifigámus cum vítiis et concupiscéntiis, fugiéntes ejus quæ in mundo est concupiscéntiæ corruptiónem, sed et vita Jesu manifestétur in corpóribus nostris et, ætérni ejus sacerdótii partícipes effécti, offerámus dona et sacrifícia pro peccátis.\par
\switchcolumn
\EN
\noindent However, we must always remember this, to wit: The expiatory value of our acts is dependent upon the Bloody Sacrifice of Christ; which same is presented bloodlessly on our altars, without pretermission. Note that in both the Unbloody and the Bloody Sacrifice, the Victim is one and the same. He that offered himself on the Cross is the very same that doth offer himself by means of our ministerial priesthood; the only difference being the manner in which the Sacrifice is made; for which reason there must needs be conjoined to the august Sacrifice of the Holy Eucharist, and act of immolation from the priests and the faithful, whereby they offer up themselves, also, to be a reasonable, holy, and living sacrifice, acceptable unto God. With this in mind Saint Cyprian dared to affirm that the Sacrifice of the Lord is not complete so far as our sanctification is concerned until our personal offerings and sacrifices are brought into union with his passion. To which end we have been given apostolic admonition: That we should always bear about in the body of the dying of the Lord Jesus; and thereby, buried with him by Baptism unto death, not only can we crucify the flesh with its affections and lusts, thus escaping the corruption that is in the world, through lust, but also thereby the life of Jesus can be made manifest in our bodies; and so, having been made partakers of Christ in his holy and eternal priesthood, we should offer both gifts and sacrifices for sins.\par
\switchcolumn*

\LA
\begin{Resp}\Rsign Confíteor tibi, Pater, Dómine cæli et terræ, quia abscondísti hæc a sapiéntibus et prudéntibus \Star{} Et revelásti ea párvulis.\end{Resp}
\begin{Resp}\Vsign Ita, Pater, quóniam sic fuit plácitum ante te.\end{Resp}
\begin{Resp}\Rsign Et revelásti ea párvulis.\end{Resp}
\switchcolumn
\EN
\begin{Resp}\Rsign I thank thee, O Father, Lord of heaven and earth, because thou hast hid these things from the wise and prudent; \Star{} Yea, thou hast revealed them unto babes.\end{Resp}
\begin{Resp}\Vsign Even so, Father, for so it seemed good in thy sight.\end{Resp}
\begin{Resp}\Rsign Yea, thou hast revealed them unto babes.\end{Resp}
\switchcolumn*

\LA
\LessonHead{Lectio VI}
\switchcolumn
\EN
\LessonHead{Lesson VI}
\switchcolumn*

\LA
\noindent Neque enim arcáni hujus sacerdótii et satisfaciéndi sacrificandíque múneris participatióne ii soli fruúntur, quibus Póntifex noster Christus Jesus adminístris útitur ad oblatiónem mundam divíno Nómini ab ortu solis usque ad occásum omni loco offeréndam; sed étiam christianórum gens univérsa, ab Apostolórum Príncipe genus eléctum, regále sacerdótium jure appelláta, debet cum pro se, tum pro toto humáno génere offérre pro peccátis, haud áliter propémodum quam sacérdos omnis ac póntifex ex homínibus assúmptus, pro homínibus constitúitur in iis quæ sunt ad Deum.\par
\switchcolumn
\EN
\noindent And note who they are that be partakers in the mysteries of such a priesthood, and in the duty of offering sacrifices and satisfaction to God; not only they who have been ordained as ministers of such sacrifices, to wit, to offer unto the divine Name in every place a pure offering, from the rising of the sun even unto the going down of the same, but also all others who are Christians; which same by the Prince of the Apostles are called, and rightly so, a chosen generation, a royal priesthood, who are to offer sacrifices for sin, not only for themselves but for all mankind, and this in much the same way as every priest and high priest taken from among men is ordained for men in the things that appertain to God.\par
\switchcolumn*

\LA
\begin{Resp}\Rsign Omnes gentes quascúmque fecísti, vénient \Star{} Et adorábunt coram te, Dómine.\end{Resp}
\begin{Resp}\Vsign Et glorificábunt nomen tuum, quóniam magnus es tu, et fáciens mirabília.\end{Resp}
\begin{Resp}\Rsign Et adorábunt coram te, Dómine.\end{Resp}
\begin{Resp}\Vsign Glória Patri, et Fílio, \Star{} et Spirítui Sancto.\end{Resp}
\begin{Resp}\Rsign Et adorábunt coram te, Dómine.\end{Resp}
\switchcolumn
\EN
\begin{Resp}\Rsign All nations whom thou hast made shall come, \Star{} And they shall worship thee, O Lord.\end{Resp}
\begin{Resp}\Vsign Yea, they shall glorify thy Name, for thou art great, and doest wondrous things.\end{Resp}
\begin{Resp}\Rsign And they shall worship thee, O Lord.\end{Resp}
\begin{Resp}\Vsign Glory be to the Father, and to the Son, \Star{} and to the Holy Ghost.\end{Resp}
\begin{Resp}\Rsign And they shall worship thee, O Lord.\end{Resp}
\switchcolumn*

\LA
\LessonHead{Lectio VII}
\switchcolumn
\EN
\LessonHead{Lesson VII}
\switchcolumn*

\LA
\LessonTitle{Léctio sancti Evangélii secúndum Joánnem}
\switchcolumn
\EN
\LessonTitle{From the Holy Gospel according to John}
\switchcolumn*

\LA
\Cite{Joannes 19:31-37}
\switchcolumn
\EN
\Cite{John 19:31-37}
\switchcolumn*

\LA
\noindent In illo témpore: Judǽi, quóniam parascéve erat, ut non remanérent in cruce córpora sábbato (erat enim magnus dies ille sábbati) rogavérunt Pilátum, ut frangeréntur eórum crura et tolleréntur. Et réliqua.\par
\switchcolumn
\EN
\noindent At that time: The Jews, because it was the Preparation, that the bodies should not remain upon the cross on the Sabbath Day, for that Sabbath Day was an high day, besought Pilate that their legs might broken, and that they might be taken away. And so on, and that which followeth.\par
\switchcolumn*

\LA
\LessonTitle{Homilía sancti Lauréntii Justiniáni Epíscopi}
\switchcolumn
\EN

\switchcolumn*

\LA
\Source{De triumpháli Christi agóne, cap. 21}
\switchcolumn
\EN
\Source{A Homily by St. Lawrence Justinian the Bishop}
\switchcolumn*

\LA
\noindent Ad Jesum ígitur cum veníssent, illúmque jam mórtuum conspexíssent, ipsíus mínime crura fregérunt, verum astántium unus, vibráta láncea, latus ejus apéruit, et contínuo exívit sanguis et aqua. Grande prorsus et inaudítum prodígium, ut de exanimáto córpore sanguis exíret et aqua. Verúmtamen, máximum, in re hac gesta, nobis vóluit Dei sapiéntia commendáre sacraméntum, sui vidélicet, et Ecclésiæ unitátem. Hujus enim spirituális cópulæ figúra præcéssit, quando de látere Adæ dormiéntis una ex costis ipsíus subtrácta narrátur, atque de eádem Eva cunctórum parens formáta est, quæ typum gerébat Ecclésiæ. Significábat tunc Spíritus Sanctus, verum et spirituálem Adam esse futúrum, qui Paracléti virtúte plasmátus, dum dormíret in cruce, de ipsíus látere, profluéntibus aqua et sánguine, speciósa sponsa sine ruga et mácula, formarétur Ecclésia.\par
\switchcolumn
\EN
\noindent But when they came to Jesus, and saw that he was dead already, they brake not his legs, but one of the soldiers with a spear pierced his side, and forthwith came there out blood and water. Truly this is a great and unheard-of wonder, that from a lifeless body should gush out blood and water! And thereby we are constrained to believe that God willed to set forth the great mystery, namely, the oneness between Christ and the Church. Consider how aforetime a figure of this spiritual union was given us in the Scriptures, where we are told that, from the side of Adam as he slept, one of his ribs was taken, wherefrom was formed Eve, the mother of us all, which same is understood to be a type of the Church. For thereby the Holy Ghost signified that there was to be a true and spiritual Adam, fashioned by the same Paraclete, from whose side, while he slept on the Cross, the Church would be formed, a comely spouse radiant with youth, that is, without spot or wrinkle, or any such thing.\par
\switchcolumn*

\LA
\begin{Resp}\Rsign Ego si exaltátus fúero a terra \Star{} Omnia traham ad meípsum.\end{Resp}
\begin{Resp}\Vsign Hoc autem dicébat signíficans qua morte esset moritúrus.\end{Resp}
\begin{Resp}\Rsign Omnia traham ad meípsum.\end{Resp}
\switchcolumn
\EN
\begin{Resp}\Rsign If I be lifted up, \Star{} I will draw all men unto me.\end{Resp}
\begin{Resp}\Vsign This he said, signifying what death he should die.\end{Resp}
\begin{Resp}\Rsign I will draw all men unto me.\end{Resp}
\switchcolumn*

\LA
\LessonHead{Lectio VIII}
\switchcolumn
\EN
\LessonHead{Lesson VIII}
\switchcolumn*

\LA
\noindent Hæc quippe sunt sacraménta ecclesiástica per quæ totum corpus Ecclésiæ ablúitur atque sanctificátur. In lavácro útique aquæ regeneratiónis, quod Christi morte sacrátur, ab origináli contagióne mundátur. In Redemptóris vero sánguine nedum ab omni culpa purgátur, verum étiam regni cæléstis illi aperítur intróitus. Ambo hæc in unum efféctum convéniunt nec unum sine áltero valet ad salútem prodésse: non enim absque baptísmi sacraménto et peccatórum remissióne, potest quis futúræ beatitúdinis hereditátem percípere. Hoc ubíque in orbe terrárum sancta mater confitétur Ecclésia, et multiplícibus Scripturárum divinárum testimóniis roborátur. Insuper et qui vidit de Christo aquam emanáre et sánguinem, testimónium perhíbuit, et verum est testimónium ejus. Nempe iste est Joánnes Apóstolus et Evangelísta, qui præcípuo amóre diléctus est a Dómino.\par
\switchcolumn
\EN
\noindent For the blood and water are the Sacraments of the Church, whereby the whole body of the Church is washed and sanctified. Certain it is that in the laver of regenerating water, (which same was consecrated by the death of Christ,) the Church is cleansed from original sin. And not only is purification given thereby, but entrance to the heavenly kingdom is also made. Both these things are done at one and the same time, nor doth one avail for salvation without the other; for no one can take unto himself the inheritance of the blessedness which is to come without the Sacrament of Baptism and the remission of sins. This truth is confessed all the world over by holy Mother Church, and confirmed by manifold testimonies in the Word of God. Moreover, he too bare witness, who saw the water and blood flow forth from the side of Christ, and his testimony is true. Now this is John the Apostle and Evangelist who was loved with an exceeding great love by the Lord.\par
\switchcolumn*

\LA
\begin{Resp}\Rsign Simus ergo imitatóres Dei \Star{} Et ambulémus in dilectióne.\end{Resp}
\begin{Resp}\Vsign Sicut et Christus diléxit nos et trádidit semetípsum pro nobis.\end{Resp}
\begin{Resp}\Rsign Et ambulémus in dilectióne.\end{Resp}
\begin{Resp}\Vsign Glória Patri, et Fílio, \Star{} et Spirítui Sancto.\end{Resp}
\begin{Resp}\Rsign Et ambulémus in dilectióne.\end{Resp}
\switchcolumn
\EN
\begin{Resp}\Rsign Be ye therefore followers of God; \Star{} And walk ye therefore in love.\end{Resp}
\begin{Resp}\Vsign For Christ also hath loved us and hath given himself for us.\end{Resp}
\begin{Resp}\Rsign And walk ye therefore in love.\end{Resp}
\begin{Resp}\Vsign Glory be to the Father, and to the Son, \Star{} and to the Holy Ghost.\end{Resp}
\begin{Resp}\Rsign And walk ye therefore in love.\end{Resp}
\switchcolumn*

\LA
\LessonHead{Lectio IX}
\switchcolumn
\EN
\LessonHead{Lesson IX}
\switchcolumn*

\LA
\noindent Profécto hæc facta sunt ut Scriptúra implerétur quæ ait: Os non comminuétis ex eo. A Dómino quippe Móysi fúerat imperátum ne in sacrifícii phase, in celebratióne immoláti agni, ullum comminuerétur os. In Dómino autem Jesu innocentíssimo agno, figúræ véritas est compléta. Nequáquam crura illíus sunt fracta, quemádmodum duórum nequam cum eo pendéntium; sed tantum illíus latus apértum, quátenus ália Scriptúra perficerétur quæ dicit: Vidébunt in quem transfixérunt. Porro vúlnerum cicatríces in suo vóluit Dóminus córpore retinére, ut, sicut eléctis incentívum devotiónis, ita et réprobis irrefragábile fíeret testimónium damnatiónis. Ita ómnia consummáta in Christo, longe ante Prophetárum oráculis promulgáta fuére, quátenus cathólica fides tam pro se quam advérsus hæreticórum munirétur erróres.\par
\switchcolumn
\EN
\noindent Verily, all these things were done that the Scriptures might be fulfilled, which say: A bone of him shall not be broken: which same was in obedience to the Lord's commandment to Moses, that in the sacrifice of the Passover, (wherein was celebrated the offering up of the lamb,) no bone should be broken. Wherefrom, when this ancient prototype was fulfilled in the most innocent Lamb, (the Lord Jesus himself,) obedience to the ancient command was also fulfilled. For in no wise were his legs broken, after the manner of what was done to the two malefactors hanging with him; but his side alone was opened, that another Scripture might be fulfilled, which saith: They shall look upon him whom they have pierced. Moreover, the Lord willed to retain the marks of the Wounds in his body, thereby to incite the devotion of the elect, and also be an unanswerable testimony of damnation to the reprobate. And so all things were consummated in Christ, which had been long before declared by the oracles of the Prophets, that the Catholic faith might be strengthened both in itself and against the errors of the hereticks.\par
\switchcolumn*

\LA
\TeDeum{Te Deum}
\switchcolumn
\EN
\TeDeum{Te Deum}
\switchcolumn*

\end{paracol}

\Office{Die V infra Octavam SSmi Cordis Jesu}{Fifth Day within the Octave of the Sacred Heart}{Semiduplex}
\addcontentsline{toc}{subsection}{Die V infra Octavam SSmi Cordis Jesu \textit{— Fifth Day within the Octave of the Sacred Heart}}
\begin{paracol}{2}
\LA
\LessonHead{Lectio I}
\switchcolumn
\EN
\LessonHead{Lesson I}
\switchcolumn*

\LA
\LessonTitle{De libro primo Regum}
\switchcolumn
\EN
\LessonTitle{Lesson from the first book of Samuel}
\switchcolumn*

\LA
\Cite{1 Reg 12:1-5}
\switchcolumn
\EN
\Cite{1 Sam 12:1-5}
\switchcolumn*

\LA
\noindent Dixit autem Sámuel ad univérsum Israël: Ecce audívi vocem vestram juxta ómnia quæ locúti estis ad me, et constítui super vos regem. Et nunc rex gráditur ante vos: ego autem sénui, et incánui: porro fílii mei vobíscum sunt: ítaque conversátus coram vobis ab adolescéntia mea usque ad hanc diem, ecce præsto sum. Loquímini de me coram Dómino, et coram christo ejus, utrum bovem cujúsquam túlerim, aut ásinum: si quémpiam calumniátus sum, si oppréssi áliquem, si de manu cujúsquam munus accépi: et contémnam illud hódie, restituámque vobis. Et dixérunt: Non es calumniátus nos, neque oppressísti, neque tulísti de manu alicújus quíppiam. Dixítque ad eos: Testis est Dóminus advérsum vos, et testis Christus ejus in die hac, quia non invenéritis in manu mea quíppiam. Et dixérunt: Testis.\par
\switchcolumn
\EN
\noindent And Samuel said to all Israel: Behold I have hearkened to your voice in all that you said to me, and have made a king over you. And now the king goeth before you: but I am old and greyheaded: and my sons are with you: having then conversed with you from my youth unto this day, behold here I am. Speak of me before the Lord, and before his anointed, whether I have taken any man's ox, or ass: If I have wronged any man, if I have oppressed any man, if I have taken a bribe at any man's hand: and I will despise it this day, and will restore it to you. And they said: Thou hast not wronged us, nor oppressed us, nor taken ought at any man's hand. And he said to them: The Lord is witness against you, and his anointed is witness this day, that you have not found any thing in my hand. And they said: He is witness.\par
\switchcolumn*

\LA
\begin{Resp}\Rsign Fériam eis pactum sempitérnum et non désinam eis benefácere et timórem meum dabo in corde eórum \Star{} Ut non recédant a me.\end{Resp}
\begin{Resp}\Vsign Et lætábor super eis cum bene eis fécero in toto Corde meo.\end{Resp}
\begin{Resp}\Rsign Ut non recédant a me.\end{Resp}
\switchcolumn
\EN
\begin{Resp}\Rsign I will make an everlasting Covenant with them, and I will not cease from doing them good, and I will put my fear in their hearts, \Star{} So that they shall not depart from me.\end{Resp}
\begin{Resp}\Vsign Yea, I will rejoice over them, to do them good with my whole Heart.\end{Resp}
\begin{Resp}\Rsign So that they shall not depart from me.\end{Resp}
\switchcolumn*

\LA
\LessonHead{Lectio II}
\switchcolumn
\EN
\LessonHead{Lesson II}
\switchcolumn*

\LA
\Cite{1 Reg 12:6-9}
\switchcolumn
\EN
\Cite{1 Sam 12:6-9}
\switchcolumn*

\LA
\noindent Et ait Sámuel ad pópulum: Dóminus, qui fecit Móysen et Aaron, et edúxit patres nostros de terra Ægýpti. Nunc ergo state, ut judício conténdam advérsum vos coram Dómino de ómnibus misericórdiis Dómini quas fecit vobíscum et cum pátribus vestris: quómodo Jacob ingréssus est in Ægýptum, et clamavérunt patres vestri ad Dóminum: et misit Dóminus Móysen et Aaron, et edúxit patres vestros de Ægýpto, et collocávit eos in loco hoc. Qui oblíti sunt Dómini Dei sui, et trádidit eos in manu Sísaræ magístri milítiæ Hasor, et in manu Philisthinórum, et in manu regis Moab: et pugnavérunt advérsum eos.\par
\switchcolumn
\EN
\noindent And Samuel said to the people: It is the Lord, who made Moses and Aaron, and brought our fathers out of the land of Egypt. Now therefore stand up, that I may plead in judgment against you before the Lord, concerning all the kindness of the Lord, which he hath shown to you, and to your fathers: How Jacob went into Egypt, and your fathers cried to the Lord: and the Lord sent Moses and Aaron, and brought your fathers out of Egypt: and made them dwell in this place. And they forgot the Lord their God, and he delivered them into the hands of Sisara, captain of the army of Hasor, and into the hands of the Philistines, and into the hand of the king of Moab, and they fought against them.\par
\switchcolumn*

\LA
\begin{Resp}\Rsign Si inimícus meus maledixísset mihi, sustinuíssem útique \Star{} Tu vero homo unánimis qui simul mecum dulces capiébas cibos.\end{Resp}
\begin{Resp}\Vsign Et si is qui me óderat super me magna locútus fuísset, abscondíssem me fórsitan ab eo.\end{Resp}
\begin{Resp}\Rsign Tu vero homo unánimis qui simul mecum dulces capiébas cibos.\end{Resp}
\switchcolumn
\EN
\begin{Resp}\Rsign It was not an open enemy that has done me this dishonour, for then I could have borne it, \Star{} But it was even thou, mine own familiar friend, who did also eat of my Bread.\end{Resp}
\begin{Resp}\Vsign Neither was it mine adversary that did magnify himself against me, for then peradventure I would have hid myself from him.\end{Resp}
\begin{Resp}\Rsign But it was even thou, mine own familiar friend, who did also eat of my Bread.\end{Resp}
\switchcolumn*

\LA
\LessonHead{Lectio III}
\switchcolumn
\EN
\LessonHead{Lesson III}
\switchcolumn*

\LA
\Cite{1 Reg 12:10-14}
\switchcolumn
\EN
\Cite{1 Sam 12:10-14}
\switchcolumn*

\LA
\noindent Póstea autem clamavérunt ad Dóminum, et dixérunt: Peccávimus, quia derelíquimus Dóminum, et servívimus Báalim et Astaroth: nunc ergo érue nos de manu inimicórum nostrórum, et serviémus tibi. Et misit Dóminus Jeróbaal, et Badan, et Jephte, et Sámuel, et éruit vos de manu inimicórum vestrórum per circúitum, et habitástis confidénter. Vidéntes autem quod Naas rex filiórum Ammon venísset advérsum vos, dixístis mihi: Nequáquam, sed rex imperábit nobis: cum Dóminus Deus vester regnáret in vobis. Nunc ergo præsto est rex vester, quem elegístis et petístis: ecce dedit vobis Dóminus regem. Si timuéritis Dóminum, et serviéritis ei, et audiéritis vocem ejus, et non exasperavéritis os Dómini, éritis et vos, et rex qui ímperat vobis, sequéntes Dóminum Deum vestrum.\par
\switchcolumn
\EN
\noindent But afterwards they cried to the Lord, and said: We have sinned, because we have forsaken the Lord, and have served Baalim and Astaroth: but now deliver us from the hand of our enemies, and we will serve thee. And the Lord sent Jerobaal, and Badan, and Jephte, and Samuel, and delivered you from the hand of your enemies round about, and you dwelt securely. But seeing that Naas king of the children of Ammon was come against you, you said to me: Nay, but a king shall reign over us: whereas the Lord your God was your king. Now therefore your king is here, whom you have chosen and desired: Behold the Lord hath given you a king. If you will fear the Lord, and serve him, and hearken to his voice, and not provoke the mouth of the Lord: then shall both you, and the king who reigneth over you, be followers of the Lord your God\par
\switchcolumn*

\LA
\begin{Resp}\Rsign Cum essémus mórtui peccátis, convivificávit nos Deus in Christo \Star{} Propter nímiam caritátem suam qua diléxit nos.\end{Resp}
\begin{Resp}\Vsign Ut osténderet in sǽculis superveniéntibus abundántes divítias grátiæ suæ.\end{Resp}
\begin{Resp}\Rsign Propter nímiam caritátem suam qua diléxit nos.\end{Resp}
\begin{Resp}\Vsign Glória Patri, et Fílio, \Star{} et Spirítui Sancto.\end{Resp}
\begin{Resp}\Rsign Propter nímiam caritátem suam qua diléxit nos.\end{Resp}
\switchcolumn
\EN
\begin{Resp}\Rsign And we, being dead in our sins, hath God quickened together with Christ, \Star{} For his great love wherewith he hath loved us.\end{Resp}
\begin{Resp}\Vsign That in the ages to come he might shew the exceeding riches of his grace.\end{Resp}
\begin{Resp}\Rsign For his great love wherewith he hath loved us.\end{Resp}
\begin{Resp}\Vsign Glory be to the Father, and to the Son, \Star{} and to the Holy Ghost.\end{Resp}
\begin{Resp}\Rsign For his great love wherewith he hath loved us.\end{Resp}
\switchcolumn*

\LA
\LessonHead{Lectio IV}
\switchcolumn
\EN
\LessonHead{Lesson IV}
\switchcolumn*

\LA
\noindent Ex lítteris Encyclicis Pii Papæ undécimi\par
\switchcolumn
\EN
\noindent According to that degree of perfection wherewith our oblation and our sacrifice do correspond to the Sacrifice of our Lord, that is to say, to the extent that we have immolated love of self and its passions, and thereby have crucified our flesh in that mystical crucifixion concerning which the Apostle wrote, in that same degree shall we gather the fruits of propitiation and expiation for ourselves and for others. For a wondrous bond doth join all the faithful unto Christ, namely, that bond which uniteth the Head with the members of the body, which is to say, the Communion of the Saints, a bond full of mystery in which we as Catholics do nevertheless verily believe. By virtue of this bond, individuals and nations are not only united the one with the other, but likewise with the Head itself, which is Christ: from whom the whole body, fitly joined together and compacted by that which every joint supplieth, according to the working in due measure of each several part, maketh increase of the body unto the building up of itself in love. This verily was the prayer which Christ Jesus himself, the Mediator between God and man, made at the time of his death: I in them, and thou in me, that they may be perfected into one.\par
\switchcolumn*

\LA
\begin{Resp}\Rsign Prope est Dóminus ómnibus invocántibus eum, \Star{} Omnibus invocántibus eum in veritáte.\end{Resp}
\begin{Resp}\Vsign Miserátor et miséricors Dóminus, pátiens et multum miséricors.\end{Resp}
\begin{Resp}\Rsign Omnibus invocántibus eum in veritáte.\end{Resp}
\switchcolumn
\EN
\begin{Resp}\Rsign The Lord is nigh unto all them that call upon him, \Star{} Yea, unto all such as call upon him faithfully.\end{Resp}
\begin{Resp}\Vsign The Lord is full of compassion and mercy, long-suffering and of great goodness.\end{Resp}
\begin{Resp}\Rsign Yea, unto all such as call upon him faithfully.\end{Resp}
\switchcolumn*

\LA
\LessonHead{Lectio V}
\switchcolumn
\EN
\LessonHead{Lesson V}
\switchcolumn*

\LA
\noindent Quemádmodum ígitur uniónem cum Christo profitétur ac firmat consecrátio, ita expiátio eámdem uniónem, et culpas detergéndo ínchoat, et Christi passiónes participándo pérficit, et víctimas pro frátribus offeréndo consúmmat. Atque id sane miseréntis Jesu consílium fuit, cum Cor nobis suum, insígnia passiónis prǽferens ac flammas amóris osténtans, patére vóluit, scílicet ut hinc infinítam peccáti malítiam conjectántes, illinc Reparatiónis caritátem infinítam admiráti, et peccátum veheméntius detestarémur et caritátis ardéntius vicem redderémus. Et vere expiatiónis potíssimum seu reparatiónis spíritus primas semper potiorésque partes hábuit in cultu Sacratíssimo Cordi Jesu exhibéndo, nihílque eo congruéntius orígini, índoli, virtúti, indústriis quæ huic religiónis formæ sunt própriæ, ut rerum memória et usus, sacra item litúrgia atque Summórum Pontíficum acta confírmant.\par
\switchcolumn
\EN
\noindent Therefore, even as our effort at consecration doth manifest and strengthen our union with Christ, so our practice of expiation (by purifying us from sins) is the beginning of such union; wherefrom our participation in the sufferings of Christ is the means of perfecting such union; and the offering which we make to him of our sacrifices for the welfare of our brethren bringeth such union to its final consummation. Now this is precisely the design of the mercy of Jesus, when he doth unveil to the gaze of mankind his Heart, surrounded by the emblems of his passion, and aflame with the Fire of Love, namely: that we, (on the one hand, perceiving the unlimited malice of sin, and on the other, filled with a knowledge of the infinite love of him who is The Reparator,) may detest sin more heartily, and substitute for it a burning love for him. And verily, the spirit of expiation or of reparation hath always played a chief part in the devotion to the most Sacred Heart of Jesus, and reparation is most consonant with the origin, nature, efficacy and particular practices of this special devotion, a fact confirmed by history and the customs of the faithful, by the sacred liturgy, and by the official documents of the Supreme Pontiffs.\par
\switchcolumn*

\LA
\begin{Resp}\Rsign Confíteor tibi, Pater, Dómine cæli et terræ, quia abscondísti hæc a sapiéntibus et prudéntibus \Star{} Et revelásti ea párvulis.\end{Resp}
\begin{Resp}\Vsign Ita, Pater, quóniam sic fuit plácitum ante te.\end{Resp}
\begin{Resp}\Rsign Et revelásti ea párvulis.\end{Resp}
\switchcolumn
\EN
\begin{Resp}\Rsign I thank thee, O Father, Lord of heaven and earth, because thou hast hid these things from the wise and prudent; \Star{} Yea, thou hast revealed them unto babes.\end{Resp}
\begin{Resp}\Vsign Even so, Father, for so it seemed good in thy sight.\end{Resp}
\begin{Resp}\Rsign Yea, thou hast revealed them unto babes.\end{Resp}
\switchcolumn*

\LA
\LessonHead{Lectio VI}
\switchcolumn
\EN
\LessonHead{Lesson VI}
\switchcolumn*

\LA
\noindent Síquidem cum se conspiciéndum Margarítæ Maríæ exhibéret Christus, caritátis suæ infinitátem prǽdicans, simul, mæréntis instar, tot tantásque sibi inústas ab ingrátis homínibus injúrias in hæc verba conquéstus est, quæ útinam in piórum ánimis insidérent nulláque unquam oblivióne deleréntur: « En Cor illud, inquit, quod tantópere hómines amávit beneficiísque ómnibus cumulávit, quodque amóri suo infiníto non tantum rédditam grátiam nullam invénit, at contra, obliviónem, negléctum, contumélias, eásque ab iis étiam illátas nonnúnquam, qui amóris peculiáris débito officióque teneréntur ».\par
\switchcolumn
\EN
\noindent Inasmuch as, when Christ revealed himself to the sight of Margaret Mary, though he then insisted on the immensity of his love, at the same time, with sorrowful mien, he grieved over the great number of horrible outrages heaped upon him by the ingratitude of mankind, he used then these words, words which should be graven on the hearts of all pious souls so as never to be forgotten by them: Behold that Heart which hath so loved men; which same hath heaped upon them so many benefits; in return for whose infinite love no gratitude is to be found; but instead cometh unto it forgetfulness, indifference, outrages; and all such things do come at times even from souls that are bound closely thereunto by the bonds of a very special love.\par
\switchcolumn*

\LA
\begin{Resp}\Rsign Omnes gentes quascúmque fecísti, vénient \Star{} Et adorábunt coram te, Dómine.\end{Resp}
\begin{Resp}\Vsign Et glorificábunt nomen tuum, quóniam magnus es tu, et fáciens mirabília.\end{Resp}
\begin{Resp}\Rsign Et adorábunt coram te, Dómine.\end{Resp}
\begin{Resp}\Vsign Glória Patri, et Fílio, \Star{} et Spirítui Sancto.\end{Resp}
\begin{Resp}\Rsign Et adorábunt coram te, Dómine.\end{Resp}
\switchcolumn
\EN
\begin{Resp}\Rsign All nations whom thou hast made shall come, \Star{} And they shall worship thee, O Lord.\end{Resp}
\begin{Resp}\Vsign Yea, they shall glorify thy Name, for thou art great, and doest wondrous things.\end{Resp}
\begin{Resp}\Rsign And they shall worship thee, O Lord.\end{Resp}
\begin{Resp}\Vsign Glory be to the Father, and to the Son, \Star{} and to the Holy Ghost.\end{Resp}
\begin{Resp}\Rsign And they shall worship thee, O Lord.\end{Resp}
\switchcolumn*

\LA
\LessonHead{Lectio VII}
\switchcolumn
\EN
\LessonHead{Lesson VII}
\switchcolumn*

\LA
\LessonTitle{Léctio sancti Evangélii secúndum Joánnem}
\switchcolumn
\EN
\LessonTitle{From the Holy Gospel according to John}
\switchcolumn*

\LA
\Cite{Joannes 19:31-37}
\switchcolumn
\EN
\Cite{John 19:31-37}
\switchcolumn*

\LA
\noindent In illo témpore: Judǽi, quóniam parascéve erat, ut non remanérent in cruce córpora sábbato (erat enim magnus dies ille sábbati) rogavérunt Pilátum, ut frangeréntur eórum crura et tolleréntur. Et réliqua.\par
\switchcolumn
\EN
\noindent At that time: The Jews, because it was the Preparation, that the bodies should not remain upon the cross on the Sabbath Day, for that Sabbath Day was an high day, besought Pilate that their legs might broken, and that they might be taken away. And so on.\par
\switchcolumn*

\LA

\switchcolumn
\EN
\LessonTitle{A Homily by St. Bernardin of Siena John continueth: One of the soldiers with a spear pierced his side, and forthwith came there out blood and water. O Love, thou that canst dissolve all things, how thou didst flow forth from our Beloved, for the sake of our redemption! In order that this thy flood of love might spread everywhere, the great firmament was rent above us, even the heights of the Heart of Jesus, which the cruel spear did not fail to pierce to its innermost abyss. And forthwith came there out blood and water. Blood for redemption flowed out, but also water for cleansing; whence the Church was formed from the side of Christ, that she might know herself ever to be the sole-beloved one of Christ, and that she might understand how greatly sin displeaseth him, as she seeth his divine blood flow forth from the God-Man, both in life and in death. For if the divine blood be shed for us, then we cost not a little.}
\switchcolumn*

\LA
\Source{Homilia sancti Bernardíni Senénsis}
\switchcolumn
\EN

\switchcolumn*

\LA
\noindent Joánnes subdit: Unus mílitum láncea latus ejus apéruit et contínuo exívit sanguis et aqua. O amor qui ómnia liquas! quómodo pro redemptióne nostra reliquísti dilectórem nostrum? Nam, ut úndique inundáret amóris dilúvium, super nos ruptæ sunt abýssi magnæ; scílicet, penetrália Cordis Jesu, quibus, ad íntima progrédiens, dira láncea non pepércit. Sanguis exívit et aqua. Sanguis in redemptiónem, sed étiam in ablutiónem aqua deflúxit; unde formáta est Ecclésia ex látere Christi, ut ætérne únicam atque diléctam a Christo se discat, et ut recognóscat quam displícuit culpa pro qua sanguis divínus ex hómine Deo vivo et mórtuo ita deflúxit. Non enim parva quantitáte constámus, si pro nobis sanguis divínus effúnditur.\par
\switchcolumn
\EN

\switchcolumn*

\LA
\begin{Resp}\Rsign Ego si exaltátus fúero a terra \Star{} Omnia traham ad meípsum.\end{Resp}
\begin{Resp}\Vsign Hoc autem dicébat signíficans qua morte esset moritúrus.\end{Resp}
\begin{Resp}\Rsign Omnia traham ad meípsum.\end{Resp}
\switchcolumn
\EN
\begin{Resp}\Rsign If I be lifted up, \Star{} I will draw all men unto me.\end{Resp}
\begin{Resp}\Vsign This he said, signifying what death he should die.\end{Resp}
\begin{Resp}\Rsign I will draw all men unto me.\end{Resp}
\switchcolumn*

\LA
\LessonHead{Lectio VIII}
\switchcolumn
\EN
\LessonHead{Lesson VIII}
\switchcolumn*

\LA
\noindent Aqua ad lítteram non cum sánguine indistíncta deflúxit. Neque enim potuísset ab insipiéntibus comprehéndi, si mixta cum sánguine defluxísset. Et forte totus sanguis deflúxit ex illo divíno córpore in signum totíus amóris effúsi, post quem humor áqueus egréssus est. Quod quidem est alto mystério factum, ut prius egrederétur ex eódem córpore rédimens prétium, deínde aqua in qua multitúdo populórum redémpta significátur. Sunt enim aquæ multæ, pópuli multi; tamen qui ad christiánam fidem pértinent unus fidélis pópulus sunt, ut non sint aquæ, sed aqua quæ manávit ex látere Christi, sicut prima Corinthiórum cápite décimo Apóstolus ait: Unus panis, et unum corpus multi sumus omnes, qui de uno pane et de uno cálice participámus. Et íterum ad Ephésios, cápite quarto, inquit: Unus Deus, una fides, unum baptísma.\par
\switchcolumn
\EN
\noindent The water flowed not forth intermingled with the blood. For if it had so done, it could not have been perceived by simple folk. And perchance all the blood flowed forth from that divine body, as a sign that all his love was poured out, after which a watery humour came forth. Verily this was a token of a profound mystery, to wit, that from one and the same body came forth first the price of our redemption amongst all peoples. For many waters may be taken to signify many peoples; yet all who belong to the Christian Faith, are one people in the Faith, so that they are not as many waters; rather they are as one stream of water, and as such did flow from the side of Christ, as saith the Apostle in the tenth chapter of the former Epistle to the Corinthians: For we being many are one bread, and one body: for we are all partakers of that one bread and one cup. And again, in the fourth chapter of Ephesians, he saith: One God, one faith, one baptism.\par
\switchcolumn*

\LA
\begin{Resp}\Rsign Simus ergo imitatóres Dei \Star{} Et ambulémus in dilectióne.\end{Resp}
\begin{Resp}\Vsign Sicut et Christus diléxit nos et trádidit semetípsum pro nobis.\end{Resp}
\begin{Resp}\Rsign Et ambulémus in dilectióne.\end{Resp}
\begin{Resp}\Vsign Glória Patri, et Fílio, \Star{} et Spirítui Sancto.\end{Resp}
\begin{Resp}\Rsign Et ambulémus in dilectióne.\end{Resp}
\switchcolumn
\EN
\begin{Resp}\Rsign Be ye therefore followers of God; \Star{} And walk ye therefore in love.\end{Resp}
\begin{Resp}\Vsign For Christ also hath loved us and hath given himself for us.\end{Resp}
\begin{Resp}\Rsign And walk ye therefore in love.\end{Resp}
\begin{Resp}\Vsign Glory be to the Father, and to the Son, \Star{} and to the Holy Ghost.\end{Resp}
\begin{Resp}\Rsign And walk ye therefore in love.\end{Resp}
\switchcolumn*

\LA
\LessonHead{Lectio IX}
\switchcolumn
\EN
\LessonHead{Lesson IX}
\switchcolumn*

\LA
\noindent Notánter tamen adverténdum est quod latus Christi apértum dícitur, non vulnerátum; quóniam próprie vulnus prætérquam in vivo córpore fíeri nequit. Ait enim Evangelísta Joánnes: Unus mílitum láncea latus ejus apéruit; ut, apérto látere, cognoscámus dilectiónem Cordis sui usque ad mortem, et ad illum ineffábilem amórem ejus ingrediámur quo ille ad nos procéssit. Accedámus ergo ad Cor ejus, Cor altum, Cor secrétum, Cor ómnia cógitans, Cor ómnia sciens, Cor díligens, immo amóre ardens; et apértam portam intelligámus saltem in amóris veheméntia; cordifórmes ingrediámur ad secrétum ab ætérno abscónditum, nunc vero in morte quasi apérto látere revelátum; quóniam apértio láteris ætérni templi apertiónem demónstrat, ubi ómnium exsisténtium consummáta est felícitas ætérna.\par
\switchcolumn
\EN
\noindent Nevertheless, it is especially deserving of our attention that the side of Christ is said to have been pierced, (that is, set open,) not wounded, since a wound cannot properly be said to be inflicted except on a living body. For John the Evangelist saith: One of the soldiers with a spear pierced his side: that through this side thus opened, we might become aware of the love of his Heart, even unto death, and that we might enter into that unutterable love of his, through the same channel whereby it came unto us. Let us draw near then to his Heart, a deep Heart, a hidden Heart, a Heart thinking of all, a Heart knowing all, a Heart loving, yea, even on fire with love. Let us also recognize at least in the vehemence of his love that the gate is open; with our hearts made like unto his, let us enter into that secret place, hidden from all eternity, but now in death revealed, as it were, through the open side; for the opening of his side is a figure of the opening of the eternal temple, where is consummated the everlasting happiness of every creature.\par
\switchcolumn*

\LA
\TeDeum{Te Deum}
\switchcolumn
\EN
\TeDeum{Te Deum}
\switchcolumn*

\end{paracol}

\Office{Die VI infra Octavam SSmi Cordis Jesu}{Sixth Day within the Octave of the Sacred Heart}{Semiduplex}
\addcontentsline{toc}{subsection}{Die VI infra Octavam SSmi Cordis Jesu \textit{— Sixth Day within the Octave of the Sacred Heart}}
\begin{paracol}{2}
\LA
\LessonHead{Lectio I}
\switchcolumn
\EN
\LessonHead{Lesson I}
\switchcolumn*

\LA
\LessonTitle{De libro primo Regum}
\switchcolumn
\EN
\LessonTitle{Lesson from the first book of Samuel}
\switchcolumn*

\LA
\Cite{1 Reg 13:1-4}
\switchcolumn
\EN
\Cite{1 Sam 13:1-4}
\switchcolumn*

\LA
\noindent Fílius uníus anni erat Saul, cum regnáre cœpísset; duóbus autem annis regnávit super Israël. Et elégit sibi Saul tria míllia de Israël. Et erant cum Saul duo míllia in Machmas et in monte Bethel, mille autem cum Jónatha in Gábaa Bénjamin. Porro céterum pópulum remísit unumquémque in tabernácula sua. Et percússit Jónathas statiónem Philisthinórum, quæ erat in Gábaa. Quod cum audíssent Philísthiim Saul cécinit búccina in omni terra dicens: Audiant Hebrǽi. Et univérsus Israël audívit hujuscémodi famam: Percússit Saul statiónem Philisthinórum, et eréxit se Israël advérsus Philísthiim; clamávit ergo pópulus post Saul in Gálgala.\par
\switchcolumn
\EN
\noindent Saul was a child of one year when he began to reign, and he reigned two years over Israel. And Saul chose him three thousand men of Israel: and two thousand were with Saul in Machmas, and in mount Bethel: and a thousand with Jonathan in Gabaa of Benjamin, and the rest of the people he sent back every man to their dwellings. And Jonathan smote the garrison of the Philistines which was in Gabaa. And when the Philistines had heard of it, Saul sounded the trumpet over all the land, saying: Let the Hebrews hear. And all Israel heard this report: Saul hath smitten the garrison of the Philistines: and Israel took courage against the Philistines. And the people were called together after Saul to Galgal.\par
\switchcolumn*

\LA
\begin{Resp}\Rsign Fériam eis pactum sempitérnum et non désinam eis benefácere et timórem meum dabo in corde eórum \Star{} Ut non recédant a me.\end{Resp}
\begin{Resp}\Vsign Et lætábor super eis cum bene eis fécero in toto Corde meo.\end{Resp}
\begin{Resp}\Rsign Ut non recédant a me.\end{Resp}
\switchcolumn
\EN
\begin{Resp}\Rsign I will make an everlasting Covenant with them, and I will not cease from doing them good, and I will put my fear in their hearts, \Star{} So that they shall not depart from me.\end{Resp}
\begin{Resp}\Vsign Yea, I will rejoice over them, to do them good with my whole Heart.\end{Resp}
\begin{Resp}\Rsign So that they shall not depart from me.\end{Resp}
\switchcolumn*

\LA
\LessonHead{Lectio II}
\switchcolumn
\EN
\LessonHead{Lesson II}
\switchcolumn*

\LA
\Cite{1 Reg 13:5-8}
\switchcolumn
\EN
\Cite{1 Sam 13:5-8}
\switchcolumn*

\LA
\noindent Et Philísthiim congregáti sunt ad præliándum contra Israël, trigínta míllia cúrruum et sex míllia équitum et réliquum vulgus, sicut aréna, quæ est in líttore maris plúrima. Et ascendéntes castrametáti sunt in Machmas ad oriéntem Betháven. Quod cum vidíssent viri Israël se in arcto pósitos, (afflíctus enim erat pópulus) abscondérunt se in spelúncis et in ábditis, in petris quoque et in antris et in cistérnis. Hebrǽi autem transiérunt Jordánem in terram Gad et Gálaad. Cumque adhuc esset Saul in Gálgala, univérsus pópulus pertérritus est, qui sequebátur eum. Et exspectávit septem diébus juxta plácitum Samuélis; et non venit Sámuel in Gálgala, dilapsúsque est pópulus ab eo.\par
\switchcolumn
\EN
\noindent The Philistines also were assembled to fight against Israel, thirty thousand chariots, and six thousand horsemen, and a multitude of people besides, like the sand on the sea shore for number. And going up they camped in Machmas at the east of Bethaven. And when the men of Israel saw that they were straitened, (for the people were distressed,) they hid themselves in caves, and in thickets, and in rocks, and in dens, and in pits. And some of the Hebrews passed over the Jordan into the land of Gad and Galaad. And when Saul was yet in Galgal, all the people that followed him were greatly afraid. And he waited seven days according to the appointment of Samuel, and Samuel came not to Galgal, and the people slipt away from him.\par
\switchcolumn*

\LA
\begin{Resp}\Rsign Si inimícus meus maledixísset mihi, sustinuíssem útique \Star{} Tu vero homo unánimis qui simul mecum dulces capiébas cibos.\end{Resp}
\begin{Resp}\Vsign Et si is qui me óderat super me magna locútus fuísset, abscondíssem me fórsitan ab eo.\end{Resp}
\begin{Resp}\Rsign Tu vero homo unánimis qui simul mecum dulces capiébas cibos.\end{Resp}
\switchcolumn
\EN
\begin{Resp}\Rsign It was not an open enemy that has done me this dishonour, for then I could have borne it, \Star{} But it was even thou, mine own familiar friend, who did also eat of my Bread.\end{Resp}
\begin{Resp}\Vsign Neither was it mine adversary that did magnify himself against me, for then peradventure I would have hid myself from him.\end{Resp}
\begin{Resp}\Rsign But it was even thou, mine own familiar friend, who did also eat of my Bread.\end{Resp}
\switchcolumn*

\LA
\LessonHead{Lectio III}
\switchcolumn
\EN
\LessonHead{Lesson III}
\switchcolumn*

\LA
\Cite{1 Reg 13:9-14}
\switchcolumn
\EN
\Cite{1 Sam 13:9-14}
\switchcolumn*

\LA
\noindent Ait ergo Saul: Afférte mihi holocáustum et pacífica. Et óbtulit holocáustum. Cumque complésset ófferens holocáustum, ecce Samuel veniébat; et egréssus est Saul óbviam ei ut salutáret eum. Locutúsque est ad eum Samuel: Quid fecísti? Respóndit Saul: Quia vidi quod pópulus dilaberétur a me, et tu non véneras juxta plácitos dies, porro Philísthiim congregáti fúerant in Machmas, dixi: Nunc descéndent Philísthiim ad me in Gálgala, et fáciem Dómini non placávi. Necessitáte compúlsus, óbtuli holocáustum. Dixítque Sámuel ad Saul: Stulte egísti, nec custodísti mandáta Dómini Dei tui, quæ præcépit tibi. Quod si non fecísses, jam nunc præparásset Dóminus regnum tuum super Israël in sempitérnum; sed nequáquam regnum tuum ultra consúrget. Quæsívit Dóminus sibi virum juxta cor suum et præcépit ei Dóminus ut esset dux super pópulum suum, eo quod non serváveris quæ præcépit Dóminus.\par
\switchcolumn
\EN
\noindent Then Saul said: Bring me the holocaust, and the peace offerings. And he offered the holocaust. And when he had made an end of offering the holocaust, behold Samuel came: and Saul went forth to meet him and salute him. And Samuel said to him: What hast thou done? Saul answered: Because I saw that the people slipt from me, and thou wast not come according to the days appointed, and the Philistines were gathered together in Machmas, I said: Now will the Philistines come down upon me to Galgal, and I have not appeased the face of the Lord. Forced by necessity, I offered the holocaust. And Samuel said to Saul: Thou hast done foolishly, and hast not kept the commandments of the Lord thy God, which he commanded thee. And if thou hadst not done thus, the Lord would now have established thy kingdom over Israel for ever. But thy kingdom shall not continue. The Lord hath sought him a man according to his own heart: and him hath the Lord commanded to be prince over his people, because thou hast not observed that which the Lord commanded.\par
\switchcolumn*

\LA
\begin{Resp}\Rsign Cum essémus mórtui peccátis, convivificávit nos Deus in Christo \Star{} Propter nímiam caritátem suam qua diléxit nos.\end{Resp}
\begin{Resp}\Vsign Ut osténderet in sǽculis superveniéntibus abundántes divítias grátiæ suæ.\end{Resp}
\begin{Resp}\Rsign Propter nímiam caritátem suam qua diléxit nos.\end{Resp}
\begin{Resp}\Vsign Glória Patri, et Fílio, \Star{} et Spirítui Sancto.\end{Resp}
\begin{Resp}\Rsign Propter nímiam caritátem suam qua diléxit nos.\end{Resp}
\switchcolumn
\EN
\begin{Resp}\Rsign And we, being dead in our sins, hath God quickened together with Christ, \Star{} For his great love wherewith he hath loved us.\end{Resp}
\begin{Resp}\Vsign That in the ages to come he might shew the exceeding riches of his grace.\end{Resp}
\begin{Resp}\Rsign For his great love wherewith he hath loved us.\end{Resp}
\begin{Resp}\Vsign Glory be to the Father, and to the Son, \Star{} and to the Holy Ghost.\end{Resp}
\begin{Resp}\Rsign For his great love wherewith he hath loved us.\end{Resp}
\switchcolumn*

\LA
\LessonHead{Lectio IV}
\switchcolumn
\EN
\LessonHead{Lesson IV}
\switchcolumn*

\LA
\Source{Ex lítteris Encyclicis Pii Papæ undécimi}
\switchcolumn
\EN
\Source{From the Encyclical Letter of Pope Pius XI}
\switchcolumn*

\LA
\noindent At enim beáte regnántem Christum in cælis qui piaculáres ritus consolári queant? Scílicet « da amántem et sentit quod dico » repónimus, Augustíni verbis usi, quæ in hunc locum aptíssime cadunt. Dei enim amantíssimus quisque, si prætériti témporis spátium respíciat, videt meditándo intuetúrque Christum pro hómine laborántem, doléntem, duríssima quæque perpetiéntem, « propter nos hómines et propter nostram salútem » tristítia, angóribus, oppróbriis pæne conféctum, immo « attrítum propter scélera nostra » ac suo nos livóre sanántem. Atque hæc ómnia eo vérius piórum meditántur ánimi, quod peccáta hóminum ac flagítia quovis témpore perpetráta in causa fuérunt cur Dei Fílius morti traderétur, eadémque nunc étiam mortem ipsam per se essent Christo illatúra, iísdem cum dolóribus mæroribúsque conjúnctam, quippe síngula passiónem Dómini, suo quodam modo renováre censeántur: « Rursus crucifigéntes sibimetípsis Fílium Dei et osténtui habéntes ».\par
\switchcolumn
\EN
\noindent But in what sense can it be said that our expiatory practices can give consolation to Christ, since he is now reigning in heavenly joy? We answer in the words of Saint Augustine (from a passage entirely appropriate to this subject): Give me a lover, and he will feel the truth of what I say. Every soul which is on fire with the love of God, if it but turns its thoughts to the past, doth in its meditation perceive and contemplate Christ suffering for mankind; afflicted by grief in the midst of sorrows endured for us men and for our salvation; almost overcome by afflictions, vexations, and reproaches; yea, bruised by our sins; with whose stripes we are healed. And devout souls will have an appreciation in proportion to their understanding of these mysteries, if they perceive that the sins and crimes of men (no matter when committed) were the real reason why the Son of God was condemned to death, and that sins committed in this present are able in some wise to cause the death of Christ, renewing as they do, the cause of that selfsame death, with its sufferings and agonies, which was accomplished on the Cross; for every sin must be said to renew, in some fashion, the Lord's passion, concerning which we read in the Scriptures: They crucify themselves the Son of God afresh, and put him to an open shame.\par
\switchcolumn*

\LA
\begin{Resp}\Rsign Prope est Dóminus ómnibus invocántibus eum, \Star{} Omnibus invocántibus eum in veritáte.\end{Resp}
\begin{Resp}\Vsign Miserátor et miséricors Dóminus, pátiens et multum miséricors.\end{Resp}
\begin{Resp}\Rsign Omnibus invocántibus eum in veritáte.\end{Resp}
\switchcolumn
\EN
\begin{Resp}\Rsign The Lord is nigh unto all them that call upon him, \Star{} Yea, unto all such as call upon him faithfully.\end{Resp}
\begin{Resp}\Vsign The Lord is full of compassion and mercy, long-suffering and of great goodness.\end{Resp}
\begin{Resp}\Rsign Yea, unto all such as call upon him faithfully.\end{Resp}
\switchcolumn*

\LA
\LessonHead{Lectio V}
\switchcolumn
\EN
\LessonHead{Lesson V}
\switchcolumn*

\LA
\noindent Quodsi propter peccáta quoque nostra, quæ futúra quidem erant at prævísa, ánima Christi tristis facta est usque ad mortem, haud dúbium quin solátii nonníhil jam tum céperit étiam e nostra item prævísa, reparatióne, cum « appáruit illi Angelus de cælo » ut Cor ejus tædio et angóribus oppréssum consolarétur. Atque ita Cor illud sacratíssimum, quod ingratórum hóminum peccátis continénter sauciátur, étiam nunc mira quidem sed vera ratióne solári póssumus ac debémus, quandóquidem, ut in sacra quoque liturgía légitur, ex ore Psaltis, Christus ipse se ab amícis suis derelíctum conquéritur: « Impropérium exspectávit Cor meum et misériam, et sustínui qui simul contristarétur et non fuit, et qui consolarétur et non invéni ».\par
\switchcolumn
\EN
\noindent And since in the agony in the Garden the soul of Christ became sorrowful even unto death, in view of our own future sins which were foreseen by him, there can be no doubt that at the same time he was in some way consoled, through the provision of our acts of reparation, when the Angel from heaven appeared unto him to comfort him in the sorrows of his Heart, which was so bowed down with weariness and grief. And so, even at this present time, we can and we ought to console, in a mystical but nonetheless real manner, that most Sacred Heart: which same is being wounded continually by the sins of thankless men. It is in this sense, namely, that Christ is grieved over his abandonment by his friends that we are accustomed to interpret the holy liturgy whenever we say in the words of the Psalmist: Thy rebuke hath broken my Heart; I am full of heaviness; I looked for some to have pity on me, but there was no man, neither found I any to comfort me.\par
\switchcolumn*

\LA
\begin{Resp}\Rsign Confíteor tibi, Pater, Dómine cæli et terræ, quia abscondísti hæc a sapiéntibus et prudéntibus \Star{} Et revelásti ea párvulis.\end{Resp}
\begin{Resp}\Vsign Ita, Pater, quóniam sic fuit plácitum ante te.\end{Resp}
\begin{Resp}\Rsign Et revelásti ea párvulis.\end{Resp}
\switchcolumn
\EN
\begin{Resp}\Rsign I thank thee, O Father, Lord of heaven and earth, because thou hast hid these things from the wise and prudent; \Star{} Yea, thou hast revealed them unto babes.\end{Resp}
\begin{Resp}\Vsign Even so, Father, for so it seemed good in thy sight.\end{Resp}
\begin{Resp}\Rsign Yea, thou hast revealed them unto babes.\end{Resp}
\switchcolumn*

\LA
\LessonHead{Lectio VI}
\switchcolumn
\EN
\LessonHead{Lesson VI}
\switchcolumn*

\LA
\noindent Accedit quod pássio Christi expiátrix renovátur et quodámmodo continuátur et adimplétur in córpore suo mystico, quod est Ecclésia. Etenim, ut rursus sancti Augustíni verbis utámur, « passus est Christus quidquid pati debúerat; jam de mensúra passiónum nihil deest. Ergo implétæ sunt passiónes, sed in cápite; restábant adhuc Christi passiónes in córpore ». Quod quidem Dóminus ipse Jesus declaráre dignátus est, cum ad Saulum « adhuc spirántem minárum et cædis in discípulos » loquens: « Ego sum, inquit, Jesus quem tu persequéris », haud obscúre signíficans, commótis in Ecclésiam insectatiónibus, ipsum divínum oppugnári ac vexári Ecclésiæ Caput. Jure ígitur meritóque Christus in córpore suo mystico adhuc pátiens, nos expiatiónis suæ sócios habére exóptat, idque étiam ipsa nostra cum eo necessitúdo póstulat; nam cum simus « corpus Christi et membra de membro », quidquid pátitur caput, ómnia cum eo membra patiántur opórtet.\par
\switchcolumn
\EN
\noindent To the foregoing we should add that the expiatory passion of Christ is renewed, and in a certain manner continued, in his mystical body, which same is the Church. In this connection, once more the words of Saint Augustine are appropriate: Christ suffered all that he needed to suffer; nothing at all is lacking in the number of his sufferings; therefore his sufferings are completed in him as the Head; but there remaineth even yet the sufferings of Christ to be endured in the body. Which truth verily the Lord Jesus himself made known, at the time that Saul was breathing out threatenings and slaughter against the disciples, to whom the Lord said: I am Jesus whom thou persecutest. By these words he plainly affirmed that persecutions visited on the Church are in reality directed against the Head of the Church. Meet it is therefore that Christ, thus still suffering in his mystical body, should desire to have us sharers in his own work of expiation, even as our own need moveth us to desire and seek union with him. And since we are the body of Christ, and members of his very flesh and bones, whatever the Head must suffer all the members of the body must needs suffer together with him.\par
\switchcolumn*

\LA
\begin{Resp}\Rsign Omnes gentes quascúmque fecísti, vénient \Star{} Et adorábunt coram te, Dómine.\end{Resp}
\begin{Resp}\Vsign Et glorificábunt nomen tuum, quóniam magnus es tu, et fáciens mirabília.\end{Resp}
\begin{Resp}\Rsign Et adorábunt coram te, Dómine.\end{Resp}
\begin{Resp}\Vsign Glória Patri, et Fílio, \Star{} et Spirítui Sancto.\end{Resp}
\begin{Resp}\Rsign Et adorábunt coram te, Dómine.\end{Resp}
\switchcolumn
\EN
\begin{Resp}\Rsign All nations whom thou hast made shall come, \Star{} And they shall worship thee, O Lord.\end{Resp}
\begin{Resp}\Vsign Yea, they shall glorify thy Name, for thou art great, and doest wondrous things.\end{Resp}
\begin{Resp}\Rsign And they shall worship thee, O Lord.\end{Resp}
\begin{Resp}\Vsign Glory be to the Father, and to the Son, \Star{} and to the Holy Ghost.\end{Resp}
\begin{Resp}\Rsign And they shall worship thee, O Lord.\end{Resp}
\switchcolumn*

\LA
\LessonHead{Lectio VII}
\switchcolumn
\EN
\LessonHead{Lesson VII}
\switchcolumn*

\LA
\LessonTitle{Léctio sancti Evangélii secúndum Joánnem}
\switchcolumn
\EN
\LessonTitle{From the Holy Gospel according to John}
\switchcolumn*

\LA
\Cite{Joannes 19:31-37}
\switchcolumn
\EN
\Cite{John 19:31-37}
\switchcolumn*

\LA
\noindent In illo témpore: Judæi, quóniam parascéve erat, ut non remanérent in cruce córpora sábbato (erat enim magnus dies ille sábbati) rogavérunt Pilátum, ut frangeréntur eórum crura et tolleréntur. Et réliqua.\par
\switchcolumn
\EN
\noindent At that time: The Jews, because it was the Preparation, that the bodies should not remain upon the cross on the Sabbath Day, for that Sabbath Day was an high day, besought Pilate that their legs might broken, and that they might be taken away. And so on.\par
\switchcolumn*

\LA

\switchcolumn
\EN
\LessonTitle{A Homily by St. Peter Canisius the Priest}
\switchcolumn*

\LA
\Source{Homilia sancti Petri Canísii Presbýteri}
\switchcolumn
\EN
\Source{Exhortationes domesticæ. Medit. 6-7}
\switchcolumn*

\LA
\noindent Diligénter tecum ánimo versa, quam ineffábilis fúerit illa cáritas qua Deus summus, in máximis Cordis angústiis et totíus mundi oppróbriis, pro te vilíssimo vermículo illam crucis acerbíssimam mortem perpéssus fúerat. Advérte ut summam Christus Servátor suis liberalitátem ómnibus exhíbuit. Aliquándo enim in médio pópuli stans, ita clamábat: Si quis sitit véniat ad me et bibat: parátum se osténdens ómnibus ómnium necessitátibus subveníre. Consideráto ut liberalíssime tibi Cordis sui pretiósum sánguinem propinávit, quando, sacro látere apérto, quidquid réliquum erat in córpore sánguinis profúdit.\par
\switchcolumn
\EN
\noindent Have thou all diligence to turn over in thy mind this matter, namely, That he, who is God over all, endured with unutterable charity, a most bitter death on the Cross, in the exceeding anguish of his Heart, whilst the whole world mocked him! And this he did for thee, who art but a paltry little worm! Meditate on the boundless generosity which Christ the Preserver manifested to all his own people. For at one time, standing in the midst of the people, he cried out: If any man thirst, let him come to me and drink; thereby shewing how ready he was to welcome everyone, and help him in his every need. Consider how freely he giveth thee to drink of his Heart's precious blood, considering that his sacred side was set open, wherefrom he poured forth whatever blood was still left in his body.\par
\switchcolumn*

\LA
\begin{Resp}\Rsign Ego si exaltátus fúero a terra \Star{} Omnia traham ad meípsum.\end{Resp}
\begin{Resp}\Vsign Hoc autem dicébat signíficans qua morte esset moritúrus.\end{Resp}
\begin{Resp}\Rsign Omnia traham ad meípsum.\end{Resp}
\switchcolumn
\EN
\begin{Resp}\Rsign If I be lifted up, \Star{} I will draw all men unto me.\end{Resp}
\begin{Resp}\Vsign This he said, signifying what death he should die.\end{Resp}
\begin{Resp}\Rsign I will draw all men unto me.\end{Resp}
\switchcolumn*

\LA
\LessonHead{Lectio VIII}
\switchcolumn
\EN
\LessonHead{Lesson VIII}
\switchcolumn*

\LA
\Source{In Evang. Dominica I post Pascha}
\switchcolumn
\EN
\Source{In Evang. Dominica I post Pascha}
\switchcolumn*

\LA
\noindent Quare, ne prorsus ingrátus sim, hos perénnes fontes donórum ac bonórum ómnium mihi ob óculos sæpe propónam, cum de illis dulcíssima exstet promíssio: Hauriétis aquas in gáudio de fóntibus Salvatóris et dicétis in die illa: Confitémini Dómino. Ad hæc ipsa ter beáta petræ non diruéndæ forámina confúgiam; in illis meum nidum firmíssimum ponam, nihil habens antíquius, prætérquam ut in meis angóribus atque perículis vúlnerum Dómini memorándo respírem.\par
\switchcolumn
\EN
\noindent And so that I may not be utterly ungrateful, often I call up before mine eyes those perennial fountains of gifts and of all good things, since from them that most sweet promise standeth out: Ye shall draw waters with joy out of the founts of the Saviour, and ye shall say in that day: Praise ye the Lord. There will I flee, to those thrice-blest holes in the rock which can never be demolished; there will I build for myself a most durable nest, holding nothing better, in all sorrows and perils that befall me, than to think on the wounds of the Lord.\par
\switchcolumn*

\LA
\begin{Resp}\Rsign Simus ergo imitatóres Dei \Star{} Et ambulémus in dilectióne.\end{Resp}
\begin{Resp}\Vsign Sicut et Christus diléxit nos et trádidit semetípsum pro nobis.\end{Resp}
\begin{Resp}\Rsign Et ambulémus in dilectióne.\end{Resp}
\begin{Resp}\Vsign Glória Patri, et Fílio, \Star{} et Spirítui Sancto.\end{Resp}
\begin{Resp}\Rsign Et ambulémus in dilectióne.\end{Resp}
\switchcolumn
\EN
\begin{Resp}\Rsign Be ye therefore followers of God; \Star{} And walk ye therefore in love.\end{Resp}
\begin{Resp}\Vsign For Christ also hath loved us and hath given himself for us.\end{Resp}
\begin{Resp}\Rsign And walk ye therefore in love.\end{Resp}
\begin{Resp}\Vsign Glory be to the Father, and to the Son, \Star{} and to the Holy Ghost.\end{Resp}
\begin{Resp}\Rsign And walk ye therefore in love.\end{Resp}
\switchcolumn*

\LA
\LessonHead{Lectio IX}
\switchcolumn
\EN
\LessonHead{Lesson IX}
\switchcolumn*

\LA
\Source{Medit. 6}
\switchcolumn
\EN
\Source{Medit. 6}
\switchcolumn*

\LA
\noindent Et tu in omni tentatióne fuge diligénter in amábile Christi Cor, ejúsque bonitátem et caritátem tibi propóne, et cum illa confer tuam vilitátem, malítiam, infidelitátem, arrogántiam. Quanta enim cáritas Christi omnes ad se convocántis: Veníte ad me omnes qui laborátis et oneráti estis, et ego refíciam vos, et sic parátum se offert et cupit, amóre nostri, ómnium et singulórum ónera sustinére! Unde magna cum fidúcia in caritátis ejus abyssum prójice peccáta tua et mox te invénies exonerátum.\par
\switchcolumn
\EN
\noindent Therefore in every trial do thou flee quickly to the lovable Heart of Christ, and call to mind his goodness and charity, setting them in contrast to thine own vileness, malice, unfaithfulness, and pride. For consider how great was the charity of Christ in inviting all men unto himself: Come unto me, all ye that travail and are heavy laden, and I will refresh you. It is with this intent that he offereth himself to us, ready and desirous, out of love for us, to carry the burdens of each and all! Cast then thy sins into the abyss of his charity, and straightway thou shalt find thyself lightened of thy load.\par
\switchcolumn*

\LA
\TeDeum{Te Deum}
\switchcolumn
\EN
\TeDeum{Te Deum}
\switchcolumn*

\end{paracol}

\Office{Die VII infra Octavam SSmi Cordis Jesu}{Seventh Day within the Octave of the Sacred Heart}{Semiduplex}
\addcontentsline{toc}{subsection}{Die VII infra Octavam SSmi Cordis Jesu \textit{— Seventh Day within the Octave of the Sacred Heart}}
\begin{paracol}{2}
\LA
\LessonHead{Lectio I}
\switchcolumn
\EN
\LessonHead{Lesson I}
\switchcolumn*

\LA
\LessonTitle{De libro primo Regum}
\switchcolumn
\EN
\LessonTitle{Lesson from the first book of Samuel}
\switchcolumn*

\LA
\Cite{1 Reg 14:6-11}
\switchcolumn
\EN
\Cite{1 Sam 14:6-11}
\switchcolumn*

\LA
\noindent Dixit autem Jónathas ad adolescéntem armígerum suum: Veni, transeámus ad statiónem incircumcisórum horum, si forte fáciat Dóminus pro nobis, quia non est Dómino diffícile salváre vel in multis vel in paucis. Dixítque ei ármiger suus: Fac ómnia, quæ placent ánimo tuo, perge quo cupis, et ero tecum ubicúmque volúeris. Et ait Jónathas: Ecce nos transímus ad viros istos. Cumque apparuérimus eis, si táliter locúti fúerint ad nos: Manéte donec veniámus ad vos; stemus in loco nostro nec ascendámus ad eos. Si autem díxerint: Ascéndite ad nos; ascendámus, quia trádidit eos Dóminus in mánibus nostris: hoc erit nobis signum. Appáruit ígitur utérque statióni Philísthiim. Dixerúntque Philísthiim: En Hebrǽi egrediúntur de cavérnis, in quibus abscónditi fúerant.\par
\switchcolumn
\EN
\noindent And Jonathan said to the young man that bore his armour: Come, let us go over to the garrison of these uncircumcised, it may be the Lord will do for us, because it is easy for the Lord to save either by many, or by few. And his armourbearer said to him: Do all that pleaseth thy mind: go whither thou wilt, and I will be with thee wheresoever thou hast a mind. And Jonathan said: Behold we will go over to these men. And when we shall be seen by them, If they shall speak thus to us: Stay till we come to you: let us stand still in our place, and not go up to them. But if they shall say: Come up to us: let us go up, because the Lord hath delivered them into our hands, this shall be a sign unto us. So both of them discovered themselves to the garrison of the Philistines: and the Philistines said: Behold the Hebrews come forth out of the holes wherein they were hid.\par
\switchcolumn*

\LA
\begin{Resp}\Rsign Fériam eis pactum sempitérnum et non désinam eis benefácere et timórem meum dabo in corde eórum \Star{} Ut non recédant a me.\end{Resp}
\begin{Resp}\Vsign Et lætábor super eis cum bene eis fécero in toto Corde meo.\end{Resp}
\begin{Resp}\Rsign Ut non recédant a me.\end{Resp}
\switchcolumn
\EN
\begin{Resp}\Rsign I will make an everlasting Covenant with them, and I will not cease from doing them good, and I will put my fear in their hearts, \Star{} So that they shall not depart from me.\end{Resp}
\begin{Resp}\Vsign Yea, I will rejoice over them, to do them good with my whole Heart.\end{Resp}
\begin{Resp}\Rsign So that they shall not depart from me.\end{Resp}
\switchcolumn*

\LA
\LessonHead{Lectio II}
\switchcolumn
\EN
\LessonHead{Lesson II}
\switchcolumn*

\LA
\Cite{1 Reg 14:12-15}
\switchcolumn
\EN
\Cite{1 Sam 14:12-15}
\switchcolumn*

\LA
\noindent Et locúti sunt viri de statióne ad Jónatham et ad armígerum ejus, dixerúntque: Ascéndite ad nos, et ostendámus vobis rem. Et ait Jónathas ad armígerum suum: Ascendámus, séquere me; trádidit enim Dóminus eos in manus Israël. Ascéndit autem Jónathas mánibus et pédibus reptans et ármiger ejus post eum. Itaque álii cadébant ante Jónatham, álios ármiger ejus interficiébat sequens eum. Et facta est plaga prima, qua percússit Jónathas et ármiger ejus, quasi vigínti virórum, in média parte júgeri, quam par boum in die aráre consuévit. Et factum est miráculum in castris per agros; sed et omnis pópulus statiónis eórum, qui íerant ad prædándum, obstúpuit, et conturbáta est terra, et áccidit quasi miráculum a Deo.\par
\switchcolumn
\EN
\noindent And the men of the garrison spoke to Jonathan, and to his armourbearer, and said: Come up to us, and we will show you a thing. And Jonathan said to his armourbearer: Let us go up, follow me: for the Lord hath delivered them into the hands of Israel. And Jonathan went up creeping on his hands and feet, and his armourbearer after him. And some fell before Jonathan, others his armourbearer slew as he followed him. And the first slaughter which Jonathan and his armourbearer made, was of about twenty men, within half an acre of land, which a yoke of oxen is wont to plough in a day. And there was a miracle in the camp, through the fields: yea and all the people of their garrison, who had gone out to plunder, were amazed, and the earth trembled: and it happened as a miracle from God.\par
\switchcolumn*

\LA
\begin{Resp}\Rsign Si inimícus meus maledixísset mihi, sustinuíssem útique \Star{} Tu vero homo unánimis qui simul mecum dulces capiébas cibos.\end{Resp}
\begin{Resp}\Vsign Et si is qui me óderat super me magna locútus fuísset, abscondíssem me fórsitan ab eo.\end{Resp}
\begin{Resp}\Rsign Tu vero homo unánimis qui simul mecum dulces capiébas cibos.\end{Resp}
\switchcolumn
\EN
\begin{Resp}\Rsign It was not an open enemy that has done me this dishonour, for then I could have borne it, \Star{} But it was even thou, mine own familiar friend, who did also eat of my Bread.\end{Resp}
\begin{Resp}\Vsign Neither was it mine adversary that did magnify himself against me, for then peradventure I would have hid myself from him.\end{Resp}
\begin{Resp}\Rsign But it was even thou, mine own familiar friend, who did also eat of my Bread.\end{Resp}
\switchcolumn*

\LA
\LessonHead{Lectio III}
\switchcolumn
\EN
\LessonHead{Lesson III}
\switchcolumn*

\LA
\Cite{1 Reg 14:16-20}
\switchcolumn
\EN
\Cite{1 Sam 14:16-20}
\switchcolumn*

\LA
\noindent Et respexérunt speculatóres Saul, qui erant in Gábaa Bénjamin, et ecce multitúdo prostráta et huc illúcque diffúgiens. Et ait Saul pópulo, qui erat cum eo: Requírite et vidéte quis abíerit ex nobis. Cumque requisíssent, repértum est non adésse Jónatham et armígerum ejus. Et ait Saul ad Achíam: Applica arcam Dei (erat enim ibi arca Dei in die illa cum fíliis Israël). Cumque loquerétur Saul ad sacerdótem, tumúltus magnus exórtus est in castris Philisthinórum, crescebátque paulátim et clárius resonábat. Et ait Saul ad sacerdótem: Cóntrahe manum tuam. Conclamávit ergo Saul, et omnis pópulus, qui erat cum eo, et venérunt usque ad locum certáminis; et ecce versus fúerat gládius uniuscujúsque ad próximum suum, et cædes magna nimis.\par
\switchcolumn
\EN
\noindent And the watchmen of Saul, who were in Gabaa of Benjamin looked, and behold a multitude overthrown, and fleeing this way and that. And Saul said to the people that were with him: Look, and see who is gone from us. And when they had sought, it was found that Jonathan and his armourbearer were not there. And Saul said to Achias: Bring the ark of the Lord. (For the ark of God was there that day with the children of Israel.) And while Saul spoke to the priest, there arose a great uproar in the camp of the Philistines: and it increased by degrees, and was heard more clearly. And Saul said to the priest: Draw in thy hand. Then Saul and all the people that were with him, shouted together, and they came to the place of the fight: and behold every man's sword was turned upon his neighbour, and there was a very great slaughter.\par
\switchcolumn*

\LA
\begin{Resp}\Rsign Cum essémus mórtui peccátis, convivificávit nos Deus in Christo \Star{} Propter nímiam caritátem suam qua diléxit nos.\end{Resp}
\begin{Resp}\Vsign Ut osténderet in sǽculis superveniéntibus abundántes divítias grátiæ suæ.\end{Resp}
\begin{Resp}\Rsign Propter nímiam caritátem suam qua diléxit nos.\end{Resp}
\begin{Resp}\Vsign Glória Patri, et Fílio, \Star{} et Spirítui Sancto.\end{Resp}
\begin{Resp}\Rsign Propter nímiam caritátem suam qua diléxit nos.\end{Resp}
\switchcolumn
\EN
\begin{Resp}\Rsign And we, being dead in our sins, hath God quickened together with Christ, \Star{} For his great love wherewith he hath loved us.\end{Resp}
\begin{Resp}\Vsign That in the ages to come he might shew the exceeding riches of his grace.\end{Resp}
\begin{Resp}\Rsign For his great love wherewith he hath loved us.\end{Resp}
\begin{Resp}\Vsign Glory be to the Father, and to the Son, \Star{} and to the Holy Ghost.\end{Resp}
\begin{Resp}\Rsign For his great love wherewith he hath loved us.\end{Resp}
\switchcolumn*

\LA
\LessonHead{Lectio IV}
\switchcolumn
\EN
\LessonHead{Lesson IV}
\switchcolumn*

\LA
\Source{Ex lítteris Encyclicis Pii Papæ undécimi}
\switchcolumn
\EN
\Source{From the Encyclical Letter of Pope Pius XI}
\switchcolumn*

\LA
\noindent Quantópere autem hujúsmodi expiatiónis seu reparatiónis necéssitas hac nostra potíssimum ætáte úrgeat, némini non maniféstum erit, qui, ut inítio díximus, hunc mundum in malígno pósitum, óculis animóque perlustráverit. Undique enim geméntium ad Nos populórum clamor adscéndit, quorum príncipes vel rectóres vere adstitérunt et convenérunt in unum advérsus Dóminum et advérsus Ecclésiam ejus. At étiam doléndum, quod inter ipsos fidéles, sánguine Agni immaculáti in baptísmo ablútos, gratiáque locupletátos, tot inveniántur cujúsvis órdinis hómines, qui incredíbili rerum divinárum ignorántia laborántes et falsis doctrínis infécti, vítiis irretítam, procul a domo Patris, vitam tradúcant, quam nec veræ fídei lumen collústrat, nec spes futúræ beatitátis deléctat, nec ardor réficit fovétque caritátis, ut sedére in ténebris et in umbra mortis vere videántur.\par
\switchcolumn
\EN
\noindent Anyone who will use his eyes and mind, if he but think of this world, whereof it is truly said: The whole world lieth in wickedness: can see how urgent, especially in these our own times, is the need for expiation or atonement. For there come to our ears from every side the cries of nations, whose rulers or governments have actually risen up, and conspired together, against the Lord, and against his Church. Nor is that other sight less sad, to wit, that even among the faithful, (washed as they are by Baptism in the blood of the Lamb without spot, and enriched by his grace,) we find so many of every station in life who are ignorant of divine things, and poisoned by false doctrine; and who do live a sinful life, far from their Father's house; without the light of the true Faith; without the joy of a hope in the future life; deprived of the strength and comfort which come with the Spirit of Love; so that it may be said of them quite truthfully: They sit in darkness and in the shadow of death.\par
\switchcolumn*

\LA
\begin{Resp}\Rsign Prope est Dóminus ómnibus invocántibus eum, \Star{} Omnibus invocántibus eum in veritáte.\end{Resp}
\begin{Resp}\Vsign Miserátor et miséricors Dóminus, pátiens et multum miséricors.\end{Resp}
\begin{Resp}\Rsign Omnibus invocántibus eum in veritáte.\end{Resp}
\switchcolumn
\EN
\begin{Resp}\Rsign The Lord is nigh unto all them that call upon him, \Star{} Yea, unto all such as call upon him faithfully.\end{Resp}
\begin{Resp}\Vsign The Lord is full of compassion and mercy, long-suffering and of great goodness.\end{Resp}
\begin{Resp}\Rsign Yea, unto all such as call upon him faithfully.\end{Resp}
\switchcolumn*

\LA
\LessonHead{Lectio V}
\switchcolumn
\EN
\LessonHead{Lesson V}
\switchcolumn*

\LA
\noindent Hisce vero malis véluti in cúmulum accédit cum eórum ignávia atque socórdia, qui, dormitántium et fugiéntium instar discipulórum, nutántes in fide, Christum angóribus oppréssum vel Sátanæ satellítibus circumvéntum mísere derelínquunt, tum eórum perfídia, qui, Judæ proditóris exémplum secúti, aut témere et sacrílege de altári libant, aut ad hóstium castra transfúgiunt. Atque ita vel invítum subit cogitátio ánimum, jam própius adventáre témpora de quibus Dóminus Noster vaticinátus est: Et quóniam abundávit iníquitas, refrigéscet cáritas multórum. Quæ quidem ómnia quotquot pie commentáti erunt fidéles, fácere non póterunt, quin, Christi perdoléntis incénsi caritáte, vehementióre stúdio suas aliorúmque culpas éxpient, Christi honórem resárciant, æternámque próvehant animárum salútem.\par
\switchcolumn
\EN
\noindent To the aforesaid accumulation of evils must be added the sloth and indifference of many of his followers, who are like unto the disciples that fled from him or slept. These are such as are not firmly rooted in the Faith, and have therefore shamefully abandoned Christ at a time when he is burdened with sorros and attacked by the hosts of Satan. And other of his followers are like unto Judas the traitor, in that they walk in the footsteps of perfidy. These are such as approach the Sacrament of the Altar with sacrilegious boldness, or even go over to the camp of the enemy. And we are therefore moved to think that now is come the hour of which our Lord prophesied: Because iniquity shall abound, the love of the many shall grow cold. If those who remain faithful on fire with love at the sufferings of Christ will but meditate on these considerations, it is unthinkable that they will do other than strive with greater zeal to expiate both their own faults and the faults of others; and that they will thus seek to make reparation for the dishonour done to Christ; and so they will be filled with zeal for the eternal salvation of souls.\par
\switchcolumn*

\LA
\begin{Resp}\Rsign Confíteor tibi, Pater, Dómine cæli et terræ, quia abscondísti hæc a sapiéntibus et prudéntibus \Star{} Et revelásti ea párvulis.\end{Resp}
\begin{Resp}\Vsign Ita, Pater, quóniam sic fuit plácitum ante te.\end{Resp}
\begin{Resp}\Rsign Et revelásti ea párvulis.\end{Resp}
\switchcolumn
\EN
\begin{Resp}\Rsign I thank thee, O Father, Lord of heaven and earth, because thou hast hid these things from the wise and prudent; \Star{} Yea, thou hast revealed them unto babes.\end{Resp}
\begin{Resp}\Vsign Even so, Father, for so it seemed good in thy sight.\end{Resp}
\begin{Resp}\Rsign Yea, thou hast revealed them unto babes.\end{Resp}
\switchcolumn*

\LA
\LessonHead{Lectio VI}
\switchcolumn
\EN
\LessonHead{Lesson VI}
\switchcolumn*

\LA
\noindent Et sane illud Apóstoli: Ubi abundávit delíctum, superabundávit grátia, áliquo pacto ad hanc quoque ætátem nostram describéndam accommodáre licet: nam, aucta ádmodum perversitáte hóminum, mirífice item, Spíritu Sancto afflánte, númerus fidélium utriúsque sexus augétur, qui alacrióre ánimo pro tot illátis injúriis divíno Cordi satisfácere student, immo étiam se ipsos Christo víctimas offérre non dúbitant. Etenim quæ usque adhuc memorávimus si quis secum ánimo réputet amánter eadémque véluti in medúllis defíxa hábeat, fíeri profécto non potest quin is non tam ab omni peccáto tamquam summo malo abhórreat atque abstíneat, quam se totum Dei voluntáti permíttat, et læsum divínæ Majestátis honórem, cum continénter orándo, tum afflictatiónibus sponte suscéptis ærumnísque, si quæ incíderint, patiénter tolerátis, tum tota demum vita hoc expiatiónis stúdio exigénda, resarcíre conténdat.\par
\switchcolumn
\EN
\noindent Surely we shall do well to apply to this our own age what the Apostle wrote: Where sin abounded, grace did much more abound. For even though the sinfulness of man doth greatly increase, by the grace of the Holy Spirit there doth also increase in number those of both sexes who most cheerfully do endeavour to make satisfaction to the divine Heart for the numerous injuries heaped thereupon; may more, they even cheerfully offer themselves as victims for sin. Verily, anyone that considereth in a spirit of love all the revelation which hath come into his mind up to this time, if he have impressed such things on the fleshly tablets of his own heart, as it were, cannot but abhor and flee sin as the greatest of all evils. Such an one will therefore offer himself wholly and completely to the will of God, and by constant prayer, willing penances, and the patient endurance of all the ills that befall him will endeavour to repair the injuries done to the Majesty of God. In a word, he will so organize his life that all things in it may be motivated by the spirit of reparation.\par
\switchcolumn*

\LA
\begin{Resp}\Rsign Omnes gentes quascúmque fecísti, vénient \Star{} Et adorábunt coram te, Dómine.\end{Resp}
\begin{Resp}\Vsign Et glorificábunt nomen tuum, quóniam magnus es tu, et fáciens mirabília.\end{Resp}
\begin{Resp}\Rsign Et adorábunt coram te, Dómine.\end{Resp}
\begin{Resp}\Vsign Glória Patri, et Fílio, \Star{} et Spirítui Sancto.\end{Resp}
\begin{Resp}\Rsign Et adorábunt coram te, Dómine.\end{Resp}
\switchcolumn
\EN
\begin{Resp}\Rsign All nations whom thou hast made shall come, \Star{} And they shall worship thee, O Lord.\end{Resp}
\begin{Resp}\Vsign Yea, they shall glorify thy Name, for thou art great, and doest wondrous things.\end{Resp}
\begin{Resp}\Rsign And they shall worship thee, O Lord.\end{Resp}
\begin{Resp}\Vsign Glory be to the Father, and to the Son, \Star{} and to the Holy Ghost.\end{Resp}
\begin{Resp}\Rsign And they shall worship thee, O Lord.\end{Resp}
\switchcolumn*

\LA
\LessonHead{Lectio VII}
\switchcolumn
\EN
\LessonHead{Lesson VII}
\switchcolumn*

\LA
\LessonTitle{Léctio sancti Evangélii secúndum Joánnem}
\switchcolumn
\EN
\LessonTitle{From the Holy Gospel according to John}
\switchcolumn*

\LA
\Cite{Joannes 19:31-37}
\switchcolumn
\EN
\Cite{John 19:31-37}
\switchcolumn*

\LA
\noindent In illo témpore: Judǽi, quóniam parascéve erat, ut non remanérent in cruce córpora sábbato (erat enim magnus dies ille sábbati) rogavérunt Pilátum, ut frangeréntur eórum crura et tolleréntur. Et réliqua.\par
\switchcolumn
\EN
\noindent At that time: The Jews, because it was the Preparation, that the bodies should not remain upon the cross on the Sabbath Day, for that Sabbath Day was an high day, besought Pilate that their legs might broken, and that they might be taken away. And so on.\par
\switchcolumn*

\LA
\LessonTitle{Homilía sancti Cyrílli Alexandríni}
\switchcolumn
\EN
\LessonTitle{A Homily by St. Cyril of Alexandria}
\switchcolumn*

\LA
\Source{Comm. in Joann. Lib. 12 Cap. 19}
\switchcolumn
\EN
\Source{Comm. in Joann. Lib. 12 Cap. 19}
\switchcolumn*

\LA
\noindent Non hæc ait beátus Evangelísta, quasi pietátem ullam efferátis et crudélibus Judǽis tríbuat, sed ut osténdat stulte eos et imperíte excoláre cúlicem, et camélum deglutíre sicut a Christo dictum est. Gravíssima enim et immánia scélera pro níhilo dúcere comperiúntur: mínima vero et exília caute ádmodum ac sollícite obsérvant, suam utrobíque pandéntes inscítiam. Quod in promptu est osténdere. Ecce enim Christo interfécto sábbati honórem magni fáciunt, et incredíbili audácia, violáto legis auctóre, pietátem erga legem præ se ferunt.\par
\switchcolumn
\EN
\noindent The blessed Evangelist doth not record these things as if they were evidence of some godliness amidst all the savage cruelty which Jewry had manifested. Rather, by them he doth shew how foolishly and ignorantly a gnat can be carefully strained out, and a camel swallowed, as Christ himself had foresaid. For we perceive that they made light of most grievous and awful crimes, the while they took careful and anxious pains to observe mere trivialities; in the which consistency they did in both things but display their ignorance, as can be readily shewn. For lo! they put Christ to death, the while they give honour to the Great Sabbath. Thus, with incredible insolence, the make a shew of reverence for the Law concerning the Sabbath, whose very Author they in so doing do dishonour.\par
\switchcolumn*

\LA
\begin{Resp}\Rsign Ego si exaltátus fúero a terra \Star{} Omnia traham ad meípsum.\end{Resp}
\begin{Resp}\Vsign Hoc autem dicébat signíficans qua morte esset moritúrus.\end{Resp}
\begin{Resp}\Rsign Omnia traham ad meípsum.\end{Resp}
\switchcolumn
\EN
\begin{Resp}\Rsign If I be lifted up, \Star{} I will draw all men unto me.\end{Resp}
\begin{Resp}\Vsign This he said, signifying what death he should die.\end{Resp}
\begin{Resp}\Rsign I will draw all men unto me.\end{Resp}
\switchcolumn*

\LA
\LessonHead{Lectio VIII}
\switchcolumn
\EN
\LessonHead{Lesson VIII}
\switchcolumn*

\LA
\noindent Magnum autem illíus præcípue sábbati diem cólere símulant, qui magnæ diéi Dóminum interemérunt, et grátiam se solis dignam flágitant, ut eórum nempe crura frangántur, intolerábili dolóre acerbiórem morte ipsa perníciem prope jam mórtuis moliéntes. Venérunt ergo mílites, et primi quidem fregérunt crura, et altérius qui crucifíxus est cum eo. Judæórum petitióni mílites obsecúti, qui símili crudelitátis furóre laborábant, duórum quidem latrónum, quippe qui vivi adhuc repérti essent, crura confríngunt. Sed, cum Jesum inclinásse caput comperíssent, et jam exspirásse putárent, frustra ejus crura confríngi exístimant, sed, cum adhuc mórtuum esse nonníhil diffidérent, láncea latus ejus perfódiunt, unde cruor aqua mixtus scatúriit, quod eulógiæ mýsticæ et sancti baptísmatis imágo quædam erat atque primítiæ.\par
\switchcolumn
\EN
\noindent The pretend to shew especial reverence for that Great Sabbath Day―they who have put to death the Lord of that great day! And thereupon they earnestly beseech a favour worthy of such folk alone, namely, That the legs of the crucified ones may be broken! And thereby was inflicted an intoleráble pain, a bitterer misfortune than death itself, on men who were already well nigh at the point of death! The Evangelist saith: Then came the soldiers and brake the legs of the first, and of the other which was crucified with him. Thus it doth appear that the soldiers who went in answer to this request were labouring under a frenzy of cruelty like to that of Jewry itself, and therefore they brake the legs of the two thieves whom they found yet alive. But when they came to Jesus, whose head was bowed, they concluded that he was dead already, so that it was useless to go to the trouble of breaking his legs. But to make quite certain that he was dead indeed, they did pierce his side with a spear. And forthwith there came out blood and water: a figure of the mystical Banquet and of holy Baptism, and also the first-fruits of the same.\par
\switchcolumn*

\LA
\begin{Resp}\Rsign Simus ergo imitatóres Dei \Star{} Et ambulémus in dilectióne.\end{Resp}
\begin{Resp}\Vsign Sicut et Christus diléxit nos et trádidit semetípsum pro nobis.\end{Resp}
\begin{Resp}\Rsign Et ambulémus in dilectióne.\end{Resp}
\begin{Resp}\Vsign Glória Patri, et Fílio, \Star{} et Spirítui Sancto.\end{Resp}
\begin{Resp}\Rsign Et ambulémus in dilectióne.\end{Resp}
\switchcolumn
\EN
\begin{Resp}\Rsign Be ye therefore followers of God; \Star{} And walk ye therefore in love.\end{Resp}
\begin{Resp}\Vsign For Christ also hath loved us and hath given himself for us.\end{Resp}
\begin{Resp}\Rsign And walk ye therefore in love.\end{Resp}
\begin{Resp}\Vsign Glory be to the Father, and to the Son, \Star{} and to the Holy Ghost.\end{Resp}
\begin{Resp}\Rsign And walk ye therefore in love.\end{Resp}
\switchcolumn*

\LA
\LessonHead{Lectio IX}
\switchcolumn
\EN
\LessonHead{Lesson IX}
\switchcolumn*

\LA
\noindent Ex iis autem quæ contigérunt, sapientíssimus Evangelísta confírmat auditóribus eum esse Christum, qui per sanctas olim prædíctus est Scriptúras; consentánea enim divínis de eo oráculis evenérunt. nec enim os ejus confráctum est, et mílitum láncea transfíxus est, secúndum Scriptúras. Spectatórem vero ejus rei ait exstitísse discípulum ipsum qui testimónium pérhibet de his, et scire se vera testári, seípsum his verbis, non álium dénotans.\par
\switchcolumn
\EN
\noindent And from these things which came to pass, the most wise Evangelist proveth to his hearers that this is the Christ, who was long ago foretold by holy Scripture; for all that happened doth agree with what was prophesied through God concerning him. And so, according to the Scriptures, not a bone of him was broken, and he was pierced with the soldier's lance. The Evangelist doth certainly say that the one who saw this thing was the same disciple that gave testimony concerning these things, and that he knew that his testimony was true, by this expression meaning not another, but himself.\par
\switchcolumn*

\LA
\TeDeum{Te Deum}
\switchcolumn
\EN
\TeDeum{Te Deum}
\switchcolumn*

\end{paracol}

\Office{In Octava SSmi Cordis Jesu}{Octave Day of the Sacred Heart}{Duplex majus}
\addcontentsline{toc}{subsection}{In Octava SSmi Cordis Jesu \textit{— Octave Day of the Sacred Heart}}
\begin{paracol}{2}
\LA
\LessonHead{Lectio I}
\switchcolumn
\EN
\LessonHead{Lesson I}
\switchcolumn*

\LA
\LessonTitle{De libro primo Regum}
\switchcolumn
\EN
\LessonTitle{Lesson from the first book of Samuel}
\switchcolumn*

\LA
\Cite{1 Reg 15:1-3}
\switchcolumn
\EN
\Cite{1 Sam 15:1-3}
\switchcolumn*

\LA
\noindent Et dixit Samuel ad Saul: Me misit Dóminus, ut úngerem te in regem super pópulum ejus Israël; nunc ergo audi vocem Dómini. Hæc dicit Dóminus exercítuum: Recénsui quæcúmque fecit Amalec Israéli, quómodo réstitit ei in via, cum ascénderet de Ægýpto. Nunc ergo vade et pércute Amalec et demolíre univérsa ejus: non parcas ei et non concupíscas ex rebus ipsíus áliquid; sed intérfice a viro usque ad mulíerem et párvulum atque lacténtem, bovem et ovem, camélum et ásinum.\par
\switchcolumn
\EN
\noindent And Samuel said to Saul: The Lord sent me to anoint thee king over his People Israel: now therefore hearken thou unto the voice of the Lord: Thus saith the Lord of hosts: I have reckoned up all that Amalec hath done to Israel: I how he opposed them in the way when they came up out of Egypt. Now therefore go, and smite Amalec, and utterly destroy all that he hath: spare him not, nor covet any thing that is his: but slay both man and woman, child and suckling, ox and sheep, camel and ass.\par
\switchcolumn*

\LA
\begin{Resp}\Rsign Fériam eis pactum sempitérnum et non désinam eis benefácere et timórem meum dabo in corde eórum \Star{} Ut non recédant a me.\end{Resp}
\begin{Resp}\Vsign Et lætábor super eis cum bene eis fécero in toto Corde meo.\end{Resp}
\begin{Resp}\Rsign Ut non recédant a me.\end{Resp}
\switchcolumn
\EN
\begin{Resp}\Rsign I will make an everlasting Covenant with them, and I will not cease from doing them good, and I will put my fear in their hearts, \Star{} So that they shall not depart from me.\end{Resp}
\begin{Resp}\Vsign Yea, I will rejoice over them, to do them good with my whole Heart.\end{Resp}
\begin{Resp}\Rsign So that they shall not depart from me.\end{Resp}
\switchcolumn*

\LA
\LessonHead{Lectio II}
\switchcolumn
\EN
\LessonHead{Lesson II}
\switchcolumn*

\LA
\Cite{1 Reg 15:4-8}
\switchcolumn
\EN
\Cite{1 Sam 15:4-8}
\switchcolumn*

\LA
\noindent Præcépit ítaque Saul pópulo, et recénsuit eos quasi agnos: ducénta míllia péditum et decem míllia virórum Juda. Cumque venísset Saul usque ad civitátem Amalec, teténdit insídias in torrénte. Dixítque Saul Cinǽo: Abíte, recédite atque descéndite ab Amalec, ne forte invólvam te cum eo; tu enim fecísti misericórdiam cum ómnibus fíliis Israël cum ascénderent de Ægýpto. Et recéssit Cinǽus de médio Amalec. Percussítque Saul Amalec ab Hévila donec vénias ad Sur, quæ est e regióne Ægýpti. Et apprehéndit Agag regem Amalec vivum; omne autem vulgus interfécit in ore gládii.\par
\switchcolumn
\EN
\noindent So Saul commanded the people, and numbered them as lambs: two hundred thousand footmen, and ten thousand of the men of Juda. And when Saul was come to the city of Amalec, he laid ambushes in the torrent. And Saul said to the Cinite: Go, depart and get ye down from Amalec: lest I destroy thee with him. For thou hast shown kindness to all the children of Israel, when they came up out of Egypt. And the Cinite departed from the midst of Amalec. And Saul smote Amalec from Hevila, until thou comest to Sur, which is over against Egypt. And he took Agag the king of Amalec alive: but all the common people he slew with the edge of the sword\par
\switchcolumn*

\LA
\begin{Resp}\Rsign Si inimícus meus maledixísset mihi, sustinuíssem útique \Star{} Tu vero homo unánimis qui simul mecum dulces capiébas cibos.\end{Resp}
\begin{Resp}\Vsign Et si is qui me óderat super me magna locútus fuísset, abscondíssem me fórsitan ab eo.\end{Resp}
\begin{Resp}\Rsign Tu vero homo unánimis qui simul mecum dulces capiébas cibos.\end{Resp}
\switchcolumn
\EN
\begin{Resp}\Rsign It was not an open enemy that has done me this dishonour, for then I could have borne it, \Star{} But it was even thou, mine own familiar friend, who did also eat of my Bread.\end{Resp}
\begin{Resp}\Vsign Neither was it mine adversary that did magnify himself against me, for then peradventure I would have hid myself from him.\end{Resp}
\begin{Resp}\Rsign But it was even thou, mine own familiar friend, who did also eat of my Bread.\end{Resp}
\switchcolumn*

\LA
\LessonHead{Lectio III}
\switchcolumn
\EN
\LessonHead{Lesson III}
\switchcolumn*

\LA
\Cite{1 Reg 15:9-11}
\switchcolumn
\EN
\Cite{1 Sam 15:9-11}
\switchcolumn*

\LA
\noindent Et pepércit Saul et pópulus Agag et óptimis grégibus óvium et armentórum et véstibus et ariétibus et univérsis, quæ pulchra erant, nec voluérunt dispérdere ea; quidquid vero vile fuit et réprobum, hoc demolíti sunt. Factum est autem verbum Dómini ad Sámuel, dicens: Pœ́nitet me quod constitúerim Saul regem, quia derelíquit me et verba mea ópere non implévit. Contristatúsque est Sámuel et clamávit ad Dóminum tota nocte.\par
\switchcolumn
\EN
\noindent And Saul and the people spared Agag and the best of the flocks of sheep and of the herds, and the garments and the rams, and all that was beautiful, and would not destroy them: but every thing that was vile and good for nothing, that they destroyed. And the word of the Lord came to Samuel, saying: It repenteth me that I have made Saul king: for he hath forsaken me, and hath not executed my commandments. And Samuel was grieved, and he cried unto the Lord all night.\par
\switchcolumn*

\LA
\begin{Resp}\Rsign Cum essémus mórtui peccátis, convivificávit nos Deus in Christo \Star{} Propter nímiam caritátem suam qua diléxit nos.\end{Resp}
\begin{Resp}\Vsign Ut osténderet in sǽculis superveniéntibus abundántes divítias grátiæ suæ.\end{Resp}
\begin{Resp}\Rsign Propter nímiam caritátem suam qua diléxit nos.\end{Resp}
\begin{Resp}\Vsign Glória Patri, et Fílio, \Star{} et Spirítui Sancto.\end{Resp}
\begin{Resp}\Rsign Propter nímiam caritátem suam qua diléxit nos.\end{Resp}
\switchcolumn
\EN
\begin{Resp}\Rsign And we, being dead in our sins, hath God quickened together with Christ, \Star{} For his great love wherewith he hath loved us.\end{Resp}
\begin{Resp}\Vsign That in the ages to come he might shew the exceeding riches of his grace.\end{Resp}
\begin{Resp}\Rsign For his great love wherewith he hath loved us.\end{Resp}
\begin{Resp}\Vsign Glory be to the Father, and to the Son, \Star{} and to the Holy Ghost.\end{Resp}
\begin{Resp}\Rsign For his great love wherewith he hath loved us.\end{Resp}
\switchcolumn*

\LA
\LessonHead{Lectio IV}
\switchcolumn
\EN
\LessonHead{Lesson IV}
\switchcolumn*

\LA
\LessonTitle{Sermo sancti Bernárdi Abbátis}
\switchcolumn
\EN
\LessonTitle{From a Sermon by St. Bernard the Abbot}
\switchcolumn*

\LA
\Source{Sermo 61 in Cantica Canticorum, nn. 3-5}
\switchcolumn
\EN
\Source{Sermo 61 in Cantica Canticorum, nn. 3-5}
\switchcolumn*

\LA
\noindent Revéra ubi tuta firmáque infírmis secúritas et réquies, nisi in vulnéribus Salvatóris? Tanto illic secúrior hábito, quanto ille poténtior est ad salvándum. Fremit mundus, premit corpus, diábolus insidiátur: non cado; fundátus enim sum supra firmam petram. Peccávi peccátum grande: turbábitur consciéntia, sed non perturbábitur, quóniam vúlnerum Dómini recordábor. Nempe vulnerátus est propter iniquitátes nostras. Quid tam ad mortem, quod non Christi morte solvátur? Si ergo in mentem vénerit tam potens tamque éfficax medicaméntum, nulla jam possum morbi malignitáte terréri.\par
\switchcolumn
\EN
\noindent For us who are so frail and weak, where is to be found a sure and certain place of abiding safety, or of everlasting rest? Where except in the Wounds of the Saviour? There alone I may dwell safely. There alone I may find a safety as great as his mighty power to save. The world may rage around me; the body may weigh me down; the devil may lay snares for me; but if I hide me there I cannot fall, for I am founded on the firm Rock. If I have committed a great sin; if my conscience is sore troubled; I will not despair, for I have always in remembrance the Wounds of the Lord. For in all truth: He was wounded for our transgressions. And there is no sin so deadly that it cannot be remítted through Christ's death. If then I keep in remembrance a remedy so powerful and efficacious, I cannot in this present life be terrified by any evil, no matter how malignant.\par
\switchcolumn*

\LA
\begin{Resp}\Rsign Prope est Dóminus ómnibus invocántibus eum, \Star{} Omnibus invocántibus eum in veritáte.\end{Resp}
\begin{Resp}\Vsign Miserátor et miséricors Dóminus, pátiens et multum miséricors.\end{Resp}
\begin{Resp}\Rsign Omnibus invocántibus eum in veritáte.\end{Resp}
\switchcolumn
\EN
\begin{Resp}\Rsign The Lord is nigh unto all them that call upon him, \Star{} Yea, unto all such as call upon him faithfully.\end{Resp}
\begin{Resp}\Vsign The Lord is full of compassion and mercy, long-suffering and of great goodness.\end{Resp}
\begin{Resp}\Rsign Yea, unto all such as call upon him faithfully.\end{Resp}
\switchcolumn*

\LA
\LessonHead{Lectio V}
\switchcolumn
\EN
\LessonHead{Lesson V}
\switchcolumn*

\LA
\noindent Ego vero fidénter quod ex me mihi deest, usúrpo mihi ex viscéribus Dómini, quóniam misericórdiæ áffluunt; nec desunt forámina, per quæ éffluant. Fodérunt manus ejus et pedes, latúsque láncea foravérunt; et per has rimas licet mihi súgere mel de petra, oleúmque de saxo duríssimo; id est gustáre et vidére quóniam suávis est Dóminus. Cogitábat cogitatiónes pacis, et ego nesciébam. Quis enim cognóvit sensum Dómini? aut quis consiliárius ejus fuit? At clavis réserans, clavus pénetrans factus est mihi ut vídeam voluntátem Dómini. Quidni vídeam per forámen? Clamat clavus, clamat vulnus, quod vere Deus sit in Christo mundum reconcílians sibi.\par
\switchcolumn
\EN
\noindent But as for me, since mercies thus abound, I take unto myself whatever is lacking in me; yea, I take it unto myself with confidence; I take it unto myself from the compassion of the Lord, with whom every kind of mercy aboundeth. For openings are not wanting, through which these mercies may flow forth. They have pierced his hands and his feet; ye, and his side too they have pierced with a spear. It is though these clefts that I am permitted to suck honey from the Rock, and oil out of the flinty Rock, and to taste and see how gracious the Lord is. His thoughts are thoughts of peace, and I knew it not. For who hath known the mind of the Lord? or who hath been his counsellor? But the nail which did pierce hath become unto me a key which doth unlock, so that the will of the Lord is set open unto me. How could I, with such an opening, do other than to see his will? For the nails cry aloud, and the Wounds speak, saying the truth: To wit, that God is in Christ reconciling the world unto himself.\par
\switchcolumn*

\LA
\begin{Resp}\Rsign Confíteor tibi, Pater, Dómine cæli et terræ, quia abscondísti hæc a sapiéntibus et prudéntibus \Star{} Et revelásti ea párvulis.\end{Resp}
\begin{Resp}\Vsign Ita, Pater, quóniam sic fuit plácitum ante te.\end{Resp}
\begin{Resp}\Rsign Et revelásti ea párvulis.\end{Resp}
\switchcolumn
\EN
\begin{Resp}\Rsign I thank thee, O Father, Lord of heaven and earth, because thou hast hid these things from the wise and prudent; \Star{} Yea, thou hast revealed them unto babes.\end{Resp}
\begin{Resp}\Vsign Even so, Father, for so it seemed good in thy sight.\end{Resp}
\begin{Resp}\Rsign Yea, thou hast revealed them unto babes.\end{Resp}
\switchcolumn*

\LA
\LessonHead{Lectio VI}
\switchcolumn
\EN
\LessonHead{Lesson VI}
\switchcolumn*

\LA
\noindent Ferrum pertránsiit ánimam ejus, et appropinquávit Cor illíus, ut non jam non sciat cómpati infirmitátibus meis. Patet arcánum Cordis per forámina córporis: patet magnum illud pietátis sacraméntum, patent víscera misericórdiæ Dei nostri, in quibus visitávit nos Oriens ex alto. Quidni víscera per vúlnera páteant? In quo enim clárius quam in vulnéribus tuis eluxísset, quod tu, Dómine, suávis et mitis, et multæ misericórdiæ? Majórem enim miseratiónem nemo habet, quam ut ánimam suam ponat quis pro addíctis morti et damnátis. Meum proínde méritum, miserátio Dómini. Non plane sum mériti inops, quámdiu ille miseratiónum non fúerit. Et si misericórdiæ Dómini ab ætérno et usque in ætérnum, ego quoque misericórdias Dómini in ætérnum cantábo.\par
\switchcolumn
\EN
\noindent The iron entered into his soul, and came nigh unto his Heart, that he might truly know compassion for mine infirmities. The secrets of his Heart lie open, through the Wounds of his body. Thus is that great mystery of love laid open, there lie open the bowels of the mercy of our God, whereby the Day-Spring from on high hath visited us. But why should not the bowels of mercy lie open through the wounds? For in what hath it appeared more clearly than in thy wounds, that thou, O Lord, art sweet and gentle, and of great mercy? For greater pity hath no man, than that a man lay down his life for those who were doomed and condemned to death. From this pity of the Lord is all my merit. I am not entirely destitute of merit, so long as he is not wanting in compassion. And if the mercies of the Lord are from eternity unto eternity, I also will sing the mercies of the Lord forever.\par
\switchcolumn*

\LA
\begin{Resp}\Rsign Omnes gentes quascúmque fecísti, vénient \Star{} Et adorábunt coram te, Dómine.\end{Resp}
\begin{Resp}\Vsign Et glorificábunt nomen tuum, quóniam magnus es tu, et fáciens mirabília.\end{Resp}
\begin{Resp}\Rsign Et adorábunt coram te, Dómine.\end{Resp}
\begin{Resp}\Vsign Glória Patri, et Fílio, \Star{} et Spirítui Sancto.\end{Resp}
\begin{Resp}\Rsign Et adorábunt coram te, Dómine.\end{Resp}
\switchcolumn
\EN
\begin{Resp}\Rsign All nations whom thou hast made shall come, \Star{} And they shall worship thee, O Lord.\end{Resp}
\begin{Resp}\Vsign Yea, they shall glorify thy Name, for thou art great, and doest wondrous things.\end{Resp}
\begin{Resp}\Rsign And they shall worship thee, O Lord.\end{Resp}
\begin{Resp}\Vsign Glory be to the Father, and to the Son, \Star{} and to the Holy Ghost.\end{Resp}
\begin{Resp}\Rsign And they shall worship thee, O Lord.\end{Resp}
\switchcolumn*

\LA
\LessonHead{Lectio VII}
\switchcolumn
\EN
\LessonHead{Lesson VII}
\switchcolumn*

\LA
\LessonTitle{Léctio sancti Evangélii secúndum Joánnem}
\switchcolumn
\EN
\LessonTitle{From the Holy Gospel according to John}
\switchcolumn*

\LA
\Cite{Joannes 19:31-37}
\switchcolumn
\EN
\Cite{John 19:31-37}
\switchcolumn*

\LA
\noindent In illo témpore: Judæi, quóniam parascéve erat, ut non remanérent in cruce córpora sábbato (erat enim magnus dies ille sábbati) rogavérunt Pilátum, ut frangeréntur eórum crura et tolleréntur. Et réliqua.\par
\switchcolumn
\EN
\noindent At that time: The Jews, because it was the Preparation, that the bodies should not remain upon the cross on the Sabbath Day, for that Sabbath Day was an high day, besought Pilate that their legs might broken, and that they might be taken away. And so on.\par
\switchcolumn*

\LA
\LessonTitle{Homilía sancti Augustíni Epíscopi}
\switchcolumn
\EN
\LessonTitle{A Homily by St. Augustine the Bishop}
\switchcolumn*

\LA
\Source{Tractatus 120 in Joánnem, nn. 2-3}
\switchcolumn
\EN
\Source{Tractatus 120 in Joannem, nn. 2-3}
\switchcolumn*

\LA
\noindent Ad Jesum autem cum veníssent, ut vidérunt eum jam mórtuum, non fregérunt ejus crura, sed unus mílitum láncea latus ejus apéruit, et contínuo exívit sanguis et aqua. Vigilánti verbo Evangelísta usus est, ut non díceret: latus ejus percússit, aut vulnerávit, aut quid áliud, sed apéruit, ut illic quodámmodo vitæ óstium panderétur, unde Sacraménta Ecclésiæ manavérunt, sine quibus ad vitam, quæ vera vita est, non intrátur. Ille sanguis in remissiónem fusus est peccatórum. Aqua illa salutáre témperat póculum, hæc et lavácrum præstat et potum. Hoc prænuntiábat quod Noë in látere arcæ óstium fácere jussus est, qua intrárent animália quæ non erant dilúvio peritúra, quibus præfigurabátur Ecclésia.\par
\switchcolumn
\EN
\noindent But when they came to Jesus, and saw that he was dead already, they brake not his legs: but one of the soldiers with a spear opened his side, and forthwith came there out blood and water. Note that the Evangelist maketh use of a word of special significance. He saith not: Penetrated his side: nor yet: Wounded: nor any other thing; but rather: Opened: that thereby in a sense the door of life might be thrown open, from whence the Sacraments of the Church have flowed forth, without which there is no entrance into the life which is the only true life. For that blood was shed for the remission of sins; and that water hath brought into being the life-giving flagon, which same is both the laver of Baptism and the cup that giveth refreshment to them that thirst. All this was announced long before, to wit, when Noah was commanded to make a door in the side of the ark, through which might enter all living creatures which were not destined to perish in the flood; and this same is a figure of the Church.\par
\switchcolumn*

\LA
\begin{Resp}\Rsign Ego si exaltátus fúero a terra \Star{} Omnia traham ad meípsum.\end{Resp}
\begin{Resp}\Vsign Hoc autem dicébat signíficans qua morte esset moritúrus.\end{Resp}
\begin{Resp}\Rsign Omnia traham ad meípsum.\end{Resp}
\switchcolumn
\EN
\begin{Resp}\Rsign If I be lifted up, \Star{} I will draw all men unto me.\end{Resp}
\begin{Resp}\Vsign This he said, signifying what death he should die.\end{Resp}
\begin{Resp}\Rsign I will draw all men unto me.\end{Resp}
\switchcolumn*

\LA
\LessonHead{Lectio VIII}
\switchcolumn
\EN
\LessonHead{Lesson VIII}
\switchcolumn*

\LA
\noindent Propter hoc prima múlier facta est de látere viri dormiéntis, et appelláta est vita matérque vivórum. Magnum quippe significávit bonum, ante magnum prævaricatiónis malum. Hic secúndus Adam, inclináto cápite, in cruce dormívit, ut inde formarétur ejus conjux, quæ de látere dormiéntis efflúxit. O mors unde mórtui revivíscunt! Quid isto sánguine múndius? Quid vúlnere isto salúbrius? Et qui vidit, inquit, testimónium perhíbuit, et verum est testimónium ejus; et ille scit quia vera dicit, ut et vos credátis. Non dixit: ut et vos sciátis, sed ut credátis. Scit enim qui vidit, cujus credat testimónio qui non vidit. Magis autem ad fidem pértinet crédere quam vidére.\par
\switchcolumn
\EN
\noindent Another figure is also to be found in the first woman. For she was made out of the side of the first man whilst he slept a deep sleep; and she was called: Life: or as it may be interpreted: The Mother of all living. Thus was indicated the great good which was later to come to pass, even before the great evil of transgression had come into being. Here, in the Gospel, the second Adam is shewn as bowing his head, and sleeping his deep sleep upon the Cross, that a bride might be formed for him out of that which came forth from his side as he slept. What a death, whereby the dead are raised anew to life! How clean and cleansing is this blood! What is more healthgiving than this wound! And he that saw it, saith he, bare record, and his record is true: and he knoweth that he saith true, that ye might believe. He said not: That ye might know; but: That ye might believe. For he knoweth, who hath seen, that he, who hath not seen, might believe his record. And believing belongeth more to the nature of faith than doth seeing.\par
\switchcolumn*

\LA
\begin{Resp}\Rsign Simus ergo imitatóres Dei \Star{} Et ambulémus in dilectióne.\end{Resp}
\begin{Resp}\Vsign Sicut et Christus diléxit nos et trádidit semetípsum pro nobis.\end{Resp}
\begin{Resp}\Rsign Et ambulémus in dilectióne.\end{Resp}
\begin{Resp}\Vsign Glória Patri, et Fílio, \Star{} et Spirítui Sancto.\end{Resp}
\begin{Resp}\Rsign Et ambulémus in dilectióne.\end{Resp}
\switchcolumn
\EN
\begin{Resp}\Rsign Be ye therefore followers of God; \Star{} And walk ye therefore in love.\end{Resp}
\begin{Resp}\Vsign For Christ also hath loved us and hath given himself for us.\end{Resp}
\begin{Resp}\Rsign And walk ye therefore in love.\end{Resp}
\begin{Resp}\Vsign Glory be to the Father, and to the Son, \Star{} and to the Holy Ghost.\end{Resp}
\begin{Resp}\Rsign And walk ye therefore in love.\end{Resp}
\switchcolumn*

\LA
\LessonHead{Lectio IX}
\switchcolumn
\EN
\LessonHead{Lesson IX}
\switchcolumn*

\LA
\noindent Duo testimónia de Scriptúris réddidit síngulis rebus quas factas fuísse narrávit. Nam quia díxerat: Ad Jesum autem cum veníssent, ut vidérunt eum jam mórtuum, non fregérunt ejus crura, ad hoc pértinet testimónium: Os non comminuétis ex eo: quod præcéptum erat eis qui celebráre Pascha jussi sunt ovis immolatióne in véteri lege, quæ Domínicæ Passiónis umbra præcésserat; unde Pascha nostrum immolátus est Christus; de quo et Isaías prophéta prædíxerat: Sicut ovis ad immolándum ductus est. Item quia subjúnxerat dicens: Sed unus mílitum láncea latus ejus apéruit, ad hoc pértinet álterum testimónium: Vidébunt in quem transfixérunt, ubi promíssus est Christus, in ea qua crucifíxus est carne ventúrus.\par
\switchcolumn
\EN
\noindent From the Scriptures he giveth two testimonies, one for each of the things which he hath recorded as having been done. To the words: When they came to Jesus, and saw that he was dead already, they brake not his legs: belongeth the testimony: A bone of him shall not be broken: which same cometh from the Mosaic injunction laid upon the Jews, who were commanded by the Old Law to celebrate the Passover by the sacrifice of a lamb, which same was a foreshadowing of the passing of Christ, whence we have the passage: Christ our Passover is sacrificed for us: concerning which the Prophet Isaiah also foretold: He is brought as a lamb to the slaughter. To the words: One of the soldiers with a spear pierced his side: belongeth the other testimony: They shall look on him whom they pierced: which same is a promise of the coming of Christ in that selfsame flesh wherein he was afterwards crucified.\par
\switchcolumn*

\LA
\TeDeum{Te Deum}
\switchcolumn
\EN
\TeDeum{Te Deum}
\switchcolumn*

\end{paracol}

\Office{Sabbato infra Hebdomadam III post Octavam Pentecostes}{Saturday of the Third Week after Pentecost}{Feria}
\addcontentsline{toc}{subsection}{Sabbato infra Hebdomadam III post Octavam Pentecostes \textit{— Saturday of the Third Week after Pentecost}}
\begin{paracol}{2}
\LA
\LessonHead{Lectio I}
\switchcolumn
\EN
\LessonHead{Lesson I}
\switchcolumn*

\LA
\LessonTitle{De libro primo Regum}
\switchcolumn
\EN
\LessonTitle{Lesson from the first book of Samuel}
\switchcolumn*

\LA
\Cite{1 Reg 16:1-3}
\switchcolumn
\EN
\Cite{1 Sam 16:1-3}
\switchcolumn*

\LA
\noindent Dixítque Dóminus ad Samuélem: Usquequo tu luges Saul, cum ego projécerim eum ne regnet super Israël? Imple cornu tuum óleo et veni, ut mittam te ad Isai Bethlehemítem; provídi enim in fíliis ejus mihi regem. Et ait Samuel: Quómodo vadam? áudiet enim Saul et interfíciet me. Et ait Dóminus: Vítulum de arménto tolles in manu tua et dices: Ad immolándum Dómino veni. Et vocábis Isai ad víctimam, et ego osténdam tibi quid fácias, et unges quemcúmque monstrávero tibi.\par
\switchcolumn
\EN
\noindent And the Lord said to Samuel. How long wilt thou mourn for Saul, whom I have rejected from reigning over Israel? fill thy horn with oil, and come, that I may send thee to Isai the Bethlehemite: for I have provided me a king among his sons. And Samuel said: How shall I go? for Saul will hear of it, and he will kill me. And the Lord said: Thou shalt take with thee a calf of the herd, and thou shalt say: I am come to sacrifice to the Lord. And thou shalt call Isai to the sacrifice, and I will show thee what thou art to do, and thou shalt anoint him whom I shall show to thee\par
\switchcolumn*

\LA
\begin{Resp}\Rsign Peccávi super númerum arénæ maris, et multiplicáta sunt peccáta mea: et non sum dignus vidére altitúdinem cæli præ multitúdine iniquitátis meæ: quóniam irritávi iram tuam, \Star{} Et malum coram te feci.\end{Resp}
\begin{Resp}\Vsign Quóniam iniquitátem meam ego cognósco: et delíctum meum contra me est semper, quia tibi soli peccávi.\end{Resp}
\begin{Resp}\Rsign Et malum coram te feci.\end{Resp}
\switchcolumn
\EN
\begin{Resp}\Rsign My sins are many, yea, they are more in number than the sands of the sea; I am not worthy to look up toward heaven because of the multitude of my iniquities; for I have provoked thee to anger; \Star{} And done evil in thy sight.\end{Resp}
\begin{Resp}\Vsign For I acknowledge my transgression, and my sin is ever before me, for against thee only have I sinned.\end{Resp}
\begin{Resp}\Rsign And done evil in thy sight.\end{Resp}
\switchcolumn*

\LA
\LessonHead{Lectio II}
\switchcolumn
\EN
\LessonHead{Lesson II}
\switchcolumn*

\LA
\Cite{1 Reg 16:4-7}
\switchcolumn
\EN
\Cite{1 Sam 15:4-7}
\switchcolumn*

\LA
\noindent Fecit ergo Samuel sicut locútus est ei Dóminus, venítque in Béthlehem, et admiráti sunt senióres civitátis occurréntes ei dixerúntque: Pacificúsne est ingréssus tuus? Et ait: Pacíficus: ad immolándum Dómino veni: sanctificámini, et veníte mecum ut ímmolem. Sanctificávit ergo Isai et fílios ejus et vocávit eos ad sacrifícium. Cumque ingréssi essent, vidit Eliab et ait: Num coram Dómino est Christus ejus? Et dixit Dóminus ad Samuélem: Ne respícias vultum ejus neque altitúdinem statúræ ejus, quóniam abjéci eum, nec juxta intúitum hóminis ego júdico; homo enim videt ea quæ parent, Dóminus autem intuétur cor.\par
\switchcolumn
\EN
\noindent Then Samuel did as the Lord had said to him. And he came to Bethlehem, and the ancients of the city wondered, and meeting him, they said: Is thy coming hither peaceable? And he said: It is peaceable: I am come to offer sacrifice to the Lord, be ye sanctified, and come with me to the sacrifice. And he sanctified Isai and his sons, and called them to the sacrifice. And when they were come in, he saw Eliab, and said: Is the Lord's anointed before him? And the Lord said to Samuel: Look not on his countenance, nor on the height of his stature: because I have rejected him, nor do I judge according to the look of man: for man seeth those things that appear, but the Lord beholdeth the heart.\par
\switchcolumn*

\LA
\begin{Resp}\Rsign Exaudísti, Dómine, oratiónem servi tui, ut ædificárem templum nómini tuo: \Star{} Bénedic et sanctífica domum istam in sempitérnum, Deus Israël.\end{Resp}
\begin{Resp}\Vsign Dómine, qui custódis pactum cum servis tuis, qui ámbulant coram te in toto corde suo.\end{Resp}
\begin{Resp}\Rsign Bénedic et sanctífica domum istam in sempitérnum, Deus Israël.\end{Resp}
\switchcolumn
\EN
\begin{Resp}\Rsign O Lord, Thou hast hearkened unto the prayer of thy servant, that I might build a temple unto thy Name, \Star{} O God of Israel, bless Thou, and hallow this house for ever.\end{Resp}
\begin{Resp}\Vsign O Lord, Who keepest covenant with thy servants that walk before thee in all their heart.\end{Resp}
\begin{Resp}\Rsign O God of Israel, bless Thou, and hallow this house for ever.\end{Resp}
\switchcolumn*

\LA
\LessonHead{Lectio III}
\switchcolumn
\EN
\LessonHead{Lesson III}
\switchcolumn*

\LA
\Cite{1 Reg 16:8-11}
\switchcolumn
\EN
\Cite{1 Sam 16:8-11}
\switchcolumn*

\LA
\noindent Et vocávit Isai Abínadab et addúxit eum coram Samuéle. Qui dixit: Nec hunc elégit Dóminus. Addúxit autem Isai Samma, de quo ait: Etiam hunc non elégit Dóminus. Addúxit ítaque Isai septem fílios suos coram Samuéle, et ait Sámuel ad Isai: Non elégit Dóminus ex istis. Dixítque Samuel ad Isai: Numquid jam compléti sunt fílii? Qui respóndit: Adhuc réliquus est párvulus et pascit oves. Et ait Sámuel ad Isai: Mitte et adduc eum.\par
\switchcolumn
\EN
\noindent And Isai called Abinadab, and brought him before Samuel. And he said: Neither hath the Lord chosen this. And Isai brought Samma, and he said of him: Neither hath the Lord chosen this. Isai therefore brought his seven sons before Samuel: and Samuel said to Isai: The Lord hath not chosen any one of these. And Samuel said to Isai: Are here all thy sons? He answered: There remaineth yet a young one, who keepeth the sheep. And Samuel said to Isai: Send, and fetch him, for we will not sit down till he come hither.\par
\switchcolumn*

\LA
\begin{Resp}\Rsign Audi, Dómine, hymnum et oratiónem, quam servus tuus orat coram te hódie, ut sint óculi tui apérti, et aures tuæ inténtæ, \Star{} Super domum istam die ac nocte.\end{Resp}
\begin{Resp}\Vsign Réspice, Dómine, de sanctuário tuo, et de excélso cælórum habitáculo.\end{Resp}
\begin{Resp}\Rsign Super domum istam die ac nocte.\end{Resp}
\begin{Resp}\Vsign Glória Patri, et Fílio, \Star{} et Spirítui Sancto.\end{Resp}
\begin{Resp}\Rsign Super domum istam die ac nocte.\end{Resp}
\switchcolumn
\EN
\begin{Resp}\Rsign Hearken, O Lord, unto the cry and to the prayer which thy servant prayeth before thee today, that thine eyes may be open and thine ears attend; \Star{} Toward this house day and night.\end{Resp}
\begin{Resp}\Vsign Look down from thine high and holy place, O Lord, even from heaven thy dwelling.\end{Resp}
\begin{Resp}\Rsign Toward this house, day and night.\end{Resp}
\begin{Resp}\Vsign Glory be to the Father, and to the Son, \Star{} and to the Holy Ghost.\end{Resp}
\begin{Resp}\Rsign Toward this house, day and night.\end{Resp}
\switchcolumn*

\end{paracol}

\Office{Dominica IV Post Pentecosten}{Fourth Sunday after Pentecost}{Semiduplex}
\addcontentsline{toc}{subsection}{Dominica IV Post Pentecosten \textit{— Fourth Sunday after Pentecost}}
\begin{paracol}{2}
\LA
\LessonHead{Lectio I}
\switchcolumn
\EN
\LessonHead{Lesson I}
\switchcolumn*

\LA
\LessonTitle{De libro primo Regum}
\switchcolumn
\EN
\LessonTitle{Lesson from the first book of Samuel}
\switchcolumn*

\LA
\Cite{1 Reg 17:1-7}
\switchcolumn
\EN
\Cite{1 Sam 17:1-7}
\switchcolumn*

\LA
\noindent Congregántes autem Philísthiim ágmina sua in prǽlium convenérunt in Socho Judæ, et castrametáti sunt inter Socho et Azéca in fínibus Dommim. Porro Saul et fílii Israël congregáti venérunt in Vallem terebínthi, et direxérunt áciem ad pugnándum contra Philísthiim. Et Philísthiim stabant super montem ex parte hac, et Israël stabat supra montem ex áltera parte, vallísque erat inter eos. Et egréssus est vir spúrius de castris Philisthinórum nómine Góliath de Geth altitúdinis sex cubitórum et palmi. Et cassis ǽrea super caput ejus, et loríca squamáta induebátur; porro pondus lorícæ ejus quinque míllia siclórum æris erat. Et ócreas ǽreas habébat in crúribus, et clípeus ǽreus tegébat húmeros ejus. Hastíle autem hastæ ejus erat quasi liciatórium texéntium; ipsum autem ferrum hastæ ejus sexcéntos siclos habébat ferri; et ármiger ejus antecedébat eum.\par
\switchcolumn
\EN
\noindent Now the Philistines gathering together their troops to battle, assembled at Socho of Juda, and camped between Socho and Azeca in the borders of Dommim. And Saul and the children of Israel being gathered together came to the valley of Terebinth, and they set the army in array to fight against the Philistines. And the Philistines stood on a mountain on the one side, and Israel stood on a mountain on the other side: and there was a valley between them. And there went out a man baseborn from the camp of the Philistines named Goliath, of Geth, whose height was six cubits and a span: And he had a helmet of brass upon his head, and he was clothed with a coat of mail with scales, and the weight of his coat of mail was five thousand sicles of brass: And he had greaves of brass on his legs, and a buckler of brass covered his shoulders. And the staff of his spear was like a weaver's beam, and the head of his spear weighed six hundred sicles of iron: and his armourbearer went before him.\par
\switchcolumn*

\LA
\begin{Resp}\Rsign Præparáte corda vestra Dómino, et servíte illi soli: \Star{} Et liberábit vos de mánibus inimicórum vestrórum.\end{Resp}
\begin{Resp}\Vsign Convertímini ad eum in toto corde vestro, et auférte deos aliénos de médio vestri.\end{Resp}
\begin{Resp}\Rsign Et liberábit vos de mánibus inimicórum vestrórum.\end{Resp}
\switchcolumn
\EN
\begin{Resp}\Rsign Prepare your hearts unto the Lord, and serve Him Only \Star{} And He will deliver you out of the hand of your enemies.\end{Resp}
\begin{Resp}\Vsign Return unto Him with all your hearts, and put away the strange gods from among you.\end{Resp}
\begin{Resp}\Rsign And He will deliver you out of the hand of your enemies.\end{Resp}
\switchcolumn*

\LA
\LessonHead{Lectio II}
\switchcolumn
\EN
\LessonHead{Lesson II}
\switchcolumn*

\LA
\Cite{1 Reg 17:8-11}
\switchcolumn
\EN
\Cite{1 Sam 17:8-11}
\switchcolumn*

\LA
\noindent Stansque clamábat advérsum phalángas Israël et dicébat eis: Quare venístis paráti ad prǽlium? Numquid ego non sum Philisthǽus, et vos servi Saul? Elígite ex vobis virum, et descéndat ad singuláre certámen: si quíverit pugnáre mecum et percússerit me, érimus vobis servi; si autem ego prævalúero et percússero eum, vos servi éritis et serviétis nobis. Et ajébat Philisthǽus: Ego exprobrávi agmínibus Israël hódie: Date mihi virum, et íneat mecum singuláre certámen. Audiens autem Saul et omnes Israëlítæ sermónes Philisthǽi hujuscémodi, stupébant et metuébant nimis.\par
\switchcolumn
\EN
\noindent And standing he cried out to the bands of Israel, and said to them: Why are you come out prepared to fight? am not I a Philistine, and you the servants of Saul? Choose out a man of you, and let him come down and fight hand to hand. If he be able to fight with me, and kill me, we will be servants to you: but if I prevail against him, and kill him, you shall be servants, and shall serve us. And the Philistine said: I have defied the bands of Israel this day: Give me a man, and let him fight with me hand to hand. And Saul and all the Israelites hearing these words of the Philistine were dismayed, and greatly afraid.\par
\switchcolumn*

\LA
\begin{Resp}\Rsign Deus ómnium exaudítor est: ipse misit Angelum suum, et tulit me de óvibus patris mei: \Star{} Et unxit me unctióne misericórdiæ suæ.\end{Resp}
\begin{Resp}\Vsign Dóminus, qui erípuit me de ore leónis, et de manu béstiæ liberávit me.\end{Resp}
\begin{Resp}\Rsign Et unxit me unctióne misericórdiæ suæ.\end{Resp}
\switchcolumn
\EN
\begin{Resp}\Rsign God, Which heareth all, even He sent His Angel, and took me from keeping my father's sheep, and \Star{} Anointed me with the oil of His mercy.\end{Resp}
\begin{Resp}\Vsign The Lord That delivered me out of the mouth of the lion, and out of the paw of the bear\end{Resp}
\begin{Resp}\Rsign And anointed me with the oil of His mercy.\end{Resp}
\switchcolumn*

\LA
\LessonHead{Lectio III}
\switchcolumn
\EN
\LessonHead{Lesson III}
\switchcolumn*

\LA
\Cite{1 Reg 17:12-16}
\switchcolumn
\EN
\Cite{1 Sam 17:12-16}
\switchcolumn*

\LA
\noindent David autem erat fílius viri Ephrathǽi, de quo supra dictum est, de Béthlehem Juda, cui nomen erat Isai, qui habébat octo fílios, et erat vir in diébus Saul senex et grandǽvus inter viros. Abiérunt autem tres fílii ejus majóres post Saul in prǽlium; et nómina trium filiórum ejus, qui perrexérunt ad bellum, Eliab primogénitus et secúndus Abínadab tertiúsque Samma; David autem erat mínimus. Tribus ergo majóribus secútis Saulem, ábiit David et revérsus est a Saul, ut pásceret gregem patris sui in Béthlehem. Procedébat vero Philisthǽus mane et véspere, et stabat quadragínta diébus.\par
\switchcolumn
\EN
\noindent Now David was the son of that Ephrathite of Bethlehem Juda before mentioned, whose name was Isai, who had eight sons, and was an old man in the days of Saul, and of great age among men. And his three eldest sons followed Saul to the battle: and the names of his three sons that went to the battle, were Eliab the firstborn, and the second Abinadab, and the third Samma. But David was the youngest. So the three eldest having followed Saul, David went, and returned from Saul, to feed his father's flock at Bethlehem. Now the Philistine came out morning and evening, and presented himself forty days.\par
\switchcolumn*

\LA
\begin{Resp}\Rsign Dóminus, qui erípuit me de ore leónis, et de manu béstiæ liberávit me, \Star{} Ipse me erípiet de mánibus inimicórum meórum.\end{Resp}
\begin{Resp}\Vsign Misit Deus misericórdiam suam et veritátem suam: ánimam meam erípuit de médio catulórum leónum.\end{Resp}
\begin{Resp}\Rsign Ipse me erípiet de mánibus inimicórum meórum.\end{Resp}
\begin{Resp}\Vsign Glória Patri, et Fílio, \Star{} et Spirítui Sancto.\end{Resp}
\begin{Resp}\Rsign Ipse me erípiet de mánibus inimicórum meórum.\end{Resp}
\switchcolumn
\EN
\begin{Resp}\Rsign The Lord That delivered me out of the mouth of the lion, and out of the paw of the bear \Star{} He will deliver me out of the hand of mine enemies.\end{Resp}
\begin{Resp}\Vsign God hath sent forth His mercy and His truth, and delivered my soul from among the lion's whelps.\end{Resp}
\begin{Resp}\Rsign He will deliver me out of the hand of mine enemies.\end{Resp}
\begin{Resp}\Vsign Glory be to the Father, and to the Son, \Star{} and to the Holy Ghost.\end{Resp}
\begin{Resp}\Rsign He will deliver me out of the hand of mine enemies.\end{Resp}
\switchcolumn*

\LA
\LessonHead{Lectio IV}
\switchcolumn
\EN
\LessonHead{Lesson IV}
\switchcolumn*

\LA
\LessonTitle{Sermo sancti Augustíni Epíscopi}
\switchcolumn
\EN
\LessonTitle{from the Sermons of St. Augustine, Bishop (of Hippo.)}
\switchcolumn*

\LA
\Source{Sermo 197 de Temp. circa med.}
\switchcolumn
\EN
\Source{Serm. 197. de Temp. circa med.}
\switchcolumn*

\LA
\noindent Stabant fílii Israël contra adversários quadragínta diébus. Quadragínta dies, propter quátuor témpora et quátuor partes orbis terræ, vitam præséntem signíficant, in qua contra Góliath vel exércitum ejus, id est, contra diábolum et ángelos ejus, Christianórum pópulus pugnáre non désinit. Nec tamen víncere posset, nisi verus David Christus cum báculo, id est, cum crucis mystério descendísset. Ante advéntum enim Christi, fratres caríssimi, solútus erat diábolus; véniens Christus fecit de eo, quod in Evangélio dictum est: Nemo potest intráre in domum fortis et vasa ejus dirípere, nisi prius alligáverit fortem. Venit ergo Christus et alligávit diábolum.\par
\switchcolumn
\EN
\noindent The children of Israel faced their enemies for forty days. These forty days, by reason of the four Seasons of the year, and of the four Continents of the globe, are a figure of this present life, during which the Christian world ceaseth not to be arrayed in battle against the devil and his angels, as it were against Goliath and the army of the Philistines. Neither can they hope to overcome him, were it not for the true David, that is, Christ, with His staff, that is, with the Mystery of His Cross. For before Christ came, my dearly beloved brethren, the devil was at large. But when Christ came, He did to him what is written in the Gospel, where it is said “How can one enter into a strong man's house, and spoil his goods, except he first bind the strong man?” Christ therefore came, and bound the devil.\par
\switchcolumn*

\LA
\begin{Resp}\Rsign Percússit Saul mille, et David decem míllia: \Star{} Quia manus Dómini erat cum illo: percússit Philisthǽum, et ábstulit oppróbrium ex Israël.\end{Resp}
\begin{Resp}\Vsign Nonne iste est David, de quo canébant in choro, dicéntes: Saul percússit mille, et David decem míllia?\end{Resp}
\begin{Resp}\Rsign Quia manus Dómini erat cum illo: percússit Philisthǽum, et ábstulit oppróbrium ex Israël.\end{Resp}
\switchcolumn
\EN
\begin{Resp}\Rsign Saul hath slain his thousands, and David his ten thousands. \Star{} Because the hand of the Lord was with him, he smote the Philistine, and took away the reproach from Israel.\end{Resp}
\begin{Resp}\Vsign Is not this David? Did they not sing one to another of him in dances, saying: Saul hath slain his thousands, and David his ten thousands?\end{Resp}
\begin{Resp}\Rsign Because the hand of the Lord was with him, he smote the Philistine, and took away the reproach from Israel.\end{Resp}
\switchcolumn*

\LA
\LessonHead{Lectio V}
\switchcolumn
\EN
\LessonHead{Lesson V}
\switchcolumn*

\LA
\noindent Sed dicit áliquis: Si alligátus est, quare adhuc tantum prǽvalet? Verum est, fratres caríssimi, quia multum prǽvalet: sed tépidis, et negligéntibus, et Deum in veritáte non timéntibus dominátur. Alligátus est enim tamquam innéxus canis caténis et néminem potest mordére, nisi eum qui se illi mortífera securitáte conjúnxerit. Jam vidéte, fratres, quam stultus est homo ille, quem canis in caténa pósitus mordet. Tu te illi per voluntátes et cupiditátes sǽculi noli conjúngere, et ille ad te non præsúmit accédere. Latráre potest, sollicitáre potest; mordére omníno non potest, nisi voléntem. Non enim cogéndo, sed suadéndo nocet; nec extórquet a nobis consénsum, sed petit.\par
\switchcolumn
\EN
\noindent But some man will say: If he is bound, why is he still so powerful? It is quite true, my dearly beloved brethren, that he is very powerful but his lordship is over the lukewarm and the careless, and such as fear not God in truth. He is chained up like a dog, and can only bite those who are such suicidal fools as to go within the length of his tether. Look you, my brethren, what a dolt a man must be who getteth himself bitten by a dog that is chained up. Let not the desires and lusts of the world draw thee within reach of him, and he will not be able to get at thee. He can bark, he can whine but he can only bite those who are willing to be bitten. He assaileth us, not by violence, but by persuasion he asketh, not seizeth, our consent.\par
\switchcolumn*

\LA
\begin{Resp}\Rsign Montes Gélboë, nec ros nec plúvia véniant super vos, \Star{} Ubi cecidérunt fortes Israël.\end{Resp}
\begin{Resp}\Vsign Omnes montes, qui estis in circúitu ejus, vísitet Dóminus: a Gélboë autem tránseat.\end{Resp}
\begin{Resp}\Rsign Ubi cecidérunt fortes Israël.\end{Resp}
\switchcolumn
\EN
\begin{Resp}\Rsign Ye mountains of Gilboa, let there be no dew, neither let there be rain upon you \Star{} For there are the mighty of Israel fallen.\end{Resp}
\begin{Resp}\Vsign All ye mountains that stand round about, the Lord look upon you but let Him pass by Gilboa\end{Resp}
\begin{Resp}\Rsign For there are the mighty of Israel fallen.\end{Resp}
\switchcolumn*

\LA
\LessonHead{Lectio VI}
\switchcolumn
\EN
\LessonHead{Lesson VI}
\switchcolumn*

\LA
\noindent Venit ergo David et invénit Judæórum pópulum contra diábolum præliántem; et cum nullus esset qui præsúmeret ad singuláre certámen accédere, ille qui figúram Christi gerébat, procéssit ad prǽlium, tulit báculum in manu sua et éxiit contra Góliath. Et in illo quidem tunc figurátum est, quod in Dómino Jesu Christo complétum est. Venit enim verus David Christus, qui contra spiritálem Góliath, id est, contra diábolum pugnatúrus crucem suam ipse portávit. Vidéte, fratres, ubi David Góliath percússerit: in fronte útique, ubi crucis signáculum non habébat. Sicut enim báculus crucis typum hábuit, ita étiam et lapis ille, de quo percússus est, Christum Dóminum figurábat.\par
\switchcolumn
\EN
\noindent David, then, came, and found the Jewish people set in battle array against the devil and since there was no one who dared to go to single combat, he, who was a type of Christ, sallied out to the battle, took his staff in his hand, and went forth against Goliath. In him was a shadow of a substance which is in Christ. Christ, the true David, when He went forth to fight against the spiritual Goliath, that is to say, against the devil, went forth bearing His Cross. Ye see, my brethren, in what part it was that David smote Goliath it was upon that forehead whereon the Cross had never been traced. And as the staff of David was a figure of the Cross of Christ, so was the stone wherewith the giant was smitten a figure of the Lord Himself.\par
\switchcolumn*

\LA
\begin{Resp}\Rsign Ego te tuli de domo patris tui, dicit Dóminus, et pósui te páscere gregem pópuli mei: \Star{} Et fui tecum in ómnibus ubicúmque ambulásti, firmans regnum tuum in ætérnum.\end{Resp}
\begin{Resp}\Vsign Fecíque tibi nomen grande, juxta nomen magnórum, qui sunt in terra: et réquiem dedi tibi ab ómnibus inimícis tuis.\end{Resp}
\begin{Resp}\Rsign Et fui tecum in ómnibus ubicúmque ambulásti, firmans regnum tuum in ætérnum.\end{Resp}
\begin{Resp}\Vsign Glória Patri, et Fílio, \Star{} et Spirítui Sancto.\end{Resp}
\begin{Resp}\Rsign Et fui tecum in ómnibus ubicúmque ambulásti, firmans regnum tuum in ætérnum.\end{Resp}
\switchcolumn
\EN
\begin{Resp}\Rsign Thus saith the Lord: I took thee out of thy father's house, and appointed thee to be ruler over My people, over Israel. \Star{} And I was with thee whithersoever thou wentest, to establish thy kingdom for ever.\end{Resp}
\begin{Resp}\Vsign And I have made thee a great name, like unto the name of the great men that are in the earth, and have caused thee to rest from all thine enemies.\end{Resp}
\begin{Resp}\Rsign And I was with thee whithersoever thou wentest, to establish thy kingdom for ever.\end{Resp}
\begin{Resp}\Vsign Glory be to the Father, and to the Son, \Star{} and to the Holy Ghost.\end{Resp}
\begin{Resp}\Rsign And I was with thee whithersoever thou wentest, to establish thy kingdom for ever.\end{Resp}
\switchcolumn*

\LA
\LessonHead{Lectio VII}
\switchcolumn
\EN
\LessonHead{Lesson VII}
\switchcolumn*

\LA
\LessonTitle{Léctio sancti Evangélii secúndum Lucam}
\switchcolumn
\EN
\LessonTitle{From the Holy Gospel according to Luke}
\switchcolumn*

\LA
\Cite{Luc 5:1-11}
\switchcolumn
\EN
\Cite{Luke 5:1-11}
\switchcolumn*

\LA
\noindent In illo témpore: Cum turbæ irrúerent in Jesum, ut audírent verbum Dei, et ipse stabat secus stagnum Genésareth. Et réliqua.\par
\switchcolumn
\EN
\noindent At that time: As the people pressed upon Jesus, to hear the word of God, He stood by the lake of Gennesareth. And so on.\par
\switchcolumn*

\LA
\LessonTitle{Homilía sancti Ambrósii Epíscopi}
\switchcolumn
\EN
\LessonTitle{Homily by St. Ambrose, Bishop of Milan.}
\switchcolumn*

\LA
\Source{Liber 4 in Lucæ cap. 5 prope finem libri}
\switchcolumn
\EN
\Source{Bk. iv. on Luke v.}
\switchcolumn*

\LA
\noindent Ubi Dóminus multis impartívit vária génera sanitátum, nec témpore, nec loco pótuit ab stúdio sanándi turba cohibéri. Vesper incúbuit, sequebántur: stagnum occúrrit, urgébant: et ídeo ascéndit in Petri navim. Hæc est illa navis, quæ adhuc secúndum Matthǽum flúctuat, secúndum Lucam replétur píscibus: ut et princípia Ecclésiæ fluctuántis, et posterióra exuberántis agnóscas. Pisces enim sunt, qui hanc enávigant vitam. Ibi adhuc discípulis Christus dormit, hic prǽcipit; dormit enim tépidis, perféctis vígilat.\par
\switchcolumn
\EN
\noindent When the Lord wrought so many works of healing, neither time nor place could restrain the people from seeking health. Evening came, and they still followed Him He went down to the lake, and they still pressed upon Him and therefore He entered into Peter's ship. This is that ship, which spiritually up to this very hour, according to the expression of Matthew, is buffeted by tempests, but still, according to Luke, is filled with fishes, this signifying, that, for a while, to labour is present to the Church, but, hereafter, it shall be to rejoice. The fishes are they which swim in the troublous waters of human life. In this ship also spiritually doth Christ, for His disciples, still sleep, and still command; for He sleepeth for the lukewarm, and watcheth for the perfect.\par
\switchcolumn*

\LA
\begin{Resp}\Rsign Peccávi super númerum arénæ maris, et multiplicáta sunt peccáta mea: et non sum dignus vidére altitúdinem cæli præ multitúdine iniquitátis meæ: quóniam irritávi iram tuam, \Star{} Et malum coram te feci.\end{Resp}
\begin{Resp}\Vsign Quóniam iniquitátem meam ego cognósco: et delíctum meum contra me est semper, quia tibi soli peccávi.\end{Resp}
\begin{Resp}\Rsign Et malum coram te feci.\end{Resp}
\switchcolumn
\EN
\begin{Resp}\Rsign My sins are many, yea, they are more in number than the sands of the sea; I am not worthy to look up toward heaven because of the multitude of my iniquities; for I have provoked thee to anger; \Star{} And done evil in thy sight.\end{Resp}
\begin{Resp}\Vsign For I acknowledge my transgression, and my sin is ever before me, for against thee only have I sinned.\end{Resp}
\begin{Resp}\Rsign And done evil in thy sight.\end{Resp}
\switchcolumn*

\LA
\LessonHead{Lectio VIII}
\switchcolumn
\EN
\LessonHead{Lesson VIII}
\switchcolumn*

\LA
\noindent Non turbátur hæc navis, in qua prudéntia navigat, abest perfidia, fides aspirat. Quemádmodum enim turbari poterat, cui præerat is, in quo Ecclésiæ firmaméntum est? Illic ergo turbátio, ubi módica fides; hic securitas, ubi perfécta diléctio. Et si aliis imperátur ut laxent rétia sua, soli tamen Petro dícitur: Duc in altum; hoc est, in profúndum disputatiónum. Quid enim tam altum, quam altitúdinem divitiárum vidére, scire Dei Fílium, et professiónem divinæ generatiónis assúmere? Quam licet mens non queat humana plene ratiónis investigatióne comprehéndere, fidei tamen plenitúdo compléctitur.\par
\switchcolumn
\EN
\noindent O fear, then, for the ship where wisdom steereth, false teaching is not known, and faith swelleth the sails. How shall she be troubled, whose Lord is Himself the Church's sure Foundation It is where faith is weak that there is fear where love is perfect, there there is safety. To many it is commanded to loose their nets, but to Peter only to “Launch out into the deep,” that is, into the depths of doctrine. What indeed is there so deep, as to gaze upon the depth of all riches, to recognise the Son of God, and to take up the confession of His Divine generation This is a thing which the mind is not able to grasp by the searchings of man's reason, but which is embraced by an hearty faith.\par
\switchcolumn*

\LA
\begin{Resp}\Rsign Duo Séraphim clamábant alter ad álterum: \Star{} Sanctus, sanctus, sanctus Dóminus Deus Sábaoth: * Plena est omnis terra glória ejus.\end{Resp}
\begin{Resp}\Vsign Tres sunt qui testimónium dant in cælo: Pater, Verbum, et Spíritus Sanctus: et hi tres unum sunt.\end{Resp}
\begin{Resp}\Rsign Sanctus, sanctus, sanctus Dóminus Deus Sábaoth.\end{Resp}
\begin{Resp}\Vsign Glória Patri, et Fílio, \Star{} et Spirítui Sancto.\end{Resp}
\begin{Resp}\Rsign Plena est omnis terra glória ejus.\end{Resp}
\switchcolumn
\EN
\begin{Resp}\Rsign One Seraph cried unto another: \Star{} Holy, Holy, Holy is the Lord God of hosts; the whole earth is full of His glory.\end{Resp}
\begin{Resp}\Vsign There are Three That bear record in heaven: the Father, the Word, and the Holy Ghost, and these Three are One.\end{Resp}
\begin{Resp}\Rsign Holy, Holy, Holy is the Lord God of hosts.\end{Resp}
\begin{Resp}\Vsign Glory be to the Father, and to the Son, \Star{} and to the Holy Ghost.\end{Resp}
\begin{Resp}\Rsign The whole earth is full of His glory.\end{Resp}
\switchcolumn*

\LA
\LessonHead{Lectio IX}
\switchcolumn
\EN
\LessonHead{Lesson IX}
\switchcolumn*

\LA
\noindent Nam, etsi non licet mihi scire quemádmodum natus sit, non licete tamen nescire quod natus sit. Seriem generatiónis ignoro; sed auctórem generatiónis agnosco. Non interfúimus cum ex Patre Dei Fílius nascerétur; sed interfúimus cum a Patre Dei Fílius dicerétur. Si Deo non crédimus, cui credémus? Omnia enim quæ crédimus, vel visu crédimus, vel audítu: visus sæpe fállitur, audítus in fide est.\par
\switchcolumn
\EN
\noindent Nor albeit, it is not given unto me to know how He was born, yet that born He was, I may not be ignorant. What the order of His generation was, I know not, but the Source of His generation I acknowledge. None hath beheld the Begetting of the Son of God by the Father, but the Church hath stood by to hear the Father testify that this is His beloved Son. (Luke iii. 22.) If we believe not God, whom shall we believe? For whatsoever we believe cometh either by sight or by hearing; sight is oftentimes deceived, but “faith cometh by hearing.” (Rom. x. 17.)\par
\switchcolumn*

\LA
\TeDeum{Te Deum}
\switchcolumn
\EN
\TeDeum{Te Deum}
\switchcolumn*

\end{paracol}

\Office{Feria secunda infra Hebdomadam IV post Octavam Pentecostes}{Monday of the Fourth Week after Pentecost}{Feria}
\addcontentsline{toc}{subsection}{Feria secunda infra Hebdomadam IV post Octavam Pentecostes \textit{— Monday of the Fourth Week after Pentecost}}
\begin{paracol}{2}
\LA
\LessonHead{Lectio I}
\switchcolumn
\EN
\LessonHead{Lesson I}
\switchcolumn*

\LA
\LessonTitle{De libro primo Regum}
\switchcolumn
\EN
\LessonTitle{Lesson from the first book of Samuel}
\switchcolumn*

\LA
\Cite{1 Reg 17:25-26}
\switchcolumn
\EN
\Cite{1 Sam 17:25-26}
\switchcolumn*

\LA
\noindent Et dixit unus quíspiam de Israël: Num vidístis virum hunc qui ascéndit? Ad exprobrándum enim Israéli ascéndit. Virum ergo, qui percússerit eum, ditábit rex divítiis magnis et fíliam suam dabit ei et domum patris ejus fáciet absque tribúto in Israël. Et ait David ad viros, qui stabant secum, dicens: Quid dábitur viro, qui percússerit Philisthǽum hunc et túlerit oppróbrium de Israël? Quis enim est hic Philisthǽus incircumcísus, qui exprobrávit ácies Dei vivéntis?\par
\switchcolumn
\EN
\noindent And some one of Israel said: Have you seen this man that is come up, for he is come up to defy Israel. And the man that shall slay him, the king will enrich with great riches, and will give him his daughter, and will make his father's house free from tribute in Israel. And David spoke to the men that stood by him, saying: What shall be given to the man that shall kill this Philistine, and shall take away the reproach from Israel? for who is this uncircumcised Philistine, that he should defy the armies of the living God?\par
\switchcolumn*

\LA
\begin{Resp}\Rsign Recordáre, Dómine, testaménti tui, et dic Angelo percutiénti: Cesset jam manus tua, \Star{} Ut non desolétur terra, et ne perdas omnem ánimam vivam.\end{Resp}
\begin{Resp}\Vsign Ego sum qui peccávi, ego qui iníque egi: isti qui oves sunt, quid fecérunt? Avertátur, óbsecro, furor tuus, Dómine, a pópulo tuo.\end{Resp}
\begin{Resp}\Rsign Ut non desolétur terra, et ne perdas omnem ánimam vivam.\end{Resp}
\switchcolumn
\EN
\begin{Resp}\Rsign Remember, O Lord, thy covenant, and say unto the destroying Angel: Stay now thine hand; \Star{} That the land be not utterly laid waste, and that thou destroy not every living soul.\end{Resp}
\begin{Resp}\Vsign Even I it is that have sinned, and done evil indeed; but these sheep, what have they done? Let thine anger, I pray thee, O Lord, be turned away from thy people.\end{Resp}
\begin{Resp}\Rsign That the land be not utterly laid waste, and that Thou destroy not every living soul.\end{Resp}
\switchcolumn*

\LA
\LessonHead{Lectio II}
\switchcolumn
\EN
\LessonHead{Lesson II}
\switchcolumn*

\LA
\Cite{1 Reg 17:31-33}
\switchcolumn
\EN
\Cite{1 Sam 17:31-33}
\switchcolumn*

\LA
\noindent Audíta sunt autem verba, quæ locútus est David, et annuntiáta in conspéctu Saul. Ad quem cum fuísset addúctus, locútus est ei: Non cóncidat cor cujúsquam in eo: ego servus tuus vadam et pugnábo advérsus Philisthǽum. Et ait Saul ad David: Non vales resistere Philisthǽo isti nec pugnáre advérsus eum, quia puer es, hic autem vir bellátor est ab adulescéntia sua.\par
\switchcolumn
\EN
\noindent And the words which David spoke were heard, and were rehearsed before Saul. And when he was brought to him, he said to him: Let not any man's heart be dismayed in him: I thy servant will go, and will fight against the Philistine. And Saul said to David: Thou art not able to withstand this Philistine, nor to fight against him: for thou art but a boy, but he is a warrior from his youth.\par
\switchcolumn*

\LA
\begin{Resp}\Rsign Exaudísti, Dómine, oratiónem servi tui, ut ædificárem templum nómini tuo: \Star{} Bénedic et sanctífica domum istam in sempitérnum, Deus Israël.\end{Resp}
\begin{Resp}\Vsign Dómine, qui custódis pactum cum servis tuis, qui ámbulant coram te in toto corde suo.\end{Resp}
\begin{Resp}\Rsign Bénedic et sanctífica domum istam in sempitérnum, Deus Israël.\end{Resp}
\switchcolumn
\EN
\begin{Resp}\Rsign O Lord, Thou hast hearkened unto the prayer of thy servant, that I might build a temple unto thy Name, \Star{} O God of Israel, bless Thou, and hallow this house for ever.\end{Resp}
\begin{Resp}\Vsign O Lord, Who keepest covenant with thy servants that walk before thee in all their heart.\end{Resp}
\begin{Resp}\Rsign O God of Israel, bless Thou, and hallow this house for ever.\end{Resp}
\switchcolumn*

\LA
\LessonHead{Lectio III}
\switchcolumn
\EN
\LessonHead{Lesson III}
\switchcolumn*

\LA
\Cite{1 Reg 17:34-36}
\switchcolumn
\EN
\Cite{1 Sam 17:34-36}
\switchcolumn*

\LA
\noindent Dixítque David ad Saul: Pascébat servus tuus patris sui gregem, et veniébat leo vel ursus et tollébat aríetem de médio gregis, et persequébar eos et percutiébam eruebámque de ore eórum; et illi consurgébant advérsum me, et apprehendébam mentum eórum et suffocábam interficiebámque eos; nam et leónem et ursum interféci ego servus tuus. Erit ígitur et Philisthǽus hic incircumcísus quasi unus ex eis.\par
\switchcolumn
\EN
\noindent And David said to Saul: thy servant kept his father's sheep, and there came a lion, or a bear, and took a ram out of the midst of the flock: And I pursued after them, and struck them, and delivered it out of their mouth: and they rose up against me, and I caught them by the throat, and I strangled and killed them. For I thy servant have killed both a lion and a bear: and this uncircumcised Philistine shall be also as one of them. I will go now, and take away the reproach of the people: for who is this uncircumcised Philistine, who hath dared to curse the army of the living God?\par
\switchcolumn*

\LA
\begin{Resp}\Rsign Audi, Dómine, hymnum et oratiónem, quam servus tuus orat coram te hódie, ut sint óculi tui apérti, et aures tuæ inténtæ, \Star{} Super domum istam die ac nocte.\end{Resp}
\begin{Resp}\Vsign Réspice, Dómine, de sanctuário tuo, et de excélso cælórum habitáculo.\end{Resp}
\begin{Resp}\Rsign Super domum istam die ac nocte.\end{Resp}
\begin{Resp}\Vsign Glória Patri, et Fílio, \Star{} et Spirítui Sancto.\end{Resp}
\begin{Resp}\Rsign Super domum istam die ac nocte.\end{Resp}
\switchcolumn
\EN
\begin{Resp}\Rsign Hearken, O Lord, unto the cry and to the prayer which thy servant prayeth before thee today, that thine eyes may be open and thine ears attend; \Star{} Toward this house day and night.\end{Resp}
\begin{Resp}\Vsign Look down from thine high and holy place, O Lord, even from heaven thy dwelling.\end{Resp}
\begin{Resp}\Rsign Toward this house, day and night.\end{Resp}
\begin{Resp}\Vsign Glory be to the Father, and to the Son, \Star{} and to the Holy Ghost.\end{Resp}
\begin{Resp}\Rsign Toward this house, day and night.\end{Resp}
\switchcolumn*

\end{paracol}

\Office{Feria tertia infra Hebdomadam IV post Octavam Pentecostes}{Tuesday of the Fourth Week after Pentecost}{Feria}
\addcontentsline{toc}{subsection}{Feria tertia infra Hebdomadam IV post Octavam Pentecostes \textit{— Tuesday of the Fourth Week after Pentecost}}
\begin{paracol}{2}
\LA
\LessonHead{Lectio I}
\switchcolumn
\EN
\LessonHead{Lesson I}
\switchcolumn*

\LA
\LessonTitle{De libro primo Regum}
\switchcolumn
\EN
\LessonTitle{Lesson from the first book of Samuel}
\switchcolumn*

\LA
\Cite{1 Reg 17:38-40}
\switchcolumn
\EN
\Cite{1 Sam 17:38-40}
\switchcolumn*

\LA
\noindent Et índuit Saul David vestiméntis suis et impósuit gáleam ǽream super caput ejus et vestívit eum loríca. Accínctus ergo David gládio ejus super vestem suam cœpit tentáre si armátus posset incédere, non enim habébat consuetúdinem. Dixítque David ad Saul: Non possum sic incédere, quia non usum hábeo. Et depósuit ea, et tulit báculum suum, quem semper habébat in mánibus, et elégit sibi quinque limpidíssimos lápides de torrénte et misit eos in peram pastorálem, quam habébat secum, et fundam manu tulit et procéssit advérsum Philisthǽum.\par
\switchcolumn
\EN
\noindent And Saul clothed David with his garments, and put a helmet of brass upon his head, and armed him with a coat of mail. And David having girded his sword upon his armour, began to try if he could walk in armour: for he was not accustomed to it. And David said to Saul: I cannot go thus, for I am not used to it. And he laid them off, And he took his staff, which he had always in his hands: and chose him five smooth stones out of the brook, and put them into the shepherd's scrip, which he had with him, and he took a sling in his hand, and went forth against the Philistine.\par
\switchcolumn*

\LA
\begin{Resp}\Rsign Dómine, si convérsus fúerit pópulus tuus, et oráverit ad sanctuárium tuum: \Star{} Tu exáudies de cælo, Dómine, et líbera eos de mánibus inimicórum suórum.\end{Resp}
\begin{Resp}\Vsign Si peccáverit in te pópulus tuus, et convérsus égerit pœniténtiam, veniénsque oráverit in isto loco.\end{Resp}
\begin{Resp}\Rsign Tu exáudies de cælo, Dómine, et líbera eos de mánibus inimicórum suórum.\end{Resp}
\switchcolumn
\EN
\begin{Resp}\Rsign Lord, when thy people shall turn again to thee, and shall pray unto thee in this house; \Star{} Then hear Thou in heaven, O Lord, and deliver them out of the hand of their enemies.\end{Resp}
\begin{Resp}\Vsign If thy people sin against thee, and turn again, and repent, and come and pray unto thee in this house.\end{Resp}
\begin{Resp}\Rsign Then hear Thou in heaven, O Lord, and deliver them out of the hand of their enemies.\end{Resp}
\switchcolumn*

\LA
\LessonHead{Lectio II}
\switchcolumn
\EN
\LessonHead{Lesson II}
\switchcolumn*

\LA
\Cite{1 Reg 17:41-46}
\switchcolumn
\EN
\Cite{1 Sam 17:41-46}
\switchcolumn*

\LA
\noindent Ibat autem Philisthǽus incédens et appropínquans advérsum David, et ármiger ejus ante eum. Cumque inspexísset Philisthǽus et vidísset David, despéxit eum; erat enim adoléscens rufus et pulcher aspéctu. Et dixit Philisthǽus ad David: Numquid ego canis sum, quod tu venis ad me cum báculo? Et maledíxit Philisthǽus David in diis suis, dixítque ad David: Veni ad me et dabo carnes tuas volatílibus cæli et béstiis terræ. Dixit autem David ad Philisthǽum: Tu venis ad me cum gládio et hasta et clípeo, ego autem vénio ad te in nómine Dómini exercítuum, Dei ágminum Israël, quibus exprobrásti hódie; et dabit te Dóminus in manu mea, et percútiam te et áuferam caput tuum a te et dabo cadávera castrórum Philísthiim hódie volatílibus cæli et béstiis terræ, ut sciat omnis terra quia est Deus in Israël.\par
\switchcolumn
\EN
\noindent And the Philistine came on, and drew nigh against David, and his armour-bearer before him. And when the Philistine looked, and beheld David, he despised him. For he was a young man, ruddy, and of a comely countenance. And the Philistine said to David: Am I a dog, that thou comest to me with a staff? And the Philistine cursed David by his gods. And he said to David: Come to me, and I will give thy flesh to the birds of the air, and to the beasts of the earth. And David said to the Philistine: Thou comest to me with a sword, and with a spear, and with a shield: but I come to thee in the name of the Lord of hosts, the God of the armies of Israel, which thou hast defied. This day, and the Lord will deliver thee into my hand, and I will slay thee, and take away thy head from thee: and I will give the carcasses of the army of the Philistines this day to the birds of the air, and to the beasts of the earth: that all the earth may know that there is a God in Israel.\par
\switchcolumn*

\LA
\begin{Resp}\Rsign Factum est, dum tólleret Dóminus Elíam per túrbinem in cælum, \Star{} Eliséus clamábat, dicens: Pater mi, pater mi, currus Israël, et auríga ejus.\end{Resp}
\begin{Resp}\Vsign Cumque pérgerent, et incedéntes sermocinaréntur, ecce currus ígneus et equi ígnei divisérunt utrúmque, et ascéndit Elías per túrbinem in cælum.\end{Resp}
\begin{Resp}\Rsign Eliséus clamábat, dicens: Pater mi, pater mi, currus Israël, et auríga ejus.\end{Resp}
\switchcolumn
\EN
\begin{Resp}\Rsign And it came to pass when the Lord would take up Elijah into heaven by a whirlwind \Star{} Elisha cried, and said: My father, my father, the chariot of Israel, and the horsemen thereof.\end{Resp}
\begin{Resp}\Vsign And as they still went on, and talked, behold, there appeared a chariot of fire, and horses of fire, and parted them both asunder, and Elijah went up by a whirlwind into heaven.\end{Resp}
\begin{Resp}\Rsign Elisha cried, and said: My father, my father, the chariot of Israel, and the horsemen thereof.\end{Resp}
\switchcolumn*

\LA
\LessonHead{Lectio III}
\switchcolumn
\EN
\LessonHead{Lesson III}
\switchcolumn*

\LA
\Cite{1 Reg 17:48-51}
\switchcolumn
\EN
\Cite{1 Sam 17:48-51}
\switchcolumn*

\LA
\noindent Cum ergo surrexísset Philisthǽus et veníret et appropinquáret contra David, festinávit David et cucúrrit ad pugnam ex advérso Philisthǽi. Et misit manum suam in peram tulítque unum lápidem et funda jecit et circumdúcens percússit Philisthǽum in fronte; et infíxus est lapis in fronte ejus, et cécidit in fáciem suam super terram. Prævaluítque David advérsum Philisthǽum in funda et lápide percussúmque Philisthǽum interfécit. Cumque gládium non habéret in manu, David cucúrrit et stetit super Philisthǽum et tulit gládium ejus et edúxit eum de vagína sua et interfécit eum, præcidítque caput ejus.\par
\switchcolumn
\EN
\noindent And when the Philistine arose and was coming, and drew nigh to meet David, David made haste, and ran to the fight to meet the Philistine. And he put his hand into his scrip, and took a stone, and cast it with the sling, and fetching it about struck the Philistine in the forehead: and the stone was fixed in his forehead, and he fell on his face upon the earth. And David prevailed over the Philistine, with a sling and a stone, and he struck, and slew the Philistine. And as David had no sword in his hand, He ran, and stood over the Philistine, and took his sword, and drew it out of the sheath, and slew him, and cut off his head. And the Philistines seeing that their champion was dead, fled away.\par
\switchcolumn*

\LA
\begin{Resp}\Rsign Ego te tuli de domo patris tui, dicit Dóminus, et pósui te páscere gregem pópuli mei: \Star{} Et fui tecum in ómnibus ubicúmque ambulásti, firmans regnum tuum in ætérnum.\end{Resp}
\begin{Resp}\Vsign Fecíque tibi nomen grande, juxta nomen magnórum, qui sunt in terra: et réquiem dedi tibi ab ómnibus inimícis tuis.\end{Resp}
\begin{Resp}\Rsign Et fui tecum in ómnibus ubicúmque ambulásti, firmans regnum tuum in ætérnum.\end{Resp}
\begin{Resp}\Vsign Glória Patri, et Fílio, \Star{} et Spirítui Sancto.\end{Resp}
\begin{Resp}\Rsign Et fui tecum in ómnibus ubicúmque ambulásti, firmans regnum tuum in ætérnum.\end{Resp}
\switchcolumn
\EN
\begin{Resp}\Rsign Thus saith the Lord: I took thee out of thy father's house, and appointed thee to be ruler over My people, over Israel. \Star{} And I was with thee whithersoever thou wentest, to establish thy kingdom for ever.\end{Resp}
\begin{Resp}\Vsign And I have made thee a great name, like unto the name of the great men that are in the earth and have caused thee to rest from all thine enemies.\end{Resp}
\begin{Resp}\Rsign And I was with thee whithersoever thou wentest, to establish thy kingdom for ever.\end{Resp}
\begin{Resp}\Vsign Glory be to the Father, and to the Son, \Star{} and to the Holy Ghost.\end{Resp}
\begin{Resp}\Rsign And I was with thee whithersoever thou wentest, to establish thy kingdom for ever.\end{Resp}
\switchcolumn*

\end{paracol}

\Office{Feria quarta infra Hebdomadam IV post Octavam Pentecostes}{Wednesday of the Fourth Week after Pentecost}{Feria}
\addcontentsline{toc}{subsection}{Feria quarta infra Hebdomadam IV post Octavam Pentecostes \textit{— Wednesday of the Fourth Week after Pentecost}}
\begin{paracol}{2}
\LA
\LessonHead{Lectio I}
\switchcolumn
\EN
\LessonHead{Lesson I}
\switchcolumn*

\LA
\LessonTitle{De libro primo Regum}
\switchcolumn
\EN
\LessonTitle{Lesson from the first book of Samuel}
\switchcolumn*

\LA
\Cite{1 Reg 18:6-8}
\switchcolumn
\EN
\Cite{1 Sam 18:6-8}
\switchcolumn*

\LA
\noindent Porro, cum reverterétur, percússo Philisthǽo, David, egréssæ sunt mulíeres de univérsis úrbibus Israël cantántes chorósque ducéntes in occúrsum Saul regis in týmpanis lætítiæ et in sistris. Et præcinébant mulíeres ludéntes atque dicéntes: Percússit Saul mille, et David decem míllia. Irátus est autem Saul nimis, et displícuit in óculis ejus sermo iste, dixítque: Dedérunt David decem míllia et mihi mille dedérunt; quid ei súperest, nisi solum regnum?\par
\switchcolumn
\EN
\noindent Now when David returned, after he slew the Philistine, the women came out of all the cities of Israel, singing and dancing, to meet king Saul, with timbrels of joy, and cornets. And the women sung as they played, and they said: Saul slew his thousands, and David his ten thousands. And Saul was exceeding angry, and this word was displeasing in his eyes, and he said: They have given David ten thousands, and to me they have given but a thousand; what can he have more but the kingdom?\par
\switchcolumn*

\LA
\begin{Resp}\Rsign Peccávi super númerum arénæ maris, et multiplicáta sunt peccáta mea: et non sum dignus vidére altitúdinem cæli præ multitúdine iniquitátis meæ: quóniam irritávi iram tuam, \Star{} Et malum coram te feci.\end{Resp}
\begin{Resp}\Vsign Quóniam iniquitátem meam ego cognósco: et delíctum meum contra me est semper, quia tibi soli peccávi.\end{Resp}
\begin{Resp}\Rsign Et malum coram te feci.\end{Resp}
\switchcolumn
\EN
\begin{Resp}\Rsign My sins are many, yea, they are more in number than the sands of the sea; I am not worthy to look up toward heaven because of the multitude of my iniquities; for I have provoked thee to anger; \Star{} And done evil in thy sight.\end{Resp}
\begin{Resp}\Vsign For I acknowledge my transgression, and my sin is ever before me, for against thee only have I sinned.\end{Resp}
\begin{Resp}\Rsign And done evil in thy sight.\end{Resp}
\switchcolumn*

\LA
\LessonHead{Lectio II}
\switchcolumn
\EN
\LessonHead{Lesson II}
\switchcolumn*

\LA
\Cite{1 Reg 18:9-13}
\switchcolumn
\EN
\Cite{1 Sam 18:9-13}
\switchcolumn*

\LA
\noindent Non rectis ergo óculis Saul aspiciébat David a die illa et deínceps. Post diem autem álteram invásit spíritus Dei malus Saul, et prophetábat in médio domus suæ; David autem psallébat manu sua sicut per síngulos dies. Tenebátque Saul lánceam et misit eam putans quod confígere posset David cum paríete; et declinávit David a fácie ejus secúndo. Et tímuit Saul David eo quod Dóminus esset cum eo, et a se recessísset. Amóvit ergo eum Saul a se et fecit eum tribúnum super mille viros; et egrediebátur et intrábat in conspéctu pópuli.\par
\switchcolumn
\EN
\noindent And Saul did not look on David with a good eye from that day and forward. And the day after the evil spirit from God came upon Saul, and he prophesied in the midst of his house. And David played with his hand as at other times. And Saul held a spear in his hand, And threw it, thinking to nail David to the wall: and David stept aside out of his presence twice. And Saul feared David, because the Lord was with him, and was departed from himself. Therefore Saul removed him from him, and made him a captain over a thousand men, and he went out and came in before the people.\par
\switchcolumn*

\LA
\begin{Resp}\Rsign Exaudísti, Dómine, oratiónem servi tui, ut ædificárem templum nómini tuo: \Star{} Bénedic et sanctífica domum istam in sempitérnum, Deus Israël.\end{Resp}
\begin{Resp}\Vsign Dómine, qui custódis pactum cum servis tuis, qui ámbulant coram te in toto corde suo.\end{Resp}
\begin{Resp}\Rsign Bénedic et sanctífica domum istam in sempitérnum, Deus Israël.\end{Resp}
\switchcolumn
\EN
\begin{Resp}\Rsign O Lord, Thou hast hearkened unto the prayer of thy servant, that I might build a temple unto thy Name, \Star{} O God of Israel, bless Thou, and hallow this house for ever.\end{Resp}
\begin{Resp}\Vsign O Lord, Who keepest covenant with thy servants that walk before thee in all their heart.\end{Resp}
\begin{Resp}\Rsign O God of Israel, bless Thou, and hallow this house for ever.\end{Resp}
\switchcolumn*

\LA
\LessonHead{Lectio III}
\switchcolumn
\EN
\LessonHead{Lesson III}
\switchcolumn*

\LA
\Cite{1 Reg 18:14-17}
\switchcolumn
\EN
\Cite{1 Sam 18:14-17}
\switchcolumn*

\LA
\noindent In ómnibus quoque viis suis David prudénter agébat, et Dóminus erat cum eo. Vidit ítaque Saul quod prudens esset nimis et cœpit cavére eum. Omnis autem Israël et Juda diligébat David; ipse enim ingrediebátur et egrediebátur ante eos. Dixítque Saul ad David: Ecce fília mea major Merob, ipsam dabo tibi uxórem; tantúmmodo esto vir fortis et præliáre bella Dómini. Saul autem reputábat dicens: Non sit manus mea in eum, sed sit super eum manus Philisthinórum.\par
\switchcolumn
\EN
\noindent And David behaved wisely in all his ways, and the Lord was with him. And Saul saw that he was exceeding prudent, and began to beware of him. But all Israel and Juda loved David, for he came in and went out before them. And Saul said to David: Behold my elder daughter Merob, her will I give thee to wife: only be a valiant man, and fight the battles of the Lord. Now Saul said within himself: Let not my hand be upon him, but let the hands of the Philistines be upon him.\par
\switchcolumn*

\LA
\begin{Resp}\Rsign Audi, Dómine, hymnum et oratiónem, quam servus tuus orat coram te hódie, ut sint óculi tui apérti, et aures tuæ inténtæ, \Star{} Super domum istam die ac nocte.\end{Resp}
\begin{Resp}\Vsign Réspice, Dómine, de sanctuário tuo, et de excélso cælórum habitáculo.\end{Resp}
\begin{Resp}\Rsign Super domum istam die ac nocte.\end{Resp}
\begin{Resp}\Vsign Glória Patri, et Fílio, \Star{} et Spirítui Sancto.\end{Resp}
\begin{Resp}\Rsign Super domum istam die ac nocte.\end{Resp}
\switchcolumn
\EN
\begin{Resp}\Rsign Hearken, O Lord, unto the cry and to the prayer which thy servant prayeth before thee today, that thine eyes may be open and thine ears attend; \Star{} Toward this house day and night.\end{Resp}
\begin{Resp}\Vsign Look down from thine high and holy place, O Lord, even from heaven thy dwelling.\end{Resp}
\begin{Resp}\Rsign Toward this house, day and night.\end{Resp}
\begin{Resp}\Vsign Glory be to the Father, and to the Son, \Star{} and to the Holy Ghost.\end{Resp}
\begin{Resp}\Rsign Toward this house, day and night.\end{Resp}
\switchcolumn*

\end{paracol}

\Office{Feria quinta infra Hebdomadam IV post Octavam Pentecostes}{Thursday of the Fourth Week after Pentecost}{Feria}
\addcontentsline{toc}{subsection}{Feria quinta infra Hebdomadam IV post Octavam Pentecostes \textit{— Thursday of the Fourth Week after Pentecost}}
\begin{paracol}{2}
\LA
\LessonHead{Lectio I}
\switchcolumn
\EN
\LessonHead{Lesson I}
\switchcolumn*

\LA
\LessonTitle{De libro primo Regum}
\switchcolumn
\EN
\LessonTitle{Lesson from the first book of Samuel}
\switchcolumn*

\LA
\Cite{1 Reg 19:1-3}
\switchcolumn
\EN
\Cite{1 Sam 19:1-3}
\switchcolumn*

\LA
\noindent Locútus est autem Saul ad Jónatham fílium suum et ad omnes servos suos ut occíderent David. Porro Jónathas fílius Saul diligébat David valde. Et indicávit Jónathas David dicens: Quærit Saul pater meus occídere te; quaprópter obsérva te, quæso, mane, et manébis clam, et abscondéris. Ego autem egrédiens stabo juxta patrem meum in agro ubicúmque fúeris; et ego loquar de te ad patrem meum, et quodcúmque vídero, nuntiábo tibi.\par
\switchcolumn
\EN
\noindent And Saul spoke to Jonathan his son and to all his servants, that they should kill David. But Jonathan the son of Saul loved David exceedingly. And Jonathan told David, saying: Saul my father seeketh to kill thee: wherefore look to thyself, I beseech thee, in the morning, and thou shalt abide in a secret place and shalt be hid. And I will go out and stand beside my father in the field where thou art: and I will speak of thee to my father, and whatsoever I shall see, I will tell thee.\par
\switchcolumn*

\LA
\begin{Resp}\Rsign Præparáte corda vestra Dómino, et servíte illi soli: \Star{} Et liberábit vos de mánibus inimicórum vestrórum.\end{Resp}
\begin{Resp}\Vsign Convertímini ad eum in toto corde vestro, et auférte deos aliénos de médio vestri.\end{Resp}
\begin{Resp}\Rsign Et liberábit vos de mánibus inimicórum vestrórum.\end{Resp}
\switchcolumn
\EN
\begin{Resp}\Rsign Prepare your hearts unto the Lord, and serve Him Only \Star{} And He will deliver you out of the hand of your enemies.\end{Resp}
\begin{Resp}\Vsign Return unto Him with all your hearts, and put away the strange gods from among you.\end{Resp}
\begin{Resp}\Rsign And He will deliver you out of the hand of your enemies.\end{Resp}
\switchcolumn*

\LA
\LessonHead{Lectio II}
\switchcolumn
\EN
\LessonHead{Lesson II}
\switchcolumn*

\LA
\Cite{1 Reg 19:4-6}
\switchcolumn
\EN
\Cite{1 Sam 19:4-6}
\switchcolumn*

\LA
\noindent Locútus est ergo Jónathas de David bona ad Saul patrem suum dixítque ad eum: Ne pecces, rex, in servum tuum David, quia non peccávit tibi, et ópera ejus bona sunt tibi valde. Et pósuit ánimam suam in manu sua et percússit Philisthǽum, et fecit Dóminus salútem magnam univérso Israéli. Vidísti et lætátus es; quare ergo peccas in sánguine innóxio interfíciens David, qui est absque culpa? Quod cum audísset Saul, placátus voce Jónathæ jurávit: Vivit Dóminus, quia non occidétur.\par
\switchcolumn
\EN
\noindent And Jonathan spoke good things of David to Saul his father: and said to him: Sin not, O king, against thy servant, David, because he hath not sinned against thee, and his works are very good towards thee. And he put his life in his hand, and slew the Philistine, and the Lord wrought great salvation for all Israel. Thou sawest it and didst rejoice. Why therefore wilt thou sin against innocent blood by killing David, who is without fault? And when Saul heard this he was appeased with the words of Jonathan, and swore: As the Lord liveth he shall not be slain.\par
\switchcolumn*

\LA
\begin{Resp}\Rsign Deus ómnium exaudítor est: ipse misit Angelum suum, et tulit me de óvibus patris mei: \Star{} Et unxit me unctióne misericórdiæ suæ.\end{Resp}
\begin{Resp}\Vsign Dóminus, qui erípuit me de ore leónis, et de manu béstiæ liberávit me.\end{Resp}
\begin{Resp}\Rsign Et unxit me unctióne misericórdiæ suæ.\end{Resp}
\switchcolumn
\EN
\begin{Resp}\Rsign God, Which heareth all, even He sent His Angel, and took me from keeping my father's sheep, and \Star{} Anointed me with the oil of His mercy.\end{Resp}
\begin{Resp}\Vsign The Lord That delivered me out of the mouth of the lion, and out of the paw of the bear\end{Resp}
\begin{Resp}\Rsign And anointed me with the oil of His mercy.\end{Resp}
\switchcolumn*

\LA
\LessonHead{Lectio III}
\switchcolumn
\EN
\LessonHead{Lesson III}
\switchcolumn*

\LA
\Cite{1 Reg 19:8-10}
\switchcolumn
\EN
\Cite{1 Sam 19:8-10}
\switchcolumn*

\LA
\noindent Motum est autem rursum bellum, et egréssus David pugnávit advérsum Philísthiim percussítque eos plaga magna, et fugérunt a fácie ejus. Et factus est spíritus Dómini malus in Saul. Sedébat autem in domo sua et tenébat lánceam; porro David psallébat manu sua. Nisúsque est Saul confígere David láncea in paríete, et declinávit David a fácie Saul; láncea autem, casso vúlnere, perláta est in paríetem. Et David fugit et salvátus est nocte illa.\par
\switchcolumn
\EN
\noindent And the war began again, and David went out and fought against the Philistines, and defeated them with a great slaughter, and they fled from his face. And the evil spirit from the Lord came upon Saul, and he sat in his house, and held a spear in his hand: and David played with his hand. And Saul endeavoured to nail David to the wall with his spear. And David slipt away out of the presence of Saul: and the spear missed him, and was fastened in the wall, and David fled and escaped that night.\par
\switchcolumn*

\LA
\begin{Resp}\Rsign Dóminus, qui erípuit me de ore leónis, et de manu béstiæ liberávit me, \Star{} Ipse me erípiet de mánibus inimicórum meórum.\end{Resp}
\begin{Resp}\Vsign Misit Deus misericórdiam suam et veritátem suam: ánimam meam erípuit de médio catulórum leónum.\end{Resp}
\begin{Resp}\Rsign Ipse me erípiet de mánibus inimicórum meórum.\end{Resp}
\begin{Resp}\Vsign Glória Patri, et Fílio, \Star{} et Spirítui Sancto.\end{Resp}
\begin{Resp}\Rsign Ipse me erípiet de mánibus inimicórum meórum.\end{Resp}
\switchcolumn
\EN
\begin{Resp}\Rsign The Lord That delivered me out of the mouth of the lion, and out of the paw of the bear \Star{} He will deliver me out of the hand of mine enemies.\end{Resp}
\begin{Resp}\Vsign God hath sent forth His mercy and His truth, and delivered my soul from among the lion's whelps.\end{Resp}
\begin{Resp}\Rsign He will deliver me out of the hand of mine enemies.\end{Resp}
\begin{Resp}\Vsign Glory be to the Father, and to the Son, \Star{} and to the Holy Ghost.\end{Resp}
\begin{Resp}\Rsign He will deliver me out of the hand of mine enemies.\end{Resp}
\switchcolumn*

\end{paracol}

\Office{Feria sexta infra Hebdomadam IV post Octavam Pentecostes}{Friday of the Fourth Week after Pentecost}{Feria}
\addcontentsline{toc}{subsection}{Feria sexta infra Hebdomadam IV post Octavam Pentecostes \textit{— Friday of the Fourth Week after Pentecost}}
\begin{paracol}{2}
\LA
\LessonHead{Lectio I}
\switchcolumn
\EN
\LessonHead{Lesson I}
\switchcolumn*

\LA
\LessonTitle{De libro primo Regum}
\switchcolumn
\EN
\LessonTitle{Lesson from the first book of Samuel}
\switchcolumn*

\LA
\Cite{1 Reg 20:1-2}
\switchcolumn
\EN
\Cite{1 Sam 20:1-2}
\switchcolumn*

\LA
\noindent Fugit autem David de Najoth, quæ est in Rámatha, veniénsque locútus est coram Jónatha: Quid feci? quæ est iníquitas mea, et quod peccátum meum in patrem tuum, quia quærit ánimam meam? Qui dixit ei: Absit, non moriéris; neque enim fáciet pater meus quidquam grande vel parvum nisi prius indicáverit mihi; hunc ergo celávit me pater meus sermónem tantúmmodo? Nequáquam erit istud.\par
\switchcolumn
\EN
\noindent But David fled from Najoth, which is in Ramatha, and came and said to Jonathan: What have I done? what is my iniquity, and what is my sin against thy father, that he seeketh my life? And he said to him: God forbid, thou shalt not die: for my father will do nothing great or little, without first telling me: hath then my father hid this word only from me? no, this shall not be.\par
\switchcolumn*

\LA
\begin{Resp}\Rsign Percússit Saul mille, et David decem míllia: \Star{} Quia manus Dómini erat cum illo: percússit Philisthǽum, et ábstulit oppróbrium ex Israël.\end{Resp}
\begin{Resp}\Vsign Nonne iste est David, de quo canébant in choro, dicéntes: Saul percússit mille, et David decem míllia?\end{Resp}
\begin{Resp}\Rsign Quia manus Dómini erat cum illo: percússit Philisthǽum, et ábstulit oppróbrium ex Israël.\end{Resp}
\switchcolumn
\EN
\begin{Resp}\Rsign Saul hath slain his thousands, and David his ten thousands. \Star{} Because the hand of the Lord was with him, he smote the Philistine, and took away the reproach from Israel.\end{Resp}
\begin{Resp}\Vsign Is not this David? Did they not sing one to another of him in dances, saying: Saul hath slain his thousands, and David his ten thousands?\end{Resp}
\begin{Resp}\Rsign Because the hand of the Lord was with him, he smote the Philistine, and took away the reproach from Israel.\end{Resp}
\switchcolumn*

\LA
\LessonHead{Lectio II}
\switchcolumn
\EN
\LessonHead{Lesson II}
\switchcolumn*

\LA
\Cite{1 Reg 20:3-4}
\switchcolumn
\EN
\Cite{1 Sam 20:3-4}
\switchcolumn*

\LA
\noindent Et jurávit rursum Davídi. Et ille ait: Scit profécto pater tuus quia invéni grátiam in óculis tuis et dicet: Nésciat hoc Jónathas, ne forte tristétur. Quinímmo vivit Dóminus, et vivit ánima tua, quia uno tantum (ut ita dicam) gradu ego morsque divídimur. Et ait Jónathas ad David: Quodcúmque díxerit mihi ánima tua, fáciam tibi.\par
\switchcolumn
\EN
\noindent And he swore again to David. And David said: thy father certainly knoweth that I have found grace in thy sight, and he will say: Let not Jonathan know this, lest he be grieved. But truly as the Lord liveth, and thy soul liveth, there is but one step (as I may say) between me and death. And Jonathan said to David: Whatsoever thy soul shall say to me, I will do for thee.\par
\switchcolumn*

\LA
\begin{Resp}\Rsign Montes Gélboë, nec ros nec plúvia véniant super vos, \Star{} Ubi cecidérunt fortes Israël.\end{Resp}
\begin{Resp}\Vsign Omnes montes, qui estis in circúitu ejus, vísitet Dóminus: a Gélboë autem tránseat.\end{Resp}
\begin{Resp}\Rsign Ubi cecidérunt fortes Israël.\end{Resp}
\switchcolumn
\EN
\begin{Resp}\Rsign Ye mountains of Gilboa, let there be no dew, neither let there be rain upon you \Star{} For there are the mighty of Israel fallen.\end{Resp}
\begin{Resp}\Vsign All ye mountains that stand round about, the Lord look upon you but let Him pass by Gilboa\end{Resp}
\begin{Resp}\Rsign For there are the mighty of Israel fallen.\end{Resp}
\switchcolumn*

\LA
\LessonHead{Lectio III}
\switchcolumn
\EN
\LessonHead{Lesson III}
\switchcolumn*

\LA
\Cite{1 Reg 20:5-7}
\switchcolumn
\EN
\Cite{1 Sam 20:5-7}
\switchcolumn*

\LA
\noindent Dixit autem David ad Jónathan: Ecce caléndæ sunt crástino, et ego ex more sedére sóleo juxta regem ad vescéndum; dimítte ergo me, ut abscóndar in agro usque ad vésperam diéi tértiæ. Si respíciens requisíerit me pater tuus, respondébis ei: Rogávit me David ut iret celériter in Béthlehem civitátem suam, quia víctimæ solémnes ibi sunt univérsis contribúlibus suis. Si díxerit: Bene, pax erit servo tuo; si autem fúerit irátus, scito quia compléta est malítia ejus.\par
\switchcolumn
\EN
\noindent And David said to Jonathan: Behold tomorrow is the new moon, and I according to custom am wont to sit beside the king to eat: let me go then that I may be hid in the field till the evening of the third day. If thy father look and inquire for me, thou shalt answer him: David asked me that he might run to Bethlehem his own city: because there are solemn sacrifices there for all his tribe. If he shall say, It is well: thy servant shall have peace: but if he be angry, know that his malice is come to its height.\par
\switchcolumn*

\LA
\begin{Resp}\Rsign Ego te tuli de domo patris tui, dicit Dóminus, et pósui te páscere gregem pópuli mei: \Star{} Et fui tecum in ómnibus ubicúmque ambulásti, firmans regnum tuum in ætérnum.\end{Resp}
\begin{Resp}\Vsign Fecíque tibi nomen grande, juxta nomen magnórum, qui sunt in terra: et réquiem dedi tibi ab ómnibus inimícis tuis.\end{Resp}
\begin{Resp}\Rsign Et fui tecum in ómnibus ubicúmque ambulásti, firmans regnum tuum in ætérnum.\end{Resp}
\begin{Resp}\Vsign Glória Patri, et Fílio, \Star{} et Spirítui Sancto.\end{Resp}
\begin{Resp}\Rsign Et fui tecum in ómnibus ubicúmque ambulásti, firmans regnum tuum in ætérnum.\end{Resp}
\switchcolumn
\EN
\begin{Resp}\Rsign Thus saith the Lord: I took thee out of thy father's house, and appointed thee to be ruler over My people, over Israel. \Star{} And I was with thee whithersoever thou wentest, to establish thy kingdom for ever.\end{Resp}
\begin{Resp}\Vsign And I have made thee a great name, like unto the name of the great men that are in the earth, and have caused thee to rest from all thine enemies.\end{Resp}
\begin{Resp}\Rsign And I was with thee whithersoever thou wentest, to establish thy kingdom for ever.\end{Resp}
\begin{Resp}\Vsign Glory be to the Father, and to the Son, \Star{} and to the Holy Ghost.\end{Resp}
\begin{Resp}\Rsign And I was with thee whithersoever thou wentest, to establish thy kingdom for ever.\end{Resp}
\switchcolumn*

\end{paracol}

\Office{Sabbato infra Hebdomadam IV post Octavam Pentecostes}{Saturday of the Fourth Week after Pentecost}{Feria}
\addcontentsline{toc}{subsection}{Sabbato infra Hebdomadam IV post Octavam Pentecostes \textit{— Saturday of the Fourth Week after Pentecost}}
\begin{paracol}{2}
\LA
\LessonHead{Lectio I}
\switchcolumn
\EN
\LessonHead{Lesson I}
\switchcolumn*

\LA
\LessonTitle{De libro primo Regum}
\switchcolumn
\EN
\LessonTitle{Lesson from the first book of Samuel}
\switchcolumn*

\LA
\Cite{1 Reg 21:1-3}
\switchcolumn
\EN
\Cite{1 Sam 21:1-3}
\switchcolumn*

\LA
\noindent Venit autem David in Nobe ad Achímelech sacerdótem. Et obstúpuit Achímelech, eo quod venísset David, et dixit ei: Quare tu solus, et nullus est tecum? Et ait David ad Achímelech sacerdótem: Rex præcépit mihi sermónem et dixit: Nemo sciat rem propter quam missus es a me, et cujúsmodi præcépta tibi déderim; nam et púeris condíxi in illum et illum locum. Nunc ergo, si quid habes ad manum, vel quinque panes, da mihi, aut quidquid invéneris.\par
\switchcolumn
\EN
\noindent And David came to Nobe to Achimelech the priest: and Achimelech was astonished at David's coming. And he said to him: Why art thou alone, and no man with thee? And David said to Achimelech the priest: The king hath commanded me a business, and said: Let no man know the thing for which thou art sent by me, and what manner of commands I have given thee: and I have appointed my servants to such and such a place. Now therefore if thou have any thing at hand, though it were but five loaves, give me, or whatsoever thou canst find.\par
\switchcolumn*

\LA
\begin{Resp}\Rsign Peccávi super númerum arénæ maris, et multiplicáta sunt peccáta mea: et non sum dignus vidére altitúdinem cæli præ multitúdine iniquitátis meæ: quóniam irritávi iram tuam, \Star{} Et malum coram te feci.\end{Resp}
\begin{Resp}\Vsign Quóniam iniquitátem meam ego cognósco: et delíctum meum contra me est semper, quia tibi soli peccávi.\end{Resp}
\begin{Resp}\Rsign Et malum coram te feci.\end{Resp}
\switchcolumn
\EN
\begin{Resp}\Rsign My sins are many, yea, they are more in number than the sands of the sea; I am not worthy to look up toward heaven because of the multitude of my iniquities; for I have provoked thee to anger; \Star{} And done evil in thy sight.\end{Resp}
\begin{Resp}\Vsign For I acknowledge my transgression, and my sin is ever before me, for against thee only have I sinned.\end{Resp}
\begin{Resp}\Rsign And done evil in thy sight.\end{Resp}
\switchcolumn*

\LA
\LessonHead{Lectio II}
\switchcolumn
\EN
\LessonHead{Lesson II}
\switchcolumn*

\LA
\Cite{1 Reg 21:4-6}
\switchcolumn
\EN
\Cite{1 Sam 21:4-6}
\switchcolumn*

\LA
\noindent Et respóndens sacérdos ad David ait illi: Non hábeo láicos panes ad manum, sed tantum panem sanctum; si mundi sunt púeri, máxime a muliéribus? Et respóndit David sacerdóti et dixit ei: Equidem, si de muliéribus ágitur, continúimus nos ab heri et nudiustértius, quando egrediebámur, et fuérunt vasa puerórum sancta. Porro via hæc pollúta est, sed et ipsa hódie sanctificábitur in vasis. Dedit ergo ei sacérdos sanctificátum panem; neque enim erat ibi panis, nisi tantum panes propositiónis, qui subláti fúerant a fácie Dómini, ut poneréntur panes cálidi.\par
\switchcolumn
\EN
\noindent And the priest answered David, saying: I have no common bread at hand, but only holy bread, if the young men be clean, especially from women? And David answered the priest, and said to him: Truly, as to what concerneth women, we have refrained ourselves from yesterday and the day before, when we came out, and the vessels of the young men were holy. Now this way is defiled, but it shall also be sanctified this day in the vessels. The priest therefore gave him hallowed bread: for there was no bread there, but only the loaves of proposition, which had been taken away from before the face of the Lord, that hot loaves might be set up.\par
\switchcolumn*

\LA
\begin{Resp}\Rsign Exaudísti, Dómine, oratiónem servi tui, ut ædificárem templum nómini tuo: \Star{} Bénedic et sanctífica domum istam in sempitérnum, Deus Israël.\end{Resp}
\begin{Resp}\Vsign Dómine, qui custódis pactum cum servis tuis, qui ámbulant coram te in toto corde suo.\end{Resp}
\begin{Resp}\Rsign Bénedic et sanctífica domum istam in sempitérnum, Deus Israël.\end{Resp}
\switchcolumn
\EN
\begin{Resp}\Rsign O Lord, Thou hast hearkened unto the prayer of thy servant, that I might build a temple unto thy Name, \Star{} O God of Israel, bless Thou, and hallow this house for ever.\end{Resp}
\begin{Resp}\Vsign O Lord, Who keepest covenant with thy servants that walk before thee in all their heart.\end{Resp}
\begin{Resp}\Rsign O God of Israel, bless Thou, and hallow this house for ever.\end{Resp}
\switchcolumn*

\LA
\LessonHead{Lectio III}
\switchcolumn
\EN
\LessonHead{Lesson III}
\switchcolumn*

\LA
\Cite{1 Reg 21:7-9}
\switchcolumn
\EN
\Cite{1 Sam 21:7-9}
\switchcolumn*

\LA
\noindent Erat autem ibi vir quidam de servis Saul in die illa intus in tabernáculo Dómini, et nomen ejus Doëg Idumǽus potentíssimus pastórum Saul. Dixit autem David ad Achímelech: Si habes hic ad manum hastam aut gládium? quia gládium meum et arma mea non tuli mecum; sermo enim regis urgébat. Et dixit sacérdos: Ecce hic gládius Góliath Philisthǽi, quem percussísti in Valle terebínthi: est involútus pállio post ephod. Si istum vis tóllere, tolle, neque enim hic est álius absque eo. Et ait David: Non est huic alter símilis: da mihi eum.\par
\switchcolumn
\EN
\noindent Now a certain man of the servants of Saul was there that day, within the tabernacle of the Lord: and his name was Doeg, an Edomite, the chiefest of Saul's herdsmen. And David said to Achimelech: Hast thou here at hand a spear, or a sword? for I brought not my own sword, nor my own weapons with me, for the king's business required haste. And the priest said: Lo, here is the sword of Goliath the Philistine whom thou slewest in the valley of Terebinth, wrapped up in a cloth behind the ephod: if thou wilt take this, take it, for here is no other but this. And David said: There is none like that, give it me.\par
\switchcolumn*

\LA
\begin{Resp}\Rsign Audi, Dómine, hymnum et oratiónem, quam servus tuus orat coram te hódie, ut sint óculi tui apérti, et aures tuæ inténtæ, \Star{} Super domum istam die ac nocte.\end{Resp}
\begin{Resp}\Vsign Réspice, Dómine, de sanctuário tuo, et de excélso cælórum habitáculo.\end{Resp}
\begin{Resp}\Rsign Super domum istam die ac nocte.\end{Resp}
\begin{Resp}\Vsign Glória Patri, et Fílio, \Star{} et Spirítui Sancto.\end{Resp}
\begin{Resp}\Rsign Super domum istam die ac nocte.\end{Resp}
\switchcolumn
\EN
\begin{Resp}\Rsign Hearken, O Lord, unto the cry and to the prayer which thy servant prayeth before thee today, that thine eyes may be open and thine ears attend; \Star{} Toward this house day and night.\end{Resp}
\begin{Resp}\Vsign Look down from thine high and holy place, O Lord, even from heaven thy dwelling.\end{Resp}
\begin{Resp}\Rsign Toward this house, day and night.\end{Resp}
\begin{Resp}\Vsign Glory be to the Father, and to the Son, \Star{} and to the Holy Ghost.\end{Resp}
\begin{Resp}\Rsign Toward this house, day and night.\end{Resp}
\switchcolumn*

\end{paracol}

\Office{Dominica V Post Pentecosten}{Fifth Sunday after Pentecost}{Semiduplex}
\addcontentsline{toc}{subsection}{Dominica V Post Pentecosten \textit{— Fifth Sunday after Pentecost}}
\begin{paracol}{2}
\LA
\LessonHead{Lectio I}
\switchcolumn
\EN
\LessonHead{Lesson I}
\switchcolumn*

\LA
\LessonTitle{Incipit liber secúndus Regum}
\switchcolumn
\EN
\LessonTitle{Lesson from the second book of Samuel}
\switchcolumn*

\LA
\Cite{2 Reg 1:1-4}
\switchcolumn
\EN
\Cite{2 Sam 1:1-4}
\switchcolumn*

\LA
\noindent Factum est autem, postquam mórtuus est Saul, ut David reverterétur a cæde Amalec et manéret in Síceleg duos dies. In die autem tértia appáruit homo véniens de castris Saul veste conscíssa et púlvere conspérsus caput et, ut venit ad David, cécidit super fáciem suam et adorávit. Dixítque ad eum David: Unde venis? Qui ait ad eum: De castris Israël fugi. Et dixit ad eum David: Quod est verbum quod factum est? Indica mihi. Qui ait: Fugit pópulus ex prǽlio, et multi corruéntes e pópulo mórtui sunt; sed et Saul et Jónathas fílius ejus interiérunt.\par
\switchcolumn
\EN
\noindent Now it came to pass, after Saul was dead, that David returned from the slaughter of the Amalecites, and abode two days in Siceleg. And on the third day, there appeared a man who came out of Saul's camp, with his garments rent, and dust strewed on his head: and when he came to David, he fell upon his face, and adored. And David said to him: From whence comest thou? And he said to him: I am fled out of the camp of Israel. And David said unto him: What is the matter that is come to pass? tell me. He said: The people are fled from the battle, and many of the people are fallen and dead: moreover Saul and Jonathan his son are slain.\par
\switchcolumn*

\LA
\begin{Resp}\Rsign Præparáte corda vestra Dómino, et servíte illi soli: \Star{} Et liberábit vos de mánibus inimicórum vestrórum.\end{Resp}
\begin{Resp}\Vsign Convertímini ad eum in toto corde vestro, et auférte deos aliénos de médio vestri.\end{Resp}
\begin{Resp}\Rsign Et liberábit vos de mánibus inimicórum vestrórum.\end{Resp}
\switchcolumn
\EN
\begin{Resp}\Rsign Prepare your hearts unto the Lord, and serve Him Only \Star{} And He will deliver you out of the hand of your enemies.\end{Resp}
\begin{Resp}\Vsign Return unto Him with all your hearts, and put away the strange gods from among you.\end{Resp}
\begin{Resp}\Rsign And He will deliver you out of the hand of your enemies.\end{Resp}
\switchcolumn*

\LA
\LessonHead{Lectio II}
\switchcolumn
\EN
\LessonHead{Lesson II}
\switchcolumn*

\LA
\Cite{2 Reg 1:5-10}
\switchcolumn
\EN
\Cite{2 Sam 1:5-10}
\switchcolumn*

\LA
\noindent Dixítque David ad adolescéntem qui nuntiábat ei: Unde scis quia mórtuus est Saul et Jónathas fílius ejus? Et ait adoléscens qui nuntiábat ei: Casu veni in montem Gélboë, et Saul incumbébat super hastam suam. Porro currus et équites appropinquábant ei, et convérsus post tergum suum vidénsque me vocávit; cui cum respondíssem: Adsum, dixit mihi: Quisnam es tu? Et ajo ad eum: Amalecítes ego sum. Et locútus est mihi: Sta super me et intérfice me, quóniam tenent me angústiæ, et adhuc tota ánima mea in me est. Stansque super eum, occídi illum, sciébam enim quod vívere non póterat post ruínam; et tuli diadéma, quod erat in cápite ejus et armíllam de brácchio illíus et áttuli ad te dóminum meum huc.\par
\switchcolumn
\EN
\noindent And David said to the young man that told him: How knowest thou that Saul and Jonathan his son, are dead? And the young man that told him, said: I came by chance upon mount Gelboe, and Saul leaned upon his spear: and the chariots and horsemen drew nigh unto him, And looking behind him, and seeing me, he called me. And I answered, Here am I. And he said to me: Who art thou? And I said to him: I am an Amalecite. And he said to me: Stand over me, and kill me: for anguish is come upon me, and as yet my whole life is in me. So standing over him, I killed him: for I knew that he could not live after the fall: and I took the diadem that was on his head, and the bracelet that was on his arm and have brought them hither to thee, my lord.\par
\switchcolumn*

\LA
\begin{Resp}\Rsign Deus ómnium exaudítor est: ipse misit Angelum suum, et tulit me de óvibus patris mei: \Star{} Et unxit me unctióne misericórdiæ suæ.\end{Resp}
\begin{Resp}\Vsign Dóminus, qui erípuit me de ore leónis, et de manu béstiæ liberávit me.\end{Resp}
\begin{Resp}\Rsign Et unxit me unctióne misericórdiæ suæ.\end{Resp}
\switchcolumn
\EN
\begin{Resp}\Rsign God, Which heareth all, even He sent His Angel, and took me from keeping my father's sheep, and \Star{} Anointed me with the oil of His mercy.\end{Resp}
\begin{Resp}\Vsign The Lord That delivered me out of the mouth of the lion, and out of the paw of the bear\end{Resp}
\begin{Resp}\Rsign And anointed me with the oil of His mercy.\end{Resp}
\switchcolumn*

\LA
\LessonHead{Lectio III}
\switchcolumn
\EN
\LessonHead{Lesson III}
\switchcolumn*

\LA
\Cite{2 Reg 1:11-15}
\switchcolumn
\EN
\Cite{2 Sam 1:11-15}
\switchcolumn*

\LA
\noindent Apprehéndens autem David, vestiménta sua scidit, omnésque viri, qui erant cum eo, et planxérunt, et flevérunt, et jejunavérunt usque ad vésperam super Saul, et super Jónathan fílium ejus, et super pópulum Dómini et super domum Israël, eo quod corruíssent gládio. Dixítque David ad júvenem qui nuntiáverat ei: Unde es tu? Qui respóndit: Fílius hóminis ádvenæ Amalecítæ ego sum. Et ait ad eum David: Quare non timuísti míttere manum tuam, ut occíderes christum Dómini? Vocánsque David unum de púeris suis ait: Accédens írrue in eum. Qui percússit illum, et mórtuus est.\par
\switchcolumn
\EN
\noindent Then David took hold of his garments and rent them, and likewise all the men that were with him. And they mourned, and wept, and fasted until evening for Saul, and for Jonathan his son, and for the people of the Lord, and for the house of Israel, because they were fallen by the sword. And David said to the young man that told him: Whence art thou? He answered: I am the son of a stranger of Amalee. David said to him: Why didst thou not fear to put out thy hand to kill the Lord's anointed? And David calling one of his servants, said: Go near and fall upon him. And he struck him so that he died.\par
\switchcolumn*

\LA
\begin{Resp}\Rsign Dóminus, qui erípuit me de ore leónis, et de manu béstiæ liberávit me, \Star{} Ipse me erípiet de mánibus inimicórum meórum.\end{Resp}
\begin{Resp}\Vsign Misit Deus misericórdiam suam et veritátem suam: ánimam meam erípuit de médio catulórum leónum.\end{Resp}
\begin{Resp}\Rsign Ipse me erípiet de mánibus inimicórum meórum.\end{Resp}
\begin{Resp}\Vsign Glória Patri, et Fílio, \Star{} et Spirítui Sancto.\end{Resp}
\begin{Resp}\Rsign Ipse me erípiet de mánibus inimicórum meórum.\end{Resp}
\switchcolumn
\EN
\begin{Resp}\Rsign The Lord That delivered me out of the mouth of the lion, and out of the paw of the bear \Star{} He will deliver me out of the hand of mine enemies.\end{Resp}
\begin{Resp}\Vsign God hath sent forth His mercy and His truth, and delivered my soul from among the lion's whelps.\end{Resp}
\begin{Resp}\Rsign He will deliver me out of the hand of mine enemies.\end{Resp}
\begin{Resp}\Vsign Glory be to the Father, and to the Son, \Star{} and to the Holy Ghost.\end{Resp}
\begin{Resp}\Rsign He will deliver me out of the hand of mine enemies.\end{Resp}
\switchcolumn*

\LA
\LessonHead{Lectio IV}
\switchcolumn
\EN
\LessonHead{Lesson IV}
\switchcolumn*

\LA
\LessonTitle{Ex libro Morálium sancti Gregórii Papæ}
\switchcolumn
\EN
\LessonTitle{From the Book of Moral (Reflections upon Job), written by Pope St. Gregory the Great.}
\switchcolumn*

\LA
\Source{Liber 4, Cap. 3 et 4}
\switchcolumn
\EN
\Source{Bk. iv. ch. 3, 4}
\switchcolumn*

\LA
\noindent Quid est quod David, qui retribuéntibus sibi mala non réddidit, cum Saul et Jónathas bello occúmberent, Gélboë móntibus maledíxit, dicens: Montes Gélboë, nec ros nec plúvia véniant super vos: neque sint agri primitiárum: quia ibi abjéctus est clýpeus fórtius, clýpeus Saul, quasi non esset unctus óleo? Quid est quod Jeremías, cum prædicatiónem suam cérneret audiéntium difficultáte præpedíri, maledíxit dicens: Maledíctus vir, qui annuntiávit patri meo, dicens: Natus est tibi puer másculus?\par
\switchcolumn
\EN
\noindent Thus was it that David, who rewarded no evil to them that did evil to him (Ps. vii. 5), when Saul and Jonathan had fallen in battle, cursed the mountains of Gilboa, saying Ye mountains of Gilboa, let there be no dew, neither let there be rain upon you, nor fields of offerings; for there the shield of the mighty is vilely cast away, the shield of Saul, as though he had not been anointed with oil. Why was it that Jeremiah, when he saw that his preaching was thrown away upon his hearers, cursed and said Cursed be the man who brought tidings to my father, saying: A man-child is born unto thee? (xx. 15.)\par
\switchcolumn*

\LA
\begin{Resp}\Rsign Percússit Saul mille, et David decem míllia: \Star{} Quia manus Dómini erat cum illo: percússit Philisthǽum, et ábstulit oppróbrium ex Israël.\end{Resp}
\begin{Resp}\Vsign Nonne iste est David, de quo canébant in choro, dicéntes: Saul percússit mille, et David decem míllia?\end{Resp}
\begin{Resp}\Rsign Quia manus Dómini erat cum illo: percússit Philisthǽum, et ábstulit oppróbrium ex Israël.\end{Resp}
\switchcolumn
\EN
\begin{Resp}\Rsign Saul hath slain his thousands, and David his ten thousands. \Star{} Because the hand of the Lord was with him, he smote the Philistine, and took away the reproach from Israel.\end{Resp}
\begin{Resp}\Vsign Is not this David? Did they not sing one to another of him in dances, saying: Saul hath slain his thousands, and David his ten thousands?\end{Resp}
\begin{Resp}\Rsign Because the hand of the Lord was with him, he smote the Philistine, and took away the reproach from Israel.\end{Resp}
\switchcolumn*

\LA
\LessonHead{Lectio V}
\switchcolumn
\EN
\LessonHead{Lesson V}
\switchcolumn*

\LA
\noindent Quid ergo montes Gélboë, Saul moriénte, deliquérunt, quátenus in eos nec ros nec plúvia cáderet et ab omni eos viriditátis gérmine senténtiæ sermo siccáret? Sed quia Gélboë interpretátur decúrsus, per Saul autem unctum et mórtuum mors nostri mediatóris exprímitur; non immérito per Gélboë montes supérba Judæórum corda signántur, quæ dum in hujus mundi desidériis défluunt, in Christi, id est, uncti se morte miscuérunt: et quia in eis unctus rex corporáliter móritur, ipsi ab omni grátiæ rore siccántur.\par
\switchcolumn
\EN
\noindent He had the mountains of Gilboa to do with the death of Saul, that they should be condemned to have dew fall on them no more, neither rain, but should wither away, barren of the green glory of the springtime? But this word Gilboa signified bubbling fountain, and the death of Saul the Anointed of God is a type of the death of our Anointed Mediator. Thus we find in the mountains of Gilboa no unfit image of the proud hearts of the Jews, which had their spring in earthly desires, and took part in the death of the Anointed Saviour. And since among them their Anointed Monarch met His death, the dew of grace is upon them no more.\par
\switchcolumn*

\LA
\begin{Resp}\Rsign Montes Gélboë, nec ros nec plúvia véniant super vos, \Star{} Ubi cecidérunt fortes Israël.\end{Resp}
\begin{Resp}\Vsign Omnes montes, qui estis in circúitu ejus, vísitet Dóminus: a Gélboë autem tránseat.\end{Resp}
\begin{Resp}\Rsign Ubi cecidérunt fortes Israël.\end{Resp}
\switchcolumn
\EN
\begin{Resp}\Rsign Ye mountains of Gilboa, let there be no dew, neither let there be rain upon you \Star{} For there are the mighty of Israel fallen.\end{Resp}
\begin{Resp}\Vsign All ye mountains that stand round about, the Lord look upon you but let Him pass by Gilboa\end{Resp}
\begin{Resp}\Rsign For there are the mighty of Israel fallen.\end{Resp}
\switchcolumn*

\LA
\LessonHead{Lectio VI}
\switchcolumn
\EN
\LessonHead{Lesson VI}
\switchcolumn*

\LA
\noindent De quibus et bene dícitur, ut agri primitiárum esse non possint. Supérbæ quippe Hebræórum mentes primitívos fructus non ferunt: quia in Redemptóris advéntu ex parte máxima in perfídia remanéntes, primórdia fídei sequi noluérunt. Sancta namque Ecclésia in primítiis suis multitúdine Géntium fœcundáta, vix in mundi fine Judǽos quos invénerit, súscipit, et extréma cólligens, eos quasi relíquias frugum ponit.\par
\switchcolumn
\EN
\noindent And well is it said of them Let there be upon you no fields of offerings. The proud minds of the Hebrews bear yet no offering. Since the coming of their Redeemer, the most part of them remain still without belief in Him, and refuse to follow the promise of their ancient faith. The Holy Church hath borne for her first-born, holy unto the Lord, a multitude of the Gentiles, and will, but in the end of the world, embrace such Jews as she then shall find, and present them as the last gatherings of her harvest.\par
\switchcolumn*

\LA
\begin{Resp}\Rsign Ego te tuli de domo patris tui, dicit Dóminus, et pósui te páscere gregem pópuli mei: \Star{} Et fui tecum in ómnibus ubicúmque ambulásti, firmans regnum tuum in ætérnum.\end{Resp}
\begin{Resp}\Vsign Fecíque tibi nomen grande, juxta nomen magnórum, qui sunt in terra: et réquiem dedi tibi ab ómnibus inimícis tuis.\end{Resp}
\begin{Resp}\Rsign Et fui tecum in ómnibus ubicúmque ambulásti, firmans regnum tuum in ætérnum.\end{Resp}
\begin{Resp}\Vsign Glória Patri, et Fílio, \Star{} et Spirítui Sancto.\end{Resp}
\begin{Resp}\Rsign Et fui tecum in ómnibus ubicúmque ambulásti, firmans regnum tuum in ætérnum.\end{Resp}
\switchcolumn
\EN
\begin{Resp}\Rsign Thus saith the Lord: I took thee out of thy father's house, and appointed thee to be ruler over My people, over Israel. \Star{} And I was with thee whithersoever thou wentest, to establish thy kingdom for ever.\end{Resp}
\begin{Resp}\Vsign And I have made thee a great name, like unto the name of the great men that are in the earth, and have caused thee to rest from all thine enemies.\end{Resp}
\begin{Resp}\Rsign And I was with thee whithersoever thou wentest, to establish thy kingdom for ever.\end{Resp}
\begin{Resp}\Vsign Glory be to the Father, and to the Son, \Star{} and to the Holy Ghost.\end{Resp}
\begin{Resp}\Rsign And I was with thee whithersoever thou wentest, to establish thy kingdom for ever.\end{Resp}
\switchcolumn*

\LA
\LessonHead{Lectio VII}
\switchcolumn
\EN
\LessonHead{Lesson VII}
\switchcolumn*

\LA
\LessonTitle{Léctio sancti Evangélii secúndum Matthǽum}
\switchcolumn
\EN
\LessonTitle{From the Holy Gospel according to Matthew}
\switchcolumn*

\LA
\Cite{Matt 5:20-24}
\switchcolumn
\EN
\Cite{Matt 5:20-24}
\switchcolumn*

\LA
\noindent In illo témpore: Dixit Jesus discípulis suis: Nisi abundáverit justítia vestra plus quam scribárum et pharisæórum, non intrábitis in regnum cælórum. Et réliqua.\par
\switchcolumn
\EN
\noindent At that time, Jesus said unto His disciples: Unless your righteousness shall exceed the righteousness of the Scribes and Pharisees, ye shall in no case enter into the kingdom of heaven. And so on.\par
\switchcolumn*

\LA
\LessonTitle{Homilía sancti Augustíni Epíscopi}
\switchcolumn
\EN
\LessonTitle{Homily by St. Augustine, Bishop of Hippo}
\switchcolumn*

\LA
\Source{Liber 1 de Sermóne Dómini in monte, cap. 9}
\switchcolumn
\EN
\Source{Bk. I on the Lord's Sermon on the Mount, ch. 9}
\switchcolumn*

\LA
\noindent Justítia pharisæórum est, ut non occídant: justítia eórum, qui intratúri sunt in regnum cælórum, ut non irascántur sine causa. Mínimum est ergo, non occídere; et qui illud sólverit, mínimus vocábitur in regno cælórum. Qui autem illud impléverit, ut non occídat, non contínuo magnus erit, et idóneus regno cælórum; sed tamen ascéndit áliquem gradum: perficiétur autem, si nec irascátur sine causa; quod si perfécerit, multo remótior erit ab homicídio. Quaprópter qui docet ut non irascámur, non solvit legem ne occidámus, et in corde, dum non iráscimur, innocéntiam custodiámus.\par
\switchcolumn
\EN
\noindent Thou shalt not kill, is of the righteousness of the Pharisees; Thou shalt not be angry with thy brother without a cause, is of the righteousness of them which shall enter into the kingdom of heaven. The least therefore is: “Thou shalt not kill, and whosoever shall break this commandment, he shall be called the least in the kingdom of heaven.” (v. 19.) But whosoever shall do it, and not kill, he is not therefore great, and meet for the kingdom of heaven; albeit, he hath risen a step; but he will have gotten farther, if he be not angry with his brother without a cause, which, if he do, he will be the farther off from manslaughter. Wherefore, He Which teacheth us that we are not to be angry without a cause, destroyeth not the law, Thou shalt not kill, but rather fulfilleth and increaseth it, making us not only to be free of the sin of outward killing, but also clean of anger within.\par
\switchcolumn*

\LA
\begin{Resp}\Rsign Peccávi super númerum arénæ maris, et multiplicáta sunt peccáta mea: et non sum dignus vidére altitúdinem cæli præ multitúdine iniquitátis meæ: quóniam irritávi iram tuam, \Star{} Et malum coram te feci.\end{Resp}
\begin{Resp}\Vsign Quóniam iniquitátem meam ego cognósco: et delíctum meum contra me est semper, quia tibi soli peccávi.\end{Resp}
\begin{Resp}\Rsign Et malum coram te feci.\end{Resp}
\switchcolumn
\EN
\begin{Resp}\Rsign My sins are many, yea, they are more in number than the sands of the sea; I am not worthy to look up toward heaven because of the multitude of my iniquities; for I have provoked thee to anger; \Star{} And done evil in thy sight.\end{Resp}
\begin{Resp}\Vsign For I acknowledge my transgression, and my sin is ever before me, for against thee only have I sinned.\end{Resp}
\begin{Resp}\Rsign And done evil in thy sight.\end{Resp}
\switchcolumn*

\LA
\LessonHead{Lectio VIII}
\switchcolumn
\EN
\LessonHead{Lesson VIII}
\switchcolumn*

\LA
\noindent Gradus ítaque sunt in istis peccátis: ut primo quisque irascátur, et eum motum retíneat corde concéptum. Jam si extórserit vocem indignántis ipsa commótio, non significántem áliquid, sed illum ánimi motum ipsa eruptióne testántem, qua feriátur ille, cui iráscitur; plus est útique, quam si surgens ira siléntio premerétur. Si vero non solum vox indignántis audiátur, sed étiam verbum, quo jam certam ejus vituperatiónem, in quem profértur, desígnet et notet; quis dúbitet, ámplius hoc esse, quam si solus indignatiónis sonus ederétur?\par
\switchcolumn
\EN
\noindent On sins of this kind there are diverse steps. First, there is the swelling feeling of anger. When this feeling appeareth in a man's heart, he keepeth it. Then the inward disturbance wringeth forth words of indignation, not themselves meaning aught, but showing the trouble of him who is provoked. And this is something more than anger kept covered under silence. Next, this audible outburst of indignation may contain direct and open reviling of him who hath roused it. And it cannot be doubted that this is something more than an empty cry of anger.\par
\switchcolumn*

\LA
\begin{Resp}\Rsign Duo Séraphim clamábant alter ad álterum: \Star{} Sanctus, sanctus, sanctus Dóminus Deus Sábaoth: * Plena est omnis terra glória ejus.\end{Resp}
\begin{Resp}\Vsign Tres sunt qui testimónium dant in cælo: Pater, Verbum, et Spíritus Sanctus: et hi tres unum sunt.\end{Resp}
\begin{Resp}\Rsign Sanctus, sanctus, sanctus Dóminus Deus Sábaoth.\end{Resp}
\begin{Resp}\Vsign Glória Patri, et Fílio, \Star{} et Spirítui Sancto.\end{Resp}
\begin{Resp}\Rsign Plena est omnis terra glória ejus.\end{Resp}
\switchcolumn
\EN
\begin{Resp}\Rsign One Seraph cried unto another: \Star{} Holy, Holy, Holy is the Lord God of hosts; the whole earth is full of His glory.\end{Resp}
\begin{Resp}\Vsign There are Three That bear record in heaven: the Father, the Word, and the Holy Ghost, and these Three are One.\end{Resp}
\begin{Resp}\Rsign Holy, Holy, Holy is the Lord God of hosts.\end{Resp}
\begin{Resp}\Vsign Glory be to the Father, and to the Son, \Star{} and to the Holy Ghost.\end{Resp}
\begin{Resp}\Rsign The whole earth is full of His glory.\end{Resp}
\switchcolumn*

\LA
\LessonHead{Lectio IX}
\switchcolumn
\EN
\LessonHead{Lesson IX}
\switchcolumn*

\LA
\noindent Vide nunc étiam tres reátus, judícii, concílii, et gehénnæ ignis. Nam in judício adhuc defensióni datur locus. In concílio autem, quamquam et judícium esse sóleat, tamen quia interésse áliquid hoc loco fatéri cogit ipsa distínctio, vidétur ad concílium pertinére senténtiæ prolátio; quando non jam cum ipso reo ágitur, utrum damnándus sit; sed inter se, qui júdicant, cónferunt quo supplício damnári opórteat, quem constat esse damnándum. Gehénna vero ignis, nec damnatiónem habet dúbiam, sicut judícium, nec damnáti pœnam, sicut concílium: in gehénna quippe ignis, certa est damnátio, et pœna damnáti.\par
\switchcolumn
\EN
\noindent Behold here the three degrees of guilt open respectively to the judgment, to the council, and to hellfire. In the judgment, there is still place for defence. In the council, albeit this also is in a sense a judgment, yet we may suppose this distinction from the judgment proper, that the council pronounceth sentence, not as the result of a trial whereat the accused is present, but as the result of a consultation among the judges, to what punishment he is to be sentenced of whom it is already established that he is guilty. When we get to hell-fire, there remaineth no longer any doubt about condemnation, as in the judgment, and no longer any doubt about sentence, as in the council. In hellfire the condemnation and the pain of him that is condemned are alike certain.\par
\switchcolumn*

\LA
\TeDeum{Te Deum}
\switchcolumn
\EN
\TeDeum{Te Deum}
\switchcolumn*

\end{paracol}

\Office{Feria secunda infra Hebdomadam V post Octavam Pentecostes}{Monday of the Fifth Week after Pentecost}{Feria}
\addcontentsline{toc}{subsection}{Feria secunda infra Hebdomadam V post Octavam Pentecostes \textit{— Monday of the Fifth Week after Pentecost}}
\begin{paracol}{2}
\LA
\LessonHead{Lectio I}
\switchcolumn
\EN
\LessonHead{Lesson I}
\switchcolumn*

\LA
\LessonTitle{De libro secúndo Regum}
\switchcolumn
\EN
\LessonTitle{Lesson from the second book of Samuel}
\switchcolumn*

\LA
\Cite{2 Reg 2:1-4}
\switchcolumn
\EN
\Cite{2 Sam 2:1-4}
\switchcolumn*

\LA
\noindent Igitur post hæc consúluit David Dóminum dicens: Num ascéndam in unam de civitátibus Juda? Et ait Dóminus ad eum: Ascénde. Dixítque David: Quo ascéndam? Et respóndit ei: In Hebron. Ascéndit ergo David et duæ uxóres ejus Achínoam Jezrahelítes et Abígail uxor Nabal Carméli; sed et viros, qui erant cum eo, duxit David síngulos cum domo sua; et mansérunt in óppidis Hebron. Venerúntque viri Juda, et unxérunt ibi David, ut regnáret super domum Juda.\par
\switchcolumn
\EN
\noindent And after these things David consulted the Lord, saying: Shall I go up into one of the cities of Juda? And the Lord said to him: Go up. And David said: Whither shall I go up? And he answered him: Into Hebron. So David went up, and his two wives, Achinoam the Jezrahelitess, and Abigail the wife of Nabal of Carmel: And the men also that were with him, David brought up every man with his household: and they abode in the towns of Hebron. And the men of Juda came, and anointed David there, to be king over the house of Juda.\par
\switchcolumn*

\LA
\begin{Resp}\Rsign Recordáre, Dómine, testaménti tui, et dic Angelo percutiénti: Cesset jam manus tua, \Star{} Ut non desolétur terra, et ne perdas omnem ánimam vivam.\end{Resp}
\begin{Resp}\Vsign Ego sum qui peccávi, ego qui iníque egi: isti qui oves sunt, quid fecérunt? Avertátur, óbsecro, furor tuus, Dómine, a pópulo tuo.\end{Resp}
\begin{Resp}\Rsign Ut non desolétur terra, et ne perdas omnem ánimam vivam.\end{Resp}
\switchcolumn
\EN
\begin{Resp}\Rsign Remember, O Lord, thy covenant, and say unto the destroying Angel: Stay now thine hand; \Star{} That the land be not utterly laid waste, and that thou destroy not every living soul.\end{Resp}
\begin{Resp}\Vsign Even I it is that have sinned, and done evil indeed; but these sheep, what have they done? Let thine anger, I pray thee, O Lord, be turned away from thy people.\end{Resp}
\begin{Resp}\Rsign That the land be not utterly laid waste, and that Thou destroy not every living soul.\end{Resp}
\switchcolumn*

\LA
\LessonHead{Lectio II}
\switchcolumn
\EN
\LessonHead{Lesson II}
\switchcolumn*

\LA
\Cite{2 Reg 2:4-7}
\switchcolumn
\EN
\Cite{2 Sam 2:4-7}
\switchcolumn*

\LA
\noindent Et nuntiátum est David, quod viri Jabes Gálaad sepelíssent Saul. Misit ergo David núntios ad viros Jabes Gálaad dixítque ad eos: Benedícti vos a Dómino, qui fecístis misericórdiam hanc cum dómino vestro Saul, et sepelístis eum. Et nunc retríbuet vobis quidem Dóminus misericórdiam et veritátem, sed et ego reddam grátiam, eo quod fecístis verbum istud. Conforténtur manus vestræ et estóte fílii fortitúdinis; licet enim mórtuus sit dóminus vester Saul, tamen me unxit domus Juda in regem sibi.\par
\switchcolumn
\EN
\noindent And it was told David, that the men of Jabes Galaad had buried Saul. David therefore sent messengers to the men of Jabes Galaad, and said to them: Blessed be you to the Lord, who have shown this mercy to your master Saul, and have buried him. And now the Lord surely will render you mercy and truth, and I also will, requite you for this good turn, because you have done this thing. Let your hands be strengthened, and be ye men of valour: for although your master Saul be dead, yet the house of Juda hath anointed me to be their king\par
\switchcolumn*

\LA
\begin{Resp}\Rsign Exaudísti, Dómine, oratiónem servi tui, ut ædificárem templum nómini tuo: \Star{} Bénedic et sanctífica domum istam in sempitérnum, Deus Israël.\end{Resp}
\begin{Resp}\Vsign Dómine, qui custódis pactum cum servis tuis, qui ámbulant coram te in toto corde suo.\end{Resp}
\begin{Resp}\Rsign Bénedic et sanctífica domum istam in sempitérnum, Deus Israël.\end{Resp}
\switchcolumn
\EN
\begin{Resp}\Rsign O Lord, Thou hast hearkened unto the prayer of thy servant, that I might build a temple unto thy Name, \Star{} O God of Israel, bless Thou, and hallow this house for ever.\end{Resp}
\begin{Resp}\Vsign O Lord, Who keepest covenant with thy servants that walk before thee in all their heart.\end{Resp}
\begin{Resp}\Rsign O God of Israel, bless Thou, and hallow this house for ever.\end{Resp}
\switchcolumn*

\LA
\LessonHead{Lectio III}
\switchcolumn
\EN
\LessonHead{Lesson III}
\switchcolumn*

\LA
\Cite{2 Reg 2:8-11}
\switchcolumn
\EN
\Cite{2 Sam 2:8-11}
\switchcolumn*

\LA
\noindent Abner autem fílius Ner, princeps exércitus Saul, tulit Isboseth fílium Saul, et circumdúxit eum per castra, regémque constítuit super Gálaad, et super Géssuri, et super Jézrahel, et super Ephraim, et super Bénjamin, et super Israël univérsum. Quadragínta annórum erat Isboseth fílius Saul cum regnáre cœpísset super Israël, et duóbus annis regnávit. Sola autem domus Juda sequebátur David. Et fuit númerus diérum, quos commorátus est David ímperans in Hebron super domum Juda, septem annórum et sex ménsium.\par
\switchcolumn
\EN
\noindent But Abner the son of Ner, general of Saul's army, took Isboseth the son of Saul, and led him about through the camp? And made him king over Galaad, and, over Gessuri, and over Jezrahel, and over Ephraim, and over Benjamin, and over all Israel. Isboseth the son of Saul was forty years old when he began to reign over, Israel, and he reigned two years: and only the house of Juda followed David. And the number of the days that David abode, reigning in Hebron over the house of Juda, was seven years and six months.\par
\switchcolumn*

\LA
\begin{Resp}\Rsign Audi, Dómine, hymnum et oratiónem, quam servus tuus orat coram te hódie, ut sint óculi tui apérti, et aures tuæ inténtæ, \Star{} Super domum istam die ac nocte.\end{Resp}
\begin{Resp}\Vsign Réspice, Dómine, de sanctuário tuo, et de excélso cælórum habitáculo.\end{Resp}
\begin{Resp}\Rsign Super domum istam die ac nocte.\end{Resp}
\begin{Resp}\Vsign Glória Patri, et Fílio, \Star{} et Spirítui Sancto.\end{Resp}
\begin{Resp}\Rsign Super domum istam die ac nocte.\end{Resp}
\switchcolumn
\EN
\begin{Resp}\Rsign Hearken, O Lord, unto the cry and to the prayer which thy servant prayeth before thee today, that thine eyes may be open and thine ears attend; \Star{} Toward this house day and night.\end{Resp}
\begin{Resp}\Vsign Look down from thine high and holy place, O Lord, even from heaven thy dwelling.\end{Resp}
\begin{Resp}\Rsign Toward this house, day and night.\end{Resp}
\begin{Resp}\Vsign Glory be to the Father, and to the Son, \Star{} and to the Holy Ghost.\end{Resp}
\begin{Resp}\Rsign Toward this house, day and night.\end{Resp}
\switchcolumn*

\end{paracol}

\Office{Feria tertia infra Hebdomadam V post Octavam Pentecostes}{Tuesday of the Fifth Week after Pentecost}{Feria}
\addcontentsline{toc}{subsection}{Feria tertia infra Hebdomadam V post Octavam Pentecostes \textit{— Tuesday of the Fifth Week after Pentecost}}
\begin{paracol}{2}
\LA
\LessonHead{Lectio I}
\switchcolumn
\EN
\LessonHead{Lesson I}
\switchcolumn*

\LA
\LessonTitle{De libro secúndo Regum}
\switchcolumn
\EN
\LessonTitle{Lesson from the second book of Samuel}
\switchcolumn*

\LA
\Cite{2 Reg 3:6-10}
\switchcolumn
\EN
\Cite{2 Sam 3:6-10}
\switchcolumn*

\LA
\noindent Cum ergo esset prǽlium inter domum Saul et domum David, Abner fílius Ner regébat domum Saul. Fúerat autem Sauli concubína nómine Respha, fília Aja. Dixítque Isboseth ad Abner: Quare ingréssus es ad concubínam patris mei? Qui irátus nimis propter verba Isboseth ait: Numquid caput canis ego sum advérsum Judam hódie, qui fécerim misericórdiam super domum Saul patris tui, et super fratres et próximos ejus, et non trádidi te in manus David? Et tu requisísti in me quod argúeres pro mulíere hódie? Hæc fáciat Deus Abner, et hæc addat ei, nisi, quómodo jurávit Dóminus David, sic fáciam cum eo, ut transferátur regnum de domo Saul, et elevétur thronus David super Israël et super Judam a Dan usque Bersabée.\par
\switchcolumn
\EN
\noindent Now while there was war between the house of Saul and the house of David, Abner the son of Ner ruled the house of Saul. And Saul had a concubine named Respha, the daughter of Aia. And Isboseth said to Abner: Why didst thou go in to my father's concubine? And he was exceedingly angry for the words of Isboseth, and said: Am I a dog's head against Juda this day, who have shown mercy to the house of Saul thy father, and to his brethren and friends, and have not delivered thee into the hands of David, and hast thou sought this day against me to charge me with a matter concerning a woman? So do God to Abner, and more also, unless as the Lord hath sworn to David, so I do to him, That the kingdom be translated from the house of Saul, and the throne of David be set up over Israel, and over Juda from Dan to Bersabee.\par
\switchcolumn*

\LA
\begin{Resp}\Rsign Dómine, si convérsus fúerit pópulus tuus, et oráverit ad sanctuárium tuum: \Star{} Tu exáudies de cælo, Dómine, et líbera eos de mánibus inimicórum suórum.\end{Resp}
\begin{Resp}\Vsign Si peccáverit in te pópulus tuus, et convérsus égerit pœniténtiam, veniénsque oráverit in isto loco.\end{Resp}
\begin{Resp}\Rsign Tu exáudies de cælo, Dómine, et líbera eos de mánibus inimicórum suórum.\end{Resp}
\switchcolumn
\EN
\begin{Resp}\Rsign Lord, when thy people shall turn again to thee, and shall pray unto thee in this house; \Star{} Then hear Thou in heaven, O Lord, and deliver them out of the hand of their enemies.\end{Resp}
\begin{Resp}\Vsign If thy people sin against thee, and turn again, and repent, and come and pray unto thee in this house.\end{Resp}
\begin{Resp}\Rsign Then hear Thou in heaven, O Lord, and deliver them out of the hand of their enemies.\end{Resp}
\switchcolumn*

\LA
\LessonHead{Lectio II}
\switchcolumn
\EN
\LessonHead{Lesson II}
\switchcolumn*

\LA
\Cite{2 Reg 3:12-16}
\switchcolumn
\EN
\Cite{2 Sam 3:12-16}
\switchcolumn*

\LA
\noindent Misit ergo Abner núntios ad David pro se dicéntes: Cujus est terra? et ut loqueréntur: Fac mecum amicítias, et erit manus mea tecum et redúcam ad te univérsum Israël. Qui ait: Optime: ego fáciam tecum amicítias, sed unam rem peto a te, dicens: Non vidébis fáciem meam, ántequam addúxeris Michol fíliam Saul; et sic vénies, et vidébis me. Misit autem David núntios ad Isboseth fílium Saul dicens: Redde uxórem meam Michol, quam despóndi mihi centum præpútiis Philísthiim. Misit ergo Isboseth et tulit eam a viro suo Pháltiel, fílio Lais. Sequebatúrque eam vir suus, plorans usque Bahúrim. Et dixit ad eum Abner: Vade et revértere. Qui revérsus est.\par
\switchcolumn
\EN
\noindent Abner therefore sent messengers to David for himself, saying: Whose is the land? and that they should say: Make a league with me, and my hand shall be with thee: and I will bring all Israel to thee. And he said: Very well: I will make a league with thee: but one thing I require of thee, saying: Thou shalt not see my face before thou bring Michol the daughter of Saul: and so thou shalt come, and see me. And David sent messengers to Isboseth the son of Saul, saying: Restore my wife Michol, whom I espoused to me for a hundred foreskins of the Philistines. And Isboseth sent, and took her from her husband Phaltiel, the son of Lais. And her husband followed her, weeping as far as Bahurim: and Abner said to him: Go and return. And he returned.\par
\switchcolumn*

\LA
\begin{Resp}\Rsign Factum est, dum tólleret Dóminus Elíam per túrbinem in cælum, \Star{} Eliséus clamábat, dicens: Pater mi, pater mi, currus Israël, et auríga ejus.\end{Resp}
\begin{Resp}\Vsign Cumque pérgerent, et incedéntes sermocinaréntur, ecce currus ígneus et equi ígnei divisérunt utrúmque, et ascéndit Elías per túrbinem in cælum.\end{Resp}
\begin{Resp}\Rsign Eliséus clamábat, dicens: Pater mi, pater mi, currus Israël, et auríga ejus.\end{Resp}
\switchcolumn
\EN
\begin{Resp}\Rsign And it came to pass when the Lord would take up Elijah into heaven by a whirlwind \Star{} Elisha cried, and said: My father, my father, the chariot of Israel, and the horsemen thereof.\end{Resp}
\begin{Resp}\Vsign And as they still went on, and talked, behold, there appeared a chariot of fire, and horses of fire, and parted them both asunder, and Elijah went up by a whirlwind into heaven.\end{Resp}
\begin{Resp}\Rsign Elisha cried, and said: My father, my father, the chariot of Israel, and the horsemen thereof.\end{Resp}
\switchcolumn*

\LA
\LessonHead{Lectio III}
\switchcolumn
\EN
\LessonHead{Lesson III}
\switchcolumn*

\LA
\Cite{2 Reg 3:17-21}
\switchcolumn
\EN
\Cite{2 Sam 3:17-21}
\switchcolumn*

\LA
\noindent Sermónem quoque íntulit Abner ad senióres Israël, dicens: Tam heri quam nudiustértius quærebátis David, ut regnáret super vos. Nunc ergo fácite, quóniam Dóminus locútus est ad David, dicens: In manu servi mei David salvábo pópulum meum Israël de manu Philísthiim, et ómnium inimicórum ejus. Locútus est autem Abner étiam ad Bénjamin. Et ábiit ut loquerétur ad David in Hebron ómnia quæ placúerant Israéli et univérso Bénjamin. Venítque ad David in Hebron cum vigínti viris, et fecit David Abner et viris ejus, qui vénerant cum eo, convívium. Et dixit Abner ad David: Surgam ut cóngregem ad te, dóminum meum regem, omnem Israël.\par
\switchcolumn
\EN
\noindent Abner also spoke to the ancients of Israel, saying: Both yesterday and the day before you sought for David that he might reign over you. Now then do it: because the Lord hath spoken to David, saying: By the hand of my servant David I will save my people Israel from the hands of the Philistines, and of all their enemies. And Abner spoke also to Benjamin. And he went to speak to David in Hebron all that seemed good to Israel, and to all Benjamin. And he came to David in Hebron with twenty men: and David made a feast for Abner, and his men that came with him. And Abner said to David: I will rise, that I may gather all Israel unto thee, my lord the king.\par
\switchcolumn*

\LA
\begin{Resp}\Rsign Ego te tuli de domo patris tui, dicit Dóminus, et pósui te páscere gregem pópuli mei: \Star{} Et fui tecum in ómnibus ubicúmque ambulásti, firmans regnum tuum in ætérnum.\end{Resp}
\begin{Resp}\Vsign Fecíque tibi nomen grande, juxta nomen magnórum, qui sunt in terra: et réquiem dedi tibi ab ómnibus inimícis tuis.\end{Resp}
\begin{Resp}\Rsign Et fui tecum in ómnibus ubicúmque ambulásti, firmans regnum tuum in ætérnum.\end{Resp}
\begin{Resp}\Vsign Glória Patri, et Fílio, \Star{} et Spirítui Sancto.\end{Resp}
\begin{Resp}\Rsign Et fui tecum in ómnibus ubicúmque ambulásti, firmans regnum tuum in ætérnum.\end{Resp}
\switchcolumn
\EN
\begin{Resp}\Rsign Thus saith the Lord: I took thee out of thy father's house, and appointed thee to be ruler over My people, over Israel. \Star{} And I was with thee whithersoever thou wentest, to establish thy kingdom for ever.\end{Resp}
\begin{Resp}\Vsign And I have made thee a great name, like unto the name of the great men that are in the earth, and have caused thee to rest from all thine enemies.\end{Resp}
\begin{Resp}\Rsign And I was with thee whithersoever thou wentest, to establish thy kingdom for ever.\end{Resp}
\begin{Resp}\Vsign Glory be to the Father, and to the Son, \Star{} and to the Holy Ghost.\end{Resp}
\begin{Resp}\Rsign And I was with thee whithersoever thou wentest, to establish thy kingdom for ever.\end{Resp}
\switchcolumn*

\end{paracol}

\Office{Feria quarta infra Hebdomadam V post Octavam Pentecostes}{Wednesday of the Fifth Week after Pentecost}{Feria}
\addcontentsline{toc}{subsection}{Feria quarta infra Hebdomadam V post Octavam Pentecostes \textit{— Wednesday of the Fifth Week after Pentecost}}
\begin{paracol}{2}
\LA
\LessonHead{Lectio I}
\switchcolumn
\EN
\LessonHead{Lesson I}
\switchcolumn*

\LA
\LessonTitle{De libro secúndo Regum}
\switchcolumn
\EN
\LessonTitle{Lesson from the second book of Samuel}
\switchcolumn*

\LA
\Cite{2 Reg 4:5-8}
\switchcolumn
\EN
\Cite{2 Sam 4:5-8}
\switchcolumn*

\LA
\noindent Veniéntes ígitur fílii Remmon Berothítæ, Rechab et Báana, ingréssi sunt, fervénte die, domum Isboseth; qui dormiébat super stratum suum merídie. Et ostiária domus purgans tríticum obdormívit. Ingréssi sunt autem domum laténter assuméntes spicas trítici, et percussérunt eum in ínguine Rechab et Báana frater ejus et fugérunt. Cum autem ingréssi fuíssent domum, ille dormiébat super lectum suum in conclávi, et percutiéntes interfecérunt eum; sublatóque cápite ejus abiérunt per viam desérti tota nocte. Et attulérunt caput Isboseth ad David in Hebron dixerúntque ad regem: Ecce caput Isboseth, fílii Saul inimíci tui, qui quærébat ánimam tuam.\par
\switchcolumn
\EN
\noindent And the sons of Remmon the Berothite, Rechab and Baana coming, went into the house of Isboseth in the heat of the day: and he was sleeping upon his bed at noon. And the doorkeeper of the house, who was cleansing wheat, was fallen asleep. And they entered into the house secretly taking ears of corn, and Rechab and Baana his brother stabbed him in the groin, and fled away. For when they came into the house, be was sleeping upon his bed in a parlour, and they struck him and killed him: and taking away his head they went off by the way of the wilderness, walking all night. And they brought the head of Isboseth to David to Hebron: and they said to the king: Behold the head of Isboseth the son of Saul thy enemy who sought thy life: and the Lord hath revenged my lord the king this day of Saul, and of his seed.\par
\switchcolumn*

\LA
\begin{Resp}\Rsign Peccávi super númerum arénæ maris, et multiplicáta sunt peccáta mea: et non sum dignus vidére altitúdinem cæli præ multitúdine iniquitátis meæ: quóniam irritávi iram tuam, \Star{} Et malum coram te feci.\end{Resp}
\begin{Resp}\Vsign Quóniam iniquitátem meam ego cognósco: et delíctum meum contra me est semper, quia tibi soli peccávi.\end{Resp}
\begin{Resp}\Rsign Et malum coram te feci.\end{Resp}
\switchcolumn
\EN
\begin{Resp}\Rsign My sins are many, yea, they are more in number than the sands of the sea; I am not worthy to look up toward heaven because of the multitude of my iniquities; for I have provoked thee to anger; \Star{} And done evil in thy sight.\end{Resp}
\begin{Resp}\Vsign For I acknowledge my transgression, and my sin is ever before me, for against thee only have I sinned.\end{Resp}
\begin{Resp}\Rsign And done evil in thy sight.\end{Resp}
\switchcolumn*

\LA
\LessonHead{Lectio II}
\switchcolumn
\EN
\LessonHead{Lesson II}
\switchcolumn*

\LA
\Cite{2 Reg 4:9-12}
\switchcolumn
\EN
\Cite{2 Sam 4:9-12}
\switchcolumn*

\LA
\noindent Respóndens autem David Rechab et Báana fratri ejus fíliis Remmon Berothítæ dixit ad eos: Vivit Dóminus, qui éruit ánimam meam de omni angústia, quóniam eum, qui annuntiáverat mihi et díxerat: Mórtuus est Saul; qui putábat se próspera nuntiáre, ténui et occídi eum in Síceleg, cui oportébat mercédem dare pro núntio: quanto magis nunc, cum hómines ímpii interfecérunt virum innóxium in domo sua super lectum suum, non quæram sánguinem ejus de manu vestra et áuferam vos de terra? Præcépit ítaque David púeris suis, et interfecérunt eos, præcidentésque manus et pedes eórum suspendérunt eos super piscínam in Hebron; caput autem Isboseth tulérunt et sepeliérunt in sepúlcro Abner in Hebron.\par
\switchcolumn
\EN
\noindent But David answered Rechab, and Baana his brother, the sons of Remmon the Berothite, and said to them: As the Lord liveth, who hath delivered my soul out of all distress, The man that told me, and said: Saul is dead, who thought he brought good tidings, I apprehended, and slew him in Siceleg, who should have been rewarded for his news. How much more now when wicked men have slain an innocent man in his own house, upon his bed, shall I not require his blood at your hand, and take you away from the earth? And David commanded his servants and they slew them: and cutting off their hands and feet, hanged them up over the pool in Hebron: but the head of Isboseth they took and buried in the sepulchre of Abner in Hebron.\par
\switchcolumn*

\LA
\begin{Resp}\Rsign Exaudísti, Dómine, oratiónem servi tui, ut ædificárem templum nómini tuo: \Star{} Bénedic et sanctífica domum istam in sempitérnum, Deus Israël.\end{Resp}
\begin{Resp}\Vsign Dómine, qui custódis pactum cum servis tuis, qui ámbulant coram te in toto corde suo.\end{Resp}
\begin{Resp}\Rsign Bénedic et sanctífica domum istam in sempitérnum, Deus Israël.\end{Resp}
\switchcolumn
\EN
\begin{Resp}\Rsign O Lord, Thou hast hearkened unto the prayer of thy servant, that I might build a temple unto thy Name, \Star{} O God of Israel, bless Thou, and hallow this house for ever.\end{Resp}
\begin{Resp}\Vsign O Lord, Who keepest covenant with thy servants that walk before thee in all their heart.\end{Resp}
\begin{Resp}\Rsign O God of Israel, bless Thou, and hallow this house for ever.\end{Resp}
\switchcolumn*

\LA
\LessonHead{Lectio III}
\switchcolumn
\EN
\LessonHead{Lesson III}
\switchcolumn*

\LA
\Cite{2 Reg 5:1-7}
\switchcolumn
\EN
\Cite{2 Sam 5:1-7}
\switchcolumn*

\LA
\noindent Et venérunt univérsæ tribus Israël ad David in Hebron dicéntes: Ecce nos os tuum et caro tua sumus. Sed et heri et nudiustértius, cum esset Saul rex super nos, tu eras edúcens et redúcens Israël; dixit autem Dóminus ad te: Tu pasces pópulum meum Israël et tu eris dux super Israël. Venérunt quoque et senióres Israël ad regem in Hebron, et percússit cum eis rex David fœdus in Hebron coram Dómino, unxerúntque David in regem super Israël. Fílius trigínta annórum erat David, cum regnáre cœpísset, et quadragínta annis regnávit. In Hebron regnávit super Judam septem annis et sex ménsibus, in Jerúsalem autem regnávit trigínta tribus annis super omnem Israël et Judam. Et ábiit rex, et omnes viri, qui erant cum eo, in Jerúsalem ad Jebusǽum habitatórem terræ, dictúmque est David ab eis: Non ingrediéris huc nisi abstúleris cæcos et claudos dicéntes: Non ingrediétur David huc. Cepit autem David arcem Sion, hæc est cívitas David.\par
\switchcolumn
\EN
\noindent Then all the tribes of Israel came to David in Hebron, saying: Behold we are thy bone and thy flesh. Moreover yesterday also and the day before, when Saul was king over us, thou wast he that did lead out and bring in Israel: and the Lord said to thee: Thou shalt feed my people Israel, and thou shalt be prince over Israel. The ancients also of Israel came to the king to Hebron, and king David made a league with them in Hebron before the Lord: and they anointed David to be king over Israel. David was thirty years old when he began to reign, and he reigned forty years. In Hebron he reigned over Juda seven years and six months: and in Jerusalem he reigned three and thirty years over all Israel and Juda. And the king and all the men that were with him went to Jerusalem to the Jebusites the inhabitants of the land: and they said to David: Thou shalt not come in hither unless thou take away the blind and the lame that say: David shall not come in hither. But David took the castle of Sion, the same is the city of David.\par
\switchcolumn*

\LA
\begin{Resp}\Rsign Audi, Dómine, hymnum et oratiónem, quam servus tuus orat coram te hódie, ut sint óculi tui apérti, et aures tuæ inténtæ, \Star{} Super domum istam die ac nocte.\end{Resp}
\begin{Resp}\Vsign Réspice, Dómine, de sanctuário tuo, et de excélso cælórum habitáculo.\end{Resp}
\begin{Resp}\Rsign Super domum istam die ac nocte.\end{Resp}
\begin{Resp}\Vsign Glória Patri, et Fílio, \Star{} et Spirítui Sancto.\end{Resp}
\begin{Resp}\Rsign Super domum istam die ac nocte.\end{Resp}
\switchcolumn
\EN
\begin{Resp}\Rsign Hearken, O Lord, unto the cry and to the prayer which thy servant prayeth before thee today, that thine eyes may be open and thine ears attend; \Star{} Toward this house day and night.\end{Resp}
\begin{Resp}\Vsign Look down from thine high and holy place, O Lord, even from heaven thy dwelling.\end{Resp}
\begin{Resp}\Rsign Toward this house, day and night.\end{Resp}
\begin{Resp}\Vsign Glory be to the Father, and to the Son, \Star{} and to the Holy Ghost.\end{Resp}
\begin{Resp}\Rsign Toward this house, day and night.\end{Resp}
\switchcolumn*

\end{paracol}

\Office{Feria quinta infra Hebdomadam V post Octavam Pentecostes}{Thursday of the Fifth Week after Pentecost}{Feria}
\addcontentsline{toc}{subsection}{Feria quinta infra Hebdomadam V post Octavam Pentecostes \textit{— Thursday of the Fifth Week after Pentecost}}
\begin{paracol}{2}
\LA
\LessonHead{Lectio I}
\switchcolumn
\EN
\LessonHead{Lesson I}
\switchcolumn*

\LA
\LessonTitle{De libro secúndo Regum}
\switchcolumn
\EN
\LessonTitle{Lesson from the second book of Samuel}
\switchcolumn*

\LA
\Cite{2 Reg 6:1-3}
\switchcolumn
\EN
\Cite{2 Sam 6:1-3}
\switchcolumn*

\LA
\noindent Congregávit autem rursum David omnes eléctos ex Israël trigínta míllia. Surrexítque pópulus et ábiit et univérsus pópulus qui erat cum eo de viris Juda, ut addúcerent arcam Dei, super quam invocátum est nomen Dómini exercítuum, sedéntis in Chérubim super eam. Et imposuérunt arcam Dei super plaustrum novum tulerúntque eam de domo Abínadab, qui erat in Gábaa; Oza autem et Ahio fílii Abínadab, minábant plaustrum novum.\par
\switchcolumn
\EN
\noindent And David again gathered together all the chosen men of Israel, thirty thousand. And David arose and went, with all the people that were with him of the men of Juda to fetch the ark of God, upon which the name of the Lord of hosts is invoked, who sitteth over it upon the cherubims. And they laid the ark of God upon a new cart: and took it out of the house of Abinadab, who was in Gabaa: and Oza, and Ahio, the sons of Abinadab, drove the new cart.\par
\switchcolumn*

\LA
\begin{Resp}\Rsign Præparáte corda vestra Dómino, et servíte illi soli: \Star{} Et liberábit vos de mánibus inimicórum vestrórum.\end{Resp}
\begin{Resp}\Vsign Convertímini ad eum in toto corde vestro, et auférte deos aliénos de médio vestri.\end{Resp}
\begin{Resp}\Rsign Et liberábit vos de mánibus inimicórum vestrórum.\end{Resp}
\switchcolumn
\EN
\begin{Resp}\Rsign Prepare your hearts unto the Lord, and serve Him Only \Star{} And He will deliver you out of the hand of your enemies.\end{Resp}
\begin{Resp}\Vsign Return unto Him with all your hearts, and put away the strange gods from among you.\end{Resp}
\begin{Resp}\Rsign And He will deliver you out of the hand of your enemies.\end{Resp}
\switchcolumn*

\LA
\LessonHead{Lectio II}
\switchcolumn
\EN
\LessonHead{Lesson II}
\switchcolumn*

\LA
\Cite{2 Reg 6:4-7}
\switchcolumn
\EN
\Cite{2 Sam 6:4-7}
\switchcolumn*

\LA
\noindent Cumque tulísset eam de domo Abínadab, qui erat in Gábaa custódiens arcam Dei, Ahio præcedébat arcam. David autem et omnis Israël ludébant coram Dómino in ómnibus lignis fabrefáctis, et cítharis et lyris et týmpanis et sistris et cýmbalis. Postquam autem venérunt ad áream Nachon, exténdit Oza manum ad arcam Dei et ténuit eam, quóniam calcitrábant boves et declinavérunt eam. Iratúsque est indignatióne Dóminus contra Ozam, et percússit eum super temeritáte: qui mórtuus est ibi juxta arcam Dei.\par
\switchcolumn
\EN
\noindent And when they had taken it out of the house of Abinadab, who was in Gabaa, Ahio having care of the ark of God went before the ark. But David and all Israel played before the Lord on all manner of instruments made of wood, on harps and lutes and timbrels and cornets and cymbals. And when they came to the floor of Nachon, Oza put forth his hand to the ark of God, and took hold of it: because the oxen kicked and made it lean aside. And the indignation of the Lord was enkindled against Oza, and he struck him for his rashness: and he died there before the ark of God.\par
\switchcolumn*

\LA
\begin{Resp}\Rsign Deus ómnium exaudítor est: ipse misit Angelum suum, et tulit me de óvibus patris mei: \Star{} Et unxit me unctióne misericórdiæ suæ.\end{Resp}
\begin{Resp}\Vsign Dóminus, qui erípuit me de ore leónis, et de manu béstiæ liberávit me.\end{Resp}
\begin{Resp}\Rsign Et unxit me unctióne misericórdiæ suæ.\end{Resp}
\switchcolumn
\EN
\begin{Resp}\Rsign God, Which heareth all, even He sent His Angel, and took me from keeping my father's sheep, and \Star{} Anointed me with the oil of His mercy.\end{Resp}
\begin{Resp}\Vsign The Lord That delivered me out of the mouth of the lion, and out of the paw of the bear\end{Resp}
\begin{Resp}\Rsign And anointed me with the oil of His mercy.\end{Resp}
\switchcolumn*

\LA
\LessonHead{Lectio III}
\switchcolumn
\EN
\LessonHead{Lesson III}
\switchcolumn*

\LA
\Cite{2 Reg 6:8-12}
\switchcolumn
\EN
\Cite{2 Sam 6:8-12}
\switchcolumn*

\LA
\noindent Contristátus est autem David eo quod percussísset Dóminus Ozam; et vocátum est nomen loci illíus: Percússio Ozæ, usque in diem hanc. Et extímuit David Dóminum in die illa, dicens: Quómodo ingrediétur ad me arca Dómini? Et nóluit divértere ad se arcam Dómini in civitátem David; sed divértit eam in domum Obédedom Gethǽi. Et habitávit arca Dómini in domo Obédedom Gethǽi tribus ménsibus, et benedíxit Dóminus Obédedom et omnem domum ejus. Nuntiatúmque est regi David quod benedixísset Dóminus Obédedom, et ómnia ejus, propter arcam Dei. Abiit ergo David et addúxit arcam Dei de domo Obédedom in civitátem David cum gáudio.\par
\switchcolumn
\EN
\noindent And David was grieved because the Lord had struck Oza, and the name of that place was called: The striking of Oza, to this day. And David was afraid of the Lord that day, saying: How shall the ark of the Lord come to me? And he would not have the ark of the Lord brought in to himself into the city of David: but he caused it to be carried into the house of Obededom the Gethite. And the ark of the Lord abode in the house of Obededom the Gethite three months: and the Lord blessed Obededom, and all his household. And it was told king David, that the Lord had blessed Obededom, and all that he had, because of the ark of God. So David went, and brought away the ark of God out of the house of Obededom into the city of David with joy. And there were with David seven choirs, and calves for victims.\par
\switchcolumn*

\LA
\begin{Resp}\Rsign Dóminus, qui erípuit me de ore leónis, et de manu béstiæ liberávit me, \Star{} Ipse me erípiet de mánibus inimicórum meórum.\end{Resp}
\begin{Resp}\Vsign Misit Deus misericórdiam suam et veritátem suam: ánimam meam erípuit de médio catulórum leónum.\end{Resp}
\begin{Resp}\Rsign Ipse me erípiet de mánibus inimicórum meórum.\end{Resp}
\begin{Resp}\Vsign Glória Patri, et Fílio, \Star{} et Spirítui Sancto.\end{Resp}
\begin{Resp}\Rsign Ipse me erípiet de mánibus inimicórum meórum.\end{Resp}
\switchcolumn
\EN
\begin{Resp}\Rsign The Lord That delivered me out of the mouth of the lion, and out of the paw of the bear \Star{} He will deliver me out of the hand of mine enemies.\end{Resp}
\begin{Resp}\Vsign God hath sent forth His mercy and His truth, and delivered my soul from among the lion's whelps.\end{Resp}
\begin{Resp}\Rsign He will deliver me out of the hand of mine enemies.\end{Resp}
\begin{Resp}\Vsign Glory be to the Father, and to the Son, \Star{} and to the Holy Ghost.\end{Resp}
\begin{Resp}\Rsign He will deliver me out of the hand of mine enemies.\end{Resp}
\switchcolumn*

\end{paracol}

\Office{Feria sexta infra Hebdomadam V post Octavam Pentecostes}{Friday of the Fifth Week after Pentecost}{Feria}
\addcontentsline{toc}{subsection}{Feria sexta infra Hebdomadam V post Octavam Pentecostes \textit{— Friday of the Fifth Week after Pentecost}}
\begin{paracol}{2}
\LA
\LessonHead{Lectio I}
\switchcolumn
\EN
\LessonHead{Lesson I}
\switchcolumn*

\LA
\LessonTitle{De libro secúndo Regum}
\switchcolumn
\EN
\LessonTitle{Lesson from the second book of Samuel}
\switchcolumn*

\LA
\Cite{2 Reg 7:4-6}
\switchcolumn
\EN
\Cite{2 Sam 7:4-6}
\switchcolumn*

\LA
\noindent Et ecce sermo Dómini ad Nathan, dicens: Vade, et lóquere ad servum meum David: Hæc dicit Dóminus: Numquid tu ædificábis mihi domum ad habitándum? Neque enim habitávi in domo ex die illa, qua edúxi fílios Israël de terra Ægýpti usque in diem hanc, sed ambulábam in tabernáculo, et in tentório.\par
\switchcolumn
\EN
\noindent But it came to pass that night, that the word of the Lord came to Nathan, saying: Go, and say to my servant David: Thus saith the Lord: Shalt thou build me a house to dwell in? Whereas I have not dwelt in a house from the day that I brought the children of Israel out of the land of Egypt even to this day: but have walked in a tabernacle, and in a tent.\par
\switchcolumn*

\LA
\begin{Resp}\Rsign Percússit Saul mille, et David decem míllia: \Star{} Quia manus Dómini erat cum illo: percússit Philisthǽum, et ábstulit oppróbrium ex Israël.\end{Resp}
\begin{Resp}\Vsign Nonne iste est David, de quo canébant in choro, dicéntes: Saul percússit mille, et David decem míllia?\end{Resp}
\begin{Resp}\Rsign Quia manus Dómini erat cum illo: percússit Philisthǽum, et ábstulit oppróbrium ex Israël.\end{Resp}
\switchcolumn
\EN
\begin{Resp}\Rsign Saul hath slain his thousands, and David his ten thousands. \Star{} Because the hand of the Lord was with him, he smote the Philistine, and took away the reproach from Israel.\end{Resp}
\begin{Resp}\Vsign Is not this David? Did they not sing one to another of him in dances, saying: Saul hath slain his thousands, and David his ten thousands?\end{Resp}
\begin{Resp}\Rsign Because the hand of the Lord was with him, he smote the Philistine, and took away the reproach from Israel.\end{Resp}
\switchcolumn*

\LA
\LessonHead{Lectio II}
\switchcolumn
\EN
\LessonHead{Lesson II}
\switchcolumn*

\LA
\Cite{2 Reg 7:7-11}
\switchcolumn
\EN
\Cite{2 Sam 7:7-11}
\switchcolumn*

\LA
\noindent Per cuncta loca, quæ transívi cum ómnibus fíliis Israël, numquid loquens locútus sum ad unum de tríbubus Israël, cui præcépi ut pásceret pópulum meum Israël, dicens: Quare non ædificástis mihi domum cédrinam? Et nunc hæc dices servo meo David: Hæc dicit Dóminus exercítuum: Ego tuli te de páscuis sequéntem greges, ut esses dux super pópulum meum Israël, et fui tecum in ómnibus ubicúmque ambulásti, et interféci univérsos inimícos tuos a fácie tua: fecíque tibi nomen grande, juxta nomen magnórum, qui sunt in terra. Et ponam locum pópulo meo Israël et plantábo eum, et habitábit sub eo et non turbábitur ámplius; nec addent fílii iniquitátis ut afflígant eum sicut prius, ex die qua constítui júdices super pópulum meum Israël. Et réquiem dabo tibi ab ómnibus inimícis tuis: prædicítque tibi Dóminus, quod domum fáciat tibi Dóminus.\par
\switchcolumn
\EN
\noindent In all the places that I have gone through with all the children of Israel, did ever I speak a word to any one of the tribes of Israel, whom I commanded to feed my people Israel, saying: Why have you not built me a house of cedar? And now thus shalt thou speak to my servant David: Thus saith the Lord of hosts: a I took thee out of the pastures from following the sheep to be ruler over my people Israel: And I have been with thee wheresoever thou hast walked, and have slain all thy enemies from before thy face: and I have made thee a great man, like unto the name of the great ones that are on the earth. And I will appoint a place for my people Israel, and I will plant them, and they shall dwell therein, and shall be disturbed no more: neither shall the children of iniquity afflict them any more as they did before, From the day that I appointed judges over my people Israel: and I will give thee rest from all thy enemies. And the Lord foretelleth to thee, that the Lord will make thee a house.\par
\switchcolumn*

\LA
\begin{Resp}\Rsign Montes Gélboë, nec ros nec plúvia véniant super vos, \Star{} Ubi cecidérunt fortes Israël.\end{Resp}
\begin{Resp}\Vsign Omnes montes, qui estis in circúitu ejus, vísitet Dóminus: a Gélboë autem tránseat.\end{Resp}
\begin{Resp}\Rsign Ubi cecidérunt fortes Israël.\end{Resp}
\switchcolumn
\EN
\begin{Resp}\Rsign Ye mountains of Gilboa, let there be no dew, neither let there be rain upon you \Star{} For there are the mighty of Israel fallen.\end{Resp}
\begin{Resp}\Vsign All ye mountains that stand round about, the Lord look upon you but let Him pass by Gilboa\end{Resp}
\begin{Resp}\Rsign For there are the mighty of Israel fallen.\end{Resp}
\switchcolumn*

\LA
\LessonHead{Lectio III}
\switchcolumn
\EN
\LessonHead{Lesson III}
\switchcolumn*

\LA
\Cite{2 Reg 7:12-17}
\switchcolumn
\EN
\Cite{2 Sam 7:12-17}
\switchcolumn*

\LA
\noindent Cumque compléti fúerint dies tui, et dormíeris cum pátribus tuis, suscitábo semen tuum post te, quod egrediétur de útero tuo, et firmábo regnum ejus. Ipse ædificábit domum nómini meo, et stabíliam thronum regni ejus usque in sempitérnum. Ego ero ei in patrem, et ipse erit mihi in fílium. Qui si iníque áliquid gésserit, árguam eum in virga virórum, et in plagis filiórum hóminum. Misericórdiam autem meam non áuferam ab eo, sicut ábstuli a Saul, quem amóvi a fácie mea. Et fidélis erit domus tua, et regnum tuum usque in ætérnum ante fáciem tuam, et thronus tuus erit firmus júgiter. Secúndum ómnia verba hæc, et juxta univérsam visiónem istam, sic locútus est Nathan ad David.\par
\switchcolumn
\EN
\noindent And when thy days shall be fulfilled, and thou shalt sleep with thy fathers, I will raise up thy seed after thee, which shall proceed out of thy bowels, and I will establish his kingdom. He shall build a house to my name, and I will establish the throne of his kingdom for ever. I will be to him a father, and he shall be to me a son: and if he commit any iniquity, I will correct him with the rod of men, and with the stripes of the children of men. But my mercy I will not take away from him, as I took it from Saul, whom I removed from before my face. And thy house shall be faithful, and thy kingdom for ever before thy face, and thy throne shall be firm for ever. According to all these words and according to all this vision, so did Nathan speak to David.\par
\switchcolumn*

\LA
\begin{Resp}\Rsign Ego te tuli de domo patris tui, dicit Dóminus, et pósui te páscere gregem pópuli mei: \Star{} Et fui tecum in ómnibus ubicúmque ambulásti, firmans regnum tuum in ætérnum.\end{Resp}
\begin{Resp}\Vsign Fecíque tibi nomen grande, juxta nomen magnórum, qui sunt in terra: et réquiem dedi tibi ab ómnibus inimícis tuis.\end{Resp}
\begin{Resp}\Rsign Et fui tecum in ómnibus ubicúmque ambulásti, firmans regnum tuum in ætérnum.\end{Resp}
\begin{Resp}\Vsign Glória Patri, et Fílio, \Star{} et Spirítui Sancto.\end{Resp}
\begin{Resp}\Rsign Et fui tecum in ómnibus ubicúmque ambulásti, firmans regnum tuum in ætérnum.\end{Resp}
\switchcolumn
\EN
\begin{Resp}\Rsign Thus saith the Lord: I took thee out of thy father's house, and appointed thee to be ruler over My people, over Israel. \Star{} And I was with thee whithersoever thou wentest, to establish thy kingdom for ever.\end{Resp}
\begin{Resp}\Vsign And I have made thee a great name, like unto the name of the great men that are in the earth, and have caused thee to rest from all thine enemies.\end{Resp}
\begin{Resp}\Rsign And I was with thee whithersoever thou wentest, to establish thy kingdom for ever.\end{Resp}
\begin{Resp}\Vsign Glory be to the Father, and to the Son, \Star{} and to the Holy Ghost.\end{Resp}
\begin{Resp}\Rsign And I was with thee whithersoever thou wentest, to establish thy kingdom for ever.\end{Resp}
\switchcolumn*

\end{paracol}

\Office{Sabbato infra Hebdomadam V post Octavam Pentecostes}{Saturday of the Fifth Week after Pentecost}{Feria}
\addcontentsline{toc}{subsection}{Sabbato infra Hebdomadam V post Octavam Pentecostes \textit{— Saturday of the Fifth Week after Pentecost}}
\begin{paracol}{2}
\LA
\LessonHead{Lectio I}
\switchcolumn
\EN
\LessonHead{Lesson I}
\switchcolumn*

\LA
\LessonTitle{De libro secúndo Regum}
\switchcolumn
\EN
\LessonTitle{Lesson from the second book of Samuel}
\switchcolumn*

\LA
\Cite{2 Reg 11:1-4}
\switchcolumn
\EN
\Cite{2 Sam 11:1-4}
\switchcolumn*

\LA
\noindent Factum est autem, verténte anno, eo témpore quo solent reges ad bella procédere, misit David Joab, et servos suos cum eo, et univérsum Israël, et vastavérunt fílios Ammon, et obsedérunt Rabba; David autem remánsit in Jerúsalem. Dum hæc ageréntur, áccidit ut súrgeret David de strato suo post merídiem et deambuláret in solário domus régiæ: vidítque mulíerem se lavántem ex advérso super solárium suum; erat autem múlier pulchra valde. Misit ergo rex, et requisívit quæ esset múlier; nuntiatúmque est ei, quod ipsa esset Bethsabée fília Elíam, uxor Uríæ Hethǽi. Missis ítaque David núntiis, tulit eam.\par
\switchcolumn
\EN
\noindent And it came to pass at the return of the year, at the time when kings go forth to war, that David sent Joab and his servants with him, and all Israel, and they spoiled the children of Ammon, and besieged Rabba: but David remained in Jerusalem. In the mean time it happened that David arose from his bed after noon, and walked upon the roof of the king's house: and he saw from the roof of his house a woman washing herself, over against him: and the woman was very beautiful. And the king sent, and inquired who the woman was. And it was told him, that she was Bethsabee the daughter of Eliam, the wife of Urias the Hethite. And David sent messengers, and took her.\par
\switchcolumn*

\LA
\begin{Resp}\Rsign Peccávi super númerum arénæ maris, et multiplicáta sunt peccáta mea: et non sum dignus vidére altitúdinem cæli præ multitúdine iniquitátis meæ: quóniam irritávi iram tuam, \Star{} Et malum coram te feci.\end{Resp}
\begin{Resp}\Vsign Quóniam iniquitátem meam ego cognósco: et delíctum meum contra me est semper, quia tibi soli peccávi.\end{Resp}
\begin{Resp}\Rsign Et malum coram te feci.\end{Resp}
\switchcolumn
\EN
\begin{Resp}\Rsign My sins are many, yea, they are more in number than the sands of the sea; I am not worthy to look up toward heaven because of the multitude of my iniquities; for I have provoked thee to anger; \Star{} And done evil in thy sight.\end{Resp}
\begin{Resp}\Vsign For I acknowledge my transgression, and my sin is ever before me, for against thee only have I sinned.\end{Resp}
\begin{Resp}\Rsign And done evil in thy sight.\end{Resp}
\switchcolumn*

\LA
\LessonHead{Lectio II}
\switchcolumn
\EN
\LessonHead{Lesson II}
\switchcolumn*

\LA
\Cite{2 Reg 11:5-11}
\switchcolumn
\EN
\Cite{2 Sam 11:5-11}
\switchcolumn*

\LA
\noindent Et revérsa est in domum suam, concépto fœtu. Mitténsque nuntiávit David et ait: Concépi. Misit autem David ad Joab dicens: Mitte ad me Uríam Hethǽum. Misítque Joab Uríam ad David. Et venit Urías ad David. Quæsivítque David quam recte ágeret Joab et pópulus, et quómodo administrarétur bellum. Et dixit David ad Uríam: Vade in domum tuam et lava pedes tuos. Et egréssus est Urías de domo regis, secutúsque est eum cibus régius. Dormívit autem Urías ante portam domus régiæ cum áliis servis dómini sui, et non descéndit ad domum suam. Nuntiatúmque est David a dicéntibus: Non ivit Urías in domum suam. Et ait David ad Uríam: Numquid non de via venísti? quare non descendísti in domum tuam? Et ait Urías ad David: Arca Dei et Israël et Juda hábitant in papiliónibus, et dóminus meus Joab et servi dómini mei super fáciem terræ manent; et ego ingrédiar domum meam, ut cómedam et bibam, et dórmiam cum uxóre mea? Per salútem tuam et per salútem ánimæ tuæ non fáciam rem hanc.\par
\switchcolumn
\EN
\noindent And she returned to her house having conceived. And she sent and told David, and said: I have conceived. And David sent to Joab, saying: Send me Urias the Hethite. And Joab sent Urias to David. And Urias came to David. And David asked how Joab did, and the people, and how the war was carried on. And David said to Urias: Go into thy house, and wash thy feet. And Urias went out from the king's house, and there went out after him a mess of meat from the king. But Urias slept before the gate of the king's house, with the other servants of his lord, and went not down to his own house. And it was told David by some that said: Urias went not to his house. And David said to Urias: Didst thou not come from thy journey? why didst thou not go down to thy house? And Urias said to David: The ark of God and Israel and Juda dwell in tents, and my lord Joab and the servants of my lord abide upon the face of the earth: and shall I go into my house, to eat and to drink, and to sleep with my wife? By thy welfare and by the welfare of thy soul I will not do this thing\par
\switchcolumn*

\LA
\begin{Resp}\Rsign Exaudísti, Dómine, oratiónem servi tui, ut ædificárem templum nómini tuo: \Star{} Bénedic et sanctífica domum istam in sempitérnum, Deus Israël.\end{Resp}
\begin{Resp}\Vsign Dómine, qui custódis pactum cum servis tuis, qui ámbulant coram te in toto corde suo.\end{Resp}
\begin{Resp}\Rsign Bénedic et sanctífica domum istam in sempitérnum, Deus Israël.\end{Resp}
\switchcolumn
\EN
\begin{Resp}\Rsign O Lord, Thou hast hearkened unto the prayer of thy servant, that I might build a temple unto thy Name, \Star{} O God of Israel, bless Thou, and hallow this house for ever.\end{Resp}
\begin{Resp}\Vsign O Lord, Who keepest covenant with thy servants that walk before thee in all their heart.\end{Resp}
\begin{Resp}\Rsign O God of Israel, bless Thou, and hallow this house for ever.\end{Resp}
\switchcolumn*

\LA
\LessonHead{Lectio III}
\switchcolumn
\EN
\LessonHead{Lesson III}
\switchcolumn*

\LA
\Cite{2 Reg 11:12-17}
\switchcolumn
\EN
\Cite{2 Sam 11:12-17}
\switchcolumn*

\LA
\noindent Ait ergo David ad Uríam: Mane hic étiam hódie, et cras dimíttam te. Mansit Urías in Jerúsalem in die illa et áltera. Et vocávit eum David ut coméderet coram se et bíberet et inebriávit eum; qui egréssus véspere dormívit in strato suo cum servis dómini sui, et in domum suam non descéndit. Factum est ergo mane, et scripsit David epístolam ad Joab: misítque per manum Uríæ, scribens in epístola: Pónite Uríam ex advérso belli, ubi fortíssimum est prǽlium, et derelínquite eum ut percússus intéreat. Igitur cum Joab obsidéret urbem, pósuit Uríam in loco ubi sciébat viros esse fortíssimos. Egressíque viri de civitáte bellábant advérsum Joab, et cecidérunt de pópulo servórum David, et mórtuus est étiam Urías Hethǽus.\par
\switchcolumn
\EN
\noindent Then David said to Urias: Tarry here today, and tomorrow I will send thee away. Urias tarried in Jerusalem that day and the next. And David called him to eat and to drink before him, and he made him drunk: and he went out in the evening, and slept on his couch with the servants of his lord, and went not down into his house. And when the morning was come, David wrote a letter to Joab: and sent it by the hand of Urias, Writing in the letter: Set ye Urias in the front of the battle, where the fight is strongest: and leave ye him, that he may be wounded and die. Wherefore as Joab was besieging the city, he put Urias in the place where he knew the bravest men were. And the men coming out of the city, fought against Joab, and there fell some of the people of the servants of David, and Urias the Hethite was killed also.\par
\switchcolumn*

\LA
\begin{Resp}\Rsign Audi, Dómine, hymnum et oratiónem, quam servus tuus orat coram te hódie, ut sint óculi tui apérti, et aures tuæ inténtæ, \Star{} Super domum istam die ac nocte.\end{Resp}
\begin{Resp}\Vsign Réspice, Dómine, de sanctuário tuo, et de excélso cælórum habitáculo.\end{Resp}
\begin{Resp}\Rsign Super domum istam die ac nocte.\end{Resp}
\begin{Resp}\Vsign Glória Patri, et Fílio, \Star{} et Spirítui Sancto.\end{Resp}
\begin{Resp}\Rsign Super domum istam die ac nocte.\end{Resp}
\switchcolumn
\EN
\begin{Resp}\Rsign Hearken, O Lord, unto the cry and to the prayer which thy servant prayeth before thee today, that thine eyes may be open and thine ears attend; \Star{} Toward this house day and night.\end{Resp}
\begin{Resp}\Vsign Look down from thine high and holy place, O Lord, even from heaven thy dwelling.\end{Resp}
\begin{Resp}\Rsign Toward this house, day and night.\end{Resp}
\begin{Resp}\Vsign Glory be to the Father, and to the Son, \Star{} and to the Holy Ghost.\end{Resp}
\begin{Resp}\Rsign Toward this house, day and night.\end{Resp}
\switchcolumn*

\end{paracol}

\Office{Dominica VI Post Pentecosten}{Sixth Sunday after Pentecost}{Semiduplex}
\addcontentsline{toc}{subsection}{Dominica VI Post Pentecosten \textit{— Sixth Sunday after Pentecost}}
\begin{paracol}{2}
\LA
\LessonHead{Lectio I}
\switchcolumn
\EN
\LessonHead{Lesson I}
\switchcolumn*

\LA
\LessonTitle{De libro secúndo Regum}
\switchcolumn
\EN
\LessonTitle{Lesson from the second book of Samuel}
\switchcolumn*

\LA
\Cite{2 Reg 12:1-4}
\switchcolumn
\EN
\Cite{2 Sam 12:1-4}
\switchcolumn*

\LA
\noindent Misit ergo Dóminus Nathan ad David: qui, cum venísset ad eum, dixit ei: Duo viri erant in civitáte una, unus dives et alter pauper. Dives habébat oves et boves plúrimos valde, pauper autem nihil habébat omníno præter ovem unam párvulam, quam émerat et nutríerat, et quæ créverat apud eum cum fíliis ejus simul de pane illíus cómedens et de cálice ejus bibens et in sinu illíus dórmiens; erátque illic sicut fília. Cum autem peregrínus quidam venísset ad dívitem, parcens ille súmere de óvibus et de bobus suis, ut exhibéret convívium peregríno illi qui vénerat ad se, tulit ovem viri páuperis et præparávit cibos hómini qui vénerat ad se.\par
\switchcolumn
\EN
\noindent And the Lord sent Nathan to David: and when he was come to him, he said to him: There were two men in one city, the one rich, and the other poor. The rich man had exceeding many sheep and oxen. But the poor man had nothing at all but one little ewe lamb, which he had bought and nourished up, and which had grown up in his house together with his children, eating of his bread, and drinking of his cup, and sleeping in his bosom: and it was unto him as a daughter. And when a certain stranger was come to the rich man, he spared to take of his own sheep and oxen, to make a feast for that stranger, who was come to him, but took the poor man's ewe, and dressed it for the man that was come to him.\par
\switchcolumn*

\LA
\begin{Resp}\Rsign Præparáte corda vestra Dómino, et servíte illi soli: \Star{} Et liberábit vos de mánibus inimicórum vestrórum.\end{Resp}
\begin{Resp}\Vsign Convertímini ad eum in toto corde vestro, et auférte deos aliénos de médio vestri.\end{Resp}
\begin{Resp}\Rsign Et liberábit vos de mánibus inimicórum vestrórum.\end{Resp}
\switchcolumn
\EN
\begin{Resp}\Rsign Prepare your hearts unto the Lord, and serve Him Only \Star{} And He will deliver you out of the hand of your enemies.\end{Resp}
\begin{Resp}\Vsign Return unto Him with all your hearts, and put away the strange gods from among you.\end{Resp}
\begin{Resp}\Rsign And He will deliver you out of the hand of your enemies.\end{Resp}
\switchcolumn*

\LA
\LessonHead{Lectio II}
\switchcolumn
\EN
\LessonHead{Lesson II}
\switchcolumn*

\LA
\Cite{2 Reg 12:5-9}
\switchcolumn
\EN
\Cite{2 Sam 12:5-9}
\switchcolumn*

\LA
\noindent Irátus autem indignatióne David advérsus hóminem illum nimis, dixit ad Nathan: Vivit Dóminus, quóniam fílius mortis est vir qui fecit hoc: ovem reddet in quádruplum, eo quod fécerit verbum istud et non pepércerit. Dixit autem Nathan ad David: Tu es ille vir. Hæc dicit Dóminus Deus Israël: Ego unxi te in regem super Israël et ego érui te de manu Saul, et dedi tibi domum dómini tui, et uxóres dómini tui in sinu tuo dedíque tibi domum Israël et Juda; et, si parva sunt ista, adíciam tibi multo majóra. Quare ergo contempsísti verbum Dómini, ut fáceres malum in conspéctu meo? Uríam Hethǽum percussísti gládio et uxórem illíus accepísti in uxórem tibi et interfecísti eum gládio filiórum Ammon.\par
\switchcolumn
\EN
\noindent And David's anger being exceedingly kindled against that man, he said to Nathan: As the Lord liveth, the man that hath done this is a child of death. He shall restore the ewe fourfold, because he did this thing, and had no pity. And Nathan said to David: Thou art the man. Thus saith the Lord the God of Israel: I anointed thee king over Israel, and I delivered thee from the hand of Saul, And gave thee thy master's house and thy master's wives into thy bosom, and gave thee the house of Israel and Juda: and if these things be little, I shall add far greater things unto thee. Why therefore hast thou despised the word of the Lord, to do evil in my sight? Thou hast killed Urias the Hethite with the sword, and hast taken his wife to be thy wife, and hast slain him with the sword of the children of Ammon.\par
\switchcolumn*

\LA
\begin{Resp}\Rsign Deus ómnium exaudítor est: ipse misit Angelum suum, et tulit me de óvibus patris mei: \Star{} Et unxit me unctióne misericórdiæ suæ.\end{Resp}
\begin{Resp}\Vsign Dóminus, qui erípuit me de ore leónis, et de manu béstiæ liberávit me.\end{Resp}
\begin{Resp}\Rsign Et unxit me unctióne misericórdiæ suæ.\end{Resp}
\switchcolumn
\EN
\begin{Resp}\Rsign God, Which heareth all, even He sent His Angel, and took me from keeping my father's sheep, and \Star{} Anointed me with the oil of His mercy.\end{Resp}
\begin{Resp}\Vsign The Lord That delivered me out of the mouth of the lion, and out of the paw of the bear\end{Resp}
\begin{Resp}\Rsign And anointed me with the oil of His mercy.\end{Resp}
\switchcolumn*

\LA
\LessonHead{Lectio III}
\switchcolumn
\EN
\LessonHead{Lesson III}
\switchcolumn*

\LA
\Cite{2 Reg 12:10-16}
\switchcolumn
\EN
\Cite{2 Sam 12:10-16}
\switchcolumn*

\LA
\noindent Quam ob rem non recédet gládius de domo tua usque in sempitérnum, eo quod despéxeris me et túleris uxórem Uríæ Hethǽi ut esset uxor tua. Itaque hæc dicit Dóminus: Ecce ego suscitábo super te malum de domo tua, et tollam uxóres tuas in óculis tuis et dabo próximo tuo, et dórmiet cum uxóribus tuis in óculis solis hujus. Tu enim fecísti abscóndite, ego autem fáciam verbum istud in conspéctu omnis Israël et in conspéctu solis. Et dixit David ad Nathan: Peccávi Dómino. Dixítque Nathan ad David: Dóminus quoque tránstulit peccátum tuum: non moriéris. Verúmtamen, quóniam blasphemáre fecísti inimícos Dómini propter verbum hoc, fílius, qui natus est tibi, morte moriétur. Et revérsus est Nathan in domum suam. Percússit quoque Dóminus párvulum, quem pepérerat uxor Uríæ David, et desperátus est. Deprecatúsque est David Dóminum pro párvulo, et jejunávit David jejúnio et ingréssus seórsum jácuit super terram.\par
\switchcolumn
\EN
\noindent Therefore the sword shall never depart from thy house, because thou hast despised me, and hast taken the wife of Urias the Hethite to be thy wife. Thus saith the Lord: Behold, I will raise up evil against thee out of thy own house, and I will take thy wives before thy eyes I and give them to thy neighhour, and he shall lie with thy wives in the sight of this sun. For thou didst it secretly: but I will do this thing in the sight of all Israel, and in the sight of the sun. And David said to Nathan: I have sinned against the Lord. And Nathan said to David: The Lord also hath taken away thy sin: thou shalt not die. Nevertheless, because thou hast given occasion to the enemies of the Lord to blaspheme, for this thing, the child that is born to thee, shall surely die. And Nathan returned to his house. The Lord also struck the child which the wife of Urias had borne to David, and his life was despaired of. And David besought the Lord for the child: and David kept a fast, and going in by himself lay upon the ground.\par
\switchcolumn*

\LA
\begin{Resp}\Rsign Dóminus, qui erípuit me de ore leónis, et de manu béstiæ liberávit me, \Star{} Ipse me erípiet de mánibus inimicórum meórum.\end{Resp}
\begin{Resp}\Vsign Misit Deus misericórdiam suam et veritátem suam: ánimam meam erípuit de médio catulórum leónum.\end{Resp}
\begin{Resp}\Rsign Ipse me erípiet de mánibus inimicórum meórum.\end{Resp}
\begin{Resp}\Vsign Glória Patri, et Fílio, \Star{} et Spirítui Sancto.\end{Resp}
\begin{Resp}\Rsign Ipse me erípiet de mánibus inimicórum meórum.\end{Resp}
\switchcolumn
\EN
\begin{Resp}\Rsign The Lord That delivered me out of the mouth of the lion, and out of the paw of the bear \Star{} He will deliver me out of the hand of mine enemies.\end{Resp}
\begin{Resp}\Vsign God hath sent forth His mercy and His truth, and delivered my soul from among the lion's whelps.\end{Resp}
\begin{Resp}\Rsign He will deliver me out of the hand of mine enemies.\end{Resp}
\begin{Resp}\Vsign Glory be to the Father, and to the Son, \Star{} and to the Holy Ghost.\end{Resp}
\begin{Resp}\Rsign He will deliver me out of the hand of mine enemies.\end{Resp}
\switchcolumn*

\LA
\LessonHead{Lectio IV}
\switchcolumn
\EN
\LessonHead{Lesson IV}
\switchcolumn*

\LA
\LessonTitle{Ex libro sancti Ambrósii Epíscopi de Apología David}
\switchcolumn
\EN
\LessonTitle{From the Book On the Defense of David written by St. Ambrose, Bishop of Milan.}
\switchcolumn*

\LA
\Source{Apolog. 1, c. 2}
\switchcolumn
\EN
\Source{i. 2}
\switchcolumn*

\LA
\noindent Unusquísque nostrum per síngulas horas quam multa delínquit! nec tamen unusquísque de plebe peccátum suum confiténdum putat. Ille rex, tantus ac potens, ne exíguo quidem moménto manére penes se delícti passus est consciéntiam; sed, præmatúra confessióne atque imménso dolóre, réddidit peccátum suum Dómino. Quem mihi nunc fácile repérias honorátum ac dívitem, qui, si arguátur alicújus culpæ reus, non moléste ferat? At ille régio clarus império, tot divínis probátus oráculis, cum a priváto hómine corriperétur quod gráviter deliquísset, non indignátus infrémuit, sed conféssus ingémuit culpæ dolóre.\par
\switchcolumn
\EN
\noindent In how many things doth each one of us transgress every hour! And nevertheless not one of all us common men thinketh it well to confess his sin. Yet that strong and great King would not suffer the acknowledgment of his iniquity to remain, even for a moment, hidden in his own heart. With eager confession and bitter sorrow, he admitted that he had sinned against the Lord. Which of you will easily find me now some honoured and wealthy person, who will not take it ill if I rebuke him for a fault whereof he is guilty? But David, amid the splendours of a throne and' the certainty of Divine revelations, when he was rebuked by one of his subjects for his grievous transgression, was not roused to anger, but contrariwise, acknowledged his sin with groans and affliction.\par
\switchcolumn*

\LA
\begin{Resp}\Rsign Percússit Saul mille, et David decem míllia: \Star{} Quia manus Dómini erat cum illo: percússit Philisthǽum, et ábstulit oppróbrium ex Israël.\end{Resp}
\begin{Resp}\Vsign Nonne iste est David, de quo canébant in choro, dicéntes: Saul percússit mille, et David decem míllia?\end{Resp}
\begin{Resp}\Rsign Quia manus Dómini erat cum illo: percússit Philisthǽum, et ábstulit oppróbrium ex Israël.\end{Resp}
\switchcolumn
\EN
\begin{Resp}\Rsign Saul hath slain his thousands, and David his ten thousands. \Star{} Because the hand of the Lord was with him, he smote the Philistine, and took away the reproach from Israel.\end{Resp}
\begin{Resp}\Vsign Is not this David? Did they not sing one to another of him in dances, saying: Saul hath slain his thousands, and David his ten thousands?\end{Resp}
\begin{Resp}\Rsign Because the hand of the Lord was with him, he smote the Philistine, and took away the reproach from Israel.\end{Resp}
\switchcolumn*

\LA
\LessonHead{Lectio V}
\switchcolumn
\EN
\LessonHead{Lesson V}
\switchcolumn*

\LA
\noindent Dénique Dóminum dolor íntimi movit afféctus, ut Nathan díceret: Quóniam pœnítuit te, et Dóminus ábstulit peccátum tuum. Matúritas ítaque véniæ, profúndam regis fuísse pœniténtiam declarávit, quæ tanti erróris offénsam tradúxerit. Alii hómines cum a sacerdótibus corripiúntur, peccátum suum íngravant, dum negáre cúpiunt aut deféndere; ibíque eórum lapsus est major, ubi sperátur corréctio. Sancti autem Dómini, qui consummáre pium certámen géstiunt, et cúrrere cursum salútis, sícubi forte ut hómines corrúerint, natúræ magis fragilitáte quam peccándi libídine, acrióres ad curréndum resúrgunt, pudóris stímulo majóra reparántes certámina; ut non solum nullum attulísse æstimétur lapsus impediméntum, sed étiam velocitátis incentíva cumulásse.\par
\switchcolumn
\EN
\noindent The heart-felt sorrow of David moved the Lord to compassion, so that Nathan said Because thou hast repented, the Lord also hath put away thy sin. The instant gift of pardon declareth the depth of the King's repentance, which was able to obtain the forgiveness of so grievous a transgression. Other men, when they be rebuked of Priests, do but aggravate the heinousness of their sins by the seeking to deny or to excuse them, and thereby make deeper their fall by means of that which should have helped them up. But the saints of the Lord who will to fight a good fight of godliness unto the end and to finish their course by saving their souls, howbeit, they may perchance have fallen like other men, have done so rather through man's weakness than through lust for iniquity, and rise more eager to go on than before. Shame goadeth them on to fly at higher things. So that not only is their fall to be reckoned to have nowise hampered them, but rather to have quickened their speed.\par
\switchcolumn*

\LA
\begin{Resp}\Rsign Montes Gélboë, nec ros nec plúvia véniant super vos, \Star{} Ubi cecidérunt fortes Israël.\end{Resp}
\begin{Resp}\Vsign Omnes montes, qui estis in circúitu ejus, vísitet Dóminus: a Gélboë autem tránseat.\end{Resp}
\begin{Resp}\Rsign Ubi cecidérunt fortes Israël.\end{Resp}
\switchcolumn
\EN
\begin{Resp}\Rsign Ye mountains of Gilboa, let there be no dew, neither let there be rain upon you \Star{} For there are the mighty of Israel fallen.\end{Resp}
\begin{Resp}\Vsign All ye mountains that stand round about, the Lord look upon you but let Him pass by Gilboa\end{Resp}
\begin{Resp}\Rsign For there are the mighty of Israel fallen.\end{Resp}
\switchcolumn*

\LA
\LessonHead{Lectio VI}
\switchcolumn
\EN
\LessonHead{Lesson VI}
\switchcolumn*

\LA
\noindent Peccávit David, quod solent reges; sed pœniténtiam gessit, flevit, ingémuit, quod non solent reges. Conféssus est culpam, obsecrávit indulgéntiam, humi stratus deplorávit ærúmnam, jejunávit, orávit, confessiónis suæ testimónium in perpétua sǽcula vulgáto dolóre transmísit. Quod erubéscunt fácere priváti, rex non erúbuit confitéri. Qui tenéntur légibus, audent suum negáre peccátum, dedignántur rogáre indulgéntiam, quam petébat qui nullis légibus tenebátur humánis. Quod peccávit, conditiónis est; quod supplicávit, correctiónis. Lapsus commúnis, sed speciális conféssio. Culpam ítaque incidísse, natúræ est: diluísse, virtútis.\par
\switchcolumn
\EN
\noindent David sinned; and so oftentimes do other kings. David repented with groaning and tears; and so do not oftentimes other kings. He admitted his guilt; he implored forgiveness; he cast himself down upon the ground, and there wept over his crime; he fasted; he prayed; by publishing his sorrow he left an everlasting witness of his acknowledgment. What meaner men blush to do, the King was not ashamed to own. They who are answerable to law are bold to deny their crimes, and too haughty to ask pardon. Not so he, though he could be haled before no earthly judgment-seat. That he sinned was a matter flowing from his nature; that he asked for pardon, his own repentance. To fall is common to all men, but his confession was his own. To transgress thus was nature; to efface his guilt, greatness.\par
\switchcolumn*

\LA
\begin{Resp}\Rsign Ego te tuli de domo patris tui, dicit Dóminus, et pósui te páscere gregem pópuli mei: \Star{} Et fui tecum in ómnibus ubicúmque ambulásti, firmans regnum tuum in ætérnum.\end{Resp}
\begin{Resp}\Vsign Fecíque tibi nomen grande, juxta nomen magnórum, qui sunt in terra: et réquiem dedi tibi ab ómnibus inimícis tuis.\end{Resp}
\begin{Resp}\Rsign Et fui tecum in ómnibus ubicúmque ambulásti, firmans regnum tuum in ætérnum.\end{Resp}
\begin{Resp}\Vsign Glória Patri, et Fílio, \Star{} et Spirítui Sancto.\end{Resp}
\begin{Resp}\Rsign Et fui tecum in ómnibus ubicúmque ambulásti, firmans regnum tuum in ætérnum.\end{Resp}
\switchcolumn
\EN
\begin{Resp}\Rsign Thus saith the Lord: I took thee out of thy father's house, and appointed thee to be ruler over My people, over Israel. \Star{} And I was with thee whithersoever thou wentest, to establish thy kingdom for ever.\end{Resp}
\begin{Resp}\Vsign And I have made thee a great name, like unto the name of the great men that are in the earth, and have caused thee to rest from all thine enemies.\end{Resp}
\begin{Resp}\Rsign And I was with thee whithersoever thou wentest, to establish thy kingdom for ever.\end{Resp}
\begin{Resp}\Vsign Glory be to the Father, and to the Son, \Star{} and to the Holy Ghost.\end{Resp}
\begin{Resp}\Rsign And I was with thee whithersoever thou wentest, to establish thy kingdom for ever.\end{Resp}
\switchcolumn*

\LA
\LessonHead{Lectio VII}
\switchcolumn
\EN
\LessonHead{Lesson VII}
\switchcolumn*

\LA
\LessonTitle{Léctio sancti Evangélii secúndum Marcum}
\switchcolumn
\EN
\LessonTitle{From the Holy Gospel according to Mark}
\switchcolumn*

\LA
\Cite{Marc 8:1-9}
\switchcolumn
\EN
\Cite{Mark 8:1-9}
\switchcolumn*

\LA
\noindent In illo témpore: Cum turba multa esset cum Jesu nec habérent quod manducárent, convocátis discípulis, ait illis: Miséreor super turbam, quia ecce jam tríduo sústinent me nec habent quod mandúcent. Et réliqua.\par
\switchcolumn
\EN
\noindent In those days, the multitude being very great, and having nothing to eat, Jesus called His disciples unto Him, and saith unto them: I have compassion on the multitude, because they have now been with Me for three days, and have nothing to eat. And so on.\par
\switchcolumn*

\LA
\LessonTitle{Homilía sancti Ambrósii Epíscopi}
\switchcolumn
\EN
\LessonTitle{Homily by St. Ambrose, Bishop of Milan.}
\switchcolumn*

\LA
\Source{Liber 6 in Lucæ cap. 9, post initium}
\switchcolumn
\EN
\Source{Bk. vi. on Luke ix.}
\switchcolumn*

\LA
\noindent Posteáquam illa, quæ Ecclésiæ typum accépit, a fluxu curáta est sánguinis, posteáquam Apóstoli ad evangelizándum regnum Dei sunt destináti, grátiæ cæléstis impartítur aliméntum. Sed quibus impartiátur, advérte. Non otiósis, non in civitáte, quasi in synagóga vel sæculári dignitáte residéntibus; sed inter desérta quæréntibus Christum. Qui enim non fastídiunt, ipsi excipiúntur a Christo, et cum ipsis lóquitur Dei Verbum, non de sæculáribus, sed de regno Dei. Et si qui corporális gerunt úlcera passiónis, his medicínam suam libénter indúlget.\par
\switchcolumn
\EN
\noindent After that woman, who is a type of the Church, was healed of the issue of blood (Luke viii. 43-48); the Lord had sent His disciples to preach the kingdom of God (ix. 2). His heavenly tenderness gave food. But consider who they were unto whom He gave it. He gave it not to such as dwell at ease, not to men in cities, not to such as sit in places of worldly splendour, but to men seeking Christ in a desert place. Such as are not given to niceness are they whom Christ receiveth, and unto whom the Word of God speaketh, not of earthly things, but of the kingdom of God (11.) And if any bear in them the running sores of fleshly passion, He healeth them.\par
\switchcolumn*

\LA
\begin{Resp}\Rsign Peccávi super númerum arénæ maris, et multiplicáta sunt peccáta mea: et non sum dignus vidére altitúdinem cæli præ multitúdine iniquitátis meæ: quóniam irritávi iram tuam, \Star{} Et malum coram te feci.\end{Resp}
\begin{Resp}\Vsign Quóniam iniquitátem meam ego cognósco: et delíctum meum contra me est semper, quia tibi soli peccávi.\end{Resp}
\begin{Resp}\Rsign Et malum coram te feci.\end{Resp}
\switchcolumn
\EN
\begin{Resp}\Rsign My sins are many, yea, they are more in number than the sands of the sea; I am not worthy to look up toward heaven because of the multitude of my iniquities; for I have provoked thee to anger; \Star{} And done evil in thy sight.\end{Resp}
\begin{Resp}\Vsign For I acknowledge my transgression, and my sin is ever before me, for against thee only have I sinned.\end{Resp}
\begin{Resp}\Rsign And done evil in thy sight.\end{Resp}
\switchcolumn*

\LA
\LessonHead{Lectio VIII}
\switchcolumn
\EN
\LessonHead{Lesson VIII}
\switchcolumn*

\LA
\noindent Cónsequens ígitur erat, ut quos a vúlnerum dolóre sanáverat, eos alimóniis spiritálibus a jejúnio liberáret. Itaque nemo cibum áccipit Christi, nisi fúerit ante sanátus; et illi qui vocántur ad cœnam, prius vocándo sanántur. Si claudus fuit, gradiéndi facultátem, ut veníret, accépit: si lúmine oculórum privátus, domum útique Dómini, nisi refúsa luce, intráre non pótuit.\par
\switchcolumn
\EN
\noindent And then it came to pass that, as He had healed them that had need of healing, He fed their hunger with ghostly meat. Thus it is that no man taketh Christ's meat, unless he be first healed, and they, that are bidden to the supper, are first cured by the invitation. The lame receive the power to walk, that they may be able to come; the blind cannot see the door of the house of the Lord, unless light be given them.\par
\switchcolumn*

\LA
\begin{Resp}\Rsign Duo Séraphim clamábant alter ad álterum: \Star{} Sanctus, sanctus, sanctus Dóminus Deus Sábaoth: * Plena est omnis terra glória ejus.\end{Resp}
\begin{Resp}\Vsign Tres sunt qui testimónium dant in cælo: Pater, Verbum, et Spíritus Sanctus: et hi tres unum sunt.\end{Resp}
\begin{Resp}\Rsign Sanctus, sanctus, sanctus Dóminus Deus Sábaoth.\end{Resp}
\begin{Resp}\Vsign Glória Patri, et Fílio, \Star{} et Spirítui Sancto.\end{Resp}
\begin{Resp}\Rsign Plena est omnis terra glória ejus.\end{Resp}
\switchcolumn
\EN
\begin{Resp}\Rsign One Seraph cried unto another: \Star{} Holy, Holy, Holy is the Lord God of hosts; the whole earth is full of His glory.\end{Resp}
\begin{Resp}\Vsign There are Three That bear record in heaven: the Father, the Word, and the Holy Ghost, and these Three are One.\end{Resp}
\begin{Resp}\Rsign Holy, Holy, Holy is the Lord God of hosts.\end{Resp}
\begin{Resp}\Vsign Glory be to the Father, and to the Son, \Star{} and to the Holy Ghost.\end{Resp}
\begin{Resp}\Rsign The whole earth is full of His glory.\end{Resp}
\switchcolumn*

\LA
\LessonHead{Lectio IX}
\switchcolumn
\EN
\LessonHead{Lesson IX}
\switchcolumn*

\LA
\noindent Ubíque ígitur mystérii ordo servátur, ut prius per remissiónem peccatórum vulnéribus medicína tribuátur, post alimónia mensæ cæléstis exúberat; quamquam nondum validióribus hæc turba reficiátur aliméntis, neque Christi córpore et sánguine jejúna solidióris fídei corda pascántur. Lacte, inquit, vos potávi, non esca; nondum enim poterátis, sed nec adhuc quidem potéstis. In modum lactis quinque sunt panes: esca autem solídior, corpus est Christus; potus veheméntior, sanguis est Dómini.\par
\switchcolumn
\EN
\noindent Everywhere is preserved the order of the Sacraments. The sinful soul is first healed by the remission of sins, and afterward is filled at the Table of the Lord albeit this multitude now present is of such as do not yet feed on those strong meats, nor pasture their starving spirits upon the Body and Blood of Christ, as do they of a manlier faith. To use the words of Paul, I have fed you with milk and not with meat, for hitherto ye were not able to bear it, neither yet now are ye able. (1 Cor. iii. 2.) The five loaves are, as it were, your milk; the stronger meat will be the Body of Christ; the more generous cup, the Blood of the Lord.\par
\switchcolumn*

\LA
\TeDeum{Te Deum}
\switchcolumn
\EN
\TeDeum{Te Deum}
\switchcolumn*

\end{paracol}

\Office{Feria secunda infra Hebdomadam VI post Octavam Pentecostes}{Monday of the Sixth Week after Pentecost}{Feria}
\addcontentsline{toc}{subsection}{Feria secunda infra Hebdomadam VI post Octavam Pentecostes \textit{— Monday of the Sixth Week after Pentecost}}
\begin{paracol}{2}
\LA
\LessonHead{Lectio I}
\switchcolumn
\EN
\LessonHead{Lesson I}
\switchcolumn*

\LA
\LessonTitle{De libro secúndo Regum}
\switchcolumn
\EN
\LessonTitle{Lesson from the second book of Samuel}
\switchcolumn*

\LA
\Cite{2 Reg 13:22-25}
\switchcolumn
\EN
\Cite{2 Sam 13:22-25}
\switchcolumn*

\LA
\noindent Porro non est locútus Absalom ad Amnon nec malum nec bonum; óderat enim Absalom Amnon, eo quod violásset Thamar sorórem suam. Factum est autem post tempus biénnii ut tonderéntur oves Absalom in Baálhasor, quæ est juxta Ephraim; et vocávit Absalom omnes fílios regis, venítque ad regem et ait ad eum: Ecce tondéntur oves servi tui; véniat, oro, rex cum servis suis ad servum suum. Dixítque rex ad Absalom: Noli, fili mi, noli rogáre ut veniámus omnes et gravémus te. Cum autem cógeret eum, et noluísset ire, benedíxit ei.\par
\switchcolumn
\EN
\noindent But Absalom spoke not to Amnon neither good nor evil: for Absalom hated Amnon because he had ravished his sister Thamar. And it came to pass after two years, that the sheep of Absalom were shorn in Baalhasor, which is near Ephraim: and Absalom invited all the king's sons: And he came to the king, and said to him: Behold thy servant's sheep are shorn. Let the king, I pray, with his servants come to his servant. And the king said to Absalom: Nay, my son, do not ask that we should all come, and be chargeable to thee. And when he pressed him, and he would not go, he blessed him.\par
\switchcolumn*

\LA
\begin{Resp}\Rsign Recordáre, Dómine, testaménti tui, et dic Angelo percutiénti: Cesset jam manus tua, \Star{} Ut non desolétur terra, et ne perdas omnem ánimam vivam.\end{Resp}
\begin{Resp}\Vsign Ego sum qui peccávi, ego qui iníque egi: isti qui oves sunt, quid fecérunt? Avertátur, óbsecro, furor tuus, Dómine, a pópulo tuo.\end{Resp}
\begin{Resp}\Rsign Ut non desolétur terra, et ne perdas omnem ánimam vivam.\end{Resp}
\switchcolumn
\EN
\begin{Resp}\Rsign Remember, O Lord, thy covenant, and say unto the destroying Angel: Stay now thine hand; \Star{} That the land be not utterly laid waste, and that thou destroy not every living soul.\end{Resp}
\begin{Resp}\Vsign Even I it is that have sinned, and done evil indeed; but these sheep, what have they done? Let thine anger, I pray thee, O Lord, be turned away from thy people.\end{Resp}
\begin{Resp}\Rsign That the land be not utterly laid waste, and that Thou destroy not every living soul.\end{Resp}
\switchcolumn*

\LA
\LessonHead{Lectio II}
\switchcolumn
\EN
\LessonHead{Lesson II}
\switchcolumn*

\LA
\Cite{2 Reg 13:26-29}
\switchcolumn
\EN
\Cite{2 Sam 13:26-29}
\switchcolumn*

\LA
\noindent Et ait Absalom: Si non vis veníre, véniat, óbsecro, nobíscum saltem Amnon frater meus. Dixítque ad eum rex: Non est necésse ut vadat tecum. Coégit ítaque Absalom eum et dimísit cum eo Amnon et univérsos fílios regis. Fecerátque Absalom convívium quasi convívium regis. Præcéperat autem Absalom púeris suis dicens: Observáte: cum temuléntus fúerit Amnon vino, et díxero vobis: Percútite eum et interfícite; nolíte timére, ego enim sum qui præcípio: roborámini et estóte viri fortes. Fecérunt ergo púeri Absalom advérsum Amnon sicut præcéperat eis Absalom. Surgentésque omnes fílii regis ascendérunt sínguli mulas suas et fugérunt.\par
\switchcolumn
\EN
\noindent And Absalom said: If thou wilt not come, at least let my brother Amnon, I beseech thee, come with us. And the king said to him: It is not necessary that he should go with thee. But Absalom pressed him, so that he let Amnon and all the king's sons go with him. And Absalom made a feast as it were the feast of a king. And Absalom had commanded his servants, saying: Take notice when Amnon shall be drunk with wine, and when I shall say to you: Strike him, and kill him, fear not: for it is I that command you: take courage, and be valiant men. And the servants of Absalom did to Amnon as Absalom had commanded them. And all the king's sons arose and got up every man upon his mule, and fled.\par
\switchcolumn*

\LA
\begin{Resp}\Rsign Exaudísti, Dómine, oratiónem servi tui, ut ædificárem templum nómini tuo: \Star{} Bénedic et sanctífica domum istam in sempitérnum, Deus Israël.\end{Resp}
\begin{Resp}\Vsign Dómine, qui custódis pactum cum servis tuis, qui ámbulant coram te in toto corde suo.\end{Resp}
\begin{Resp}\Rsign Bénedic et sanctífica domum istam in sempitérnum, Deus Israël.\end{Resp}
\switchcolumn
\EN
\begin{Resp}\Rsign O Lord, Thou hast hearkened unto the prayer of thy servant, that I might build a temple unto thy Name, \Star{} O God of Israel, bless Thou, and hallow this house for ever.\end{Resp}
\begin{Resp}\Vsign O Lord, Who keepest covenant with thy servants that walk before thee in all their heart.\end{Resp}
\begin{Resp}\Rsign O God of Israel, bless Thou, and hallow this house for ever.\end{Resp}
\switchcolumn*

\LA
\LessonHead{Lectio III}
\switchcolumn
\EN
\LessonHead{Lesson III}
\switchcolumn*

\LA
\Cite{2 Reg 13:30-34}
\switchcolumn
\EN
\Cite{2 Sam 13:30-34}
\switchcolumn*

\LA
\noindent Cumque adhuc pérgerent in itínere, fama pervénit ad David dicens: Percússit Absalom omnes fílios regis, et non remánsit ex eis saltem unus. Surréxit ítaque rex et scidit vestiménta sua et cécidit super terram, et omnes servi illíus, qui assistébant ei, scidérunt vestiménta sua. Respóndens autem Jónadab fílius Sémmaa fratris David dixit: Ne ǽstimet dóminus meus rex quod omnes púeri fílii regis occísi sint; Amnon solus mórtuus est, quóniam in ore Absalom erat pósitus ex die qua oppréssit Thamar sorórem ejus. Nunc ergo ne ponat dóminus meus rex super cor suum verbum istud dicens: Omnes fílii regis occísi sunt, quóniam Amnon solus mórtuus est. Fugit autem Absalom.\par
\switchcolumn
\EN
\noindent And while they were yet in the way, a rumour came to David, saying: Absalom hath slain all the king's sons, and there is not one of them left. Then the king rose up, and rent his garments: and fell upon the ground, and all his servants, that stood about him, rent their garments. But Jonadab the son of Semmaa David's brother answering, said: Let not my lord the king think that all the king's sons are slain: Amnon only is dead, for he was appointed by the mouth of Absalom from the day that he ravished his sister Thamar. Now therefore let not my lord the king take this thing into his heart, saying: All the king's sons are slain: for Amnon only is dead. But Absalom fled away:\par
\switchcolumn*

\LA
\begin{Resp}\Rsign Audi, Dómine, hymnum et oratiónem, quam servus tuus orat coram te hódie, ut sint óculi tui apérti, et aures tuæ inténtæ, \Star{} Super domum istam die ac nocte.\end{Resp}
\begin{Resp}\Vsign Réspice, Dómine, de sanctuário tuo, et de excélso cælórum habitáculo.\end{Resp}
\begin{Resp}\Rsign Super domum istam die ac nocte.\end{Resp}
\begin{Resp}\Vsign Glória Patri, et Fílio, \Star{} et Spirítui Sancto.\end{Resp}
\begin{Resp}\Rsign Super domum istam die ac nocte.\end{Resp}
\switchcolumn
\EN
\begin{Resp}\Rsign Hearken, O Lord, unto the cry and to the prayer which thy servant prayeth before thee today, that thine eyes may be open and thine ears attend; \Star{} Toward this house day and night.\end{Resp}
\begin{Resp}\Vsign Look down from thine high and holy place, O Lord, even from heaven thy dwelling.\end{Resp}
\begin{Resp}\Rsign Toward this house, day and night.\end{Resp}
\begin{Resp}\Vsign Glory be to the Father, and to the Son, \Star{} and to the Holy Ghost.\end{Resp}
\begin{Resp}\Rsign Toward this house, day and night.\end{Resp}
\switchcolumn*

\end{paracol}

\Office{Feria tertia infra Hebdomadam VI post Octavam Pentecostes}{Tuesday of the Sixth Week after Pentecost}{Feria}
\addcontentsline{toc}{subsection}{Feria tertia infra Hebdomadam VI post Octavam Pentecostes \textit{— Tuesday of the Sixth Week after Pentecost}}
\begin{paracol}{2}
\LA
\LessonHead{Lectio I}
\switchcolumn
\EN
\LessonHead{Lesson I}
\switchcolumn*

\LA
\LessonTitle{De libro secúndo Regum}
\switchcolumn
\EN
\LessonTitle{Lesson from the second book of Samuel}
\switchcolumn*

\LA
\Cite{2 Reg 14:4-7}
\switchcolumn
\EN
\Cite{2 Sam 14:4-7}
\switchcolumn*

\LA
\noindent Cum ingréssa fuísset múlier Thecuítis ad regem, cécidit coram eo super terram et adorávit et dixit: Serva me, rex. Et ait ad eam rex: Quid causæ habes? Quæ respóndit: Heu, múlier vídua ego sum; mórtuus est enim vir meus. Et ancíllæ tuæ erant duo fílii, qui rixáti sunt advérsum se in agro, nullúsque erat qui eos prohibére posset, et percússit alter álterum et interfécit eum. Et ecce consúrgens univérsa cognátio advérsum ancíllam tuam dicit: Trade eum qui percússit fratrem suum, ut occidámus eum pro ánima fratris sui quem interfécit, et deleámus herédem. Et quærunt exstínguere scintíllam meam quæ relícta est, ut non supérsit viro meo nomen, et relíquiæ super terram.\par
\switchcolumn
\EN
\noindent And when the woman of Thecua was come in to the king, she fell before him upon the ground, and worshipped, and said: Save me, O king. And the king said to her: What is the matter with thee? She answered: Alas, I am a widow woman: for my husband is dead. And thy handmaid had two sons: and they quarrelled with each other in the field, and there was none to part them: and the one struck the other, and slew him. And behold the whole kindred rising against thy handmaid, saith: Deliver him that hath slain his brother, that we may kill him for the life of his brother, whom he slew, and that we may destroy the heir: and they seek to quench my spark which is left, and will leave my husband no name, nor remainder upon the earth.\par
\switchcolumn*

\LA
\begin{Resp}\Rsign Dómine, si convérsus fúerit pópulus tuus, et oráverit ad sanctuárium tuum: \Star{} Tu exáudies de cælo, Dómine, et líbera eos de mánibus inimicórum suórum.\end{Resp}
\begin{Resp}\Vsign Si peccáverit in te pópulus tuus, et convérsus égerit pœniténtiam, veniénsque oráverit in isto loco.\end{Resp}
\begin{Resp}\Rsign Tu exáudies de cælo, Dómine, et líbera eos de mánibus inimicórum suórum.\end{Resp}
\switchcolumn
\EN
\begin{Resp}\Rsign Lord, when thy people shall turn again to thee, and shall pray unto thee in this house; \Star{} Then hear Thou in heaven, O Lord, and deliver them out of the hand of their enemies.\end{Resp}
\begin{Resp}\Vsign If thy people sin against thee, and turn again, and repent, and come and pray unto thee in this house.\end{Resp}
\begin{Resp}\Rsign Then hear Thou in heaven, O Lord, and deliver them out of the hand of their enemies.\end{Resp}
\switchcolumn*

\LA
\LessonHead{Lectio II}
\switchcolumn
\EN
\LessonHead{Lesson II}
\switchcolumn*

\LA
\Cite{2 Reg 14:10-14}
\switchcolumn
\EN
\Cite{2 Sam 14:10-14}
\switchcolumn*

\LA
\noindent Et ait rex: Qui contradíxerit tibi, adduc eum ad me, et ultra non addet ut tangat te. Quæ ait: Recordétur rex Dómini Dei sui, ut non multiplicéntur próximi sánguinis ad ulciscéndum, et nequáquam interfíciant fílium meum. Qui ait: Vivit Dóminus, quia non cadet de capíllis fílii tui super terram. Dixit ergo múlier: Loquátur ancílla tua ad dóminum meum regem verbum. Et ait: Lóquere. Dixítque múlier: Quare cogitásti hujuscémodi rem contra pópulum Dei et locútus est rex verbum istud, ut peccet et non redúcat ejéctum suum? Omnes mórimur et quasi aquæ dilábimur in terram, quæ non revertúntur, nec vult Deus períre ánimam, sed retráctat cógitans ne pénitus péreat qui abjéctus est.\par
\switchcolumn
\EN
\noindent And the king said: If any one shall say ought against thee, bring him to me, and be shall not touch thee any more. And she said: Let the king remember the Lord his God, that the next of kin be not multiplied to take revenge, and that they may not kill my son. And he said: As the Lord liveth, there shall not one hair of thy son fall to the earth. Then the woman said: Let thy handmaid speak one word to my lord the king. And he said: Speak. And the woman said: Why hast thou thought such a thing against the people of God, and why hath the king spoken this word, to sin, and not bring home again his own exile? We all die, and like waters that return no more, we fall down into the earth: neither will God have a soul to perish, but recalleth, meaning that he that is cast off should not altogether perish.\par
\switchcolumn*

\LA
\begin{Resp}\Rsign Factum est, dum tólleret Dóminus Elíam per túrbinem in cælum, \Star{} Eliséus clamábat, dicens: Pater mi, pater mi, currus Israël, et auríga ejus.\end{Resp}
\begin{Resp}\Vsign Cumque pérgerent, et incedéntes sermocinaréntur, ecce currus ígneus et equi ígnei divisérunt utrúmque, et ascéndit Elías per túrbinem in cælum.\end{Resp}
\begin{Resp}\Rsign Eliséus clamábat, dicens: Pater mi, pater mi, currus Israël, et auríga ejus.\end{Resp}
\switchcolumn
\EN
\begin{Resp}\Rsign And it came to pass when the Lord would take up Elijah into heaven by a whirlwind \Star{} Elisha cried, and said: My father, my father, the chariot of Israel, and the horsemen thereof.\end{Resp}
\begin{Resp}\Vsign And as they still went on, and talked, behold, there appeared a chariot of fire, and horses of fire, and parted them both asunder, and Elijah went up by a whirlwind into heaven.\end{Resp}
\begin{Resp}\Rsign Elisha cried, and said: My father, my father, the chariot of Israel, and the horsemen thereof.\end{Resp}
\switchcolumn*

\LA
\LessonHead{Lectio III}
\switchcolumn
\EN
\LessonHead{Lesson III}
\switchcolumn*

\LA
\Cite{2 Reg 14:19-21}
\switchcolumn
\EN
\Cite{2 Sam 14:19-21}
\switchcolumn*

\LA
\noindent Et ait rex: Numquid manus Joab tecum est in ómnibus istis? Respóndit múlier et ait: Per salútem ánimæ tuæ, dómine mi rex, nec ad sinístram nec ad déxteram est ex ómnibus his quæ locútus est dóminus meus rex; servus enim tuus Joab ipse præcépit mihi et ipse pósuit in os ancíllæ tuæ ómnia verba hæc. Ut vérterem figúram sermónis hujus, servus tuus Joab præcépit istud; tu autem, dómine mi rex, sápiens es, sicut habet sapiéntiam Angelus Dei, ut intéllegas ómnia super terram. Et ait rex ad Joab: Ecce placátus feci verbum tuum; vade ergo, et révoca púerum Absalom.\par
\switchcolumn
\EN
\noindent And the king said: Is not the hand of Joab with thee in all this? The woman answered, and said: By the health of thy soul, my lord, O king, it is neither on the left hand, nor on the right, in all these things which my lord the king hath spoken: for thy servant Joab, he commanded me, and he put all these words into the mouth of thy handmaid. That I should come about with this form of speech, thy servant Joab, commanded this: but thou, my lord, O king, art wise, according to the wisdom of an angel of God, to understand all things upon earth. And the king said to Joab: Behold I am appeased and have granted thy request: Go therefore and fetch back the boy Absalom\par
\switchcolumn*

\LA
\begin{Resp}\Rsign Ego te tuli de domo patris tui, dicit Dóminus, et pósui te páscere gregem pópuli mei: \Star{} Et fui tecum in ómnibus ubicúmque ambulásti, firmans regnum tuum in ætérnum.\end{Resp}
\begin{Resp}\Vsign Fecíque tibi nomen grande, juxta nomen magnórum, qui sunt in terra: et réquiem dedi tibi ab ómnibus inimícis tuis.\end{Resp}
\begin{Resp}\Rsign Et fui tecum in ómnibus ubicúmque ambulásti, firmans regnum tuum in ætérnum.\end{Resp}
\begin{Resp}\Vsign Glória Patri, et Fílio, \Star{} et Spirítui Sancto.\end{Resp}
\begin{Resp}\Rsign Et fui tecum in ómnibus ubicúmque ambulásti, firmans regnum tuum in ætérnum.\end{Resp}
\switchcolumn
\EN
\begin{Resp}\Rsign Thus saith the Lord: I took thee out of thy father's house, and appointed thee to be ruler over My people, over Israel. \Star{} And I was with thee whithersoever thou wentest, to establish thy kingdom for ever.\end{Resp}
\begin{Resp}\Vsign And I have made thee a great name, like unto the name of the great men that are in the earth, and have caused thee to rest from all thine enemies.\end{Resp}
\begin{Resp}\Rsign And I was with thee whithersoever thou wentest, to establish thy kingdom for ever.\end{Resp}
\begin{Resp}\Vsign Glory be to the Father, and to the Son, \Star{} and to the Holy Ghost.\end{Resp}
\begin{Resp}\Rsign And I was with thee whithersoever thou wentest, to establish thy kingdom for ever.\end{Resp}
\switchcolumn*

\end{paracol}

\Office{Feria quarta infra Hebdomadam VI post Octavam Pentecostes}{Wednesday of the Sixth Week after Pentecost}{Feria}
\addcontentsline{toc}{subsection}{Feria quarta infra Hebdomadam VI post Octavam Pentecostes \textit{— Wednesday of the Sixth Week after Pentecost}}
\begin{paracol}{2}
\LA
\LessonHead{Lectio I}
\switchcolumn
\EN
\LessonHead{Lesson I}
\switchcolumn*

\LA
\LessonTitle{De libro secúndo Regum}
\switchcolumn
\EN
\LessonTitle{Lesson from the second book of Samuel}
\switchcolumn*

\LA
\Cite{2 Reg 15:1-3}
\switchcolumn
\EN
\Cite{2 Sam 15:1-3}
\switchcolumn*

\LA
\noindent Igitur post hæc fecit sibi Absalom currus et équites et quinquagínta viros, qui præcéderent eum. Et mane consúrgens Absalom stabat juxta intróitum portæ et omnem virum, qui habébat negótium ut veníret ad regis judícium, vocábat Absalom ad se et dicébat: De qua civitáte es tu? Qui respóndens ajébat: Ex una tribu Israël ego sum servus tuus. Respondebátque ei Absalom: Vidéntur mihi sermónes tui boni et justi; sed non est qui te áudiat constitútus a rege.\par
\switchcolumn
\EN
\noindent Now after these things Absalom made himself chariots, and horsemen, and fifty men to run before him. And Absalom rising up early stood by the entrance of the gate, and when any man had business to come to the king's judgment, Absalom called him to him, and said: Of what city art thou? He answered, and said: thy servant is of such a tribe of Israel. And Absalom answered him: thy words seem to me good and just. But there is no man appointed by the king to hear thee.\par
\switchcolumn*

\LA
\begin{Resp}\Rsign Peccávi super númerum arénæ maris, et multiplicáta sunt peccáta mea: et non sum dignus vidére altitúdinem cæli præ multitúdine iniquitátis meæ: quóniam irritávi iram tuam, \Star{} Et malum coram te feci.\end{Resp}
\begin{Resp}\Vsign Quóniam iniquitátem meam ego cognósco: et delíctum meum contra me est semper, quia tibi soli peccávi.\end{Resp}
\begin{Resp}\Rsign Et malum coram te feci.\end{Resp}
\switchcolumn
\EN
\begin{Resp}\Rsign My sins are many, yea, they are more in number than the sands of the sea; I am not worthy to look up toward heaven because of the multitude of my iniquities; for I have provoked thee to anger; \Star{} And done evil in thy sight.\end{Resp}
\begin{Resp}\Vsign For I acknowledge my transgression, and my sin is ever before me, for against thee only have I sinned.\end{Resp}
\begin{Resp}\Rsign And done evil in thy sight.\end{Resp}
\switchcolumn*

\LA
\LessonHead{Lectio II}
\switchcolumn
\EN
\LessonHead{Lesson II}
\switchcolumn*

\LA
\Cite{2 Reg 15:3-6}
\switchcolumn
\EN
\Cite{2 Sam 15:3-6}
\switchcolumn*

\LA
\noindent Dicebátque Absalom: Quis me constítuat júdicem super terram, ut ad me véniant omnes qui habent negótium et juste júdicem? Sed et, cum accéderet ad eum homo ut salutáret illum, extendébat manum suam et apprehéndens osculabátur eum. Faciebátque hoc omni Israël veniénti ad judícium ut audirétur a rege, et sollicitábat corda virórum Israël.\par
\switchcolumn
\EN
\noindent O that they would make me judge over the land, that all that have business might come to me, that I might do them justice. Moreover when any man came to him to salute him, he put forth his hand, and took him, and kissed him. And this he did to all Israel that came for judgment, to be heard by the king, and he enticed the hearts of the men of Israel.\par
And Absalom said:\par
\switchcolumn*

\LA
\begin{Resp}\Rsign Exaudísti, Dómine, oratiónem servi tui, ut ædificárem templum nómini tuo: \Star{} Bénedic et sanctífica domum istam in sempitérnum, Deus Israël.\end{Resp}
\begin{Resp}\Vsign Dómine, qui custódis pactum cum servis tuis, qui ámbulant coram te in toto corde suo.\end{Resp}
\begin{Resp}\Rsign Bénedic et sanctífica domum istam in sempitérnum, Deus Israël.\end{Resp}
\switchcolumn
\EN
\begin{Resp}\Rsign O Lord, Thou hast hearkened unto the prayer of thy servant, that I might build a temple unto thy Name, \Star{} O God of Israel, bless Thou, and hallow this house for ever.\end{Resp}
\begin{Resp}\Vsign O Lord, Who keepest covenant with thy servants that walk before thee in all their heart.\end{Resp}
\begin{Resp}\Rsign O God of Israel, bless Thou, and hallow this house for ever.\end{Resp}
\switchcolumn*

\LA
\LessonHead{Lectio III}
\switchcolumn
\EN
\LessonHead{Lesson III}
\switchcolumn*

\LA
\Cite{2 Reg 15:7-10}
\switchcolumn
\EN
\Cite{2 Sam 15:7-10}
\switchcolumn*

\LA
\noindent Post quadragínta autem annos, dixit Absalom ad regem David: Vadam et reddam vota mea, quæ vovi Dómino, in Hebron. Vovens enim vovit servus tuus, cum esset in Gessur Sýriæ, dicens: Si redúxerit me Dóminus in Jerúsalem, sacrificábo Dómino. Dixítque ei rex David: Vade in pace. Et surréxit et ábiit in Hebron. Misit autem Absalom exploratóres in univérsas tribus Israël dicens: Statim ut audiéritis clangórem búccinæ, dícite: Regnávit Absalom in Hebron.\par
\switchcolumn
\EN
\noindent And after forty years, Absalom said to king David: Let me go, and pay my vows which I have vowed to the Lord in Hebron. For thy servant made a vow, when he was in Gessur of Syria, saying: If the Lord shall bring me again into Jerusalem I will offer sacrifice to the Lord. And king David said to him: Go in peace. And he arose, and went to Hebron. And Absalom sent spies into all the tribes of Israel, saying: As soon as you shall hear the sound of the trumpet, say ye: Absalom reigneth in Hebron.\par
\switchcolumn*

\LA
\begin{Resp}\Rsign Audi, Dómine, hymnum et oratiónem, quam servus tuus orat coram te hódie, ut sint óculi tui apérti, et aures tuæ inténtæ, \Star{} Super domum istam die ac nocte.\end{Resp}
\begin{Resp}\Vsign Réspice, Dómine, de sanctuário tuo, et de excélso cælórum habitáculo.\end{Resp}
\begin{Resp}\Rsign Super domum istam die ac nocte.\end{Resp}
\begin{Resp}\Vsign Glória Patri, et Fílio, \Star{} et Spirítui Sancto.\end{Resp}
\begin{Resp}\Rsign Super domum istam die ac nocte.\end{Resp}
\switchcolumn
\EN
\begin{Resp}\Rsign Hearken, O Lord, unto the cry and to the prayer which thy servant prayeth before thee today, that thine eyes may be open and thine ears attend; \Star{} Toward this house day and night.\end{Resp}
\begin{Resp}\Vsign Look down from thine high and holy place, O Lord, even from heaven thy dwelling.\end{Resp}
\begin{Resp}\Rsign Toward this house, day and night.\end{Resp}
\begin{Resp}\Vsign Glory be to the Father, and to the Son, \Star{} and to the Holy Ghost.\end{Resp}
\begin{Resp}\Rsign Toward this house, day and night.\end{Resp}
\switchcolumn*

\end{paracol}

\Office{Feria quinta infra Hebdomadam VI post Octavam Pentecostes}{Thursday of the Sixth Week after Pentecost}{Feria}
\addcontentsline{toc}{subsection}{Feria quinta infra Hebdomadam VI post Octavam Pentecostes \textit{— Thursday of the Sixth Week after Pentecost}}
\begin{paracol}{2}
\LA
\LessonHead{Lectio I}
\switchcolumn
\EN
\LessonHead{Lesson I}
\switchcolumn*

\LA
\LessonTitle{De libro secúndo Regum}
\switchcolumn
\EN
\LessonTitle{Lesson from the second book of Samuel}
\switchcolumn*

\LA
\Cite{2 Reg 15:13-15}
\switchcolumn
\EN
\Cite{2 Sam 15:13-15}
\switchcolumn*

\LA
\noindent Venit ígitur núntius ad David dicens: Toto corde univérsus Israël séquitur Absalom. Et ait David servis suis qui erant cum eo in Jerúsalem: Súrgite, fugiámus; neque enim erit nobis effúgium a fácie Absalom. Festináte égredi, ne forte véniens óccupet nos et impéllat super nos ruínam et percútiat civitátem in ore gládii. Dixerúntque servi regis ad eum: Omnia, quæcúmque præcéperit dóminus noster rex, libénter exsequémur servi tui.\par
\switchcolumn
\EN
\noindent And there came a messenger to David, saying: All Israel with their whole heart followeth Absalom. And David said to his servants, that were with him in Jerusalem: Arise and let us flee: for we shall not escape else from the face of Absalom: make haste to go out, lest he come and overtake us, and bring ruin upon us, and smite the city with the edge of the sword. And the king's servants said to him: Whatsoever our lord the king shall command, we thy servants will willingly execute.\par
\switchcolumn*

\LA
\begin{Resp}\Rsign Præparáte corda vestra Dómino, et servíte illi soli: \Star{} Et liberábit vos de mánibus inimicórum vestrórum.\end{Resp}
\begin{Resp}\Vsign Convertímini ad eum in toto corde vestro, et auférte deos aliénos de médio vestri.\end{Resp}
\begin{Resp}\Rsign Et liberábit vos de mánibus inimicórum vestrórum.\end{Resp}
\switchcolumn
\EN
\begin{Resp}\Rsign Prepare your hearts unto the Lord, and serve Him Only \Star{} And He will deliver you out of the hand of your enemies.\end{Resp}
\begin{Resp}\Vsign Return unto Him with all your hearts, and put away the strange gods from among you.\end{Resp}
\begin{Resp}\Rsign And He will deliver you out of the hand of your enemies.\end{Resp}
\switchcolumn*

\LA
\LessonHead{Lectio II}
\switchcolumn
\EN
\LessonHead{Lesson II}
\switchcolumn*

\LA
\Cite{2 Reg 15:16-18}
\switchcolumn
\EN
\Cite{2 Sam 15:16-18}
\switchcolumn*

\LA
\noindent Egréssus est ergo rex et univérsa domus ejus pédibus suis; et derelíquit rex decem mulíeres concubínas ad custodiéndam domum. Egressúsque rex et omnis Israël pédibus suis stetit procul a domo. Et univérsi servi ejus ambulábant juxta eum; et legiónes Ceréthi et Pheléthi et omnes Gethǽi pugnatóres válidi sexcénti viri, qui secúti eum fúerant de Geth, pédites præcedébant regem.\par
\switchcolumn
\EN
\noindent And the king went forth, and all his household on foot: and the king left ten women his concubines to keep the house: And the king going forth and all Israel on foot, stood afar off from the house: And all his servants walked by him, and the bands of the Cerethi, and the Phelethi, and all the Gethites, valiant warriors, six hundred men who had followed him from Geth on foot, went before the king.\par
\switchcolumn*

\LA
\begin{Resp}\Rsign Deus ómnium exaudítor est: ipse misit Angelum suum, et tulit me de óvibus patris mei: \Star{} Et unxit me unctióne misericórdiæ suæ.\end{Resp}
\begin{Resp}\Vsign Dóminus, qui erípuit me de ore leónis, et de manu béstiæ liberávit me.\end{Resp}
\begin{Resp}\Rsign Et unxit me unctióne misericórdiæ suæ.\end{Resp}
\switchcolumn
\EN
\begin{Resp}\Rsign God, Which heareth all, even He sent His Angel, and took me from keeping my father's sheep, and \Star{} Anointed me with the oil of His mercy.\end{Resp}
\begin{Resp}\Vsign The Lord That delivered me out of the mouth of the lion, and out of the paw of the bear\end{Resp}
\begin{Resp}\Rsign And anointed me with the oil of His mercy.\end{Resp}
\switchcolumn*

\LA
\LessonHead{Lectio III}
\switchcolumn
\EN
\LessonHead{Lesson III}
\switchcolumn*

\LA
\Cite{2 Reg 15:19-20}
\switchcolumn
\EN
\Cite{2 Sam 15:19-20}
\switchcolumn*

\LA
\noindent Dixit autem rex ad Ethái Gethǽum: Cur venis nobíscum? Revértere et hábita cum rege, quia peregrínus es et egréssus es de loco tuo. Heri venísti, et hódie compélleris nobíscum égredi? Ego autem vadam quo itúrus sum; revértere et reduc tecum fratres tuos, et Dóminus fáciet tecum misericórdiam et veritátem, quia ostendísti grátiam et fidem.\par
\switchcolumn
\EN
\noindent And the king said to Ethai the Gethite: Why comest thou with us? return and dwell with the king, for thou art a stranger, and art come out of thy own place. Yesterday thou camest, and today shalt thou be forced to go forth with us? but I shall go whither I am going: return thou, and take back thy brethren with thee, and the Lord will show thee mercy, and truth, because thou hast shown grace and fidelity.\par
\switchcolumn*

\LA
\begin{Resp}\Rsign Dóminus, qui erípuit me de ore leónis, et de manu béstiæ liberávit me, \Star{} Ipse me erípiet de mánibus inimicórum meórum.\end{Resp}
\begin{Resp}\Vsign Misit Deus misericórdiam suam et veritátem suam: ánimam meam erípuit de médio catulórum leónum.\end{Resp}
\begin{Resp}\Rsign Ipse me erípiet de mánibus inimicórum meórum.\end{Resp}
\begin{Resp}\Vsign Glória Patri, et Fílio, \Star{} et Spirítui Sancto.\end{Resp}
\begin{Resp}\Rsign Ipse me erípiet de mánibus inimicórum meórum.\end{Resp}
\switchcolumn
\EN
\begin{Resp}\Rsign The Lord That delivered me out of the mouth of the lion, and out of the paw of the bear \Star{} He will deliver me out of the hand of mine enemies.\end{Resp}
\begin{Resp}\Vsign God hath sent forth His mercy and His truth, and delivered my soul from among the lion's whelps.\end{Resp}
\begin{Resp}\Rsign He will deliver me out of the hand of mine enemies.\end{Resp}
\begin{Resp}\Vsign Glory be to the Father, and to the Son, \Star{} and to the Holy Ghost.\end{Resp}
\begin{Resp}\Rsign He will deliver me out of the hand of mine enemies.\end{Resp}
\switchcolumn*

\end{paracol}

\Office{Feria sexta infra Hebdomadam VI post Octavam Pentecostes}{Friday of the Sixth Week after Pentecost}{Feria}
\addcontentsline{toc}{subsection}{Feria sexta infra Hebdomadam VI post Octavam Pentecostes \textit{— Friday of the Sixth Week after Pentecost}}
\begin{paracol}{2}
\LA
\LessonHead{Lectio I}
\switchcolumn
\EN
\LessonHead{Lesson I}
\switchcolumn*

\LA
\LessonTitle{De libro secúndo Regum}
\switchcolumn
\EN
\LessonTitle{Lesson from the second book of Samuel}
\switchcolumn*

\LA
\Cite{2 Reg 16:5-8}
\switchcolumn
\EN
\Cite{2 Sam 16:5-8}
\switchcolumn*

\LA
\noindent Venit ergo rex David usque Bahúrim, et ecce egrediebátur inde vir de cognatióne domus Saul nómine Sémei fílius Gera; procedebátque egrédiens et maledicébat, mittebátque lápides contra David et contra univérsos servos regis David. Omnis autem pópulus et univérsi bellatóres a dextro et a sinístro látere regis incedébant. Ita autem loquebátur Sémei, cum maledíceret regi: Egrédere, egrédere, vir sánguinum et vir Bélial: réddidit tibi Dóminus univérsum sánguinem domus Saul, quóniam invasísti regnum pro eo.\par
\switchcolumn
\EN
\noindent And king David came as far as Bahurim: and behold there came out from thence a man of the kindred of the house of Saul named Semei, the son of Gera, and coming out he cursed as he went on, And he threw stones at David, and at all the servants of king David: and all the people, and all the warriors walked on the right, and on the left side of the king. And thus said Semei when he cursed the king: Come out, come out, thou man of blood, and thou man of Belial. The Lord hath repaid thee for all the blood of the house of Saul: because thou hast usurped the kingdom in his stead, and the Lord hath given the kingdom into the hand of Absalom thy son: and behold thy evils press upon thee, because thou art a man of blood.\par
\switchcolumn*

\LA
\begin{Resp}\Rsign Percússit Saul mille, et David decem míllia: \Star{} Quia manus Dómini erat cum illo: percússit Philisthǽum, et ábstulit oppróbrium ex Israël.\end{Resp}
\begin{Resp}\Vsign Nonne iste est David, de quo canébant in choro, dicéntes: Saul percússit mille, et David decem míllia?\end{Resp}
\begin{Resp}\Rsign Quia manus Dómini erat cum illo: percússit Philisthǽum, et ábstulit oppróbrium ex Israël.\end{Resp}
\switchcolumn
\EN
\begin{Resp}\Rsign Saul hath slain his thousands, and David his ten thousands. \Star{} Because the hand of the Lord was with him, he smote the Philistine, and took away the reproach from Israel.\end{Resp}
\begin{Resp}\Vsign Is not this David? Did they not sing one to another of him in dances, saying: Saul hath slain his thousands, and David his ten thousands?\end{Resp}
\begin{Resp}\Rsign Because the hand of the Lord was with him, he smote the Philistine, and took away the reproach from Israel.\end{Resp}
\switchcolumn*

\LA
\LessonHead{Lectio II}
\switchcolumn
\EN
\LessonHead{Lesson II}
\switchcolumn*

\LA
\Cite{2 Reg 16:9-10}
\switchcolumn
\EN
\Cite{2 Sam 16:9-10}
\switchcolumn*

\LA
\noindent Dixit autem Abísai fílius Sárviæ regi: Quare maledícit canis hic mórtuus dómino meo regi? vadam et amputábo caput ejus. Et ait rex: Quid mihi et vobis est, fílii Sárviæ? Dimíttite eum ut maledícat; Dóminus enim præcépit ei ut maledíceret David, et quis est qui áudeat dícere quare sic fécerit?\par
\switchcolumn
\EN
\noindent And Abisai the son of Sarvia said to the king: Why should this dead dog curse my lord the king? I will go, and cut off his head. And the king said: What have I to do with you, ye sons of Sarvia? Let him alone and let him curse: for the Lord hath bid him curse David: and who is he that shall dare say, why hath he done so?\par
\switchcolumn*

\LA
\begin{Resp}\Rsign Montes Gélboë, nec ros nec plúvia véniant super vos, \Star{} Ubi cecidérunt fortes Israël.\end{Resp}
\begin{Resp}\Vsign Omnes montes, qui estis in circúitu ejus, vísitet Dóminus: a Gélboë autem tránseat.\end{Resp}
\begin{Resp}\Rsign Ubi cecidérunt fortes Israël.\end{Resp}
\switchcolumn
\EN
\begin{Resp}\Rsign Ye mountains of Gilboa, let there be no dew, neither let there be rain upon you \Star{} For there are the mighty of Israel fallen.\end{Resp}
\begin{Resp}\Vsign All ye mountains that stand round about, the Lord look upon you but let Him pass by Gilboa\end{Resp}
\begin{Resp}\Rsign For there are the mighty of Israel fallen.\end{Resp}
\switchcolumn*

\LA
\LessonHead{Lectio III}
\switchcolumn
\EN
\LessonHead{Lesson III}
\switchcolumn*

\LA
\Cite{2 Reg 16:11-12}
\switchcolumn
\EN
\Cite{2 Sam 16:11-12}
\switchcolumn*

\LA
\noindent Et ait rex Abísai et univérsis servis suis: Ecce fílius meus, qui egréssus est de útero meo, quærit ánimam meam; quanto magis nunc fílius Jémini? Dimíttite eum ut maledícat juxta præcéptum Dómini; si forte respíciat Dóminus afflictiónem meam et reddat mihi Dóminus bonum pro maledictióne hac hodiérna.\par
\switchcolumn
\EN
\noindent And the king said to Abisai, and to all his servants: Behold my son, who came forth from my bowels, seeketh my life: how much more now a son of Jemini? let him alone that he may curse as the Lord hath bidden him. Perhaps the Lord may look upon my affliction, and the Lord may render me good for the cursing of this day.\par
\switchcolumn*

\LA
\begin{Resp}\Rsign Ego te tuli de domo patris tui, dicit Dóminus, et pósui te páscere gregem pópuli mei: \Star{} Et fui tecum in ómnibus ubicúmque ambulásti, firmans regnum tuum in ætérnum.\end{Resp}
\begin{Resp}\Vsign Fecíque tibi nomen grande, juxta nomen magnórum, qui sunt in terra: et réquiem dedi tibi ab ómnibus inimícis tuis.\end{Resp}
\begin{Resp}\Rsign Et fui tecum in ómnibus ubicúmque ambulásti, firmans regnum tuum in ætérnum.\end{Resp}
\begin{Resp}\Vsign Glória Patri, et Fílio, \Star{} et Spirítui Sancto.\end{Resp}
\begin{Resp}\Rsign Et fui tecum in ómnibus ubicúmque ambulásti, firmans regnum tuum in ætérnum.\end{Resp}
\switchcolumn
\EN
\begin{Resp}\Rsign Thus saith the Lord: I took thee out of thy father's house, and appointed thee to be ruler over My people, over Israel. \Star{} And I was with thee whithersoever thou wentest, to establish thy kingdom for ever.\end{Resp}
\begin{Resp}\Vsign And I have made thee a great name, like unto the name of the great men that are in the earth, and have caused thee to rest from all thine enemies.\end{Resp}
\begin{Resp}\Rsign And I was with thee whithersoever thou wentest, to establish thy kingdom for ever.\end{Resp}
\begin{Resp}\Vsign Glory be to the Father, and to the Son, \Star{} and to the Holy Ghost.\end{Resp}
\begin{Resp}\Rsign And I was with thee whithersoever thou wentest, to establish thy kingdom for ever.\end{Resp}
\switchcolumn*

\end{paracol}

\Office{Sabbato infra Hebdomadam VI post Octavam Pentecostes}{Saturday of the Sixth Week after Pentecost}{Feria}
\addcontentsline{toc}{subsection}{Sabbato infra Hebdomadam VI post Octavam Pentecostes \textit{— Saturday of the Sixth Week after Pentecost}}
\begin{paracol}{2}
\LA
\LessonHead{Lectio I}
\switchcolumn
\EN
\LessonHead{Lesson I}
\switchcolumn*

\LA
\LessonTitle{De libro secúndo Regum}
\switchcolumn
\EN
\LessonTitle{Lesson from the second book of Samuel}
\switchcolumn*

\LA
\Cite{2 Reg 18:6-8}
\switchcolumn
\EN
\Cite{2 Sam 18:6-8}
\switchcolumn*

\LA
\noindent Egréssus est pópulus in campum contra Israël, et factum est prǽlium in saltu Ephraim. Et cæsus est ibi pópulus Israël ab exércitu David, fáctaque est plaga magna in die illa vigínti míllium. Fuit autem ibi prǽlium dispérsum super fáciem omnis terræ, et multo plures erant, quos saltus consúmpserat de pópulo, quam hi quos voráverat gládius in die illa.\par
\switchcolumn
\EN
\noindent So the people went out into the field against Israel and the battle was fought in the forest of Ephraim. And the people of Israel were defeated there by David's army, and a great slaughter was made that day of twenty thousand men. And the battle there was scattered over the face of all the country, and there were many more of the people whom the forest consumed, than whom the sword devoured that day.\par
\switchcolumn*

\LA
\begin{Resp}\Rsign Peccávi super númerum arénæ maris, et multiplicáta sunt peccáta mea: et non sum dignus vidére altitúdinem cæli præ multitúdine iniquitátis meæ: quóniam irritávi iram tuam, \Star{} Et malum coram te feci.\end{Resp}
\begin{Resp}\Vsign Quóniam iniquitátem meam ego cognósco: et delíctum meum contra me est semper, quia tibi soli peccávi.\end{Resp}
\begin{Resp}\Rsign Et malum coram te feci.\end{Resp}
\switchcolumn
\EN
\begin{Resp}\Rsign My sins are many, yea, they are more in number than the sands of the sea; I am not worthy to look up toward heaven because of the multitude of my iniquities; for I have provoked thee to anger; \Star{} And done evil in thy sight.\end{Resp}
\begin{Resp}\Vsign For I acknowledge my transgression, and my sin is ever before me, for against thee only have I sinned.\end{Resp}
\begin{Resp}\Rsign And done evil in thy sight.\end{Resp}
\switchcolumn*

\LA
\LessonHead{Lectio II}
\switchcolumn
\EN
\LessonHead{Lesson II}
\switchcolumn*

\LA
\Cite{2 Reg 18:9-12}
\switchcolumn
\EN
\Cite{2 Sam 18:9-12}
\switchcolumn*

\LA
\noindent Accidit autem ut occúrreret Absalom servis David sedens mulo; cumque ingréssus fuísset mulus subter condénsam quercum et magnam, adhǽsit caput ejus quércui et, illo suspénso inter cælum et terram, mulus cui inséderat, pertransívit. Vidit autem hoc quíspiam, et nuntiávit Joab dicens: Vidi Absalom pendére de quercu. Et ait Joab viro qui nuntiáverat ei: Si vidísti, quare non confodísti eum cum terra et ego dedíssem tibi decem argénti siclos et unum bálteum? Qui dixit ad Joab: Si appénderes in mánibus meis mille argénteos, nequáquam mítterem manum meam in fílium regis; audiéntibus enim nobis, præcépit rex tibi, et Abísai, et Ethái, dicens: Custodíte mihi púerum Absalom.\par
\switchcolumn
\EN
\noindent And it happened that Absalom met the servants of David, riding on a mule: and as the mule went under a thick and large oak, his head stuck in the oak: and while he hung between the heaven and the earth, the mule on which he rode passed on. And one saw this and told Joab, saying: I saw Absalom hanging upon an oak. And Joab said to the man that told him: If thou sawest him, why didst thou not stab him to the ground, and I would have given thee ten sicles of silver, and belt? And he said to Joab: If thou wouldst have paid down in my hands a thousand pieces of silver, I would not lay my hands upon the king's son: for in our hearing the king charged thee, and Abisai, and Ethai, saying: Save me the boy Absalom.\par
\switchcolumn*

\LA
\begin{Resp}\Rsign Exaudísti, Dómine, oratiónem servi tui, ut ædificárem templum nómini tuo: \Star{} Bénedic et sanctífica domum istam in sempitérnum, Deus Israël.\end{Resp}
\begin{Resp}\Vsign Dómine, qui custódis pactum cum servis tuis, qui ámbulant coram te in toto corde suo.\end{Resp}
\begin{Resp}\Rsign Bénedic et sanctífica domum istam in sempitérnum, Deus Israël.\end{Resp}
\switchcolumn
\EN
\begin{Resp}\Rsign O Lord, Thou hast hearkened unto the prayer of thy servant, that I might build a temple unto thy Name, \Star{} O God of Israel, bless Thou, and hallow this house for ever.\end{Resp}
\begin{Resp}\Vsign O Lord, Who keepest covenant with thy servants that walk before thee in all their heart.\end{Resp}
\begin{Resp}\Rsign O God of Israel, bless Thou, and hallow this house for ever.\end{Resp}
\switchcolumn*

\LA
\LessonHead{Lectio III}
\switchcolumn
\EN
\LessonHead{Lesson III}
\switchcolumn*

\LA
\Cite{2 Reg 18:14-17}
\switchcolumn
\EN
\Cite{2 Sam 18:14-17}
\switchcolumn*

\LA
\noindent Et ait Joab: Non sicut tu vis, sed aggrédiar eum coram te. Tulit ergo tres lánceas in manu sua et infíxit eas in corde Absalom; cumque adhuc palpitáret hærens in quercu, cucurrérunt decem júvenes armígeri Joab et percutiéntes interfecérunt eum. Cécinit autem Joab búccina et retínuit pópulum, ne persequerétur fugiéntem Israël, volens párcere multitúdini. Et tulérunt Absalom et projecérunt eum in saltu in fóveam grandem et comportavérunt super eum acérvum lápidum magnum nimis.\par
\switchcolumn
\EN
\noindent And Joab said: Not as thou wilt, but will set upon him in thy sight. So he took three lances in his hand, and thrust them into the heart of Absalom: and whilst he yet panted for life, sticking on the oak, Ten young men, armourbearers of Joab, ran up, and striking him slew him. And Joab sounded the trumpet, and kept back the people from pursuing after Israel in their flight, being willing to spare the multitude. And they took Absalom, and cast him into a great pit in the forest, and they laid an exceeding great heap of stones upon him: but all Israel fled to their own dwellings.\par
\switchcolumn*

\LA
\begin{Resp}\Rsign Audi, Dómine, hymnum et oratiónem, quam servus tuus orat coram te hódie, ut sint óculi tui apérti, et aures tuæ inténtæ, \Star{} Super domum istam die ac nocte.\end{Resp}
\begin{Resp}\Vsign Réspice, Dómine, de sanctuário tuo, et de excélso cælórum habitáculo.\end{Resp}
\begin{Resp}\Rsign Super domum istam die ac nocte.\end{Resp}
\begin{Resp}\Vsign Glória Patri, et Fílio, \Star{} et Spirítui Sancto.\end{Resp}
\begin{Resp}\Rsign Super domum istam die ac nocte.\end{Resp}
\switchcolumn
\EN
\begin{Resp}\Rsign Hearken, O Lord, unto the cry and to the prayer which thy servant prayeth before thee today, that thine eyes may be open and thine ears attend; \Star{} Toward this house day and night.\end{Resp}
\begin{Resp}\Vsign Look down from thine high and holy place, O Lord, even from heaven thy dwelling.\end{Resp}
\begin{Resp}\Rsign Toward this house, day and night.\end{Resp}
\begin{Resp}\Vsign Glory be to the Father, and to the Son, \Star{} and to the Holy Ghost.\end{Resp}
\begin{Resp}\Rsign Toward this house, day and night.\end{Resp}
\switchcolumn*

\end{paracol}

\Office{Dominica VII Post Pentecosten}{Seventh Sunday after Pentecost}{Semiduplex}
\addcontentsline{toc}{subsection}{Dominica VII Post Pentecosten \textit{— Seventh Sunday after Pentecost}}
\begin{paracol}{2}
\LA
\LessonHead{Lectio I}
\switchcolumn
\EN
\LessonHead{Lesson I}
\switchcolumn*

\LA
\LessonTitle{Incipit liber tértius Regum}
\switchcolumn
\EN
\LessonTitle{Lesson from the first book of Kings}
\switchcolumn*

\LA
\Cite{3 Reg 1:1-4}
\switchcolumn
\EN
\Cite{1 Kgs 1:1-4}
\switchcolumn*

\LA
\noindent Et rex David senúerat habebátque ætátis plúrimos dies: cumque operirétur véstibus, non calefiébat. Dixérunt ergo ei servi sui: Quærámus dómino nostro regi adolescéntulam vírginem, et stet coram rege et fóveat eum dormiátque in sinu suo et calefáciat dóminum nostrum regem. Quæsiérunt ígitur adolescéntulam speciósam in ómnibus fínibus Israël et invenérunt Abísag Sunamítidem et adduxérunt eam ad regem. Erat autem puélla pulchra nimis dormiebátque cum rege et ministrábat ei; rex vero non cognóvit eam.\par
\switchcolumn
\EN
\noindent Now king David was old, and advanced in years: and when he was covered with clothes, he was not warm. His servants therefore said to him: Let us seek for our lord the king, a young virgin, and let her stand before the king, and cherish him, and sleep in his bosom, and warm our lord the king. So they sought a beautiful young woman in all the coasts of Israel, and they found Abisag a Sunamitess, and brought her to the king. And the damsel was exceeding beautiful, and she slept with the king: and served him, but the king did not know her.\par
\switchcolumn*

\LA
\begin{Resp}\Rsign Præparáte corda vestra Dómino, et servíte illi soli: \Star{} Et liberábit vos de mánibus inimicórum vestrórum.\end{Resp}
\begin{Resp}\Vsign Convertímini ad eum in toto corde vestro, et auférte deos aliénos de médio vestri.\end{Resp}
\begin{Resp}\Rsign Et liberábit vos de mánibus inimicórum vestrórum.\end{Resp}
\switchcolumn
\EN
\begin{Resp}\Rsign Prepare your hearts unto the Lord, and serve Him Only \Star{} And He will deliver you out of the hand of your enemies.\end{Resp}
\begin{Resp}\Vsign Return unto Him with all your hearts, and put away the strange gods from among you.\end{Resp}
\begin{Resp}\Rsign And He will deliver you out of the hand of your enemies.\end{Resp}
\switchcolumn*

\LA
\LessonHead{Lectio II}
\switchcolumn
\EN
\LessonHead{Lesson II}
\switchcolumn*

\LA
\Cite{3 Reg 1:5-8}
\switchcolumn
\EN
\Cite{1 Kgs 1:5-8}
\switchcolumn*

\LA
\noindent Adonías autem fílius Haggith elevabátur, dicens: Ego regnábo. Fecítque sibi currus et équites et quinquagínta viros, qui cúrrerent ante eum. Nec corrípuit eum pater suus aliquándo dicens: Quare hoc fecísti? Erat autem et ipse pulcher valde, secúndus natu post Absalom. Et sermo ei cum Joab fílio Sárviæ et cum Abíathar sacerdóte, qui adjuvábant partes Adoníæ. Sadoc vero sacérdos, et Banájas fílius Jójadæ, et Nathan prophéta, et Sémei et Rei, et robur exércitus David non erat cum Adonía.\par
\switchcolumn
\EN
\noindent And Adonias the son of Haggith exalted himself, saying: I will be king. And he made himself chariots and horsemen, and fifty men to run before him. Neither did his father rebuke him at any time, saying: Why hast thou done this? And he also was very beautiful, the next in birth after Absalom. And he conferred with Joab the son of Sarvia, and with Abiathar the priest, who furthered Adonias's side. But Sadoc the priest, and Banaias the son of Joiada, and Nathan the prophet, and Semei, and Rei, and the strength of David's army was not with Adonias.\par
\switchcolumn*

\LA
\begin{Resp}\Rsign Deus ómnium exaudítor est: ipse misit Angelum suum, et tulit me de óvibus patris mei: \Star{} Et unxit me unctióne misericórdiæ suæ.\end{Resp}
\begin{Resp}\Vsign Dóminus, qui erípuit me de ore leónis, et de manu béstiæ liberávit me.\end{Resp}
\begin{Resp}\Rsign Et unxit me unctióne misericórdiæ suæ.\end{Resp}
\switchcolumn
\EN
\begin{Resp}\Rsign God, Which heareth all, even He sent His Angel, and took me from keeping my father's sheep, and \Star{} Anointed me with the oil of His mercy.\end{Resp}
\begin{Resp}\Vsign The Lord That delivered me out of the mouth of the lion, and out of the paw of the bear\end{Resp}
\begin{Resp}\Rsign And anointed me with the oil of His mercy.\end{Resp}
\switchcolumn*

\LA
\LessonHead{Lectio III}
\switchcolumn
\EN
\LessonHead{Lesson III}
\switchcolumn*

\LA
\Cite{3 Reg 1:11-15}
\switchcolumn
\EN
\Cite{1 Kgs 1:11-15}
\switchcolumn*

\LA
\noindent Dixit ítaque Nathan ad Bethsabée matrem Salomónis: Num audísti quod regnáverit Adonías fílius Haggith, et dóminus noster David hoc ignórat? Nunc ergo veni, áccipe consílium a me et salva ánimam tuam filiíque tui Salomónis. Vade et ingrédere ad regem David et dic ei: Nonne tu, dómine mi rex, jurásti mihi ancíllæ tuæ dicens: Sálomon fílius tuus regnábit post me et ipse sedébit in sólio meo? Quare ergo regnat Adonías? Et, adhuc ibi te loquénte cum rege, ego véniam post te et complébo sermónes tuos. Ingréssa est ítaque Bethsabée ad regem in cubículum.\par
\switchcolumn
\EN
\noindent And Nathan said to Bethsabee the mother of Solomon: Hast thou not heard that Adonias the son of Haggith reigneth, and our lord David knoweth it not? Now then come, take my counsel and save thy life, and the life of thy son Solomon. Go, and get thee in to king David, and say to him: Didst not thou, my lord O king, swear to me thy handmaid, saying: Solomon thy son shall reign after me, and he shall sit on my throne? why then doth Adonias reign? And while thou art yet speaking there with the king, I will come in after thee, and will fill up thy words. So Bethsabee went in to the king into the chamber.\par
\switchcolumn*

\LA
\begin{Resp}\Rsign Dóminus, qui erípuit me de ore leónis, et de manu béstiæ liberávit me, \Star{} Ipse me erípiet de mánibus inimicórum meórum.\end{Resp}
\begin{Resp}\Vsign Misit Deus misericórdiam suam et veritátem suam: ánimam meam erípuit de médio catulórum leónum.\end{Resp}
\begin{Resp}\Rsign Ipse me erípiet de mánibus inimicórum meórum.\end{Resp}
\begin{Resp}\Vsign Glória Patri, et Fílio, \Star{} et Spirítui Sancto.\end{Resp}
\begin{Resp}\Rsign Ipse me erípiet de mánibus inimicórum meórum.\end{Resp}
\switchcolumn
\EN
\begin{Resp}\Rsign The Lord That delivered me out of the mouth of the lion, and out of the paw of the bear \Star{} He will deliver me out of the hand of mine enemies.\end{Resp}
\begin{Resp}\Vsign God hath sent forth His mercy and His truth, and delivered my soul from among the lion's whelps.\end{Resp}
\begin{Resp}\Rsign He will deliver me out of the hand of mine enemies.\end{Resp}
\begin{Resp}\Vsign Glory be to the Father, and to the Son, \Star{} and to the Holy Ghost.\end{Resp}
\begin{Resp}\Rsign He will deliver me out of the hand of mine enemies.\end{Resp}
\switchcolumn*

\LA
\LessonHead{Lectio IV}
\switchcolumn
\EN
\LessonHead{Lesson IV}
\switchcolumn*

\LA
\LessonTitle{Ex Epístola sancti Hierónymi Presbýteri ad Nepotiánum}
\switchcolumn
\EN
\LessonTitle{From the Epistle written to Nepotian by St. Jerome, Priest at Bethlehem.}
\switchcolumn*

\LA
\Source{Epist. 2, tom. 1}
\switchcolumn
\EN
\Source{Ep. ii. tom. i.}
\switchcolumn*

\LA
\noindent David annos natus septuagínta, bellicósus quondam vir, senectúte frigescénte, non póterat calefíeri. Quǽritur ítaque puélla de univérsis fínibus Israël Abísag Sunamítis, quæ cum rege dormíret, et seníle corpus calefáceret. Quæ est ista Sunamítis, uxor et virgo, tam fervens, ut frígidum calefáceret; tam sancta, ut caléntem ad libídinem non provocáret? Expónat sapientíssimus Sálomon patris sui delícias, et pacíficus bellatóris viri narret ampléxus: Pósside sapiéntiam, pósside intelligéntiam. Ne obliviscáris, et ne declináveris a verbis oris mei: neque derelínquas illam, et apprehéndet te: ama illam, et servábit te. Princípium sapiéntiæ, pósside sapiéntiam, et in omni possessióne tua pósside intelligéntiam: circúmda illam, et exaltábit te; honóra illam, et amplexábitur te, ut det cápiti tuo corónam gratiárum. Coróna quoque deliciárum próteget te.\par
\switchcolumn
\EN
\noindent Then David, who had once been a man of war, was seventy years old, the chill of old age came upon him, and he could get no heat. So they sought out for him throughout all the coasts of Israel Abishag the Shunamite, to sleep with the king and to warm his aged body. Who is this Shunamite, wife and yet virgin, so hot, that she could heat the chilly, so holy, that her warmth provoked him not to lust. Let Solomon the Wise explain his father's enjoyment, and the “Peaceful One” tell of the warrior's embraces. “Get wisdom, get understanding, forget it not, neither decline from the words of my mouth, forsake her not, and she shall preserve thee love her, and she shall keep thee. Wisdom is the principal thing; therefore, get wisdom and with all thy getting, get understanding. Exalt her, and she shall promote thee. Honour her, and she shall embrace thee, and shall give to thine head an ornament of grace. She shall compass thee like acrown of delights.” (Prov. iv. 5-9.)\par
\switchcolumn*

\LA
\begin{Resp}\Rsign Percússit Saul mille, et David decem míllia: \Star{} Quia manus Dómini erat cum illo: percússit Philisthǽum, et ábstulit oppróbrium ex Israël.\end{Resp}
\begin{Resp}\Vsign Nonne iste est David, de quo canébant in choro, dicéntes: Saul percússit mille, et David decem míllia?\end{Resp}
\begin{Resp}\Rsign Quia manus Dómini erat cum illo: percússit Philisthǽum, et ábstulit oppróbrium ex Israël.\end{Resp}
\switchcolumn
\EN
\begin{Resp}\Rsign Saul hath slain his thousands, and David his ten thousands. \Star{} Because the hand of the Lord was with him, he smote the Philistine, and took away the reproach from Israel.\end{Resp}
\begin{Resp}\Vsign Is not this David? Did they not sing one to another of him in dances, saying: Saul hath slain his thousands, and David his ten thousands?\end{Resp}
\begin{Resp}\Rsign Because the hand of the Lord was with him, he smote the Philistine, and took away the reproach from Israel.\end{Resp}
\switchcolumn*

\LA
\LessonHead{Lectio V}
\switchcolumn
\EN
\LessonHead{Lesson V}
\switchcolumn*

\LA
\noindent Omnes pene virtútes córporis mutántur in sénibus, et crescénte sola sapiéntia, decréscunt cétera: jejúnia, vigíliæ, chaméuniæ, id est, super paviméntum dormitiónes, huc illúcque discúrsus, peregrinórum suscéptio, defénsio páuperum, instántia oratiónum et perseverántia, visitátio languéntium, labor mánuum unde præbeántur eleemósynæ. Et, ne sermónem lóngius prótraham, cuncta exercéntur, fracto córpore, minóra fiunt.\par
\switchcolumn
\EN
\noindent In old men almost all the powers of the body become weakened, and while wisdom only is increasing, all things else beside wisdom fail. Then faileth strength for fasting, for watching, for “chameuniae,” (that is, sleeping on the floor,) for wandering hither and thither, for receiving strangers, for defending the poor, for instance and constancy in prayer, for visiting the sick, for that work with the hands whence alms are given. I need not treat of this with long talk, but, in short, when the body is broken down, all the works of the body wax enfeebled.\par
\switchcolumn*

\LA
\begin{Resp}\Rsign Montes Gélboë, nec ros nec plúvia véniant super vos, \Star{} Ubi cecidérunt fortes Israël.\end{Resp}
\begin{Resp}\Vsign Omnes montes, qui estis in circúitu ejus, vísitet Dóminus: a Gélboë autem tránseat.\end{Resp}
\begin{Resp}\Rsign Ubi cecidérunt fortes Israël.\end{Resp}
\switchcolumn
\EN
\begin{Resp}\Rsign Ye mountains of Gilboa, let there be no dew, neither let there be rain upon you \Star{} For there are the mighty of Israel fallen.\end{Resp}
\begin{Resp}\Vsign All ye mountains that stand round about, the Lord look upon you but let Him pass by Gilboa\end{Resp}
\begin{Resp}\Rsign For there are the mighty of Israel fallen.\end{Resp}
\switchcolumn*

\LA
\LessonHead{Lectio VI}
\switchcolumn
\EN
\LessonHead{Lesson VI}
\switchcolumn*

\LA
\noindent Nec hoc dico, quod in juvénibus et adhuc solidióris ætátis, his dumtáxat, qui labóre et ardentíssimo stúdio, vitæ quoque sanctimónia, et oratiónis ad Dóminum Jesum frequéntia, sciéntiam consecúti sunt, frígeat sapiéntia, quæ in plerísque sénibus ætáte marcéscit; sed quod adolescéntia multa córporis bella sustíneat, et inter incentíva vitiórum et carnis titillatiónes, quasi ignis in lignis virídibus suffocétur, ut suum non possit explicáre fulgórem. Senéctus vero rursus eórum qui adolescéntiam suam honéstis ártibus instruxérunt, et in lege Dómini meditáti sunt die ac nocte, ætáte fit dóctior, usu trítior, procéssu témporis sapiéntior, et véterum studiórum dulcíssimos fructus metit.\par
\switchcolumn
\EN
\noindent Don't either do I say, on the other hand, that wisdom, which in many old men drivelleth into second childhood, is weak, or wanting in such of the young and stout, as win knowledge by work and earnest study, by holiness of life and instancy of prayer to the Lord Jesus, but this I do say, that the more spiritual faculties have in youth many wrestlings with the body to go through, and that, what with violent provocations to vice, and what with the sensual ticklings of the flesh, they are apt to be smothered like fire among green wood, and not able to blaze forth in all their brightness. But when old age cometh upon them, who have spent their youth in acquiring sound knowledge, and have meditated in the law of the Lord day and night, it hath this effect on them, to make them more learned by their increased years, more experienced by constant use, more wise, through the advance of time and, in short, doth offer them the rich harvest of their past diligence.\par
\switchcolumn*

\LA
\begin{Resp}\Rsign Ego te tuli de domo patris tui, dicit Dóminus, et pósui te páscere gregem pópuli mei: \Star{} Et fui tecum in ómnibus ubicúmque ambulásti, firmans regnum tuum in ætérnum.\end{Resp}
\begin{Resp}\Vsign Fecíque tibi nomen grande, juxta nomen magnórum, qui sunt in terra: et réquiem dedi tibi ab ómnibus inimícis tuis.\end{Resp}
\begin{Resp}\Rsign Et fui tecum in ómnibus ubicúmque ambulásti, firmans regnum tuum in ætérnum.\end{Resp}
\begin{Resp}\Vsign Glória Patri, et Fílio, \Star{} et Spirítui Sancto.\end{Resp}
\begin{Resp}\Rsign Et fui tecum in ómnibus ubicúmque ambulásti, firmans regnum tuum in ætérnum.\end{Resp}
\switchcolumn
\EN
\begin{Resp}\Rsign Thus saith the Lord: I took thee out of thy father's house, and appointed thee to be ruler over My people, over Israel. \Star{} And I was with thee whithersoever thou wentest, to establish thy kingdom for ever.\end{Resp}
\begin{Resp}\Vsign And I have made thee a great name, like unto the name of the great men that are in the earth, and have caused thee to rest from all thine enemies.\end{Resp}
\begin{Resp}\Rsign And I was with thee whithersoever thou wentest, to establish thy kingdom for ever.\end{Resp}
\begin{Resp}\Vsign Glory be to the Father, and to the Son, \Star{} and to the Holy Ghost.\end{Resp}
\begin{Resp}\Rsign And I was with thee whithersoever thou wentest, to establish thy kingdom for ever.\end{Resp}
\switchcolumn*

\LA
\LessonHead{Lectio VII}
\switchcolumn
\EN
\LessonHead{Lesson VII}
\switchcolumn*

\LA
\LessonTitle{Léctio sancti Evangélii secúndum Matthǽum}
\switchcolumn
\EN
\LessonTitle{From the Holy Gospel according to Matthew}
\switchcolumn*

\LA
\Cite{Matt 7:15-21}
\switchcolumn
\EN
\Cite{Matt 7:15-21}
\switchcolumn*

\LA
\noindent In illo témpore: Dixit Jesus discípulis suis: Atténdite a falsis prophétis, qui véniunt ad vos in vestiméntis óvium, intrínsecus autem sunt lupi rapáces. Et réliqua.\par
\switchcolumn
\EN
\noindent At that time, Jesus said unto His disciples: Beware of false prophets, which come to you in sheep's clothing, but inwardly they are ravening wolves. And so on.\par
\switchcolumn*

\LA
\LessonTitle{Homilía sancti Hilárii Epíscopi}
\switchcolumn
\EN
\LessonTitle{Homily by St. Hilary, Bishop of Poitiers.}
\switchcolumn*

\LA
\Source{Comment. in Matth. can. 6}
\switchcolumn
\EN
\Source{Comment. on Matth. ch. vi.}
\switchcolumn*

\LA
\noindent Blandiménta verbórum et mansuetúdinis simulatiónem ádmonet fructu operatiónis expéndi oportére; ut non qualem quis verbis réferat, sed qualem se rebus effíciat, spectémus; quia in multis vestítu óvium rábies lupína contégitur. Ergo ut spinæ uvas, ut tríbuli ficus non génerant, et ut iníquæ árbores utília poma non áfferunt; ita ne in istis quidem consístere docet boni óperis efféctum, et idcírco omnes cognoscéndos esse de frúctibus. Regnum enim cælórum sola verbórum offícia non óbtinent; neque qui díxerit: Dómine, Dómine, heres illíus erit.\par
\switchcolumn
\EN
\noindent The Lord here warneth us that we must rate the worth of soft words and seeming meekness, by the fruits which they that manifest such things bring forth in their works, and that we should look, in order to see what a man is, not at his professions, but at his deeds. For there are many in whom sheep's clothing is but a mask to hide wolfish ravening. But “Do men gather grapes of thorns, or figs of thistles? Even so, every good tree bringeth forth good fruit, but a corrupt tree bringeth forth evil fruit.” Thus, the Lord teacheth us, is it with men also; evil men bring not forth good fruits, and hereby are we to know them. Lip-service alone winneth not the kingdom of heaven, nor is every one that saith unto Christ: “Lord, Lord,” an heir thereof.\par
\switchcolumn*

\LA
\begin{Resp}\Rsign Peccávi super númerum arénæ maris, et multiplicáta sunt peccáta mea: et non sum dignus vidére altitúdinem cæli præ multitúdine iniquitátis meæ: quóniam irritávi iram tuam, \Star{} Et malum coram te feci.\end{Resp}
\begin{Resp}\Vsign Quóniam iniquitátem meam ego cognósco: et delíctum meum contra me est semper, quia tibi soli peccávi.\end{Resp}
\begin{Resp}\Rsign Et malum coram te feci.\end{Resp}
\switchcolumn
\EN
\begin{Resp}\Rsign My sins are many, yea, they are more in number than the sands of the sea; I am not worthy to look up toward heaven because of the multitude of my iniquities; for I have provoked thee to anger; \Star{} And done evil in thy sight.\end{Resp}
\begin{Resp}\Vsign For I acknowledge my transgression, and my sin is ever before me, for against thee only have I sinned.\end{Resp}
\begin{Resp}\Rsign And done evil in thy sight.\end{Resp}
\switchcolumn*

\LA
\LessonHead{Lectio VIII}
\switchcolumn
\EN
\LessonHead{Lesson VIII}
\switchcolumn*

\LA
\noindent Quid enim mériti est Dómino dícere, Dómine? Numquid Dóminus non erit, nisi fúerit dictus a nobis? Et quæ offícii sánctitas est nóminis nuncupátio, cum cæléstis regni iter obediéntia pótius voluntátis Dei, non nuncupátio, repertúra sit? Multi mihi dicent in illa die: Dómine, Dómine, nonne in tuo nómine prophetávimus? Etiam nunc pseudoprophetárum frauduléntiam et hypocritárum simulaménta condémnat, qui glóriam sibi ex verbi virtúte præsúmunt, in doctrínæ prophetía, et dæmoniórum fuga, et istiúsmodi óperum virtútibus.\par
\switchcolumn
\EN
\noindent What use is there in calling the Lord, Lord? Would He not be Lord all the same, whether or not we called Him so What holiness is there in this ascription of a name, when the true way to enter into the kingdom of heaven is to do the will of our Father, Who is in heaven? “Many will say to Me in that day: Lord, Lord, have we not prophesied in thy Name?” Already here doth the Lord rebuke the deceit of the false prophets, and the feigning of the hypocrites, who take glory to themselves because of the power of their words, their prophesying in teaching, their casting out of devils, and such-like mighty works.\par
\switchcolumn*

\LA
\begin{Resp}\Rsign Duo Séraphim clamábant alter ad álterum: \Star{} Sanctus, sanctus, sanctus Dóminus Deus Sábaoth: * Plena est omnis terra glória ejus.\end{Resp}
\begin{Resp}\Vsign Tres sunt qui testimónium dant in cælo: Pater, Verbum, et Spíritus Sanctus: et hi tres unum sunt.\end{Resp}
\begin{Resp}\Rsign Sanctus, sanctus, sanctus Dóminus Deus Sábaoth.\end{Resp}
\begin{Resp}\Vsign Glória Patri, et Fílio, \Star{} et Spirítui Sancto.\end{Resp}
\begin{Resp}\Rsign Plena est omnis terra glória ejus.\end{Resp}
\switchcolumn
\EN
\begin{Resp}\Rsign One Seraph cried unto another: \Star{} Holy, Holy, Holy is the Lord God of hosts; the whole earth is full of His glory.\end{Resp}
\begin{Resp}\Vsign There are Three That bear record in heaven: the Father, the Word, and the Holy Ghost, and these Three are One.\end{Resp}
\begin{Resp}\Rsign Holy, Holy, Holy is the Lord God of hosts.\end{Resp}
\begin{Resp}\Vsign Glory be to the Father, and to the Son, \Star{} and to the Holy Ghost.\end{Resp}
\begin{Resp}\Rsign The whole earth is full of His glory.\end{Resp}
\switchcolumn*

\LA
\LessonHead{Lectio IX}
\switchcolumn
\EN
\LessonHead{Lesson IX}
\switchcolumn*

\LA
\noindent Atque hinc sibi regnum cælórum pollicéntur: quasi vero eórum áliquid próprium sit, quæ loquúntur aut fáciunt, et non ómnia virtus Dei invocáta perfíciat; cum doctrínæ sciéntiam léctio áfferat, dæmónia Christi nomen exágitet. De nostro ígitur est beáta illa ætérnitas promerénda, præstandúmque est áliquid ex próprio, ut bonum velímus, malum omne vitémus, totóque afféctu præcéptis cæléstibus obtemperémus, ac tálibus offíciis cógniti Deo simus, agamúsque pótius quod vult, quam quod potest gloriémur; repúdians eos ac repéllens, quos a cognitióne sua, ópera iniquitátis avérterint.\par
\switchcolumn
\EN
\noindent Because of all these things they promised unto themselves that they shall enter into the kingdom of heaven as though in their words and works any good thing were their own, and not all the mighty working of that God upon Whom they call, since reading bringeth knowledge of doctrine, and the Name of Christ driveth out devils. That which is needed on our part to win that blessed eternity, that of our own which we must give, is to will to do right, to turn away from all evil, to obey with our whole heart the commandments laid on us from heaven, and so to become the friends of God. It should be ours rather to do God's will, than to boast of God's power. And we must put off from us and thrust away such as are by their wicked works already estranged from His friendship.\par
\switchcolumn*

\LA
\TeDeum{Te Deum}
\switchcolumn
\EN
\TeDeum{Te Deum}
\switchcolumn*

\end{paracol}

\Office{Feria secunda infra Hebdomadam VII post Octavam Pentecostes}{Monday of the Seventh Week after Pentecost}{Feria}
\addcontentsline{toc}{subsection}{Feria secunda infra Hebdomadam VII post Octavam Pentecostes \textit{— Monday of the Seventh Week after Pentecost}}
\begin{paracol}{2}
\LA
\LessonHead{Lectio I}
\switchcolumn
\EN
\LessonHead{Lesson I}
\switchcolumn*

\LA
\LessonTitle{De libro tértio Regum}
\switchcolumn
\EN
\LessonTitle{Lesson from the first book of Kings}
\switchcolumn*

\LA
\Cite{3 Reg 1:28-31}
\switchcolumn
\EN
\Cite{1 Kgs 1:28-31}
\switchcolumn*

\LA
\noindent Et respóndit rex David dicens: Vocáte ad me Bethsabée. Quæ cum fuísset ingréssa coram rege et stetísset ante eum, jurávit rex, et ait: Vivit Dóminus, qui éruit ánimam meam de omni angústia, quia, sicut jurávi tibi per Dóminum Deum Israël dicens: Sálomon fílius tuus regnábit post me et ipse sedébit super sólium meum pro me, sic fáciam hódie. Summissóque Bethsabée in terram vultu, adorávit regem, dicens: Vivat dóminus meus David in ætérnum.\par
\switchcolumn
\EN
\noindent And King David answered and said: Call to me Bethsabee. And when she was come in to the king, and stood before him, The king swore and said: As the Lord liveth, who hath delivered my soul out of all distress, Even as I swore to thee by the Lord the God of Israel, saying: Solomon thy son shall reign after me, and he shall sit upon my throne in my stead, so will I do this day. And Bethsabee bowing with her face to the earth worshipped the king, saying: May my lord David live for ever\par
\switchcolumn*

\LA
\begin{Resp}\Rsign Recordáre, Dómine, testaménti tui, et dic Angelo percutiénti: Cesset jam manus tua, \Star{} Ut non desolétur terra, et ne perdas omnem ánimam vivam.\end{Resp}
\begin{Resp}\Vsign Ego sum qui peccávi, ego qui iníque egi: isti qui oves sunt, quid fecérunt? Avertátur, óbsecro, furor tuus, Dómine, a pópulo tuo.\end{Resp}
\begin{Resp}\Rsign Ut non desolétur terra, et ne perdas omnem ánimam vivam.\end{Resp}
\switchcolumn
\EN
\begin{Resp}\Rsign Remember, O Lord, thy covenant, and say unto the destroying Angel: Stay now thine hand; \Star{} That the land be not utterly laid waste, and that thou destroy not every living soul.\end{Resp}
\begin{Resp}\Vsign Even I it is that have sinned, and done evil indeed; but these sheep, what have they done? Let thine anger, I pray thee, O Lord, be turned away from thy people.\end{Resp}
\begin{Resp}\Rsign That the land be not utterly laid waste, and that Thou destroy not every living soul.\end{Resp}
\switchcolumn*

\LA
\LessonHead{Lectio II}
\switchcolumn
\EN
\LessonHead{Lesson II}
\switchcolumn*

\LA
\Cite{3 Reg 1:32-35}
\switchcolumn
\EN
\Cite{1 Kgs 1:32-35}
\switchcolumn*

\LA
\noindent Dixit quoque rex David: Vocáte mihi Sadoc sacerdótem et Nathan prophétam et Banájam fílium Jójadæ. Qui cum ingréssi fuíssent coram rege, dixit ad eos: tóllite vobíscum servos Dómini vestri, et impónite Salomónem fílium meum super mulam meam, et dúcite eum in Gihon, et ungat eum ibi Sadoc sacérdos et Nathan prophéta in regem super Israël. Et canétis búccina, atque dicétis: Vivat rex Sálomon. Et ascendétis post eum, et véniet, et sedébit super sólium meum, et ipse regnábit pro me.\par
\switchcolumn
\EN
\noindent King David also said: Call me Sadoc the priest, and Nathan the prophet, and Banaias the son of Joiada. And when they were come in before the king, He said to them: Take with you the servants of your lord, and set my son Solomon upon my mule: and bring him to Gihon. And let Sadoc the priest, and Nathan the prophet anoint him there king over Israel: and you shall sound the trumpet, and shall say: God save king Solomon. And you shall come up after him, and he shall come, and shall sit upon my throne, and he shall reign in my stead: and I will appoint him to be ruler over Israel, and over Juda.\par
\switchcolumn*

\LA
\begin{Resp}\Rsign Exaudísti, Dómine, oratiónem servi tui, ut ædificárem templum nómini tuo: \Star{} Bénedic et sanctífica domum istam in sempitérnum, Deus Israël.\end{Resp}
\begin{Resp}\Vsign Dómine, qui custódis pactum cum servis tuis, qui ámbulant coram te in toto corde suo.\end{Resp}
\begin{Resp}\Rsign Bénedic et sanctífica domum istam in sempitérnum, Deus Israël.\end{Resp}
\switchcolumn
\EN
\begin{Resp}\Rsign O Lord, Thou hast hearkened unto the prayer of thy servant, that I might build a temple unto thy Name, \Star{} O God of Israel, bless Thou, and hallow this house for ever.\end{Resp}
\begin{Resp}\Vsign O Lord, Who keepest covenant with thy servants that walk before thee in all their heart.\end{Resp}
\begin{Resp}\Rsign O God of Israel, bless Thou, and hallow this house for ever.\end{Resp}
\switchcolumn*

\LA
\LessonHead{Lectio III}
\switchcolumn
\EN
\LessonHead{Lesson III}
\switchcolumn*

\LA
\Cite{3 Reg 1:38-40}
\switchcolumn
\EN
\Cite{1 Kgs 1:38-40}
\switchcolumn*

\LA
\noindent Descéndit ergo Sadoc sacérdos, et Nathan prophéta, et Banájas fílius Jójadæ et Ceréthi, et Pheléthi: et imposuérunt Salomónem super mulam regis David, et adduxérunt eum in Gihon. Sumpsítque Sadoc sacérdos cornu ólei de tabernáculo, et unxit Salomónem: et cecinérunt búccina, et dixit omnis pópulus: Vivat rex Sálomon. Et ascéndit univérsa multitúdo post eum, et pópulus canéntium tíbiis, et lætántium gáudio magno: et insónuit terra a clamóre eórum.\par
\switchcolumn
\EN
\noindent So Sadoc the priest, and Nathan the prophet went down, and Banaias the son of Joiada, and the Cerethi, and Phelethi: and they set Solomon upon the mule of king David, and brought him to Gihon. And Sadoc the priest took a horn of oil out of the tabernacle, and anointed Solomon: and they sounded the trumpet, and all the people said: God save king Solomon. And all the multitude went up after him, and the people played with pipes, and rejoiced with a great joy, and the earth rang with the noise of their cry.\par
\switchcolumn*

\LA
\begin{Resp}\Rsign Audi, Dómine, hymnum et oratiónem, quam servus tuus orat coram te hódie, ut sint óculi tui apérti, et aures tuæ inténtæ, \Star{} Super domum istam die ac nocte.\end{Resp}
\begin{Resp}\Vsign Réspice, Dómine, de sanctuário tuo, et de excélso cælórum habitáculo.\end{Resp}
\begin{Resp}\Rsign Super domum istam die ac nocte.\end{Resp}
\begin{Resp}\Vsign Glória Patri, et Fílio, \Star{} et Spirítui Sancto.\end{Resp}
\begin{Resp}\Rsign Super domum istam die ac nocte.\end{Resp}
\switchcolumn
\EN
\begin{Resp}\Rsign Hearken, O Lord, unto the cry and to the prayer which thy servant prayeth before thee today, that thine eyes may be open and thine ears attend; \Star{} Toward this house day and night.\end{Resp}
\begin{Resp}\Vsign Look down from thine high and holy place, O Lord, even from heaven thy dwelling.\end{Resp}
\begin{Resp}\Rsign Toward this house, day and night.\end{Resp}
\begin{Resp}\Vsign Glory be to the Father, and to the Son, \Star{} and to the Holy Ghost.\end{Resp}
\begin{Resp}\Rsign Toward this house, day and night.\end{Resp}
\switchcolumn*

\end{paracol}

\Office{Feria tertia infra Hebdomadam VII post Octavam Pentecostes}{Tuesday of the Seventh Week after Pentecost}{Feria}
\addcontentsline{toc}{subsection}{Feria tertia infra Hebdomadam VII post Octavam Pentecostes \textit{— Tuesday of the Seventh Week after Pentecost}}
\begin{paracol}{2}
\LA
\LessonHead{Lectio I}
\switchcolumn
\EN
\LessonHead{Lesson I}
\switchcolumn*

\LA
\LessonTitle{De libro tértio Regum}
\switchcolumn
\EN
\LessonTitle{Lesson from the first book of Kings}
\switchcolumn*

\LA
\Cite{3 Reg 2:1-4}
\switchcolumn
\EN
\Cite{1 Kgs 2:1-4}
\switchcolumn*

\LA
\noindent Appropinquavérunt autem dies David ut morerétur, præcepítque Salomóni fílio suo, dicens: Ego ingrédior viam univérsæ terræ: confortáre, et esto vir, et obsérva custódias Dómini Dei tui, ut ámbules in viis ejus: ut custódias cæremónias ejus, et præcépta ejus, et judícia, et testimónia, sicut scriptum est in lege Móysi; ut intéllegas univérsa quæ facis, et quocúmque te vérteris; ut confírmet Dóminus sermónes suos, quos locútus est de me dicens: Si custodíerint fílii tui vias suas, et ambuláverint coram me in veritáte in omni corde suo et in omni ánima sua, non auferétur tibi vir de sólio Israël.\par
\switchcolumn
\EN
\noindent And the days of David drew nigh that he should die, and he charged his son Solomon, saying: I am going the way of all flesh: take thou courage, and show thyself a man. And keep the charge of the Lord thy God, to walk in his ways, and observe his ceremonies, and his precepts, and judgments, and testimonies, as it is written in the law of Moses: that thou mayest understand all thou dost, and whithersoever thou shalt turn thyself: That the Lord may confirm his words, which he hath spoken of me, saying: If thy children shall take heed to their ways, and shall walk before me in truth, with all their heart, and with all their soul, there shall not be taken away from thee a man on the throne of Israel.\par
\switchcolumn*

\LA
\begin{Resp}\Rsign Dómine, si convérsus fúerit pópulus tuus, et oráverit ad sanctuárium tuum: \Star{} Tu exáudies de cælo, Dómine, et líbera eos de mánibus inimicórum suórum.\end{Resp}
\begin{Resp}\Vsign Si peccáverit in te pópulus tuus, et convérsus égerit pœniténtiam, veniénsque oráverit in isto loco.\end{Resp}
\begin{Resp}\Rsign Tu exáudies de cælo, Dómine, et líbera eos de mánibus inimicórum suórum.\end{Resp}
\switchcolumn
\EN
\begin{Resp}\Rsign Lord, when thy people shall turn again to thee, and shall pray unto thee in this house; \Star{} Then hear Thou in heaven, O Lord, and deliver them out of the hand of their enemies.\end{Resp}
\begin{Resp}\Vsign If thy people sin against thee, and turn again, and repent, and come and pray unto thee in this house.\end{Resp}
\begin{Resp}\Rsign Then hear Thou in heaven, O Lord, and deliver them out of the hand of their enemies.\end{Resp}
\switchcolumn*

\LA
\LessonHead{Lectio II}
\switchcolumn
\EN
\LessonHead{Lesson II}
\switchcolumn*

\LA
\Cite{3 Reg 2:5-6}
\switchcolumn
\EN
\Cite{1 Kgs 2:5-6}
\switchcolumn*

\LA
\noindent Tu quoque nosti quæ fécerit mihi Joab fílius Sárviæ, quæ fécerit duóbus princípibus exércitus Israël, Abner fílio Ner, et Amasæ fílio Jether: quos occídit, et effúdit sánguinem belli in pace, et pósuit cruórem prǽlii in bálteo suo, qui erat circa lumbos ejus, et in calceaménto suo, quod erat in pédibus ejus. Fácies ergo juxta sapiéntiam tuam, et non dedúces canítiem ejus pacífice ad ínferos.\par
\switchcolumn
\EN
\noindent Thou knowest also what Joab the son of Sarvia hath done to me, what he did to the two captains of the army of Israel, to Abner the son of Ner, and to Amasa the son of Jether: whom he slew, and shed the blood of war in peace, and put the blood of war on his girdle that was about his loins, and in his shoes that were on his feet. Do therefore according to thy wisdom, and let not his hoary head go down to hell in peace.\par
\switchcolumn*

\LA
\begin{Resp}\Rsign Factum est, dum tólleret Dóminus Elíam per túrbinem in cælum, \Star{} Eliséus clamábat, dicens: Pater mi, pater mi, currus Israël, et auríga ejus.\end{Resp}
\begin{Resp}\Vsign Cumque pérgerent, et incedéntes sermocinaréntur, ecce currus ígneus et equi ígnei divisérunt utrúmque, et ascéndit Elías per túrbinem in cælum.\end{Resp}
\begin{Resp}\Rsign Eliséus clamábat, dicens: Pater mi, pater mi, currus Israël, et auríga ejus.\end{Resp}
\switchcolumn
\EN
\begin{Resp}\Rsign And it came to pass when the Lord would take up Elijah into heaven by a whirlwind \Star{} Elisha cried, and said: My father, my father, the chariot of Israel, and the horsemen thereof.\end{Resp}
\begin{Resp}\Vsign And as they still went on, and talked, behold, there appeared a chariot of fire, and horses of fire, and parted them both asunder, and Elijah went up by a whirlwind into heaven.\end{Resp}
\begin{Resp}\Rsign Elisha cried, and said: My father, my father, the chariot of Israel, and the horsemen thereof.\end{Resp}
\switchcolumn*

\LA
\LessonHead{Lectio III}
\switchcolumn
\EN
\LessonHead{Lesson III}
\switchcolumn*

\LA
\Cite{3 Reg 2:7-9}
\switchcolumn
\EN
\Cite{1 Kgs 2:7-9}
\switchcolumn*

\LA
\noindent Sed et fíliis Berzellái Galaadítis reddes grátiam, erúntque comedéntes in mensa tua; occurérunt enim mihi quando fugiébam a fácie Absalom fratris tui. Habes quoque apud te Sémei fílium Gera fílii Jémini de Bahúrim, qui maledíxit mihi maledictióne péssima, quando ibam ad castra; sed quia descéndit mihi in occúrsum, cum transírem Jordánem, et jurávi ei per Dóminum, dicens: Non te interfíciam gládio. Tu noli pati eum esse innóxium.\par
\switchcolumn
\EN
\noindent But show kindness to the sons of Berzellai the Galaadite, and let them eat at thy table: for they met me when I fled from the face of Absalom thy brother. Thou hast also with thee Semei the son of Gera the son of Jemini of Bahurim, who cursed me with a grievous curse, when I went to the camp: but because he came down to meet me when I passed over the Jordan, and I swore to him by the Lord, saying: I will not kill thee with a sword: Do not thou hold him guiltless. But thou art a wise man, and knowest what to do with him, and thou shalt bring down his grey hairs with blood to hell.\par
\switchcolumn*

\LA
\begin{Resp}\Rsign Ego te tuli de domo patris tui, dicit Dóminus, et pósui te páscere gregem pópuli mei: \Star{} Et fui tecum in ómnibus ubicúmque ambulásti, firmans regnum tuum in ætérnum.\end{Resp}
\begin{Resp}\Vsign Fecíque tibi nomen grande, juxta nomen magnórum, qui sunt in terra: et réquiem dedi tibi ab ómnibus inimícis tuis.\end{Resp}
\begin{Resp}\Rsign Et fui tecum in ómnibus ubicúmque ambulásti, firmans regnum tuum in ætérnum.\end{Resp}
\begin{Resp}\Vsign Glória Patri, et Fílio, \Star{} et Spirítui Sancto.\end{Resp}
\begin{Resp}\Rsign Et fui tecum in ómnibus ubicúmque ambulásti, firmans regnum tuum in ætérnum.\end{Resp}
\switchcolumn
\EN
\begin{Resp}\Rsign Thus saith the Lord: I took thee out of thy father's house, and appointed thee to be ruler over My people, over Israel. \Star{} And I was with thee whithersoever thou wentest, to establish thy kingdom for ever.\end{Resp}
\begin{Resp}\Vsign And I have made thee a great name, like unto the name of the great men that are in the earth, and have caused thee to rest from all thine enemies.\end{Resp}
\begin{Resp}\Rsign And I was with thee whithersoever thou wentest, to establish thy kingdom for ever.\end{Resp}
\begin{Resp}\Vsign Glory be to the Father, and to the Son, \Star{} and to the Holy Ghost.\end{Resp}
\begin{Resp}\Rsign And I was with thee whithersoever thou wentest, to establish thy kingdom for ever.\end{Resp}
\switchcolumn*

\end{paracol}

\Office{Feria quarta infra Hebdomadam VII post Octavam Pentecostes}{Wednesday of the Seventh Week after Pentecost}{Feria}
\addcontentsline{toc}{subsection}{Feria quarta infra Hebdomadam VII post Octavam Pentecostes \textit{— Wednesday of the Seventh Week after Pentecost}}
\begin{paracol}{2}
\LA
\LessonHead{Lectio I}
\switchcolumn
\EN
\LessonHead{Lesson I}
\switchcolumn*

\LA
\LessonTitle{De libro tértio Regum}
\switchcolumn
\EN
\LessonTitle{Lesson from the first book of Kings}
\switchcolumn*

\LA
\Cite{3 Reg 3:5-6}
\switchcolumn
\EN
\Cite{1 Kgs 3:5-6}
\switchcolumn*

\LA
\noindent Appáruit autem Dóminus Salomóni per sómnium nocte, dicens: Póstula quod vis ut dem tibi. Et ait Sálomon: Tu fecísti cum servo tuo David patre meo misericórdiam magnam, sicut ambulávit in conspéctu tuo in veritáte et justítia, et recto corde tecum; custodísti ei misericórdiam tuam grandem, et dedísti ei fílium sedéntem super thronum ejus, sicut est hódie.\par
\switchcolumn
\EN
\noindent And the Lord appeared to Solomon in a dream by night, saying: Ask what thou wilt that I should give thee. And Solomon said: Thou hast shown great mercy to thy servant David my father, even as he walked before thee in truth, and justice, and an upright heart with thee: and thou hast kept thy great mercy for him, and hast given him a son to sit on his throne, as it is this day.\par
\switchcolumn*

\LA
\begin{Resp}\Rsign Peccávi super númerum arénæ maris, et multiplicáta sunt peccáta mea: et non sum dignus vidére altitúdinem cæli præ multitúdine iniquitátis meæ: quóniam irritávi iram tuam, \Star{} Et malum coram te feci.\end{Resp}
\begin{Resp}\Vsign Quóniam iniquitátem meam ego cognósco: et delíctum meum contra me est semper, quia tibi soli peccávi.\end{Resp}
\begin{Resp}\Rsign Et malum coram te feci.\end{Resp}
\switchcolumn
\EN
\begin{Resp}\Rsign My sins are many, yea, they are more in number than the sands of the sea; I am not worthy to look up toward heaven because of the multitude of my iniquities; for I have provoked thee to anger; \Star{} And done evil in thy sight.\end{Resp}
\begin{Resp}\Vsign For I acknowledge my transgression, and my sin is ever before me, for against thee only have I sinned.\end{Resp}
\begin{Resp}\Rsign And done evil in thy sight.\end{Resp}
\switchcolumn*

\LA
\LessonHead{Lectio II}
\switchcolumn
\EN
\LessonHead{Lesson II}
\switchcolumn*

\LA
\Cite{3 Reg 3:7-9}
\switchcolumn
\EN
\Cite{1 Kgs 3:7-9}
\switchcolumn*

\LA
\noindent Et nunc, Dómine Deus, tu regnáre fecísti servum tuum pro David patre meo: Ego autem sum puer párvulus et ignórans egréssum et intróitum meum; et servus tuus in médio est pópuli quem elegísti, pópuli infiníti, qui numerári et supputári non potest præ multitúdine. Dabis ergo servo tuo cor dócile, ut pópulum tuum judicáre possit et discérnere inter bonum et malum. Quis enim póterit judicáre pópulum istum, pópulum tuum hunc multum?\par
\switchcolumn
\EN
\noindent And now, O Lord God, thou hast made thy servant king instead of David my father: and I am but a child, and know not how to go out and come in. And thy servant is in the midst of the people which thou hast chosen, an immense people, which cannot be numbered nor counted for multitude. Give therefore to thy servant an understanding heart, to judge thy people, and discern between good and evil. For who shall be able to judge this people, thy people which is so numerous?\par
\switchcolumn*

\LA
\begin{Resp}\Rsign Exaudísti, Dómine, oratiónem servi tui, ut ædificárem templum nómini tuo: \Star{} Bénedic et sanctífica domum istam in sempitérnum, Deus Israël.\end{Resp}
\begin{Resp}\Vsign Dómine, qui custódis pactum cum servis tuis, qui ámbulant coram te in toto corde suo.\end{Resp}
\begin{Resp}\Rsign Bénedic et sanctífica domum istam in sempitérnum, Deus Israël.\end{Resp}
\switchcolumn
\EN
\begin{Resp}\Rsign O Lord, Thou hast hearkened unto the prayer of thy servant, that I might build a temple unto thy Name, \Star{} O God of Israel, bless Thou, and hallow this house for ever.\end{Resp}
\begin{Resp}\Vsign O Lord, Who keepest covenant with thy servants that walk before thee in all their heart.\end{Resp}
\begin{Resp}\Rsign O God of Israel, bless Thou, and hallow this house for ever.\end{Resp}
\switchcolumn*

\LA
\LessonHead{Lectio III}
\switchcolumn
\EN
\LessonHead{Lesson III}
\switchcolumn*

\LA
\Cite{3 Reg 3:10-13}
\switchcolumn
\EN
\Cite{1 Kgs 3:10-13}
\switchcolumn*

\LA
\noindent Plácuit ergo sermo coram Dómino, quod Sálomon postulásset hujuscémodi rem. Et dixit Dóminus Salomóni: Quia postulásti verbum hoc, et non petísti tibi dies multos, nec divítias, aut ánimas inimicórum tuórum, sed postulásti tibi sapiéntiam ad discernéndum judícium: ecce feci tibi secúndum sermónes tuos, et dedi tibi cor sápiens et intéllegens, in tantum ut nullus ante símilis tui fúerit, nec post te surrectúrus sit. Sed et hæc, quæ non postulásti, dedi tibi, divítias scílicet et glóriam, ut nemo fúerit símilis tui in régibus cunctis retro diébus.\par
\switchcolumn
\EN
\noindent And the word was pleasing to the Lord that Solomon had asked such a thing. And the Lord said to Solomon: Because thou hast asked this thing, and hast not asked for thyself long life or riches, nor the lives of thy enemies, but hast asked for thyself wisdom to discern judgment, Behold I have done for thee according to thy words, and have given thee a wise and understanding heart, insomuch that there hath been no one like thee before thee, nor shall arise after thee. Yea and the things also which thou didst not ask, I have given thee: to wit riches and glory, as that no one hath been like thee among the kings in all days heretofore.\par
\switchcolumn*

\LA
\begin{Resp}\Rsign Audi, Dómine, hymnum et oratiónem, quam servus tuus orat coram te hódie, ut sint óculi tui apérti, et aures tuæ inténtæ, \Star{} Super domum istam die ac nocte.\end{Resp}
\begin{Resp}\Vsign Réspice, Dómine, de sanctuário tuo, et de excélso cælórum habitáculo.\end{Resp}
\begin{Resp}\Rsign Super domum istam die ac nocte.\end{Resp}
\begin{Resp}\Vsign Glória Patri, et Fílio, \Star{} et Spirítui Sancto.\end{Resp}
\begin{Resp}\Rsign Super domum istam die ac nocte.\end{Resp}
\switchcolumn
\EN
\begin{Resp}\Rsign Hearken, O Lord, unto the cry and to the prayer which thy servant prayeth before thee today, that thine eyes may be open and thine ears attend; \Star{} Toward this house day and night.\end{Resp}
\begin{Resp}\Vsign Look down from thine high and holy place, O Lord, even from heaven thy dwelling.\end{Resp}
\begin{Resp}\Rsign Toward this house, day and night.\end{Resp}
\begin{Resp}\Vsign Glory be to the Father, and to the Son, \Star{} and to the Holy Ghost.\end{Resp}
\begin{Resp}\Rsign Toward this house, day and night.\end{Resp}
\switchcolumn*

\end{paracol}

\Office{Feria quinta infra Hebdomadam VII post Octavam Pentecostes}{Thursday of the Seventh Week after Pentecost}{Feria}
\addcontentsline{toc}{subsection}{Feria quinta infra Hebdomadam VII post Octavam Pentecostes \textit{— Thursday of the Seventh Week after Pentecost}}
\begin{paracol}{2}
\LA
\LessonHead{Lectio I}
\switchcolumn
\EN
\LessonHead{Lesson I}
\switchcolumn*

\LA
\LessonTitle{De libro tértio Regum}
\switchcolumn
\EN
\LessonTitle{Lesson from the first book of Kings}
\switchcolumn*

\LA
\Cite{3 Reg 4:21-24}
\switchcolumn
\EN
\Cite{1 Kgs 4:21-24}
\switchcolumn*

\LA
\noindent Sálomon autem erat in dicióne sua, habens ómnia regna a flúmine terræ Philísthiim usque ad términum Ægýpti: offeréntium sibi múnera, et serviéntium ei cunctis diébus vitæ ejus. Erat autem cibus Salomónis per dies síngulos trigínta cori símilæ, et sexagínta cori farínæ, decem boves pingues, et vigínti boves pascuáles, et centum aríetes, excépta venatióne cervórum, capreárum, atque bubalórum, et ávium altílium. Ipse enim obtinébat omnem regiónem, quæ erat trans flumen, a Thaphsa usque ad Gazam, et cunctos reges illárum regiónum: et habébat pacem ex omni parte in circúitu.\par
\switchcolumn
\EN
\noindent And Solomon had under him all the kingdoms from the river to the land of the Philistines, even to the border of Egypt: and they brought him presents, and served him, all the days of his life. And the provision of Solomon for each day was thirty measures of fine flour, and threescore measures of meal, Ten fat oxen and twenty out of the pastures, and a hundred rams, besides venison of harts, roes, and buffles, and fatted fowls. For he had all the country which was beyond the river, from Thaphsa to Gazan, and all the kings of those countries: and he had peace on every side round about.\par
\switchcolumn*

\LA
\begin{Resp}\Rsign Præparáte corda vestra Dómino, et servíte illi soli: \Star{} Et liberábit vos de mánibus inimicórum vestrórum.\end{Resp}
\begin{Resp}\Vsign Convertímini ad eum in toto corde vestro, et auférte deos aliénos de médio vestri.\end{Resp}
\begin{Resp}\Rsign Et liberábit vos de mánibus inimicórum vestrórum.\end{Resp}
\switchcolumn
\EN
\begin{Resp}\Rsign Prepare your hearts unto the Lord, and serve Him Only \Star{} And He will deliver you out of the hand of your enemies.\end{Resp}
\begin{Resp}\Vsign Return unto Him with all your hearts, and put away the strange gods from among you.\end{Resp}
\begin{Resp}\Rsign And He will deliver you out of the hand of your enemies.\end{Resp}
\switchcolumn*

\LA
\LessonHead{Lectio II}
\switchcolumn
\EN
\LessonHead{Lesson II}
\switchcolumn*

\LA
\Cite{3 Reg 4:25-29}
\switchcolumn
\EN
\Cite{1 Kgs 4:25-29}
\switchcolumn*

\LA
\noindent Habitabátque Juda et Israël absque timóre ullo, unusquísque sub vite sua et sub ficu sua, a Dan usque Bersabée, cunctis diébus Salomónis. Et habébat Sálomon quadragínta míllia præsépia equórum currílium, et duódecim míllia equéstrium. Nutriebántque eos supradícti regis præfécti; sed et necessária mensæ regis Salomónis cum ingénti cura præbébant in témpore suo. Hórdeum quoque, et páleas equórum et jumentórum, deferébant in locum, ubi erat rex, juxta constitútum sibi. Dedit quoque Deus sapiéntiam Salomóni et prudéntiam multam nimis et latitúdinem cordis quasi arénam quæ est in líttore maris.\par
\switchcolumn
\EN
\noindent And Juda and Israel dwelt without any fear, every one under his vine, and under his fig tree, from Dan to Bersabee, all the days of Solomon. And Solomon had forty thousand stalls of chariot horses, and twelve thousand for the saddle. And the foresaid governors of the king fed them: and they furnished the necessaries also for king Solomon's table, with great care in their time. They brought barley also and straw for the horses, and beasts, to the place where the king was, according as it was appointed them. And God gave to Solomon wisdom and understanding exceeding much, and largeness of heart as the sand that is on the sea shore.\par
\switchcolumn*

\LA
\begin{Resp}\Rsign Deus ómnium exaudítor est: ipse misit Angelum suum, et tulit me de óvibus patris mei: \Star{} Et unxit me unctióne misericórdiæ suæ.\end{Resp}
\begin{Resp}\Vsign Dóminus, qui erípuit me de ore leónis, et de manu béstiæ liberávit me.\end{Resp}
\begin{Resp}\Rsign Et unxit me unctióne misericórdiæ suæ.\end{Resp}
\switchcolumn
\EN
\begin{Resp}\Rsign God, Which heareth all, even He sent His Angel, and took me from keeping my father's sheep, and \Star{} Anointed me with the oil of His mercy.\end{Resp}
\begin{Resp}\Vsign The Lord That delivered me out of the mouth of the lion, and out of the paw of the bear\end{Resp}
\begin{Resp}\Rsign And anointed me with the oil of His mercy.\end{Resp}
\switchcolumn*

\LA
\LessonHead{Lectio III}
\switchcolumn
\EN
\LessonHead{Lesson III}
\switchcolumn*

\LA
\Cite{3 Reg 4:30-34}
\switchcolumn
\EN
\Cite{1 Kgs 4:30-34}
\switchcolumn*

\LA
\noindent Et præcedébat sapiéntia Salomónis sapiéntiam ómnium Orientálium et Ægyptiórum; et erat sapiéntior cunctis homínibus: sapiéntior Ethan Ezrahíta, et Heman, et Chalcol, et Dorda fíliis Mahol: et erat nominátus in univérsis géntibus per circúitum. Locútus est quoque Sálomon tria míllia parábolas: et fuérunt cármina ejus quinque et mille. Et disputávit super lignis a cedro, quæ est in Líbano, usque ad hyssópum quæ egrediétur de paríete; et disséruit de juméntis, et volúcribus, et reptílibus, et píscibus. Et veniébant de cunctis pópulis ad audiéndam sapiéntiam Salomónis, et ab univérsis régibus terræ, qui audiébant sapiéntiam ejus.\par
\switchcolumn
\EN
\noindent And the wisdom of Solomon surpassed the wisdom of all the Orientals, and of the Egyptians, And he was wiser than all men: wiser than Ethan the Ezrahite, and Heman, and Chalcol, and Dorda the sons of Mahol, and he was renowned in all nations round about. Solomon also spoke three thousand parables: and his poems were a thousand and five. And he treated about trees from the cedar that is in Libanus, unto the hyssop that cometh out of the wall: and he discoursed of beasts, and of fowls, and of creeping things, and of fishes. And they came from all nations to hear the wisdom of Solomon, and from all the kings of the earth, who heard of his wisdom.\par
\switchcolumn*

\LA
\begin{Resp}\Rsign Dóminus, qui erípuit me de ore leónis, et de manu béstiæ liberávit me, \Star{} Ipse me erípiet de mánibus inimicórum meórum.\end{Resp}
\begin{Resp}\Vsign Misit Deus misericórdiam suam et veritátem suam: ánimam meam erípuit de médio catulórum leónum.\end{Resp}
\begin{Resp}\Rsign Ipse me erípiet de mánibus inimicórum meórum.\end{Resp}
\begin{Resp}\Vsign Glória Patri, et Fílio, \Star{} et Spirítui Sancto.\end{Resp}
\begin{Resp}\Rsign Ipse me erípiet de mánibus inimicórum meórum.\end{Resp}
\switchcolumn
\EN
\begin{Resp}\Rsign The Lord That delivered me out of the mouth of the lion, and out of the paw of the bear \Star{} He will deliver me out of the hand of mine enemies.\end{Resp}
\begin{Resp}\Vsign God hath sent forth His mercy and His truth, and delivered my soul from among the lion's whelps.\end{Resp}
\begin{Resp}\Rsign He will deliver me out of the hand of mine enemies.\end{Resp}
\begin{Resp}\Vsign Glory be to the Father, and to the Son, \Star{} and to the Holy Ghost.\end{Resp}
\begin{Resp}\Rsign He will deliver me out of the hand of mine enemies.\end{Resp}
\switchcolumn*

\end{paracol}

\Office{Feria sexta infra Hebdomadam VII post Octavam Pentecostes}{Friday of the Seventh Week after Pentecost}{Feria}
\addcontentsline{toc}{subsection}{Feria sexta infra Hebdomadam VII post Octavam Pentecostes \textit{— Friday of the Seventh Week after Pentecost}}
\begin{paracol}{2}
\LA
\LessonHead{Lectio I}
\switchcolumn
\EN
\LessonHead{Lesson I}
\switchcolumn*

\LA
\LessonTitle{De libro tértio Regum}
\switchcolumn
\EN
\LessonTitle{Lesson from the first book of Kings}
\switchcolumn*

\LA
\Cite{3 Reg 5:1-4}
\switchcolumn
\EN
\Cite{1 Kgs 5:1-4}
\switchcolumn*

\LA
\noindent Misit quoque Hiram, rex Tyri, servos suos ad Salomónem; audívit enim quod ipsum unxíssent regem pro patre ejus; quia amícus fúerat Hiram David omni témpore. Misit autem Sálomon ad Hiram, dicens: Tu scis voluntátem David patris mei, et quia non potúerit ædificáre domum nómini Dómini Dei sui propter bella imminéntia per circúitum, donec daret Dóminus eos sub vestígio pedum ejus. Nunc autem réquiem dedit Dóminus Deus meus mihi per circúitum, et non est satan neque occúrsus malus.\par
\switchcolumn
\EN
\noindent And Hiram king of Tyre sent his servants to Solomon: for he heard that they had anointed him king in the room of his father: for Hiram had always been David's friend. And Solomon sent to Hiram, saying: Thou knowest the will of David my father, and that he could not build a house to the name of the Lord his God, because of the wars that were round about him, until the Lord put them under the soles of his feet. But now the Lord my God hath given me rest round about: and there is no adversary nor evil occurrence.\par
\switchcolumn*

\LA
\begin{Resp}\Rsign Percússit Saul mille, et David decem míllia: \Star{} Quia manus Dómini erat cum illo: percússit Philisthǽum, et ábstulit oppróbrium ex Israël.\end{Resp}
\begin{Resp}\Vsign Nonne iste est David, de quo canébant in choro, dicéntes: Saul percússit mille, et David decem míllia?\end{Resp}
\begin{Resp}\Rsign Quia manus Dómini erat cum illo: percússit Philisthǽum, et ábstulit oppróbrium ex Israël.\end{Resp}
\switchcolumn
\EN
\begin{Resp}\Rsign Saul hath slain his thousands, and David his ten thousands. \Star{} Because the hand of the Lord was with him, he smote the Philistine, and took away the reproach from Israel.\end{Resp}
\begin{Resp}\Vsign Is not this David? Did they not sing one to another of him in dances, saying: Saul hath slain his thousands, and David his ten thousands?\end{Resp}
\begin{Resp}\Rsign Because the hand of the Lord was with him, he smote the Philistine, and took away the reproach from Israel.\end{Resp}
\switchcolumn*

\LA
\LessonHead{Lectio II}
\switchcolumn
\EN
\LessonHead{Lesson II}
\switchcolumn*

\LA
\Cite{3 Reg 5:5-6}
\switchcolumn
\EN
\Cite{1 Kgs 5:5-6}
\switchcolumn*

\LA
\noindent Quam ob rem cógito ædificáre templum nómini Dómini Dei mei, sicut locútus est Dóminus David, patri meo, dicens: Fílius tuus, quem dabo pro te super sólium tuum, ipse ædificábit domum nómini meo. Prǽcipe ígitur ut præcídant mihi servi tui cedros de Líbano, et servi mei sint cum servis tuis; mercédem autem servórum tuórum dabo tibi quamcúmque petíeris; scis enim quómodo non est in pópulo meo vir qui nóverit ligna cǽdere sicut Sidónii.\par
\switchcolumn
\EN
\noindent Wherefore I purpose to build a temple to the name of the Lord my God, as the Lord spoke to David my father, saying: my son, whom I will set upon the throne in thy place, he shall build a house to my name. Give orders therefore that thy servants cut me down cedar trees out of Libanus, and let my servants be with thy servants: and I will give thee the hire of thy servants whatsoever thou wilt ask, for thou knowest how there is not among my people a man that has skill to hew wood like to the Sidonians.\par
\switchcolumn*

\LA
\begin{Resp}\Rsign Montes Gélboë, nec ros nec plúvia véniant super vos, \Star{} Ubi cecidérunt fortes Israël.\end{Resp}
\begin{Resp}\Vsign Omnes montes, qui estis in circúitu ejus, vísitet Dóminus: a Gélboë autem tránseat.\end{Resp}
\begin{Resp}\Rsign Ubi cecidérunt fortes Israël.\end{Resp}
\switchcolumn
\EN
\begin{Resp}\Rsign Ye mountains of Gilboa, let there be no dew, neither let there be rain upon you \Star{} For there are the mighty of Israel fallen.\end{Resp}
\begin{Resp}\Vsign All ye mountains that stand round about, the Lord look upon you but let Him pass by Gilboa\end{Resp}
\begin{Resp}\Rsign For there are the mighty of Israel fallen.\end{Resp}
\switchcolumn*

\LA
\LessonHead{Lectio III}
\switchcolumn
\EN
\LessonHead{Lesson III}
\switchcolumn*

\LA
\Cite{3 Reg 5:7-9}
\switchcolumn
\EN
\Cite{1 Kgs 5:7-9}
\switchcolumn*

\LA
\noindent Cum ergo audísset Hiram verba Salomónis, lætátus est valde, et ait: Benedíctus Dóminus Deus hódie, qui dedit David fílium sapientíssimum super pópulum hunc plúrimum. Et misit Hiram ad Salomónem, dicens: Audívi quæcúmque mandásti mihi: ego fáciam omnem voluntátem tuam in lignis cédrinis et abiégnis. Servi mei depónent ea de Líbano ad mare, et ego compónam ea in rátibus in mari usque ad locum quem significáveris mihi; et applicábo ea ibi, et tu tolles ea præbebísque necessária mihi ut detur cibus dómui meæ.\par
\switchcolumn
\EN
\noindent Now when Hiram had heard the words of Solomon, he rejoiced exceedingly, and said: Blessed be the Lord God this day, who hath given to David a very wise son over this numerous people. And Hiram sent to Solomon, saying: I have heard all thou hast desired of me: and I will do all thy desire concerning cedar trees, and fir trees. My servants shall bring them down from Libanus to the sea: and I will put them together in floats in the sea, and convey them to the place, which thou shalt signify to me; and will land them there, and thou shalt receive them: and thou shalt allow me necessaries, to furnish food for my household.\par
\switchcolumn*

\LA
\begin{Resp}\Rsign Ego te tuli de domo patris tui, dicit Dóminus, et pósui te páscere gregem pópuli mei: \Star{} Et fui tecum in ómnibus ubicúmque ambulásti, firmans regnum tuum in ætérnum.\end{Resp}
\begin{Resp}\Vsign Fecíque tibi nomen grande, juxta nomen magnórum, qui sunt in terra: et réquiem dedi tibi ab ómnibus inimícis tuis.\end{Resp}
\begin{Resp}\Rsign Et fui tecum in ómnibus ubicúmque ambulásti, firmans regnum tuum in ætérnum.\end{Resp}
\begin{Resp}\Vsign Glória Patri, et Fílio, \Star{} et Spirítui Sancto.\end{Resp}
\begin{Resp}\Rsign Et fui tecum in ómnibus ubicúmque ambulásti, firmans regnum tuum in ætérnum.\end{Resp}
\switchcolumn
\EN
\begin{Resp}\Rsign Thus saith the Lord: I took thee out of thy father's house, and appointed thee to be ruler over My people, over Israel. \Star{} And I was with thee whithersoever thou wentest, to establish thy kingdom for ever.\end{Resp}
\begin{Resp}\Vsign And I have made thee a great name, like unto the name of the great men that are in the earth, and have caused thee to rest from all thine enemies.\end{Resp}
\begin{Resp}\Rsign And I was with thee whithersoever thou wentest, to establish thy kingdom for ever.\end{Resp}
\begin{Resp}\Vsign Glory be to the Father, and to the Son, \Star{} and to the Holy Ghost.\end{Resp}
\begin{Resp}\Rsign And I was with thee whithersoever thou wentest, to establish thy kingdom for ever.\end{Resp}
\switchcolumn*

\end{paracol}

\Office{Sabbato infra Hebdomadam VII post Octavam Pentecostes}{Saturday of the Seventh Week after Pentecost}{Feria}
\addcontentsline{toc}{subsection}{Sabbato infra Hebdomadam VII post Octavam Pentecostes \textit{— Saturday of the Seventh Week after Pentecost}}
\begin{paracol}{2}
\LA
\LessonHead{Lectio I}
\switchcolumn
\EN
\LessonHead{Lesson I}
\switchcolumn*

\LA
\LessonTitle{De libro tértio Regum}
\switchcolumn
\EN
\LessonTitle{Lesson from the first book of Kings}
\switchcolumn*

\LA
\Cite{3 Reg 7:51; 8:1-2}
\switchcolumn
\EN
\Cite{1 Kgs 7:51; 8:1-2}
\switchcolumn*

\LA
\noindent Et perfécit omne opus, quod faciébat Sálomon in domo Dómini, et íntulit quæ sanctificáverat David pater suus, argéntum, et aurum, et vasa, reposuítque in thesáuris domus Dómini. Tunc congregáti sunt omnes majóres natu Israël cum princípibus tríbuum et duces familiárum filiórum Israël ad regem Salomónem in Jerúsalem, ut deférrent arcam fœ́deris Dómini de civitáte David, id est de Sion. Convenítque ad regem Salomónem univérsus Israël in mense Ethánim in solémni die (ipse est mensis séptimus).\par
\switchcolumn
\EN
\noindent And Solomon finished all the work that he made in the house of the Lord, and brought in the things that David his father had dedicated, the silver and the gold, and the vessels, and laid them up in the treasures of the house of the Lord. Then all the ancients of Israel with the princes of the tribes, and the heads of the families of the children of Israel were assembled to king Solomon in Jerusalem: that they might carry the ark of the covenant of the Lord out of the city of David, that is, out of Sion. And all Israel assembled themselves to king Solomon on the festival day in the month of Ethanim, the same is the seventh month.\par
\switchcolumn*

\LA
\begin{Resp}\Rsign Peccávi super númerum arénæ maris, et multiplicáta sunt peccáta mea: et non sum dignus vidére altitúdinem cæli præ multitúdine iniquitátis meæ: quóniam irritávi iram tuam, \Star{} Et malum coram te feci.\end{Resp}
\begin{Resp}\Vsign Quóniam iniquitátem meam ego cognósco: et delíctum meum contra me est semper, quia tibi soli peccávi.\end{Resp}
\begin{Resp}\Rsign Et malum coram te feci.\end{Resp}
\switchcolumn
\EN
\begin{Resp}\Rsign My sins are many, yea, they are more in number than the sands of the sea; I am not worthy to look up toward heaven because of the multitude of my iniquities; for I have provoked thee to anger; \Star{} And done evil in thy sight.\end{Resp}
\begin{Resp}\Vsign For I acknowledge my transgression, and my sin is ever before me, for against thee only have I sinned.\end{Resp}
\begin{Resp}\Rsign And done evil in thy sight.\end{Resp}
\switchcolumn*

\LA
\LessonHead{Lectio II}
\switchcolumn
\EN
\LessonHead{Lesson II}
\switchcolumn*

\LA
\Cite{3 Reg 8:3-7}
\switchcolumn
\EN
\Cite{1 Kgs 8:3-7}
\switchcolumn*

\LA
\noindent Venerúntque cuncti senes de Israël, et tulérunt arcam sacerdótes, et portavérunt arcam Dómini, et tabernáculum fœ́deris et ómnia vasa sanctuárii, quæ erant in tabernáculo, et ferébant ea sacerdótes et levítæ. Rex autem Sálomon et omnis multitúdo Israël, quæ convénerat ad eum, gradiebátur cum illo ante arcam, et immolábant oves et boves absque æstimatióne et número. Et intulérunt sacerdótes arcam fœ́deris Dómini in locum suum, in oráculum templi, in Sanctum sanctórum, subter alas Chérubim; síquidem Chérubim expandébant alas super locum arcæ, et protegébant arcam, et vectes ejus désuper.\par
\switchcolumn
\EN
\noindent And all the ancients of Israel came, and the priests took up the ark, And carried the ark of the Lord, and the tabernacle of the covenant, and all the vessels of the sanctuary, that were in the tabernacle: and the priests and the Levites carried them. And king Solomon, and all the multitude of Israel, that were assembled unto him went with him before the ark, and they sacrificed sheep and oxen that could not be counted or numbered. And the priests brought in the ark of the covenant of the Lord into its place, into the oracle of the temple, into the holy of holies under the wings of the cherubims. For the cherubims spread forth their wings over the place of the ark, and covered the ark, and the staves thereof above.\par
\switchcolumn*

\LA
\begin{Resp}\Rsign Exaudísti, Dómine, oratiónem servi tui, ut ædificárem templum nómini tuo: \Star{} Bénedic et sanctífica domum istam in sempitérnum, Deus Israël.\end{Resp}
\begin{Resp}\Vsign Dómine, qui custódis pactum cum servis tuis, qui ámbulant coram te in toto corde suo.\end{Resp}
\begin{Resp}\Rsign Bénedic et sanctífica domum istam in sempitérnum, Deus Israël.\end{Resp}
\switchcolumn
\EN
\begin{Resp}\Rsign O Lord, Thou hast hearkened unto the prayer of thy servant, that I might build a temple unto thy Name, \Star{} O God of Israel, bless Thou, and hallow this house for ever.\end{Resp}
\begin{Resp}\Vsign O Lord, Who keepest covenant with thy servants that walk before thee in all their heart.\end{Resp}
\begin{Resp}\Rsign O God of Israel, bless Thou, and hallow this house for ever.\end{Resp}
\switchcolumn*

\LA
\LessonHead{Lectio III}
\switchcolumn
\EN
\LessonHead{Lesson III}
\switchcolumn*

\LA
\Cite{3 Reg 8:9-12}
\switchcolumn
\EN
\Cite{1 Kgs 8:9-12}
\switchcolumn*

\LA
\noindent In arca autem non erat áliud nisi duæ tábulæ lapídeæ, quas posúerat in ea Móyses in Horeb, quando pépigit Dóminus fœdus cum fíliis Israël, cum egrederéntur de terra Ægýpti. Factum est autem, cum exíssent sacerdótes de sanctuário, nébula implévit domum Dómini, et non póterant sacerdótes stare et ministráre propter nébulam; impléverat enim glória Dómini domum Dómini. Tunc ait Sálomon: Dóminus dixit ut habitáret in nébula.\par
\switchcolumn
\EN
\noindent Now in the ark there was nothing else but the two tables of stone, which Moses put there at Horeb, when the Lord made a covenant with the children of Israel, when they came out of the land of Egypt. And it came to pass, when the priests were come out of the sanctuary, that a cloud filled the house of the Lord, And the priests could not stand to minister because of the cloud: for the glory of the Lord had filled the house of the Lord. Then Solomon said: The Lord said that he would dwell in a cloud.\par
\switchcolumn*

\LA
\begin{Resp}\Rsign Audi, Dómine, hymnum et oratiónem, quam servus tuus orat coram te hódie, ut sint óculi tui apérti, et aures tuæ inténtæ, \Star{} Super domum istam die ac nocte.\end{Resp}
\begin{Resp}\Vsign Réspice, Dómine, de sanctuário tuo, et de excélso cælórum habitáculo.\end{Resp}
\begin{Resp}\Rsign Super domum istam die ac nocte.\end{Resp}
\begin{Resp}\Vsign Glória Patri, et Fílio, \Star{} et Spirítui Sancto.\end{Resp}
\begin{Resp}\Rsign Super domum istam die ac nocte.\end{Resp}
\switchcolumn
\EN
\begin{Resp}\Rsign Hearken, O Lord, unto the cry and to the prayer which thy servant prayeth before thee today, that thine eyes may be open and thine ears attend; \Star{} Toward this house day and night.\end{Resp}
\begin{Resp}\Vsign Look down from thine high and holy place, O Lord, even from heaven thy dwelling.\end{Resp}
\begin{Resp}\Rsign Toward this house, day and night.\end{Resp}
\begin{Resp}\Vsign Glory be to the Father, and to the Son, \Star{} and to the Holy Ghost.\end{Resp}
\begin{Resp}\Rsign Toward this house, day and night.\end{Resp}
\switchcolumn*

\end{paracol}

\Office{Dominica VIII Post Pentecosten}{Eighth Sunday after Pentecost}{Semiduplex}
\addcontentsline{toc}{subsection}{Dominica VIII Post Pentecosten \textit{— Eighth Sunday after Pentecost}}
\begin{paracol}{2}
\LA
\LessonHead{Lectio I}
\switchcolumn
\EN
\LessonHead{Lesson I}
\switchcolumn*

\LA
\LessonTitle{De libro tértio Regum}
\switchcolumn
\EN
\LessonTitle{Lesson from the first book of Kings}
\switchcolumn*

\LA
\Cite{3 Reg 9:1-5}
\switchcolumn
\EN
\Cite{1 Kgs 9:1-5}
\switchcolumn*

\LA
\noindent Factum est autem cum perfecísset Sálomon ædifícium domus Dómini et ædifícium regis et omne quod optáverat et volúerat fácere, appáruit ei Dóminus secúndo, sicut apparúerat ei in Gábaon. Dixítque Dóminus ad eum: Exaudívi oratiónem tuam et deprecatiónem tuam, quam deprecátus es coram me; sanctificávi domum hanc, quam ædificásti, ut pónerem nomen meum ibi in sempitérnum; et erunt óculi mei et cor meum ibi cunctis diébus. Tu quoque, si ambuláveris coram me sicut ambulávit pater tuus in simplicitáte cordis et in æquitáte, et féceris ómnia quæ præcépi tibi et legítima mea et judícia mea serváveris, ponam thronum regni tui super Israël in sempitérnum, sicut locútus sum David, patri tuo dicens: Non auferétur vir de génere tuo de sólio Israël.\par
\switchcolumn
\EN
\noindent And it came to pass when Solomon had finished the building of the house of the Lord, and the king's house, and all that he desired, and was pleased to do, That the Lord appeared to him the second time, as he had appeared to him in Gabaon. And the Lord said to him: I have heard thy prayer and thy supplication, which thou hast made before me: I have sanctified this house, which thou hast built, to put my name there for ever, and my eyes and my heart shall be there always. And if thou wilt walk before me, as thy father walked, in simplicity of heart, and in uprightness: and wilt do all that I have commanded thee, and wilt keep my ordinances and my judgments, I will establish the throne of thy kingdom over Israel for ever, as I promised David thy father, saying: There shall not fail a man of thy race upon the throne of Israel.\par
\switchcolumn*

\LA
\begin{Resp}\Rsign Præparáte corda vestra Dómino, et servíte illi soli: \Star{} Et liberábit vos de mánibus inimicórum vestrórum.\end{Resp}
\begin{Resp}\Vsign Convertímini ad eum in toto corde vestro, et auférte deos aliénos de médio vestri.\end{Resp}
\begin{Resp}\Rsign Et liberábit vos de mánibus inimicórum vestrórum.\end{Resp}
\switchcolumn
\EN
\begin{Resp}\Rsign Prepare your hearts unto the Lord, and serve Him Only \Star{} And He will deliver you out of the hand of your enemies.\end{Resp}
\begin{Resp}\Vsign Return unto Him with all your hearts, and put away the strange gods from among you.\end{Resp}
\begin{Resp}\Rsign And He will deliver you out of the hand of your enemies.\end{Resp}
\switchcolumn*

\LA
\LessonHead{Lectio II}
\switchcolumn
\EN
\LessonHead{Lesson II}
\switchcolumn*

\LA
\Cite{3 Reg 9:6-9}
\switchcolumn
\EN
\Cite{1 Kgs 9:6-9}
\switchcolumn*

\LA
\noindent Si autem aversióne avérsi fuéritis vos et fílii vestri non sequéntes me nec custodiéntes mandáta mea et cæremónias meas quas propósui vobis, sed abiéritis et coluéritis deos aliénos et adoravéritis eos; áuferam Israël de superfície terræ quam dedi eis, et templum quod sanctificávi nómini meo proíciam a conspéctu meo, erítque Israël in provérbium et in fábulam cunctis pópulis, et domus hæc erit in exémplum: omnis qui transíerit per eam stupébit et sibilábit et dicet: Quare fecit Dóminus sic terræ huic et dómui huic? Et respondébunt: Quia dereliquérunt Dóminum Deum suum, qui edúxit patres eórum de terra Ægýpti, et secúti sunt deos aliénos et adoravérunt eos et coluérunt eos; idcírco indúxit Dóminus super eos omne malum hoc.\par
\switchcolumn
\EN
\noindent But if you and your children revolting shall turn away from following me, and will not keep my commandments, and my ceremonies, which I have set before you, but will go and worship strange gods, and adore them: I will take away Israel from the face of the land which I have given them; and the temple which I have sanctified to my name, I will cast out of my sight; and Israel shall be a proverb, and a byword among all people. And this house shall be made an example of: every one that shall pass by it, shall be astonished, and shall hiss, and say: Why hath the Lord done thus to this land, and to this house: And they shall answer: Because they forsook the Lord their God, who brought their fathers out of the land of Egypt, and followed strange gods, and adored them, and worshipped them: therefore hath the Lord brought upon them all this evil.\par
\switchcolumn*

\LA
\begin{Resp}\Rsign Deus ómnium exaudítor est: ipse misit Angelum suum, et tulit me de óvibus patris mei: \Star{} Et unxit me unctióne misericórdiæ suæ.\end{Resp}
\begin{Resp}\Vsign Dóminus, qui erípuit me de ore leónis, et de manu béstiæ liberávit me.\end{Resp}
\begin{Resp}\Rsign Et unxit me unctióne misericórdiæ suæ.\end{Resp}
\switchcolumn
\EN
\begin{Resp}\Rsign God, Which heareth all, even He sent His Angel, and took me from keeping my father's sheep, and \Star{} Anointed me with the oil of His mercy.\end{Resp}
\begin{Resp}\Vsign The Lord That delivered me out of the mouth of the lion, and out of the paw of the bear\end{Resp}
\begin{Resp}\Rsign And anointed me with the oil of His mercy.\end{Resp}
\switchcolumn*

\LA
\LessonHead{Lectio III}
\switchcolumn
\EN
\LessonHead{Lesson III}
\switchcolumn*

\LA
\Cite{3 Reg 9:10-14}
\switchcolumn
\EN
\Cite{1 Kgs 9:10-14}
\switchcolumn*

\LA
\noindent Explétis autem annis vigínti postquam ædificáverat Sálomon duas domos, id est domus Dómini et domum regis (Hiram rege Tyri præbénte Salomóni ligna cédrina et abiégna et aurum juxta omne quod opus habúerat), tunc dedit Sálomon Hiram vigínti óppida in terra Galilǽæ. Et egréssus est Hiram de Tyro ut vidéret óppida, quæ déderat ei Sálomon, et non placuérunt ei, et ait: Hǽcine sunt civitátes, quas dedísti mihi, frater? Et appellávit eas terram Chabul usque in diem hanc. Misit quoque Hiram ad regem Salomónem centum vigínti talénta auri.\par
\switchcolumn
\EN
\noindent And when twenty years were ended after Solomon had built the two houses, that is, the house of the Lord, and the house of the king, (Hiram the king of Tyre furnishing Solomon with cedar trees and fir trees, and gold according to all he had need of.) then Solomon gave Hiram twenty cities in the land of Galilee. And Hiram came out of Tyre, to see the towns which Solomon had given him, and they pleased him not, And he said: Are these the cities which thou hast given me, brother? And he called them the land of Chabul, unto this day. And Hiram sent to king Solomon a hundred and twenty talents of gold.\par
\switchcolumn*

\LA
\begin{Resp}\Rsign Dóminus, qui erípuit me de ore leónis, et de manu béstiæ liberávit me, \Star{} Ipse me erípiet de mánibus inimicórum meórum.\end{Resp}
\begin{Resp}\Vsign Misit Deus misericórdiam suam et veritátem suam: ánimam meam erípuit de médio catulórum leónum.\end{Resp}
\begin{Resp}\Rsign Ipse me erípiet de mánibus inimicórum meórum.\end{Resp}
\begin{Resp}\Vsign Glória Patri, et Fílio, \Star{} et Spirítui Sancto.\end{Resp}
\begin{Resp}\Rsign Ipse me erípiet de mánibus inimicórum meórum.\end{Resp}
\switchcolumn
\EN
\begin{Resp}\Rsign The Lord That delivered me out of the mouth of the lion, and out of the paw of the bear \Star{} He will deliver me out of the hand of mine enemies.\end{Resp}
\begin{Resp}\Vsign God hath sent forth His mercy and His truth, and delivered my soul from among the lion's whelps.\end{Resp}
\begin{Resp}\Rsign He will deliver me out of the hand of mine enemies.\end{Resp}
\begin{Resp}\Vsign Glory be to the Father, and to the Son, \Star{} and to the Holy Ghost.\end{Resp}
\begin{Resp}\Rsign He will deliver me out of the hand of mine enemies.\end{Resp}
\switchcolumn*

\LA
\LessonHead{Lectio IV}
\switchcolumn
\EN
\LessonHead{Lesson IV}
\switchcolumn*

\LA
\LessonTitle{Ex libro sancti Augustíni Epíscopi de Civitáte Dei.}
\switchcolumn
\EN
\LessonTitle{From the Book upon “The City of God” written by St. Augustine, Bishop of Hippo}
\switchcolumn*

\LA
\Source{Liber 17. cap. 8. sub medio.}
\switchcolumn
\EN
\Source{Bk. xvii, ch. 8}
\switchcolumn*

\LA
\noindent Facta est quidem nonnúlla imágo rei futúræ étiam in Salomóne, in eo quod templum ædificávit, et pacem hábuit secúndum nomen suum (Sálomon quippe pacíficus est Latíne), et in exórdio regni sui mirabíliter laudábilis fuit. Sed eádem sua persóna per umbram futúri prænuntiábat étiam ipse Christum Dóminum nostrum, non exhibébat. Unde quædam de illo ita scripta sunt, quasi de ipso ista prædícta sint, dum Scriptúra sancta étiam rebus gestis prophétans, quodámmodo in eo figúram delíneat futurórum.\par
\switchcolumn
\EN
\noindent Things which were then still to come were in a certain manner imagined in Solomon, who built the temple, who had that peace which his name implieth (for the name Solomon signifieth “the Peaceful One”) and who, at the beginning of his reign, was marvellously praiseworthy. By these things he foreshadowed in his own person, though he set not forth with his mouth, our Lord Christ. Thus are some things written of Solomon which are, as it were, things written concerning Christ, the Holy Scripture in this way, by giving the history of things past, prophesying all the while of things to come.\par
\switchcolumn*

\LA
\begin{Resp}\Rsign Percússit Saul mille, et David decem míllia: \Star{} Quia manus Dómini erat cum illo: percússit Philisthǽum, et ábstulit oppróbrium ex Israël.\end{Resp}
\begin{Resp}\Vsign Nonne iste est David, de quo canébant in choro, dicéntes: Saul percússit mille, et David decem míllia?\end{Resp}
\begin{Resp}\Rsign Quia manus Dómini erat cum illo: percússit Philisthǽum, et ábstulit oppróbrium ex Israël.\end{Resp}
\switchcolumn
\EN
\begin{Resp}\Rsign Saul hath slain his thousands, and David his ten thousands. \Star{} Because the hand of the Lord was with him, he smote the Philistine, and took away the reproach from Israel.\end{Resp}
\begin{Resp}\Vsign Is not this David? Did they not sing one to another of him in dances, saying: Saul hath slain his thousands, and David his ten thousands?\end{Resp}
\begin{Resp}\Rsign Because the hand of the Lord was with him, he smote the Philistine, and took away the reproach from Israel.\end{Resp}
\switchcolumn*

\LA
\LessonHead{Lectio V}
\switchcolumn
\EN
\LessonHead{Lesson V}
\switchcolumn*

\LA
\noindent Nam præter libros divínæ históriæ, ubi regnásse narrátur, Psalmus étiam septuagésimus primus título nóminis ejus inscríptus est. In quo tam multa dicúntur, quæ omníno ei conveníre non possunt, Dómino autem Christo aptíssima perspicuitáte convéniunt: ut evidénter appáreat, quod in illo figúra qualiscúmque adumbráta sit, in isto autem ipsa véritas præsentáta.\par
\switchcolumn
\EN
\noindent Nor besides the books of the Divine history wherein his reign is recorded, the 71st Psalm is superscribed with his name. In this Psalm are many things which cannot suit him, but are most clearly applicable to the Lord Christ, thus showing that Solomon was, as it were, a shadowy figure cast before, of that which was afterwards revealed in very truth in the Person of Christ.\par
\switchcolumn*

\LA
\begin{Resp}\Rsign Montes Gélboë, nec ros nec plúvia véniant super vos, \Star{} Ubi cecidérunt fortes Israël.\end{Resp}
\begin{Resp}\Vsign Omnes montes, qui estis in circúitu ejus, vísitet Dóminus: a Gélboë autem tránseat.\end{Resp}
\begin{Resp}\Rsign Ubi cecidérunt fortes Israël.\end{Resp}
\switchcolumn
\EN
\begin{Resp}\Rsign Ye mountains of Gilboa, let there be no dew, neither let there be rain upon you \Star{} For there are the mighty of Israel fallen.\end{Resp}
\begin{Resp}\Vsign All ye mountains that stand round about, the Lord look upon you but let Him pass by Gilboa\end{Resp}
\begin{Resp}\Rsign For there are the mighty of Israel fallen.\end{Resp}
\switchcolumn*

\LA
\LessonHead{Lectio VI}
\switchcolumn
\EN
\LessonHead{Lesson VI}
\switchcolumn*

\LA
\noindent Notum est enim quibus términis regnum conclúsum fúerit Salomónis: et tamen in eo Psalmo légitur, ut ália táceam: Dominábitur a mari usque ad mare, et a flúmine usque ad términos orbis terræ; quod in Christo vidémus impléri. A flúmine quippe dominándi sumpsit exórdium, ubi baptizátus a Joánne, eódem monstránte, cœpit agnósci a discípulis, qui eum non solum magístrum, verum étiam Dóminum appellavérunt.\par
\switchcolumn
\EN
\noindent Without going into the rest, I may say that it is known what were the limits of Solomon's dominions, and yet in that Psalm we read: “He shall have dominion also from sea to sea, and from the river unto the ends of the earth.” This we see fulfilled in Christ. It was from the river, that is, from His Baptism by John in Jordan, that He began His assumption of dominion, for then it was when John bare testimony unto Him that His disciples began to acknowledge Him, calling Him, not Master only, but Lord.\par
\switchcolumn*

\LA
\begin{Resp}\Rsign Ego te tuli de domo patris tui, dicit Dóminus, et pósui te páscere gregem pópuli mei: \Star{} Et fui tecum in ómnibus ubicúmque ambulásti, firmans regnum tuum in ætérnum.\end{Resp}
\begin{Resp}\Vsign Fecíque tibi nomen grande, juxta nomen magnórum, qui sunt in terra: et réquiem dedi tibi ab ómnibus inimícis tuis.\end{Resp}
\begin{Resp}\Rsign Et fui tecum in ómnibus ubicúmque ambulásti, firmans regnum tuum in ætérnum.\end{Resp}
\begin{Resp}\Vsign Glória Patri, et Fílio, \Star{} et Spirítui Sancto.\end{Resp}
\begin{Resp}\Rsign Et fui tecum in ómnibus ubicúmque ambulásti, firmans regnum tuum in ætérnum.\end{Resp}
\switchcolumn
\EN
\begin{Resp}\Rsign Thus saith the Lord: I took thee out of thy father's house, and appointed thee to be ruler over My people, over Israel. \Star{} And I was with thee whithersoever thou wentest, to establish thy kingdom for ever.\end{Resp}
\begin{Resp}\Vsign And I have made thee a great name, like unto the name of the great men that are in the earth, and have caused thee to rest from all thine enemies.\end{Resp}
\begin{Resp}\Rsign And I was with thee whithersoever thou wentest, to establish thy kingdom for ever.\end{Resp}
\begin{Resp}\Vsign Glory be to the Father, and to the Son, \Star{} and to the Holy Ghost.\end{Resp}
\begin{Resp}\Rsign And I was with thee whithersoever thou wentest, to establish thy kingdom for ever.\end{Resp}
\switchcolumn*

\LA
\LessonHead{Lectio VII}
\switchcolumn
\EN
\LessonHead{Lesson VII}
\switchcolumn*

\LA
\LessonTitle{Léctio sancti Evangélii secúndum Lucam}
\switchcolumn
\EN
\LessonTitle{From the Holy Gospel according to Luke}
\switchcolumn*

\LA
\Cite{Luc 16:1-9}
\switchcolumn
\EN
\Cite{Luke 16:1-9}
\switchcolumn*

\LA
\noindent In illo témpore: Dixit Jesus discípulis suis parábolam hanc: Homo quidam erat dives, qui habébat víllicum, et hic diffamátus est apud illum quasi dissipásset bona ipsíus. Et réliqua.\par
\switchcolumn
\EN
\noindent At that time, Jesus spake this parable unto His disciples: There was a certain rich man, which had a steward and the same was accused unto him that he had wasted his goods. And so on.\par
\switchcolumn*

\LA
\LessonTitle{Homilía sancti Hierónymi Presbýteri}
\switchcolumn
\EN
\LessonTitle{Homily by St. Jerome, Priest at Bethlehem}
\switchcolumn*

\LA
\Source{Epistola 151 ad Algasium, quæst. 6, tom. 3}
\switchcolumn
\EN
\Source{Letter 151 to Algasia}
\switchcolumn*

\LA
\noindent Si dispensátor iníqui mammónæ, dómini voce laudátur, quod de re iníqua sibi justítiam præparárit; et passus dispéndia dóminus laudat dispensatóris prudéntiam, quod advérsus dóminum quidem fraudulénter, sed pro se prudénter égerit: quanto magis Christus, qui nullum damnum sustinére potest, et pronus est ad cleméntiam, laudábit discípulos suos, si in eos qui creditúri sibi sunt, misericórdes fúerint?\par
\switchcolumn
\EN
\noindent The lord commended the unjust steward, because he had done wisely though wickedly. The lord, although himself defrauded by it, could not but praise the shrewdness of his dishonest servant, because he had cheated him with profit to himself. How much more will our Master Christ, Who is above any defrauding by us, and is Himself the Great Forgiver, praise us if we win a blessing from Him by dealing indulgently with those who are to believe in Him?\par
\switchcolumn*

\LA
\begin{Resp}\Rsign Peccávi super númerum arénæ maris, et multiplicáta sunt peccáta mea: et non sum dignus vidére altitúdinem cæli præ multitúdine iniquitátis meæ: quóniam irritávi iram tuam, \Star{} Et malum coram te feci.\end{Resp}
\begin{Resp}\Vsign Quóniam iniquitátem meam ego cognósco: et delíctum meum contra me est semper, quia tibi soli peccávi.\end{Resp}
\begin{Resp}\Rsign Et malum coram te feci.\end{Resp}
\switchcolumn
\EN
\begin{Resp}\Rsign My sins are many, yea, they are more in number than the sands of the sea; I am not worthy to look up toward heaven because of the multitude of my iniquities; for I have provoked thee to anger; \Star{} And done evil in thy sight.\end{Resp}
\begin{Resp}\Vsign For I acknowledge my transgression, and my sin is ever before me, for against thee only have I sinned.\end{Resp}
\begin{Resp}\Rsign And done evil in thy sight.\end{Resp}
\switchcolumn*

\LA
\LessonHead{Lectio VIII}
\switchcolumn
\EN
\LessonHead{Lesson VIII}
\switchcolumn*

\LA
\noindent Dénique post parábolam íntulit: Et ego vobis dico, Fácite vobis amícos de iníquo mammóna. Mammóna autem non Hebræórum, sed Syrórum lingua divítiæ nuncupántur, quod de iniquitáte colléctæ sint. Si ergo iníquitas bene dispensáta vértitur in justítiam; quanto magis sermo divínus, in quo nulla est iníquitas, qui et Apóstolis créditus est, si bene fúerit dispensátus, dispensatóres suos levábit in cælum?\par
\switchcolumn
\EN
\noindent After this parable the Lord saith: “Make to yourselves friends of the mammon of unrighteousness.” This word “mammon” is a Syriac, not an Hebrew, word, signifying ill-gotten gains. If then even ill-gotten gains can be so used by such as have them as to profit them, how much more can they who, like the Apostles, are “stewards of the mysteries of God,” those true and blameless riches, how much more can they profit themselves, (1 Cor. iv. 1,) even everlastingly, by their right use of them?\par
\switchcolumn*

\LA
\begin{Resp}\Rsign Duo Séraphim clamábant alter ad álterum: \Star{} Sanctus, sanctus, sanctus Dóminus Deus Sábaoth: * Plena est omnis terra glória ejus.\end{Resp}
\begin{Resp}\Vsign Tres sunt qui testimónium dant in cælo: Pater, Verbum, et Spíritus Sanctus: et hi tres unum sunt.\end{Resp}
\begin{Resp}\Rsign Sanctus, sanctus, sanctus Dóminus Deus Sábaoth.\end{Resp}
\begin{Resp}\Vsign Glória Patri, et Fílio, \Star{} et Spirítui Sancto.\end{Resp}
\begin{Resp}\Rsign Plena est omnis terra glória ejus.\end{Resp}
\switchcolumn
\EN
\begin{Resp}\Rsign One Seraph cried unto another: \Star{} Holy, Holy, Holy is the Lord God of hosts; the whole earth is full of His glory.\end{Resp}
\begin{Resp}\Vsign There are Three That bear record in heaven: the Father, the Word, and the Holy Ghost, and these Three are One.\end{Resp}
\begin{Resp}\Rsign Holy, Holy, Holy is the Lord God of hosts.\end{Resp}
\begin{Resp}\Vsign Glory be to the Father, and to the Son, \Star{} and to the Holy Ghost.\end{Resp}
\begin{Resp}\Rsign The whole earth is full of His glory.\end{Resp}
\switchcolumn*

\LA
\LessonHead{Lectio IX}
\switchcolumn
\EN
\LessonHead{Lesson IX}
\switchcolumn*

\LA
\noindent Quam ob rem séquitur: Qui fidélis est in mínimo, hoc est, in carnálibus; et in multis fidélis erit, hoc est, in spirituálibus. Qui autem in parvo iníquus est, ut non det frátribus ad uténdum, quod a Deo pro ómnibus est creátum; iste et in spirituáli pecúnia dividénda iníquus erit, ut non pro necessitáte, sed pro persónis doctrínam Dómini dívidat. Si autem, inquit, carnáles divítia, quæ labúntur, non bene dispensátis; veras æternásque divítias doctrínæ Dei quis credet vobis?\par
\switchcolumn
\EN
\noindent Therefore it is immediately written: “He that is faithful in that which is least, that is to say, in bodily things, is faithful also in much that is to say, in spiritual things.” “And he that is unjust in the least”, that is to say, by not giving to his needy brother succour of those things which are needful for the body, and which God hath made for all men, such an one is unjust also in much that is to say, he will deal out spiritual things unfairly, this to one and that to another, and not according to their true spiritual needs. “If therefore,” saith the Lord, “ye have not been faithful in the “use of earthly riches which pass away, “who will commit to your trust the true and abiding riches,” that is, the spiritual riches of the word of God?\par
\switchcolumn*

\LA
\TeDeum{Te Deum}
\switchcolumn
\EN
\TeDeum{Te Deum}
\switchcolumn*

\end{paracol}

\Office{Feria secunda infra Hebdomadam VIII post Octavam Pentecostes}{Monday of the Eighth Week after Pentecost}{Feria}
\addcontentsline{toc}{subsection}{Feria secunda infra Hebdomadam VIII post Octavam Pentecostes \textit{— Monday of the Eighth Week after Pentecost}}
\begin{paracol}{2}
\LA
\LessonHead{Lectio I}
\switchcolumn
\EN
\LessonHead{Lesson I}
\switchcolumn*

\LA
\LessonTitle{De libro tértio Regum}
\switchcolumn
\EN
\LessonTitle{Lesson from the first book of Kings}
\switchcolumn*

\LA
\Cite{3 Reg 10:1-3}
\switchcolumn
\EN
\Cite{1 Kgs 10:1-3}
\switchcolumn*

\LA
\noindent Sed et regína Saba audíta fama Salomónis in nómine Dómini, venit tentáre eum in ænigmátibus. Et ingréssa Jerúsalem multo cum comitátu et divítiis, camélis portántibus arómata, et aurum infinítum nimis, et gemmas pretiósas, venit ad regem Salomónem, et locúta est ei univérsa quæ habébat in corde suo. Et dócuit eam Sálomon ómnia verba quæ proposúerat; non fuit sermo qui regem posset latére, et non respondéret ei.\par
\switchcolumn
\EN
\noindent And the queen of Saba, having heard of the fame of Solomon in the name of the Lord, came to try him with hard questions. And entering into Jerusalem with a great train, and riches, and camels that carried spices, and an immense quantity of gold, and precious stones, she came to king Solomon, and spoke to him all that she had in her heart. And Solomon informed her of all the things she proposed to him: there was not any word the king was ignorant of, and which he could not answer her\par
\switchcolumn*

\LA
\begin{Resp}\Rsign Recordáre, Dómine, testaménti tui, et dic Angelo percutiénti: Cesset jam manus tua, \Star{} Ut non desolétur terra, et ne perdas omnem ánimam vivam.\end{Resp}
\begin{Resp}\Vsign Ego sum qui peccávi, ego qui iníque egi: isti qui oves sunt, quid fecérunt? Avertátur, óbsecro, furor tuus, Dómine, a pópulo tuo.\end{Resp}
\begin{Resp}\Rsign Ut non desolétur terra, et ne perdas omnem ánimam vivam.\end{Resp}
\switchcolumn
\EN
\begin{Resp}\Rsign Remember, O Lord, thy covenant, and say unto the destroying Angel: Stay now thine hand; \Star{} That the land be not utterly laid waste, and that thou destroy not every living soul.\end{Resp}
\begin{Resp}\Vsign Even I it is that have sinned, and done evil indeed; but these sheep, what have they done? Let thine anger, I pray thee, O Lord, be turned away from thy people.\end{Resp}
\begin{Resp}\Rsign That the land be not utterly laid waste, and that Thou destroy not every living soul.\end{Resp}
\switchcolumn*

\LA
\LessonHead{Lectio II}
\switchcolumn
\EN
\LessonHead{Lesson II}
\switchcolumn*

\LA
\Cite{3 Reg 10:4-7}
\switchcolumn
\EN
\Cite{1 Kgs 10:4-7}
\switchcolumn*

\LA
\noindent Videns autem regína Saba omnem sapiéntiam Salomónis et domum quam ædificáverat, et cibos mensæ ejus, et habitácula servórum, et órdines ministrántium, vestésque eórum, et pincérnas, et holocáusta quæ offerébat in domo Dómini, non habébat ultra spíritum; dixítque ad regem: Verus est sermo, quem audívi in terra mea super sermónibus tuis, et super sapiéntia tua: et non credébam narrántibus mihi, donec ipsa veni, et vidi óculis meis, et probávi quod média pars mihi nuntiáta non fúerit. Major est sapiéntia et ópera tua, quam rumor quem audívi.\par
\switchcolumn
\EN
\noindent And when the queen of Saba saw all the wisdom of Solomon, and the house which he had built, And the meat of his table, and the apartments of his servants, and the order of his ministers, and their apparel, and the cupbearers, and the holocausts, which he offered in the house of the Lord: she had no longer any spirit in her, And she said to the king: The report is true, which I heard in my own country, Concerning thy words, and concerning thy wisdom. And I did not believe them that told me, till I came myself, and saw with my own eyes, and have found that the half hath not been told me: thy wisdom and thy works, exceed the fame which I heard.\par
\switchcolumn*

\LA
\begin{Resp}\Rsign Exaudísti, Dómine, oratiónem servi tui, ut ædificárem templum nómini tuo: \Star{} Bénedic et sanctífica domum istam in sempitérnum, Deus Israël.\end{Resp}
\begin{Resp}\Vsign Dómine, qui custódis pactum cum servis tuis, qui ámbulant coram te in toto corde suo.\end{Resp}
\begin{Resp}\Rsign Bénedic et sanctífica domum istam in sempitérnum, Deus Israël.\end{Resp}
\switchcolumn
\EN
\begin{Resp}\Rsign O Lord, Thou hast hearkened unto the prayer of thy servant, that I might build a temple unto thy Name, \Star{} O God of Israel, bless Thou, and hallow this house for ever.\end{Resp}
\begin{Resp}\Vsign O Lord, Who keepest covenant with thy servants that walk before thee in all their heart.\end{Resp}
\begin{Resp}\Rsign O God of Israel, bless Thou, and hallow this house for ever.\end{Resp}
\switchcolumn*

\LA
\LessonHead{Lectio III}
\switchcolumn
\EN
\LessonHead{Lesson III}
\switchcolumn*

\LA
\Cite{3 Reg 10:8-11}
\switchcolumn
\EN
\Cite{1 Kgs 10:8-11}
\switchcolumn*

\LA
\noindent Beáti viri tui, et beáti servi tui, qui stant coram te semper et áudiunt sapiéntiam tuam. Sit Dóminus Deus tuus benedíctus, cui complacuísti, et pósuit te super thronum Israël, eo quod diléxerit Dóminus Israël in sempitérnum, et constítuit te regem ut fáceres judícium et justítiam. Dedit ergo regi centum vigínti talénta auri, et arómata multa nimis, et gemmas pretiósas. Non sunt alláta ultra arómata tam multa, quam ea quæ dedit regína Saba regi Salomóni. Sed et classis Hiram, quæ portábat aurum de Ophir, áttulit ex Ophir ligna thyína multa nimis, et gemmas pretiósas.\par
\switchcolumn
\EN
\noindent Blessed are thy men, and blessed are thy servants, who stand before thee always, and hear thy wisdom. Blessed be the Lord thy God, whom thou hast pleased, and who hath set thee upon the throne of Israel, because the Lord hath loved Israel for ever, and hath appointed thee king, to do judgment and justice. And she gave the king a hundred and twenty talents of gold, and of spices a very great store, and precious stones: there was brought no more such abundance of spices as these which the queen of Saba gave to king Solomon. The navy also of Hiram, which brought gold from Ophir, brought from Ophir great plenty of thyine trees, and precious stones.\par
\switchcolumn*

\LA
\begin{Resp}\Rsign Audi, Dómine, hymnum et oratiónem, quam servus tuus orat coram te hódie, ut sint óculi tui apérti, et aures tuæ inténtæ, \Star{} Super domum istam die ac nocte.\end{Resp}
\begin{Resp}\Vsign Réspice, Dómine, de sanctuário tuo, et de excélso cælórum habitáculo.\end{Resp}
\begin{Resp}\Rsign Super domum istam die ac nocte.\end{Resp}
\begin{Resp}\Vsign Glória Patri, et Fílio, \Star{} et Spirítui Sancto.\end{Resp}
\begin{Resp}\Rsign Super domum istam die ac nocte.\end{Resp}
\switchcolumn
\EN
\begin{Resp}\Rsign Hearken, O Lord, unto the cry and to the prayer which thy servant prayeth before thee today, that thine eyes may be open and thine ears attend; \Star{} Toward this house day and night.\end{Resp}
\begin{Resp}\Vsign Look down from thine high and holy place, O Lord, even from heaven thy dwelling.\end{Resp}
\begin{Resp}\Rsign Toward this house, day and night.\end{Resp}
\begin{Resp}\Vsign Glory be to the Father, and to the Son, \Star{} and to the Holy Ghost.\end{Resp}
\begin{Resp}\Rsign Toward this house, day and night.\end{Resp}
\switchcolumn*

\end{paracol}

\Office{Feria tertia infra Hebdomadam VIII post Octavam Pentecostes}{Tuesday of the Eighth Week after Pentecost}{Feria}
\addcontentsline{toc}{subsection}{Feria tertia infra Hebdomadam VIII post Octavam Pentecostes \textit{— Tuesday of the Eighth Week after Pentecost}}
\begin{paracol}{2}
\LA
\LessonHead{Lectio I}
\switchcolumn
\EN
\LessonHead{Lesson I}
\switchcolumn*

\LA
\LessonTitle{De libro tértio Regum}
\switchcolumn
\EN
\LessonTitle{Lesson from the first book of Kings}
\switchcolumn*

\LA
\Cite{3 Reg 11:1-4}
\switchcolumn
\EN
\Cite{1 Kgs 11:1-4}
\switchcolumn*

\LA
\noindent Rex autem Sálomon adamávit mulíeres alienígenas multas, fíliam quoque pharaónis, et Moabítidas, et Ammonítidas, Idumǽas, et Sidónias, et Hethǽas: de géntibus super quibus dixit Dóminus fíliis Israël: Non ingrediémini ad eas, neque de illis ingrediéntur ad vestras; certíssime enim avértent corda vestra ut sequámini deos eárum. His ítaque copulátus est Sálomon ardentíssimo amóre: fuerúntque ei uxóres quasi regínæ septingéntæ, et concubínæ trecéntæ; et avertérunt mulíeres cor ejus. Cumque jam esset senex, depravátum est cor ejus per mulíeres, ut sequerétur deos aliénos; nec erat cor ejus perféctum cum Dómino Deo suo, sicut cor David patris ejus.\par
\switchcolumn
\EN
\noindent Of the nations concerning which the Lord said to the children of Israel: You shall not go in unto them, neither shall any of them come in to yours: for they will most certainly turn away your heart to follow their gods. And to these was Solomon joined with a most ardent love. And he had seven hundred wives as queens, and three hundred concubines: and the women turned away his heart. And when he was now old, his heart was turned away by women to follow strange gods: and his heart was not perfect with the Lord his God, as was the heart of David his father.\par
And king Solomon loved many strange women besides the daughter of Pharao, and women of Moab, and of Ammon, and of Edom, and of Sidon, and of the Hethites:\par
\switchcolumn*

\LA
\begin{Resp}\Rsign Dómine, si convérsus fúerit pópulus tuus, et oráverit ad sanctuárium tuum: \Star{} Tu exáudies de cælo, Dómine, et líbera eos de mánibus inimicórum suórum.\end{Resp}
\begin{Resp}\Vsign Si peccáverit in te pópulus tuus, et convérsus égerit pœniténtiam, veniénsque oráverit in isto loco.\end{Resp}
\begin{Resp}\Rsign Tu exáudies de cælo, Dómine, et líbera eos de mánibus inimicórum suórum.\end{Resp}
\switchcolumn
\EN
\begin{Resp}\Rsign Lord, when thy people shall turn again to thee, and shall pray unto thee in this house; \Star{} Then hear Thou in heaven, O Lord, and deliver them out of the hand of their enemies.\end{Resp}
\begin{Resp}\Vsign If thy people sin against thee, and turn again, and repent, and come and pray unto thee in this house.\end{Resp}
\begin{Resp}\Rsign Then hear Thou in heaven, O Lord, and deliver them out of the hand of their enemies.\end{Resp}
\switchcolumn*

\LA
\LessonHead{Lectio II}
\switchcolumn
\EN
\LessonHead{Lesson II}
\switchcolumn*

\LA
\Cite{3 Reg 11:5-8}
\switchcolumn
\EN
\Cite{1 Kgs 11:5-8}
\switchcolumn*

\LA
\noindent Sed colébat Sálomon Astárthen deam Sidoniórum, et Moloch idólum Ammonitárum. Fecítque Sálomon quod non placúerat coram Dómino, et non adimplévit ut sequerétur Dóminum sicut David pater ejus. Tunc ædificávit Sálomon fanum Chamos idólo Moab in monte, qui est contra Jerúsalem, et Moloch idólo filiórum Ammon. Atque in hunc modum fecit univérsis uxóribus suis alienígenis, quæ adolébant thura, et immolábant diis suis.\par
\switchcolumn
\EN
\noindent But Solomon worshipped Astarthe the goddess of the Sidonians, and Moloch the idol of the ammonites. And Solomon did that which was not pleasing before the Lord, and did not fully follow the Lord, as David his father. Then Solomon built a temple for Chamos the idol of Moab, on the hill that is over against Jerusalem, and for Moloch the idol of the children of Ammon. And he did in this manner for all his wives that were strangers, who burnt incense, and offered sacrifice to their gods.\par
\switchcolumn*

\LA
\begin{Resp}\Rsign Factum est, dum tólleret Dóminus Elíam per túrbinem in cælum, \Star{} Eliséus clamábat, dicens: Pater mi, pater mi, currus Israël, et auríga ejus.\end{Resp}
\begin{Resp}\Vsign Cumque pérgerent, et incedéntes sermocinaréntur, ecce currus ígneus et equi ígnei divisérunt utrúmque, et ascéndit Elías per túrbinem in cælum.\end{Resp}
\begin{Resp}\Rsign Eliséus clamábat, dicens: Pater mi, pater mi, currus Israël, et auríga ejus.\end{Resp}
\switchcolumn
\EN
\begin{Resp}\Rsign And it came to pass when the Lord would take up Elijah into heaven by a whirlwind \Star{} Elisha cried, and said: My father, my father, the chariot of Israel, and the horsemen thereof.\end{Resp}
\begin{Resp}\Vsign And as they still went on, and talked, behold, there appeared a chariot of fire, and horses of fire, and parted them both asunder, and Elijah went up by a whirlwind into heaven.\end{Resp}
\begin{Resp}\Rsign Elisha cried, and said: My father, my father, the chariot of Israel, and the horsemen thereof.\end{Resp}
\switchcolumn*

\LA
\LessonHead{Lectio III}
\switchcolumn
\EN
\LessonHead{Lesson III}
\switchcolumn*

\LA
\Cite{3 Reg 11:9-12}
\switchcolumn
\EN
\Cite{1 Kgs 11:9-12}
\switchcolumn*

\LA
\noindent Igitur irátus est Dóminus Salomóni, quod avérsa esset mens ejus a Dómino Deo Israël, qui apparúerat ei secúndo, et præcéperat de verbo hoc ne sequerétur deos aliénos: et non custodívit quæ mandávit ei Dóminus. Dixit ítaque Dóminus Salomóni: Quia habuísti hoc apud te, et non custodísti pactum meum, et præcépta mea quæ mandávi tibi, disrúmpens scindam regnum tuum, et dabo illud servo tuo. Verúmtamen in diébus tuis non fáciam, propter David patrem tuum.\par
\switchcolumn
\EN
\noindent And the Lord was angry with Solomon, because his mind was turned away from the Lord the God of Israel, who had appeared to him twice, And had commanded him concerning this thing, that he should not follow strange gods: but he kept not the things which the Lord commanded him. The Lord therefore said to Solomon: Because thou hast done this, and hast not kept my covenant, and my precepts, which I have commanded thee, I will divide and rend thy kingdom, and will give it to thy servant. Nevertheless in thy days I will not do it, for David thy father's sake: but I will rend it out of the hand of thy son\par
\switchcolumn*

\LA
\begin{Resp}\Rsign Ego te tuli de domo patris tui, dicit Dóminus, et pósui te páscere gregem pópuli mei: \Star{} Et fui tecum in ómnibus ubicúmque ambulásti, firmans regnum tuum in ætérnum.\end{Resp}
\begin{Resp}\Vsign Fecíque tibi nomen grande, juxta nomen magnórum, qui sunt in terra: et réquiem dedi tibi ab ómnibus inimícis tuis.\end{Resp}
\begin{Resp}\Rsign Et fui tecum in ómnibus ubicúmque ambulásti, firmans regnum tuum in ætérnum.\end{Resp}
\begin{Resp}\Vsign Glória Patri, et Fílio, \Star{} et Spirítui Sancto.\end{Resp}
\begin{Resp}\Rsign Et fui tecum in ómnibus ubicúmque ambulásti, firmans regnum tuum in ætérnum.\end{Resp}
\switchcolumn
\EN
\begin{Resp}\Rsign Thus saith the Lord: I took thee out of thy father's house, and appointed thee to be ruler over My people, over Israel. \Star{} And I was with thee whithersoever thou wentest, to establish thy kingdom for ever.\end{Resp}
\begin{Resp}\Vsign And I have made thee a great name, like unto the name of the great men that are in the earth, and have caused thee to rest from all thine enemies.\end{Resp}
\begin{Resp}\Rsign And I was with thee whithersoever thou wentest, to establish thy kingdom for ever.\end{Resp}
\begin{Resp}\Vsign Glory be to the Father, and to the Son, \Star{} and to the Holy Ghost.\end{Resp}
\begin{Resp}\Rsign And I was with thee whithersoever thou wentest, to establish thy kingdom for ever.\end{Resp}
\switchcolumn*

\end{paracol}

\Office{Feria quarta infra Hebdomadam VIII post Octavam Pentecostes}{Wednesday of the Eighth Week after Pentecost}{Feria}
\addcontentsline{toc}{subsection}{Feria quarta infra Hebdomadam VIII post Octavam Pentecostes \textit{— Wednesday of the Eighth Week after Pentecost}}
\begin{paracol}{2}
\LA
\LessonHead{Lectio I}
\switchcolumn
\EN
\LessonHead{Lesson I}
\switchcolumn*

\LA
\LessonTitle{De libro tértio Regum}
\switchcolumn
\EN
\LessonTitle{Lesson from the first book of Kings}
\switchcolumn*

\LA
\Cite{3 Reg 11:26-28}
\switchcolumn
\EN
\Cite{1 Kgs 11:26-28}
\switchcolumn*

\LA
\noindent Jeróboam quoque fílius Nabat Ephrathǽus de Saréda servus Salomónis, cujus mater erat nómine Sarva, múlier vídua, levávit manum contra regem. Et hæc causa rebelliónis advérsus eum, quia Sálomon ædificávit Mello, et coæquávit voráginem civitátis David patris sui. Erat autem Jeróboam vir fortis et potens; vidénsque Sálomon adolescéntem bonæ índolis et indústrium, constitúerat eum præféctum super tribúta univérsæ domus Joseph.\par
\switchcolumn
\EN
\noindent Jeroboam also the son of Nabat an Ephrathite of Sareda, a servant of Solomon, whose mother was named Sarua, a widow woman, lifted up his hand against the king. And this is the cause of his rebellion against him, for Solomon built Mello, and filled up the breach of the city of David his father. And Jeroboam was a valiant and mighty man: and Solomon seeing him a young man ingenious and industrious, made him chief over the tributes of all the house of Joseph.\par
\switchcolumn*

\LA
\begin{Resp}\Rsign Peccávi super númerum arénæ maris, et multiplicáta sunt peccáta mea: et non sum dignus vidére altitúdinem cæli præ multitúdine iniquitátis meæ: quóniam irritávi iram tuam, \Star{} Et malum coram te feci.\end{Resp}
\begin{Resp}\Vsign Quóniam iniquitátem meam ego cognósco: et delíctum meum contra me est semper, quia tibi soli peccávi.\end{Resp}
\begin{Resp}\Rsign Et malum coram te feci.\end{Resp}
\switchcolumn
\EN
\begin{Resp}\Rsign My sins are many, yea, they are more in number than the sands of the sea; I am not worthy to look up toward heaven because of the multitude of my iniquities; for I have provoked thee to anger; \Star{} And done evil in thy sight.\end{Resp}
\begin{Resp}\Vsign For I acknowledge my transgression, and my sin is ever before me, for against thee only have I sinned.\end{Resp}
\begin{Resp}\Rsign And done evil in thy sight.\end{Resp}
\switchcolumn*

\LA
\LessonHead{Lectio II}
\switchcolumn
\EN
\LessonHead{Lesson II}
\switchcolumn*

\LA
\Cite{3 Reg 11:29-31}
\switchcolumn
\EN
\Cite{1 Kgs 11:29-31}
\switchcolumn*

\LA
\noindent Factum est ígitur in témpore illo, ut Jeróboam egrederétur de Jerúsalem, et inveníret eum Ahías Silónites prophéta in via opértus pállio novo; erant autem duo tantum in agro. Apprehendénsque Ahías pállium suum novum, quo coopértus erat, scidit in duódecim partes, et ait ad Jeróboam: Tolle tibi decem scissúras; hæc enim dicit Dóminus Deus Israël: Ecce ego scindam regnum de manu Salomónis, et dabo tibi decem tribus.\par
\switchcolumn
\EN
\noindent So it came to pass at that time, that Jeroboam went out of Jerusalem, and the prophet Ahias the Silonite, clad with a new garment, found him in the way: and they two were alone in the field. And Ahias taking his new garment, wherewith he was clad, divided it into twelve parts: And he said to Jeroboam: Take to thee ten pieces: for thus saith the Lord the God of Israel: Behold I will rend the kingdom out of the hand of Solomon, and will give thee ten tribes.\par
\switchcolumn*

\LA
\begin{Resp}\Rsign Exaudísti, Dómine, oratiónem servi tui, ut ædificárem templum nómini tuo: \Star{} Bénedic et sanctífica domum istam in sempitérnum, Deus Israël.\end{Resp}
\begin{Resp}\Vsign Dómine, qui custódis pactum cum servis tuis, qui ámbulant coram te in toto corde suo.\end{Resp}
\begin{Resp}\Rsign Bénedic et sanctífica domum istam in sempitérnum, Deus Israël.\end{Resp}
\switchcolumn
\EN
\begin{Resp}\Rsign O Lord, Thou hast hearkened unto the prayer of thy servant, that I might build a temple unto thy Name, \Star{} O God of Israel, bless Thou, and hallow this house for ever.\end{Resp}
\begin{Resp}\Vsign O Lord, Who keepest covenant with thy servants that walk before thee in all their heart.\end{Resp}
\begin{Resp}\Rsign O God of Israel, bless Thou, and hallow this house for ever.\end{Resp}
\switchcolumn*

\LA
\LessonHead{Lectio III}
\switchcolumn
\EN
\LessonHead{Lesson III}
\switchcolumn*

\LA
\Cite{3 Reg 11:40-43}
\switchcolumn
\EN
\Cite{1 Kgs 11:40-43}
\switchcolumn*

\LA
\noindent Vóluit ergo Sálomon interfícere Jeróboam, qui surréxit, et aufúgit in Ægýptum ad Sesac regem Ægýpti, et fuit in Ægýpto usque ad mortem Salomónis. Réliquum autem verbórum Salomónis, et ómnia quæ fecit, et sapiéntia ejus, ecce univérsa scripta sunt in libro verbórum diérum Salomónis. Dies autem, quos regnávit Sálomon in Jerúsalem super omnem Israël, quadragínta anni sunt. Dormivítque Sálomon cum pátribus suis, et sepúltus est in civitáte David patris sui: regnavítque Róboam fílius ejus pro eo.\par
\switchcolumn
\EN
\noindent Solomon therefore sought to kill Jeroboam: but he arose, and fled into Egypt to Sesac the king of Egypt, and was in Egypt till the death of Solomon. And the rest of the words of Solomon, and all that he did, and his wisdom: behold they are all written in the book of the words of the days of Solomon. And the days that Solomon reigned in Jerusalem over all Israel, were forty years. And Solomon slept with his fathers, and was buried in the city of David his father, and Roboam his son reigned in his stead.\par
\switchcolumn*

\LA
\begin{Resp}\Rsign Audi, Dómine, hymnum et oratiónem, quam servus tuus orat coram te hódie, ut sint óculi tui apérti, et aures tuæ inténtæ, \Star{} Super domum istam die ac nocte.\end{Resp}
\begin{Resp}\Vsign Réspice, Dómine, de sanctuário tuo, et de excélso cælórum habitáculo.\end{Resp}
\begin{Resp}\Rsign Super domum istam die ac nocte.\end{Resp}
\begin{Resp}\Vsign Glória Patri, et Fílio, \Star{} et Spirítui Sancto.\end{Resp}
\begin{Resp}\Rsign Super domum istam die ac nocte.\end{Resp}
\switchcolumn
\EN
\begin{Resp}\Rsign Hearken, O Lord, unto the cry and to the prayer which thy servant prayeth before thee today, that thine eyes may be open and thine ears attend; \Star{} Toward this house day and night.\end{Resp}
\begin{Resp}\Vsign Look down from thine high and holy place, O Lord, even from heaven thy dwelling.\end{Resp}
\begin{Resp}\Rsign Toward this house, day and night.\end{Resp}
\begin{Resp}\Vsign Glory be to the Father, and to the Son, \Star{} and to the Holy Ghost.\end{Resp}
\begin{Resp}\Rsign Toward this house, day and night.\end{Resp}
\switchcolumn*

\end{paracol}

\Office{Feria quinta infra Hebdomadam VIII post Octavam Pentecostes}{Thursday of the Eighth Week after Pentecost}{Feria}
\addcontentsline{toc}{subsection}{Feria quinta infra Hebdomadam VIII post Octavam Pentecostes \textit{— Thursday of the Eighth Week after Pentecost}}
\begin{paracol}{2}
\LA
\LessonHead{Lectio I}
\switchcolumn
\EN
\LessonHead{Lesson I}
\switchcolumn*

\LA
\LessonTitle{De libro tértio Regum}
\switchcolumn
\EN
\LessonTitle{Lesson from the first book of Kings}
\switchcolumn*

\LA
\Cite{3 Reg 12:1-5}
\switchcolumn
\EN
\Cite{1 Kgs 12:1-5}
\switchcolumn*

\LA
\noindent Venit autem Róboam in Sichem; illuc enim congregátus erat omnis Israël ad constituéndum eum regem. At vero Jeróboam fílius Nabat, cum adhuc esset in Ægýpto prófugus a fácie regis Salomónis, audíta morte ejus, revérsus est de Ægýpto. Miserúntque et vocavérunt eum. Venit ergo Jeróboam, et omnis multitúdo Israël, et locúti sunt ad Róboam, dicéntes: Pater tuus duríssimum jugum impósuit nobis; tu ítaque nunc immínue páululum de império patris tui duríssimo, et de jugo gravíssimo quod impósuit nobis, et serviémus tibi. Qui ait eis: Ite usque ad tértium diem, et revertímini ad me.\par
\switchcolumn
\EN
\noindent And Roboam went to Sichem: for thither were all Israel come together to make him king. But Jeroboam the son of Nabat, who was yet in Egypt, a fugitive from the face of king Solomon, hearing of his death, returned out of Egypt. And they sent and called him: and Jeroboam came, and all the multitude of Israel, and they spoke to Roboam, saying: thy father laid a grievous yoke upon us: now therefore do thou take off a little of the grievous service of thy father, and of his most heavy yoke, which he put upon us, and we will serve thee. And he said to them: Go till the third day, and come to me again.\par
\switchcolumn*

\LA
\begin{Resp}\Rsign Præparáte corda vestra Dómino, et servíte illi soli: \Star{} Et liberábit vos de mánibus inimicórum vestrórum.\end{Resp}
\begin{Resp}\Vsign Convertímini ad eum in toto corde vestro, et auférte deos aliénos de médio vestri.\end{Resp}
\begin{Resp}\Rsign Et liberábit vos de mánibus inimicórum vestrórum.\end{Resp}
\switchcolumn
\EN
\begin{Resp}\Rsign Prepare your hearts unto the Lord, and serve Him Only \Star{} And He will deliver you out of the hand of your enemies.\end{Resp}
\begin{Resp}\Vsign Return unto Him with all your hearts, and put away the strange gods from among you.\end{Resp}
\begin{Resp}\Rsign And He will deliver you out of the hand of your enemies.\end{Resp}
\switchcolumn*

\LA
\LessonHead{Lectio II}
\switchcolumn
\EN
\LessonHead{Lesson II}
\switchcolumn*

\LA
\Cite{3 Reg 12:5-8}
\switchcolumn
\EN
\Cite{1 Kgs 12:5-8}
\switchcolumn*

\LA
\noindent Cumque abiísset pópulus, íniit consílium rex Róboam cum senióribus, qui assistébant coram Salomóne patre ejus cum adhuc víveret, et ait: Quod datis mihi consílium ut respóndeam pópulo huic? Qui dixérunt ei: Si hódie obedíeris pópulo huic et servíeris et petitióni eórum césseris locutúsque fúeris ad eos verba lénia, erunt tibi servi cunctis diébus. Qui derelíquit consílium senum, quod déderant ei, at adhíbuit adolescéntes, qui nutríti fúerant cum eo, et assistébant illi.\par
\switchcolumn
\EN
\noindent And when the people was gone, King Roboam took counsel with the old men, that stood before Solomon his father while he yet lived, and he said: What counsel do you give me, that I may answer this people? They said to him: If thou wilt yield to this people today, and condescend to them, and grant their petition, and wilt speak gentle words to them, they will be thy servants always. But he left the counsel of the old men, which they had given him, and consulted with the young men, that had been brought up with him, and stood before him.\par
\switchcolumn*

\LA
\begin{Resp}\Rsign Deus ómnium exaudítor est: ipse misit Angelum suum, et tulit me de óvibus patris mei: \Star{} Et unxit me unctióne misericórdiæ suæ.\end{Resp}
\begin{Resp}\Vsign Dóminus, qui erípuit me de ore leónis, et de manu béstiæ liberávit me.\end{Resp}
\begin{Resp}\Rsign Et unxit me unctióne misericórdiæ suæ.\end{Resp}
\switchcolumn
\EN
\begin{Resp}\Rsign God, Which heareth all, even He sent His Angel, and took me from keeping my father's sheep, and \Star{} Anointed me with the oil of His mercy.\end{Resp}
\begin{Resp}\Vsign The Lord That delivered me out of the mouth of the lion, and out of the paw of the bear\end{Resp}
\begin{Resp}\Rsign And anointed me with the oil of His mercy.\end{Resp}
\switchcolumn*

\LA
\LessonHead{Lectio III}
\switchcolumn
\EN
\LessonHead{Lesson III}
\switchcolumn*

\LA
\Cite{3 Reg 12:13-16}
\switchcolumn
\EN
\Cite{1 Kgs 12:13-16}
\switchcolumn*

\LA
\noindent Respondítque rex pópulo dura, derelícto consílio seniórum quod ei déderant, et locútus est eis secúndum consílium júvenum dicens: Pater meus aggravávit jugum vestrum, ego autem addam jugo vestro; pater meus cécidit vos flagéllis, ego autem cædam vos scorpiónibus. Et non acquiévit rex pópulo, quóniam aversátus fúerat eum Dóminus, ut suscitáret verbum suum, quod locútus fúerat in manu Ahíæ Silonítæ, ad Jeróboam fílium Nabat. Videns ítaque pópulus quod noluísset eos audíre rex, respóndit ei, dicens: Quæ nobis pars in David? vel quæ heréditas in fílio Isai?\par
\switchcolumn
\EN
\noindent And the king answered the people roughly, leaving the counsel of the old men, which they had given him, And he spoke to them according to the counsel of the young men, saying: My father made your yoke heavy, but I will add to your yoke: my father beat you with whips, but I will beat you with scorpions. And the king condescended not to the people: for the Lord was turned away from him, to make good his word, which he had spoken in the hand of Ahias the Silonite, to Jeroboam the son of Nabat. Then the people seeing that the king would not hearken to them, answered him, saying: What portion have we in David? or what inheritance in the son of Isai? Go home to thy dwellings, O Israel, now David look to thy own house. So Israel departed to their dwellings.\par
\switchcolumn*

\LA
\begin{Resp}\Rsign Dóminus, qui erípuit me de ore leónis, et de manu béstiæ liberávit me, \Star{} Ipse me erípiet de mánibus inimicórum meórum.\end{Resp}
\begin{Resp}\Vsign Misit Deus misericórdiam suam et veritátem suam: ánimam meam erípuit de médio catulórum leónum.\end{Resp}
\begin{Resp}\Rsign Ipse me erípiet de mánibus inimicórum meórum.\end{Resp}
\begin{Resp}\Vsign Glória Patri, et Fílio, \Star{} et Spirítui Sancto.\end{Resp}
\begin{Resp}\Rsign Ipse me erípiet de mánibus inimicórum meórum.\end{Resp}
\switchcolumn
\EN
\begin{Resp}\Rsign The Lord That delivered me out of the mouth of the lion, and out of the paw of the bear \Star{} He will deliver me out of the hand of mine enemies.\end{Resp}
\begin{Resp}\Vsign God hath sent forth His mercy and His truth, and delivered my soul from among the lion's whelps.\end{Resp}
\begin{Resp}\Rsign He will deliver me out of the hand of mine enemies.\end{Resp}
\begin{Resp}\Vsign Glory be to the Father, and to the Son, \Star{} and to the Holy Ghost.\end{Resp}
\begin{Resp}\Rsign He will deliver me out of the hand of mine enemies.\end{Resp}
\switchcolumn*

\end{paracol}

\Office{Feria sexta infra Hebdomadam VIII post Octavam Pentecostes}{Friday of the Eighth Week after Pentecost}{Feria}
\addcontentsline{toc}{subsection}{Feria sexta infra Hebdomadam VIII post Octavam Pentecostes \textit{— Friday of the Eighth Week after Pentecost}}
\begin{paracol}{2}
\LA
\LessonHead{Lectio I}
\switchcolumn
\EN
\LessonHead{Lesson I}
\switchcolumn*

\LA
\LessonTitle{De libro tértio Regum}
\switchcolumn
\EN
\LessonTitle{Lesson from the first book of Kings}
\switchcolumn*

\LA
\Cite{3 Reg 14:5-6}
\switchcolumn
\EN
\Cite{1 Kgs 14:5-6}
\switchcolumn*

\LA
\noindent Dixit autem Dóminus ad Ahíam: Ecce uxor Jeróboam ingréditur ut cónsulat te super fílio suo, qui ægrótat: hæc et hæc loquéris ei. Cum ergo illa intráret et dissimuláret se esse quæ erat, audívit Ahías sónitum pedum ejus introëúntis per óstium, et ait: Ingrédere, uxor Jeróboam: quare áliam te esse símulas? Ego autem missus sum ad te durus núntius.\par
\switchcolumn
\EN
\noindent And the Lord said to Ahias: Behold the wife of Jeroboam cometh in, to consult thee concerning her son that is sick: thus and thus shalt thou speak to her. So when she was coming in, and made as if she were another woman, Ahias heard the sound of her feet coming in at the door, and said: Come in, thou wife of Jeroboam: why dost thou feign thyself to be another? But I am sent to thee with heavy tidings.\par
\switchcolumn*

\LA
\begin{Resp}\Rsign Percússit Saul mille, et David decem míllia: \Star{} Quia manus Dómini erat cum illo: percússit Philisthǽum, et ábstulit oppróbrium ex Israël.\end{Resp}
\begin{Resp}\Vsign Nonne iste est David, de quo canébant in choro, dicéntes: Saul percússit mille, et David decem míllia?\end{Resp}
\begin{Resp}\Rsign Quia manus Dómini erat cum illo: percússit Philisthǽum, et ábstulit oppróbrium ex Israël.\end{Resp}
\switchcolumn
\EN
\begin{Resp}\Rsign Saul hath slain his thousands, and David his ten thousands. \Star{} Because the hand of the Lord was with him, he smote the Philistine, and took away the reproach from Israel.\end{Resp}
\begin{Resp}\Vsign Is not this David? Did they not sing one to another of him in dances, saying: Saul hath slain his thousands, and David his ten thousands?\end{Resp}
\begin{Resp}\Rsign Because the hand of the Lord was with him, he smote the Philistine, and took away the reproach from Israel.\end{Resp}
\switchcolumn*

\LA
\LessonHead{Lectio II}
\switchcolumn
\EN
\LessonHead{Lesson II}
\switchcolumn*

\LA
\Cite{3 Reg 14:7-9}
\switchcolumn
\EN
\Cite{1 Kgs 14:7-9}
\switchcolumn*

\LA
\noindent Vade et dic Jeróboam: Hæc dicit Dóminus Deus Israël: Quia exaltávi te de médio pópuli, et dedi te ducem super pópulum meum Israël, et scidi regnum domus David, et dedi illud tibi, et non fuísti sicut servus meus David, qui custodívit mandáta mea, et secútus est me in toto corde suo, fáciens quod plácitum esset in conspéctu meo; sed operátus es mala super omnes, qui fuérunt ante te, et fecísti tibi deos aliénos et conflátiles, ut me ad iracúndiam provocáres; me autem projecísti post corpus tuum.\par
\switchcolumn
\EN
\noindent Go, and tell Jeroboam: Thus saith the Lord the God of Israel: Forasmuch as I exalted thee from among the people, and made thee prince over my people Israel: And rent the kingdom away from the house of David, and gave it to thee, and thou hast not been as my servant David, who kept my commandments, and followed me with all his heart, doing that which was well pleasing in my sight: But hast done evil above all that were before thee, and hast made thee strange gods and molten gods, to provoke me to anger, and hast cast me behind thy back\par
\switchcolumn*

\LA
\begin{Resp}\Rsign Montes Gélboë, nec ros nec plúvia véniant super vos, \Star{} Ubi cecidérunt fortes Israël.\end{Resp}
\begin{Resp}\Vsign Omnes montes, qui estis in circúitu ejus, vísitet Dóminus: a Gélboë autem tránseat.\end{Resp}
\begin{Resp}\Rsign Ubi cecidérunt fortes Israël.\end{Resp}
\switchcolumn
\EN
\begin{Resp}\Rsign Ye mountains of Gilboa, let there be no dew, neither let there be rain upon you \Star{} For there are the mighty of Israel fallen.\end{Resp}
\begin{Resp}\Vsign All ye mountains that stand round about, the Lord look upon you but let Him pass by Gilboa\end{Resp}
\begin{Resp}\Rsign For there are the mighty of Israel fallen.\end{Resp}
\switchcolumn*

\LA
\LessonHead{Lectio III}
\switchcolumn
\EN
\LessonHead{Lesson III}
\switchcolumn*

\LA
\Cite{3 Reg 14:10-12}
\switchcolumn
\EN
\Cite{1 Kgs 14:10-12}
\switchcolumn*

\LA
\noindent Idcírco ecce ego indúcam mala super domum Jeróboam, et percútiam de Jeróboam mingéntem ad paríetem, et clausum, et novíssimum in Israël; et mundábo relíquias domus Jeróboam, sicut mundári solet fimus usque ad purum. Qui mórtui fúerint de Jeróboam in civitáte, cómedent eos canes; qui autem mórtui fúerint in agro, vorábunt eos aves cæli: quia Dóminus locútus est. Tu ígitur surge, et vade in domum tuam, et in ipso intróitu pedum tuórum in urbem, moriétur puer.\par
\switchcolumn
\EN
\noindent Therefore behold I will bring evils upon the house of Jeroboam, and will cut of from Jeroboam him that pisseth against the wall, and him that is shut up, and the last in Israel: and I will sweep away the remnant of the house of Jeroboam, as dung is swept away till all be clean. Them that shall die of Jeroboam in the city, the dogs shall eat: and them that shall die in the field, the birds of the air shall devour: for the Lord hath spoken it. Arise thou therefore, and go to thy house: and when thy feet shall be entering into the city, the child shall die,\par
\switchcolumn*

\LA
\begin{Resp}\Rsign Ego te tuli de domo patris tui, dicit Dóminus, et pósui te páscere gregem pópuli mei: \Star{} Et fui tecum in ómnibus ubicúmque ambulásti, firmans regnum tuum in ætérnum.\end{Resp}
\begin{Resp}\Vsign Fecíque tibi nomen grande, juxta nomen magnórum, qui sunt in terra: et réquiem dedi tibi ab ómnibus inimícis tuis.\end{Resp}
\begin{Resp}\Rsign Et fui tecum in ómnibus ubicúmque ambulásti, firmans regnum tuum in ætérnum.\end{Resp}
\begin{Resp}\Vsign Glória Patri, et Fílio, \Star{} et Spirítui Sancto.\end{Resp}
\begin{Resp}\Rsign Et fui tecum in ómnibus ubicúmque ambulásti, firmans regnum tuum in ætérnum.\end{Resp}
\switchcolumn
\EN
\begin{Resp}\Rsign Thus saith the Lord: I took thee out of thy father's house, and appointed thee to be ruler over My people, over Israel. \Star{} And I was with thee whithersoever thou wentest, to establish thy kingdom for ever.\end{Resp}
\begin{Resp}\Vsign And I have made thee a great name, like unto the name of the great men that are in the earth, and have caused thee to rest from all thine enemies.\end{Resp}
\begin{Resp}\Rsign And I was with thee whithersoever thou wentest, to establish thy kingdom for ever.\end{Resp}
\begin{Resp}\Vsign Glory be to the Father, and to the Son, \Star{} and to the Holy Ghost.\end{Resp}
\begin{Resp}\Rsign And I was with thee whithersoever thou wentest, to establish thy kingdom for ever.\end{Resp}
\switchcolumn*

\end{paracol}

\Office{Sabbato infra Hebdomadam VIII post Octavam Pentecostes}{Saturday of the Eighth Week after Pentecost}{Feria}
\addcontentsline{toc}{subsection}{Sabbato infra Hebdomadam VIII post Octavam Pentecostes \textit{— Saturday of the Eighth Week after Pentecost}}
\begin{paracol}{2}
\LA
\LessonHead{Lectio I}
\switchcolumn
\EN
\LessonHead{Lesson I}
\switchcolumn*

\LA
\LessonTitle{De libro tértio Regum}
\switchcolumn
\EN
\LessonTitle{Lesson from the first book of Kings}
\switchcolumn*

\LA
\Cite{3 Reg 18:21-22}
\switchcolumn
\EN
\Cite{1 Kgs 18:21-22}
\switchcolumn*

\LA
\noindent Accédens autem Elías ad omnem pópulum ait: Usquequo claudicátis in duas partes? Si Dóminus est Deus, sequímini eum; si autem Baal, sequímini illum. Et non respóndit ei pópulus verbum. Et ait rursus Elías ad pópulum: Ego remánsi prophéta Dómini solus, prophétæ autem Baal quadringénti et quinquagínta viri sunt.\par
\switchcolumn
\EN
\noindent And Elias coming to all the people, said: How long do you halt between two sides? if the Lord be God, follow him: but if Baal, then follow him. And the people did not answer him a word. And Elias said again to the people: I only remain a prophet of the Lord: but the prophets of Baal are four hundred and fifty men.\par
\switchcolumn*

\LA
\begin{Resp}\Rsign Peccávi super númerum arénæ maris, et multiplicáta sunt peccáta mea: et non sum dignus vidére altitúdinem cæli præ multitúdine iniquitátis meæ: quóniam irritávi iram tuam, \Star{} Et malum coram te feci.\end{Resp}
\begin{Resp}\Vsign Quóniam iniquitátem meam ego cognósco: et delíctum meum contra me est semper, quia tibi soli peccávi.\end{Resp}
\begin{Resp}\Rsign Et malum coram te feci.\end{Resp}
\switchcolumn
\EN
\begin{Resp}\Rsign My sins are many, yea, they are more in number than the sands of the sea; I am not worthy to look up toward heaven because of the multitude of my iniquities; for I have provoked thee to anger; \Star{} And done evil in thy sight.\end{Resp}
\begin{Resp}\Vsign For I acknowledge my transgression, and my sin is ever before me, for against thee only have I sinned.\end{Resp}
\begin{Resp}\Rsign And done evil in thy sight.\end{Resp}
\switchcolumn*

\LA
\LessonHead{Lectio II}
\switchcolumn
\EN
\LessonHead{Lesson II}
\switchcolumn*

\LA
\Cite{3 Reg 18:23-24}
\switchcolumn
\EN
\Cite{1 Kgs 18:23-24}
\switchcolumn*

\LA
\noindent Dentur nobis duo boves; et illi éligant sibi bovem unum, et in frusta cædéntes ponant super ligna, ignem autem non suppónant; et ego fáciam bovem álterum et impónam super ligna, ignem autem non suppónam. Invocáte nómina deórum vestrórum, et ego invocábo nomen Dómini mei; et Deus, qui exaudíerit per ignem, ipse sit Deus. Respóndens omnis pópulus ait: Optima proposítio.\par
\switchcolumn
\EN
\noindent Let two bullocks be given us, and let them choose one bullock for themselves, and cut it in pieces and lay it upon wood, but put no fire under: and I will dress the other bullock, and lay it on wood, and put no fire under it. Call ye on the names of your gods, and I will call on the name of my Lord: and the God that shall answer by fire, let him be God. And all the people answering said: A very good proposal.\par
\switchcolumn*

\LA
\begin{Resp}\Rsign Exaudísti, Dómine, oratiónem servi tui, ut ædificárem templum nómini tuo: \Star{} Bénedic et sanctífica domum istam in sempitérnum, Deus Israël.\end{Resp}
\begin{Resp}\Vsign Dómine, qui custódis pactum cum servis tuis, qui ámbulant coram te in toto corde suo.\end{Resp}
\begin{Resp}\Rsign Bénedic et sanctífica domum istam in sempitérnum, Deus Israël.\end{Resp}
\switchcolumn
\EN
\begin{Resp}\Rsign O Lord, Thou hast hearkened unto the prayer of thy servant, that I might build a temple unto thy Name, \Star{} O God of Israel, bless Thou, and hallow this house for ever.\end{Resp}
\begin{Resp}\Vsign O Lord, Who keepest covenant with thy servants that walk before thee in all their heart.\end{Resp}
\begin{Resp}\Rsign O God of Israel, bless Thou, and hallow this house for ever.\end{Resp}
\switchcolumn*

\LA
\LessonHead{Lectio III}
\switchcolumn
\EN
\LessonHead{Lesson III}
\switchcolumn*

\LA
\Cite{3 Reg 18:25-27}
\switchcolumn
\EN
\Cite{1 Kgs 18:25-27}
\switchcolumn*

\LA
\noindent Dixit ergo Elías prophétis Baal: Elígite vobis bovem unum et fácite primi, quia vos plures estis, et invocáte nómina deórum vestrórum, ignémque non supponátis. Qui, cum tulíssent bovem quem déderat eis, fecérunt, et invocábant nomen Baal de mane usque ad merídiem, dicéntes: Baal, exáudi nos. Et non erat vox, nec qui respondéret. Transiliebántque altáre quod fécerant. Cumque esset jam merídies, illudébat illis Elías, dicens: Clamáte voce majóre.\par
\switchcolumn
\EN
\noindent Then Elias said to the prophets of Baal: Choose you one bullock and dress it first, because you are many: and call on the names of your gods, but put no fire under. And they took the bullock which he gave them, and dressed it: and they called on the name of Baal from morning even till noon, saying: O Baal, hear us. But there was no voice, nor any that answered: and they leaped over the altar that they had made. And when it was now noon, Elias jested at them, saying: Cry with a louder voice: for he is a God, and perhaps he is talking, or is in an inn, or on a journey, or perhaps he is asleep, and must be awaked.\par
\switchcolumn*

\LA
\begin{Resp}\Rsign Audi, Dómine, hymnum et oratiónem, quam servus tuus orat coram te hódie, ut sint óculi tui apérti, et aures tuæ inténtæ, \Star{} Super domum istam die ac nocte.\end{Resp}
\begin{Resp}\Vsign Réspice, Dómine, de sanctuário tuo, et de excélso cælórum habitáculo.\end{Resp}
\begin{Resp}\Rsign Super domum istam die ac nocte.\end{Resp}
\begin{Resp}\Vsign Glória Patri, et Fílio, \Star{} et Spirítui Sancto.\end{Resp}
\begin{Resp}\Rsign Super domum istam die ac nocte.\end{Resp}
\switchcolumn
\EN
\begin{Resp}\Rsign Hearken, O Lord, unto the cry and to the prayer which thy servant prayeth before thee today, that thine eyes may be open and thine ears attend; \Star{} Toward this house day and night.\end{Resp}
\begin{Resp}\Vsign Look down from thine high and holy place, O Lord, even from heaven thy dwelling.\end{Resp}
\begin{Resp}\Rsign Toward this house, day and night.\end{Resp}
\begin{Resp}\Vsign Glory be to the Father, and to the Son, \Star{} and to the Holy Ghost.\end{Resp}
\begin{Resp}\Rsign Toward this house, day and night.\end{Resp}
\switchcolumn*

\end{paracol}

\Office{Dominica IX Post Pentecosten}{Ninth Sunday after Pentecost}{Semiduplex}
\addcontentsline{toc}{subsection}{Dominica IX Post Pentecosten \textit{— Ninth Sunday after Pentecost}}
\begin{paracol}{2}
\LA
\RefLine{Lectiones I–VI: de Scriptura occurrente.}
\switchcolumn
\EN
\RefLine{Lessons I–VI: from the Scripture of the day.}
\switchcolumn*

\LA
\LessonHead{Lectio I}
\switchcolumn
\EN
\LessonHead{Lesson I}
\switchcolumn*

\LA
\LessonTitle{Incipit liber quartus Regum}
\switchcolumn
\EN
\LessonTitle{Lesson from the second book of Kings}
\switchcolumn*

\LA
\Cite{4 Reg 1:1-4}
\switchcolumn
\EN
\Cite{2 Kgs 1:1-4}
\switchcolumn*

\LA
\noindent Prævaricátus est autem Moab in Israël, postquam mórtuus est Achab. Cecidítque Ochozías per cancéllos cœnáculi sui, quod habébat in Samaría, et ægrotávit; misítque núntios dicens ad eos: Ite, consúlite Beélzebub deum Accaron, utrum vívere queam de infirmitáte mea hac. Angelus autem Dómini locútus est ad Elíam Thesbíten dicens: Surge et ascénde in occúrsum nuntiórum regis Samaríæ et dices ad eos: Numquid non est Deus in Israël, ut eátis ad consuléndum Beélzebub deum Accaron? Quam ob rem hæc dicit Dóminus: De léctulo, super quem ascendísti, non descéndes, sed morte moriéris.\par
\switchcolumn
\EN
\noindent And Moab rebelled against Israel, after the death of Achab. And Ochozias fell through the lattices of his upper chamber which he had in Samaria, and was sick: and he sent messengers, saying to them: Go, consult Beelzebub, the god of Accaron, whether I shall recover of this my illness. And an angel of the Lord spoke to Elias the Thesbite, saying: Arise, and go up to meet the messengers of the king of Samaria, and say to them: Is there not a God in Israel, that ye go to consult Beelzebub the god of Accaron? Wherefore thus saith the Lord: From the bed, on which thou art gone up, thou shalt not come down, but thou shalt surely die.\par
\switchcolumn*

\LA
\begin{Resp}\Rsign Præparáte corda vestra Dómino, et servíte illi soli: \Star{} Et liberábit vos de mánibus inimicórum vestrórum.\end{Resp}
\begin{Resp}\Vsign Convertímini ad eum in toto corde vestro, et auférte deos aliénos de médio vestri.\end{Resp}
\begin{Resp}\Rsign Et liberábit vos de mánibus inimicórum vestrórum.\end{Resp}
\switchcolumn
\EN
\begin{Resp}\Rsign Prepare your hearts unto the Lord, and serve Him Only \Star{} And He will deliver you out of the hand of your enemies.\end{Resp}
\begin{Resp}\Vsign Return unto Him with all your hearts, and put away the strange gods from among you.\end{Resp}
\begin{Resp}\Rsign And He will deliver you out of the hand of your enemies.\end{Resp}
\switchcolumn*

\LA
\LessonHead{Lectio II}
\switchcolumn
\EN
\LessonHead{Lesson II}
\switchcolumn*

\LA
\Cite{4 Reg 1:4-6}
\switchcolumn
\EN
\Cite{2 Kgs 1:4-6}
\switchcolumn*

\LA
\noindent Et ábiit Elías. Reversíque sunt núntii ad Ochozíam. Qui dixit eis: Quare revérsi estis? At illi respondérunt ei: vir occúrrit nobis et dixit ad nos: Ite et revertímini ad regem qui misit vos, et dicétis ei: Hæc dicit Dóminus: Numquid quia non erat Deus in Israël, mittis ut consulátur Beélzebub deus Accaron? Idcírco de léctulo, super quem ascendísti, non descéndes, sed morte moriéris.\par
\switchcolumn
\EN
\noindent And Elias went away. And the messengers turned back to Ochozias. And he said to them: Why are you come back? But they answered him: A man met us, and said to us: Go, and return to the king, that sent you, and you shall say to him: Thus saith the Lord: Is it because there was no God in Israel that thou sendest to Beelzebub the god of Accaron? Therefore thou shalt not come down from the bed, on which thou art gone up, but then shalt surely die.\par
\switchcolumn*

\LA
\begin{Resp}\Rsign Deus ómnium exaudítor est: ipse misit Angelum suum, et tulit me de óvibus patris mei: \Star{} Et unxit me unctióne misericórdiæ suæ.\end{Resp}
\begin{Resp}\Vsign Dóminus, qui erípuit me de ore leónis, et de manu béstiæ liberávit me.\end{Resp}
\begin{Resp}\Rsign Et unxit me unctióne misericórdiæ suæ.\end{Resp}
\switchcolumn
\EN
\begin{Resp}\Rsign God, Which heareth all, even He sent His Angel, and took me from keeping my father's sheep, and \Star{} Anointed me with the oil of His mercy.\end{Resp}
\begin{Resp}\Vsign The Lord That delivered me out of the mouth of the lion, and out of the paw of the bear\end{Resp}
\begin{Resp}\Rsign And anointed me with the oil of His mercy.\end{Resp}
\switchcolumn*

\LA
\LessonHead{Lectio III}
\switchcolumn
\EN
\LessonHead{Lesson III}
\switchcolumn*

\LA
\Cite{4 Reg 1:7-10}
\switchcolumn
\EN
\Cite{2 Kgs 1:7-10}
\switchcolumn*

\LA
\noindent Qui dixit eis: Cujus figúræ et hábitus est vir ille, qui occúrrit vobis et locútus est verba hæc? At illi dixérunt: Vir pilósus et zona pellícea accínctus rénibus. Qui ait: Elías Thesbítes est. Misítque ad eum quinquagenárium príncipem et quinquagínta qui erant sub eo. Qui ascéndit ad eum sedentíque in vértice montis ait: Homo Dei, rex præcépit ut descéndas. Respondénsque Elías dixit quinquagenário: Si homo Dei sum, descéndat ignis de cælo, et dévoret te et quinquagínta tuos. Descéndit ítaque ignis de cælo, et devorávit eum et quinquagínta qui erant cum eo.\par
\switchcolumn
\EN
\noindent And he said to them: What manner of man was he who met you, and spoke these words? But they said: A hairy man with a girdle of leather about his loins. And he said: It is Elias the Thesbite. And he sent to him a captain of fifty, and the fifty men that were under him. And he went up to him, and as he was sitting on the top of a hill, said to him: Man of God, the king hath commanded that thou come down. And Elias answering, said to the captain of fifty: If I be a man of God, let fire come down from heaven, and consume thee, and thy fifty. And there came down fire from heaven, and consumed him, and the fifty that were with him.\par
\switchcolumn*

\LA
\begin{Resp}\Rsign Dóminus, qui erípuit me de ore leónis, et de manu béstiæ liberávit me, \Star{} Ipse me erípiet de mánibus inimicórum meórum.\end{Resp}
\begin{Resp}\Vsign Misit Deus misericórdiam suam et veritátem suam: ánimam meam erípuit de médio catulórum leónum.\end{Resp}
\begin{Resp}\Rsign Ipse me erípiet de mánibus inimicórum meórum.\end{Resp}
\begin{Resp}\Vsign Glória Patri, et Fílio, \Star{} et Spirítui Sancto.\end{Resp}
\begin{Resp}\Rsign Ipse me erípiet de mánibus inimicórum meórum.\end{Resp}
\switchcolumn
\EN
\begin{Resp}\Rsign The Lord That delivered me out of the mouth of the lion, and out of the paw of the bear \Star{} He will deliver me out of the hand of mine enemies.\end{Resp}
\begin{Resp}\Vsign God hath sent forth His mercy and His truth, and delivered my soul from among the lion's whelps.\end{Resp}
\begin{Resp}\Rsign He will deliver me out of the hand of mine enemies.\end{Resp}
\begin{Resp}\Vsign Glory be to the Father, and to the Son, \Star{} and to the Holy Ghost.\end{Resp}
\begin{Resp}\Rsign He will deliver me out of the hand of mine enemies.\end{Resp}
\switchcolumn*

\LA
\LessonHead{Lectio IV}
\switchcolumn
\EN
\LessonHead{Lesson IV}
\switchcolumn*

\LA
\LessonTitle{Sermo sancti Augustíni Epíscopi}
\switchcolumn
\EN
\LessonTitle{From the sermons of St. Augustine, Bishop of Hippo}
\switchcolumn*

\LA
\Source{Sermo 201 de Tempore}
\switchcolumn
\EN
\Source{201st for the Season.}
\switchcolumn*

\LA
\noindent In lectiónibus quæ nobis diébus istis recitántur, fratres caríssimi, frequénter admónui, ut non sequámur lítteram occidéntem, et vivificántem spíritum relinquámur. Sic enim Apóstolus ait: Líttera enim occídit, spíritus vivíficat. Si enim hoc tantum vólumus intellégere quod sonat in líttera, aut parvam, aut prope nullam ædificatiónem de divínis lectiónibus capiémus. Illa enim ómnia quæ recitántur, typus erant et imágo futurórum. In Judǽis enim figuráta, in nobis, grátia Dei donánte, compléta sunt.\par
\switchcolumn
\EN
\noindent Dearly beloved brethren, in the Lessons which are now being read to us day by day, I have often warned you that we must not follow the deathful letter, to the abandonment of the quickening spirit. For it is thus that the Apostle saith: “The letter killeth, but the spirit giveth life.” (2 Cor. iii. 6.) If we will understand only the plain meaning of the letter, we shall get little or no edification from our readings in the Divine Scriptures. All those things whereof we hear were types and images of things.\par
\switchcolumn*

\LA
\begin{Resp}\Rsign Percússit Saul mille, et David decem míllia: \Star{} Quia manus Dómini erat cum illo: percússit Philisthǽum, et ábstulit oppróbrium ex Israël.\end{Resp}
\begin{Resp}\Vsign Nonne iste est David, de quo canébant in choro, dicéntes: Saul percússit mille, et David decem míllia?\end{Resp}
\begin{Resp}\Rsign Quia manus Dómini erat cum illo: percússit Philisthǽum, et ábstulit oppróbrium ex Israël.\end{Resp}
\switchcolumn
\EN
\begin{Resp}\Rsign Saul hath slain his thousands, and David his ten thousands. \Star{} Because the hand of the Lord was with him, he smote the Philistine, and took away the reproach from Israel.\end{Resp}
\begin{Resp}\Vsign Is not this David? Did they not sing one to another of him in dances, saying: Saul hath slain his thousands, and David his ten thousands?\end{Resp}
\begin{Resp}\Rsign Because the hand of the Lord was with him, he smote the Philistine, and took away the reproach from Israel.\end{Resp}
\switchcolumn*

\LA
\LessonHead{Lectio V}
\switchcolumn
\EN
\LessonHead{Lesson V}
\switchcolumn*

\LA
\noindent Beátus enim Elías typum hábuit Dómini Salvatóris. Sicut enim Elías a Judǽis persecutiónem passus est; ita et verus Elías Dóminus noster ab ipsis Judǽis reprobátus est et contémptus. Elías relíquit gentem suam; et Christus deséruit synagógam. Elías ábiit in desértum; et Christus venit in mundum. Elías in desérto corvis ministrántibus pascebátur; et Christus in desérto mundi hujus Géntium fide refícitur.\par
\switchcolumn
\EN
\noindent The Blessed Elias was a type of the Lord our Saviour. Just as Elias was rejected by the Jews, so was the true Elias, even our Lord, rejected and despised by the same Jews. Elias went away out of his own country, and Christ hath left the synagogue Elias went into the desert, and Christ hath come into the world. Elias, when he was in the desert, was fed by ravens, and Christ in the desert of this world is comforted by the faith of the Gentiles.\par
\switchcolumn*

\LA
\begin{Resp}\Rsign Montes Gélboë, nec ros nec plúvia véniant super vos, \Star{} Ubi cecidérunt fortes Israël.\end{Resp}
\begin{Resp}\Vsign Omnes montes, qui estis in circúitu ejus, vísitet Dóminus: a Gélboë autem tránseat.\end{Resp}
\begin{Resp}\Rsign Ubi cecidérunt fortes Israël.\end{Resp}
\switchcolumn
\EN
\begin{Resp}\Rsign Ye mountains of Gilboa, let there be no dew, neither let there be rain upon you \Star{} For there are the mighty of Israel fallen.\end{Resp}
\begin{Resp}\Vsign All ye mountains that stand round about, the Lord look upon you but let Him pass by Gilboa\end{Resp}
\begin{Resp}\Rsign For there are the mighty of Israel fallen.\end{Resp}
\switchcolumn*

\LA
\LessonHead{Lectio VI}
\switchcolumn
\EN
\LessonHead{Lesson VI}
\switchcolumn*

\LA
\noindent Corvi enim illi qui beáto Elíæ, jubénte Dómino, ministrábant, Géntium pópulum figurábant. Proptérea et de géntium Ecclésia dícitur: Nigra sum et formósa, fília Jerúsalem. Unde est Ecclésia nigra et formósa? Nigra per natúram, formósa per grátiam. Unde nigra? Ecce in iniquitátibus concéptus sum, et in delíctis péperit me mater mea. Unde formósa? Aspérges me hyssópo, et mundábor: lavábis me, et super nivem dealbábor.\par
\switchcolumn
\EN
\noindent For the ravens which, at the command of the Lord, ministered unto Elias, were a type of the flock of Gentiles. Wherefore also it is said for the Gentile Church: “I am black, but comely, O daughters of Jerusalem!” (Cant. i. 4.) Why is the Church black but comely She is black by nature, but comely by grace. Why is she black by nature? Because she must needs own: “Behold, I was shapen in iniquity, and in sin did my mother conceive me.” (Ps. 1. 7.) Why is she comely by? Sprinkle me with hyssop, and I shall be clean wash me, and I shall be whiter than snow.\par
\switchcolumn*

\LA
\begin{Resp}\Rsign Ego te tuli de domo patris tui, dicit Dóminus, et pósui te páscere gregem pópuli mei: \Star{} Et fui tecum in ómnibus ubicúmque ambulásti, firmans regnum tuum in ætérnum.\end{Resp}
\begin{Resp}\Vsign Fecíque tibi nomen grande, juxta nomen magnórum, qui sunt in terra: et réquiem dedi tibi ab ómnibus inimícis tuis.\end{Resp}
\begin{Resp}\Rsign Et fui tecum in ómnibus ubicúmque ambulásti, firmans regnum tuum in ætérnum.\end{Resp}
\begin{Resp}\Vsign Glória Patri, et Fílio, \Star{} et Spirítui Sancto.\end{Resp}
\begin{Resp}\Rsign Et fui tecum in ómnibus ubicúmque ambulásti, firmans regnum tuum in ætérnum.\end{Resp}
\switchcolumn
\EN
\begin{Resp}\Rsign Thus saith the Lord: I took thee out of thy father's house, and appointed thee to be ruler over My people, over Israel. \Star{} And I was with thee whithersoever thou wentest, to establish thy kingdom for ever.\end{Resp}
\begin{Resp}\Vsign And I have made thee a great name, like unto the name of the great men that are in the earth, and have caused thee to rest from all thine enemies.\end{Resp}
\begin{Resp}\Rsign And I was with thee whithersoever thou wentest, to establish thy kingdom for ever.\end{Resp}
\begin{Resp}\Vsign Glory be to the Father, and to the Son, \Star{} and to the Holy Ghost.\end{Resp}
\begin{Resp}\Rsign And I was with thee whithersoever thou wentest, to establish thy kingdom for ever.\end{Resp}
\switchcolumn*

\LA
\LessonHead{Lectio VII}
\switchcolumn
\EN
\LessonHead{Lesson VII}
\switchcolumn*

\LA
\LessonTitle{Léctio sancti Evangélii secúndum Lucam}
\switchcolumn
\EN
\LessonTitle{From the Holy Gospel according to Luke}
\switchcolumn*

\LA
\Cite{Luc 19:41-47}
\switchcolumn
\EN
\Cite{Luke 19:41-17}
\switchcolumn*

\LA
\noindent In illo témpore: Cum appropinquáret Jesus Jerúsalem, videns civitátem flevit super illam dicens: Quia si cognovísses et tu, et quidem in hac die tua, quæ ad pacem tibi! Nunc autem abscóndita sunt ab óculis tuis. Et réliqua.\par
\switchcolumn
\EN
\noindent At that time: When Jesus was come near to Jerusalem, He beheld the city, and wept over it, saying: If thou hadst known, even thou, at least in this thy day, the things which belong unto thy peace! but now they are hid from thine eyes. And so on.\par
\switchcolumn*

\LA
\LessonTitle{Homilía sancti Gregórii Papæ}
\switchcolumn
\EN
\LessonTitle{Homily by Pope St. Gregory the Great.}
\switchcolumn*

\LA
\Source{Homilia 39 in Evangelia}
\switchcolumn
\EN
\Source{39 on the Gospels.}
\switchcolumn*

\LA
\noindent Quod a flente Dómino illa Jerosolymórum subvérsio descríbitur, quæ a Vespasiáno et Tito Románis princípibus facta est, nullus qui históriam eversiónis ejúsdem legit, ignórat. Románi étenim príncipes denuntiántur, cum dícitur: Quia vénient dies in te, et circúmdabunt te inimíci tui vallo. Hoc quoque quod ádditur: Non relínquent in te lápidem super lápidem: étiam jam ipsa ejúsdem civitátis transmigrátio testátur; quia dum nunc in eo loco constrúcta est, ubi extra portam fúerat Dóminus crucifíxus, prior illa Jerúsalem, ut dícitur, fúnditus est evérsa.\par
\switchcolumn
\EN
\noindent No man that hath read the history of the destruction of Jerusalem by the Roman Princes Vespasian and Titus, can be ignorant that it was of that destruction that the Lord spoke when He wept over the ruin of the city. It is these Princes that are pointed at where it is said “For the days shall come upon thee that thine enemies shall cast a trench about thee.” The truth of what followeth: “They shall not leave in thee one stone upon another” is even now fulfilled in the change of site of the city, which hath been re-built round about that place without the gates, where the Lord was crucified, while the ancient city hath been, as I am told, rooted up from the very foundations.\par
\switchcolumn*

\LA
\begin{Resp}\Rsign Dómine, Pater et Deus vitæ meæ, ne derelínquas me in cogitátu malígno: extolléntiam oculórum meórum ne déderis mihi, et desidérium malígnum avérte a me, Dómine; aufer a me concupiscéntiam, \Star{} Et ánimo irreverénti et infruníto ne tradas me, Dómine.\end{Resp}
\begin{Resp}\Vsign Ne derelínquas me, Dómine, ne accréscant ignorántiæ meæ, nec multiplicéntur delícta mea.\end{Resp}
\begin{Resp}\Rsign Et ánimo irreverénti et infruníto ne tradas me, Dómine.\end{Resp}
\switchcolumn
\EN
\begin{Resp}\Rsign O Lord, Father and God of my life, leave me not to evil counsels; give me not a proud look, but turn away from me an haughty mind, O Lord turn away from me concupiscence, \Star{} And give me not over unto an impudent and froward mind, O Lord!\end{Resp}
\begin{Resp}\Vsign Leave me not, O Lord, lest mine ignorance increase, and my sins abound.\end{Resp}
\begin{Resp}\Rsign And give me not over unto an impudent and froward mind, O Lord.\end{Resp}
\switchcolumn*

\LA
\LessonHead{Lectio VIII}
\switchcolumn
\EN
\LessonHead{Lesson VIII}
\switchcolumn*

\LA
\noindent Cui ex qua culpa eversiónis suæ pœna fúerit illáta, subjúngitur: Eo quod non cognóveris tempus visitatiónis tuæ. Creátor quippe hóminum per incarnatiónis suæ mystérium hanc visitáre dignátus est; sed ipsa timóris et amóris ílius recordáta non est. Unde étiam per prophétam in increpatióne cordis humáni aves cæli ad testimónium deducúntur, dum dícitur: Milvus in cælo cognóvit tempus suum, turtur et hirúndo et cicónia custodiérunt tempus advéntus sui; pópulus autem meus non cognóvit judícium Dómini.\par
\switchcolumn
\EN
\noindent What the sin of Jerusalem was which brought upon her the punishment of this destruction, we find written after: “Because thou knewest not the time of thy visitation.” The Maker of men, through the mystery of His Incarnation, was pleased to visit her, but she remembered not to fear and to love Him. Hence also the Prophet Jeremiah, rebuking the hardness of man's heart, calleth the birds of the air to testify against it, saying “The stork in the heaven knoweth her appointed time and the turtle, and the swallow, and the crane, observe the time of their coming but my people know not the judgment of the Lord.” (viii. 7.)\par
\switchcolumn*

\LA
\begin{Resp}\Rsign Duo Séraphim clamábant alter ad álterum: \Star{} Sanctus, sanctus, sanctus Dóminus Deus Sábaoth: * Plena est omnis terra glória ejus.\end{Resp}
\begin{Resp}\Vsign Tres sunt qui testimónium dant in cælo: Pater, Verbum, et Spíritus Sanctus: et hi tres unum sunt.\end{Resp}
\begin{Resp}\Rsign Sanctus, sanctus, sanctus Dóminus Deus Sábaoth.\end{Resp}
\begin{Resp}\Vsign Glória Patri, et Fílio, \Star{} et Spirítui Sancto.\end{Resp}
\begin{Resp}\Rsign Plena est omnis terra glória ejus.\end{Resp}
\switchcolumn
\EN
\begin{Resp}\Rsign One Seraph cried unto another: \Star{} Holy, Holy, Holy is the Lord God of hosts; the whole earth is full of His glory.\end{Resp}
\begin{Resp}\Vsign There are Three That bear record in heaven: the Father, the Word, and the Holy Ghost, and these Three are One.\end{Resp}
\begin{Resp}\Rsign Holy, Holy, Holy is the Lord God of hosts.\end{Resp}
\begin{Resp}\Vsign Glory be to the Father, and to the Son, \Star{} and to the Holy Ghost.\end{Resp}
\begin{Resp}\Rsign The whole earth is full of His glory.\end{Resp}
\switchcolumn*

\LA
\LessonHead{Lectio IX}
\switchcolumn
\EN
\LessonHead{Lesson IX}
\switchcolumn*

\LA
\noindent Flevit étenim prius Redémptor ruínam pérfidæ civitátis, quam ipsa sibi cívitas non cognoscébat esse ventúram. Cui a flente Dómino recte dícitur: Quia si cognovísses, et tu; subáudi, fleres; quæ modo, quæ nescis quod ímminet, exsúltas. Unde et súbditur: Et quidem in hac die tua, quæ ad pacem tibi. Cum enim carnis se voluptátibus daret, et ventúra mala non prospíceret; in die sua, quæ ad pacem esse ei póterant, habébat.\par
\switchcolumn
\EN
\noindent The Saviour wept over the ruin of the unfaithful city, while she herself as yet knew not that it was coming. If thou hadst known, said He, even thou and we may understand Him to have meant thou wouldest thyself have wept, in place of making merry as thou now dost, knowing not what hangeth over thee. And hence He saith farther: “at least in this thy day, the things which belong unto thy peace.” While she was giving herself up to fleshly pleasures, and casting no look ahead upon coming sorrows, she had still for a day in her power the things which might have brought unto her peace.\par
\switchcolumn*

\LA
\TeDeum{Te Deum}
\switchcolumn
\EN
\TeDeum{Te Deum}
\switchcolumn*

\end{paracol}

\Office{Feria secunda infra Hebdomadam IX post Octavam Pentecostes}{Monday of the Ninth Week after Pentecost}{Feria}
\addcontentsline{toc}{subsection}{Feria secunda infra Hebdomadam IX post Octavam Pentecostes \textit{— Monday of the Ninth Week after Pentecost}}
\begin{paracol}{2}
\LA
\LessonHead{Lectio I}
\switchcolumn
\EN
\LessonHead{Lesson I}
\switchcolumn*

\LA
\LessonTitle{De libro quarto Regum}
\switchcolumn
\EN
\LessonTitle{Lesson from the second book of Kings}
\switchcolumn*

\LA
\Cite{4 Reg 2:5-7}
\switchcolumn
\EN
\Cite{2 Kgs 2:5-7}
\switchcolumn*

\LA
\noindent Accessérunt fílii prophetárum qui erant in Jéricho, ad Eliséum, et dixérunt ei: Numquid nosti quia Dóminus hódie tollet dóminum tuum a te? Et ait: Et ego novi: siléte. Dixit autem ei Elías: Sede hic, quia Dóminus misit me usque ad Jordánem. Qui ait: Vivit Dóminus, et vivit ánima tua, quia non derelínquam te. Iérunt ígitur ambo páriter. Et quinquagínta viri de fíliis prophetárum secúti sunt eos, qui et stetérunt e contra longe; illi autem ambo stabant super Jordánem.\par
\switchcolumn
\EN
\noindent The sons of the prophets that were at Jericho, came to Eliseus, and said to him: Dost thou know that this day the Lord will take away thy master from thee? And he said: I also know it: hold your peace. And Elias said to him: Stay here, because the Lord hath sent me as far as the Jordan. And he said: As the Lord liveth, and as thy soul liveth, I will not leave thee; and they two went on together, And fifty men of the sons of the prophets followed them, and stood in sight at a distance: but they two stood by the Jordan.\par
\switchcolumn*

\LA
\begin{Resp}\Rsign Recordáre, Dómine, testaménti tui, et dic Angelo percutiénti: Cesset jam manus tua, \Star{} Ut non desolétur terra, et ne perdas omnem ánimam vivam.\end{Resp}
\begin{Resp}\Vsign Ego sum qui peccávi, ego qui iníque egi: isti qui oves sunt, quid fecérunt? Avertátur, óbsecro, furor tuus, Dómine, a pópulo tuo.\end{Resp}
\begin{Resp}\Rsign Ut non desolétur terra, et ne perdas omnem ánimam vivam.\end{Resp}
\switchcolumn
\EN
\begin{Resp}\Rsign Remember, O Lord, thy covenant, and say unto the destroying Angel: Stay now thine hand; \Star{} That the land be not utterly laid waste, and that thou destroy not every living soul.\end{Resp}
\begin{Resp}\Vsign Even I it is that have sinned, and done evil indeed; but these sheep, what have they done? Let thine anger, I pray thee, O Lord, be turned away from thy people.\end{Resp}
\begin{Resp}\Rsign That the land be not utterly laid waste, and that Thou destroy not every living soul.\end{Resp}
\switchcolumn*

\LA
\LessonHead{Lectio II}
\switchcolumn
\EN
\LessonHead{Lesson II}
\switchcolumn*

\LA
\Cite{4 Reg 2:8-10}
\switchcolumn
\EN
\Cite{2 Kgs 2:8-10}
\switchcolumn*

\LA
\noindent Tulítque Elías pállium suum et invólvit illud et percússit aquas, quæ divísæ sunt in utrámque partem, et transiérunt ambo per siccum. Cumque transíssent, Elías dixit ad Eliséum: Póstula quod vis ut fáciam tibi, ántequam tollar a te. Dixítque Eliséus: Obsecro ut fiat in me duplex spíritus tuus. Qui respóndit: Rem diffícilem postulásti; áttamen si víderis me, quando tollar a te, erit tibi quod petísti; si autem non víderis, non erit.\par
\switchcolumn
\EN
\noindent And Elias took his mantle and folded it together, and struck the waters, and they were divided hither and thither, and they both passed over on dry ground. And when they were gone over, Elias said to Eliseus: Ask what thou wilt have me to do for thee, before I be taken away from thee. And Eliseus said: I beseech thee that in me may be thy double spirit. And he answered: Thou hast asked a hard thing: nevertheless if thou see me when I am taken from thee, thou shalt have what thou hast asked: but if thou see me not, thou shalt not have it.\par
\switchcolumn*

\LA
\begin{Resp}\Rsign Exaudísti, Dómine, oratiónem servi tui, ut ædificárem templum nómini tuo: \Star{} Bénedic et sanctífica domum istam in sempitérnum, Deus Israël.\end{Resp}
\begin{Resp}\Vsign Dómine, qui custódis pactum cum servis tuis, qui ámbulant coram te in toto corde suo.\end{Resp}
\begin{Resp}\Rsign Bénedic et sanctífica domum istam in sempitérnum, Deus Israël.\end{Resp}
\switchcolumn
\EN
\begin{Resp}\Rsign O Lord, Thou hast hearkened unto the prayer of thy servant, that I might build a temple unto thy Name, \Star{} O God of Israel, bless Thou, and hallow this house for ever.\end{Resp}
\begin{Resp}\Vsign O Lord, Who keepest covenant with thy servants that walk before thee in all their heart.\end{Resp}
\begin{Resp}\Rsign O God of Israel, bless Thou, and hallow this house for ever.\end{Resp}
\switchcolumn*

\LA
\LessonHead{Lectio III}
\switchcolumn
\EN
\LessonHead{Lesson III}
\switchcolumn*

\LA
\Cite{4 Reg 2:11-13}
\switchcolumn
\EN
\Cite{2 Kgs 2:11-13}
\switchcolumn*

\LA
\noindent Cumque pérgerent, et incedéntes sermocinaréntur, ecce currus ígneus, et equi ígnei divisérunt utrúmque; et ascéndit Elías per túrbinem in cælum. Eliséus autem vidébat et clamábat: Pater mi, pater mi, currus Israël et auríga ejus. Et non vidit eum ámplius. Apprehendítque vestiménta sua et scidit illa in duas partes. Et levávit pállium Elíæ, quod cecíderat ei. Reversúsque stetit super ripam Jordánis.\par
\switchcolumn
\EN
\noindent And as they went on, walking and talking together, behold a fiery chariot, and fiery horses parted them both asunder: and Elias went up by a whirlwind into heaven. And Eliseus saw him, and cried: My father, my father, the chariot of Israel, and the driver thereof. And he saw him no more: and he took hold of his own garments, and rent them in two pieces. And he took up the mantle of Elias, that fell from him: and going back, he stood upon the bank of the Jordan,\par
\switchcolumn*

\LA
\begin{Resp}\Rsign Audi, Dómine, hymnum et oratiónem, quam servus tuus orat coram te hódie, ut sint óculi tui apérti, et aures tuæ inténtæ, \Star{} Super domum istam die ac nocte.\end{Resp}
\begin{Resp}\Vsign Réspice, Dómine, de sanctuário tuo, et de excélso cælórum habitáculo.\end{Resp}
\begin{Resp}\Rsign Super domum istam die ac nocte.\end{Resp}
\begin{Resp}\Vsign Glória Patri, et Fílio, \Star{} et Spirítui Sancto.\end{Resp}
\begin{Resp}\Rsign Super domum istam die ac nocte.\end{Resp}
\switchcolumn
\EN
\begin{Resp}\Rsign Hearken, O Lord, unto the cry and to the prayer which thy servant prayeth before thee today, that thine eyes may be open and thine ears attend; \Star{} Toward this house day and night.\end{Resp}
\begin{Resp}\Vsign Look down from thine high and holy place, O Lord, even from heaven thy dwelling.\end{Resp}
\begin{Resp}\Rsign Toward this house, day and night.\end{Resp}
\begin{Resp}\Vsign Glory be to the Father, and to the Son, \Star{} and to the Holy Ghost.\end{Resp}
\begin{Resp}\Rsign Toward this house, day and night.\end{Resp}
\switchcolumn*

\end{paracol}

\Office{Feria tertia infra Hebdomadam IX post Octavam Pentecostes}{Tuesday of the Ninth Week after Pentecost}{Feria}
\addcontentsline{toc}{subsection}{Feria tertia infra Hebdomadam IX post Octavam Pentecostes \textit{— Tuesday of the Ninth Week after Pentecost}}
\begin{paracol}{2}
\LA
\LessonHead{Lectio I}
\switchcolumn
\EN
\LessonHead{Lesson I}
\switchcolumn*

\LA
\LessonTitle{De libro quarto Regum}
\switchcolumn
\EN
\LessonTitle{Lesson from the second book of Kings}
\switchcolumn*

\LA
\Cite{4 Reg 3:6-9}
\switchcolumn
\EN
\Cite{2 Kgs 3:6-9}
\switchcolumn*

\LA
\noindent Egréssus est ígitur rex Joram in die illa de Samaría, et recénsuit univérsum Israël. Misítque ad Jósaphat regem Juda, dicens: Rex Moab recéssit a me; veni mecum contra eum ad prǽlium. Qui respóndit: Ascéndam: qui meus est, tuus est; pópulus meus, pópulus tuus, et equi mei, equi tui. Dixítque: Per quam viam ascendémus? At ille respóndit: Per desértum Idumǽæ. Perrexérunt ígitur rex Israël, et rex Juda, et rex Edom, et circuiérunt per viam septem diérum; nec erat aqua exercítui et juméntis quæ sequebántur eos.\par
\switchcolumn
\EN
\noindent And king Joram went out that day from Samaria, and mustered all Israel. And he sent to Josaphat king of Juda, saying: The king of Moab is revolted from me, come with me against him to battle. And he answered: I will come up: he that is mine, is thine: my people, thy people: and my horses, thy horses. And he said: Which way shall we go up? But he answered: By the desert of Edom. So the king of Israel, and the king of Juda, and the king of Edom went, and they fetched a compass of seven days' journey, and there was no water for the army, and for the beasts, that followed them.\par
\switchcolumn*

\LA
\begin{Resp}\Rsign Dómine, si convérsus fúerit pópulus tuus, et oráverit ad sanctuárium tuum: \Star{} Tu exáudies de cælo, Dómine, et líbera eos de mánibus inimicórum suórum.\end{Resp}
\begin{Resp}\Vsign Si peccáverit in te pópulus tuus, et convérsus égerit pœniténtiam, veniénsque oráverit in isto loco.\end{Resp}
\begin{Resp}\Rsign Tu exáudies de cælo, Dómine, et líbera eos de mánibus inimicórum suórum.\end{Resp}
\switchcolumn
\EN
\begin{Resp}\Rsign Lord, when thy people shall turn again to thee, and shall pray unto thee in this house; \Star{} Then hear Thou in heaven, O Lord, and deliver them out of the hand of their enemies.\end{Resp}
\begin{Resp}\Vsign If thy people sin against thee, and turn again, and repent, and come and pray unto thee in this house.\end{Resp}
\begin{Resp}\Rsign Then hear Thou in heaven, O Lord, and deliver them out of the hand of their enemies.\end{Resp}
\switchcolumn*

\LA
\LessonHead{Lectio II}
\switchcolumn
\EN
\LessonHead{Lesson II}
\switchcolumn*

\LA
\Cite{4 Reg 3:10-13}
\switchcolumn
\EN
\Cite{2 Kgs 3:10-13}
\switchcolumn*

\LA
\noindent Dixítque rex Israël: Heu! Heu! Heu! congregávit nos Dóminus tres reges, ut tráderet in manus Moab. Et ait Jósaphat: Estne hic prophéta Dómini, ut deprecémur Dóminum per eum? Et respóndit unus de servis regis Israël: Est hic Eliséus fílius Saphat, qui fundébat aquam super manus Elíæ. Et ait Jósaphat: Est apud eum sermo Dómini. Descendítque ad eum rex Israël, et Jósaphat rex Juda, et rex Edom. Dixit autem Eliséus ad regem Israël: Quid mihi et tibi est? Vade ad prophétas patris tui et matris tuæ.\par
\switchcolumn
\EN
\noindent And the king of Israel said: Alas, alas, alas, the Lord hath gathered us three kings together, to deliver us into the hands of Moab! And Josaphat said: Is there not here a prophet of the Lord, that we may beseech the Lord by him? And one of the servants of the king of Israel answered: Here is Eliseus the son of Saphat, who poured water on the hands of Elias. And Josaphat said: The word of the Lord is with him. And the king of Israel, and Josaphat king of Juda, and the king of Edom went down to him. And Eliseus said to the king of Israel: What have I to do with thee? go to the prophets of thy father, and thy mother.\par
\switchcolumn*

\LA
\begin{Resp}\Rsign Factum est, dum tólleret Dóminus Elíam per túrbinem in cælum, \Star{} Eliséus clamábat, dicens: Pater mi, pater mi, currus Israël, et auríga ejus.\end{Resp}
\begin{Resp}\Vsign Cumque pérgerent, et incedéntes sermocinaréntur, ecce currus ígneus et equi ígnei divisérunt utrúmque, et ascéndit Elías per túrbinem in cælum.\end{Resp}
\begin{Resp}\Rsign Eliséus clamábat, dicens: Pater mi, pater mi, currus Israël, et auríga ejus.\end{Resp}
\switchcolumn
\EN
\begin{Resp}\Rsign And it came to pass when the Lord would take up Elijah into heaven by a whirlwind \Star{} Elisha cried, and said: My father, my father, the chariot of Israel, and the horsemen thereof.\end{Resp}
\begin{Resp}\Vsign And as they still went on, and talked, behold, there appeared a chariot of fire, and horses of fire, and parted them both asunder, and Elijah went up by a whirlwind into heaven.\end{Resp}
\begin{Resp}\Rsign Elisha cried, and said: My father, my father, the chariot of Israel, and the horsemen thereof.\end{Resp}
\switchcolumn*

\LA
\LessonHead{Lectio III}
\switchcolumn
\EN
\LessonHead{Lesson III}
\switchcolumn*

\LA
\Cite{4 Reg 3:13-18}
\switchcolumn
\EN
\Cite{2 Kgs 3:13-18}
\switchcolumn*

\LA
\noindent Et ait illi rex Israël: Quare congregávit Dóminus tres reges hos, ut tráderet eos in manus Moab? Dixítque ad eum Eliséus: Vivit Dóminus exercítuum, in cujus conspéctu sto, quod si non vultum Jósaphat regis Judæ erubéscerem, non attendíssem quidem te, nec respexíssem. Nunc autem addúcite mihi psaltem. Cumque cáneret psaltes, facta est super eum manus Dómini, et ait: Hæc dicit Dóminus: Fácite álveum torréntis hujus fossas et fossas; hæc enim dicit Dóminus: Non vidébitis ventum, neque plúviam, et álveus iste replébitur aquis, et bibétis vos, et famíliæ vestræ et juménta vestra. Parúmque est hoc in conspéctu Dómini; ínsuper tradet étiam Moab in manus vestras.\par
\switchcolumn
\EN
\noindent And the king of Israel said to him: Why hath the Lord gathered together these three kings, to deliver them into the hands of Moab? And Eliseus said to him: As the Lord of hosts liveth, in whose sight I stand, if I did not reverence the face of Josaphat king of Juda, I would not have hearkened to thee, nor looked on thee. But now bring me hither a minstrel. And when the minstrel played, the hand of the Lord came upon him, and he said: Thus saith the Lord: Make the channel of this torrent full of ditches. For thus saith the Lord: You shall not see wind, nor rain: and yet this channel shall be filled with waters, and you shall drink, you and your families, and your beasts. And this is a small thing in the sight of the Lord: moreover he will deliver also Moab into your hands.\par
\switchcolumn*

\LA
\begin{Resp}\Rsign Ego te tuli de domo patris tui, dicit Dóminus, et pósui te páscere gregem pópuli mei: \Star{} Et fui tecum in ómnibus ubicúmque ambulásti, firmans regnum tuum in ætérnum.\end{Resp}
\begin{Resp}\Vsign Fecíque tibi nomen grande, juxta nomen magnórum, qui sunt in terra: et réquiem dedi tibi ab ómnibus inimícis tuis.\end{Resp}
\begin{Resp}\Rsign Et fui tecum in ómnibus ubicúmque ambulásti, firmans regnum tuum in ætérnum.\end{Resp}
\begin{Resp}\Vsign Glória Patri, et Fílio, \Star{} et Spirítui Sancto.\end{Resp}
\begin{Resp}\Rsign Et fui tecum in ómnibus ubicúmque ambulásti, firmans regnum tuum in ætérnum.\end{Resp}
\switchcolumn
\EN
\begin{Resp}\Rsign Thus saith the Lord: I took thee out of thy father's house, and appointed thee to be ruler over My people, over Israel. \Star{} And I was with thee whithersoever thou wentest, to establish thy kingdom for ever.\end{Resp}
\begin{Resp}\Vsign And I have made thee a great name, like unto the name of the great men that are in the earth and have caused thee to rest from all thine enemies.\end{Resp}
\begin{Resp}\Rsign And I was with thee whithersoever thou wentest, to establish thy kingdom for ever.\end{Resp}
\begin{Resp}\Vsign Glory be to the Father, and to the Son, \Star{} and to the Holy Ghost.\end{Resp}
\begin{Resp}\Rsign And I was with thee whithersoever thou wentest, to establish thy kingdom for ever.\end{Resp}
\switchcolumn*

\end{paracol}

\Office{Feria quarta infra Hebdomadam IX post Octavam Pentecostes}{Wednesday of the Ninth Week after Pentecost}{Feria}
\addcontentsline{toc}{subsection}{Feria quarta infra Hebdomadam IX post Octavam Pentecostes \textit{— Wednesday of the Ninth Week after Pentecost}}
\begin{paracol}{2}
\LA
\LessonHead{Lectio I}
\switchcolumn
\EN
\LessonHead{Lesson I}
\switchcolumn*

\LA
\LessonTitle{De libro quarto Regum}
\switchcolumn
\EN
\LessonTitle{Lesson from the second book of Kings}
\switchcolumn*

\LA
\Cite{4 Reg 4:1-4}
\switchcolumn
\EN
\Cite{2 Kgs 4:1-4}
\switchcolumn*

\LA
\noindent Múlier autem quædam de uxóribus prophetárum clamábat ad Eliséum, dicens: Servus tuus vir meus mórtuus est, et tu nosti quia servus tuus fuit timens Dóminum; et ecce créditor venit ut tollat duos fílios meos ad serviéndum sibi. Cui dixit Eliséus: Quid vis ut fáciam tibi? Dic mihi, quid habes in domo tua? At illa respóndit: Non hábeo ancílla tua quidquam in domo mea, nisi parum ólei, quo ungar. Cui ait: Vade, pete mútuo ab ómnibus vicínis tuis vasa vácua non pauca; et ingrédere et claude óstium tuum, cum intrínsecus fúeris tu et fílii tui, et mitte inde in ómnia vasa hæc; et, cum plena fúerint, tolles.\par
\switchcolumn
\EN
\noindent Now a certain woman of the wives of the prophets cried to Eliseus, saying: thy servant my husband is dead, and thou knowest that thy servant was one that feared God, and behold the creditor is come to take away my two sons to serve him. And Eliseus said to her: What wilt thou have me to do for thee? Tell me, what hast thou in thy house? And she answered: I thy handmaid have nothing in my house but a little oil, to anoint me. And he said to her: Go, borrow of all thy neighbours empty vessels not a few. And go in, and shut thy door, when thou art within, and thy sons: and pour out thereof into all those vessels: and when they are full take them away\par
\switchcolumn*

\LA
\begin{Resp}\Rsign Peccávi super númerum arénæ maris, et multiplicáta sunt peccáta mea: et non sum dignus vidére altitúdinem cæli præ multitúdine iniquitátis meæ: quóniam irritávi iram tuam, \Star{} Et malum coram te feci.\end{Resp}
\begin{Resp}\Vsign Quóniam iniquitátem meam ego cognósco: et delíctum meum contra me est semper, quia tibi soli peccávi.\end{Resp}
\begin{Resp}\Rsign Et malum coram te feci.\end{Resp}
\switchcolumn
\EN
\begin{Resp}\Rsign My sins are many, yea, they are more in number than the sands of the sea; I am not worthy to look up toward heaven because of the multitude of my iniquities; for I have provoked thee to anger; \Star{} And done evil in thy sight.\end{Resp}
\begin{Resp}\Vsign For I acknowledge my transgression, and my sin is ever before me, for against thee only have I sinned.\end{Resp}
\begin{Resp}\Rsign And done evil in thy sight.\end{Resp}
\switchcolumn*

\LA
\LessonHead{Lectio II}
\switchcolumn
\EN
\LessonHead{Lesson II}
\switchcolumn*

\LA
\Cite{4 Reg 4:5-10}
\switchcolumn
\EN
\Cite{2 Kgs 4:5-10}
\switchcolumn*

\LA
\noindent Ivit ítaque múlier et clausit óstium super se et super fílios suos; illi offerébant vasa, et illa infundébat. Cumque plena fuíssent vasa, dixit ad fílium suum: affer mihi adhuc vas. Et ille respóndit: Non hábeo. Stetítque óleum. Venit autem illa et indicávit hómini Dei. Et ille, Vade, inquit, vende óleum et redde creditóri tuo; tu autem et fílii tui vívite de réliquo. Facta est autem quædam dies, et transíbat Eliséus per Sunam. Erat autem ibi múlier magna, quæ ténuit eum et ut coméderet panem; cumque frequénter inde transíret, divertébat ad eam ut coméderet panem. Quæ dixit ad virum suum: Animadvérto quod vir Dei sanctus est iste, qui transit per nos frequénter: faciámus ergo ei cenáculum parvum, et ponámus ei in eo léctulum, et mensam, et sellam, et candelábrum, ut, cum vénerit ad nos, máneat ibi.\par
\switchcolumn
\EN
\noindent So the woman went, and shut the door upon her, and upon her sons: they brought her the vessels, and she poured in. And when the vessels were full, she said to her son: Bring me yet a vessel. And he answered: I have no more. And the oil stood. And she came, and told the man of God. And he said: Go, sell the oil, and pay thy creditor: and thou and thy sons live of the rest. And there was a day when Eliseus passed by Sunam: now there was a great woman there, who detained him to eat bread; and as he passed often that way, he turned into her house to eat bread. And she said to her husband: I perceive that this is a holy man of God, who often passeth by us. Let us therefore make him a little chamber, and put a little bed in it for him, and a table, and a stool, and a candlestick, that when he cometh to us, he may abide there.\par
\switchcolumn*

\LA
\begin{Resp}\Rsign Exaudísti, Dómine, oratiónem servi tui, ut ædificárem templum nómini tuo: \Star{} Bénedic et sanctífica domum istam in sempitérnum, Deus Israël.\end{Resp}
\begin{Resp}\Vsign Dómine, qui custódis pactum cum servis tuis, qui ámbulant coram te in toto corde suo.\end{Resp}
\begin{Resp}\Rsign Bénedic et sanctífica domum istam in sempitérnum, Deus Israël.\end{Resp}
\switchcolumn
\EN
\begin{Resp}\Rsign O Lord, Thou hast hearkened unto the prayer of thy servant, that I might build a temple unto thy Name, \Star{} O God of Israel, bless Thou, and hallow this house for ever.\end{Resp}
\begin{Resp}\Vsign O Lord, Who keepest covenant with thy servants that walk before thee in all their heart.\end{Resp}
\begin{Resp}\Rsign O God of Israel, bless Thou, and hallow this house for ever.\end{Resp}
\switchcolumn*

\LA
\LessonHead{Lectio III}
\switchcolumn
\EN
\LessonHead{Lesson III}
\switchcolumn*

\LA
\Cite{4 Reg 4:11-17}
\switchcolumn
\EN
\Cite{2 Kgs 4:11-17}
\switchcolumn*

\LA
\noindent Facta est ergo dies quædam, et véniens divértit in cenáculum et requiévit ibi. Dixítque ad Giézi púerum suum: Voca Sunamítidem istam. Qui, cum vocásset eam et illa stetísset coram eo, dixit ad púerum suum: Lóquere ad eam: Ecce, sédule in ómnibus ministrásti nobis, quid vis ut fáciam tibi? Numquid habes negótium, et vis ut loquar regi, sive príncipi milítiæ? Quæ respóndit: In médio pópuli mei hábito. Et ait: Quid ergo vult ut fáciam ei? Dixítque Giézi: Ne quæras, fílium enim non habet, et vir ejus senex est. Præcépit ítaque ut vocáret eam; quæ cum vocáta fuísset, et stetísset ante óstium, dixit ad eam: In témpore isto et in hac eádem hora, si vita comes fúerit, habébis in útero fílium. At illa respóndit: Noli, quæso, dómine mi, vir Dei, noli mentíri ancíllæ tuæ. Et concépit múlier, et péperit fílium in témpore, et in hora eádem, qua díxerat Eliséus.\par
\switchcolumn
\EN
\noindent Now there was a certain day when he came and turned in to the chamber, and rested there. And he said to Giezi his servant Call this Sunamitess. And when he had called her, and she stood before him, He said to his servant: Say to her Behold thou hast diligently served us in all things, what wilt thou have me to do for thee? hast thou any business, and wilt thou that I speak to the king, or to the general of the army? And she answered: I dwell in the midst of my own people. And he said: What will she then that I do for her? And Giezi said: Do not ask, for she hath no son, and her husband is old. Then he bid him call her: And when she was called, and stood before the door. He said to her: At this time, and this same hour, if life accompany, thou shalt have a son in thy womb. But she answered: Do not, I beseech thee, my lord, thou man of God, do not lie to thy handmaid. And the woman conceived, and brought forth a son in the time, and at the same hour, that Eliseus had said\par
\switchcolumn*

\LA
\begin{Resp}\Rsign Audi, Dómine, hymnum et oratiónem, quam servus tuus orat coram te hódie, ut sint óculi tui apérti, et aures tuæ inténtæ, \Star{} Super domum istam die ac nocte.\end{Resp}
\begin{Resp}\Vsign Réspice, Dómine, de sanctuário tuo, et de excélso cælórum habitáculo.\end{Resp}
\begin{Resp}\Rsign Super domum istam die ac nocte.\end{Resp}
\begin{Resp}\Vsign Glória Patri, et Fílio, \Star{} et Spirítui Sancto.\end{Resp}
\begin{Resp}\Rsign Super domum istam die ac nocte.\end{Resp}
\switchcolumn
\EN
\begin{Resp}\Rsign Hearken, O Lord, unto the cry and to the prayer which thy servant prayeth before thee today, that thine eyes may be open and thine ears attend; \Star{} Toward this house day and night.\end{Resp}
\begin{Resp}\Vsign Look down from thine high and holy place, O Lord, even from heaven thy dwelling.\end{Resp}
\begin{Resp}\Rsign Toward this house, day and night.\end{Resp}
\begin{Resp}\Vsign Glory be to the Father, and to the Son, \Star{} and to the Holy Ghost.\end{Resp}
\begin{Resp}\Rsign Toward this house, day and night.\end{Resp}
\switchcolumn*

\end{paracol}

\Office{Feria quinta infra Hebdomadam IX post Octavam Pentecostes}{Thursday of the Ninth Week after Pentecost}{Feria}
\addcontentsline{toc}{subsection}{Feria quinta infra Hebdomadam IX post Octavam Pentecostes \textit{— Thursday of the Ninth Week after Pentecost}}
\begin{paracol}{2}
\LA
\LessonHead{Lectio I}
\switchcolumn
\EN
\LessonHead{Lesson I}
\switchcolumn*

\LA
\LessonTitle{De libro quarto Regum}
\switchcolumn
\EN
\LessonTitle{Lesson from the second book of Kings}
\switchcolumn*

\LA
\Cite{4 Reg 6:24-27}
\switchcolumn
\EN
\Cite{2 Kgs 6:24-27}
\switchcolumn*

\LA
\noindent Congregávit Bénadad rex Sýriæ univérsum exércitum suum et ascéndit et obsidébat Samaríam. Factáque est fames magna in Samaría, et támdiu obséssa est, donec venumdarétur caput ásini octogínta argénteis, et quarta pars cabi stércoris columbárum quinque argénteis. Cumque rex Israël transíret per murum, múlier quædam exclamávit ad eum, dicens: Salva me, dómine mi rex. Qui ait: Non te salvat Dóminus, unde te possum salváre? de área, vel de torculári?\par
\switchcolumn
\EN
\noindent And it came to pass after these things, that Benadad king of Syria gathered together all his army, and went up, and besieged Samaria. And there was a great famine in Samaria: and so long did the siege continue, till the head of an ass was sold for fourscore pieces of silver, and the fourth part of a cabe of pigeon's dung, for five pieces of silver. And as the king of Israel was passing by the wall, a certain woman cried out to him, saying: Save me, my lord O king. And he said: If the Lord doth not save thee, how can I save thee? out of the barnfloor, or out of the winepress?\par
\switchcolumn*

\LA
\begin{Resp}\Rsign Præparáte corda vestra Dómino, et servíte illi soli: \Star{} Et liberábit vos de mánibus inimicórum vestrórum.\end{Resp}
\begin{Resp}\Vsign Convertímini ad eum in toto corde vestro, et auférte deos aliénos de médio vestri.\end{Resp}
\begin{Resp}\Rsign Et liberábit vos de mánibus inimicórum vestrórum.\end{Resp}
\switchcolumn
\EN
\begin{Resp}\Rsign Prepare your hearts unto the Lord, and serve Him Only \Star{} And He will deliver you out of the hand of your enemies.\end{Resp}
\begin{Resp}\Vsign Return unto Him with all your hearts, and put away the strange gods from among you.\end{Resp}
\begin{Resp}\Rsign And He will deliver you out of the hand of your enemies.\end{Resp}
\switchcolumn*

\LA
\LessonHead{Lectio II}
\switchcolumn
\EN
\LessonHead{Lesson II}
\switchcolumn*

\LA
\Cite{4 Reg 6:27-32}
\switchcolumn
\EN
\Cite{2 Kgs 6:27-32}
\switchcolumn*

\LA
\noindent Dixítque ad eam rex: Quid tibi vis? Quæ respóndit: Múlier ista dixit mihi: Da fílium tuum ut comedámus eum hódie, et fílium meum comedémus cras. Cóximus ergo fílium meum, et comédimus; dixíque ei die áltera: Da fílium tuum, ut comedámus eum; quæ abscóndit fílium suum. Quod cum audísset rex, scidit vestiménta sua, et transíbat per murum; vidítque omnis pópulus cilícium, quo vestítus erat ad carnem intrínsecus. Et ait rex: Hæc mihi fáciat Deus et hæc addat, si stéterit caput Eliséi fílii Saphat super ipsum hódie. Eliséus autem sedébat in domo sua, et senes sedébant cum eo.\par
\switchcolumn
\EN
\noindent And the king said to her: What aileth thee? And she answered: This woman said to me: Give thy son, that we may eat him today, and we will eat my son tomorrow. So we boiled my son, and ate him. And I said to her on the next day: Give thy son that we may eat him. And she hath hid her son. When the king heard this, he rent his garments, and passed by upon the wall. And all the people saw the haircloth which he wore within next to his flesh. And the king said: May God do so and so to me, and may he add more, if the head of Eliseus the son of Saphat shall stand on him this day. But Eliseus sat in his house, and the ancients sat with him.\par
\switchcolumn*

\LA
\begin{Resp}\Rsign Deus ómnium exaudítor est: ipse misit Angelum suum, et tulit me de óvibus patris mei: \Star{} Et unxit me unctióne misericórdiæ suæ.\end{Resp}
\begin{Resp}\Vsign Dóminus, qui erípuit me de ore leónis, et de manu béstiæ liberávit me.\end{Resp}
\begin{Resp}\Rsign Et unxit me unctióne misericórdiæ suæ.\end{Resp}
\switchcolumn
\EN
\begin{Resp}\Rsign God, Which heareth all, even He sent His Angel, and took me from keeping my father's sheep, and \Star{} Anointed me with the oil of His mercy.\end{Resp}
\begin{Resp}\Vsign The Lord That delivered me out of the mouth of the lion, and out of the paw of the bear\end{Resp}
\begin{Resp}\Rsign And anointed me with the oil of His mercy.\end{Resp}
\switchcolumn*

\LA
\LessonHead{Lectio III}
\switchcolumn
\EN
\LessonHead{Lesson III}
\switchcolumn*

\LA
\Cite{4 Reg 6:32-33; 7:1}
\switchcolumn
\EN
\Cite{2 Kgs 6:32-33; 7:1}
\switchcolumn*

\LA
\noindent Præmísit ítaque virum; et ántequam veníret núntius ille, dixit ad senes: Numquid scitis quod míserit fílius homicídæ hic, ut præcidátur caput meum? Vidéte ergo, cum vénerit núntius, cláudite óstium, et non sinátis eum introíre; ecce enim sónitus pedum dómini ejus post eum est. Adhuc illo loquénte eis, appáruit núntius qui veniébat ad eum. Et ait: Ecce, tantum malum a Dómino est; quid ámplius exspectábo a Dómino? Dixit autem Eliséus: Audíte verbum Dómini: hæc dicit Dóminus: In témpore hoc cras módius símilæ uno statére erit, et duo módii hórdei statére uno, in porta Samaríæ.\par
\switchcolumn
\EN
\noindent So he sent a man before: and before that messenger came, he said to the ancients: Do you know that this son of a murderer hath sent to cut off my head? Look then, when the messenger shall come, shut the door, and suffer him not to come in: for behold the sound of his master's feet is behind him. While he was yet speaking to them, the messenger appeared who was coming to him. And he said: Behold, so great an evil is from the Lord: what shall I look for more from the Lord?\par
\switchcolumn*

\LA
\begin{Resp}\Rsign Dóminus, qui erípuit me de ore leónis, et de manu béstiæ liberávit me, \Star{} Ipse me erípiet de mánibus inimicórum meórum.\end{Resp}
\begin{Resp}\Vsign Misit Deus misericórdiam suam et veritátem suam: ánimam meam erípuit de médio catulórum leónum.\end{Resp}
\begin{Resp}\Rsign Ipse me erípiet de mánibus inimicórum meórum.\end{Resp}
\begin{Resp}\Vsign Glória Patri, et Fílio, \Star{} et Spirítui Sancto.\end{Resp}
\begin{Resp}\Rsign Ipse me erípiet de mánibus inimicórum meórum.\end{Resp}
\switchcolumn
\EN
\begin{Resp}\Rsign The Lord That delivered me out of the mouth of the lion, and out of the paw of the bear \Star{} He will deliver me out of the hand of mine enemies.\end{Resp}
\begin{Resp}\Vsign God hath sent forth His mercy and His truth, and delivered my soul from among the lion's whelps.\end{Resp}
\begin{Resp}\Rsign He will deliver me out of the hand of mine enemies.\end{Resp}
\begin{Resp}\Vsign Glory be to the Father, and to the Son, \Star{} and to the Holy Ghost.\end{Resp}
\begin{Resp}\Rsign He will deliver me out of the hand of mine enemies.\end{Resp}
\switchcolumn*

\end{paracol}

\Office{Feria sexta infra Hebdomadam IX post Octavam Pentecostes}{Friday of the Ninth Week after Pentecost}{Feria}
\addcontentsline{toc}{subsection}{Feria sexta infra Hebdomadam IX post Octavam Pentecostes \textit{— Friday of the Ninth Week after Pentecost}}
\begin{paracol}{2}
\LA
\LessonHead{Lectio I}
\switchcolumn
\EN
\LessonHead{Lesson I}
\switchcolumn*

\LA
\LessonTitle{De libro quarto Regum}
\switchcolumn
\EN
\LessonTitle{Lesson from the second book of Kings}
\switchcolumn*

\LA
\Cite{4 Reg 8:1-3}
\switchcolumn
\EN
\Cite{2 Kgs 8:1-3}
\switchcolumn*

\LA
\noindent Eliséus autem locútus est ad mulíerem, cujus vívere fécerat fílium, dicens: Surge, vade tu et domus tua et peregrináre ubicúmque repéreris; vocávit enim Dóminus famem, et véniet super terram septem annis. Quæ surréxit, et fecit juxta verbum hóminis Dei; et vadens cum domo sua, peregrináta est in terra Philísthiim diébus multis. Cumque finíti essent anni septem, revérsa est múlier de terra Philísthiim: et egréssa est ut interpelláret regem pro domo sua, et pro agris suis.\par
\switchcolumn
\EN
\noindent And Eliseus spoke to the woman, whose son he had restored to life, saying: Arise, and go thou and thy household, and sojourn wheresoever thou canst find: for the Lord hath exiled a famine, and it shall come upon the land seven years. And she arose, and did according to the word of the man of God: and going with her household, she sojourned in the land of the Philistines many days. And when the seven years were ended, the woman returned out of the land of the Philistines, and she went forth to speak to the king for her house, and for her lands.\par
\switchcolumn*

\LA
\begin{Resp}\Rsign Percússit Saul mille, et David decem míllia: \Star{} Quia manus Dómini erat cum illo: percússit Philisthǽum, et ábstulit oppróbrium ex Israël.\end{Resp}
\begin{Resp}\Vsign Nonne iste est David, de quo canébant in choro, dicéntes: Saul percússit mille, et David decem míllia?\end{Resp}
\begin{Resp}\Rsign Quia manus Dómini erat cum illo: percússit Philisthǽum, et ábstulit oppróbrium ex Israël.\end{Resp}
\switchcolumn
\EN
\begin{Resp}\Rsign Saul hath slain his thousands, and David his ten thousands. \Star{} Because the hand of the Lord was with him, he smote the Philistine, and took away the reproach from Israel.\end{Resp}
\begin{Resp}\Vsign Is not this David? Did they not sing one to another of him in dances, saying: Saul hath slain his thousands, and David his ten thousands?\end{Resp}
\begin{Resp}\Rsign Because the hand of the Lord was with him, he smote the Philistine, and took away the reproach from Israel.\end{Resp}
\switchcolumn*

\LA
\LessonHead{Lectio II}
\switchcolumn
\EN
\LessonHead{Lesson II}
\switchcolumn*

\LA
\Cite{4 Reg 8:4-6}
\switchcolumn
\EN
\Cite{2 Kgs 8:4-6}
\switchcolumn*

\LA
\noindent Rex autem loquebátur cum Giézi púero viri Dei dicens: Narra mihi ómnia magnália, quæ fecit Eliséus. Cumque ille narráret regi quómodo mórtuum suscitásset, appáruit múlier, cujus vivificáverat fílium, clamans ad regem pro domo sua, et pro agris suis. Dixítque Giézi: Dómine mi rex, hæc est múlier, et hic est fílius ejus quem suscitávit Eliséus. Et interrogávit rex mulíerem, quæ narrávit ei, dedítque ei rex eunúchum unum dicens: Restítue ei ómnia quæ sua sunt, et univérsos réditus agrórum, a die qua relíquit terram usque ad præsens.\par
\switchcolumn
\EN
\noindent And the king talked with Giezi, the servant of the man of God, saying: Tell me all the great things that Eliseus hath done. And when he was telling the king how he had raised one dead to life, the woman appeared, whose son he had restored to life, crying to the king for her house, and her lands. And Giezi said: My lord O king, this is the woman, and this is her son, whom Eliseus raised to life. And the king asked the woman: and she told him. And the king appointed her an eunuch, saying: Restore her all that is hers, and all the revenues of the lands, from the day that she left the land, to this present.\par
\switchcolumn*

\LA
\begin{Resp}\Rsign Montes Gélboë, nec ros nec plúvia véniant super vos, \Star{} Ubi cecidérunt fortes Israël.\end{Resp}
\begin{Resp}\Vsign Omnes montes, qui estis in circúitu ejus, vísitet Dóminus: a Gélboë autem tránseat.\end{Resp}
\begin{Resp}\Rsign Ubi cecidérunt fortes Israël.\end{Resp}
\switchcolumn
\EN
\begin{Resp}\Rsign Ye mountains of Gilboa, let there be no dew, neither let there be rain upon you \Star{} For there are the mighty of Israel fallen.\end{Resp}
\begin{Resp}\Vsign All ye mountains that stand round about, the Lord look upon you but let Him pass by Gilboa\end{Resp}
\begin{Resp}\Rsign For there are the mighty of Israel fallen.\end{Resp}
\switchcolumn*

\LA
\LessonHead{Lectio III}
\switchcolumn
\EN
\LessonHead{Lesson III}
\switchcolumn*

\LA
\Cite{4 Reg 8:7-10}
\switchcolumn
\EN
\Cite{2 Kgs 8:7-10}
\switchcolumn*

\LA
\noindent Venit quoque Eliséus Damáscum, et Bénadad rex Sýriæ ægrotábat. Nuntiaverúntque ei dicéntes: Venit vir Dei huc. Et ait rex ad Házaël: Tolle tecum múnera et vade in occúrsum viri Dei et cónsule Dóminum per eum dicens: Si evádere pótero de infirmitáte mea hac? Ivit ígitur Házaël in occúrsum ejus. Cumque stetísset coram eo, ait: Fílius tuus Bénadad rex Sýriæ misit me ad te, dicens: Si sanári pótero de infirmitáte mea hac? Dixítque ei Eliséus: Vade, dic ei: Sanáberis. Porro osténdit mihi Dóminus quia morte moriétur.\par
\switchcolumn
\EN
\noindent Eliseus also came to Damascus, and Benadad king of Syria was sick: and they told him, saying: The man of God is come hither. And the king said to Hazael: Take with thee presents, and go to meet the man of God, and consult the Lord by him, saying: Can I recover of this my illness? And Hazael went to meet him, taking with him presents, and all the good things of Damascus, the burdens of forty camels. And when he stood before him, he said: thy son Benadad the king of Syria hath sent me to thee, saying: Can I recover of this my illness? And Eliseus said to him: Go tell him: Thou shalt recover: but the Lord hath shown me that he shall surely die.\par
\switchcolumn*

\LA
\begin{Resp}\Rsign Ego te tuli de domo patris tui, dicit Dóminus, et pósui te páscere gregem pópuli mei: \Star{} Et fui tecum in ómnibus ubicúmque ambulásti, firmans regnum tuum in ætérnum.\end{Resp}
\begin{Resp}\Vsign Fecíque tibi nomen grande, juxta nomen magnórum, qui sunt in terra: et réquiem dedi tibi ab ómnibus inimícis tuis.\end{Resp}
\begin{Resp}\Rsign Et fui tecum in ómnibus ubicúmque ambulásti, firmans regnum tuum in ætérnum.\end{Resp}
\begin{Resp}\Vsign Glória Patri, et Fílio, \Star{} et Spirítui Sancto.\end{Resp}
\begin{Resp}\Rsign Et fui tecum in ómnibus ubicúmque ambulásti, firmans regnum tuum in ætérnum.\end{Resp}
\switchcolumn
\EN
\begin{Resp}\Rsign Thus saith the Lord: I took thee out of thy father's house, and appointed thee to be ruler over My people, over Israel. \Star{} And I was with thee whithersoever thou wentest, to establish thy kingdom for ever.\end{Resp}
\begin{Resp}\Vsign And I have made thee a great name, like unto the name of the great men that are in the earth, and have caused thee to rest from all thine enemies.\end{Resp}
\begin{Resp}\Rsign And I was with thee whithersoever thou wentest, to establish thy kingdom for ever.\end{Resp}
\begin{Resp}\Vsign Glory be to the Father, and to the Son, \Star{} and to the Holy Ghost.\end{Resp}
\begin{Resp}\Rsign And I was with thee whithersoever thou wentest, to establish thy kingdom for ever.\end{Resp}
\switchcolumn*

\end{paracol}

\Office{Sabbato secunda infra Hebdomadam IX post Octavam Pentecostes}{Saturday of the Ninth Week after Pentecost}{Feria}
\addcontentsline{toc}{subsection}{Sabbato secunda infra Hebdomadam IX post Octavam Pentecostes \textit{— Saturday of the Ninth Week after Pentecost}}
\begin{paracol}{2}
\LA
\LessonHead{Lectio I}
\switchcolumn
\EN
\LessonHead{Lesson I}
\switchcolumn*

\LA
\LessonTitle{De libro quarto Regum}
\switchcolumn
\EN
\LessonTitle{Lesson from the second book of Kings}
\switchcolumn*

\LA
\Cite{4 Reg 9:1-5}
\switchcolumn
\EN
\Cite{2 Kgs 9:1-5}
\switchcolumn*

\LA
\noindent Eliséus autem prophétes vocávit unum de fíliis prophetárum et ait illi: Accínge lumbos tuos et tolle lentículam ólei hanc in manu tua et vade in Ramoth Gálaad. Cumque véneris illuc, vidébis Jehu fílium Jósaphat fílii Namsi; et ingréssus suscitábis eum de médio fratrum suórum, et introdúces in intérius cubículum. Tenénsque lentículam ólei fundes super caput ejus, et dices: Hæc dicit Dóminus: Unxi te regem super Israël. Aperiésque óstium, et fúgies, et non ibi subsístes. Abiit ergo adoléscens puer prophétæ in Ramoth Gálaad et ingréssus est illuc: ecce autem príncipes exércitus sedébant, et ait: Verbum mihi ad te, o princeps. Dixítque Jehu: Ad quem ex ómnibus nobis? At ille dixit: Ad te, o princeps.\par
\switchcolumn
\EN
\noindent And Eliseus the prophet called one of the sons of the prophets, slid said to him: Gird up thy loins, and take this little bottle of oil in thy hand, and go to Ramoth Galaad. And when thou art come thither, thou shalt see Jehu the son of Josaphat the son of Namsi: and going in thou shalt make him rise up from amongst his brethren, and carry him into an inner chamber. Then taking the little bottle of oil, thou shalt pour it on his head, and shalt say: Thus saith the Lord: I have anointed thee king over Israel. And thou shalt open the door and flee, and shalt not stay there. So the young man, the servant of the prophet, went awry to Ramoth Galaad, And went in thither: and behold the captains of the army were sitting: and he said: I have a word to thee, O prince. And Jehu said: Unto whom of us all? And he said: To thee, O prince.\par
\switchcolumn*

\LA
\begin{Resp}\Rsign Peccávi super númerum arénæ maris, et multiplicáta sunt peccáta mea: et non sum dignus vidére altitúdinem cæli præ multitúdine iniquitátis meæ: quóniam irritávi iram tuam, \Star{} Et malum coram te feci.\end{Resp}
\begin{Resp}\Vsign Quóniam iniquitátem meam ego cognósco: et delíctum meum contra me est semper, quia tibi soli peccávi.\end{Resp}
\begin{Resp}\Rsign Et malum coram te feci.\end{Resp}
\switchcolumn
\EN
\begin{Resp}\Rsign My sins are many, yea, they are more in number than the sands of the sea; I am not worthy to look up toward heaven because of the multitude of my iniquities; for I have provoked thee to anger; \Star{} And done evil in thy sight.\end{Resp}
\begin{Resp}\Vsign For I acknowledge my transgression, and my sin is ever before me, for against thee only have I sinned.\end{Resp}
\begin{Resp}\Rsign And done evil in thy sight.\end{Resp}
\switchcolumn*

\LA
\LessonHead{Lectio II}
\switchcolumn
\EN
\LessonHead{Lesson II}
\switchcolumn*

\LA
\Cite{4 Reg 9:6-10}
\switchcolumn
\EN
\Cite{2 Kgs 9:6-10}
\switchcolumn*

\LA
\noindent Et surréxit et ingréssus est cubículum; at ille fudit óleum super caput ejus et ait: Hæc dicit Dóminus Deus Israël: Unxi te regem super pópulum Dómini Israël, et percúties domum Achab dómini tui, et ulcíscar sánguinem servórum meórum prophetárum, et sánguinem ómnium servórum Dómini de manu Jézabel. Perdámque omnem domum Achab: et interfíciam de Achab mingéntem ad paríetem, et clausum et novíssimum in Israël. Et dabo domum Achab sicut domum Jeróboam fílii Nabat, et sicut domum Báasa fílii Ahía. Jézabel quoque cómedent canes in agro Jézrahel, nec erit qui sepéliat eam. Aperuítque óstium, et fugit.\par
\switchcolumn
\EN
\noindent And he arose, and went into the chamber: and he poured the oil upon his head, and said: Thus saith the Lord God of Israel: I have anointed thee king over Israel, the people of the Lord. And thou shalt cut off the house of Achab thy master, and I will revenge the blood of my servants the prophets, and the blood of all the servants of the Lord at the hand of Jezabel. And I will destroy all the house of Achab, and I will cut off from Achab him that pisseth against the well, and him that is shut up, and the meanest in Israel. And I will make the house of Achab like the house of Jeroboam the son of Nabat, and like the house of Baasa the son of Ahias. And the dogs shall eat Jezabel in the field of Jezrahel, and there shall be no one to bury her. And he opened the door and fled.\par
\switchcolumn*

\LA
\begin{Resp}\Rsign Exaudísti, Dómine, oratiónem servi tui, ut ædificárem templum nómini tuo: \Star{} Bénedic et sanctífica domum istam in sempitérnum, Deus Israël.\end{Resp}
\begin{Resp}\Vsign Dómine, qui custódis pactum cum servis tuis, qui ámbulant coram te in toto corde suo.\end{Resp}
\begin{Resp}\Rsign Bénedic et sanctífica domum istam in sempitérnum, Deus Israël.\end{Resp}
\switchcolumn
\EN
\begin{Resp}\Rsign O Lord, Thou hast hearkened unto the prayer of thy servant, that I might build a temple unto thy Name, \Star{} O God of Israel, bless Thou, and hallow this house for ever.\end{Resp}
\begin{Resp}\Vsign O Lord, Who keepest covenant with thy servants that walk before thee in all their heart.\end{Resp}
\begin{Resp}\Rsign O God of Israel, bless Thou, and hallow this house for ever.\end{Resp}
\switchcolumn*

\LA
\LessonHead{Lectio III}
\switchcolumn
\EN
\LessonHead{Lesson III}
\switchcolumn*

\LA
\Cite{4 Reg 9:11-13}
\switchcolumn
\EN
\Cite{2 Kgs 9:11-13}
\switchcolumn*

\LA
\noindent Jehu autem egréssus est ad servos dómini sui, qui dixérunt ei: Recte ne sunt ómnia? quid venit insánus iste ad te? Qui ait eis: Nostis hóminem, et quid locútus sit. At illi respondérunt: Falsum est, sed magis narra nobis. Qui ait eis: Hæc et hæc locútus est mihi, et ait: Hæc dicit Dóminus: Unxi te regem super Israël. Festinavérunt ítaque, et unusquísque tollens pállium suum posuérunt sub pédibus ejus in similitúdinem tribunális, et cecinérunt tuba, atque dixérunt: Regnávit Jehu.\par
\switchcolumn
\EN
\noindent Then Jehu went forth to the servants of his lord: and they said to him: Are all things well? why came this mad man to thee? And he said to them: You know the man, and what he said. But they answered: It is false, but rather do thou tell us. And he said to them: Thus and thus did he speak to me: and he said: Thus saith the Lord: I have anointed thee king over Israel. Then they made haste and taking every man his garment laid it under his feet, after the manner of a judgment seat, and they sounded the trumpet, and said: Jehu is king.\par
\switchcolumn*

\LA
\begin{Resp}\Rsign Audi, Dómine, hymnum et oratiónem, quam servus tuus orat coram te hódie, ut sint óculi tui apérti, et aures tuæ inténtæ, \Star{} Super domum istam die ac nocte.\end{Resp}
\begin{Resp}\Vsign Réspice, Dómine, de sanctuário tuo, et de excélso cælórum habitáculo.\end{Resp}
\begin{Resp}\Rsign Super domum istam die ac nocte.\end{Resp}
\begin{Resp}\Vsign Glória Patri, et Fílio, \Star{} et Spirítui Sancto.\end{Resp}
\begin{Resp}\Rsign Super domum istam die ac nocte.\end{Resp}
\switchcolumn
\EN
\begin{Resp}\Rsign Hearken, O Lord, unto the cry and to the prayer which thy servant prayeth before thee today, that thine eyes may be open and thine ears attend; \Star{} Toward this house day and night.\end{Resp}
\begin{Resp}\Vsign Look down from thine high and holy place, O Lord, even from heaven thy dwelling.\end{Resp}
\begin{Resp}\Rsign Toward this house, day and night.\end{Resp}
\begin{Resp}\Vsign Glory be to the Father, and to the Son, \Star{} and to the Holy Ghost.\end{Resp}
\begin{Resp}\Rsign Toward this house, day and night.\end{Resp}
\switchcolumn*

\end{paracol}

\Office{Dominica X Post Pentecosten}{Tenth Sunday after Pentecost}{Semiduplex}
\addcontentsline{toc}{subsection}{Dominica X Post Pentecosten \textit{— Tenth Sunday after Pentecost}}
\begin{paracol}{2}
\LA
\RefLine{Lectiones I–VI: de Scriptura occurrente.}
\switchcolumn
\EN
\RefLine{Lessons I–VI: from the Scripture of the day.}
\switchcolumn*

\LA
\LessonHead{Lectio I}
\switchcolumn
\EN
\LessonHead{Lesson I}
\switchcolumn*

\LA
\LessonTitle{De libro quarto Regum}
\switchcolumn
\EN
\LessonTitle{Lesson from the second book of Kings}
\switchcolumn*

\LA
\Cite{4 Reg 9:29-34}
\switchcolumn
\EN
\Cite{2 Kgs 9:29-34}
\switchcolumn*

\LA
\noindent Anno undécimo Joram fílii Achab regnávit Ochozías super Judam, venítque Jehu in Jézrahel. Porro Jézabel, intróitu ejus audíto, depínxit óculos suos stíbio et ornávit caput suum et respéxit per fenéstram ingrediéntem Jehu per portam et ait: Numquid pax potest esse Zambri, qui interfécit dóminum suum? Levavítque Jehu fáciem suam ad fenéstram et ait: Quæ est ista? Et inclinavérunt se ad eum duo vel tres eunúchi. At ille dixit eis: Præcipitáte eam deórsum; et præcipitavérunt eam, aspersúsque est sánguine páries, et equórum úngulæ conculcavérunt eam. Cumque introgréssus esset ut coméderet biberétque, ait: Ite et vidéte maledíctam illam et sepelíte eam, quia fília regis est.\par
\switchcolumn
\EN
\noindent In the eleventh year of Joram the son of Achab, Ochozias reigned over Juda, And Jehu came into Jezrahel. But Jezabel hearing of his coming in, painted her face with stibic stone, and adorned her head, and looked out of a window At Jehu coming in at the gate, and said: Can there be peace for Zambri, that hath killed his master? And Jehu lifted up his face to the window, and said: Who is this? And two or three eunuchs bowed down to him. And he said to them: Throw her down headlong: and they threw her down, and the wall was sprinkled with her blood, and the hoofs of the horses trod upon her. And when he was come in, to eat, and to drink, he said: Go, and see after that cursed woman, and bury her: because she is a king's daughter.\par
\switchcolumn*

\LA
\begin{Resp}\Rsign Præparáte corda vestra Dómino, et servíte illi soli: \Star{} Et liberábit vos de mánibus inimicórum vestrórum.\end{Resp}
\begin{Resp}\Vsign Convertímini ad eum in toto corde vestro, et auférte deos aliénos de médio vestri.\end{Resp}
\begin{Resp}\Rsign Et liberábit vos de mánibus inimicórum vestrórum.\end{Resp}
\switchcolumn
\EN
\begin{Resp}\Rsign Prepare your hearts unto the Lord, and serve Him Only \Star{} And He will deliver you out of the hand of your enemies.\end{Resp}
\begin{Resp}\Vsign Return unto Him with all your hearts, and put away the strange gods from among you.\end{Resp}
\begin{Resp}\Rsign And He will deliver you out of the hand of your enemies.\end{Resp}
\switchcolumn*

\LA
\LessonHead{Lectio II}
\switchcolumn
\EN
\LessonHead{Lesson II}
\switchcolumn*

\LA
\Cite{4 Reg 9:35-37; 10:1-3}
\switchcolumn
\EN
\Cite{2 Kgs 9:35-37; 10:1-3}
\switchcolumn*

\LA
\noindent Cumque issent ut sepelírent eam, non invenérunt nisi calváriam et pedes et summas manus. Reversíque nuntiavérunt ei. Et ait Jehu: Sermo Dómini est, quem locútus est per servum suum Elíam Thesbíten dicens: In agro Jézrahel cómedent canes carnes Jézabel, et erunt carnes Jézabel sicut stercus super fáciem terræ in agro Jézrahel, ita ut prætereúntes dicant: Hǽccine est illa Jézabel? Erant autem Achab septuagínta fílii in Samaría. Scripsit ergo Jehu lítteras et misit in Samaríam ad optimátes civitátis et ad majóres natu et ad nutrícios Achab, dicens: Statim ut accepéritis lítteras has, qui habétis fílios dómini vestri et currus et equos et civitátes firmas et arma, elígite meliórem et eum qui vobis placúerit de fíliis dómini vestri et eum pónite super sólium patris sui et pugnáte pro domo dómini vestri.\par
\switchcolumn
\EN
\noindent And when they went to bury her, they found nothing but the skull, and the feet, and the extremities of her hands. And coming back they told him. And Jehu said: It is the word of the Lord, which he spoke by his servant Elias the Thesbite, saying: In the field of Jezrahel the dogs shall eat the flesh of Jezabel, And the flesh of Jezabel shall be as dung upon the face of the earth in the field of Jezrahel, so that they who pass by shall say: Is this that same Jezabel? And Achab had seventy sons in Samaria: so Jehu wrote letters, and sent to Samaria, to the chief men of the city, and to the ancients, and to them that brought up Achab's children, saying: As soon as you receive these letters, ye that have your master's sons, and chariots, and horses, and fenced cities, and armour, Choose the best, and him that shall please you most of your master's sons, and set him on his father's throne, and fight for the house of your master.\par
\switchcolumn*

\LA
\begin{Resp}\Rsign Deus ómnium exaudítor est: ipse misit Angelum suum, et tulit me de óvibus patris mei: \Star{} Et unxit me unctióne misericórdiæ suæ.\end{Resp}
\begin{Resp}\Vsign Dóminus, qui erípuit me de ore leónis, et de manu béstiæ liberávit me.\end{Resp}
\begin{Resp}\Rsign Et unxit me unctióne misericórdiæ suæ.\end{Resp}
\switchcolumn
\EN
\begin{Resp}\Rsign God, Which heareth all, even He sent His Angel, and took me from keeping my father's sheep, and \Star{} Anointed me with the oil of His mercy.\end{Resp}
\begin{Resp}\Vsign The Lord That delivered me out of the mouth of the lion, and out of the paw of the bear\end{Resp}
\begin{Resp}\Rsign And anointed me with the oil of His mercy.\end{Resp}
\switchcolumn*

\LA
\LessonHead{Lectio III}
\switchcolumn
\EN
\LessonHead{Lesson III}
\switchcolumn*

\LA
\Cite{4 Reg 10:4-7}
\switchcolumn
\EN
\Cite{2 Kgs 10:4-7}
\switchcolumn*

\LA
\noindent Timuérunt illi veheménter et dixérunt: Ecce duo reges non potuérunt stare coram eo, et quómodo nos valébimus resístere? Misérunt ergo præpósiti domus et præfécti civitátis et majóres natu et nutrícii ad Jehu dicéntes: Servi tui sumus: quæcúmque jússeris faciémus, nec constituémus nobis regem: quæcúmque tibi placent, fac. Rescrípsit autem eis lítteras secúndo dicens: Si mei estis et obedítis mihi, tóllite cápita filiórum dómini vestri et veníte ad me hac eádem hora cras in Jézrahel. Porro fílii regis, septuagínta viri, apud optimátes civitátis nutriebántur. Cumque veníssent lítteræ ad eos, tulérunt fílios regis et occidérunt septuagínta viros et posuérunt cápita eórum in cóphinis et misérunt ad eum in Jézrahel.\par
\switchcolumn
\EN
\noindent But they were exceedingly afraid, and said: Behold two kings could not stand before him, and how shall we be able to resist? Therefore the overseers of the house, and the rulers of the city, and the ancients, and the tutors sent to Jehu, saying: We are thy servants, whatsoever thou shalt command us we will do, neither will we make us a king: do thou all that pleaseth thee. And he wrote letters the second time to them, saying: If you be mine, and will obey me, take the heads of the sons of your master, and come to me to Jezrahel by tomorrow this time. Now the king's sons, being seventy men, were brought up with the chief men of the city. And when the letters came to them, they took the king's sons, and slew seventy persons, and put their heads in baskets, and sent them to him to Jezrahel.\par
\switchcolumn*

\LA
\begin{Resp}\Rsign Dóminus, qui erípuit me de ore leónis, et de manu béstiæ liberávit me, \Star{} Ipse me erípiet de mánibus inimicórum meórum.\end{Resp}
\begin{Resp}\Vsign Misit Deus misericórdiam suam et veritátem suam: ánimam meam erípuit de médio catulórum leónum.\end{Resp}
\begin{Resp}\Rsign Ipse me erípiet de mánibus inimicórum meórum.\end{Resp}
\begin{Resp}\Vsign Glória Patri, et Fílio, \Star{} et Spirítui Sancto.\end{Resp}
\begin{Resp}\Rsign Ipse me erípiet de mánibus inimicórum meórum.\end{Resp}
\switchcolumn
\EN
\begin{Resp}\Rsign The Lord That delivered me out of the mouth of the lion, and out of the paw of the bear \Star{} He will deliver me out of the hand of mine enemies.\end{Resp}
\begin{Resp}\Vsign God hath sent forth His mercy and His truth, and delivered my soul from among the lion's whelps.\end{Resp}
\begin{Resp}\Rsign He will deliver me out of the hand of mine enemies.\end{Resp}
\begin{Resp}\Vsign Glory be to the Father, and to the Son, \Star{} and to the Holy Ghost.\end{Resp}
\begin{Resp}\Rsign He will deliver me out of the hand of mine enemies.\end{Resp}
\switchcolumn*

\LA
\LessonHead{Lectio IV}
\switchcolumn
\EN
\LessonHead{Lesson IV}
\switchcolumn*

\LA
\LessonTitle{Sermo sancti Joánnis Chrysóstomi}
\switchcolumn
\EN
\LessonTitle{From the Sermons of St. John Chrysostom, Patriarch of Constantinople.}
\switchcolumn*

\LA
\Source{Homilia 25 in Epistola ad Romanos}
\switchcolumn
\EN
\Source{Hom. 25 on the Epistle to the Romans.}
\switchcolumn*

\LA
\noindent Non putémus nos excusatiónem habitúros, si quando delictórum sócios invenérimus: nam istud, supplícium magis augébit. Quandóquidem et serpens magis punítus est quam múlier, quemádmodum et múlier plus quam vir. Et Jézabel majóres pœnas dedit, quam Achab víneæ raptor; ipsa quippe univérsum istud negótium texúerat, regíque lapsus occasiónem déderat. Igitur et tu quoque, si réliquis perditiónis causa fúeris, gravióra patiéris quam qui per te subvérsi sunt. Neque enim peccáre tantum in se perditiónis habet, quantum quod réliqui ad peccándum inducúntur.\par
\switchcolumn
\EN
\noindent Let us not dream that we are ourselves to be held less guilty, when we find that we have not been alone in sin. On the contrary, such fellowship addeth to our punishment. The serpent was more heavily punished than Eve, and Eve than Adam, and Jezebel suffered more than Ahab, who took the vineyard of Naboth. She it had been that planned the whole matter, and opened the way for her husband's crime. Even so thou also, who shalt have caused another's perdition, shalt suffer more grievously than shall they, whom thou hast ruined. Since, for a man to commit sin himself is less wicked than to lead others into sin.\par
\switchcolumn*

\LA
\begin{Resp}\Rsign Percússit Saul mille, et David decem míllia: \Star{} Quia manus Dómini erat cum illo: percússit Philisthǽum, et ábstulit oppróbrium ex Israël.\end{Resp}
\begin{Resp}\Vsign Nonne iste est David, de quo canébant in choro, dicéntes: Saul percússit mille, et David decem míllia?\end{Resp}
\begin{Resp}\Rsign Quia manus Dómini erat cum illo: percússit Philisthǽum, et ábstulit oppróbrium ex Israël.\end{Resp}
\switchcolumn
\EN
\begin{Resp}\Rsign Saul hath slain his thousands, and David his ten thousands. \Star{} Because the hand of the Lord was with him, he smote the Philistine, and took away the reproach from Israel.\end{Resp}
\begin{Resp}\Vsign Is not this David? Did they not sing one to another of him in dances, saying: Saul hath slain his thousands, and David his ten thousands?\end{Resp}
\begin{Resp}\Rsign Because the hand of the Lord was with him, he smote the Philistine, and took away the reproach from Israel.\end{Resp}
\switchcolumn*

\LA
\LessonHead{Lectio V}
\switchcolumn
\EN
\LessonHead{Lesson V}
\switchcolumn*

\LA
\noindent Itaque, si quando peccántes vidérimus, non solum non impellámus, sed et extrahámus ex ipso malítiæ bárathro, ne et aliénæ perditiónis pœnas demus. Recordémur quoque perpétuo terríbilis illíus tribunális, flúminis ígnei, vinculórum insolubílium, profundárum tenebrárum, stridórum déntium, venenosíque vermis. Sed dices: Benígnus est Deus. Ergo hæc ómnia verba sunt, et neque punítur dives ille Lázari contémptor, neque fátuæ vírgines a sponso rejiciúntur? Ergo qui Christum non pavérunt, in ignem diábolo præparátum non abíbunt? Ergo qui sórdibus est véstibus, non períbit, manus ac pedes vinctus? Qui centum denários a consérvo suo exégit, non tradétur tortóribus? Quod de mœchis dictum est, nimírum quod vermis eórum non moriétur, et ignis eórum non exstinguétur, verum non erit?\par
\switchcolumn
\EN
\noindent If, therefore, we should see others sinning, let us not only not help them, but let us do what in us lieth to draw them out of the bottomless pit of destruction, lest we should suffer as accomplices in their trespass. Let our memory never forget that right awful judgment-seat, the river of fire, the chains that can never be unlocked, the darkness that cannot be pierced, the sound of teeth gnashing, the deadly worm. But thou sayest God is good. Are then all these things but idle words? Is there no punishment for the rich man which giveth no heed to Lazarus? (Luke xvi. 20-26.) Doth the bridegroom open to the foolish virgins the door of the marriage-chamber? (Matth. xxv. 1- 12.) They that have denied to Christ the necessaries of life, are they not to depart from Him into everlasting fire, prepared for the devil and his angels? (Matth. xxv. 41-46.) The man that cometh in to the marriage-supper, not having a wedding garment, shall he, or shall he not, be bound hand and foot, and taken away, and cast into outer darkness? (Matth. xxii. 11-13.) The servant that hath no compassion on his fellow-servant, which oweth him an hundred pence, shall he, or shall he not, be delivered to the tormentors? (Matth. xviii. 23-35.) It is said, concerning such as commit adultery, that their worm dieth not and their fire is not quenched? (Mark ix. 44.) Is that not true?\par
\switchcolumn*

\LA
\begin{Resp}\Rsign Montes Gélboë, nec ros nec plúvia véniant super vos, \Star{} Ubi cecidérunt fortes Israël.\end{Resp}
\begin{Resp}\Vsign Omnes montes, qui estis in circúitu ejus, vísitet Dóminus: a Gélboë autem tránseat.\end{Resp}
\begin{Resp}\Rsign Ubi cecidérunt fortes Israël.\end{Resp}
\switchcolumn
\EN
\begin{Resp}\Rsign Ye mountains of Gilboa, let there be no dew, neither let there be rain upon you \Star{} For there are the mighty of Israel fallen.\end{Resp}
\begin{Resp}\Vsign All ye mountains that stand round about, the Lord look upon you but let Him pass by Gilboa\end{Resp}
\begin{Resp}\Rsign For there are the mighty of Israel fallen.\end{Resp}
\switchcolumn*

\LA
\LessonHead{Lectio VI}
\switchcolumn
\EN
\LessonHead{Lesson VI}
\switchcolumn*

\LA
\noindent Sed minátur ista tantúmmodo Deus? Utique ínquies. Et unde, dic quæso, tantam rem audes públice loqui, atque ex te ipso ferre senténtiam? Ego quippe et ex iis quæ dixit Deus, et ex iis quæ fecit, contrárium probáre pótero. Quod si propter futúra non credis, vel saltem propter ea quæ jam facta sunt, crede. Non enim certe minæ sunt et verba tantúmmodo, quæ facta sunt et in opus ipsum exiérunt. Quis ígitur totum orbem indúcto dilúvio stagnávit, ac grave illud naufrágium, omnimodámque géneris nostri perditiónem effécit? Quis deínde fúlmina illa, télaque flammántia super terram Sodomórum dimísit? Quis univérsum Ægýpti exércitum in mare demérsit? Quis synagógam Abíron combússit? Quis septuagínta illa míllia propter Davídis peccátum uno témporis moménto peste occídit? Nonne hæc ómnia et ália Deus illis íntulit?\par
\switchcolumn
\EN
\noindent But these are perhaps only threats on God's part, no doubt, quoth thou. I ask thee: How darest thou say such a thing out loud, and deliver this judgment from thine own imagining? For sooth, I can prove to thee, from the things which God hath done, that thou art wrong. If thou wilt not believe for things to come, at least believe for things past. Of them at least it cannot be said that they are nought but threats and mere words, for they have happened, and actually been realised in fact. Who was He which brought in a great flood, until the whole land was standing water, and our whole race perished, save eight persons? Who was He which rained upon Sodom brimstone and fire out of heaven? (Gen. xix. 24.) Who was He which overthrew all the host of Egypt in the Red Sea? (Ex. xiv. 27.) Who was He which sent out a fire and consumed them that were of the faction of Abiram? (Num. xvi. 35.) Who was He which sent a pestilence upon Israel, because David had sinned, and, from the morning even to the time appointed, there died of the people seventy thousand men? (2 Kgs xxiv.) Was it not God, and none other, Which brought upon them all these things, and more also?\par
\switchcolumn*

\LA
\begin{Resp}\Rsign Ego te tuli de domo patris tui, dicit Dóminus, et pósui te páscere gregem pópuli mei: \Star{} Et fui tecum in ómnibus ubicúmque ambulásti, firmans regnum tuum in ætérnum.\end{Resp}
\begin{Resp}\Vsign Fecíque tibi nomen grande, juxta nomen magnórum, qui sunt in terra: et réquiem dedi tibi ab ómnibus inimícis tuis.\end{Resp}
\begin{Resp}\Rsign Et fui tecum in ómnibus ubicúmque ambulásti, firmans regnum tuum in ætérnum.\end{Resp}
\begin{Resp}\Vsign Glória Patri, et Fílio, \Star{} et Spirítui Sancto.\end{Resp}
\begin{Resp}\Rsign Et fui tecum in ómnibus ubicúmque ambulásti, firmans regnum tuum in ætérnum.\end{Resp}
\switchcolumn
\EN
\begin{Resp}\Rsign Thus saith the Lord: I took thee out of thy father's house, and appointed thee to be ruler over My people, over Israel. \Star{} And I was with thee whithersoever thou wentest, to establish thy kingdom for ever.\end{Resp}
\begin{Resp}\Vsign And I have made thee a great name, like unto the name of the great men that are in the earth, and have caused thee to rest from all thine enemies.\end{Resp}
\begin{Resp}\Rsign And I was with thee whithersoever thou wentest, to establish thy kingdom for ever.\end{Resp}
\begin{Resp}\Vsign Glory be to the Father, and to the Son, \Star{} and to the Holy Ghost.\end{Resp}
\begin{Resp}\Rsign And I was with thee whithersoever thou wentest, to establish thy kingdom for ever.\end{Resp}
\switchcolumn*

\LA
\LessonHead{Lectio VII}
\switchcolumn
\EN
\LessonHead{Lesson VII}
\switchcolumn*

\LA
\LessonTitle{Léctio sancti Evangélii secúndum Lucam}
\switchcolumn
\EN
\LessonTitle{From the Holy Gospel according to Luke}
\switchcolumn*

\LA
\Cite{Luc 18:9-14}
\switchcolumn
\EN
\Cite{Luke 18:9-14}
\switchcolumn*

\LA
\noindent In illo témpore: Dixit Jesus ad quosdam, qui in se confidébant tamquam justi et aspernabántur céteros, parábolam istam: Duo hómines ascendérunt in templum ut orárent: unus pharisǽus, et alter publicánus. Et réliqua.\par
\switchcolumn
\EN
\noindent At that time: Jesus spoke this parable unto certain which trusted in themselves that they were righteous, and despised others: Two men went up into the Temple to pray, the one a Pharisee, and the other a publican. And so on.\par
\switchcolumn*

\LA
\LessonTitle{Homilía sancti Augustíni Epíscopi}
\switchcolumn
\EN
\LessonTitle{Homily by St. Augustine, Bishop of Hippo}
\switchcolumn*

\LA
\Source{Sermo 36 de Verbis Domini, circa medium}
\switchcolumn
\EN
\Source{Serm 36 of the Word of the Lord}
\switchcolumn*

\LA
\noindent Díceret saltem pharisǽus: Non sum sicut multi hómines. Quid est, céteri hómines, nisi omnes præter ipsum? Ego, inquit, justus sum, céteri peccatóres. Non sum sicut céteri hómines, injústi, raptóres, adúlteri. Et ecce tibi ex vicíno publicáno majóris tumóris occásio: Sicut, inquit, publicánus iste. Ego, inquit, solus sum: iste de céteris est. Non sum, inquit, talis, qualis iste, per justítias meas, quibus iníquus non sum.\par
\switchcolumn
\EN
\noindent The Pharisee might at least have said: “I am not as many men are.” But what meaneth “other men”? All other men except himself. “I,” said he, “am righteous; others are sinners.” “I am not as other men are, extortioners, unjust, adulterers,” and then he took occasion, from the neighborhood of the publican, to plume himself “or even,” quoth he, “as this publican.” “I am alone,” he thought, “that publican” is one of the others. Mine own righteousness maketh the gulf between me and the wicked, such as he is.\par
\switchcolumn*

\LA
\begin{Resp}\Rsign Dómine, Pater et Deus vitæ meæ, ne derelínquas me in cogitátu malígno: extolléntiam oculórum meórum ne déderis mihi, et desidérium malígnum avérte a me, Dómine; aufer a me concupiscéntiam, \Star{} Et ánimo irreverénti et infruníto ne tradas me, Dómine.\end{Resp}
\begin{Resp}\Vsign Ne derelínquas me, Dómine, ne accréscant ignorántiæ meæ, nec multiplicéntur delícta mea.\end{Resp}
\begin{Resp}\Rsign Et ánimo irreverénti et infruníto ne tradas me, Dómine.\end{Resp}
\switchcolumn
\EN
\begin{Resp}\Rsign O Lord, Father and God of my life, leave me not to evil counsels; give me not a proud look, but turn away from me an haughty mind, O Lord turn away from me concupiscence, \Star{} And give me not over unto an impudent and froward mind, O Lord!\end{Resp}
\begin{Resp}\Vsign Leave me not, O Lord, lest mine ignorance increase, and my sins abound.\end{Resp}
\begin{Resp}\Rsign And give me not over unto an impudent and froward mind, O Lord.\end{Resp}
\switchcolumn*

\LA
\LessonHead{Lectio VIII}
\switchcolumn
\EN
\LessonHead{Lesson VIII}
\switchcolumn*

\LA

\switchcolumn
\EN
\LessonTitle{“I fast twice in the week; I give tithes of all that I possess.” If we look in his prayer to find what he went to the Temple to pray to God for, we shall find nothing. He went up to pray, but his prayer was not a request of anything from God, but a glorification of himself. It was little enough not to pray to God, but he also glorified himself and despised his neighbour. But the publican stood afar off and yet drew nigh to God. Self-knowledge bade him keep at a distance, but his earnestness made him close. The publican stood afar off, but the Lord was at hand to hear him.}
\switchcolumn*

\LA
\noindent Jejúno bis in sábbato, décimas do ómnium quæ possídeo. Quid rogáverit Deum, quære in verbis ejus, nihil invénies. Ascéndit oráre: nóluit Deum rogáre, sed se laudáre. Parum est, non Deum rogáre, sed se laudáre; ínsuper et rogánti insultáre. Publicánus autem de longínquo stabat et Deo tamen ipse appropinquábat: cordis consciéntia eum removébat, píetas applicábat. Publicánus autem de longínquo stabat, sed Dóminus eum de propínquo attendébat.\par
\switchcolumn
\EN

\switchcolumn*

\LA
\begin{Resp}\Rsign Duo Séraphim clamábant alter ad álterum: \Star{} Sanctus, sanctus, sanctus Dóminus Deus Sábaoth: * Plena est omnis terra glória ejus.\end{Resp}
\begin{Resp}\Vsign Tres sunt qui testimónium dant in cælo: Pater, Verbum, et Spíritus Sanctus: et hi tres unum sunt.\end{Resp}
\begin{Resp}\Rsign Sanctus, sanctus, sanctus Dóminus Deus Sábaoth.\end{Resp}
\begin{Resp}\Vsign Glória Patri, et Fílio, \Star{} et Spirítui Sancto.\end{Resp}
\begin{Resp}\Rsign Plena est omnis terra glória ejus.\end{Resp}
\switchcolumn
\EN
\begin{Resp}\Rsign One Seraph cried unto another: \Star{} Holy, Holy, Holy is the Lord God of hosts; the whole earth is full of His glory.\end{Resp}
\begin{Resp}\Vsign There are Three That bear record in heaven: the Father, the Word, and the Holy Ghost, and these Three are One.\end{Resp}
\begin{Resp}\Rsign Holy, Holy, Holy is the Lord God of hosts.\end{Resp}
\begin{Resp}\Vsign Glory be to the Father, and to the Son, \Star{} and to the Holy Ghost.\end{Resp}
\begin{Resp}\Rsign The whole earth is full of His glory.\end{Resp}
\switchcolumn*

\LA
\LessonHead{Lectio IX}
\switchcolumn
\EN
\LessonHead{Lesson IX}
\switchcolumn*

\LA

\switchcolumn
\EN
\LessonTitle{“Though the Lord be high, yet hath He respect unto the lowly” but the proud, such as was this Pharisee, “He knoweth afar off.” (Ps. cxxxvii. 6.) He knoweth the proud, all the same, but they are afar off from Him. Consider now the lowliness of the publican. It was not only that he stood afar off, but “he would not lift up so much as his eyes unto heaven”. He looked carefully, lest he should look up, he dared not to lift up his eyes unto heaven. Self-knowledge kept him down, though hope raised him up. Consider again, how that he “smote upon his breast.” He afflicted himself, and therefore the Lord had compassion upon his acknowledgment of guilt.” He smote upon his breast, saying Lord, be merciful to me a sinner.” Hearken here to a prayer and wonderest thou that when the sinner remembereth, God forgetteth.}
\switchcolumn*

\LA
\noindent Excélsus enim Dóminus, et humília réspicit; excélsos autem, qualis erat ille pharisǽus, a longe cognóscit. Excélsa quidem Deus a longe cognóscit, sed non ignóscit. Adhuc audi humilitátem publicáni. Parum est, quia de longínquo stabat: nec óculos suos ad cælum levábat: ut aspicerétur, non aspiciébat; respícere sursum non audébat: premébat consciéntia, spes sublevábat. Adhuc audi: Percutiébat pectus suum. Pœnas a se ipso exigébat; proptérea Dóminus confiténti parcébat. Percutiébat pectus suum, dicens: Dómine, propítius esto mihi peccatóri. Ecce qui rogat. Quid miráris, si Deus ignóscit, quando ipse agnóscit?\par
\switchcolumn
\EN

\switchcolumn*

\LA
\TeDeum{Te Deum}
\switchcolumn
\EN
\TeDeum{Te Deum}
\switchcolumn*

\end{paracol}

\Office{Feria secunda infra Hebdomadam X post Octavam Pentecostes}{Monday of the Tenth Week after Pentecost}{Feria}
\addcontentsline{toc}{subsection}{Feria secunda infra Hebdomadam X post Octavam Pentecostes \textit{— Monday of the Tenth Week after Pentecost}}
\begin{paracol}{2}
\LA
\LessonHead{Lectio I}
\switchcolumn
\EN
\LessonHead{Lesson I}
\switchcolumn*

\LA
\LessonTitle{De libro quarto Regum}
\switchcolumn
\EN
\LessonTitle{Lesson from the second book of Kings}
\switchcolumn*

\LA
\Cite{4 Reg 11:1-3}
\switchcolumn
\EN
\Cite{2 Kgs 11:1-3}
\switchcolumn*

\LA
\noindent Athalía vero mater Ochozíæ videns mórtuum fílium suum surréxit et interfécit omne semen régium. Tollens autem Jósaba fília regis Joram, soror Ochozíæ, Joas fílium Ochozíæ, furáta est eum de médio filiórum regis qui interficiebántur, et nutrícem ejus de triclínio: et abscóndit eum a fácie Athalíæ, ut non interficerétur. Erátque cum ea sex annis clam in domo Dómini; porro Athalía regnávit super terram.\par
\switchcolumn
\EN
\noindent And Athalia the mother of Ochozias seeing that her son was dead, arose, and slew all the royal seed. But Josaba the daughter of king Joram, sister of Ochozias, took Joas the son of Ochozias, and stole him from among the king's sons that were slain, out of the bedchamber with his nurse: and hid him from the face of Athalia, so that he was not slain. And he was with her six years hid in the house of the Lord. And Athalia reigned over the land.\par
\switchcolumn*

\LA
\begin{Resp}\Rsign Recordáre, Dómine, testaménti tui, et dic Angelo percutiénti: Cesset jam manus tua, \Star{} Ut non desolétur terra, et ne perdas omnem ánimam vivam.\end{Resp}
\begin{Resp}\Vsign Ego sum qui peccávi, ego qui iníque egi: isti qui oves sunt, quid fecérunt? Avertátur, óbsecro, furor tuus, Dómine, a pópulo tuo.\end{Resp}
\begin{Resp}\Rsign Ut non desolétur terra, et ne perdas omnem ánimam vivam.\end{Resp}
\switchcolumn
\EN
\begin{Resp}\Rsign Remember, O Lord, thy covenant, and say unto the destroying Angel: Stay now thine hand; \Star{} That the land be not utterly laid waste, and that thou destroy not every living soul.\end{Resp}
\begin{Resp}\Vsign Even I it is that have sinned, and done evil indeed; but these sheep, what have they done? Let thine anger, I pray thee, O Lord, be turned away from thy people.\end{Resp}
\begin{Resp}\Rsign That the land be not utterly laid waste, and that Thou destroy not every living soul.\end{Resp}
\switchcolumn*

\LA
\LessonHead{Lectio II}
\switchcolumn
\EN
\LessonHead{Lesson II}
\switchcolumn*

\LA
\Cite{4 Reg 11:4-7}
\switchcolumn
\EN
\Cite{2 Kgs 11:4-7}
\switchcolumn*

\LA
\noindent Anno autem séptimo misit Jójada, et assúmens centuriónes et mílites, introdúxit ad se in templum Dómini, pepigítque cum eis fœdus; et adjúrans eos in domo Dómini, osténdit eis fílium regis: et præcépit illis, dicens: Iste est sermo quem fácere debétis: tértia pars vestrum intróëat sábbato, et obsérvet excúbias domus regis; tértia autem pars sit ad portam Sur, et tértia pars sit ad portam quæ est post habitáculum scutariórum, et custodiétis excúbias domus Messa. Duæ vero partes e vobis, omnes egrediéntes sábbato, custódiant excúbias domus Dómini circa regem.\par
\switchcolumn
\EN
\noindent And in the seventh year Joiada sent, and taking the centurions and the soldiers, brought them in to him into the temple of the Lord, and made a covenant with them: and taking an oath of them in the house of the Lord, showed them the king's son: And he commanded them, saying: This is the thing that you must do: Let a third part of you go in on the sabbath, and keep the watch of the king's house. And let a third part be at the gate of Sur: and let a third part be at the gate behind the dwelling of the shieldbearers: and you shall keep the watch of the house of Messa. But let two parts of you, all that go forth on the sabbath, keep the watch of the house of the Lord about the king.\par
\switchcolumn*

\LA
\begin{Resp}\Rsign Exaudísti, Dómine, oratiónem servi tui, ut ædificárem templum nómini tuo: \Star{} Bénedic et sanctífica domum istam in sempitérnum, Deus Israël.\end{Resp}
\begin{Resp}\Vsign Dómine, qui custódis pactum cum servis tuis, qui ámbulant coram te in toto corde suo.\end{Resp}
\begin{Resp}\Rsign Bénedic et sanctífica domum istam in sempitérnum, Deus Israël.\end{Resp}
\switchcolumn
\EN
\begin{Resp}\Rsign O Lord, Thou hast hearkened unto the prayer of thy servant, that I might build a temple unto thy Name, \Star{} O God of Israel, bless Thou, and hallow this house for ever.\end{Resp}
\begin{Resp}\Vsign O Lord, Who keepest covenant with thy servants that walk before thee in all their heart.\end{Resp}
\begin{Resp}\Rsign O God of Israel, bless Thou, and hallow this house for ever.\end{Resp}
\switchcolumn*

\LA
\LessonHead{Lectio III}
\switchcolumn
\EN
\LessonHead{Lesson III}
\switchcolumn*

\LA
\Cite{4 Reg 11:9-12}
\switchcolumn
\EN
\Cite{2 Kgs 11:9-12}
\switchcolumn*

\LA
\noindent Et assuméntes sínguli viros suos, qui ingrediebántur sábbatum, cum his qui egrediebántur sábbato, venérunt ad Jójadam sacerdótem, qui dedit eis hastas et arma regis David, quæ erant in domo Dómini. Et stetérunt sínguli habéntes arma in manu sua, a parte templi déxtera usque ad partem sinístram altáris et ædis, circum regem. Produxítque fílium regis, et pósuit super eum diadéma et testimónium: fecerúntque eum regem, et unxérunt: et plaudéntes manu, dixérunt: Vivat rex.\par
\switchcolumn
\EN
\noindent And the centurions did according to all things that Joiada the priest had commanded them: and taking every one their men, that went in on the sabbath, with them that went out on the sabbath, came to Joiada the priest. And he gave them the spears, and the arms of king David, which were in the house of the Lord. And they stood having every one their weapons in their hands, from the right side of the temple, unto the left side of the altar, and of the temple, about the king. And he brought forth the king's son, and put the diadem upon him, and the testimony: and they made him king, and anointed him: and clapping their hands. they said, God save the king.\par
\switchcolumn*

\LA
\begin{Resp}\Rsign Audi, Dómine, hymnum et oratiónem, quam servus tuus orat coram te hódie, ut sint óculi tui apérti, et aures tuæ inténtæ, \Star{} Super domum istam die ac nocte.\end{Resp}
\begin{Resp}\Vsign Réspice, Dómine, de sanctuário tuo, et de excélso cælórum habitáculo.\end{Resp}
\begin{Resp}\Rsign Super domum istam die ac nocte.\end{Resp}
\begin{Resp}\Vsign Glória Patri, et Fílio, \Star{} et Spirítui Sancto.\end{Resp}
\begin{Resp}\Rsign Super domum istam die ac nocte.\end{Resp}
\switchcolumn
\EN
\begin{Resp}\Rsign Hearken, O Lord, unto the cry and to the prayer which thy servant prayeth before thee today, that thine eyes may be open and thine ears attend; \Star{} Toward this house day and night.\end{Resp}
\begin{Resp}\Vsign Look down from thine high and holy place, O Lord, even from heaven thy dwelling.\end{Resp}
\begin{Resp}\Rsign Toward this house, day and night.\end{Resp}
\begin{Resp}\Vsign Glory be to the Father, and to the Son, \Star{} and to the Holy Ghost.\end{Resp}
\begin{Resp}\Rsign Toward this house, day and night.\end{Resp}
\switchcolumn*

\end{paracol}

\Office{Feria tertia infra Hebdomadam X post Octavam Pentecostes}{Tuesday of the Tenth Week after Pentecost}{Feria}
\addcontentsline{toc}{subsection}{Feria tertia infra Hebdomadam X post Octavam Pentecostes \textit{— Tuesday of the Tenth Week after Pentecost}}
\begin{paracol}{2}
\LA
\LessonHead{Lectio I}
\switchcolumn
\EN
\LessonHead{Lesson I}
\switchcolumn*

\LA
\LessonTitle{De libro quarto Regum}
\switchcolumn
\EN
\LessonTitle{Lesson from the second book of Kings}
\switchcolumn*

\LA
\Cite{4 Reg 12:1-3}
\switchcolumn
\EN
\Cite{2 Kgs 12:1-3}
\switchcolumn*

\LA
\noindent Anno séptimo Jehu, regnávit Joas: et quadragínta annis regnávit in Jerúsalem. Nomen matris ejus Sébia de Bersabée. Fecítque Joas rectum coram Dómino cunctis diébus quibus dócuit eum Jójada sacérdos; verúmtamen excélsa non ábstulit; adhuc enim pópulus immolábat, et adolébat in excélsis incénsum.\par
\switchcolumn
\EN
\noindent In the seventh year of Jehu Joas began to reign: and he reigned forty years in Jerusalem. The name of his mother was Sebia of Bersabee. And Joas did that which was right before the Lord all the days that Joiada the priest taught him. But yet he took not away the high places: for the people still sacrificed and burnt incense in the high places.\par
\switchcolumn*

\LA
\begin{Resp}\Rsign Dómine, si convérsus fúerit pópulus tuus, et oráverit ad sanctuárium tuum: \Star{} Tu exáudies de cælo, Dómine, et líbera eos de mánibus inimicórum suórum.\end{Resp}
\begin{Resp}\Vsign Si peccáverit in te pópulus tuus, et convérsus égerit pœniténtiam, veniénsque oráverit in isto loco.\end{Resp}
\begin{Resp}\Rsign Tu exáudies de cælo, Dómine, et líbera eos de mánibus inimicórum suórum.\end{Resp}
\switchcolumn
\EN
\begin{Resp}\Rsign Lord, when thy people shall turn again to thee, and shall pray unto thee in this house; \Star{} Then hear Thou in heaven, O Lord, and deliver them out of the hand of their enemies.\end{Resp}
\begin{Resp}\Vsign If thy people sin against thee, and turn again, and repent, and come and pray unto thee in this house.\end{Resp}
\begin{Resp}\Rsign Then hear Thou in heaven, O Lord, and deliver them out of the hand of their enemies.\end{Resp}
\switchcolumn*

\LA
\LessonHead{Lectio II}
\switchcolumn
\EN
\LessonHead{Lesson II}
\switchcolumn*

\LA
\Cite{4 Reg 12:4-5}
\switchcolumn
\EN
\Cite{2 Kgs 12:4-5}
\switchcolumn*

\LA
\noindent Dixítque Joas ad sacerdótes: Omnem pecúniam sanctórum, quæ illáta fúerit in templum Dómini a prætereúntibus, quæ offértur pro prétio ánimæ, et quam sponte et arbítrio cordis sui ínferunt in templum Dómini, accípiant illam sacerdótes juxta órdinem suum, et instáurent sartatécta domus, si quid necessárium víderint instauratióne.\par
\switchcolumn
\EN
\noindent And Joas said to the priests: O All the money of the sanctified things, which is brought into the temple of the Lord by those that pass, which is offered for the price of a soul, and which of their own accord, and of their own free heart they bring into the temple of the Lord: Let the priests take it according to their order, and repair the house, wheresoever they shall see any thing that wanteth repairing.\par
\switchcolumn*

\LA
\begin{Resp}\Rsign Factum est, dum tólleret Dóminus Elíam per túrbinem in cælum, \Star{} Eliséus clamábat, dicens: Pater mi, pater mi, currus Israël, et auríga ejus.\end{Resp}
\begin{Resp}\Vsign Cumque pérgerent, et incedéntes sermocinaréntur, ecce currus ígneus et equi ígnei divisérunt utrúmque, et ascéndit Elías per túrbinem in cælum.\end{Resp}
\begin{Resp}\Rsign Eliséus clamábat, dicens: Pater mi, pater mi, currus Israël, et auríga ejus.\end{Resp}
\switchcolumn
\EN
\begin{Resp}\Rsign And it came to pass when the Lord would take up Elijah into heaven by a whirlwind \Star{} Elisha cried, and said: My father, my father, the chariot of Israel, and the horsemen thereof.\end{Resp}
\begin{Resp}\Vsign And as they still went on, and talked, behold, there appeared a chariot of fire, and horses of fire, and parted them both asunder, and Elijah went up by a whirlwind into heaven.\end{Resp}
\begin{Resp}\Rsign Elisha cried, and said: My father, my father, the chariot of Israel, and the horsemen thereof.\end{Resp}
\switchcolumn*

\LA
\LessonHead{Lectio III}
\switchcolumn
\EN
\LessonHead{Lesson III}
\switchcolumn*

\LA
\Cite{4 Reg 12:6-8}
\switchcolumn
\EN
\Cite{2 Kgs 12:6-8}
\switchcolumn*

\LA
\noindent Igitur usque ad vigésimum tértium annum regis Joas, non instauravérunt sacerdótes sartatécta templi. Vocavítque rex Joas Jójadam pontíficem et sacerdótes, dicens eis: Quare sartatécta non instaurátis templi? Nolíte ergo ámplius accípere pecúniam juxta órdinem vestrum; sed ad instauratiónem templi réddite eam. Prohibitíque sunt sacerdótes ultra accípere pecúniam a pópulo, et instauráre sartatécta domus.\par
\switchcolumn
\EN
\noindent Now till the three and twentieth year of king Joas, the priests did not make the repairs of the temple. And king Joas called Joiada the high priest and the priests, saying to them: Why do you not repair the temple? Take you therefore money no more according to your order, but restore it for the repairing of the temple. And the priests were forbidden to take any more money of the people, and to make the repairs of the house.\par
\switchcolumn*

\LA
\begin{Resp}\Rsign Ego te tuli de domo patris tui, dicit Dóminus, et pósui te páscere gregem pópuli mei: \Star{} Et fui tecum in ómnibus ubicúmque ambulásti, firmans regnum tuum in ætérnum.\end{Resp}
\begin{Resp}\Vsign Fecíque tibi nomen grande, juxta nomen magnórum, qui sunt in terra: et réquiem dedi tibi ab ómnibus inimícis tuis.\end{Resp}
\begin{Resp}\Rsign Et fui tecum in ómnibus ubicúmque ambulásti, firmans regnum tuum in ætérnum.\end{Resp}
\begin{Resp}\Vsign Glória Patri, et Fílio, \Star{} et Spirítui Sancto.\end{Resp}
\begin{Resp}\Rsign Et fui tecum in ómnibus ubicúmque ambulásti, firmans regnum tuum in ætérnum.\end{Resp}
\switchcolumn
\EN
\begin{Resp}\Rsign Thus saith the Lord: I took thee out of thy father's house, and appointed thee to be ruler over My people, over Israel. \Star{} And I was with thee whithersoever thou wentest, to establish thy kingdom for ever.\end{Resp}
\begin{Resp}\Vsign And I have made thee a great name, like unto the name of the great men that are in the earth and have caused thee to rest from all thine enemies.\end{Resp}
\begin{Resp}\Rsign And I was with thee whithersoever thou wentest, to establish thy kingdom for ever.\end{Resp}
\begin{Resp}\Vsign Glory be to the Father, and to the Son, \Star{} and to the Holy Ghost.\end{Resp}
\begin{Resp}\Rsign And I was with thee whithersoever thou wentest, to establish thy kingdom for ever.\end{Resp}
\switchcolumn*

\end{paracol}

\Office{Feria quarta infra Hebdomadam X post Octavam Pentecostes}{Wednesday of the Tenth Week after Pentecost}{Feria}
\addcontentsline{toc}{subsection}{Feria quarta infra Hebdomadam X post Octavam Pentecostes \textit{— Wednesday of the Tenth Week after Pentecost}}
\begin{paracol}{2}
\LA
\LessonHead{Lectio I}
\switchcolumn
\EN
\LessonHead{Lesson I}
\switchcolumn*

\LA
\LessonTitle{De libro quarto Regum}
\switchcolumn
\EN
\LessonTitle{Lesson from the second book of Kings}
\switchcolumn*

\LA
\Cite{4 Reg 13:14-17}
\switchcolumn
\EN
\Cite{2 Kgs 13:14-17}
\switchcolumn*

\LA
\noindent Eliséus autem ægrotábat infirmitáte, qua et mórtuus est; descendítque ad eum Joas rex Israël, et flebat coram eo, dicebátque: Pater mi, pater mi, currus Israël et auríga ejus. Et ait illi Eliséus: Affer arcum et sagíttas. Cumque attulísset ad eum arcum et sagíttas, dixit ad regem Israël: Pone manum tuam super arcum. Et, cum posuísset ille manum suam, superpósuit Eliséus manus suas mánibus regis, et ait: Aperi fenéstram orientálem. Cumque aperuísset, dixit Eliséus: Jace sagíttam. Et jecit. Et ait Eliséus: Sagítta salútis Dómini, et sagítta salútis contra Sýriam; percutiésque Sýriam in Aphec, donec consúmas eam.\par
\switchcolumn
\EN
\noindent Now Eliseus was sick of the illness whereof he died: and Joas king of Israel went down to him, and wept before him, and said: O my father, my father, the chariot of Israel and the guider thereof. And Eliseus said to him: Bring a bow and arrows. And when he had brought him a bow, and arrows, He said to the king of Israel: Put thy hand upon the bow. And when he had put his hand, Eliseus put his hands over the king's hands, And said: Open the window to the east. And when he had opened it, Eliseus said: Shoot an arrow. And he shot. And Eliseus said: The arrow of the Lord's deliverance, and the arrow of the deliverance from Syria: and thou shalt strike the Syrians in Aphec, till thou consume them.\par
\switchcolumn*

\LA
\begin{Resp}\Rsign Peccávi super númerum arénæ maris, et multiplicáta sunt peccáta mea: et non sum dignus vidére altitúdinem cæli præ multitúdine iniquitátis meæ: quóniam irritávi iram tuam, \Star{} Et malum coram te feci.\end{Resp}
\begin{Resp}\Vsign Quóniam iniquitátem meam ego cognósco: et delíctum meum contra me est semper, quia tibi soli peccávi.\end{Resp}
\begin{Resp}\Rsign Et malum coram te feci.\end{Resp}
\switchcolumn
\EN
\begin{Resp}\Rsign My sins are many, yea, they are more in number than the sands of the sea; I am not worthy to look up toward heaven because of the multitude of my iniquities; for I have provoked thee to anger; \Star{} And done evil in thy sight.\end{Resp}
\begin{Resp}\Vsign For I acknowledge my transgression, and my sin is ever before me, for against thee only have I sinned.\end{Resp}
\begin{Resp}\Rsign And done evil in thy sight.\end{Resp}
\switchcolumn*

\LA
\LessonHead{Lectio II}
\switchcolumn
\EN
\LessonHead{Lesson II}
\switchcolumn*

\LA
\Cite{4 Reg 13:18-20}
\switchcolumn
\EN
\Cite{2 Kgs 13:18-20}
\switchcolumn*

\LA
\noindent Et ait: Tolle sagíttas. Qui cum tulísset, rursum dixit ei: Pércute jáculo terram. Et, cum percussísset tribus vícibus, et stetísset, irátus est vir Dei contra eum, et ait: Si percussísses quínquies aut séxies sive sépties, percussísses Sýriam usque ad consumptiónem; nunc autem tribus vícibus percúties eam. Mórtuus est ergo Eliséus, et sepeliérunt eum. Latrúnculi autem de Moab venérunt in terram in ipso anno.\par
\switchcolumn
\EN
\noindent And he said: Take the arrows. And when he had taken them, he said to him: Strike with an arrow upon the ground. And he struck three times and stood still. And the man of God was angry with him, and said: If thou hadst smitten five or six or seven times, thou hadst smitten Syria even to utter destruction: but now three times shalt thou smite it. And Eliseus died, and they buried him. And the rovers from Moab came into the land the same year.\par
\switchcolumn*

\LA
\begin{Resp}\Rsign Exaudísti, Dómine, oratiónem servi tui, ut ædificárem templum nómini tuo: \Star{} Bénedic et sanctífica domum istam in sempitérnum, Deus Israël.\end{Resp}
\begin{Resp}\Vsign Dómine, qui custódis pactum cum servis tuis, qui ámbulant coram te in toto corde suo.\end{Resp}
\begin{Resp}\Rsign Bénedic et sanctífica domum istam in sempitérnum, Deus Israël.\end{Resp}
\switchcolumn
\EN
\begin{Resp}\Rsign O Lord, Thou hast hearkened unto the prayer of thy servant, that I might build a temple unto thy Name, \Star{} O God of Israel, bless Thou, and hallow this house for ever.\end{Resp}
\begin{Resp}\Vsign O Lord, Who keepest covenant with thy servants that walk before thee in all their heart.\end{Resp}
\begin{Resp}\Rsign O God of Israel, bless Thou, and hallow this house for ever.\end{Resp}
\switchcolumn*

\LA
\LessonHead{Lectio III}
\switchcolumn
\EN
\LessonHead{Lesson III}
\switchcolumn*

\LA
\Cite{4 Reg 13:21; 13:24-25}
\switchcolumn
\EN
\Cite{2 Kgs 13:21; 13:24-25}
\switchcolumn*

\LA
\noindent Quidam autem sepeliéntes hóminem, vidérunt latrúnculos, et projecérunt cadáver in sepúlcro Eliséi. Quod cum tetigísset ossa Eliséi, revíxit homo et stetit super pedes suos. Mórtuus est autem Házaël rex Sýriæ, et regnávit Bénadad fílius ejus pro eo. Porro Joas fílius Jóachaz tulit urbes de manu Bénadad fílii Házaël, quas túlerat de manu Jóachaz patris sui jure prǽlii; tribus vícibus percússit eum Joas, et réddidit civitátes Israël.\par
\switchcolumn
\EN
\noindent And some that were burying a man, saw the rovers, and cast the body into the sepulchre of Eliseus. And when it had touched the bones of Eliseus, the man came to life, and stood upon his feet. And Hazael king of Syria died, and Benadad his son reigned in his stead. Now Joas the son of Joachaz, took the cities out of the hand of Benadad, the son of Hazael, which he had taken out of the hand of Joachaz his father by war, three times did Joas beat him, and he restored the cities to Israel.\par
\switchcolumn*

\LA
\begin{Resp}\Rsign Audi, Dómine, hymnum et oratiónem, quam servus tuus orat coram te hódie, ut sint óculi tui apérti, et aures tuæ inténtæ, \Star{} Super domum istam die ac nocte.\end{Resp}
\begin{Resp}\Vsign Réspice, Dómine, de sanctuário tuo, et de excélso cælórum habitáculo.\end{Resp}
\begin{Resp}\Rsign Super domum istam die ac nocte.\end{Resp}
\begin{Resp}\Vsign Glória Patri, et Fílio, \Star{} et Spirítui Sancto.\end{Resp}
\begin{Resp}\Rsign Super domum istam die ac nocte.\end{Resp}
\switchcolumn
\EN
\begin{Resp}\Rsign Hearken, O Lord, unto the cry and to the prayer which thy servant prayeth before thee today, that thine eyes may be open and thine ears attend; \Star{} Toward this house day and night.\end{Resp}
\begin{Resp}\Vsign Look down from thine high and holy place, O Lord, even from heaven thy dwelling.\end{Resp}
\begin{Resp}\Rsign Toward this house, day and night.\end{Resp}
\begin{Resp}\Vsign Glory be to the Father, and to the Son, \Star{} and to the Holy Ghost.\end{Resp}
\begin{Resp}\Rsign Toward this house, day and night.\end{Resp}
\switchcolumn*

\end{paracol}

\Office{Feria quinta infra Hebdomadam X post Octavam Pentecostes}{Thursday of the Tenth Week after Pentecost}{Feria}
\addcontentsline{toc}{subsection}{Feria quinta infra Hebdomadam X post Octavam Pentecostes \textit{— Thursday of the Tenth Week after Pentecost}}
\begin{paracol}{2}
\LA
\LessonHead{Lectio I}
\switchcolumn
\EN
\LessonHead{Lesson I}
\switchcolumn*

\LA
\LessonTitle{De libro quarto Regum}
\switchcolumn
\EN
\LessonTitle{Lesson from the second book of Kings}
\switchcolumn*

\LA
\Cite{4 Reg 17:6-9}
\switchcolumn
\EN
\Cite{2 Kgs 17:6-9}
\switchcolumn*

\LA
\noindent Anno autem nono Osée, cepit rex Assyriórum Samaríam et tránstulit Israël in Assyrios: posuítque eos in Hala et in Habor juxta fluvium Gozan, in civitátibus Medórum. Factum est enim, cum peccássent fílii Israël Dómino Deo suo, qui edúxerat eos de terra Ægýpti, de manu pharaónis regis Ægýpti, coluérunt deos alienos. Et ambulavérunt juxta ritum Géntium, quas consúmpserat Dóminus in conspéctu filiórum Israël et regum Israël, quia simíliter fécerant. Et offendérunt fílii Israël verbis non rectis Dóminum Deum suum: et ædificavérunt sibi excélsa in cunctis úrbibus suis.\par
\switchcolumn
\EN
\noindent And in the ninth year of Osee, the king of the Assyrians took Samaria, and carried Israel away to Assyria: and he placed them in Hala and Habor by the river of Gozan, in the cities of the Medes. For so it was that the children of Israel had sinned against the Lord their God, who brought them out of the land of Egypt, from under the hand of Pharao king of Egypt, and they worshipped strange gods. And they walked according to the way of the nations which the Lord had destroyed in the sight of the children of Israel and of the kings of Israel: because they had done in like manner. And the children of Israel offended the Lord their God with things that were not right: and built them high places in all their cities from the tower of the watchmen to the fenced city.\par
\switchcolumn*

\LA
\begin{Resp}\Rsign Præparáte corda vestra Dómino, et servíte illi soli: \Star{} Et liberábit vos de mánibus inimicórum vestrórum.\end{Resp}
\begin{Resp}\Vsign Convertímini ad eum in toto corde vestro, et auférte deos aliénos de médio vestri.\end{Resp}
\begin{Resp}\Rsign Et liberábit vos de mánibus inimicórum vestrórum.\end{Resp}
\switchcolumn
\EN
\begin{Resp}\Rsign Prepare your hearts unto the Lord, and serve Him Only \Star{} And He will deliver you out of the hand of your enemies.\end{Resp}
\begin{Resp}\Vsign Return unto Him with all your hearts, and put away the strange gods from among you.\end{Resp}
\begin{Resp}\Rsign And He will deliver you out of the hand of your enemies.\end{Resp}
\switchcolumn*

\LA
\LessonHead{Lectio II}
\switchcolumn
\EN
\LessonHead{Lesson II}
\switchcolumn*

\LA
\Cite{4 Reg 17:13-15}
\switchcolumn
\EN
\Cite{2 Kgs 17:13-15}
\switchcolumn*

\LA
\noindent Et testificátus est Dóminus in Israël et in Juda per manum ómnium prophetárum et vidéntium, dicens: Revertímini a viis vestris pessimis, et custodíte præcépta mea et cæremónias, juxta omnem legem, quam præcépi pátribus vestris, et sicut misi ad vos in manu servórum meórum prophetárum. Qui non audiérunt, sed induravérunt cervícem suam juxta cervícem patrum suórum, qui noluérunt obedíre Dómino Deo suo; et abjecérunt legítima ejus, et pactum quod pépigit cum pátribus eórum, et testificatiónes, quibus contestátus est eos, secutíque sunt vanitátes, et vane egérunt.\par
\switchcolumn
\EN
\noindent And the Lord testified to them in Israel and in Juda by the hand of all the prophets and seers, saying: Return from your wicked ways, and keep my precepts, and ceremonies, according to all the law which I commanded your fathers: and as I have sent to you in the hand of my servants the prophets. And they hearkened not, but hardened their necks like to the neck of their fathers, who would not obey the Lord their God. And they rejected his ordinances and the covenant that he made with their fathers, and the testimonies which he testified against them: and they followed vanities, and acted vainly: and they followed the nations that were round about them, concerning which the Lord had commanded them that they should not do as they did.\par
\switchcolumn*

\LA
\begin{Resp}\Rsign Deus ómnium exaudítor est: ipse misit Angelum suum, et tulit me de óvibus patris mei: \Star{} Et unxit me unctióne misericórdiæ suæ.\end{Resp}
\begin{Resp}\Vsign Dóminus, qui erípuit me de ore leónis, et de manu béstiæ liberávit me.\end{Resp}
\begin{Resp}\Rsign Et unxit me unctióne misericórdiæ suæ.\end{Resp}
\switchcolumn
\EN
\begin{Resp}\Rsign God, Which heareth all, even He sent His Angel, and took me from keeping my father's sheep, and \Star{} Anointed me with the oil of His mercy.\end{Resp}
\begin{Resp}\Vsign The Lord That delivered me out of the mouth of the lion, and out of the paw of the bear\end{Resp}
\begin{Resp}\Rsign And anointed me with the oil of His mercy.\end{Resp}
\switchcolumn*

\LA
\LessonHead{Lectio III}
\switchcolumn
\EN
\LessonHead{Lesson III}
\switchcolumn*

\LA
\Cite{4 Reg 17:18-21}
\switchcolumn
\EN
\Cite{2 Kgs 17:18-21}
\switchcolumn*

\LA
\noindent Iratúsque est Dóminus veheménter Israéli, et ábstulit eos a conspéctu suo; et non remánsit nisi tribus Juda tantúmmodo. Sed nec ipse Juda custodívit mandáta Dómini Dei sui: verum ambulávit in erróribus Israël, quos operátus fuerat. Projecítque Dóminus omne semen Israël, et afflíxit eos, et trádidit eos in manu diripiéntium, donec proíceret eos a fácie sua: ex eo jam témpore, quo scissus est Israël a domo David, et constituérunt sibi regem Jeróboam fílium Nabat.\par
\switchcolumn
\EN
\noindent And the Lord was very angry with Israel, and removed them from his sight, and there remained only the tribe of Juda. But neither did Juda itself keep the commandments of the Lord their God: but they walked in the errors of Israel, which they had wrought. And the Lord cast off all the seed of Israel, and afflicted them and delivered them into the hand of spoilers, till he cast them away from his face: Even from that time, when Israel was rent from the house of David, and made Jeroboam son of Nabat their king: for Jeroboam separated Israel from the Lord, and made them commit a great sin.\par
\switchcolumn*

\LA
\begin{Resp}\Rsign Dóminus, qui erípuit me de ore leónis, et de manu béstiæ liberávit me, \Star{} Ipse me erípiet de mánibus inimicórum meórum.\end{Resp}
\begin{Resp}\Vsign Misit Deus misericórdiam suam et veritátem suam: ánimam meam erípuit de médio catulórum leónum.\end{Resp}
\begin{Resp}\Rsign Ipse me erípiet de mánibus inimicórum meórum.\end{Resp}
\begin{Resp}\Vsign Glória Patri, et Fílio, \Star{} et Spirítui Sancto.\end{Resp}
\begin{Resp}\Rsign Ipse me erípiet de mánibus inimicórum meórum.\end{Resp}
\switchcolumn
\EN
\begin{Resp}\Rsign The Lord That delivered me out of the mouth of the lion, and out of the paw of the bear \Star{} He will deliver me out of the hand of mine enemies.\end{Resp}
\begin{Resp}\Vsign God hath sent forth His mercy and His truth, and delivered my soul from among the lion's whelps.\end{Resp}
\begin{Resp}\Rsign He will deliver me out of the hand of mine enemies.\end{Resp}
\begin{Resp}\Vsign Glory be to the Father, and to the Son, \Star{} and to the Holy Ghost.\end{Resp}
\begin{Resp}\Rsign He will deliver me out of the hand of mine enemies.\end{Resp}
\switchcolumn*

\end{paracol}

\Office{Feria sexta infra Hebdomadam X post Octavam Pentecostes}{Friday of the Tenth Week after Pentecost}{Feria}
\addcontentsline{toc}{subsection}{Feria sexta infra Hebdomadam X post Octavam Pentecostes \textit{— Friday of the Tenth Week after Pentecost}}
\begin{paracol}{2}
\LA
\LessonHead{Lectio I}
\switchcolumn
\EN
\LessonHead{Lesson I}
\switchcolumn*

\LA
\LessonTitle{De libro quarto Regum}
\switchcolumn
\EN
\LessonTitle{Lesson from the second book of Kings}
\switchcolumn*

\LA
\Cite{4 Reg 17:21-23}
\switchcolumn
\EN
\Cite{2 Kgs 17:21-23}
\switchcolumn*

\LA
\noindent Separávit Jeróboam Israël a Dómino, et peccáre eos fecit peccátum magnum. Et ambulavérunt fílii Israël in univérsis peccátis Jeróboam quæ fécerat, et non recessérunt ab eis, úsquequo Dóminus auférret Israël a fácie sua, sicut locútus fúerat in manu ómnium servórum suórum prophetárum; translatúsque est Israël de terra sua in Assýrios, usque in diem hanc.\par
\switchcolumn
\EN
\noindent Jeroboam separated Israel from the Lord, and made them commit a great sin. And the children of Israel walked in all the sins of Jeroboam, which he had done: and they departed not from them, Till the Lord removed Israel from his face, as he had spoken in the hand of all his servants the prophets: and Israel was carried away out of their land to Assyria, unto this day.\par
\switchcolumn*

\LA
\begin{Resp}\Rsign Percússit Saul mille, et David decem míllia: \Star{} Quia manus Dómini erat cum illo: percússit Philisthǽum, et ábstulit oppróbrium ex Israël.\end{Resp}
\begin{Resp}\Vsign Nonne iste est David, de quo canébant in choro, dicéntes: Saul percússit mille, et David decem míllia?\end{Resp}
\begin{Resp}\Rsign Quia manus Dómini erat cum illo: percússit Philisthǽum, et ábstulit oppróbrium ex Israël.\end{Resp}
\switchcolumn
\EN
\begin{Resp}\Rsign Saul hath slain his thousands, and David his ten thousands. \Star{} Because the hand of the Lord was with him, he smote the Philistine, and took away the reproach from Israel.\end{Resp}
\begin{Resp}\Vsign Is not this David? Did they not sing one to another of him in dances, saying: Saul hath slain his thousands, and David his ten thousands?\end{Resp}
\begin{Resp}\Rsign Because the hand of the Lord was with him, he smote the Philistine, and took away the reproach from Israel.\end{Resp}
\switchcolumn*

\LA
\LessonHead{Lectio II}
\switchcolumn
\EN
\LessonHead{Lesson II}
\switchcolumn*

\LA
\Cite{4 Reg 17:24-25}
\switchcolumn
\EN
\Cite{2 Kgs 17:24-25}
\switchcolumn*

\LA
\noindent Addúxit autem rex Assyriórum de Babylóne et de Cutha et de Avah et de Emath et de Sephárvaim, et collocávit eos in civitátibus Samaríæ pro fíliis Israël: qui possedérunt Samaríam, et habitavérunt in úrbibus ejus. Cumque ibi habitáre cœpíssent, non timébant Dóminum, et immísit in eos Dóminus leónes, qui interficiébant eos.\par
\switchcolumn
\EN
\noindent And the king of the Assyrians brought people from Babylon, and from Cutha, and from Avah, and from Emath, and from Sepharvaim: and placed them in the cities of Samaria instead of the children of Israel: and they possessed Samaria, and dwelt in the cities thereof. And when they began to dwell there, they feared not the Lord: and the Lord sent lions among them, which killed them.\par
\switchcolumn*

\LA
\begin{Resp}\Rsign Montes Gélboë, nec ros nec plúvia véniant super vos, \Star{} Ubi cecidérunt fortes Israël.\end{Resp}
\begin{Resp}\Vsign Omnes montes, qui estis in circúitu ejus, vísitet Dóminus: a Gélboë autem tránseat.\end{Resp}
\begin{Resp}\Rsign Ubi cecidérunt fortes Israël.\end{Resp}
\switchcolumn
\EN
\begin{Resp}\Rsign Ye mountains of Gilboa, let there be no dew, neither let there be rain upon you \Star{} For there are the mighty of Israel fallen.\end{Resp}
\begin{Resp}\Vsign All ye mountains that stand round about, the Lord look upon you but let Him pass by Gilboa\end{Resp}
\begin{Resp}\Rsign For there are the mighty of Israel fallen.\end{Resp}
\switchcolumn*

\LA
\LessonHead{Lectio III}
\switchcolumn
\EN
\LessonHead{Lesson III}
\switchcolumn*

\LA
\Cite{4 Reg 17:26-27}
\switchcolumn
\EN
\Cite{2 Kgs 17:26-27}
\switchcolumn*

\LA
\noindent Nuntiatúmque est regi Assyriórum et dictum: gentes, quas transtulísti, et habitáre fecísti in civitátibus Samaríæ, ignórant legítima Dei terræ: et immísit in eos Dóminus leónes, et ecce interfíciunt eos, eo quod ignórent ritum Dei terræ. Præcépit autem rex Assyriórum, dicens: Dúcite illuc unum de sacerdótibus, quos inde captívos adduxístis; et vadat, et hábitet cum eis: et dóceat eos legítima Dei terræ.\par
\switchcolumn
\EN
\noindent And it was told the king of the Assyrians, and it was said: The nations which thou hast removed, and made to dwell in the cities of Samaria, know not the ordinances of the God of the land: and the Lord hath sent lions among them: and behold they kill them, because they know not the manner of the God of the land. And the king of the Assyrians commanded, saying: Carry thither one of the priests whom you brought from thence captive, and let him go, and dwell with them: and let him teach them the ordinances of the God of the land.\par
\switchcolumn*

\LA
\begin{Resp}\Rsign Ego te tuli de domo patris tui, dicit Dóminus, et pósui te páscere gregem pópuli mei: \Star{} Et fui tecum in ómnibus ubicúmque ambulásti, firmans regnum tuum in ætérnum.\end{Resp}
\begin{Resp}\Vsign Fecíque tibi nomen grande, juxta nomen magnórum, qui sunt in terra: et réquiem dedi tibi ab ómnibus inimícis tuis.\end{Resp}
\begin{Resp}\Rsign Et fui tecum in ómnibus ubicúmque ambulásti, firmans regnum tuum in ætérnum.\end{Resp}
\begin{Resp}\Vsign Glória Patri, et Fílio, \Star{} et Spirítui Sancto.\end{Resp}
\begin{Resp}\Rsign Et fui tecum in ómnibus ubicúmque ambulásti, firmans regnum tuum in ætérnum.\end{Resp}
\switchcolumn
\EN
\begin{Resp}\Rsign Thus saith the Lord: I took thee out of thy father's house, and appointed thee to be ruler over My people, over Israel. \Star{} And I was with thee whithersoever thou wentest, to establish thy kingdom for ever.\end{Resp}
\begin{Resp}\Vsign And I have made thee a great name, like unto the name of the great men that are in the earth, and have caused thee to rest from all thine enemies.\end{Resp}
\begin{Resp}\Rsign And I was with thee whithersoever thou wentest, to establish thy kingdom for ever.\end{Resp}
\begin{Resp}\Vsign Glory be to the Father, and to the Son, \Star{} and to the Holy Ghost.\end{Resp}
\begin{Resp}\Rsign And I was with thee whithersoever thou wentest, to establish thy kingdom for ever.\end{Resp}
\switchcolumn*

\end{paracol}

\Office{Sabbato infra Hebdomadam X post Octavam Pentecostes}{Saturday of the Tenth Week after Pentecost}{Feria}
\addcontentsline{toc}{subsection}{Sabbato infra Hebdomadam X post Octavam Pentecostes \textit{— Saturday of the Tenth Week after Pentecost}}
\begin{paracol}{2}
\LA
\LessonHead{Lectio I}
\switchcolumn
\EN
\LessonHead{Lesson I}
\switchcolumn*

\LA
\LessonTitle{De libro quarto Regum}
\switchcolumn
\EN
\LessonTitle{Lesson from the second book of Kings}
\switchcolumn*

\LA
\Cite{4 Reg 18:1-5}
\switchcolumn
\EN
\Cite{2 Kgs 18:1-5}
\switchcolumn*

\LA
\noindent Anno tértio Osée fílii Ela regis Israël, regnávit Ezechías fílius Achaz regis Juda. Vigínti quinque annórum erat cum regnáre cœpísset, et vigínti novem annis regnávit in Jerúsalem: Nomen matris ejus Abi fília Zacharíæ. Fecítque quod erat bonum coram Dómino, juxta ómnia quæ fécerat David pater ejus. Ipse dissipávit excélsa, et contrívit státuas, et succídit lucos, confregítque serpéntem ǽneum, quem fécerat Móyses: síquidem usque ad illud tempus fílii Israël adolébant ei incénsum: vocavítque nomen ejus Nohéstan. In Dómino Deo Israël sperávit.\par
\switchcolumn
\EN
\noindent In the third year of Osee the son of Ela king of Israel, reigned Ezechias the son of Achaz king of Juda. He was five and twenty years old when he began to reign: and he reigned nine and twenty years in Jerusalem: the name of his mother was Abi the daughter of Zacharias. And he did that which was good before the Lord, according to all that David his father had done. He destroyed the high places, and broke the statues in pieces, and cut down the groves, and broke the brazen serpent, which Moses had made: for till that time the children of Israel burnt incense to it: and he called its name Nohestan. He trusted in the Lord the God of Israel:\par
\switchcolumn*

\LA
\begin{Resp}\Rsign Peccávi super númerum arénæ maris, et multiplicáta sunt peccáta mea: et non sum dignus vidére altitúdinem cæli præ multitúdine iniquitátis meæ: quóniam irritávi iram tuam, \Star{} Et malum coram te feci.\end{Resp}
\begin{Resp}\Vsign Quóniam iniquitátem meam ego cognósco: et delíctum meum contra me est semper, quia tibi soli peccávi.\end{Resp}
\begin{Resp}\Rsign Et malum coram te feci.\end{Resp}
\switchcolumn
\EN
\begin{Resp}\Rsign My sins are many, yea, they are more in number than the sands of the sea; I am not worthy to look up toward heaven because of the multitude of my iniquities; for I have provoked thee to anger; \Star{} And done evil in thy sight.\end{Resp}
\begin{Resp}\Vsign For I acknowledge my transgression, and my sin is ever before me, for against thee only have I sinned.\end{Resp}
\begin{Resp}\Rsign And done evil in thy sight.\end{Resp}
\switchcolumn*

\LA
\LessonHead{Lectio II}
\switchcolumn
\EN
\LessonHead{Lesson II}
\switchcolumn*

\LA
\Cite{4 Reg 18:5-8}
\switchcolumn
\EN
\Cite{2 Kgs 18:5-8}
\switchcolumn*

\LA
\noindent Itaque post eum non fuit símilis ei de cunctis régibus Juda, sed neque in his qui ante eum fuérunt. Et adhǽsit Dómino, et non recéssit a vestígiis ejus, fecítque mandáta ejus, quæ præcéperat Dóminus Móysi. Unde et erat Dóminus cum eo, et in cunctis ad quæ procedébat, sapiénter se agébat. Rebellávit quoque contra regem Assyriórum, et non servívit ei. Ipse percússit Philisthǽos usque ad Gazam, et omnes términos eórum, a turre custódum usque ad civitátem munítam.\par
\switchcolumn
\EN
\noindent so that after him there was none like him among all the kings of Juda, nor any of them that were before him: And he stuck to the Lord, and departed not from his steps, but kept his commandments, which the Lord commanded Moses. Wherefore the Lord also was with him, and in all things, to which he went forth, he behaved himself wisely. And he rebelled against the king of the Assyrians, and served him not. He smote the Philistines as far as Gaza, and all their borders, from the tower of the watchmen to the fenced city.\par
\switchcolumn*

\LA
\begin{Resp}\Rsign Exaudísti, Dómine, oratiónem servi tui, ut ædificárem templum nómini tuo: \Star{} Bénedic et sanctífica domum istam in sempitérnum, Deus Israël.\end{Resp}
\begin{Resp}\Vsign Dómine, qui custódis pactum cum servis tuis, qui ámbulant coram te in toto corde suo.\end{Resp}
\begin{Resp}\Rsign Bénedic et sanctífica domum istam in sempitérnum, Deus Israël.\end{Resp}
\switchcolumn
\EN
\begin{Resp}\Rsign O Lord, Thou hast hearkened unto the prayer of thy servant, that I might build a temple unto thy Name, \Star{} O God of Israel, bless Thou, and hallow this house for ever.\end{Resp}
\begin{Resp}\Vsign O Lord, Who keepest covenant with thy servants that walk before thee in all their heart.\end{Resp}
\begin{Resp}\Rsign O God of Israel, bless Thou, and hallow this house for ever.\end{Resp}
\switchcolumn*

\LA
\LessonHead{Lectio III}
\switchcolumn
\EN
\LessonHead{Lesson III}
\switchcolumn*

\LA
\Cite{4 Reg 18:9-12}
\switchcolumn
\EN
\Cite{2 Kgs 18:9-12}
\switchcolumn*

\LA
\noindent Anno quarto regis Ezechíæ, qui erat annus séptimus Osée fílii Ela regis Israël, ascéndit Salmánasar rex Assyriórum in Samaríam, et oppugnávit eam, et cepit, nam post annos tres, anno sexto Ezechíæ, id est nono anno Osée regis Israël, capta est Samaría. Et tránstulit rex Assyriórum Israël in Assýrios, collocavítque eos in Hala et in Habor flúviis Gozan, in civitátibus Medórum, quia non audiérunt vocem Dómini Dei sui, sed prætergréssi sunt pactum ejus; ómnia, quæ præcéperat Móyses servus Dómini, non audiérunt, neque fecérunt.\par
\switchcolumn
\EN
\noindent In the fourth year of king Ezechias, which was the seventh year of Osee the son of Ela king of Israel, Salmanasar king of the Assyrians came up to Samaria, and besieged it, And took it. For after three years, in the sixth year of Ezechias, that is, in the ninth year of Osee king of Israel, Samaria was taken: And the king of the Assyrians carried away Israel into Assyria, and placed them in Hale, and in Habor by the rivers of Gozan in the cities of the Medes: Because they hearkened not to the voice of the Lord their God, but transgressed his covenant: all that Moses the servant of the Lord commanded, they would not hear nor do.\par
\switchcolumn*

\LA
\begin{Resp}\Rsign Audi, Dómine, hymnum et oratiónem, quam servus tuus orat coram te hódie, ut sint óculi tui apérti, et aures tuæ inténtæ, \Star{} Super domum istam die ac nocte.\end{Resp}
\begin{Resp}\Vsign Réspice, Dómine, de sanctuário tuo, et de excélso cælórum habitáculo.\end{Resp}
\begin{Resp}\Rsign Super domum istam die ac nocte.\end{Resp}
\begin{Resp}\Vsign Glória Patri, et Fílio, \Star{} et Spirítui Sancto.\end{Resp}
\begin{Resp}\Rsign Super domum istam die ac nocte.\end{Resp}
\switchcolumn
\EN
\begin{Resp}\Rsign Hearken, O Lord, unto the cry and to the prayer which thy servant prayeth before thee today, that thine eyes may be open and thine ears attend; \Star{} Toward this house day and night.\end{Resp}
\begin{Resp}\Vsign Look down from thine high and holy place, O Lord, even from heaven thy dwelling.\end{Resp}
\begin{Resp}\Rsign Toward this house, day and night.\end{Resp}
\begin{Resp}\Vsign Glory be to the Father, and to the Son, \Star{} and to the Holy Ghost.\end{Resp}
\begin{Resp}\Rsign Toward this house, day and night.\end{Resp}
\switchcolumn*

\end{paracol}

\Office{Dominica XI Post Pentecosten}{Eleventh Sunday after Pentecost}{Semiduplex}
\addcontentsline{toc}{subsection}{Dominica XI Post Pentecosten \textit{— Eleventh Sunday after Pentecost}}
\begin{paracol}{2}
\LA
\RefLine{Lectiones I–VI: de Scriptura occurrente.}
\switchcolumn
\EN
\RefLine{Lessons I–VI: from the Scripture of the day.}
\switchcolumn*

\LA
\LessonHead{Lectio I}
\switchcolumn
\EN
\LessonHead{Lesson I}
\switchcolumn*

\LA
\LessonTitle{De libro quarto Regum}
\switchcolumn
\EN
\LessonTitle{Lesson from the second book of Kings}
\switchcolumn*

\LA
\Cite{4 Reg 20:1-3}
\switchcolumn
\EN
\Cite{2 Kgs 20:1-3}
\switchcolumn*

\LA
\noindent In diébus illis ægrotáret Ezechías usque ad mortem, et venit ad eum Isaías fílius Amos prophéta dixítque ei: Hæc dicit Dóminus Deus: Prǽcipe dómui tuæ, moriéris enim tu et non vives. Qui convértit fáciem suam ad paríetem et orávit Dóminum, dicens: Obsecro, Dómine: meménto, quæso, quómodo ambuláverim coram te in veritáte et in corde perfécto et quod plácitum est coram te fécerim. Flevit ítaque Ezechías fletu magno.\par
\switchcolumn
\EN
\noindent In those days Ezechias was sick unto death: and Isaias the son of Amos the prophet came and said to him: Thus saith the Lord God: Give charge concerning thy house, for thou shalt die, and not live. And he turned his face to the wall, and prayed to the Lord, saying: I beseech thee, O Lord, remember how I have walked before thee in truth, and with a perfect heart, and have done that which is pleasing before thee. And Ezechias wept with much weeping.\par
\switchcolumn*

\LA
\begin{Resp}\Rsign Præparáte corda vestra Dómino, et servíte illi soli: \Star{} Et liberábit vos de mánibus inimicórum vestrórum.\end{Resp}
\begin{Resp}\Vsign Convertímini ad eum in toto corde vestro, et auférte deos aliénos de médio vestri.\end{Resp}
\begin{Resp}\Rsign Et liberábit vos de mánibus inimicórum vestrórum.\end{Resp}
\switchcolumn
\EN
\begin{Resp}\Rsign Prepare your hearts unto the Lord, and serve Him Only \Star{} And He will deliver you out of the hand of your enemies.\end{Resp}
\begin{Resp}\Vsign Return unto Him with all your hearts, and put away the strange gods from among you.\end{Resp}
\begin{Resp}\Rsign And He will deliver you out of the hand of your enemies.\end{Resp}
\switchcolumn*

\LA
\LessonHead{Lectio II}
\switchcolumn
\EN
\LessonHead{Lesson II}
\switchcolumn*

\LA
\Cite{4 Reg 20:4-7}
\switchcolumn
\EN
\Cite{2 Kgs 20:4-7}
\switchcolumn*

\LA
\noindent Et ántequam egrederétur Isaías médiam partem átrii, factus est sermo Dómini ad eum, dicens: Revértere et dic Ezechíæ duci pópuli mei: Hæc dicit Dóminus Deus David patris tui: Audívi oratiónem tuam et vidi lácrimas tuas et ecce sanávi te: die tértio ascéndes templum Dómini, et addam diébus tuis quíndecim annos; sed et de manu regis Assyriórum liberábo te et civitátem hanc et prótegam urbem istam propter me et propter David servum meum. Dixítque Isaías: Afférte massam ficórum. Quam cum attulíssent et posuíssent super ulcus ejus, curátus est.\par
\switchcolumn
\EN
\noindent And before Isaias was gone out of the middle of the court, the word of the Lord came to him, saying: Go back, and tell Ezechias the captain of my people: Thus saith the Lord the God of David thy father: I have heard thy prayer, and I have seen thy tears: and behold I have healed thee; on the third day thou shalt go up to the temple of the Lord. And I will add to thy days fifteen years: and I will deliver thee and this city out of the hand of the king of the Assyrians, and I will protect this city for my own sake, and for David my servant's sake. And Isaias said: Bring me a lump of figs. And when they had brought it, and laid it upon his boil. he was healed.\par
\switchcolumn*

\LA
\begin{Resp}\Rsign Deus ómnium exaudítor est: ipse misit Angelum suum, et tulit me de óvibus patris mei: \Star{} Et unxit me unctióne misericórdiæ suæ.\end{Resp}
\begin{Resp}\Vsign Dóminus, qui erípuit me de ore leónis, et de manu béstiæ liberávit me.\end{Resp}
\begin{Resp}\Rsign Et unxit me unctióne misericórdiæ suæ.\end{Resp}
\switchcolumn
\EN
\begin{Resp}\Rsign God, Which heareth all, even He sent His Angel, and took me from keeping my father's sheep, and \Star{} Anointed me with the oil of His mercy.\end{Resp}
\begin{Resp}\Vsign The Lord That delivered me out of the mouth of the lion, and out of the paw of the bear\end{Resp}
\begin{Resp}\Rsign And anointed me with the oil of His mercy.\end{Resp}
\switchcolumn*

\LA
\LessonHead{Lectio III}
\switchcolumn
\EN
\LessonHead{Lesson III}
\switchcolumn*

\LA
\Cite{4 Reg 20:8-11}
\switchcolumn
\EN
\Cite{2 Kgs 20:8-11}
\switchcolumn*

\LA
\noindent Díxerat autem Ezechías ad Isaíam: Quod erit signum quia Dóminus me sanábit, et quia ascensúrus sum die tértia templum Dómini? Cui ait Isaías: Hoc erit signum a Dómino, quod factúrus sit Dóminus sermónem quem locútus est: vis ut ascéndat umbra decem líneis, an ut revertátur tótidem grádibus? Et ait Ezechías: Fácile est umbram créscere decem líneis, nec hoc volo ut fiat; sed ut revertátur retrórsum decem grádibus. Invocávit ítaque Isaías prophéta Dóminum, et redúxit umbram per líneas, quibus jam descénderat in horológio Achaz, retrórsum decem grádibus.\par
\switchcolumn
\EN
\noindent And Ezechias had said to Isaias: What shall be the sign that the Lord will heal me, and that I shall go up to the temple of the Lord the third day? And Isaias said to him: This shall be the sign from the Lord, that the Lord will do the word which he hath spoken: Wilt thou that the shadow go forward ten lines, or that it go back so many degrees? And Ezechias said: It is an easy matter for the shadow to go forward ten lines: and I do not desire that this be done, but let it return back ten degrees. And Isaias the prophet called upon the Lord, and he brought the shadow ten degrees backwards by the lines, by which it had already gone down in the dial of Achaz.\par
\switchcolumn*

\LA
\begin{Resp}\Rsign Dóminus, qui erípuit me de ore leónis, et de manu béstiæ liberávit me, \Star{} Ipse me erípiet de mánibus inimicórum meórum.\end{Resp}
\begin{Resp}\Vsign Misit Deus misericórdiam suam et veritátem suam: ánimam meam erípuit de médio catulórum leónum.\end{Resp}
\begin{Resp}\Rsign Ipse me erípiet de mánibus inimicórum meórum.\end{Resp}
\begin{Resp}\Vsign Glória Patri, et Fílio, \Star{} et Spirítui Sancto.\end{Resp}
\begin{Resp}\Rsign Ipse me erípiet de mánibus inimicórum meórum.\end{Resp}
\switchcolumn
\EN
\begin{Resp}\Rsign The Lord That delivered me out of the mouth of the lion, and out of the paw of the bear \Star{} He will deliver me out of the hand of mine enemies.\end{Resp}
\begin{Resp}\Vsign God hath sent forth His mercy and His truth, and delivered my soul from among the lion's whelps.\end{Resp}
\begin{Resp}\Rsign He will deliver me out of the hand of mine enemies.\end{Resp}
\begin{Resp}\Vsign Glory be to the Father, and to the Son, \Star{} and to the Holy Ghost.\end{Resp}
\begin{Resp}\Rsign He will deliver me out of the hand of mine enemies.\end{Resp}
\switchcolumn*

\LA
\LessonHead{Lectio IV}
\switchcolumn
\EN
\LessonHead{Lesson IV}
\switchcolumn*

\LA
\LessonTitle{Ex Expositióne sancti Hierónymi Presbýteri in Isaíam Prophétam}
\switchcolumn
\EN
\LessonTitle{From the Exposition of the Prophet Isaiah written by St. Jerome, Priest at Bethlehem.}
\switchcolumn*

\LA
\Source{Lib. 11 in Isaíæ cap. 38}
\switchcolumn
\EN
\Source{Bk. ii. on Is. xxxviii.}
\switchcolumn*

\LA
\noindent Ne elevarétur cor Ezechíæ post incredíbiles triúmphos et de média captivitáte victóriam, infirmitáte córporis sui visitátur, et audit se esse moritúrum; ut convérsus ad Dóminum flectat senténtiam ejus. Quod quidem et in Jona prophéta légimus, et in comminatiónibus contra David: quæ dicúntur futúra, nec facta sunt, non Deo mutánte senténtiam, sed provocánte humánum genus ad notítiam sui. Dóminus enim pœ́nitens est super malítiis. Convertítque Ezechías fáciem suam ad paríetem, quia ad templum ire non póterat. Ad paríetem autem templi, juxta quod Sálomon palátium exstrúxerat; vel absolúte ad paríetem, ne lácrimas suas assidéntibus ostentáre viderétur.\par
\switchcolumn
\EN
\noindent Lest the heart of Hezekiah should be puffed up by his strange and unlooked for triumphs, and by his victory when he was but a prisoner, he was visited by bodily weakness, and told that he was to die that he might betake himself to the Lord, and turn Him from carrying out the sentence. We read of a like case in the history of the Prophet Jonah (ii. 4, 10.) And in regard to the threatening made against David (2 Kgs (Sam.) xxiv. 16) - when punishments were foretold which were not brought to pass. This is not because that God is a Being capable of changing His mind, but because He willeth to mankind to know Him, how that “He [is gracious and merciful, slow to anger, and of great kindness, and] repenteth Him of the evil.” (Joel ii. 13.) Hezekiah turned his face unto the wall, not being able to go up to the Temple. This may either mean that he turned towards the wall of the Temple, hard by which Solomon had built a palace, or simply, that he turned his face to the wall, so as not to parade his tears before his attendants.\par
\switchcolumn*

\LA
\begin{Resp}\Rsign Percússit Saul mille, et David decem míllia: \Star{} Quia manus Dómini erat cum illo: percússit Philisthǽum, et ábstulit oppróbrium ex Israël.\end{Resp}
\begin{Resp}\Vsign Nonne iste est David, de quo canébant in choro, dicéntes: Saul percússit mille, et David decem míllia?\end{Resp}
\begin{Resp}\Rsign Quia manus Dómini erat cum illo: percússit Philisthǽum, et ábstulit oppróbrium ex Israël.\end{Resp}
\switchcolumn
\EN
\begin{Resp}\Rsign Saul hath slain his thousands, and David his ten thousands. \Star{} Because the hand of the Lord was with him, he smote the Philistine, and took away the reproach from Israel.\end{Resp}
\begin{Resp}\Vsign Is not this David? Did they not sing one to another of him in dances, saying: Saul hath slain his thousands, and David his ten thousands?\end{Resp}
\begin{Resp}\Rsign Because the hand of the Lord was with him, he smote the Philistine, and took away the reproach from Israel.\end{Resp}
\switchcolumn*

\LA
\LessonHead{Lectio V}
\switchcolumn
\EN
\LessonHead{Lesson V}
\switchcolumn*

\LA
\noindent Audiénsque se esse moritúrum, non precátur vitam et annos plúrimos, sed in Dei judício, quid velit præstáre, dimíttit. Nóverat enim idcírco Deo placuísse Salomónem, quod annos vitæ non petíerit amplióres; sed itúrus ad Dóminum narrat ópera sua, quómodo ambuláverit coram eo in veritáte et in corde perfécto. Felix consciéntia, quæ afflictiónis témpore bonórum óperum recordátur. Beáti enim mundo corde: quóniam ipsi Deum vidébunt. Et quómodo álibi scríbitur: Quis gloriábitur purum habére se cor? Quod ita sólvitur: perfectiónem cordis in eo nunc dici, quod idóla destrúxerit, templi valvas aperúerit, serpéntem ǽneum comminúerit, et cétera fécerit quæ Scriptúra commémorat.\par
\switchcolumn
\EN
\noindent Having been told that he was about to die, he prayed not for life and many years, but left it to God to do as in His good judgment He was pleased to will. He knew how this had pleased God on the part of Solomon. (3 (1) Kings iii. 11.) So, when he betook him to the Lord, he only made mention of his works, how he had walked before Him in truth, and with a perfect heart. Happy is he whose conscience in the hour of affliction can assure him of good works. Yea, “blessed are the pure in heart, for they shall see God” (Matth. v. 8.) It is indeed written in another place “Who can say, I have made my heart clean, I am pure from my sin?” (Prov. xx. 9.) How then could Hezekiah say that he had walked with a perfect heart But the answer is, that by this is meant that he had destroyed the idols, opened the doors of the Temple, broken in pieces the brazen serpent, and done the rest of the things whereof the Scripture maketh mention.\par
\switchcolumn*

\LA
\begin{Resp}\Rsign Montes Gélboë, nec ros nec plúvia véniant super vos, \Star{} Ubi cecidérunt fortes Israël.\end{Resp}
\begin{Resp}\Vsign Omnes montes, qui estis in circúitu ejus, vísitet Dóminus: a Gélboë autem tránseat.\end{Resp}
\begin{Resp}\Rsign Ubi cecidérunt fortes Israël.\end{Resp}
\switchcolumn
\EN
\begin{Resp}\Rsign Ye mountains of Gilboa, let there be no dew, neither let there be rain upon you \Star{} For there are the mighty of Israel fallen.\end{Resp}
\begin{Resp}\Vsign All ye mountains that stand round about, the Lord look upon you but let Him pass by Gilboa\end{Resp}
\begin{Resp}\Rsign For there are the mighty of Israel fallen.\end{Resp}
\switchcolumn*

\LA
\LessonHead{Lectio VI}
\switchcolumn
\EN
\LessonHead{Lesson VI}
\switchcolumn*

\LA

\switchcolumn
\EN
\LessonTitle{“And Hezekiah wept sore.” He had then no children, and it seemed as though the promise which God had made unto David, [that “his seed should endure for ever, and his throne as the sun before” Him, (Ps. lxxxviii. 29,)] was about to fail in his own death. It is written that “Manasseh was twelve years old when he began to reign,” (xxi. 1) - whence it is evident that Hezekiah begat him not till after three years of his new lease of life. Sore therefore wept he, when all hope was torn from him that the Messiah should spring from his seed. Others again remark that he wept sore, since death terrifieth sometimes even the saints, since they know not what sentence is about to be pronounced upon them, and what place shall be allotted them in the inscrutable judgment.}
\switchcolumn*

\LA
\noindent Flevit autem fletu magno, propter promissiónem Dómini ad David, quam vidébat in sua morte peritúram. Eo enim témpore Ezechías fílios non habébat; nam post mortem ejus Manásses cum duódecim esset annórum, regnáre cœpit in Judǽa: ex quo perspícuum est, post tértium annum concéssæ vitæ Manássen esse generátum. Ergo iste omnis est fletus, quod desperábat Christum de suo sémine nascitúrum. Alii ásserunt, quamvis sanctos viros morte terréri, propter incértum judícii et ignoratiónem senténtiæ Dei, quam sedem habitúri sint.\par
\switchcolumn
\EN

\switchcolumn*

\LA
\begin{Resp}\Rsign Ego te tuli de domo patris tui, dicit Dóminus, et pósui te páscere gregem pópuli mei: \Star{} Et fui tecum in ómnibus ubicúmque ambulásti, firmans regnum tuum in ætérnum.\end{Resp}
\begin{Resp}\Vsign Fecíque tibi nomen grande, juxta nomen magnórum, qui sunt in terra: et réquiem dedi tibi ab ómnibus inimícis tuis.\end{Resp}
\begin{Resp}\Rsign Et fui tecum in ómnibus ubicúmque ambulásti, firmans regnum tuum in ætérnum.\end{Resp}
\begin{Resp}\Vsign Glória Patri, et Fílio, \Star{} et Spirítui Sancto.\end{Resp}
\begin{Resp}\Rsign Et fui tecum in ómnibus ubicúmque ambulásti, firmans regnum tuum in ætérnum.\end{Resp}
\switchcolumn
\EN
\begin{Resp}\Rsign Thus saith the Lord: I took thee out of thy father's house, and appointed thee to be ruler over My people, over Israel. \Star{} And I was with thee whithersoever thou wentest, to establish thy kingdom for ever.\end{Resp}
\begin{Resp}\Vsign And I have made thee a great name, like unto the name of the great men that are in the earth, and have caused thee to rest from all thine enemies.\end{Resp}
\begin{Resp}\Rsign And I was with thee whithersoever thou wentest, to establish thy kingdom for ever.\end{Resp}
\begin{Resp}\Vsign Glory be to the Father, and to the Son, \Star{} and to the Holy Ghost.\end{Resp}
\begin{Resp}\Rsign And I was with thee whithersoever thou wentest, to establish thy kingdom for ever.\end{Resp}
\switchcolumn*

\LA
\LessonHead{Lectio VII}
\switchcolumn
\EN
\LessonHead{Lesson VII}
\switchcolumn*

\LA
\LessonTitle{Léctio sancti Evangélii secúndum Marcum}
\switchcolumn
\EN
\LessonTitle{From the Holy Gospel according to Mark}
\switchcolumn*

\LA
\Cite{Marc 7:31-37}
\switchcolumn
\EN
\Cite{Mark 7:31-37}
\switchcolumn*

\LA
\noindent In illo témpore: Exiens Jesus de fínibus Tyri venit per Sidónem ad mare Galilǽæ inter médios fines Decapóleos. Et réliqua.\par
\switchcolumn
\EN
\noindent At that time: Jesus, departing from the coasts of Tyre and Sidon, came unto the sea of Galilee, through the midst of the coasts of Decapolis. And so on.\par
\switchcolumn*

\LA
\LessonTitle{Homilía sancti Gregórii Papæ}
\switchcolumn
\EN
\LessonTitle{Homily by Pope St. Gregory the Great}
\switchcolumn*

\LA
\Source{Homilia 10, liber 1 in Ezech., ante medium}
\switchcolumn
\EN
\Source{Hom. x. Bk. i. on Ezekiel}
\switchcolumn*

\LA
\noindent Quid est quod creátor ómnium Deus, cum surdum et mutum sanáre voluísset, in aures illíus suos dígitos misit, et éxspuens linguam ejus tétigit? Quid per dígitos Redemptóris, nisi dona Sancti Spíritus designántur? Unde cum in álio loco ejecísset dæmónium, dixit: Si in dígito Dei eício dæmónia, profécto pervénit in vos regnum Dei. Qua de re per Evangelístam álium dixísse descríbitur: Si ego in Spíritu Dei eício dǽmones, ígitur pervénit in vos regnum Dei. Ex quo utróque loco collígitur, quia dígitus Spíritus vocátur. Dígitos ergo in aurículas míttere, est per dona Spíritus Sancti mentem surdi ad obediéndum aperíre.\par
\switchcolumn
\EN
\noindent What signifieth it that when God, the Maker of all, would heal a deaf and dumb man, “He put His Fingers into his ears, and He spit, and touched his tongue.” What is figured by the Fingers of the Redeemer but the gifts of the Holy Ghost? Hence it is written in another place (Luke xi. 20) that after He had cast out an evil spirit, He said: “If I with the finger of God cast out devils, no doubt the kingdom of God is come upon you.” Which words are thus given by another Evangelist (Matth. xii. 28): “If I cast out devils by the Spirit of God, then the kingdom of God is come unto you.” By setting these two passages together we see that the Spirit is called the Finger. For our Lord, then, to put His Fingers into the deaf man's ears was by the gift of the Holy Spirit to enlighten his dark mind unto obedience.\par
\switchcolumn*

\LA
\begin{Resp}\Rsign Dómine, Pater et Deus vitæ meæ, ne derelínquas me in cogitátu malígno: extolléntiam oculórum meórum ne déderis mihi, et desidérium malígnum avérte a me, Dómine; aufer a me concupiscéntiam, \Star{} Et ánimo irreverénti et infruníto ne tradas me, Dómine.\end{Resp}
\begin{Resp}\Vsign Ne derelínquas me, Dómine, ne accréscant ignorántiæ meæ, nec multiplicéntur delícta mea.\end{Resp}
\begin{Resp}\Rsign Et ánimo irreverénti et infruníto ne tradas me, Dómine.\end{Resp}
\switchcolumn
\EN
\begin{Resp}\Rsign O Lord, Father and God of my life, leave me not to evil counsels; give me not a proud look, but turn away from me an haughty mind, O Lord turn away from me concupiscence, \Star{} And give me not over unto an impudent and froward mind, O Lord!\end{Resp}
\begin{Resp}\Vsign Leave me not, O Lord, lest mine ignorance increase, and my sins abound.\end{Resp}
\begin{Resp}\Rsign And give me not over unto an impudent and froward mind, O Lord.\end{Resp}
\switchcolumn*

\LA
\LessonHead{Lectio VIII}
\switchcolumn
\EN
\LessonHead{Lesson VIII}
\switchcolumn*

\LA
\noindent Quid est vero, quod éxspuens linguam ejus tétigit? Salíva nobis est ex ore Redemptóris, accépta sapiéntia in elóquio divíno. Salíva quippe ex cápite défluit in ore. Ea ergo sapiéntia, quæ ipse est, dum lingua nostra tángitur, mox ad prædicatiónis verba formátur. Qui suspíciens in cælum, ingémuit: non quod ipse necessárium gémitum habéret, qui dabat quod postulábat; sed nos ad eum gémere, qui cælo prǽsidet, dócuit: ut et aures nostræ per donum Spíritus Sancti aperíri, et lingua per salívam oris, id est, per sciéntiam divínæ locutiónis, solvi débeat ad verba prædicatiónis.\par
\switchcolumn
\EN
\noindent That signifieth it also that “He spit and touched his tongue.” We receive spittle out of the Redeemer's mouth upon our tongues when we receive wisdom to speak God's truth. Spittle is a secretion of the head which floweth into the mouth. And so, that wisdom, which is Himself, the great Head of His Church, as soon as it hath touched our tongue, doth straightway take the form of preaching. “And looking up to heaven, He sighed,” not that He had any need to sigh, Who gave whatsoever He asked, but that He was fain to teach us to look up and sigh toward Him Whose throne is in heaven, confessing our need, that our ears should be opened by the gift of the Holy Spirit, and our tongue loosed by the spittle of our Saviour's Mouth, that is, by knowledge of His Divine Word, before we can use it to preach to others.\par
\switchcolumn*

\LA
\begin{Resp}\Rsign Duo Séraphim clamábant alter ad álterum: \Star{} Sanctus, sanctus, sanctus Dóminus Deus Sábaoth: * Plena est omnis terra glória ejus.\end{Resp}
\begin{Resp}\Vsign Tres sunt qui testimónium dant in cælo: Pater, Verbum, et Spíritus Sanctus: et hi tres unum sunt.\end{Resp}
\begin{Resp}\Rsign Sanctus, sanctus, sanctus Dóminus Deus Sábaoth.\end{Resp}
\begin{Resp}\Vsign Glória Patri, et Fílio, \Star{} et Spirítui Sancto.\end{Resp}
\begin{Resp}\Rsign Plena est omnis terra glória ejus.\end{Resp}
\switchcolumn
\EN
\begin{Resp}\Rsign One Seraph cried unto another: \Star{} Holy, Holy, Holy is the Lord God of hosts; the whole earth is full of His glory.\end{Resp}
\begin{Resp}\Vsign There are Three That bear record in heaven: the Father, the Word, and the Holy Ghost, and these Three are One.\end{Resp}
\begin{Resp}\Rsign Holy, Holy, Holy is the Lord God of hosts.\end{Resp}
\begin{Resp}\Vsign Glory be to the Father, and to the Son, \Star{} and to the Holy Ghost.\end{Resp}
\begin{Resp}\Rsign The whole earth is full of His glory.\end{Resp}
\switchcolumn*

\LA
\LessonHead{Lectio IX}
\switchcolumn
\EN
\LessonHead{Lesson IX}
\switchcolumn*

\LA

\switchcolumn
\EN
\LessonTitle{“And He said unto him: Ephphatha, that is, be opened. And straightway his ears were opened, and the string of his tongue was loosed.” Herein we must remark the command, “Be opened” was addressed to the deaf ears, but the tongue also was immediately loosed. Just so, when the ears of a man's heart have been opened to learn the obedience of faith, the string of his tongue also is thereupon loosed, that he may exhort others to do the good things which himself doth. It is well added “And he spake plain.” He only doth well preach obedience to others who hath first learnt himself to obey.}
\switchcolumn*

\LA
\noindent Cui mox, Ephphétha, id est, Adaperíre dícitur: et statim apértæ sunt aures ejus, et solútum est vínculum linguæ ejus. Qua in re notándum est, quia propter clausas aures dictum est, Adaperíre. Sed cui aures cordis ad obediéndum apértæ fúerint, ex subsequénti procul dúbio étiam linguæ ejus vínculum sólvitur; ut bona quæ ipse fécerit, étiam faciénda áliis loquátur. Ubi bene ádditur: Et loquebátur recte. Ille enim recte lóquitur, qui prius obediéndo fécerit quæ loquéndo ádmonet esse faciénda.\par
\switchcolumn
\EN

\switchcolumn*

\LA
\TeDeum{Te Deum}
\switchcolumn
\EN
\TeDeum{Te Deum}
\switchcolumn*

\end{paracol}

\Office{Feria secunda infra Hebdomadam XI post Octavam Pentecostes}{Monday of the Eleventh Week after Pentecost}{Feria}
\addcontentsline{toc}{subsection}{Feria secunda infra Hebdomadam XI post Octavam Pentecostes \textit{— Monday of the Eleventh Week after Pentecost}}
\begin{paracol}{2}
\LA
\LessonHead{Lectio I}
\switchcolumn
\EN
\LessonHead{Lesson I}
\switchcolumn*

\LA
\LessonTitle{De libro quarto Regum}
\switchcolumn
\EN
\LessonTitle{Lesson from the second book of Kings}
\switchcolumn*

\LA
\Cite{4 Reg 22:1-5}
\switchcolumn
\EN
\Cite{2 Kgs 22:1-5}
\switchcolumn*

\LA
\noindent Octo annórum erat Josías cum regnáre cœpísset : trigínta et uno anno regnávit in Jerúsalem. Nomen matris ejus Idida fília Hadaía de Bésecath. Fecítque quod plácitum erat coram Dómino et ambulávit per omnes vias David patris sui: non declinávit ad déxteram, sive ad sinistram. Anno autem octavo décimo regis Josíæ misit rex Saphan fílium Aslía fílii Méssulam scribam templi Dómini, dicens ei: Vade ad Helcíam sacerdotem magnum, ut conflétur pecúnia, quæ illáta est in templum Dómini, quam collegérunt janitores templi a pópulo, detúrque febris per præpósitos domus Dómini; qui et distríbuant eam his qui operántur in templo Domini.\par
\switchcolumn
\EN
\noindent Josias was eight years old when he began to reign: he reigned one and thirty years in Jerusalem: the name of his mother was Idida, the daughter of Hadaia, of Besecath. And he did that which was right in the sight of the Lord, and walked in all the ways of David his father: he turned not aside to the right hand, or to the left. And in the eighteenth year of king Josias, the king sent Saphan the son of Assia, the son of Messulam, the scribe of the temple of the Lord, saying to him: Go to Helcias the high priest, that the money may be put together which is brought into the temple of the Lord, which the doorkeepers of the temple have gathered of the people. And let it be given to the workmen by the overseers of the house of the Lord: and lot them distribute it to those that work in the temple of the Lord, to repair the temple:\par
\switchcolumn*

\LA
\begin{Resp}\Rsign Recordáre, Dómine, testaménti tui, et dic Angelo percutiénti: Cesset jam manus tua, \Star{} Ut non desolétur terra, et ne perdas omnem ánimam vivam.\end{Resp}
\begin{Resp}\Vsign Ego sum qui peccávi, ego qui iníque egi: isti qui oves sunt, quid fecérunt? Avertátur, óbsecro, furor tuus, Dómine, a pópulo tuo.\end{Resp}
\begin{Resp}\Rsign Ut non desolétur terra, et ne perdas omnem ánimam vivam.\end{Resp}
\switchcolumn
\EN
\begin{Resp}\Rsign Remember, O Lord, thy covenant, and say unto the destroying Angel: Stay now thine hand; \Star{} That the land be not utterly laid waste, and that thou destroy not every living soul.\end{Resp}
\begin{Resp}\Vsign Even I it is that have sinned, and done evil indeed; but these sheep, what have they done? Let thine anger, I pray thee, O Lord, be turned away from thy people.\end{Resp}
\begin{Resp}\Rsign That the land be not utterly laid waste, and that Thou destroy not every living soul.\end{Resp}
\switchcolumn*

\LA
\LessonHead{Lectio II}
\switchcolumn
\EN
\LessonHead{Lesson II}
\switchcolumn*

\LA
\Cite{4 Reg 22:8-10}
\switchcolumn
\EN
\Cite{2 Kgs 22:8-10}
\switchcolumn*

\LA
\noindent Dixit autem Helcías pontifex ad Saphan scribam: Librum legis réperi in domo Dómini. Dedítque Helcías volúmen Saphan, qui et legit illud. Venit quoque Saphan scriba ad regem, et renuntiávit ei quod præceperat, et ait: Conflavérunt servi tui pecúniam, quæ repérta est in domo Dómini, et dedérunt ut distribuerétur fabris a præfectis óperum templi Dómini. Narrávit quoque Saphan scriba regi, dicens: Librum dedit mihi Helcías sacérdos.\par
\switchcolumn
\EN
\noindent And Helcias the high priest said to Saphan the scribe: I have found the book of the law in the house of the Lord: and Helcias gave the book to Saphan, and he read it. And Saphan the scribe came to the king, and brought him word again concerning that which he had commanded, and said: thy servants have gathered together the money that was found in the house of the Lord, and they have given it to be distributed to the workmen, by the overseers of the works of the temple of the Lord. And Saphan the scribe told the king, saying: Helcias the priest hath delivered to me a book.\par
\switchcolumn*

\LA
\begin{Resp}\Rsign Exaudísti, Dómine, oratiónem servi tui, ut ædificárem templum nómini tuo: \Star{} Bénedic et sanctífica domum istam in sempitérnum, Deus Israël.\end{Resp}
\begin{Resp}\Vsign Dómine, qui custódis pactum cum servis tuis, qui ámbulant coram te in toto corde suo.\end{Resp}
\begin{Resp}\Rsign Bénedic et sanctífica domum istam in sempitérnum, Deus Israël.\end{Resp}
\switchcolumn
\EN
\begin{Resp}\Rsign O Lord, Thou hast hearkened unto the prayer of thy servant, that I might build a temple unto thy Name, \Star{} O God of Israel, bless Thou, and hallow this house for ever.\end{Resp}
\begin{Resp}\Vsign O Lord, Who keepest covenant with thy servants that walk before thee in all their heart.\end{Resp}
\begin{Resp}\Rsign O God of Israel, bless Thou, and hallow this house for ever.\end{Resp}
\switchcolumn*

\LA
\LessonHead{Lectio III}
\switchcolumn
\EN
\LessonHead{Lesson III}
\switchcolumn*

\LA
\Cite{4 Reg 22:10-13}
\switchcolumn
\EN
\Cite{2 Kgs 22:10-13}
\switchcolumn*

\LA
\noindent Quem cum legísset Saphan coram rege, et audísset rex verba libri legis Dómini, scidit vestiménta sua. Et præcépit Helcíæ sacerdóti, et Ahícam fílio Saphan, et Achobor fílio Micha, et Saphan scribæ, et Asaíæ servo regis, dicens: Ite et consúlite Dóminum super me, et super pópulo, et super omni Juda, de verbis volúminis istius, quod invéntum est; magna enim ira Dómini succénsa est contra nos, quia non audiérunt patres nostri verba libri hujus, ut fácerent omne quod scriptum est nobis.\par
\switchcolumn
\EN
\noindent And when Saphan had read it before the king, And the king had heard the words of the law of the Lord, he rent his garments. And he commanded Helcias the priest, and Ahicam the son of Saphan, and Achobor the son of Micha, and Saphan the scribe, and Asaia the king's servant, saying: Go and consult the Lord for me, and for the people, and for all Juda, concerning the words of this book which is found: for the great wrath of the Lord is kindled against us, because our fathers have not hearkened to the words of this book, to do all that is written for us.\par
\switchcolumn*

\LA
\begin{Resp}\Rsign Audi, Dómine, hymnum et oratiónem, quam servus tuus orat coram te hódie, ut sint óculi tui apérti, et aures tuæ inténtæ, \Star{} Super domum istam die ac nocte.\end{Resp}
\begin{Resp}\Vsign Réspice, Dómine, de sanctuário tuo, et de excélso cælórum habitáculo.\end{Resp}
\begin{Resp}\Rsign Super domum istam die ac nocte.\end{Resp}
\begin{Resp}\Vsign Glória Patri, et Fílio, \Star{} et Spirítui Sancto.\end{Resp}
\begin{Resp}\Rsign Super domum istam die ac nocte.\end{Resp}
\switchcolumn
\EN
\begin{Resp}\Rsign Hearken, O Lord, unto the cry and to the prayer which thy servant prayeth before thee today, that thine eyes may be open and thine ears attend; \Star{} Toward this house day and night.\end{Resp}
\begin{Resp}\Vsign Look down from thine high and holy place, O Lord, even from heaven thy dwelling.\end{Resp}
\begin{Resp}\Rsign Toward this house, day and night.\end{Resp}
\begin{Resp}\Vsign Glory be to the Father, and to the Son, \Star{} and to the Holy Ghost.\end{Resp}
\begin{Resp}\Rsign Toward this house, day and night.\end{Resp}
\switchcolumn*

\end{paracol}

\Office{Feria tertia infra Hebdomadam XI post Octavam Pentecostes}{Tuesday of the Eleventh Week after Pentecost}{Feria}
\addcontentsline{toc}{subsection}{Feria tertia infra Hebdomadam XI post Octavam Pentecostes \textit{— Tuesday of the Eleventh Week after Pentecost}}
\begin{paracol}{2}
\LA
\LessonHead{Lectio I}
\switchcolumn
\EN
\LessonHead{Lesson I}
\switchcolumn*

\LA
\LessonTitle{De libro quarto Regum}
\switchcolumn
\EN
\LessonTitle{Lesson from the second book of Kings}
\switchcolumn*

\LA
\Cite{4 Reg 23:2-3}
\switchcolumn
\EN
\Cite{2 Kgs 23:2-3}
\switchcolumn*

\LA
\noindent Ascéndit rex templum Dómini, et omnes viri Juda, universíque qui habitábant in Jerúsalem cum eo, sacerdótes et prophétæ, et omnis pópulus a parvo usque ad magnum: legítque, cunctis audiéntibus, omnia verba libri fœ́deris, qui invéntus est in domo Domini. Stetítque rex super gradum: et fœ́dus percússit coram Dómino, ut ambulárent post Dóminum, et custodírent præcépta ejus, et testimónia, et cæremónias in omni corde, et in tota ánima, et suscitárent verba fœ́deris hujus, quæ scripta erant in libro illo: Acquievítque pópulus pacto.\par
\switchcolumn
\EN
\noindent And the king went up to the temple of the Lord, and all the men of Juda, and all the inhabitants of Jerusalem with him, the priests and the prophets, and all the people both little and great: and in the hearing of them all he read all the words of the book of the covenant, which was found in the house of the Lord. And the king stood upon the step: and made a covenant with the Lord, to walk after the Lord, and to keep his commandments, and his testimonies and his ceremonies, with all their heart, and with all their soul, and to perform the words of this covenant, which were written in that book: and the people agreed to the covenant.\par
\switchcolumn*

\LA
\begin{Resp}\Rsign Dómine, si convérsus fúerit pópulus tuus, et oráverit ad sanctuárium tuum: \Star{} Tu exáudies de cælo, Dómine, et líbera eos de mánibus inimicórum suórum.\end{Resp}
\begin{Resp}\Vsign Si peccáverit in te pópulus tuus, et convérsus égerit pœniténtiam, veniénsque oráverit in isto loco.\end{Resp}
\begin{Resp}\Rsign Tu exáudies de cælo, Dómine, et líbera eos de mánibus inimicórum suórum.\end{Resp}
\switchcolumn
\EN
\begin{Resp}\Rsign Lord, when thy people shall turn again to thee, and shall pray unto thee in this house; \Star{} Then hear Thou in heaven, O Lord, and deliver them out of the hand of their enemies.\end{Resp}
\begin{Resp}\Vsign If thy people sin against thee, and turn again, and repent, and come and pray unto thee in this house.\end{Resp}
\begin{Resp}\Rsign Then hear Thou in heaven, O Lord, and deliver them out of the hand of their enemies.\end{Resp}
\switchcolumn*

\LA
\LessonHead{Lectio II}
\switchcolumn
\EN
\LessonHead{Lesson II}
\switchcolumn*

\LA
\Cite{4 Reg 23:4-5}
\switchcolumn
\EN
\Cite{2 Kgs 23:4-5}
\switchcolumn*

\LA
\noindent Et præcépit rex Helcíæ pontífici, et sacerdótibus secúndi órdinis, et janitóribus, ut proícerent de templo Dómini omnia vasa, quæ facta fúerant Baal, et in luco, et univérsæ milítiæ cæli: et combússit ea foris Jerúsalem in conválle Cedron, et tulit púlverem eorum in Bethel. Et delévit arúspices quos posúerant reges Juda ad sacrificándum in excélsis per civitátes Juda, et in circúitu Jerúsalem: et eos qui adolébant incénsum Baal, et soli, et lunæ, et duódecim signis, et omni milítiæ cæli.\par
\switchcolumn
\EN
\noindent And the king commanded Helcias the high priest, and the priests of the second order, and the doorkeepers, to cast out of the temple of the Lord all the vessels that had been made for Baal, and for the grove, and for all the host of heaven: and he burnt them without Jerusalem in the valley of Cedron, and he carried the ashes of them to Bethel. And he destroyed the soothsayers, whom the kings of Juda had appointed to sacrifice in the high places in the cities of Juda, and round about Jerusalem: them also that burnt incense to Baal, and to the sun, and to the moon, and to the twelve signs, and to all the host of heaven.\par
\switchcolumn*

\LA
\begin{Resp}\Rsign Factum est, dum tólleret Dóminus Elíam per túrbinem in cælum, \Star{} Eliséus clamábat, dicens: Pater mi, pater mi, currus Israël, et auríga ejus.\end{Resp}
\begin{Resp}\Vsign Cumque pérgerent, et incedéntes sermocinaréntur, ecce currus ígneus et equi ígnei divisérunt utrúmque, et ascéndit Elías per túrbinem in cælum.\end{Resp}
\begin{Resp}\Rsign Eliséus clamábat, dicens: Pater mi, pater mi, currus Israël, et auríga ejus.\end{Resp}
\switchcolumn
\EN
\begin{Resp}\Rsign And it came to pass when the Lord would take up Elijah into heaven by a whirlwind \Star{} Elisha cried, and said: My father, my father, the chariot of Israel, and the horsemen thereof.\end{Resp}
\begin{Resp}\Vsign And as they still went on, and talked, behold, there appeared a chariot of fire, and horses of fire, and parted them both asunder, and Elijah went up by a whirlwind into heaven.\end{Resp}
\begin{Resp}\Rsign Elisha cried, and said: My father, my father, the chariot of Israel, and the horsemen thereof.\end{Resp}
\switchcolumn*

\LA
\LessonHead{Lectio III}
\switchcolumn
\EN
\LessonHead{Lesson III}
\switchcolumn*

\LA
\Cite{4 Reg 23:6-8}
\switchcolumn
\EN
\Cite{2 Kgs 23:6-8}
\switchcolumn*

\LA
\noindent Et efférri fecit lucum de domo Dómini foras Jerúsalem in conválle Cedron, et combússit eum ibi, et redégit in púlverem, et projécit super sepúlcra vulgi. Destrúxit quoque ædículas effeminatórum quæ erant in domo Dómini, pro quibus mulíeres texébant quasi domúnculas luci. Congregavítque omnes sacerdótes de civitátibus Juda, et contaminávit excélsa ubi sacrificábant sacerdótes, de Gábaa usque Bersabée.\par
\switchcolumn
\EN
\noindent And he caused the grove to be carried out from the house of the Lord without Jerusalem to the valley of Cedron, and he burnt it there, and reduced it to dust, and cast the dust upon the graves of the common people. He destroyed also the pavilions of the effeminate, which were in the house of the Lord, for which the women wove as it were little dwellings for the grove. And he gathered together all the priests out of the cities of Juda: and he defiled the high places, where the priests offered sacrifice, from Gabaa to Bersabee: and he broke down the altars of the gates that were in the entering in of the gate of Josue governor of tile city, which was on the left hand of the gate of the city.\par
\switchcolumn*

\LA
\begin{Resp}\Rsign Ego te tuli de domo patris tui, dicit Dóminus, et pósui te páscere gregem pópuli mei: \Star{} Et fui tecum in ómnibus ubicúmque ambulásti, firmans regnum tuum in ætérnum.\end{Resp}
\begin{Resp}\Vsign Fecíque tibi nomen grande, juxta nomen magnórum, qui sunt in terra: et réquiem dedi tibi ab ómnibus inimícis tuis.\end{Resp}
\begin{Resp}\Rsign Et fui tecum in ómnibus ubicúmque ambulásti, firmans regnum tuum in ætérnum.\end{Resp}
\begin{Resp}\Vsign Glória Patri, et Fílio, \Star{} et Spirítui Sancto.\end{Resp}
\begin{Resp}\Rsign Et fui tecum in ómnibus ubicúmque ambulásti, firmans regnum tuum in ætérnum.\end{Resp}
\switchcolumn
\EN
\begin{Resp}\Rsign Thus saith the Lord: I took thee out of thy father's house, and appointed thee to be ruler over My people, over Israel. \Star{} And I was with thee whithersoever thou wentest, to establish thy kingdom for ever.\end{Resp}
\begin{Resp}\Vsign And I have made thee a great name, like unto the name of the great men that are in the earth and have caused thee to rest from all thine enemies.\end{Resp}
\begin{Resp}\Rsign And I was with thee whithersoever thou wentest, to establish thy kingdom for ever.\end{Resp}
\begin{Resp}\Vsign Glory be to the Father, and to the Son, \Star{} and to the Holy Ghost.\end{Resp}
\begin{Resp}\Rsign And I was with thee whithersoever thou wentest, to establish thy kingdom for ever.\end{Resp}
\switchcolumn*

\end{paracol}

\Office{Feria quarta infra Hebdomadam XI post Octavam Pentecostes}{Wednesday of the Eleventh Week after Pentecost}{Feria}
\addcontentsline{toc}{subsection}{Feria quarta infra Hebdomadam XI post Octavam Pentecostes \textit{— Wednesday of the Eleventh Week after Pentecost}}
\begin{paracol}{2}
\LA
\LessonHead{Lectio I}
\switchcolumn
\EN
\LessonHead{Lesson I}
\switchcolumn*

\LA
\LessonTitle{De libro quarto Regum}
\switchcolumn
\EN
\LessonTitle{Lesson from the second book of Kings}
\switchcolumn*

\LA
\Cite{4 Reg 23:24-26}
\switchcolumn
\EN
\Cite{2 Kgs 23:24-26}
\switchcolumn*

\LA
\noindent Pythónes, et haríolos, et figúras idolórum, et immundítias, et abominatiónes, quæ fúerant in terra Juda et Jerúsalem, ábstulit Josías, ut statúeret verba legis, quæ scripta sunt in libro, quem invenit Helcías sacérdos in templo Dómini. Símilis illi non fuit ante eum rex, qui reverterétur ad Dóminum in omni corde suo, et in tota ánima sua, et in univérsa virtúte sua juxta omnem legem Móysi: neque post eum surréxit símilis illi. Verúmtamen non est avérsus Dóminus ab ira furóris sui magni, quo irátus est furor ejus contra Judam propter irritatiónes, quibus provocáverat eum Manásses.\par
\switchcolumn
\EN
\noindent Moreover the diviners by spirits, and soothsayers, and the figures of idols, and the uncleannesses, and the abominations, that had been in the land of Juda, and Jerusalem, Josias took away: that he might perform the words of the law, that were written in the book which Helcias the priest had found in the temple of the Lord. There was no king before him like unto him, that returned to the Lord with all his heart, and with all his soul, and with ail his strength, according to all the law of Moses: neither after him did there arise any like him. But yet the Lord turned not away from the wrath of his great indignation, wherewith his anger was kindled against Juda: because of the provocations, wherewith Manasses had provoked him\par
\switchcolumn*

\LA
\begin{Resp}\Rsign Peccávi super númerum arénæ maris, et multiplicáta sunt peccáta mea: et non sum dignus vidére altitúdinem cæli præ multitúdine iniquitátis meæ: quóniam irritávi iram tuam, \Star{} Et malum coram te feci.\end{Resp}
\begin{Resp}\Vsign Quóniam iniquitátem meam ego cognósco: et delíctum meum contra me est semper, quia tibi soli peccávi.\end{Resp}
\begin{Resp}\Rsign Et malum coram te feci.\end{Resp}
\switchcolumn
\EN
\begin{Resp}\Rsign My sins are many, yea, they are more in number than the sands of the sea; I am not worthy to look up toward heaven because of the multitude of my iniquities; for I have provoked thee to anger; \Star{} And done evil in thy sight.\end{Resp}
\begin{Resp}\Vsign For I acknowledge my transgression, and my sin is ever before me, for against thee only have I sinned.\end{Resp}
\begin{Resp}\Rsign And done evil in thy sight.\end{Resp}
\switchcolumn*

\LA
\LessonHead{Lectio II}
\switchcolumn
\EN
\LessonHead{Lesson II}
\switchcolumn*

\LA
\Cite{4 Reg 23:27-30}
\switchcolumn
\EN
\Cite{2 Kgs 23:27-30}
\switchcolumn*

\LA
\noindent Dixit ítaque Dóminus: Etiam Judam áuferam a fácie mea, sicut ábstuli Israël, et proíciam civitátem hanc, quam elégi, Jerúsalem, et domum de qua dixi: Erit nomen meum ibi. Réliqua autem sermónum Josíæ, et univérsa quæ fecit, nonne hæc scripta sunt in libro verbórum diérum regum Juda? In diébus ejus ascéndit phárao Néchao rex Ægýpti contra regem Assyriórum ad flumen Euphráten, et abiit Josías rex in occúrsum ejus: et occísus est in Magéddo, cum vidísset eum. Et portavérunt eum servi sui mórtuum de Magéddo: et pertulérunt in Jerúsalem, et sepeliérunt eum in sepúlcro suo.\par
\switchcolumn
\EN
\noindent And the Lord said: I will remove Juda also from before my face, as I have removed Israel: and I will cast off this city Jerusalem, which I chose, and the house, of which I said: My name shall be there. 28 Now the rest of the acts of Josias, and all that he did, are they not written in the book of the words of the days of the kings of Juda? In his days Pharao Nechao king of Egypt went up against the king of Assyria to the river Euphrates: and king Josias went to meet him: and was slain at Mageddo, when he had seen him. And his servants carried him dead from Mageddo: and they brought him to Jerusalem, and buried him in Iris own sepulchre.\par
\switchcolumn*

\LA
\begin{Resp}\Rsign Exaudísti, Dómine, oratiónem servi tui, ut ædificárem templum nómini tuo: \Star{} Bénedic et sanctífica domum istam in sempitérnum, Deus Israël.\end{Resp}
\begin{Resp}\Vsign Dómine, qui custódis pactum cum servis tuis, qui ámbulant coram te in toto corde suo.\end{Resp}
\begin{Resp}\Rsign Bénedic et sanctífica domum istam in sempitérnum, Deus Israël.\end{Resp}
\begin{Resp}\Vsign Glória Patri, et Fílio, \Star{} et Spirítui Sancto.\end{Resp}
\begin{Resp}\Rsign Bénedic et sanctífica domum istam in sempitérnum, Deus Israël.\end{Resp}
\switchcolumn
\EN
\begin{Resp}\Rsign O Lord, Thou hast hearkened unto the prayer of thy servant, that I might build a temple unto thy Name, \Star{} O God of Israel, bless Thou, and hallow this house for ever.\end{Resp}
\begin{Resp}\Vsign O Lord, Who keepest covenant with thy servants that walk before thee in all their heart.\end{Resp}
\begin{Resp}\Rsign O God of Israel, bless Thou, and hallow this house for ever.\end{Resp}
\begin{Resp}\Vsign Glory be to the Father, and to the Son, \Star{} and to the Holy Ghost.\end{Resp}
\begin{Resp}\Rsign O God of Israel, bless Thou, and hallow this house for ever.\end{Resp}
\switchcolumn*

\end{paracol}

\Office{Feria quinta infra Hebdomadam XI post Octavam Pentecostes}{Thursday of the Eleventh Week after Pentecost}{Feria}
\addcontentsline{toc}{subsection}{Feria quinta infra Hebdomadam XI post Octavam Pentecostes \textit{— Thursday of the Eleventh Week after Pentecost}}
\begin{paracol}{2}
\LA
\LessonHead{Lectio I}
\switchcolumn
\EN
\LessonHead{Lesson I}
\switchcolumn*

\LA
\LessonTitle{De libro quarto Regum}
\switchcolumn
\EN
\LessonTitle{Lesson from the second book of Kings}
\switchcolumn*

\LA
\Cite{4 Reg 23:36-37; 24:1}
\switchcolumn
\EN
\Cite{2 Kgs 23:36-37; 24:1}
\switchcolumn*

\LA
\noindent Vigínti quinque annórum erat Jóakim cum regnáre cœpísset, et úndecim annis regnávit in Jerúsalem. Nomen matris ejus Zébida filia Phadaía de Ruma. Et fecit malum coram Dómino juxta omnia quæ fécerant patres ejus. In diébus ejus ascéndit Nabuchodónosor rex Babylónis, et factus est ei Jóakim servus tribus annis: et rursum rebellávit contra eum.\par
\switchcolumn
\EN
\noindent Joakim was five and twenty years old when he began to reign: and he reigned eleven years in Jerusalem: the name of his mother was Zebida the daughter of Phadaia of Ruma. And he did evil before the Lord according to all that his fathers had done. In his days Nabuchodonosor king of Babylon came up, and Joakim became his servant three years: then again he rebelled against him.\par
\switchcolumn*

\LA
\begin{Resp}\Rsign Præparáte corda vestra Dómino, et servíte illi soli: \Star{} Et liberábit vos de mánibus inimicórum vestrórum.\end{Resp}
\begin{Resp}\Vsign Convertímini ad eum in toto corde vestro, et auférte deos aliénos de médio vestri.\end{Resp}
\begin{Resp}\Rsign Et liberábit vos de mánibus inimicórum vestrórum.\end{Resp}
\switchcolumn
\EN
\begin{Resp}\Rsign Prepare your hearts unto the Lord, and serve Him Only \Star{} And He will deliver you out of the hand of your enemies.\end{Resp}
\begin{Resp}\Vsign Return unto Him with all your hearts, and put away the strange gods from among you.\end{Resp}
\begin{Resp}\Rsign And He will deliver you out of the hand of your enemies.\end{Resp}
\switchcolumn*

\LA
\LessonHead{Lectio II}
\switchcolumn
\EN
\LessonHead{Lesson II}
\switchcolumn*

\LA
\Cite{4 Reg 24:2-4}
\switchcolumn
\EN
\Cite{2 Kgs 24:2-4}
\switchcolumn*

\LA
\noindent Immisítque ei Dóminus latrúnculos Chaldæórum, et latrúnculos Sýriæ, et latrúnculos Moab, et latrúnculos filiórum Ammon: et immísit eos in Judam ut dispérderent eum, juxta verbum Dómini quod locútus fúerat per servos suos prophétas. Factum est autem hoc per verbum Dómini contra Judam, ut auférret eum coram se propter peccáta Manásse univérsa quæ fecit, et propter sánguinem innóxium quem effúdit, et implévit Jerúsalem cruóre innocéntium: et ob hanc rem nóluit Dóminus propitiári.\par
\switchcolumn
\EN
\noindent And the Lord sent against him the rovers of the Chaldees, and the rovers of Syria, and the rovers of Moab, and the rovers of the children of Ammon: and he sent them against Juda, to destroy it, according to the word of the Lord, which he had spoken by his servants the prophets. And this came by the word of the Lord against Juda, to remove them from before him for all the sins of Manasses which he did. And for the innocent blood that he shed, filling Jerusalem with innocent blood: and therefore the Lord would not be appeased\par
\switchcolumn*

\LA
\begin{Resp}\Rsign Deus ómnium exaudítor est: ipse misit Angelum suum, et tulit me de óvibus patris mei: \Star{} Et unxit me unctióne misericórdiæ suæ.\end{Resp}
\begin{Resp}\Vsign Dóminus, qui erípuit me de ore leónis, et de manu béstiæ liberávit me.\end{Resp}
\begin{Resp}\Rsign Et unxit me unctióne misericórdiæ suæ.\end{Resp}
\switchcolumn
\EN
\begin{Resp}\Rsign God, Which heareth all, even He sent His Angel, and took me from keeping my father's sheep, and \Star{} Anointed me with the oil of His mercy.\end{Resp}
\begin{Resp}\Vsign The Lord That delivered me out of the mouth of the lion, and out of the paw of the bear\end{Resp}
\begin{Resp}\Rsign And anointed me with the oil of His mercy.\end{Resp}
\switchcolumn*

\LA
\LessonHead{Lectio III}
\switchcolumn
\EN
\LessonHead{Lesson III}
\switchcolumn*

\LA
\Cite{4 Reg 24:5-7}
\switchcolumn
\EN
\Cite{2 Kgs 24:5-7}
\switchcolumn*

\LA
\noindent Réliqua autem sermónum Jóakim, et univérsa quæ fecit, nonne hæc scripta sunt in libro sermónum diérum regum Juda? Et dormívit Joakim cum pátribus suis: Et regnávit Jóachin fílius ejus pro eo. Et ultra non áddidit rex Ægýpti ut egrederétur de terra sua; túlerat enim rex Babylónis, a rivo Ægýpti usque ad flúvium Euphráten, ómnia quæ fúerant regis Ægýpti.\par
\switchcolumn
\EN
\noindent But the rest of the acts of Joakim, and all that he did, are they not written in the book of the words of the days of the kings of Juda? And Joakim slept with his fathers: And Joachin his son reigned in his stead. And the king of Egypt came not again any more out of his own country: for the king of Babylon had taken all that had belonged to the king of Egypt, from the river of Egypt, unto the river Euphrates\par
\switchcolumn*

\LA
\begin{Resp}\Rsign Dóminus, qui erípuit me de ore leónis, et de manu béstiæ liberávit me, \Star{} Ipse me erípiet de mánibus inimicórum meórum.\end{Resp}
\begin{Resp}\Vsign Misit Deus misericórdiam suam et veritátem suam: ánimam meam erípuit de médio catulórum leónum.\end{Resp}
\begin{Resp}\Rsign Ipse me erípiet de mánibus inimicórum meórum.\end{Resp}
\begin{Resp}\Vsign Glória Patri, et Fílio, \Star{} et Spirítui Sancto.\end{Resp}
\begin{Resp}\Rsign Ipse me erípiet de mánibus inimicórum meórum.\end{Resp}
\switchcolumn
\EN
\begin{Resp}\Rsign The Lord That delivered me out of the mouth of the lion, and out of the paw of the bear \Star{} He will deliver me out of the hand of mine enemies.\end{Resp}
\begin{Resp}\Vsign God hath sent forth His mercy and His truth, and delivered my soul from among the lion's whelps.\end{Resp}
\begin{Resp}\Rsign He will deliver me out of the hand of mine enemies.\end{Resp}
\begin{Resp}\Vsign Glory be to the Father, and to the Son, \Star{} and to the Holy Ghost.\end{Resp}
\begin{Resp}\Rsign He will deliver me out of the hand of mine enemies.\end{Resp}
\switchcolumn*

\end{paracol}

\Office{Sabbato infra Hebdomadam XI post Octavam Pentecostes}{Saturday of the Eleventh Week after Pentecost}{Feria}
\addcontentsline{toc}{subsection}{Sabbato infra Hebdomadam XI post Octavam Pentecostes \textit{— Saturday of the Eleventh Week after Pentecost}}
\begin{paracol}{2}
\LA
\LessonHead{Lectio I}
\switchcolumn
\EN
\LessonHead{Lesson I}
\switchcolumn*

\LA
\LessonTitle{De libro quarto Regum}
\switchcolumn
\EN
\LessonTitle{Lesson from the second book of Kings}
\switchcolumn*

\LA
\Cite{4 Reg 24:18-20; 25:1-3}
\switchcolumn
\EN
\Cite{2 Kgs 24:18-20; 25:1-3}
\switchcolumn*

\LA
\noindent Vigésimum et primum annum ætátis habébat Sedecías cum regnáre cœpísset, et úndecim annis regnávit in Jerúsalem. Nomen matris ejus erat Amítal fília Jeremíæ de Lobna. Et fecit malum coram Dómino, juxta ómnia quæ fecerat Jóakim; irascebátur enim Dóminus contra Jerúsalem et contra Judam, donec proíceret eos a fácie sua: recessítque Sedecías a rege Babylónis. Factum est autem anno nono regni ejus, mense décimo, décima die mensis, venit Nabuchodónosor rex Babylónis, ipse et omnis exércitus ejus, in Jerúsalem, et circumdedérunt eam: et exstruxérunt in circúitu ejus munitiónes; et clausa est cívitas atque valláta usque ad undécimum annum regis Sedecíæ, nona die mensis: prævaluítque fames in civitáte, nec erat panis pópulo terræ.\par
\switchcolumn
\EN
\noindent Sedecias was one and twenty years old when he began to reign, and he reigned eleven years in Jerusalem: the name of his mother was Amital, the daughter of Jeremias of Lobna. And he did evil before the Lord, according to all that Joakim had done. For the Lord was angry against Jerusalem and against Juda, till he cast them out from his face: and Sedecias revolted from the king of Babylon. And it came to pass in the ninth year of his reign, in the tenth month, the tenth day of the month, that Nabuchodonosor king of Babylon came, he and all his army against Jerusalem: and they surrounded it: and raised works round about it. And the city was shut up and besieged till the eleventh year of king Sedecias, The ninth day of the month: and a famine prevailed in the city, and there was no bread for the people of the land.\par
\switchcolumn*

\LA
\begin{Resp}\Rsign Peccávi super númerum arénæ maris, et multiplicáta sunt peccáta mea: et non sum dignus vidére altitúdinem cæli præ multitúdine iniquitátis meæ: quóniam irritávi iram tuam, \Star{} Et malum coram te feci.\end{Resp}
\begin{Resp}\Vsign Quóniam iniquitátem meam ego cognósco: et delíctum meum contra me est semper, quia tibi soli peccávi.\end{Resp}
\begin{Resp}\Rsign Et malum coram te feci.\end{Resp}
\switchcolumn
\EN
\begin{Resp}\Rsign My sins are many, yea, they are more in number than the sands of the sea; I am not worthy to look up toward heaven because of the multitude of my iniquities; for I have provoked thee to anger; \Star{} And done evil in thy sight.\end{Resp}
\begin{Resp}\Vsign For I acknowledge my transgression, and my sin is ever before me, for against thee only have I sinned.\end{Resp}
\begin{Resp}\Rsign And done evil in thy sight.\end{Resp}
\switchcolumn*

\LA
\LessonHead{Lectio II}
\switchcolumn
\EN
\LessonHead{Lesson II}
\switchcolumn*

\LA
\Cite{4 Reg 25:4-7}
\switchcolumn
\EN
\Cite{2 Kgs 25:4-7}
\switchcolumn*

\LA
\noindent Et interrúpta est cívitas: et omnes viri bellatóres nocte fugérunt per viam portæ quæ est inter dúplicem murum ad hortum regis. Porro Chaldǽi obsidébant in circúitu civitátem. Fugit ítaque Sedecías per viam quæ ducit ad campéstria solitúdinis. Et persecútus est exércitus Chaldæórum regem, comprehendítque eum in planítie Jéricho: et omnes bellatóres qui erant cum eo, dispérsi sunt, et reliquérunt eum. Apprehénsum ergo regem duxérunt ad regem Babylónis in Réblatha: qui locútus est cum eo judícium. Fílios autem Sedecíæ occídit coram eo, et óculos ejus effódit, vinxítque eum caténis, et addúxit in Babylónem.\par
\switchcolumn
\EN
\noindent And a breach was made into the city: and all the men of war fled in the night between the two walls by the king's garden, (now the Chaldees besieged the city round about,) and Sedecias fled by the way that leadeth to the plains of the wilderness. And the army of the Chaldees pursued after the king, and overtook him in the plains of Jericho: and all the warriors that were with him were scattered, and left him: So they took the king, and brought him to the king of Babylon to Reblatha, and he gave judgment upon him. And he slew the sons of Sedecias before his face, and he put out his eyes, and bound him with chains, and brought him to Babylon.\par
\switchcolumn*

\LA
\begin{Resp}\Rsign Exaudísti, Dómine, oratiónem servi tui, ut ædificárem templum nómini tuo: \Star{} Bénedic et sanctífica domum istam in sempitérnum, Deus Israël.\end{Resp}
\begin{Resp}\Vsign Dómine, qui custódis pactum cum servis tuis, qui ámbulant coram te in toto corde suo.\end{Resp}
\begin{Resp}\Rsign Bénedic et sanctífica domum istam in sempitérnum, Deus Israël.\end{Resp}
\switchcolumn
\EN
\begin{Resp}\Rsign O Lord, Thou hast hearkened unto the prayer of thy servant, that I might build a temple unto thy Name, \Star{} O God of Israel, bless Thou, and hallow this house for ever.\end{Resp}
\begin{Resp}\Vsign O Lord, Who keepest covenant with thy servants that walk before thee in all their heart.\end{Resp}
\begin{Resp}\Rsign O God of Israel, bless Thou, and hallow this house for ever.\end{Resp}
\switchcolumn*

\LA
\LessonHead{Lectio III}
\switchcolumn
\EN
\LessonHead{Lesson III}
\switchcolumn*

\LA
\Cite{4 Reg 25:8-13}
\switchcolumn
\EN
\Cite{2 Kgs 25:8-13}
\switchcolumn*

\LA
\noindent Mense quinto, séptima die mensis, ipse est annus nonus décimus regis Babylónis, venit Nabuzárdan princeps exércitus, servus regis Babylónis, in Jerúsalem. Et succendit domum Domini, et domum regis: et domos Jerúsalem, omnémque domum combússit igni. Et muros Jerúsalem in circúitu destrúxit omnis exércitus Chaldæórum, qui erat cum príncipe mílitum. Réliquam autem pópuli partem quæ remánserat in civitáte, et pérfugas qui transfúgerant ad regem Babylónis, et réliquum vulgus tránstulit Nabuzárdan princeps milítiæ; et de paupéribus terræ relíquit vinitóres et agrícolas. Colúmnas autem ǽreas quæ erant in templo Dómini, et bases, et mare ǽreum quod erat in domo Dómini, confregérunt Chaldǽi, et transtulérunt æs omne in Babylónem.\par
\switchcolumn
\EN
\noindent In the fifth month, the seventh day of the month, that is, the nineteenth year of the king of Babylon, came Nabuzardan commander of the army, a servant of the king of Babylon, into Jerusalem. And he burnt the house of the Lord, and the king's house, and the houses of Jerusalem, and every house he burnt with fire. And all the army of the Chaldees, which was with the commander of the troops, broke down the walls of Jerusalem round about. And Nabuzardan the commander of the army, carried away the rest of the people that remained in the city, and the fugitives that had gone over to the king of Babylon, and the remnant of the common people. But of the poor of the land he left some dressers of vines and husbandmen. And the pillars of brass that were in the temple of the Lord, and the bases, and the sea of brass which was in the house of the Lord, the Chaldees broke in pieces, and carried all the brass of them to Babylon\par
\switchcolumn*

\LA
\begin{Resp}\Rsign Audi, Dómine, hymnum et oratiónem, quam servus tuus orat coram te hódie, ut sint óculi tui apérti, et aures tuæ inténtæ, \Star{} Super domum istam die ac nocte.\end{Resp}
\begin{Resp}\Vsign Réspice, Dómine, de sanctuário tuo, et de excélso cælórum habitáculo.\end{Resp}
\begin{Resp}\Rsign Super domum istam die ac nocte.\end{Resp}
\begin{Resp}\Vsign Glória Patri, et Fílio, \Star{} et Spirítui Sancto.\end{Resp}
\begin{Resp}\Rsign Super domum istam die ac nocte.\end{Resp}
\switchcolumn
\EN
\begin{Resp}\Rsign Hearken, O Lord, unto the cry and to the prayer which thy servant prayeth before thee today, that thine eyes may be open and thine ears attend; \Star{} Toward this house day and night.\end{Resp}
\begin{Resp}\Vsign Look down from thine high and holy place, O Lord, even from heaven thy dwelling.\end{Resp}
\begin{Resp}\Rsign Toward this house, day and night.\end{Resp}
\begin{Resp}\Vsign Glory be to the Father, and to the Son, \Star{} and to the Holy Ghost.\end{Resp}
\begin{Resp}\Rsign Toward this house, day and night.\end{Resp}
\switchcolumn*

\end{paracol}

\Office{Dominica XII Post Pentecosten}{Twelfth Sunday after Pentecost}{Semiduplex}
\addcontentsline{toc}{subsection}{Dominica XII Post Pentecosten \textit{— Twelfth Sunday after Pentecost}}
\begin{paracol}{2}
\LA
\RefLine{Lectiones I–VI: de Scriptura occurrente.}
\switchcolumn
\EN
\RefLine{Lessons I–VI: from the Scripture of the day.}
\switchcolumn*

\LA
\LessonHead{Lectio VII}
\switchcolumn
\EN
\LessonHead{Lesson VII}
\switchcolumn*

\LA
\LessonTitle{Léctio sancti Evangélii secúndum Lucam}
\switchcolumn
\EN
\LessonTitle{From the Holy Gospel according to Luke}
\switchcolumn*

\LA
\Cite{Luc 10:23-37}
\switchcolumn
\EN
\Cite{Luke 10:23-37}
\switchcolumn*

\LA
\noindent In illo témpore: Dixit Jesus discípulis suis: Beáti óculi qui vident quæ vos vidétis; dico enim vobis quod multi prophétæ et reges voluérunt vidére quæ vos vidétis, et non vidérunt. Et réliqua.\par
\switchcolumn
\EN
\noindent At that time, Jesus said unto His disciples: Blessed are the eyes which see the things that ye see. For I tell you that many prophets and kings have desired to see those things which ye see, and have not seen them. And so on.\par
\switchcolumn*

\LA
\LessonTitle{Homilía sancti Bedæ Venerábilis Presbýteri}
\switchcolumn
\EN
\LessonTitle{Homily by the Venerable Bede, Priest (at Jarrow.)}
\switchcolumn*

\LA
\Source{Liber 3, cap. 43 in Lucæ 10}
\switchcolumn
\EN
\Source{Bk. iii. ch. 43 on Luke x.}
\switchcolumn*

\LA
\noindent Non óculi scribárum et pharisæórum, qui corpus tantum Dómini vidére; sed illi beáti óculi, qui ejus possunt cognóscere sacraménta, de quibus dícitur: Et revelásti ea párvulis. Beáti óculi parvulórum, quibus et se et Patrem Fílius reveláre dignátur. Abraham exsultávit ut vidéret diem Christi; et vidit, et gavísus est. Isaías quoque, et Michǽas, et multi álii prophétæ vidérunt glóriam Dómini, qui et proptérea Vidéntes sunt appelláti; sed hi omnes, a longe aspiciéntes et salutántes, per spéculum et in ænígmate vidérunt.\par
\switchcolumn
\EN
\noindent Blessed were the eyes not of Scribes and Pharisees, which saw but the Body of the Lord, but those eyes, eyes blessed indeed, which were able to see those things whereof it is written “Thou hast hid these things from the wise and prudent, and hast revealed them unto babes.” Blessed are the eyes of those little ones unto whom it seemeth good in the eyes of the Son to reveal Himself and the Father also. Abraham rejoiced to see the day of Christ and he saw it, and was glad. (John viii. 56.) Isaiah, and Micah, and many among the Prophets, saw the glory of the Lord, wherefore also they be called Seers, but all they beheld it and hailed it afar off, seeing but as through a glass, darkly. (1 Cor. xiii. 12.)\par
\switchcolumn*

\LA
\begin{Resp}\Rsign Dómine, Pater et Deus vitæ meæ, ne derelínquas me in cogitátu malígno: extolléntiam oculórum meórum ne déderis mihi, et desidérium malígnum avérte a me, Dómine; aufer a me concupiscéntiam, \Star{} Et ánimo irreverénti et infruníto ne tradas me, Dómine.\end{Resp}
\begin{Resp}\Vsign Ne derelínquas me, Dómine, ne accréscant ignorántiæ meæ, nec multiplicéntur delícta mea.\end{Resp}
\begin{Resp}\Rsign Et ánimo irreverénti et infruníto ne tradas me, Dómine.\end{Resp}
\switchcolumn
\EN
\begin{Resp}\Rsign O Lord, Father and God of my life, leave me not to evil counsels; give me not a proud look, but turn away from me an haughty mind, O Lord turn away from me concupiscence, \Star{} And give me not over unto an impudent and froward mind, O Lord!\end{Resp}
\begin{Resp}\Vsign Leave me not, O Lord, lest mine ignorance increase, and my sins abound.\end{Resp}
\begin{Resp}\Rsign And give me not over unto an impudent and froward mind, O Lord.\end{Resp}
\switchcolumn*

\LA
\LessonHead{Lectio VIII}
\switchcolumn
\EN
\LessonHead{Lesson VIII}
\switchcolumn*

\LA
\noindent Apóstoli autem, in præsentiárum habéntes Dóminum, convescentésque ei, et quæcúmque voluíssent interrogándo discéntes, nequáquam per Angelos aut várias visiónum spécies opus habébant docéri. Quos vero Lucas multos prophétas et reges dicit, Matthǽus apértius prophétas et justos appéllat. Ipse sunt enim reges magni; quia tentatiónum suárum mótibus non consentiéndo succúmbere, sed regéndo præésse novérunt.\par
\switchcolumn
\EN
\noindent Otherwise were the Apostles, who saw the Lord face to Face, eating with Him, and learning from Him by asking whatsoever they listed. For them there was no need to be taught by Angels, or the shifting fabric of visions. They whom Luke doth call Prophets and kings, Matthew nameth as “Prophets and righteous men” (xiii. 17.) Righteous men are indeed mighty kings, who know how to lord it over their own rebellious temptations, instead of falling under them to become their slaves.\par
\switchcolumn*

\LA
\begin{Resp}\Rsign Duo Séraphim clamábant alter ad álterum: \Star{} Sanctus, sanctus, sanctus Dóminus Deus Sábaoth: * Plena est omnis terra glória ejus.\end{Resp}
\begin{Resp}\Vsign Tres sunt qui testimónium dant in cælo: Pater, Verbum, et Spíritus Sanctus: et hi tres unum sunt.\end{Resp}
\begin{Resp}\Rsign Sanctus, sanctus, sanctus Dóminus Deus Sábaoth.\end{Resp}
\begin{Resp}\Vsign Glória Patri, et Fílio, \Star{} et Spirítui Sancto.\end{Resp}
\begin{Resp}\Rsign Plena est omnis terra glória ejus.\end{Resp}
\switchcolumn
\EN
\begin{Resp}\Rsign One Seraph cried unto another: \Star{} Holy, Holy, Holy is the Lord God of hosts; the whole earth is full of His glory.\end{Resp}
\begin{Resp}\Vsign There are Three That bear record in heaven: the Father, the Word, and the Holy Ghost, and these Three are One.\end{Resp}
\begin{Resp}\Rsign Holy, Holy, Holy is the Lord God of hosts.\end{Resp}
\begin{Resp}\Vsign Glory be to the Father, and to the Son, \Star{} and to the Holy Ghost.\end{Resp}
\begin{Resp}\Rsign The whole earth is full of His glory.\end{Resp}
\switchcolumn*

\LA
\LessonHead{Lectio IX}
\switchcolumn
\EN
\LessonHead{Lesson IX}
\switchcolumn*

\LA

\switchcolumn
\EN
\LessonTitle{“And, behold, a certain lawyer stood up, and tempted Him, saying Master, what shall I do to inherit eternal life?” This lawyer, who stood up to ask the Lord a tempting question touching eternal life, took the subject of his asking, as I think, from the words which the Lord had just uttered, when He said “Rejoice, because your names are written in heaven”. But his attempt was a proof of the truth of that which the Lord immediately added “I thank thee, O Father, Lord of heaven and earth, that Thou hast hid these things from the wise and prudent, and hast revealed them unto babes!”}
\switchcolumn*

\LA
\noindent Et ecce quidam legisperítus surréxit, tentans eum et dicens: Magíster, quid faciéndo vitam ætérnam possidébo? Legisperítus, qui de vita ætérna Dóminum tentans intérrogat, occasiónem, ut reor, tentándi de ipsis Dómini sermónibus sumpsit, ubi ait: Gaudéte autem quod nómina vestra scripta sunt in cælis. Sed ipsa sua tentatióne declárat quam vera sit illa Dómini conféssio, qua Patri lóquitur: Quod abscondísti hæc a sapiéntibus, et prudéntibus, et revelásti ea párvulis.\par
\switchcolumn
\EN

\switchcolumn*

\LA
\TeDeum{Te Deum}
\switchcolumn
\EN
\TeDeum{Te Deum}
\switchcolumn*

\end{paracol}

\Office{Dominica XIII Post Pentecosten}{Thirteenth Sunday after Pentecost}{Semiduplex}
\addcontentsline{toc}{subsection}{Dominica XIII Post Pentecosten \textit{— Thirteenth Sunday after Pentecost}}
\begin{paracol}{2}
\LA
\RefLine{Lectiones I–VI: de Scriptura occurrente.}
\switchcolumn
\EN
\RefLine{Lessons I–VI: from the Scripture of the day.}
\switchcolumn*

\LA
\LessonHead{Lectio VII}
\switchcolumn
\EN
\LessonHead{Lesson VII}
\switchcolumn*

\LA
\LessonTitle{Léctio sancti Evangélii secúndum Lucam}
\switchcolumn
\EN
\LessonTitle{From the Holy Gospel according to Luke}
\switchcolumn*

\LA
\Cite{Luc 17:11-19}
\switchcolumn
\EN
\Cite{Luke 17:11-19}
\switchcolumn*

\LA
\noindent In illo témpore: Dum iret Jesus in Jerúsalem, transíbat per médiam Samaríam et Galilǽam. Et cum ingrederétur quoddam castéllum, occurrérunt ei decem viri leprósi. Et réliqua.\par
\switchcolumn
\EN
\noindent It came to pass, as Jesus went to Jerusalem, that He passed through the midst of Samaria and Galilee. And, as He entered into a certain village, there met Him ten men that were lepers. And so on.\par
\switchcolumn*

\LA
\LessonTitle{Homilía sancti Augustíni Epíscopi}
\switchcolumn
\EN
\LessonTitle{Homily by St. Augustine, Bishop of Hippo}
\switchcolumn*

\LA
\Source{Lib. 2 quæst. Evang. cap. 40}
\switchcolumn
\EN
\Source{Bk. ii, Gospel Questions, ch. 40}
\switchcolumn*

\LA
\noindent De decem leprósis, quos Dóminus ita mundávit, cum ait: Ite, osténdite vos sacerdótibus; quæri potest, cur eos ad sacerdótes míserit, ut cum irent, mundaréntur. Nullum enim eórum, quibus hæc corporália benefícia prǽstitit, invenítur misísse ad sacerdótes, nisi leprósos. Nam et illum a lepra mundáverat, cui dixit: Vade, osténde te sacerdótibus, et offer pro te sacrifícium, quod præcépit Móyses, in testimónium illis. Quæréndum ígitur est, quid ipsa lepra signíficet: non enim sanáti, sed mundáti dicúntur, qui ea caruérunt. Colóris quippe vítium est, non valetúdinis, aut integritátis sénsuum atque membrórum.\par
\switchcolumn
\EN
\noindent The ten lepers “lifted up their voices and said: Jesus, Master, have mercy on us. And when He saw them, He said unto them: Go, show yourselves unto the Priests. And it came to pass that, as they went, they were cleansed.” Question: why did the Lord send them unto the Priests, that, as they went, they might be cleansed? Lepers were the only class among those upon whose bodies He worked mercy, whom we find that He sent unto the Priests. It is written in another place that He said to a leper whom He had cleansed: “Go, and show thyself to the Priest, and offer for thy cleansing according as Moses commanded, for a testimony unto them” (Luke v. 14, Lev. xiv. seq.) We ask then, of what leprosy was a type, whereof they that were ridded were called, not “healed,” but “cleansed.” It is a disease which doth first appear in the skin, but destroyeth not immediately the strength, nor the use of feeling and the limbs.\par
\switchcolumn*

\LA
\begin{Resp}\Rsign Quis mihi tríbuat, ut in inférno prótegas me et abscóndas me, donec pertránseat furor tuus, Dómine, nisi tu, qui solus es Deus? \Star{} Et constítuas mihi tempus, in quo recordéris mei?\end{Resp}
\begin{Resp}\Vsign Numquid sicut dies hóminis dies tui, ut quæras iniquitátem meam; cum sit nemo, qui de manu tua possit erúere.\end{Resp}
\begin{Resp}\Rsign Et constítuas mihi tempus, in quo recordéris mei?\end{Resp}
\switchcolumn
\EN
\begin{Resp}\Rsign O that Thou wouldest hide me in the grave; that Thou wouldest keep me secret, until thy wrath be past, even thine, O Lord, Thou That alone art God; \Star{} That Thou wouldest appoint me a set time, and remember me!\end{Resp}
\begin{Resp}\Vsign Are thy days as the days of man, that Thou inquirest after mine iniquity, and there is none that can deliver out of thine hand.\end{Resp}
\begin{Resp}\Rsign That Thou wouldest appoint me a set time, and remember me!\end{Resp}
\switchcolumn*

\LA
\LessonHead{Lectio VIII}
\switchcolumn
\EN
\LessonHead{Lesson VIII}
\switchcolumn*

\LA
\noindent Leprósi ergo non absúrde intélligi possunt, qui sciéntiam veræ fídei non habéntes, várias doctrínas profiténtur erróris. Non enim abscóndunt imperítiam suam; sed pro summa perítia próferunt in lucem, et jactántia sermónis osténtant. Nulla porro falsa doctrína est, quæ non áliqua vera intermísceat. Vera ergo falsis inordináte permíxta, in una disputatióne vel narratióne hóminis, tamquam in uníus córporis colóre apparéntia, signíficant lepram, tamquam veris falsísque colórum variántem atque maculántem.\par
\switchcolumn
\EN
\noindent The lepers, therefore, we may not absurdly suppose such to be figured as have not the knowledge of the true faith, but do show forth diverse-coloured teachings of error. They hide not their witlessness, but do use all such wit as they have to make it manifest, and proclaim it in high-sounding phrases. There is no false doctrine but hath some truth mixed up with it. A man's discourse then, with some truths in it unequally mingled with falsehoods, and all confounded in one mass, is like to the body of one that is stricken with leprosy, whereon all manner of foul colours do appear in this and that place along with the true colour of skin.\par
\switchcolumn*

\LA
\begin{Resp}\Rsign Duo Séraphim clamábant alter ad álterum: \Star{} Sanctus, sanctus, sanctus Dóminus Deus Sábaoth: * Plena est omnis terra glória ejus.\end{Resp}
\begin{Resp}\Vsign Tres sunt qui testimónium dant in cælo: Pater, Verbum, et Spíritus Sanctus: et hi tres unum sunt.\end{Resp}
\begin{Resp}\Rsign Sanctus, sanctus, sanctus Dóminus Deus Sábaoth.\end{Resp}
\begin{Resp}\Vsign Glória Patri, et Fílio, \Star{} et Spirítui Sancto.\end{Resp}
\begin{Resp}\Rsign Plena est omnis terra glória ejus.\end{Resp}
\switchcolumn
\EN
\begin{Resp}\Rsign One Seraph cried unto another: \Star{} Holy, Holy, Holy is the Lord God of hosts; the whole earth is full of His glory.\end{Resp}
\begin{Resp}\Vsign There are Three That bear record in heaven: the Father, the Word, and the Holy Ghost, and these Three are One.\end{Resp}
\begin{Resp}\Rsign Holy, Holy, Holy is the Lord God of hosts.\end{Resp}
\begin{Resp}\Vsign Glory be to the Father, and to the Son, \Star{} and to the Holy Ghost.\end{Resp}
\begin{Resp}\Rsign The whole earth is full of His glory.\end{Resp}
\switchcolumn*

\LA
\LessonHead{Lectio IX}
\switchcolumn
\EN
\LessonHead{Lesson IX}
\switchcolumn*

\LA
\noindent Hi autem tam vitándi sunt Ecclésiæ, ut, si fíeri potest, lóngius remóti, magno clamóre Christum interpéllent; sicut isti decem stetérunt a longe, et levavérunt vocem, dicéntes: Jesu præcéptor, miserére nostri. Nam et quod præceptórem vocant, quo nómine néscio utrum quisquam Dóminum interpelláverit pro medicína corporáli; satis puto significáre, lepram falsam esse doctrínam, quam bonus præcéptor abstérgit.\par
\switchcolumn
\EN
\noindent Such men as these are banished out of the walls of the Church, to the end that haply when they stand afar off they may lift up their voices and cry to Christ for pardon, just as those ten men that were lepers, which stood afar off, outside the village, lifted up their voices and said “Jesus, Master, have mercy on us.” That they styled Him Master, by which title I know not if any besought the Lord for bodily healing, I think doth sufficiently show that leprosy signifieth false doctrine, whereof the Good Master doth cleanse us.\par
\switchcolumn*

\LA
\TeDeum{Te Deum}
\switchcolumn
\EN
\TeDeum{Te Deum}
\switchcolumn*

\end{paracol}

\Office{Dominica XIV Post Pentecosten}{Fourteenth Sunday after Pentecost}{Semiduplex}
\addcontentsline{toc}{subsection}{Dominica XIV Post Pentecosten \textit{— Fourteenth Sunday after Pentecost}}
\begin{paracol}{2}
\LA
\RefLine{Lectiones I–VI: de Scriptura occurrente.}
\switchcolumn
\EN
\RefLine{Lessons I–VI: from the Scripture of the day.}
\switchcolumn*

\LA
\LessonHead{Lectio VII}
\switchcolumn
\EN
\LessonHead{Lesson VII}
\switchcolumn*

\LA
\LessonTitle{Léctio sancti Evangélii secúndum Matthǽum}
\switchcolumn
\EN
\LessonTitle{From the Holy Gospel according to Matthew}
\switchcolumn*

\LA
\Cite{Matt 6:24-33}
\switchcolumn
\EN
\Cite{Matt 6:24-33}
\switchcolumn*

\LA
\noindent In illo témpore: Dixit Jesus discípulis suis: Nemo potest duóbus dóminis servíre. Et réliqua.\par
\switchcolumn
\EN
\noindent At that time, Jesus said unto His disciples: No man can serve two masters. And so on.\par
\switchcolumn*

\LA
\LessonTitle{Homilía sancti Augustíni Epíscopi}
\switchcolumn
\EN
\LessonTitle{Homily by St. Augustine, Bishop of Hippo.}
\switchcolumn*

\LA
\Source{Lib. 2 de Sermone Dómini in monte, cap. 14}
\switchcolumn
\EN
\Source{Bk. ii. on the Lord's Sermon on the Mount, ch. xiv.}
\switchcolumn*

\LA
\noindent Nemo potest duóbus dóminis servíre. Ad hanc ipsam intentiónem referéndum est quod consequénter expónit, dicens: Aut enim unum ódio habébit, et álterum díliget; aut álterum patiétur, et álterum contémnet. Quæ verba diligénter consideránda sunt: nam, qui sint duo dómini, deínceps osténdit, cum dicit: Non potéstis Deo servíre, et mammónæ. Mammóna apud Hebrǽos divítiæ appellári dicúntur. Cóngruit et Púnicum nomen; nam lucrum Púnice mammon dícitur.\par
\switchcolumn
\EN
\noindent “No man can serve two masters,” and this is further explained “for either he will hate the one, and love the other; or else he will hold to the one, and despise the other.” These words we ought carefully to weigh, for the Lord showeth straightway who be the two masters whom we have choice of: “Ye cannot serve God and Mammon.” Mammon is a term which the Hebrews are said to use for riches. It is also a Carthaginian word for the Punic for “gain” is “mammon.”\par
\switchcolumn*

\LA
\begin{Resp}\Rsign Quis mihi tríbuat, ut in inférno prótegas me et abscóndas me, donec pertránseat furor tuus, Dómine, nisi tu, qui solus es Deus? \Star{} Et constítuas mihi tempus, in quo recordéris mei?\end{Resp}
\begin{Resp}\Vsign Numquid sicut dies hóminis dies tui, ut quæras iniquitátem meam; cum sit nemo, qui de manu tua possit erúere.\end{Resp}
\begin{Resp}\Rsign Et constítuas mihi tempus, in quo recordéris mei?\end{Resp}
\switchcolumn
\EN
\begin{Resp}\Rsign O that Thou wouldest hide me in the grave; that Thou wouldest keep me secret, until thy wrath be past, even thine, O Lord, Thou That alone art God; \Star{} That Thou wouldest appoint me a set time, and remember me!\end{Resp}
\begin{Resp}\Vsign Are thy days as the days of man, that Thou inquirest after mine iniquity, and there is none that can deliver out of thine hand.\end{Resp}
\begin{Resp}\Rsign That Thou wouldest appoint me a set time, and remember me!\end{Resp}
\switchcolumn*

\LA
\LessonHead{Lectio VIII}
\switchcolumn
\EN
\LessonHead{Lesson VIII}
\switchcolumn*

\LA
\noindent Sed qui servit mammónæ, illi útique servit, qui rebus istis terrénis mérito suæ perversitátis præpósitus, magistrátus hujus sǽculi a Dómino dícitur. Aut enim unum ódio habébit homo, et álterum díliget, id est, Deum; aut álterum patiétur, et álterum contémnet. Patiétur enim durum et perniciósum dóminum, quisquis servit mammónæ; sua enim cupiditáte implicátus, súbditur diábolo, et non eum díligit. Quis enim est qui díligat diábolum? sed tamen pátitur.\par
\switchcolumn
\EN
\noindent He which serveth mammon, serveth that evil one who hath perversely chosen to be lord of these earthly things, and is called by the Lord “the prince of this world.” (John xiv. 30.) Of these two masters, either a man will hate the one and love the other, that is God or he will hold to the one and despise the other. He which serveth mammon holdeth to an hard and destroying master, for he is led captive by his lust, and sold a slave to the devil, and him loveth no man is there any man that loveth the devil And yet there be that hold to him.\par
\switchcolumn*

\LA
\begin{Resp}\Rsign Duo Séraphim clamábant alter ad álterum: \Star{} Sanctus, sanctus, sanctus Dóminus Deus Sábaoth: * Plena est omnis terra glória ejus.\end{Resp}
\begin{Resp}\Vsign Tres sunt qui testimónium dant in cælo: Pater, Verbum, et Spíritus Sanctus: et hi tres unum sunt.\end{Resp}
\begin{Resp}\Rsign Sanctus, sanctus, sanctus Dóminus Deus Sábaoth.\end{Resp}
\begin{Resp}\Vsign Glória Patri, et Fílio, \Star{} et Spirítui Sancto.\end{Resp}
\begin{Resp}\Rsign Plena est omnis terra glória ejus.\end{Resp}
\switchcolumn
\EN
\begin{Resp}\Rsign One Seraph cried unto another: \Star{} Holy, Holy, Holy is the Lord God of hosts; the whole earth is full of His glory.\end{Resp}
\begin{Resp}\Vsign There are Three That bear record in heaven: the Father, the Word, and the Holy Ghost, and these Three are One.\end{Resp}
\begin{Resp}\Rsign Holy, Holy, Holy is the Lord God of hosts.\end{Resp}
\begin{Resp}\Vsign Glory be to the Father, and to the Son, \Star{} and to the Holy Ghost.\end{Resp}
\begin{Resp}\Rsign The whole earth is full of His glory.\end{Resp}
\switchcolumn*

\LA
\LessonHead{Lectio IX}
\switchcolumn
\EN
\LessonHead{Lesson IX}
\switchcolumn*

\LA

\switchcolumn
\EN
\LessonTitle{“Therefore, I say unto you, Take no thought for your life, what ye shall eat, or what ye shall drink; nor yet for your body, what ye shall put on”, lest, albeit such things are not idle, but needful to be sought after, yet the seeking for things even needful should divide the heart and our intention should be corrupted when we do something as it were mercifully that is, lest, when we would seem to be seeking another's good, it should be profit to ourselves, rather than benefit to him, that we seek and therefore we seem not to ourselves to sin, because we would seek things not idle, but needful.}
\switchcolumn*

\LA
\noindent Ideo, inquit, dico vobis, non habére sollicitúdinem ánimæ vestræ quid edátis, neque córpori vestro quid induátis; ne forte, quamvis jam supérflua non quærántur, propter ipsa necessária cor duplicétur, et ad ipsa conquirénda, nostra detorqueátur inténtio, cum áliquid quasi misericórditer operámur: id est, ut cum consúlere alícui vidéri vólumus, nostrum emoluméntum ibi pótius, quam illíus utilitátem attendámus; et ídeo nobis non videámur peccáre, quia non supérflua, sed necessária sunt, quæ cónsequi vólumus.\par
\switchcolumn
\EN

\switchcolumn*

\LA
\TeDeum{Te Deum}
\switchcolumn
\EN
\TeDeum{Te Deum}
\switchcolumn*

\end{paracol}

\Office{Dominica XV Post Pentecosten}{Fifteenth Sunday after Pentecost}{Semiduplex}
\addcontentsline{toc}{subsection}{Dominica XV Post Pentecosten \textit{— Fifteenth Sunday after Pentecost}}
\begin{paracol}{2}
\LA
\RefLine{Lectiones I–VI: de Scriptura occurrente.}
\switchcolumn
\EN
\RefLine{Lessons I–VI: from the Scripture of the day.}
\switchcolumn*

\LA
\LessonHead{Lectio VII}
\switchcolumn
\EN
\LessonHead{Lesson VII}
\switchcolumn*

\LA
\LessonTitle{Léctio sancti Evangélii secúndum Lucam}
\switchcolumn
\EN
\LessonTitle{From the Holy Gospel according to Luke}
\switchcolumn*

\LA
\Cite{Luc 7:11-16}
\switchcolumn
\EN
\Cite{Luke 7:11-16}
\switchcolumn*

\LA
\noindent In illo témpore: Ibat Jesus in civitátem, quæ vocátur Naim; et ibant cum eo discípuli ejus, et turba copiósa. Et réliqua.\par
\switchcolumn
\EN
\noindent At that time: Jesus went into a city called Nain and His disciples went with Him, and much people. And so on.\par
\switchcolumn*

\LA
\LessonTitle{Homilía sancti Augustíni Epíscopi}
\switchcolumn
\EN
\LessonTitle{Homily by St. Augustine, Bishop of Hippo}
\switchcolumn*

\LA
\Source{Sermo 44 de verbis Domini, circa initium}
\switchcolumn
\EN
\Source{44th Discourse on the Words of the Lord}
\switchcolumn*

\LA
\noindent De júvene illo resuscitáto gavísa est mater vídua; de homínibus in spíritu cotídie suscitátis gaudet mater Ecclésia. Ille quidem mórtuus erat córpore; illi autem mente. Illíus mors visíbilis visibíliter plangebátur; illórum mors invisíbilis nec quærebátur, nec videbátur. Quæsívit ille, qui nóverat mórtuos. Ille solus nóverat mórtuos, qui póterat fácere vivos. Nisi enim ad mórtuos suscitándos venísset, Apóstolus non díceret: Surge, qui dormis, et exsúrge a mórtuis, et illuminábit te Christus.\par
\switchcolumn
\EN
\noindent That her son was called again to life was the joy of that widowed mother; that souls of men are every day called to life is the joy of our Mother the Church. He was dead in body they have been dead in mind. His death was outward, and was outwardly bewailed; their inward. Death hath been neither mourned for nor seen. But He hath sought for them, Who hath seen that they are dead, and He only hath seen that they are dead, Who hath been able to make them alive. If He had not come to raise the dead, the Apostle had not said: “Awake, thou that sleepest, and arise from the dead, and Christ shall give thee light.” (Eph. v. 14.)\par
\switchcolumn*

\LA
\begin{Resp}\Rsign Quis mihi tríbuat, ut in inférno prótegas me et abscóndas me, donec pertránseat furor tuus, Dómine, nisi tu, qui solus es Deus? \Star{} Et constítuas mihi tempus, in quo recordéris mei?\end{Resp}
\begin{Resp}\Vsign Numquid sicut dies hóminis dies tui, ut quæras iniquitátem meam; cum sit nemo, qui de manu tua possit erúere.\end{Resp}
\begin{Resp}\Rsign Et constítuas mihi tempus, in quo recordéris mei?\end{Resp}
\switchcolumn
\EN
\begin{Resp}\Rsign O that Thou wouldest hide me in the grave; that Thou wouldest keep me secret, until thy wrath be past, even thine, O Lord, Thou That alone art God; \Star{} That Thou wouldest appoint me a set time, and remember me!\end{Resp}
\begin{Resp}\Vsign Are thy days as the days of man, that Thou inquirest after mine iniquity, and there is none that can deliver out of thine hand.\end{Resp}
\begin{Resp}\Rsign That Thou wouldest appoint me a set time, and remember me!\end{Resp}
\switchcolumn*

\LA
\LessonHead{Lectio VIII}
\switchcolumn
\EN
\LessonHead{Lesson VIII}
\switchcolumn*

\LA
\noindent Tres autem mórtuos invenímus a Dómino resuscitátos visibíliter, míllia invisibíliter. Quot autem mórtuos visibíliter suscitáverit, quis novit? Non enim ómnia, quæ fecit, scripta sunt. Joánnes hoc dixit: Multa ália fecit Jesus, quæ si scripta essent, árbitror totum mundum non posse libros cápere. Multi ergo sunt álii sine dúbio suscitáti, sed non tres frustra commemoráti. Dóminus enim noster Jesus Christus ea quæ faciébat corporáliter, étiam spiritáliter volébat intélligi. Neque enim tantum mirácula propter mirácula faciébat; sed ut illa, quæ faciébat, mira essent vidéntibus, vera essent intelligéntibus.\par
\switchcolumn
\EN
\noindent We find written how the Lord raised from the dead three persons visibly, but thousands invisibly. But how many they may have been whom He raised visibly, who knoweth For all the things which He did are not written. John saith thus: “There are also many other things which Jesus did, the which, if they should be written every one, I suppose that even the world itself could not contain the books that should be written.” (xxi. 25.) There were then, doubtless, many more raised to life, but it is not meaningless that three are recorded. For our Lord Jesus Christ hath willed that those things which He did carnally, we should understand also spiritually. He worked not miracles only for the sake of working wonders, but that His works might be at once wonderful to them that beheld, and true to them that understand them.\par
\switchcolumn*

\LA
\begin{Resp}\Rsign Duo Séraphim clamábant alter ad álterum: \Star{} Sanctus, sanctus, sanctus Dóminus Deus Sábaoth: * Plena est omnis terra glória ejus.\end{Resp}
\begin{Resp}\Vsign Tres sunt qui testimónium dant in cælo: Pater, Verbum, et Spíritus Sanctus: et hi tres unum sunt.\end{Resp}
\begin{Resp}\Rsign Sanctus, sanctus, sanctus Dóminus Deus Sábaoth.\end{Resp}
\begin{Resp}\Vsign Glória Patri, et Fílio, \Star{} et Spirítui Sancto.\end{Resp}
\begin{Resp}\Rsign Plena est omnis terra glória ejus.\end{Resp}
\switchcolumn
\EN
\begin{Resp}\Rsign One Seraph cried unto another: \Star{} Holy, Holy, Holy is the Lord God of hosts; the whole earth is full of His glory.\end{Resp}
\begin{Resp}\Vsign There are Three That bear record in heaven: the Father, the Word, and the Holy Ghost, and these Three are One.\end{Resp}
\begin{Resp}\Rsign Holy, Holy, Holy is the Lord God of hosts.\end{Resp}
\begin{Resp}\Vsign Glory be to the Father, and to the Son, \Star{} and to the Holy Ghost.\end{Resp}
\begin{Resp}\Rsign The whole earth is full of His glory.\end{Resp}
\switchcolumn*

\LA
\LessonHead{Lectio IX}
\switchcolumn
\EN
\LessonHead{Lesson IX}
\switchcolumn*

\LA
\noindent Quemádmodum qui videt lítteras in códice óptime scripto, et non novit légere, laudat quidem antiquárii manum, admírans ápicum pulchritúdinem; sed quid sibi velint, quid índicent illi ápices, nescit, et est óculis laudátor, mente non cógnitor. Alius autem et laudat artifícium, et capit intelléctum: ille útique, qui non solum vidére quod commúne est ómnibus, potest, sed étiam légere; quod qui non dídicit, non potest. Ita qui vidérunt Christi mirácula, et non intellexérunt quid sibi vellent, et quid intelligéntibus quodámmodo innúerent, miráti sunt tantum quia facta sunt; álii vero et facta miráti, et intellécta assecúti. Tales nos in schola Christi esse debémus.\par
\switchcolumn
\EN
\noindent Even as one that looketh upon a scroll right fairly written, and knoweth not how to read therein, praiseth the hand of the old scribe when he seeth the beauty of the points, but what it saith, what those points mean, he knoweth not, and praiseth by the eye, without understanding by the mind, and as, on the other hand, he that can not only gaze on it, as can all men, but also can read it, praiseth the penmanship, and catcheth the sense likewise, which the unlearned cannot do even so, there were some that saw the miracles which Christ did, and understood not what they meant, nor what they, as it were, hinted to such as did understand them, and these only marvelled to see them wrought. And other some there were which saw the works, and marvelled, and understood them, and profited by them. And it is as these last that we ought to be in the school of Christ.\par
\switchcolumn*

\LA
\TeDeum{Te Deum}
\switchcolumn
\EN
\TeDeum{Te Deum}
\switchcolumn*

\end{paracol}

\SeasonTitle{Dominicæ et Feriæ mensis Augusti}{Sundays and Ferias of August}

\addcontentsline{toc}{section}{Dominicæ et Feriæ mensis Augusti \textit{— Sundays and Ferias of August}}

\Office{Dominica I Augusti}{First Sunday of August}{}
\addcontentsline{toc}{subsection}{Dominica I Augusti \textit{— First Sunday of August}}
\begin{paracol}{2}
\LA
\LessonHead{Lectio I}
\switchcolumn
\EN
\LessonHead{Lesson I}
\switchcolumn*

\LA
\LessonTitle{Incípiunt Parábolæ Salomónis}
\switchcolumn
\EN
\LessonTitle{Lesson from the book of Proverbs}
\switchcolumn*

\LA
\Cite{Prov 1:1-6}
\switchcolumn
\EN
\Cite{Prov 1:1-6}
\switchcolumn*

\LA
\noindent Parábolæ Salomónis, fílii David regis Israël, ad sciéndam sapiéntiam et disciplínam, ad intelligénda verba prudéntiæ et suscipiéndam eruditiónem doctrínæ, justítiam et judícium et æquitátem, ut detur párvulis astútia, adolescénti sciéntia et intelléctus. Audiens sápiens sapiéntior erit, et intéllegens gubernácula possidébit: animadvértet parábolam et interpretatiónem, verba sapiéntum et ænígmata eórum.\par
\switchcolumn
\EN
\noindent The parables of Solomon, the son of David, king of Israel. To know wisdom, and instruction: To understand the words of prudence: and to receive the instruction of doctrine, justice, and judgment, and equity: To give subtilty to little ones, to the young man knowledge and understanding. A wise man shall hear and shall be wiser: and he that understandeth, shall possess governments. He shall understand a parable, and the interpretation, the words of the wise, and their mysterious sayings.\par
\switchcolumn*

\LA
\begin{Resp}\Rsign In princípio Deus ántequam terram fáceret, priúsquam abýssos constitúeret, priúsquam prodúceret fontes aquárum, \Star{} Antequam montes collocaréntur, ante omnes colles generávit me Dóminus.\end{Resp}
\begin{Resp}\Vsign Quando præparábat cælos, áderam, cum eo cuncta compónens.\end{Resp}
\begin{Resp}\Rsign Antequam montes collocaréntur, ante omnes colles generávit me Dóminus.\end{Resp}
\switchcolumn
\EN
\begin{Resp}\Rsign God possessed me in the beginning, before He made the earth, before He created the depths, before He caused the fountains of water to spring. \Star{} Before the mountains were settled, before there were any hills, did the Lord beget me.\end{Resp}
\begin{Resp}\Vsign When He prepared the heavens, I was there with Him, ordering all things.\end{Resp}
\begin{Resp}\Rsign Before the mountains were settled, before there were any hills, did the Lord beget me.\end{Resp}
\switchcolumn*

\LA
\LessonHead{Lectio II}
\switchcolumn
\EN
\LessonHead{Lesson II}
\switchcolumn*

\LA
\Cite{Prov 1:7-14}
\switchcolumn
\EN
\Cite{Prov 1:7-14}
\switchcolumn*

\LA
\noindent Timor Dómini princípium sapiéntiæ. Sapiéntiam atque doctrínam stulti despíciunt. Audi, fili mi, disciplínam patris tui et ne dimíttas legem matris tuæ, ut addátur grátia cápiti tuo et torques collo tuo. Fili mi, si te lactáverint peccatóres, ne acquiéscas eis: si díxerint: Veni nobíscum, insidiémur sánguini, abscondámus tendículas contra insóntem frustra, deglutiámus eum sicut inférnus vivéntem et íntegrum quasi descendéntem in lacum; omnem pretiósam substántiam reperiémus, implébimus domos nostras spóliis: sortem mitte nobíscum, marsúpium unum sit ómnium nostrum.\par
\switchcolumn
\EN
\noindent The fear of the Lord is the beginning of wisdom. Fools despise wisdom and instruction. My son, hear the instruction of thy father, and forsake not the law of thy mother: That grace may be added to thy head, and a chain of gold to thy neck. My son, if sinners shall entice thee, consent not to them. If they shall say: Come with us, let us lie in wait for blood, let us hide snares for the innocent without cause: Let us swallow him up alive like hell, and whole as one that goeth down into the pit. We shall find all precious substance, we shall fill our houses with spoils. Cast in thy lot with us, let us all have one purse.\par
\switchcolumn*

\LA
\begin{Resp}\Rsign Gyrum cæli circuívi sola, et in flúctibus maris ambulávi, in omni gente et in omni pópulo primátum ténui: \Star{} Superbórum et sublímium colla própria virtúte calcávi.\end{Resp}
\begin{Resp}\Vsign Ego in altíssimis hábito, et thronus meus in colúmna nubis.\end{Resp}
\begin{Resp}\Rsign Superbórum et sublímium colla própria virtúte calcávi.\end{Resp}
\switchcolumn
\EN
\begin{Resp}\Rsign I alone compassed the circuit of heaven, and walked on the waves of the sea. In every nation and in every people, I held the first place. \Star{} In the greatness of my strength have I trodden under my feet the necks of such as be haughty and proud.\end{Resp}
\begin{Resp}\Vsign I dwell in the highest places, and my throne is in a cloudy pillar.\end{Resp}
\begin{Resp}\Rsign In the greatness of my strength have I trodden under my feet the necks of such as be haughty and proud.\end{Resp}
\switchcolumn*

\LA
\LessonHead{Lectio III}
\switchcolumn
\EN
\LessonHead{Lesson III}
\switchcolumn*

\LA
\Cite{Prov 1:15-19}
\switchcolumn
\EN
\Cite{Prov 1:15-19}
\switchcolumn*

\LA
\noindent Fili mi, ne ámbules cum eis: próhibe pedem tuum a sémitis eórum; pedes enim illórum ad malum currunt, et festínant ut effúndant sánguinem. Frustra autem jácitur rete ante óculos pennatórum. Ipsi quoque contra sánguinem suum insidiántur et moliúntur fraudes contra ánimas suas. Sic sémitæ omnis avári: ánimas possidéntium rápiunt.\par
\switchcolumn
\EN
\noindent My son, walk not thou with them, restrain thy foot from their paths. For their feet run to evil, and make haste to shed blood. But a net is spread in vain before the eyes of them that have wings. And they themselves lie in wait for their own blood, and practise deceits against their own souls. So the wage of every covetous man destroy the souls of the possessors.\par
\switchcolumn*

\LA
\begin{Resp}\Rsign Emítte, Dómine, sapiéntiam de sede magnitúdinis tuæ, ut mecum sit et mecum labóret: \Star{} Ut sciam, quid accéptum sit coram te omni témpore.\end{Resp}
\begin{Resp}\Vsign Da mihi, Dómine, sédium tuárum assistrícem sapiéntiam.\end{Resp}
\begin{Resp}\Rsign Ut sciam quid accéptum sit coram te omni témpore.\end{Resp}
\begin{Resp}\Vsign Glória Patri, et Fílio, \Star{} et Spirítui Sancto.\end{Resp}
\begin{Resp}\Rsign Ut sciam quid accéptum sit coram te omni témpore.\end{Resp}
\switchcolumn
\EN
\begin{Resp}\Rsign O send out wisdom from the throne of thy glory, O Lord, to be with me, and to labour with me, \Star{} That I may know at all times what is pleasing unto thee.\end{Resp}
\begin{Resp}\Vsign Give me wisdom, O Lord, that sitteth by thy throne.\end{Resp}
\begin{Resp}\Rsign That I may know at all times what is pleasing unto thee.\end{Resp}
\begin{Resp}\Vsign Glory be to the Father, and to the Son, \Star{} and to the Holy Ghost.\end{Resp}
\begin{Resp}\Rsign That I may know at all times what is pleasing unto thee.\end{Resp}
\switchcolumn*

\LA
\LessonHead{Lectio IV}
\switchcolumn
\EN
\LessonHead{Lesson IV}
\switchcolumn*

\LA
\LessonTitle{Ex Tractatu sancti Ambrósii Epíscopi in Psalmum centésimum décimum octavum}
\switchcolumn
\EN
\LessonTitle{From the Treatise of St. Ambrose, Bishop of Milan, upon the 18-th Psalm.}
\switchcolumn*

\LA
\Source{Sermo 5, n. 36-37}
\switchcolumn
\EN
\Source{Sermon v. 6}
\switchcolumn*

\LA
\noindent Inítium esse sapiéntiæ timórem Dómini, dicit prophéta. Quid est autem inítium sapiéntiæ, nisi sǽculo renuntiáre? Quia sápere sæculária, stultítia est. Dénique sapiéntiam hujus mundi, stultítiam esse apud Deum, Apóstolus dicit. Sed et ipse timor Dómini, nisi secúndum sciéntiam sit, nihil prodest, immo obest plúrimum. Síquidem Judǽi habent zelum Dei; sed quia non habent secúndum sciéntiam, in ipso zelo et timóre majórem cóntrahunt divinitátis offénsam. Quod circumcídunt infántulos suos, quod sábbata custódiunt, timórem Dei habent; sed quia nesciunt legem spiritálem esse, circumcídunt corpus, non cor suum.\par
\switchcolumn
\EN
\noindent The Prophet saith that “the fear of the Lord is the beginning of wisdom.” And what is the first act of wisdom but to renounce the world since to love the things of the world is folly. So indeed saith the Apostle “The wisdom of this world is foolishness with God.” (i Cor. iii. 19.) But the very fear of the Lord itself is useless, nay, harmful, if it be not according to knowledge. The Jews have a truly fervent zeal for God, but since they have not knowledge, their very zeal and fear do cause them to do things contrary to God's will. That they circumcise their children, that they keep holy the Sabbath-Day, showeth how they fear the Lord, but knowing not the spiritual meaning of the Law, they circumcise the body and not the heart.\par
\switchcolumn*

\LA
\begin{Resp}\Rsign Da mihi, Dómine, sédium tuárum assistrícem sapiéntiam, et noli me reprobáre a púeris tuis: \Star{} Quóniam servus tuus sum ego, et fílius ancíllæ tuæ.\end{Resp}
\begin{Resp}\Vsign Mitte illam de sede magnitúdinis tuæ, ut mecum sit et mecum labóret.\end{Resp}
\begin{Resp}\Rsign Quóniam servus tuus sum ego, et fílius ancíllæ tuæ.\end{Resp}
\switchcolumn
\EN
\begin{Resp}\Rsign Give me wisdom, O Lord, that sitteth by thy throne, and reject me not from among thy children. \Star{} For I am thy servant and son of thine handmaid.\end{Resp}
\begin{Resp}\Vsign O send her out from the throne of thy glory, to be with me and to labour with me.\end{Resp}
\begin{Resp}\Rsign For I am thy servant and son of thine handmaid.\end{Resp}
\switchcolumn*

\LA
\LessonHead{Lectio V}
\switchcolumn
\EN
\LessonHead{Lesson V}
\switchcolumn*

\LA
\noindent Et quid de Judǽis dico? Sunt étiam in nobis qui habent timórem Dei, sed non secúndum sciéntiam, statuéntes durióra præcépta, quæ non possit humána condítio sustinére. Timor in eo est, quia vidéntur sibi consúlere disciplínæ, opus virtútis exígere; sed inscítia in eo est, quia non compatiúntur natúræ, non ǽstimant possibilitátem. Non sit ergo irrationábilis timor. Etenim vera sapiéntia a timóre Dei íncipit, nec est sapiéntia spiritális sine timóre Dei; ita timor sine sapiéntia esse non debet.\par
\switchcolumn
\EN
\noindent But wherefore should I speak of Jews There are those among ourselves who have the fear of God, but not according to knowledge, and set up hard ordinances which the weakness of man is not able to bear. They fear God in this, that they seem to themselves to be looking to discipline, and to be enforcing the practice of godliness, but they lack knowledge, in that they feel not for the weakness of nature, nor consider whether a thing can, or cannot be done. Let not then the fear of God be unreasonable. True wisdom beginneth with the fear of God, neither is it spiritual wisdom without the fear of God, but neither ought the fear of God to be without wisdom.\par
\switchcolumn*

\LA
\begin{Resp}\Rsign Inítium sapiéntiæ timor Dómini: \Star{} Intelléctus bonus ómnibus faciéntibus eum: laudátio ejus manet in sǽculum sǽculi.\end{Resp}
\begin{Resp}\Vsign Diléctio illíus custódia legum est: quia omnis sapiéntia timor Dómini.\end{Resp}
\begin{Resp}\Rsign Intelléctus bonus ómnibus faciéntibus eum: laudátio ejus manet in sǽculum sǽculi.\end{Resp}
\switchcolumn
\EN
\begin{Resp}\Rsign The fear of the Lord is the beginning of wisdom. \Star{} A good understanding have all they that do His commandments. His praise endureth for ever.\end{Resp}
\begin{Resp}\Vsign Love is the keeping of her laws, for all wisdom is the fear of the Lord.\end{Resp}
\begin{Resp}\Rsign A good understanding have all they that do His commandments. His praise endureth for ever.\end{Resp}
\switchcolumn*

\LA
\LessonHead{Lectio VI}
\switchcolumn
\EN
\LessonHead{Lesson VI}
\switchcolumn*

\LA
\noindent Basis quædam verbi est timor sanctus. Sicut enim simulácrum áliquod in basi statúitur, et tunc majórem habet grátiam, cum in basi státua fúerit collocáta, standíque áccipit firmitátem: ita verbum Dei in timóre sancto mélius statúitur, fórtius radicátur in péctore timéntis Dóminum; ne labátur verbum de corde viri, ne véniant vólucres et áuferant illud de incuriósi et dissimulántis afféctu.\par
\switchcolumn
\EN
\noindent Holy fear is the foundation of all good instruction. Just as a statue is set up upon a pedestal, and thereby receiveth both beauty and strength, even so doth it become the word of God to be set forth based upon an holy fear, and it is in the heart of him that feareth that it getteth the firmest root, even an home wherefrom it droppeth not, neither do the fowls of the air come and carry it away, as from the heart of him that is careless and deceiving.\par
\switchcolumn*

\LA
\begin{Resp}\Rsign Verbum iníquum et dolósum longe fac a me, Dómine: \Star{} Divítias et paupertátem ne déderis mihi, sed tantum víctui meo tríbue necessária.\end{Resp}
\begin{Resp}\Vsign Duo rogávi te, ne déneges mihi ántequam móriar.\end{Resp}
\begin{Resp}\Rsign Divítias et paupertátem ne déderis mihi, sed tantum víctui meo tríbue necessária.\end{Resp}
\begin{Resp}\Vsign Glória Patri, et Fílio, \Star{} et Spirítui Sancto.\end{Resp}
\begin{Resp}\Rsign Divítias et paupertátem ne déderis mihi, sed tantum víctui meo tríbue necessária.\end{Resp}
\switchcolumn
\EN
\begin{Resp}\Rsign Lord, remove far from me vanity and lies. \Star{} Give me neither poverty nor riches, but feed me with food convenient for me.\end{Resp}
\begin{Resp}\Vsign Two things have I required of thee; deny me them not before I die.\end{Resp}
\begin{Resp}\Rsign Give me neither poverty nor riches, but feed me with food convenient for me.\end{Resp}
\begin{Resp}\Vsign Glory be to the Father, and to the Son, \Star{} and to the Holy Ghost.\end{Resp}
\begin{Resp}\Rsign Give me neither poverty nor riches, but feed me with food convenient for me.\end{Resp}
\switchcolumn*

\end{paracol}

\Office{Feria II infra Hebdomadam I Augusti}{Monday of the first week of August}{}
\addcontentsline{toc}{subsection}{Feria II infra Hebdomadam I Augusti \textit{— Monday of the first week of August}}
\begin{paracol}{2}
\LA
\LessonHead{Lectio I}
\switchcolumn
\EN
\LessonHead{Lesson I}
\switchcolumn*

\LA
\LessonTitle{De Parábolis Salomónis}
\switchcolumn
\EN
\LessonTitle{Lesson from the book of Proverbs}
\switchcolumn*

\LA
\Cite{Prov 3:1-6}
\switchcolumn
\EN
\Cite{Prov 3:1-6}
\switchcolumn*

\LA
\noindent Fili mi, ne obliviscáris legis meæ, et præcépta mea cor tuum custódiat; longitúdinem enim diérum et annos vitæ et pacem appónent tibi. Misericórdia et véritas te non déserant: circúmda eas gútturi tuo et descríbe in tábulis cordis tui, et invénies grátiam et disciplínam bonam coram Deo et homínibus. Habe fidúciam in Dómino ex toto corde tuo et ne innitáris prudéntiæ tuæ; in ómnibus viis tuis cógita illum, et ipse díriget gressus tuos.\par
\switchcolumn
\EN
\noindent My son, forget not my law, and let thy heart keep my commandments. For they shall add to thee length of days, and years of life and peace. Let not mercy and truth leave thee, put them about thy neck, and write them in the tables of thy heart: And thou shalt find grace and good understanding before God and men. Have confidence in the Lord with all thy heart, and lean not upon thy own prudence. In all thy ways think on him, and he will direct thy steps.\par
\switchcolumn*

\LA
\begin{Resp}\Rsign Ne derelínquas me, Dómine, pater et dominátor vitæ meæ, ut non córruam in conspéctu adversariórum meórum: \Star{} Ne gáudeat de me inimícus meus.\end{Resp}
\begin{Resp}\Vsign Apprehénde arma et scutum et exsúrge in adjutórium mihi.\end{Resp}
\begin{Resp}\Rsign Ne gáudeat de me inimícus meus.\end{Resp}
\switchcolumn
\EN
\begin{Resp}\Rsign O Lord, Father and Governor of my life, leave me not, lest I fall before mine adversaries, \Star{} and mine enemy rejoice over me.\end{Resp}
\begin{Resp}\Vsign Take hold of shield and buckler, and stand up for mine help.\end{Resp}
\begin{Resp}\Rsign Lest mine enemy rejoice over me.\end{Resp}
\switchcolumn*

\LA
\LessonHead{Lectio II}
\switchcolumn
\EN
\LessonHead{Lesson II}
\switchcolumn*

\LA
\Cite{Prov 3:7-10}
\switchcolumn
\EN
\Cite{Prov 3:7-10}
\switchcolumn*

\LA
\noindent Ne sis sápiens apud temetípsum, time Deum et recéde a malo; sánitas quippe erit umbilíco tuo et irrigátio óssium tuórum. Honóra Dóminum de tua substántia et de primítiis ómnium frugum tuárum da ei; et implebúntur hórrea tua saturitáte, et vino torculária tua redundábunt.\par
\switchcolumn
\EN
\noindent Be not wise in thy own conceit: fear God, and depart from evil: For it shall be health to thy navel, and moistening to thy bones. Honour the Lord with thy substance, and give him of the first of all thy fruits: And thy barns shall be filled with abundance, and thy presses shall run over with wine.\par
\switchcolumn*

\LA
\begin{Resp}\Rsign Magna enim sunt judícia tua, Dómine, et inenarrabília verba tua: \Star{} Magnificásti pópulum tuum et honorásti.\end{Resp}
\begin{Resp}\Vsign Transtulísti illos per Mare Rubrum et transvexísti eos per aquam nímiam.\end{Resp}
\begin{Resp}\Rsign Magnificásti pópulum tuum et honorásti.\end{Resp}
\switchcolumn
\EN
\begin{Resp}\Rsign Great are thy judgments, O Lord, and thy words cannot be expressed. \Star{} Thou didst make thy people mighty and honourable.\end{Resp}
\begin{Resp}\Vsign Thou broughtest them through the Red Sea, and leddest them through much water.\end{Resp}
\begin{Resp}\Rsign Thou didst make thy people mighty and honourable.\end{Resp}
\switchcolumn*

\LA
\LessonHead{Lectio III}
\switchcolumn
\EN
\LessonHead{Lesson III}
\switchcolumn*

\LA
\Cite{Prov 3:11-15}
\switchcolumn
\EN
\Cite{Prov 3:11-15}
\switchcolumn*

\LA
\noindent Disciplínam Dómini, fili mi, ne abícias, nec defícias cum ab eo corríperis; quem enim díligit Dóminus córripit, et quasi pater in fílio cómplacet sibi. Beátus homo qui invénit sapiéntiam et qui áffluit prudéntia. Mélior est acquisítio ejus negotiatióne argénti, et auri primi et puríssimi fructus ejus; pretiósior est cunctis ópibus, et ómnia quæ desiderántur, huic non valent comparári.\par
\switchcolumn
\EN
\noindent My son, reject not the correction of the Lord: and do not faint when thou art chastised by him: For whom the Lord loveth, he chastiseth: and as a father in the son he pleaseth himself. Blessed is the man that findeth wisdom and is rich in prudence: The purchasing thereof is better than the merchandise of silver, and her fruit than the chiefest and purest gold: She is more precious than all riches: and all the things that are desired, are not to be compared with her.\par
\switchcolumn*

\LA
\begin{Resp}\Rsign Quæ sunt in corde hóminum, óculi tui vident, Dómine, et in libro tuo ómnia scribéntur: \Star{} Homo videt in fácie, Deus autem in corde.\end{Resp}
\begin{Resp}\Vsign Omnia enim corda scrutátur, et univérsas méntium cogitatiónes intéllegit.\end{Resp}
\begin{Resp}\Rsign Homo videt in fácie, Deus autem in corde.\end{Resp}
\begin{Resp}\Vsign Glória Patri, et Fílio, \Star{} et Spirítui Sancto.\end{Resp}
\begin{Resp}\Rsign Homo videt in fácie, Deus autem in corde.\end{Resp}
\switchcolumn
\EN
\begin{Resp}\Rsign Lord, thine eyes behold all that is in the heart of man, and in thy book are they all written. \Star{} Man looketh on the outward appearance, but God looketh on the heart.\end{Resp}
\begin{Resp}\Vsign For He searcheth all hearts, and understandeth all the imaginations of the thoughts.\end{Resp}
\begin{Resp}\Rsign Man looketh on the outward appearance, but God looketh on the heart.\end{Resp}
\begin{Resp}\Vsign Glory be to the Father, and to the Son, \Star{} and to the Holy Ghost.\end{Resp}
\begin{Resp}\Rsign Man looketh on the outward appearance, but God looketh on the heart.\end{Resp}
\switchcolumn*

\end{paracol}

\Office{Feria III infra Hebdomadam I Augusti}{Tuesday of the first week of August}{}
\addcontentsline{toc}{subsection}{Feria III infra Hebdomadam I Augusti \textit{— Tuesday of the first week of August}}
\begin{paracol}{2}
\LA
\LessonHead{Lectio I}
\switchcolumn
\EN
\LessonHead{Lesson I}
\switchcolumn*

\LA
\LessonTitle{De Parábolis Salomónis}
\switchcolumn
\EN
\LessonTitle{Lesson from the book of Proverbs}
\switchcolumn*

\LA
\Cite{Prov 5:1-6}
\switchcolumn
\EN
\Cite{Prov 5:1-6}
\switchcolumn*

\LA
\noindent Fili mi, atténde ad sapiéntiam meam, et prudéntiæ meæ inclína aurem tuam, ut custódias cogitatiónes, et disciplínam lábia tua consérvent. Ne atténdas falláciæ mulíeris; favus enim distíllans lábia meretrícis, et nitídius óleo guttur ejus; novíssima autem illíus amára quasi absýnthium et acúta quasi gládius biceps; pedes ejus descéndunt in mortem, et ad ínferos gressus illíus pénetrant; per sémitam vitæ non ámbulant, vagi sunt gressus ejus et investigábiles.\par
\switchcolumn
\EN
\noindent My son, attend to my wisdom, and incline thy ear to my prudence. That thou mayst keep thoughts, and thy lips may preserve instruction. Mind not the deceit of a woman. For the lips of a harlot are like a honeycomb dropping, and her throat is smoother than oil. But her end is bitter as wormwood, and sharp as a two-edged sword. Her feet go down into death, and her steps go in as far as hell. They walk not by the path of life, her steps are wandering, and unaccountable.\par
\switchcolumn*

\LA
\begin{Resp}\Rsign Præbe, fili, cor mihi, et óculi tui vias meas custódiant: \Star{} Ut addátur grátia cápiti tuo.\end{Resp}
\begin{Resp}\Vsign Atténde, fili mi, sapiéntiam meam et ad elóquium meum inclína aurem tuam.\end{Resp}
\begin{Resp}\Rsign Ut addátur grátia cápiti tuo.\end{Resp}
\switchcolumn
\EN
\begin{Resp}\Rsign My son, give me thine heart, and let thine eyes observe my ways. \Star{} For they shall be an ornament of grace unto thine head.\end{Resp}
\begin{Resp}\Vsign My son, attend unto my wisdom, and incline thine ear unto my sayings.\end{Resp}
\begin{Resp}\Rsign For they shall be an ornament of grace unto thine head.\end{Resp}
\switchcolumn*

\LA
\LessonHead{Lectio II}
\switchcolumn
\EN
\LessonHead{Lesson II}
\switchcolumn*

\LA
\Cite{Prov 5:7-13}
\switchcolumn
\EN
\Cite{Prov 5:7-13}
\switchcolumn*

\LA
\noindent Nunc ergo, fili mi, audi me, et ne recédas a verbis oris mei. Longe fac ab ea viam tuam et ne appropínques fóribus domus ejus; ne des aliénis honórem tuum, et annos tuos crudéli; ne forte impleántur extránei víribus tuis, et labóres tui sint in domo aliéna, et gemas in novíssimis, quando consúmpseris carnes tuas et corpus tuum, et dicas: Cur detestátus sum disciplínam, et increpatiónibus non acquiévit cor meum; nec audívi vocem docéntium me, et magístris non inclinávi aurem meam?\par
\switchcolumn
\EN
\noindent Now therefore, my son, hear me, and depart not from the words of my mouth. Remove thy way far from her, and come not nigh the doors of her house. Give not thy honour to strangers, and thy years to the cruel. Lest strangers be filled with thy strength, and thy labours be in another man's house, And thou mourn it the last, when thou shalt have spent thy flesh and thy body, and say: Why have I hated instruction, and my heart consented not to reproof, And have not heard the voice of them that taught me, and have not inclined my ear to masters?\par
\switchcolumn*

\LA
\begin{Resp}\Rsign Inítium sapiéntiæ timor Dómini: \Star{} Intelléctus bonus ómnibus faciéntibus eum: laudátio ejus manet in sǽculum sǽculi.\end{Resp}
\begin{Resp}\Vsign Diléctio illíus custódia legum est: quia omnis sapiéntia timor Dómini.\end{Resp}
\begin{Resp}\Rsign Intelléctus bonus ómnibus faciéntibus eum: laudátio ejus manet in sǽculum sǽculi.\end{Resp}
\switchcolumn
\EN
\begin{Resp}\Rsign The fear of the Lord is the beginning of wisdom. \Star{} A good understanding have all they that do His commandments. His praise endureth for ever.\end{Resp}
\begin{Resp}\Vsign Love is the keeping of her laws, for all wisdom is the fear of the Lord.\end{Resp}
\begin{Resp}\Rsign A good understanding have all they that do His commandments. His praise endureth for ever.\end{Resp}
\switchcolumn*

\LA
\LessonHead{Lectio III}
\switchcolumn
\EN
\LessonHead{Lesson III}
\switchcolumn*

\LA
\Cite{Prov 5:20-23}
\switchcolumn
\EN
\Cite{Prov 5:20-23}
\switchcolumn*

\LA
\noindent Quare sedúceris, fili mi, ab aliéna, et fovéris in sinu altérius? Réspicit Dóminus vias hóminis, et omnes gressus ejus consíderat; iniquitátes suæ cápiunt ímpium, et fúnibus peccatórum suórum constríngitur; ipse moriétur, quia non hábuit disciplínam, et in multitúdine stultítiæ suæ decipiétur.\par
\switchcolumn
\EN
\noindent Why art thou seduced, my son, by a strange woman, and art cherished in the bosom of another? The Lord beholdeth the ways of man, and considereth all his steps. His own iniquities catch the wicked, and he is fast bound with the ropes of his own sins. He shall die, because he hath not received instruction, and in the multitude of his folly he shall be deceived.\par
\switchcolumn*

\LA
\begin{Resp}\Rsign Verbum iníquum et dolósum longe fac a me, Dómine: \Star{} Divítias et paupertátem ne déderis mihi, sed tantum víctui meo tríbue necessária.\end{Resp}
\begin{Resp}\Vsign Duo rogávi te, ne déneges mihi ántequam móriar.\end{Resp}
\begin{Resp}\Rsign Divítias et paupertátem ne déderis mihi, sed tantum víctui meo tríbue necessária.\end{Resp}
\begin{Resp}\Vsign Glória Patri, et Fílio, \Star{} et Spirítui Sancto.\end{Resp}
\begin{Resp}\Rsign Divítias et paupertátem ne déderis mihi, sed tantum víctui meo tríbue necessária.\end{Resp}
\switchcolumn
\EN
\begin{Resp}\Rsign Lord, remove far from me vanity and lies. \Star{} Give me neither poverty nor riches, but feed me with food convenient for me.\end{Resp}
\begin{Resp}\Vsign Two things have I required of thee; deny me them not before I die.\end{Resp}
\begin{Resp}\Rsign Give me neither poverty nor riches, but feed me with food convenient for me.\end{Resp}
\begin{Resp}\Vsign Glory be to the Father, and to the Son, \Star{} and to the Holy Ghost.\end{Resp}
\begin{Resp}\Rsign Give me neither poverty nor riches, but feed me with food convenient for me.\end{Resp}
\switchcolumn*

\end{paracol}

\Office{Feria IV infra Hebdomadam I Augusti}{Wednesday of the first week of August}{}
\addcontentsline{toc}{subsection}{Feria IV infra Hebdomadam I Augusti \textit{— Wednesday of the first week of August}}
\begin{paracol}{2}
\LA
\LessonHead{Lectio I}
\switchcolumn
\EN
\LessonHead{Lesson I}
\switchcolumn*

\LA
\LessonTitle{De Parábolis Salomónis}
\switchcolumn
\EN
\LessonTitle{Lesson from the book of Proverbs}
\switchcolumn*

\LA
\Cite{Prov 8:1-6}
\switchcolumn
\EN
\Cite{Prov 8:1-6}
\switchcolumn*

\LA
\noindent Numquid non sapiéntia clámitat, et prudéntia dat vocem suam? In summis excelsísque vertícibus supra viam, in médiis sémitis stans, juxta portas civitátis, in ipsis fóribus lóquitur, dicens: O viri, ad vos clámito, et vox mea ad fílios hóminum; intellégite, párvuli, astútiam, et, insipiéntes, animadvérte; audíte, quóniam de rebus magnis locutúra sum, et aperiéntur lábia mea, ut recta prǽdicent.\par
\switchcolumn
\EN
\noindent Doth not wisdom cry aloud, and prudence put forth her voice? Standing in the top of the highest places by the way, in the midst of the paths. Beside the gates of the city, in the very doors she speaketh, saying: O ye men, to you I call, and my voice is to the sons of men. O little ones, understand subtilty, and ye unwise, take notice. Hear, for I will speak of great things: and my lips shall be opened to preach right things.\par
\switchcolumn*

\LA
\begin{Resp}\Rsign Dómine, Pater et Deus vitæ meæ, ne derelínquas me in cogitátu malígno: extolléntiam oculórum meórum ne déderis mihi, et desidérium malígnum avérte a me, Dómine; aufer a me concupiscéntiam, \Star{} Et ánimo irreverénti et infruníto ne tradas me, Dómine.\end{Resp}
\begin{Resp}\Vsign Ne derelínquas me, Dómine, ne accréscant ignorántiæ meæ, nec multiplicéntur delícta mea.\end{Resp}
\begin{Resp}\Rsign Et ánimo irreverénti et infruníto ne tradas me, Dómine.\end{Resp}
\switchcolumn
\EN
\begin{Resp}\Rsign O Lord, Father and God of my life, leave me not to evil counsels; give me not a proud look, but turn away from me an haughty mind, O Lord turn away from me concupiscence, \Star{} And give me not over unto an impudent and froward mind, O Lord!\end{Resp}
\begin{Resp}\Vsign Leave me not, O Lord, lest mine ignorance increase, and my sins abound.\end{Resp}
\begin{Resp}\Rsign And give me not over unto an impudent and froward mind, O Lord.\end{Resp}
\switchcolumn*

\LA
\LessonHead{Lectio II}
\switchcolumn
\EN
\LessonHead{Lesson II}
\switchcolumn*

\LA
\Cite{Prov 8:7-11}
\switchcolumn
\EN
\Cite{Prov 8:7-11}
\switchcolumn*

\LA
\noindent Veritátem meditábitur guttur meum, et lábia mea detestabúntur ímpium; justi sunt omnes sermónes mei, non est in eis pravum quid, neque pervérsum: recti sunt intellegéntibus, et æqui inveniéntibus sciéntiam. Accípite disciplínam meam, et non pecúniam, doctrínam magis quam aurum elígite; mélior est enim sapiéntia cunctis pretiosíssimis, et omne desiderábile ei non potest comparári.\par
\switchcolumn
\EN
\noindent My mouth shall meditate truth, and my lips shall hate wickedness. All my words are just, there is nothing wicked nor perverse in them. They are right to them that understand, and just to them that find knowledge. Receive my instruction, and not money: choose knowledge rather than gold. For wisdom is better than all the most precious things: and whatsoever may be desired cannot be compared to it.\par
\switchcolumn*

\LA
\begin{Resp}\Rsign Magna enim sunt judícia tua, Dómine, et inenarrabília verba tua: \Star{} Magnificásti pópulum tuum et honorásti.\end{Resp}
\begin{Resp}\Vsign Transtulísti illos per Mare Rubrum et transvexísti eos per aquam nímiam.\end{Resp}
\begin{Resp}\Rsign Magnificásti pópulum tuum et honorásti.\end{Resp}
\switchcolumn
\EN
\begin{Resp}\Rsign Great are thy judgments, O Lord, and thy words cannot be expressed. \Star{} Thou didst make thy people mighty and honourable.\end{Resp}
\begin{Resp}\Vsign Thou broughtest them through the Red Sea, and leddest them through much water.\end{Resp}
\begin{Resp}\Rsign Thou didst make thy people mighty and honourable.\end{Resp}
\switchcolumn*

\LA
\LessonHead{Lectio III}
\switchcolumn
\EN
\LessonHead{Lesson III}
\switchcolumn*

\LA
\Cite{Prov 8:12-17}
\switchcolumn
\EN
\Cite{Prov 8:12-17}
\switchcolumn*

\LA
\noindent Ego sapiéntia hábito in consílio, et erudítis intérsum cogitatiónibus; timor Dómini odit malum, arrogántiam, et supérbiam, et viam pravam, et os bilíngue detéstor; meum est consílium et ǽquitas, mea est prudéntia, mea est fortitúdo; per me reges regnant, et legum conditóres justa decérnunt; per me príncipes ímperant, et poténtes decérnunt justítiam; ego diligéntes me díligo, et qui mane vígilant ad me, invénient me.\par
\switchcolumn
\EN
\noindent I wisdom dwell in counsel, and am present in learned thoughts. The fear of the Lord hateth evil: I hate arrogance, and pride, and every wicked way, and a mouth with a double tongue. Counsel and equity is mine, prudence is mine, strength is mine. By me kings reign, and lawgivers decree just things, By me princes rule, and the mighty decree justice. I love them that love me: and they that in the morning early watch for me, shall find me.\par
\switchcolumn*

\LA
\begin{Resp}\Rsign Quæ sunt in corde hóminum, óculi tui vident, Dómine, et in libro tuo ómnia scribéntur: \Star{} Homo videt in fácie, Deus autem in corde.\end{Resp}
\begin{Resp}\Vsign Omnia enim corda scrutátur, et univérsas méntium cogitatiónes intéllegit.\end{Resp}
\begin{Resp}\Rsign Homo videt in fácie, Deus autem in corde.\end{Resp}
\begin{Resp}\Vsign Glória Patri, et Fílio, \Star{} et Spirítui Sancto.\end{Resp}
\begin{Resp}\Rsign Homo videt in fácie, Deus autem in corde.\end{Resp}
\switchcolumn
\EN
\begin{Resp}\Rsign Lord, thine eyes behold all that is in the heart of man, and in thy book are they all written. \Star{} Man looketh on the outward appearance, but God looketh on the heart.\end{Resp}
\begin{Resp}\Vsign For He searcheth all hearts, and understandeth all the imaginations of the thoughts.\end{Resp}
\begin{Resp}\Rsign Man looketh on the outward appearance, but God looketh on the heart.\end{Resp}
\begin{Resp}\Vsign Glory be to the Father, and to the Son, \Star{} and to the Holy Ghost.\end{Resp}
\begin{Resp}\Rsign Man looketh on the outward appearance, but God looketh on the heart.\end{Resp}
\switchcolumn*

\end{paracol}

\Office{Feria V infra Hebdomadam I Augusti}{Thursday of the first week of August}{}
\addcontentsline{toc}{subsection}{Feria V infra Hebdomadam I Augusti \textit{— Thursday of the first week of August}}
\begin{paracol}{2}
\LA
\LessonHead{Lectio I}
\switchcolumn
\EN
\LessonHead{Lesson I}
\switchcolumn*

\LA
\LessonTitle{De Parábolis Salomónis}
\switchcolumn
\EN
\LessonTitle{Lesson from the book of Proverbs}
\switchcolumn*

\LA
\Cite{Prov 10:1-5}
\switchcolumn
\EN
\Cite{Prov 10:1-5}
\switchcolumn*

\LA
\noindent Fílius sápiens lætíficat patrem, fílius vero stultus mæstítia est matris suæ. Nihil próderunt thesáuri impietátis, justítia vero liberábit a morte. Non afflíget Dóminus fame ánimam justi, et insídias impiórum subvértet. Egestátem operáta est manus remíssa, manus autem fórtium divítias parat. Qui nítitur mendáciis, hic pascit ventos, idem autem ipse séquitur aves volántes. Qui cóngregat in messe, fílius sápiens est; qui autem stertit æstáte, fílius confusiónis.\par
\switchcolumn
\EN
\noindent A wise son maketh the father glad: but a foolish son is the sorrow of his mother. Treasures of wickedness shall profit nothing: but justice shall deliver from death. The Lord will not afflict the soul of the just with famine, and he will disappoint the deceitful practices of the wicked. The slothful hand hath wrought poverty: but the hand of the industrious getteth riches. He that trusteth to lies feedeth the winds: and the same runneth after birds that fly away. He that gathered in the harvest is a wise son: but he that snorteth in the summer, is the son of confusion.\par
\switchcolumn*

\LA
\begin{Resp}\Rsign In princípio Deus ántequam terram fáceret, priúsquam abýssos constitúeret, priúsquam prodúceret fontes aquárum, \Star{} Antequam montes collocaréntur, ante omnes colles generávit me Dóminus.\end{Resp}
\begin{Resp}\Vsign Quando præparábat cælos, áderam, cum eo cuncta compónens.\end{Resp}
\begin{Resp}\Rsign Antequam montes collocaréntur, ante omnes colles generávit me Dóminus.\end{Resp}
\switchcolumn
\EN
\begin{Resp}\Rsign God possessed me in the beginning, before He made the earth, before He created the depths, before He caused the fountains of water to spring. \Star{} Before the mountains were settled, before there were any hills, did the Lord beget me.\end{Resp}
\begin{Resp}\Vsign When He prepared the heavens, I was there with Him, ordering all things.\end{Resp}
\begin{Resp}\Rsign Before the mountains were settled, before there were any hills, did the Lord beget me.\end{Resp}
\switchcolumn*

\LA
\LessonHead{Lectio II}
\switchcolumn
\EN
\LessonHead{Lesson II}
\switchcolumn*

\LA
\Cite{Prov 10:6-10}
\switchcolumn
\EN
\Cite{Prov 10:6-10}
\switchcolumn*

\LA
\noindent Benedíctio Dómini super caput justi, os autem impiórum óperit iníquitas. Memória justi cum láudibus, et nomen impiórum putréscet. Sápiens corde præcépta súscipit, stultus cǽditur lábiis. Qui ámbulat simplíciter, ámbulat confidénter; qui autem deprávat vias suas, maniféstus erit. Qui ánnuit óculo dabit dolórem, et stultus lábiis verberábitur.\par
\switchcolumn
\EN
\noindent The blessing of the Lord is upon the head of the just: but iniquity covereth the mouth of the wicked. The memory of the just is with praises: and the name of the wicked shall rot. The wise of heart receiveth precepts: a fool is beaten with lips. He that walketh sincerely, walketh confidently: but he that perverteth his ways, shall be manifest. He that winketh with the eye shall cause sorrow: and the foolish in lips shall be beaten.\par
\switchcolumn*

\LA
\begin{Resp}\Rsign Gyrum cæli circuívi sola, et in flúctibus maris ambulávi, in omni gente et in omni pópulo primátum ténui: \Star{} Superbórum et sublímium colla própria virtúte calcávi.\end{Resp}
\begin{Resp}\Vsign Ego in altíssimis hábito, et thronus meus in colúmna nubis.\end{Resp}
\begin{Resp}\Rsign Superbórum et sublímium colla própria virtúte calcávi.\end{Resp}
\switchcolumn
\EN
\begin{Resp}\Rsign I alone compassed the circuit of heaven, and walked on the waves of the sea. In every nation and in every people, I held the first place. \Star{} In the greatness of my strength have I trodden under my feet the necks of such as be haughty and proud.\end{Resp}
\begin{Resp}\Vsign I dwell in the highest places, and my throne is in a cloudy pillar.\end{Resp}
\begin{Resp}\Rsign In the greatness of my strength have I trodden under my feet the necks of such as be haughty and proud.\end{Resp}
\switchcolumn*

\LA
\LessonHead{Lectio III}
\switchcolumn
\EN
\LessonHead{Lesson III}
\switchcolumn*

\LA
\Cite{Prov 10:11-16}
\switchcolumn
\EN
\Cite{Prov 10:11-16}
\switchcolumn*

\LA
\noindent Vena vitæ os justi, et os impiórum óperit iniquitátem. Odium súscitat rixas, et univérsa delícta óperit cáritas. In lábiis sapiéntis invenítur sapiéntia, et virga in dorso ejus qui índiget corde. Sapiéntes abscóndunt sciéntiam, os autem stulti confusióni próximum est. Substántia dívitis urbs fortitúdinis ejus, pavor páuperum egéstas eórum. Opus justi ad vitam, fructus autem ímpii ad peccátum.\par
\switchcolumn
\EN
\noindent The mouth of the just is a vein of life: and the mouth of the wicked covereth iniquity. Hatred stirreth up strifes: and charity covereth all sins. In the lips of the wise is wisdom found: and a rod on the back of him that wanteth sense. Wise men lay up knowledge: but the mouth of the fool is next to confusion. The substance of a rich man is the city of his strength: the fear of the poor is their poverty. The work of the just is unto life: but the fruit of the wicked, unto sin.\par
\switchcolumn*

\LA
\begin{Resp}\Rsign Emítte, Dómine, sapiéntiam de sede magnitúdinis tuæ, ut mecum sit et mecum labóret: \Star{} Ut sciam, quid accéptum sit coram te omni témpore.\end{Resp}
\begin{Resp}\Vsign Da mihi, Dómine, sédium tuárum assistrícem sapiéntiam.\end{Resp}
\begin{Resp}\Rsign Ut sciam quid accéptum sit coram te omni témpore.\end{Resp}
\begin{Resp}\Vsign Glória Patri, et Fílio, \Star{} et Spirítui Sancto.\end{Resp}
\begin{Resp}\Rsign Ut sciam quid accéptum sit coram te omni témpore.\end{Resp}
\switchcolumn
\EN
\begin{Resp}\Rsign O send out wisdom from the throne of thy glory, O Lord, to be with me, and to labour with me, \Star{} That I may know at all times what is pleasing unto thee.\end{Resp}
\begin{Resp}\Vsign Give me wisdom, O Lord, that sitteth by thy throne.\end{Resp}
\begin{Resp}\Rsign That I may know at all times what is pleasing unto thee.\end{Resp}
\begin{Resp}\Vsign Glory be to the Father, and to the Son, \Star{} and to the Holy Ghost.\end{Resp}
\begin{Resp}\Rsign That I may know at all times what is pleasing unto thee.\end{Resp}
\switchcolumn*

\end{paracol}

\Office{Feria VI infra Hebdomadam I Augusti}{Friday of the first week of August}{}
\addcontentsline{toc}{subsection}{Feria VI infra Hebdomadam I Augusti \textit{— Friday of the first week of August}}
\begin{paracol}{2}
\LA
\LessonHead{Lectio I}
\switchcolumn
\EN
\LessonHead{Lesson I}
\switchcolumn*

\LA
\LessonTitle{De Parábolis Salomónis}
\switchcolumn
\EN
\LessonTitle{Lesson from the book of Proverbs}
\switchcolumn*

\LA
\Cite{Prov 14:1-5}
\switchcolumn
\EN
\Cite{Prov 14:1-5}
\switchcolumn*

\LA
\noindent Sápiens múlier ædíficat domum suam, insípiens exstrúctam quoque mánibus déstruet. Ambulans recto itínere et timens Deum despícitur ab eo qui infámi gráditur via. In ore stulti virga supérbiæ, lábia autem sapiéntium custódiunt eos. Ubi non sunt boves, præsépe vácuum est; ubi autem plúrimæ ségetes, ibi manifésta est fortitúdo bovis. Testis fidélis non mentítur, profert autem mendácium dolósus testis.\par
\switchcolumn
\EN
\noindent A wise woman buildeth her house: but the foolish will pull down with her hands that also which is built. He that walketh in the right way, and feareth God, is despised by him that goeth by an infamous way. In the mouth of a fool is the rod of pride: but the lips of the wise preserve them. Where there are no oxen, the crib is empty: but where there is much corn, there the strength of the ox is manifest. A faithful witness will not lie: but a deceitful witness uttereth a lie.\par
\switchcolumn*

\LA
\begin{Resp}\Rsign Da mihi, Dómine, sédium tuárum assistrícem sapiéntiam, et noli me reprobáre a púeris tuis: \Star{} Quóniam servus tuus sum ego, et fílius ancíllæ tuæ.\end{Resp}
\begin{Resp}\Vsign Mitte illam de sede magnitúdinis tuæ, ut mecum sit et mecum labóret.\end{Resp}
\begin{Resp}\Rsign Quóniam servus tuus sum ego, et fílius ancíllæ tuæ.\end{Resp}
\switchcolumn
\EN
\begin{Resp}\Rsign Give me wisdom, O Lord, that sitteth by thy throne, and reject me not from among thy children. \Star{} For I am thy servant and son of thine handmaid.\end{Resp}
\begin{Resp}\Vsign O send her out from the throne of thy glory, to be with me and to labour with me.\end{Resp}
\begin{Resp}\Rsign For I am thy servant and son of thine handmaid.\end{Resp}
\switchcolumn*

\LA
\LessonHead{Lectio II}
\switchcolumn
\EN
\LessonHead{Lesson II}
\switchcolumn*

\LA
\Cite{Prov 14:6-11}
\switchcolumn
\EN
\Cite{Prov 14:6-11}
\switchcolumn*

\LA
\noindent Quærit derísor sapiéntiam et non ínvenit, doctrína prudéntium fácilis. Vade contra virum stultum, et nescit lábia prudéntiæ. Sapiéntia cállidi est intellégere viam suam, et imprudéntia stultórum errans. Stultus illúdet peccátum, et inter justos morábitur grátia. Cor, quod novit amaritúdinem ánimæ suæ, in gáudio ejus non miscébitur extráneus. Domus impiórum delébitur, tabernácula vero justórum germinábunt.\par
\switchcolumn
\EN
\noindent A scorner seeketh wisdom, and findeth it not: the learning of the wise is easy. Go against a foolish man, and he knoweth not the lips of prudence. The wisdom of a. discreet man is to understand his way: and the imprudence of fools erreth. A fool will laugh at sin, but among the just grace shall abide. The heart that knoweth the bitterness of his own soul, in his joy the stranger shall not intermeddle. The house of the wicked shall be destroyed: but the tabernacles of the just shall flourish.\par
\switchcolumn*

\LA
\begin{Resp}\Rsign Inítium sapiéntiæ timor Dómini: \Star{} Intelléctus bonus ómnibus faciéntibus eum: laudátio ejus manet in sǽculum sǽculi.\end{Resp}
\begin{Resp}\Vsign Diléctio illíus custódia legum est: quia omnis sapiéntia timor Dómini.\end{Resp}
\begin{Resp}\Rsign Intelléctus bonus ómnibus faciéntibus eum: laudátio ejus manet in sǽculum sǽculi.\end{Resp}
\switchcolumn
\EN
\begin{Resp}\Rsign The fear of the Lord is the beginning of wisdom. \Star{} A good understanding have all they that do His commandments. His praise endureth for ever.\end{Resp}
\begin{Resp}\Vsign Love is the keeping of her laws, for all wisdom is the fear of the Lord.\end{Resp}
\begin{Resp}\Rsign A good understanding have all they that do His commandments. His praise endureth for ever.\end{Resp}
\switchcolumn*

\LA
\LessonHead{Lectio III}
\switchcolumn
\EN
\LessonHead{Lesson III}
\switchcolumn*

\LA
\Cite{Prov 14:12-16}
\switchcolumn
\EN
\Cite{Prov 14:12-16}
\switchcolumn*

\LA
\noindent Est via, quæ vidétur hómini justa, novíssima autem ejus dedúcunt ad mortem. Risus dolóre miscébitur, et extréma gáudii luctus óccupat. Viis suis replébitur stultus, et super eum erit vir bonus. Innocens credit omni verbo, astútus consíderat gressus suos. Fílio dolóso nihil erit boni, servo autem sapiénti prósperi erunt actus et dirigétur via ejus. Sápiens timet et declínat a malo, stultus tránsilit et confídit.\par
\switchcolumn
\EN
\noindent There is a way which seemeth just to a man: but the ends thereof lead to death. Laughter shall be mingled with sorrow, and mourning taketh hold of the end of joy. A fool shall be filled with his own ways, and the good man shall be above him. The innocent believeth every word: the discreet man considereth his steps. No good shall come to the deceitful son: but the wise servant shall prosper in his dealings, and his way shall be made straight. A wise man feareth and declineth from evil: the fool leapeth over and is confident.\par
\switchcolumn*

\LA
\begin{Resp}\Rsign Verbum iníquum et dolósum longe fac a me, Dómine: \Star{} Divítias et paupertátem ne déderis mihi, sed tantum víctui meo tríbue necessária.\end{Resp}
\begin{Resp}\Vsign Duo rogávi te, ne déneges mihi ántequam móriar.\end{Resp}
\begin{Resp}\Rsign Divítias et paupertátem ne déderis mihi, sed tantum víctui meo tríbue necessária.\end{Resp}
\begin{Resp}\Vsign Glória Patri, et Fílio, \Star{} et Spirítui Sancto.\end{Resp}
\begin{Resp}\Rsign Divítias et paupertátem ne déderis mihi, sed tantum víctui meo tríbue necessária.\end{Resp}
\switchcolumn
\EN
\begin{Resp}\Rsign Lord, remove far from me vanity and lies. \Star{} Give me neither poverty nor riches, but feed me with food convenient for me.\end{Resp}
\begin{Resp}\Vsign Two things have I required of thee; deny me them not before I die.\end{Resp}
\begin{Resp}\Rsign Give me neither poverty nor riches, but feed me with food convenient for me.\end{Resp}
\begin{Resp}\Vsign Glory be to the Father, and to the Son, \Star{} and to the Holy Ghost.\end{Resp}
\begin{Resp}\Rsign Give me neither poverty nor riches, but feed me with food convenient for me.\end{Resp}
\switchcolumn*

\end{paracol}

\Office{Sabbato infra Hebdomadam I Augusti}{Saturday of the first week of August}{}
\addcontentsline{toc}{subsection}{Sabbato infra Hebdomadam I Augusti \textit{— Saturday of the first week of August}}
\begin{paracol}{2}
\LA
\LessonHead{Lectio I}
\switchcolumn
\EN
\LessonHead{Lesson I}
\switchcolumn*

\LA
\LessonTitle{De Parábolis Salomónis}
\switchcolumn
\EN
\LessonTitle{Lesson from the book of Proverbs}
\switchcolumn*

\LA
\Cite{Prov 16:1-5}
\switchcolumn
\EN
\Cite{Prov 16:1-5}
\switchcolumn*

\LA
\noindent Hóminis est ánimam præparáre, et Dómini gubernáre linguam. Omnes viæ hóminis patent óculis ejus: spirítuum ponderátor est Dóminus. Revéla Dómino ópera tua, et dirigéntur cogitatiónes tuæ. Univérsa propter semetípsum operátus est Dóminus, ímpium quoque ad diem malum. Abominátio Dómini est omnis árrogans, étiam si manus ad manum fúerit, non est ínnocens.\par
\switchcolumn
\EN
\noindent It is the part of man to prepare the soul: and of the Lord to govern the tongue. All the ways of a man are open to his eyes: the Lord is the weigher of spirits. Lay open thy works to the Lord: and thy thoughts shall be directed. The Lord hath made all things for himself: the wicked also for the evil day. Every proud man is an abomination to the Lord: though hand should be joined to hand, he is not innocent.\par
\switchcolumn*

\LA
\begin{Resp}\Rsign Dómine, Pater et Deus vitæ meæ, ne derelínquas me in cogitátu malígno: extolléntiam oculórum meórum ne déderis mihi, et desidérium malígnum avérte a me, Dómine; aufer a me concupiscéntiam, \Star{} Et ánimo irreverénti et infruníto ne tradas me, Dómine.\end{Resp}
\begin{Resp}\Vsign Ne derelínquas me, Dómine, ne accréscant ignorántiæ meæ, nec multiplicéntur delícta mea.\end{Resp}
\begin{Resp}\Rsign Et ánimo irreverénti et infruníto ne tradas me, Dómine.\end{Resp}
\switchcolumn
\EN
\begin{Resp}\Rsign O Lord, Father and God of my life, leave me not to evil counsels; give me not a proud look, but turn away from me an haughty mind, O Lord turn away from me concupiscence, \Star{} And give me not over unto an impudent and froward mind, O Lord!\end{Resp}
\begin{Resp}\Vsign Leave me not, O Lord, lest mine ignorance increase, and my sins abound.\end{Resp}
\begin{Resp}\Rsign And give me not over unto an impudent and froward mind, O Lord.\end{Resp}
\switchcolumn*

\LA
\LessonHead{Lectio II}
\switchcolumn
\EN
\LessonHead{Lesson II}
\switchcolumn*

\LA
\Cite{Prov 16:5-9}
\switchcolumn
\EN
\Cite{Prov 16:5-9}
\switchcolumn*

\LA
\noindent Inítium viæ bonæ fácere justítiam, accépta est autem apud Deum magis quam immoláre hóstias. Misericórdia et veritáte redímitur iníquitas, et in timóre Dómini declinátur a malo. Cum placúerint Dómino viæ hóminis, inimícos quoque ejus convértet ad pacem. Mélius est parum cum justítia, quam multi fructus cum iniquitáte. Cor hóminis dispónit viam suam, sed Dómini est dirígere gressus ejus.\par
\switchcolumn
\EN
\noindent The beginning of a good way is to do justice; and this is more acceptable with God, than to offer sacrifices. By mercy and truth iniquity is redeemed: and by the fear of the Lord men depart from evil. When the ways of man shall please the Lord, he will convert even his enemies to peace. Better is a little with justice, than great revenues with iniquity. The heart of man disposeth his way: but the Lord must direct his steps.\par
\switchcolumn*

\LA
\begin{Resp}\Rsign Magna enim sunt judícia tua, Dómine, et inenarrabília verba tua: \Star{} Magnificásti pópulum tuum et honorásti.\end{Resp}
\begin{Resp}\Vsign Transtulísti illos per Mare Rubrum et transvexísti eos per aquam nímiam.\end{Resp}
\begin{Resp}\Rsign Magnificásti pópulum tuum et honorásti.\end{Resp}
\switchcolumn
\EN
\begin{Resp}\Rsign Great are thy judgments, O Lord, and thy words cannot be expressed. \Star{} Thou didst make thy people mighty and honourable.\end{Resp}
\begin{Resp}\Vsign Thou broughtest them through the Red Sea, and leddest them through much water.\end{Resp}
\begin{Resp}\Rsign Thou didst make thy people mighty and honourable.\end{Resp}
\switchcolumn*

\LA
\LessonHead{Lectio III}
\switchcolumn
\EN
\LessonHead{Lesson III}
\switchcolumn*

\LA
\Cite{Prov 16:10-15}
\switchcolumn
\EN
\Cite{Prov 16:10-15}
\switchcolumn*

\LA
\noindent Divinátio in lábiis regis, in judício non errábit os ejus. Pondus et statéra judícia Dómini sunt, et ópera ejus omnes lápides sácculi. Abominábiles regi qui agunt ímpie, quóniam justítia firmátur sólium. Volúntas regum lábia justa, qui recta lóquitur diligétur. Indignátio regis núntii mortis, et vir sápiens placábit eam. In hilaritáte vultus regis vita, et cleméntia ejus quasi imber serótinus.\par
\switchcolumn
\EN
\noindent Divination is in the lips of the king, his mouth shall not err in judgment. Weight and balance are judgments of the Lord: and his work all the weights of the bag. They that act wickedly are abominable to the king: for the throne is established by justice. Just lips are the delight of kings: he that speaketh right things shall be loved. The wrath of a king is as messengers of death: and the wise man will pacify it. In the cheerfulness of the king's countenance is life: and his clemency is like the latter rain.\par
\switchcolumn*

\LA
\begin{Resp}\Rsign Quæ sunt in corde hóminum, óculi tui vident, Dómine, et in libro tuo ómnia scribéntur: \Star{} Homo videt in fácie, Deus autem in corde.\end{Resp}
\begin{Resp}\Vsign Omnia enim corda scrutátur, et univérsas méntium cogitatiónes intéllegit.\end{Resp}
\begin{Resp}\Rsign Homo videt in fácie, Deus autem in corde.\end{Resp}
\begin{Resp}\Vsign Glória Patri, et Fílio, \Star{} et Spirítui Sancto.\end{Resp}
\begin{Resp}\Rsign Homo videt in fácie, Deus autem in corde.\end{Resp}
\switchcolumn
\EN
\begin{Resp}\Rsign Lord, thine eyes behold all that is in the heart of man, and in thy book are they all written. \Star{} Man looketh on the outward appearance, but God looketh on the heart.\end{Resp}
\begin{Resp}\Vsign For He searcheth all hearts, and understandeth all the imaginations of the thoughts.\end{Resp}
\begin{Resp}\Rsign Man looketh on the outward appearance, but God looketh on the heart.\end{Resp}
\begin{Resp}\Vsign Glory be to the Father, and to the Son, \Star{} and to the Holy Ghost.\end{Resp}
\begin{Resp}\Rsign Man looketh on the outward appearance, but God looketh on the heart.\end{Resp}
\switchcolumn*

\end{paracol}

\Office{Dominica II Augusti}{Second Sunday of August}{}
\addcontentsline{toc}{subsection}{Dominica II Augusti \textit{— Second Sunday of August}}
\begin{paracol}{2}
\LA
\LessonHead{Lectio I}
\switchcolumn
\EN
\LessonHead{Lesson I}
\switchcolumn*

\LA
\LessonTitle{Incipit liber Ecclesiástes}
\switchcolumn
\EN
\LessonTitle{Lesson from the book of Ecclesiastes}
\switchcolumn*

\LA
\Cite{Eccl 1:1-7}
\switchcolumn
\EN
\Cite{Eccl 1:1-7}
\switchcolumn*

\LA
\noindent Verba Ecclesiástæ, fílii David, regis Jerúsalem. Vánitas vanitátum, dixit Ecclesiástes; vánitas vanitátum, et ómnia vánitas. Quid habet ámplius homo de univérso labóre suo quo labórat sub sole? Generátio prǽterit, et generátio ádvenit; terra autem in ætérnum stat. Oritur sol et óccidit et ad locum suum revértitur; ibíque renáscens gyrat per merídiem et fléctitur ad aquilónem. Lustrans univérsa in circúitu pergit spíritus et in círculos suos revértitur. Omnia flúmina intrant in mare, et mare non redúndat; ad locum unde éxeunt flúmina revertúntur, ut íterum fluant.\par
\switchcolumn
\EN
\noindent The words of Ecclesiastes, the son of David, king of Jerusalem. Vanity of vanities, said Ecclesiastes vanity of vanities, and all is vanity. What hath a man more of all his labour, that he taketh under the sun? One generation passeth away, and another generation cometh: but the earth standeth for ever. The sun riseth, and goeth down, and returneth to his place: and there rising again, Maketh his round by the south, and turneth again to the north: the spirit goeth forward surveying all places round about, and returneth to his circuits. All the rivers run into the sea, yet the sea doth not overflow: unto the place from whence the rivers come, they return, to flow again.\par
\switchcolumn*

\LA
\begin{Resp}\Rsign In princípio Deus ántequam terram fáceret, priúsquam abýssos constitúeret, priúsquam prodúceret fontes aquárum, \Star{} Antequam montes collocaréntur, ante omnes colles generávit me Dóminus.\end{Resp}
\begin{Resp}\Vsign Quando præparábat cælos, áderam, cum eo cuncta compónens.\end{Resp}
\begin{Resp}\Rsign Antequam montes collocaréntur, ante omnes colles generávit me Dóminus.\end{Resp}
\switchcolumn
\EN
\begin{Resp}\Rsign God possessed me in the beginning, before He made the earth, before He created the depths, before He caused the fountains of water to spring. \Star{} Before the mountains were settled, before there were any hills, did the Lord beget me.\end{Resp}
\begin{Resp}\Vsign When He prepared the heavens, I was there with Him, ordering all things.\end{Resp}
\begin{Resp}\Rsign Before the mountains were settled, before there were any hills, did the Lord beget me.\end{Resp}
\switchcolumn*

\LA
\LessonHead{Lectio II}
\switchcolumn
\EN
\LessonHead{Lesson II}
\switchcolumn*

\LA
\Cite{Eccl 1:8-11}
\switchcolumn
\EN
\Cite{Eccl 1:8-11}
\switchcolumn*

\LA
\noindent Cunctæ res diffíciles: non potest eas homo explicáre sermóne. Non saturátur óculus visu, nec auris audítu implétur. Quid est quod fuit? Ipsum quod futúrum est. Quid est quod factum est? Ipsum quod faciéndum est. Nihil sub sole novum, nec valet quisquam dícere: Ecce hoc recens est; jam enim præcéssit in sǽculis quæ fuérunt ante nos. Non est priórum memória; sed nec eórum quidem quæ póstea futúra sunt, erit recordátio apud eos qui futúri sunt in novíssimo.\par
\switchcolumn
\EN
\noindent All things are hard: man cannot explain them by word. The eye is not filled with seeing, neither is the ear filled with hearing. What is it that hath been? the same thing that shall be. What is it that hath been done? the same that shall be done. Nothing under the sun is new, neither is any man able to say: Behold this is new: for it hath already gone before in the ages that were before us. There is no remembrance of former things: nor indeed of those things which hereafter are to come, shall there be any remembrance with them that shall be in the latter end.\par
\switchcolumn*

\LA
\begin{Resp}\Rsign Gyrum cæli circuívi sola, et in flúctibus maris ambulávi, in omni gente et in omni pópulo primátum ténui: \Star{} Superbórum et sublímium colla própria virtúte calcávi.\end{Resp}
\begin{Resp}\Vsign Ego in altíssimis hábito, et thronus meus in colúmna nubis.\end{Resp}
\begin{Resp}\Rsign Superbórum et sublímium colla própria virtúte calcávi.\end{Resp}
\switchcolumn
\EN
\begin{Resp}\Rsign I alone compassed the circuit of heaven, and walked on the waves of the sea. In every nation and in every people, I held the first place. \Star{} In the greatness of my strength have I trodden under my feet the necks of such as be haughty and proud.\end{Resp}
\begin{Resp}\Vsign I dwell in the highest places, and my throne is in a cloudy pillar.\end{Resp}
\begin{Resp}\Rsign In the greatness of my strength have I trodden under my feet the necks of such as be haughty and proud.\end{Resp}
\switchcolumn*

\LA
\LessonHead{Lectio III}
\switchcolumn
\EN
\LessonHead{Lesson III}
\switchcolumn*

\LA
\Cite{Eccl 1:12-17}
\switchcolumn
\EN
\Cite{Eccl 1:12-17}
\switchcolumn*

\LA
\noindent Ego Ecclesiástes fui rex Israël in Jerúsalem; et propósui in ánimo meo quǽrere et investigáre sapiénter de ómnibus quæ fiunt sub sole. Hanc occupatiónem péssimam dedit Deus fíliis hóminum, ut occuparéntur in ea. Vidi cuncta quæ fiunt sub sole, et ecce univérsa vánitas et afflíctio spíritus. Pervérsi diffícile corrigúntur, et stultórum infinítus est númerus. Locútus sum in corde meo dicens: Ecce magnus efféctus sum et præcéssi omnes sapiéntia qui fuérunt ante me in Jerúsalem; et mens mea contempláta est multa sapiénter, et dídici. Dedíque cor meum ut scirem prudéntiam atque doctrínam, errorésque et stultítiam.\par
\switchcolumn
\EN
\noindent I Ecclesiastes was king over Israel in Jerusalem, And I proposed in my mind to seek and search out wisely concerning all things that are done under the sun. This painful occupation hath God given to the children of men, to be exercised therein. I have seen all things that are done under the sun, and behold all is vanity, and vexation of spirit. The perverse are hard to be corrected, and the number of fools is infinite. I have spoken in my heart, saying: Behold I am become great, and have gone beyond all in wisdom, that were before me in Jerusalem: and my mind hath contemplated many things wisely, and I have learned. And I have given my heart to know prudence, and learning, and errors, and folly: and I have perceived that in these also there was labour, and vexation of spirit.\par
\switchcolumn*

\LA
\begin{Resp}\Rsign Emítte, Dómine, sapiéntiam de sede magnitúdinis tuæ, ut mecum sit et mecum labóret: \Star{} Ut sciam, quid accéptum sit coram te omni témpore.\end{Resp}
\begin{Resp}\Vsign Da mihi, Dómine, sédium tuárum assistrícem sapiéntiam.\end{Resp}
\begin{Resp}\Rsign Ut sciam quid accéptum sit coram te omni témpore.\end{Resp}
\begin{Resp}\Vsign Glória Patri, et Fílio, \Star{} et Spirítui Sancto.\end{Resp}
\begin{Resp}\Rsign Ut sciam quid accéptum sit coram te omni témpore.\end{Resp}
\switchcolumn
\EN
\begin{Resp}\Rsign O send out wisdom from the throne of thy glory, O Lord, to be with me, and to labour with me, \Star{} That I may know at all times what is pleasing unto thee.\end{Resp}
\begin{Resp}\Vsign Give me wisdom, O Lord, that sitteth by thy throne.\end{Resp}
\begin{Resp}\Rsign That I may know at all times what is pleasing unto thee.\end{Resp}
\begin{Resp}\Vsign Glory be to the Father, and to the Son, \Star{} and to the Holy Ghost.\end{Resp}
\begin{Resp}\Rsign That I may know at all times what is pleasing unto thee.\end{Resp}
\switchcolumn*

\LA
\LessonHead{Lectio IV}
\switchcolumn
\EN
\LessonHead{Lesson IV}
\switchcolumn*

\LA
\LessonTitle{Sermo sancti Joánnis Chrysóstomi}
\switchcolumn
\EN
\LessonTitle{From the Sermons of St. John Chrysostom, Patriarch of Constantinople.}
\switchcolumn*

\LA
\Source{Sermo contra concubinarios, in fine, tomo 5}
\switchcolumn
\EN
\Source{Sermon against concubinage}
\switchcolumn*

\LA
\noindent Sálomon cum sæculárium rerum concupiscéntia tenerétur, magnas eas et admirándas putábat, multúmque in eis labóris et sollicitúdinis insumébat, magníficas ædificándo domos, copiósum coacervándo aurum, congregándo cantórum choros, vária génera ministrórum mensæ et popínæ, quæréndo ánimæ suæ voluptátem ab hortórum et córporum formosórum grátia, et omnem, ut ita dicam, oblectatiónis et refrigérii viam sectándo.\par
\switchcolumn
\EN
\noindent While Solomon was given up to the lust of the world, he deemed the same a great and noble pursuit, and expended thereon great labour and care. He built magnificent palaces, he heaped up gold in plenty, he gathered together choirs of singers, and all sorts of servants to minister to the luxury of his table and of his fare. He sought enjoyment for his heart from the charm of gardens and of fair bodies. In short, he gave himself up to the study of all kinds of pleasure and recreation.\par
\switchcolumn*

\LA
\begin{Resp}\Rsign Da mihi, Dómine, sédium tuárum assistrícem sapiéntiam, et noli me reprobáre a púeris tuis: \Star{} Quóniam servus tuus sum ego, et fílius ancíllæ tuæ.\end{Resp}
\begin{Resp}\Vsign Mitte illam de sede magnitúdinis tuæ, ut mecum sit et mecum labóret.\end{Resp}
\begin{Resp}\Rsign Quóniam servus tuus sum ego, et fílius ancíllæ tuæ.\end{Resp}
\switchcolumn
\EN
\begin{Resp}\Rsign Give me wisdom, O Lord, that sitteth by thy throne, and reject me not from among thy children. \Star{} For I am thy servant and son of thine handmaid.\end{Resp}
\begin{Resp}\Vsign O send her out from the throne of thy glory, to be with me and to labour with me.\end{Resp}
\begin{Resp}\Rsign For I am thy servant and son of thine handmaid.\end{Resp}
\switchcolumn*

\LA
\LessonHead{Lectio V}
\switchcolumn
\EN
\LessonHead{Lesson V}
\switchcolumn*

\LA
\noindent At ubi inde ad se revérsus, et quasi ex umbrósa quadam abýsso ad lumen veræ sapiéntiæ respícere váluit, tunc sublímem illam et cælis dignam emísit vocem: Vánitas vanitátum, dicens, et ómnia vánitas. Hanc et vos, et hac sublimiórem, si voluéritis, efferétis senténtiam de intempestíva hac voluptáte, si aliquantísper a mala consuetúdine vos sequestravéritis.\par
\switchcolumn
\EN
\noindent But when he came to himself again, and was once more able, as it were, out of that dark pit, to look upon the light of true wisdom, he uttered that saying, so high, so worthy of heaven “Vanity of vanities; all is vanity.” And ye also, if ever ye will shake yourselves clear of your debasing habit, will utter this cry, and an higher cry than this, as ye turn from your untimely indulgences.\par
\switchcolumn*

\LA
\begin{Resp}\Rsign Inítium sapiéntiæ timor Dómini: \Star{} Intelléctus bonus ómnibus faciéntibus eum: laudátio ejus manet in sǽculum sǽculi.\end{Resp}
\begin{Resp}\Vsign Diléctio illíus custódia legum est: quia omnis sapiéntia timor Dómini.\end{Resp}
\begin{Resp}\Rsign Intelléctus bonus ómnibus faciéntibus eum: laudátio ejus manet in sǽculum sǽculi.\end{Resp}
\switchcolumn
\EN
\begin{Resp}\Rsign The fear of the Lord is the beginning of wisdom. \Star{} A good understanding have all they that do His commandments. His praise endureth for ever.\end{Resp}
\begin{Resp}\Vsign Love is the keeping of her laws, for all wisdom is the fear of the Lord.\end{Resp}
\begin{Resp}\Rsign A good understanding have all they that do His commandments. His praise endureth for ever.\end{Resp}
\switchcolumn*

\LA
\LessonHead{Lectio VI}
\switchcolumn
\EN
\LessonHead{Lesson VI}
\switchcolumn*

\LA
\noindent Quamvis autem a Salomóne sǽculis superióribus non tam multa sapiéntiæ exigebátur diligéntia: neque enim delícias lex vetus prohibébat, neque áliis frui supervácuis dicébat vanum: áttamen et sic se habéntibus rebus, in ipsis contuéri licébit, quam viles et vanitáti obnóxiæ res sint. Nos vero ad majórem vocáti vitam, et ad excelléntius fastígium ascéndimus, et in majóribus exercémur palǽstris: et quid áliud, quam quod, sicut supérnæ virtútes intellectuáles et incorpóreæ illæ, vitam institúere jubémur?\par
\switchcolumn
\EN
\noindent The ages that had rolled before the time of Solomon had not left to his own so precious an inheritance of wisdom as those which have preceded us have left to us; the old law did not forbid these indulgences, nor pronounce it folly to enjoy other idle luxuries and yet, even with matters so, we can see how low, how worthless, such things be. As for us, we are called to a higher life, we ascend to a nobler stand-point, and brace ourselves in a manlier school and why, but because we are bidden to strive for a life like the life of the spiritual and bodiless powers.\par
\switchcolumn*

\LA
\begin{Resp}\Rsign Verbum iníquum et dolósum longe fac a me, Dómine: \Star{} Divítias et paupertátem ne déderis mihi, sed tantum víctui meo tríbue necessária.\end{Resp}
\begin{Resp}\Vsign Duo rogávi te, ne déneges mihi ántequam móriar.\end{Resp}
\begin{Resp}\Rsign Divítias et paupertátem ne déderis mihi, sed tantum víctui meo tríbue necessária.\end{Resp}
\begin{Resp}\Vsign Glória Patri, et Fílio, \Star{} et Spirítui Sancto.\end{Resp}
\begin{Resp}\Rsign Divítias et paupertátem ne déderis mihi, sed tantum víctui meo tríbue necessária.\end{Resp}
\switchcolumn
\EN
\begin{Resp}\Rsign Lord, remove far from me vanity and lies. \Star{} Give me neither poverty nor riches, but feed me with food convenient for me.\end{Resp}
\begin{Resp}\Vsign Two things have I required of thee; deny me them not before I die.\end{Resp}
\begin{Resp}\Rsign Give me neither poverty nor riches, but feed me with food convenient for me.\end{Resp}
\begin{Resp}\Vsign Glory be to the Father, and to the Son, \Star{} and to the Holy Ghost.\end{Resp}
\begin{Resp}\Rsign Give me neither poverty nor riches, but feed me with food convenient for me.\end{Resp}
\switchcolumn*

\end{paracol}

\Office{Feria II infra Hebdomadam II Augusti}{Monday of the second week of August}{}
\addcontentsline{toc}{subsection}{Feria II infra Hebdomadam II Augusti \textit{— Monday of the second week of August}}
\begin{paracol}{2}
\LA
\LessonHead{Lectio I}
\switchcolumn
\EN
\LessonHead{Lesson I}
\switchcolumn*

\LA
\LessonTitle{De libro Ecclesiástæ}
\switchcolumn
\EN
\LessonTitle{Lesson from the book of Ecclesiastes}
\switchcolumn*

\LA
\Cite{Eccl 2:1-4}
\switchcolumn
\EN
\Cite{Eccl 2:1-4}
\switchcolumn*

\LA
\noindent Dixi ego in corde meo: Vadam et áffluam delíciis et fruar bonis; et vidi quod hoc quoque esset vánitas. Risum reputávi errórem et gáudio dixi: Quid frustra decíperis? Cogitávi in corde meo abstráhere a vino carnem meam, ut ánimum meum transférrem ad sapiéntiam, devitarémque stultítiam, donec vidérem quid esset útile fíliis hóminum, quo facto opus est sub sole número diérum vitæ suæ. Magnificávi ópera mea, ædificávi mihi domos et plantávi víneas.\par
\switchcolumn
\EN
\noindent I said in my heart: I will go, and abound with delights, and enjoy good things. And I saw that this also was vanity. Laughter I counted error: and to mirth I said: Why art thou vainly deceived? I thought in my heart, to withdraw my flesh from wine, that I might turn my mind to wisdom, and might avoid folly, till I might see what was profitable for the children of men: and what they ought to do under the sun, all the days of their life. I made me great works, I built me houses, and planted vineyards.\par
\switchcolumn*

\LA
\begin{Resp}\Rsign Ne derelínquas me, Dómine, pater et dominátor vitæ meæ, ut non córruam in conspéctu adversariórum meórum: \Star{} Ne gáudeat de me inimícus meus.\end{Resp}
\begin{Resp}\Vsign Apprehénde arma et scutum et exsúrge in adjutórium mihi.\end{Resp}
\begin{Resp}\Rsign Ne gáudeat de me inimícus meus.\end{Resp}
\switchcolumn
\EN
\begin{Resp}\Rsign O Lord, Father and Governor of my life, leave me not, lest I fall before mine adversaries, \Star{} and mine enemy rejoice over me.\end{Resp}
\begin{Resp}\Vsign Take hold of shield and buckler, and stand up for mine help.\end{Resp}
\begin{Resp}\Rsign Lest mine enemy rejoice over me.\end{Resp}
\switchcolumn*

\LA
\LessonHead{Lectio II}
\switchcolumn
\EN
\LessonHead{Lesson II}
\switchcolumn*

\LA
\Cite{Eccl 2:7-9}
\switchcolumn
\EN
\Cite{Eccl 2:7-9}
\switchcolumn*

\LA
\noindent Possédi servos et ancíllas multámque famíliam hábui, arménta quoque et magnos óvium greges, ultra omnes qui fuérunt ante me in Jerúsalem; coacervávi mihi argéntum et aurum et substántias regum ac provinciárum; feci mihi cantóres et cantatríces et delícias filiórum hóminum, scyphos et úrceos in ministério ad vina fundénda, et supergréssus sum ópibus omnes qui ante me fuérunt in Jerúsalem: sapiéntia quoque perseverávit mecum.\par
\switchcolumn
\EN
\noindent I got me menservants, and maidservants, and had a great family: and herds of oxen, and great flocks of sheep, above all that were before me in Jerusalem: I heaped together for myself silver and gold, and the wealth of kings, and provinces: I made me singing men, and singing women, and the delights of the sons of men, cups and vessels to serve to pour out wine: And I surpassed in riches all that were before me in Jerusalem: my wisdom also remained with me.\par
\switchcolumn*

\LA
\begin{Resp}\Rsign Magna enim sunt judícia tua, Dómine, et inenarrabília verba tua: \Star{} Magnificásti pópulum tuum et honorásti.\end{Resp}
\begin{Resp}\Vsign Transtulísti illos per Mare Rubrum et transvexísti eos per aquam nímiam.\end{Resp}
\begin{Resp}\Rsign Magnificásti pópulum tuum et honorásti.\end{Resp}
\switchcolumn
\EN
\begin{Resp}\Rsign Great are thy judgments, O Lord, and thy words cannot be expressed. \Star{} Thou didst make thy people mighty and honourable.\end{Resp}
\begin{Resp}\Vsign Thou broughtest them through the Red Sea, and leddest them through much water.\end{Resp}
\begin{Resp}\Rsign Thou didst make thy people mighty and honourable.\end{Resp}
\switchcolumn*

\LA
\LessonHead{Lectio III}
\switchcolumn
\EN
\LessonHead{Lesson III}
\switchcolumn*

\LA
\Cite{Eccl 2:10-11}
\switchcolumn
\EN
\Cite{Eccl 2:10-11}
\switchcolumn*

\LA
\noindent Et ómnia, quæ desideravérunt óculi mei, non negávi eis, nec prohíbui cor meum, quin omni voluptáte fruerétur et oblectáret se in his quæ præparáveram; et hanc ratus sum partem meam si úterer labóre meo. Cumque me convertíssem ad univérsa ópera, quæ fécerant manus meæ, et ad labóres in quibus frustra sudáveram, vidi in ómnibus vanitátem et afflictiónem ánimi, et nihil permanére sub sole.\par
\switchcolumn
\EN
\noindent And whatsoever my eyes desired, I refused them not: and I withheld not my heart from enjoying every pleasure, and delighting itself in the things which I had prepared: and esteemed this my portion, to make use of my own labour. And when I turned myself to all the works which my hands had wrought, and to the labours wherein I had laboured in vain, I saw in all things vanity, and vexation of mind, and that nothing was lasting under the sun.\par
\switchcolumn*

\LA
\begin{Resp}\Rsign Quæ sunt in corde hóminum, óculi tui vident, Dómine, et in libro tuo ómnia scribéntur: \Star{} Homo videt in fácie, Deus autem in corde.\end{Resp}
\begin{Resp}\Vsign Omnia enim corda scrutátur, et univérsas méntium cogitatiónes intéllegit.\end{Resp}
\begin{Resp}\Rsign Homo videt in fácie, Deus autem in corde.\end{Resp}
\begin{Resp}\Vsign Glória Patri, et Fílio, \Star{} et Spirítui Sancto.\end{Resp}
\begin{Resp}\Rsign Homo videt in fácie, Deus autem in corde.\end{Resp}
\switchcolumn
\EN
\begin{Resp}\Rsign Lord, thine eyes behold all that is in the heart of man, and in thy book are they all written. \Star{} Man looketh on the outward appearance, but God looketh on the heart.\end{Resp}
\begin{Resp}\Vsign For He searcheth all hearts, and understandeth all the imaginations of the thoughts.\end{Resp}
\begin{Resp}\Rsign Man looketh on the outward appearance, but God looketh on the heart.\end{Resp}
\begin{Resp}\Vsign Glory be to the Father, and to the Son, \Star{} and to the Holy Ghost.\end{Resp}
\begin{Resp}\Rsign Man looketh on the outward appearance, but God looketh on the heart.\end{Resp}
\switchcolumn*

\end{paracol}

\Office{Feria III infra Hebdomadam II Augusti}{Tuesday of the second week of August}{}
\addcontentsline{toc}{subsection}{Feria III infra Hebdomadam II Augusti \textit{— Tuesday of the second week of August}}
\begin{paracol}{2}
\LA
\LessonHead{Lectio I}
\switchcolumn
\EN
\LessonHead{Lesson I}
\switchcolumn*

\LA
\LessonTitle{De libro Ecclesiástæ}
\switchcolumn
\EN
\LessonTitle{Lesson from the book of Ecclesiastes}
\switchcolumn*

\LA
\Cite{Eccl 3:1-8}
\switchcolumn
\EN
\Cite{Eccl 3:1-8}
\switchcolumn*

\LA
\noindent Omnia tempus habent, et suis spátiis tránseunt univérsa sub cælo. Tempus nascéndi et tempus moriéndi, tempus plantándi, et tempus evelléndi quod plantátum est; tempus occidéndi et tempus sanándi; tempus destruéndi, et tempus ædificándi; tempus flendi et tempus ridéndi, tempus plangéndi, et tempus saltándi; tempus spargéndi lápides, et tempus colligéndi, tempus amplexándi et tempus longe fíeri ab ampléxibus, tempus acquiréndi, et tempus perdéndi; tempus custodiéndi, et tempus abiciéndi; tempus scindéndi, et tempus consuéndi; tempus tacéndi, et tempus loquéndi; tempus dilectiónis, et tempus ódii; tempus belli, et tempus pacis.\par
\switchcolumn
\EN
\noindent All things have their season, and in their times all things pass under heaven. A time to be born and a time to die. A time to plant, and a time to pluck up that which is planted. A time to kill, and a time to heal. A time to destroy, and a time to build. A time to weep, and a time to laugh. A time to mourn, and a time to dance. A time to scatter stones, and a time to gather. A time to embrace, and a time to be far from embraces. A time to get, and a time to lose. A time to keep, and a time to cast away. A time to rend, and a time to sew. A time to keep silence, and a time to speak. A time of love, and a time of hatred. A time of war, and a time of peace.\par
\switchcolumn*

\LA
\begin{Resp}\Rsign Præbe, fili, cor mihi, et óculi tui vias meas custódiant: \Star{} Ut addátur grátia cápiti tuo.\end{Resp}
\begin{Resp}\Vsign Atténde, fili mi, sapiéntiam meam et ad elóquium meum inclína aurem tuam.\end{Resp}
\begin{Resp}\Rsign Ut addátur grátia cápiti tuo.\end{Resp}
\switchcolumn
\EN
\begin{Resp}\Rsign My son, give me thine heart, and let thine eyes observe my ways. \Star{} For they shall be an ornament of grace unto thine head.\end{Resp}
\begin{Resp}\Vsign My son, attend unto my wisdom, and incline thine ear unto my sayings.\end{Resp}
\begin{Resp}\Rsign For they shall be an ornament of grace unto thine head.\end{Resp}
\switchcolumn*

\LA
\LessonHead{Lectio II}
\switchcolumn
\EN
\LessonHead{Lesson II}
\switchcolumn*

\LA
\Cite{Eccl 3:9-13}
\switchcolumn
\EN
\Cite{Eccl 3:9-13}
\switchcolumn*

\LA
\noindent Quid habet ámplius homo de labóre suo? Vidi afflictiónem, quam dedit Deus fíliis hóminum, ut distendántur in ea. Cuncta fecit bona in témpore suo et mundum trádidit disputatióni eórum, ut non invéniat homo opus quod operátus est Deus ab inítio usque ad finem. Et cognóvi quod non esset mélius nisi lætári et fácere bene in vita sua; omnis enim homo qui cómedit et bibit, et videt bonum de labóre suo, hoc donum Dei est.\par
\switchcolumn
\EN
\noindent What hath man more of his labour? I have seen the trouble, which God hath given the sons of men to be exercised in it. He hath made all things good in their time, and hath delivered the world to their consideration, so that man cannot find out the work which God hath made from the beginning to the end. And I have known that there was no better thing than to rejoice, and to do well in this life. For every man that eateth and drinketh, and seeth good of his labour, this is the gift of God.\par
\switchcolumn*

\LA
\begin{Resp}\Rsign Inítium sapiéntiæ timor Dómini: \Star{} Intelléctus bonus ómnibus faciéntibus eum: laudátio ejus manet in sǽculum sǽculi.\end{Resp}
\begin{Resp}\Vsign Diléctio illíus custódia legum est: quia omnis sapiéntia timor Dómini.\end{Resp}
\begin{Resp}\Rsign Intelléctus bonus ómnibus faciéntibus eum: laudátio ejus manet in sǽculum sǽculi.\end{Resp}
\switchcolumn
\EN
\begin{Resp}\Rsign The fear of the Lord is the beginning of wisdom. \Star{} A good understanding have all they that do His commandments. His praise endureth for ever.\end{Resp}
\begin{Resp}\Vsign Love is the keeping of her laws, for all wisdom is the fear of the Lord.\end{Resp}
\begin{Resp}\Rsign A good understanding have all they that do His commandments. His praise endureth for ever.\end{Resp}
\switchcolumn*

\LA
\LessonHead{Lectio III}
\switchcolumn
\EN
\LessonHead{Lesson III}
\switchcolumn*

\LA
\Cite{Eccl 3:14-17}
\switchcolumn
\EN
\Cite{Eccl 3:14-17}
\switchcolumn*

\LA
\noindent Dídici quod ómnia ópera, quæ fecit Deus, persevérent in perpétuum; non póssumus eis quidquam áddere nec auférre, quæ fecit Deus ut timeátur. Quod factum est ipsum pérmanet, quæ futúra sunt jam fuérunt, et Deus instáurat quod ábiit. Vidi sub sole in loco judícii impietátem et in loco justítiæ iniquitátem: et dixi in corde meo: Justum et ímpium judicábit Deus, et tempus omnis rei tunc erit.\par
\switchcolumn
\EN
\noindent I have learned that all the works which God hath made, continue for ever: we cannot add any thing, nor take away from those things which God hath made that he may be feared. That which hath been made, the same continueth: the things that shall be, have already been: and God restoreth that which is past. I saw under the sun in the place of judgment wickedness, and in the place of justice iniquity. And I said in my heart: God shall judge both the just and the wicked, and then shall be the time of every thing.\par
\switchcolumn*

\LA
\begin{Resp}\Rsign Verbum iníquum et dolósum longe fac a me, Dómine: \Star{} Divítias et paupertátem ne déderis mihi, sed tantum víctui meo tríbue necessária.\end{Resp}
\begin{Resp}\Vsign Duo rogávi te, ne déneges mihi ántequam móriar.\end{Resp}
\begin{Resp}\Rsign Divítias et paupertátem ne déderis mihi, sed tantum víctui meo tríbue necessária.\end{Resp}
\begin{Resp}\Vsign Glória Patri, et Fílio, \Star{} et Spirítui Sancto.\end{Resp}
\begin{Resp}\Rsign Divítias et paupertátem ne déderis mihi, sed tantum víctui meo tríbue necessária.\end{Resp}
\switchcolumn
\EN
\begin{Resp}\Rsign Lord, remove far from me vanity and lies. \Star{} Give me neither poverty nor riches, but feed me with food convenient for me.\end{Resp}
\begin{Resp}\Vsign Two things have I required of thee; deny me them not before I die.\end{Resp}
\begin{Resp}\Rsign Give me neither poverty nor riches, but feed me with food convenient for me.\end{Resp}
\begin{Resp}\Vsign Glory be to the Father, and to the Son, \Star{} and to the Holy Ghost.\end{Resp}
\begin{Resp}\Rsign Give me neither poverty nor riches, but feed me with food convenient for me.\end{Resp}
\switchcolumn*

\end{paracol}

\Office{Feria IV infra Hebdomadam II Augusti}{Wednesday of the second week of August}{}
\addcontentsline{toc}{subsection}{Feria IV infra Hebdomadam II Augusti \textit{— Wednesday of the second week of August}}
\begin{paracol}{2}
\LA
\LessonHead{Lectio I}
\switchcolumn
\EN
\LessonHead{Lesson I}
\switchcolumn*

\LA
\LessonTitle{De libro Ecclesiástæ}
\switchcolumn
\EN
\LessonTitle{Lesson from the book of Ecclesiastes}
\switchcolumn*

\LA
\Cite{Eccl 4:1-4}
\switchcolumn
\EN
\Cite{Eccl 4:1-4}
\switchcolumn*

\LA
\noindent Verti me ad ália, et vidi calúmnias, quæ sub sole gerúntur, et lácrimas innocéntium, et néminem consolatórem, nec posse resístere eórum violéntiæ cunctórum auxílio destitútos. Et laudávi magis mórtuos quam vivéntes; et feliciórem utróque judicávi qui necdum natus est nec vidit mala quæ sub sole fiunt. Rursum contemplátus sum omnes labóres hóminum, et indústrias animadvérti patére invídiæ próximi; et in hoc ergo vánitas et cura supérflua est.\par
\switchcolumn
\EN
\noindent I turned myself to other things, and I saw the oppressions that are done under the sun, and the tears of the innocent, and they had no comforter; and they were not able to resist their violence, being destitute of help from any. And I praised the dead rather than the living: And I judged him happier than them both, that is not yet born, nor hath seen the evils that are done under the sun. Again I considered all the labours of men, and I remarked that their industries are exposed to the envy of their neighhour: so in this also there is vanity, and fruitless care.\par
\switchcolumn*

\LA
\begin{Resp}\Rsign Dómine, Pater et Deus vitæ meæ, ne derelínquas me in cogitátu malígno: extolléntiam oculórum meórum ne déderis mihi, et desidérium malígnum avérte a me, Dómine; aufer a me concupiscéntiam, \Star{} Et ánimo irreverénti et infruníto ne tradas me, Dómine.\end{Resp}
\begin{Resp}\Vsign Ne derelínquas me, Dómine, ne accréscant ignorántiæ meæ, nec multiplicéntur delícta mea.\end{Resp}
\begin{Resp}\Rsign Et ánimo irreverénti et infruníto ne tradas me, Dómine.\end{Resp}
\switchcolumn
\EN
\begin{Resp}\Rsign O Lord, Father and God of my life, leave me not to evil counsels; give me not a proud look, but turn away from me an haughty mind, O Lord turn away from me concupiscence, \Star{} And give me not over unto an impudent and froward mind, O Lord!\end{Resp}
\begin{Resp}\Vsign Leave me not, O Lord, lest mine ignorance increase, and my sins abound.\end{Resp}
\begin{Resp}\Rsign And give me not over unto an impudent and froward mind, O Lord.\end{Resp}
\switchcolumn*

\LA
\LessonHead{Lectio II}
\switchcolumn
\EN
\LessonHead{Lesson II}
\switchcolumn*

\LA
\Cite{Eccl 4:5-8}
\switchcolumn
\EN
\Cite{Eccl 4:5-8}
\switchcolumn*

\LA
\noindent Stultus cómplicat manus suas et cómedit carnes suas, dicens: Mélior est pugíllus cum réquie, quam plena útraque manus cum labóre et afflictióne ánimi. Consíderans, réperi et áliam vanitátem sub sole. Unus est, et secúndum non habet, non fílium, non fratrem, et tamen laboráre non cessat, nec satiántur óculi ejus divítiis, nec recógitat, dicens: Cui labóro, et fraudo ánimam meam bonis? In hoc quoque vánitas est et afflíctio péssima.\par
\switchcolumn
\EN
\noindent The fool foldeth his hands together, and eateth his own flesh, saying: Better is a handful with rest, than both hands full with labour, and vexation of mind. Considering I found also another vanity under the sun: There is but one, and he hath not a second, no child, no brother, and yet he ceaseth not to labour, neither are his eyes satisfied with riches, neither doth he reflect, saying: For whom do I labour, and defraud my soul of good things? in this also is vanity, and a grievous vexation.\par
\switchcolumn*

\LA
\begin{Resp}\Rsign Magna enim sunt judícia tua, Dómine, et inenarrabília verba tua: \Star{} Magnificásti pópulum tuum et honorásti.\end{Resp}
\begin{Resp}\Vsign Transtulísti illos per Mare Rubrum et transvexísti eos per aquam nímiam.\end{Resp}
\begin{Resp}\Rsign Magnificásti pópulum tuum et honorásti.\end{Resp}
\switchcolumn
\EN
\begin{Resp}\Rsign Great are thy judgments, O Lord, and thy words cannot be expressed. \Star{} Thou didst make thy people mighty and honourable.\end{Resp}
\begin{Resp}\Vsign Thou broughtest them through the Red Sea, and leddest them through much water.\end{Resp}
\begin{Resp}\Rsign Thou didst make thy people mighty and honourable.\end{Resp}
\switchcolumn*

\LA
\LessonHead{Lectio III}
\switchcolumn
\EN
\LessonHead{Lesson III}
\switchcolumn*

\LA
\Cite{Eccl 4:9-13}
\switchcolumn
\EN
\Cite{Eccl 4:9-13}
\switchcolumn*

\LA
\noindent Mélius est ergo duos esse simul quam unum; habent enim emoluméntum societátis suæ. Si unus cecíderit, ab áltero fulciétur. Væ soli, quia, cum cecíderit, non habet sublevántem se. Et si dormíerint duo, fovebúntur mútuo: unus quómodo calefíet? Et, si quíspiam prævalúerit contra unum, duo resístunt ei; funículus triplex diffícile rúmpitur. Mélior est puer pauper et sápiens, rege sene et stulto, qui nescit prævidére in pósterum.\par
\switchcolumn
\EN
\noindent It is better therefore that two should be together, than one: for they have the advantage of their society: If one fall he shall be supported by the other: woe to him that is alone, for when he falleth, he hath none to lift him up. And if two lie together, they shall warm one another: how shall one alone be warmed? And if a man prevail against one, two shall withstand him: a threefold cord is not easily broken. Better is a child that is poor and wise, than a king that is old and foolish, who knoweth not to foresee for hereafter.\par
\switchcolumn*

\LA
\begin{Resp}\Rsign Quæ sunt in corde hóminum, óculi tui vident, Dómine, et in libro tuo ómnia scribéntur: \Star{} Homo videt in fácie, Deus autem in corde.\end{Resp}
\begin{Resp}\Vsign Omnia enim corda scrutátur, et univérsas méntium cogitatiónes intéllegit.\end{Resp}
\begin{Resp}\Rsign Homo videt in fácie, Deus autem in corde.\end{Resp}
\begin{Resp}\Vsign Glória Patri, et Fílio, \Star{} et Spirítui Sancto.\end{Resp}
\begin{Resp}\Rsign Homo videt in fácie, Deus autem in corde.\end{Resp}
\switchcolumn
\EN
\begin{Resp}\Rsign Lord, thine eyes behold all that is in the heart of man, and in thy book are they all written. \Star{} Man looketh on the outward appearance, but God looketh on the heart.\end{Resp}
\begin{Resp}\Vsign For He searcheth all hearts, and understandeth all the imaginations of the thoughts.\end{Resp}
\begin{Resp}\Rsign Man looketh on the outward appearance, but God looketh on the heart.\end{Resp}
\begin{Resp}\Vsign Glory be to the Father, and to the Son, \Star{} and to the Holy Ghost.\end{Resp}
\begin{Resp}\Rsign Man looketh on the outward appearance, but God looketh on the heart.\end{Resp}
\switchcolumn*

\end{paracol}

\Office{Feria V infra Hebdomadam II Augusti}{Thursday of the second week of August}{}
\addcontentsline{toc}{subsection}{Feria V infra Hebdomadam II Augusti \textit{— Thursday of the second week of August}}
\begin{paracol}{2}
\LA
\LessonHead{Lectio I}
\switchcolumn
\EN
\LessonHead{Lesson I}
\switchcolumn*

\LA
\LessonTitle{De libro Ecclesiástæ}
\switchcolumn
\EN
\LessonTitle{Lesson from the book of Ecclesiastes}
\switchcolumn*

\LA
\Cite{Eccl 5:1-4}
\switchcolumn
\EN
\Cite{Eccl 5:1-4}
\switchcolumn*

\LA
\noindent Ne témere quid loquáris, neque cor tuum sit velox ad proferéndum sermónem coram Deo. Deus enim in cælo, et tu super terram; idcírco sint pauci sermónes tui. Multas curas sequúntur sómnia, et in multis sermónibus inveniétur stultítia. Si quid vovísti Deo, ne moréris réddere; dísplicet enim ei infidélis et stulta promíssio; sed quodcúmque vóveris redde; multóque mélius est non vovére, quam post votum promíssa non réddere.\par
\switchcolumn
\EN
\noindent Speak not any thing rashly, and let not thy heart be hasty to utter a word before God. For God is in heaven, and thou upon earth: therefore let thy words be few. Dreams follow many cares: and in many words shall be found folly. If thou hast vowed any thing to God, defer not to pay it: for an unfaithful and foolish promise displeaseth him: but whatsoever thou hast vowed, pay it. And it is much better not to vow, than after a vow not to perform the things promised.\par
\switchcolumn*

\LA
\begin{Resp}\Rsign In princípio Deus ántequam terram fáceret, priúsquam abýssos constitúeret, priúsquam prodúceret fontes aquárum, \Star{} Antequam montes collocaréntur, ante omnes colles generávit me Dóminus.\end{Resp}
\begin{Resp}\Vsign Quando præparábat cælos, áderam, cum eo cuncta compónens.\end{Resp}
\begin{Resp}\Rsign Antequam montes collocaréntur, ante omnes colles generávit me Dóminus.\end{Resp}
\switchcolumn
\EN
\begin{Resp}\Rsign God possessed me in the beginning, before He made the earth, before He created the depths, before He caused the fountains of water to spring. \Star{} Before the mountains were settled, before there were any hills, did the Lord beget me.\end{Resp}
\begin{Resp}\Vsign When He prepared the heavens, I was there with Him, ordering all things.\end{Resp}
\begin{Resp}\Rsign Before the mountains were settled, before there were any hills, did the Lord beget me.\end{Resp}
\switchcolumn*

\LA
\LessonHead{Lectio II}
\switchcolumn
\EN
\LessonHead{Lesson II}
\switchcolumn*

\LA
\Cite{Eccl 5:5-8}
\switchcolumn
\EN
\Cite{Eccl 5:5-8}
\switchcolumn*

\LA
\noindent Ne déderis os tuum ut peccáre fácias carnem tuam; neque dicas coram Angelo: Non est providéntia; ne forte irátus Deus contra sermónes tuos díssipet cuncta ópera mánuum tuárum. Ubi multa sunt sómnia, plúrimæ sunt vanitátes et sermónes innúmeri; tu vero Deum time. Si víderis calúmnias egenórum et violénta judícia et subvérti justítiam in província, non miréris super hoc negótio; quia excélso excélsior est álius, et super hos quoque eminentióres sunt álii; et ínsuper univérsæ terræ rex ímperat serviénti.\par
\switchcolumn
\EN
\noindent Give not thy mouth to cause thy flesh to sin: and say not before the angel: There is no providence: lest God be angry at thy words, and destroy all the works of thy hands. Where there are many dreams, there are many vanities, and words without number: but do thou fear God. If thou shalt see the oppressions of the poor, and violent judgments, and justice perverted in the province, wonder not at this matter: for he that is high hath another higher, and there are others still higher than these: Moreover there is the king that reigneth over all the land subject to him.\par
\switchcolumn*

\LA
\begin{Resp}\Rsign Gyrum cæli circuívi sola, et in flúctibus maris ambulávi, in omni gente et in omni pópulo primátum ténui: \Star{} Superbórum et sublímium colla própria virtúte calcávi.\end{Resp}
\begin{Resp}\Vsign Ego in altíssimis hábito, et thronus meus in colúmna nubis.\end{Resp}
\begin{Resp}\Rsign Superbórum et sublímium colla própria virtúte calcávi.\end{Resp}
\switchcolumn
\EN
\begin{Resp}\Rsign I alone compassed the circuit of heaven, and walked on the waves of the sea. In every nation and in every people, I held the first place. \Star{} In the greatness of my strength have I trodden under my feet the necks of such as be haughty and proud.\end{Resp}
\begin{Resp}\Vsign I dwell in the highest places, and my throne is in a cloudy pillar.\end{Resp}
\begin{Resp}\Rsign In the greatness of my strength have I trodden under my feet the necks of such as be haughty and proud.\end{Resp}
\switchcolumn*

\LA
\LessonHead{Lectio III}
\switchcolumn
\EN
\LessonHead{Lesson III}
\switchcolumn*

\LA
\Cite{Eccl 5:9-13}
\switchcolumn
\EN
\Cite{Eccl 5:9-13}
\switchcolumn*

\LA
\noindent Avárus non implébitur pecúnia, et qui amat divítias fructum non cápiet ex eis: et hoc ergo vánitas. Ubi multæ sunt opes, multi et qui cómedunt eas. Et quid prodest possessóri, nisi quod cernit divítias óculis suis? Dulcis est somnus operánti, sive parum sive multum cómedat; satúritas autem dívitis non sinit eum dormíre. Est et ália infírmitas péssima, quam vidi sub sole: divítiæ conservátæ in malum dómini sui. Péreunt enim in afflictióne péssima: generávit fílium qui in summa egestáte erit.\par
\switchcolumn
\EN
\noindent A covetous man shall not be satisfied with money: and he that loveth riches shall reap no fruit from them: so this also is vanity. Where there are great riches, there are also many to eat them. And what doth it profit the owner, but that he seeth the riches with his eyes? Sleep is sweet to a labouring man, whether he eat little or much: but the fulness of the rich will not suffer him to sleep. There is also another grievous evil, which I have seen under the sun: riches kept to the hurt of the owner. For they are lost with very great affliction: he hath begotten a son, who shall be in extremity of want.\par
\switchcolumn*

\LA
\begin{Resp}\Rsign Emítte, Dómine, sapiéntiam de sede magnitúdinis tuæ, ut mecum sit et mecum labóret: \Star{} Ut sciam, quid accéptum sit coram te omni témpore.\end{Resp}
\begin{Resp}\Vsign Da mihi, Dómine, sédium tuárum assistrícem sapiéntiam.\end{Resp}
\begin{Resp}\Rsign Ut sciam quid accéptum sit coram te omni témpore.\end{Resp}
\begin{Resp}\Vsign Glória Patri, et Fílio, \Star{} et Spirítui Sancto.\end{Resp}
\begin{Resp}\Rsign Ut sciam quid accéptum sit coram te omni témpore.\end{Resp}
\switchcolumn
\EN
\begin{Resp}\Rsign O send out wisdom from the throne of thy glory, O Lord, to be with me, and to labour with me, \Star{} That I may know at all times what is pleasing unto thee.\end{Resp}
\begin{Resp}\Vsign Give me wisdom, O Lord, that sitteth by thy throne.\end{Resp}
\begin{Resp}\Rsign That I may know at all times what is pleasing unto thee.\end{Resp}
\begin{Resp}\Vsign Glory be to the Father, and to the Son, \Star{} and to the Holy Ghost.\end{Resp}
\begin{Resp}\Rsign That I may know at all times what is pleasing unto thee.\end{Resp}
\switchcolumn*

\end{paracol}

\Office{Feria VI infra Hebdomadam II Augusti}{Friday of the second week of August}{}
\addcontentsline{toc}{subsection}{Feria VI infra Hebdomadam II Augusti \textit{— Friday of the second week of August}}
\begin{paracol}{2}
\LA
\LessonHead{Lectio I}
\switchcolumn
\EN
\LessonHead{Lesson I}
\switchcolumn*

\LA
\LessonTitle{De libro Ecclesiástæ}
\switchcolumn
\EN
\LessonTitle{Lesson from the book of Ecclesiastes}
\switchcolumn*

\LA
\Cite{Eccl 6:1-2}
\switchcolumn
\EN
\Cite{Eccl 6:1-2}
\switchcolumn*

\LA
\noindent Est et áliud malum quod vidi sub sole, et quidem frequens apud hómines: vir, cui dedit Deus divítias et substántiam et honórem, et nihil deest ánimæ suæ, ex ómnibus quæ desíderat; nec tríbuit ei potestátem Deus ut cómedat ex eo, sed homo extráneus vorábit illud: hoc vánitas et miséria magna est.\par
\switchcolumn
\EN
\noindent There is also another evil, which I have seen under the sun, and that frequent among men: A man to whom God hath given riches, and substance, and honour, and his soul wanteth nothing of all that he desireth: yet God doth not give him power to eat thereof, but a stranger shall eat it up. This is vanity and a great misery.\par
\switchcolumn*

\LA
\begin{Resp}\Rsign Da mihi, Dómine, sédium tuárum assistrícem sapiéntiam, et noli me reprobáre a púeris tuis: \Star{} Quóniam servus tuus sum ego, et fílius ancíllæ tuæ.\end{Resp}
\begin{Resp}\Vsign Mitte illam de sede magnitúdinis tuæ, ut mecum sit et mecum labóret.\end{Resp}
\begin{Resp}\Rsign Quóniam servus tuus sum ego, et fílius ancíllæ tuæ.\end{Resp}
\switchcolumn
\EN
\begin{Resp}\Rsign Give me wisdom, O Lord, that sitteth by thy throne, and reject me not from among thy children. \Star{} For I am thy servant and son of thine handmaid.\end{Resp}
\begin{Resp}\Vsign O send her out from the throne of thy glory, to be with me and to labour with me.\end{Resp}
\begin{Resp}\Rsign For I am thy servant and son of thine handmaid.\end{Resp}
\switchcolumn*

\LA
\LessonHead{Lectio II}
\switchcolumn
\EN
\LessonHead{Lesson II}
\switchcolumn*

\LA
\Cite{Eccl 6:3-6}
\switchcolumn
\EN
\Cite{Eccl 6:3-6}
\switchcolumn*

\LA
\noindent Si genúerit quíspiam centum líberos et víxerit multos annos et plures dies ætátis habúerit, et ánima illíus non utátur bonis substántiæ suæ sepulturáque cáreat; de hoc ego pronúntio, quod mélior illo sit abortívus. Frustra enim venit et pergit ad ténebras, et oblivióne delébitur nomen ejus. Non vidit solem neque cognóvit distántiam boni et mali, étiam si duóbus míllibus annis víxerit, et non fúerit perfrúitus bonis.\par
\switchcolumn
\EN
\noindent If a man beget a hundred children, and live many years, and attain to a great age, and his soul make no use of the goods of his substance, and he be without burial: of this man I pronounce, that the untimely born is better than he. For he came in vain, and goeth to darkness, and his name shall be wholly forgotten. He hath not seen the sun, nor known the distance of good and evil: Although he lived two thousand years, and hath not enjoyed good things.\par
\switchcolumn*

\LA
\begin{Resp}\Rsign Inítium sapiéntiæ timor Dómini: \Star{} Intelléctus bonus ómnibus faciéntibus eum: laudátio ejus manet in sǽculum sǽculi.\end{Resp}
\begin{Resp}\Vsign Diléctio illíus custódia legum est: quia omnis sapiéntia timor Dómini.\end{Resp}
\begin{Resp}\Rsign Intelléctus bonus ómnibus faciéntibus eum: laudátio ejus manet in sǽculum sǽculi.\end{Resp}
\begin{Resp}\Vsign Glória Patri, et Fílio, \Star{} et Spirítui Sancto.\end{Resp}
\begin{Resp}\Rsign Intelléctus bonus ómnibus faciéntibus eum: laudátio ejus manet in sǽculum sǽculi.\end{Resp}
\switchcolumn
\EN
\begin{Resp}\Rsign The fear of the Lord is the beginning of wisdom. \Star{} A good understanding have all they that do His commandments. His praise endureth for ever.\end{Resp}
\begin{Resp}\Vsign Love is the keeping of her laws, for all wisdom is the fear of the Lord.\end{Resp}
\begin{Resp}\Rsign A good understanding have all they that do His commandments. His praise endureth for ever.\end{Resp}
\begin{Resp}\Vsign Glory be to the Father, and to the Son, \Star{} and to the Holy Ghost.\end{Resp}
\begin{Resp}\Rsign A good understanding have all they that do His commandments. His praise endureth for ever.\end{Resp}
\switchcolumn*

\LA
\LessonHead{Lectio III}
\switchcolumn
\EN
\LessonHead{Lesson III}
\switchcolumn*

\LA
\Cite{Eccl 6:6-9}
\switchcolumn
\EN
\Cite{Eccl 6:6-9}
\switchcolumn*

\LA
\noindent Nonne ad unum locum próperant ómnia? Omnis labor hóminis in ore ejus; sed ánima ejus non implébitur. Quid habet ámplius sápiens a stulto? et quid pauper, nisi ut pergat illuc ubi est vita? Mélius est vidére quod cúpias, quam desideráre quod néscias; sed et hoc vánitas est, et præsúmptio spíritus.\par
\switchcolumn
\EN
\noindent do not all make haste to one place? All the labour of man is for his mouth, but his soul shall not be filled. What hath the wise man more than the fool? and what the poor man, but to go thither, where there is life? Better it is to see what thou mayst desire, than to desire that which thou canst not know. But this also is vanity, and presumption of spirit.\par
\switchcolumn*

\LA
\begin{Resp}\Rsign Verbum iníquum et dolósum longe fac a me, Dómine: \Star{} Divítias et paupertátem ne déderis mihi, sed tantum víctui meo tríbue necessária.\end{Resp}
\begin{Resp}\Vsign Duo rogávi te, ne déneges mihi ántequam móriar.\end{Resp}
\begin{Resp}\Rsign Divítias et paupertátem ne déderis mihi, sed tantum víctui meo tríbue necessária.\end{Resp}
\begin{Resp}\Vsign Glória Patri, et Fílio, \Star{} et Spirítui Sancto.\end{Resp}
\begin{Resp}\Rsign Divítias et paupertátem ne déderis mihi, sed tantum víctui meo tríbue necessária.\end{Resp}
\switchcolumn
\EN
\begin{Resp}\Rsign Lord, remove far from me vanity and lies. \Star{} Give me neither poverty nor riches, but feed me with food convenient for me.\end{Resp}
\begin{Resp}\Vsign Two things have I required of thee; deny me them not before I die.\end{Resp}
\begin{Resp}\Rsign Give me neither poverty nor riches, but feed me with food convenient for me.\end{Resp}
\begin{Resp}\Vsign Glory be to the Father, and to the Son, \Star{} and to the Holy Ghost.\end{Resp}
\begin{Resp}\Rsign Give me neither poverty nor riches, but feed me with food convenient for me.\end{Resp}
\switchcolumn*

\end{paracol}

\Office{Sabbato infra Hebdomadam II Augusti}{Saturday of the second week of August}{}
\addcontentsline{toc}{subsection}{Sabbato infra Hebdomadam II Augusti \textit{— Saturday of the second week of August}}
\begin{paracol}{2}
\LA
\LessonHead{Lectio I}
\switchcolumn
\EN
\LessonHead{Lesson I}
\switchcolumn*

\LA
\LessonTitle{De libro Ecclesiástæ}
\switchcolumn
\EN
\LessonTitle{Lesson from the book of Ecclesiastes}
\switchcolumn*

\LA
\Cite{Eccl 7:1-3}
\switchcolumn
\EN
\Cite{Eccl 7:1-3}
\switchcolumn*

\LA
\noindent Quid necésse est hómini majóra se quǽrere, cum ignóret quid condúcat sibi in vita sua, número diérum peregrinatiónis suæ et témpore quod velut umbra prǽterit? Aut quis ei póterit indicáre quid post eum futúrum sub sole sit? Mélius est nomen bonum quam unguénta pretiósa, et dies mortis die nativitátis. Mélius est ire ad domum luctus quam ad domum convívii; in illa enim finis cunctórum admonétur hóminum, et vivens cógitat quid futúrum sit.\par
\switchcolumn
\EN
\noindent What needeth a man to seek things that are above him, whereas he knoweth not what is profitable for him in his life, in all the days of his pilgrimage, and the time that passeth like a shadow? Or who can tell him what shall be after him under the sun? A good name is better than precious ointments: and the day of death than the day of one's birth. It is better to go to the house of mourning, than to the house of feasting: for in that we are put in mind of the end of all, and the living thinketh what is to come.\par
\switchcolumn*

\LA
\begin{Resp}\Rsign Dómine, Pater et Deus vitæ meæ, ne derelínquas me in cogitátu malígno: extolléntiam oculórum meórum ne déderis mihi, et desidérium malígnum avérte a me, Dómine; aufer a me concupiscéntiam, \Star{} Et ánimo irreverénti et infruníto ne tradas me, Dómine.\end{Resp}
\begin{Resp}\Vsign Ne derelínquas me, Dómine, ne accréscant ignorántiæ meæ, nec multiplicéntur delícta mea.\end{Resp}
\begin{Resp}\Rsign Et ánimo irreverénti et infruníto ne tradas me, Dómine.\end{Resp}
\switchcolumn
\EN
\begin{Resp}\Rsign O Lord, Father and God of my life, leave me not to evil counsels; give me not a proud look, but turn away from me an haughty mind, O Lord turn away from me concupiscence, \Star{} And give me not over unto an impudent and froward mind, O Lord!\end{Resp}
\begin{Resp}\Vsign Leave me not, O Lord, lest mine ignorance increase, and my sins abound.\end{Resp}
\begin{Resp}\Rsign And give me not over unto an impudent and froward mind, O Lord.\end{Resp}
\switchcolumn*

\LA
\LessonHead{Lectio II}
\switchcolumn
\EN
\LessonHead{Lesson II}
\switchcolumn*

\LA
\Cite{Eccl 7:4-9}
\switchcolumn
\EN
\Cite{Eccl 7:4-9}
\switchcolumn*

\LA
\noindent Mélior est ira risu, quia per tristítiam vultus corrígitur ánimus delinquéntis. Cor sapiéntium ubi tristítia est, et cor stultórum ubi lætítia. Mélius est a sapiénte córripi, quam stultórum adulatióne décipi; quia sicut sónitus spinárum ardéntium sub olla, sic risus stulti; sed et hoc vánitas. Calúmnia contúrbat sapiéntem et perdet robur cordis illíus. Mélior est finis oratiónis quam princípium, mélior est pátiens arrogánte.\par
\switchcolumn
\EN
\noindent Anger is better than laughter: because by the sadness of the countenance the mind of the offender is corrected. The heart of the wise is where there is mourning, and the heart of fools where there is mirth. It is better to be rebuked by a wise man, than to be deceived by the flattery of fools. For as the crackling of thorns burning under a pot, so is the laughter of a fool: now this also is vanity. Oppression troubleth the wise, and shall destroy the strength of his heart. Better is the end of a speech than the beginning. Better is the patient man than the presumptuous.\par
\switchcolumn*

\LA
\begin{Resp}\Rsign Magna enim sunt judícia tua, Dómine, et inenarrabília verba tua: \Star{} Magnificásti pópulum tuum et honorásti.\end{Resp}
\begin{Resp}\Vsign Transtulísti illos per Mare Rubrum et transvexísti eos per aquam nímiam.\end{Resp}
\begin{Resp}\Rsign Magnificásti pópulum tuum et honorásti.\end{Resp}
\switchcolumn
\EN
\begin{Resp}\Rsign Great are thy judgments, O Lord, and thy words cannot be expressed. \Star{} Thou didst make thy people mighty and honourable.\end{Resp}
\begin{Resp}\Vsign Thou broughtest them through the Red Sea, and leddest them through much water.\end{Resp}
\begin{Resp}\Rsign Thou didst make thy people mighty and honourable.\end{Resp}
\switchcolumn*

\LA
\LessonHead{Lectio III}
\switchcolumn
\EN
\LessonHead{Lesson III}
\switchcolumn*

\LA
\Cite{Eccl 7:11-14}
\switchcolumn
\EN
\Cite{Eccl 7:11-14}
\switchcolumn*

\LA
\noindent Ne dicas: Quid putas causæ est quod prióra témpora melióra fuére quam nunc sunt? stulta enim est hujuscémodi interrogátio. Utílior est sapiéntia cum divítiis et magis prodest vidéntibus solem. Sicut enim prótegit sapiéntia, sic prótegit pecúnia; hoc autem plus habet erudítio et sapiéntia, quod vitam tríbuunt possessóri suo. Consídera ópera Dei, quod nemo possit corrígere quem ille despéxerit.\par
\switchcolumn
\EN
\noindent Say not: What thinkest thou is the cause that former times were better than they are now? for this manner of question is foolish. Wisdom with riches is more profitable, and bringeth more advantage to them that see the sun. For as wisdom is a defence, so money is a defence: but learning and wisdom excel in this, that they give life to him that possesseth them. Consider the works of God, that no man can correct whom he hath despised.\par
\switchcolumn*

\LA
\begin{Resp}\Rsign Quæ sunt in corde hóminum, óculi tui vident, Dómine, et in libro tuo ómnia scribéntur: \Star{} Homo videt in fácie, Deus autem in corde.\end{Resp}
\begin{Resp}\Vsign Omnia enim corda scrutátur, et univérsas méntium cogitatiónes intéllegit.\end{Resp}
\begin{Resp}\Rsign Homo videt in fácie, Deus autem in corde.\end{Resp}
\begin{Resp}\Vsign Glória Patri, et Fílio, \Star{} et Spirítui Sancto.\end{Resp}
\begin{Resp}\Rsign Homo videt in fácie, Deus autem in corde.\end{Resp}
\switchcolumn
\EN
\begin{Resp}\Rsign Lord, thine eyes behold all that is in the heart of man, and in thy book are they all written. \Star{} Man looketh on the outward appearance, but God looketh on the heart.\end{Resp}
\begin{Resp}\Vsign For He searcheth all hearts, and understandeth all the imaginations of the thoughts.\end{Resp}
\begin{Resp}\Rsign Man looketh on the outward appearance, but God looketh on the heart.\end{Resp}
\begin{Resp}\Vsign Glory be to the Father, and to the Son, \Star{} and to the Holy Ghost.\end{Resp}
\begin{Resp}\Rsign Man looketh on the outward appearance, but God looketh on the heart.\end{Resp}
\switchcolumn*

\end{paracol}

\Office{Dominica III Augusti}{Third Sunday of August}{}
\addcontentsline{toc}{subsection}{Dominica III Augusti \textit{— Third Sunday of August}}
\begin{paracol}{2}
\LA
\LessonHead{Lectio I}
\switchcolumn
\EN
\LessonHead{Lesson I}
\switchcolumn*

\LA
\LessonTitle{Incipit liber Sapiéntiæ}
\switchcolumn
\EN
\LessonTitle{Lesson from the book of Wisdom}
\switchcolumn*

\LA
\Cite{Sap 1:1-4}
\switchcolumn
\EN
\Cite{Wis 1:1-4}
\switchcolumn*

\LA
\noindent Dilígite justítiam, qui judicátis terram, sentíte de Dómino in bonitáte et in simplicitáte cordis quǽrite illum; quóniam invenítur ab his qui non tentant illum, appáret autem eis qui fidem habent in illum. Pervérsæ enim cogitatiónes séparant a Deo, probáta autem virtus córripit insipiéntes; quóniam in malévolam ánimam non introíbit sapiéntia, nec habitábit in córpore súbdito peccátis.\par
\switchcolumn
\EN
\noindent Love justice, you that are the judges of the earth. Think of the Lord in goodness, and seek him in simplicity of heart. For he is found by them that tempt him not: and he showeth himself to them that have faith in him. For perverse thoughts separate from God: and his power, when it is tried, reproveth the unwise: For wisdom will not enter into a malicious soul, nor dwell in a body subject to sins.\par
\switchcolumn*

\LA
\begin{Resp}\Rsign In princípio Deus ántequam terram fáceret, priúsquam abýssos constitúeret, priúsquam prodúceret fontes aquárum, \Star{} Antequam montes collocaréntur, ante omnes colles generávit me Dóminus.\end{Resp}
\begin{Resp}\Vsign Quando præparábat cælos, áderam, cum eo cuncta compónens.\end{Resp}
\begin{Resp}\Rsign Antequam montes collocaréntur, ante omnes colles generávit me Dóminus.\end{Resp}
\switchcolumn
\EN
\begin{Resp}\Rsign God possessed me in the beginning, before He made the earth, before He created the depths, before He caused the fountains of water to spring. \Star{} Before the mountains were settled, before there were any hills, did the Lord beget me.\end{Resp}
\begin{Resp}\Vsign When He prepared the heavens, I was there with Him, ordering all things.\end{Resp}
\begin{Resp}\Rsign Before the mountains were settled, before there were any hills, did the Lord beget me.\end{Resp}
\switchcolumn*

\LA
\LessonHead{Lectio II}
\switchcolumn
\EN
\LessonHead{Lesson II}
\switchcolumn*

\LA
\Cite{Sap 1:5-8}
\switchcolumn
\EN
\Cite{Wis 1:5-8}
\switchcolumn*

\LA
\noindent Spíritus enim sanctus disciplínæ effúgiet fictum et áuferet se a cogitatiónibus quæ sunt sine intelléctu et corripiétur a superveniénte iniquitáte. Benígnus est enim spíritus sapiéntiæ et non liberábit malédicum a lábiis suis; quóniam renum illíus testis est Deus et cordis illíus scrutátor est verus, et linguæ ejus audítor: quóniam spíritus Dómini replévit orbem terrárum, et hoc quod cóntinet ómnia, sciéntiam habet vocis. Propter hoc qui lóquitur iníqua non potest latére, nec prætériet illum corrípiens judícium.\par
\switchcolumn
\EN
\noindent For the Holy Spirit of discipline will flee from the deceitful, and will withdraw himself from thoughts that are without understanding, and he shall not abide when iniquity cometh in. For the spirit of wisdom is benevolent, and will not acquit the evil speaker from his lips: for God is witness of his reins, and he is a true searcher of his heart, and a hearer of his tongue. For the spirit of the Lord hath filled the whole world: and that, which containeth all things, hath knowledge of the voice. Therefore he that speaketh unjust things cannot be hid, neither shall the chastising judgment pass him by.\par
\switchcolumn*

\LA
\begin{Resp}\Rsign Gyrum cæli circuívi sola, et in flúctibus maris ambulávi, in omni gente et in omni pópulo primátum ténui: \Star{} Superbórum et sublímium colla própria virtúte calcávi.\end{Resp}
\begin{Resp}\Vsign Ego in altíssimis hábito, et thronus meus in colúmna nubis.\end{Resp}
\begin{Resp}\Rsign Superbórum et sublímium colla própria virtúte calcávi.\end{Resp}
\switchcolumn
\EN
\begin{Resp}\Rsign I alone compassed the circuit of heaven, and walked on the waves of the sea. In every nation and in every people, I held the first place. \Star{} In the greatness of my strength have I trodden under my feet the necks of such as be haughty and proud.\end{Resp}
\begin{Resp}\Vsign I dwell in the highest places, and my throne is in a cloudy pillar.\end{Resp}
\begin{Resp}\Rsign In the greatness of my strength have I trodden under my feet the necks of such as be haughty and proud.\end{Resp}
\switchcolumn*

\LA
\LessonHead{Lectio III}
\switchcolumn
\EN
\LessonHead{Lesson III}
\switchcolumn*

\LA
\Cite{Sap 1:9-11}
\switchcolumn
\EN
\Cite{Wis 1:9-11}
\switchcolumn*

\LA
\noindent In cogitatiónibus enim ímpii interrogátio erit, sermónum autem illíus audítio ad Deum véniet ad correptiónem iniquitátum illíus; quóniam auris zeli audit ómnia, et tumúltus murmuratiónum non abscondétur. Custodíte ergo vos a murmuratióne, quæ nihil prodest, et a detractióne párcite linguæ, quóniam sermo obscúrus in vácuum non ibit, os autem quod mentítur occídit ánimam.\par
\switchcolumn
\EN
\noindent For inquisition shall be made into the thoughts of the ungodly: and the hearing of his words shall come to God, to the chastising of his iniquities. For the ear of jealousy heareth all things, and the tumult of murmuring shall not be hid. Keep yourselves therefore from murmuring, which profiteth nothing, and refrain your tongue from detraction, for an obscure speech shall not go for nought: and the mouth that belieth, killeth the soul.\par
\switchcolumn*

\LA
\begin{Resp}\Rsign Emítte, Dómine, sapiéntiam de sede magnitúdinis tuæ, ut mecum sit et mecum labóret: \Star{} Ut sciam, quid accéptum sit coram te omni témpore.\end{Resp}
\begin{Resp}\Vsign Da mihi, Dómine, sédium tuárum assistrícem sapiéntiam.\end{Resp}
\begin{Resp}\Rsign Ut sciam quid accéptum sit coram te omni témpore.\end{Resp}
\begin{Resp}\Vsign Glória Patri, et Fílio, \Star{} et Spirítui Sancto.\end{Resp}
\begin{Resp}\Rsign Ut sciam quid accéptum sit coram te omni témpore.\end{Resp}
\switchcolumn
\EN
\begin{Resp}\Rsign O send out wisdom from the throne of thy glory, O Lord, to be with me, and to labour with me, \Star{} That I may know at all times what is pleasing unto thee.\end{Resp}
\begin{Resp}\Vsign Give me wisdom, O Lord, that sitteth by thy throne.\end{Resp}
\begin{Resp}\Rsign That I may know at all times what is pleasing unto thee.\end{Resp}
\begin{Resp}\Vsign Glory be to the Father, and to the Son, \Star{} and to the Holy Ghost.\end{Resp}
\begin{Resp}\Rsign That I may know at all times what is pleasing unto thee.\end{Resp}
\switchcolumn*

\LA
\LessonHead{Lectio IV}
\switchcolumn
\EN
\LessonHead{Lesson IV}
\switchcolumn*

\LA
\LessonTitle{Ex libro Officiórum sancti Ambrósii Epíscopi}
\switchcolumn
\EN
\LessonTitle{From the Book of St. Ambrose, Bishop of Milan, “On Offices.”}
\switchcolumn*

\LA
\Source{Lib. 1. Cap. 28 et 29}
\switchcolumn
\EN
\Source{Bk. i. c. 28}
\switchcolumn*

\LA
\noindent Magnus justítiæ splendor, quæ áliis pótius nata quam sibi, communitátem et societátem nostram ádjuvat, excelsitátem tenet, ut suo judício ómnia subjécta hábeat, opem áliis ferat, pecúniam cónferat, offícia non ábnuat, perícula suscípiat aliéna. Quis non cúperet hanc virtútis arcem tenére, nisi prima avarítia infirmáret atque inflécteret tantæ virtútis vigórem? Etenim dum augére opes, aggregáre pecúnias, occupáre terras possessiónibus cúpimus, præstáre divítiis; justítiæ formam exúimus, beneficéntiam commúnem amíttimus.\par
\switchcolumn
\EN
\noindent Great is the glory of justice. She liveth for others rather than for herself. By her our commonwealth and fellowship are holpen. She holdeth such a pre-eminence that all things are subject unto her judgment. She helpeth others. She giveth wealth. She refuseth not to labour. She taketh upon her the dangers of others. Who would not desire to hold this castle of power and courage, if the covetousness of our first parents had not weakened and distorted the strength of our nerve But so it is, that, while we are fain to increase wealth, to put by money, to add lands to our possessions, or to make show of our abundance, we put off the image of justice, and lose charity toward our brethren.\par
\switchcolumn*

\LA
\begin{Resp}\Rsign Da mihi, Dómine, sédium tuárum assistrícem sapiéntiam, et noli me reprobáre a púeris tuis: \Star{} Quóniam servus tuus sum ego, et fílius ancíllæ tuæ.\end{Resp}
\begin{Resp}\Vsign Mitte illam de sede magnitúdinis tuæ, ut mecum sit et mecum labóret.\end{Resp}
\begin{Resp}\Rsign Quóniam servus tuus sum ego, et fílius ancíllæ tuæ.\end{Resp}
\switchcolumn
\EN
\begin{Resp}\Rsign Give me wisdom, O Lord, that sitteth by thy throne, and reject me not from among thy children. \Star{} For I am thy servant and son of thine handmaid.\end{Resp}
\begin{Resp}\Vsign O send her out from the throne of thy glory, to be with me and to labour with me.\end{Resp}
\begin{Resp}\Rsign For I am thy servant and son of thine handmaid.\end{Resp}
\switchcolumn*

\LA
\LessonHead{Lectio V}
\switchcolumn
\EN
\LessonHead{Lesson V}
\switchcolumn*

\LA
\noindent Quanta autem justítia sit, ex hoc intélligi potest, quod nec locis, nec persónis, nec tempóribus excípitur, quæ étiam hóstibus reservátur: ut si constitútus sit cum hoste aut locus aut dies prǽlio, advérsus justítiam putétur aut loco præveníre aut témpore. Interest enim utrum áliquis pugna áliqua et conflíctu gravi capiátur, an superióre grátia, vel áliquo evéntu. Si ergo in bello justítia valet, quanto magis in pace servánda est?\par
\switchcolumn
\EN
\noindent How far-spreading is the field of justice appeareth by this, that there is excepted therefrom no place, person, or time, nay, she hath to do even as regards enemies, for if one be agreed with his enemy of a certain place, or day for battle, it should be deemed unjust to fall on him beforehand, at some other place, or time. For it is a very different thing, whether one get the better of another in a hard fight, or by skill, or by accident. If therefore in war justice hath place, how much more is she to be observed in time of peace\par
\switchcolumn*

\LA
\begin{Resp}\Rsign Inítium sapiéntiæ timor Dómini: \Star{} Intelléctus bonus ómnibus faciéntibus eum: laudátio ejus manet in sǽculum sǽculi.\end{Resp}
\begin{Resp}\Vsign Diléctio illíus custódia legum est: quia omnis sapiéntia timor Dómini.\end{Resp}
\begin{Resp}\Rsign Intelléctus bonus ómnibus faciéntibus eum: laudátio ejus manet in sǽculum sǽculi.\end{Resp}
\switchcolumn
\EN
\begin{Resp}\Rsign The fear of the Lord is the beginning of wisdom. \Star{} A good understanding have all they that do His commandments. His praise endureth for ever.\end{Resp}
\begin{Resp}\Vsign Love is the keeping of her laws, for all wisdom is the fear of the Lord.\end{Resp}
\begin{Resp}\Rsign A good understanding have all they that do His commandments. His praise endureth for ever.\end{Resp}
\switchcolumn*

\LA
\LessonHead{Lectio VI}
\switchcolumn
\EN
\LessonHead{Lesson VI}
\switchcolumn*

\LA
\noindent Fundaméntum ergo est justítiæ fides. Justórum enim corda meditántur fidem: et qui se justus accúsat, justítiam supra fidem cóllocat. Nam tunc justítia ejus appáret, si vera fateátur. Dénique et Dóminus per Isaíam: Ecce, inquit, mitto lápidem in fundaméntum Sion: id est, Christum in fundaméntum Ecclésiæ. Fides enim ómnium Christus: Ecclésia autem quædam forma justítiæ est, commúne jus ómnium: in commúne orat, in commúne operátur, in commúne tentátur. Dénique qui seípsum sibi ábnegat, ipse justus, ipse dignus Christo est. Ideo et Paulus fundaméntum pósuit Christum, ut supra eum ópera justítiæ locarémus, quia fides fundaméntum est.\par
\switchcolumn
\EN
\noindent Honour is the foundation of justice. The thoughts in the hearts of just men are honourable thoughts and when the just man accuseth himself, it is honour that bringeth him to that just deed. Then is his justice made manifest by his honourable avowal. The Lord saith by Isaiah “Behold, I lay in Zion a foundation-stone” (xxviii. 16), that is to say, He giveth Christ unto the Church to be her foundation. Christ is the true honour for all men, and the Church is as it were a figure of justice, being a commonwealth wherein all have rights, and which worketh as one, and suffereth as one. Whosoever denieth himself, the same is just, and worthy of Christ. Therefore also Paul saith “Other foundation can no man lay than that is laid, which is Jesus Christ” (I Cor. iii. 11), and upon that foundation is it, that every building of justice must be raised. For the spirit of Christ is the true spirit of honour which is the foundation whereon justice resteth.\par
\switchcolumn*

\LA
\begin{Resp}\Rsign Verbum iníquum et dolósum longe fac a me, Dómine: \Star{} Divítias et paupertátem ne déderis mihi, sed tantum víctui meo tríbue necessária.\end{Resp}
\begin{Resp}\Vsign Duo rogávi te, ne déneges mihi ántequam móriar.\end{Resp}
\begin{Resp}\Rsign Divítias et paupertátem ne déderis mihi, sed tantum víctui meo tríbue necessária.\end{Resp}
\begin{Resp}\Vsign Glória Patri, et Fílio, \Star{} et Spirítui Sancto.\end{Resp}
\begin{Resp}\Rsign Divítias et paupertátem ne déderis mihi, sed tantum víctui meo tríbue necessária.\end{Resp}
\switchcolumn
\EN
\begin{Resp}\Rsign Lord, remove far from me vanity and lies. \Star{} Give me neither poverty nor riches, but feed me with food convenient for me.\end{Resp}
\begin{Resp}\Vsign Two things have I required of thee; deny me them not before I die.\end{Resp}
\begin{Resp}\Rsign Give me neither poverty nor riches, but feed me with food convenient for me.\end{Resp}
\begin{Resp}\Vsign Glory be to the Father, and to the Son, \Star{} and to the Holy Ghost.\end{Resp}
\begin{Resp}\Rsign Give me neither poverty nor riches, but feed me with food convenient for me.\end{Resp}
\switchcolumn*

\end{paracol}

\Office{Feria II infra Hebdomadam III Augusti}{Monday of the third week of August}{}
\addcontentsline{toc}{subsection}{Feria II infra Hebdomadam III Augusti \textit{— Monday of the third week of August}}
\begin{paracol}{2}
\LA
\LessonHead{Lectio I}
\switchcolumn
\EN
\LessonHead{Lesson I}
\switchcolumn*

\LA
\LessonTitle{De libro Sapiéntiæ}
\switchcolumn
\EN
\LessonTitle{Lesson from the book of Wisdom}
\switchcolumn*

\LA
\Cite{Sap 3:1-6}
\switchcolumn
\EN
\Cite{Wis 3:1-6}
\switchcolumn*

\LA
\noindent Justórum autem ánimæ in manu Dei sunt, et non tanget illos torméntum mortis. Visi sunt óculis insipiéntium mori, et æstimáta est afflíctio éxitus illórum, et quod a nobis est iter extermínium; illi autem sunt in pace; et si coram homínibus torménta passi sunt, spes illórum immortalitáte plena est. In paucis vexáti in multis bene disponéntur, quóniam Deus tentávit eos et invénit illos dignos se. Tamquam aurum in fornáce probávit illos, et quasi holocáusti hóstiam accépit illos, et in témpore erit respéctus illórum.\par
\switchcolumn
\EN
\noindent But the souls of the just are in the hand of God, and the torment of death shall not touch them. In the sight of the unwise they seemed to die: and their departure was taken for misery: And their going away from us, for utter destruction: but they are in peace. And though in the sight of men they suffered torments, their hope is full of immortality. Afflicted in few things, in many they shall be well rewarded: because God hath tried them, and found them worthy of himself. As gold in the furnace he hath proved them, and as a victim of a holocaust he hath received them, and in time there shall be respect had to them.\par
\switchcolumn*

\LA
\begin{Resp}\Rsign Ne derelínquas me, Dómine, pater et dominátor vitæ meæ, ut non córruam in conspéctu adversariórum meórum: \Star{} Ne gáudeat de me inimícus meus.\end{Resp}
\begin{Resp}\Vsign Apprehénde arma et scutum et exsúrge in adjutórium mihi.\end{Resp}
\begin{Resp}\Rsign Ne gáudeat de me inimícus meus.\end{Resp}
\switchcolumn
\EN
\begin{Resp}\Rsign O Lord, Father and Governor of my life, leave me not, lest I fall before mine adversaries, \Star{} and mine enemy rejoice over me.\end{Resp}
\begin{Resp}\Vsign Take hold of shield and buckler, and stand up for mine help.\end{Resp}
\begin{Resp}\Rsign Lest mine enemy rejoice over me.\end{Resp}
\switchcolumn*

\LA
\LessonHead{Lectio II}
\switchcolumn
\EN
\LessonHead{Lesson II}
\switchcolumn*

\LA
\Cite{Sap 3:7-11}
\switchcolumn
\EN
\Cite{Wis 3:7-11}
\switchcolumn*

\LA
\noindent Fulgébunt justi et tamquam scintíllæ in arundinéto discúrrent; judicábunt natiónes et dominabúntur pópulis, et regnábit Dóminus illórum in perpétuum. Qui confídunt in illo intéllegent veritátem, et fidéles in dilectióne acquiéscent illi, quóniam donum et pax est eléctis ejus. Impii autem, secúndum quæ cogitavérunt, correptiónem habébunt, qui neglexérunt justum et a Dómino recessérunt. Sapiéntiam enim et disciplínam qui ábicit infélix est; et vácua est spes illórum, et labóres sine fructu, et inutília ópera eórum.\par
\switchcolumn
\EN
\noindent The just shall shine, and shall run to and fro like sparks among the reeds. They shall judge nations, and rule over people, and their Lord shall reign for ever. They that trust in him, shall understand the truth: and they that are faithful in love shall rest in him: for grace and peace is to his elect. But the wicked shall be punished according to their own devices: who have neglected the just, and have revolted from the Lord. For he that rejecteth wisdom, and discipline, is unhappy: and their hope is vain, and their labours without fruit, and their works unprofitable.\par
\switchcolumn*

\LA
\begin{Resp}\Rsign Magna enim sunt judícia tua, Dómine, et inenarrabília verba tua: \Star{} Magnificásti pópulum tuum et honorásti.\end{Resp}
\begin{Resp}\Vsign Transtulísti illos per Mare Rubrum et transvexísti eos per aquam nímiam.\end{Resp}
\begin{Resp}\Rsign Magnificásti pópulum tuum et honorásti.\end{Resp}
\begin{Resp}\Vsign Glória Patri, et Fílio, \Star{} et Spirítui Sancto.\end{Resp}
\begin{Resp}\Rsign Magnificásti pópulum tuum et honorásti.\end{Resp}
\switchcolumn
\EN
\begin{Resp}\Rsign Great are thy judgments, O Lord, and thy words cannot be expressed. \Star{} Thou didst make thy people mighty and honourable.\end{Resp}
\begin{Resp}\Vsign Thou broughtest them through the Red Sea, and leddest them through much water.\end{Resp}
\begin{Resp}\Rsign Thou didst make thy people mighty and honourable.\end{Resp}
\begin{Resp}\Vsign Glory be to the Father, and to the Son, \Star{} and to the Holy Ghost.\end{Resp}
\begin{Resp}\Rsign Thou didst make thy people mighty and honourable.\end{Resp}
\switchcolumn*

\LA
\LessonHead{Lectio III}
\switchcolumn
\EN
\LessonHead{Lesson III}
\switchcolumn*

\LA
\Cite{Sap 5:16-21}
\switchcolumn
\EN
\Cite{Wis 5:16-21}
\switchcolumn*

\LA
\noindent Justi autem in perpétuum vivent, et apud Dóminum est merces eórum et cogitátio illórum apud Altíssimum. Ideo accípient regnum decóris et diadéma speciéi de manu Dómini; quóniam déxtera sua teget eos et brácchio sancto suo deféndet illos. Accípiet armatúram zelus illíus et armábit creatúram ad ultiónem inimicórum. Induet pro thoráce justítiam et accípiet pro gálea judícium certum, sumet scutum inexpugnábile æquitátem. Acuet autem duram iram in lánceam, et pugnábit cum illo orbis terrárum contra insensátos.\par
\switchcolumn
\EN
\noindent But the just shall live for evermore: and their reward is with the Lord, and the care of them with the most High. Therefore shall they receive a kingdom of glory, and a crown of beauty at the hand of the Lord: for with his right hand he will cover them, and with his holy arm he will defend them. And his zeal will take armour, and he will arm the creature for the revenge of his enemies. He will put on justice as a breastplate, and will take true judgment instead of a helmet. He will take equity for an invincible shield: And he will sharpen his severe wrath for a spear, and the whole world shall fight with him against the unwise.\par
\switchcolumn*

\LA
\begin{Resp}\Rsign Quæ sunt in corde hóminum, óculi tui vident, Dómine, et in libro tuo ómnia scribéntur: \Star{} Homo videt in fácie, Deus autem in corde.\end{Resp}
\begin{Resp}\Vsign Omnia enim corda scrutátur, et univérsas méntium cogitatiónes intéllegit.\end{Resp}
\begin{Resp}\Rsign Homo videt in fácie, Deus autem in corde.\end{Resp}
\begin{Resp}\Vsign Glória Patri, et Fílio, \Star{} et Spirítui Sancto.\end{Resp}
\begin{Resp}\Rsign Homo videt in fácie, Deus autem in corde.\end{Resp}
\switchcolumn
\EN
\begin{Resp}\Rsign Lord, thine eyes behold all that is in the heart of man, and in thy book are they all written. \Star{} Man looketh on the outward appearance, but God looketh on the heart.\end{Resp}
\begin{Resp}\Vsign For He searcheth all hearts, and understandeth all the imaginations of the thoughts.\end{Resp}
\begin{Resp}\Rsign Man looketh on the outward appearance, but God looketh on the heart.\end{Resp}
\begin{Resp}\Vsign Glory be to the Father, and to the Son, \Star{} and to the Holy Ghost.\end{Resp}
\begin{Resp}\Rsign Man looketh on the outward appearance, but God looketh on the heart.\end{Resp}
\switchcolumn*

\end{paracol}

\Office{Feria III infra Hebdomadam III Augusti}{Tuesday of the third week of August}{}
\addcontentsline{toc}{subsection}{Feria III infra Hebdomadam III Augusti \textit{— Tuesday of the third week of August}}
\begin{paracol}{2}
\LA
\LessonHead{Lectio I}
\switchcolumn
\EN
\LessonHead{Lesson I}
\switchcolumn*

\LA
\LessonTitle{De libro Sapiéntiæ}
\switchcolumn
\EN
\LessonTitle{Lesson from the book of Wisdom}
\switchcolumn*

\LA
\Cite{Sap 6:1-5}
\switchcolumn
\EN
\Cite{Wis 6:1-5}
\switchcolumn*

\LA
\noindent Mélior est sapiéntia quam vires, et vir prudens quam fortis. Audíte ergo, reges, et intellégite; díscite, júdices fínium terræ; præbéte aures, vos qui continétis multitúdines et placétis vobis in turbis natiónum; quóniam data est a Dómino potéstas vobis, et virtus ab Altíssimo, qui interrogábit ópera vestra et cogitatiónes scrutábitur; quóniam, cum essétis minístri regni illíus, non recte judicástis nec custodístis legem justítiæ neque secúndum voluntátem Dei ambulástis.\par
\switchcolumn
\EN
\noindent Wisdom is better than strength, and a wise man is better than a strong man. Hear therefore, ye kings, and understand: learn, ye that are judges of the ends of the earth. Give ear, you that rule the people, and that please yourselves in multitudes of nations: For power is given you by the Lord, and strength by the most High, who will examine your works, and search out your thoughts: Because being ministers of his kingdom, you have not judged rightly, nor kept the law of justice, nor walked according to the will of God.\par
\switchcolumn*

\LA
\begin{Resp}\Rsign Præbe, fili, cor mihi, et óculi tui vias meas custódiant: \Star{} Ut addátur grátia cápiti tuo.\end{Resp}
\begin{Resp}\Vsign Atténde, fili mi, sapiéntiam meam et ad elóquium meum inclína aurem tuam.\end{Resp}
\begin{Resp}\Rsign Ut addátur grátia cápiti tuo.\end{Resp}
\switchcolumn
\EN
\begin{Resp}\Rsign My son, give me thine heart, and let thine eyes observe my ways. \Star{} For they shall be an ornament of grace unto thine head.\end{Resp}
\begin{Resp}\Vsign My son, attend unto my wisdom, and incline thine ear unto my sayings.\end{Resp}
\begin{Resp}\Rsign For they shall be an ornament of grace unto thine head.\end{Resp}
\switchcolumn*

\LA
\LessonHead{Lectio II}
\switchcolumn
\EN
\LessonHead{Lesson II}
\switchcolumn*

\LA
\Cite{Sap 6:6-9}
\switchcolumn
\EN
\Cite{Wis 6:6-9}
\switchcolumn*

\LA
\noindent Horrénde et cito apparébit vobis, quóniam judícium duríssimum his qui præsunt fiet. Exíguo enim concéditur misericórdia; poténtes autem poténter torménta patiéntur; non enim súbtrahet persónam cujúsquam Deus nec verébitur magnitúdinem cujúsquam; quóniam pusíllum et magnum ipse fecit, et æquáliter cura est illi de ómnibus; fortióribus autem fórtior instat cruciátio.\par
\switchcolumn
\EN
\noindent Horribly and speedily will he appear to you: for a most severe judgment shall be for them that bear rule. For to him that is little, mercy is granted: but the mighty shall be mightily tormented. For God will not except any man's person, neither will he stand in awe of any man's greatness: for he made the little and the great, and he hath equally care of all. But a greater punishment is ready for the more mighty.\par
\switchcolumn*

\LA
\begin{Resp}\Rsign Inítium sapiéntiæ timor Dómini: \Star{} Intelléctus bonus ómnibus faciéntibus eum: laudátio ejus manet in sǽculum sǽculi.\end{Resp}
\begin{Resp}\Vsign Diléctio illíus custódia legum est: quia omnis sapiéntia timor Dómini.\end{Resp}
\begin{Resp}\Rsign Intelléctus bonus ómnibus faciéntibus eum: laudátio ejus manet in sǽculum sǽculi.\end{Resp}
\switchcolumn
\EN
\begin{Resp}\Rsign The fear of the Lord is the beginning of wisdom. \Star{} A good understanding have all they that do His commandments. His praise endureth for ever.\end{Resp}
\begin{Resp}\Vsign Love is the keeping of her laws, for all wisdom is the fear of the Lord.\end{Resp}
\begin{Resp}\Rsign A good understanding have all they that do His commandments. His praise endureth for ever.\end{Resp}
\switchcolumn*

\LA
\LessonHead{Lectio III}
\switchcolumn
\EN
\LessonHead{Lesson III}
\switchcolumn*

\LA
\Cite{Sap 6:10-13}
\switchcolumn
\EN
\Cite{Wis 6:10-13}
\switchcolumn*

\LA
\noindent Ad vos ergo, reges, sunt hi sermónes mei, ut discátis sapiéntiam et non excidátis. Qui enim custodíerint justa juste, justificabúntur; et, qui didícerint ista, invénient quid respóndeant. Concupíscite ergo sermónes meos, dilígite illos et habébitis disciplínam. Clara est, et quæ numquam marcéscit, sapiéntia: et fácile vidétur ab his qui díligunt eam, et invenítur ab his qui quærunt illam.\par
\switchcolumn
\EN
\noindent To you, therefore, O kings, are these my words, that you may learn wisdom, and not fall from it. For they that have kept just things justly, shall be justified: and they that have learned these things, shall find what to answer. Covet ye therefore my words, and love them, and you shall have instruction. Wisdom is glorious, and never fadeth away, and is easily seen by them that love her, and is found by them that seek her.\par
\switchcolumn*

\LA
\begin{Resp}\Rsign Verbum iníquum et dolósum longe fac a me, Dómine: \Star{} Divítias et paupertátem ne déderis mihi, sed tantum víctui meo tríbue necessária.\end{Resp}
\begin{Resp}\Vsign Duo rogávi te, ne déneges mihi ántequam móriar.\end{Resp}
\begin{Resp}\Rsign Divítias et paupertátem ne déderis mihi, sed tantum víctui meo tríbue necessária.\end{Resp}
\begin{Resp}\Vsign Glória Patri, et Fílio, \Star{} et Spirítui Sancto.\end{Resp}
\begin{Resp}\Rsign Divítias et paupertátem ne déderis mihi, sed tantum víctui meo tríbue necessária.\end{Resp}
\switchcolumn
\EN
\begin{Resp}\Rsign Lord, remove far from me vanity and lies. \Star{} Give me neither poverty nor riches, but feed me with food convenient for me.\end{Resp}
\begin{Resp}\Vsign Two things have I required of thee; deny me them not before I die.\end{Resp}
\begin{Resp}\Rsign Give me neither poverty nor riches, but feed me with food convenient for me.\end{Resp}
\begin{Resp}\Vsign Glory be to the Father, and to the Son, \Star{} and to the Holy Ghost.\end{Resp}
\begin{Resp}\Rsign Give me neither poverty nor riches, but feed me with food convenient for me.\end{Resp}
\switchcolumn*

\end{paracol}

\Office{Feria IV infra Hebdomadam III Augusti}{Wednesday of the third week of August}{}
\addcontentsline{toc}{subsection}{Feria IV infra Hebdomadam III Augusti \textit{— Wednesday of the third week of August}}
\begin{paracol}{2}
\LA
\LessonHead{Lectio I}
\switchcolumn
\EN
\LessonHead{Lesson I}
\switchcolumn*

\LA
\LessonTitle{De libro Sapiéntiæ}
\switchcolumn
\EN
\LessonTitle{Lesson from the book of Wisdom}
\switchcolumn*

\LA
\Cite{Sap 7:1-6}
\switchcolumn
\EN
\Cite{Wis 7:1-6}
\switchcolumn*

\LA
\noindent Sum quidem et ego mortális homo, símilis ómnibus, et ex génere terréni illíus, qui prior factus est, et in ventre matris figurátus sum caro, decem ménsium témpore coagulátus sum in sánguine: ex sémine hóminis et delectaménto somni conveniénte; et ego natus accépi commúnem áërem et in simíliter factam décidi terram et primam vocem símilem ómnibus emísi plorans; in involuméntis nutrítus sum, et curis magnis. Nemo enim ex régibus áliud hábuit nativitátis inítium: unus ergo intróitus est ómnibus ad vitam, et símilis éxitus.\par
\switchcolumn
\EN
\noindent I myself also am a mortal man, like all others, and of the race of him, that was first made of the earth, and in the womb of my mother I was fashioned to be flesh. In the time of ten months I was compacted in blood, of the seed of man, and the pleasure of sleep concurring. And being born I drew in the common air, and fell upon the earth, that is made alike, and the first voice which I uttered was crying, as all others do. I was nursed in swaddling clothes, and with great cares. For none of the kings had any other beginning of birth. For all men have one entrance into life, and the like going out.\par
\switchcolumn*

\LA
\begin{Resp}\Rsign Dómine, Pater et Deus vitæ meæ, ne derelínquas me in cogitátu malígno: extolléntiam oculórum meórum ne déderis mihi, et desidérium malígnum avérte a me, Dómine; aufer a me concupiscéntiam, \Star{} Et ánimo irreverénti et infruníto ne tradas me, Dómine.\end{Resp}
\begin{Resp}\Vsign Ne derelínquas me, Dómine, ne accréscant ignorántiæ meæ, nec multiplicéntur delícta mea.\end{Resp}
\begin{Resp}\Rsign Et ánimo irreverénti et infruníto ne tradas me, Dómine.\end{Resp}
\switchcolumn
\EN
\begin{Resp}\Rsign O Lord, Father and God of my life, leave me not to evil counsels; give me not a proud look, but turn away from me an haughty mind, O Lord turn away from me concupiscence, \Star{} And give me not over unto an impudent and froward mind, O Lord!\end{Resp}
\begin{Resp}\Vsign Leave me not, O Lord, lest mine ignorance increase, and my sins abound.\end{Resp}
\begin{Resp}\Rsign And give me not over unto an impudent and froward mind, O Lord.\end{Resp}
\switchcolumn*

\LA
\LessonHead{Lectio II}
\switchcolumn
\EN
\LessonHead{Lesson II}
\switchcolumn*

\LA
\Cite{Sap 7:7-10}
\switchcolumn
\EN
\Cite{Wis 7:7-10}
\switchcolumn*

\LA
\noindent Propter hoc optávi, et datus est mihi sensus; et invocávi, et venit in me spíritus sapiéntiæ; et præpósui illam regnis et sédibus, et divítias nihil esse duxi in comparatióne illíus. Nec comparávi illi lápidem pretiósum, quóniam omne aurum in comparatióne illíus aréna est exígua, et tamquam lutum æstimábitur argéntum in conspéctu illíus. Super salútem et spéciem diléxi illam et propósui pro luce habére illam, quóniam inexstinguíbile est lumen illíus.\par
\switchcolumn
\EN
\noindent Wherefore I wished, and understanding was given me: and I called upon God, and the spirit of wisdom came upon me: And I preferred her before kingdoms and thrones, and esteemed riches nothing in comparison of her. Neither did I compare unto her any precious stone: for all gold in comparison of her, is as a little sand, and silver in respect to her shall be counted as clay. I loved her above health and beauty, and chose to have her instead of light: for her light cannot be put out.\par
\switchcolumn*

\LA
\begin{Resp}\Rsign Magna enim sunt judícia tua, Dómine, et inenarrabília verba tua: \Star{} Magnificásti pópulum tuum et honorásti.\end{Resp}
\begin{Resp}\Vsign Transtulísti illos per Mare Rubrum et transvexísti eos per aquam nímiam.\end{Resp}
\begin{Resp}\Rsign Magnificásti pópulum tuum et honorásti.\end{Resp}
\switchcolumn
\EN
\begin{Resp}\Rsign Great are thy judgments, O Lord, and thy words cannot be expressed. \Star{} Thou didst make thy people mighty and honourable.\end{Resp}
\begin{Resp}\Vsign Thou broughtest them through the Red Sea, and leddest them through much water.\end{Resp}
\begin{Resp}\Rsign Thou didst make thy people mighty and honourable.\end{Resp}
\switchcolumn*

\LA
\LessonHead{Lectio III}
\switchcolumn
\EN
\LessonHead{Lesson III}
\switchcolumn*

\LA
\Cite{Sap 7:11-14}
\switchcolumn
\EN
\Cite{Wis 7:11-14}
\switchcolumn*

\LA
\noindent Venérunt autem mihi ómnia bona páriter cum illa, et innumerábilis honéstas per manus illíus; et lætátus sum in ómnibus, quóniam antecedébat me ista sapiéntia, et ignorábam quóniam horum ómnium mater est. Quam sine fictióne dídici et sine invídia commúnico et honestátem illíus non abscóndo. Infinítus enim thesáurus est homínibus; quo qui usi sunt, partícipes facti sunt amicítiæ Dei, propter disciplínæ dona commendáti.\par
\switchcolumn
\EN
\noindent Now all good things came to me together with her, and innumerable riches through her hands, And I rejoiced in all these: for this wisdom went before me, and I knew not that she was the mother of them all. Which I have learned without guile, and communicate without envy, and her riches I hide not. For she is an infinite treasure to men! which they that use, become the friends of God, being commended for the gift of discipline.\par
\switchcolumn*

\LA
\begin{Resp}\Rsign Quæ sunt in corde hóminum, óculi tui vident, Dómine, et in libro tuo ómnia scribéntur: \Star{} Homo videt in fácie, Deus autem in corde.\end{Resp}
\begin{Resp}\Vsign Omnia enim corda scrutátur, et univérsas méntium cogitatiónes intéllegit.\end{Resp}
\begin{Resp}\Rsign Homo videt in fácie, Deus autem in corde.\end{Resp}
\begin{Resp}\Vsign Glória Patri, et Fílio, \Star{} et Spirítui Sancto.\end{Resp}
\begin{Resp}\Rsign Homo videt in fácie, Deus autem in corde.\end{Resp}
\switchcolumn
\EN
\begin{Resp}\Rsign Lord, thine eyes behold all that is in the heart of man, and in thy book are they all written. \Star{} Man looketh on the outward appearance, but God looketh on the heart.\end{Resp}
\begin{Resp}\Vsign For He searcheth all hearts, and understandeth all the imaginations of the thoughts.\end{Resp}
\begin{Resp}\Rsign Man looketh on the outward appearance, but God looketh on the heart.\end{Resp}
\begin{Resp}\Vsign Glory be to the Father, and to the Son, \Star{} and to the Holy Ghost.\end{Resp}
\begin{Resp}\Rsign Man looketh on the outward appearance, but God looketh on the heart.\end{Resp}
\switchcolumn*

\end{paracol}

\Office{Feria V infra Hebdomadam III Augusti}{Thursday of the third week of August}{}
\addcontentsline{toc}{subsection}{Feria V infra Hebdomadam III Augusti \textit{— Thursday of the third week of August}}
\begin{paracol}{2}
\LA
\LessonHead{Lectio I}
\switchcolumn
\EN
\LessonHead{Lesson I}
\switchcolumn*

\LA
\LessonTitle{De libro Sapiéntiæ}
\switchcolumn
\EN
\LessonTitle{Lesson from the book of Wisdom}
\switchcolumn*

\LA
\Cite{Sap 9:13-19}
\switchcolumn
\EN
\Cite{Wis 9:13-19}
\switchcolumn*

\LA
\noindent Quis hóminum póterit scire consílium Dei? aut quis póterit cogitáre quid velit Deus? Cogitatiónes enim mortálium tímidæ, et incértæ providéntiæ nostræ; corpus enim quod corrúmpitur ággravat ánimam, et terréna inhabitátio déprimit sensum multa cogitántem. Et diffícile æstimámus quæ in terra sunt, et quæ in prospéctu sunt invenímus cum labóre; quæ autem in cælis sunt quis investigábit? Sensum autem tuum quis sciet, nisi tu déderis sapiéntiam, et míseris Spíritum Sanctum tuum de altíssimis, et sic corréctæ sint sémitæ eórum qui sunt in terris, et quæ tibi placent didícerint hómines? Nam per sapiéntiam sanáti sunt quicúmque placuérunt tibi, Dómine, a princípio.\par
\switchcolumn
\EN
\noindent For who among men is he that can know the counsel of God? or who can think what the will of God is? For the thoughts of mortal men are fearful, and our counsels uncertain. For the corruptible body is a load upon the soul, and the earthly habitation presseth down the mind that museth upon many things. And hardly do we guess aright at things that are upon earth: and with labour do we find the things that are before us. But the things that are in heaven, who shall search out? And who shall know thy thought, except thou give wisdom, and send thy Holy Spirit from above: And so the ways of them that are upon earth may be corrected, and men may learn the things that please thee? For by wisdom they were healed, whosoever have pleased thee, O Lord, from the beginning.\par
\switchcolumn*

\LA
\begin{Resp}\Rsign In princípio Deus ántequam terram fáceret, priúsquam abýssos constitúeret, priúsquam prodúceret fontes aquárum, \Star{} Antequam montes collocaréntur, ante omnes colles generávit me Dóminus.\end{Resp}
\begin{Resp}\Vsign Quando præparábat cælos, áderam, cum eo cuncta compónens.\end{Resp}
\begin{Resp}\Rsign Antequam montes collocaréntur, ante omnes colles generávit me Dóminus.\end{Resp}
\switchcolumn
\EN
\begin{Resp}\Rsign God possessed me in the beginning, before He made the earth, before He created the depths, before He caused the fountains of water to spring. \Star{} Before the mountains were settled, before there were any hills, did the Lord beget me.\end{Resp}
\begin{Resp}\Vsign When He prepared the heavens, I was there with Him, ordering all things.\end{Resp}
\begin{Resp}\Rsign Before the mountains were settled, before there were any hills, did the Lord beget me.\end{Resp}
\switchcolumn*

\LA
\LessonHead{Lectio II}
\switchcolumn
\EN
\LessonHead{Lesson II}
\switchcolumn*

\LA
\Cite{Sap 10:1-5}
\switchcolumn
\EN
\Cite{Wis 10:1-5}
\switchcolumn*

\LA
\noindent Hæc illum, qui primus formátus est a Deo pater orbis terrárum, cum solus esset creátus, custodívit et edúxit illum a delícto suo, et dedit illi virtútem continéndi ómnia. Ab hac ut recéssit injústus in ira sua, per iram homicídii fratérni depériit. Propter quem, cum aqua deléret terram, sanávit íterum sapiéntia per contemptíbile lignum justum gubérnans. Hæc, et in consénsu nequítiæ cum se natiónes contulíssent, scivit justum et conservávit sine queréla Deo, et in fílii misericórdia fortem custodívit.\par
\switchcolumn
\EN
\noindent She preserved him, that was first formed by God the father of the world, when he was created alone, And she brought him out of his sin, and gave him power to govern all things. But when the unjust went away from her in his anger, he perished by the fury wherewith he murdered his brother. For whose cause, when water destroyed the earth, wisdom healed it again, directing the course of the just by contemptible wood. Moreover when the nations had conspired together to consent to wickedness, she knew the just, and preserved him without blame to God, and kept him strong against the compassion for his son.\par
\switchcolumn*

\LA
\begin{Resp}\Rsign Gyrum cæli circuívi sola, et in flúctibus maris ambulávi, in omni gente et in omni pópulo primátum ténui: \Star{} Superbórum et sublímium colla própria virtúte calcávi.\end{Resp}
\begin{Resp}\Vsign Ego in altíssimis hábito, et thronus meus in colúmna nubis.\end{Resp}
\begin{Resp}\Rsign Superbórum et sublímium colla própria virtúte calcávi.\end{Resp}
\switchcolumn
\EN
\begin{Resp}\Rsign I alone compassed the circuit of heaven, and walked on the waves of the sea. In every nation and in every people, I held the first place. \Star{} In the greatness of my strength have I trodden under my feet the necks of such as be haughty and proud.\end{Resp}
\begin{Resp}\Vsign I dwell in the highest places, and my throne is in a cloudy pillar.\end{Resp}
\begin{Resp}\Rsign In the greatness of my strength have I trodden under my feet the necks of such as be haughty and proud.\end{Resp}
\switchcolumn*

\LA
\LessonHead{Lectio III}
\switchcolumn
\EN
\LessonHead{Lesson III}
\switchcolumn*

\LA
\Cite{Sap 10:6-9}
\switchcolumn
\EN
\Cite{Wis 10:6-9}
\switchcolumn*

\LA
\noindent Hæc justum a pereúntibus ímpiis liberávit fugiéntem, descendénte igne in Pentápolim, quibus in testimónium nequítiæ fumigabúnda constat desérta terra, et incérto témpore fructus habéntes árbores, et incredíbilis ánimæ memória stans figméntum salis. Sapiéntiam enim prætereúntes, non tantum in hoc lapsi sunt ut ignorárent bona, sed et insipiéntiæ suæ reliquérunt homínibus memóriam, ut in his quæ peccavérunt nec latére potuíssent. Sapiéntia autem hos qui se obsérvant a dolóribus liberávit.\par
\switchcolumn
\EN
\noindent She delivered the just man who fled from the wicked that were perishing, when the fire came down upon Pentapolis: Whose land for a testimony of their wickedness is desolate, and smoketh to this day, and the trees bear fruits that ripen not, and a standing pillar of salt is a monument of an incredulous soul. For regarding not wisdom, they did not only slip in this, that they were ignorant of good things, but they left also unto men a memorial of their folly, so that in the things in which they sinned, they could not so much as lie hid. But wisdom hath delivered from sorrow them that attend upon her.\par
\switchcolumn*

\LA
\begin{Resp}\Rsign Emítte, Dómine, sapiéntiam de sede magnitúdinis tuæ, ut mecum sit et mecum labóret: \Star{} Ut sciam, quid accéptum sit coram te omni témpore.\end{Resp}
\begin{Resp}\Vsign Da mihi, Dómine, sédium tuárum assistrícem sapiéntiam.\end{Resp}
\begin{Resp}\Rsign Ut sciam quid accéptum sit coram te omni témpore.\end{Resp}
\begin{Resp}\Vsign Glória Patri, et Fílio, \Star{} et Spirítui Sancto.\end{Resp}
\begin{Resp}\Rsign Ut sciam quid accéptum sit coram te omni témpore.\end{Resp}
\switchcolumn
\EN
\begin{Resp}\Rsign O send out wisdom from the throne of thy glory, O Lord, to be with me, and to labour with me, \Star{} That I may know at all times what is pleasing unto thee.\end{Resp}
\begin{Resp}\Vsign Give me wisdom, O Lord, that sitteth by thy throne.\end{Resp}
\begin{Resp}\Rsign That I may know at all times what is pleasing unto thee.\end{Resp}
\begin{Resp}\Vsign Glory be to the Father, and to the Son, \Star{} and to the Holy Ghost.\end{Resp}
\begin{Resp}\Rsign That I may know at all times what is pleasing unto thee.\end{Resp}
\switchcolumn*

\end{paracol}

\Office{Feria VI infra Hebdomadam III Augusti}{Friday of the third week of August}{}
\addcontentsline{toc}{subsection}{Feria VI infra Hebdomadam III Augusti \textit{— Friday of the third week of August}}
\begin{paracol}{2}
\LA
\LessonHead{Lectio I}
\switchcolumn
\EN
\LessonHead{Lesson I}
\switchcolumn*

\LA
\LessonTitle{De libro Sapiéntiæ}
\switchcolumn
\EN
\LessonTitle{Lesson from the book of Wisdom}
\switchcolumn*

\LA
\Cite{Sap 13:1-3}
\switchcolumn
\EN
\Cite{Wis 13:1-3}
\switchcolumn*

\LA
\noindent Vani autem sunt omnes hómines, in quibus non subest sciéntia Dei; et de his, quæ vidéntur bona, non potuérunt intellégere eum qui est, neque opéribus attendéntes agnovérunt quis esset ártifex; sed aut ignem aut spíritum aut citátum áërem aut gyrum stellárum aut nímiam aquam aut solem et lunam rectóres orbis terrárum deos putavérunt. Quorum si spécie delectáti, deos putavérunt, sciant quanto his dominátor eórum speciósior est; speciéi enim generátor hæc ómnia constítuit.\par
\switchcolumn
\EN
\noindent But all men are vain, in whom there is not the knowledge of God: and who by these good things that are seen, could not understand him that is, neither by attending to the works have acknowledged who was the workman: But have imagined either the fire, or the wind, or the swift air, or the circle of the stars, or the great water, or the sun and moon, to be the gods that rule the world. With whose beauty, if they, being delighted, took them to be gods: let them know how much the Lord of them is more beautiful than they: for the first author of beauty made all those things.\par
\switchcolumn*

\LA
\begin{Resp}\Rsign Da mihi, Dómine, sédium tuárum assistrícem sapiéntiam, et noli me reprobáre a púeris tuis: \Star{} Quóniam servus tuus sum ego, et fílius ancíllæ tuæ.\end{Resp}
\begin{Resp}\Vsign Mitte illam de sede magnitúdinis tuæ, ut mecum sit et mecum labóret.\end{Resp}
\begin{Resp}\Rsign Quóniam servus tuus sum ego, et fílius ancíllæ tuæ.\end{Resp}
\switchcolumn
\EN
\begin{Resp}\Rsign Give me wisdom, O Lord, that sitteth by thy throne, and reject me not from among thy children. \Star{} For I am thy servant and son of thine handmaid.\end{Resp}
\begin{Resp}\Vsign O send her out from the throne of thy glory, to be with me and to labour with me.\end{Resp}
\begin{Resp}\Rsign For I am thy servant and son of thine handmaid.\end{Resp}
\switchcolumn*

\LA
\LessonHead{Lectio II}
\switchcolumn
\EN
\LessonHead{Lesson II}
\switchcolumn*

\LA
\Cite{Sap 13:4-7}
\switchcolumn
\EN
\Cite{Wis 13:4-7}
\switchcolumn*

\LA
\noindent Aut, si virtútem et ópera eórum miráti sunt, intéllegant ab illis, quóniam qui hæc fecit, fórtior est illis; a magnitúdine enim speciéi et creatúræ cognoscibíliter póterit Creátor horum vidéri. Sed tamen adhuc in his minor est queréla, et hi enim fortásse errant, Deum quæréntes et voléntes inveníre. Etenim, cum in opéribus illíus converséntur, inquírunt et persuásum habent quóniam bona sunt quæ vidéntur.\par
\switchcolumn
\EN
\noindent Or if they admired their power and their effects, let them understand by them, that he that made them, is mightier than they: For by the greatness of the beauty, and of the creature, the creator of them may be seen, so as to be known thereby. But yet as to these they are less to be blamed. For they perhaps err, seeking God, and desirous to find him. For being conversant among his works, they search: and they are persuaded that the things are good which are seen.\par
\switchcolumn*

\LA
\begin{Resp}\Rsign Inítium sapiéntiæ timor Dómini: \Star{} Intelléctus bonus ómnibus faciéntibus eum: laudátio ejus manet in sǽculum sǽculi.\end{Resp}
\begin{Resp}\Vsign Diléctio illíus custódia legum est: quia omnis sapiéntia timor Dómini.\end{Resp}
\begin{Resp}\Rsign Intelléctus bonus ómnibus faciéntibus eum: laudátio ejus manet in sǽculum sǽculi.\end{Resp}
\switchcolumn
\EN
\begin{Resp}\Rsign The fear of the Lord is the beginning of wisdom. \Star{} A good understanding have all they that do His commandments. His praise endureth for ever.\end{Resp}
\begin{Resp}\Vsign Love is the keeping of her laws, for all wisdom is the fear of the Lord.\end{Resp}
\begin{Resp}\Rsign A good understanding have all they that do His commandments. His praise endureth for ever.\end{Resp}
\switchcolumn*

\LA
\LessonHead{Lectio III}
\switchcolumn
\EN
\LessonHead{Lesson III}
\switchcolumn*

\LA
\Cite{Sap 13:8-10}
\switchcolumn
\EN
\Cite{Wis 13:8-10}
\switchcolumn*

\LA
\noindent Iterum autem nec his debet ignósci. Si enim tantum potuérunt scire, ut possent æstimáre sǽculum, quómodo hujus Dóminum non facílius invenérunt? Infelíces autem sunt, et inter mórtuos spes illórum est, qui appellavérunt deos ópera mánuum hóminum, aurum et argéntum, artis inventiónem et similitúdines animálium aut lápidem inútilem opus manus antíquæ.\par
\switchcolumn
\EN
\noindent But then again they are not to be pardoned. For if they were able to know so much as to make a judgment of the world: how did they not more easily find out the Lord thereof? But unhappy are they, and their hope is among the dead, who have called gods the works of the hands of men, gold and silver, the inventions of art, and the resemblances of beasts, or an unprofitable stone the work of an ancient hand.\par
\switchcolumn*

\LA
\begin{Resp}\Rsign Verbum iníquum et dolósum longe fac a me, Dómine: \Star{} Divítias et paupertátem ne déderis mihi, sed tantum víctui meo tríbue necessária.\end{Resp}
\begin{Resp}\Vsign Duo rogávi te, ne déneges mihi ántequam móriar.\end{Resp}
\begin{Resp}\Rsign Divítias et paupertátem ne déderis mihi, sed tantum víctui meo tríbue necessária.\end{Resp}
\begin{Resp}\Vsign Glória Patri, et Fílio, \Star{} et Spirítui Sancto.\end{Resp}
\begin{Resp}\Rsign Divítias et paupertátem ne déderis mihi, sed tantum víctui meo tríbue necessária.\end{Resp}
\switchcolumn
\EN
\begin{Resp}\Rsign Lord, remove far from me vanity and lies. \Star{} Give me neither poverty nor riches, but feed me with food convenient for me.\end{Resp}
\begin{Resp}\Vsign Two things have I required of thee; deny me them not before I die.\end{Resp}
\begin{Resp}\Rsign Give me neither poverty nor riches, but feed me with food convenient for me.\end{Resp}
\begin{Resp}\Vsign Glory be to the Father, and to the Son, \Star{} and to the Holy Ghost.\end{Resp}
\begin{Resp}\Rsign Give me neither poverty nor riches, but feed me with food convenient for me.\end{Resp}
\switchcolumn*

\end{paracol}

\Office{Sabbato infra Hebdomadam III Augusti}{Saturday of the third week of August}{}
\addcontentsline{toc}{subsection}{Sabbato infra Hebdomadam III Augusti \textit{— Saturday of the third week of August}}
\begin{paracol}{2}
\LA
\LessonHead{Lectio I}
\switchcolumn
\EN
\LessonHead{Lesson I}
\switchcolumn*

\LA
\LessonTitle{De libro Sapiéntiæ}
\switchcolumn
\EN
\LessonTitle{Lesson from the book of Wisdom}
\switchcolumn*

\LA
\Cite{Sap 15:1-3}
\switchcolumn
\EN
\Cite{Wis 15:1-3}
\switchcolumn*

\LA
\noindent Tu autem, Deus noster, suávis et verus es, pátiens et in misericórdia dispónens ómnia. Etenim, si peccavérimus, tui sumus sciéntes magnitúdinem tuam; et, si non peccavérimus, scimus quóniam apud te sumus computáti. Nosse enim te, consummáta justítia est, et scire justítiam et virtútem tuam, radix est immortalitátis.\par
\switchcolumn
\EN
\noindent But thou, our God, art gracious and true, patient, and ordering all things in mercy. For if we sin, we are thine, knowing thy greatness: and if we sin not, we know that we are counted with thee. For to know thee is perfect justice: and to know thy justice, and thy power, is the root of immortality.\par
\switchcolumn*

\LA
\begin{Resp}\Rsign Dómine, Pater et Deus vitæ meæ, ne derelínquas me in cogitátu malígno: extolléntiam oculórum meórum ne déderis mihi, et desidérium malígnum avérte a me, Dómine; aufer a me concupiscéntiam, \Star{} Et ánimo irreverénti et infruníto ne tradas me, Dómine.\end{Resp}
\begin{Resp}\Vsign Ne derelínquas me, Dómine, ne accréscant ignorántiæ meæ, nec multiplicéntur delícta mea.\end{Resp}
\begin{Resp}\Rsign Et ánimo irreverénti et infruníto ne tradas me, Dómine.\end{Resp}
\switchcolumn
\EN
\begin{Resp}\Rsign O Lord, Father and God of my life, leave me not to evil counsels; give me not a proud look, but turn away from me an haughty mind, O Lord turn away from me concupiscence, \Star{} And give me not over unto an impudent and froward mind, O Lord!\end{Resp}
\begin{Resp}\Vsign Leave me not, O Lord, lest mine ignorance increase, and my sins abound.\end{Resp}
\begin{Resp}\Rsign And give me not over unto an impudent and froward mind, O Lord.\end{Resp}
\switchcolumn*

\LA
\LessonHead{Lectio II}
\switchcolumn
\EN
\LessonHead{Lesson II}
\switchcolumn*

\LA
\Cite{Sap 15:4-6}
\switchcolumn
\EN
\Cite{Wis 15:4-6}
\switchcolumn*

\LA
\noindent Non enim in errórem indúxit nos hóminum malæ artis excogitátio, nec umbra pictúræ labor sine fructu, effígies sculpta per vários colóres, cujus aspéctus insensáto dat concupiscéntiam, et díligit mórtuæ imáginis effígiem sine ánima. Malórum amatóres digni sunt qui spem hábeant in tálibus, et qui fáciunt illos et qui díligunt, et qui colunt.\par
\switchcolumn
\EN
\noindent For the invention of mischievous men hath not deceived us, nor the shadow of a picture, a fruitless labour, a graven figure with diverse colours, The sight whereof enticeth the fool to lust after it, and he loveth the lifeless figure of a dead image. The lovers of evil things deserve to have no better things to trust in, both they that make them, and they that love them, and they that worship them.\par
\switchcolumn*

\LA
\begin{Resp}\Rsign Magna enim sunt judícia tua, Dómine, et inenarrabília verba tua: \Star{} Magnificásti pópulum tuum et honorásti.\end{Resp}
\begin{Resp}\Vsign Transtulísti illos per Mare Rubrum et transvexísti eos per aquam nímiam.\end{Resp}
\begin{Resp}\Rsign Magnificásti pópulum tuum et honorásti.\end{Resp}
\switchcolumn
\EN
\begin{Resp}\Rsign Great are thy judgments, O Lord, and thy words cannot be expressed. \Star{} Thou didst make thy people mighty and honourable.\end{Resp}
\begin{Resp}\Vsign Thou broughtest them through the Red Sea, and leddest them through much water.\end{Resp}
\begin{Resp}\Rsign Thou didst make thy people mighty and honourable.\end{Resp}
\switchcolumn*

\LA
\LessonHead{Lectio III}
\switchcolumn
\EN
\LessonHead{Lesson III}
\switchcolumn*

\LA
\Cite{Sap 15:7-8}
\switchcolumn
\EN
\Cite{Wis 15:7-8}
\switchcolumn*

\LA
\noindent Sed et fígulus mollem terram premens laborióse fingit ad usus nostros unumquódque vas; et de eódem luto fingit, quæ munda sunt in usum vasa, et simíliter quæ his sunt contrária; horum autem vasórum quis sit usus judex est fígulus. Et cum labóre vano deum fingit de eódem luto ille, qui paulo ante de terra factus fúerat, et post pusíllum redúcit se unde accéptus est, repetítus ánimæ débitum quam habébat.\par
\switchcolumn
\EN
\noindent The potter also tempering soft earth, with labour fashioneth every vessel for our service, and of the same clay he maketh both vessels that are for clean uses, and likewise such as serve to the contrary: but what is the use of these vessels, the potter is the judge. And of the same clay by a vain labour he maketh a god: he who a little before was made of earth himself, and a little after returneth to the same out of which he was taken, when his life which was lent him shall be called for again.\par
\switchcolumn*

\LA
\begin{Resp}\Rsign Quæ sunt in corde hóminum, óculi tui vident, Dómine, et in libro tuo ómnia scribéntur: \Star{} Homo videt in fácie, Deus autem in corde.\end{Resp}
\begin{Resp}\Vsign Omnia enim corda scrutátur, et univérsas méntium cogitatiónes intéllegit.\end{Resp}
\begin{Resp}\Rsign Homo videt in fácie, Deus autem in corde.\end{Resp}
\begin{Resp}\Vsign Glória Patri, et Fílio, \Star{} et Spirítui Sancto.\end{Resp}
\begin{Resp}\Rsign Homo videt in fácie, Deus autem in corde.\end{Resp}
\switchcolumn
\EN
\begin{Resp}\Rsign Lord, thine eyes behold all that is in the heart of man, and in thy book are they all written. \Star{} Man looketh on the outward appearance, but God looketh on the heart.\end{Resp}
\begin{Resp}\Vsign For He searcheth all hearts, and understandeth all the imaginations of the thoughts.\end{Resp}
\begin{Resp}\Rsign Man looketh on the outward appearance, but God looketh on the heart.\end{Resp}
\begin{Resp}\Vsign Glory be to the Father, and to the Son, \Star{} and to the Holy Ghost.\end{Resp}
\begin{Resp}\Rsign Man looketh on the outward appearance, but God looketh on the heart.\end{Resp}
\switchcolumn*

\end{paracol}

\Office{Dominica IV Augusti}{Fourth Sunday of August}{}
\addcontentsline{toc}{subsection}{Dominica IV Augusti \textit{— Fourth Sunday of August}}
\begin{paracol}{2}
\LA
\LessonHead{Lectio I}
\switchcolumn
\EN
\LessonHead{Lesson I}
\switchcolumn*

\LA
\LessonTitle{Incipit liber Ecclesiástici}
\switchcolumn
\EN
\LessonTitle{Lesson from the book of Ecclesiasticus}
\switchcolumn*

\LA
\Cite{Sir 1:1-5}
\switchcolumn
\EN
\Cite{Sir 1:1-5}
\switchcolumn*

\LA
\noindent Omnis sapiéntia a Dómino Deo est et cum illo fuit semper et est ante ævum. Arénam maris et plúviæ guttas, et dies sǽculi, quis dinumerávit? Altitúdinem cæli, et latitúdinem terræ, et profúndum abýssi, quis diménsus est? Sapiéntiam Dei præcedéntem ómnia quis investigávit? Prior ómnium creáta est sapiéntia, et intelléctus prudéntiæ ab ævo. Fons sapiéntiæ Verbum Dei in excélsis, et ingréssus illíus mandáta ætérna.\par
\switchcolumn
\EN
\noindent All wisdom is from the Lord God, and hath been always with him, and is before all time. Who hath numbered the sand of the sea, and the drops of rain, and the days of the world? Who hath measured the height of heaven, and the breadth of the earth, and the depth of the abyss? Who hath searched out the wisdom of God that goeth before all things? Wisdom hath been created before all things, and the understanding of prudence from everlasting. The word of God on high is the fountain of wisdom, and her ways are everlasting commandments.\par
\switchcolumn*

\LA
\begin{Resp}\Rsign In princípio Deus ántequam terram fáceret, priúsquam abýssos constitúeret, priúsquam prodúceret fontes aquárum, \Star{} Antequam montes collocaréntur, ante omnes colles generávit me Dóminus.\end{Resp}
\begin{Resp}\Vsign Quando præparábat cælos, áderam, cum eo cuncta compónens.\end{Resp}
\begin{Resp}\Rsign Antequam montes collocaréntur, ante omnes colles generávit me Dóminus.\end{Resp}
\switchcolumn
\EN
\begin{Resp}\Rsign God possessed me in the beginning, before He made the earth, before He created the depths, before He caused the fountains of water to spring. \Star{} Before the mountains were settled, before there were any hills, did the Lord beget me.\end{Resp}
\begin{Resp}\Vsign When He prepared the heavens, I was there with Him, ordering all things.\end{Resp}
\begin{Resp}\Rsign Before the mountains were settled, before there were any hills, did the Lord beget me.\end{Resp}
\switchcolumn*

\LA
\LessonHead{Lectio II}
\switchcolumn
\EN
\LessonHead{Lesson II}
\switchcolumn*

\LA
\Cite{Sir 1:6-10}
\switchcolumn
\EN
\Cite{Sir 1:6-10}
\switchcolumn*

\LA
\noindent Radix sapiéntiæ cui reveláta est? et astútias illíus quis agnóvit? Disciplína sapiéntiæ cui reveláta est et manifestáta? et multiplicatiónem ingréssus illíus quis intelléxit? Unus est altíssimus, Creátor omnípotens, et Rex potens et metuéndus nimis, sedens super thronum illíus et dóminans, Deus. Ipse creávit illam in Spíritu Sancto, et vidit, et dinumerávit, et mensus est. Et effúdit illam super ómnia ópera sua, et super omnem carnem, secúndum datum suum, et prǽbuit illam diligéntibus se.\par
\switchcolumn
\EN
\noindent To whom hath the root of wisdom been revealed, and who hath known her wise counsels? To whom hath the discipline of wisdom been revealed and made manifest? and who hath understood the multiplicity of her steps? There is one most high Creator Almighty, and a powerful king, and greatly to be feared, who sitteth upon his throne, and is the God of dominion. He created her in the Holy Ghost, and saw her, and numbered her, and measured her. And he poured her out upon all his works, and upon all flesh according to his gift, and hath given her to them that love him.\par
\switchcolumn*

\LA
\begin{Resp}\Rsign Gyrum cæli circuívi sola, et in flúctibus maris ambulávi, in omni gente et in omni pópulo primátum ténui: \Star{} Superbórum et sublímium colla própria virtúte calcávi.\end{Resp}
\begin{Resp}\Vsign Ego in altíssimis hábito, et thronus meus in colúmna nubis.\end{Resp}
\begin{Resp}\Rsign Superbórum et sublímium colla própria virtúte calcávi.\end{Resp}
\switchcolumn
\EN
\begin{Resp}\Rsign I alone compassed the circuit of heaven, and walked on the waves of the sea. In every nation and in every people, I held the first place. \Star{} In the greatness of my strength have I trodden under my feet the necks of such as be haughty and proud.\end{Resp}
\begin{Resp}\Vsign I dwell in the highest places, and my throne is in a cloudy pillar.\end{Resp}
\begin{Resp}\Rsign In the greatness of my strength have I trodden under my feet the necks of such as be haughty and proud.\end{Resp}
\switchcolumn*

\LA
\LessonHead{Lectio III}
\switchcolumn
\EN
\LessonHead{Lesson III}
\switchcolumn*

\LA
\Cite{Sir 1:11-15}
\switchcolumn
\EN
\Cite{Sir 1:11-16}
\switchcolumn*

\LA
\noindent Timor Dómini glória, et gloriátio, et lætítia, et coróna exsultatiónis. Timor Dómini delectábit cor, et dabit lætítiam, et gáudium, et longitúdinem diérum. Timénti Dóminum bene erit in extrémis, et in die defunctiónis suæ benedicétur. Diléctio Dei honorábilis sapiéntia; quibus autem apparúerit in visu díligunt eam in visióne, et in agnitióne magnálium suórum. Inítium sapiéntiæ timor Dómini: et cum fidélibus in vulva concreátus est.\par
\switchcolumn
\EN
\noindent The fear of the Lord is honour, and glory, and gladness, and a crown of joy. The fear of the Lord shall delight the heart, and shall give joy, and gladness, and length of days. With him that feareth the Lord, it shall go well in the latter end, and in the day of his death he shall be blessed. The love of God is honourable wisdom. And they to whom she shall show herself love her by the sight, and by the knowledge of her great works. The fear of the Lord is the beginning of wisdom, and was created with the faithful in the womb, it walketh with chosen women, and is known with the just and faithful.\par
\switchcolumn*

\LA
\begin{Resp}\Rsign Emítte, Dómine, sapiéntiam de sede magnitúdinis tuæ, ut mecum sit et mecum labóret: \Star{} Ut sciam, quid accéptum sit coram te omni témpore.\end{Resp}
\begin{Resp}\Vsign Da mihi, Dómine, sédium tuárum assistrícem sapiéntiam.\end{Resp}
\begin{Resp}\Rsign Ut sciam quid accéptum sit coram te omni témpore.\end{Resp}
\begin{Resp}\Vsign Glória Patri, et Fílio, \Star{} et Spirítui Sancto.\end{Resp}
\begin{Resp}\Rsign Ut sciam quid accéptum sit coram te omni témpore.\end{Resp}
\switchcolumn
\EN
\begin{Resp}\Rsign O send out wisdom from the throne of thy glory, O Lord, to be with me, and to labour with me, \Star{} That I may know at all times what is pleasing unto thee.\end{Resp}
\begin{Resp}\Vsign Give me wisdom, O Lord, that sitteth by thy throne.\end{Resp}
\begin{Resp}\Rsign That I may know at all times what is pleasing unto thee.\end{Resp}
\begin{Resp}\Vsign Glory be to the Father, and to the Son, \Star{} and to the Holy Ghost.\end{Resp}
\begin{Resp}\Rsign That I may know at all times what is pleasing unto thee.\end{Resp}
\switchcolumn*

\LA
\LessonHead{Lectio IV}
\switchcolumn
\EN
\LessonHead{Lesson IV}
\switchcolumn*

\LA
\LessonTitle{Ex libro Morálium sancti Gregórii Papæ}
\switchcolumn
\EN
\LessonTitle{From the Book of Moral Reflections on Job, written by Pope St. Gregory the Great.}
\switchcolumn*

\LA
\Source{Liber 1, cap. 10 in cap. 1 Job}
\switchcolumn
\EN
\Source{Bk. i. 10 on Job i.}
\switchcolumn*

\LA
\noindent Sunt nonnúlli qui vitam suam négligunt et, dum transitória áppetunt, dum ætérna vel non intélligunt, vel intellécta contémnunt, nec dolórem séntiunt, nec habére consílium sciunt; cumque supérna quæ amisérunt, non consíderant, esse se (heu míseri) in bonis felíces putant. Nequáquam enim ad veritátis lucem, cui cónditi fúerant, mentis óculos érigunt; nequáquam ad contemplatiónem pátriæ ætérnæ desidérii áciem tendunt; sed semetípsos in his, ad quæ projécti sunt, deseréntes, vice pátriæ díligunt exsílium quod patiúntur, et in cæcitáte quam tólerant, quasi in claritáte lúminis exsúltant.\par
\switchcolumn
\EN
\noindent Come there are who are careless concerning their true life, greedy of the things which pass away, but as to the things which are eternal, either understand them not, or, understanding them, holding them to be but of little moment they feel no sorrow, nor know how to take wise advice and, in forgetfulness of the heavenly possessions which they have lost, they deem themselves (alas, poor wretches) happy in their goods. They lift not up their eyes to the light of truth for which they were created no keen desire ever maketh them to cast a longing look toward the everlasting Fatherland. Leaving alone the chief end for which they were made, they fix their affections upon the exile which they are enduring, instead of upon their home, and make merry in the blindness which they are suffering, as though it were glorious day-light.\par
\switchcolumn*

\LA
\begin{Resp}\Rsign Da mihi, Dómine, sédium tuárum assistrícem sapiéntiam, et noli me reprobáre a púeris tuis: \Star{} Quóniam servus tuus sum ego, et fílius ancíllæ tuæ.\end{Resp}
\begin{Resp}\Vsign Mitte illam de sede magnitúdinis tuæ, ut mecum sit et mecum labóret.\end{Resp}
\begin{Resp}\Rsign Quóniam servus tuus sum ego, et fílius ancíllæ tuæ.\end{Resp}
\switchcolumn
\EN
\begin{Resp}\Rsign Give me wisdom, O Lord, that sitteth by thy throne, and reject me not from among thy children. \Star{} For I am thy servant and son of thine handmaid.\end{Resp}
\begin{Resp}\Vsign O send her out from the throne of thy glory, to be with me and to labour with me.\end{Resp}
\begin{Resp}\Rsign For I am thy servant and son of thine handmaid.\end{Resp}
\switchcolumn*

\LA
\LessonHead{Lectio V}
\switchcolumn
\EN
\LessonHead{Lesson V}
\switchcolumn*

\LA
\noindent At contra, electórum mentes, dum transitória cuncta nulla esse conspíciunt, ad quæ sint cónditæ exquírunt: cumque eórum satisfactióni nihil extra Deum súfficit, ipsa inquisitiónis exercitatióne fatigáta illórum cogitátio, in Conditóris sui spe et contemplatióne requiéscit, supérnis intérseri cívibus áppetit; et unusquísque eórum adhuc in mundo córpore pósitus, mente tamen extra mundum surgit: ærúmnam exsílii, quam tólerat, deplórat, et ad sublímem pátriam incessántibus se amóris stímulis éxcitat. Cum ergo dolens videt, quam sit ætérnum quod pérdidit, ínvenit salúbre consílium, temporále hoc despícere quod percúrrit: et quo magis crescit consílii sciéntia ut peritúra déserat, eo augétur dolor, quod necdum ad mansúra pertíngat.\par
\switchcolumn
\EN
\noindent But, on the other hand, the understandings of the elect, while they apprehend the things which pass away, perceive them to be indeed nothings, and work towards grasping the true end to which they were created, and since nothing outside God satisfieth them, their thought, wearied by the intensity of speculation, findeth rest in the hope for, and the contemplation of, their Maker. They are fain to take their place among the citizens above, and each one of them, although still placed in the world as concerns his body, doth yet in heart and mind ascend above the world. They bemoan the hardships of the exile which they are enduring, and rouse themselves by the constant pricking of their love, to look to their Fatherland above. When therefore such an one seeth with grief that by sin he hath lost an eternal inheritance, he findeth this healthy counsel, to reckon but lightly the things of time through which he is passing, and as the riper groweth his wise course that he hath chosen, to let be these perishing things, the deeper groweth his sorrow that he hath not yet attained unto the things which endure.\par
\switchcolumn*

\LA
\begin{Resp}\Rsign Inítium sapiéntiæ timor Dómini: \Star{} Intelléctus bonus ómnibus faciéntibus eum: laudátio ejus manet in sǽculum sǽculi.\end{Resp}
\begin{Resp}\Vsign Diléctio illíus custódia legum est: quia omnis sapiéntia timor Dómini.\end{Resp}
\begin{Resp}\Rsign Intelléctus bonus ómnibus faciéntibus eum: laudátio ejus manet in sǽculum sǽculi.\end{Resp}
\switchcolumn
\EN
\begin{Resp}\Rsign The fear of the Lord is the beginning of wisdom. \Star{} A good understanding have all they that do His commandments. His praise endureth for ever.\end{Resp}
\begin{Resp}\Vsign Love is the keeping of her laws, for all wisdom is the fear of the Lord.\end{Resp}
\begin{Resp}\Rsign A good understanding have all they that do His commandments. His praise endureth for ever.\end{Resp}
\switchcolumn*

\LA
\LessonHead{Lectio VI}
\switchcolumn
\EN
\LessonHead{Lesson VI}
\switchcolumn*

\LA
\noindent Intuéndum quoque est, quod nullus dolor mentis sit in actióne præcipitatiónis. Qui enim sine consíliis vivunt, qui seípsos rerum evéntibus præcípites déserunt, nullo ínterim cogitatiónum dolóre fatigántur. Nam qui solérter in vitæ consílio figit mentem, caute sese in omni actióne circumspiciéndo consíderat; et ne ex re quæ ágitur, repentínus finis adversúsque surrípiat, hunc prius mólliter pósito pede cogitatiónis palpat: pensat, ne ab his quæ agénda sunt, præpédiat formído; ne in his quæ differénda sunt, præcipitátio impéllat; ne prava per concupiscéntiam apérto bello súperent; ne recta per inánem glóriam insidiándo supplántent.\par
\switchcolumn
\EN
\noindent We must also realise that they who are headlong in their courses, feel not sorrow of heart. They that live without thought, who leave themselves recklessly to the guidance of events, escape the weariness of thought. He that ordereth his life by prudent consideration, looketh carefully around him before each thing that he doth, and, like a man, that before advancing on an uncertain way, trieth the ground with his foot, so he taketh thought beforehand, lest some sudden and evil thing should happen to him he considereth whether that which he hath a mind to do is not forbidden to him by caution, whether he be not too hasty about things which were better put off to another season, lest evil should overcome him by open attack upon his lusts, or even good undo him by the in-bringing of vain glory.\par
\switchcolumn*

\LA
\begin{Resp}\Rsign Inítium sapiéntiæ timor Dómini: \Star{} Intelléctus bonus ómnibus faciéntibus eum: laudátio ejus manet in sǽculum sǽculi.\end{Resp}
\begin{Resp}\Vsign Diléctio illíus custódia legum est: quia omnis sapiéntia timor Dómini.\end{Resp}
\begin{Resp}\Rsign Intelléctus bonus ómnibus faciéntibus eum: laudátio ejus manet in sǽculum sǽculi.\end{Resp}
\begin{Resp}\Vsign Glória Patri, et Fílio, \Star{} et Spirítui Sancto.\end{Resp}
\begin{Resp}\Rsign Intelléctus bonus ómnibus faciéntibus eum: laudátio ejus manet in sǽculum sǽculi.\end{Resp}
\switchcolumn
\EN
\begin{Resp}\Rsign The fear of the Lord is the beginning of wisdom. \Star{} A good understanding have all they that do His commandments. His praise endureth for ever.\end{Resp}
\begin{Resp}\Vsign Love is the keeping of her laws, for all wisdom is the fear of the Lord.\end{Resp}
\begin{Resp}\Rsign A good understanding have all they that do His commandments. His praise endureth for ever.\end{Resp}
\begin{Resp}\Vsign Glory be to the Father, and to the Son, \Star{} and to the Holy Ghost.\end{Resp}
\begin{Resp}\Rsign A good understanding have all they that do His commandments. His praise endureth for ever.\end{Resp}
\switchcolumn*

\end{paracol}

\Office{Feria II infra Hebdomadam IV Augusti}{Monday of the fourth week of August}{}
\addcontentsline{toc}{subsection}{Feria II infra Hebdomadam IV Augusti \textit{— Monday of the fourth week of August}}
\begin{paracol}{2}
\LA
\LessonHead{Lectio I}
\switchcolumn
\EN
\LessonHead{Lesson I}
\switchcolumn*

\LA
\LessonTitle{De libro Ecclesiástici}
\switchcolumn
\EN
\LessonTitle{Lesson from the book of Ecclesiasticus}
\switchcolumn*

\LA
\Cite{Sir 1:22-26}
\switchcolumn
\EN
\Cite{Sir 1:22-26}
\switchcolumn*

\LA
\noindent Coróna sapiéntiæ timor Dómini, replens pacem et salútis fructum; et vidit et dinumerávit eam: útraque autem sunt dona Dei. Sciéntiam et intelléctum prudéntiæ sapiéntia compartiétur et glóriam tenéntium se exáltat. Radix sapiéntiæ est timére Dóminum, et rami illíus longǽvi. In thesáuris sapiéntiæ intelléctus et sciéntiæ religiósitas, exsecrátio autem peccatóribus sapiéntia.\par
\switchcolumn
\EN
\noindent The fear of the Lord is a crown of wisdom, filling up peace and the fruit of salvation: And it hath seen, and numbered her: but both are the gifts of God. Wisdom shall distribute knowledge, and understanding of prudence: and exalteth the glory of them that hold her. The root of wisdom is to fear the Lord: and the branches thereof are longlived. In the treasures of wisdom is understanding, and religiousness of knowledge: but to sinners wisdom is an abomination\par
\switchcolumn*

\LA
\begin{Resp}\Rsign Ne derelínquas me, Dómine, pater et dominátor vitæ meæ, ut non córruam in conspéctu adversariórum meórum: \Star{} Ne gáudeat de me inimícus meus.\end{Resp}
\begin{Resp}\Vsign Apprehénde arma et scutum et exsúrge in adjutórium mihi.\end{Resp}
\begin{Resp}\Rsign Ne gáudeat de me inimícus meus.\end{Resp}
\switchcolumn
\EN
\begin{Resp}\Rsign O Lord, Father and Governor of my life, leave me not, lest I fall before mine adversaries, \Star{} and mine enemy rejoice over me.\end{Resp}
\begin{Resp}\Vsign Take hold of shield and buckler, and stand up for mine help.\end{Resp}
\begin{Resp}\Rsign Lest mine enemy rejoice over me.\end{Resp}
\switchcolumn*

\LA
\LessonHead{Lectio II}
\switchcolumn
\EN
\LessonHead{Lesson II}
\switchcolumn*

\LA
\Cite{Sir 1:27-33}
\switchcolumn
\EN
\Cite{Sir 1:27-33}
\switchcolumn*

\LA
\noindent Timor Dómini expéllit peccátum; nam qui sine timóre est non póterit justificári; iracúndia enim animositátis illíus subvérsio illíus est. Usque in tempus sustinébit pátiens, et póstea reddítio jucunditátis. Bonus sensus usque in tempus abscóndet verba illíus, et lábia multórum enarrábunt sensum illíus. In thesáuris sapiéntiæ significátio disciplínæ, exsecrátio autem peccatóri cultúra Dei. Fili, concupíscens sapiéntiam consérva justítiam, et Deus præbébit illam tibi.\par
\switchcolumn
\EN
\noindent The fear of the Lord driveth out sin: For he that is without fear, cannot be justified: for the wrath of his high spirits is his ruin. A patient man shall bear for a time, and afterwards joy shall be restored to him. A good understanding will hide his words for a time, and the lips of many shall declare his wisdom. In the treasures of wisdom is the signification of discipline: But the worship of God is an abomination to a sinner. Son, if thou desire wisdom, keep justice, and God will give her to thee\par
\switchcolumn*

\LA
\begin{Resp}\Rsign Magna enim sunt judícia tua, Dómine, et inenarrabília verba tua: \Star{} Magnificásti pópulum tuum et honorásti.\end{Resp}
\begin{Resp}\Vsign Transtulísti illos per Mare Rubrum et transvexísti eos per aquam nímiam.\end{Resp}
\begin{Resp}\Rsign Magnificásti pópulum tuum et honorásti.\end{Resp}
\switchcolumn
\EN
\begin{Resp}\Rsign Great are thy judgments, O Lord, and thy words cannot be expressed. \Star{} Thou didst make thy people mighty and honourable.\end{Resp}
\begin{Resp}\Vsign Thou broughtest them through the Red Sea, and leddest them through much water.\end{Resp}
\begin{Resp}\Rsign Thou didst make thy people mighty and honourable.\end{Resp}
\switchcolumn*

\LA
\LessonHead{Lectio III}
\switchcolumn
\EN
\LessonHead{Lesson III}
\switchcolumn*

\LA
\Cite{Sir 1:34-40}
\switchcolumn
\EN
\Cite{Sir 1:34-40}
\switchcolumn*

\LA
\noindent Sapiéntia enim et disciplína timor Dómini, et quod beneplácitum est illi, fides et mansuetúdo, et adimplébit thesáuros illíus. Ne sis incredíbilis timóri Dómini, et ne accésseris ad illum dúplici corde. Ne fúeris hypócrita in conspéctu hóminum et non scandalizéris in lábiis tuis. Atténde in illis, ne forte cadas et addúcas ánimæ tuæ inhonoratiónem, et revélet Deus abscónsa tua et in médio synagógæ elídat te; quóniam accessísti malígne ad Dóminum, et cor tuum plenum est dolo et fallácia.\par
\switchcolumn
\EN
\noindent For the fear of the Lord is wisdom and discipline: and that which is agreeable to him, Is faith, and meekness: and he will fill up his treasures. Be not incredulous to the fear of the Lord: and come not to him with a double heart. Be not a hypocrite in the sight of men, and let not thy lips be a stumblingblock to thee. Watch over them, lest thou fall, and bring dishonour upon thy soul, And God discover thy secrets, and cast thee down in the midst of the congregation. Because thou camest to the Lord wickedly, and thy heart is full of guile and deceit.\par
\switchcolumn*

\LA
\begin{Resp}\Rsign Quæ sunt in corde hóminum, óculi tui vident, Dómine, et in libro tuo ómnia scribéntur: \Star{} Homo videt in fácie, Deus autem in corde.\end{Resp}
\begin{Resp}\Vsign Omnia enim corda scrutátur, et univérsas méntium cogitatiónes intéllegit.\end{Resp}
\begin{Resp}\Rsign Homo videt in fácie, Deus autem in corde.\end{Resp}
\begin{Resp}\Vsign Glória Patri, et Fílio, \Star{} et Spirítui Sancto.\end{Resp}
\begin{Resp}\Rsign Homo videt in fácie, Deus autem in corde.\end{Resp}
\switchcolumn
\EN
\begin{Resp}\Rsign Lord, thine eyes behold all that is in the heart of man, and in thy book are they all written. \Star{} Man looketh on the outward appearance, but God looketh on the heart.\end{Resp}
\begin{Resp}\Vsign For He searcheth all hearts, and understandeth all the imaginations of the thoughts.\end{Resp}
\begin{Resp}\Rsign Man looketh on the outward appearance, but God looketh on the heart.\end{Resp}
\begin{Resp}\Vsign Glory be to the Father, and to the Son, \Star{} and to the Holy Ghost.\end{Resp}
\begin{Resp}\Rsign Man looketh on the outward appearance, but God looketh on the heart.\end{Resp}
\switchcolumn*

\end{paracol}

\Office{Feria III infra Hebdomadam IV Augusti}{Tuesday of the fourth week of August}{}
\addcontentsline{toc}{subsection}{Feria III infra Hebdomadam IV Augusti \textit{— Tuesday of the fourth week of August}}
\begin{paracol}{2}
\LA
\LessonHead{Lectio I}
\switchcolumn
\EN
\LessonHead{Lesson I}
\switchcolumn*

\LA
\LessonTitle{De libro Ecclesiástici}
\switchcolumn
\EN
\LessonTitle{Lesson from the book of Ecclesiasticus}
\switchcolumn*

\LA
\Cite{Sir 2:1-3}
\switchcolumn
\EN
\Cite{Sir 2:1-3}
\switchcolumn*

\LA
\noindent Fili, accédens ad servitútem Dei sta in justítia et timóre, et prǽpara ánimam tuam ad tentatiónem. Déprime cor tuum et sústine, inclína aurem tuam et súscipe verba intelléctus et ne festínes in témpore obductiónis. Sústine sustentatiónes Dei, conjúngere Deo et sústine, ut crescat in novíssimo vita tua.\par
\switchcolumn
\EN
\noindent Son, when thou comest to the service of God, stand in justice and in fear, and prepare thy soul for temptation. Humble thy heart, and endure: incline thy ear, and receive the words of understanding: and make not haste in the time of clouds. Wait on God with patience: join thyself to God, and endure, that thy life may be increased in the latter end.\par
\switchcolumn*

\LA
\begin{Resp}\Rsign Præbe, fili, cor mihi, et óculi tui vias meas custódiant: \Star{} Ut addátur grátia cápiti tuo.\end{Resp}
\begin{Resp}\Vsign Atténde, fili mi, sapiéntiam meam et ad elóquium meum inclína aurem tuam.\end{Resp}
\begin{Resp}\Rsign Ut addátur grátia cápiti tuo.\end{Resp}
\switchcolumn
\EN
\begin{Resp}\Rsign My son, give me thine heart, and let thine eyes observe my ways. \Star{} For they shall be an ornament of grace unto thine head.\end{Resp}
\begin{Resp}\Vsign My son, attend unto my wisdom, and incline thine ear unto my sayings.\end{Resp}
\begin{Resp}\Rsign For they shall be an ornament of grace unto thine head.\end{Resp}
\switchcolumn*

\LA
\LessonHead{Lectio II}
\switchcolumn
\EN
\LessonHead{Lesson II}
\switchcolumn*

\LA
\Cite{Sir 2:4-6}
\switchcolumn
\EN
\Cite{Sir 2:4-6}
\switchcolumn*

\LA
\noindent Omne quod tibi applícitum fúerit áccipe et in dolóre sústine et in humilitáte tua patiéntiam habe; quóniam in igne probátur aurum et argéntum, hómines vero receptíbiles in camíno humiliatiónis. Crede Deo, et recuperábit te, et dírige viam tuam et spera in illum: serva timórem illíus et in illo veterásce.\par
\switchcolumn
\EN
\noindent Take all that shall be brought upon thee: and in thy sorrow endure, and in thy humiliation keep patience. For gold and silver are tried in the fire, but acceptable men in the furnace of humiliation. Believe God, and he will recover thee: and direct thy way, and trust in him. Keep his fear, and grow old therein.\par
\switchcolumn*

\LA
\begin{Resp}\Rsign Inítium sapiéntiæ timor Dómini: \Star{} Intelléctus bonus ómnibus faciéntibus eum: laudátio ejus manet in sǽculum sǽculi.\end{Resp}
\begin{Resp}\Vsign Diléctio illíus custódia legum est: quia omnis sapiéntia timor Dómini.\end{Resp}
\begin{Resp}\Rsign Intelléctus bonus ómnibus faciéntibus eum: laudátio ejus manet in sǽculum sǽculi.\end{Resp}
\switchcolumn
\EN
\begin{Resp}\Rsign The fear of the Lord is the beginning of wisdom. \Star{} A good understanding have all they that do His commandments. His praise endureth for ever.\end{Resp}
\begin{Resp}\Vsign Love is the keeping of her laws, for all wisdom is the fear of the Lord.\end{Resp}
\begin{Resp}\Rsign A good understanding have all they that do His commandments. His praise endureth for ever.\end{Resp}
\switchcolumn*

\LA
\LessonHead{Lectio III}
\switchcolumn
\EN
\LessonHead{Lesson III}
\switchcolumn*

\LA
\Cite{Sir 2:7-12}
\switchcolumn
\EN
\Cite{Sir 2:7-12}
\switchcolumn*

\LA
\noindent Metuéntes Dóminum, sustinéte misericórdiam ejus et non deflectátis ab illo, ne cadátis. Qui timétis Dóminum, crédite illi, et non evacuábitur merces vestra. Qui timétis Dóminum, speráte in illum, et in oblectatiónem véniet vobis misericórdia. Qui timétis Dóminum dilígite illum, et illuminabúntur corda vestra. Respícite, fílii, natiónes hóminum et scitóte quia nullus sperávit in Dómino et confúsus est. Quis enim permánsit in mandátis ejus et derelíctus est? aut quis invocávit eum, et despéxit illum?\par
\switchcolumn
\EN
\noindent Ye that fear the Lord, wait for his mercy: and go not aside from him, lest ye fall. Ye that fear the Lord, believe him: and your reward shall not be made void. Ye that fear the Lord, hope in him: and mercy shall come to you for your delight. Ye that fear the Lord, love him, and your hearts shall be enlightened. My children behold the generations of men: and know ye that no one hath hoped in the Lord, and hath been confounded. For who hath continued in his commandment, and hath been forsaken? or who hath called upon him, and he despised him?\par
\switchcolumn*

\LA
\begin{Resp}\Rsign Verbum iníquum et dolósum longe fac a me, Dómine: \Star{} Divítias et paupertátem ne déderis mihi, sed tantum víctui meo tríbue necessária.\end{Resp}
\begin{Resp}\Vsign Duo rogávi te, ne déneges mihi ántequam móriar.\end{Resp}
\begin{Resp}\Rsign Divítias et paupertátem ne déderis mihi, sed tantum víctui meo tríbue necessária.\end{Resp}
\begin{Resp}\Vsign Glória Patri, et Fílio, \Star{} et Spirítui Sancto.\end{Resp}
\begin{Resp}\Rsign Divítias et paupertátem ne déderis mihi, sed tantum víctui meo tríbue necessária.\end{Resp}
\switchcolumn
\EN
\begin{Resp}\Rsign Lord, remove far from me vanity and lies. \Star{} Give me neither poverty nor riches, but feed me with food convenient for me.\end{Resp}
\begin{Resp}\Vsign Two things have I required of thee; deny me them not before I die.\end{Resp}
\begin{Resp}\Rsign Give me neither poverty nor riches, but feed me with food convenient for me.\end{Resp}
\begin{Resp}\Vsign Glory be to the Father, and to the Son, \Star{} and to the Holy Ghost.\end{Resp}
\begin{Resp}\Rsign Give me neither poverty nor riches, but feed me with food convenient for me.\end{Resp}
\switchcolumn*

\end{paracol}

\Office{Feria IV infra Hebdomadam IV Augusti}{Wednesday of the fourth week of August}{}
\addcontentsline{toc}{subsection}{Feria IV infra Hebdomadam IV Augusti \textit{— Wednesday of the fourth week of August}}
\begin{paracol}{2}
\LA
\LessonHead{Lectio I}
\switchcolumn
\EN
\LessonHead{Lesson I}
\switchcolumn*

\LA
\LessonTitle{De libro Ecclesiástici}
\switchcolumn
\EN
\LessonTitle{Lesson from the book of Ecclesiasticus}
\switchcolumn*

\LA
\Cite{Sir 3:1-4}
\switchcolumn
\EN
\Cite{Sir 3:1-4}
\switchcolumn*

\LA
\noindent Fílii sapiéntiæ ecclésia justórum, et nátio illórum obediéntia et diléctio. Judícium patris audíte, fílii, et sic fácite ut salvi sitis. Deus enim honorávit patrem in fíliis, et judícium matris exquírens firmávit in fílios. Qui díligit Deum exorábit pro peccátis et continébit se ab illis et in oratióne diérum exaudiétur.\par
\switchcolumn
\EN
\noindent The sons of wisdom are the church of the just: and their generation, obedience and love. Children, hear the judgment of your father, and so do that you may be saved. For God hath made the father honourable to the children: and seeking the judgment of the mothers, hath confirmed it upon the children. He that loveth God, shall obtain pardon for his sins by prayer, and shall refrain himself from them, and shall be heard in the prayer of days.\par
\switchcolumn*

\LA
\begin{Resp}\Rsign Dómine, Pater et Deus vitæ meæ, ne derelínquas me in cogitátu malígno: extolléntiam oculórum meórum ne déderis mihi, et desidérium malígnum avérte a me, Dómine; aufer a me concupiscéntiam, \Star{} Et ánimo irreverénti et infruníto ne tradas me, Dómine.\end{Resp}
\begin{Resp}\Vsign Ne derelínquas me, Dómine, ne accréscant ignorántiæ meæ, nec multiplicéntur delícta mea.\end{Resp}
\begin{Resp}\Rsign Et ánimo irreverénti et infruníto ne tradas me, Dómine.\end{Resp}
\switchcolumn
\EN
\begin{Resp}\Rsign O Lord, Father and God of my life, leave me not to evil counsels; give me not a proud look, but turn away from me an haughty mind, O Lord turn away from me concupiscence, \Star{} And give me not over unto an impudent and froward mind, O Lord!\end{Resp}
\begin{Resp}\Vsign Leave me not, O Lord, lest mine ignorance increase, and my sins abound.\end{Resp}
\begin{Resp}\Rsign And give me not over unto an impudent and froward mind, O Lord.\end{Resp}
\switchcolumn*

\LA
\LessonHead{Lectio II}
\switchcolumn
\EN
\LessonHead{Lesson II}
\switchcolumn*

\LA
\Cite{Sir 3:5-8}
\switchcolumn
\EN
\Cite{Sir 3:5-8}
\switchcolumn*

\LA
\noindent Et, sicut qui thesaurízat, ita et qui honoríficat matrem suam; qui honórat patrem suum jucundábitur in fíliis et in die oratiónis suæ exaudiétur. Qui honórat patrem suum vita vivet longióre, et qui obédit patri refrigerábit matrem. Qui timet Dóminum honórat paréntes, et quasi dóminis sérviet his qui se genuérunt.\par
\switchcolumn
\EN
\noindent And he that honoureth his mother is as one that layeth up a treasure. He that honoureth his father shall have joy in his own children, and in the day of his prayer he shall be heard. He that honoureth his father shall enjoy a long life: and he that obeyeth the father, shall be a comfort to his mother. He that feareth the Lord, honoureth his parents, and will serve them as his masters that brought him into the world.\par
\switchcolumn*

\LA
\begin{Resp}\Rsign Magna enim sunt judícia tua, Dómine, et inenarrabília verba tua: \Star{} Magnificásti pópulum tuum et honorásti.\end{Resp}
\begin{Resp}\Vsign Transtulísti illos per Mare Rubrum et transvexísti eos per aquam nímiam.\end{Resp}
\begin{Resp}\Rsign Magnificásti pópulum tuum et honorásti.\end{Resp}
\switchcolumn
\EN
\begin{Resp}\Rsign Great are thy judgments, O Lord, and thy words cannot be expressed. \Star{} Thou didst make thy people mighty and honourable.\end{Resp}
\begin{Resp}\Vsign Thou broughtest them through the Red Sea, and leddest them through much water.\end{Resp}
\begin{Resp}\Rsign Thou didst make thy people mighty and honourable.\end{Resp}
\switchcolumn*

\LA
\LessonHead{Lectio III}
\switchcolumn
\EN
\LessonHead{Lesson III}
\switchcolumn*

\LA
\Cite{Sir 3:9-13}
\switchcolumn
\EN
\Cite{Sir 3:9-13}
\switchcolumn*

\LA
\noindent In ópere, et sermóne, et omni patiéntia, honóra patrem tuum, ut supervéniat tibi benedíctio ab eo, et benedíctio illíus in novíssimo máneat. Benedíctio patris firmat domos filiórum, maledíctio autem matris eradícat fundaménta. Ne gloriéris in contumélia patris tui, non enim est tibi glória ejus confúsio. Glória enim hóminis ex honóre patris sui, et dédecus fílii pater sine honóre.\par
\switchcolumn
\EN
\noindent Honour thy father, in work and word, and all patience, That a blessing may come upon thee from him, and his blessing may remain in the latter end. The father's blessing establisheth the houses of the children: but the mother's curse rooteth up the foundation. Glory not in the dishonour of thy father: for his shame is no glory to thee. For the glory of a man is from the honour of his father, and a father without honour is the disgrace of the son.\par
\switchcolumn*

\LA
\begin{Resp}\Rsign Quæ sunt in corde hóminum, óculi tui vident, Dómine, et in libro tuo ómnia scribéntur: \Star{} Homo videt in fácie, Deus autem in corde.\end{Resp}
\begin{Resp}\Vsign Omnia enim corda scrutátur, et univérsas méntium cogitatiónes intéllegit.\end{Resp}
\begin{Resp}\Rsign Homo videt in fácie, Deus autem in corde.\end{Resp}
\begin{Resp}\Vsign Glória Patri, et Fílio, \Star{} et Spirítui Sancto.\end{Resp}
\begin{Resp}\Rsign Homo videt in fácie, Deus autem in corde.\end{Resp}
\switchcolumn
\EN
\begin{Resp}\Rsign Lord, thine eyes behold all that is in the heart of man, and in thy book are they all written. \Star{} Man looketh on the outward appearance, but God looketh on the heart.\end{Resp}
\begin{Resp}\Vsign For He searcheth all hearts, and understandeth all the imaginations of the thoughts.\end{Resp}
\begin{Resp}\Rsign Man looketh on the outward appearance, but God looketh on the heart.\end{Resp}
\begin{Resp}\Vsign Glory be to the Father, and to the Son, \Star{} and to the Holy Ghost.\end{Resp}
\begin{Resp}\Rsign Man looketh on the outward appearance, but God looketh on the heart.\end{Resp}
\switchcolumn*

\end{paracol}

\Office{Feria V infra Hebdomadam IV Augusti}{Thursday of the fourth week of August}{}
\addcontentsline{toc}{subsection}{Feria V infra Hebdomadam IV Augusti \textit{— Thursday of the fourth week of August}}
\begin{paracol}{2}
\LA
\LessonHead{Lectio I}
\switchcolumn
\EN
\LessonHead{Lesson I}
\switchcolumn*

\LA
\LessonTitle{De libro Ecclesiástici}
\switchcolumn
\EN
\LessonTitle{Lesson from the book of Ecclesiasticus}
\switchcolumn*

\LA
\Cite{Sir 3:22-26}
\switchcolumn
\EN
\Cite{Sir 3:22-26}
\switchcolumn*

\LA
\noindent Altióra te ne quæsíeris, et fortióra te ne scrutátus fúeris; sed, quæ præcépit tibi Deus, illa cógita semper, et in plúribus opéribus ejus ne fúeris curiósus; non est enim tibi necessárium ea quæ abscóndita sunt vidére óculis tuis. In supervácuis rebus noli scrutári multiplíciter, et in plúribus opéribus ejus non eris curiósus; plúrima enim super sensum hóminum osténsa sunt tibi; multos quoque supplantávit suspício illórum, et in vanitáte detínuit sensus illórum.\par
\switchcolumn
\EN
\noindent Seek not the things that are too high for thee, and search not into things above thy ability: but the things that God hath commanded thee, think on them always, and in many of his works be not curious. For it is not necessary for thee to see with thy eyes those things that are hid. In unnecessary matters be not over curious, and in many of his works thou shalt not be inquisitive. For many things are shown to thee above the understanding of men. And the suspicion of them hath deceived many, and hath detained their minds in vanity.\par
\switchcolumn*

\LA
\begin{Resp}\Rsign In princípio Deus ántequam terram fáceret, priúsquam abýssos constitúeret, priúsquam prodúceret fontes aquárum, \Star{} Antequam montes collocaréntur, ante omnes colles generávit me Dóminus.\end{Resp}
\begin{Resp}\Vsign Quando præparábat cælos, áderam, cum eo cuncta compónens.\end{Resp}
\begin{Resp}\Rsign Antequam montes collocaréntur, ante omnes colles generávit me Dóminus.\end{Resp}
\switchcolumn
\EN
\begin{Resp}\Rsign God possessed me in the beginning, before He made the earth, before He created the depths, before He caused the fountains of water to spring. \Star{} Before the mountains were settled, before there were any hills, did the Lord beget me.\end{Resp}
\begin{Resp}\Vsign When He prepared the heavens, I was there with Him, ordering all things.\end{Resp}
\begin{Resp}\Rsign Before the mountains were settled, before there were any hills, did the Lord beget me.\end{Resp}
\switchcolumn*

\LA
\LessonHead{Lectio II}
\switchcolumn
\EN
\LessonHead{Lesson II}
\switchcolumn*

\LA
\Cite{Sir 3:27-30}
\switchcolumn
\EN
\Cite{Sir 3:27-30}
\switchcolumn*

\LA
\noindent Cor durum habébit male in novíssimo, et qui amat perículum in illo períbit. Cor ingrédiens duas vias non habébit succéssus, et pravus corde in illis scandalizábitur. Cor nequam gravábitur in dolóribus, et peccátor adíciet ad peccándum. Synagógæ superbórum non erit sánitas, frutex enim peccáti radicábitur in illis, et non intellégitur.\par
\switchcolumn
\EN
\noindent A hard heart shall fear evil at the last: and he that loveth danger shall perish in it. A heart that goeth two ways shall not have success, and the perverse of heart shall be scandalized therein. A wicked heart shall be laden with sorrows, and the sinner will add sin to sin. The congregation of the proud shall not be healed: for the plant of wickedness shall take root in them, and it shall not be perceived.\par
\switchcolumn*

\LA
\begin{Resp}\Rsign Gyrum cæli circuívi sola, et in flúctibus maris ambulávi, in omni gente et in omni pópulo primátum ténui: \Star{} Superbórum et sublímium colla própria virtúte calcávi.\end{Resp}
\begin{Resp}\Vsign Ego in altíssimis hábito, et thronus meus in colúmna nubis.\end{Resp}
\begin{Resp}\Rsign Superbórum et sublímium colla própria virtúte calcávi.\end{Resp}
\switchcolumn
\EN
\begin{Resp}\Rsign I alone compassed the circuit of heaven, and walked on the waves of the sea. In every nation and in every people, I held the first place. \Star{} In the greatness of my strength have I trodden under my feet the necks of such as be haughty and proud.\end{Resp}
\begin{Resp}\Vsign I dwell in the highest places, and my throne is in a cloudy pillar.\end{Resp}
\begin{Resp}\Rsign In the greatness of my strength have I trodden under my feet the necks of such as be haughty and proud.\end{Resp}
\switchcolumn*

\LA
\LessonHead{Lectio III}
\switchcolumn
\EN
\LessonHead{Lesson III}
\switchcolumn*

\LA
\Cite{Sir 3:31-34}
\switchcolumn
\EN
\Cite{Sir 3:31-34}
\switchcolumn*

\LA
\noindent Cor sapiéntis intellégitur in sapiéntia, et auris bona áudiet cum omni concupiscéntia sapiéntiam. Sápiens cor et intellegíbile abstinébit se a peccátis et in opéribus justítiæ succéssus habébit. Ignem ardéntem exstínguit aqua, et eleemósyna resístit peccátis; et Deus prospéctor est ejus qui reddit grátiam, méminit ejus in pósterum, et in témpore casus sui invéniet firmaméntum.\par
\switchcolumn
\EN
\noindent The heart of the wise is understood in wisdom, and a good ear will hear wisdom with all desire. A wise heart, and which hath understanding, will abstain from sins, and in the works of justice shall have success. Water quencheth a flaming fire, and alms resisteth sins: And God provideth for him that showeth favour: he remembereth him afterwards, and in the time of his fall he shall find a sure stay.\par
\switchcolumn*

\LA
\begin{Resp}\Rsign Emítte, Dómine, sapiéntiam de sede magnitúdinis tuæ, ut mecum sit et mecum labóret: \Star{} Ut sciam, quid accéptum sit coram te omni témpore.\end{Resp}
\begin{Resp}\Vsign Da mihi, Dómine, sédium tuárum assistrícem sapiéntiam.\end{Resp}
\begin{Resp}\Rsign Ut sciam quid accéptum sit coram te omni témpore.\end{Resp}
\begin{Resp}\Vsign Glória Patri, et Fílio, \Star{} et Spirítui Sancto.\end{Resp}
\begin{Resp}\Rsign Ut sciam quid accéptum sit coram te omni témpore.\end{Resp}
\switchcolumn
\EN
\begin{Resp}\Rsign O send out wisdom from the throne of thy glory, O Lord, to be with me, and to labour with me, \Star{} That I may know at all times what is pleasing unto thee.\end{Resp}
\begin{Resp}\Vsign Give me wisdom, O Lord, that sitteth by thy throne.\end{Resp}
\begin{Resp}\Rsign That I may know at all times what is pleasing unto thee.\end{Resp}
\begin{Resp}\Vsign Glory be to the Father, and to the Son, \Star{} and to the Holy Ghost.\end{Resp}
\begin{Resp}\Rsign That I may know at all times what is pleasing unto thee.\end{Resp}
\switchcolumn*

\end{paracol}

\Office{Feria VI infra Hebdomadam IV Augusti}{Friday of the fourth week of August}{}
\addcontentsline{toc}{subsection}{Feria VI infra Hebdomadam IV Augusti \textit{— Friday of the fourth week of August}}
\begin{paracol}{2}
\LA
\LessonHead{Lectio I}
\switchcolumn
\EN
\LessonHead{Lesson I}
\switchcolumn*

\LA
\LessonTitle{De libro Ecclesiástici}
\switchcolumn
\EN
\LessonTitle{Lesson from the book of Ecclesiasticus}
\switchcolumn*

\LA
\Cite{Sir 4:1-4}
\switchcolumn
\EN
\Cite{Sir 4:1-4}
\switchcolumn*

\LA
\noindent Fili, eleemósynam páuperis ne defráudes et óculos tuos ne transvértas a páupere. Animam esuriéntem ne despéxeris et non exásperes páuperem in inópia sua. Cor ínopis ne afflíxeris et non prótrahas datum angustiánti. Rogatiónem contribuláti ne abícias, et non avértas fáciem tuam ab egéno.\par
\switchcolumn
\EN
\noindent Son, defraud not the poor of alms, and turn not away thy eyes from the poor. Despise not the hungry soul: and provoke not the poor in his want. Afflict not the heart of the needy, and defer not to give to him that is in distress. Reject not the petition of the afflicted: and turn not away thy face from the needy.\par
\switchcolumn*

\LA
\begin{Resp}\Rsign Da mihi, Dómine, sédium tuárum assistrícem sapiéntiam, et noli me reprobáre a púeris tuis: \Star{} Quóniam servus tuus sum ego, et fílius ancíllæ tuæ.\end{Resp}
\begin{Resp}\Vsign Mitte illam de sede magnitúdinis tuæ, ut mecum sit et mecum labóret.\end{Resp}
\begin{Resp}\Rsign Quóniam servus tuus sum ego, et fílius ancíllæ tuæ.\end{Resp}
\switchcolumn
\EN
\begin{Resp}\Rsign Give me wisdom, O Lord, that sitteth by thy throne, and reject me not from among thy children. \Star{} For I am thy servant and son of thine handmaid.\end{Resp}
\begin{Resp}\Vsign O send her out from the throne of thy glory, to be with me and to labour with me.\end{Resp}
\begin{Resp}\Rsign For I am thy servant and son of thine handmaid.\end{Resp}
\switchcolumn*

\LA
\LessonHead{Lectio II}
\switchcolumn
\EN
\LessonHead{Lesson II}
\switchcolumn*

\LA
\Cite{Sir 4:5-7}
\switchcolumn
\EN
\Cite{Sir 4:5-7}
\switchcolumn*

\LA
\noindent Ab ínope ne avértas óculos tuos propter iram: et non relínquas quæréntibus tibi retro maledícere; maledicéntis enim tibi in amaritúdine ánimæ, exaudiétur deprecátio illíus; exáudiet autem eum qui fecit illum. Congregatióni páuperum affábilem te fácito: et presbýtero humília ánimam tuam, et magnáto humília caput tuum.\par
\switchcolumn
\EN
\noindent Turn not away thy eyes from the poor for fear of anger: and leave not to them that ask of thee to curse thee behind thy back. For the prayer of him that curseth thee in the bitterness of his soul, shall be heard, for he that made him will hear him. Make thyself affable to the congregation of the poor, and humble thy soul to the ancient, and bow thy head to a great man\par
\switchcolumn*

\LA
\begin{Resp}\Rsign Inítium sapiéntiæ timor Dómini: \Star{} Intelléctus bonus ómnibus faciéntibus eum: laudátio ejus manet in sǽculum sǽculi.\end{Resp}
\begin{Resp}\Vsign Diléctio illíus custódia legum est: quia omnis sapiéntia timor Dómini.\end{Resp}
\begin{Resp}\Rsign Intelléctus bonus ómnibus faciéntibus eum: laudátio ejus manet in sǽculum sǽculi.\end{Resp}
\switchcolumn
\EN
\begin{Resp}\Rsign The fear of the Lord is the beginning of wisdom. \Star{} A good understanding have all they that do His commandments. His praise endureth for ever.\end{Resp}
\begin{Resp}\Vsign Love is the keeping of her laws, for all wisdom is the fear of the Lord.\end{Resp}
\begin{Resp}\Rsign A good understanding have all they that do His commandments. His praise endureth for ever.\end{Resp}
\switchcolumn*

\LA
\LessonHead{Lectio III}
\switchcolumn
\EN
\LessonHead{Lesson III}
\switchcolumn*

\LA
\Cite{Sir 4:8-11}
\switchcolumn
\EN
\Cite{Sir 4:8-11}
\switchcolumn*

\LA
\noindent Declína páuperi sine tristítia aurem tuam et redde débitum tuum et respónde illi pacífica in mansuetúdine. Líbera eum, qui injúriam pátitur, de manu supérbi, et non ácide feras in ánima tua. In judicándo esto pupíllis miséricors ut pater, et pro viro matri illórum; et eris tu velut fílius Altíssimi obédiens, et miserébitur tui magis quam mater.\par
\switchcolumn
\EN
\noindent Bow down thy ear cheerfully to the poor, and pay what thou owest, and answer him peaceable words with mildness. Deliver him that suffereth wrong out of the hand of the proud: and be not fainthearted in thy soul. In judging be merciful to the fatherless as a father, and as a husband to their mother. And thou shalt be as the obedient son of the most High, and he will have mercy on thee more than a mother.\par
\switchcolumn*

\LA
\begin{Resp}\Rsign Verbum iníquum et dolósum longe fac a me, Dómine: \Star{} Divítias et paupertátem ne déderis mihi, sed tantum víctui meo tríbue necessária.\end{Resp}
\begin{Resp}\Vsign Duo rogávi te, ne déneges mihi ántequam móriar.\end{Resp}
\begin{Resp}\Rsign Divítias et paupertátem ne déderis mihi, sed tantum víctui meo tríbue necessária.\end{Resp}
\begin{Resp}\Vsign Glória Patri, et Fílio, \Star{} et Spirítui Sancto.\end{Resp}
\begin{Resp}\Rsign Divítias et paupertátem ne déderis mihi, sed tantum víctui meo tríbue necessária.\end{Resp}
\switchcolumn
\EN
\begin{Resp}\Rsign Lord, remove far from me vanity and lies. \Star{} Give me neither poverty nor riches, but feed me with food convenient for me.\end{Resp}
\begin{Resp}\Vsign Two things have I required of thee; deny me them not before I die.\end{Resp}
\begin{Resp}\Rsign Give me neither poverty nor riches, but feed me with food convenient for me.\end{Resp}
\begin{Resp}\Vsign Glory be to the Father, and to the Son, \Star{} and to the Holy Ghost.\end{Resp}
\begin{Resp}\Rsign Give me neither poverty nor riches, but feed me with food convenient for me.\end{Resp}
\switchcolumn*

\end{paracol}

\Office{Sabbato infra Hebdomadam IV Augusti}{Saturday of the fourth week of August}{}
\addcontentsline{toc}{subsection}{Sabbato infra Hebdomadam IV Augusti \textit{— Saturday of the fourth week of August}}
\begin{paracol}{2}
\LA
\LessonHead{Lectio I}
\switchcolumn
\EN
\LessonHead{Lesson I}
\switchcolumn*

\LA
\LessonTitle{De libro Ecclesiástici}
\switchcolumn
\EN
\LessonTitle{Lesson from the book of Ecclesiasticus}
\switchcolumn*

\LA
\Cite{Sir 4:23-28}
\switchcolumn
\EN
\Cite{Sir 4:23-28}
\switchcolumn*

\LA
\noindent Fili, consérva tempus et devíta a malo. Pro ánima tua ne confundáris dícere verum; est enim confúsio addúcens peccátum, et est confúsio addúcens glóriam et grátiam. Ne accípias fáciem advérsus fáciem tuam nec advérsus ánimam tuam mendácium. Ne revereáris próximum tuum in casu suo, nec retíneas verbum in témpore salútis. Non abscóndas sapiéntiam tuam in decóre suo.\par
\switchcolumn
\EN
\noindent Son, observe the time, and fly from evil. For thy soul be not ashamed to say the truth. For there is a shame that bringeth sin, and there is a shame that bringeth glory and grace. Accept no person against thy own person, nor against thy soul a lie. Reverence not thy neighbour in his fall: And refrain not to speak in the time of salvation. Hide not thy wisdom in her beauty.\par
\switchcolumn*

\LA
\begin{Resp}\Rsign Dómine, Pater et Deus vitæ meæ, ne derelínquas me in cogitátu malígno: extolléntiam oculórum meórum ne déderis mihi, et desidérium malígnum avérte a me, Dómine; aufer a me concupiscéntiam, \Star{} Et ánimo irreverénti et infruníto ne tradas me, Dómine.\end{Resp}
\begin{Resp}\Vsign Ne derelínquas me, Dómine, ne accréscant ignorántiæ meæ, nec multiplicéntur delícta mea.\end{Resp}
\begin{Resp}\Rsign Et ánimo irreverénti et infruníto ne tradas me, Dómine.\end{Resp}
\switchcolumn
\EN
\begin{Resp}\Rsign O Lord, Father and God of my life, leave me not to evil counsels; give me not a proud look, but turn away from me an haughty mind, O Lord turn away from me concupiscence, \Star{} And give me not over unto an impudent and froward mind, O Lord!\end{Resp}
\begin{Resp}\Vsign Leave me not, O Lord, lest mine ignorance increase, and my sins abound.\end{Resp}
\begin{Resp}\Rsign And give me not over unto an impudent and froward mind, O Lord.\end{Resp}
\switchcolumn*

\LA
\LessonHead{Lectio II}
\switchcolumn
\EN
\LessonHead{Lesson II}
\switchcolumn*

\LA
\Cite{Sir 4:29-32}
\switchcolumn
\EN
\Cite{Sir 4:29-32}
\switchcolumn*

\LA
\noindent In lingua enim sapiéntia dignóscitur, et sensus et sciéntia et doctrína in verbo sensáti, et firmaméntum in opéribus justítiæ. Non contradícas verbo veritátis ullo modo, et de mendácio ineruditiónis tuæ confúndere. Non confundáris confitéri peccáta tua, et ne subícias te omni hómini pro peccáto. Noli resístere contra fáciem poténtis, nec conéris contra ictum flúvii.\par
\switchcolumn
\EN
\noindent For by the tongue wisdom is discerned: and understanding, and knowledge, and learning by the word of the wise, and steadfastness in the works of justice. In nowise speak against the truth, but be ashamed of the lie of thy ignorance. Be not ashamed to confess thy sins, but submit not thyself to every man for sin. Resist not against the face of the mighty, and do not strive against the stream of the river\par
\switchcolumn*

\LA
\begin{Resp}\Rsign Magna enim sunt judícia tua, Dómine, et inenarrabília verba tua: \Star{} Magnificásti pópulum tuum et honorásti.\end{Resp}
\begin{Resp}\Vsign Transtulísti illos per Mare Rubrum et transvexísti eos per aquam nímiam.\end{Resp}
\begin{Resp}\Rsign Magnificásti pópulum tuum et honorásti.\end{Resp}
\switchcolumn
\EN
\begin{Resp}\Rsign Great are thy judgments, O Lord, and thy words cannot be expressed. \Star{} Thou didst make thy people mighty and honourable.\end{Resp}
\begin{Resp}\Vsign Thou broughtest them through the Red Sea, and leddest them through much water.\end{Resp}
\begin{Resp}\Rsign Thou didst make thy people mighty and honourable.\end{Resp}
\switchcolumn*

\LA
\LessonHead{Lectio III}
\switchcolumn
\EN
\LessonHead{Lesson III}
\switchcolumn*

\LA
\Cite{Sir 4:33-36}
\switchcolumn
\EN
\Cite{Sir 4:33-36}
\switchcolumn*

\LA
\noindent Pro justítia agonizáre pro ánima tua, et usque ad mortem certa pro justítia, et Deus expugnábit pro te inimícos tuos. Noli citátus esse in lingua tua, et inútilis, et remíssus in opéribus tuis. Noli esse sicut leo in domo tua evértens domésticos tuos, et ópprimens subjéctos tibi. Non sit porrécta manus tua ad accipiéndum et ad dandum collécta.\par
\switchcolumn
\EN
\noindent Strive for justice for thy soul, and even unto death fight for justice, and God will overthrow thy enemies for thee. Be not hasty in thy tongue: and slack and remiss in thy works. Be not as a lion in thy house, terrifying them of thy household, and oppressing them that are under thee. Let not thy hand be stretched out to receive, and shut when thou shouldst give.\par
\switchcolumn*

\LA
\begin{Resp}\Rsign Quæ sunt in corde hóminum, óculi tui vident, Dómine, et in libro tuo ómnia scribéntur: \Star{} Homo videt in fácie, Deus autem in corde.\end{Resp}
\begin{Resp}\Vsign Omnia enim corda scrutátur, et univérsas méntium cogitatiónes intéllegit.\end{Resp}
\begin{Resp}\Rsign Homo videt in fácie, Deus autem in corde.\end{Resp}
\begin{Resp}\Vsign Glória Patri, et Fílio, \Star{} et Spirítui Sancto.\end{Resp}
\begin{Resp}\Rsign Homo videt in fácie, Deus autem in corde.\end{Resp}
\switchcolumn
\EN
\begin{Resp}\Rsign Lord, thine eyes behold all that is in the heart of man, and in thy book are they all written. \Star{} Man looketh on the outward appearance, but God looketh on the heart.\end{Resp}
\begin{Resp}\Vsign For He searcheth all hearts, and understandeth all the imaginations of the thoughts.\end{Resp}
\begin{Resp}\Rsign Man looketh on the outward appearance, but God looketh on the heart.\end{Resp}
\begin{Resp}\Vsign Glory be to the Father, and to the Son, \Star{} and to the Holy Ghost.\end{Resp}
\begin{Resp}\Rsign Man looketh on the outward appearance, but God looketh on the heart.\end{Resp}
\switchcolumn*

\end{paracol}

\Office{Dominica V Augusti}{Fifth Sunday of August}{}
\addcontentsline{toc}{subsection}{Dominica V Augusti \textit{— Fifth Sunday of August}}
\begin{paracol}{2}
\LA
\LessonHead{Lectio I}
\switchcolumn
\EN
\LessonHead{Lesson I}
\switchcolumn*

\LA
\LessonTitle{De libro Ecclesiástici}
\switchcolumn
\EN
\LessonTitle{Lesson from the book of Ecclesiasticus}
\switchcolumn*

\LA
\Cite{Sir 5:1-5}
\switchcolumn
\EN
\Cite{Sir 5:1-5}
\switchcolumn*

\LA
\noindent Noli atténdere ad possessiónes iníquas, et ne díxeris: Est mihi suffíciens vita; nihil enim próderit in témpore vindíctæ et obductiónis. Ne sequáris in fortitúdine tua concupiscéntiam cordis tui, et ne díxeris: Quómodo pótui? aut quis me subíciet propter facta mea? Deus enim víndicans vindicábit. Ne díxeris: Peccávi, et quid mihi áccidit triste? Altíssimus enim est pátiens rédditor. De propitiáto peccáto noli esse sine metu, neque adícias peccátum super peccátum.\par
\switchcolumn
\EN
\noindent Set not thy heart upon unjust possessions, and say not: I have enough to live on: for it shall be of no service in the time of vengeance and darkness. Follow not in thy strength the desires of thy heart: And say not: How mighty am I? and who shall bring me under for my deeds? for God will surely take revenge. Say not: I have sinned, and what harm hath befallen me? for the most High is a patient rewarder. Be not without fear about sin forgiven, and add not sin upon sin:\par
\switchcolumn*

\LA
\begin{Resp}\Rsign In princípio Deus ántequam terram fáceret, priúsquam abýssos constitúeret, priúsquam prodúceret fontes aquárum, \Star{} Antequam montes collocaréntur, ante omnes colles generávit me Dóminus.\end{Resp}
\begin{Resp}\Vsign Quando præparábat cælos, áderam, cum eo cuncta compónens.\end{Resp}
\begin{Resp}\Rsign Antequam montes collocaréntur, ante omnes colles generávit me Dóminus.\end{Resp}
\switchcolumn
\EN
\begin{Resp}\Rsign God possessed me in the beginning, before He made the earth, before He created the depths, before He caused the fountains of water to spring. \Star{} Before the mountains were settled, before there were any hills, did the Lord beget me.\end{Resp}
\begin{Resp}\Vsign When He prepared the heavens, I was there with Him, ordering all things.\end{Resp}
\begin{Resp}\Rsign Before the mountains were settled, before there were any hills, did the Lord beget me.\end{Resp}
\switchcolumn*

\LA
\LessonHead{Lectio II}
\switchcolumn
\EN
\LessonHead{Lesson II}
\switchcolumn*

\LA
\Cite{Sir 5:6-11}
\switchcolumn
\EN
\Cite{Sir 5:6-11}
\switchcolumn*

\LA
\noindent Et ne dicas: Miserátio Dómini magna est, multitúdinis peccatórum meórum miserébitur; misericórdia enim et ira ab illo cito próximant, et in peccatóres réspicit ira illíus. Non tardes convérti ad Dóminum et ne dífferas de die in diem; súbito enim véniet ira illíus et in témpore vindíctæ dispérdet te. Noli ánxius esse in divítiis injústis; non enim próderunt tibi in die obductiónis et vindíctæ. Non véntiles te in omnem ventum et non eas in omnem viam; sic enim omnis peccátor probátur in dúplici lingua.\par
\switchcolumn
\EN
\noindent And say not: The mercy of the Lord is great, he will have mercy on the multitude of my sins. For mercy and wrath quickly come from him, and his wrath looketh upon sinners. Delay not to be converted to the Lord, and defer it not from day to day. For his wrath shall come on a sudden, and in the time of vengeance he will destroy thee. Be not anxious for goods unjustly gotten: for they shall not profit thee in the day of calamity and revenge. Winnow not with every wind, and go not into every way: for so is every sinner proved by a double tongue.\par
\switchcolumn*

\LA
\begin{Resp}\Rsign Gyrum cæli circuívi sola, et in flúctibus maris ambulávi, in omni gente et in omni pópulo primátum ténui: \Star{} Superbórum et sublímium colla própria virtúte calcávi.\end{Resp}
\begin{Resp}\Vsign Ego in altíssimis hábito, et thronus meus in colúmna nubis.\end{Resp}
\begin{Resp}\Rsign Superbórum et sublímium colla própria virtúte calcávi.\end{Resp}
\switchcolumn
\EN
\begin{Resp}\Rsign I alone compassed the circuit of heaven, and walked on the waves of the sea. In every nation and in every people, I held the first place. \Star{} In the greatness of my strength have I trodden under my feet the necks of such as be haughty and proud.\end{Resp}
\begin{Resp}\Vsign I dwell in the highest places, and my throne is in a cloudy pillar.\end{Resp}
\begin{Resp}\Rsign In the greatness of my strength have I trodden under my feet the necks of such as be haughty and proud.\end{Resp}
\switchcolumn*

\LA
\LessonHead{Lectio III}
\switchcolumn
\EN
\LessonHead{Lesson III}
\switchcolumn*

\LA
\Cite{Sir 5:12-16}
\switchcolumn
\EN
\Cite{Sir 5:12-16}
\switchcolumn*

\LA
\noindent Esto firmus in via Dómini, et in veritáte sensus tui, et sciéntia; et prosequátur te verbum pacis et justítiæ. Esto mansuétus ad audiéndum verbum, ut intéllegas, et cum sapiéntia próferas respónsum verum. Si est tibi intelléctus, respónde próximo; sin autem, sit manus tua super os tuum, ne capiáris in verbo indisciplináto et confundáris. Honor et glória in sermóne sensáti; lingua vero imprudéntis subvérsio est ipsíus. Non appelléris susúrro, et lingua tua ne capiáris et confundáris.\par
\switchcolumn
\EN
\noindent Be steadfast in the way of the Lord, and in the truth of thy judgment, and in knowledge, and let the word of peace and justice keep with thee. Be meek to hear the word, that thou mayst understand: and return a true answer with wisdom. If thou have understanding, answer thy neighbour: but if not, let thy hand be upon thy mouth, lest thou be surprised in an unskillful word, and be confounded. Honour and glory is in the word of the wise, but the tongue of the fool is his ruin. Be not called a whisperer, and be not taken in thy tongue, and confounded.\par
\switchcolumn*

\LA
\begin{Resp}\Rsign Emítte, Dómine, sapiéntiam de sede magnitúdinis tuæ, ut mecum sit et mecum labóret: \Star{} Ut sciam, quid accéptum sit coram te omni témpore.\end{Resp}
\begin{Resp}\Vsign Da mihi, Dómine, sédium tuárum assistrícem sapiéntiam.\end{Resp}
\begin{Resp}\Rsign Ut sciam quid accéptum sit coram te omni témpore.\end{Resp}
\begin{Resp}\Vsign Glória Patri, et Fílio, \Star{} et Spirítui Sancto.\end{Resp}
\begin{Resp}\Rsign Ut sciam quid accéptum sit coram te omni témpore.\end{Resp}
\switchcolumn
\EN
\begin{Resp}\Rsign O send out wisdom from the throne of thy glory, O Lord, to be with me, and to labour with me, \Star{} That I may know at all times what is pleasing unto thee.\end{Resp}
\begin{Resp}\Vsign Give me wisdom, O Lord, that sitteth by thy throne.\end{Resp}
\begin{Resp}\Rsign That I may know at all times what is pleasing unto thee.\end{Resp}
\begin{Resp}\Vsign Glory be to the Father, and to the Son, \Star{} and to the Holy Ghost.\end{Resp}
\begin{Resp}\Rsign That I may know at all times what is pleasing unto thee.\end{Resp}
\switchcolumn*

\LA
\LessonHead{Lectio IV}
\switchcolumn
\EN
\LessonHead{Lesson IV}
\switchcolumn*

\LA
\LessonTitle{Sermo sancti Joánnis Chrysóstomi}
\switchcolumn
\EN
\LessonTitle{From the Sermons of St. John Chrysostom, Patriarch of Constantinople.}
\switchcolumn*

\LA
\Source{Homilia 22 in 2 Cor. 10 in Moralibus.}
\switchcolumn
\EN
\Source{Homilia 22}
\switchcolumn*

\LA
\noindent Ne tardes convérti ad Dóminum, et ne dífferas de die in diem; nescis enim quid paritúra sit superventúra dies. Perículum enim et metus est in defferéndo; salus vero certa ac secúra, si nulla sit dilátio. Virtútem ígitur cole: sic enim, licet júvenis moriáris, secúre discésseris: quod si ad senectútem pervéneris, cum multa facilitáte et nulla moléstia e vita discédes; duplicémque habébis festivitátem, et quod a vitæ malítia abstinúeris, et quod virtútem colúeris. Ne dicas: Erit tempus, quando convérti licébit: verba enim hæc Deum valde exásperant.\par
\switchcolumn
\EN
\noindent Tarry not to turn thee to the Lord, yea, defer not thy repentance from day to day for thou knowest not what the morrow may bring forth. In delay there is danger and terror, but where there is no delay, there health is safe and secure. Live well then, and, then, however young thou diest, thou wilt die safely; and if thou come to old age thou wilt depart without vexation or trouble; and thou wilt have a double happiness, in that thou wilt be leaving all the evils of life, and in that thou hast lived well. Say not: There will be a time meet for repentance for such words as these do greatly rouse the anger of God.\par
\switchcolumn*

\LA
\begin{Resp}\Rsign Da mihi, Dómine, sédium tuárum assistrícem sapiéntiam, et noli me reprobáre a púeris tuis: \Star{} Quóniam servus tuus sum ego, et fílius ancíllæ tuæ.\end{Resp}
\begin{Resp}\Vsign Mitte illam de sede magnitúdinis tuæ, ut mecum sit et mecum labóret.\end{Resp}
\begin{Resp}\Rsign Quóniam servus tuus sum ego, et fílius ancíllæ tuæ.\end{Resp}
\switchcolumn
\EN
\begin{Resp}\Rsign Give me wisdom, O Lord, that sitteth by thy throne, and reject me not from among thy children. \Star{} For I am thy servant and son of thine handmaid.\end{Resp}
\begin{Resp}\Vsign O send her out from the throne of thy glory, to be with me and to labour with me.\end{Resp}
\begin{Resp}\Rsign For I am thy servant and son of thine handmaid.\end{Resp}
\switchcolumn*

\LA
\LessonHead{Lectio V}
\switchcolumn
\EN
\LessonHead{Lesson V}
\switchcolumn*

\LA
\noindent Cur nam, cum ipse tibi ætérna sǽcula promísit, tu in præsénti vita laboráre non vis, quæ parva et momentánea est; sed sic ignávus ac dissolútus agis, quasi hac breviórem áliam quamdam inquíras? Nonne illæ quotidiánæ comessatiónes, nonne illæ mensæ, nonne scorta illa, nonne theátra illa, nonne divítiæ illæ testántur inexplébilem malítiæ concupiscéntiam? Cógita bene, quod quóties scortátus es, tóties condemnásti teípsum; peccátum enim ita se habet, ut mox atque patrátum fúerit, senténtiam ferat judex.\par
\switchcolumn
\EN
\noindent He hath promised thee eternal ages, and thou willest not to work in this present life, which is so short and so fleeting. Dost thou so idly and loosely carry thyself, as though the life for which thou seekest were a shorter life than this Do not daily feastings, daily gluttonies, daily uncleanness, shows, and riches bear witness to the undying nature of sinful cravings Think it well over, that as often as thou dost commit uncleanness, thou dost damn thyself for this is the nature of sin, as soon as it is committed, the Judge's sentence is uttered.\par
\switchcolumn*

\LA
\begin{Resp}\Rsign Inítium sapiéntiæ timor Dómini: \Star{} Intelléctus bonus ómnibus faciéntibus eum: laudátio ejus manet in sǽculum sǽculi.\end{Resp}
\begin{Resp}\Vsign Diléctio illíus custódia legum est: quia omnis sapiéntia timor Dómini.\end{Resp}
\begin{Resp}\Rsign Intelléctus bonus ómnibus faciéntibus eum: laudátio ejus manet in sǽculum sǽculi.\end{Resp}
\switchcolumn
\EN
\begin{Resp}\Rsign The fear of the Lord is the beginning of wisdom. \Star{} A good understanding have all they that do His commandments. His praise endureth for ever.\end{Resp}
\begin{Resp}\Vsign Love is the keeping of her laws, for all wisdom is the fear of the Lord.\end{Resp}
\begin{Resp}\Rsign A good understanding have all they that do His commandments. His praise endureth for ever.\end{Resp}
\switchcolumn*

\LA
\LessonHead{Lectio VI}
\switchcolumn
\EN
\LessonHead{Lesson VI}
\switchcolumn*

\LA
\noindent Inebriátus es? ventri indulsísti? rapuísti? siste jam gradum: verte te in divérsum: confitére Deo grátiam, quod non in médiis peccátis te ábstulit: ne quære áliud tibi prorogári tempus ut male operéris. Multi dum male ac vitióse víverent, súbito periérunt et in maniféstam damnatiónem abiérunt: time ne idem tibi áccidat. Sed multis, inquis, dedit Deus spátium, ut in última senécta confiteréntur. Quid ígitur? numquid et tibi dábitur? Fortásse dabit, inquis. Cur dicis, fortásse? Cóntigit aliquóties? Cógita quod de ánima delíberas: proínde étiam de contrário cógita, et dic: Quid autem, si non det? Quid autem, si det, inquis? Esto, dat quidem ipse: verúmtamen hoc illo cértius et utílius.\par
\switchcolumn
\EN
\noindent Hast thou been drunken Hast thou over-eaten thyself? Hast thou stolen Stop, and turn back thank the goodness of God, that He hath not taken thee away in the midst of thy sins seek not more time wherein to commit iniquity. Many have they been who have perished suddenly, in the midst of bad and vicious lives, and have gone away to manifest damnation have a fear lest the same thing befall thee. But, thou sayest, they have been many to whom God hath given time, and they have been to confession in their old age. What then Is that a proof that it will be given to thee Perchance, sayest thou. Why sayest thou, Perchance It befalleth sometimes. Bethink thee that it is of thy soul thou art considering. Look at it the other way, and say What if it be not given But, sayest thou, and what if it be May it be so it is true it is among His gifts but nevertheless, this is the safer and the better way.\par
\switchcolumn*

\LA
\begin{Resp}\Rsign Verbum iníquum et dolósum longe fac a me, Dómine: \Star{} Divítias et paupertátem ne déderis mihi, sed tantum víctui meo tríbue necessária.\end{Resp}
\begin{Resp}\Vsign Duo rogávi te, ne déneges mihi ántequam móriar.\end{Resp}
\begin{Resp}\Rsign Divítias et paupertátem ne déderis mihi, sed tantum víctui meo tríbue necessária.\end{Resp}
\begin{Resp}\Vsign Glória Patri, et Fílio, \Star{} et Spirítui Sancto.\end{Resp}
\begin{Resp}\Rsign Divítias et paupertátem ne déderis mihi, sed tantum víctui meo tríbue necessária.\end{Resp}
\switchcolumn
\EN
\begin{Resp}\Rsign Lord, remove far from me vanity and lies. \Star{} Give me neither poverty nor riches, but feed me with food convenient for me.\end{Resp}
\begin{Resp}\Vsign Two things have I required of thee; deny me them not before I die.\end{Resp}
\begin{Resp}\Rsign Give me neither poverty nor riches, but feed me with food convenient for me.\end{Resp}
\begin{Resp}\Vsign Glory be to the Father, and to the Son, \Star{} and to the Holy Ghost.\end{Resp}
\begin{Resp}\Rsign Give me neither poverty nor riches, but feed me with food convenient for me.\end{Resp}
\switchcolumn*

\end{paracol}

\Office{Feria II infra Hebdomadam V Augusti}{Monday of the fifth week of August}{}
\addcontentsline{toc}{subsection}{Feria II infra Hebdomadam V Augusti \textit{— Monday of the fifth week of August}}
\begin{paracol}{2}
\LA
\LessonHead{Lectio I}
\switchcolumn
\EN
\LessonHead{Lesson I}
\switchcolumn*

\LA
\LessonTitle{De libro Ecclesiástici}
\switchcolumn
\EN
\LessonTitle{Lesson from the book of Ecclesiasticus}
\switchcolumn*

\LA
\Cite{Sir 7:1-5}
\switchcolumn
\EN
\Cite{Sir 7:1-5}
\switchcolumn*

\LA
\noindent Noli fácere mala, et non te apprehéndent: discéde ab iníquo, et defícient mala abs te. Fili, non sémines mala in sulcis injustítiæ et non metes ea in séptuplum. Noli quǽrere a dómino ducátum, neque a rege cáthedram honóris. Non te justífices ante Deum, quóniam ágnitor cordis ipse est: et penes regem noli velle vidéri sápiens.\par
\switchcolumn
\EN
\noindent Do no evils, and no evils shall lay hold of thee. Depart from the unjust, and evils shall depart from thee. My son, sow not evils in the furrows of injustice, and thou shalt not reap them sevenfold. Seek not of the Lord a pre-eminence, nor of the king the seat of honour. Justify not thyself before God, for he knoweth the heart: and desire not to appear wise before the king.\par
\switchcolumn*

\LA
\begin{Resp}\Rsign Ne derelínquas me, Dómine, pater et dominátor vitæ meæ, ut non córruam in conspéctu adversariórum meórum: \Star{} Ne gáudeat de me inimícus meus.\end{Resp}
\begin{Resp}\Vsign Apprehénde arma et scutum et exsúrge in adjutórium mihi.\end{Resp}
\begin{Resp}\Rsign Ne gáudeat de me inimícus meus.\end{Resp}
\switchcolumn
\EN
\begin{Resp}\Rsign O Lord, Father and Governor of my life, leave me not, lest I fall before mine adversaries, \Star{} and mine enemy rejoice over me.\end{Resp}
\begin{Resp}\Vsign Take hold of shield and buckler, and stand up for mine help.\end{Resp}
\begin{Resp}\Rsign Lest mine enemy rejoice over me.\end{Resp}
\switchcolumn*

\LA
\LessonHead{Lectio II}
\switchcolumn
\EN
\LessonHead{Lesson II}
\switchcolumn*

\LA
\Cite{Sir 7:6-10}
\switchcolumn
\EN
\Cite{Sir 7:6-10}
\switchcolumn*

\LA
\noindent Noli quǽrere fíeri judex, nisi váleas virtúte irrúmpere iniquitátes; ne forte extiméscas fáciem poténtis et ponas scándalum in æquitáte tua. Non pecces in multitúdinem civitátis nec te immíttas in pópulum: neque álliges duplícia peccáta, nec enim in uno eris immúnis. Noli esse pusillánimis in ánimo tuo, exoráre et fácere eleemósynam ne despícias.\par
\switchcolumn
\EN
\noindent Seek not to be made a judge, unless thou have strength enough to extirpate iniquities: lest thou fear the person of the powerful, and lay a stumblingblock for thy integrity. Offend not against the multitude of a city, neither cast thyself in upon the people, Nor bind sin to sin: for even in one thou shalt not be unpunished. Be not fainthearted in thy mind: Neglect not to pray, and to give alms.\par
\switchcolumn*

\LA
\begin{Resp}\Rsign Magna enim sunt judícia tua, Dómine, et inenarrabília verba tua: \Star{} Magnificásti pópulum tuum et honorásti.\end{Resp}
\begin{Resp}\Vsign Transtulísti illos per Mare Rubrum et transvexísti eos per aquam nímiam.\end{Resp}
\begin{Resp}\Rsign Magnificásti pópulum tuum et honorásti.\end{Resp}
\switchcolumn
\EN
\begin{Resp}\Rsign Great are thy judgments, O Lord, and thy words cannot be expressed. \Star{} Thou didst make thy people mighty and honourable.\end{Resp}
\begin{Resp}\Vsign Thou broughtest them through the Red Sea, and leddest them through much water.\end{Resp}
\begin{Resp}\Rsign Thou didst make thy people mighty and honourable.\end{Resp}
\switchcolumn*

\LA
\LessonHead{Lectio III}
\switchcolumn
\EN
\LessonHead{Lesson III}
\switchcolumn*

\LA
\Cite{Sir 7:11-15}
\switchcolumn
\EN
\Cite{Sir 7:11-15}
\switchcolumn*

\LA
\noindent Ne dicas: In multitúdine múnerum meórum respíciet Deus et, offerénte me Deo altíssimo, múnera mea suscípiet. Non irrídeas hóminem in amaritúdine ánimæ; est enim qui humíliat et exáltat, circumspéctor Deus. Noli aráre mendácium advérsus fratrem tuum, neque in amícum simíliter fácias. Noli velle mentíri omne mendácium; assidúitas enim illíus non est bona. Noli verbósus esse in multitúdine presbyterórum, et non íteres verbum in oratióne tua.\par
\switchcolumn
\EN
\noindent Say not: God will have respect to the multitude of my gifts, and when I offer to the most high God, he will accept my offerings. Laugh no man to scorn in the bitterness of his soul: for there is one that humbleth and exalteth, God who seeth all. Devise not a lie against thy brother: neither do the like against thy friend. Be not willing to make any manner of lie: for the custom thereof is not good. Be not full of words in a multitude of ancients, and repeat not the word in thy prayer.\par
\switchcolumn*

\LA
\begin{Resp}\Rsign Quæ sunt in corde hóminum, óculi tui vident, Dómine, et in libro tuo ómnia scribéntur: \Star{} Homo videt in fácie, Deus autem in corde.\end{Resp}
\begin{Resp}\Vsign Omnia enim corda scrutátur, et univérsas méntium cogitatiónes intéllegit.\end{Resp}
\begin{Resp}\Rsign Homo videt in fácie, Deus autem in corde.\end{Resp}
\begin{Resp}\Vsign Glória Patri, et Fílio, \Star{} et Spirítui Sancto.\end{Resp}
\begin{Resp}\Rsign Homo videt in fácie, Deus autem in corde.\end{Resp}
\switchcolumn
\EN
\begin{Resp}\Rsign Lord, thine eyes behold all that is in the heart of man, and in thy book are they all written. \Star{} Man looketh on the outward appearance, but God looketh on the heart.\end{Resp}
\begin{Resp}\Vsign For He searcheth all hearts, and understandeth all the imaginations of the thoughts.\end{Resp}
\begin{Resp}\Rsign Man looketh on the outward appearance, but God looketh on the heart.\end{Resp}
\begin{Resp}\Vsign Glory be to the Father, and to the Son, \Star{} and to the Holy Ghost.\end{Resp}
\begin{Resp}\Rsign Man looketh on the outward appearance, but God looketh on the heart.\end{Resp}
\switchcolumn*

\end{paracol}

\Office{Feria III infra Hebdomadam V Augusti}{Tuesday of the fifth week of August}{}
\addcontentsline{toc}{subsection}{Feria III infra Hebdomadam V Augusti \textit{— Tuesday of the fifth week of August}}
\begin{paracol}{2}
\LA
\LessonHead{Lectio I}
\switchcolumn
\EN
\LessonHead{Lesson I}
\switchcolumn*

\LA
\LessonTitle{De libro Ecclesiástici}
\switchcolumn
\EN
\LessonTitle{Lesson from the book of Ecclesiasticus}
\switchcolumn*

\LA
\Cite{Sir 10:1-5}
\switchcolumn
\EN
\Cite{Sir 10:1-5}
\switchcolumn*

\LA
\noindent Judex sápiens judicábit pópulum suum, et principátus sensáti stábilis erit. Secúndum júdicem pópuli sic et minístri ejus, et qualis rector est civitátis, tales et inhabitántes in ea. Rex insípiens perdet pópulum suum, et civitátes inhabitabúntur per sensum poténtium. In manu Dei potéstas terræ, et útilem rectórem suscitábit in tempus super illam. In manu Dei prospéritas hóminis, et super fáciem scribæ impónet honórem suum.\par
\switchcolumn
\EN
\noindent A wise judge shall judge his people, and the government of a prudent man shall be steady. As the judge of the people is himself, so also are his ministers: and what manner of man the ruler of a city is, such also are they that dwell therein. An unwise king shall be the ruin of his people: and cities shall be inhabited through the prudence of the rulers. The power of the earth is in the hand of God, and in his time he will raise up a profitable ruler over it. The prosperity of man is in the hand of God, and upon the person of the scribe he shall lay his honour.\par
\switchcolumn*

\LA
\begin{Resp}\Rsign Præbe, fili, cor mihi, et óculi tui vias meas custódiant: \Star{} Ut addátur grátia cápiti tuo.\end{Resp}
\begin{Resp}\Vsign Atténde, fili mi, sapiéntiam meam et ad elóquium meum inclína aurem tuam.\end{Resp}
\begin{Resp}\Rsign Ut addátur grátia cápiti tuo.\end{Resp}
\switchcolumn
\EN
\begin{Resp}\Rsign My son, give me thine heart, and let thine eyes observe my ways. \Star{} For they shall be an ornament of grace unto thine head.\end{Resp}
\begin{Resp}\Vsign My son, attend unto my wisdom, and incline thine ear unto my sayings.\end{Resp}
\begin{Resp}\Rsign For they shall be an ornament of grace unto thine head.\end{Resp}
\switchcolumn*

\LA
\LessonHead{Lectio II}
\switchcolumn
\EN
\LessonHead{Lesson II}
\switchcolumn*

\LA
\Cite{Sir 10:6-10}
\switchcolumn
\EN
\Cite{Sir 10:6-10}
\switchcolumn*

\LA
\noindent Omnis injúriæ próximi ne memíneris et nihil agas in opéribus injúriæ. Odíbilis coram Deo est et homínibus supérbia, et exsecrábilis omnis iníquitas géntium. Regnum a gente in gentem transfértur propter injustítias et injúrias et contumélias et divérsos dolos. Aváro autem nihil est sceléstius. Quid supérbit terra et cinis? Nihil est iníquius quam amáre pecúniam; hic enim et ánimam suam venálem habet, quóniam in vita sua projécit íntima sua.\par
\switchcolumn
\EN
\noindent Pride is hateful before God and men: and all iniquity of nations is execrable. A kingdom is translated from one people to another, because of injustices, and wrongs, and injuries, and diverse deceits. But nothing is more wicked than the covetous man. Why is earth and ashes proud? There is not a more wicked thing than to love money: for such a one setteth even his own soul to sale: because while he liveth he hath cast away his bowels.\par
Remember not any injury done thee by thy neighbour, and do thou nothing by deeds of injury.\par
\switchcolumn*

\LA
\begin{Resp}\Rsign Inítium sapiéntiæ timor Dómini: \Star{} Intelléctus bonus ómnibus faciéntibus eum: laudátio ejus manet in sǽculum sǽculi.\end{Resp}
\begin{Resp}\Vsign Diléctio illíus custódia legum est: quia omnis sapiéntia timor Dómini.\end{Resp}
\begin{Resp}\Rsign Intelléctus bonus ómnibus faciéntibus eum: laudátio ejus manet in sǽculum sǽculi.\end{Resp}
\switchcolumn
\EN
\begin{Resp}\Rsign The fear of the Lord is the beginning of wisdom. \Star{} A good understanding have all they that do His commandments. His praise endureth for ever.\end{Resp}
\begin{Resp}\Vsign Love is the keeping of her laws, for all wisdom is the fear of the Lord.\end{Resp}
\begin{Resp}\Rsign A good understanding have all they that do His commandments. His praise endureth for ever.\end{Resp}
\switchcolumn*

\LA
\LessonHead{Lectio III}
\switchcolumn
\EN
\LessonHead{Lesson III}
\switchcolumn*

\LA
\Cite{Sir 10:11-16}
\switchcolumn
\EN
\Cite{Sir 10:11-16}
\switchcolumn*

\LA
\noindent Omnis potentátus brevis vita: languor prolíxior gravat médicum, brevem languórem præcídit médicus; sic et rex hódie est, et cras moriétur: cum enim moriétur homo, hereditábit serpéntes, et béstias, et vermes. Inítium supérbiæ hóminis apostatáre a Deo; quóniam ab eo qui fecit illum, recéssit cor ejus, quóniam inítium omnis peccáti est supérbia. Qui tenúerit illam adimplébitur maledíctis, et subvértet eum in finem. Proptérea exhonorávit Dóminus convéntus malórum, et destrúxit eos usque in finem.\par
\switchcolumn
\EN
\noindent All power is of short life. A long sickness is troublesome to the physician. The physician cutteth off it short sickness: so also a king is today, and tomorrow he shall die. For when a man shall die, he shall inherit serpents, and beasts, and worms. The beginning of the pride of man, is to fall off from God: Because his heart is departed from him that made him: for pride is the beginning of all sin: be that holdeth it, shall be filled with maledictions, and it shall ruin him in the end. Therefore hath the Lord disgraced the assemblies of the wicked, and hath utterly destroyed them.\par
\switchcolumn*

\LA
\begin{Resp}\Rsign Verbum iníquum et dolósum longe fac a me, Dómine: \Star{} Divítias et paupertátem ne déderis mihi, sed tantum víctui meo tríbue necessária.\end{Resp}
\begin{Resp}\Vsign Duo rogávi te, ne déneges mihi ántequam móriar.\end{Resp}
\begin{Resp}\Rsign Divítias et paupertátem ne déderis mihi, sed tantum víctui meo tríbue necessária.\end{Resp}
\begin{Resp}\Vsign Glória Patri, et Fílio, \Star{} et Spirítui Sancto.\end{Resp}
\begin{Resp}\Rsign Divítias et paupertátem ne déderis mihi, sed tantum víctui meo tríbue necessária.\end{Resp}
\switchcolumn
\EN
\begin{Resp}\Rsign Lord, remove far from me vanity and lies. \Star{} Give me neither poverty nor riches, but feed me with food convenient for me.\end{Resp}
\begin{Resp}\Vsign Two things have I required of thee; deny me them not before I die.\end{Resp}
\begin{Resp}\Rsign Give me neither poverty nor riches, but feed me with food convenient for me.\end{Resp}
\begin{Resp}\Vsign Glory be to the Father, and to the Son, \Star{} and to the Holy Ghost.\end{Resp}
\begin{Resp}\Rsign Give me neither poverty nor riches, but feed me with food convenient for me.\end{Resp}
\switchcolumn*

\end{paracol}

\Office{Feria IV infra Hebdomadam V Augusti}{Wednesday of the fifth week of August}{}
\addcontentsline{toc}{subsection}{Feria IV infra Hebdomadam V Augusti \textit{— Wednesday of the fifth week of August}}
\begin{paracol}{2}
\LA
\LessonHead{Lectio I}
\switchcolumn
\EN
\LessonHead{Lesson I}
\switchcolumn*

\LA
\LessonTitle{De libro Ecclesiástici}
\switchcolumn
\EN
\LessonTitle{Lesson from the book of Ecclesiasticus}
\switchcolumn*

\LA
\Cite{Sir 13:1-6}
\switchcolumn
\EN
\Cite{Sir 13:1-6}
\switchcolumn*

\LA
\noindent Qui tetígerit picem, inquinábitur ab ea; et, qui communicáverit supérbo, índuet supérbiam. Pondus super se tollet qui honestióri se commúnicat, et ditióri te ne sócius fúeris. Quid communicábit cácabus ad ollam? quando enim se collíserint, confringétur. Dives injúste egit et fremet, pauper autem læsus tacébit. Si largítus fúeris, assúmet te, et, si non habúeris, derelínquet te. Si habes, convívet tecum, et evacuábit te: et ipse non dolébit super te.\par
\switchcolumn
\EN
\noindent He that toucheth pitch, shall be defiled with it: and he that hath fellowship with the proud, shall put on pride. He shall take a burden upon him that hath fellowship with one more honourable than himself. And have no fellowship with one that is richer than thyself. What agreement shall the earthen pot have with the kettle? for if they knock one against the other, it shall be broken. The rich man hath done wrong, and yet he will fume: but the poor is wronged and must hold his peace. If thou give, he will make use of thee: and if thou have nothing, he will forsake thee. If thou have any thing, he will live with thee, and will make thee bare, and he will not be sorry for thee.\par
\switchcolumn*

\LA
\begin{Resp}\Rsign Dómine, Pater et Deus vitæ meæ, ne derelínquas me in cogitátu malígno: extolléntiam oculórum meórum ne déderis mihi, et desidérium malígnum avérte a me, Dómine; aufer a me concupiscéntiam, \Star{} Et ánimo irreverénti et infruníto ne tradas me, Dómine.\end{Resp}
\begin{Resp}\Vsign Ne derelínquas me, Dómine, ne accréscant ignorántiæ meæ, nec multiplicéntur delícta mea.\end{Resp}
\begin{Resp}\Rsign Et ánimo irreverénti et infruníto ne tradas me, Dómine.\end{Resp}
\switchcolumn
\EN
\begin{Resp}\Rsign O Lord, Father and God of my life, leave me not to evil counsels; give me not a proud look, but turn away from me an haughty mind, O Lord turn away from me concupiscence, \Star{} And give me not over unto an impudent and froward mind, O Lord!\end{Resp}
\begin{Resp}\Vsign Leave me not, O Lord, lest mine ignorance increase, and my sins abound.\end{Resp}
\begin{Resp}\Rsign And give me not over unto an impudent and froward mind, O Lord.\end{Resp}
\switchcolumn*

\LA
\LessonHead{Lectio II}
\switchcolumn
\EN
\LessonHead{Lesson II}
\switchcolumn*

\LA
\Cite{Sir 13:9-15}
\switchcolumn
\EN
\Cite{Sir 13:9-15}
\switchcolumn*

\LA
\noindent Humiliáre Deo et exspécta manus ejus; atténde, ne sedúctus in stultítiam humiliéris. Noli esse húmilis in sapiéntia tua, ne humiliátus in stultítiam seducáris. Advocátus a potentióre discéde, ex hoc enim magis te advocábit. Ne ímprobus sis, ne impingáris, et ne longe sis ab eo, ne eas in obliviónem. Ne retíneas ex æquo loqui cum illo nec credas multis verbis illíus; ex multa enim loquéla tentábit te, et subrídens interrogábit te de abscónditis tuis. Immítis ánimus illíus conservábit verba tua: et non parcet de malítia, et de vínculis.\par
\switchcolumn
\EN
\noindent Humble thyself to God, and wait for his hands. Beware that thou be not deceived Into folly, and be humbled. Be not lowly in thy wisdom, lest being humbled thou be deceived into folly. If thou be invited by one that is mightier, withdraw thyself: for so he will invite thee the more. Be not troublesome to him, lest thou be put back: and keep not far from him, lest thou be forgotten. Affect not to speak with him as an equal: and believe not his many words: for by much talk he will sift thee, and smiling will examine thee concerning thy secrets. His cruel mind will lay up thy words: and he will not spare to do thee hurt, and to cast thee into prison.\par
\switchcolumn*

\LA
\begin{Resp}\Rsign Magna enim sunt judícia tua, Dómine, et inenarrabília verba tua: \Star{} Magnificásti pópulum tuum et honorásti.\end{Resp}
\begin{Resp}\Vsign Transtulísti illos per Mare Rubrum et transvexísti eos per aquam nímiam.\end{Resp}
\begin{Resp}\Rsign Magnificásti pópulum tuum et honorásti.\end{Resp}
\switchcolumn
\EN
\begin{Resp}\Rsign Great are thy judgments, O Lord, and thy words cannot be expressed. \Star{} Thou didst make thy people mighty and honourable.\end{Resp}
\begin{Resp}\Vsign Thou broughtest them through the Red Sea, and leddest them through much water.\end{Resp}
\begin{Resp}\Rsign Thou didst make thy people mighty and honourable.\end{Resp}
\switchcolumn*

\LA
\LessonHead{Lectio III}
\switchcolumn
\EN
\LessonHead{Lesson III}
\switchcolumn*

\LA
\Cite{Sir 13:16-22}
\switchcolumn
\EN
\Cite{Sir 13:16-22}
\switchcolumn*

\LA
\noindent Cave tibi, et atténde diligénter audítui tuo, quóniam cum subversióne tua ámbulas; áudiens vero illa, quasi in somnis vide, et vigilábis. Omni vita tua dílige Deum, et ínvoca illum in salúte tua. Omne ánimal díligit símile sibi, sic et omnis homo próximum sibi; omnis caro ad símilem sibi conjungétur, et omnis homo símili sui sociábitur. Si communicábit lupus agno aliquándo, sic peccátor justo. Quæ communicátio sancto hómini ad canem? aut quæ pars díviti ad páuperem?\par
\switchcolumn
\EN
\noindent Take heed to thyself, and attend diligently to what thou hearest: for thou walkest in danger of thy ruin. When thou hearest those things, see as it were in sleep, and thou shalt awake. Love God all thy life, and call upon him for thy salvation. Every beast loveth its like: so also every man him that is nearest to himself. All flesh shall consort with the like to itself, and every man shall associate himself to his like. If the wolf shall at any time have fellowship with the lamb, so the sinner with the just. What fellowship hath a holy man with a dog, or what part hath the rich with the poor?\par
\switchcolumn*

\LA
\begin{Resp}\Rsign Quæ sunt in corde hóminum, óculi tui vident, Dómine, et in libro tuo ómnia scribéntur: \Star{} Homo videt in fácie, Deus autem in corde.\end{Resp}
\begin{Resp}\Vsign Omnia enim corda scrutátur, et univérsas méntium cogitatiónes intéllegit.\end{Resp}
\begin{Resp}\Rsign Homo videt in fácie, Deus autem in corde.\end{Resp}
\begin{Resp}\Vsign Glória Patri, et Fílio, \Star{} et Spirítui Sancto.\end{Resp}
\begin{Resp}\Rsign Homo videt in fácie, Deus autem in corde.\end{Resp}
\switchcolumn
\EN
\begin{Resp}\Rsign Lord, thine eyes behold all that is in the heart of man, and in thy book are they all written. \Star{} Man looketh on the outward appearance, but God looketh on the heart.\end{Resp}
\begin{Resp}\Vsign For He searcheth all hearts, and understandeth all the imaginations of the thoughts.\end{Resp}
\begin{Resp}\Rsign Man looketh on the outward appearance, but God looketh on the heart.\end{Resp}
\begin{Resp}\Vsign Glory be to the Father, and to the Son, \Star{} and to the Holy Ghost.\end{Resp}
\begin{Resp}\Rsign Man looketh on the outward appearance, but God looketh on the heart.\end{Resp}
\switchcolumn*

\end{paracol}

\Office{Feria V infra Hebdomadam V Augusti}{Thursday of the fifth week of August}{}
\addcontentsline{toc}{subsection}{Feria V infra Hebdomadam V Augusti \textit{— Thursday of the fifth week of August}}
\begin{paracol}{2}
\LA
\LessonHead{Lectio I}
\switchcolumn
\EN
\LessonHead{Lesson I}
\switchcolumn*

\LA
\LessonTitle{De libro Ecclesiástici}
\switchcolumn
\EN
\LessonTitle{Lesson from the book of Ecclesiasticus}
\switchcolumn*

\LA
\Cite{Sir 14:1-5}
\switchcolumn
\EN
\Cite{Sir 14:1-5}
\switchcolumn*

\LA
\noindent Beátus vir qui non est lapsus verbo ex ore suo, et non est stimulátus in tristítia delícti; felix qui non hábuit ánimi sui tristítiam, et non éxcidit a spe sua. Viro cúpido et tenáci sine ratióne est substántia, et hómini lívido ad quid aurum? Qui acérvat ex ánimo suo injúste áliis cóngregat, et in bonis illíus álius luxuriábitur. Qui sibi nequam est, cui álii bonus erit? et non jucundábitur in bonis suis.\par
\switchcolumn
\EN
\noindent Blessed is the man that hath not slipped by a word out of his mouth, and is not pricked with the remorse of sin. Happy is he that hath had no sadness of his mind, and who is not fallen from his hope. Riches are not comely for a covetous man and a niggard, and what should an envious man do with gold? He that gathereth together by wronging his own soul, gathereth for others, and another will squander away his goods in rioting. He that is evil to himself, to whom will he be good? and he shall not take pleasure in his goods.\par
\switchcolumn*

\LA
\begin{Resp}\Rsign In princípio Deus ántequam terram fáceret, priúsquam abýssos constitúeret, priúsquam prodúceret fontes aquárum, \Star{} Antequam montes collocaréntur, ante omnes colles generávit me Dóminus.\end{Resp}
\begin{Resp}\Vsign Quando præparábat cælos, áderam, cum eo cuncta compónens.\end{Resp}
\begin{Resp}\Rsign Antequam montes collocaréntur, ante omnes colles generávit me Dóminus.\end{Resp}
\switchcolumn
\EN
\begin{Resp}\Rsign God possessed me in the beginning, before He made the earth, before He created the depths, before He caused the fountains of water to spring. \Star{} Before the mountains were settled, before there were any hills, did the Lord beget me.\end{Resp}
\begin{Resp}\Vsign When He prepared the heavens, I was there with Him, ordering all things.\end{Resp}
\begin{Resp}\Rsign Before the mountains were settled, before there were any hills, did the Lord beget me.\end{Resp}
\switchcolumn*

\LA
\LessonHead{Lectio II}
\switchcolumn
\EN
\LessonHead{Lesson II}
\switchcolumn*

\LA
\Cite{Sir 14:6-10}
\switchcolumn
\EN
\Cite{Sir 14:6-10}
\switchcolumn*

\LA
\noindent Qui sibi ínvidet, nihil est illo néquius, et hæc reddítio est malítiæ illíus; et, si bene fécerit, ignoránter et non volens facit: et in novíssimo maniféstat malítiam suam. Nequam est óculus lívidi: et avértens fáciem suam, et despíciens ánimam suam. Insatiábilis óculus cúpidi in parte iniquitátis: non satiábitur, donec consúmat arefáciens ánimam suam. Oculus malus ad mala, et non satiábitur pane, sed índigens et in tristítia erit super mensam suam.\par
\switchcolumn
\EN
\noindent There is none worse than he that envieth himself, and this is the reward of his wickedness: And if he do good, he doth it ignorantly, and unwillingly: and at the last he discovereth his wickedness. The eye of the envious is wicked: and he turneth away his face, and despiseth his own soul. The eye of the covetous man is insatiable in his portion of iniquity: he will not be satisfied till he consume his own soul, drying it up. An evil eye is towards evil things: and he shall not have his fill of bread, but shall be needy and pensive at his own table.\par
\switchcolumn*

\LA
\begin{Resp}\Rsign Gyrum cæli circuívi sola, et in flúctibus maris ambulávi, in omni gente et in omni pópulo primátum ténui: \Star{} Superbórum et sublímium colla própria virtúte calcávi.\end{Resp}
\begin{Resp}\Vsign Ego in altíssimis hábito, et thronus meus in colúmna nubis.\end{Resp}
\begin{Resp}\Rsign Superbórum et sublímium colla própria virtúte calcávi.\end{Resp}
\switchcolumn
\EN
\begin{Resp}\Rsign I alone compassed the circuit of heaven, and walked on the waves of the sea. In every nation and in every people, I held the first place. \Star{} In the greatness of my strength have I trodden under my feet the necks of such as be haughty and proud.\end{Resp}
\begin{Resp}\Vsign I dwell in the highest places, and my throne is in a cloudy pillar.\end{Resp}
\begin{Resp}\Rsign In the greatness of my strength have I trodden under my feet the necks of such as be haughty and proud.\end{Resp}
\switchcolumn*

\LA
\LessonHead{Lectio III}
\switchcolumn
\EN
\LessonHead{Lesson III}
\switchcolumn*

\LA
\Cite{Sir 14:11-17}
\switchcolumn
\EN
\Cite{Sir 14:11-17}
\switchcolumn*

\LA
\noindent Fili, si habes, bénefac tecum et Deo dignas oblatiónes offer. Memor esto, quóniam mors non tardat, et testaméntum inferórum, quia demonstrátum est tibi; testaméntum enim hujus mundi morte moriétur. Ante mortem bénefac amíco tuo et secúndum vires tuas expórrigens da páuperi. Non defraudéris a die bono, et partícula boni doni non te prætéreat. Nonne áliis relínques dolóres et labóres tuos in divisióne sortis? Da et áccipe et justífica ánimam tuam: ante óbitum tuum operáre justítiam; quóniam non est apud ínferos inveníre cibum.\par
\switchcolumn
\EN
\noindent My son, if thou have any thing, do good to thyself, and offer to God worthy offerings. Remember that death is not slow, and that the covenant of hell hath been shown to thee: for the covenant of this world shall surely die. Do good to thy friend before thou die, and according to thy ability, stretching out thy hand give to the poor. Defraud not thyself of the good day, and let not the part of a good gift overpass thee. Shalt thou not leave to others to divide by lot thy sorrows and labours? Give and take, and justify thy soul. Before thy death work justice: for in hell there is no finding food.\par
\switchcolumn*

\LA
\begin{Resp}\Rsign Emítte, Dómine, sapiéntiam de sede magnitúdinis tuæ, ut mecum sit et mecum labóret: \Star{} Ut sciam, quid accéptum sit coram te omni témpore.\end{Resp}
\begin{Resp}\Vsign Da mihi, Dómine, sédium tuárum assistrícem sapiéntiam.\end{Resp}
\begin{Resp}\Rsign Ut sciam quid accéptum sit coram te omni témpore.\end{Resp}
\begin{Resp}\Vsign Glória Patri, et Fílio, \Star{} et Spirítui Sancto.\end{Resp}
\begin{Resp}\Rsign Ut sciam quid accéptum sit coram te omni témpore.\end{Resp}
\switchcolumn
\EN
\begin{Resp}\Rsign O send out wisdom from the throne of thy glory, O Lord, to be with me, and to labour with me, \Star{} That I may know at all times what is pleasing unto thee.\end{Resp}
\begin{Resp}\Vsign Give me wisdom, O Lord, that sitteth by thy throne.\end{Resp}
\begin{Resp}\Rsign That I may know at all times what is pleasing unto thee.\end{Resp}
\begin{Resp}\Vsign Glory be to the Father, and to the Son, \Star{} and to the Holy Ghost.\end{Resp}
\begin{Resp}\Rsign That I may know at all times what is pleasing unto thee.\end{Resp}
\switchcolumn*

\end{paracol}

\Office{Feria VI infra Hebdomadam V Augusti}{Friday of the fifth week of August}{}
\addcontentsline{toc}{subsection}{Feria VI infra Hebdomadam V Augusti \textit{— Friday of the fifth week of August}}
\begin{paracol}{2}
\LA
\LessonHead{Lectio I}
\switchcolumn
\EN
\LessonHead{Lesson I}
\switchcolumn*

\LA
\LessonTitle{De libro Ecclesiástici}
\switchcolumn
\EN
\LessonTitle{Lesson from the book of Ecclesiasticus}
\switchcolumn*

\LA
\Cite{Sir 21:1-5}
\switchcolumn
\EN
\Cite{Sir 21:1-5}
\switchcolumn*

\LA
\noindent Fili, peccásti? non adícias íterum, sed et de prístinis deprecáre, ut tibi dimittántur. Quasi a fácie cólubri fuge peccáta et, si accésseris ad illa, suscípient te. Dentes leónis dentes ejus interficiéntes ánimas hóminum, quasi rhomphǽa bis acúta omnis iníquitas, plagæ illíus non est sánitas. Objurgátio et injúriæ annullábunt substántiam, et domus quæ nimis lócuples est, annullábitur supérbia: sic substántia supérbi eradicábitur.\par
\switchcolumn
\EN
\noindent My son, hast thou sinned? do so no more: but for thy former sins also pray that they may be forgiven thee. Flee from sins as from the face of a serpent: for if thou comest near them, they will take hold of thee. The teeth thereof are the teeth of a lion, killing the souls of men. All iniquity is like a two-edged sword, there is no remedy for the wound thereof. Injuries and wrongs will waste riches: and the house that is very rich shall be brought to nothing by pride: so the substance of the proud shall be rooted out.\par
\switchcolumn*

\LA
\begin{Resp}\Rsign Da mihi, Dómine, sédium tuárum assistrícem sapiéntiam, et noli me reprobáre a púeris tuis: \Star{} Quóniam servus tuus sum ego, et fílius ancíllæ tuæ.\end{Resp}
\begin{Resp}\Vsign Mitte illam de sede magnitúdinis tuæ, ut mecum sit et mecum labóret.\end{Resp}
\begin{Resp}\Rsign Quóniam servus tuus sum ego, et fílius ancíllæ tuæ.\end{Resp}
\switchcolumn
\EN
\begin{Resp}\Rsign Give me wisdom, O Lord, that sitteth by thy throne, and reject me not from among thy children. \Star{} For I am thy servant and son of thine handmaid.\end{Resp}
\begin{Resp}\Vsign O send her out from the throne of thy glory, to be with me and to labour with me.\end{Resp}
\begin{Resp}\Rsign For I am thy servant and son of thine handmaid.\end{Resp}
\switchcolumn*

\LA
\LessonHead{Lectio II}
\switchcolumn
\EN
\LessonHead{Lesson II}
\switchcolumn*

\LA
\Cite{Sir 21:6-10}
\switchcolumn
\EN
\Cite{Sir 21:6-10}
\switchcolumn*

\LA
\noindent Deprecátio páuperis ex ore usque ad aures ejus pervéniet, et judícium festináto advéniet illi. Qui odit correptiónem, vestígium est peccatóris; et, qui timet Deum, convertétur ad cor suum. Notus a longe potens lingua audáci, et sensátus scit labi se ab ipso. Qui ædíficat domum suam impéndiis aliénis, quasi qui cólligit lápides suos in híeme. Stuppa collécta synagóga peccántium, et consummátio illórum flamma ignis.\par
\switchcolumn
\EN
\noindent The prayer out of the mouth of the poor shall reach the ears of God, and judgment shall come for him speedily. He that hateth to be reproved walketh in the trace of a sinner: and he that feareth God will turn to his own heart. He that is mighty by a bold tongue is known afar off, but a wise man knoweth to slip by him. He that buildeth his house at other men's charges, is as he that gathereth himself stones to build in the winter. The congregation of sinners is like tow heaped together, and the end of them is a flame of fire.\par
\switchcolumn*

\LA
\begin{Resp}\Rsign Inítium sapiéntiæ timor Dómini: \Star{} Intelléctus bonus ómnibus faciéntibus eum: laudátio ejus manet in sǽculum sǽculi.\end{Resp}
\begin{Resp}\Vsign Diléctio illíus custódia legum est: quia omnis sapiéntia timor Dómini.\end{Resp}
\begin{Resp}\Rsign Intelléctus bonus ómnibus faciéntibus eum: laudátio ejus manet in sǽculum sǽculi.\end{Resp}
\switchcolumn
\EN
\begin{Resp}\Rsign The fear of the Lord is the beginning of wisdom. \Star{} A good understanding have all they that do His commandments. His praise endureth for ever.\end{Resp}
\begin{Resp}\Vsign Love is the keeping of her laws, for all wisdom is the fear of the Lord.\end{Resp}
\begin{Resp}\Rsign A good understanding have all they that do His commandments. His praise endureth for ever.\end{Resp}
\switchcolumn*

\LA
\LessonHead{Lectio III}
\switchcolumn
\EN
\LessonHead{Lesson III}
\switchcolumn*

\LA
\Cite{Sir 21:11-16}
\switchcolumn
\EN
\Cite{Sir 21:11-16}
\switchcolumn*

\LA
\noindent Via peccántium complanáta lapídibus, et in fine illórum ínferi et ténebræ et pœnæ. Qui custódit justítiam, continébit sensum ejus. Consummátio timóris Dei, sapiéntia et sensus. Non erudiétur qui non est sápiens in bono; est autem sapiéntia quæ abúndat in malo, et non est sensus ubi est amaritúdo. Sciéntia sapiéntis tamquam inundátio abundábit, et consílium illíus sicut fons vitæ pérmanet.\par
\switchcolumn
\EN
\noindent The way of sinners is made plain with stones, and in their end is hell, and darkness, and pains. He that keepeth justice shall get the understanding thereof. The perfection of the fear of God is wisdom and understanding. He that is not wise in good, will not be taught. But there is a wisdom that aboundeth in evil: and there is no understanding where there is bitterness. The knowledge of a wise man shall abound like a flood, and his counsel continueth like a fountain of life\par
\switchcolumn*

\LA
\begin{Resp}\Rsign Verbum iníquum et dolósum longe fac a me, Dómine: \Star{} Divítias et paupertátem ne déderis mihi, sed tantum víctui meo tríbue necessária.\end{Resp}
\begin{Resp}\Vsign Duo rogávi te, ne déneges mihi ántequam móriar.\end{Resp}
\begin{Resp}\Rsign Divítias et paupertátem ne déderis mihi, sed tantum víctui meo tríbue necessária.\end{Resp}
\begin{Resp}\Vsign Glória Patri, et Fílio, \Star{} et Spirítui Sancto.\end{Resp}
\begin{Resp}\Rsign Divítias et paupertátem ne déderis mihi, sed tantum víctui meo tríbue necessária.\end{Resp}
\switchcolumn
\EN
\begin{Resp}\Rsign Lord, remove far from me vanity and lies. \Star{} Give me neither poverty nor riches, but feed me with food convenient for me.\end{Resp}
\begin{Resp}\Vsign Two things have I required of thee; deny me them not before I die.\end{Resp}
\begin{Resp}\Rsign Give me neither poverty nor riches, but feed me with food convenient for me.\end{Resp}
\begin{Resp}\Vsign Glory be to the Father, and to the Son, \Star{} and to the Holy Ghost.\end{Resp}
\begin{Resp}\Rsign Give me neither poverty nor riches, but feed me with food convenient for me.\end{Resp}
\switchcolumn*

\end{paracol}

\Office{Sabbato infra Hebdomadam V Augusti}{Saturday of the fifth week of August}{}
\addcontentsline{toc}{subsection}{Sabbato infra Hebdomadam V Augusti \textit{— Saturday of the fifth week of August}}
\begin{paracol}{2}
\LA
\LessonHead{Lectio I}
\switchcolumn
\EN
\LessonHead{Lesson I}
\switchcolumn*

\LA
\LessonTitle{De libro Ecclesiástici}
\switchcolumn
\EN
\LessonTitle{Lesson from the book of Ecclesiasticus}
\switchcolumn*

\LA
\Cite{Sir 32:1-5}
\switchcolumn
\EN
\Cite{Sir 32:1-5}
\switchcolumn*

\LA
\noindent Rectórem te posuérunt? noli extólli: esto in illis quasi unus ex ipsis. Curam illórum habe, et sic consíde, et omni cura tua explícita recúmbe: ut lætéris propter illos et ornaméntum grátiæ accípias corónam, et dignatiónem consequáris corrogatiónis. Lóquere, major natu; decet enim te primum verbum diligénti sciéntia, et non impédias músicam.\par
\switchcolumn
\EN
\noindent Have they made thee ruler? be not lifted up: be among them as one of them. Have care of them, and so sit down, and when thou hast acquitted thyself of all thy charge, take thy place: That thou mayst rejoice for them, and receive a crown as an ornament of grace, and get the honour of the contribution. Speak, thou that art elder: for it becometh thee, To speak the first word with careful knowledge, and hinder not music.\par
\switchcolumn*

\LA
\begin{Resp}\Rsign Dómine, Pater et Deus vitæ meæ, ne derelínquas me in cogitátu malígno: extolléntiam oculórum meórum ne déderis mihi, et desidérium malígnum avérte a me, Dómine; aufer a me concupiscéntiam, \Star{} Et ánimo irreverénti et infruníto ne tradas me, Dómine.\end{Resp}
\begin{Resp}\Vsign Ne derelínquas me, Dómine, ne accréscant ignorántiæ meæ, nec multiplicéntur delícta mea.\end{Resp}
\begin{Resp}\Rsign Et ánimo irreverénti et infruníto ne tradas me, Dómine.\end{Resp}
\switchcolumn
\EN
\begin{Resp}\Rsign O Lord, Father and God of my life, leave me not to evil counsels; give me not a proud look, but turn away from me an haughty mind, O Lord turn away from me concupiscence, \Star{} And give me not over unto an impudent and froward mind, O Lord!\end{Resp}
\begin{Resp}\Vsign Leave me not, O Lord, lest mine ignorance increase, and my sins abound.\end{Resp}
\begin{Resp}\Rsign And give me not over unto an impudent and froward mind, O Lord.\end{Resp}
\switchcolumn*

\LA
\LessonHead{Lectio II}
\switchcolumn
\EN
\LessonHead{Lesson II}
\switchcolumn*

\LA
\Cite{Sir 32:6-11}
\switchcolumn
\EN
\Cite{Sir 32:6-11}
\switchcolumn*

\LA
\noindent Ubi audítus non est, non effúndas sermónem, et importúne noli extólli in sapiéntia tua. Gémmula carbúnculi in ornaménto auri, et comparátio musicórum in convívio vini; sicut in fabricatióne auri signum est smarágdi, sic númerus musicórum in jucúndo et moderáto vino. Audi tacens, et pro reveréntia accédet tibi bona grátia. Adoléscens, lóquere in tua causa vix. Si bis interrogátus fúeris, hábeat caput respónsum tuum.\par
\switchcolumn
\EN
\noindent Where there is no hearing, pour not out words, and be not lifted up out of season with thy wisdom. A concert of music in a banquet of wine is as a carbuncle set in gold. As a signet of an emerald in a work of gold: so is the melody of music with pleasant and moderate wine. Hear in silence, and for thy reverence good grace shall come to thee. Young man, scarcely speak in thy own cause. If thou be asked twice, let thy answer be short.\par
\switchcolumn*

\LA
\begin{Resp}\Rsign Magna enim sunt judícia tua, Dómine, et inenarrabília verba tua: \Star{} Magnificásti pópulum tuum et honorásti.\end{Resp}
\begin{Resp}\Vsign Transtulísti illos per Mare Rubrum et transvexísti eos per aquam nímiam.\end{Resp}
\begin{Resp}\Rsign Magnificásti pópulum tuum et honorásti.\end{Resp}
\switchcolumn
\EN
\begin{Resp}\Rsign Great are thy judgments, O Lord, and thy words cannot be expressed. \Star{} Thou didst make thy people mighty and honourable.\end{Resp}
\begin{Resp}\Vsign Thou broughtest them through the Red Sea, and leddest them through much water.\end{Resp}
\begin{Resp}\Rsign Thou didst make thy people mighty and honourable.\end{Resp}
\switchcolumn*

\LA
\LessonHead{Lectio III}
\switchcolumn
\EN
\LessonHead{Lesson III}
\switchcolumn*

\LA
\Cite{Sir 32:12-17}
\switchcolumn
\EN
\Cite{Sir 32:12-17}
\switchcolumn*

\LA
\noindent In multis esto quasi ínscius, et audi tacens simul et quærens. In médio magnatórum non præsúmas et, ubi sunt senes, non multum loquáris. Ante grándinem præíbit coruscátio, et ante verecúndiam præíbit grátia et pro reveréntia accédet tibi bona grátia. Et hora surgéndi non te trices: præcúrre autem prior in domum tuam, et illic avocáre, et illic lude, et age conceptiónes tuas et non in delíctis et verbo supérbo; et super his ómnibus benedícito Dóminum, qui fecit te, et inebriántem te ab ómnibus bonis suis.\par
\switchcolumn
\EN
\noindent In many things be as if thou wert ignorant, and hear in silence and withal seeking. In the company of great men take not upon thee: and when the ancients are present, speak not much. Before a storm goeth lightning: and before shamefacedness goeth favour: and for thy reverence good grace shall come to thee. And at the time of rising be not slack: but be first to run home to thy house, and there withdraw thyself, and there take thy pastime. And do what thou hast a mind, but not in sin or proud speech. And for all these things bless the Lord, that made thee, and that replenisheth thee with all his good things.\par
\switchcolumn*

\LA
\begin{Resp}\Rsign Quæ sunt in corde hóminum, óculi tui vident, Dómine, et in libro tuo ómnia scribéntur: \Star{} Homo videt in fácie, Deus autem in corde.\end{Resp}
\begin{Resp}\Vsign Omnia enim corda scrutátur, et univérsas méntium cogitatiónes intéllegit.\end{Resp}
\begin{Resp}\Rsign Homo videt in fácie, Deus autem in corde.\end{Resp}
\begin{Resp}\Vsign Glória Patri, et Fílio, \Star{} et Spirítui Sancto.\end{Resp}
\begin{Resp}\Rsign Homo videt in fácie, Deus autem in corde.\end{Resp}
\switchcolumn
\EN
\begin{Resp}\Rsign Lord, thine eyes behold all that is in the heart of man, and in thy book are they all written. \Star{} Man looketh on the outward appearance, but God looketh on the heart.\end{Resp}
\begin{Resp}\Vsign For He searcheth all hearts, and understandeth all the imaginations of the thoughts.\end{Resp}
\begin{Resp}\Rsign Man looketh on the outward appearance, but God looketh on the heart.\end{Resp}
\begin{Resp}\Vsign Glory be to the Father, and to the Son, \Star{} and to the Holy Ghost.\end{Resp}
\begin{Resp}\Rsign Man looketh on the outward appearance, but God looketh on the heart.\end{Resp}
\switchcolumn*

\end{paracol}

\PartTitle{Proprium de Sanctis}{Proper of Saints}

\addcontentsline{toc}{chapter}{Proprium de Sanctis \textit{— Proper of Saints}}

\SeasonTitle{Mensis Majus}{May}

\addcontentsline{toc}{section}{Mensis Majus \textit{— May}}

\par\needspace{10\baselineskip}\addvspace{1.2em}{\centering\small\textsc{Die 16 Maii} · \textsc{16 May}\par}
\Office{S. Ubaldi Episcopi et Confessoris}{St. Ubald, Bishop and Confessor}{Semiduplex}
\addcontentsline{toc}{subsection}{Die 16 Maii · S. Ubaldi Episcopi et Confessoris \textit{— St. Ubald, Bishop and Confessor}}
\begin{paracol}{2}
\LA
\RefLine{Lectiones I–III: de Scriptura occurrente.}
\switchcolumn
\EN
\RefLine{Lessons I–III: from the Scripture of the day.}
\switchcolumn*

\LA
\RefLine{Lectiones VII–IX: ex Commúni unius Confessoris Pontificis (C4).}
\switchcolumn
\EN
\RefLine{Lessons VII–IX: from the Common of a Confessor Bishop (C4).}
\switchcolumn*

\LA
\LessonHead{Lectio IV}
\switchcolumn
\EN
\LessonHead{Lesson IV}
\switchcolumn*

\LA
\noindent Ubaldus Eugubii in Umbria nobili genere natus, a primis annis pietáte et litteris egregie est institútus; jamque adoléscens, ut uxórem duceret, sæpe tentatus, numquam tamen a proposito servandæ virginitátis recéssit. Sacerdos effectus, patrimónium suum paupéribus et ecclésiis distríbuit, et canonicórum regulárium ordinis sancti Augustíni institutum suscípiens, illud in pátriam tránstulit, atque in eo aliquámdiu sanctíssime vixit. Cujus sanctitátis opinióne evulgata, ab Honorio secundo summo Pontifice ecclésiæ Eugubinæ invítus præfícitur, et episcopalis consecratiónis munere decorátur.\par
\switchcolumn
\EN
\noindent Ubald was born of a noble family at Gubbio in Umbria, and well established in godliness and learning from his earliest years. When he was a young man, it was often proposed to him to marry, but he never abandoned his determination to preserve his virginity. After that he was ordained Priest he divided his inheritance among the poor and Churches, and embraced the Institute of Canons Regular of St. Augustine. This Institute he brought to Gubbio, and for some time led therein a most holy life. When the fame of his saintliness had got noised abroad, Pope Honorius II set him, contrary to his own wishes, over the Church of Gubbio, and he was honoured with consecration as Bishop by the hands of the said Pope himself, (in the year of our Lord 1129.)\par
\switchcolumn*

\LA
\begin{Resp}\Rsign Invéni David servum meum, óleo sancto meo unxi eum: \Star{} Manus enim mea auxiliábitur ei, allelúja.\end{Resp}
\begin{Resp}\Vsign Nihil profíciet inimícus in eo, et fílius iniquitátis non nocébit ei.\end{Resp}
\begin{Resp}\Rsign Manus enim mea auxiliábitur ei, allelúja.\end{Resp}
\switchcolumn
\EN
\begin{Resp}\Rsign I have found David My servant, with My holy oil have I anointed him \Star{} For My hand shall help him, alleluia.\end{Resp}
\begin{Resp}\Vsign The enemy shall prevail nothing against him, nor the son of wickedness afflict him.\end{Resp}
\begin{Resp}\Rsign For My hand shall help him, alleluia.\end{Resp}
\switchcolumn*

\LA
\LessonHead{Lectio V}
\switchcolumn
\EN
\LessonHead{Lesson V}
\switchcolumn*

\LA
\noindent Ad suam ítaque revertens ecclésiam, cum de consueta vivéndi ratióne nihil ádmodum immutasset, in omni tamen virtútum genere eo magis eminére cœpit, quo efficacius aliórum étiam salútem verbo et exemplo procuraret, factus forma gregis ex animo. Nam victu parco, vestítu moderáto, léctulo áspero et paupérrimo, crucis mortificatiónem jugiter in suo corpore circumferebat, dum inexplébili oratiónis studio spíritum quotídie recrearet. Hinc admirábilem illam mansuetúdinem est adeptus, qua gravíssimas injurias et contumelias non modo æquanimiter tulit, verum étiam mirifico dilectiónis afféctu persecutores suos omni benignitátis testimonio complectebátur.\par
\switchcolumn
\EN
\noindent When Ubald came to live as Bishop in Gubbio, he changed his way of life in no wise from that which he had led before, but his virtues began to be more eminent because his word and example were now more able to benefit his neighbours, to whom the shepherd of their souls was a pattern, not by outward showing only, but from his heart. He ate little, dressed simply, and slept upon a hard and very poor bed. He “always bore in the body the dying of the Lord Jesus,” (2 Cor. iv. 10,) while he daily fed his soul in unceasing and earnest prayer. Hence he acquired such wonderful meekness, that when he was most grievously wronged and insulted he not only took it patiently, but, by a strange impulse of love for them, embraced his persecutors with every proof of affection.\par
\switchcolumn*

\LA
\begin{Resp}\Rsign Pósui adjutórium super poténtem, et exaltávi eléctum de plebe mea: \Star{} Manus enim mea auxiliábitur ei, allelúja.\end{Resp}
\begin{Resp}\Vsign Invéni David servum meum, óleo sancto meo unxi eum.\end{Resp}
\begin{Resp}\Rsign Manus enim mea auxiliábitur ei, allelúja.\end{Resp}
\switchcolumn
\EN
\begin{Resp}\Rsign I have laid help upon one that is mighty, and have exalted one chosen out of My people \Star{} For My hand shall help him, alleluia.\end{Resp}
\begin{Resp}\Vsign I have found David My servant, with My holy oil have I anointed him.\end{Resp}
\begin{Resp}\Rsign For My hand shall help him, alleluia.\end{Resp}
\switchcolumn*

\LA
\LessonHead{Lectio VI}
\switchcolumn
\EN
\LessonHead{Lesson VI}
\switchcolumn*

\LA
\noindent Biennio, antequam ex hac vita migraret, cum diutinis afflicetarétur infirmitátibus, inter acerbíssimos corporis cruciátus velut aurum in fornace purgátum, Deo grátias indesinenter agebat. Adveniénte autem sacro Pentecostes die, cum multis annis ecclésiam sibi commissam cumma cum laude gubernasset, sanctis opéribus ac miraculis clarus quiévit in pace. Quem Cælestinus Papa tertius in Sanctórum númerum retulit. Ejus virtus præcipue in effugandis spirítibus immundis elucet. Corpus vero, per tot sæcula incorruptum, magna fidelium veneratióne in pátria colitur, quam non semel a præsénti discrimine liberávit.\par
\switchcolumn
\EN
\noindent Just the space of two years before Ubald passed away from this present life, he was tried as gold in the furnace, by grievous bodily weakness, and, day after day, amid the sharpest sufferings, he never ceased patiently to give God thanks. He rested in peace on the sacred day of Pentecost, (in the year 1160,) having for many years governed with great praise the Church which had been entrusted to him, and glorious for good works and miracles. Pope Celestine III. numbered him with the Saints. His strength is most chiefly shown in the casting out of evil spirits. His body hath remained without corruption for all these ages, and is reverenced greatly in his native town by Christ's faithful people. To them he hath more than once shown himself good at need.\par
\switchcolumn*

\LA
\begin{Resp}\Rsign Iste est, qui ante Deum magnas virtútes operátus est, et omnis terra doctrína ejus repléta est: \Star{} Ipse intercédat pro peccátis ómnium populórum, allelúja.\end{Resp}
\begin{Resp}\Vsign Iste est, qui contémpsit vitam mundi, et pervénit ad cæléstia regna.\end{Resp}
\begin{Resp}\Rsign Ipse intercédat pro peccátis ómnium populórum, allelúja.\end{Resp}
\begin{Resp}\Vsign Glória Patri, et Fílio, \Star{} et Spirítui Sancto.\end{Resp}
\begin{Resp}\Rsign Ipse intercédat pro peccátis ómnium populórum, allelúja.\end{Resp}
\switchcolumn
\EN
\begin{Resp}\Rsign This is he which wrought great wonders before God, and the whole earth is full of his teaching \Star{} May he pray for all people, that their sins may be forgiven unto them, alleluia.\end{Resp}
\begin{Resp}\Vsign This is he which loved not his life in this world, and hath attained unto the kingdom of heaven.\end{Resp}
\begin{Resp}\Rsign May he pray for all people, that their sins may be forgiven unto them, alleluia.\end{Resp}
\begin{Resp}\Vsign Glory be to the Father, and to the Son, \Star{} and to the Holy Ghost.\end{Resp}
\begin{Resp}\Rsign May he pray for all people, that their sins may be forgiven unto them, alleluia.\end{Resp}
\switchcolumn*

\LA
\LessonHead{Lectio IX (pro commemoratione)}
\switchcolumn
\EN
\LessonHead{Lesson IX (when commemorated)}
\switchcolumn*

\LA
\LessonTitle{Commemoratio: S. Ubaldi Episcopi et Confessoris}
\switchcolumn
\EN
\LessonTitle{Commemoration of: St. Ubald, Bishop and Confessor}
\switchcolumn*

\LA
\noindent Ubáldus, Eugúbii in Umbria nóbili génere natus, a primis annis pietáte et lítteris egrégie est institútus; jamque adoléscens, ut uxórem dúceret, sæpe tentátus, nunquam tamen a propósito servándæ virginitátis recéssit. Sacérdos efféctus, patrimónium suum paupéribus et ecclésiis distríbuit. Canonicórum regulárium órdinis sancti Augustíni institútum suscípiens, illud in pátriam tránstulit. Ab Honório secúndo summo Pontífice ecclésiæ Eugubínæ invítus præfícitur, et episcopális consecratiónis múnere decorátur. Factus forma gregis ex ánimo, de consuéta vivéndi ratióne nihil ádmodum immutávit, et in omni virtútum génere enítuit. Diútius infirmitátibus afflíctus, Deo indesinénter grátias agébat. Cum multis annis ecclésiam sibi commíssam summa cum laude gubernásset, sanctis opéribus et miráculis clarus, quiévit in pace.\par
\switchcolumn
\EN
\noindent St. Ubald was born of a noble family at Gubbio in Umbria, and was well trained, from his earliest years, in devotion and in letters. As a young man, he was often urged to marry, but he never gave up his determination to preserve his virginity. When he was made a priest, he distributed his patrimony to the poor and to churches. He entered the Institute of the Canons Regular of St. Augustine and brought it to his own country. Against his will, he was put over the Church of Gubbio by Pope Honorius II and given episcopal consecration. Making himself with his whole heart a model for his flock, he changed nothing in his way of life and became illustrious for every kind of virtue. Although he suffered from various infirmities for a long time, he never ceased to give thanks to God. When he had governed the Church entrusted to him for many years in a way deserving of the highest praise, he fell asleep in peace, famous for his works and miracles.\par
\switchcolumn*

\LA
\TeDeum{Te Deum}
\switchcolumn
\EN
\TeDeum{Te Deum}
\switchcolumn*

\end{paracol}

\par\needspace{10\baselineskip}\addvspace{1.2em}{\centering\small\textsc{Die 17 Maii} · \textsc{17 May}\par}
\Office{S. Paschalis Baylon Confessoris}{St. Pascal Baylon Confessor}{Duplex}
\addcontentsline{toc}{subsection}{Die 17 Maii · S. Paschalis Baylon Confessoris \textit{— St. Pascal Baylon Confessor}}
\begin{paracol}{2}
\LA
\RefLine{Lectiones I–III: de Scriptura occurrente.}
\switchcolumn
\EN
\RefLine{Lessons I–III: from the Scripture of the day.}
\switchcolumn*

\LA
\RefLine{Lectiones VII–IX: ex Commúni Confessoris non pontificis (C5).}
\switchcolumn
\EN
\RefLine{Lessons VII–IX: from the Common of a Confessor not a Bishop (C5).}
\switchcolumn*

\LA
\LessonHead{Lectio IV}
\switchcolumn
\EN
\LessonHead{Lesson IV}
\switchcolumn*

\LA
\noindent Paschalis Baylon, paupéribus piisque paréntibus in oppido Turris Formosæ, Seguntinæ diœcesis, in Aragónia natus, a téneris annis plura dedit futuræ sanctitátis indicia. Sortítus ánimam bonam ac rerum cæléstium appríme studiosam, puerítiam atque adolescéntiam in gregis custódia transegit; quam ille vivéndi ratiónem ideo præcipue diligebat, quod humilitáti fovendæ ac innocéntiæ conservandæ in primis utilem atque opportunam judicáret. Erat in victu modicus, in oratióne assiduus, tantáque apud coævos et socios florebat auctoritate et grátia, ut eórum lites compónens, erróres corrigens, ignorantiam erudiens ac desidiam éxcitans, velut ómnium parens et magister máximo studio colerétur ac amarétur: beátus étiam tum a plerisque appellatus.\par
\switchcolumn
\EN
\noindent Paschal Baylon was the son of poor and godly parents, in the town of Torre Hermosa, and Diocese of Sagunta in Aragon, (in the year of our Lord 1540.) From his childhood he gave indications of a holy life. He was naturally of a good disposition, and very wishful to learn about heavenly things. His boyhood and youth he passed in the occupation of a shepherd. This way of life pleased him well, because he thought it one useful and fitted to nourish lowliness and keep innocency. He ate little, and was instant in prayer. He had great weight and favour with his fellows and neighbours, whose quarrels he healed, corrected their mistakes, enlightened their ignorance, and roused them from idleness. They all greatly honoured and loved him, as though he were their father and teacher, and even then many called him “Beato,” that is “the Blessed.”\par
\switchcolumn*

\LA
\begin{Resp}\Rsign Honéstum fecit illum Dóminus, et custodívit eum ab inimícis, et a seductóribus tutávit illum: \Star{} Et dedit illi claritátem ætérnam, allelúja.\end{Resp}
\begin{Resp}\Vsign Justum dedúxit Dóminus per vias rectas, et osténdit illi regnum Dei.\end{Resp}
\begin{Resp}\Rsign Et dedit illi claritátem ætérnam, allelúja.\end{Resp}
\switchcolumn
\EN
\begin{Resp}\Rsign The Lord made him honourable, and defended him from his enemies, and kept him safe from those that lay in wait for him; \Star{} And gave him perpetual glory, alleluia.\end{Resp}
\begin{Resp}\Vsign He went down with him into the pit, and left him not in bonds.\end{Resp}
\begin{Resp}\Rsign And gave him perpetual glory, alleluia.\end{Resp}
\switchcolumn*

\LA
\LessonHead{Lectio V}
\switchcolumn
\EN
\LessonHead{Lesson V}
\switchcolumn*

\LA
\noindent Qui vero in sæculo, terra nempe desérta et inaquósa, adeo feliciter adoléverat, flos convallium, plantátus in domo Dómini, mirum ubíque sparsit sanctitátis odórem. Igitur Paschalis arrepto vitæ severioris instituto, atque in ordine Minórum strictioris observantiæ Discalceatórum cooptatus, exsultávit ut gigas ad curréndam viam suam; totumque se Dómino excoléndum tradens, dies noctesque cogitabat, qua se ratióne magis ei magisque conformaret. Ita factum est brevi, ut eum, tamquam seraphicæ perfectiónis exemplar, ipsi quoque provectiores imitándum sibi propónerent. Ipse autem in humili serviéntium gradu constitútus, se velut ómnium peripséma réputans, ardua quæque et abjecta domus ministeria, véluti jure quodam peculiari sibi débita, summa cum hilaritate suscipiebat, et exercebat humilitate ac patiéntia pari. Carnem spiritui quandoque reluctari nitentem jugi maceratióne afflictávit, atque in servitútem redégit; spíritum vero assidua sui abnegatióne ferventiórem in dies ad anterióra extendebat.\par
\switchcolumn
\EN
\noindent In a world which was to him “a dry land, where no water is” the vallies, “planted in the House of the Lord” (Ps. xci. 14,) whose strange sweetness spread all around. When he took upon him an harder life, by entering the Institute of barefooted Grey Friars, of the strict Observance, “he rejoiced as a strong man to run a race” (Ps. xviii. 6,) and gave himself up altogether to serve the Lord, thinking by day and by night only how he might attain more and more to have that mind in him which was also in Christ Jesus (Phil. ii. 5.) And so it came to pass in a little while, that his very elders set him before them for their model, as a pattern of a man seeking to be perfect in the path of the Seraphic Order. Paschal himself held the lowly place of a lay brother, and deemed himself “the off-scouring of all things” (1. Cor. iv. 13.) He took most cheerfully, and discharged with the greatest humility and patience, the hardest and meanest work of the house, as though such were his peculiar right. His flesh would sometimes rebel against his spirit, but he broke it under the yoke of mortification, and brought it into subjection. Day by day the spirit of self-denial waxed stronger in him, and “forgetting those things which were behind, he reached forth unto those things which were before” (Phil. iii. 13.)\par
\switchcolumn*

\LA
\begin{Resp}\Rsign Amávit eum Dóminus, et ornávit eum: stolam glóriæ índuit eum, \Star{} Et ad portas paradísi coronávit eum, allelúja.\end{Resp}
\begin{Resp}\Vsign Índuit eum Dóminus lorícam fídei, et ornávit eum.\end{Resp}
\begin{Resp}\Rsign Et ad portas paradísi coronávit eum, allelúja.\end{Resp}
\switchcolumn
\EN
\begin{Resp}\Rsign The Lord loved him and beautified him; He clothed him with a robe of glory, \Star{} And crowned him at the gates of Paradise, alleluia.\end{Resp}
\begin{Resp}\Vsign The Lord hath put on him the breast-plate of faith, and hath adorned him.\end{Resp}
\begin{Resp}\Rsign And crowned him at the gates of Paradise, alleluia.\end{Resp}
\switchcolumn*

\LA
\LessonHead{Lectio VI}
\switchcolumn
\EN
\LessonHead{Lesson VI}
\switchcolumn*

\LA
\noindent Deíparam Vírginem, cujus clientelæ se ab ineunte ætate dicáverat, tamquam matrem quotidiánis colebat obsequiis atque filiali exorábat fiducia. Porro erga sanctíssimum Eucharistiæ sacraméntum difficile dictu est quam ardénti tenerétur devotiónis afféctu; quem defunctus étiam in cadávere retinére visus est, dum, jacens in féretro, ad sacræ Hostiæ elevatiónem bis óculos reserávit et clausit, magna ómnium, qui áderant, admiratióne. Ejusdem veritátem inter hæreticos publice palamque proféssus, multa et gravia ob eam causam perpéssus est; crebro étiam ad necem petítus, sed singulari Dei providéntia ab impiórum mánibus ereptus. Sæpe inter orándum ómnibus destituebátur sensibus, dulcique languebat amoris deliquio: quo témpore cælestem illam sciéntiam hausisse créditus est, qua, homo rudis et illiteratus, de mystériis fidei difficillimis respondére, atque áliquot étiam libros conscríbere pótuit. Denique meritis plenus, eádem qua prædixerat hora, feliciter migrávit ad Dóminum, anno salútis millesimo quingentésimo nonagesimo secundo, sexto décimo Kalendas Junii, eodem, quo natus fuerat, festo Pentecostes recurrente, annum agens secúndum supra quinquagesimum. Quibus aliisque virtútibus insignem, ac miraculis tam in vita quam post mortem clarum, Paulus quintus Pontifex maximus illum Beátum appellávit; Alexander autem octavus Sanctórum catálogo adscripsit; tandem Leo décimus tertius peculiarem cœtuum eucharisticórum, item societátum ómnium a sanctíssima Eucharistia, sive quæ in posterum futuræ sunt, patronum cælestem declarávit et constítuit.\par
\switchcolumn
\EN
\noindent The Virgin Mother of God he had vowed himself when he was but a little lad, and he paid her every day the services of a son, and trusted her as a mother. It is hard to tell how intense was the love which bound him to the Most Holy Sacrament of the Eucharist, a love which seemed literally stronger than death, for when his dead body was found lying on the bier, its eyes opened and shut twice when the Sacred Host was lifted up, to the amazement of all that were there. When he was among heretics, he suffered much and grievously at their hands for plainly and openly telling the truth touching this Sacrament they often sought after him to murder him, but by the singular Providence of God he was delivered from those wicked men. When he was at prayer he often became utterly insensible, and his soul fainted away with the love of God. During these trances it was believed that he received directly from heaven that knowledge which he had, and which enabled him, although a man altogether rough and unlettered, to answer the hardest questions upon the mysteries of the faith, and even to write some books. At last, full of good works, he joyfully passed away to be ever with the Lord, at the hour foretold by himself, on the Feast of Pentecost, the 17th day of May, in the year of salvation 1592, on which day also he had been born fifty -two years before. Illustrious for the graces above mentioned, and for the miracles which he worked both during his life and after his death, he was named Blessed by Pope Paul V., and Alexander VIII. enrolled him among the Saints.\par
\switchcolumn*

\LA
\begin{Resp}\Rsign Iste homo perfécit ómnia quæ locútus est ei Deus, et dixit ad eum: Ingrédere in réquiem meam: \Star{} Quia te vidi justum coram me ex ómnibus géntibus, allelúja.\end{Resp}
\begin{Resp}\Vsign Iste est, qui contémpsit vitam mundi, et pervénit ad cæléstia regna.\end{Resp}
\begin{Resp}\Rsign Quia te vidi justum coram me ex ómnibus géntibus, allelúja.\end{Resp}
\begin{Resp}\Vsign Glória Patri, et Fílio, \Star{} et Spirítui Sancto.\end{Resp}
\begin{Resp}\Rsign Quia te vidi justum coram me ex ómnibus géntibus, allelúja.\end{Resp}
\switchcolumn
\EN
\begin{Resp}\Rsign This is he which did according unto all that God commanded him; and God said unto him: Enter thou into My rest, \Star{} For thee have I seen righteous before Me among all people, alleluia.\end{Resp}
\begin{Resp}\Vsign This is he which loved not his life in this world, and is come unto an everlasting kingdom.\end{Resp}
\begin{Resp}\Rsign For thee have I seen righteous before Me among all people, alleluia.\end{Resp}
\begin{Resp}\Vsign Glory be to the Father, and to the Son, \Star{} and to the Holy Ghost.\end{Resp}
\begin{Resp}\Rsign For thee have I seen righteous before Me among all people, alleluia.\end{Resp}
\switchcolumn*

\LA
\LessonHead{Lectio IX (pro commemoratione)}
\switchcolumn
\EN
\LessonHead{Lesson IX (when commemorated)}
\switchcolumn*

\LA
\LessonTitle{Commemoratio: S. Paschalis Baylon Confessoris}
\switchcolumn
\EN
\LessonTitle{Commemoration of: St. Pascal Baylon Confessor}
\switchcolumn*

\LA
\noindent Paschális Baylon, paupéribus piísque paréntibus in óppido Turris Formósæ in Aragónia natus puerítiam atque adolescéntiam in gregis custódia transégit. Póstea severióris vitæ institútum ampléxus, et órdini fratrum Minórum adscríptus, júgiter cogitábat, qua se ratióne magis magísque Christo crucifíxo conformáret. Deíparam Vírginem, cujus clientélæ se ab ineúnte ætáte dicáverat, tamquam matrem filiálibus et cotidiánis obséquiis colébat. Erga Eucharístiam ténero et assíduo flagrávit devotiónis afféctu, quem defúnctus étiam retinére visus est, dum jacens in féretro, ad sacræ Hóstiæ elevatiónem bis óculos reserávit et clausit, magna ómnium, qui áderant, admiratióne. Méritis plenus migrávit ad Dóminum, anno millésimo quingentésimo nonagésimo secúndo. Eum Leo décimus tértius peculiárem cœ́tuum eucharisticórum, societatúmque ómnium a sanctíssima Eucharístia patrónum cæléstem declarávit et constítuit.\par
\switchcolumn
\EN
\noindent Paschal Baylon was born of poor and devout parents in the town of Torre Hermosa in Aragon, and spent his boyhood and youth in herding sheep. When he had embraced an austere way of life by entering the Friars Minor, he thought ceaselessly of how he might conform ever more closely to Christ crucified. With daily filial devotion he honoured as mother the Virgin Mother of God, to whom he had consecrated himself from his earliest years. He burned with a tender and steady love for the Eucharist, a love which he manifested even in death, when, lying on his bier, he opened and closed his eyes twice at the elevation of the sacred Host, to the astonishment of all present. Full of merits, he went to the Lord in the year 1592. Leo XIII declared and appointed him the special heavenly patron of all Eucharistic congresses and of all associations in honour of the most holy Eucharist.\par
\switchcolumn*

\LA
\TeDeum{Te Deum}
\switchcolumn
\EN
\TeDeum{Te Deum}
\switchcolumn*

\end{paracol}

\par\needspace{10\baselineskip}\addvspace{1.2em}{\centering\small\textsc{Die 18 Maii} · \textsc{18 May}\par}
\Office{S. Venantii Martyris}{St. Venantius, Martyr}{Duplex}
\addcontentsline{toc}{subsection}{Die 18 Maii · S. Venantii Martyris \textit{— St. Venantius, Martyr}}
\begin{paracol}{2}
\LA
\RefLine{Lectiones I–III: de Scriptura occurrente.}
\switchcolumn
\EN
\RefLine{Lessons I–III: from the Scripture of the day.}
\switchcolumn*

\LA
\RefLine{Lectiones VII–IX: ex Commúni unius Martyris Tempore Paschali (C2ap).}
\switchcolumn
\EN
\RefLine{Lessons VII–IX: from the Common of a Single Martyr in Paschaltide (C2ap).}
\switchcolumn*

\LA
\LessonHead{Lectio IV}
\switchcolumn
\EN
\LessonHead{Lesson IV}
\switchcolumn*

\LA
\noindent Venantius Camers, quindecim annos natus, cum christianæ religiónis accusarétur apud Antíochum, qui sub Decio imperatóre Camerino præerat, in porta civitátis præsidi se obtulit. Quem ille pollicitatiónibus ac terróribus diu tentátum, flagris cædi et vinculis astringi jussit; sed is mirabíliter ab Angelo solutus, lampadibus póstea aduritur, atque inverso ore, fumo supposito, suspénditur. Ejus constantiam in tormentis demirátus Anastasius Cornicularius, et quod eum ab Angelo íterum solutum cándida veste supra fumum ambulántem vidísset, in Christum credidit; et a beato Porphyrio presbytero cum família baptizatus, paulo post martyrii palmam cum eodem promeruit.\par
\switchcolumn
\EN
\noindent Venantius was a lad of Camerino (in the neighbourhood of Ancona,) who at fifteen years of age was accused of Christianity before Antiochus, Praefect of Camerino under the Emperor Decius. Venantius therefore appeared before Antiochus at the gate of the city, and when the Praefect had striven with him for a long while, by promises and threats, he commanded him to be scourged and thrown into irons, but an Angel loosed his bonds. He was afterwards scarified with lamps, and hung head downwards in smoke. Anastasius the trumpeter was amazed at his hardiness under suffering, and when it appeared to him that the Martyr was a second time loosed by an Angel, and was walking in white raiment on the smoke, he believed in Christ, and was baptized, with all his house, by the blessed Priest Porphyry, and a little while after they both together earned the palm of martyrdom.\par
\switchcolumn*

\LA
\begin{Resp}\Rsign Lux perpétua lucébit Sanctis tuis, Dómine, \Star{} Et ætérnitas témporum, allelúja, allelúja.\end{Resp}
\begin{Resp}\Vsign Lætítia sempitérna erit super cápita eórum: gáudium et exsultatiónem obtinébunt.\end{Resp}
\begin{Resp}\Rsign Et ætérnitas témporum, allelúja, allelúja.\end{Resp}
\switchcolumn
\EN
\begin{Resp}\Rsign Eternal light will shine over your saints, O Lord, \Star{} And the eternity of the times, alleluia, alleluia.\end{Resp}
\begin{Resp}\Vsign Everlasting joy shall be upon their heads, they shall obtain joy and gladness\end{Resp}
\begin{Resp}\Rsign And the eternity of the times, alleluia, alleluia.\end{Resp}
\switchcolumn*

\LA
\LessonHead{Lectio V}
\switchcolumn
\EN
\LessonHead{Lesson V}
\switchcolumn*

\LA
\noindent At Venantius præsidi sistitur, et ab eo íterum frustra tentátus ut Christi fidem deséreret, in carcerem conjícitur; quo Attalus præco mittitur, qui ei dicat se quoque christianum fuísse et ei nómini proptérea renuntiásse, quod cognovísset inane esse fidei comméntum, quo Christiáni præséntibus se ábdicant ob vanam futurórum spem. Verum nobilis Christi athleta, callidi hostis insidias non ignorans, diaboli ministrum a se pénitus rejécit. Quare ad præsidem íterum adducto omnes contusi sunt dentes maxillæque confractæ; atque, ita cæsus, in sterquilinium dejícitur. Sed inde ab Angelo quoque ereptus, rursus stetit ante judicem; qui, Venantio adhuc loquente, e tribunali cécidit, et in ea voce, Verus est Venantii Deus, nostros deos destrúite, exclámans exspirávit.\par
\switchcolumn
\EN
\noindent Now Venantius stood before the Praefect, and when he had again vainly tempted him to give up his faith in Christ, he cast him into prison, and sent unto him Attalus the crier. Attalus told him how that he also had been a Christian, but had denied that name, seeing it was a foolish faith which made Christians to throw away things present for a groundless hope of things to come. But Christ's brave champion, well knowing the wiles of our subtle enemy, drove the devil's servant from his presence. When he appeared again before the Praefect, his teeth and jaws were broken, and so mangled he was cast out upon a dunghill. But thence also an Angel delivered him, and he stood again before the judge. And there while Venantius was yet speaking, the judge fell from off the judgment-seat, and when he had cried with a loud voice, “Venantius, his God is true, take away our gods,” he died.\par
\switchcolumn*

\LA
\begin{Resp}\Rsign In servis suis, allelúja, \Star{} Consolábitur Deus, allelúja.\end{Resp}
\begin{Resp}\Vsign Judicábit Dóminus pópulum suum, et in servis suis.\end{Resp}
\begin{Resp}\Rsign Consolábitur Deus, allelúja.\end{Resp}
\switchcolumn
\EN
\begin{Resp}\Rsign His servants, alleluia \Star{} God will comfort, alleluia, alleluia.\end{Resp}
\begin{Resp}\Vsign The Lord will judge His people and, His servants,\end{Resp}
\begin{Resp}\Rsign God will comfort, alleluia, alleluia.\end{Resp}
\switchcolumn*

\LA
\LessonHead{Lectio VI}
\switchcolumn
\EN
\LessonHead{Lesson VI}
\switchcolumn*

\LA
\noindent Quod cum præsidi nuntiátum esset, extemplo Venántium leónibus óbjici jussit; qui, naturali feritáte omissa, ad ejus se pedes abjecérunt. Interim ille pópulum Christi fidem edocébat. Quare inde amotus, íterum in carcerem traditur. Cumque postridie præsidi referret Porphyrius, se per visum noctu pópulos, quos Venantius aqua tingebat, claríssima luce fulgéntes, ipsum vero præsidem obscuríssima caligine operátum vidisse; præses, ira incénsus, eum illico cápite plecti imperat, deínde Venántium per loca vépribus et cárduis cónsita trahi usque ad vésperam. Is, cum semiánimis relíctus esset, mane se íterum præsidi præsentávit; cujus jussu statim e rupe præcipitátur. Sed inde étiam divinitus ereptus, denuo per loca aspera ad mille passus trahitur; ubi, milítibus siti æstuántibus, in próxima convalle, ex lápide, in quo et genuum formam relíquit, sicut étiam nunc in ejus ecclésia vidére licet, crucis signo a Venantio facto, aquæ manarunt. Eo miraculo plures permoti in Christum credidérunt, quos omnes præses eo loci una cum Venantio cápite feriri jussit. Fulgura et terræmotus eo témpore ita magni fuére, ut præses aufugeret; qui paucis tamen post diébus, divinam haud valens effugere justítiam, turpíssimam mortem oppétiit. Christiáni interim Venantii et aliórum córpora honorifico loco sepeliérunt: quæ Camerini, in ecclésia Venantio dicata, cóndita adhuc sunt.\par
\switchcolumn
\EN
\noindent Then they told the President of it, he commanded Venantius to be straightway thrown to the lions. But the beasts were not wild to him, and lay down at his feet. And meanwhile he taught the Christian faith to the people. So they took him away from thence and cast him once more into prison. The next day Porphyry came to the President, and told him how that he had seen in a vision of the night Venantius sprinkling certain ones with water, and they that were sprinkled shone with a marvellous light, and the President himself hidden in deep darkness. Then the President was moved to great anger and commanded forthwith to behead Porphyry. As for Venantius, he bade them drag him about in rough places, full of briars and thistles, until the evening. When it was over, he was left half dead, but in the morning he stood for the last time before the President, who commanded to cast him down from a steep rock. It pleased God that this should not kill him, and he was haled again through rough places for about a mile. There the soldiers were athirst, and Venantius, by the sign of the Cross, made waters to flow from a stone in a gulley hard by. This is that stone whereon also he left the imprint of his knees, and which can be seen to this day in his Church. By this wonder many were moved to believe in Christ and the President commanded them all, and Venantius with them, to be beheaded in the same place where they were. When it was done there were great lightnings and earthquakes, so that the President fled, but he could not fly from the judgment of God, and but a few days thereafter he died a most shameful death. Meanwhile the Christians took the bodies of Venantius and the others, and buried them in an honourable place, wherein they lie to this day, under the Church at Camerino which is dedicated to Venantius.\par
\switchcolumn*

\LA
\begin{Resp}\Rsign Fíliæ Jerúsalem, veníte et vidéte Mártyres cum corónis, quibus coronávit eos Dóminus \Star{} In die solemnitátis et lætítiæ, allelúja.\end{Resp}
\begin{Resp}\Vsign Quóniam confortávit seras portárum tuárum, benedíxit fílios tuos in te.\end{Resp}
\begin{Resp}\Rsign In die solemnitátis et lætítiæ, allelúja.\end{Resp}
\begin{Resp}\Vsign Glória Patri, et Fílio, \Star{} et Spirítui Sancto.\end{Resp}
\begin{Resp}\Rsign In die solemnitátis et lætítiæ, allelúja.\end{Resp}
\switchcolumn
\EN
\begin{Resp}\Rsign Come forth, O ye daughters of Jerusalem, and behold the Martyrs with the crowns wherewith the Lord crowned them; \Star{} In the day of His feasting, and of His gladness. Alleluia.\end{Resp}
\begin{Resp}\Vsign For He hath strengthened the bars of thy gates He hath blessed thy children within thee.\end{Resp}
\begin{Resp}\Rsign In the day of His feasting, and of His gladness. Alleluia.\end{Resp}
\begin{Resp}\Vsign Glory be to the Father, and to the Son, \Star{} and to the Holy Ghost.\end{Resp}
\begin{Resp}\Rsign In the day of His feasting, and of His gladness. Alleluia.\end{Resp}
\switchcolumn*

\LA
\LessonHead{Lectio IX (pro commemoratione)}
\switchcolumn
\EN
\LessonHead{Lesson IX (when commemorated)}
\switchcolumn*

\LA
\LessonTitle{Commemoratio: S. Venantii Martyris}
\switchcolumn
\EN
\LessonTitle{Commemoration of: St. Venantius, Martyr}
\switchcolumn*

\LA
\noindent Venántius Camers, quíndecim annos natus, christiánæ religiónis accusátus apud Antíochum, qui sub Décio imperatóre Cameríno prǽerat, in porta civitátis prǽsidi se óbtulit. Quem ille, pollicitatiónibus ac terróribus diu tentátum, flagris cædi et vínculis astríngi jussit; e quibus mirabíliter ab Angelo solútus, lampádibus póstea adúritur, atque invérso ore, fumo suppósito, suspénditur. Eídem, ad prǽsidem íterum addúcto, omnes contúsi sunt dentes maxillǽque confráctæ, atque ita cæsus in sterquilínium deícitur. Sed inde ab Angelo eréptus, rursus stetit ante júdicem; qui, Venántio adhuc loquénte, e tribunáli cécidit et in ea voce, Verus est Venántii Deus, nostros deos destrúite, exclámans exspirávit. Demum, post nova et exquisíta torménta, una cum áliis decem gloriósi certáminis cursum cervícibus abscíssis implévit. Quorum córpora christiáni honorífico loco sepeliérunt, quæ Cameríni, in ecclésia Venántio dicáta, cóndita sunt.\par
\switchcolumn
\EN
\noindent At the age of fifteen, Venantius of Camerino was denounced for his Christian religion to Antiochus, who was Prefect of Camerino under the Emperor Decius. Venantius presented himself to the prefect at the city gates. For a long time the prefect tempted him by means of promises and threats, then commanded that he be beaten and chained. Miraculously freed by an angel, Venantius was then burned with torches and suspended face down over a smoking fire. Led back again to the governor, he had all his teeth and jaws broken and, thus mutilated, he was thrown into a pit of dung. Rescued from this pit by an Angel, he stood once more before the judge, who, even as Venantius was speaking to him, fell from his tribunal, crying out, “Venantius, his God is true, take away our gods!” and expired. At length, after new and exquisite torments, Venantius was beheaded, along with ten others, and so finished the course of his glorious struggle. The Christians gave honourable burial to the bodies of these martyrs, who now rest in Camerino, in the church dedicated to Venantius.\par
\switchcolumn*

\LA
\TeDeum{Te Deum}
\switchcolumn
\EN
\TeDeum{Te Deum}
\switchcolumn*

\end{paracol}

\par\needspace{10\baselineskip}\addvspace{1.2em}{\centering\small\textsc{Die 19 Maii} · \textsc{19 May}\par}
\Office{S. Petri Celestini Papæ et Confessoris}{St. Peter Celestine, Pope and Confessor}{Duplex}
\addcontentsline{toc}{subsection}{Die 19 Maii · S. Petri Celestini Papæ et Confessoris \textit{— St. Peter Celestine, Pope and Confessor}}
\begin{paracol}{2}
\LA
\RefLine{Lectiones I–III: de Scriptura occurrente.}
\switchcolumn
\EN
\RefLine{Lessons I–III: from the Scripture of the day.}
\switchcolumn*

\LA
\RefLine{Lectiones VII–VIII: ex Commúni Summorum Pontificum Confessorum (C4b).}
\switchcolumn
\EN
\RefLine{Lessons VII–VIII: from the Common of Popes (C4b).}
\switchcolumn*

\LA
\RefLine{Lectio IX: de S. Pudentianæ Virginis (05-19o).}
\switchcolumn
\EN
\RefLine{Lesson IX: of St. Pudentiana, Virgin (05-19o).}
\switchcolumn*

\LA
\LessonHead{Lectio IV}
\switchcolumn
\EN
\LessonHead{Lesson IV}
\switchcolumn*

\LA
\noindent Petrus, a nomine quo Pontifex est appellatus, Cælestinus dictus, honestis catholicisque parentibus Aeserniæ in Samnitibus natus, adolescentiam vix ingressus, ut ánimum a mundi illecebris custodiret, in solitudinem secessit. Ibi contemplationibus mentem mitriens, corpus in servitutem redigens, ferream catenam ad nudam carnem adhibebat. Congregationem, quæ postea Cælestinorum dicta est, sub regula sancti Benedicti instituit. Hinc quasi lucerna supra candelabrum posita, cum abscondi nequiret, (Romana Ecclesia diu viduata pastore) in Petri cathedram ignorans et absens ascitus, magna novitatis admiratione non minus quam repentino gaudio cunctos affecit. Cum autem in pontificatus sublimitate collocatus, variis distentus curis, assuitis incumbere meditationibus vix posse cognosceret, oneri pariter et honori voluntarie cessit. Indeque priscam vitæ rationem repetens, obdormivit in Domino, ejusque pretiosam mortem crux præfulgens in aëre ante cubiculi ostium reddidit amplius gloriosam. Miraculis multis tam vivens, quam post obitum claruit; quibus rite examinatis, Clemens quintus anno postquam decessit undecimo, sanctorum numero adscripsit.\par
\switchcolumn
\EN
\noindent This Peter, who is called Peter Celestine, because when he became Pope he did so under the title of Celestine V, was the son of respectable Catholic parents, and was born at Isernia in Apulia, (about the year of grace 1221.) He was hardly entered on boyhood, when he withdrew into a desert, in order to keep his soul safe from the snares of the world. In solitude he fed his mind with heavenly meditation, and brought his body into subjection, even by wearing an iron chain next to his bare flesh. He founded, under the Rule of St. Benedict, that congregation which was afterwards known as the Celestine. His light, as of a candle set upon a candlestick, could not be kept hidden, and after the Church of Rome had for a long while been widowed of a shepherd, he was chosen without his knowledge and in his absence, to fill the chair of Peter. The news of his election filled himself with as great amazement, as it did all others with sudden joy. When, however, he was seated in the exalted place of the Papal dignity, he found that the many cares by which he was beset made it well-nigh impossible for him to give himself to his accustomed meditations (after four months,) of his own free will he resigned the burden and the honour together (on the 13th day of December, 1294); and, while he sought to return to his old way of life, (on the 19th day of May, 1296,) he fell asleep in the Lord. How precious his death was in His sight was gloriously manifested by a Cross which appeared shining in the air before the door of the cell. He was illustrious for miracles both during his life and after his death, and when these had been duly investigated, Clement V., in the eleventh year after his departure hence, enrolled his name among those of the Saints.\par
\switchcolumn*

\LA
\begin{Resp}\Rsign Invéni David servum meum, óleo sancto meo unxi eum: \Star{} Manus enim mea auxiliábitur ei, allelúja.\end{Resp}
\begin{Resp}\Vsign Nihil profíciet inimícus in eo, et fílius iniquitátis non nocébit ei.\end{Resp}
\begin{Resp}\Rsign Manus enim mea auxiliábitur ei, allelúja.\end{Resp}
\switchcolumn
\EN
\begin{Resp}\Rsign I have found David My servant, with My holy oil have I anointed him \Star{} For My hand shall help him, alleluia.\end{Resp}
\begin{Resp}\Vsign The enemy shall prevail nothing against him, nor the son of wickedness afflict him.\end{Resp}
\begin{Resp}\Rsign For My hand shall help him, alleluia.\end{Resp}
\switchcolumn*

\LA
\LessonHead{Lectio V}
\switchcolumn
\EN
\LessonHead{Lesson V}
\switchcolumn*

\LA
\LessonTitle{Ex libro Moralium sancti Gregorii Papæ.}
\switchcolumn
\EN
\LessonTitle{From the Book of Moral Reflections upon Job, written by Pope St. Gregory the Great}
\switchcolumn*

\LA
\Source{Lib. 10. cap. 16. in cap. 12. Job.}
\switchcolumn
\EN
\Source{(Bk. x Chap. xvi on Job xii)}
\switchcolumn*

\LA
\noindent Deridétur justi simplícitas. Hujus mundi sapiéntia est, cor machinatiónibus tégere, sensum verbis veláre: quæ falsa sunt, vera osténdere; quæ vera sunt, falsa demonstráre. Hæc nimírum prudéntia usu a juvénibus scitur hæc a púeris prétio díscitur; hanc qui sciunt, céteros despiciéndo supérbiunt: hanc qui nésciunt, subjécti et tímidi in áliis mirántur: quia ab eis hæc éadem duplicitátis iníquitas nómine palliáta dilígitur, dum mentis pervérsitas urbánitas vocátur. Hæc sibi obsequéntibus præcipit honórum cúlmina quærere, adépta temporális glóriæ vanitáte gaudére, irrogáta ab áliis mala multiplícius réddere: cum vires súppetunt, nullis resisténtibus cédere: cum virtútis possibílitas deest, quidquid explére per malítiam non valent, hoc in pacífica bonitáte simuláre.\par
\switchcolumn
\EN
\noindent The simplicity of the righteous is made a subject of derision. The wisdom of this world hideth our true feelings by artifice, and useth language to conceal our thoughts this is the wisdom which demonstrated the truth of falsehood, and showeth the falsehood of the truth. This kind of shrewdness the young acquire by practice, and children pay for the learning it. Those who are good at this look down upon their neighbours those who are bad at it are humble and timid, and wonder at it in others they regard this astuteness too, wrong though it be, with wistful admiration, under softened epithets. Un-straightforwardness is called good breeding. The principles of the world teach those who entertain them, to try and rise to distinction, and when they have attained the bubble of glory which is so soon to pass away, to feel it sweet to have at their feet them on whom they may wreak rich revenge. These principles teach a man, as long as he is strong enough, to give way to nobody else, and, if he hath no chance by force, to try and attain his object by diplomacy.\par
\switchcolumn*

\LA
\begin{Resp}\Rsign Pósui adjutórium super poténtem, et exaltávi eléctum de plebe mea: \Star{} Manus enim mea auxiliábitur ei, allelúja.\end{Resp}
\begin{Resp}\Vsign Invéni David servum meum, óleo sancto meo unxi eum.\end{Resp}
\begin{Resp}\Rsign Manus enim mea auxiliábitur ei, allelúja.\end{Resp}
\switchcolumn
\EN
\begin{Resp}\Rsign I have laid help upon one that is mighty, and have exalted one chosen out of My people \Star{} For My hand shall help him, alleluia.\end{Resp}
\begin{Resp}\Vsign I have found David My servant, with My holy oil have I anointed him.\end{Resp}
\begin{Resp}\Rsign For My hand shall help him, alleluia.\end{Resp}
\switchcolumn*

\LA
\LessonHead{Lectio VI}
\switchcolumn
\EN
\LessonHead{Lesson VI}
\switchcolumn*

\LA
\noindent At contra, sapiéntia justórum est: nil per ostensiónem fíngere, sensum verbis aperíre, vera ut sunt dilígere, falsa devitáre; bona gratis exhibére, mala libéntius toleráre quam fácere; nullam injúriæ ultiónem quærere, pro veritáte contuméliam lucrum putáre. Sed hæc justórum simplícitas deridétur; quia ab hujus mundi sapiéntibus puritátis virtus fatúitas créditur. Omne enim quod innocénter ágitur, ab eis proculdúbio stultum putátur; et quidquid in ópere véritas ápprobat, carnáli sapiéntiæ fátuum sonat. Quid namque stúltius vidétur mundo, quam mentem verbis osténdere, nil cállida machinatióne simuláre, nullas injúriis contumélias réddere, pro maledicéntibus oráre, paupertátem quærere, posséssa relínquere, rapiénti non resístere, percutiénti álteram maxíllam præbére?\par
\switchcolumn
\EN
\noindent The wisdom of the righteous is the contrary of all this. They seek to avoid deception, to give their thoughts a clear expression in their words, to love the truth because it is the truth, to avoid falsehood, and rather to suffer than to inflict evil. Such are they who seek not to avenge themselves for wrong, and deem it gain to be despised for the truth's sake. This their simplicity is made a subject of derision, for such as are wise in this world believe the purity of their virtue to be simple foolery. Whatsoever is done innocently, they consider without doubt stupid. Such works as the truth approveth are idiotic, when tried by carnal standards of wisdom. After all, what stupider thing is there in this world than to express our real thoughts in our words, to keep nothing quiet by skilful tact, to repay no injuries, to pray for them which curse us, to seek poverty, to give up property, to strive not with such as take from us, to turn the other cheek to the smiter\par
\switchcolumn*

\LA
\begin{Resp}\Rsign Iste est, qui ante Deum magnas virtútes operátus est, et omnis terra doctrína ejus repléta est: \Star{} Ipse intercédat pro peccátis ómnium populórum, allelúja.\end{Resp}
\begin{Resp}\Vsign Iste est, qui contémpsit vitam mundi, et pervénit ad cæléstia regna.\end{Resp}
\begin{Resp}\Rsign Ipse intercédat pro peccátis ómnium populórum, allelúja.\end{Resp}
\begin{Resp}\Vsign Glória Patri, et Fílio, \Star{} et Spirítui Sancto.\end{Resp}
\begin{Resp}\Rsign Ipse intercédat pro peccátis ómnium populórum, allelúja.\end{Resp}
\switchcolumn
\EN
\begin{Resp}\Rsign This is he which wrought great wonders before God, and the whole earth is full of his teaching \Star{} May he pray for all people, that their sins may be forgiven unto them, alleluia.\end{Resp}
\begin{Resp}\Vsign This is he which loved not his life in this world, and hath attained unto the kingdom of heaven.\end{Resp}
\begin{Resp}\Rsign May he pray for all people, that their sins may be forgiven unto them, alleluia.\end{Resp}
\begin{Resp}\Vsign Glory be to the Father, and to the Son, \Star{} and to the Holy Ghost.\end{Resp}
\begin{Resp}\Rsign May he pray for all people, that their sins may be forgiven unto them, alleluia.\end{Resp}
\switchcolumn*

\LA
\LessonHead{Lectio IX (pro commemoratione)}
\switchcolumn
\EN
\LessonHead{Lesson IX (when commemorated)}
\switchcolumn*

\LA
\LessonTitle{Commemoratio: S. Petri Celestini Papæ et Confessoris}
\switchcolumn
\EN
\LessonTitle{Commemoration of: St. Peter Celestine, Pope and Confessor}
\switchcolumn*

\LA
\noindent Petrus, a nómine quo Póntifex est appellátus, Cælestínus dictus, honéstis catholicísque paréntibus Æsérniæ in Samnítibus natus, adolescéntiam vix ingréssus, ut ánimum a mundi illécebris custodíret, in solitúdinem secéssit. Ibi contemplatiónibus mentem nútriens, corpus in servitútem rédigens, férream caténam ad nudam carnem adhibébat. Congregatiónem, quæ póstea, Cælestinórum dicta est, sub régula sancti Benedícti instítuit. Hinc, quasi lucérna supra candelábrum pósita, cum abscóndi nequíret, (Romána Ecclésia diu viduáta pastóre) in Petri Cáthedram ignórans et absens adscítus, magna novitátis admiratióne non minus quam repentíno gáudio cunctos affécit. Cum autem in pontificátus sublimitáte collocátus, váriis disténtus curis, assuétis incúmbere meditatiónibus vix posse cognósceret, óneri páriter et honóri voluntárie cessit. Indéque priscam vitæ ratiónem répetens, obdormívit in Dómino, ejúsque pretiósam mortem crux præfúlgens in áëre ante cubículi óstium réddidit ámplius gloriósam. Miráculis multis tam vivens quam post óbitum cláruit; quibus rite examinátis, Clemens quintus anno postquam decéssit undécimo, eum Sanctórum número adscrípsit.\par
\switchcolumn
\EN
\noindent Peter, who is called Peter Celestine, because when he became Pope he did so under the title of Celestine V, was the son of respectable Catholic parents, and was born at Isernia in Apulia. He was hardly entered on boyhood, when he withdrew into a desert, in order to keep his soul safe from the snares of the world. In solitude he fed his mind with heavenly meditation, and brought his body into subjection, even by wearing an iron chain next to his bare flesh. He founded, under the Rule of St. Benedict, that congregation which was afterwards known as the Celestines. His light, as of a candle set upon a candlestick, could not be kept hidden, and after the Church of Rome had for a long while been widowed of a shepherd, he was chosen without his knowledge and in his absence, to fill the chair of Peter. The news of his election filled himself with as great amazement, as it did all others with sudden joy. When, however, he was seated in the exalted place of the Papal dignity, he found that the many cares by which he was beset made it well-nigh impossible for him to give himself to his accustomed meditations; of his own free will, he resigned the burden and the honor together; and, while he sought to return to his old way of life, he fell asleep in the Lord. How precious his death was in his sight was gloriously manifested by a Cross which appeared shining in the air before the door of the cell. He was illustrious for miracles both during his life and after his death, and when these had been duly investigated, Clement V, in the eleventh year after his departure hence, enrolled his name among those of the Saints.\par
\switchcolumn*

\LA
\TeDeum{Te Deum}
\switchcolumn
\EN
\TeDeum{Te Deum}
\switchcolumn*

\end{paracol}

\par\needspace{10\baselineskip}\addvspace{1.2em}{\centering\small\textsc{Die 19 Maii} · \textsc{19 May}\par}
\Office{S. Pudentianæ Virginis}{St. Pudentiana, Virgin}{Simplex}
\addcontentsline{toc}{subsection}{Die 19 Maii · S. Pudentianæ Virginis \textit{— St. Pudentiana, Virgin}}
\begin{paracol}{2}
\LA
\LessonHead{Lectio IX (pro commemoratione)}
\switchcolumn
\EN
\LessonHead{Lesson IX (when commemorated)}
\switchcolumn*

\LA
\LessonTitle{Commemoratio: S. Pudentianæ Virginis}
\switchcolumn
\EN
\LessonTitle{Commemoration of: St. Pudentiana, Virgin}
\switchcolumn*

\LA
\noindent Pudentiana virgo, Pudentis Romani filia, parentibus orbata, cum admirabili pietate christianam religionem coleret, una cum sorore Praxede pecuniam ex vendito patrimonio redactam pauperibus distribuit, seque jejuniis et orationibus dedit. Cujus étiam opera, tota ejus familia, in qua erant nonaginta sex homines, a Pio Pontifice baptizata est. Quod autem ab Antonino imperatore sancitum erat, ne Christiani publice sacrificia facerent, Pius Pontifex in sedibus Pudentianæ cum Christianis sacra celebrabat. Quibus illa benigne acceptis, quæ ad vitam necessaria essent, suppeditabat Itaque in his christianæ pietatis officiis migravit e vita, et in sepulcro patris ad cemeterium Priscillæ via Salaria sepulta est, decimo quarto Kalendas Junii.\par
\switchcolumn
\EN
\noindent The maiden Pudentiana was the orphan daughter of Pudens the Roman Senator. She was a Christian of eminent godliness. She with her sister Praxedes distributed to the poor the money which they obtained by the sale of their inheritance. She gave herself continually to fasting and prayer. By her care the whole of the household, being ninety-six persons, were baptized by Pope Pius. Whereas the Emperor Antonine had forbidden the Christians to offer sacrifice in public, Pope Pius used to meet with them in Pudentiana's house, to celebrate the holy rites. She was a gracious hostess to them, and ministered to them in such things as are needful for the body. She thus busied herself in works of Christian godliness until she passed from this present life to a better. She was buried in her father's sepulchre in the cemetery of Priscilla on the Salarian Way upon the 19th day of May.\par
\switchcolumn*

\LA
\TeDeum{Te Deum}
\switchcolumn
\EN
\TeDeum{Te Deum}
\switchcolumn*

\end{paracol}

\par\needspace{10\baselineskip}\addvspace{1.2em}{\centering\small\textsc{Die 20 Maii} · \textsc{20 May}\par}
\Office{S. Bernardini Senensis Confessoris}{St. Bernardine of Siena, Confessor}{Semiduplex}
\addcontentsline{toc}{subsection}{Die 20 Maii · S. Bernardini Senensis Confessoris \textit{— St. Bernardine of Siena, Confessor}}
\begin{paracol}{2}
\LA
\RefLine{Lectiones I–III: de Scriptura occurrente.}
\switchcolumn
\EN
\RefLine{Lessons I–III: from the Scripture of the day.}
\switchcolumn*

\LA
\RefLine{Lectiones VII–IX: ex Commúni Confessoris non pontificis (C5).}
\switchcolumn
\EN
\RefLine{Lessons VII–IX: from the Common of a Confessor not a Bishop (C5).}
\switchcolumn*

\LA
\LessonHead{Lectio IV}
\switchcolumn
\EN
\LessonHead{Lesson IV}
\switchcolumn*

\LA
\noindent Bernardinus Albizesca, nobili Senénsi família ortus, ab ineunte ætate non obscura sanctitátis dedit indícia; nam a piis paréntibus honeste educatus, neglectis puerílibus ludis, inter prima grammaticæ stúdia pietátis opéribus ánimum inténdit, jejuniis, oratióni, et beatíssimæ Vírginis cultui præcipue addictus. Misericórdia vero in páuperes fuit insígnis. Quæ quidem ómnia procedénte témpore quo mélius posset excólere, eórum número adscribi vóluit, qui Senis in hospitali domo beátæ Maríæ de Scala Deo inserviunt, unde complures sanctitáte celebres viri prodiérunt. Ibi corporis afflictatióne et ægrotántium cura, dum atrox pestiléntia grassarétur, incredibili caritate sese exercuit. Inter ceteras autem virtútes castitátem, egregia forma repugnante, sanctíssime custodívit adeo ut eo præsénte nemo umquam, ne impudentíssimus quidem, verbum minus honestum proferre auderet.\par
\switchcolumn
\EN
\noindent This Bernardine was born of the noble family of the Albizeschi, in the Republic of Siena, (on the 8th of September, in the year 1380.) His saintliness began to manifest itself from his earliest years. He was well brought up by a godly father and mother, and even when he was being taught the first rudiments of worldly learning, he used to give up his play-time to occupy himself with devout works, being much drawn to fasting, prayer, and the devotion to the most Blessed Virgin. He abounded likewise in tenderness for the poor. As time went on, that he might the more entirely do these things, it was his will to enroll himself among those who work in the Hospital of Blessed Mary, called “of the Ladder,” at Sienna. There, during the raging of an horrible distemper, he laboured with marvellous charity and great bodily suffering, in serving the sick. In bodily presence he was a very goodly person, but, with all his other virtues, he kept ever so holy a guard over his purity, that it soon came to pass that no one, however shameless, dared to say an unseemly word in his presence.\par
\switchcolumn*

\LA
\begin{Resp}\Rsign Honéstum fecit illum Dóminus, et custodívit eum ab inimícis, et a seductóribus tutávit illum: \Star{} Et dedit illi claritátem ætérnam, allelúja.\end{Resp}
\begin{Resp}\Vsign Justum dedúxit Dóminus per vias rectas, et osténdit illi regnum Dei.\end{Resp}
\begin{Resp}\Rsign Et dedit illi claritátem ætérnam, allelúja.\end{Resp}
\switchcolumn
\EN
\begin{Resp}\Rsign The Lord made him honourable, and defended him from his enemies, and kept him safe from those that lay in wait for him; \Star{} And gave him perpetual glory, alleluia.\end{Resp}
\begin{Resp}\Vsign He went down with him into the pit, and left him not in bonds.\end{Resp}
\begin{Resp}\Rsign And gave him perpetual glory, alleluia.\end{Resp}
\switchcolumn*

\LA
\LessonHead{Lectio V}
\switchcolumn
\EN
\LessonHead{Lesson V}
\switchcolumn*

\LA
\noindent Gravi morbo tentatus, eoque ad quatuor menses patientíssime toleráto, demum incolumis, de religiosæ vitæ instituto capesséndo deliberare cœpit. Quo ut sibi viam muniret, ædiculam in extrema urbe conduxit; in quam cum sese abdidísset, asperrimam omni ex parte vitam trahebat, Deum assidue orans, ut, quid sibi sequéndum esset, osténderet. Quare divinitus factus est, ut beáti Francisci ordinem præ ceteris optaret, in quo humilitate, patiéntia aliisque religiosi hóminis virtútibus excelluit. Id cum cœnobii rector animadverteret, jamque antea Bernardini doctrinam et sacrárum litterárum perítiam perspectam haberet, prædicándi onus eidem impósuit: quo humillime suscepto, cum se minus idoneum agnosceret ob vocis exilitátem ac raucitátem, Dei ope implorata, non sine miraculo ejusmodi impediménto liberátus est.\par
\switchcolumn
\EN
\noindent He suffered a severe sickness, and when, after bearing it with the utmost patience, he recovered his health, he began to think of embracing some institute of the religious life. To make his way sure, he built a little hut in the outskirts of the city, where he hid himself and led a life of hardships of all kinds, continuing instant in prayer to God that He would be pleased to make clear to him what path he should follow. And so it came to pass by God's will that he chose the Order of Blessed Francis. In that Order he shone a bright instance of lowliness, long-suffering, and every other grace of a religious man. When the superior of his convent saw this, and had already considered what his teaching and knowledge of sacred learning were, he laid on Bernardine the duty of preaching. This the Saint humbly accepted, and finding that his usefulness was much impaired by his having a shrill, harsh voice, he betook him to implore the help of God, Who was pleased, not without a miracle, to free him from this drawback.\par
\switchcolumn*

\LA
\begin{Resp}\Rsign Amávit eum Dóminus, et ornávit eum: stolam glóriæ índuit eum, \Star{} Et ad portas paradísi coronávit eum, allelúja.\end{Resp}
\begin{Resp}\Vsign Índuit eum Dóminus lorícam fídei, et ornávit eum.\end{Resp}
\begin{Resp}\Rsign Et ad portas paradísi coronávit eum, allelúja.\end{Resp}
\switchcolumn
\EN
\begin{Resp}\Rsign The Lord loved him and beautified him; He clothed him with a robe of glory, \Star{} And crowned him at the gates of Paradise, alleluia.\end{Resp}
\begin{Resp}\Vsign The Lord hath put on him the breast-plate of faith, and hath adorned him.\end{Resp}
\begin{Resp}\Rsign And crowned him at the gates of Paradise, alleluia.\end{Resp}
\switchcolumn*

\LA
\LessonHead{Lectio VI}
\switchcolumn
\EN
\LessonHead{Lesson VI}
\switchcolumn*

\LA
\noindent Cumque ea témpora vitiis criminibúsque redundarent, et cruentis factiónibus in Italia, divina humanaque ómnia permixta essent, Bernardinus urbes atque oppida concursans, in nómine Jesu, quem semper in ore et in péctore gerebat, collapsam pietátem moresque verbo et exemplo magna ex parte restituit. Quo factum est, ut præcláræ civitátes eum sibi episcopum a summo Pontifice postularent; quod ille munus invicta humilitate constantíssime rejécit. Denique vir Dei imménsis labóribus exhaustus, multis magnisque editis miraculis, libris étiam pie docteque conscriptis, cum vixísset annos sex ac sexagínta, in urbe Aquila in Vestinis beato fine quiévit. Quem novis in dies coruscántem signis, anno post obitum sexto, Nicolaus quintus Pontifex maximus in Sanctórum númerum retulit.\par
\switchcolumn
\EN
\noindent Those were times fruitful in vices and crimes and the bloody civil wars which raged in Italy confounded all things Divine and human. Bernardine went through the cities and towns, and, in the Name of Jesus, that Name which he ever bore upon his lips and in his heart, he prevailed in great measure by his word and example, in setting up falling godliness and morality. Illustrious cities demanded him from the Pope as their Bishop, but this was an honour which his unconquerable humility caused him always steadily to refuse. At last the man of God, after untold labours, the working of many and great miracles, and the writing of godly and learned books, in the 67th year of his age, at Aquila in the Abruzzi, rested in a blessed death, (upon the 20th day of May 1444.) As the fame of new signs and wonders increased day by day, Pope Nicholas V., in the sixth year after his death, added his name to the roll of the Saints.\par
\switchcolumn*

\LA
\begin{Resp}\Rsign Iste homo perfécit ómnia quæ locútus est ei Deus, et dixit ad eum: Ingrédere in réquiem meam: \Star{} Quia te vidi justum coram me ex ómnibus géntibus, allelúja.\end{Resp}
\begin{Resp}\Vsign Iste est, qui contémpsit vitam mundi, et pervénit ad cæléstia regna.\end{Resp}
\begin{Resp}\Rsign Quia te vidi justum coram me ex ómnibus géntibus, allelúja.\end{Resp}
\begin{Resp}\Vsign Glória Patri, et Fílio, \Star{} et Spirítui Sancto.\end{Resp}
\begin{Resp}\Rsign Quia te vidi justum coram me ex ómnibus géntibus, allelúja.\end{Resp}
\switchcolumn
\EN
\begin{Resp}\Rsign This is he which did according unto all that God commanded him; and God said unto him: Enter thou into My rest, \Star{} For thee have I seen righteous before Me among all people, alleluia.\end{Resp}
\begin{Resp}\Vsign This is he which loved not his life in this world, and is come unto an everlasting kingdom.\end{Resp}
\begin{Resp}\Rsign For thee have I seen righteous before Me among all people, alleluia.\end{Resp}
\begin{Resp}\Vsign Glory be to the Father, and to the Son, \Star{} and to the Holy Ghost.\end{Resp}
\begin{Resp}\Rsign For thee have I seen righteous before Me among all people, alleluia.\end{Resp}
\switchcolumn*

\LA
\LessonHead{Lectio IX (pro commemoratione)}
\switchcolumn
\EN
\LessonHead{Lesson IX (when commemorated)}
\switchcolumn*

\LA
\LessonTitle{Commemoratio: S. Bernardini Senensis Confessoris}
\switchcolumn
\EN
\LessonTitle{Commemoration of: St. Bernardine of Siena, Confessor}
\switchcolumn*

\LA
\noindent Bernardínus Albizésca, nóbili Senénsi família ortus, inter prima grammáticæ stúdia, negléctis puerílibus ludis, pietátis opéribus ánimum inténdit, beátæ Vírginis cúltui præcípue addíctus. Caritáte et misericórdia in páuperes insígnis, eórum servítio Senis in hospitáli beátæ Maríæ de Scala se mancipávit. De capesséndo religiósæ vitæ institúto delíberans, Deo sic disponénte, beáti Francísci órdinem præ céteris optávit, in quo humilitáte, patiéntia, aliísque religiósi hóminis virtútibus excélluit. Prædicándi ónere a superióribus suscépto, cum se minus idóneum agnósceret ob vocis exilitátem et raucitátem, Dei ope imploráta, prodigióse ab hoc impediménto liberátus est. Urbes atque óppida concúrsans, in nómine Jesu, quem semper in ore et péctore gerébat, cívium ubíque discórdias exstínxit, et collápsam pietátem morésque verbo et exémplo magna ex parte restítuit. Libros pie doctéque conscrípsit. Plenus méritis et miráculis clarus, annos natus sex ac sexagínta, in urbe Aquila in Vestínis, beáto fine quiévit.\par
\switchcolumn
\EN
\noindent St. Bernardine Albizeschi was born of a noble family of Siena. Even in his first years in school, he turned away from children's games and applied himself to exercises of devotion, especially those honouring the Blessed Virgin. Outstanding for his charity and mercy to the poor, he gave himself over to serving them at the hospital Santa Maria della Scala in Siena. When he came to think of entering the religious life, divine Providence led him to choose, in preference to others, the Franciscan Order, where he excelled in humility, patience, and the other virtues of a religious. His superiors imposed on him the duty of preaching; and, although he knew that his voice was too weak and too hoarse for a preacher's, he accepted the charge and implored the help of God. As a result, he was marvellously freed of his handicap. Going about the towns and villages in the Name of Jesus, which was always carried on his lips and in his heart, he everywhere put an end to the dissensions of the citizens, and by his word and example did much to restore their piety and morality, which had fallen to a low level. He also wrote devout and learned books. At the age of sixty-six, rich in merit and famous for his miracles, he died a happy death at the town of L'Aquila in the Abruzzi.\par
\switchcolumn*

\LA
\TeDeum{Te Deum}
\switchcolumn
\EN
\TeDeum{Te Deum}
\switchcolumn*

\end{paracol}

\par\needspace{10\baselineskip}\addvspace{1.2em}{\centering\small\textsc{Die 25 Maii} · \textsc{25 May}\par}
\Office{S. Gregorii VII Papæ et Confessoris}{St. Gregory VII, Pope and Confessor}{Duplex}
\addcontentsline{toc}{subsection}{Die 25 Maii · S. Gregorii VII Papæ et Confessoris \textit{— St. Gregory VII, Pope and Confessor}}
\begin{paracol}{2}
\LA
\RefLine{Lectiones I–III: de Scriptura occurrente.}
\switchcolumn
\EN
\RefLine{Lessons I–III: from the Scripture of the day.}
\switchcolumn*

\LA
\RefLine{Lectiones VII–VIII: ex Commúni Summorum Pontificum Confessorum (C4b).}
\switchcolumn
\EN
\RefLine{Lessons VII–VIII: from the Common of Popes (C4b).}
\switchcolumn*

\LA
\RefLine{Lectio IX: de S. Urbanis Papæ et Martyris (05-25o).}
\switchcolumn
\EN
\RefLine{Lesson IX: of St. Urban, Pope and Martyr (05-25o).}
\switchcolumn*

\LA
\LessonHead{Lectio IV}
\switchcolumn
\EN
\LessonHead{Lesson IV}
\switchcolumn*

\LA
\noindent Gregórius Papa septimus, antea Hildebrandus, Suánæ in Etruria natus, doctrina, sanctitáte, omnique virtútum genere cum primis nobilis, mirifice universam Dei illustrávit Ecclésiam. Cum párvulus ad fabri ligna edolántis pedes, jam litterárum ínscius, lúderet, ex rejectis tamen segmentis illa Davidici eleménta oraculi, Dominábitur a mari usque ad mare, casu formasse narrátur: manum púeri ductante Númine, quo significarétur ejus fore amplíssimam in mundo auctoritátem. Romam deínde profectus, sub protectióne sancti Petri educátus est. Juvenis, Ecclésiæ libertátem a laicis oppressam ac depravatos ecclesiasticórum mores vehementius dolens, in Cluniacénsi monasterio, ubi sub regula sancti Benedicti austerioris vitæ observántia eo témpore maxime vigebat, monachi habitum induens, tanto pietátis ardore divinæ majestáti deserviébat, ut a sanctis ejusdem cœnobii pátribus prior sit eléctus. Sed, divina providéntia majora de eo disponente, in salútem plurimórum Cluníaco edúctus Hildebrandus, abbas primum monasterii sancti Pauli extra muros Urbis eléctus, ac postmodum Romanæ Ecclésiæ cardinalis creatus, sub Summis Pontificibus Leone nono, Victore secundo, Stephano nono, Nicolao secundo et Alexandro secundo, præcipuis munéribus et legatiónibus perfunctus est ; sanctíssimi et puríssimi consílii vir a beato Petro Damiáno nuncupatus. A Victore Papa secundo legátus a látere in Gálliam missus, Lugduni episcopum, simoníaca labe infectum, ad sui críminis confessiónem miraculo adegit. Berengárium in concílio Turonénsi ad iteratam hæresis abjuratiónem cómpulit. Cadalói quoque schisma sua virtúte compressit.\par
\switchcolumn
\EN
\noindent Hildebrand, who reigned as Pope under the name of Gregory VII, was born at Saona in Tuscany. By his teaching, by his holiness, and by his graces of all kinds, he was a noble light of the Church, whose brightness hath shone throughout all lands. There is a story to the effect that when he was a little child without any schooling, he was playing at the feet of a carpenter who was planing wood, and that God guided his hand to arrange the shavings which fell into the form of letters, making the inspired words of David, He shall have dominion from sea to sea, (Ps. lxxi. 8,) a fore-shadowing, as it were, of that wide lordship over the earth which was afterwards his. He was taken to Rome, and brought up under the shelter of St. Peter. As a young man he bitterly sorrowed over the oppression of the freedom of the Church by the laity, and over the corruption of the clergy themselves. He took the habit of a monk in the Abbey of Cluny, which was then in all the glory of the severest observance of the Rule of St. Benedict. There he served God's majesty with such warmth of earnestness that the saintly fathers of the convent chose him to be their Prior. But the Providence of God had greater things in store for him, whereby to make him a source of health to many, and he was brought away from Cluny. He was first elected Abbot of the monastery of St. Paul-without-the-walls at Rome, and afterwards created a Cardinal of the Roman Church. Under the Popes Leo IX, Victor II, Stephen IX, Nicholas II, and Alexander II, he discharged great offices of trust, and the duties of a Legate, and Blessed Peter Damian, speaking of him at this time, calleth him a man of most holy and honest thoughts. When Pope Victor II. sent him as his Legate into France, he, by a miracle, forced the Bishop of Lyons, who was befouled by the pollution of simony, to acknowledge his sin in the Council of Tours he wrung from Berenger a second abjuration of his heresy and he prevailed against the schism of Cadolaus, and strangled it.\par
\switchcolumn*

\LA
\begin{Resp}\Rsign Invéni David servum meum, óleo sancto meo unxi eum: \Star{} Manus enim mea auxiliábitur ei, allelúja.\end{Resp}
\begin{Resp}\Vsign Nihil profíciet inimícus in eo, et fílius iniquitátis non nocébit ei.\end{Resp}
\begin{Resp}\Rsign Manus enim mea auxiliábitur ei, allelúja.\end{Resp}
\switchcolumn
\EN
\begin{Resp}\Rsign I have found David My servant, with My holy oil have I anointed him \Star{} For My hand shall help him, alleluia.\end{Resp}
\begin{Resp}\Vsign The enemy shall prevail nothing against him, nor the son of wickedness afflict him.\end{Resp}
\begin{Resp}\Rsign For My hand shall help him, alleluia.\end{Resp}
\switchcolumn*

\LA
\LessonHead{Lectio V}
\switchcolumn
\EN
\LessonHead{Lesson V}
\switchcolumn*

\LA
\noindent Mortuo Alexandro secundo, invitus et mærens, unánimi ómnium consénsu, décimo Kalendas Maji anno Christi millesimo septuagesimo tertio, Summus Pontifex eléctus, sicut sol effulsit in domo Dei. Nam, potens opere et sermóne, ecclesiásticæ disciplinæ reparandæ, fidei propagandæ, libertáti Ecclésiæ restituéndæ, exstirpandis erróribus et corruptélis tanto studio incúbuit, ut ex Apostolórum ætate nullus Pontificum fuísse tradátur, qui majores pro Ecclésia Dei labóres molestiasque pertulerit, aut qui pro ejus libertáte acrius pugnaverit. Aliquot provincias a simoníaca labe expurgávit. Contra Henrici imperatóris ímpios conatus, fortis per ómnia athleta, impávidus permansit, seque pro muro dómui Israël pónere non tímuit ; ac eumdem Henricum, in profúndum malórum prolapsum, fidelium communióne regnosque privávit, atque súbditos pópulos fide ei data liberávit.\par
\switchcolumn
\EN
\noindent After the death of Alexander II, Hildebrand, against his own will and to his own grief, was, on the 22nd day of April, in the year of Christ 1073, chosen Pope by one common consent of all. Reigning as Gregory VII., he was as the sun shining upon the Temple of the Most High. (Ecclus. 1. 7.) Mighty both in word and deed, he toiled for the restoration of Ecclesiastical discipline, for the spread of the Faith, for the defence of the freedom of the Church, for the suppression of error and corruption, so that since the time of the Apostles there is said never to have been a Pope who bore more labour and trouble for the sake of God's Church, or contended more manfully for her liberties. He purged diverse provinces of the pollution of simony. Like a brave soldier he withstood without dread the unrighteous contendings of the Emperor Henry IV, against whom he shrank not from setting himself as a wall of defence for the house of Israel. And when the said Henry fell into the depths of sin he cut him off from the communion of the faithful, and from his kingdom, and loosed the nations that were subject to him from their sworn allegiance.\par
\switchcolumn*

\LA
\begin{Resp}\Rsign Pósui adjutórium super poténtem, et exaltávi eléctum de plebe mea: \Star{} Manus enim mea auxiliábitur ei, allelúja.\end{Resp}
\begin{Resp}\Vsign Invéni David servum meum, óleo sancto meo unxi eum.\end{Resp}
\begin{Resp}\Rsign Manus enim mea auxiliábitur ei, allelúja.\end{Resp}
\switchcolumn
\EN
\begin{Resp}\Rsign I have laid help upon one that is mighty, and have exalted one chosen out of My people \Star{} For My hand shall help him, alleluia.\end{Resp}
\begin{Resp}\Vsign I have found David My servant, with My holy oil have I anointed him.\end{Resp}
\begin{Resp}\Rsign For My hand shall help him, alleluia.\end{Resp}
\switchcolumn*

\LA
\LessonHead{Lectio VI}
\switchcolumn
\EN
\LessonHead{Lesson VI}
\switchcolumn*

\LA
\noindent Dum Missárum solemnia perágeret, visa est viris piis columba e cælo delapsa, humero ejus dextro ínsidens, alis exténsis caput ejus velare quo significátum est, Spíritus Sancti afflatu, non humanæ prudéntiæ ratiónibus ipsum duci in Ecclésiæ regímine. Cum ab iníqui Henrici exercitu Romæ gravi obsidióne premerétur, excitátum ab hostibus incendium signo crucis exstinxit. De ejus manu tandem a Roberto Guiscardo duce Northmanno ereptus, Cassinum se cóntulit ; atque inde Salernum ad dedicandam ecclésiam sancti Matthæi Apóstoli contendit. Cum aliquándo in ea civitáte sermónem habuísset ad pópulum, ærumnis confectus in morbum íncidit, quo se interitúrum præscívit. Postrema moriéntis Gregórii verba fuére: Diléxi justítiam et odívi iniquitátem, proptérea morior in exsílio. Innumerabília sunt, quæ vel fortiter sustinuit, vel multis coactis in Urbe Synodis sapiénter constítuit ; vir vere sanctus, criminum vindex, et acerrimus Ecclésiæ defensor. Exactis ítaque in pontificatu annis duodecim, migrávit in cælum anno salútis millesimo octogesimo quinto, plúribus in vita et post mortem miraculis clarus ; ejusque sacrum corpus in cathedrali basilica Salernitana est honorifice cónditum.\par
\switchcolumn
\EN
\noindent While he was celebrating solemn Mass, godly men saw a dove descend from heaven* perch upon his right shoulder, and spread out its wings so as to veil his head, a testimony that it was not by reasonings of man's wisdom, but by the teachings of the Holy Ghost, that he was guided in his rule over the Church. When the armies of the infamous Henry encompassed Rome, and hedged her in on every side, a great fire which the enemy had raised became extinct, when Gregory made the sign of the Cross towards it. The Norman Duke, Robert Guiscard, at length delivered Gregory from the hand of Henry, and he departed from Rome, first to the Abbey of Monte Cassino, and thence onward to Salerno, to dedicate the Church of St. Matthew the Apostle at that place. While he was preaching to the people there, on a certain day he was smitten with grievous pains, and fell into a sickness whereof he foresaw that he should never be healed. As he lay on his death-bed, Gregory's last words were I have loved righteousness and hated iniquity, and therefore I am dying in exile. He was a man really holy, a visitor of sin, and a most leal soldier of the Church. It is past reckoning how many sufferings he manfully bore, and how much he wisely ordained in many Councils, which he gathered together in Rome. He had been Pope twelve years, when, (on the 25th day of May,) in the year of salvation 1085, he went hence to be ever with the Lord. Both during his life and after his death he was marked by signs and wonders not a few. His holy body was honourably buried in the Cathedral Church of Salerno.\par
\switchcolumn*

\LA
\begin{Resp}\Rsign Iste est, qui ante Deum magnas virtútes operátus est, et omnis terra doctrína ejus repléta est: \Star{} Ipse intercédat pro peccátis ómnium populórum, allelúja.\end{Resp}
\begin{Resp}\Vsign Iste est, qui contémpsit vitam mundi, et pervénit ad cæléstia regna.\end{Resp}
\begin{Resp}\Rsign Ipse intercédat pro peccátis ómnium populórum, allelúja.\end{Resp}
\begin{Resp}\Vsign Glória Patri, et Fílio, \Star{} et Spirítui Sancto.\end{Resp}
\begin{Resp}\Rsign Ipse intercédat pro peccátis ómnium populórum, allelúja.\end{Resp}
\switchcolumn
\EN
\begin{Resp}\Rsign This is he which wrought great wonders before God, and the whole earth is full of his teaching \Star{} May he pray for all people, that their sins may be forgiven unto them, alleluia.\end{Resp}
\begin{Resp}\Vsign This is he which loved not his life in this world, and hath attained unto the kingdom of heaven.\end{Resp}
\begin{Resp}\Rsign May he pray for all people, that their sins may be forgiven unto them, alleluia.\end{Resp}
\begin{Resp}\Vsign Glory be to the Father, and to the Son, \Star{} and to the Holy Ghost.\end{Resp}
\begin{Resp}\Rsign May he pray for all people, that their sins may be forgiven unto them, alleluia.\end{Resp}
\switchcolumn*

\LA
\LessonHead{Lectio IX (pro commemoratione)}
\switchcolumn
\EN
\LessonHead{Lesson IX (when commemorated)}
\switchcolumn*

\LA
\LessonTitle{Commemoratio: S. Gregorii VII Papæ et Confessoris}
\switchcolumn
\EN
\LessonTitle{Commemoration of: St. Gregory VII, Pope and Confessor}
\switchcolumn*

\LA
\noindent Gregórius Papa séptimus, ántea Hildebrándus, Suánæ in Etrúria natus, doctrína, sanctitáte, omníque virtútum génere cum primis nóbilis, mirífice univérsam Dei Ecclésiam illustrávit. Adoléscens religiósi hábitum índuit in Cluniacénsi monastério, tantóque pietátis ardóre Deo desérviit, ut a sanctis ejúsdem cœnóbii pátribus prior fúerit eléctus. Póstea factus abbas monastérii sancti Pauli extra muros Urbis, ac deínde creátus Románæ Ecclésiæ cardinális, sub Summis Pontifícibus Leóne nono, Victóre secúndo, Stéphano nono, Nicoláo secúndo et Alexándro secúndo, præcípuis munéribus et legatiónibus perfúnctus est. Mórtuo Alexándro secúndo, unanímiter Summus Póntifex eléctus, ecclesiásticæ libertátis propugnátor ac defénsor acérrimus éxstitit; quaprópter multa passus, et Roma discédere coáctus est. Postréma ejus moriéntis verba fuére: Diléxi justítiam et odívi iniquitátem, proptérea mórior in exsílio. Migrávit in cælum, anno salútis millésimo octogésimo quinto, ejúsque corpus in cathedráli basílica Salernitána honorífice cónditum est.\par
\switchcolumn
\EN
\noindent Pope Gregory VII, the former Hildebrand, was born near Soana in Tuscany. As noble as any of the nobility in learning, in holiness and in every kind of virtue, he was a shining light to the whole Church of God. As a young man, he donned the religious habit at the monastery of Cluny, and served God with such zeal and devotion that he was chosen Prior by the holy religious of that monastery. Later, he was made Abbot of the monastery of St. Paul-outside-the-Walls, and then Cardinal of the Roman Church, performing noteworthy services and missions under Popes Leo IX, Victor II, Stephen IX, Nicholas II and Alexander II. At the death of Alexander, he was unanimously elected Pope, and stood out as a most zealous promoter and defender of the freedom of the Church, for which he suffered many things, even having to leave Rome. His last words, as he lay dying, were: “I have loved righteousness and hated iniquity, and therefore I am dying in exile.” He went to heaven in year of salvation 1085, and his body was buried with honour in the Cathedral of Salerno.\par
\switchcolumn*

\LA
\TeDeum{Te Deum}
\switchcolumn
\EN
\TeDeum{Te Deum}
\switchcolumn*

\end{paracol}

\par\needspace{10\baselineskip}\addvspace{1.2em}{\centering\small\textsc{Die 25 Maii} · \textsc{25 May}\par}
\Office{S. Urbanis Papæ et Martyris}{St. Urban, Pope and Martyr}{Simplex}
\addcontentsline{toc}{subsection}{Die 25 Maii · S. Urbanis Papæ et Martyris \textit{— St. Urban, Pope and Martyr}}
\begin{paracol}{2}
\LA
\LessonHead{Lectio IX (pro commemoratione)}
\switchcolumn
\EN
\LessonHead{Lesson IX (when commemorated)}
\switchcolumn*

\LA
\LessonTitle{Commemoratio: S. Urbanis Papæ et Martyris}
\switchcolumn
\EN

\switchcolumn*

\LA

\switchcolumn
\EN
\Source{For St. Urban, Pope and Martyr}
\switchcolumn*

\LA
\noindent Urbanus Romanus, Alexandro Sevéro imperatóre, doctrina et vitæ sanctitáte multos ad Christi fidem convértit; in illis Valerianum, beátæ Cæcíliæ sponsum, et Tiburtium, Valeriáni fratrem, qui póstea martyrium forti animo subiérunt. Hic de bonis Ecclésiæ attributis scripsit his verbis: Ipsæ res fidelium, quæ Dómino offerúntur, non debent in alios usus quam ecclesiásticos et christianórum fratrum, vel indigéntium, converti; quia vota sunt fidelium, et pretia peccatórum, ac patrimónia páuperum. Sedit annos sex, menses septem, dies quatuor: ac martyrio coronátus, sepúltus est in cœmeterio Prætextati, octavo Kalendas Junii. Ordinatiónibus quinque habitis mense Decembri, creávit presbyteros novem, diaconos quinque, episcopos per diversa loca octo.\par
\switchcolumn
\EN
\noindent This Urban was a Roman, who, in the reign of the Emperor Alexander Severus, by his teaching and holy life, brought many to believe in Christ. Among others was Valerian, the husband of the blessed Cecily, and Tiburtius, the brother of Valerian, both of whom afterwards bravely underwent martyrdom. It was Urban l who wrote the following words\par
concerning the property of the Church: Those things which His faithful ones\par
make offering of unto the Lord, must never be turned to any other use than those\par
of the Church, or of our Christian brethren, or of the poor. They are the\par
free-will offerings of faithful believers, the trespass offerings of sinners,\par
and the inheritance of the poor. He sat in the chair of Peter six years, seven\par
months, and four days, and being crowned with martyrdom, was buried in the\par
cemetery of Praetextatus, on the 25th day of May. He held five ordinations in\par
December, wherein he ordained nine Priests, five Deacons, and eight Bishops for\par
diverse places.\par
\switchcolumn*

\LA
\TeDeum{Te Deum}
\switchcolumn
\EN
\TeDeum{Te Deum}
\switchcolumn*

\end{paracol}

\par\needspace{10\baselineskip}\addvspace{1.2em}{\centering\small\textsc{Die 26 Maii} · \textsc{26 May}\par}
\Office{S. Philippi Neri Confessoris}{St. Philip Neri, Confessor}{Duplex}
\addcontentsline{toc}{subsection}{Die 26 Maii · S. Philippi Neri Confessoris \textit{— St. Philip Neri, Confessor}}
\begin{paracol}{2}
\LA
\RefLine{Lectiones I–III: de Scriptura occurrente.}
\switchcolumn
\EN
\RefLine{Lessons I–III: from the Scripture of the day.}
\switchcolumn*

\LA
\RefLine{Lectiones VII–VIII: ex Commúni Confessoris non pontificis (C5).}
\switchcolumn
\EN
\RefLine{Lessons VII–VIII: from the Common of a Confessor not a Bishop (C5).}
\switchcolumn*

\LA
\RefLine{Lectio IX: de S. Eleutherii Papæ et Martyris (05-26o).}
\switchcolumn
\EN
\RefLine{Lesson IX: of St. Eleutherius, Pope and Martyr (05-26o).}
\switchcolumn*

\LA
\LessonHead{Lectio IV}
\switchcolumn
\EN
\LessonHead{Lesson IV}
\switchcolumn*

\LA
\noindent Philippus Nerius, piis honestisque paréntibus Floréntiæ natus, ab ipsa ineunte ætate non obscura dedit futuræ sanctitátis indicia. Adoléscens, ampla patrui hereditáte dimissa, Romam se contulit; ubi philosophía ac sacris litteris eruditus, totum se Christo dicávit. Ea fuit abstinéntia, ut sæpe jejunus triduum permanserit. Vigiliis et oratiónibus intentus, septem Urbis Ecclésias frequenter visitans, apud cœmetérium Callísti in cæléstium rerum contemplatióne pernoctare consuevit. Sacerdos ex obediéntia factus, in animárum salúte procuranda totus fuit; et in confessiónibus audiéndis ad extremum usque diem perseverans, innumeros penne fílios Christo péperit; quos verbi Dei quotidiáno pabulo, sacramentórum frequéntia, oratiónis assiduitate, aliisque piis exercitatiónibus enutriri cupiens, Oratórii congregatiónem instituit.\par
\switchcolumn
\EN
\noindent Philip Neri was born of godly and respectable parents at Florence, (on the 23rd day of July, in the year of grace 1515) From his earliest childhood he gave signs of the holiness of life which he afterwards attained. As a young man he gave up the rich inheritance which would have come to him from his uncle, and went to Rome, where, in the study of philosophy and theology, he gave himself altogether to Christ. His self-control was such that he sometimes fasted from all food for three days at a time. He was instant in watching and prayer, and during the frequent pilgrimages which he made to the seven churches of Rome, it was his custom to remain all night in prayer to God in the Catacomb of Kallistus. (On the 23rd of May, 1551) he became a Priest in obedience to the advice of his Confessor, and afterwards made the salvation of souls the one object of his existence; he heard confessions with unwearied tenderness until his dying day, and became the spiritual father in Christ of many sons, whom it was his beloved work to feed day by day upon the Word of God, upon the often receiving of the Sacraments, upon instant prayer, and upon other godly works: to the which end he founded the Congregation of the Oratory.\par
\switchcolumn*

\LA
\begin{Resp}\Rsign Honéstum fecit illum Dóminus, et custodívit eum ab inimícis, et a seductóribus tutávit illum: \Star{} Et dedit illi claritátem ætérnam, allelúja.\end{Resp}
\begin{Resp}\Vsign Justum dedúxit Dóminus per vias rectas, et osténdit illi regnum Dei.\end{Resp}
\begin{Resp}\Rsign Et dedit illi claritátem ætérnam, allelúja.\end{Resp}
\switchcolumn
\EN
\begin{Resp}\Rsign The Lord made him honourable, and defended him from his enemies, and kept him safe from those that lay in wait for him; \Star{} And gave him perpetual glory, alleluia.\end{Resp}
\begin{Resp}\Vsign He went down with him into the pit, and left him not in bonds.\end{Resp}
\begin{Resp}\Rsign And gave him perpetual glory, alleluia.\end{Resp}
\switchcolumn*

\LA
\LessonHead{Lectio V}
\switchcolumn
\EN
\LessonHead{Lesson V}
\switchcolumn*

\LA
\noindent Caritate Dei vulnerátus languebat jugiter, tantoque cor ejus æstuábat ardore, ut, cum intra fines suos contineri non posset, illíus sinum, confractis atque elátis duabus cóstulis, mirabíliter Dóminus ampliaverit. Sacrum vero fáciens aut ferventius orans, in áëra quandoque sublatus, mira undique luce fulgere visus fuit. Egenos et páuperes omni caritátis officio prosequebátur: dignus, qui et Angelo in spécie páuperis eleemosynam erogaret; et, dum egéntibus noctu panem deferret, in fóveam lapsus, inde páriter ab Angelo incolumis eriperétur. Humilitáti addictus, ab honóribus semper abhorruit, atque ecclesiásticas dignitates, étiam primárias, non semel ultro delatas, constantíssime recusávit.\par
\switchcolumn
\EN
\noindent He was full of the love of God, and his heart was so hot therewith, that it became straitened in its place, and the Lord was pleased to ease him by(the gristle which joined) the fourth and fifth ribs (on his left side) breaking, and so allowing more play to the internal organs. Sometimes, when he was saying Mass, or more intent than usual in prayer, he was seen to be raised off the ground, and become, in a strange manner, all shining. He was ever ready to succour the poor and needy with kindly services, in which works God was pleased to make him meet once to give alms to an Angel, and again, when once by night he was carrying bread to the hungry, and was fallen into a pit, an Angel drew him out unhurt. He longed to be lowly, and always shrank from honours, and from dignities in the Church, whereof several of the highest were diverse times offered to him, but he always firmly refused them.\par
\switchcolumn*

\LA
\begin{Resp}\Rsign Amávit eum Dóminus, et ornávit eum: stolam glóriæ índuit eum, \Star{} Et ad portas paradísi coronávit eum, allelúja.\end{Resp}
\begin{Resp}\Vsign Índuit eum Dóminus lorícam fídei, et ornávit eum.\end{Resp}
\begin{Resp}\Rsign Et ad portas paradísi coronávit eum, allelúja.\end{Resp}
\switchcolumn
\EN
\begin{Resp}\Rsign The Lord loved him and beautified him; He clothed him with a robe of glory, \Star{} And crowned him at the gates of Paradise, alleluia.\end{Resp}
\begin{Resp}\Vsign The Lord hath put on him the breast-plate of faith, and hath adorned him.\end{Resp}
\begin{Resp}\Rsign And crowned him at the gates of Paradise, alleluia.\end{Resp}
\switchcolumn*

\LA
\LessonHead{Lectio VI}
\switchcolumn
\EN
\LessonHead{Lesson VI}
\switchcolumn*

\LA
\noindent Prophetíæ dono fuit illustris, et in animórum sensibus penetrandis mirifice enituit. Virginitátem perpetuo illibatam servávit; idque assecutus est, ut eos qui puritátem cólerent, ex odore, qui vero secus, ex fœtóre dignosceret. Abséntibus interdum appáruit, iisque periclitántibus opem tulit. Ægrotos plurimos et morti próximos sanitáti restituit. Mortuum quoque ad vitam revocávit. Cælestium spirituum et ipsíus Deiparæ Vírginis frequenter fuit apparitióne dignátus, ac plurimórum ánimas splendore circumfusas in cælum conscéndere vidit. Denique, anno salútis millesimo quingentésimo nonagesimo quinto, octavo Kalendas Junias, in quem diem inciderat festum Corporis Christi, Sacro maxima spíritus exsultatióne peracto, ceterisque functiónibus expletis, post mediam noctem, qua prædixerat hora, octogenarius obdormívit in Dómino. Quem Gregórius décimus quintus, miraculis clarum, in Sanctórum númerum rétulit.\par
\switchcolumn
\EN
\noindent He was illustrious for the gift of prophecy, and had a marked and very wonderful power of reading the thoughts of men's souls. He ever kept his own virginity undefiled, and could distinguish those that were pureminded by a sort of sweet savour, and the unclean, on the contrary, by a kind of stench. He sometimes appeared in double to persons at a distance, and brought them help when they were in peril. He healed many that were sick and dying. He also raised one dead man to life. He was honoured by seeing several times heavenly spirits, and likewise the Virgin Mother of God herself. He saw the souls of diverse persons, radiant with glory, ascend to heaven. In the year of salvation 1595 the Feast of the Body of Christ fell upon the 25th day of May. Philip, on that day, said Mass with extraordinary gladness of spirit, and performed the other religious works of the day, and after the hour of midnight, at the time he had himself foretold, he fell asleep in the Lord, in the eighty-second year of his life. Gregory XV., finding that God had glorified him by many miracles, enrolled his name among those of the saints.\par
\switchcolumn*

\LA
\begin{Resp}\Rsign Iste homo perfécit ómnia quæ locútus est ei Deus, et dixit ad eum: Ingrédere in réquiem meam: \Star{} Quia te vidi justum coram me ex ómnibus géntibus, allelúja.\end{Resp}
\begin{Resp}\Vsign Iste est, qui contémpsit vitam mundi, et pervénit ad cæléstia regna.\end{Resp}
\begin{Resp}\Rsign Quia te vidi justum coram me ex ómnibus géntibus, allelúja.\end{Resp}
\begin{Resp}\Vsign Glória Patri, et Fílio, \Star{} et Spirítui Sancto.\end{Resp}
\begin{Resp}\Rsign Quia te vidi justum coram me ex ómnibus géntibus, allelúja.\end{Resp}
\switchcolumn
\EN
\begin{Resp}\Rsign This is he which did according unto all that God commanded him; and God said unto him: Enter thou into My rest, \Star{} For thee have I seen righteous before Me among all people, alleluia.\end{Resp}
\begin{Resp}\Vsign This is he which loved not his life in this world, and is come unto an everlasting kingdom.\end{Resp}
\begin{Resp}\Rsign For thee have I seen righteous before Me among all people, alleluia.\end{Resp}
\begin{Resp}\Vsign Glory be to the Father, and to the Son, \Star{} and to the Holy Ghost.\end{Resp}
\begin{Resp}\Rsign For thee have I seen righteous before Me among all people, alleluia.\end{Resp}
\switchcolumn*

\LA
\LessonHead{Lectio IX (pro commemoratione)}
\switchcolumn
\EN
\LessonHead{Lesson IX (when commemorated)}
\switchcolumn*

\LA
\LessonTitle{Commemoratio: S. Philippi Neri Confessoris}
\switchcolumn
\EN
\LessonTitle{Commemoration of: St. Philip Neri, Confessor}
\switchcolumn*

\LA
\noindent Philíppus Nérius, piis honestísque paréntibus Floréntiæ natus, ampla pátrui hereditáte dimíssa, Romam se cóntulit; ubi, philosophía ac sacris lítteris erudítus, totum se Christo dicávit. Sacérdos ex obediéntia factus, in animárum salúte procuránda totus fuit, et in confessiónibus audiéndis ad extrémum usque diem persevérans, innúmeros pene fílios Christo péperit; quos verbi Dei cotidiáno pábulo, sacramentórum frequéntia, oratiónis assiduitáte aliísque piis exercitatiónibus enutríri cúpiens, Oratórii congregatiónem instítuit. Caritáte Dei vulnerátum tanto cor ejus æstuábat ardóre, ut, cum intra fines suos continéri non posset, sinum, confráctis atque elátis duábus cóstulis, mirabíliter Dóminus ampliáverit. Prophetíæ dono fuit illústris, et in animórum sénsibus penetrándis mirífice enítuit. Virginitátem perpétuo illibátam servávit; idque assecútus est, ut eos, qui puritátem cólerent, ex odóre, qui vero secus, ex fœtóre dignósceret. Anno salútis millésimo quingentésimo nonagésimo quinto, octogenárius obdormívit in Dómino.\par
\switchcolumn
\EN
\noindent Philip Neri was born in Florence of good and devout parents. Giving up a large inheritance from his uncle, he went to Rome, where he studied philosophy and the sacred sciences and dedicated himself wholly to Christ. He became a priest out of obedience and gave himself up completely to the saving of souls. Through hearing confessions, in which he persevered to the last day of his life, he brought forth innumerable sons for Christ. Desiring to nourish them with the daily bread of God's word, with frequent reception of the sacraments, with constant prayer, and with other exercises of piety, he founded the Congregation of the Oratory. His heart was wounded by the love of God, burning with such ardour that it could only be contained within his breast because the Lord miraculously enlarged the breast by breaking two of his ribs, and forming an arch over the heart. Philip was famed for the gift of prophecy and for his wonderful penetration of the thoughts of men's hearts. He kept his virginity always intact; and he had the gift of distinguishing those who cultivated purity by a good odour, and those who did not by a stench. At the age of eighty, in the year of salvation 1595, he fell asleep in the Lord.\par
\switchcolumn*

\LA
\TeDeum{Te Deum}
\switchcolumn
\EN
\TeDeum{Te Deum}
\switchcolumn*

\end{paracol}

\par\needspace{10\baselineskip}\addvspace{1.2em}{\centering\small\textsc{Die 26 Maii} · \textsc{26 May}\par}
\Office{S. Eleutherii Papæ et Martyris}{St. Eleutherius, Pope and Martyr}{Simplex}
\addcontentsline{toc}{subsection}{Die 26 Maii · S. Eleutherii Papæ et Martyris \textit{— St. Eleutherius, Pope and Martyr}}
\begin{paracol}{2}
\LA
\LessonHead{Lectio IX (pro commemoratione)}
\switchcolumn
\EN
\LessonHead{Lesson IX (when commemorated)}
\switchcolumn*

\LA
\LessonTitle{Commemoratio: S. Eleutherii Papæ et Martyris}
\switchcolumn
\EN
\LessonTitle{Commemoration of: St. Eleutherius, Pope and Martyr}
\switchcolumn*

\LA
\noindent Eleutherius, Nicopoli in Græcia natus, Aniceti Pontificis diaconus, Commodo imperatóre, præfuit Ecclésiæ. Huic, inítio pontificatus, supplices litteræ venérunta Lucio Britannórum rege, ut se ac suos in Christianórum númerum reciperet. Quam ob rem Fugatium et Damianum, doctos et pios viros, misit in Britanniam, per quos rex et réliqui fidem susciperent. Hoc Pontifice, Irenæus Polycarpi discipulus, Romam veniens, ab eo benígne acceptus est. Quo témpore summa pace et quiete fruebátur Ecclésia Dei; ac per totum orbem terrárum, maxime Romæ, fides propagabátur. Vixit Eleutherius in pontificatu annos quindecim, dies viginti tres. Fecit ordinatiónes tres mense Decembri, quibus creávit presbyteros duodecim, diaconos octo, episcopos per diversa loca quindecim: sepultúsque est in Vaticano prope corpus sancti Petri.\par
\switchcolumn
\EN
\noindent Eleutherius was a Greek by race, and was born at Nicopolis, a city of Epirus. His father's name was Abundius. He was a Priest of the holy Roman Church. In the year of our Lord 179, during the reign of the Emperor Marcus Aurelius Augustus, after the death of Soter, he was chosen Bishop of Rome by the votes of all the clergy. He discharged the duties of this office excellently, and with all praise, for fifteen years and twenty-three days. He held three ordinations in the month of December, wherein he ordained twelve Priests, eight Deacons, and fifteen Bishops for diverse places. He was consulted by the church of Lyons by letter concerning certain questions. The bearer of these letters was that most learned Irenaeus. Him he hospitably welcomed, and to him he opened the traditions of the Apostles which the Church of Rome had kept pure. He condemned the superstitious dry-meat system of the Montanists. He laid down excellent rules as to the right form of church sentences. When Marcion and Valentine had oftentimes relapsed he cast them out of the Church. In his days the Church enjoyed the utmost peace, and he brought many even of the chiefest of Rome to believe in Christ. He received letters by messengers from Lleurwg, Prince of the Britons (of Morganwg,) praying him for ministers of the Word of God, and he sent unto him Fagan and Dyfan, Priests of the Roman Church, through whose hands the Prince himself, with his whole household and nearly all his subjects, were born again in the sacred washing of regeneration. At length, when he had done all these things and others for the worship of God, Eleutherius died a holy death upon the 28th day of May, (in the year of our Lord 192,) in the reign of the Emperor Commodus, and was buried upon the Vatican Mount.\par
\switchcolumn*

\LA
\TeDeum{Te Deum}
\switchcolumn
\EN
\TeDeum{Te Deum}
\switchcolumn*

\end{paracol}

\par\needspace{10\baselineskip}\addvspace{1.2em}{\centering\small\textsc{Die 27 Maii} · \textsc{27 May}\par}
\Office{S. Bedæ Venerabilis Confessoris et Ecclesiæ Doctoris}{St. Bede the Venerable, Confessor and Doctor of the Church}{Duplex}
\addcontentsline{toc}{subsection}{Die 27 Maii · S. Bedæ Venerabilis Confessoris et Ecclesiæ Doctoris \textit{— St. Bede the Venerable, Confessor and Doctor of the Church}}
\begin{paracol}{2}
\LA
\RefLine{Lectiones I–III: de Scriptura occurrente.}
\switchcolumn
\EN
\RefLine{Lessons I–III: from the Scripture of the day.}
\switchcolumn*

\LA
\RefLine{Lectio IX: de S. Joannis Papæ et Martyris (05-27o).}
\switchcolumn
\EN
\RefLine{Lesson IX: of St. John I, Pope and Martyr (05-27o).}
\switchcolumn*

\LA
\LessonHead{Lectio IV}
\switchcolumn
\EN
\LessonHead{Lesson IV}
\switchcolumn*

\LA
\noindent Beda presbyter, Girvi in Britanniæ et Scotiæ finibus ortus, septennis sancto Benedicto Biscopio abbati Wiremuthensi educandus traditur. Monachus deínde factus, vitam sic instituit, ut, dum se artium et doctrinarum studiis totum impenderet, nihil umquam de regulari disciplina remitteret. Nullum fuit doctrinæ genus, in quo non esset diligentissime versatus; sed præcipua illi cura divinarum Scripturarum meditatio, quarum sententiam ut plenius assequeretur, Græci Hebraicique sermonis notitiam est adeptus. Trigesimo ætatis anno abbatis sui jussu sacerdos initiatus, statim, suasore Acca Hagulstadensi episcopo sacros explanare libros aggressus est: in quo sanctorum Patrum doctrinis adeo inhæsit, ut nihil proferret nisi illorum judicio comprobatum, eorumdem étiam fere verbis usus. Otium perosus semper, ex lectione ad orationem transibat ac vicissim ex oratione ad lectionem: in qua adeo animo inflammabatur, ut sæpe inter legendum et docendum lacrimis perfunderetur. Ne autem rerum fluxarum curis distraheretur, delatum abbatis munus constantissime detrectavit.\par
\switchcolumn
\EN
\noindent Bede, a priest, was born at Jarrow, on the borders of England and Scotland. At the age of seven years he was placed under the care of holy Benedict Biscop, Abbot of Wearmouth, to be educated. Thereafter he became a monk, and so ordered his life that, whilst he should devote himself wholly to the study of the sciences and of doctrine, he might in nothing relax the discipline of his Order. There was no branch of learning in which he was not most thoroughly versed, but his chief care was the study of Holy Scriptures; and that he might the better understand them he acquired a knowledge of the Greek and Hebrew tongues. When he was thirty years of age he was ordained priest at the command of his Abbot, and immediately, on the advice of Acca, Bishop of Hexham, undertook the work of expounding the Sacred Books. In his interpretations he so strictly adhered to the teaching of the holy Fathers that he would advance nothing which was not approved by their judgment, nay, had the warrant of their very words. He ever hated sloth, and by habitually passing from reading to prayer, and in turn from prayer to reading, he so inflamed his soul that often amid his reading and teaching he was bathed in tears. Lest also his mind should be distracted by the cares of transitory things, he never would take the office of Abbot when it was offered to him.\par
\switchcolumn*

\LA
\begin{Resp}\Rsign Honéstum fecit illum Dóminus, et custodívit eum ab inimícis, et a seductóribus tutávit illum: \Star{} Et dedit illi claritátem ætérnam, allelúja.\end{Resp}
\begin{Resp}\Vsign Justum dedúxit Dóminus per vias rectas, et osténdit illi regnum Dei.\end{Resp}
\begin{Resp}\Rsign Et dedit illi claritátem ætérnam, allelúja.\end{Resp}
\switchcolumn
\EN
\begin{Resp}\Rsign The Lord made him honourable, and defended him from his enemies, and kept him safe from those that lay in wait for him; \Star{} And gave him perpetual glory, alleluia.\end{Resp}
\begin{Resp}\Vsign He went down with him into the pit, and left him not in bonds.\end{Resp}
\begin{Resp}\Rsign And gave him perpetual glory, alleluia.\end{Resp}
\switchcolumn*

\LA
\LessonHead{Lectio V}
\switchcolumn
\EN
\LessonHead{Lesson V}
\switchcolumn*

\LA
\noindent Scientiæ ac pietatis laude Bedæ nomen sic brevi claruit, ut sanctus Sergius Papa de eo Romam arcessendo cogitaverit; quo difficillimis scilicet, quæ de rebus sacris exortæ erant, quæstionibus definiendis conferret operam. Emendandis fidelium moribus, fidei vindicandæ atque asserendæ libros plures conscripsit, quibus tantam sui apud omnes opinionem fecit, ut illum sanctus Bonifatius episcopus et martyr, Ecclesiæ lumen prædicaverit; Lanfrancus, Anglorum doctorem; concilium Aquisgranense, doctorem admirabilem dixerit. Quin ejus scripta eo adhuc vivente, publice in ecclesiis legebantur. Quod cum fieret, quoniam ipsum sanctum minime appellare liceret, Venerabilis titulo efferebant; qui deínde veluti proprius secutis étiam temporibus semper habitus est. Ejus autem doctrinæ eo vis efficacior erat, quod vitæ sanctimonia religiosisque virtutibus confirmabatur. Quam ob rem discipulos, quos multos et egregios imbuendos habuit, studio et exemplo non litteris modo atque scientiis, sed étiam sanctitate fecit insignes.\par
\switchcolumn
\EN
\noindent The name of Bede soon became so famous for learning and piety that St. Sergius the Pope thought of calling him to Rome, where, certainly, he might have helped to solve the very difficult questions which had then arisen concerning sacred things. He wrote many books for the bettering of the lives of the faithful, and defending and extending of the faith. By those he gained everywhere such a reputation that the holy martyr Bishop Boniface styled him a Light of the Church; Lanfranc called him The Teacher of the English, and the Council of Aix-la-Chapelle The Admirable Doctor. But as his writings were publicly read in the churches during his life, and as it was not allowable to call him already a saint, they named him The Venerable, a title which in all times after has remained peculiarly his. The power of his teaching was the greater also, in that it was attested by a holy life and the graces of religious observance. In this way, by his earnestness and example, his disciples, who were many and distinguished, were made eminent, not only in letters and the sciences, but in personal holiness.\par
\switchcolumn*

\LA
\begin{Resp}\Rsign Amávit eum Dóminus, et ornávit eum: stolam glóriæ índuit eum, \Star{} Et ad portas paradísi coronávit eum, allelúja.\end{Resp}
\begin{Resp}\Vsign Índuit eum Dóminus lorícam fídei, et ornávit eum.\end{Resp}
\begin{Resp}\Rsign Et ad portas paradísi coronávit eum, allelúja.\end{Resp}
\switchcolumn
\EN
\begin{Resp}\Rsign The Lord loved him and beautified him; He clothed him with a robe of glory, \Star{} And crowned him at the gates of Paradise, alleluia.\end{Resp}
\begin{Resp}\Vsign The Lord hath put on him the breast-plate of faith, and hath adorned him.\end{Resp}
\begin{Resp}\Rsign And crowned him at the gates of Paradise, alleluia.\end{Resp}
\switchcolumn*

\LA
\LessonHead{Lectio VI}
\switchcolumn
\EN
\LessonHead{Lesson VI}
\switchcolumn*

\LA
\noindent Aetate demum et laboribus fractus, gravi morbo correptus est. Quo cum amplius quinquaginta dies detentus esset, consuetum orandi morem Scripturasque interpretandi non intercepit; eo namque tempore Evangelium Joannis in popularium suorum usum Anglice vertit. Cum autem in Ascensionis præludio instare sibi mortem persentiret, supremis Ecclesiæ sacramentis muniri voluit; tum, sodales amplexatus, atque humi super cilicio stratus, cum illa verba ingeminaret, Gloria Patri, et Filio, et Spiritui Sancto, obdormivit in Domino. Ejus corpus suavissimum uti fertur, spirans odorem sepultum est in monasterio Girvensi, ac postea Dunclinum cum sancti Cuthberti, reliquiis translatum. Eum tamquam Doctorem a Benedictinis aliisque religiosis familiis ac diœcesibus cultum, Leo decimus tertius Pontifex maximus ex sacrorum Rituum Congregationis consulto, universalis Ecclesiæ Doctorem declaravit, et festo ipsius die Missam et Officium de Doctoribus ab omnibus recitari decrevit.\par
\switchcolumn
\EN
\noindent Broken at length by age and labour, he was seized by a grievous illness. Though he suffered under it for more than seven weeks, he ceased not from his prayers and his interpreting of the Scriptures; for at that time he was turning the Gospel of John into English for the use of his people. But when, on the Eve of the Ascension, he perceived that death was coming upon him, he desired to be fortified with the last sacraments of the Church: then, after he had embraced his companions, and was laid on a piece of sackcloth on the ground, he repeated the words, Glory be to the Father, and to the Son, and to the Holy Ghost, and fell asleep in the Lord. His body, very sweet, as it is related, breathing sweet odour, was buried in the monastery of Jarrow, and afterwards was translated to Durham with the relics of St. Cuthbert. Bede, who was already a Doctor among the Benedictines, and in other religious Orders, and venerated in certain dioceses, was declared by Pope Leo XIII., after consulting with the Congregation of Sacred Rites, to be a Doctor of the universal Church; and the Mass and Office for Doctors was ordered to be recited by all on his feast-day.\par
\switchcolumn*

\LA
\begin{Resp}\Rsign Iste homo perfécit ómnia quæ locútus est ei Deus, et dixit ad eum: Ingrédere in réquiem meam: \Star{} Quia te vidi justum coram me ex ómnibus géntibus, allelúja.\end{Resp}
\begin{Resp}\Vsign Iste est, qui contémpsit vitam mundi, et pervénit ad cæléstia regna.\end{Resp}
\begin{Resp}\Rsign Quia te vidi justum coram me ex ómnibus géntibus, allelúja.\end{Resp}
\begin{Resp}\Vsign Glória Patri, et Fílio, \Star{} et Spirítui Sancto.\end{Resp}
\begin{Resp}\Rsign Quia te vidi justum coram me ex ómnibus géntibus, allelúja.\end{Resp}
\switchcolumn
\EN
\begin{Resp}\Rsign This is he which did according unto all that God commanded him; and God said unto him: Enter thou into My rest, \Star{} For thee have I seen righteous before Me among all people, alleluia.\end{Resp}
\begin{Resp}\Vsign This is he which loved not his life in this world, and is come unto an everlasting kingdom.\end{Resp}
\begin{Resp}\Rsign For thee have I seen righteous before Me among all people, alleluia.\end{Resp}
\begin{Resp}\Vsign Glory be to the Father, and to the Son, \Star{} and to the Holy Ghost.\end{Resp}
\begin{Resp}\Rsign For thee have I seen righteous before Me among all people, alleluia.\end{Resp}
\switchcolumn*

\LA
\LessonHead{Lectio VII}
\switchcolumn
\EN
\LessonHead{Lesson VII}
\switchcolumn*

\LA
\LessonTitle{Léctio sancti Evangélii secúndum Matthǽum}
\switchcolumn
\EN
\LessonTitle{From the Holy Gospel according to Matthew}
\switchcolumn*

\LA
\Cite{Matt 5:13-19}
\switchcolumn
\EN
\Cite{Matt 5:13-19}
\switchcolumn*

\LA
\noindent In illo témpore: Dixit Jesus discipulis suis: Vos estis sal terræ. Quod si sal evanuerit, in quo salietur? Et réliqua.\par
\switchcolumn
\EN
\noindent AT that time Jesus said to His disciples: Ye are the salt of the earth. But if the salt have lost its savour, wherewith shall it be salted? And so on.\par
\switchcolumn*

\LA
\LessonTitle{Homilía sancti Bedæ Venerábilis Presbýteri}
\switchcolumn
\EN
\LessonTitle{Homily of the Venerable Bede, Priest.}
\switchcolumn*

\LA
\Source{In Evang. Vos estis sal terra}
\switchcolumn
\EN
\Source{On the Gospel}
\switchcolumn*

\LA
\noindent In terra, humana natura; in sale, sapientia verbis significatur. Salis enim natura, terra efficitur infructuosa; unde quasdam urbes legimus, victorum ira, sale seminatas. Et hoc convenit apostolicæ doctrinæ, ut sale sapientiæ compescat in terra humanæ carnis luxum sæculi aut fœditatem vitiorum germinare. Quod si sal evanuerit, in quo salietur? Id est, si vos, per quos condiendi sunt populi, propter metum persecutionum, aut terrorem, amiseritis regna cælorum, extra Ecclesiam positi, inimicorum opprobria sustinetis non dubium. Vos estis lux mundi: id est, vos, quia vera luce illuminati estis, lux eis qui in mundo sunt, esse debetis. Non potest civitas abscondi supra montem posita: id est, apostolica doctrina super Christum fundata, sive Ecclesia super Christum ex multis gentibus fidei unitate constructa et caritatis bitumine conglutinata; quæ sit tuta intrantibus; et laboriosa adeuntibus, habitatores custodit, et omnes inimicos secludit.\par
\switchcolumn
\EN
\noindent The Gospel saith, Ye are the salt of the earth. In these words the earth signifies human nature, and the salt signifies wisdom. Salt, verily, by its nature renders the earth unfruitful. Hence we read of cities, which in the anger of their victors were sown with salt. And hereto agreeth the teaching of the Apostle that by the salt of wisdom the lust of this world is restrained in the earth of human flesh, lest the foulness of vice should sprout up. But what if the salt shall have lost its savour? That is to say If you, by whom the people are to be seasoned, are, on account of fear of persecution, or terror, you should lose the kingdom of heaven, placed outside the Church, there is no doubt that you will incur the taunts of the enemy. Ye are the light of the world that is to say You, because ye are enlightened by the true light, ought to be the light of them who are in the world. A city set on an hill cannot be hid; that is to say The Apostles' teaching, founded upon Christ; in other words, the Church built upon Christ, out of many nations, in the unity of the faith, and bound together with the cement of love; to those who enter it, a place of safety; to those who go up to it, toilsome; the guardian of those who dwell in it, and excluding every enemy.\par
\switchcolumn*

\LA
\begin{Resp}\Rsign Iste est qui ante Deum magnas virtútes operátus est, et de omni corde suo laudávit Dóminum: \Star{} Ipse intercédat pro peccátis ómnium populórum, allelúja.\end{Resp}
\begin{Resp}\Vsign Ecce homo sine queréla, verus Dei cultor, ábstinens se ab omni ópere malo, et pérmanens in innocéntia sua.\end{Resp}
\begin{Resp}\Rsign Ipse intercédat pro peccátis ómnium populórum, allelúja.\end{Resp}
\switchcolumn
\EN
\begin{Resp}\Rsign This is he which wrought great wonders before God, and praised the Lord with all his heart. \Star{} May he pray for all people, that their sins may be forgiven unto them, alleluia.\end{Resp}
\begin{Resp}\Vsign Behold a man without blame, a worshipper of God in truth, keeping himself clean from every evil work, and abiding still in his innocency.\end{Resp}
\begin{Resp}\Rsign May he pray for all people, that their sins may be forgiven unto them, alleluia.\end{Resp}
\switchcolumn*

\LA
\LessonHead{Lectio VIII}
\switchcolumn
\EN
\LessonHead{Lesson VIII}
\switchcolumn*

\LA
\noindent Neque accendunt lucernam et ponunt eam sub modio, sed super candelabrum. Sub modio ergo lucernam ponit quisquis lucem doctrinæ commodis temporalibus obscurat et tegit; super candelabrum vero, qui se ita ministerio Dei subjicit, ut superior sit doctrina veritatis quam servitus corporis. Aliter Salvator accendit lucernam, qui humanæ testam naturæ flamma suæ divinitatis implevit; et hanc super candelabrum, id est Ecclesiam, posuit, quod in frontibus nostris fidem suæ incarnationis fixit. Quæ lucerna non potuit sub modio poni, id est sub mensura legis includi; nec in sola Judǽa, sed in universo illuxit orbe.\par
\switchcolumn
\EN
\noindent It either doth any man light a candle and put it under a bushel; but upon a candlestick. So he who puts the light under the bushel is he who for his own temporal ends would hide and tamper with the light of doctrine; but upon the candlestick he places it who follows the ministry of God in order that the teaching of the truth may be accounted a greater thing than the service of the body. In another aspect, the Saviour lighted the candle when He filled our mortal body with the flame of the God-head; and He placed it on a candlestick, that is the Church; for He fixed the faith of His incarnation upon our foreheads. Which light cannot be placed under a bushel; that is to say, it cannot be included within the measures of the law, nor in Judea alone, but has lightened the whole earth.\par
\switchcolumn*

\LA
\begin{Resp}\Rsign In médio Ecclésiæ apéruit os ejus, \Star{} Et implévit eum Dóminus spíritu sapiéntiæ et intelléctus, allelúja.\end{Resp}
\begin{Resp}\Vsign Jucunditátem et exsultatiónem thesaurizávit super eum.\end{Resp}
\begin{Resp}\Rsign Et implévit eum Dóminus spíritu sapiéntiæ et intelléctus, allelúja.\end{Resp}
\begin{Resp}\Vsign Glória Patri, et Fílio, \Star{} et Spirítui Sancto.\end{Resp}
\begin{Resp}\Rsign Et implévit eum Dóminus spíritu sapiéntiæ et intelléctus, allelúja.\end{Resp}
\switchcolumn
\EN
\begin{Resp}\Rsign In the midst of the congregation did the Lord open his mouth. \Star{} And filled him with the spirit of wisdom and understanding, alleluia.\end{Resp}
\begin{Resp}\Vsign He made him rich with joy and gladness.\end{Resp}
\begin{Resp}\Rsign And filled him with the spirit of wisdom and understanding, alleluia.\end{Resp}
\begin{Resp}\Vsign Glory be to the Father, and to the Son, \Star{} and to the Holy Ghost.\end{Resp}
\begin{Resp}\Rsign And filled him with the spirit of wisdom and understanding, alleluia.\end{Resp}
\switchcolumn*

\LA
\LessonHead{Lectio IX (pro commemoratione)}
\switchcolumn
\EN
\LessonHead{Lesson IX (when commemorated)}
\switchcolumn*

\LA
\LessonTitle{Commemoratio: S. Bedæ Venerabilis Confessoris et Ecclesiæ Doctoris}
\switchcolumn
\EN
\LessonTitle{Commemoration of: St. Bede the Venerable, Confessor and Doctor of the Church}
\switchcolumn*

\LA
\noindent Beda présbyter, Girvi in Británniæ et Scótiæ fínibus ortus est. Mónachus factus, vitam sic instítuit, ut, dum se ártium et doctrinárum stúdiis totum impénderet, nihil umquam de regulári disciplína remítteret. Nullum fuit doctrínæ genus, in quo non esset diligentíssime versátus; sed præcípua illi cura fuit divinárum Scripturárum meditátio, ita ut, sacerdótio initiátus, sacros explanáre libros aggréssus sit; in quo sanctórum Patrum doctrínis ádeo inhǽsit, ut nihil proférret nisi illórum judício comprobátum, eorúmdem étiam fere verbis usus. Otium perósus semper, ex lectióne ad oratiónem transíbat, ac vicíssim ex oratióne ad lectiónem. Emendándis fidélium móribus, fídei vindicándæ atque asseréndæ libros plures conscrípsit, quibus tantam sui apud omnes opiniónem fecit, ut ejus scripta, eo adhuc vivénte públice in ecclésiis legeréntur. Ætáte demum et labóribus fractus, pie obdormívit in Dómino. Eum Leo décimus tértius universális Ecclésiæ Doctórem declarávit.\par
\switchcolumn
\EN
\noindent Bede the priest was born at Jarrow, on the borders of England and Scotland. When a monk, he so arranged his life as to devote himself completely to the study of the liberal arts and sacred doctrine, without in any way relaxing the discipline of the Rule. There was no kind of learning in which he was not thoroughly versed; but his special interest was the study of the Scriptures; and when he was made a priest, he undertook the task of explaining the holy books. In doing so, he adhered to the teaching of the holy Fathers so closely that he would say nothing not already approved by their judgment, and he even made use of their very words. Abhorring laziness, he would go straight from reading to prayer and from prayer to reading. To raise the level of morality among Christians and to defend and spread the faith, he wrote many books, which gained him such a reputation with everyone that his writings were publicly read in churches during his own lifetime. At length, worn out with age and labours, he fell asleep peacefully in the Lord. Leo XIII declared him a Doctor of the universal Church.\par
\switchcolumn*

\LA
\TeDeum{Te Deum}
\switchcolumn
\EN
\TeDeum{Te Deum}
\switchcolumn*

\end{paracol}

\par\needspace{10\baselineskip}\addvspace{1.2em}{\centering\small\textsc{Die 27 Maii} · \textsc{27 May}\par}
\Office{S. Joannis Papæ et Martyris}{St. John I, Pope and Martyr}{Simplex}
\addcontentsline{toc}{subsection}{Die 27 Maii · S. Joannis Papæ et Martyris \textit{— St. John I, Pope and Martyr}}
\begin{paracol}{2}
\LA
\LessonHead{Lectio IX (pro commemoratione)}
\switchcolumn
\EN
\LessonHead{Lesson IX (when commemorated)}
\switchcolumn*

\LA
\LessonTitle{Commemoratio: S. Joannis Papæ et Martyris}
\switchcolumn
\EN
\LessonTitle{Commemoration of: St. John I, Pope and Martyr}
\switchcolumn*

\LA
\noindent Joannes Etruscus, Justino seniore imperatore, rexit Ecclesiam: ad quem profectus est Constantinopolim auxilii causa, quod Theodoricus rex hæreticus divexabat Italiam. Cujus étiam iter Deus miraculis illustravit, Nam, cum ei nobilis vir ad Corinthum, equum, quo ejus uxor mansueto utebatur, itineris causa commodasset, factum est, ut, domino postea remissus equus ita ferox evaderet, ut fremitu et totius corporis agitatione semper deinceps dominam expulerit: tamquam indignaretur mulierem recipere, ex quo sedisset in eo Jesu Christi Vicarius. Quam ob rem illi equum Pontifici donaverunt. Sed illud majus miraculum, quod Constantinopoli, in aditu portæ Aureæ inspectante frequentissimo populo, qui una cum imperatore Pontifici honoris causa occurrerat, cæco lumen restituit. Ad cujus pedes prostratus étiam imperator eum veneratus est. Rebus cum imperatore compositis, in Italiam rediit, statimque epistolam scripsit ad omnes Italiæ episcopos, jubens eos Arianorum ecclesias ad catholicum ritum consecrare, illud subjungens: Quia et nos, quando fuimus Constantinopoli tam pro religione catholica quam pro regis Theodorici causa, quascumque illis in partibus eorum ecclesias reperire potuimus, catholicas eas consecravimus. Quod iniquissimo animo ferens Theodoricus, dolo accersitum Pontificem Ravennam in carcerem conjecit; ubi, squalore inediaque afflictus, paucis diebus cessit e vita, cum sedisset annos duos, menses novem, dies quatuordecim, ordinatis eo tempore episcopis quindecim. Paulo post moritur Theodoricus: quem quidam eremita, ut scribit sanctus Gregorius, vidit inter Joánnem Pontificem et Symmachum patricium, quem idem occiderat, demergi in ignem Liparitanum; ut videlicet illi, quibus mortem attulerat, tamquam judices essent ejus interitus. Joannis corpus Ravenna Romam portatum est, et in basilica sancti Petri sepultum.\par
\switchcolumn
\EN
\noindent Pope John I was a Tuscan, who ruled the Church during the reign of the Emperor Justinian. He went to Constantinople to get help from Justinian in the troubles which the heretic King Theodoric was then causing in Italy. It pleased the Lord to mark this journey with wonders. A certain nobleman at Corinth lent to the Pope for his journey a very quiet horse on which his own wife was used to ride. But when the horse was returned to his owner he was found become so vicious, that by his restiveness and plunging he was always throwing off his mistress, as though he were not content to carry the lady after having carried the Vicar of Jesus Christ. When the nobleman and his wife found the beast to be thus worthless, they gave him for a present to the Pope. But a thing much more marvellous was that when the Pope, accompanied by the Emperor, and under the gaze of an immense multitude of people, who had come forth with Justinian to do him honour, was at the entering in of the Golden Gate of Constantinople, he gave sight to a blind man. Even the Emperor fell at his feet to show him respect. When he had arranged his business with Justinian he returned into Italy, and forthwith sent out a letter to all the Bishops of Italy, bidding them hallow for Catholic worship the churches of the Arians, and adding these words: We Ourselves when We were at Constantinople on some matters pertaining to the Catholic Religion and others pertaining to the King Theodoric, hallowed as Catholic all their Churches which We were able to find in those parts. Theodoric took this rule very ill, and, having enticed John by fraud to come to Ravenna, he cast him into prison, wherein, in a few days, he died of filth and hunger. He had sat in the chair of Peter two years, nine months, and fourteen days, within which time he had ordained fifteen Bishops. A little while afterward Theodoric also died. St. Gregory writeth that a certain hermit saw him between Pope John and Symmachus the Patrician, whom he had likewise slain, going down into the fiery crater of Lipari, as though they who had been his victims were become the judges of his punishment. The body of John was carried from Ravenna to Rome, and there buried in the Church of St. Peter.\par
\switchcolumn*

\LA
\TeDeum{Te Deum}
\switchcolumn
\EN
\TeDeum{Te Deum}
\switchcolumn*

\end{paracol}

\par\needspace{10\baselineskip}\addvspace{1.2em}{\centering\small\textsc{Die 28 Maii} · \textsc{28 May}\par}
\Office{S. Augustini Episcopi et Confessoris}{St. Augustine of Canterbury, Bishop and Confessor}{Duplex}
\addcontentsline{toc}{subsection}{Die 28 Maii · S. Augustini Episcopi et Confessoris \textit{— St. Augustine of Canterbury, Bishop and Confessor}}
\begin{paracol}{2}
\LA
\RefLine{Lectiones I–III: de Scriptura occurrente.}
\switchcolumn
\EN
\RefLine{Lessons I–III: from the Scripture of the day.}
\switchcolumn*

\LA
\LessonHead{Lectio IV}
\switchcolumn
\EN
\LessonHead{Lesson IV}
\switchcolumn*

\LA
\noindent Augustinus Romæ in Lateranensi cœnobio monachus, a Gregorio Magno cum secus monachis fere quadraginta in Angliam missus est anno quingentesimo nonagesimo septimo, ut gentes illas ad Christum converteret. Erat eo tempore rex Ethelbertus in Cantio potentissimus, qui audita adventus Augustíni causa, eum cum sociis Cantuariam, sui regni metropolim, invitavit; ibique manendi et Christum prædicandi facultatem eidem liberaliter concessit. Quare sanctus vir prope Cantuariam oratorium exstruxit, ubi ipse aliquamdiu consedit, atque apostolicam vivendi rationem cum suis æmulatus est.\par
\switchcolumn
\EN
\noindent Augustine, the first Archbishop of Canterbury, and the Apostle of the English, was sent into England by blessed Gregory, and came thither in the year 597. At that time there was in Kent a most mighty king named Ethelbert, whose power reached even to the Humber. When this King had heard wherefore the holy man was come, he received him kindly, and bade him and his companions, who were all monks, to come to his own capital city of Canterbury being struck with astonishment at the perfect blamelessness of their lives, and the power of the heavenly doctrine which they preached, and which God confirmed with signs following.\par
\switchcolumn*

\LA
\begin{Resp}\Rsign Invéni David servum meum, óleo sancto meo unxi eum: \Star{} Manus enim mea auxiliábitur ei, allelúja.\end{Resp}
\begin{Resp}\Vsign Nihil profíciet inimícus in eo, et fílius iniquitátis non nocébit ei.\end{Resp}
\begin{Resp}\Rsign Manus enim mea auxiliábitur ei, allelúja.\end{Resp}
\switchcolumn
\EN
\begin{Resp}\Rsign I have found David My servant, with My holy oil have I anointed him \Star{} For My hand shall help him, alleluia.\end{Resp}
\begin{Resp}\Vsign The enemy shall prevail nothing against him, nor the son of wickedness afflict him.\end{Resp}
\begin{Resp}\Rsign For My hand shall help him, alleluia.\end{Resp}
\switchcolumn*

\LA
\LessonHead{Lectio V}
\switchcolumn
\EN
\LessonHead{Lesson V}
\switchcolumn*

\LA
\noindent Cælestis doctrina prædicatione plurimis firmata miraculis, ac vitai exemplo sic insulanos illos demulsit, ut eorum plerosque ad christianam fidem perduxerit, ac demum regem ipsum, quem cum innumero suorum comitatu sacro fonte lustravit, summa cum lætitia Berthæ regiæ uxoris, quæ christiana erat. Olim in Natali Domini, cum decem millibus et amplius baptismum in alveo fluminis Eboraci contulisset, quotquot ex iis morbo aliquo affecti erant, cum animæ salute, corporis quoque sanitatem recepisse memoriæ proditum est. Jussu Gregorii ordinatus episcopus, sedem Cantuariæ instituit in ecclesia Salvatoris a se erecta, in qua monachos operis sui subsidiarios collocavit; et sancti Petri monasterium, quod postea et a suo nomine dictum est, in suburbanis construxit. Idem Gregorius usum pallii cum facultate ecclesiasticæ hierarchiæ in Anglia instituendas ei concessit, quo novam étiam operariorum manum misit, nempe Mellitum, Justum, Paulinum, et Rufinianum.\par
\switchcolumn
\EN
\noindent They drew nigh to the city in solemn procession, singing the Litany, and bearing before them for their standard a silver cross and a picture of the Lord our Saviour painted on a panel. Hard by the city, upon the east side, there was a Church builded of old time in honour of St. Martin, and wherein the Queen, who was a Christian, was used to pray. There they first began to meet together, to sing, to pray, to celebrate Masses, to preach, and to baptize, until the King was turned to the faith, and the most part of his people were led by his example, (but not his authority,) to take the name of Christian, for he had learnt from his teachers and his own soul's physicians, that men are to be drawn, and not driven to heaven. And now Augustine, being ordained Archbishop of the English and of Britain, lest he should leave untravailed any part of the Lord's vineyard, asked from the Apostolic See a new band of labourers, Mellitus, Justus, Paulinus, and Rufinian.\par
\switchcolumn*

\LA
\begin{Resp}\Rsign Pósui adjutórium super poténtem, et exaltávi eléctum de plebe mea: \Star{} Manus enim mea auxiliábitur ei, allelúja.\end{Resp}
\begin{Resp}\Vsign Invéni David servum meum, óleo sancto meo unxi eum.\end{Resp}
\begin{Resp}\Rsign Manus enim mea auxiliábitur ei, allelúja.\end{Resp}
\switchcolumn
\EN
\begin{Resp}\Rsign I have laid help upon one that is mighty, and have exalted one chosen out of My people \Star{} For My hand shall help him, alleluia.\end{Resp}
\begin{Resp}\Vsign I have found David My servant, with My holy oil have I anointed him.\end{Resp}
\begin{Resp}\Rsign For My hand shall help him, alleluia.\end{Resp}
\switchcolumn*

\LA
\LessonHead{Lectio VI}
\switchcolumn
\EN
\LessonHead{Lesson VI}
\switchcolumn*

\LA
\noindent Dispositis ejus ecclesiæ rebus, synodum habuit Augustinus cum episcopis atque doctoribus veterum Britonum, qui in Paschæ celebratione aliusque ritibus ab Ecclesia Romana jamdudum dissidebant. Sed cum eos neque Apostolicæ Sedis auctoritate, neque miraculis movere posset ut dissidio cessarent, prophetico spiritu eis excidium prænuntiavit. Denique maximis pro Christo exantlutis laboribus, miraculis clarus, cum Mellitum Londinensi ecclesiæ præfecisset, Justum Roffensi, suæ Laurentium, in cælum migravit septimo Kalendas Junias, Ethelberto regnante, ac sepultus est in monasterio sancti Petri, quod exinde Cantuariensium Antistitum et aliquot regum conditorium fuit. Ejus cultum ferventi studio prosecutæ sunt Anglorum gentes, ac Leo decimus tertius Pontifex Maximus ejus Oificium et Missam ad universam extendit Ecclesiam.\par
\switchcolumn
\EN
\noindent Having arranged the affairs of his church, Augustine held a synod with the bishops and doctors of the ancient Britons, who had long been at variance with the Roman Church in the celebration of Easter and other rites. But since he could not move them, either by the authority of the apostolic see or by miracles, to put an end to these variations, in a prophetic spirit he foretold their ruin. At length, after having endured many difficulties for Christ, and having become noted for miracles, when he had placed Mellitus in charge of the church of London, Justus of that of Rochester, and Laurence in charge of his own church, he passed to heaven on the 26th day of May, in the reign of Ethelbert, and was buried in the monastery of St. Peter, which thereafter became the burying-place of the bishops of Canterbury and of some kings. The English people honoured his memory with fervent zeal; and the Supreme Pontiff Leo XIII extended his Office and Mass to the universal Church.\par
\switchcolumn*

\LA
\begin{Resp}\Rsign Iste est, qui ante Deum magnas virtútes operátus est, et omnis terra doctrína ejus repléta est: \Star{} Ipse intercédat pro peccátis ómnium populórum, allelúja.\end{Resp}
\begin{Resp}\Vsign Iste est, qui contémpsit vitam mundi, et pervénit ad cæléstia regna.\end{Resp}
\begin{Resp}\Rsign Ipse intercédat pro peccátis ómnium populórum, allelúja.\end{Resp}
\begin{Resp}\Vsign Glória Patri, et Fílio, \Star{} et Spirítui Sancto.\end{Resp}
\begin{Resp}\Rsign Ipse intercédat pro peccátis ómnium populórum, allelúja.\end{Resp}
\switchcolumn
\EN
\begin{Resp}\Rsign This is he which wrought great wonders before God, and the whole earth is full of his teaching \Star{} May he pray for all people, that their sins may be forgiven unto them, alleluia.\end{Resp}
\begin{Resp}\Vsign This is he which loved not his life in this world, and hath attained unto the kingdom of heaven.\end{Resp}
\begin{Resp}\Rsign May he pray for all people, that their sins may be forgiven unto them, alleluia.\end{Resp}
\begin{Resp}\Vsign Glory be to the Father, and to the Son, \Star{} and to the Holy Ghost.\end{Resp}
\begin{Resp}\Rsign May he pray for all people, that their sins may be forgiven unto them, alleluia.\end{Resp}
\switchcolumn*

\LA
\LessonHead{Lectio VII}
\switchcolumn
\EN
\LessonHead{Lesson VII}
\switchcolumn*

\LA
\LessonTitle{Léctio sancti Evangélii secúndum Lucam.}
\switchcolumn
\EN
\LessonTitle{From the Holy Gospel according to Luke}
\switchcolumn*

\LA
\Cite{Luc 10:1-9}
\switchcolumn
\EN
\Cite{Luke 10:1-9}
\switchcolumn*

\LA
\noindent In illo témpore: Designavit Dominus et alios septuaginta duos: et misit illos binos ante faciem suam, in omnem civitatem et locum, quo erat ipse venturus. Et reliqua.\par
\switchcolumn
\EN
\noindent At that time, the Lord appointed other seventy-two also, and sent them two and two before His face into every city and place, whither He Himself would come. And so on.\par
\switchcolumn*

\LA
\LessonTitle{Homilia sancti Gregorii Papæ.}
\switchcolumn
\EN
\LessonTitle{Homily by Pope St. Gregory (the Great)}
\switchcolumn*

\LA
\Source{Homilia 17. in Evang.}
\switchcolumn
\EN
\Source{17th Homily on the Gospels}
\switchcolumn*

\LA
\noindent Dominus et Salvator noster, fratres carissimi, aliquando nos sermonibus, aliquando vero operibus admonet. Ipsa etenim facta ejus præcepta sunt: quia dum aliquid tacitus facit, quid agere debeamus innotescit. Ecce enim binos in prædicationem discipulos mittit: quia duo sunt præcepta caritatis, Dei videlicet amor, et proximi: et minus quam inter duos caritas haberi non potest. Nemo enim proprie ad semetipsum habere caritatem dicitur: sed dilectio in alterum tendit, ut caritas esse possit.\par
\switchcolumn
\EN
\noindent Dearly beloved brethren, our Lord and Saviour doth sometimes admonish us by words, and sometimes by works. Yea, His very works do themselves teach us for that which He doth silently His example still moveth us to copy. Behold how He sendeth forth His disciples to preach by two and two since there are two commandments to love, that is, a commandment to love God, and a commandment to love our neighbour and where there are not two, the one, being alone, hath not whereon to do the Lord's commandment. And no man can properly be said to love himself: for love tendeth outward toward our neighbour, if it be thelove whereto the Gospel doth oblige us.\par
\switchcolumn*

\LA
\begin{Resp}\Rsign Amávit eum Dóminus, et ornávit eum: stolam glóriæ índuit eum, \Star{} Et ad portas paradísi coronávit eum, allelúja.\end{Resp}
\begin{Resp}\Vsign Induit eum Dóminus lorícam fídei, et ornávit eum.\end{Resp}
\begin{Resp}\Rsign Et ad portas paradísi coronávit eum, allelúja.\end{Resp}
\switchcolumn
\EN
\begin{Resp}\Rsign The Lord loved him and beautified him He clothed him with a robe of glory \Star{} And crowned him at the gates of Paradise, alleluia.\end{Resp}
\begin{Resp}\Vsign The Lord hath put on him the breast-plate of faith, and hath adorned him.\end{Resp}
\begin{Resp}\Rsign And crowned him at the gates of Paradise, alleluia.\end{Resp}
\switchcolumn*

\LA
\LessonHead{Lectio VIII}
\switchcolumn
\EN
\LessonHead{Lesson VIII}
\switchcolumn*

\LA
\noindent Ecce enim binos ad prædicandum discipulos Dominus mittit: quatenus hoc nobis tacitus innuat, quia qui caritatem erga alterum non habet, prædicationis officium suscipere nullatenus debet. Bene autem dicitur, quia misit eos ante faciem suam in omnem civitatem et locum, quo erat ipse venturus. Prædicatores enim suos Dominus sequitur: quia prædicatio prævenit, et tunc ad mentis nostræ habitaculum Dominus venit, quando verba exhortationis præcurrunt: atque per hoc veritas in mente suscipitur.\par
\switchcolumn
\EN
\noindent Behold, the Lord sendeth forth His disciples to preach by two and two and thus doing, He doth silently teach us that whosoever loveth not his neighbour, such an one it behoveth not to take upon him the office of a preacher. Well also is it said that He sent them before His face into every city and place whither He Himself would come. The Lord followeth His preachers first cometh preaching, and then the Lord Himself cometh to the house of our mind, whither the word of exhortation hath come before and so cometh the truth into our mind.\par
\switchcolumn*

\LA
\begin{Resp}\Rsign Sint lumbi vestri præcíncti, et lucérnæ ardéntes in mánibus vestris: \Star{} Et vos símiles homínibus exspectántibus dóminum suum, quando revertátur a núptiis, allelúja.\end{Resp}
\begin{Resp}\Vsign Vigiláte ergo, quia nescítis qua hora Dóminus vester ventúrus sit.\end{Resp}
\begin{Resp}\Rsign Et vos símiles homínibus exspectántibus dóminum suum, quando revertátur a núptiis, allelúja.\end{Resp}
\begin{Resp}\Vsign Glória Patri, et Fílio, \Star{} et Spirítui Sancto.\end{Resp}
\begin{Resp}\Rsign Et vos símiles homínibus exspectántibus dóminum suum, quando revertátur a núptiis, allelúja.\end{Resp}
\switchcolumn
\EN
\begin{Resp}\Rsign Let your loins be girded about, and your lights burning; \Star{} And ye yourselves like unto men that wait for their lord, when he will return from the wedding, alleluia.\end{Resp}
\begin{Resp}\Vsign Watch therefore, for ye know not what hour your Lord doth come.\end{Resp}
\begin{Resp}\Rsign And ye yourselves like unto men that wait for their lord, when he will return from the wedding, alleluia.\end{Resp}
\begin{Resp}\Vsign Glory be to the Father, and to the Son, \Star{} and to the Holy Ghost.\end{Resp}
\begin{Resp}\Rsign And ye yourselves like unto men that wait for their lord, when he will return from the wedding, alleluia.\end{Resp}
\switchcolumn*

\LA
\LessonHead{Lectio IX}
\switchcolumn
\EN
\LessonHead{Lesson IX}
\switchcolumn*

\LA
\noindent Hinc namque eisdem prædicatoribus Isaias dicit: Parate viam Domini, rectas facite semitas Dei nostri. Hinc filiis Psalmista ait: Iter facite ei, qui ascendit super occasum. Super occasum namque Dominus ascendit: quia unde in passione occubuit, inde majorem suam gloriam resurgendo manifestavit. Super occasum videlicet ascendit; quia mortem quam pertulit, resurgendo calcavit. Ei ergo qui ascendit super occasum, iter facimus, cum nos ejus gloriam vestris mentibus prædicamus, ut eas et ipse post veniens, per amoris sui præsentiam illustret.\par
\switchcolumn
\EN
\noindent Therefore to preachers saith Isaiah: “Prepare ye the way of the Lord, make straight an highway for our God.” (xl. 3.) And again the Psalmist saith: “Spread a path before Him That rideth upon the West.” (lxvii. 4.) The Lord rideth upon the West above that from which in death He veiled His glory hath He royally exalted that glory that excelleth, even the glory of His rising again. He rideth upon the West, Who, being risen again from the dead, is throned high above the death to which He bowed. Before Him, therefore, That rideth upon the West, we spread a path, when we set forth His glory before the eyes of your mind, to the end that He Himself may come after, and Himself enlighten the same your minds by His presence and His love.\par
\switchcolumn*

\LA
\TeDeum{Te Deum}
\switchcolumn
\EN
\TeDeum{Te Deum}
\switchcolumn*

\LA
\LessonHead{Lectio IX (pro commemoratione)}
\switchcolumn
\EN
\LessonHead{Lesson IX (when commemorated)}
\switchcolumn*

\LA
\LessonTitle{Commemoratio: S. Augustini Episcopi et Confessoris}
\switchcolumn
\EN
\LessonTitle{Commemoration of: St. Augustine of Canterbury, Bishop and Confessor}
\switchcolumn*

\LA
\noindent Augustínus, Romæ in Lateranénsi cœnóbio mónachus, a Gregório Magno cum sóciis mónachis fere quadragínta in Angliam missus est, anno quingentésimo nonagésimo séptimo. A rege Ethelbérto Cantuáriam, ejus regni metrópolim, invitátus cum sóciis, prope eam oratórium exstrúxit. Cæléstis doctrínæ prædicatióne plerósque insulános ac regem ipsum ad christiánam fidem perdúxit, summa cum lætítia Berthæ, régiæ uxóris, quæ Christiána erat. Jussu Gregórii ordinátus epíscopus, Sedem Cantuariénsem instítuit, et ab eódem Pontífice usum pállii cum facultáte hierarchíæ in Anglia instituéndæ obtínuit. Máximis demum pro Christo exantlátis labóribus, cum Mellítum Londinénsi ecclésiæ præfecísset, Justum Roffénsi, suæ Lauréntium, in cælum migrávit séptimo kaléndas júnias, et sepúltus est in monastério sancti Petri, quod exínde Cantuariénsium antístitum et áliquot regum conditórium fuit.\par
\switchcolumn
\EN
\noindent Augustine, a monk of the Lateran monastery in Rome, was sent by Gregory the Great in 597 to England with about forty monks as his companions. They were invited by King Ethelbert to Canterbury, the chief city of the kingdom, and they built an oratory nearby. Through preaching the doctrine of heaven, Augustine brought many of the islanders and the king himself to the Christian faith, to the great joy of the king's wife, Bertha, who was a Christian. By order of Pope Gregory, Augustine was ordained bishop and founded the see of Canterbury; by the same Pontiff he was granted the use of the pallium and the right to organize the hierarchy of England. At length, after suffering great hardships for Christ, having set Mellitus over the Church of London, Justus over that of Rochester, and Lawrence over his own Church, he made his journey to heaven on the 26th day of May. He was buried in the monastery of St. Peter, which then became the burial place of bishops of Canterbury and of several kings.\par
\switchcolumn*

\LA
\TeDeum{Te Deum}
\switchcolumn
\EN
\TeDeum{Te Deum}
\switchcolumn*

\end{paracol}

\par\needspace{10\baselineskip}\addvspace{1.2em}{\centering\small\textsc{Die 29 Maii} · \textsc{29 May}\par}
\Office{S. Mariæ Magdalenæ de Pazzis Virginis}{St. Mary Magdalene de Pazzi, Virgin}{Semiduplex}
\addcontentsline{toc}{subsection}{Die 29 Maii · S. Mariæ Magdalenæ de Pazzis Virginis \textit{— St. Mary Magdalene de Pazzi, Virgin}}
\begin{paracol}{2}
\LA
\RefLine{Lectiones I–III: de Scriptura occurrente.}
\switchcolumn
\EN
\RefLine{Lessons I–III: from the Scripture of the day.}
\switchcolumn*

\LA
\RefLine{Lectiones VII–IX: ex Commúni Virginum tantum (C6a).}
\switchcolumn
\EN
\RefLine{Lessons VII–IX: from the Common of a Virgin not a Martyr (C6a).}
\switchcolumn*

\LA
\LessonHead{Lectio IV}
\switchcolumn
\EN
\LessonHead{Lesson IV}
\switchcolumn*

\LA
\noindent Maria Magdalena, illustrióri Pazziórum genere Floréntiæ nata, fere ab incunabulis iter perfectiónis arrípuit. Decennis perpetuam virginitátem vovit, susceptoque habitu in monastério sanctæ Maríæ Angelórum, ordinis Carmelitárum, se ómnium virtútum exemplar exhíbuit. Adeo casta fuit, ut quidquid puritátem lǽdere potest, pénitus ignoraverit. Quinquénnium, Deo jubénte, solo pane et aqua transegit, exceptis diébus Domínicis, quibus cibis quadragesimálibus vescebátur. Corpus suum cilício, flagellis, frígore, inedia, vigiliis, nuditate, atque omni pœnárum genere cruciábat.\par
\switchcolumn
\EN
\noindent This Mary Magdalen was born of the noble Florentine family of the Pazzi, (on the 2nd day of April, in the year of Christ 1566.) She was hardly out of her cradle when she set her feet in the path of perfection. At ten years of age she made a vow of perpetual virginity, and (at fifteen) took the habit of the Order of Mount Carmel, in the convent of Saint Mary of the Angels. In that sisterhood she was in all ways a pattern to all. She was pure to that degree, that she did not even know of the existence of anything which can hurt modesty. For the space of five years, by the command of God, she lived upon nothing but bread and water, the Lord's Day only excepted, in which she used the food which is taken in Lent. She chastised her body with hair-cloth, scourging, cold, hunger, watching, nakedness, and all manner of hardships.\par
\switchcolumn*

\LA
\begin{Resp}\Rsign Propter veritátem, et mansuetúdinem, et justítiam: \Star{} Et dedúcet te mirabíliter déxtera tua, allelúja.\end{Resp}
\begin{Resp}\Vsign Spécie tua et pulchritúdine tua inténde, próspere procéde, et regna.\end{Resp}
\begin{Resp}\Rsign Et dedúcet te mirabíliter déxtera tua, allelúja.\end{Resp}
\switchcolumn
\EN
\begin{Resp}\Rsign Because of truth, and meekness, and righteousness; \Star{} And thy right hand shall lead thee wonderfully, alleluia.\end{Resp}
\begin{Resp}\Vsign In thy comeliness, and thy beauty, go forward, fare prosperously, and reign.\end{Resp}
\begin{Resp}\Rsign And thy right hand shall lead thee wonderfully, alleluia.\end{Resp}
\switchcolumn*

\LA
\LessonHead{Lectio V}
\switchcolumn
\EN
\LessonHead{Lesson V}
\switchcolumn*

\LA
\noindent Tanto igne divini amoris æstuábat, ut, ei feréndo impar, ingesta aqua pectus refrigerare cogerétur. Extra sensus frequenter rapta, diuturnas et admirábiles éxtases passa est, in quibus et arcana cæléstia penetrávit, et exímiis a Deo grátiis illustráta fuit. His autem muníta longum certamen a princípibus tenebrárum sustinuit, árida, desoláta, ab ómnibus derelícta, variisque tentatiónibus vexáta; Deo sic permittente, ut invíctæ patiéntiæ ac profundíssimæ humilitátis exémplar præbéret.\par
\switchcolumn
\EN
\noindent The love of God was so hot within her, that she was sometimes fain to bathe her breast with cold water to allay the agitation. She was oftentimes rapt in the spirit, and that most marvellously, for whole days at a time, during which trances she saw things hidden and heavenly, and was enlightened of God with great gifts. But after all these things she had a stern tussling with the prince of the darkness of this world, while God allowed her spirit to remain dry, deserted, abandoned by all, and tormented with diverse temptations. And all that while she remained an example of unconquered patience and the deepest lowly-mindedness.\par
\switchcolumn*

\LA
\begin{Resp}\Rsign Dilexísti justítiam, et odísti iniquitátem: \Star{} Proptérea unxit te Deus, Deus tuus, óleo lætítiæ, allelúja.\end{Resp}
\begin{Resp}\Vsign Propter veritátem, et mansuetúdinem, et justítiam.\end{Resp}
\begin{Resp}\Rsign Proptérea unxit te Deus, Deus tuus, óleo lætítiæ, allelúja.\end{Resp}
\switchcolumn
\EN
\begin{Resp}\Rsign Thou hast loved righteousness, and hated iniquity; \Star{} Therefore God, thy God, hath anointed thee with the oil of gladness, alleluia.\end{Resp}
\begin{Resp}\Vsign Because of truth, and meekness, and righteousness.\end{Resp}
\begin{Resp}\Rsign Therefore God, thy God, hath anointed thee with the oil of gladness, alleluia.\end{Resp}
\switchcolumn*

\LA
\LessonHead{Lectio VI}
\switchcolumn
\EN
\LessonHead{Lesson VI}
\switchcolumn*

\LA
\noindent Caritáte erga próximum singuláriter enítuit; nam sæpe noctes ducebat insomnes, vel obeúndis sorórum ministériis, vel inserviéndo infirmis occupata, quarum aliquándo úlcera lambens sanávit. Infidélium et peccatórum perditiónem amare deflens, se ad quælibet pro illórum salúte torménta parátam offerébat. Multis ante obitum annis, univérsis cæli deliciis, quibus copiose affluébat, heróica virtúte renuntians, illud frequénter in ore habébat: Pati, non mori. Tandem longa et gravíssima infirmitáte exháusta, transívit ad Sponsum die vigésima quinta Maji, anno millésimo sexcentésimo septimo, expleto anno quadragésimo primo ætátis suæ. Eam, multis in vita et post mortem miráculis claram, Clemens nonus sanctárum Vírginum número adscrípsit: cujus corpus in præséntem diem incorrúptum conservátur.\par
\switchcolumn
\EN
\noindent She was very remarkable for her tender love toward her neighbours. Sometimes she went whole nights without sleep, while she was working for the service of the sisters, Or waiting upon the sick. She sometimes healed sores even by licking them. That there should be unbelievers and sinners perishing caused her bitter weeping, and she offered herself to God to suffer for their conversion whatsoever He chose. For many years, therefore, before her death, her mighty charity towards Others, made her freely to give up that heavenly joy of spirit, wherewith she had once so overflowed. She had often in her mouth the words: To suffer, not to die. At length, in the forty-second year of her age, on the 25th day of May, in the year 1607, after a long and grievous sickness, the Bridegroom came, and she entered with Him into the marriage-chamber. Clement IX, finding that God had glorified her by many miracles, both during her life and after her death, enrolled her name among those of the Holy Virgins. Her body, up to the present day, has never shown the least sign of corruption.\par
\switchcolumn*

\LA
\begin{Resp}\Rsign Afferéntur Regi vírgines post eam, próximæ ejus; \Star{} Afferéntur tibi in lætítia et exsultatióne, allelúja.\end{Resp}
\begin{Resp}\Vsign Spécie tua et pulchritúdine tua; inténde, próspere procéde, et regna.\end{Resp}
\begin{Resp}\Rsign Afferéntur tibi in lætítia et exsultatióne, allelúja.\end{Resp}
\begin{Resp}\Vsign Glória Patri, et Fílio, \Star{} et Spirítui Sancto.\end{Resp}
\begin{Resp}\Rsign Afferéntur tibi in lætítia et exsultatióne, allelúja.\end{Resp}
\switchcolumn
\EN
\begin{Resp}\Rsign After her shall virgins be brought unto the King, her fellows \Star{} They shall be brought unto thee with gladness and rejoicing, alleluia.\end{Resp}
\begin{Resp}\Vsign In thy comeliness and thy beauty, go forward, fare prosperously, and reign.\end{Resp}
\begin{Resp}\Rsign They shall be brought unto thee with gladness and rejoicing, alleluia.\end{Resp}
\begin{Resp}\Vsign Glory be to the Father, and to the Son, \Star{} and to the Holy Ghost.\end{Resp}
\begin{Resp}\Rsign They shall be brought unto thee with gladness and rejoicing, alleluia.\end{Resp}
\switchcolumn*

\LA
\LessonHead{Lectio IX (pro commemoratione)}
\switchcolumn
\EN
\LessonHead{Lesson IX (when commemorated)}
\switchcolumn*

\LA
\LessonTitle{Commemoratio: S. Mariæ Magdalenæ de Pazzis Virginis}
\switchcolumn
\EN
\LessonTitle{Commemoration of: St. Mary Magdalene de Pazzi, Virgin}
\switchcolumn*

\LA
\noindent María Magdaléna Floréntiæ illústri Pazziórum génere nata, fere ab incunábulis iter perfectiónis arrípuit. Decénnis perpétuam virginitátem vovit, susceptóque hábitu in monastério sorórum Carmelitárum, se ómnium virtútum exémplar prǽbuit. Adeo casta fuit, ut, quidquid puritátem lǽdere posset, pénitus ignoráverit. Tanto igne divíni amóris æstuábat, ut, ei feréndo impar, ingésta aqua pectus refrigeráre cogerétur. Caritáte erga próximum excélluit; nam sæpe noctes ducébat insómnes, vel obeúndis sorórum ministériis, vel inserviéndo infírmis occupáta, quarum aliquándo úlcera lambens sanávit. Illud frequénter in ore habébat: Pati, non mori. Tandem, diútino et gravi morbo exháusta, transívit ad Sponsum, anno millésimo sexcentésimo séptimo, expléto ætátis suæ quadragésimo primo. Eam Clemens nonus sanctárum Vírginum número adscrípsit.\par
\switchcolumn
\EN
\noindent Mary Magdalene was born of the famous family of the Pazzi in Florence, and almost from her cradle entered on the way of perfection. At ten years old, she made a vow of perpetual virginity, and when she had taken the habit in the monastery of the Carmelite sisters, she showed herself to be exemplary in all virtues. She was so chaste that she did not even know of anything that could harm purity. She burned with such a fire of divine love that she was unable to bear it, and had to cool her breast by pouring water over it. She was notable for her love of neighbor, often passing sleepless nights either in doing the work of the sisters, or in serving the sick, whose ulcers she sometimes cured by licking them. She frequently repeated this saying: “To suffer, not to die.” At length, worn out by a long and serious illness, she went to the Bridegroom in the year 1607, having completed her forty-first year. Clement IX numbered her among the Holy Virgins.\par
\switchcolumn*

\LA
\TeDeum{Te Deum}
\switchcolumn
\EN
\TeDeum{Te Deum}
\switchcolumn*

\end{paracol}

\par\needspace{10\baselineskip}\addvspace{1.2em}{\centering\small\textsc{Die 30 Maii} · \textsc{30 May}\par}
\Office{S. Felicis I Papæ et Martyris}{St. Felix, Pope and Martyr}{Simplex}
\addcontentsline{toc}{subsection}{Die 30 Maii · S. Felicis I Papæ et Martyris \textit{— St. Felix, Pope and Martyr}}
\begin{paracol}{2}
\LA
\RefLine{Lectiones I–II: de Scriptura occurrente.}
\switchcolumn
\EN
\RefLine{Lessons I–II: from the Scripture of the day.}
\switchcolumn*

\LA
\LessonHead{Lectio III (pro commemoratione)}
\switchcolumn
\EN
\LessonHead{Lesson III (when commemorated)}
\switchcolumn*

\LA
\noindent Felix Románus, patre Constántio, Aureliáno imperatóre præfuit Ecclesiæ. Constítuit ut Missa supra memórias et sepúlcra Mártyrum celebrarétur. Qui cum mense Decémbri habuísset ordinatiónes duas, et creásset presbýteros novem, diáconos quinque, epíscopos per divérsa loca quinque, martyrio coronátus, via Aurelia sepelítur in basílica quam a se ædificátam dedicárat. Vixit in pontificátu annos duos, menses quátuor, dies vigínti novem.\par
\switchcolumn
\EN
\noindent Pope Felix I was a Roman who ruled the Church in the days of the Emperor Aurelian. His father's name was Constantius. His is the ordinance which commands that Mass should be celebrated on the monuments and graves of martyrs. He held two December ordinations, wherein he ordained nine Priests, five Deacons, and five Bishops for diverse places. Having finished his testimony he was buried upon the Aurelian Way, in the Church which he had himself built and dedicated. He lived as Pope two years, four months, and twenty-nine days.\par
\switchcolumn*

\LA
\TeDeum{Te Deum}
\switchcolumn
\EN
\TeDeum{Te Deum}
\switchcolumn*

\end{paracol}

\par\needspace{10\baselineskip}\addvspace{1.2em}{\centering\small\textsc{Die 31 Maii} · \textsc{31 May}\par}
\Office{Beatæ Mariæ Virginis Reginæ}{Beatæ Mariæ Virginis Reginæ}{Duplex II classis}
\addcontentsline{toc}{subsection}{Die 31 Maii · Beatæ Mariæ Virginis Reginæ \textit{— Beatæ Mariæ Virginis Reginæ}}
\begin{paracol}{2}
\LA
\RefLine{Lectio IX: de Scriptura occurrente.}
\switchcolumn
\EN
\RefLine{Lesson IX: from the Scripture of the day.}
\switchcolumn*

\LA
\LessonHead{Lectio I}
\switchcolumn
\EN
\LessonHead{Lesson I}
\switchcolumn*

\LA
\LessonTitle{De libro Ecclesiástici}
\switchcolumn
\EN
\LessonTitle{Lesson from the book of Ecclesiasticus}
\switchcolumn*

\LA
\Cite{Sir 24:5-11}
\switchcolumn
\EN
\Cite{Sir 24:5-11}
\switchcolumn*

\LA
\noindent Ego ex ore Altíssimi prodívi, primogénita ante omnem creatúram: ego feci in cælis ut orirétur lumen indefíciens, et sicut nébula texi omnem terram: ego in altíssimis habitávi et thronus meus in colúmna nubis. Gyrum cæli circuívi sola, et profúndum abýssi penetrávi, in flúctibus maris ambulávi, et in omni terra steti: et in omni pópulo, et in omni gente primátum hábui: et ómnium excelléntium et humílium corda virtúte calcávi: et in his ómnibus réquiem quæsívi, et in hereditáte Dómini morábor.\par
\switchcolumn
\EN
\noindent I came out of the mouth of the most High, the firstborn before all creatures: I made that in the heavens there should rise light that never faileth, and as a cloud I covered all the earth: I dwelt in the highest places, and my throne is in a pillar of a cloud. I alone have compassed the circuit of heaven, and have penetrated into the bottom of the deep, and have walked in the waves of the sea, And have stood in all the earth: and in every people, And in every nation I have had the chief rule: And by my power I have trodden under my feet the hearts of all the high and low: and in all these I sought rest, and I shall abide in the inheritance of the Lord.\par
\switchcolumn*

\LA
\begin{Resp}\Rsign Beáta es, María, quæ credidísti Dómino: perfécta sunt in te quæ dicta sunt tibi. \Star{} Ecce exaltáta es super choros Angelórum ad cæléstia regna.\end{Resp}
\begin{Resp}\Vsign Ave Maria, grátia plena, Dominus tecum.\end{Resp}
\begin{Resp}\Rsign Ecce exaltáta es super choros Angelorum ad cæléstia regna.\end{Resp}
\switchcolumn
\EN
\begin{Resp}\Rsign Blessed art thou, Mary, because thou hast believed the Lord: \Star{} The things promised thee have been accomplished in thee.\end{Resp}
\begin{Resp}\Vsign Hail, Mary, full of grace, the Lord is with thee.\end{Resp}
\begin{Resp}\Rsign The things promised thee have been accomplished in thee.\end{Resp}
\switchcolumn*

\LA
\LessonHead{Lectio II}
\switchcolumn
\EN
\LessonHead{Lesson II}
\switchcolumn*

\LA
\Cite{Sir 24:14-16}
\switchcolumn
\EN
\Cite{Sir 24:14-16}
\switchcolumn*

\LA
\noindent Ab inítio, et ante sǽcula creáta sum, et usque ad futúrum sǽculum non désinam, et in habitatióne sancta coram ipso ministrávi. Et sic in Sion firmáta sum, et in civitáte sanctificáta simíliter requiévi, et in Jerúsalem potéstas mea. Et radicávi in pópulo honorificáto, et in parte Dei mei heréditas illíus, et in plenitúdine sanctórum deténtio mea.\par
\switchcolumn
\EN
\noindent From the beginning, and before the world, was I created, and unto the world to come I shall not cease to be, and in the holy dwelling place I have ministered before him. And so was I established in Sion, and in the holy city likewise I rested, and my power was in Jerusalem. And I took root in an honourable people, and in the portion of my God his inheritance, and my abode is in the full assembly of saints.\par
\switchcolumn*

\LA
\begin{Resp}\Rsign Regálem dignitátem Vírginis Maríæ recolámus. \Star{} Quia cum Christo regnat in ætérnum. (Allelúja.)\end{Resp}
\begin{Resp}\Vsign Glóriam Regínæ nostræ celebrémus.\end{Resp}
\begin{Resp}\Rsign Quia cum Christo regnat in ætérnum. (Allelúja.)\end{Resp}
\switchcolumn
\EN
\begin{Resp}\Rsign Let us reflect upon the royal dignity of the Virgin Mary, \Star{} For with Christ she reigneth forever.\end{Resp}
\begin{Resp}\Vsign Let us proclaim the glory of our Queen.\end{Resp}
\begin{Resp}\Rsign For with Christ she reigneth forever.\end{Resp}
\switchcolumn*

\LA
\LessonHead{Lectio III}
\switchcolumn
\EN
\LessonHead{Lesson III}
\switchcolumn*

\LA
\Cite{Sir 24:24-30}
\switchcolumn
\EN
\Cite{Sir 24:24-30}
\switchcolumn*

\LA
\noindent Ego mater pulchræ dilectiónis, et timóris, et agnitiónis, et sanctæ spei. In me grátia omnis viæ et veritátis: in me omnis spes vitæ et virtútis. Transíte ad me, omnes qui concupíscitis me, et a generatiónibus meis implémini; spíritus enim meus super mel dulcis, et heréditas mea super mel et favum. Memória mea in generatiónes sæculórum. Qui edunt me, adhuc esúrient, et qui bibunt me, adhuc sítient. qui audit me non confundétur, et qui operántur in me non peccábunt: qui elúcidant me, vitam ætérnam habébunt.\par
\switchcolumn
\EN
\noindent I am the mother of fair love, and of fear, and of knowledge, and of holy hope. In me is all grace of the way and of the truth, in me is all hope of life and of virtue. Come over to me, all ye that desire me, and be filled with my fruits. For my spirit is sweet above honey, and my inheritance above honey and the honeycomb. My memory is unto everlasting generations. They that eat me, shall yet hunger: and they that drink me, shall yet thirst. He that hearkeneth to me, shall not be confounded: and they that work by me, shall not sin.\par
\switchcolumn*

\LA
\begin{Resp}\Rsign Elégit eam Deus, et præelégit eam; \Star{} Corónam glóriæ pósuit super caput ejus. (Allelúja.)\end{Resp}
\begin{Resp}\Vsign Et in tabernáculo suo habitáre fecit eam.\end{Resp}
\begin{Resp}\Rsign Corónam glóriæ pósuit super caput ejus. (Allelúja.)\end{Resp}
\begin{Resp}\Vsign Glória Patri, et Fílio, \Star{} et Spirítui Sancto.\end{Resp}
\begin{Resp}\Rsign Corónam glóriæ pósuit super caput ejus. (Allelúja.)\end{Resp}
\switchcolumn
\EN
\begin{Resp}\Rsign God hath chosen her, and preferred her; \Star{} He hath placed a crown of glory on her head.\end{Resp}
\begin{Resp}\Vsign And he maketh her dwell in his tabernacle.\end{Resp}
\begin{Resp}\Rsign He hath placed a crown of glory on her head.\end{Resp}
\begin{Resp}\Vsign Glory be to the Father, and to the Son, \Star{} and to the Holy Ghost.\end{Resp}
\begin{Resp}\Rsign He hath placed a crown of glory on her head.\end{Resp}
\switchcolumn*

\LA
\LessonHead{Lectio IV}
\switchcolumn
\EN
\LessonHead{Lesson IV}
\switchcolumn*

\LA
\LessonTitle{Sermo sancti Petri Canísii Presbýteri}
\switchcolumn
\EN
\LessonTitle{Sermon of St. Peter Canisius, Priest}
\switchcolumn*

\LA
\Source{De Maria Deipara Virgine incomparabili, lib. 15, c. 13}
\switchcolumn
\EN
\Source{On the Incomparable Virgin Mary, Mother of God}
\switchcolumn*

\LA
\noindent Cur beatíssimam Vírginem Maríam Regínæ nómine, Damascénum, Athanásium aliósque secúti, non compellémus, cujus et pater David rex ínclitus, et fílius Rex regum Dominúsque dominántium sine fine ímperans, laudem in Scriptúris præstantíssimam tenent? Regína est ínsuper, si cum illis conferátur, quibus, véluti régibus, cæléste regnum cum Christo, Rege summo, cóntigit, útpote illíus coherédibus, et in eódem véluti throno, ut Scriptúra lóquitur, cum illo collocátis. Regína est étiam nulli electórum secúnda, sed simul Angelis et homínibus tanto præláta dígnius, quo nihil illa sublímius ac sánctius esse potest, quæ sola cum Deo Patre Fílium habet commúnem, et quæ supra se Deum et Christum tantum, infra se vero réliqua videt ómnia.\par
\switchcolumn
\EN
\noindent If we follow St. John Damascene, St. Athanasius and others, are we not forced to call Mary Queen, since her father David receiveth the highest praise in Scripture as a renowned king, and her son as King of kings and Lord of lords, reigning forever? She is Queen, moreover, when compared with the Saints who reign like kings in the heavenly kingdom, co-heirs with Christ, the great King, placed on the same throne with him, as the Scripture saith. And as Queen she is second to none of the elect, but in dignity is raised so high above both Angels and men that nothing can be higher or holier than she, who alone hath the same Son as God the Father, and who seeth above her only God and Christ, and below her creatures other than herself.\par
\switchcolumn*

\LA
\begin{Resp}\Rsign Súscipe verbum, Virgo María, quod tibi a Dómino transmíssum est: \Star{} Ecce concípies et páries Deum páriter et hóminem. (Allelúja.)\end{Resp}
\begin{Resp}\Vsign Et Regína vocáberis super omnes gentes.\end{Resp}
\begin{Resp}\Rsign Ecce concípies et páries Deum páriter et hóminem. (Allelúja.)\end{Resp}
\switchcolumn
\EN
\begin{Resp}\Rsign Receive the word, O Virgin Mary, sent to thee by the Lord: \Star{} Behold, thou wilt conceive and bring forth God as well as man.\end{Resp}
\begin{Resp}\Vsign And thou wilt be called Queen over all the nations.\end{Resp}
\begin{Resp}\Rsign Behold, thou wilt conceive and bring forth God as well as man.\end{Resp}
\switchcolumn*

\LA
\LessonHead{Lectio V}
\switchcolumn
\EN
\LessonHead{Lesson V}
\switchcolumn*

\LA
\noindent Magnus Athanásius perspícue dixit: Maríam non modo Deíparam, sed étiam Regínam et Dóminam próprie veréque censéri, quandóquidem Christus ex ipsa matre Vírgine natus, Deus et Dominus idémque Rex máneat. De hac ígitur Regína dictum illud Psalmógraphi interpretátur: Adstitit Regína a dextris tuis in vestítu deauráto. Non ergo solum cæli, sed et cælórum Regína recte dicitur María, útpote mater Regis Angelórum, Regísque cælórum et amíca et sponsa. A te vero, augustíssima Regína, eadémque fidíssima Mater María, quam sine fructu pie nullus implórat, et cui mortáles omnes beneficiórum memória sempitérna devíncti sunt, étiam atque étiam reverénter oro et óbsecro, ut hoc qualecúmque observántiæ erga te meæ testimónium ratum et gratum habére velis, utque obláti múneris exiguitátem studiósa offeréntis voluntáte metíri ac tuo præpoténti Fílio comprobáre dignéris.\par
\switchcolumn
\EN
\noindent The great Athanasius said clearly: Mary is not only the Mother of God, but also can be properly and truly called Queen and Lady, since in fact the Christ who was born of the Virgin Mother is God and Lord and also King. It is to this Queen, therefore, that the Psalmist's words are applied: The Queen taketh her place at thy right hand in garments of gold. Thus Mary is rightly called Queen, not only of heaven, but also of the heavens, as the Mother of the King of Angels, and as the Bride and beloved of the King of the heavens. O Mary, most august Queen and most faithful Mother, to whom no one prayeth in vain who prayeth devoutly, and to whom all mortal men are bound by the enduring memory of so many benefits, again and again reverently I beseech thee to accept and be pleased with every evidence of my devotion to thee, to value the poor gift I offer according to the zeal with which it is offered, and to recommend it to thine all-powerful Son.\par
\switchcolumn*

\LA
\begin{Resp}\Rsign Ecce pósitus est hic in ruínam et resurrectiónem multórum, \Star{} Et tuam ipsíus ánimam pertransíbit gládius. (Allelúja.)\end{Resp}
\begin{Resp}\Vsign Ave, Christi Mater, passiónis sócia, totíus mundi Regína.\end{Resp}
\begin{Resp}\Rsign Et tuam ipsíus ánimam pertransíbit gládius. (Allelúja.)\end{Resp}
\switchcolumn
\EN
\begin{Resp}\Rsign Behold, this Child is destined for the fall and the rise of many, \Star{} And thine own soul a sword shall pierce.\end{Resp}
\begin{Resp}\Vsign Hail, Mother of Christ, sharing in his passion, Queen of all the world.\end{Resp}
\begin{Resp}\Rsign And thine own soul a sword shall pierce.\end{Resp}
\switchcolumn*

\LA
\LessonHead{Lectio VI}
\switchcolumn
\EN
\LessonHead{Lesson VI}
\switchcolumn*

\LA
\LessonTitle{Ex Lítteris Encýclicis Pii Papæ duodécimi}
\switchcolumn
\EN
\LessonTitle{From the Encyclical Letter of Pope Pius XII}
\switchcolumn*

\LA
\Source{Encyclica « Ad cæli Regínam », diei 11 Octobris 1954}
\switchcolumn
\EN
\Source{Ad cæli Reginam, diei 11 Octobris 1954}
\switchcolumn*

\LA
\noindent E Christiánæ vetustátis monuméntis, e litúrgicis précibus, ex índito christiáno pópulo religiónis sensu, ex opéribus arte conféctis, úndique collégimus voces quæ ásserunt Deíparam Vírginem regáli dignitáte præstáre; ratiónes étiam, quas sancta theología ex divínæ fídei thesáuro deducéndo ástruit, eándem veritátem prorsus confirmáre argúimus. Tot ex allátis testimóniis quasi latíssime résonans concéntus effícitur, qui extóllit régii honóris præcélsum fastígium Dei hominúmque Matris, cui cuncta creáta subsunt, quæ est: « exaltáta super choros Angelórum ad cæléstia regna ». Cum vero, matúro ponderatóque consílio, persuásum Nobis habeámus magna oritúra esse Ecclésiæ emoluménta, si quasi in suo candelábro rutilántior lucérna pósita, illa sólide probáta véritas maniféstior ómnibus refúlgeat, Apostólica Nostra Potestáte decérnimus et institúimus festum Maríæ Regínæ, quod toto terrárum orbe quotánnis die trigésimo primo mensis maji est celebrándum.\par
\switchcolumn
\EN
\noindent From the documents of ancient Christianity, from the prayers of the liturgy, from the innate religious sense of the Christian people, from works of art, from all sides we gather witnesses which assert that the Virgin Mother of God excelleth in queenly dignity. And we have set forth the reasons which sacred theology deduces from the treasury of divine faith to confirm the same truth. All these witnesses form a sort of chorus, proclaiming far and wide the supreme queenly honour granted to the Mother of God and man, who is above all created things and exalted over the choirs of Angels to reign in heaven. Thus it is that after mature and thoughtful consideration we have been persuaded that great benefits would flow to the Church if, like a light that illumines more brightly when placed in its stand, this solidly proved truth were to shine out more clearly to all, and so, by Our Apostolic Authority, we decree and institute the feast of Mary, Queen, which is to be celebrated every year on the thirty-first day of May throughout the world.\par
\switchcolumn*

\LA
\begin{Resp}\Rsign Signum magnum appáruit in cælo: \Star{} Múlier amícta sole, et luna sub pédibus ejus, et in cápite ejus coróna stellárum duódecim. (Allelúja.)\end{Resp}
\begin{Resp}\Vsign Cujus Fílius regnat in ætérnum.\end{Resp}
\begin{Resp}\Rsign Múlier amícta sole, et luna sub pédibus ejus, et in cápite ejus coróna stellárum duódecim. (Allelúja.)\end{Resp}
\begin{Resp}\Vsign Glória Patri, et Fílio, \Star{} et Spirítui Sancto.\end{Resp}
\begin{Resp}\Rsign Múlier amícta sole, et luna sub pédibus ejus, et in cápite ejus coróna stellárum duódecim. (Allelúja.)\end{Resp}
\switchcolumn
\EN
\begin{Resp}\Rsign A great sign appeared in heaven: \Star{} A woman clothed with the sun, and the moon was under her feet, and upon her head a crown of twelve stars.\end{Resp}
\begin{Resp}\Vsign Whose Son reigneth forever.\end{Resp}
\begin{Resp}\Rsign A woman clothed with the sun, and the moon was under her feet, and upon her head a crown of twelve stars.\end{Resp}
\begin{Resp}\Vsign Glory be to the Father, and to the Son, \Star{} and to the Holy Ghost.\end{Resp}
\begin{Resp}\Rsign A woman clothed with the sun, and the moon was under her feet, and upon her head a crown of twelve stars.\end{Resp}
\switchcolumn*

\LA
\LessonHead{Lectio VII}
\switchcolumn
\EN
\LessonHead{Lesson VII}
\switchcolumn*

\LA
\LessonTitle{Léctio sancti Evangélii secúndum Lucam}
\switchcolumn
\EN
\LessonTitle{From the Holy Gospel according to Luke}
\switchcolumn*

\LA
\Cite{Luc 1:26-33}
\switchcolumn
\EN
\Cite{Luke 1:26-33}
\switchcolumn*

\LA
\noindent In illo témpore: Missus est Angelus Gábriel a Deo in civitátem Galilǽæ, cui nomen Názareth, ad Vírginem desponsátam viro, cui nomen erat Joseph de domo David, et nomen Vírginis María. Et réliqua.\par
\switchcolumn
\EN
\noindent In that time: In the sixth month, the angel Gabriel was sent from God into a city of Galilee, called Nazareth, to a virgin espoused to a man whose name was Joseph, of the house of David; and the virgin's name was Mary. And so on.\par
\switchcolumn*

\LA
\LessonTitle{Homilía sancti Bonaventúræ Epíscopi}
\switchcolumn
\EN
\LessonTitle{Homily of St. Bonaventure, Bishop}
\switchcolumn*

\LA
\Source{Sermo de regia dignitate Beatæ Mariæ Virginis}
\switchcolumn
\EN
\Source{Sermon on the Royal Dignity of the Blessed Virgin Mary}
\switchcolumn*

\LA
\noindent Beáta María Virgo, summi Regis mater est per generósam conceptiónem, secúndum quod dícitur in verbo sibi dicto ab Angelo: Ecce inquit, concípies et páries fílium; et póstea: Dabit ei Dóminus sedem David patris ejus, et regnábit in domo Jacob in ætérnum, et regni ejus non erit finis. Ac si apérte dicat: Concípies et páries fílium Regem, in regáli sólio æternáliter residéntem, ac per hoc tamquam Mater Regis regnábis, et ut Regína in regáli sólio residébis. Si enim decet filium honórem matri dare, decet ut ei commúnicet thronum regálem; unde Virgo María, quia concépit eum, qui habet in fémore scriptum: Rex regum et Dóminus dominántium, statim ex quo concépit Fílium Dei, Regína fuit, non solum terræ, sed étiam cæli, quod designátum est in Apocalýpsi ubi dícitur: Signum magnum appáruit in cælo: múlier amícta sole, et luna sub pédibus ejus, et in cápite ejus coróna stellárum duódecim.\par
\switchcolumn
\EN
\noindent The Blessed Virgin Mary is the Mother of the great King by reason of a noble kind of conception, according to the message given her by the Angel. Behold, he said, thou shalt conceive and shalt bring forth a son; and again, the Lord God will give him the throne of David his father, and he shall be king over the house of Jacob forever, and of his kingdom there shall be no end. This is as if to say in so many words, Thou wilt conceive and bear a son who is King, eternally reigning on the royal throne, and as Queen thou wilt be seated on the royal throne. For if it becometh a son to give honour to his mother, it is also fitting that he share his royal throne with her; and so the Virgin Mary, because she conceived him on whose thigh was written, King of kings and Lord of lords, was Queen not only of earth but also of heaven as soon as she conceived the Son of God. This is indicated in the Apocalypse where it saith: A great sign appeared in heaven: a woman clothed with the sun, and the moon was under her feet, and upon her head a crown of twelve stars.\par
\switchcolumn*

\LA
\begin{Resp}\Rsign Ecce Virgo concípiet et páriet Fílium, \Star{} Et vocábitur nomen ejus Admirábilis, Deus, Fortis. (Allelúja.)\end{Resp}
\begin{Resp}\Vsign Super sólium David et super regnum ejus sedébit in ætérnum.\end{Resp}
\begin{Resp}\Rsign Et vocábitur nomen ejus Admirábilis, Deus, Fortis. (Allelúja.)\end{Resp}
\switchcolumn
\EN
\begin{Resp}\Rsign Behold a Virgin shall conceive and bear a Son, \Star{} And his name shall be called Wonderful, God, the Mighty One.\end{Resp}
\begin{Resp}\Vsign He shall sit upon the throne of David, and reign over his kingdom forever.\end{Resp}
\begin{Resp}\Rsign And his name shall be called Wonderful, God, the Mighty One.\end{Resp}
\switchcolumn*

\LA
\LessonHead{Lectio VIII}
\switchcolumn
\EN
\LessonHead{Lesson VIII}
\switchcolumn*

\LA
\noindent María Regína est præclaríssima quantum ad glóriam, quod bene desígnat Prophéta in psalmo qui speciáliter est de Christo et Vírgine María, ubi primo dícitur de Christo: Sedes tua, Deus, in sǽculum sǽculi, et paulo post de Vírgine: Adstitit regína a dextris tuis, hoc est in potióribus bonis, quod quidem dictum est quantum ad glóriam mentis. Séquitur: In vestítu deauráto, in quo exprímitur vestítus gloriósæ immortalitátis, quem décuit habére Vírginem in sua Assumptióne. Absit enim ut vestiméntum illud quo opértus est Christus, quod étiam fuit perfécte sanctificátum in terra per Verbum incarnátum, fiat cibus vérmium. Sicut décuit Christum dare grátiam Matri suæ pleníssimam in sua Conceptióne, sic décuit tribúere pleníssimam glóriam in ipsíus Matris Assumptióne. Et ídeo tenéndum est quod Virgo, gloriósa in ánima et córpore, sédeat juxta Fílium.\par
\switchcolumn
\EN
\noindent Mary the Queen outshineth all others in glory, as the Prophet clearly showeth in the Psalm which particularly concerns Christ and the Virgin Mary. It first saith of Christ: Thy throne, O God, standeth forever and ever, and shortly thereafter of the Virgin: The Queen taketh her place at thy right hand, that is, in the position of highest blessedness, for it referreth to glory of soul. It continueth: In garments of gold, by which is meant the clothing of glorious immortality which was proper to the Virgin in her Assumption. For it could not be that the garment which clothed Christ, the garment completely sanctified on earth by the incarnate Word, should be the food of worms. As it was fitting for Christ to grant the fullness of grace to his Mother at her Conception, so was it fitting that he grant her the fullness of glory at her Assumption. And so we are to hold that the Virgin, glorious in soul and body, is enthroned next to her Son.\par
\switchcolumn*

\LA
\begin{Resp}\Rsign Exaltáta es, sancta Dei Génetrix, \Star{} Super choros Angelórum ad cæléstia regna. (Allelúja.)\end{Resp}
\begin{Resp}\Vsign Intercéde pro nobis ad Dóminum Jesum Christum.\end{Resp}
\begin{Resp}\Rsign Super choros Angelórum ad cæléstia regna. (Allelúja.)\end{Resp}
\begin{Resp}\Vsign Glória Patri, et Fílio, \Star{} et Spirítui Sancto.\end{Resp}
\begin{Resp}\Rsign Super choros Angelórum ad cæléstia regna. (Allelúja.)\end{Resp}
\switchcolumn
\EN
\begin{Resp}\Rsign Thou hast been exalted, holy Mother of God, \Star{} Over the choirs of Angels, to reign in heaven.\end{Resp}
\begin{Resp}\Vsign Intercede for us with the Lord Jesus Christ.\end{Resp}
\begin{Resp}\Rsign Over the choirs of Angels, to reign in heaven.\end{Resp}
\begin{Resp}\Vsign Glory be to the Father, and to the Son, \Star{} and to the Holy Ghost.\end{Resp}
\begin{Resp}\Rsign Over the choirs of Angels, to reign in heaven.\end{Resp}
\switchcolumn*

\LA
\LessonHead{Lectio IX}
\switchcolumn
\EN
\LessonHead{Lesson IX}
\switchcolumn*

\LA
\noindent María Regína est et dispensátrix grátiæ, quod designátum fuit in libro Esther, ubi dícitur: Parvus fons qui crevit in flúvium, et in lucem et solem convérsus est. Virgo Maria sub figura Esther comparátur diffusióni fontis et lúminis, propter diffusiónem grátiæ quantum ad dúplicem usum, actiónis licet et contemplatiónis. Grátia enim Dei, quæ est medicatíva géneris humáni, per ipsam ad nos descéndit quasi per aquædúctum, quia ad ipsam Vírginem pértinet dispensátio grátiæ non per modum princípii, sed per modum mériti. Mérito ígitur Virgo María est excellentíssima Regína in comparatiónem ad plebem, cum ímpetrat véniam, súperat pugnam, et distríbuit grátiam, et consequénter perdúcit ad glóriam.\par
\switchcolumn
\EN
\noindent Mary the Queen is also the distributrix of grace. This is indicated in the book of Esther, where it is said: The little spring which grew into a river and was turned into a light and into the sun. The Virgin Mary, under the type of Esther, is compared to the outpouring of a spring and of light, because of the diffusion of grace for two uses, that is, for action and for contemplation. For the grace of God, which is a healing for the human race, descendeth to us through her as if through an aqueduct, since the dispensing of grace is attributed to the Virgin not as to its beginning, but because of her position through merit. By position the Virgin Mary is a most excellent Queen towards her people: she obtaineth forgiveness, overcometh strife, distributeth grace, and thereby she leadeth them to glory.\par
\switchcolumn*

\LA
\TeDeum{Te Deum}
\switchcolumn
\EN
\TeDeum{Te Deum}
\switchcolumn*

\end{paracol}

\SeasonTitle{Mensis Junius}{June}

\addcontentsline{toc}{section}{Mensis Junius \textit{— June}}

\par\needspace{10\baselineskip}\addvspace{1.2em}{\centering\small\textsc{Die 1 Junii} · \textsc{1 June}\par}
\Office{S. Angelæ Mericiæ Virginis}{St. Angela Merici, Virgin}{Duplex}
\addcontentsline{toc}{subsection}{Die 1 Junii · S. Angelæ Mericiæ Virginis \textit{— St. Angela Merici, Virgin}}
\begin{paracol}{2}
\LA
\RefLine{Lectiones I–III: de Scriptura occurrente.}
\switchcolumn
\EN
\RefLine{Lessons I–III: from the Scripture of the day.}
\switchcolumn*

\LA
\RefLine{Lectiones VII–IX: ex Commúni Virginum tantum (C6a).}
\switchcolumn
\EN
\RefLine{Lessons VII–IX: from the Common of a Virgin not a Martyr (C6a).}
\switchcolumn*

\LA
\LessonHead{Lectio IV}
\switchcolumn
\EN
\LessonHead{Lesson IV}
\switchcolumn*

\LA
\noindent Angela Merícia, Decentiáni, Veronénsis diœcésis óppido ad lacum Benácum in dicióne Venéta, piis orta paréntibus, a prima ætáte virginitátis lílium, quod perpétuo serváre statúerat, sédula sepsit. Ab omni mulíebri ornátu abhórrens, egrégiam vultus formam pulchrámque cæsáriem studióse fœdávit, ut cælésti dumtáxat animárum Sponso placéret. In ipso autem adolescéntiæ flore paréntibus orbáta, austerióris vitæ desidério in desértum locum aufúgere tentávit; sed, ab avúnculo prohíbita, novit præstáre domi, quod in solitúdine non lícuit. Cilício ac flagéllis frequénter usa, carnem nónnisi infírma valetúdine, vinum in Nativitátis et Resurrectiónis Domínicæ tantum celebritáte; complúres vero dies, nihil omníno degustávit. Oratióni dédita, brevíssimum humi carpébat somnum. Dæmónem vero sub lucéntis ángeli forma sibi illúdere conántem agnóvit prótinus, et conjécit in fugam. Tandem patérnis bonis abdicátis, et hábitum ac régulam tértii órdinis sancti Francísci ampléxa, evangélicam paupertátem virginitátis laudi conjúnxit.\par
\switchcolumn
\EN
\noindent Angela Merici was born of godly parents at Decenzano on the western shore of the Lake of Garda, in the diocese of Verona and territory of Venice, (on the 21st day of March, about the year of grace 1474.) From her earliest years she carefully guarded the lily of her virginity, with the intention of keeping it for ever unbroken. She had no taste for women's finery, and purposely marred the exceeding comeliness of her body and her sightly hair, as seeking to appear beautiful only in the eyes of Him Who is the Lover of souls. At ten years of age she lost both her father and mother, and thereafter, being fain to take upon her a life of greater hardness, she essayed to retire into a desert place apart, but this her uncle forbade her to do, and she learnt how to practice at home what she was not allowed to attempt in the wilderness. She often used hair-cloth and scourging never ate flesh-meat, except when she was sick drank wine only on the Feast-days of Christmas and Easter; and many a day took nothing at all. She was instant in prayer. What little sleep she took, she took lying on the ground. The devil strove to beguile her, appearing under the form of an angel of light, but she quickly detected him and put him to flight. At length she added to the glory of virginity that poverty which is commended in the Gospel she gave up all that she had, and adopted the dress and rule of the Third Order of St Francis.\par
\switchcolumn*

\LA
\begin{Resp}\Rsign Propter veritátem, et mansuetúdinem, et justítiam: \Star{} Et dedúcet te mirabíliter déxtera tua, allelúja.\end{Resp}
\begin{Resp}\Vsign Spécie tua et pulchritúdine tua inténde, próspere procéde, et regna.\end{Resp}
\begin{Resp}\Rsign Et dedúcet te mirabíliter déxtera tua, allelúja.\end{Resp}
\switchcolumn
\EN
\begin{Resp}\Rsign Because of truth, and meekness, and righteousness; \Star{} And thy right hand shall lead thee wonderfully, alleluia.\end{Resp}
\begin{Resp}\Vsign In thy comeliness, and thy beauty, go forward, fare prosperously, and reign.\end{Resp}
\begin{Resp}\Rsign And thy right hand shall lead thee wonderfully, alleluia.\end{Resp}
\switchcolumn*

\LA
\LessonHead{Lectio V}
\switchcolumn
\EN
\LessonHead{Lesson V}
\switchcolumn*

\LA
\noindent Nullum pietátis offícium erga próximos omíttens, paupéribus, quidquid sibi ex mendicáto victu superésset, largiebátur; libénter ministrábat ægrótis, plúraque cum magna sanctitátis fama peragrávit loca, ut vel solátio esset afflíctis, vel reis véniam impetráret, vel infénsos ínvicem reconciliáret ánimos, vel e vitiórum cœno sceléstos revocáret. Angelórum pane, quem únice esuriébat, frequentíssime refécta, tanta caritátis vi ferebátur in Deum, ut sæpius extra sensus raperétur. Sacra Palæstínæ loca, summa cum religióne, obívit; quo in itínere et visum, quem ad Cydónias appúlsa oras amíserat, eódem regréssa recuperávit, et barbarórum captivitátem ac naufrágium ímminens divínitus evásit. Romam dénique, firmam Ecclésiæ petram veneratúra et amplíssimæ jubilæi véniæ percúpida, sedénte Cleménte séptimo, accéssit; quam summus Póntifex allocútus, ejúsdem sanctimóniam suspéxit et commendávit summópere; nec ab Urbe ipsam abíre ante permísit, quam álio cælitus vocátam agnóvit.\par
\switchcolumn
\EN
\noindent She left undone no service of kindness which she was able to do to her neighbors. If there remained anything over of the food which was given in alms to herself, she gave that to the poor. She cheerfully waited upon the sick. She journeyed about, with a great reputation for holiness, comforting the afflicted, asking forgiveness for the guilty, reconciling the angry, and recalling the wicked from evil. Her only hunger was for the bread of Angels, and she took the Same right often, and then arose in her vehemence of love bearing her towards God, which oftentimes made her beside herself. She made a pilgrimage, with intense feeling, to the Holy Places in Palestine, during which journey she lost her sight at Canea in Crete on her way out, and recovered it at the same place on her way home. In this journey also, God saved her from being made prisoner by the unbelievers and from shipwreck. She went to Rome, (in 1525,) at once to pray at the immovable Rock of the Church, and to gain the abundant pardons of the Jubilee. Pope Clement VII conversed with her, was edified by her holiness, and highly commended her neither would he let her leave Rome, until he knew that God was calling her elsewhere.\par
\switchcolumn*

\LA
\begin{Resp}\Rsign Dilexísti justítiam, et odísti iniquitátem: \Star{} Proptérea unxit te Deus, Deus tuus, óleo lætítiæ, allelúja.\end{Resp}
\begin{Resp}\Vsign Propter veritátem, et mansuetúdinem, et justítiam.\end{Resp}
\begin{Resp}\Rsign Proptérea unxit te Deus, Deus tuus, óleo lætítiæ, allelúja.\end{Resp}
\switchcolumn
\EN
\begin{Resp}\Rsign Thou hast loved righteousness, and hated iniquity; \Star{} Therefore God, thy God, hath anointed thee with the oil of gladness, alleluia.\end{Resp}
\begin{Resp}\Vsign Because of truth, and meekness, and righteousness.\end{Resp}
\begin{Resp}\Rsign Therefore God, thy God, hath anointed thee with the oil of gladness, alleluia.\end{Resp}
\switchcolumn*

\LA
\LessonHead{Lectio VI}
\switchcolumn
\EN
\LessonHead{Lesson VI}
\switchcolumn*

\LA
\noindent Bríxiam ítaque, ubi domum ad sanctæ Afræ templum condúxit, revérsa, novam ibi vírginum societátem, sicut cælésti voce ac visióne mandátum sibi fúerat, sub certa disciplína sanctísque vivéndi régulis constítuit, quam sanctæ Ursulæ invíctæ vírginum ducis patrocínio ac nómine insignívit: eam vero perénnem futúram morti próxima prædíxit. Tandem prope septuagenária, dives méritis evolávit in cælum sexto Kaléndas Februárii anni millésimi quingentésimi et quadragésimi. Cujus cadáver per ipsos trigínta dies inhumátum, flexíbile ac vivo simíllimum perseverávit. Demum in sanctæ Afræ templo inter céteras, quibus illud abúndat, Sanctórum relíquias repósito, plúrima ad ejus sepúlcrum agi statim cœpére mirácula; quorum fama late diffúsa, non Bríxiæ modo et Decentiáni, sed álibi étiam vulgo cœpit nuncupári beáta, ejúsque imágo aris impóni: immo sanctus ipse Cárolus Borromæus non multis post annis dignam, quæ ab apostólica Sede in sanctárum Vírginum album referrétur, Bríxiæ palam asséruit. Cultum vero illi jámdiu a pópulis exhíbitum, et et tum locórum Ordináriis probátum, tum plúribus étiam summórum Pontíficum indúltis munítum, Clemens Papa décimus tértius solémni decréto ratum hábuit ac confirmávit. Eam tandem, novis miráculis rite probátis insígnem, Pius Papa séptimus, solémni canonizatióne in Vaticána basílica perácta, die vigésima quarta Maji anno millésimo octingentésimo séptimo, sanctárum Vírginum catálogo adscrípsit.\par
\switchcolumn
\EN
\noindent She went back to Brescia, and there hired an house near the Church of St. Afra, in which house, in obedience to a vision and command from heaven, she founded a new Order of religious women, constituted under certain rules and holy regulations of life. This Order she put under the name and patronage of St Ursula, the fearless leader of maidens. When Angela was near to death, she foretold that this Order will never cease. She was well-nigh three score and ten years of age, and full of good works, when, in the (night between the) 27th (and the 28th days) of January, in the year 1540, she winged her flight heavenward. Her dead body lay unburied thirty days, supple and life-like. It was laid at last in the Church of St. Afra, where sleep so many more of God's holy children. Diverse miracles forthwith began to be worked at her grave. The fame of these being noised about, she began to be commonly called Blessed, and that not only at Brescia and Decenzano and pictures of her were put over Altars. Not many years afterward, holy Charles Borromeo said openly at Brescia, that she was one whose name the Apostolic See might well enroll among those of holy virgins. The reverence which had of a long time been shown to her memory was approved by the local Ordinaries, confirmed by diverse Papal Indults, and solemnly ratified and established by decree of Pope Clement XIII. As she continued famous for new and proved miracles, Pope Pius VII, at the solemn canonization held in the Vatican Basilica, upon the 24th day of May, in the year 1807, added her name to the list of holy maids.\par
\switchcolumn*

\LA
\begin{Resp}\Rsign Afferéntur Regi vírgines post eam, próximæ ejus; \Star{} Afferéntur tibi in lætítia et exsultatióne, allelúja.\end{Resp}
\begin{Resp}\Vsign Spécie tua et pulchritúdine tua; inténde, próspere procéde, et regna.\end{Resp}
\begin{Resp}\Rsign Afferéntur tibi in lætítia et exsultatióne, allelúja.\end{Resp}
\begin{Resp}\Vsign Glória Patri, et Fílio, \Star{} et Spirítui Sancto.\end{Resp}
\begin{Resp}\Rsign Afferéntur tibi in lætítia et exsultatióne, allelúja.\end{Resp}
\switchcolumn
\EN
\begin{Resp}\Rsign After her shall virgins be brought unto the King, her fellows \Star{} They shall be brought unto thee with gladness and rejoicing, alleluia.\end{Resp}
\begin{Resp}\Vsign In thy comeliness and thy beauty, go forward, fare prosperously, and reign.\end{Resp}
\begin{Resp}\Rsign They shall be brought unto thee with gladness and rejoicing, alleluia.\end{Resp}
\begin{Resp}\Vsign Glory be to the Father, and to the Son, \Star{} and to the Holy Ghost.\end{Resp}
\begin{Resp}\Rsign They shall be brought unto thee with gladness and rejoicing, alleluia.\end{Resp}
\switchcolumn*

\LA
\LessonHead{Lectio IX (pro commemoratione)}
\switchcolumn
\EN
\LessonHead{Lesson IX (when commemorated)}
\switchcolumn*

\LA
\LessonTitle{Commemoratio: S. Angelæ Mericiæ Virginis}
\switchcolumn
\EN
\LessonTitle{Commemoration of: St. Angela Merici, Virgin}
\switchcolumn*

\LA
\noindent Angela Merícia, piis orta paréntibus, a prima ætáte magna virtútum specímina dedit, cilício ac flagéllis frequénter usa, et oratióni indesinénter dédita. Patérnis bonis abdicátis, ac régulam tértii órdinis sancti Francísci ampléxa, evangélicam paupertátem virginitátis laudi conjúnxit, nullúmque pietátis offícium erga próximos omísit. Sacra Eucharístia frequentíssime refécta, tanta caritátis vi ferebátur in Deum, ut sǽpius extra sensus raperétur. Bríxiæ novam vírginum societátem sub certa disciplína sanctísque vivéndi régulis constítuit, quam sanctæ Ursulæ patrocínio ac nómine insignívit. Tandem prope septuagenária evolávit in cælum, anno Dómini millésimo quingentésimo quadragésimo, sexto kaléndas februárii. Cultum illi jámdiu exhíbitum Clemens Papa décimus tértius solémni decréto ratum hábuit et confirmávit. Pius vero Papa séptimus sanctárum Vírginum catálogo eam adscrípsit.\par
\switchcolumn
\EN
\noindent Angela Merici was born of devout parents and from her earliest years gave evidence of great virtue, frequently using sackcloth and scourges, and praying unceasingly. She renounced her inheritance and embraced the rule of the Third Order of St. Francis. Uniting evangelical poverty with the glory of virginity, she withheld no service of kindness to her neighbour. She was frequently refreshed by the Holy Eucharist, and was carried up to God with such force of love as often to be rapt out of her senses. At Brescia, she founded a new society of Virgins, under a fixed discipline and a rule of life conducive to holiness, which she put under the patronage and name of St. Ursula. At last, when she was almost seventy years old, she went to heaven, in the year of the Lord 1540, on January 27th. Pope Clement XIII by a solemn decree ratified and confirmed the veneration which had already been offered her for a long time; and Pope Pius VII enrolled her in the list of holy Virgins.\par
\switchcolumn*

\LA
\TeDeum{Te Deum}
\switchcolumn
\EN
\TeDeum{Te Deum}
\switchcolumn*

\end{paracol}

\par\needspace{10\baselineskip}\addvspace{1.2em}{\centering\small\textsc{Die 2 Junii} · \textsc{2 June}\par}
\Office{Ss. Marcellini, Petri, atque Erasmi, Episcopi, Martyrum}{Sts. Marcellinus, Peter, and Erasmus, Martyrs}{Simplex}
\addcontentsline{toc}{subsection}{Die 2 Junii · Ss. Marcellini, Petri, atque Erasmi, Episcopi, Martyrum \textit{— Sts. Marcellinus, Peter, and Erasmus, Martyrs}}
\begin{paracol}{2}
\LA
\RefLine{Lectiones I–II: de Scriptura occurrente.}
\switchcolumn
\EN
\RefLine{Lessons I–II: from the Scripture of the day.}
\switchcolumn*

\LA
\LessonHead{Lectio III (pro commemoratione)}
\switchcolumn
\EN
\LessonHead{Lesson III (when commemorated)}
\switchcolumn*

\LA
\noindent Petrus exorcísta, Diocletiáno imperatóre, Romæ a Seréno júdice propter christiánæ fídei confessiónem missus in cárcerem, Paulínam Artémii, qui cárceri præerat, filiam a dæmone agitátam liberávit. Quo facto et paréntes puéllæ cum tota família, et vicínos qui ad rei novitátem concúrrerant, Jesu Christo conciliátos ad Marcellínum presbyterum addúxit, a quo omnes baptizáti sunt. Quod ubi rescívit Serénus, Petrum et Marcellínum ad se vocátos aspérius objúrgat, et ad verbórum acerbitátem minas ac terróres adjúngit, nisi Christo renúntient. Cui cum Marcellínus christiána libertáte respondéret, pugnis contúsum, et a Petro sejúnctum, nudum inclúdit in cárcerem stratum vitri fragméntis, sine cibo ac sine lúmine. Petrum item constríngi ímperat arctíssimis vínculis. Sed cum utríque ex torméntis fides et ánimus crésceret, constánti confessióne, et abscísso cápite, illústre testimónium Jesu Christo dedérunt. Erásmus epíscopus, imperatóribus Diocletiáno et Maximiáno, in Campánia plumbátis et fústibus cæsus, resína quoque, súlphure, plumbo liquefácto, et fervénti pice, cera, oleóque perfúsus, inde tamen ínteger et inviolátus evásit. Quo miráculo multi se ad Christi fidem convertérunt. Verum is, íterum detrúsus in cárcerem, constríctus férreis gravissimísque vínculis, inde ab Angelo mirabíliter eréptus est. Deínde Fórmiis a Maximiáno váriis afféctus supplíciis, tunicáque ærea candénti indútus, illa étiam torménta, divína virtúte superávit. Dénique, plúrimis et in fide confirmátis, et ad fidem convérsis, insígnem martyrii palmam adéptus est.\par
\switchcolumn
\EN
\noindent This Peter was an exorcist, whom, in the reign of the Emperor Diocletian, Serenus the Judge cast into prison at Rome because he confessed the Christian faith. He there set free Paulina, the daughter of Artemius, the keeper of the prison, from an evil spirit which tormented her. Upon this, Artemius and his wife and all their house, with their neighbours who had run together to see the strange thing, would fain be made friends with Jesus Christ. Peter therefore brought them to Marcellinus, the Priest, who baptized them all. When Serenus heard of it, he called Peter and Marcellinus before him, and sharply rebuked them, adding to his bitter words threats and terrors, unless they would deny Christ. Marcellinus answered him with Christian boldness, whereupon he caused him to be buffeted, separated him from Peter and shut him up naked in a prison strewn with broken glass, without either food or light. Peter also he straitly confined. But when both of them were found to wax faithfuller and braver in their bonds, they were beheaded, unshaken in their testimony, and confessing Jesus Christ gloriously by their blood. Elmo was a Bishop in Campania who, (in the year 303,) in the reign of the Emperors Diocletian and Maximian was beaten with clubs and whips loaded with lead, and afterwards anointed with melted pitch, sulphur, and lead, and boiling resin, wax, and oil. From all this he came forth whole and sound which wonder turned many to believe in Christ. He was remanded again to prison, and straitly bound in heavy iron fetters. But from these he was wondrously delivered by an angel. At last, at Formi, Maximian caused him to be subjected to diverse torments, and in the end being clad in a coat of redhot brass the power of God made him to be more than conqueror in this thing also, and to grasp the palm-branch of a glorious testimony, whereby he strengthened many in the faith and turned many to it.\par
\switchcolumn*

\LA
\TeDeum{Te Deum}
\switchcolumn
\EN
\TeDeum{Te Deum}
\switchcolumn*

\end{paracol}

\par\needspace{10\baselineskip}\addvspace{1.2em}{\centering\small\textsc{Die 4 Junii} · \textsc{4 June}\par}
\Office{S. Francisci Caracciolo Confessoris}{St. Francis Caracciolo, Confessor}{Duplex}
\addcontentsline{toc}{subsection}{Die 4 Junii · S. Francisci Caracciolo Confessoris \textit{— St. Francis Caracciolo, Confessor}}
\begin{paracol}{2}
\LA
\RefLine{Lectiones I–III: de Scriptura occurrente.}
\switchcolumn
\EN
\RefLine{Lessons I–III: from the Scripture of the day.}
\switchcolumn*

\LA
\RefLine{Lectiones VII–IX: ex Commúni Confessoris non pontificis (C5).}
\switchcolumn
\EN
\RefLine{Lessons VII–IX: from the Common of a Confessor not a Bishop (C5).}
\switchcolumn*

\LA
\LessonHead{Lectio IV}
\switchcolumn
\EN
\LessonHead{Lesson IV}
\switchcolumn*

\LA
\noindent Francíscus, dictus ántea Ascánius, ex nóbili família Carácciolo in óppido sanctæ Maríæ de Villa in Aprútio ortus, a primis annis exímio enítuit pietátis cultu. Adoléscens gráviter ægrótans státuit sese prorsus Dei proximíque mancipáre servítio. Neápolim proféctus, sacerdótio initiátus sacróque adscríptus sodalítio, contemplatióni lucrandísque animábus se totum devóvit, ac extrémo supplício damnátis hortatórem se præbuit assíduum. Cóntigit autem, ut epístolum álteri destinátum ei per errórem redderétur, quo a piíssimis viris Joánne Augustíno Adórno et Fabrício Carácciolo ad novi religiósi institúti fundatiónem vocabátur. Rei novitáte captus et divínæ voluntátis demirátus consília, álacri ánimo sese illis adjúnxit. Cónditis autem in Camaldulénsium erémo, quo secésserant, novi órdinis légibus, inde Romam simul profécti confirmatiónem a Xysto quinto impetrárunt, qui eósdem cléricos reguláres Minóres appellári vóluit, áddito ad tria consuéta áltero de non ambiéndis dignitátibus voto.\par
\switchcolumn
\EN
\noindent This Francis, whose worldly name was Ascanius, was one of the noble family of Caracciolo. He was born in the town of Santa Maria della Villa, in the Abruzzi, \textasciitilde{}(on the 13th day of October, in the year of grace 1563.) From his earliest years he showed great marks of godliness. When he was a young man he had a severe illness, and on his recovery determined to serve God only, and bade farewell to the world. He betook himself to Naples, where he was ordained Priest, enrolled himself in a devout guild, and gave himself up altogether to seek after God, and to gain souls for Him, in which work he showed himself an unwearied comforter to such prisoners as were condemned to death. It came to pass that those two great servants of God, John Augustine Adorno and Fabricius Caracciolo, wrote a letter to a certain person, wherein they exhorted him to found a new religious Institute. This letter came by a mistake to be delivered to Francis Caracciolo. The newness of the idea and the strange ways of God's Providence took possession of his mind, and he joyfully added himself to their company. They withdrew themselves to the wilderness of the Camaldolese hermits (near Naples,) and there concerted the Rule of the New Order. Thence they went together to Rome, and obtained the confirmation of their work from Sixtus V, who was pleased that they should be called The Lesser Clerks Regular, since they add to the three accustomed vows (of Poverty, Chastity, and Obedience,) a fourth, binding themselves not to seek preferment in the Church.\par
\switchcolumn*

\LA
\begin{Resp}\Rsign Honéstum fecit illum Dóminus, et custodívit eum ab inimícis, et a seductóribus tutávit illum: \Star{} Et dedit illi claritátem ætérnam, allelúja.\end{Resp}
\begin{Resp}\Vsign Justum dedúxit Dóminus per vias rectas, et osténdit illi regnum Dei.\end{Resp}
\begin{Resp}\Rsign Et dedit illi claritátem ætérnam, allelúja.\end{Resp}
\switchcolumn
\EN
\begin{Resp}\Rsign The Lord made him honourable, and defended him from his enemies, and kept him safe from those that lay in wait for him; \Star{} And gave him perpetual glory, alleluia.\end{Resp}
\begin{Resp}\Vsign He went down with him into the pit, and left him not in bonds.\end{Resp}
\begin{Resp}\Rsign And gave him perpetual glory, alleluia.\end{Resp}
\switchcolumn*

\LA
\LessonHead{Lectio V}
\switchcolumn
\EN
\LessonHead{Lesson V}
\switchcolumn*

\LA
\noindent Solémni emíssa professióne, ob singulárem ejus in divum Francíscum Assisinátem cultum, Francísci nomen assúmpsit. Adórno biénnio post vita functo, ipse toti religióni, quamquam invítus, præfícitur; quo in múnere virtútum ómnium præclára præbuit exémpla. Institúti amplificándi studiosíssimus, id assíduis oratiónibus, lácrimis, et jugi córporis maceratióne eníxe a Deo postulábat. Quam ob rem tértio in Hispániam se cóntulit, peregríni hábitu indútus, victúmque ostiátim mendícans. In itínere aspérrima quæque perpéssus, Omnipoténtis auxílium mirum in modum expértus, navim quam conscénderat, ab imminénti naufrágio oratiónis præsídio servávit incólumem. Ut in regnis illis voti compos fíeret, plúrimum laborávit; sed, ejus sanctitátis fama prælucénte, amplissimáque catholicórum regum Philíppi secúndi et Philíppi tértii munificéntia, adversariórum conátibus singulári ánimi fortitúdine superátis, plura sui órdinis domicília fundávit: quod pari evéntu per Itáliam præstitit.\par
\switchcolumn
\EN
\noindent Ascanius Caracciolo, moved by a special love and devotion he had to the holy Francis of Assisi, took, when he made his solemn profession, the name of Francis. After two years, John Adorno departed this life, and Francis, against his own will, was made Head of the Order. In this office he shone a burning light of grace. Devoted to the prosperity of the Institute, he earnestly sought the blessing of God upon it, by constant prayer, by tears, and by stern treatment of his own body. In this work, he thrice travelled into Spain in the guise of a pilgrim, and begging his bread from door to door. In these his journeys he suffered very great hardships, and was most wonderfully helped by the Almighty, especially one while when he was on shipboard and the ship nigh to perish, but for the work of his prayers. He toiled hard in those countries to attain ' his wishes, but through the widespread fame of his holy life, and the noble generosity of the Most Catholic Kings Philip II. and Philip III., he overcame with his brave perseverance the opposition of all that withstood him, and founded several houses of his Order. This he was able to do in Italy also.\par
\switchcolumn*

\LA
\begin{Resp}\Rsign Amávit eum Dóminus, et ornávit eum: stolam glóriæ índuit eum, \Star{} Et ad portas paradísi coronávit eum, allelúja.\end{Resp}
\begin{Resp}\Vsign Índuit eum Dóminus lorícam fídei, et ornávit eum.\end{Resp}
\begin{Resp}\Rsign Et ad portas paradísi coronávit eum, allelúja.\end{Resp}
\switchcolumn
\EN
\begin{Resp}\Rsign The Lord loved him and beautified him; He clothed him with a robe of glory, \Star{} And crowned him at the gates of Paradise, alleluia.\end{Resp}
\begin{Resp}\Vsign The Lord hath put on him the breast-plate of faith, and hath adorned him.\end{Resp}
\begin{Resp}\Rsign And crowned him at the gates of Paradise, alleluia.\end{Resp}
\switchcolumn*

\LA
\LessonHead{Lectio VI}
\switchcolumn
\EN
\LessonHead{Lesson VI}
\switchcolumn*

\LA
\noindent Humilitáte ádeo excélluit, ut Romam véniens, in páuperum hospítio recéptus, se lepróso sociáverit, et ecclesiásticas dignitátes a Paulo quinto sibi oblátas constantíssime recusáverit. Illibátam perpétuo servávit virginitátem, effrontésque mulíeres, ejus castimóniæ insidiántes, Christo lucrifécit. Erga diviníssimum Eucharístiæ mystérium ardénti æstuans amóre, noctes pene íntegras in ejus adoratióne insómnes ducébat: quod pium exercítium, véluti sui órdinis tésseram, in eo perpétuo servándum constítuit. Deíparæ Vírginis cultum impénse fovit. In próximum exímia exársit caritáte. Prophetíæ dono et córdium scrutatióne ditátus fuit. Quadragésimum quartum ætátis suæ annum agens, dum in sacra Lauretána, æde in oratióne persísteret, sibi vitæ finem imminére cognóvit. Aprútium statim defléxit, et in óppido Agnóni apud alúmnos sancti Philíppi Nerii letháli febre corréptus, sacraméntis Ecclésiæ devotíssime suscéptis, prídie Nonas Júnii anni millésimi sexcentésimi octávi, in pervigílio festi Córporis Christi, placidíssime obdormívit in Dómino. Sacrum ejus corpus, Neápolim delátum, in ecclésia sanctæ Maríæ Majóris, ubi prima sui órdinis jécerat fundaménta, honorífice cónditum fuit. Eum póstea, miráculis clarum, Clemens decimus quartus Póntifex máximus solémni ritu inter Beátos; Pius vero séptimus Póntifex máximus, novis fulgéntem signis, anno millésimo octingentésimo séptimo, Sanctórum albo adscrípsit.\par
\switchcolumn
\EN
\noindent There was a great pattern of lowliness, so that when he came to Rome he betook himself to an almshouse, and chose a leper for his familiar friend. Paul V. offered him diverse honours in the Church, but he firmly refused them all. He preserved his purity unspotted, and when certain shameless women set themselves to attack his chastity, he took the occasion to gain over their souls for Christ. Toward God's great mystery of the Eucharist he was drawn with passionate tenderness, and would pass almost whole nights without sleep, simply adoring It. This godly custom he established in his Order, to be kept up therein for ever, the peculiar mark thereof. He was a great encourager of the worship of the Maiden Mother of God. He was hot with strong love for his neighbour. He was gifted with prophecy, and the discerning of spirits. In the forty-fourth year of his age he was continuing long in prayer in the Holy House of Loretto, when it was made known to him that the end of his earthly life was at hand. He straightway took his way to the Abruzzi, and was there seized with illness while he was with the disciples of St. Philip Neri, in the town of Agnone. He received with great devotion the Sacraments of the Church, and then, upon the 4th day of June, being the Eve of the Feast of the Body of Christ, in the year 1608, he very peacefully fell asleep in the Lord. His sacred body was carried to Naples, and there honourably buried in the Church of St. Mary the Greater, where he had laid the first foundations of his Order. As he became distinguished for miracles Pope Clement XIV enrolled his name, with solemn pomp, among those of the Blessed, and Pope Pius VII, in the year 1807, finding his mighty works continue, added it to the list of the Saints.\par
\switchcolumn*

\LA
\begin{Resp}\Rsign Iste homo perfécit ómnia quæ locútus est ei Deus, et dixit ad eum: Ingrédere in réquiem meam: \Star{} Quia te vidi justum coram me ex ómnibus géntibus, allelúja.\end{Resp}
\begin{Resp}\Vsign Iste est, qui contémpsit vitam mundi, et pervénit ad cæléstia regna.\end{Resp}
\begin{Resp}\Rsign Quia te vidi justum coram me ex ómnibus géntibus, allelúja.\end{Resp}
\begin{Resp}\Vsign Glória Patri, et Fílio, \Star{} et Spirítui Sancto.\end{Resp}
\begin{Resp}\Rsign Quia te vidi justum coram me ex ómnibus géntibus, allelúja.\end{Resp}
\switchcolumn
\EN
\begin{Resp}\Rsign This is he which did according unto all that God commanded him; and God said unto him: Enter thou into My rest, \Star{} For thee have I seen righteous before Me among all people, alleluia.\end{Resp}
\begin{Resp}\Vsign This is he which loved not his life in this world, and is come unto an everlasting kingdom.\end{Resp}
\begin{Resp}\Rsign For thee have I seen righteous before Me among all people, alleluia.\end{Resp}
\begin{Resp}\Vsign Glory be to the Father, and to the Son, \Star{} and to the Holy Ghost.\end{Resp}
\begin{Resp}\Rsign For thee have I seen righteous before Me among all people, alleluia.\end{Resp}
\switchcolumn*

\LA
\LessonHead{Lectio IX (pro commemoratione)}
\switchcolumn
\EN
\LessonHead{Lesson IX (when commemorated)}
\switchcolumn*

\LA
\LessonTitle{Commemoratio: S. Francisci Caracciolo Confessoris}
\switchcolumn
\EN
\LessonTitle{Commemoration of: St. Francis Caracciolo, Confessor}
\switchcolumn*

\LA
\noindent Francíscus, dictus ántea Ascánius, ex nóbili família Carácciolo, in óppido sanctæ Maríæ de Villa in Aprútio ortus est. Adoléscens gráviter ægrótans státuit sese prorsus Dei proximíque mancipáre servítio. Neápolim proféctus et sacerdótio initiátus, contemplatióni lucrandísque animábus se totum devóvit, ac extrémo supplício damnátis hortatórem se prǽbuit assíduum. Joánni Augustíno Adórno et Fabrício Carácciolo, mira Dei dispositióne, adjúnctus, Clericórum regulárium Minórum órdinem instítuit, áddito ad tria consuéta áltero de non ambiéndis dignitátibus voto; quem, post óbitum Adórni, sanctíssime rexit, et summo stúdio per Hispániam et Itáliam propagávit. Erga sanctíssimæ Eucharístiæ sacraméntum tanto æstuábat afféctu, ut noctes pene íntegras in ejus adoratióne impénderet; quod pium exercítium, véluti sui órdinis tésseram, perpétuo in eo servándum constítuit. Tandem prophetíæ dono et córdium scrutatióne ditátus, quadragésimum quartum annum agens, in óppido Agnóni in Aprútio letháli febre corréptus, in Dómino obdormívit prídie nonas júnii, anno millésimo sexcentésimo octávo. Sacrum ejus corpus, Neápolim delátum, in ecclésia sui órdinis cónditum est.\par
\switchcolumn
\EN
\noindent Francis, originally called Ascanio, came from the noble family of the Caracciolo, in the city of Santa Maria della Villa in the Abruzzi. During a severe illness in his youth, he resolved to give himself up completely to the service of God and neighbour. He went to Naples and was ordained priest, after which he devoted himself wholly to contemplation and the winning of souls; and he was always on hand to comfort prisoners condemned to death. By a wonderful disposition of God, he joined with John Augustine Adorno and Fabrizio Caracciolo in founding the Order of Clerks Regular Minor, adding to the three usual vows another: not to seek dignities. After the death of Adorno, Francis ruled the Order in a most holy way, and with the greatest zeal promoted it throughout Spain and Italy. He burned with such love for the Most Blessed Sacrament that he would spend almost the whole night in adoring it; and he established adoration of the Blessed Sacrament as a practice to be kept for ever in his Order, the peculiar mark thereof. He was enriched with the gifts of prophecy and the reading of hearts. At length, in his forty-fourth year, he was taken with a deadly fever in the town of Agnone, in the Abruzzi, and fell asleep in the Lord on the 4th day of June, 1608. His holy body was taken to Naples and buried in the church of his Order.\par
\switchcolumn*

\LA
\TeDeum{Te Deum}
\switchcolumn
\EN
\TeDeum{Te Deum}
\switchcolumn*

\end{paracol}

\par\needspace{10\baselineskip}\addvspace{1.2em}{\centering\small\textsc{Die 5 Junii} · \textsc{5 June}\par}
\Office{S. Bonifatii Episcopi et Martyris}{St. Boniface, Bishop and Martyr}{Duplex}
\addcontentsline{toc}{subsection}{Die 5 Junii · S. Bonifatii Episcopi et Martyris \textit{— St. Boniface, Bishop and Martyr}}
\begin{paracol}{2}
\LA
\RefLine{Lectiones I–III: de Scriptura occurrente.}
\switchcolumn
\EN
\RefLine{Lessons I–III: from the Scripture of the day.}
\switchcolumn*

\LA
\LessonHead{Lectio IV}
\switchcolumn
\EN
\LessonHead{Lesson IV}
\switchcolumn*

\LA
\noindent Bonifátius, ántea Winfrídus appellátus, apud Anglos natus est exeúnte sæculo séptimo, et ab ipsa infántia mundum aversátus, vitam monásticam in votis hábuit. Cum ejus pater ánimum sæculi illécebris permutáre frustra tentásset, monastérium ingréditu, et sub beáti Wolphárdi disciplína ómnium virtútum ac scientiárum génere imbúitur. Annum agens trigésimum sacerdótio insignítur, ac verbi divíni prædicátor assíduus, magno animárum lucro hoc in múnere versátur. Attamen, regnum Christi adaugére desíderans, contínuo flebat ingéntem multitúdinem barbarórum, qui ignorántiæ ténebris immérsi dæmoni famulabántur. Qui quidem animárum zelus cum in dies inexstinguíbili ardóre accrésceret, divíno númine per lácrimas et oratiónes exploráto, facultátem a monastérii præpósito obtínuit ad Germánicas oras proficiscéndi.\par
\switchcolumn
\EN
\noindent Winfried, afterwards called Boniface, was an Englishman, and born in England, towards the end of the seventh century. From his very childhood, he turned away from the world, and set his heart upon becoming a monk. His father tried in vain to turn him from his wishes by the beguilements of the world, and he entered a Monastery, where the Blessed Wolphard instructed him in all godliness and diverse kinds of learning. At the age of twenty-nine years he was ordained Priest, and became an unwearied preacher of the Word of God, wherein he had a gift which he used with great gain of souls. Nevertheless, his great desire was to spread the kingdom of Christ, and he continually bewailed the vast number of savages who were plunged in the darkness of ignorance and were the servants of the devil. This zealous love of souls increased in him in intensity day by day, till nothing would serve him, but, having implored the blessing of God by tears and prayers, and obtained authority from the head of his monastery, to set forth for the coast of Germany.\par
\switchcolumn*

\LA
\begin{Resp}\Rsign Lux perpétua lucébit Sanctis tuis, Dómine, \Star{} Et ætérnitas témporum, allelúja, allelúja.\end{Resp}
\begin{Resp}\Vsign Lætítia sempitérna erit super cápita eórum: gáudium et exsultatiónem obtinébunt.\end{Resp}
\begin{Resp}\Rsign Et ætérnitas témporum, allelúja, allelúja.\end{Resp}
\switchcolumn
\EN
\begin{Resp}\Rsign Eternal light will shine over your saints, O Lord, \Star{} And the eternity of the times, alleluia, alleluia.\end{Resp}
\begin{Resp}\Vsign Everlasting joy shall be upon their heads, they shall obtain joy and gladness\end{Resp}
\begin{Resp}\Rsign And the eternity of the times, alleluia, alleluia.\end{Resp}
\switchcolumn*

\LA
\LessonHead{Lectio V}
\switchcolumn
\EN
\LessonHead{Lesson V}
\switchcolumn*

\LA
\noindent Ex Anglia duóbus cum sóciis navem solvens, Dorestádium in Frísiæ óppidum venit. Cum autem bellum gravíssimum inter Frísonum regem Radbódum et Cárolum Martéllum exarsísset, sine fructu Evangélium prædicávit. Quaprópter in Angliam revérsus, ad suum redívit monastérium, cui invítus præfícitur. Post elápsum biénnium, ex consénsu epíscopi Vintoniénsis, munus abdicávit, et Romam proféctus est, ut apostólica auctoritáte ad Gentílium conversiónem delegarétur. Cum ad Urbem pervenísset, a Gregório secúndo benígne excípitur, pro Winfrído Bonifátius a Pontífice nominátur. In Germániam diréctus, Thuríngiæ Saxoniæque pópulis Christum annuntiávit. Cum intérea Radbódus, Frísiæ rex ac infestíssimus christiáni nóminis hostis, occubuísset, Bonifátius ad Frísones rédiit, ubi sancti Willibrórdi sócius per triénnium tanto cum fructu Evangélium prædicávit, ut, destrúctis idolórum simulácris, innúmeræ vero Deo ecclésiæ excitaréntur.\par
\switchcolumn
\EN
\noindent He set sail from England with two companions (in the year 716) and reached the town of Dorestadt in Friesland. A great war being then raging between Radbod, King of the Frieslanders, and Charles Martel, Winfrid preached the Gospel in vain. He went back to England, and betook himself again to his Monastery, whereof he was, against his own will, chosen to be the head. After two years he obtained the consent of the Bishop of Winchester to resign his office, and (in 719) went to Rome, to seek an Apostolic commission to preach to the heathen. When he arrived at the city he was courteously welcomed by Gregory II., who changed his name from Winfrid to Boniface. He departed thence to Germany, and preached Christ to the tribes in Thuringia and Saxony. Radbod, King of Friesland, who bitterly hated the Christian name, being dead, Boniface went a second time among the Frieslanders, and there, with his comrade St. Willibrord, preached the Gospel for three years with so much fruit, that the idols were hewn down, and countless churches arose to the true God.\par
\switchcolumn*

\LA
\begin{Resp}\Rsign In servis suis, allelúja, \Star{} Consolábitur Deus, allelúja.\end{Resp}
\begin{Resp}\Vsign Judicábit Dóminus pópulum suum, et in servis suis.\end{Resp}
\begin{Resp}\Rsign Consolábitur Deus, allelúja.\end{Resp}
\switchcolumn
\EN
\begin{Resp}\Rsign His servants, alleluia \Star{} God will comfort, alleluia, alleluia.\end{Resp}
\begin{Resp}\Vsign The Lord will judge His people and, His servants,\end{Resp}
\begin{Resp}\Rsign God will comfort, alleluia, alleluia.\end{Resp}
\switchcolumn*

\LA
\LessonHead{Lectio VI}
\switchcolumn
\EN
\LessonHead{Lesson VI}
\switchcolumn*

\LA
\noindent A sancto Willibrórdo ad episcopále munus expetítus, illud detrectávit ut prómpius infidélium salúti instáret. In Germániam proféctus, plura Hassórum míllia a dæmonis superstitióne avocávit. A Gregório Pontífice Romam evocátus, post insígnem fídei professiónem epíscopus consecrátur. Exínde ad Germános redux, Hássiam et Thuríngiam ab idololatríæ relíquiis pénitus expurgávit. Tanta propter mérita Bonifátius a Gregório tértio ad dignitátem archiepiscopálem evéhitur, et tértio Romam proféctus a summo Pontífice Sedis apostólicæ legátus constitúitur. Qua insignítus auctoritáte quátuor episcopátus instítuit, et várias synodos celebrávit, inter quas concílium Leptinénse memorábile est, apud Belgas in Cameracénsi diœcési celebrátum, quo quidem témpore ad fidem in Bélgio adaugéndam egrégie cóntulit. A Zacharía Papa creátus Mogúntius archiepíscopus, ipso Pontífice jubénte, Pipínum in regem Francórum unxit. Post mortem sancti Willibrórdi Ultrajecténsem ecclésiam gubernándum suscépit, primo per Eóbanum, deínde per seípsum, dum ab ecclésia Moguntína absolútus Ultrajécti resédit. Frisónibus ad idololatríam relápsis, Evangélium prædicáre rursus aggréditur, cumque offício pastoráli occuparétur, a bárbaris et ímpiis homínibus juxta Bornam flúvium cum Eóbano coëpíscopo multísque áliis cruénta cæde perémptus martyrii palma condecorátur. Corpus sancti Bonifátii Mogúntiam translátum, et, ut ipse vivens petíerat, in Fuldénsi monastério, quod exstrúxerat, recónditum fuit, ubi multis miráculis incláruit. Pius autem nonus Póntifex máximus, ejus Offícium et Missam ad univérsam Ecclésiam exténdit.\par
\switchcolumn
\EN
\noindent Willibrord urged upon him to take the office of a Bishop, but he deferred to seek it, that he might the more instantly toil for the salvation of the unbelievers. Advancing into Germany, he reclaimed thousands of the Hessians from devil-worship. Pope Gregory sent for him to Rome, (whither he came in 723,) and after hearing a noble profession of his faith, consecrated him a Bishop. He again returned to Germany, and thoroughly purged Hesse and Thuringia from all remains of idolatry. On account of such great works Gregory III. advanced Boniface to the dignity of an Archbishop, and on the occasion of a third journey to Rome, (in 738,) he was invested by the Sovereign Pontiff with the powers of Legate of the Apostolic See. As such, he founded (the) four Bishoprics (of Erfurt, Paderborn, Wurtzburg, and Eichstadt,) and held diverse Synods, among which is especially to be remembered that of Lessines, held in Belgium, in the diocese of Cambrai, wherein he made his strongest endeavours to spread the Faith among the Belgians. By Pope Zacharias, he was named Archbishop of Maintz, and by command of the same Pope, he anointed Pepin to be King of the Franks. After the death of St. Willibrord, he undertook the government of the Church of Utrecht, at first through Eoban but he afterwards was released from the care of the Church of Maintz and established his see at Utrecht. The Frieslanders having again fallen back into idolatry, he once more betook himself to preach the Gospel among them, and while he was busied in this duty, he grasped the crown of martyrdom, being murdered by some ungodly savages, along with his fellow-Bishop Eoban, and many others, in a bloody massacre near the River Born, (on the th day of June, in the year of our Lord 755, and of his own age the 75th.) In accordance with the wish expressed by himself during life the body of St. Boniface was carried to Maintz, and buried in the monastery of Fulda, of which he had been the founder, and where God has gloriously honoured it by the working of many signs and wonders. Pope Pius IX. ordered the Office and Mass in bis memory to be used throughout the whole Church.\par
\switchcolumn*

\LA
\begin{Resp}\Rsign Fíliæ Jerúsalem, veníte et vidéte Mártyres cum corónis, quibus coronávit eos Dóminus \Star{} In die solemnitátis et lætítiæ, allelúja.\end{Resp}
\begin{Resp}\Vsign Quóniam confortávit seras portárum tuárum, benedíxit fílios tuos in te.\end{Resp}
\begin{Resp}\Rsign In die solemnitátis et lætítiæ, allelúja.\end{Resp}
\begin{Resp}\Vsign Glória Patri, et Fílio, \Star{} et Spirítui Sancto.\end{Resp}
\begin{Resp}\Rsign In die solemnitátis et lætítiæ, allelúja.\end{Resp}
\switchcolumn
\EN
\begin{Resp}\Rsign Come forth, O ye daughters of Jerusalem, and behold the Martyrs with the crowns wherewith the Lord crowned them; \Star{} In the day of His feasting, and of His gladness. Alleluia.\end{Resp}
\begin{Resp}\Vsign For He hath strengthened the bars of thy gates He hath blessed thy children within thee.\end{Resp}
\begin{Resp}\Rsign In the day of His feasting, and of His gladness. Alleluia.\end{Resp}
\begin{Resp}\Vsign Glory be to the Father, and to the Son, \Star{} and to the Holy Ghost.\end{Resp}
\begin{Resp}\Rsign In the day of His feasting, and of His gladness. Alleluia.\end{Resp}
\switchcolumn*

\LA
\LessonHead{Lectio VII}
\switchcolumn
\EN
\LessonHead{Lesson VII}
\switchcolumn*

\LA
\LessonTitle{Léctio sancti Evangélii secúndum Matthǽum.}
\switchcolumn
\EN
\LessonTitle{From the Holy Gospel according to Matthew}
\switchcolumn*

\LA
\Cite{Matt 5:1-12}
\switchcolumn
\EN
\Cite{Matt 5:1-12}
\switchcolumn*

\LA
\noindent In illo témpore: Videns Jesus turbas, ascéndit in montem, et cum sedísset, accessérunt ad eum discípuli ejus. Et réliqua.\par
\switchcolumn
\EN
\noindent At that time Jesus, seeing the multitudes, went up into a mountain, and, when He was set, His disciples came unto Him. And so on.\par
\switchcolumn*

\LA
\LessonTitle{Homilia sancti Augustíni Epíscopi.}
\switchcolumn
\EN
\LessonTitle{Homily by St. Augustine, Bishop (of Hippo.)}
\switchcolumn*

\LA
\Source{Liber 1 de sermone Dómini in monte, cap. 2. et 3.}
\switchcolumn
\EN
\Source{Bk. i. on the Lord's Sermon.}
\switchcolumn*

\LA
\noindent Beáti mundo corde; quóniam ipsi Deum vidébunt. Quam ergo stulti sunt, qui Deum istis exterióribus óculis quærunt, cum corde videátur, sicut álibi scriptum est: Et in simplicitáte cordis quærite illum. Hoc est enim mundum cor, quod est simplex cor. Et quemádmodum lumen hoc vidéri non potest, nisi óculis mundis; ita nec Deus vidétur, nisi mundum sit illud quo vidéri potest. Beáti pacífici; quóniam ipsi fílii Dei vocabúntur. In pace perféctio est, ubi nihil repúgnat; et ídeo fílii Dei pacífici, quóniam nihil in his resístit Deo, et útique fílii similitúdinem patris habére debent.\par
\switchcolumn
\EN
\noindent Blessed are the pure in heart, for they shall see God. What fools then be they that seek God with their outward eyes, since it is in the heart that He is seen, as it is written elsewhere In simplicity of heart seek Him. (Wisd. i. i.) A simple heart is a pure heart. And even as we cannot see this earthly light, unless the eyes be open, so cannot God be seen, unless that be open which alone can perceive Him. Blessed are the peacemakers, for they shall be called the children of God. The perfection of peace is the absence of contrariety, and the peacemakers are called the children of God because they offer no contrariety against the will of God. As beseemeth children, they have their Father's likeness.\par
\switchcolumn*

\LA
\begin{Resp}\Rsign Ego sum vitis vera, et vos pálmites: \Star{} Qui manet in me, et ego in eo, hic fert fructum multum, allelúja, allelúja.\end{Resp}
\begin{Resp}\Vsign Sicut diléxit me Pater, et ego diléxi vos.\end{Resp}
\begin{Resp}\Rsign Qui manet in me, et ego in eo, hic fert fructum multum, allelúja, allelúja.\end{Resp}
\switchcolumn
\EN
\begin{Resp}\Rsign I am the vine: you the branches. \Star{} He that abideth in me, and I in him, the same beareth much fruit, alleluia, alleluia.\end{Resp}
\begin{Resp}\Vsign As the Father hath loved me, I also have loved you.\end{Resp}
\begin{Resp}\Rsign He that abideth in me, and I in him, the same beareth much fruit, alleluia, alleluia.\end{Resp}
\switchcolumn*

\LA
\LessonHead{Lectio VIII}
\switchcolumn
\EN
\LessonHead{Lesson VIII}
\switchcolumn*

\LA
\noindent Pacífici autem in semetípsis sunt, qui omnes ánimi sui motus componéntes, et subjiciéntes ratióni, id est menti et spirítui, carnalésque concupiscéntias habéntes edómitas, fiunt regnum Dei. In quo ita sunt ordináta ómnia, ut id quod est in hómine præcípuum et excéllens, hoc ímperet, céteris non reluctántibus, quæ sunt nobis bestiísque commúnia; atque idípsum quod excélluit in hómine, id est mens et rátio, subjiciátur potióri, quod est ipsa véritas, unigénitus Fílius Dei. Neque enim imperáre inferióribus potest, nisi superióri se ipse subjíciat. Et hæc est pax, quæ datur in terra homínibus bonæ voluntátis; hæc vita consummáti perfectíque sapiéntis.\par
\switchcolumn
\EN
\noindent They are peacemakers in themselves, who order all the movements of their own mind in obedience to reason, that is, to their intellect and soul, and so doing, and taming the lusts of the flesh, become a kingdom for God. In such kingdom all things are so ordered, that the chiefest and noblest part of man ruleth without contention over those lower things which we have in common with beasts. And just in the same way, must that nobler part of man, that is to say, intellect and reason, needs be put in subjection to what is above it, namely, Truth, the Only begotten Son of God. He only can rule well who hath learnt to obey. And this ordering is that peace which is given on earth to men of good will this is the life of whomsoever is thoroughly and perfectly wise.\par
\switchcolumn*

\LA
\begin{Resp}\Rsign Cándidi facti sunt Nazarǽi ejus, allelúja: splendórem Deo dedérunt, allelúja: \Star{} Et sicut lac coaguláti sunt, allelúja, allelúja.\end{Resp}
\begin{Resp}\Vsign Candidióres nive, nitidióres lacte, rubicundióres ébore antíquo, sapphíro pulchrióres.\end{Resp}
\begin{Resp}\Rsign Et sicut lac coaguláti sunt, allelúja, allelúja.\end{Resp}
\begin{Resp}\Vsign Glória Patri, et Fílio, \Star{} et Spirítui Sancto.\end{Resp}
\begin{Resp}\Rsign Et sicut lac coaguláti sunt, allelúja, allelúja.\end{Resp}
\switchcolumn
\EN
\begin{Resp}\Rsign Her Nazarites are become pure, Alleluia: they reflect the glory of God, Alleluia. \Star{} They are whiter than milk. Alleluia, Alleluia.\end{Resp}
\begin{Resp}\Vsign They are purer than snow, they are whiter than milk, they are more ruddy in body than coral, their polishing is of sapphire.\end{Resp}
\begin{Resp}\Rsign They are whiter than milk. Alleluia, Alleluia.\end{Resp}
\begin{Resp}\Vsign Glory be to the Father, and to the Son, \Star{} and to the Holy Ghost.\end{Resp}
\begin{Resp}\Rsign They are whiter than milk. Alleluia, Alleluia.\end{Resp}
\switchcolumn*

\LA
\LessonHead{Lectio IX}
\switchcolumn
\EN
\LessonHead{Lesson IX}
\switchcolumn*

\LA
\noindent De hujúsmodi regno pacatíssimo et ordinatíssimo missus est foras princeps hujus sæculi, qui pervérsis inordinatísque dominátur. Hac pace intrínsecus constitúta atque firmáta, quascúmque persecutiónes ille, qui foras missus est, forínsecus concitáverit, auget glóriam, quæ secúndum Deum est; non áliquid in illo ædifício labefáctans, sed deficiéntibus máchinis suis innotéscere fáciens, quanta fírmitas intus exstrúcta sit. Ideo séquitur: Beáti, qui persecutiónem patiúntur propter justítiam; quóniam ipsórum est regnum cælórum.\par
\switchcolumn
\EN
\noindent From this most peaceful and most orderly kingdom is cast forth the prince of this world, whose rule is over the contentious and disorderly. When once this peace hath been proclaimed and established within, whatsoever wars he that is without can raise, can but heap more glory upon that glory which is according to God, for nothing of the castle will yield before him, but the yielding of his own engines will witness how strong be its ramparts. And therefore cometh next Blessed are they which are persecuted for righteousness' sake, for theirs is the kingdom of heaven.\par
\switchcolumn*

\LA
\TeDeum{Te Deum}
\switchcolumn
\EN
\TeDeum{Te Deum}
\switchcolumn*

\LA
\LessonHead{Lectio IX (pro commemoratione)}
\switchcolumn
\EN
\LessonHead{Lesson IX (when commemorated)}
\switchcolumn*

\LA
\LessonTitle{Commemoratio: S. Bonifatii Episcopi et Martyris}
\switchcolumn
\EN
\LessonTitle{Commemoration of: St. Boniface, Bishop and Martyr}
\switchcolumn*

\LA
\noindent Bonifátius, ántea Winfrídus appellátus, apud Anglos natus est exeúnte sǽculo séptimo. Monastérium ingréssus et sacerdótio auctus, magno animárum lucro in prædicatóris múnere est versátus. Zelo augéndæ fídei accénsus, apud Frísones Evangélium prædicávit. In Angliam revérsus, cum per biénnium monastério sanctíssime præfuísset, superióris múnere abdicáto, Romam se cóntulit, ubi a Gregório secúndo Bonifátii nomen accépit, et in Germániam missus, Thuríngiæ Saxoniǽque pópulis Christum annuntiávit. Ad Frísones revérsus, cum sancto Willibrórdo, magno fructu Evangélium prædicávit. Mox Romam accersítus, episcopáli dignitáte insignítur, et in Germániam íterum proféctus, Hássiam et Thuríngiam ab idololatríæ relíquiis pénitus expurgávit. Sedis apostólicæ legátus creátus et Moguntínus archiepíscopus, plures eréxit et, per se vel per discípulos, administrávit ecclésias. Frisónibus demum ad idololatríam relápsis Evangélium prædicáre rursus aggréssus, cum Eóbano coëpíscopo multísque áliis, juxta Bornam flúvium cruénta cæde perémptus, martýrii palmam accépit. Ejus corpus in Fuldénsi monastério cónditum est.\par
\switchcolumn
\EN
\noindent Boniface, originally called Winfrid, was born in England towards the end of the seventh century. After entering a monastery and becoming a priest, he showed great skill in winning souls through preaching. Burning with zeal to spread the faith, he preached the Gospel among the Frisians. Then he returned to England, where he ruled his monastery for two years in a most holy manner. Having resigned the office of Superior, he went to Rome, where he received from Gregory II the name of Boniface and the commission to proclaim Christ to the peoples of Thuringia and Saxony. With holy Willibrord, he returned to the Frisians and preached the Gospel with great fruit. Soon he was summoned to Rome and invested with the episcopal dignity; after which, he set out once more for Germany. There he rid Hesse and Thuringia of almost the last vestiges of idolatry. He was made apostolic delegate and Archbishop of Mainz, and he built many churches, and administered them either personally or through his disciples. At length, he went back once again to the Frisians, who had lapsed into idolatry, to preach the Gospel to them. There, with Eobanus his fellow bishop and many others, he was killed in a bloody massacre near the River Born and received the crown of martyrdom. His body lieth in the monastery of Fulda.\par
\switchcolumn*

\LA
\TeDeum{Te Deum}
\switchcolumn
\EN
\TeDeum{Te Deum}
\switchcolumn*

\end{paracol}

\par\needspace{10\baselineskip}\addvspace{1.2em}{\centering\small\textsc{Die 6 Junii} · \textsc{6 June}\par}
\Office{S. Norberti Episcopi et Confessoris}{St. Norbert, Bishop and Confessor}{Duplex}
\addcontentsline{toc}{subsection}{Die 6 Junii · S. Norberti Episcopi et Confessoris \textit{— St. Norbert, Bishop and Confessor}}
\begin{paracol}{2}
\LA
\RefLine{Lectiones I–III: de Scriptura occurrente.}
\switchcolumn
\EN
\RefLine{Lessons I–III: from the Scripture of the day.}
\switchcolumn*

\LA
\RefLine{Lectiones VII–IX: ex Commúni unius Confessoris Pontificis (C4).}
\switchcolumn
\EN
\RefLine{Lessons VII–IX: from the Common of a Confessor Bishop (C4).}
\switchcolumn*

\LA
\LessonHead{Lectio IV}
\switchcolumn
\EN
\LessonHead{Lesson IV}
\switchcolumn*

\LA
\noindent Norbértus, nobilíssimis paréntibus natus, adoléscens liberálibus disciplínis erudítus, in ipsa póstea imperatóris aula, spretis mundi illécebris, ecclesiásticæ milítiæ adscríbi vóluit. Sacris initiátus, rejéctis móllibus ac spléndidis véstibus, pellícea melóte indútus, prædicatióni verbi Dei se totum dedit. Abdicátis ecclesiásticis provéntibus satis amplis, et património in páuperes erogáto, semel in die sub vésperam solo cibo quadragesimáli utens, nudísque pédibus et lácera veste sub brumáli rigóre incédens, miræ austeritátis vitam est aggréssus. Potens ígitur ópere et sermóne, innúmeros hæréticos ad fidem, peccatóres ad pœniténtiam, dissidéntes ad pacem et concórdiam revocávit.\par
\switchcolumn
\EN
\noindent Norbert, born in the year 1080 of parents of the highest rank, thoroughly educated in his youth in worldly knowledge, and a member of the Imperial court, turned his back upon the glory of the world, and chose rather to enlist himself as a soldier of the Church. Being ordained Priest, he laid aside all soft and showy raiment, clad himself in a coat of skins, and made the preaching of the Word of God the one object of his life. He had the right to rich revenues of the Church but these he renounced and to an ample fortune from his father; but this he gave to the poor. He ate only once a day, and that in the evening, and then his meal was of the fare of Lent. His life was one of singular hardness, and he was used even in the depth of winter to go out with bare feet and ragged garments. Hence came that mighty power of his words and deeds, whereby he was enabled to turn countless heretics to the true faith, sinners to repentance, and enemies to peace and brotherly love.\par
\switchcolumn*

\LA
\begin{Resp}\Rsign Invéni David servum meum, óleo sancto meo unxi eum: \Star{} Manus enim mea auxiliábitur ei.\end{Resp}
\begin{Resp}\Vsign Nihil profíciet inimícus in eo, et fílius iniquitátis non nocébit ei.\end{Resp}
\begin{Resp}\Rsign Manus enim mea auxiliábitur ei.\end{Resp}
\switchcolumn
\EN
\begin{Resp}\Rsign I have found David My servant, with My holy oil have I anointed him \Star{} For My hand shall help him.\end{Resp}
\begin{Resp}\Vsign The enemy shall prevail nothing against him, nor the son of wickedness afflict him.\end{Resp}
\begin{Resp}\Rsign For My hand shall help him.\end{Resp}
\switchcolumn*

\LA
\LessonHead{Lectio V}
\switchcolumn
\EN
\LessonHead{Lesson V}
\switchcolumn*

\LA
\noindent Cum Laudúni esset, ab epíscopo rogátus ne a sua diœcési discéderet, desértum in ea locum, qui Præmonstrátus dicebátur, sibi delégit; ibíque, trédecim sóciis aggregátis, Præmonstraténsem órdinem instítuit, divínitus accépta per visum régula a sancto Augustíno. Cum vero ejus fama sanctitátis in dies magis augerétur, ac plúrimi ad eum quotídie discípuli convenírent, idem ordo ab Honório secúndo aliísque summis Pontifícibus confirmátus, ac plúribus ab eo monastériis ædificátis, mirífice propagátus est.\par
\switchcolumn
\EN
\noindent Being one while at Laon, the Bishop besought him not to leave his diocese, and he therefore made choice of a wilderness at the place called Prémontré, whither he withdrew himself with thirteen disciples, and thus founded the Order of the Praemonstratensian Canons, whereof he, by the will of God, received the Rule, in a vision, from St. Augustine. When, however, the fame of his holy life became every day more and more noised abroad, and great numbers sought to become his disciples, and the Order had been approved by Honorius II, and other Popes, many more monasteries were built by him, and the Institute wonderfully extended.\par
\switchcolumn*

\LA
\begin{Resp}\Rsign Pósui adjutórium super poténtem, et exaltávi eléctum de plebe mea: \Star{} Manus enim mea auxiliábitur ei.\end{Resp}
\begin{Resp}\Vsign Invéni David servum meum, óleo sancto meo unxi eum.\end{Resp}
\begin{Resp}\Rsign Manus enim mea auxiliábitur ei.\end{Resp}
\switchcolumn
\EN
\begin{Resp}\Rsign I have laid help upon one that is mighty, and have exalted one chosen out of My people \Star{} For My hand shall help him.\end{Resp}
\begin{Resp}\Vsign I have found David My servant, with My holy oil have I anointed him.\end{Resp}
\begin{Resp}\Rsign For My hand shall help him.\end{Resp}
\switchcolumn*

\LA
\LessonHead{Lectio VI}
\switchcolumn
\EN
\LessonHead{Lesson VI}
\switchcolumn*

\LA
\noindent Antvérpiam accersítus, in ea urbe Tanchelíni nefáriam hæresim profligávit. Prophético spíritu et miráculis cláruit. Archiepíscopus tandem (licet relúctans) Magdeburgénsis creátus, ecclesiásticam disciplínam, præsértim cælibátum, constánter propugnávit. Rhemis in concílio Innocéntium secúndum egrégie adjúvit, et Romam cum áliis epíscopis proféctus, schisma Petri Leónis compréssit. Postrémo vir Dei, méritis et Spíritu Sancto plenus, Magdebúrgi obdormívit in Dómino, anno salútis millésimo centésimo trigésimo quarto, die sexta Júnii.\par
\switchcolumn
\EN
\noindent Being called to Antwerpen, he there gave the death-blow to the shameful heresy of Tanchelm. He was remarkable for the spirit of prophecy and for the gift of miracles. He was created (albeit he would rather not have had it so) Archbishop of Magdeburg, and as such he was a strong upholder of the discipline of the Church, especially contending against the marriage of the clergy. At a Council held at Reims he was a great help to Innocent II, and went with some Other Bishops to Rome, where they stamped out the schism of Peter Leoni. It was at last at Magdeburg that this man of God, full of good works and of the Holy Ghost, fell asleep in the Lord, on the 6th day of June, in the year of salvation 1134.\par
\switchcolumn*

\LA
\begin{Resp}\Rsign Iste est, qui ante Deum magnas virtútes operátus est, et omnis terra doctrína ejus repléta est: \Star{} Ipse intercédat pro peccátis ómnium populórum.\end{Resp}
\begin{Resp}\Vsign Iste est, qui contémpsit vitam mundi, et pervénit ad cæléstia regna.\end{Resp}
\begin{Resp}\Rsign Ipse intercédat pro peccátis ómnium populórum.\end{Resp}
\begin{Resp}\Vsign Glória Patri, et Fílio, \Star{} et Spirítui Sancto.\end{Resp}
\begin{Resp}\Rsign Ipse intercédat pro peccátis ómnium populórum.\end{Resp}
\switchcolumn
\EN
\begin{Resp}\Rsign This is he which wrought great wonders before God, and the whole earth is full of his teaching \Star{} May he pray for all people, that their sins may be forgiven unto them\end{Resp}
\begin{Resp}\Vsign This is he which loved not his life in this world, and hath attained unto the kingdom of heaven.\end{Resp}
\begin{Resp}\Rsign May he pray for all people, that their sins may be forgiven unto them\end{Resp}
\begin{Resp}\Vsign Glory be to the Father, and to the Son, \Star{} and to the Holy Ghost.\end{Resp}
\begin{Resp}\Rsign May he pray for all people, that their sins may be forgiven unto them\end{Resp}
\switchcolumn*

\LA
\LessonHead{Lectio IX (pro commemoratione)}
\switchcolumn
\EN
\LessonHead{Lesson IX (when commemorated)}
\switchcolumn*

\LA
\LessonTitle{Commemoratio: S. Norberti Episcopi et Confessoris}
\switchcolumn
\EN
\LessonTitle{Commemoration of: St. Norbert, Bishop and Confessor}
\switchcolumn*

\LA
\noindent Norbértus, nobilíssimis paréntibus natus, adoléscens liberálibus disciplínis erudítus, in ipsa póstea imperatóris aula, spretis mundi illécebris, ecclesiásticæ milítiæ adscríbi vóluit. Sacris initiátus, prædicatióni verbi Dei se totum dedit. Innúmeros hæréticos ad fidem, peccatóres ad pœniténtiam, dissidéntes ad pacem et concórdiam revocávit. Desértum locum, qui Præmonstrátus dicebátur, in Laudunénsi diœcési sibi delégit; ibíque, trédecim sóciis aggregátis, Præmonstraténsem órdinem instítuit, qui mirífice propagátus est. Archiepíscopus Magdeburgénsis, licet relúctans, creátus, ecclesiásticam disciplínam, cælibátum præsértim, constánter propugnávit. Rhemis in concílio Innocéntium secúndum egrégie adjúvit, et Romam cum áliis epíscopis proféctus, schisma Petri Leónis compréssit. Magdebúrgi obdormívit in Dómino, anno salútis millésimo centésimo trigésimo quarto, die sexta júnii.\par
\switchcolumn
\EN
\noindent Norbert was born of most noble parents and, as a young man, studied the liberal arts. Then, while serving in the court of the Emperor himself, he spurned the seductions of the world and decided to enroll among the soldiers of the Church. After receiving holy orders he devoted himself entirely to preaching the word of God. He brought back innumerable heretics to the faith, sinners to penance, quarrellers to peace and concord. He retired to a desert place called Prémontré in the diocese of Laon; and there, with thirteen companions, he founded the Praemonstratensian Order, which spread in a marvelous way. Against his will, he was made Archbishop of Magdeburg, and constantly defended ecclesiastical discipline and especially celibacy. At the Council of Reims, he was a strong champion of Innocent II; and going to Rome with other bishops, he put an end to the schism of Pierleone. He fell asleep in the Lord at Magdeburg on the 6th day of June in the year of salvation 1134.\par
\switchcolumn*

\LA
\TeDeum{Te Deum}
\switchcolumn
\EN
\TeDeum{Te Deum}
\switchcolumn*

\end{paracol}

\par\needspace{10\baselineskip}\addvspace{1.2em}{\centering\small\textsc{Die 9 Junii} · \textsc{9 June}\par}
\Office{Ss. Primi et Feliciani Martyrum}{Sts. Primus and Felician, Martyrs}{Simplex}
\addcontentsline{toc}{subsection}{Die 9 Junii · Ss. Primi et Feliciani Martyrum \textit{— Sts. Primus and Felician, Martyrs}}
\begin{paracol}{2}
\LA
\RefLine{Lectiones I–II: de Scriptura occurrente.}
\switchcolumn
\EN
\RefLine{Lessons I–II: from the Scripture of the day.}
\switchcolumn*

\LA
\LessonHead{Lectio III (pro commemoratione)}
\switchcolumn
\EN
\LessonHead{Lesson III (when commemorated)}
\switchcolumn*

\LA
\noindent Primus et Felicianus fratres, in persecutióne Diocletiáni et Maximiáni accusáti christianæ religiónis, in víncula conjiciúntur; quibus soluti, inde eripiúntur ab Angelo. Mox ad prætorem adducti, cum christianam fidem acerrime tueréntur, alter ab altero distracti sunt; ac primum varie tentáta est constántia Feliciáni. Sed, cum suasores impietátis se posse quidquam verbis proficere desperárent, affixis stipiti mánibus ejus et pédibus, ipsum sine cibo et potu inde triduum pendentem reliquérunt. Postridie ejus diéi, prætor vocátum ad se Primum sic affátur: Vides quanto sit prudentior, quam tu, frater tuus, qui obsecutus imperatóribus, apud ipsos est honoratus? Quem si tu quoque imitari volueris, particeps eris ejus honoris et gratiæ. Cui Primus: Qui factum sit fratri meo, cognóvi ex Angelo. Utinam, quemádmodum sum cum eo voluntáte conjunctíssimus, sic ab eodem ne martyrio disjungar. Quo dicto excanduit prætor, et ad ceteros cruciatus, quibus Primum affecit, præsénte jam Feliciáno, liquátum igne plumbum in os ejus jussit infundi. Mox utrumque perduci imperat in theatrum, in eosque immitti duos leónis; qui, prostráti ad eórum genua, cápite et cauda ipsis blandiebántur. Ad id spectaculum cum ámplius duódecim míllia hóminum convenissent, quingénti cum suis familiis christianam religiónem suscepérunt. Quibus rebus permotus prætor, eos secúri pércuti jussit.\par
\switchcolumn
\EN
\noindent Primus and Felician were two brothers who were accused of Christianity during the persecution by Diocletian and Maximian, and thrown into irons, which an angel broke, and so freed their limbs. In the presence of the Praetor they most earnestly clave to the profession of their faith, and were immediately parted one from the other. Felician's was the steadfastness which was first tried in diverse ways. They, however, that strove to argue him into sin, when they found that words availed nothing, fastened his hands and feet to a post, and left him to hang there three days without food or drink. On the fourth day the Praetor called Primus before him, and said to him Seest thou how much thy brother is wiser than thou He hath obeyed the Emperors, and they have made him honourable. Thou hast only to follow his example to be made partaker of his honours and favours. Primus answered him: What hath befallen my brother I know, for an angel hath told me. God grant that, seeing I have the same will that he hath, I may not be divided from him in uplifting of testimony. These words raised the wrath of the Praetor, and to the torments which he had already inflicted on Primus, he added this also, that he had boiling lead put into his mouth, compelling his brother Felician to be present and see it done. After that, he had them led into the theatre and two lions let loose upon them, in the presence of about twelve thousand people who were gathered together to see the show. The lions only fawned upon the knees of the Saints, making friends with them with motions of their heads and tails. This exhibition turned five hundred persons and their households to Christ. The Praetor, then, moved beyond all endurance by what had passed, caused Primus and Felician to be beheaded.\par
\switchcolumn*

\LA
\TeDeum{Te Deum}
\switchcolumn
\EN
\TeDeum{Te Deum}
\switchcolumn*

\end{paracol}

\par\needspace{10\baselineskip}\addvspace{1.2em}{\centering\small\textsc{Die 10 Junii} · \textsc{10 June}\par}
\Office{S. Margaritæ Reginæ Viduæ}{St. Margaret, Queen and Widow}{Semiduplex}
\addcontentsline{toc}{subsection}{Die 10 Junii · S. Margaritæ Reginæ Viduæ \textit{— St. Margaret, Queen and Widow}}
\begin{paracol}{2}
\LA
\RefLine{Lectiones I–III: de Scriptura occurrente.}
\switchcolumn
\EN
\RefLine{Lessons I–III: from the Scripture of the day.}
\switchcolumn*

\LA
\RefLine{Lectiones VII–IX: ex Commúni non Virginum non Martyrum (C7a).}
\switchcolumn
\EN
\RefLine{Lessons VII–IX: from the Common of Widows (C7a).}
\switchcolumn*

\LA
\LessonHead{Lectio IV}
\switchcolumn
\EN
\LessonHead{Lesson IV}
\switchcolumn*

\LA
\noindent Margarita, Scotórum regína, paterno Angliæ regum, materno Cæsárum sánguine claríssima, illustrior adhuc fuit christiana virtúte. Hæc in Hungária nata, ubi pater tunc témporis exsulabat, post exactam summa cum pietáte puerílem ætátem, una cum genitore, qui a sancto Eduardo patruo, Anglórum rege, ad paterni regni fastigium vocabátur, in Angliam venit. Mox, alternante paréntum fortuna, ex Angliæ littore solvens, vi tempestátis expulsa, seu verius divinæ providéntiæ consílio deducta est in oram marítimam Scotiæ. Ibi cum ex matris imperio Malchólmo tertio Scotórum regi, egregiis ejus dótibus capto, nupsísset, sanctimoniæ ac pietátis opéribus, trigínta quibus regnávit annis, toti regno mirifice prófuit.\par
\switchcolumn
\EN
\noindent Margaret, Queen of Scots, was most noble by birth, uniting in herself, from her father the blood of the Kings of England and from her mother the blood of the Caesars, but her greatest nobleness was in her brave Christian life. She was born in Hungary, where her father was then an exile, (in the year 1046,) and had passed a religious childhood, when her uncle Edward, the holy King of England, recalled him to his own royal home, and she came to England with him (in 1054.) A few years after, upon the ruin of her family, she was escaping from England by sea, when the violence of the weather, or, to speak more truly, the Providence of God, caused that the ship should take refuge upon the coast of Scotland. There her extraordinary graces of mind and body so attracted King Malcolm III., that by the advice of his mother, he took her to wife (in 1070,) and of Scotland she deserved exceedingly well for the thirty years of her reign, by the holiness of her life and the abundance of her works of mercy.\par
\switchcolumn*

\LA
\begin{Resp}\Rsign Propter veritátem, et mansuetúdinem, et justítiam: \Star{} Et dedúcet te mirabíliter déxtera tua.\end{Resp}
\begin{Resp}\Vsign Spécie tua et pulchritúdine tua inténde, próspere procéde, et regna.\end{Resp}
\begin{Resp}\Rsign Et dedúcet te mirabíliter déxtera tua.\end{Resp}
\switchcolumn
\EN
\begin{Resp}\Rsign Because of truth, and meekness, and righteousness; \Star{} And thy right hand shall lead thee wonderfully.\end{Resp}
\begin{Resp}\Vsign In thy comeliness, and thy beauty, go forward, fare prosperously, and reign.\end{Resp}
\begin{Resp}\Rsign And thy right hand shall lead thee wonderfully.\end{Resp}
\switchcolumn*

\LA
\LessonHead{Lectio V}
\switchcolumn
\EN
\LessonHead{Lesson V}
\switchcolumn*

\LA
\noindent Inter regales delicias corpus afflictatiónibus ac vigiliis mácerans, magnam noctis partem piis precatiónibus extrahebat. Præter alia jejunia, quæ idéntidem usurpabat, íntegros quadragínta dies ante Natalícia festa tanta cum severitáte jejunare consuevit, ut ne in gravíssimis quidem dolóribus intermiserit. Divino cultui addictíssima, templa pluria et cœnobia partim ex íntegro excitávit, partim resarcívit, et sacra supelléctili ac largo censu ditávit. Regem cónjugem ad meliórem frugem et ad simília suis exercitatiónibus ópera saluberrimo exemplo traduxit, liberosque omnes tam sancte et feliciter educávit, ut eórum plerique, quemádmodum et Agatha mater, et Christina soror, sanctíssimum vitæ genus ampléxi sint. Universi demum regni felicitáti cónsulens, a vitiis ómnibus, quæ furtim irrepserant, pópulos expurgávit, eisque mores christiana pietáte dignos restituit.\par
\switchcolumn
\EN
\noindent In the midst of kingly dainties, she afflicted her body with hardships and watching, using to spend great part of the night in earnest prayer. Besides other fasts which she imposed upon herself, it was her custom to observe one of forty days before Christmas, concerning which fast she was so rigid, that she would not relax it even under sharp suffering. She took great delight in the public worship of God, and founded or renewed a great number of Churches and convents, which she enriched at great cost with sacred furniture. Her healthy example drew the King her husband to habits of sobriety, and to imitate her in her good works. To all her children she had the happiness of giving a godly education, and several of them, like her mother Agatha and her sister Christina, led notable holy lives. The happiness of the whole kingdom was the object for which she constantly strove, and she successfully rooted out all the vices which had stealthily crept in, and established among the people a standard of living worthy of Christians.\par
\switchcolumn*

\LA
\begin{Resp}\Rsign Dilexísti justítiam, et odísti iniquitátem: \Star{} Proptérea unxit te Deus, Deus tuus, óleo lætítiæ.\end{Resp}
\begin{Resp}\Vsign Propter veritátem, et mansuetúdinem, et justítiam.\end{Resp}
\begin{Resp}\Rsign Proptérea unxit te Deus, Deus tuus, óleo lætítiæ.\end{Resp}
\switchcolumn
\EN
\begin{Resp}\Rsign Thou hast loved righteousness, and hated iniquity; \Star{} Therefore God, thy God, hath anointed thee with the oil of gladness.\end{Resp}
\begin{Resp}\Vsign Because of truth, and meekness, and righteousness.\end{Resp}
\begin{Resp}\Rsign Therefore God, thy God, hath anointed thee with the oil of gladness.\end{Resp}
\switchcolumn*

\LA
\LessonHead{Lectio VI}
\switchcolumn
\EN
\LessonHead{Lesson VI}
\switchcolumn*

\LA
\noindent Nihil tamen æque in illa mirábile fuit ac flagrantíssima caritas erga próximos, præsertim egénos, quorum numerosis grégibus non modo stipem áffatim suppeditare, verum étiam trecentis quotídie materna benignitate dapes præbére, flexis genibus in morem ancíllæ ministrare, regiis mánibus pedes ablúere, et pressis étiam ósculis úlcera fovére, solemne hábuit. His porro aliisque piis sumptibus non regias tantum vestes et preiosa monília distraxit, sed ipsum non semel exhausit ærárium. Tolerátis demum ad patiéntiæ miraculum acerbíssimis dolóribus, ánimam semestri corporis ægrotatióne purgatam Auctori suo, sextodecimo Kalendas Decembris, réddidit. Quo témporis momento fácies ejus, diuturni morbi macie ac pallore fœdata, insólita quadam venustáte reflóruit. Miris étiam post mortem prodigiis clara, et Clementis décimi auctoritate in Scotiæ patronam accepta, ubíque terrárum religiosíssime cólitur.\par
\switchcolumn
\EN
\noindent The most remarkable feature of her life was the tenderness of her charity toward her neighbour, especially the needy. Of these she would not only order whole flocks to be relieved, but was accustomed to give dinner to three hundred of them every day, treating them with the tenderness of a mother, and waiting upon them on her knees like a maidservant. She held it one of the privileges of her rank to wash their feet with her own Royal hands, and to dress their sores, which latter she would even kiss. To meet the expenses of her charities she sold not only her queenly raiment and her precious jewels, but more than once exhausted her funds entirely. Purified by grievous suffering, which she bore with marvelous patience during an illness of six months, she resigned her soul into the hands of Him Who had created it, upon the 11th day of June, (1093.) At the moment of death, the bystanders saw her poor worn face, pale and disfigured by continual suffering, flush again with a beauty to which it had long been unused. After her death she became illustrious on account of great signs and wonders. With the approval of Clement X., she was chosen Patroness of Scotland, and her memory is held in profound reverence throughout the whole earth.\par
\switchcolumn*

\LA
\begin{Resp}\Rsign Fallax grátia, et vana est pulchritúdo: \Star{} Múlier timens Dóminum, ipsa laudábitur.\end{Resp}
\begin{Resp}\Vsign Date ei de fructu mánuum suárum, et laudent eam in portis ópera ejus.\end{Resp}
\begin{Resp}\Rsign Múlier timens Dóminum, ipsa laudábitur.\end{Resp}
\begin{Resp}\Vsign Glória Patri, et Fílio, \Star{} et Spirítui Sancto.\end{Resp}
\begin{Resp}\Rsign Múlier timens Dóminum, ipsa laudábitur.\end{Resp}
\switchcolumn
\EN
\begin{Resp}\Rsign Favour is deceitful, and beauty is vain \Star{} A woman that feareth God she shall be praised.\end{Resp}
\begin{Resp}\Vsign Give her of the fruit of her hands, and let her own works praise her in the gates.\end{Resp}
\begin{Resp}\Rsign A woman that feareth God she shall be praised.\end{Resp}
\begin{Resp}\Vsign Glory be to the Father, and to the Son, \Star{} and to the Holy Ghost.\end{Resp}
\begin{Resp}\Rsign A woman that feareth God she shall be praised.\end{Resp}
\switchcolumn*

\LA
\LessonHead{Lectio IX (pro commemoratione)}
\switchcolumn
\EN
\LessonHead{Lesson IX (when commemorated)}
\switchcolumn*

\LA
\LessonTitle{Commemoratio: S. Margaritæ Reginæ Viduæ}
\switchcolumn
\EN
\LessonTitle{Commemoration of: St. Margaret, Queen and Widow}
\switchcolumn*

\LA
\noindent Margaríta, ex régia Anglórum stirpe in Hungária nata, post exáctam summa cum pietáte puerítiam, una cum genitóre, qui a sancto Eduárdo pátruo, Anglórum rege, ad patérni regni fastígium vocabátur, in Angliam, dein in Scótiam venit. Ibi, cum ex matris império Scotórum regi Malchólmo tértio nupsísset, sanctimóniæ et pietátis opéribus annis trigínta toti regno prófuit. Máxima erat in ea vitæ austéritas et flagrantíssimum erga próximos caritátis stúdium, præsértim in egénos; pro quibus aléndis non semel exháusit ærárium. Demum acerbis dolóribus et diutúrno morbo patientíssime tolerátis, ánimam Deo réddidit sextodécimo caléndas decémbris. Quo témporis moménto fácies ejus, mácie ac pallóre fœdáta, insólita quadam venustáte reflóruit. Cleméntis décimi auctoritáte in Scótiæ patrónam accépta, ubíque terrárum religiosíssime cólitur.\par
\switchcolumn
\EN
\noindent Margaret, of the royal house of England, was born in Hungary and spent her childhood there as an unusually devout and pious girl. When her father was called to high office in his own country by his uncle, St. Edward, King of England, she went to England and then to Scotland. There, upon instructions from her mother, she married King Malcolm III. The country was blessed by her holy life and by her deeds of charity for the next thirty years. The austerity of her life was exceedingly great, and her charity towards her neighbour most ardent and zealous, especially for those in need, for whom she not infrequently exhausted the treasury. At length, having most patiently endured bitter sorrows and long illness, she rendered her soul to God on the 16th day of November. At the moment of her death her features, emaciated and pale, bloomed again with unusual beauty. On the authority of Clement X she was chosen Patroness of Scotland, and is honoured with great devotion throughout the world.\par
\switchcolumn*

\LA
\TeDeum{Te Deum}
\switchcolumn
\EN
\TeDeum{Te Deum}
\switchcolumn*

\end{paracol}

\par\needspace{10\baselineskip}\addvspace{1.2em}{\centering\small\textsc{Die 11 Junii} · \textsc{11 June}\par}
\Office{S. Barnabæ Apostoli}{St. Barnabas the Apostle}{Duplex majus}
\addcontentsline{toc}{subsection}{Die 11 Junii · S. Barnabæ Apostoli \textit{— St. Barnabas the Apostle}}
\begin{paracol}{2}
\LA
\LessonHead{Lectio I}
\switchcolumn
\EN
\LessonHead{Lesson I}
\switchcolumn*

\LA
\LessonTitle{De Actibus Apostolórum}
\switchcolumn
\EN
\LessonTitle{From the Acts of the Apostles}
\switchcolumn*

\LA
\Cite{Acts 13:43-47}
\switchcolumn
\EN
\Cite{Acts 13:43-47}
\switchcolumn*

\LA
\noindent Cum dimíssa esset synagóga, secúti sunt multi Judæórum et coléntium advenárum, Paulum et Bárnabam; qui loquéntes suadébant eis ut permanérent in grátia Dei. Sequénti vero sábbato pene univérsa cívitas convénit audíre verbum Dei. Vidéntes autem turbas Judǽi, repléti sunt zelo, et contradicébant his, quæ a Paulo dicebántur, blasphemántes. Tunc constánter Paulus et Bárnabas dixérunt: Vobis oportébat primum loqui verbum Dei; sed, quóniam repéllitis illud, et indígnos vos judicátis ætérnæ vitæ, ecce convértimur ad gentes; sic enim præcépit nobis Dóminus: Pósui te in lucem Géntium, ut sis in salútem usque ad extrémum terræ.\par
\switchcolumn
\EN
\noindent And when the synagogue was broken up, many of the Jews, and of the strangers who served God, followed Paul and Barnabas: who speaking to them, persuaded them to continue in the grace of God. But the next sabbath day, the whole city almost came together, to hear the word of God. And the Jews seeing the multitudes, were filled with envy, and contradicted those things which were said by Paul, blaspheming. Then Paul and Barnabas said boldly: To you it behoved us first to speak the word of God: but because you reject it, and judge yourselves unworthy of eternal life, behold we turn to the Gentiles. For so the Lord hath commanded us: I have set thee to be the light of the Gentiles; that thou mayest be for salvation unto the utmost part of the earth.\par
\switchcolumn*

\LA
\begin{Resp}\Rsign Ecce ego mitto vos sicut oves in médio lupórum, dicit Dóminus: \Star{} Estóte ergo prudéntes sicut serpéntes, et símplices sicut colúmbæ.\end{Resp}
\begin{Resp}\Vsign Dum lucem habétis, crédite in lucem, ut fílii lucis sitis.\end{Resp}
\begin{Resp}\Rsign Estóte ergo prudéntes sicut serpéntes, et símplices sicut colúmbæ.\end{Resp}
\switchcolumn
\EN
\begin{Resp}\Rsign Behold, I send you as sheep in the midst of wolves, said the Lord: \Star{} Be ye, therefore, wise as serpents and simple as doves.\end{Resp}
\begin{Resp}\Vsign Whilst you have the light, believe in the light, that you may be the children of light.\end{Resp}
\begin{Resp}\Rsign Be ye therefore wise as serpents and simple as doves.\end{Resp}
\switchcolumn*

\LA
\LessonHead{Lectio II}
\switchcolumn
\EN
\LessonHead{Lesson II}
\switchcolumn*

\LA
\Cite{Act 13:48-52}
\switchcolumn
\EN
\Cite{Acts 13:48-52}
\switchcolumn*

\LA
\noindent Audiéntes autem gentes, gavísæ sunt, et glorificábant verbum Dómini: et credidérunt quotquot erant præordináti ad vitam ætérnam. Disseminabátur autem verbum Dómini per univérsam regiónem. Judǽi autem concitavérunt mulíeres religiósas et honéstas, et primos civitátis, et excitavérunt persecutiónem in Paulum et Bárnabam: et ejecérunt eos de fínibus suis. At illi, excússo púlvere pedum in eos, venérunt Icónium. Discípuli quoque replebántur gáudio, et Spíritu Sancto.\par
\switchcolumn
\EN
\noindent And the Gentiles hearing it, were glad, and glorified the word of the Lord: and as many as were ordained to life everlasting, believed. And the word of the Lord was published throughout the whole country. But the Jews stirred up religious and honourable women, and the chief men of the city, and raised persecution against Paul and Barnabas: and cast them out of their coasts. But they, shaking off the dust of their feet against them, came to Iconium. And the disciples were filled with joy and with the Holy Ghost.\par
\switchcolumn*

\LA
\begin{Resp}\Rsign Tóllite jugum meum super vos, dicit Dóminus, et díscite a me, quia mitis sum et húmilis corde: \Star{} Jugum enim meum suáve est, et onus meum leve.\end{Resp}
\begin{Resp}\Vsign Et inveniétis réquiem animábus vestris.\end{Resp}
\begin{Resp}\Rsign Jugum enim meum suáve est, et onus meum leve.\end{Resp}
\switchcolumn
\EN
\begin{Resp}\Rsign Take up my yoke upon you, said the Lord, and learn of me, because I am meek, and humble of heart. \Star{} For my yoke is sweet and my burden light.\end{Resp}
\begin{Resp}\Vsign And you shall find rest to your souls.\end{Resp}
\begin{Resp}\Rsign For my yoke is sweet and my burden light.\end{Resp}
\switchcolumn*

\LA
\LessonHead{Lectio III}
\switchcolumn
\EN
\LessonHead{Lesson III}
\switchcolumn*

\LA
\Cite{Act 14:1-3}
\switchcolumn
\EN
\Cite{Acts 14:1-3}
\switchcolumn*

\LA
\noindent Factum est autem Icónii, ut simul introírent in synagógam Judæórum, et loqueréntur, ita ut créderet Judæórum et Græcórum copiósa multitúdo. Qui vero incréduli fuérunt Judǽi, suscitavérunt et ad iracúndiam concitavérunt ánimas Géntium advérsus fratres. Multo ígitur témpore demoráti sunt, fiduciáliter agéntes in Dómino, testimónium perhibénte verbo grátiæ suæ, dante signa et prodígia fíeri per manus eórum.\par
\switchcolumn
\EN
\noindent And it came to pass in Iconium, that they entered together into the synagogue of the Jews, and so spoke that a very great multitude both of the Jews and of the Greeks did believe. But the unbelieving Jews stirred up and incensed the minds of the Gentiles against the brethren. A long time therefore they abode there, dealing confidently in the Lord, who gave testimony to the word of his grace, granting signs and wonders to be done by their hands.\par
\switchcolumn*

\LA
\begin{Resp}\Rsign Dum stetéritis ante reges et prǽsides, nolíte cogitáre quómodo aut quid loquámini: \Star{} Dábitur enim vobis in illa hora, quid loquámini.\end{Resp}
\begin{Resp}\Vsign Non enim vos estis qui loquímini: sed Spíritus Patris vestri, qui lóquitur in vobis.\end{Resp}
\begin{Resp}\Rsign Dábitur enim vobis in illa hora, quid loquámini.\end{Resp}
\begin{Resp}\Vsign Glória Patri, et Fílio, \Star{} et Spirítui Sancto.\end{Resp}
\begin{Resp}\Rsign Dábitur enim vobis in illa hora, quid loquámini.\end{Resp}
\switchcolumn
\EN
\begin{Resp}\Rsign But when they shall deliver you up to the judges, take no thought how or what to speak: \Star{} for it shall be given you in that hour what to speak:\end{Resp}
\begin{Resp}\Vsign For it is not you that speak, but the spirit of your Father that speaketh in you.\end{Resp}
\begin{Resp}\Rsign For it shall be given you in that hour what to speak.\end{Resp}
\begin{Resp}\Vsign Glory be to the Father, and to the Son, \Star{} and to the Holy Ghost.\end{Resp}
\begin{Resp}\Rsign For it shall be given you in that hour what to speak.\end{Resp}
\switchcolumn*

\LA
\LessonHead{Lectio IV}
\switchcolumn
\EN
\LessonHead{Lesson IV}
\switchcolumn*

\LA
\noindent Bárnabas Levítes, Cýprius génere, qui et Joseph, cum Paulo Géntium Apóstolus ordinátus est ad prædicándum Jesu Christi Evangélium. Is, agro véndito quem habébat, redáctam ex eo pecúniam áttulit Apóstolis. Missus autem Antiochíam prædicatiónis causa, cum ibi multos ad Christi Dómini fidem convérsos esse comperísset, incredibíliter lætátus, eos hortabátur ut in Christi fide permanérent. Qua cohortatióne multum proficiébat, quod ab ómnibus vir bonus et Spíritu Sancto plenus habebátur.\par
\switchcolumn
\EN
\noindent Joseph, who by the Apostles was surnamed Barnabas, (which is, being interpreted, the Son of Consolation,) a Levite and of the country of Cyprus, having land, sold it, and brought the money, and laid it at the Apostles' feet. (Acts iv. 36, 37.) When Paul, after his conversion, was come to Jerusalem, the disciples were all afraid of him, but Barnabas took him, and brought him to the Apostles, (ix. 26, 27,) When tidings that a great number believed and turned unto the Lord at Antioch came unto the ears of the Church which was at Jerusalem, they sent forth Barnabas that he should go as far as Antioch. Who, when he came, and had seen the grace of God, was glad, and exhorted them all that with purpose of heart they would cleave unto the Lord. For he was a good man, and full of the Holy Ghost, and of faith, and much people was added unto the Lord. (xi. 21-24.)\par
\switchcolumn*

\LA
\begin{Resp}\Rsign Vidi conjúnctos viros, habéntes spléndidas vestes, et Ángelus Dómini locútus est ad me, dicens: \Star{} Isti sunt viri sancti facti amíci Dei.\end{Resp}
\begin{Resp}\Vsign Vidi Ángelum Dei fortem, volántem per médium cælum, voce magna clamántem et dicéntem.\end{Resp}
\begin{Resp}\Rsign Isti sunt viri sancti facti amíci Dei.\end{Resp}
\switchcolumn
\EN
\begin{Resp}\Rsign I saw men standing together, clad in shining raiment, and the Angel of the Lord spoke unto me, saying: \Star{} These men are holy, for they are the friends of God.\end{Resp}
\begin{Resp}\Vsign I saw a strong Angel of God fly into the midst of heaven, saying with a loud voice:\end{Resp}
\begin{Resp}\Rsign These men are holy, for they are the friends of God.\end{Resp}
\switchcolumn*

\LA
\LessonHead{Lectio V}
\switchcolumn
\EN
\LessonHead{Lesson V}
\switchcolumn*

\LA
\noindent Proféctus inde Tarsum ut quæreret Paulum, cum eo Antiochíam venit. In ejus urbis ecclésia annum commoráti, christiánæ fídei et vitæ illis homínibus præcépta dedérunt: ubi étiam Jesu Christi cultóres primum Christiáni sunt appelláti. Discípuli autem Pauli et Bárnabæ, suis facultátibus Christiános, qui in Judǽa erant, sustentábant, eo mitténtes pecúniam per Paulum et Bárnabam. Qui perfúncti illo caritátis offício, adhíbito Joánne, cui cognómen erat Marcus, rediérunt Antiochíam.\par
\switchcolumn
\EN
\noindent Then departed Barnabas to Tarsus for to seek Paul, and, when he had found him, he brought him unto Antioch. And it came to pass that a whole year they assembled themselves with the Church, and taught much people. And the disciples were called Christians first in Antioch. And in these days came prophets from Jerusalem unto Antioch. And there stood up one of them, named Agabus, and signified, by the Spirit, that there should be great dearth throughout all the world which came to pass in the days of Claudius Caesar. Then the disciples, every man according to his ability, determined to send relief unto the brethren which dwelt in Judea, which also they did, and sent it to the elders by the hands of Barnabas and Paul. (xi. 25-30.) And Barnabas and Paul returned from Jerusalem, when they had fulfilled their ministry, and took with them John, whose surname was Mark. (xii. 25.)\par
\switchcolumn*

\LA
\begin{Resp}\Rsign Beáti estis, cum maledíxerint vobis hómines, et persecúti vos fúerint, et díxerint omne malum advérsum vos, mentiéntes, propter me: \Star{} Gaudéte et exsultáte, quóniam merces vestra copiósa est in cælis.\end{Resp}
\begin{Resp}\Vsign Cum vos óderint hómines, et cum separáverint vos, et exprobráverint, et ejécerint nomen vestrum tamquam malum propter Fílium hóminis.\end{Resp}
\begin{Resp}\Rsign Gaudéte et exsultáte, quóniam merces vestra copiósa est in cælis.\end{Resp}
\switchcolumn
\EN
\begin{Resp}\Rsign Blessed are ye when men shall revile you, and persecute you, and shall say all manner of evil against you falsely, for My sake; \Star{} Rejoice, and be exceeding glad, for great is your reward in heaven.\end{Resp}
\begin{Resp}\Vsign When men shall hate you, and when they shall separate you from their company, and shall reproach you, and cast out your name as evil, for the Son of Man's sake.\end{Resp}
\begin{Resp}\Rsign Rejoice, and be exceeding glad, for great is your reward in heaven.\end{Resp}
\switchcolumn*

\LA
\LessonHead{Lectio VI}
\switchcolumn
\EN
\LessonHead{Lesson VI}
\switchcolumn*

\LA
\noindent Cum autem Antiochíæ in Ecclésia, cum céteris prophétis et doctóribus, Paulus et Bárnabas in jejúnio et oratióne Dómino deservírent, dixit Spíritus Sanctus: Segregáte mihi Saulum et Bárnabam in opus, ad quod assúmpsi eos. Tunc jejunántes et orántes, imponentésque eis manus, dimisérunt illos. Itaque Seleucíam venérunt, inde in Cyprum; ac multas prætérea urbes regionésque, prædicántes Evangélium summa cum audéntium utilitáte, peragrárunt. Postrémo Bárnabas digréssus a Paulo, una cum Joánne, qui cognominátus est Marcus, navigávit in Cyprum; ibíque círciter séptimum Nerónis annum, tértio Idus Júnii ad apostólici múneris laudem, martyrii corónam adjúnxit. Ejus corpus, Zenóne imperatóre, repértum est in ínsula Cypro; ad cujus pectus erat Evangélium Matthæi, Bárnabæ manu conscríptum.\par
\switchcolumn
\EN
\noindent Now there were in the Church that was at Antioch, certain Prophets and teachers and, as Paul and Barnabas, together with them, ministered to the Lord and fasted, the Holy Ghost said Separate Me Barnabas and Saul for the work whereunto I have called them. And when they had fasted and prayed, and laid their hands on them, they sent them away. So they, being sent forth by the Holy Ghost, departed unto Seleucia and from thence they sailed to Cyprus (xiii. 1-4) in the which island, and in many other cities and countries, they journeyed about, preaching the Gospel with great gain to them that heard them. Nevertheless, at last, Paul and Barnabas departed asunder one from the other. And so Barnabas took Mark and sailed unto Cyprus, (xv. 39,) once more. And there it was that upon a certain nth of June, in or about the seventh year of the reign of Nero, Barnabas crowned the dignity of the Apostolate with the glory of martyrdom. During the reign of the Emperor Zeno, his body was found in its grave in Cyprus on his breast lay a copy of the Gospel according to Matthew, written by the hand of Barnabas himself.\par
\switchcolumn*

\LA
\begin{Resp}\Rsign Isti sunt triumphatóres et amíci Dei, qui contemnéntes jussa príncipum, meruérunt prǽmia ætérna: \Star{} Modo coronántur, et accípiunt palmam.\end{Resp}
\begin{Resp}\Vsign Isti sunt qui venérunt ex magna tribulatióne, et lavérunt stolas suas in sánguine Agni.\end{Resp}
\begin{Resp}\Rsign Modo coronántur, et accípiunt palmam.\end{Resp}
\begin{Resp}\Vsign Glória Patri, et Fílio, \Star{} et Spirítui Sancto.\end{Resp}
\begin{Resp}\Rsign Modo coronántur, et accípiunt palmam.\end{Resp}
\switchcolumn
\EN
\begin{Resp}\Rsign These are they which have conquered, and are become the friends of God, who recked not of the commandments of princes, and earned the everlasting reward. \Star{} And now have they crowns on their heads, and palms in their hands.\end{Resp}
\begin{Resp}\Vsign These are they which came out of great tribulation, and have washed their robes in the blood of the Lamb.\end{Resp}
\begin{Resp}\Rsign And now have they crowns on their heads, and palms in their hands.\end{Resp}
\begin{Resp}\Vsign Glory be to the Father, and to the Son, \Star{} and to the Holy Ghost.\end{Resp}
\begin{Resp}\Rsign And now have they crowns on their heads, and palms in their hands.\end{Resp}
\switchcolumn*

\LA
\LessonHead{Lectio VII}
\switchcolumn
\EN
\LessonHead{Lesson VII}
\switchcolumn*

\LA
\LessonTitle{Léctio sancti Evangélii secúndum Matthǽum.}
\switchcolumn
\EN
\LessonTitle{From the Holy Gospel according to Matthew}
\switchcolumn*

\LA
\Cite{Matt 10:16-22}
\switchcolumn
\EN
\Cite{Matt 10:6-22}
\switchcolumn*

\LA
\noindent In illo témpore: Dixit Jesus discípulis suis: Ecce ego mitto vos sicut oves in médio lupórum. Et réliqua.\par
\switchcolumn
\EN
\noindent At that time: Jesus said unto His disciples Behold, I send you forth as sheep in the midst of wolves. And so on.\par
\switchcolumn*

\LA
\LessonTitle{Homilía sancti Joánnis Chrysóstomi.}
\switchcolumn
\EN
\LessonTitle{Homily by St. John Chrysostom, Patriarch of Constantinople}
\switchcolumn*

\LA
\Source{Homilia 34 in Matthæum, post initium}
\switchcolumn
\EN
\Source{34/A on Matthew.}
\switchcolumn*

\LA
\noindent Cum Dóminus omnem sollicitúdinem a discipulórum córdibus ejécerit, et ostensióne signórum armáverit, atque, ab ómnibus negótiis sæculáribus alienátos et ab omni temporálium rerum cura liberátos, férreos quodámmodo atque adamántinos fécerit, tum dénique eventúra illis advérsa prædicit. Multa enim ex hac prædicatióne futurárum rerum cómmoda consequebántur. Primum, ut ejus præsciéntiæ vim edíscerent. Deínde, ut nemo suspicarétur, ex Magístri infirmitáte tam grávia mala descéndere. Prætérea, ne, qui ea passúri erant, súbito ac inopináto rerum evéntu perturbaréntur. Dénique, ne, cum ista sub ipsum passiónis suæ tempus audírent, nímium commoveréntur.\par
\switchcolumn
\EN
\noindent When the Lord had cleared the minds of His disciples of all care, and had armed them by showing forth His mighty works, had estranged them from all business of this world, and freed them from all anxiety touching the things of time, moulding them into a frame of iron-like, nay, diamondlike, hardness, then at length He told them of the contendings against the which they were afterward to wrestle. By this foretelling of things to come they were much holpen. First, they learnt the power of His fore-knowledge. Then, they were guarded against all suspicion that these great sorrows flowed from faultiness in their Master. Again, the future sufferers were made safe from all trouble of being taken unawares. Lastly, seeing that they heard these things at a time nigh to His own suffering, they were not over troubled.\par
\switchcolumn*

\LA
\begin{Resp}\Rsign Isti sunt qui vivéntes in carne, plantavérunt Ecclésiam sánguine suo: \Star{} Cálicem Dómini bibérunt, et amíci Dei facti sunt.\end{Resp}
\begin{Resp}\Vsign In omnem terram exívit sonus eórum, et in fines orbis terræ verba eórum.\end{Resp}
\begin{Resp}\Rsign Cálicem Dómini bibérunt, et amíci Dei facti sunt.\end{Resp}
\switchcolumn
\EN
\begin{Resp}\Rsign These are they who while yet they lived in the flesh, planted the Church in their own blood; \Star{} They drank of the Lord's cup, and became the friends of God.\end{Resp}
\begin{Resp}\Vsign Their sound is gone out through all the earth, and their words to the ends of the world.\end{Resp}
\begin{Resp}\Rsign They drank of the Lord's cup, and became the friends of God.\end{Resp}
\switchcolumn*

\LA
\LessonHead{Lectio VIII}
\switchcolumn
\EN
\LessonHead{Lesson VIII}
\switchcolumn*

\LA
\noindent Jam vero, ut intélligant novum hoc esse belli genus et insólitum præliándi morem, cum illos nudos mítteret, una indútos túnica, sine cálceis, absque virga et absque zona et pera, et ab excipiéntibus ali jubéret; non fecit hic dicéndi finem, sed inexplicábilem virtútem suam próferens, Etiam sic eúntes, inquit, mansuetúdinem tamen óvium osténdite, quamvis ad lupos itúri, nec simplíciter ad lupos, sed étiam in médio lupórum: (neque vero óvium tantum mansuetúdinem habére jubet, sed étiam colúmbæ simplicitátem); sic enim virtútem meam máxime osténdam, cum ab óvibus lupi superabúntur; et quamvis illæ sint in médio lupórum, et innúmeris móribus laceréntur, non modo non consúmptæ fúerint, verum étiam illos in sui natúram transmutáverint.\par
\switchcolumn
\EN
\noindent And now, that they may understand how that this is a new kind of warfare, and an unaccustomed manner of contending, when He sendeth them forth unarmed, providing neither gold, nor silver, nor brass in their purses nor scrip for their journey, neither two coats, neither shoes, nor yet staves, (x. 9, 10,) left to the hospitality of whosoever would receive them, He maketh not here an end to His discourse, but, in manifestation of His unspeakable power, He biddeth them, so going, to show forth the meekness of sheep, seeing they were about going unto wolves neither simply unto wolves, but in the very midst of wolves. Neither is it only the meekness of sheep which He biddeth them have, but also the harmlessness of doves, that He might so much the more gloriously display His power, when the sheep overcame the wolves. These are the sheep which albeit they abide in the midst of wolves, and are mangled by many a bite, not only are not destroyed, but do gradually make the wolves change their nature, and become sheep themselves.\par
\switchcolumn*

\LA
\begin{Resp}\Rsign Isti sunt viri sancti, quos elégit Dóminus in caritáte non ficta, et dedit illis glóriam sempitérnam: \Star{} Quorum doctrína fulget Ecclésia, ut sole luna.\end{Resp}
\begin{Resp}\Vsign Sancti per fidem vicérunt regna: operáti sunt justítiam.\end{Resp}
\begin{Resp}\Rsign Quorum doctrína fulget Ecclésia, ut sole luna.\end{Resp}
\begin{Resp}\Vsign Glória Patri, et Fílio, \Star{} et Spirítui Sancto.\end{Resp}
\begin{Resp}\Rsign Quorum doctrína fulget Ecclésia, ut sole luna.\end{Resp}
\switchcolumn
\EN
\begin{Resp}\Rsign These men are saints, whom the Lord hath chosen in love unfeigned, and hath given them glory everlasting. These are they \Star{} By the light of whose teaching the Church is glorified, even as the moon is glorified by the light of the sun.\end{Resp}
\begin{Resp}\Vsign The saints through faith subdued kingdoms, wrought righteousness.\end{Resp}
\begin{Resp}\Rsign By the light of whose teaching the Church is glorified, even as the moon is glorified by the light of the sun.\end{Resp}
\begin{Resp}\Vsign Glory be to the Father, and to the Son, \Star{} and to the Holy Ghost.\end{Resp}
\begin{Resp}\Rsign By the light of whose teaching the Church is glorified, even as the moon is glorified by the light of the sun.\end{Resp}
\switchcolumn*

\LA
\LessonHead{Lectio IX}
\switchcolumn
\EN
\LessonHead{Lesson IX}
\switchcolumn*

\LA
\noindent Majus certe atque admirabílius est mentem adversariórum commutáre, et ánimum in divérsum transférre, quam illos occídere; præsértim cum duódecim tantum essent, et lupis plenus esset orbis univérsus. Erubescámus ígitur, qui, longe divérsa faciéntes, tamquam lupi in adversários rúimus. Nam, quámdiu oves fuérimus, víncimur; tunc enim a nobis pastóris auxílium recédit, qui non lupos, sed oves pascit.\par
\switchcolumn
\EN
\noindent Beyond all doubt it is a greater and more marvellous thing to change the minds of enemies, and to turn their thoughts round, than to kill them more especially when the work is to be done by only twelve sheep, and the whole world is full of the wolves. Shame then upon us, whose deeds are so contrary, and who rather run like wolves upon our enemies. For so long as we are sheep we conquer, yea, though a thousand wolves be gathered round about us, we overcome, and are the conquerors but if we become wolves ourselves, then are we conquered. For then doth the Shepherd's help forsake us, Who feedeth not wolves but sheep.\par
\switchcolumn*

\LA
\TeDeum{Te Deum}
\switchcolumn
\EN
\TeDeum{Te Deum}
\switchcolumn*

\LA
\LessonHead{Lectio IX (pro commemoratione)}
\switchcolumn
\EN
\LessonHead{Lesson IX (when commemorated)}
\switchcolumn*

\LA
\LessonTitle{Commemoratio: S. Barnabæ Apostoli}
\switchcolumn
\EN
\LessonTitle{Commemoration of: St. Barnabas the Apostle}
\switchcolumn*

\LA
\noindent Bárnabas Levítes, Cýprius génere, cum Paulo Apóstolo ordinátus est ad prædicándum Jesu Christi Evangélium. Is, agro véndito quem habébat, redáctam ex eo pecúniam áttulit Apóstolis. Missus Antiochíam prædicatiónis causa, multos ibi ad Christi fidem convérsos suis hortatiónibus confirmávit. Inde proféctus cum eódem Paulo, multas urbes regionésque summa cum audiéntium utilitáte peragrávit. Postrémo, digréssus a Paulo, una cum Joánne, qui cognominátus est Marcus, navigávit in Cyprum, ibíque, séptimo Nerónis anno, ad apostólici múneris laudem martýrii corónam adjúnxit.\par
\switchcolumn
\EN
\noindent Barnabas the Levite, of the people of Cyprus, was sent with Paul the Apostle to preach the Gospel of Jesus Christ. He sold a field he owned and brought the money from it to the Apostles. Sent to Antioch to preach, he gave new strength by his exhortations to many who had been converted to belief in Christ. Thence he set forth with St. Paul and travelled through many cities and countries, to the great benefit of his hearers. Finally he left Paul, and with John, surnamed Mark, he took ship to Cyprus. There, in the seventh year of Nero's reign, he added to the dignity of the apostolic office the crown of martyrdom.\par
\switchcolumn*

\LA
\TeDeum{Te Deum}
\switchcolumn
\EN
\TeDeum{Te Deum}
\switchcolumn*

\end{paracol}

\par\needspace{10\baselineskip}\addvspace{1.2em}{\centering\small\textsc{Die 12 Junii} · \textsc{12 June}\par}
\Office{S. Joannis a S. Facundo Confessoris}{St. John of St. Facundus, Confessor}{Duplex}
\addcontentsline{toc}{subsection}{Die 12 Junii · S. Joannis a S. Facundo Confessoris \textit{— St. John of St. Facundus, Confessor}}
\begin{paracol}{2}
\LA
\RefLine{Lectiones I–III: de Scriptura occurrente.}
\switchcolumn
\EN
\RefLine{Lessons I–III: from the Scripture of the day.}
\switchcolumn*

\LA
\RefLine{Lectiones VII–VIII: ex Commúni Confessoris non pontificis (C5).}
\switchcolumn
\EN
\RefLine{Lessons VII–VIII: from the Common of a Confessor not a Bishop (C5).}
\switchcolumn*

\LA
\RefLine{Lectio IX: de Ss. Basilidis, Cyrini, Naboris and Nazarii Martyrum (06-12o).}
\switchcolumn
\EN
\RefLine{Lesson IX: of Sts. Basilides, Cyrinus, Nabor and Nazarius, Martyrs (06-12o).}
\switchcolumn*

\LA
\LessonHead{Lectio IV}
\switchcolumn
\EN
\LessonHead{Lesson IV}
\switchcolumn*

\LA
\noindent Joánnem Sahaguni in Hispania nobili genere natum, parentes cum diu prole caruissent, piis operibus, et orationibus a Deo impetrarunt. Ab ineunte ætate egregium futuræ sanctitatis specimen dedit: nam e loco superiori ad ceteros pueros crebro verba faciebat, quibus eos ad virtutem, et Dei cultum hortabatur, eorumque dissidia componebat. In patria monachis sancti Facundi ordinis sancti Benedicti primis litterarum rudimentis imbuendus traditur. Dum iis operam daret, curavit pater, ut parochus ecclesiam administraret: quod munus juvenis nullis rationibus adduci potuit, ut retineret. Inter familiares episcopi Burgensis adscriptus, ob spectatam ipsius probitatem intimus ei fuit, ab eoque presbyter, et canonicus factus multis beneficiis auctus est. Sed relicta aula episcopi ut Deo quietius serviret, omnibus ecclesiæ proventibus abdicatis se cuidam sacello addixit, ubi Sacrum quotidie faciebat, ac de rebus divinis magna cum auditorum ædificatione frequenter concionabatur.\par
\switchcolumn
\EN
\noindent John Gonzalez was born, the offspring of a noble race, at San Fagondez in Spain, (on Midsummer Day in the year of grace 1430.) His father and mother after long childlessness, obtained him from God by prayers and good works. From his earliest years he gave clear signs of his after holiness of life. He was used to climb up upon an high place to preach to the other little boys, and to exhort them to be good and to worship God, and he made it his work to reconcile their quarrels. While he was still at home he was given in charge to the monks of the Order of Saint Benedict, at the village of San Fagondez, to teach him his first lessons. While he was thus busied, his father obtained for him the benefice of the Parish, but no persuasions could induce him to keep this preferment. He became one of the household of the Bishop of Burgos, and that Prelate, seeing his uprightness, took him into his counsels, ordained him Priest, and made him a Canon, heaping upon him many kindnesses. However, that he might serve God the more quietly, he left the Bishop's Palace, resigned all his Church income, and betook him to a certain Chapel wherein he celebrated the Holy Liturgy every day, and oftentimes preached concerning the things of God, with great profit to all that heard him.\par
\switchcolumn*

\LA
\begin{Resp}\Rsign Honéstum fecit illum Dóminus, et custodívit eum ab inimícis, et a seductóribus tutávit illum: \Star{} Et dedit illi claritátem ætérnam.\end{Resp}
\begin{Resp}\Vsign Justum dedúxit Dóminus per vias rectas, et osténdit illi regnum Dei.\end{Resp}
\begin{Resp}\Rsign Et dedit illi claritátem ætérnam.\end{Resp}
\switchcolumn
\EN
\begin{Resp}\Rsign The Lord made him honourable, and defended him from his enemies, and kept him safe from those that lay in wait for him; \Star{} And gave him perpetual glory.\end{Resp}
\begin{Resp}\Vsign He went down with him into the pit, and left him not in bonds.\end{Resp}
\begin{Resp}\Rsign And gave him perpetual glory.\end{Resp}
\switchcolumn*

\LA
\LessonHead{Lectio V}
\switchcolumn
\EN
\LessonHead{Lesson V}
\switchcolumn*

\LA
\noindent Postea studiorum causa Salmanticam profectus, in celebre collegium divi Bartholomæi cooptatus, sacerdotis munus ita exercuit, ut simul optatis studiis incumberet, et in sacris étiam concionibus assidue versaretur. Cum vero in gravissimum morbum incidisset, arctioris disciplinæ voto se obstrinxit, quod ut redderet, cum prius cuidam pauperi pene nudo ex duabus, quas tantum habebat vestes, meliorem dedisset, ad cœnobium sancti Augustíni severiori disciplina tum maxime florens se contulit: in quo admissus, obedientia, animi demissione, vigiliis ac oratione provectiores anteibat. Triclinii cura cum ipsi demandata esset, vini doliolum ipso attingente omnibus monachis per annum abunde suffecit. Exacto tirocinii anno, præfecti jussu munus concionandi suscepit. Salmanticæ id temporis adeo cruentis factionibus divina, humanaque omnia permixta erant, ut singulis propemodum horis cædes fierent, et omnium ordinum, ac præsertim nobilium sanguine non viæ solum, et fora, sed templa étiam redundarent.\par
\switchcolumn
\EN
\noindent He went later to Salamanca to study, and there being taken into the celebrated College of St. Bartholomew, he did his priestly office, so that he was at once constant to the studies he desired and busy with sermons. Here he had a severe illness, and vowed to take up a sterner way of living. In fulfillment of this vow, he gave to an half-naked beggar the better of the two garments which were all that he had, and then went to a Convent of the friars of St. Augustine, which was then in the richest bloom of rigid discipline. Being admitted therein, he surpassed the most advanced in obedience, lowliness, watchings, and prayer. At the time that he had charge of the table, one keg of wine abundantly sufficed in his hands for all the friars, throughout an whole year. After his year of novitiate, he undertook the duty of preacher at the command of his Superior. At that time, owing to bloody feuds, all things human and divine at Salamanca were in such utter confusion, that murders were committed almost every hour, and the streets and squares, and the very churches, flowed with the blood of all classes, especially of the nobility.\par
\switchcolumn*

\LA
\begin{Resp}\Rsign Amávit eum Dóminus, et ornávit eum: stolam glóriæ índuit eum, \Star{} Et ad portas paradísi coronávit eum.\end{Resp}
\begin{Resp}\Vsign Índuit eum Dóminus lorícam fídei, et ornávit eum.\end{Resp}
\begin{Resp}\Rsign Et ad portas paradísi coronávit eum.\end{Resp}
\switchcolumn
\EN
\begin{Resp}\Rsign The Lord loved him and beautified him; He clothed him with a robe of glory, \Star{} And crowned him at the gates of Paradise.\end{Resp}
\begin{Resp}\Vsign The Lord hath put on him the breast-plate of faith, and hath adorned him.\end{Resp}
\begin{Resp}\Rsign And crowned him at the gates of Paradise.\end{Resp}
\switchcolumn*

\LA
\LessonHead{Lectio VI}
\switchcolumn
\EN
\LessonHead{Lesson VI}
\switchcolumn*

\LA
\noindent At Joannes, tum concionibus, tum privatis colloquiis civium animos demulcens, ad tranquillitatem urbem reduxit. Virum principem graviter offendit, quod illius in subditos sævitiam increpasset. Qua de causa equites duos immisit, qui eum in itinere confoderent, jamque ad ipsum propinquaverant: cum, stupore divinitus immisso, simul cum equis immobiles steterunt, donec ad pedes sancti viri provoluti, sceleris veniam precarentur. Ipse quoque princeps repentino terrore perculsus, jam de salute desperaverat, cum revocato Joanne, facti pœnitens incolumitati redditus est. Factiosi étiam homines, cum eum fustibus peterent, brachiis diriguere, nec ante redditæ vires quam delicti veniam precarentur. Christum Dominum, dum Sacrum faceret, præsentem contueri, atque ex ipso divinitatis fonte cælestia mysteria haurire solitus. Abdita cordis inspicere, ac futura raro eventu præsagire frequens illi fuit, fratrisque filiam septennem mortuam excitavit. Denique mortis die prænuntiato, et Ecclesiæ sacramentis devotissime susceptis, extremum diem clausit, multis ante et post obitum miraculis gloriosus. Quibus rite probatis Alexander octavus Sanctorum numero eum adscripsit.\par
\switchcolumn
\EN
\noindent It was John, who by public preaching and private conversations, softened the hearts of the citizens so that the town was restored to peace. He grievously offended one of the nobles by rebuking him for his cruelty toward his vassals. This man sent two knights to murder him on the road. They had already come nigh him when God sent a terror upon them, so that they and their horses stood still, until they cast themselves down before the feet of the Saint, imploring his forgiveness for their sin. The Prince himself, also, smitten with a sudden dread, despaired of his salvation, till he had sent for John, who, finding him repent of his deed, restored him to soundness. Some quarrelsome men, likewise, who were fain to give him a cudgelling, found their arms stiffen, nor would their strength come back till they had asked his pardon for their wickedness. Oftentimes when he was celebrating the Holy Liturgy, the Presence of the Lord Christ became sensibly manifest to him, and he drank in things heavenly from their Divine Well-head Himself. Oftentimes also he could see the secrets of men's hearts, and foretell strange things to come. He raised from the dead his own niece, aged seven years. He foretold the day of his own death, and prepared himself by receiving most devoutly the Sacraments of the Church, (and then fell asleep in the Lord, upon the 11th day of June, in the year 1475.) God glorified him by many miracles, both before and after his death. These being duly proved, Alexander VIII. numbered him among the Saints. the Christians took them, and buried them honourably.\par
\switchcolumn*

\LA
\begin{Resp}\Rsign Iste homo perfécit ómnia quæ locútus est ei Deus, et dixit ad eum: Ingrédere in réquiem meam: \Star{} Quia te vidi justum coram me ex ómnibus géntibus.\end{Resp}
\begin{Resp}\Vsign Iste est, qui contémpsit vitam mundi, et pervénit ad cæléstia regna.\end{Resp}
\begin{Resp}\Rsign Quia te vidi justum coram me ex ómnibus géntibus.\end{Resp}
\begin{Resp}\Vsign Glória Patri, et Fílio, \Star{} et Spirítui Sancto.\end{Resp}
\begin{Resp}\Rsign Quia te vidi justum coram me ex ómnibus géntibus.\end{Resp}
\switchcolumn
\EN
\begin{Resp}\Rsign This is he which did according unto all that God commanded him; and God said unto him: Enter thou into My rest, \Star{} For thee have I seen righteous before Me among all people.\end{Resp}
\begin{Resp}\Vsign This is he which loved not his life in this world, and is come unto an everlasting kingdom.\end{Resp}
\begin{Resp}\Rsign For thee have I seen righteous before Me among all people.\end{Resp}
\begin{Resp}\Vsign Glory be to the Father, and to the Son, \Star{} and to the Holy Ghost.\end{Resp}
\begin{Resp}\Rsign For thee have I seen righteous before Me among all people.\end{Resp}
\switchcolumn*

\LA
\LessonHead{Lectio IX (pro commemoratione)}
\switchcolumn
\EN
\LessonHead{Lesson IX (when commemorated)}
\switchcolumn*

\LA
\LessonTitle{Commemoratio: S. Joannis a S. Facundo Confessoris}
\switchcolumn
\EN
\LessonTitle{Commemoration of: St. John of St. Facundus, Confessor}
\switchcolumn*

\LA
\noindent Joánnem, Sahagúni in Hispánia, nóbili génere natum, paréntes cum diu prole caruíssent, piis opéribus et oratiónibus a Deo impetrárunt. Ab ineúnte ætáte futúræ sanctitátis indícia prǽbuit. Présbyter ordinátus, ut Deo quiétius servíret, omnes ecclesiásticos provéntus, quibus mérito auctus fúerat, sponte dimísit. Salmánticæ, cum in gravíssimum morbum incidísset, arctióris disciplínæ voto se obstrínxit, quod ut rédderet, ad cœnóbium sancti Augustíni, severióri disciplína tum máxime florens, se cóntulit; in quo admíssus, virtútibus omnis provectióres anteíbat. Salmanticénses cives, cruéntis factiónibus exagitátos, tum conciónibus tum privátis collóquiis ac vitæ sanctitáte, ad tranquillitátem redúxit, non semel a præsénti discrímine divínitus liberátus. Christum Dóminum, dum Sacrum fáceret, præséntem contuéri, ábdita cordis inspícere, ac futúra præsagíre frequens illi fuit. Dénique, mortis die prænuntiáto, sanctíssime ex hac vita migrávit, multis ante et post óbitum miráculis gloriósus. Quibus rite probátis, Alexánder octávus Sanctórum número eum adscrípsit.\par
\switchcolumn
\EN
\noindent John, born of a noble family at Sahagun [St. Facundus] in Spain, was granted by God to his parents, who had long been childless, in answer to their good works and prayers. From his earliest years, he gave signs of his future holiness. When ordained priest, he renounced, of his own accord, all the ecclesiastical benefits which had been given him, that he might serve God in greater tranquillity. When he incurred a serious illness in Salamanca, he bound himself by vow to observe a severer discipline. To do so, he went to the monastery of St. Augustine, where there flourished the greatest severity of discipline. As a religious, he excelled the most advanced monks in all virtues. Through his public talks and private conversations, as well as the holiness of his life, he brought back to peaceful living the citizens of Salamanca, who had been disturbed by bloody factions. In the course of this work, he was not infrequently saved from imminent death by divine power. The Lord Christ often appeared to him while he was celebrating Mass. Often also he could divine the secrets of hearts and foretell the future. At length, having predicted the day of his death, he departed this life in a most holy way, glorified by many miracles both before and after his death. These miracles were duly proved, and Alexander VIII enrolled him in the number of the Saints.\par
\switchcolumn*

\LA
\TeDeum{Te Deum}
\switchcolumn
\EN
\TeDeum{Te Deum}
\switchcolumn*

\end{paracol}

\par\needspace{10\baselineskip}\addvspace{1.2em}{\centering\small\textsc{Die 12 Junii} · \textsc{12 June}\par}
\Office{Ss. Basilidis, Cyrini, Naboris and Nazarii Martyrum}{Sts. Basilides, Cyrinus, Nabor and Nazarius, Martyrs}{Simplex}
\addcontentsline{toc}{subsection}{Die 12 Junii · Ss. Basilidis, Cyrini, Naboris and Nazarii Martyrum \textit{— Sts. Basilides, Cyrinus, Nabor and Nazarius, Martyrs}}
\begin{paracol}{2}
\LA
\LessonHead{Lectio IX (pro commemoratione)}
\switchcolumn
\EN
\LessonHead{Lesson IX (when commemorated)}
\switchcolumn*

\LA
\LessonTitle{Commemoratio: Ss. Basilidis, Cyrini, Naboris and Nazarii Martyrum}
\switchcolumn
\EN
\LessonTitle{Commemoration of: Sts. Basilides, Cyrinus, Nabor and Nazarius, Martyrs}
\switchcolumn*

\LA
\noindent Basilides, Cyrinus, Nabor et Nazarius, Romani milites, nobiles genere, et virtute illustres, christiana religione suscepta, cum Christum Dei Filium, Diocletiano imperatore, prædicarent, ab Aurelio præfecto Urbis comprehensi, et ut diis sacra facerent admoniti, ejus jussa contemnentes, missi sunt in carcerem. Quibus orantibus, cum subito clarissima lux oborta omnium oculis, qui ibidem essent, carcerem collustrasset: illo cælesti splendore commotus Marcellus custodiæ præpositus, multique alii, Christo Domino crediderunt. Verum postea e carcere emissi, ab imperatore Maximiano, cum ejus étiam neglecto imperio, unum Christum Deum et Dominum in ore haberent, scorpionibus cruciati iterum conjiciuntur in vincula: unde septimo die educti, et ante pedes imperatoris constituti, perstiterunt in irrisione inanium deorum, Jesum Christum Deum constantissime confitentes. Quam ob rem damnati, securi feriuntur. Quorum corpora feris objecta, nec ab illis tacta a Christianis honorifice sepulta sunt.\par
\switchcolumn
\EN
\noindent Basilides, Cyrinus, Nabor, and Nazarius were Roman soldiers, of illustrious birth, and distinguished gallantry. Having embraced the Christian Religion, and being found publishing that Christ was the Son of God, they were arrested by Aurelius, Praefect of Rome under the Emperor Diocletian. As they despised his orders to sacrifice to the gods, they were committed to prison. While they were at prayer there, a brilliant light broke forth before the eyes of all that were there, and shone in all the prison. Marcellinus the keeper of the prison and many others were moved by this heavenly glory to believe in the Lord ChristBasilides, Cyrinus, Nabor, and Nazarius were afterwards discharged out of the prison. However, in the reign of the Emperor Maximian, when they set light by his commands also, and had ever in their mouth that there is but one Christ, one God, and one Lord, they were tormented with whips loaded with metal, and again cast into chains. Thence, on the seventh day, they were brought out, and set before the Emperor, and there still persisted in mocking at the foolish idols, and declaring that Jesus Christ is God. They were accordingly condemned to death and beheaded. Their bodies were given to wild beasts to eat, but, as the creatures would not touch them,\par
\switchcolumn*

\LA
\TeDeum{Te Deum}
\switchcolumn
\EN
\TeDeum{Te Deum}
\switchcolumn*

\end{paracol}

\par\needspace{10\baselineskip}\addvspace{1.2em}{\centering\small\textsc{Die 13 Junii} · \textsc{13 June}\par}
\Office{S. Antonii de Padua Confessoris et Ecclesiæ Doctoris}{St. Anthony of Padua, Confessor}{Duplex}
\addcontentsline{toc}{subsection}{Die 13 Junii · S. Antonii de Padua Confessoris et Ecclesiæ Doctoris \textit{— St. Anthony of Padua, Confessor}}
\begin{paracol}{2}
\LA
\RefLine{Lectiones I–III: de Scriptura occurrente.}
\switchcolumn
\EN
\RefLine{Lessons I–III: from the Scripture of the day.}
\switchcolumn*

\LA
\RefLine{Lectiones VII–IX: ex Commúni Doctoris non Pontificis (C5a).}
\switchcolumn
\EN
\RefLine{Lessons VII–IX: from the Common of a Doctor not a Bishop (C5a).}
\switchcolumn*

\LA
\LessonHead{Lectio IV}
\switchcolumn
\EN
\LessonHead{Lesson IV}
\switchcolumn*

\LA
\noindent Antónius, Ulyssipóne in Lusitánia honéstis ortus paréntibus, et ab iis pie educátus, adoléscens, institútum canonicórum regulárium suscépit. Sed, cum córpora beatórum quinque Mártyrum fratrum Minórum Conímbriam transferéntur, qui paulo ante apud Marróchium pro Christi fide passi erant, martýrii desidério incénsus ad Franciscánum órdinem transívit. Mox eódem ardóre impúlsus, ad Saracénos ire perréxit; sed, advérsa valetúdine afflíctus et redíre coáctus, cum navi ad Hispániæ líttora ténderet, ventórum vi in Sicíliam delátus est.\par
\switchcolumn
\EN
\noindent Ferdinand de Bullones, afterwards called Anthony, was born of decent parents at Lisbon in Portugal, (on the Feast of the Assumption, in the year of grace 1195.) They gave him a godly training, and while he was still a young man, he joined an Institute of Canons Regular. However, when the bodies of the five holy martyred Friars Minor, who had just suffered in Morocco for Christ's sake, were brought to Coimbra, the desire to be himself a martyr took a strong hold upon him, and in 1220 he left the Canons Regular and became a Franciscan. The same yearning led him to attempt to go among the Saracens, but he fell sick on the way, and, being obliged to turn back, the ship in which he had embarked for Spain was driven by stress of weather to Sicily.\par
\switchcolumn*

\LA
\begin{Resp}\Rsign Honéstum fecit illum Dóminus, et custodívit eum ab inimícis, et a seductóribus tutávit illum: \Star{} Et dedit illi claritátem ætérnam.\end{Resp}
\begin{Resp}\Vsign Justum dedúxit Dóminus per vias rectas, et osténdit illi regnum Dei.\end{Resp}
\begin{Resp}\Rsign Et dedit illi claritátem ætérnam.\end{Resp}
\switchcolumn
\EN
\begin{Resp}\Rsign The Lord made him honourable, and defended him from his enemies, and kept him safe from those that lay in wait for him; \Star{} And gave him perpetual glory.\end{Resp}
\begin{Resp}\Vsign He went down with him into the pit, and left him not in bonds.\end{Resp}
\begin{Resp}\Rsign And gave him perpetual glory.\end{Resp}
\switchcolumn*

\LA
\LessonHead{Lectio V}
\switchcolumn
\EN
\LessonHead{Lesson V}
\switchcolumn*

\LA
\noindent Assísium e Sicília ad capítulum generále venit: inde in erémum montis Pauli in Æmília secéssit, ubi divínis contemplatiónibus, jejúniis et vigíliis diu vacávit. Póstea, sacris ordínibus initiátus et ad prædicándum Evangélium missus, dicéndi sapiéntia et cópia tantum profécit, tantámque sui admiratiónem commóvit, ut eum summus Póntifex aliquándo concionántem áudiens, arcam Testaménti appellárit. In primis vero hǽreses summa vi profligávit, ideóque perpétuus hæreticórum málleus est vocátus.\par
\switchcolumn
\EN
\noindent From Sicily he came to Assisi to attend the General Chapter of his Order, and thence withdrew himself to the Hermitage of Monte Paolo near Bologna, where he gave himself up for a long while to consideration of the things of God, to fastings, and to watchings. Being afterwards ordained Priest and sent to preach the Gospel, his wisdom and fluency were very marked, and drew on him such admiration of men, that the Pope, once hearing him preach, called him The Ark of the Covenant. One of his chief points was to expend all his strength in attacking heresies, whence he gained the name of the Heretics' everlasting Hammer.\par
\switchcolumn*

\LA
\begin{Resp}\Rsign Amávit eum Dóminus, et ornávit eum: stolam glóriæ índuit eum, \Star{} Et ad portas paradísi coronávit eum.\end{Resp}
\begin{Resp}\Vsign Índuit eum Dóminus lorícam fídei, et ornávit eum.\end{Resp}
\begin{Resp}\Rsign Et ad portas paradísi coronávit eum.\end{Resp}
\switchcolumn
\EN
\begin{Resp}\Rsign The Lord loved him and beautified him; He clothed him with a robe of glory, \Star{} And crowned him at the gates of Paradise.\end{Resp}
\begin{Resp}\Vsign The Lord hath put on him the breast-plate of faith, and hath adorned him.\end{Resp}
\begin{Resp}\Rsign And crowned him at the gates of Paradise.\end{Resp}
\switchcolumn*

\LA
\LessonHead{Lectio VI}
\switchcolumn
\EN
\LessonHead{Lesson VI}
\switchcolumn*

\LA
\noindent Primus ex suo ordine, ob doctrínæ præstántiam, Bonóniæ et álibi sacras lítteras est interpretátus, fratrúmque suórum stúdiis præfuit. Multis vero peragrátis provínciis, anno ante óbitum Patávium venit, ubi illústria sanctitátis suæ monuménta relíquit. Dénique magnis labóribus pro glória Dei perfúnctus, méritis et miráculis clarus obdormívit in Dómino Idibus Júnii, anno salútis millésimo ducentésimo trigésimo primo. Quem Gregórius nonus, Póntifex máximus, sanctórum Confessórum número adscrípsit, et Pius duodécimus, ex Sacrórum Rítuum Congregatiónis consúlto, universális Ecclésiæ Doctórem declarávit.\par
\switchcolumn
\EN
\noindent He was the first of his Order who, on account of his excellent gift of teaching, publicly lectured at Bologna on the interpretation of Holy Scripture, and directed the studies of his brethren. He traveled through many provinces. The year before his death he came to Padua, where he left some remarkable records of his holy life. After having undergone much toil for the glory of God, full of good works and miracles, he fell asleep in the Lord upon the 13th day of June, in the year of salvation 1231. Pope Gregory IX. enrolled his name among those of the Holy Confessors.\par
\switchcolumn*

\LA
\begin{Resp}\Rsign Iste homo perfécit ómnia quæ locútus est ei Deus, et dixit ad eum: Ingrédere in réquiem meam: \Star{} Quia te vidi justum coram me ex ómnibus géntibus.\end{Resp}
\begin{Resp}\Vsign Iste est, qui contémpsit vitam mundi, et pervénit ad cæléstia regna.\end{Resp}
\begin{Resp}\Rsign Quia te vidi justum coram me ex ómnibus géntibus.\end{Resp}
\begin{Resp}\Vsign Glória Patri, et Fílio, \Star{} et Spirítui Sancto.\end{Resp}
\begin{Resp}\Rsign Quia te vidi justum coram me ex ómnibus géntibus.\end{Resp}
\switchcolumn
\EN
\begin{Resp}\Rsign This is he which did according unto all that God commanded him; and God said unto him: Enter thou into My rest, \Star{} For thee have I seen righteous before Me among all people.\end{Resp}
\begin{Resp}\Vsign This is he which loved not his life in this world, and is come unto an everlasting kingdom.\end{Resp}
\begin{Resp}\Rsign For thee have I seen righteous before Me among all people.\end{Resp}
\begin{Resp}\Vsign Glory be to the Father, and to the Son, \Star{} and to the Holy Ghost.\end{Resp}
\begin{Resp}\Rsign For thee have I seen righteous before Me among all people.\end{Resp}
\switchcolumn*

\LA
\LessonHead{Lectio IX (pro commemoratione)}
\switchcolumn
\EN
\LessonHead{Lesson IX (when commemorated)}
\switchcolumn*

\LA
\LessonTitle{Commemoratio: S. Antonii de Padua Confessoris et Ecclesiæ Doctoris}
\switchcolumn
\EN
\LessonTitle{Commemoration of: St. Anthony of Padua, Confessor}
\switchcolumn*

\LA
\noindent Antónius, Ulyssipóne in Lusitánia honéstis piísque ortus paréntibus, adoléscens institútum canonicórum regulárium suscépit; sed, martýrii desidério incénsus, ad Franciscánum órdinem transívit. Ad Saracénos missus, et advérsa valetúdine redíre coáctus, vi ventórum in Sicíliam delátus est. Mox sacris ordínibus initiátus et prædicatóris múnere fungens, tantam sui admiratiónem commóvit, ut ad sacras lítteras interpretándas Bonóniæ et álibi vocátus sit, fratrum suórum stúdiis fúerit præféctus, et arca Testaménti atque hǽresum málleus merúerit appellári. Multis autem peragrátis provínciis, anno ante óbitum Patávium venit, ubi illústria sanctitátis suæ monuménta relíquit. Méritis et miráculis clarus obdormívit in Dómino idibus júnii, anno salútis millésimo ducentésimo trigésimo primo, ætátis suæ trigésimo sexto, et a summo Pontífice Pio duodécimo universális Ecclésiæ Doctor fuit declarátus.\par
\switchcolumn
\EN
\noindent Anthony came of good and devout parents at Lisbon in Portugal. As a young man, he entered the Institute of Canons Regular, but, fired with the desire for martyrdom, he changed to the Franciscan Order. He was sent to the Saracens, but was forced to return because of ill-health. On the return voyage, strong winds blew the ship off course to Sicily. Anthony was soon given holy orders and took up the work of preaching, in which he aroused so much admiration that he was called to interpret Sacred Scriptures at Bologna and other places. He was put in charge of this brethren's studies, and was deservedly called the Ark of the Covenant. and Hammer of Heretics. After many travels, he came, a year before his death, to Padua, where left shining memories of his holiness. Famous for both merits and miracles, he fell asleep in the Lord on the 13th day of June in the year of salvation 1231, at the age of thirty-six. He was declared a Doctor of the universal Church by Pope Pius XII.\par
\switchcolumn*

\LA
\TeDeum{Te Deum}
\switchcolumn
\EN
\TeDeum{Te Deum}
\switchcolumn*

\end{paracol}

\par\needspace{10\baselineskip}\addvspace{1.2em}{\centering\small\textsc{Die 14 Junii} · \textsc{14 June}\par}
\Office{S. Basilii Magni, Episcopis Confessoris et Ecclesiæ Doctoris}{St. Basil the Great, Confessor and Doctor of the Church}{Duplex}
\addcontentsline{toc}{subsection}{Die 14 Junii · S. Basilii Magni, Episcopis Confessoris et Ecclesiæ Doctoris \textit{— St. Basil the Great, Confessor and Doctor of the Church}}
\begin{paracol}{2}
\LA
\RefLine{Lectiones I–III: de Scriptura occurrente.}
\switchcolumn
\EN
\RefLine{Lessons I–III: from the Scripture of the day.}
\switchcolumn*

\LA
\LessonHead{Lectio IV}
\switchcolumn
\EN
\LessonHead{Lesson IV}
\switchcolumn*

\LA
\noindent Basilíus, nobilis Cáppadox, Athénis una cum Gregório Nazianzéno, ejus amicíssimo, sæculáribus lítteris, deínde in monastério sacris mirabíliter erudítus, eum brevi cursum fecit ad omnem doctrínæ et morum excelléntiam, ut inde Magni cognómen invénerit. Is ad prædicándum Jesu Christi Evangélium in Pontum accersítus, eam provínciam, a christiánis institútis aberrántem, ad viam salútis revocávit. Mox ab Eusébio Cæsaréæ epíscopo ad erudiéndam eam civitátem adjútor adhibétur: in cujus locum póstea succéssit. Is Fílium Patri consubstantiálem esse in primis deféndit, ac Valéntem imperatórem, sibi irátum, miráculis ádeo flexit, ut, incumbéntem ad voluntátem ejiciéndi ipsum in exsílium, a senténtia discédere coégerit.\par
\switchcolumn
\EN
\noindent This Basil was a noble Cappadocian v who studied earthly learning at Athens, in company with Gregory of Nazianzus, to whom he was united in a warm and tender friendship. He afterwards studied things sacred in a monastery, where he quickly attained an eminent degree of excellence in doctrine and life, whereby he gained to himself the surname of the Great. He was called to Pontus to preach the Gospel of Christ Jesus, and brought back into the way of salvation that country which before had been wandering astray from the rules of Christian discipline. He was shortly united as coadjutor to Eusebius, Bishop of Cassarea, for the edification of that city, and afterwards became his successor in the see. One of his greatest labours was to maintain that the Son is of one Substance with the Father, and when the Emperor Valens, moved to wrath against him, was willing to send him into exile, he so bent him by dint of the miracles which he worked that he forced him to forego his intention.\par
\switchcolumn*

\LA
\begin{Resp}\Rsign Invéni David servum meum, óleo sancto meo unxi eum: \Star{} Manus enim mea auxiliábitur ei.\end{Resp}
\begin{Resp}\Vsign Nihil profíciet inimícus in eo, et fílius iniquitátis non nocébit ei.\end{Resp}
\begin{Resp}\Rsign Manus enim mea auxiliábitur ei.\end{Resp}
\switchcolumn
\EN
\begin{Resp}\Rsign I have found David My servant, with My holy oil have I anointed him \Star{} For My hand shall help him.\end{Resp}
\begin{Resp}\Vsign The enemy shall prevail nothing against him, nor the son of wickedness afflict him.\end{Resp}
\begin{Resp}\Rsign For My hand shall help him.\end{Resp}
\switchcolumn*

\LA
\LessonHead{Lectio V}
\switchcolumn
\EN
\LessonHead{Lesson V}
\switchcolumn*

\LA
\noindent Nam et Valéntis sella, in qua factúrus decrétum de ejiciéndo e civitáte Basilío, sedére volébat, confrácta est. Et tribus ab eo cálamis adhíbitis ad scribéndam exsílii legem, nullus eórum réddidit atraméntum; et, cum nihilóminus in propósito scribéndi ímpium decrétum persísteret, ipsíus déxtera, dissolútis nervis, tota contrémuit. His commótus Valens chartam utráque manu conscídit. Ea autem nocte, quæ ad deliberándum Basilío data est, Valéntis uxor íntimis est cruciáta dolóribus, et únicus fílius in gravem morbum íncidit. Quibus ille pertérritus, iniquitátem suam recognóscens, Basilíum accérsit, quo præsénte, puer cœpit convaléscere; verum, vocátis a Valénte ad viséndum púerum hæréticis, paulo post móritur.\par
\switchcolumn
\EN
\noindent The chair upon which Valens sat down, in order to sign the decree of Basil's ejectment from the city, broke down under him, and three pens which he took one after the other to sign the edict of banishment, all would not write and when nevertheless he remained firm to write the ungodly order, his right hand shook. Valens was so frightened at these omens, that he tore the paper in two. During the night which was allowed to Basil to make up his mind, Valens' wife had a severe stomach-ache, and their only son was taken seriously ill. These things alarmed Valens so much that he acknowledged his wickedness, and sent for Basil, during whose visit the child began to get better. However, when Valens sent for some heretics to see it, it presently died.\par
\switchcolumn*

\LA
\begin{Resp}\Rsign Pósui adjutórium super poténtem, et exaltávi eléctum de plebe mea: \Star{} Manus enim mea auxiliábitur ei.\end{Resp}
\begin{Resp}\Vsign Invéni David servum meum, óleo sancto meo unxi eum.\end{Resp}
\begin{Resp}\Rsign Manus enim mea auxiliábitur ei.\end{Resp}
\switchcolumn
\EN
\begin{Resp}\Rsign I have laid help upon one that is mighty, and have exalted one chosen out of My people \Star{} For My hand shall help him.\end{Resp}
\begin{Resp}\Vsign I have found David My servant, with My holy oil have I anointed him.\end{Resp}
\begin{Resp}\Rsign For My hand shall help him.\end{Resp}
\switchcolumn*

\LA
\LessonHead{Lectio VI}
\switchcolumn
\EN
\LessonHead{Lesson VI}
\switchcolumn*

\LA
\noindent Abstinéntia et continéntia fuit admirábili; una túnica conténtus erat: in jejúnio servándo diligentíssimus, in oratióne assíduus, in qua sæpe totam noctem consumébat. Virginitátem perpétuo cóluit. Monastériis exstrúctis, ita monachórum institútum temperávit, ut solitáriæ atque actuósæ vitæ utilitátes præcláre simul conjúngeret. Multa erudíte scripsit; ac nemo, teste Gregório Nazianzéno, sacræ Scriptúræ libros vérius aut ubérius explicávit. Obiit Kaléndis Januárii, cum tantum spíritu vivens, præter ossa et pellem, nulla prætérea córporis parte constáre viderétur.\par
\switchcolumn
\EN
\noindent The abstinence and self-control of Basil were truly wonderful. He was content to wear nothing but one single garment. In observance of fasting he was most earnest, and so instant in prayer, that he would oftentimes pass the whole night therein. His virginity he kept always unsullied. He built monasteries, wherein he so adapted the institution of monasticism, that he exquisitely united for the inmates the advantages of the contemplative and of the active life. He was the author of many learned writings, and, according to the witness of Gregory of Nazianzus, no one hath ever composed more faithful and edifying explanations of the books of the Holy Scripture. He died upon the 1st day of January, (in the year of our Lord 379,) at which time so essentially spiritual was his life, that his body showed nothing but skin and bones.\par
\switchcolumn*

\LA
\begin{Resp}\Rsign Iste est, qui ante Deum magnas virtútes operátus est, et omnis terra doctrína ejus repléta est: \Star{} Ipse intercédat pro peccátis ómnium populórum.\end{Resp}
\begin{Resp}\Vsign Iste est, qui contémpsit vitam mundi, et pervénit ad cæléstia regna.\end{Resp}
\begin{Resp}\Rsign Ipse intercédat pro peccátis ómnium populórum.\end{Resp}
\begin{Resp}\Vsign Glória Patri, et Fílio, \Star{} et Spirítui Sancto.\end{Resp}
\begin{Resp}\Rsign Ipse intercédat pro peccátis ómnium populórum.\end{Resp}
\switchcolumn
\EN
\begin{Resp}\Rsign This is he which wrought great wonders before God, and the whole earth is full of his teaching \Star{} May he pray for all people, that their sins may be forgiven unto them\end{Resp}
\begin{Resp}\Vsign This is he which loved not his life in this world, and hath attained unto the kingdom of heaven.\end{Resp}
\begin{Resp}\Rsign May he pray for all people, that their sins may be forgiven unto them\end{Resp}
\begin{Resp}\Vsign Glory be to the Father, and to the Son, \Star{} and to the Holy Ghost.\end{Resp}
\begin{Resp}\Rsign May he pray for all people, that their sins may be forgiven unto them\end{Resp}
\switchcolumn*

\LA
\LessonHead{Lectio VII}
\switchcolumn
\EN
\LessonHead{Lesson VII}
\switchcolumn*

\LA
\LessonTitle{Léctio sancti Evangélii secúndum Lucam}
\switchcolumn
\EN
\LessonTitle{From the Holy Gospel according to Luke}
\switchcolumn*

\LA
\Cite{Luc 14:26-35}
\switchcolumn
\EN
\Cite{Luke 14:26-35}
\switchcolumn*

\LA
\noindent In illo témpore : Dixit Jesus turbis: Si quis venit ad me, et non odit patrem suum, et matrem, et uxórem, et fílios, et fratres, et soróres, adhuc autem et ánimam suam, non potest meus esse discípulus. Et réliqua.\par
\switchcolumn
\EN
\noindent At that time Jesus said unto the multitudes If any man come to Me, and hate not his father, and mother, and wife, and children, and brethren, and sisters, yea, and his own life also, he cannot be My disciple. And so on.\par
\switchcolumn*

\LA
\LessonTitle{Homilía sancti Basilíi Epíscopi}
\switchcolumn
\EN
\LessonTitle{Homily by St. Basil, Archbishop of Caesarea}
\switchcolumn*

\LA
\Source{Lib. Regul. fusius explic. ad interrog. 8}
\switchcolumn
\EN
\Source{Paraphrase of his Rule.}
\switchcolumn*

\LA
\noindent Perfécta quidem renuntiátio in eo consístit, ut id assequámur, ne ad ipsíus étiam vitæ affectiónem propénsi simus, et respónsum mortis habeámus, ut non simus fidéntes in nobis ipsis. Hujúsmodi autem renuntiátio inítium sumit ab alienatióne rerum externárum, véluti a possessiónibus, ab ináni glória, a vivéndi consuetúdine, a rerum inutílium amóre; quemádmodum étiam suo exémplo nobis ostendérunt sancti Dómini nostri discípuli, Jacóbus quidem et Joánnes, relícto patre Zebedǽo et ipsa quoque navícula, de qua omnis illórum victus rátio pendébat; Matthǽus vero, cum ab ipso telónio surréxit ac Dóminum secútus est.\par
\switchcolumn
\EN
\noindent This is perfect self-renunciation, when we attain to indifference as regards our own lives, and wring from death himself the confession that our trust is not in our own strength. The first step toward this crown of abnegation is to estrange ourselves from outward things, such as property, public reputation, habits of life, and affection for things unnecessary, whereof the immediate disciples of our Holy Lord have left us a fine example James, for instance, and John, who left their father Zebedee and the boats which were their only means of getting their daily bread or Matthew, who got up from the receipt of custom, and straightway followed the Lord.\par
\switchcolumn*

\LA
\begin{Resp}\Rsign Amávit eum Dóminus, et ornávit eum: stolam glóriæ índuit eum, \Star{} Et ad portas paradísi coronávit eum.\end{Resp}
\begin{Resp}\Vsign Induit eum Dóminus lorícam fídei, et ornávit eum.\end{Resp}
\begin{Resp}\Rsign Et ad portas paradísi coronávit eum.\end{Resp}
\switchcolumn
\EN
\begin{Resp}\Rsign The Lord loved him and beautified him He clothed him with a robe of glory \Star{} And crowned him at the gates of Paradise.\end{Resp}
\begin{Resp}\Vsign The Lord hath put on him the breast-plate of faith, and hath adorned him.\end{Resp}
\begin{Resp}\Rsign And crowned him at the gates of Paradise.\end{Resp}
\switchcolumn*

\LA
\LessonHead{Lectio VIII}
\switchcolumn
\EN
\LessonHead{Lesson VIII}
\switchcolumn*

\LA
\noindent Sed quid opus est nostris ratiónibus aut sanctórum virórum exémplis id quod dícimus confirmáre, cum ipsa Dómini verba in médium líceat afférre, iísque ipsis religiósam ac Deum timéntem ánimam commovére, quibus ille perspícue et sine controvérsia protestátur, dicens: Sic ígitur quicúmque ex vobis non renuntiáverit ómnibus quæ póssidet, non potest meus esse discípulus? Et álio in loco, cum prius dixísset: Si vis perféctus esse, vade, et vende ómnia quæ habes, et da paupéribus; póstea subjúnxit: Veni, séquere me.\par
\switchcolumn
\EN
\noindent But what need have we for our own arguments, or for the examples of holy men to confirm what we say, when we are able to cite the very words of the Lord Himself, and by them to move any earnest soul that loveth God Such were they unto whom He plainly and unhesitatingly declared Whosoever he be of you that forsaketh not all that he hath, he cannot be My disciple. And again, in another place, when He had said If thou wilt be perfect, go, and sell that thou hast, and give to the poor the completion of the sentence was, And come and follow Me. \textasciitilde{}(Matth. xix. 19.)\par
\switchcolumn*

\LA
\begin{Resp}\Rsign In médio Ecclésiæ apéruit os ejus, \Star{} Et implévit eum Dóminus spíritu sapiéntiæ et intelléctus.\end{Resp}
\begin{Resp}\Vsign Jucunditátem et exsultatiónem thesaurizávit super eum.\end{Resp}
\begin{Resp}\Rsign Et implévit eum Dóminus spíritu sapiéntiæ et intelléctus.\end{Resp}
\begin{Resp}\Vsign Glória Patri, et Fílio, \Star{} et Spirítui Sancto.\end{Resp}
\begin{Resp}\Rsign Et implévit eum Dóminus spíritu sapiéntiæ et intelléctus.\end{Resp}
\switchcolumn
\EN
\begin{Resp}\Rsign In the midst of the congregation did the Lord open his mouth. \Star{} And filled him with the spirit of wisdom and understanding.\end{Resp}
\begin{Resp}\Vsign He made him rich with joy and gladness.\end{Resp}
\begin{Resp}\Rsign And filled him with the spirit of wisdom and understanding.\end{Resp}
\begin{Resp}\Vsign Glory be to the Father, and to the Son, \Star{} and to the Holy Ghost.\end{Resp}
\begin{Resp}\Rsign And filled him with the spirit of wisdom and understanding.\end{Resp}
\switchcolumn*

\LA
\LessonHead{Lectio IX}
\switchcolumn
\EN
\LessonHead{Lesson IX}
\switchcolumn*

\LA
\noindent Est ígitur renuntiátio, quemádmodum docúimus, vinculórum terrénæ hujus ac temporális vitæ solútio, atque ab humánis negótiis liberátio, per quam ad ineúndam viam, qua ad Deum pervenítur, aptióres et promptióres effícimur; et expedíta rátio ad acquisitiónem usúmque rerum, quæ super aurum et lápidem pretiósum multum longe sunt pretiosióres. Et in summa, cordis humáni ad cæléstem conversatiónem translátio, ita ut dícere líceat: Nostra conversátio in cælis est; et (quod máximum est) inítium unde ad Christi similitúdinem evádimus, qui cum dives esset, propter nos pauper est factus.\par
\switchcolumn
\EN
\noindent This then, as we have taught, is self-renunciation to unlock the chains of this earthly life, which passeth away, and to set oneself free from the business of men, and so to make ourselves lither and meeter to enter on that path which leadeth to God, and let our reason be more unhampered to gain and to use those things which are far more precious than gold or precious stones. (Ps. xviii. 11.) In short, it is to have our heart in heaven and not on earth, so as to be able to say Our conversation is in heaven. (Phil. iii. 20.) And (which is the great thing) this is the first step towards the attaining to be like Christ, who, though He was rich, yet for our sakes He became poor. (2 Cor. viii. 9.)\par
\switchcolumn*

\LA
\TeDeum{Te Deum}
\switchcolumn
\EN
\TeDeum{Te Deum}
\switchcolumn*

\LA
\LessonHead{Lectio IX (pro commemoratione)}
\switchcolumn
\EN
\LessonHead{Lesson IX (when commemorated)}
\switchcolumn*

\LA
\LessonTitle{Commemoratio: S. Basilii Magni, Episcopis Confessoris et Ecclesiæ Doctoris}
\switchcolumn
\EN
\LessonTitle{Commemoration of: St. Basil the Great, Confessor and Doctor of the Church}
\switchcolumn*

\LA
\noindent Basilíus, nóbilis Cáppadox, Athénis una cum Gregório Nazianzéno, ejus amicíssimo, sæculáribus lítteris, deínde in monastério sacris mirabíliter erudítus, eum brevi cursum fecit ad omnem doctrínæ et morum excelléntiam, ut inde Magni cognómen invénerit. Ad prædicándum Jesu Christi Evangélium in Pontum accersítus, eam provínciam ad viam salútis revocávit; mox ab Eusébio Cæsaréæ epíscopo ad erudiéndam eam civitátem adjútor adhibétur, in cujus locum póstea succéssit. Is Fílium Patri consubstantiálem esse in primis deféndit, ac Valéntem imperatórem, sibi irátum et exsílium minitántem, miráculis ádeo flexit, ut a senténtia discédere coégerit. Abstinéntia et continéntia fuit admirábili; in oratióne assíduus, in ea sæpe totam noctem consumébat. Monastériis exstrúctis, ita monachórum institútum temperávit, ut solitáriæ atque actuósæ vitæ utilitátis præcláre simul conjúngeret. Multa erudíte scripsit; ac nemo, teste Gregório Nazianzéno, sacræ Scriptúræ libros vérius aut ubérius explicávit. Obiit kaléndis januárii.\par
\switchcolumn
\EN
\noindent Basil, a Cappodocian nobleman, studied profane letters at Athens together with his close friend, Gregory of Nazianzus, and took his sacred studies in a monastery. Becoming marvelously proficient in both, he soon attained such excellence in learning and in his way of life that from then on he was given the name of The Great. Summoned to preach the Gospel of Christ Jesus in Pontus, he called that province back to the way of salvation. Soon he was asked by Eusebius, Bishop of Caesarea, to aid him in teaching; and he succeeded Eusebius as bishop. Basil was among the first to defend the consubstantiality of the Son with the Father; and by his miracles he caused Emperor Valens, who was angry with him and threatening him with exile, to give up any such intentions. Basil's abstinence and continence were marvelled at; and he was constant in prayer, often spending the whole night in it. He built monasteries, ordering the monastic life so that it would best combine the advantages of the solitary life with those of the active life. He wrote many learned books; and, as Gregory of Nazianzus testifieth, no one hath explained the books of Holy Scripture more truly and fruitfully. He died on the 1st day of January.\par
\switchcolumn*

\LA
\TeDeum{Te Deum}
\switchcolumn
\EN
\TeDeum{Te Deum}
\switchcolumn*

\end{paracol}

\par\needspace{10\baselineskip}\addvspace{1.2em}{\centering\small\textsc{Die 15 Junii} · \textsc{15 June}\par}
\Office{Ss. Viti, Modesti atque Crescentiæ Martyrum}{Sts. Vitus, Modestus and Crescentia, Martyrs}{Simplex}
\addcontentsline{toc}{subsection}{Die 15 Junii · Ss. Viti, Modesti atque Crescentiæ Martyrum \textit{— Sts. Vitus, Modestus and Crescentia, Martyrs}}
\begin{paracol}{2}
\LA
\RefLine{Lectiones I–II: de Scriptura occurrente.}
\switchcolumn
\EN
\RefLine{Lessons I–II: from the Scripture of the day.}
\switchcolumn*

\LA
\LessonHead{Lectio III (pro commemoratione)}
\switchcolumn
\EN
\LessonHead{Lesson III (when commemorated)}
\switchcolumn*

\LA
\noindent Vitus admodum puer inscio patre baptizatus est: quod cum ille rescivisset, nihil prætermisit, quo filium a christiana religione removeret. Qua in voluntate permanentem, Valeriano judici verberibus castigandum tradidit. Sed nihilominus in sententia persistens, patri redditus est. Sed dum eum pater gravius punire cogitat, Vitus Angeli monitu, comitibus Modesto et Crescentia, ejus educatoribus, migrat in alienas terras: ibique eam sanctitatis laudem adeptus est, ut ejus fama ad Diocletianum perlata, ipsum imperator accerseret, ut filium suum a dæmone vexatum liberaret: Quo liberato, cum ei amplissimis præmiis ingratus imperator, ut deos coleret, persuadere non potuisset, una cum Modesto et Crescentia vinculis constrictum mittit in carcerem. Quos ubi constantiores esse comperit, demitti jubet in ingens vas liquato plumbo, ferventi resina ac pice plenum: in quo cum trium Hebræorum puerorum more divinos hymnos canerent, inde erepti leoni objiciuntur: qui prostérnens se, eorum pedes lambebat. Quare inflammatus ira imperator, quod multitudinem videbat miraculo commoveri, eos in catasta sterni jubet et ita cædi eorum membra, atque ossa divelli. Quo tempore tonitrua, fulgura, magnique terræmotus fuere, quibus templa deorum corruerunt, et multi oppressi sunt. Eorum reliquias Florentia nobilis femina, unguentis conditas, honorifice sepelivit.\par
\switchcolumn
\EN
\noindent This Vitus was a child who was baptized without his father's knowledge. When his father had found it out, he used his best endeavours to dissuade his son from the Christian religion, but as he found him persistent in it, he handed him over to Valerian the judge to be whipped. But as he still remained as unshaken as before, he was given back to his father. But while his father was turning over in his mind to what severe discipline to subject him, Vitus, being warned by an Angel, fled out of the country, in company with his foster-parents Modestus and Crescentia. In his new home he gained great praise for holiness, so that the fame of it came to Diocletian, which Emperor sent for him to deliver his own child which was vexed with a devil. Him Vitus delivered, but when the Emperor found that with all his great gifts he could not bring him to worship the gods, he had the ingratitude to cast him and Modestus and Crescentia into prison, binding them in fetters. But when they were found in their prison more faithful than ever to their confession, the Emperor commanded them to be thrown into a great vessel full of melted lead, resin, and pitch. Therein these three, like the three Holy Children in the burning fiery furnace, sang praise to God and upon that they were haled forth and cast to a lion, but he lay down before them, and licked their feet. Then the Emperor, being filled with fury, more especially because he saw that the multitude that looked on were stirred up at the miracle, commanded Vitus, Modestus, and Crescentia to be stretched upon a block and their limbs crushed, and their bones rent one from the other. While as they were dying there came great thunderings, and lightnings, and earthquakes, so that temples of the gods fell down, and many men were killed. As for that which remained of the Martyrs, the noble lady Florence took it, and embalmed it with spices, and honourably buried it.\par
\switchcolumn*

\LA
\TeDeum{Te Deum}
\switchcolumn
\EN
\TeDeum{Te Deum}
\switchcolumn*

\end{paracol}

\par\needspace{10\baselineskip}\addvspace{1.2em}{\centering\small\textsc{Die 18 Junii} · \textsc{18 June}\par}
\Office{S. Ephræm Syri Confessoris et Ecclesiæ Doctoris}{St. Ephraem of Syria, Confessor and Doctor of the Church}{Duplex}
\addcontentsline{toc}{subsection}{Die 18 Junii · S. Ephræm Syri Confessoris et Ecclesiæ Doctoris \textit{— St. Ephraem of Syria, Confessor and Doctor of the Church}}
\begin{paracol}{2}
\LA
\RefLine{Lectiones I–III: de Scriptura occurrente.}
\switchcolumn
\EN
\RefLine{Lessons I–III: from the Scripture of the day.}
\switchcolumn*

\LA
\RefLine{Lectio IX: de SS. Marci et Marcelliani Martyrum (06-18t).}
\switchcolumn
\EN
\RefLine{Lesson IX: of Sts. Mark and Marcellian, Martyrs (06-18t).}
\switchcolumn*

\LA
\LessonHead{Lectio IV}
\switchcolumn
\EN
\LessonHead{Lesson IV}
\switchcolumn*

\LA
\noindent Ephræm genere Syrus, Nisibeno patre natus est. Adhuc juvenis ad sanctum Jacobum episcopum se contulit, a quo baptizatus, brevi ita sanctitate et doctrina profecit, ut in schola Nisibi, Mesopotamiæ urbe, florente magister fuerit constitutus. Post Jacobi episcopi mortem, Nisibi a Persis capta, Edessam profectus est; ubi primum in monte inter monachos consedit, deínde, ut plurimos ad se confluentes homines vitaret, vitam duxit eremiticam. Edessenæ ecclesiæ diaconus ordinatus, et ob humilitatem sacerdotium recusans, omnium virtutum splendore enituit, et pietatem et religionem vera sapientiæ professione sibi comparare sategit. Spem omnem in solo Deo defixam habens, quævis humana ac transitoria contemnens, divina ac sempiterna assidue concupiscebat.\par
\switchcolumn
\EN
\noindent Ephraem was of Syrian descent, and son of a citizen of Nisibis. While yet a young man he went to the holy bishop James, by whom he was baptized, and he soon made such progress in holiness and learning as to be appointed master of a flourishing school at Nisibis, a city of Mesopotamia. After the death of the bishop James, Nisibis was captured by the Persians, and Ephraem went to Edessa. Here he settled first on the mountain among the monks, and then, that he might avoid the great numbers of men who flocked to him, he adopted the eremitical life. He was ordained deacon of the Church of Edessa, but refused the priesthood out of humility. He was conspicuous with the splendour of every virtue and strove to acquire piety and religion by professing true wisdom. He placed all his hope in God alone, despised all human and transitory things, and always longed for the divine and eternal.\par
\switchcolumn*

\LA
\begin{Resp}\Rsign Honéstum fecit illum Dóminus, et custodívit eum ab inimícis, et a seductóribus tutávit illum: \Star{} Et dedit illi claritátem ætérnam.\end{Resp}
\begin{Resp}\Vsign Justum dedúxit Dóminus per vias rectas, et osténdit illi regnum Dei.\end{Resp}
\begin{Resp}\Rsign Et dedit illi claritátem ætérnam.\end{Resp}
\switchcolumn
\EN
\begin{Resp}\Rsign The Lord made him honourable, and defended him from his enemies, and kept him safe from those that lay in wait for him; \Star{} And gave him perpetual glory.\end{Resp}
\begin{Resp}\Vsign He went down with him into the pit, and left him not in bonds.\end{Resp}
\begin{Resp}\Rsign And gave him perpetual glory.\end{Resp}
\switchcolumn*

\LA
\LessonHead{Lectio V}
\switchcolumn
\EN
\LessonHead{Lesson V}
\switchcolumn*

\LA
\noindent Cæsaream Cappadociæ, divino ductus Spiritu, cum petiisset, ipsum ibi os Ecclesiæ Basilium vidit, et uterque mutua consuetudine opportunum in modum usus est. Ad innumeros errores refellendos, qui, tunc temporis grassantes, Ecclesiam Dei divexabant, atque ad mysteria Domini nostri Jesu Christi sedulo illustranda, plurimas edidit lucubrationes, Syro sermone compositas, et fere omnes in linguam Græcam versas; atque, teste sancto Hieronymo, ipse ad tantum venit claritudinem, ut, post lectionem Scripturarum, publice in quibusdam ecclesiis ejus scripta recitarentur.\par
\switchcolumn
\EN
\noindent When, led by the Spirit of God, he went to Caesarea in Cappadocia, there he saw Basil, that mouthpiece of the Church, and both enjoyed mutual companionship in a suitable manner. In order to refute the countless errors which were rife at that time, and which were troubling the Church of God, and in order to expound zealously the divine mysteries of our Lord Jesus Christ, he wrote many studies in Syrian, almost all of which have been translated into Greek. St. Jerome beareth witness that he attained such fame, that his writings were read publicly in certain churches after the reading from the Scriptures.\par
\switchcolumn*

\LA
\begin{Resp}\Rsign Amávit eum Dóminus, et ornávit eum: stolam glóriæ índuit eum, \Star{} Et ad portas paradísi coronávit eum.\end{Resp}
\begin{Resp}\Vsign Índuit eum Dóminus lorícam fídei, et ornávit eum.\end{Resp}
\begin{Resp}\Rsign Et ad portas paradísi coronávit eum.\end{Resp}
\switchcolumn
\EN
\begin{Resp}\Rsign The Lord loved him and beautified him; He clothed him with a robe of glory, \Star{} And crowned him at the gates of Paradise.\end{Resp}
\begin{Resp}\Vsign The Lord hath put on him the breast-plate of faith, and hath adorned him.\end{Resp}
\begin{Resp}\Rsign And crowned him at the gates of Paradise.\end{Resp}
\switchcolumn*

\LA
\LessonHead{Lectio VI}
\switchcolumn
\EN
\LessonHead{Lesson VI}
\switchcolumn*

\LA
\noindent Universa illius opera, tam splendido doctrinæ lumine referta effecerunt, ut idem Sanctus, adhuc vivens, tamquam Ecclesiæ Doctor, magno honore habitus fuerit. Metrica quoque cantica composuit in laudem beatissimæ Virginis Mariæ ac Sanctorum: quam ob causam a Syris Spiritus Sancti cithara merito fuit appellatus. In mirifica ac pia devotione erga eamdem Virginem immaculatam primum excelluit. Meritis plenus, Edessæ, in Mesopotamia, decimo quarto Kalendas Julii, decessit sub Valente principe: eumque, instántibus plúribus Sanctæ Románæ Ecclésiæ Cardinálibus, Patriárchis, Archiepíscopis, Epíscopis, Abbátibus, et religiósis famíliis, Benedíctus Papa décimus quintus, ex Sacrórum Rítuum Congregatiónis consúlto, universális Ecclésiæ Doctórem declarávit.\par
\switchcolumn
\EN
\noindent His works taken as a whole, so infused with the bright light of learning, brought it about that this holy man, while yet alive, was held in great honour, and was even considered a Doctor of the Church. He also composed songs in verse, in honour of the Most Blessed Virgin Mary, and of the Saints, and for this reason he was appropriately named by the Syrians the Harp of the Holy Ghost. He was noted for his great and tender devotion towards the Immaculate Virgin. He died, rich in merits, at Edessa in Mesopotamia on the 18th day of June in the reign of Valens. Pope Benedict XV, at the instance of many Cardinals of the Holy Roman Church, Patriarchs, Archbishops, Bishops, Abbots, and religious communities, declared him by a decree of the Congregation of Sacred Rites to be a Doctor of the universal Church.\par
\switchcolumn*

\LA
\begin{Resp}\Rsign Iste homo perfécit ómnia quæ locútus est ei Deus, et dixit ad eum: Ingrédere in réquiem meam: \Star{} Quia te vidi justum coram me ex ómnibus géntibus.\end{Resp}
\begin{Resp}\Vsign Iste est, qui contémpsit vitam mundi, et pervénit ad cæléstia regna.\end{Resp}
\begin{Resp}\Rsign Quia te vidi justum coram me ex ómnibus géntibus.\end{Resp}
\begin{Resp}\Vsign Glória Patri, et Fílio, \Star{} et Spirítui Sancto.\end{Resp}
\begin{Resp}\Rsign Quia te vidi justum coram me ex ómnibus géntibus.\end{Resp}
\switchcolumn
\EN
\begin{Resp}\Rsign This is he which did according unto all that God commanded him; and God said unto him: Enter thou into My rest, \Star{} For thee have I seen righteous before Me among all people.\end{Resp}
\begin{Resp}\Vsign This is he which loved not his life in this world, and is come unto an everlasting kingdom.\end{Resp}
\begin{Resp}\Rsign For thee have I seen righteous before Me among all people.\end{Resp}
\begin{Resp}\Vsign Glory be to the Father, and to the Son, \Star{} and to the Holy Ghost.\end{Resp}
\begin{Resp}\Rsign For thee have I seen righteous before Me among all people.\end{Resp}
\switchcolumn*

\LA
\LessonHead{Lectio VII}
\switchcolumn
\EN
\LessonHead{Lesson VII}
\switchcolumn*

\LA
\LessonTitle{Léctio sancti Evangélii secúndum Matthǽum.}
\switchcolumn
\EN
\LessonTitle{From the Holy Gospel according to Matthew}
\switchcolumn*

\LA
\Cite{Matt 5:13-19}
\switchcolumn
\EN
\Cite{Matt 5:13-19}
\switchcolumn*

\LA
\noindent In illo témpore: Dixit Jesus discípulis suis: Vos estis sal terræ. Quod si sal evanúerit, in quo saliétur? Et réliqua.\par
\switchcolumn
\EN
\noindent At that time: Jesus said unto His disciples: Ye are the salt of the earth. But if the salt have lost his savour, wherewith shall it be salted. And so on.\par
\switchcolumn*

\LA
\LessonTitle{Homilia sancti Ephræm Syri, Diaconi}
\switchcolumn
\EN
\LessonTitle{Homily of St. Ephrem of Syria Deacon}
\switchcolumn*

\LA
\Source{Sermo de vita et exercitatione monastica}
\switchcolumn
\EN
\Source{Sermo de vita et exercitatione monastica}
\switchcolumn*

\LA
\noindent Præclarum est bonum inchoare atque perficere, et gratum Deo esse et utilem proximo, ipsique summo ac dulcissimo rectori nostro Christo Jesu placere, qui ait: Vos estis sal terræ, et columna cælorum. Labor afflictionis tuæ, dilectissime, tamquam somnus est; porro laboris requies inenarrabilis atque inæstimabilis. Attende ergo tibi ipsi sollicite, ne utrumque pariter amittas, dum neutrum plene persequeris, præsentem scilicet sempiternamque lætitiam. Stude potius perfectam virtutem consequi, ornatam atque insignitam omnibus quæ diligit Deus. Hanc si assequaris, numquam irritabis Deum, neque proximum tuum violabis.\par
\switchcolumn
\EN
\noindent Clearly it is a good thing to begin and to accomplish, to be acceptable to God and useful to one's neighbour, to please our most high and gracious ruler, Christ Jesus, who saith: Ye are the salt of the earth and the pillars of the heavens. The burden of your affliction, dearly beloved, is as a sleep; then there followeth unspeakable and priceless rest from labour. Therefore, watch thyself carefully, so that while thou followest after neither, wholeheartedly, thou shouldst not lose both the present and the eternal joy. Study rather to attain to the perfect virtue, adorned and stamped by all that God loveth. If thou dost strive after this, thou wilt never either anger God nor injure thy neighbour.\par
\switchcolumn*

\LA
\begin{Resp}\Rsign Iste est qui ante Deum magnas virtútes operátus est, et de omni corde suo laudávit Dóminum: \Star{} Ipse intercédat pro peccátis ómnium populórum.\end{Resp}
\begin{Resp}\Vsign Ecce homo sine queréla, verus Dei cultor, ábstinens se ab omni ópere malo, et pérmanens in innocéntia sua.\end{Resp}
\begin{Resp}\Rsign Ipse intercédat pro peccátis ómnium populórum.\end{Resp}
\switchcolumn
\EN
\begin{Resp}\Rsign This is he which wrought great wonders before God, and praised the Lord with all his heart. \Star{} May he pray for all people, that their sins may be forgiven unto them.\end{Resp}
\begin{Resp}\Vsign Behold a man without blame, a worshipper of God in truth, keeping himself clean from every evil work, and abiding still in his innocency.\end{Resp}
\begin{Resp}\Rsign May he pray for all people, that their sins may be forgiven unto them.\end{Resp}
\switchcolumn*

\LA
\LessonHead{Lectio VIII}
\switchcolumn
\EN
\LessonHead{Lesson VIII}
\switchcolumn*

\LA
\noindent Porro virtus ista, unica uniusque speciei dicitur, variarum virtutum in se ipsa habens pulchritudinem. Diadema regium absque pretiosis lapidibus candentibusque margaritis connecti texique non potest; ita et hæc unica virtus sine variarum fulgore virtutum constare nequit. Est enim profecto simillima diademati regio. Nam, ut illi, si lapis unus aut margarita defuerit, in regio capite lucere pleniter nequit; ita et hæc unica virtus, nisi virtutum ceterarum honore conseritur, perfecta virtus non appellatur. Similis item est pretiosissimis epulis, exquisitissimis condimentis præparatis, sed sale carentibus. Sicut enim pretiosi illi cibi sine sale comedi nequeunt; ita et ista virtus uniformis, si variarum virtutum gloria et honore decoretur, absit autem Dei proximique dilectio, vilis prorsus atque contemptibilis est.\par
\switchcolumn
\EN
\noindent Moreover this virtue is called special and singular, having within itself the beauty of different virtues. We cannot have a royal diadem without precious stones and gleaming pearls arranged and fitted together; so likewise this single virtue cannot remain without the splendour of other different virtues. For it is, indeed, most like a kingly crown. For as in the latter case, if one stone or one pearl be missing it cannot shine perfectly upon the royal head; so, then, this special virtue cannot be called a perfect virtue, unless it is worthily connected with other virtues. Again, it is like unto very rich food, furnished with exquisite seasonings, but lacking salt. For as those rich dishes cannot be eaten without salt, so this simple virtue may be adorned with the glory and honour of different virtues, but if a man lack the love of God and of his neighbour, he is wholly worthless and contemptible.\par
\switchcolumn*

\LA
\begin{Resp}\Rsign In médio Ecclésiæ apéruit os ejus, \Star{} Et implévit eum Dóminus spíritu sapiéntiæ et intelléctus.\end{Resp}
\begin{Resp}\Vsign Jucunditátem et exsultatiónem thesaurizávit super eum.\end{Resp}
\begin{Resp}\Rsign Et implévit eum Dóminus spíritu sapiéntiæ et intelléctus.\end{Resp}
\begin{Resp}\Vsign Glória Patri, et Fílio, \Star{} et Spirítui Sancto.\end{Resp}
\begin{Resp}\Rsign Et implévit eum Dóminus spíritu sapiéntiæ et intelléctus.\end{Resp}
\switchcolumn
\EN
\begin{Resp}\Rsign In the midst of the congregation did the Lord open his mouth. \Star{} And filled him with the spirit of wisdom and understanding.\end{Resp}
\begin{Resp}\Vsign He made him rich with joy and gladness.\end{Resp}
\begin{Resp}\Rsign And filled him with the spirit of wisdom and understanding.\end{Resp}
\begin{Resp}\Vsign Glory be to the Father, and to the Son, \Star{} and to the Holy Ghost.\end{Resp}
\begin{Resp}\Rsign And filled him with the spirit of wisdom and understanding.\end{Resp}
\switchcolumn*

\LA
\LessonHead{Lectio IX (pro commemoratione)}
\switchcolumn
\EN
\LessonHead{Lesson IX (when commemorated)}
\switchcolumn*

\LA
\LessonTitle{Commemoratio: S. Ephræm Syri Confessoris et Ecclesiæ Doctoris}
\switchcolumn
\EN
\LessonTitle{Commemoration of: St. Ephraem of Syria, Confessor and Doctor of the Church}
\switchcolumn*

\LA
\noindent Ephræm, génere Syrus, Nisibéno patre natus est. Adhuc júvenis ad sanctum Jacóbum epíscopum se cóntulit, a quo baptizátus, brevi ita sanctitáte et doctrína profécit, ut in schola Nísibi, Mesopotámiæ urbe, florénte magíster fúerit constitútus. Edessénæ ecclésiæ diáconus ordinátus et ob humilitátem sacerdótium recúsans, ómnium virtútum splendóre enítuit, et pietátem et religiónem vera sapiéntiæ professióne sibi comparáre satégit. Univérsa illíus ópera, tam spléndido doctrínæ lúmine reférta, effecérunt, ut idem Sanctus, adhuc vivens, tamquam Ecclésiæ Doctor, magno honóre hábitus fúerit. In mirífica ac pia devotióne erga Vírginem Immaculátam primum excélluit. Méritis plenus, Edéssæ in Mesopotámia, décimo quarto kaléndas júlii, decéssit sub Valénte príncipe: eúmque Benedíctus Papa décimus quintus, ex Sacrórum Rítuum Congregatiónis consúlto, universális Ecclésiæ Doctórem declarávit.\par
\switchcolumn
\EN
\noindent Ephraem was of Syrian stock, his father being a citizen of Nisibis. While he was still young, he went to the holy bishop James to be baptized. In a short time he advanced so much in holiness and learning that he was made master of a flourishing school at Nisibis, a city in Mesopotamia. Ordained deacon of the Church of Edessa, and refusing the priesthood out of humility, he shone with the splendour of all virtues, and sought to acquire devotion and religion by the profession of true wisdom. All his works, illuminated with the bright light of learning, caused this Saint to be treated with great honour as a Doctor of the Church even in his lifetime. He excelled, above all, in a wonderful and loving devotion to the Immaculate Virgin. Rich in merits, he died on the 18th day of June at Edessa in Mesopotamia, under the Emperor Valens. Pope Benedict XV, after consulting the Congregation of Sacred Rites, declared him a Doctor of the Universal Church.\par
\switchcolumn*

\LA
\TeDeum{Te Deum}
\switchcolumn
\EN
\TeDeum{Te Deum}
\switchcolumn*

\end{paracol}

\par\needspace{10\baselineskip}\addvspace{1.2em}{\centering\small\textsc{Die 18 Junii} · \textsc{18 June}\par}
\Office{SS. Marci et Marcelliani Martyrum}{Sts. Mark and Marcellian, Martyrs}{Simplex}
\addcontentsline{toc}{subsection}{Die 18 Junii · SS. Marci et Marcelliani Martyrum \textit{— Sts. Mark and Marcellian, Martyrs}}
\begin{paracol}{2}
\LA
\LessonHead{Lectio IX (pro commemoratione)}
\switchcolumn
\EN
\LessonHead{Lesson IX (when commemorated)}
\switchcolumn*

\LA
\LessonTitle{Commemoratio: SS. Marci et Marcelliani Martyrum}
\switchcolumn
\EN
\LessonTitle{Commemoration of: Sts. Mark and Marcellian, Martyrs}
\switchcolumn*

\LA
\noindent Marcus et Marcellianus fratres Romani, propter christianam fidem a Fabiano duce comprehensi, ad stipitem alligati sunt, pedibus clavis confixis. Ad quos cum ita loqueretur judex: Resipiscite miseri, et vos ipsos ab his cruciatibus eripite: responderunt: Numquam tam jucunde epulati sumus, quam hæc libenter Jesu Christi causa perferimus, in cujus amóre nunc fixi esse cœpimus: utinam tamdiu nos hæc pati sinat, quamdiu hoc corruptibili corpore vestiti erimus. Qui diem noctemque in tormentis divinas laudes canentes, denique telis transfixi, ad martyrii gloriam pervenerunt. Quorum corpora via Ardeatina sepulta sunt.\par
\switchcolumn
\EN
\noindent Mark and Marcellian were two brothers, Romans, who were arrested by Duke Fabian for believing in Christ, and fastened to a beam, to which their feet were nailed. The Judge said to them Wretched creatures, do think for a moment, and free yourselves from such suffering. But they answered him We have never enjoyed any dinner so much as we do what we are now undergoing here for Jesus Christ's sake. We have got ourselves a little fast to His love now. Would that He would let us suffer this as long as we are clad in this corruptible body. Still suffering, they for a day and a night sang the praises of God continually, and in the end were thrust through with darts, and so attained the glory of Martyrdom. Their bodies are buried upon the Way to Ardea.\par
\switchcolumn*

\LA
\TeDeum{Te Deum}
\switchcolumn
\EN
\TeDeum{Te Deum}
\switchcolumn*

\end{paracol}

\par\needspace{10\baselineskip}\addvspace{1.2em}{\centering\small\textsc{Die 19 Junii} · \textsc{19 June}\par}
\Office{S. Julianæ de Falconeriis Virginis}{St. Juliana Falconeri, Virgin}{Duplex}
\addcontentsline{toc}{subsection}{Die 19 Junii · S. Julianæ de Falconeriis Virginis \textit{— St. Juliana Falconeri, Virgin}}
\begin{paracol}{2}
\LA
\RefLine{Lectiones I–III: de Scriptura occurrente.}
\switchcolumn
\EN
\RefLine{Lessons I–III: from the Scripture of the day.}
\switchcolumn*

\LA
\RefLine{Lectiones VII–VIII: ex Commúni Virginum tantum (C6a).}
\switchcolumn
\EN
\RefLine{Lessons VII–VIII: from the Common of a Virgin not a Martyr (C6a).}
\switchcolumn*

\LA
\RefLine{Lectio IX: de Ss. Gervasii et Protasii Martyrum (06-19o).}
\switchcolumn
\EN
\RefLine{Lesson IX: of Ss. Gervase and Protase, Martyrs (06-19o).}
\switchcolumn*

\LA
\LessonHead{Lectio IV}
\switchcolumn
\EN
\LessonHead{Lesson IV}
\switchcolumn*

\LA
\noindent Juliana, ex nobili Falconéria família, Claríssimo patre, qui templum Deiparæ ab Angelo salutátæ ære suo magnifice a fundamentis Floréntiæ, ut nunc visitur, eréxit, matre Reguardata, ambóbus jam senescéntibus ac ad id tempus sterílibus, nata est anno millesimo ducentésimo septuagesimo. Ab incunabulis non éxiguum futuræ sanctitátis specimen dedit ; vagiéntibus quippe labris, suavíssima Jesu et Maríæ nómina ultro proferre audíta est. Puerítiam postmodum ingressa totam se christiánis virtútibus mancipávit, in quibus adeo excelluit, ut beátus Aléxius patruus, cujus institútis ac exemplis instruebátur, matri dicere non dubitaverit, ipsam non feminam peperisse, sed angelum. Nam ita modesto vultu animoque ab omni vel levíssima erroris macula pura fuit, ut óculos vel levíssima, erroris macula pura fuit, ut óculos numquam in toto vitæ cursu ad hóminis fáciem intuéndam erexerit, auditoque peccáti vocabulo contremúerit. Expleto nondum décimo quinto ætátis suæ anno, re familiari, licet opulenta, terrenisque posthabitis nuptiis, Deo virginitátem in mánibus divi Philippi Benitii solemniter vovit, ab eoque ómnium prima religiósum Mantellatárum habitum, ut dicunt, sumpsit.\par
\switchcolumn
\EN
\noindent Juliana was a daughter of the noble family of the Falconieri, and was born in the year 1270. Her father was the same who at his own costs so splendidly built from the foundations the Church of Our Lady of the Annunciation as it now standeth at Florence. Her mother's name was Reguardata. They were both well stricken in years, and, until the birth of Juliana, had been childless. From her very cradle she gave tokens of the holiness of life to which she afterwards attained. And from the murmuring of her baby lips was caught the sweet sound of the names of Jesus and Mary. As she entered on her girlhood, she delivered herself up entirely to the pursuit of Christian godliness, and so excellently shone therein, that her uncle, the Blessed Alexius, scrupled not to tell her mother that she had given birth to an Angel rather than to a woman. So modest was her carriage, and so clean her soul from the lightest speck of indiscretion, that she never in her whole life stared a man in the face, and that the very mention of sin made her shiver, and when the story of a grievous crime was told her, she dropped down nearly fainting. Before she had finished her fifteenth year, she renounced her inheritance, although a rich one, and all prospect of an earthly marriage, and made to God a vow of virginity, before holy Philip Benizi, from whom she was the first to receive the religious habit of what are called the mantled nuns.\par
\switchcolumn*

\LA
\begin{Resp}\Rsign Propter veritátem, et mansuetúdinem, et justítiam: \Star{} Et dedúcet te mirabíliter déxtera tua.\end{Resp}
\begin{Resp}\Vsign Spécie tua et pulchritúdine tua inténde, próspere procéde, et regna.\end{Resp}
\begin{Resp}\Rsign Et dedúcet te mirabíliter déxtera tua.\end{Resp}
\switchcolumn
\EN
\begin{Resp}\Rsign Because of truth, and meekness, and righteousness; \Star{} And thy right hand shall lead thee wonderfully.\end{Resp}
\begin{Resp}\Vsign In thy comeliness, and thy beauty, go forward, fare prosperously, and reign.\end{Resp}
\begin{Resp}\Rsign And thy right hand shall lead thee wonderfully.\end{Resp}
\switchcolumn*

\LA
\LessonHead{Lectio V}
\switchcolumn
\EN
\LessonHead{Lesson V}
\switchcolumn*

\LA
\noindent Julianæ exemplum secutæ sunt plurimæ ex nobilióribus familiis feminæ, ac mater ipsa fíliæ sese religiose instituéndam dedit, ita ut, aucto paulatim numero, ordinem Mantellatárum institúerit, ac illi pie vivéndi leges summa prudéntia ac sanctitáte tradíderit. Ejus virtútes cum optime perspectas divus Benitius haberet, morti próximus nulli mélius quam Julianæ, non feminas tantum, sed et totum Servórum ordinem, cujus propagator et moderátor exstiterat, commendátum vóluit. Verum ipsa demisse semper de se cogitábat ; et, cum ceterárum esset magistra, in re quaque domestica, licet vili, soróribus famulabátur. Assiduitate orándi íntegras insumebat dies in éxtasim sæpíssime rapta, et, si reliquum, in sedandis civium dissidiis, criminosis a via iniquitátis retrahendis, ac inserviéndis impendebat ægrotis ; quorum quandoque saniem ex ulcéribus manántem, admoto ore lambens, eos sanitáti restituébat. Corpus suum flagris, nodosis funículis, férreis cingulis, vigiliis, humi nudæ cubando, térere sólita fuit. Parcíssimo cibo, et hoc vili, quatuor hebdomadæ diebus, et reliquis duobus solo Angelórum pane contenta, excepto die sábbati, quo pane solo et aqua nutriebátur.\par
\switchcolumn
\EN
\noindent She put herself for instruction under her daughter. Thus in a little while their number increased, and she became the foundress of the Order of Mantled nuns, to whom she gave a rule of life full of wisdom and godliness. Holy Philip Benizi, having thorough knowledge of her excellence, chose her above all living to whom at his death to leave the care not of the women only but of the whole of the Order of the Servants of the Blessed Virgin Mary, of which he had been the propagator and director. Juliana, who deemed ever lowly of herself, even when she was the mistress of the others, ministered to her sisters in the meanest offices of the work of the house. She passed whole days in incessant prayer, and was often rapt in spirit, and the remainder of her time she toiled to make peace among the citizens, who were at variance together, to recall transgressors from the ways of iniquity, and to nurse the sick, to cure whom she would sometimes even use her tongue to remove the matter that ran from their sores. It was her custom to afflict her own body with whips, knotted cords, iron girdles, watching, and sleeping upon the ground. Upon Mondays, Tuesdays, Wednesdays, and Thursdays, she ate very sparingly some unpalatable food, upon Fridays she took nothing except the Bread of Angels, and upon Saturdays, besides the Holy Communion, only bread and water.\par
\switchcolumn*

\LA
\begin{Resp}\Rsign Dilexísti justítiam, et odísti iniquitátem: \Star{} Proptérea unxit te Deus, Deus tuus, óleo lætítiæ.\end{Resp}
\begin{Resp}\Vsign Propter veritátem, et mansuetúdinem, et justítiam.\end{Resp}
\begin{Resp}\Rsign Proptérea unxit te Deus, Deus tuus, óleo lætítiæ.\end{Resp}
\switchcolumn
\EN
\begin{Resp}\Rsign Thou hast loved righteousness, and hated iniquity; \Star{} Therefore God, thy God, hath anointed thee with the oil of gladness.\end{Resp}
\begin{Resp}\Vsign Because of truth, and meekness, and righteousness.\end{Resp}
\begin{Resp}\Rsign Therefore God, thy God, hath anointed thee with the oil of gladness.\end{Resp}
\switchcolumn*

\LA
\LessonHead{Lectio VI}
\switchcolumn
\EN
\LessonHead{Lesson VI}
\switchcolumn*

\LA
\noindent Dura hujusmodi vivéndi ratióne in stómachi morbum íncidit, quo ingravescente, cum septuagesimum ætátis annum ágeret, ad extremum vitæ spatium redacta est. Diuturnæ valetúdinis incómmoda hílari vultu constantique animo pértulit ; de uno tantum cónqueri audíta est, quod, cum cibum cápere ac retinére nullo modo posset, ab Eucharistica mensa ob sacraménti reveréntiam arcerétur. Verum, his in angustiis constituta, sacerdotem rogávit ut allátum divinum Panem, quem ore sumere nequíbat, péctori saltem exterius admovéret. Precibus illíus morem gessit sacerdos ; et mirum! eodem témporis momento divinus Panis dispáruit, et Juliana sereno ac ridénti vultu exspirávit. Res supra fidem tamdiu fuit, donec virgineum de more curarétur corpus ; invénta enim est circa sinistrum péctoris latus carni véluti sigillo impressa forma hostiæ, quæ Christi crucifixi effigiem repræsentábat. Hujus prodigii fama ceterorúmque miraculórum non Floréntiæ tantum, sed totius christiáni orbis veneratiónem illi conciliávit, ac per quatuor prope íntegra sæcula adeo aucta est, ut tandem Benedíctus Papa décimus tertius in ejus celebritate Offícium proprium recitari ab universo ordine beátæ Maríæ Vírginis Servórum jússerit. Clemens vero duodecimus, munificentíssimus ejusdem ordinis protéctor, novis in dies miraculis coruscántem, sanctárum Vírginum catálogo adscripsit.\par
\switchcolumn
\EN
\noindent The self-inflicted hardships of her life brought upon her a disease of the stomach, whereby, when she was seventy years of age, she was brought to the point of death. She bore the daily sufferings of her illness with a smiling face and a brave heart. The only thing of which she was heard to complain was that, her stomach being so weak that she could not keep down any food, she was withheld by reverence for the Sacrament from drawing near to the Lord's Table. Finding herself in these straits she begged the Priest to bring the Bread of God, and, as she dared not take It into her mouth, to put It as near as possible to her heart. The Priest did as she wished, and, to the amazement of all present, the Divine Bread at once disappeared from sight, and at the same instant a smile of joyous peace crossed the face of Juliana, and she gave up the ghost. All were confounded until the virgin body was being laid out after death in the accustomed manner. Then there was found upon the left side of the bosom a mark like the stamp of a seal, reproducing the form of the Sacred Host, the mould of which was one of those that bear a figure of Christ crucified. The noise of this and other wonders got for Juliana a reverence not only from Florence, but from all parts of the Christian world, which so increased through the course of four hundred years, that Pope Benedict XIII. commanded an office in her honour to be said by the whole Order of Servants of the Blessed Virgin Mary, and Clement XII., a munificent Protector of the same Order, finding new signs and wonders shedding lustre upon her memory every day, numbered her among Holy Virgins.\par
\switchcolumn*

\LA
\begin{Resp}\Rsign Afferéntur Regi vírgines post eam, próximæ ejus; \Star{} Afferéntur tibi in lætítia et exsultatióne.\end{Resp}
\begin{Resp}\Vsign Spécie tua et pulchritúdine tua; inténde, próspere procéde, et regna.\end{Resp}
\begin{Resp}\Rsign Afferéntur tibi in lætítia et exsultatióne.\end{Resp}
\begin{Resp}\Vsign Glória Patri, et Fílio, \Star{} et Spirítui Sancto.\end{Resp}
\begin{Resp}\Rsign Afferéntur tibi in lætítia et exsultatióne.\end{Resp}
\switchcolumn
\EN
\begin{Resp}\Rsign After her shall virgins be brought unto the King, her fellows \Star{} They shall be brought unto thee with gladness and rejoicing.\end{Resp}
\begin{Resp}\Vsign In thy comeliness and thy beauty, go forward, fare prosperously, and reign.\end{Resp}
\begin{Resp}\Rsign They shall be brought unto thee with gladness and rejoicing.\end{Resp}
\begin{Resp}\Vsign Glory be to the Father, and to the Son, \Star{} and to the Holy Ghost.\end{Resp}
\begin{Resp}\Rsign They shall be brought unto thee with gladness and rejoicing.\end{Resp}
\switchcolumn*

\LA
\LessonHead{Lectio IX (pro commemoratione)}
\switchcolumn
\EN
\LessonHead{Lesson IX (when commemorated)}
\switchcolumn*

\LA
\LessonTitle{Commemoratio: S. Julianæ de Falconeriis Virginis}
\switchcolumn
\EN
\LessonTitle{Commemoration of: St. Juliana Falconeri, Virgin}
\switchcolumn*

\LA
\noindent Juliána, ex nóbili Falconéria família, ab incunábulis, vagiéntibus lábiis suavíssima Jesu et Maríæ nómina ultro proférre audíta est. Expléto nondum décimo quinto ætátis anno, re familiári, licet opulénta, terrenísque núptiis posthábitis, Deo virginitátem in mánibus divi Philíppi Benítii solémniter vovit, et ab eo ómnium prima religiósum Mantellatárum, quas vocant, hábitum sumpsit. Cum vero ejus exémplum plúrimæ ex nobilióribus féminis sequeréntur, et mater fíliæ sese instituéndam dedísset, Juliána órdinem Mantellatárum instítuit. Mira humilitáte, assíduo oratiónis stúdio, singulári abstinéntia excélluit. Cum ob advérsam valetúdinem cibum cápere ac retinére nullo modo posset, ideóque ab Eucharística mensa arcerétur, sacerdótem rogávit, ut allátum divínum Panem, quem ore súmere nequíbat, péctori saltem extérius admovéret. Quod cum sacérdos præstitísset, íllico divínus Panis dispáruit, et Juliána ridénti vultu exspirávit.\par
\switchcolumn
\EN
\noindent When Juliana, of the noble family of the Falconieri, was still in her cradle, her baby lips were heard to utter, without any prompting, the sweet names of Jesus and Mary. Before she was fifteen years old, she renounced a rich inheritance and an earthly wedding and took a solemn vow of virginity in the presence of St. Philip Benizi. She was the first to receive from him the habit of the religious called the Mantellates. When many noble ladies followed her example, and even her mother gave herself over to her daughter to be instructed in the religious life, Juliana founded the Order of the Mantellate Nuns. She excelled in a wonderful humility, a constant zeal for prayer and an amazing abstinence. When her health failed so that she could take and retain no food at all, and was therefore kept from the Eucharistic table, she asked the Priest to place the divine Bread on her breast, since she could not receive it with her mouth. When he did so, the holy Bread disappeared at once, and Juliana, smiling, departed this life.\par
\switchcolumn*

\LA
\TeDeum{Te Deum}
\switchcolumn
\EN
\TeDeum{Te Deum}
\switchcolumn*

\end{paracol}

\par\needspace{10\baselineskip}\addvspace{1.2em}{\centering\small\textsc{Die 19 Junii} · \textsc{19 June}\par}
\Office{Ss. Gervasii et Protasii Martyrum}{Ss. Gervase and Protase, Martyrs}{Simplex}
\addcontentsline{toc}{subsection}{Die 19 Junii · Ss. Gervasii et Protasii Martyrum \textit{— Ss. Gervase and Protase, Martyrs}}
\begin{paracol}{2}
\LA
\LessonHead{Lectio IX (pro commemoratione)}
\switchcolumn
\EN
\LessonHead{Lesson IX (when commemorated)}
\switchcolumn*

\LA
\LessonTitle{Commemoratio: Ss. Gervasii et Protasii Martyrum}
\switchcolumn
\EN
\LessonTitle{Commemoration of: Ss. Gervase and Protase, Martyrs}
\switchcolumn*

\LA
\noindent Gervasius et Protasius, Vitalis et Valériæ fílii, quorum pater Ravennæ, mater Mediolani pro Christi Dómini fide martyrium subiérunt, distributo paupéribus patrimonio, domesticos servos libertáte donarunt. Quo facto Gentílium sacerdótes immane in illos conceptum ódium habebant. Quare, cum Astasius comes in bellum proficisci vellet, hanc occasiónem perdéndi pios fratres se nactos esse putavérunt. Itaque Astasio persuadent se a diis admonitos esse, nullo modo eum in bello victórem futurum, nisi Gervasio et Protasio coactis Christum negare, eosdem ad sacra diis facienda compélleret. Quod cum illi detestaréntur, Astasius imperávit Gervasium tamdiu cædi, dum inter verbera exspiraret; Protasium fustibus contusum secúri pércuti jubet. Quorum córpora Philippus Christi servus clam sústulit et in suis ædibus sepelívit: quæ póstea sanctus Ambrosius, Dei mónitu invénta, in loco sacro et insígni collocanda curávit. Passi sunt Mediolani décimo tertio Kalendas Julii.\par
\switchcolumn
\EN
\noindent Gervase and Protase were the sons of Vitalis and Valeria, both of whom testified even unto death for the Lord Christ's sake, the father at Ravenna, and the mother at Milan. After the victory of their parents, Gervase and Protase gave all the inheritance to the poor, and set free their slaves. This act of theirs stirred up against them a savage hatred on the part of the heathen priests, and when the Count Astasius was about setting forth to war, they believed they had got a good occasion for the destruction of the two godly brethren. They persuaded Astasius that their gods had revealed to them that he had no chance of conquering in the war, unless he had first made Gervase and Protase to deny Christ and to offer sacrifice to the gods. Being commanded so to do, they flatly refused, and Astasius then ordered Gervase to be lashed until he died between the stripes, and Protase to be cudgelled and beheaded. A servant of Christ named Philip took away their dead bodies by stealth, and buried them in his own house, and, in after times, St. Ambrose, being warned of God, found them, and bestowed them in an hallowed and honourable place. They suffered at Milan upon the 19th day of June.\par
\switchcolumn*

\LA
\TeDeum{Te Deum}
\switchcolumn
\EN
\TeDeum{Te Deum}
\switchcolumn*

\end{paracol}

\par\needspace{10\baselineskip}\addvspace{1.2em}{\centering\small\textsc{Die 20 Junii} · \textsc{20 June}\par}
\Office{S. Silverii Papæ et Martyris}{St. Silverius, Pope and Martyr}{Simplex}
\addcontentsline{toc}{subsection}{Die 20 Junii · S. Silverii Papæ et Martyris \textit{— St. Silverius, Pope and Martyr}}
\begin{paracol}{2}
\LA
\RefLine{Lectiones I–II: de Scriptura occurrente.}
\switchcolumn
\EN
\RefLine{Lessons I–II: from the Scripture of the day.}
\switchcolumn*

\LA
\LessonHead{Lectio III (pro commemoratione)}
\switchcolumn
\EN
\LessonHead{Lesson III (when commemorated)}
\switchcolumn*

\LA
\noindent Silverius Campanus post Agapitum proxime Pontifex creatus est: cujus doctrina et sanctitas illuxit in insectandis hæreticis, et constantis animi magnitudo perspecta est in tuendo judicio Agapiti. Nam Anthimum, quem, quia Eutychianam hæresim defendebat, Agapitus ab episcopatu Constantinopolitano deposuerat, cum a Theodora Augusta sæpissime rogatus esset, restituere noluit. Quam ob rem irata mulier mandat Belisario ut Silverium mittat in exsilium. Qui exsulavit in insula Pontia, unde his verbis scripsisse fertur ad Amatorem episcopum: Sustentor pane tribulationis et aqua angustiæ; nec tamen dimisi, aut dimitto officium meum. Et sane, brevi incommodis ærumnisque confectus, obdormivit in Domino, duodecimo Kalendas Julii; Cujus corpus Romam delatum, et in basilica Vaticana depositum, multis miraculis illustratum fuit. Præfuit Ecclesiæ annos tres et amplius, creatis mense Decembri presbyteris tredecim diaconis quinque, episcopis per diversa loca decem et novem.\par
\switchcolumn
\EN
\noindent This Silverius was a native of Campania, and succeeded Agapitus in the Papacy \textasciitilde{}(in the year 536.) His orthodoxy and holiness shone brightest in his onslaughts upon heretics, and he showed admirable firmness in upholding a sentence by Agapitus. Agapitus had deposed Anthimus from the Patriarchate of Constantinople for defending Save the heresy of Eutyches; and Silverius would never allow of his restoration, although the Empress Theodora repeatedly asked him to do so. The woman was enraged at him on this account, and ordered Bellisarius to send Silverius into exile. He was accordingly banished to the island of Ponza, whence he is said to have written these words to Bishop Amator I am fed upon the bread of tribulation and the water of affliction, but nevertheless I have not given up, and I will not give up, doing my duty. But sickness and the hardships of his exile soon broke his strength, and he fell asleep in the Lord upon the 20th day of June (in the year of grace 538.) Many miracles shed a lustre upon his grave. He ruled the Church for more than three years, and ordained in the month of December thirteen Priests, five Deacons, and nineteen Bishops for diverse Sees.\par
\switchcolumn*

\LA
\TeDeum{Te Deum}
\switchcolumn
\EN
\TeDeum{Te Deum}
\switchcolumn*

\end{paracol}

\par\needspace{10\baselineskip}\addvspace{1.2em}{\centering\small\textsc{Die 21 Junii} · \textsc{21 June}\par}
\Office{S. Aloisii Gonzagæ Confessoris}{St. Aloysius Gonzaga, Confessor}{Duplex}
\addcontentsline{toc}{subsection}{Die 21 Junii · S. Aloisii Gonzagæ Confessoris \textit{— St. Aloysius Gonzaga, Confessor}}
\begin{paracol}{2}
\LA
\RefLine{Lectiones I–III: de Scriptura occurrente.}
\switchcolumn
\EN
\RefLine{Lessons I–III: from the Scripture of the day.}
\switchcolumn*

\LA
\LessonHead{Lectio IV}
\switchcolumn
\EN
\LessonHead{Lesson IV}
\switchcolumn*

\LA
\noindent Aloisius, Ferdinandi Gonzagæ Castellionis Stiverorum marchionis filius, festinato propter vitæ periculum baptismo, prius cælo quam terris nasci visus, primam illam gratiam tam constanter retinuit, ut in ea confirmatus crederetur. A primo rationis usu, quo se Deo statim obtulit, vitam duxit quotidie sanctiorem. Novennis Florentiæ ante aram beatæ Virginis, quam parentis loco semper habuit, perpetuam virginitatem vovit: eamque, insigni Dei beneficio, nulla mentis aut corporis pugna tentatam servavit. Reliquas animi perturbationes cœpit ætate illa tam fortiter comprimere, ut ne primo quidem earum motu deínde incitaretur. Sensus etiam, oculos præcipue, ita cohibuit, ut non modo illos numquam in faciem intenderit Mariæ Austriacæ, quam plures annos inter honorarios Hispaniarum principis ephebos fere quotidie salutavit; sed a matris étiam vultu contineret: homo propterea sine carne, aut angelus in carne merito appellatus.\par
\switchcolumn
\EN
\noindent Aloysius, eldest son of Ferdinand Gonzaga, Marquess of Castiglione, was so hurriedly baptized on account of danger that he seemed to be born to heaven almost before he was born to earth, and he so faithfully kept that his first grace that he seemed to have been confirmed therein. From his first use of reason, which he employed to offer himself to God, he led a life more holy day by day. At Florence, when he was nine years old, he made a vow of perpetual virginity before the Altar of the Blessed Virgin, upon whom he always looked as in the place of a mother to him, and by a remarkable mercy, from God, he kept this vow wholly and without the slightest impure temptation, either of mind or body, during his whole life. As for any other uprisings of the soul, he began at that age to check them so sternly, that he was never more pricked by even their earliest movements. His senses, and especially, his eye-sight, he so mortified, that he never once looked upon the face of Mary of Austria, whom, when he was for several years one of the Pages of honour of the King of Spain, he saluted almost every day and he even denied himself in part, the pleasure of looking on the face of his own mother. He might indeed have been justly called a fleshless man, or an infleshed angel.\par
\switchcolumn*

\LA
\begin{Resp}\Rsign Honéstum fecit illum Dóminus, et custodívit eum ab inimícis, et a seductóribus tutávit illum: \Star{} Et dedit illi claritátem ætérnam.\end{Resp}
\begin{Resp}\Vsign Justum dedúxit Dóminus per vias rectas, et osténdit illi regnum Dei.\end{Resp}
\begin{Resp}\Rsign Et dedit illi claritátem ætérnam.\end{Resp}
\switchcolumn
\EN
\begin{Resp}\Rsign The Lord made him honourable, and defended him from his enemies, and kept him safe from those that lay in wait for him; \Star{} And gave him perpetual glory.\end{Resp}
\begin{Resp}\Vsign He went down with him into the pit, and left him not in bonds.\end{Resp}
\begin{Resp}\Rsign And gave him perpetual glory.\end{Resp}
\switchcolumn*

\LA
\LessonHead{Lectio V}
\switchcolumn
\EN
\LessonHead{Lesson V}
\switchcolumn*

\LA
\noindent Adjecit sensuum custodiæ corporis cruciatum. Tria singulis hebdomadis jejunia, eaque plerumque modico pane et aqua tolerabat, quamquam perpetuum fuísse per id tempus ipsius jejunium videri potest, cum ejus prandia ferme vix unciam æquarent. Sæpe étiam ter in die se funibus aut catenis cruentabat: flagella quandoque canum loris, cilicia equorum calcaribus supplevit. Mollem lectulum clam injectis asserum fragmentis asperabat, eo étiam ut citius ad orandum excitaretur, magnam quippe noctis partem, summa étiam hieme, solo tectus indusio, positis humi genibus, vel præ languore jacens ac pronus, in cælestium contemplatione traducebat. Interdiu quoque tres, quatuor, quinque horas in ea perstabat immotus, donec unam saltem animo nusquam distracto percurrisset. Cujus constantiæ præmium fuit stabilitas mentis inter orandum alio non vagantis, immo perpetua velut extasi in Deo defixæ. Ei demum ut unice adhæreret, victo post triennale acerrimum certamen patre, et aviti principatus jure in fratrem translato, societati Jesu, ad quam cælesti voce Matriti fuerat accitus, Romæ se adjunxit.\par
\switchcolumn
\EN
\noindent With this fettering of the senses he added torture of the body. He kept three days as fasts in every week, and that mostly upon a little bread and water. But indeed he as it were fasted every day, for he hardly ever took so much as an ounce weight of food at breakfast. Often also, even thrice in one day, he would lash himself to flowing of blood with cords, or prick himself with spiked chains. He sometimes used a dog-whip, instead of a scourge, and the rowels of spurs instead of hair-cloth. He privately filled his soft bed with pieces of broken plates, that he might find it easier to wake to pray. He passed great part of the night, clad only in a shirt even in the depth of winter, kneeling on the ground, or lying flat on his face when too weak and weary to remain upright, busied with heavenly thoughts. Sometimes he would keep himself thus for three, four, or five hours, until he had spent at least one without any movement of body or any wandering of mind. Such perseverance obtained for him the reward of being able to keep his understanding quite concentrated in prayer without distraction, as though rapt in God in an unbroken extasy. Desiring to give himself up to Him alone, he overcame, after a strong opposition for three years, the objections of his father, procured the transfer to his brother of his right to the Marquessate, and (on the 25th of November, 1585,) joined at Rome the Society of Jesus, to which he had been called by a voice from heaven when he was at Madrid.\par
\switchcolumn*

\LA
\begin{Resp}\Rsign Amávit eum Dóminus, et ornávit eum: stolam glóriæ índuit eum, \Star{} Et ad portas paradísi coronávit eum.\end{Resp}
\begin{Resp}\Vsign Índuit eum Dóminus lorícam fídei, et ornávit eum.\end{Resp}
\begin{Resp}\Rsign Et ad portas paradísi coronávit eum.\end{Resp}
\switchcolumn
\EN
\begin{Resp}\Rsign The Lord loved him and beautified him; He clothed him with a robe of glory, \Star{} And crowned him at the gates of Paradise.\end{Resp}
\begin{Resp}\Vsign The Lord hath put on him the breast-plate of faith, and hath adorned him.\end{Resp}
\begin{Resp}\Rsign And crowned him at the gates of Paradise.\end{Resp}
\switchcolumn*

\LA
\LessonHead{Lectio VI}
\switchcolumn
\EN
\LessonHead{Lesson VI}
\switchcolumn*

\LA
\noindent In tirocinio ipso virtutum omnium magister haberi cœpit, Exactissima in eo erat legum étiam minimarum custodia, mundi contemptus singularis implacabile odium sui Dei vero amor tam ardens, ut corpus étiam sensim absumeret, Jussus propterea mentem a divinis rebus tantisper avertere, occurrentem sibi ubique Deum irrito conatu fugiebat. Mira étiam proximos caritate amplexus, in publicis, quibus alacriter ministrabat, nosocomiis, contagiosam luem traxit, Qua lente consumptus, die quem prædixerat, undecimo Kalendas Julii, ætatis anno quarto et vigesimo jam inchoato, cum antea flagellis cædi, atque humi stratus mori postulasset, migravit in cælum, Ibi eum sancta Maria Magdalena de Pazzis tanta frui gloria, Deo monstrante, vidit, quantam vix esse in cælo credidisset; ipsumque sanctimonia insignem, et caritate Martyrem incognitum fuísse prædicavit, Multis étiam magnisque claruit miraculis, Quibus rite probatis, Benedictus decimus tertius Sanctorum fastis angelicum juvenum adscripsit, atque innocentiæ et castitatis exemplar simul et patronum studiosæ præsertim juventuti dedit.\par
\switchcolumn
\EN
\noindent In his very Noviciate he began to be held a master of all godliness. His obedience to even the most trifling rules was absolutely exact, his indifference to the world extraordinary, and his hatred of self implacable. His love of God was so keen that it gradually undermined his bodily strength. Being commanded to give his mind some rest from thinking unceasingly of God, he struggled vainly to distract himself from Him Who met him everywhere. From tender love toward his neighbour, he joyfully ministered to the sick in the public hospitals, (during the great distemper at Rome in 1591,) and in the exercise of this charity he caught a deadly disease. This sickness slowly wore him away, and soon after he had entered on the 24th year of his age, upon the 21st day of June, a day which he had himself foretold, after entreating that he might be scourged, and laid upon the ground to die, he passed away to heaven. What the glory is which he there enjoyeth holy Mary Magdalen de' Pazzi was enabled, by the revelation of God, to behold, and she declared that it was such as she had hardly believed existed even in heaven, and that his holiness and love were so great that she should call him an unknown martyr of charity. On earth God glorified him by many great miracles. These being duly proved, Benedict XIII inserted the name of this angellad in the Kalendar of the Saints, and commended him to all young scholars both as a pattern of innocency and purity, and as a patron.\par
\switchcolumn*

\LA
\begin{Resp}\Rsign Iste homo perfécit ómnia quæ locútus est ei Deus, et dixit ad eum: Ingrédere in réquiem meam: \Star{} Quia te vidi justum coram me ex ómnibus géntibus.\end{Resp}
\begin{Resp}\Vsign Iste est, qui contémpsit vitam mundi, et pervénit ad cæléstia regna.\end{Resp}
\begin{Resp}\Rsign Quia te vidi justum coram me ex ómnibus géntibus.\end{Resp}
\begin{Resp}\Vsign Glória Patri, et Fílio, \Star{} et Spirítui Sancto.\end{Resp}
\begin{Resp}\Rsign Quia te vidi justum coram me ex ómnibus géntibus.\end{Resp}
\switchcolumn
\EN
\begin{Resp}\Rsign This is he which did according unto all that God commanded him; and God said unto him: Enter thou into My rest, \Star{} For thee have I seen righteous before Me among all people.\end{Resp}
\begin{Resp}\Vsign This is he which loved not his life in this world, and is come unto an everlasting kingdom.\end{Resp}
\begin{Resp}\Rsign For thee have I seen righteous before Me among all people.\end{Resp}
\begin{Resp}\Vsign Glory be to the Father, and to the Son, \Star{} and to the Holy Ghost.\end{Resp}
\begin{Resp}\Rsign For thee have I seen righteous before Me among all people.\end{Resp}
\switchcolumn*

\LA
\LessonHead{Lectio VII}
\switchcolumn
\EN
\LessonHead{Lesson VII}
\switchcolumn*

\LA
\LessonTitle{Léctio sancti Evangélii secúndum Matthǽum}
\switchcolumn
\EN
\LessonTitle{From the Holy Gospel according to Matthew}
\switchcolumn*

\LA
\Cite{Matt 22:29-40}
\switchcolumn
\EN
\Cite{Matt 22:29-40}
\switchcolumn*

\LA
\noindent In illo témpore: Respondens Jesus ait Sadducæis: Erratis nescientes Scripturas, neque virtutem Dei. In resurrectione enim neque nubent, neque nubentur: sed erunt sicut Angeli Dei in cælo. Et réliqua.\par
\switchcolumn
\EN
\noindent At that time, Jesus answered and said unto the Sadducees: Ye do err, not knowing the Scriptures, nor the power of God for in the resurrection they neither marry nor are given in marriage, but are as the angels of God in heaven. And so on.\par
\switchcolumn*

\LA
\LessonTitle{Homilia sancti Joannis Chrysostomi}
\switchcolumn
\EN
\LessonTitle{Homily by St. John Chrysostom, Patriarch of Constantinople}
\switchcolumn*

\LA
\Source{Liber de Virginitate}
\switchcolumn
\EN
\Source{Bk on Virginity.}
\switchcolumn*

\LA
\noindent Virginitas est bona: id ego quoque fateor. Atqui nuptiis étiam melior et istud tibi assentior: ac si libet, illud adjungam, tanto nuptiis eam præstare, quanto cælum terræ, quanto hominibus Angeli antecellunt: ac, si quid præterea addendum est, étiam magis. Nam si neque nubunt Angeli, neque uxorem ducunt, non étiam carne et sanguine coagmentati sunt: in terris præterea non commorantur, non cupiditatum aut libidinum perturbationibus sunt obnoxii; non cibi indigent, aut potus; non sunt ejusmodi, ut eos dulcis sonus, aut cantus mollis, aut præclara species possit allicere; nulla denique ejus generis illecebra capiuntur.\par
\switchcolumn
\EN
\noindent This I say, that virginity is good. And in this I agree likewise, that it is better than marriage. And I will even add, that it is as much more excellent than marriage, as heaven is more noble than earth, or Angels than men, and indeed, if I must say more, even more so. For if Angels neither marry nor are given in marriage, at least, they are not creatures of flesh and blood, they dwell not upon earth, they are exposed to no restless troublings of desire or lust, they need neither meat nor drink, they are not such that sweet sound, or soft song, or the delight of beauty can charm them there is nothing of this sort to take hold on them and draw them away.\par
\switchcolumn*

\LA
\begin{Resp}\Rsign Iste est qui ante Deum magnas virtútes operátus est, et de omni corde suo laudávit Dóminum: \Star{} Ipse intercédat pro peccátis ómnium populórum.\end{Resp}
\begin{Resp}\Vsign Ecce homo sine queréla, verus Dei cultor, ábstinens se ab omni ópere malo, et pérmanens in innocéntia sua.\end{Resp}
\begin{Resp}\Rsign Ipse intercédat pro peccátis ómnium populórum.\end{Resp}
\switchcolumn
\EN
\begin{Resp}\Rsign This is he which wrought great wonders before God, and praised the Lord with all his heart. \Star{} May he pray for all people, that their sins may be forgiven unto them.\end{Resp}
\begin{Resp}\Vsign Behold a man without blame, a worshipper of God in truth, keeping himself clean from every evil work, and abiding still in his innocency.\end{Resp}
\begin{Resp}\Rsign May he pray for all people, that their sins may be forgiven unto them.\end{Resp}
\switchcolumn*

\LA
\LessonHead{Lectio VIII}
\switchcolumn
\EN
\LessonHead{Lesson VIII}
\switchcolumn*

\LA
\noindent At humánum genus cum natura beatis illis mentibus inferius sit, omni vi studioque contendit, ut, quoad ejus fieri potest, illas assequatur. Quomodo? Non nubunt Angeli: at neque étiam virgo. Assistentes illi semper ad Deum, eidem inserviunt; et istud ipsum virgo. Quod si virgines, quamdiu corporis onere deprimuntur, quemadmodum Angeli in cælum nequeunt ascendere, illud eo vel maximo solatio compensant, quod modo spiritu et corpore sancti sint, cæli Regem recipiunt. Videsne virginitatis præstantiam? quomodo terrarum incolas sic afficiat, ut qui corpore vestiti sunt, eos incorporeis mentibus exæquet?\par
\switchcolumn
\EN
\noindent Due the human nature which striveth its utmost to follow them, is not so exalted as that of these blessed intelligences. How Angels marry not nor are given in marriage neither doth a virgin. Angels stand ever before God, and serve Him and so doth a virgin. But if a virgin, still weighed down with this body, and unable, like the Angels, to ascend to heaven, doth make it his one great comfort here to be holy in body and in spirit, and to open his heart for a home for the King of heaven dost thou not see wherein a virgin is higher than an Angel The excellence of virginity in men over virginity in Angels lieth in this, that it maketh them which are yet earth-dwellers and body-burdened equal to intelligences unshackled by bodies.\par
\switchcolumn*

\LA
\begin{Resp}\Rsign Sint lumbi vestri præcíncti, et lucérnæ ardéntes in mánibus vestris: \Star{} Et vos símiles homínibus exspectántibus dóminum suum, quando revertátur a núptiis.\end{Resp}
\begin{Resp}\Vsign Vigiláte ergo, quia nescítis qua hora Dóminus vester ventúrus sit.\end{Resp}
\begin{Resp}\Rsign Et vos símiles homínibus exspectántibus dóminum suum, quando revertátur a núptiis.\end{Resp}
\begin{Resp}\Vsign Glória Patri, et Fílio, \Star{} et Spirítui Sancto.\end{Resp}
\begin{Resp}\Rsign Et vos símiles homínibus exspectántibus dóminum suum, quando revertátur a núptiis.\end{Resp}
\switchcolumn
\EN
\begin{Resp}\Rsign Let your loins be girded about, and your lights burning; \Star{} And ye yourselves like unto men that wait for their lord, when he will return from the wedding.\end{Resp}
\begin{Resp}\Vsign Watch therefore, for ye know not what hour your Lord doth come.\end{Resp}
\begin{Resp}\Rsign And ye yourselves like unto men that wait for their lord, when he will return from the wedding.\end{Resp}
\begin{Resp}\Vsign Glory be to the Father, and to the Son, \Star{} and to the Holy Ghost.\end{Resp}
\begin{Resp}\Rsign And ye yourselves like unto men that wait for their lord, when he will return from the wedding.\end{Resp}
\switchcolumn*

\LA
\LessonHead{Lectio IX}
\switchcolumn
\EN
\LessonHead{Lesson IX}
\switchcolumn*

\LA
\noindent Qua enim, quæso, re differebant ab Angelis Elias, Eliseus, Joannes, veri hi virginitatis amatores? Nulla, nisi quod mortali natura constabant. Nam cetera si quis diligenter inquirat, hi nihilo minus affecti reperientur, quam beatæ illæ mentes: et idipsum quo inferiore conditione videntur esse, in magna est eorum laude ponendum. Ut enim terrarum incolæ, et ii qui essent mortali natura, possent ad illam virtutem vi et contentione pervenire; vide quanta eos fortitudine, quanta vitæ ratione præditos fuísse oporteat.\par
\switchcolumn
\EN
\noindent In what respect, I ask, differed Elijah, Elisha, and John, those great lovers of virginity, from Angels In nothing, except that their faithfulness was exercised in a dying body. For the rest, if we look carefully, their minds were no otherwise than those of the blessed spirits, and their crown of glory is this that they attained the same honour under conditions less favourable. For consider of what manliness, of what superiority of reason over feeling they must have been possessed, to enable them bravely to fight their way, earth-dwellers and dying creatures as they were, to the bright summit of grace which was theirs.\par
\switchcolumn*

\LA
\TeDeum{Te Deum}
\switchcolumn
\EN
\TeDeum{Te Deum}
\switchcolumn*

\LA
\LessonHead{Lectio IX (pro commemoratione)}
\switchcolumn
\EN
\LessonHead{Lesson IX (when commemorated)}
\switchcolumn*

\LA
\LessonTitle{Commemoratio: S. Aloisii Gonzagæ Confessoris}
\switchcolumn
\EN
\LessonTitle{Commemoration of: St. Aloysius Gonzaga, Confessor}
\switchcolumn*

\LA
\noindent Aloísius, Ferdinándi Gonzágæ Castelliónis Stiverórum marchiónis fílius, festináto propter vitæ perículum baptísmo, prius cælo quam terris nasci visus, primam illam grátiam tam constánter retínuit, ut in ea confirmátus crederétur. Novénnis Floréntiæ ante aram beátæ Vírginis, quam paréntis loco semper hábuit, perpétuam virginitátem vovit; eámque, insígni Dei benefício, nulla mentis aut córporis pugna tentátam servávit, homo proptérea sine carne, aut ángelus in carne mérito appellátus. Avíti principátus jure in fratrem transmísso, societáti Jesu Romæ se adjúnxit. In tirocínio ipso ómnium virtútum magíster habéri cœpit. In eo Dei amor erat tam ardens, ut corpus étiam sensim absúmeret. Mira étiam próximos caritáte ampléxus, in públicis, quibus alácriter ministrábat, nosocomíis contagiósam luem traxit. Qua lente consúmptus, undécimo kaléndas júlii, ætátis anno quarto et vigésimo jam inchoáto, migrávit in cælum. Quem Benedíctus décimus tértius inter Sanctos rétulit, atque innocéntiæ et castitátis exémplar simul et patrónum studiósæ præsértim juventúti dedit.\par
\switchcolumn
\EN
\noindent Aloysius, son of Ferdinand Gonzaga, Marquis of Castiglióne delle Stiviere, was in danger of death while he was being born. He was therefore baptized without delay, so that it seemed he was born to heaven even before he was born to earth. He retained this first grace so faithfully that he was believed to have been confirmed in it. When he was nine years old, he took a vow of virginity at Florence before the altar of the Blessed Virgin, whom he always thought of as his mother. By a singular blessing of God, he kept this vow without any rebellion of mind or body so that he was deservedly called a man without a body or an angel in the flesh. He handed over the right of succession to his brother and joined the Society of Jesus in Rome. Even in the novitiate he began to be considered a master of all the virtues. So ardent was the love of God in him that he would be rapt out of his body. Possessed by a wonderful charity for his neighbour, he zealously served in the public hospitals, and as a result he contacted a contagious fever. After slowly wasting away, he went to heaven on the 21st day of June, having just entered his twenty-fourth year. Benedict XIII enrolled him among the Saints and gave him to students as both a model of innocency and charity and their heavenly Patron.\par
\switchcolumn*

\LA
\TeDeum{Te Deum}
\switchcolumn
\EN
\TeDeum{Te Deum}
\switchcolumn*

\end{paracol}

\par\needspace{10\baselineskip}\addvspace{1.2em}{\centering\small\textsc{Die 22 Junii} · \textsc{22 June}\par}
\Office{S. Paulini Episcopi et Confessoris}{St. Paulinus of Nola, Bishop and Confessor}{Duplex}
\addcontentsline{toc}{subsection}{Die 22 Junii · S. Paulini Episcopi et Confessoris \textit{— St. Paulinus of Nola, Bishop and Confessor}}
\begin{paracol}{2}
\LA
\RefLine{Lectiones I–III: de Scriptura occurrente.}
\switchcolumn
\EN
\RefLine{Lessons I–III: from the Scripture of the day.}
\switchcolumn*

\LA
\LessonHead{Lectio IV}
\switchcolumn
\EN
\LessonHead{Lesson IV}
\switchcolumn*

\LA
\Source{Breve Pii X diei 18 Sept. 1908}
\switchcolumn
\EN
\Source{Breve Pii X diei 18 Sept. 1908}
\switchcolumn*

\LA
\noindent Póntius Merópius Anícius Paulínus, anno reparátæ salútis trecentésimo quinquagésimo tértio, a claríssima cívium Romanórum família, Burdígalæ in Aquitánia natus, acri fuit ingénio ac móribus suávibus. Ausónio magístro, eloquéntiæ ac poéseos laude excélluit. Prænóbilis ac ditíssimus, honórum cursum ingréssus, florénti ætáte, senatória dignitáte potítus est. Dein Itáliam pétiit consul, et Campániam provínciam nactus, sedem Nolæ státuit. Hic divíno lúmine tactus, ob cæléstia signa, quæ Felícis presbyteri Mártyris sepúlcrum illustrábant, veræ Christi fídei, quam jam ánimo cogitábat, impénsius adhærére cœpit. Fasces ígitur ac secúrim nulla cæde maculátam depósuit, et revérsus in Gálliam, váriis ærúmnis ac magnis terra maríque labóribus jactátus, óculo cápitur; sed a beáto Martíno Turonénsi epíscopo sanitáti restitútus, lustrálibus baptísmatis aquis a beáto Delphíno Burdigalénsi antístite ablúitur.\par
\switchcolumn
\EN
\noindent Pontius Meropius Anicius Paulinus was born of a most illustrious family of Roman citizens at Bordeaux, in Aquitaine, in the year of restored salvation 353, and was of keen intelligence and graceful manners. With Ausonius as his master, he was honourably distinguished for eloquence and poetry. Being of the higher nobility, and very wealthy, he began a career of honours in the flower of his youth, and attained senatorial dignity. Then he went to Italy as consul, where he obtained the province of Campánia and fixed his residence at Nola. here he was struck, as by a ray of the divine light, by the heavenly miracles which were making illustrious the tomb of Felix, Priest and Martyr, and began to adhere more earnestly to the true faith of Christ, which he had long been revolving in his mind. And so he resigned the consular fasces and axe, never defiled by blood, and returned to Gaul, and after suffering various hardships and many toils both by land and by sea he received an injury to the eye; but being restored to health by the blessed Martin, Bishop of Tours, he was washed in the lustral waters of baptism by the blessed Delphinus, Bishop of Bordeaux.\par
\switchcolumn*

\LA
\begin{Resp}\Rsign Invéni David servum meum, óleo sancto meo unxi eum: \Star{} Manus enim mea auxiliábitur ei.\end{Resp}
\begin{Resp}\Vsign Nihil profíciet inimícus in eo, et fílius iniquitátis non nocébit ei.\end{Resp}
\begin{Resp}\Rsign Manus enim mea auxiliábitur ei.\end{Resp}
\switchcolumn
\EN
\begin{Resp}\Rsign I have found David My servant, with My holy oil have I anointed him \Star{} For My hand shall help him.\end{Resp}
\begin{Resp}\Vsign The enemy shall prevail nothing against him, nor the son of wickedness afflict him.\end{Resp}
\begin{Resp}\Rsign For My hand shall help him.\end{Resp}
\switchcolumn*

\LA
\LessonHead{Lectio V}
\switchcolumn
\EN
\LessonHead{Lesson V}
\switchcolumn*

\LA
\noindent Divítiis quibus abundábat spretis, bona véndidit pretiúmque paupéribus distríbuit, et uxórem linquens Therásiam, mutáta pátria et ruptis vínculis carnis, in Hispániam secéssit, venerándam secútus ac toto sibi pretiosiórem orbe Christi paupériem. Barcinóne dum Sacris devóte astáret, solémni die Domínicæ Nativitátis, repentíno admirátæ plebis tumúltu corréptus, ac frustra relúctans, a Lampídio epíscopo présbyter ordinátur. Inde redit in Itáliam, et Nolæ, quo sancti Felícis religióne ductus fúerat, penes illíus sepúlcrum monastérium cóndidit, et adscítis sóciis cœnobíticam vitam aggréditur. Hic vir, jam senatória et consulári dignitáte præclárus, stultítiam crucis ampléxus, toto fere orbe admiránte, vili indútus túnica, vigílias inter ac jejúnia, in assídua cæléstium rerum contemplatióne dies noctésque defíxus manébat. Sed, percrescénte sanctimóniæ fama, ad Nolánum episcopátum evéhitur, atque eódem in pastoráli múnere obeúndo, miránda pietátis, sapiéntiæ ac potíssimum caritátis exémpla relíquit.\par
\switchcolumn
\EN
\noindent He despised his abundant wealth, sold his property, and gave the money to the poor, and left his wife Therasia, changed his country, broke all natural ties, and retired to Spain, conforming to the more precious poverty of Christ, which he valued more than the whole world. While he was devoutly assisting at Mass in Barcelona on the feast of the Lord's Nativity, he was suddenly seized by force by an admiring crowd, and in spite of his reluctance, ordained priest by the Bishop Lampidius. Then he returned to Italy, and at Nola, where he had been drawn by devotion to St. Felix, he built a monastery near his tomb, and entered upon the monastic life with some companions. This man, already notable for his senatorial and consular dignity, embraced the folly of the cross to the admiration of almost the whole world, and, clothed in a mean garment, in vigils and in fastings, remained for days and nights fixed in constant contemplation of heavenly things. But, as the fame of his sanctity spread, he was elevated to the see of Nola, and entering upon that pastoral office, he gave a wonderful example of piety, wisdom, and above all, of charity.\par
\switchcolumn*

\LA
\begin{Resp}\Rsign Pósui adjutórium super poténtem, et exaltávi eléctum de plebe mea: \Star{} Manus enim mea auxiliábitur ei.\end{Resp}
\begin{Resp}\Vsign Invéni David servum meum, óleo sancto meo unxi eum.\end{Resp}
\begin{Resp}\Rsign Manus enim mea auxiliábitur ei.\end{Resp}
\switchcolumn
\EN
\begin{Resp}\Rsign I have laid help upon one that is mighty, and have exalted one chosen out of My people \Star{} For My hand shall help him.\end{Resp}
\begin{Resp}\Vsign I have found David My servant, with My holy oil have I anointed him.\end{Resp}
\begin{Resp}\Rsign For My hand shall help him.\end{Resp}
\switchcolumn*

\LA
\LessonHead{Lectio VI}
\switchcolumn
\EN
\LessonHead{Lesson VI}
\switchcolumn*

\LA
\noindent Hæc inter, sapiéntia reférta, de religióne ac fide pertractántia edíderat scripta, sæpe étiam númeris indúlgens concínnis carmínibus Sanctórum acta concelebráverat, summam christiáni poétæ famam adéptus. Quotquot sanctitáte ac doctrína præstantíssimi viri eo témpore erant, tot sibi amicítia atque admiratióne devínxit. Quamplúrimi ad eum, ceu ad christiánæ perfectiónis magístrum, undequáquam confluébant. Vastáta a Gothis Campánia, facultátem omnem, ne relíctis quidem sibi rebus ad vitam necessáriis, in aléndos páuperes et captívos rediméndos cóntulit. Póstea vero, Vándalis eásdem regiónes infestántibus, cum ab eo pósceret vídua ut fílium sibi redímeret ab hóstibus captum; consúmptis bonis ómnibus, in offício pietátis, se ipsum pro illo tradit in servitútem, atque in víncula conjéctus in Africam rápitur. Tandem, non sine præsénti Dei ope, libertáte donátus et Nolam revérsus, diléctum ovíle bonus pastor revísit; ibíque annum agens septuagésimum octávum ætátis suæ, placidíssimo éxitu obdormívit in Dómino. Corpus, prope sancti Felícis sepúlcrum cónditum, póstea, Longobardórum témpore, Benevéntum, atque, Ottóne tértio imperatóre, Romam ad basílicam sancti Bartholomæi ad ínsulam Tiberínam translátum fuit. Pius vero Papa décimus jussit sacras Paulíni exúvias Nolæ restítui, et festum ipsíus ad ritum dúplicem pro univérsa Ecclésia evéxit.\par
\switchcolumn
\EN
\noindent Meanwhile he had produced writings full of wisdom, treating of religion and faith, and often also as a recreation, he had celebrated the deeds of the Saints in numerous elegant poems, attaining the highest fame as a Christian poet. All the men living at that time who were pre-eminent for holiness and learning, he attached to himself with the bonds of friendship and admiration. Very many flocked to him from all parts of the country, as if to a master of Christian perfection. When Campánia was laid waste by the Goths, he devoted all his resources to feeding the poor, and ransoming captives, without even leaving for himself the necessities of life. And after that, when the Vandals invaded the same region, a widow begged him to ransom her son who had been captured by the enemy; and as he had spent all his resources in works of piety, he sold himself into slavery in the young man's place, and was taken in chains to Africa. At length he was set at liberty, not without the evident assistance of God, and, returning to Nola, the good shepherd once more saw his beloved flock. There, in a most peaceful end, he fell asleep in the Lord in the seventy-eighth year of his age. His body was buried near the tomb of St. Felix, and later, in the time of the Lombards, was translated to Benevento and, under the Emperor Otto III, to Rome, where it was laid in the basilica of St. Bartholomew on an island in the Tiber. But Pope Pius X ordered the sacred relics of Paulinus to be restored to Nola, and he raised his feast to the rite of a double for the universal Church.\par
\switchcolumn*

\LA
\begin{Resp}\Rsign Iste est, qui ante Deum magnas virtútes operátus est, et omnis terra doctrína ejus repléta est: \Star{} Ipse intercédat pro peccátis ómnium populórum.\end{Resp}
\begin{Resp}\Vsign Iste est, qui contémpsit vitam mundi, et pervénit ad cæléstia regna.\end{Resp}
\begin{Resp}\Rsign Ipse intercédat pro peccátis ómnium populórum.\end{Resp}
\begin{Resp}\Vsign Glória Patri, et Fílio, \Star{} et Spirítui Sancto.\end{Resp}
\begin{Resp}\Rsign Ipse intercédat pro peccátis ómnium populórum.\end{Resp}
\switchcolumn
\EN
\begin{Resp}\Rsign This is he which wrought great wonders before God, and the whole earth is full of his teaching \Star{} May he pray for all people, that their sins may be forgiven unto them\end{Resp}
\begin{Resp}\Vsign This is he which loved not his life in this world, and hath attained unto the kingdom of heaven.\end{Resp}
\begin{Resp}\Rsign May he pray for all people, that their sins may be forgiven unto them\end{Resp}
\begin{Resp}\Vsign Glory be to the Father, and to the Son, \Star{} and to the Holy Ghost.\end{Resp}
\begin{Resp}\Rsign May he pray for all people, that their sins may be forgiven unto them\end{Resp}
\switchcolumn*

\LA
\LessonHead{Lectio VII}
\switchcolumn
\EN
\LessonHead{Lesson VII}
\switchcolumn*

\LA
\LessonTitle{Léctio sancti Evangélii secúndum Lucam}
\switchcolumn
\EN
\LessonTitle{The Lesson is taken from the Holy Gospel according to Luke}
\switchcolumn*

\LA
\Cite{Luc 12:32-34}
\switchcolumn
\EN

\switchcolumn*

\LA

\switchcolumn
\EN
\Source{Chap. 12, 32-34}
\switchcolumn*

\LA
\noindent In illo témpore: Dixit Jesus discípulis suis: Nolíte timére, pusíllus grex, quia complácuit Patri vestro dare vobis regnum. Et réliqua.\par
\switchcolumn
\EN
\noindent At that time: Jesus said unto his disciples: Fear not, little flock; for it is your Father's good pleasure to give you the kingdom. And so on, and that which followeth.\par
\switchcolumn*

\LA
\LessonTitle{Homilía sancti Paulíni Epíscopi}
\switchcolumn
\EN
\LessonTitle{A Homily by St. Paulinus the Bishop}
\switchcolumn*

\LA
\Source{Sermo, alias Epistola 34 de Gazophylacio}
\switchcolumn
\EN
\Source{Sermo, alias Epistola 34 de Gazophylacio}
\switchcolumn*

\LA
\noindent Potúerat, dilectíssimi, Dóminus omnípotens æque univérsos dívites fácere, ut nemo indigéret áltero; sed infinítæ bonitátis consílio sic parávit miséricors et miserátor Dóminus, ut tuam in illis mentem probet. Fecit míserum, ut agnósceret misericórdem; fecit ínopem, ut exercéret opuléntum. Matéria divitiárum tibi est fratérna paupértas, si intélligas super egénum et páuperem, nec tibi tantum hábeas quod accepísti; quia ídeo et illíus partem tibi in hoc sæculo cóntulit Deus, ut tibi debéret quidquid de suis donis tuo voluntário afféctu indigéntibus obtulísses, ac te vicíssim in ætérna die de illíus parte ditáret. Per ipsos enim nunc áccipit Christus, et tunc pro ipsis repéndet.\par
\switchcolumn
\EN
\noindent The Lord who is omnipotent, dearly beloved, might have made all men equally rich, so that no one need ask anything from another. In his infinite goodness, however, the merciful and gracious Lord hath planned otherwise, in order to prove your disposition in these matters. He hath made misery, that he might discern mercy; he hath made the needy, that he might make use of the rich. For your brother's poverty is your material of riches, if ye do not understand concerning the needy and the poor, and do not consider as your own what wealth ye have received. For God hath bestowed upon you your brother's portion in this world, in order that ye should offer of your own willing affection something of his gifts to those in need, and that he may enrich you in your turn with that portion in eternity. For now Christ receiveth through them, and hereafter he will repay for them.\par
\switchcolumn*

\LA
\begin{Resp}\Rsign Amávit eum Dóminus, et ornávit eum: stolam glóriæ índuit eum, \Star{} Et ad portas paradísi coronávit eum.\end{Resp}
\begin{Resp}\Vsign Induit eum Dóminus lorícam fídei, et ornávit eum.\end{Resp}
\begin{Resp}\Rsign Et ad portas paradísi coronávit eum.\end{Resp}
\switchcolumn
\EN
\begin{Resp}\Rsign The Lord loved him and beautified him He clothed him with a robe of glory \Star{} And crowned him at the gates of Paradise.\end{Resp}
\begin{Resp}\Vsign The Lord hath put on him the breast-plate of faith, and hath adorned him.\end{Resp}
\begin{Resp}\Rsign And crowned him at the gates of Paradise.\end{Resp}
\switchcolumn*

\LA
\LessonHead{Lectio VIII}
\switchcolumn
\EN
\LessonHead{Lesson VIII}
\switchcolumn*

\LA
\noindent Réfice esuriéntem ánimam, et non timébis in die malo ab ira superventúra. Beátus enim (inquit), qui intélligit super egénum et páuperem, in die malo liberábit eum Dóminus. Operáre ígitur et éxcole hanc regiónem terræ tuæ, frater, ut gérminet tibi frugem fértilem, plenam ádipe fruménti, magno cum fœnore centésimum tibi fructum multiplicáti séminis afferéntem. In hujus vel possessiónis vel negotiatiónis appetítum et stúdium sancta et salutáris est avarítia; nam talis cupíditas, quæ regnum cæléste merétur et bonum perénne desíderat, radix bonórum est. Tales ígitur divítias concupíscite, et hujúsmodi possidéte patrimónium, quod in centénos fructus vobis créditor pénsitet, ut vestros quoque vobíscum bonis perénnibus augeátis herédes. Posséssio enim hæc vere magna et pretiósa est, quæ possessórem suum non cúmulo sæculári ónerat, sed réditu ditat ætérno.\par
\switchcolumn
\EN
\noindent Refresh the hungry soul, and ye will not fear the wrath to come in the evil day. For, saith he, Blessed is he that understandeth concerning the needy and poor; the Lord will deliver him in the evil day. Therefore, brother, work and cultivate this parcel of thy land, that it may bear for you a rich crop, full of the fatness of wheat, bringing to you with great interest the fruit of the seed an hundred times multiplied. In the desire and pursuit of this business or this possession there is a holy and salutary avarice; for such desire which doth merit the kingdom of heaven, and doth long for eternal goods, is the root of all good. Therefore covet such riches, and take possession of this kind of patrimony, that the creditor may weigh out to you the fruit increased an hundredfold, and ye and your heirs may abound in everlasting good things. For this possession is great and precious, which doth not burden the owner with earthly wealth, but enricheth him with an eternal reward.\par
\switchcolumn*

\LA
\begin{Resp}\Rsign Sint lumbi vestri præcíncti, et lucérnæ ardéntes in mánibus vestris: \Star{} Et vos símiles homínibus exspectántibus dóminum suum, quando revertátur a núptiis.\end{Resp}
\begin{Resp}\Vsign Vigiláte ergo, quia nescítis qua hora Dóminus vester ventúrus sit.\end{Resp}
\begin{Resp}\Rsign Et vos símiles homínibus exspectántibus dóminum suum, quando revertátur a núptiis.\end{Resp}
\begin{Resp}\Vsign Glória Patri, et Fílio, \Star{} et Spirítui Sancto.\end{Resp}
\begin{Resp}\Rsign Et vos símiles homínibus exspectántibus dóminum suum, quando revertátur a núptiis.\end{Resp}
\switchcolumn
\EN
\begin{Resp}\Rsign Let your loins be girded about, and your lights burning; \Star{} And ye yourselves like unto men that wait for their lord, when he will return from the wedding.\end{Resp}
\begin{Resp}\Vsign Watch therefore, for ye know not what hour your Lord doth come.\end{Resp}
\begin{Resp}\Rsign And ye yourselves like unto men that wait for their lord, when he will return from the wedding.\end{Resp}
\begin{Resp}\Vsign Glory be to the Father, and to the Son, \Star{} and to the Holy Ghost.\end{Resp}
\begin{Resp}\Rsign And ye yourselves like unto men that wait for their lord, when he will return from the wedding.\end{Resp}
\switchcolumn*

\LA
\LessonHead{Lectio IX}
\switchcolumn
\EN
\LessonHead{Lesson IX}
\switchcolumn*

\LA
\noindent Verum, dilectíssimi, non solum ut bona ætérna quærátis, sed ut mala innúmera vitáre mereámini, præsénti sollicitúdine et sédula operatióne justítiæ providéte. Magno enim adjutório atque præsídio nobis opus est, et multárum atque indeficiéntium oratiónum patrocíniis indigémus. Adversárius enim noster non quiéscit, et in nostrum pérvigil hostis intéritum óbsidet omnes vias nostras. Multæ prætérea nobis in hoc sæculo cruces, innúmera discrímina, morbórum labes, fébrium ignes et dolórum tela grassántur in ánimas, cupiditátum faces accendúntur; ubíque præténti latent láquei, úndique stricti horrent gládii, inter insídias et pugnas vita transígitur, et per ignes dolóso cíneri suppósitos ambulámus. Igitur, priúsquam in áliquam tantárum ægritúdinum labem casu vel mérito actus incúrras, festína médico suscéptus et carus fíeri, ut in témpore necessitátis parátum hábeas remédium salútis. Aliud est, quando tu solus oras pro te; et áliud, quando multitúdo pro te apud Deum trépidat.\par
\switchcolumn
\EN
\noindent Truly, dearly beloved, by present carefulness and constant just dealing, provide not only that ye make seek eternal good things, but that ye may deserve to escape innumerable evils. For we require strong helps and protection, and stand in need of many and unceasing prayers for our defence. For our adversary does not rest, and the ever-watchful enemy besetteth all our ways in order to work our ruin. Moreover, in this world we have many crosses, very numerous struggles, pestilent diseases, fires of fever and stabs of pain do violence to our souls. The flames of lust are kindled; hidden snares are set everywhere, on all sides there bristle drawn swords, life is passed in the midst of ambushes and combats, and we walk through fires craftily buried under ashes. Therefore before ye should by your misfortune or your fault rush into any one of such great calamities, hasten to become acceptable and dear to the physician, so that in time of need ye may have ready the remedy of salvation. It is one thing when ye alone pray for yourself; and quite another when a multitude entreateth for you before God.\par
\switchcolumn*

\LA
\TeDeum{Te Deum}
\switchcolumn
\EN
\TeDeum{Te Deum}
\switchcolumn*

\LA
\LessonHead{Lectio IX (pro commemoratione)}
\switchcolumn
\EN
\LessonHead{Lesson IX (when commemorated)}
\switchcolumn*

\LA
\LessonTitle{Commemoratio: S. Paulini Episcopi et Confessoris}
\switchcolumn
\EN
\LessonTitle{Commemoration of: St. Paulinus of Nola, Bishop and Confessor}
\switchcolumn*

\LA
\noindent Paulínus, anno reparátæ salútis trecentésimo quinquagésimo tértio a claríssima cívium Romanórum família Burdígalæ natus, senatória dignitáte potítus est. Nolæ consul renuntiátus, divíno lúmine tactus consulátum dimísit, et, Burdígalam revérsus, a beáto Delphíno baptizátus est. Dein, bonórum quibus abundábat prétio paupéribus distribúto, in Hispániam secéssit, ubi présbyter ordinátur. Nolam cum rediísset, penes sancti Felícis sepúlcrum monastérium cóndidit, sociísque adscítis, cœnobíticam vitam eámque arctíssimam est aggréssus. Percrebrescénte ejus sanctimóniæ fama, ad Nolánum episcopátum evéhitur, quo in múnere admiránda pietátis, patiéntiæ ac potíssimum caritátis exémpla relíquit. Multa ad sacram doctrínam pertinéntia scripsit, atque ínsuper eloquéntiæ et poéseos laude excélluit. Vastáta a Gothis Campánia, facultátes omnes in aléndos páuperes et captívos rediméndos cóntulit. Póstea vero, Vándalis eásdem regiónes infestántibus, cum nil ámplius erogándum habéret, se ipsum pro fílio cujúsdam víduæ in servitútem trádidit, et in Africam ductus est. Tandem, Dei ope libertáte donátus, Nolæ placidíssimo éxitu obdormívit in Dómino.\par
\switchcolumn
\EN
\noindent Paulinus, born in the year of restored salvation 353 of a very famous Roman family of Bordeaux, acquired the dignity of senator. He was made consul of Nola, but a divine light prompted him to renounce the consulship and return to Bordeaux, where he was baptized by St. Delphinus. Then, giving to the poor the large sum obtained by the sale of his goods, he went to Spain, where he was ordained priest. When he returned to Nola, he built a monastery near the grave of St. Felix and, with the companions who joined him, undertook a most austere cenobitical life. As the fame of his holiness grew, he was elevated to the bishopric of Nola, in which office he left an example of wonderful devotion, patience, and especially charity. He wrote many works on sacred doctrine and also gained a reputation for eloquence and poetry. When Campania was devastated by the Goths, he used all his goods to feed the poor and redeem captives. And later, when the Vandals infested the same region and he had nothing more to give, he gave himself into slavery for the son of a widow, and was taken to Africa. At length, restored to liberty, by the hand of God, he died a peaceful death in the Lord at Nola.\par
\switchcolumn*

\LA
\TeDeum{Te Deum}
\switchcolumn
\EN
\TeDeum{Te Deum}
\switchcolumn*

\end{paracol}

\par\needspace{10\baselineskip}\addvspace{1.2em}{\centering\small\textsc{Die 23 Junii} · \textsc{23 June}\par}
\Office{In Vigilia S. Joannis Baptistæ}{Vigil of the Nativity of St. John the Baptist}{Simplex}
\addcontentsline{toc}{subsection}{Die 23 Junii · In Vigilia S. Joannis Baptistæ \textit{— Vigil of the Nativity of St. John the Baptist}}
\begin{paracol}{2}
\LA
\LessonHead{Lectio I}
\switchcolumn
\EN
\LessonHead{Lesson I}
\switchcolumn*

\LA
\LessonTitle{Léctio sancti Evangélii secúndum Lucam}
\switchcolumn
\EN
\LessonTitle{Continuation of the Holy Gospel according to Luke}
\switchcolumn*

\LA
\Cite{Luc 1:5-17}
\switchcolumn
\EN
\Cite{Luke 1:5-17}
\switchcolumn*

\LA
\noindent Fuit in diébus Heródis, regis Judǽæ, sacérdos quidam, nómine Zacharías, de vice Abía et uxor illíus de filiábus Aaron, et nomen ejus Elísabeth. Et réliqua.\par
\switchcolumn
\EN
\noindent There was, in the days of Herod the King of Judaea, a certain Priest named Zacharias, of the course of Abia; and his wife was of the daughters of Aaron, and her name was Elizabeth. And so on.\par
\switchcolumn*

\LA
\LessonTitle{Homilía sancti Ambrósii Epíscopi}
\switchcolumn
\EN
\LessonTitle{Homily by St. Ambrose, Bishop of Milan}
\switchcolumn*

\LA
\Source{Liber 1 in Lucam}
\switchcolumn
\EN
\Source{Bk. i. on Luke}
\switchcolumn*

\LA
\noindent Docet nos Scriptúra divína non solum mores in iis qui prædicábiles sunt, sed étiam paréntes oportére laudári; ut véluti transmíssa immaculátæ puritátis heréditas, in iis quos vólumus laudáre, præcéllat. Quæ enim ália inténtio hoc loco sancti Evangelístæ, nisi ut sanctus Joánnes Baptísta nobilitétur paréntibus, miráculis, móribus, múnere, passióne? Sic étiam sancti Sámuel mater Anna laudátur; sic Isaac a paréntibus nobilitátem pietátis accépit, quam pósteris derelíquit. Sacérdos ítaque Zacharías, nec solum sacérdos, sed étiam de vice Abía, id est, nóbilis inter superióres famílias.\par
\switchcolumn
\EN
\noindent The Divine Scriptures teach us that we are behoven to praise the lives, not only of those concerning whom we are to speak honourably, but the lives also of their fathers, so as to show that that which we will praise in our subjects was in them a gift inherited from the bright purity of the source from which they came. What other meaning can the holy Evangelist have had in this place but to glorify St. John the Baptist, as well for having been the offspring of such parents, as for his miracles, his life, his gifts, and his sufferings? So likewise is praise ascribed to Hannah, the mother of Samuel so also did Isaac draw from his parents that noble godliness which he in his turn bequeathed to his children. Thus it is told not only that Zacharias was a Priest, but a Priest of the course of Abia, that is to say, of a family noble among the noblest.\par
\switchcolumn*

\LA
\begin{Resp}\Rsign Peccávi super númerum arénæ maris, et multiplicáta sunt peccáta mea: et non sum dignus vidére altitúdinem cæli præ multitúdine iniquitátis meæ: quóniam irritávi iram tuam, \Star{} Et malum coram te feci.\end{Resp}
\begin{Resp}\Vsign Quóniam iniquitátem meam ego cognósco: et delíctum meum contra me est semper, quia tibi soli peccávi.\end{Resp}
\begin{Resp}\Rsign Et malum coram te feci.\end{Resp}
\switchcolumn
\EN
\begin{Resp}\Rsign My sins are many, yea, they are more in number than the sands of the sea; I am not worthy to look up toward heaven because of the multitude of my iniquities; for I have provoked thee to anger; \Star{} And done evil in thy sight.\end{Resp}
\begin{Resp}\Vsign For I acknowledge my transgression, and my sin is ever before me, for against thee only have I sinned.\end{Resp}
\begin{Resp}\Rsign And done evil in thy sight.\end{Resp}
\switchcolumn*

\LA
\LessonHead{Lectio II}
\switchcolumn
\EN
\LessonHead{Lesson II}
\switchcolumn*

\LA
\noindent Et uxor, inquit, illi de filiábus Aaron. Non solum ígitur a paréntibus, sed étiam a majóribus sancti Joánnis nobílitas propagátur, non sæculári potestáte sublímis, sed religiónis successióne venerábilis. Tales enim majóres habére débuit prænúntius Christi; ut non repénte concéptam, sed a majóribus accéptam et ipso infúsam jure natúræ, prædicáre fidem Domínici viderétur advéntus. Erant, inquit, ambo justi ante Deum, incedéntes in ómnibus mandátis et justificatiónibus Dómini sine queréla. Quid ad hoc réferunt qui, peccátis suis solátia præferéntes, sine peccátis frequéntibus hóminem putant esse non posse; et utúntur versículo, quia scriptum est in Job: Nemo mundus a sorde, nec si uníus diéi vita ejus sit in terra?\par
\switchcolumn
\EN
\noindent And his wife was of the daughters of Aaron. Thus we see that the noble blood of St. John was inherited not only from parents, but from an ancient ancestry, not illustrious indeed by worldly power, but worshipful for the tradition of a sacred succession. Such were the forefathers whom it well became the Fore-runner of the Christ to have, that it might manifestly fall to his lot, not as a sudden gift, but as an heir-loom, to preach belief in the coming of the Lord. And they were both righteous before God, walking in all the commandments and ordinances of the Lord, blameless. What do they make of this text who, to take them some consolation for their own sins, hold that man cannot exist without oftentimes sinning, and quote to that end that which is written in Job: Not one is clean, even though his life on the earth be but one day?\par
\switchcolumn*

\LA
\begin{Resp}\Rsign Exaudísti, Dómine, oratiónem servi tui, ut ædificárem templum nómini tuo: \Star{} Bénedic et sanctífica domum istam in sempitérnum, Deus Israël.\end{Resp}
\begin{Resp}\Vsign Dómine, qui custódis pactum cum servis tuis, qui ámbulant coram te in toto corde suo.\end{Resp}
\begin{Resp}\Rsign Bénedic et sanctífica domum istam in sempitérnum, Deus Israël.\end{Resp}
\switchcolumn
\EN
\begin{Resp}\Rsign O Lord, Thou hast hearkened unto the prayer of thy servant, that I might build a temple unto thy Name, \Star{} O God of Israel, bless Thou, and hallow this house for ever.\end{Resp}
\begin{Resp}\Vsign O Lord, Who keepest covenant with thy servants that walk before thee in all their heart.\end{Resp}
\begin{Resp}\Rsign O God of Israel, bless Thou, and hallow this house for ever.\end{Resp}
\switchcolumn*

\LA
\LessonHead{Lectio III}
\switchcolumn
\EN
\LessonHead{Lesson III}
\switchcolumn*

\LA
\noindent Quibus respondéndum est, prius ut quid sit hóminem sine peccáto esse defíniant: utrum nunquam omníno peccásse, an desiísse peccáre. Si enim hoc putant sine peccáto esse, nunquam omníno peccásse; et ipse conséntio. Omnes enim peccavérunt, et egent glória Dei. Sin autem eum qui véterem errórem corréxerit, et in eam se vitæ transformáverit qualitátem, ut témperet a peccáto, negant abstinére a delíctis; non possum in eórum conveníre senténtiam, cum legámus, quia Sic Dóminus diléxit Ecclésiam, ut exhíbeat ipsam sibi gloriósam, et non habéntem máculam, aut rugam, aut áliquid ejúsmodi; sed ut sit sancta et immaculáta.\par
\switchcolumn
\EN
\noindent To such we must reply by asking them first to tell us what they mean by a man without sin whether it be one who hath never sinned, or one who hath ceased to sin. If they mean by a man without sin one who hath never sinned, I myself agree in their position, for all have sinned and come short of the glory of God. (Rom. iii. 23.) But if they mean to deny that he who hath reformed his old crooked ways, and changed his life for a new one, on purpose to avoid sin, cannot avoid sin, I am not able to subscribe to their opinion while I read that Christ loved the Church and gave Himself for it, that He might sanctify and cleanse it with the washing of water by the word, that He might present it to Himself a glorious Church, not having spot, or wrinkle, or any such thing but that it should be holy and without blemish. (Eph. v. 25-27.)\par
\switchcolumn*

\LA
\begin{Resp}\Rsign Audi, Dómine, hymnum et oratiónem, quam servus tuus orat coram te hódie, ut sint óculi tui apérti, et aures tuæ inténtæ, \Star{} Super domum istam die ac nocte.\end{Resp}
\begin{Resp}\Vsign Réspice, Dómine, de sanctuário tuo, et de excélso cælórum habitáculo.\end{Resp}
\begin{Resp}\Rsign Super domum istam die ac nocte.\end{Resp}
\begin{Resp}\Vsign Glória Patri, et Fílio, \Star{} et Spirítui Sancto.\end{Resp}
\begin{Resp}\Rsign Super domum istam die ac nocte.\end{Resp}
\switchcolumn
\EN
\begin{Resp}\Rsign Hearken, O Lord, unto the cry and to the prayer which thy servant prayeth before thee today, that thine eyes may be open and thine ears attend; \Star{} Toward this house day and night.\end{Resp}
\begin{Resp}\Vsign Look down from thine high and holy place, O Lord, even from heaven thy dwelling.\end{Resp}
\begin{Resp}\Rsign Toward this house, day and night.\end{Resp}
\begin{Resp}\Vsign Glory be to the Father, and to the Son, \Star{} and to the Holy Ghost.\end{Resp}
\begin{Resp}\Rsign Toward this house, day and night.\end{Resp}
\switchcolumn*

\end{paracol}

\par\needspace{10\baselineskip}\addvspace{1.2em}{\centering\small\textsc{Die 24 Junii} · \textsc{24 June}\par}
\Office{In Nativitate S. Joannis Baptistæ}{Nativity of St. John the Baptist}{Duplex I classis cum Octava communi}
\addcontentsline{toc}{subsection}{Die 24 Junii · In Nativitate S. Joannis Baptistæ \textit{— Nativity of St. John the Baptist}}
\begin{paracol}{2}
\LA
\RefLine{Lectio IX: de Scriptura occurrente.}
\switchcolumn
\EN
\RefLine{Lesson IX: from the Scripture of the day.}
\switchcolumn*

\LA
\LessonHead{Lectio I}
\switchcolumn
\EN
\LessonHead{Lesson I}
\switchcolumn*

\LA
\LessonTitle{Incipit liber Jeremíæ Prophétæ}
\switchcolumn
\EN
\LessonTitle{Lesson from the book of Jeremias}
\switchcolumn*

\LA
\Cite{Jer 1:1-5}
\switchcolumn
\EN
\Cite{Jer 1:1-5}
\switchcolumn*

\LA
\noindent Verba Jeremíæ fílii Helcíæ, de sacerdótibus, qui fuérunt in Anathoth, in terra Bénjamin. Quod factum est verbum Dómini ad eum in diébus Josíæ fílii Amon regis Juda, in tertiodécimo anno regni ejus. Et factum est in diébus Jóakim fílii Josíæ regis Juda, usque ad consummatiónem undécimi anni Sedecíæ fílii Josíæ regis Juda, usque ad transmigratiónem Jerúsalem, in mense quinto. Et factum est verbum Dómini ad me, dicens: Priúsquam te formárem in útero, novi te: et ántequam exíres de vulva, sanctificávi te, et prophétam in géntibus dedi te.\par
\switchcolumn
\EN
\noindent The words of Jeremias the son of Helcias, of the priests that were in Anathoth, in the land of Benjamin. The word of the Lord which came to him in the days of Josias the son of Amon king of Juda, in the thirteenth year of his reign. And which came to him in the days of Joakim the son of Josias king of Juda, unto the end of the eleventh year of Sedecias the son of Josias king of Juda, even unto the carrying away of Jerusalem captive, in the fifth month. And the word of the Lord came to me, saying: Before I formed thee in the bowels of thy mother, I knew thee: and before thou camest forth out of the womb, I sanctified thee, and made thee a prophet unto the nations.\par
\switchcolumn*

\LA
\begin{Resp}\Rsign Fuit homo missus a Deo, cui nomen erat Joánnes: \Star{} Hic venit in testimónium, ut testimónium perhibéret de lúmine, et paráret Dómino plebem perféctam.\end{Resp}
\begin{Resp}\Vsign Erat Joánnes in desérto prǽdicans baptísmum pœniténtiæ.\end{Resp}
\begin{Resp}\Rsign Hic venit in testimónium, ut testimónium perhibéret de lúmine, et paráret Dómino plebem perféctam.\end{Resp}
\switchcolumn
\EN
\begin{Resp}\Rsign There was a man sent from God, whose name was John. \Star{} The same came for a witness, to bear witness of the Light, to make ready a people prepared for the Lord.\end{Resp}
\begin{Resp}\Vsign John was in the wilderness, preaching the baptism of repentance.\end{Resp}
\begin{Resp}\Rsign The same came for a witness, to bear witness of the Light, to make ready a people prepared for the Lord.\end{Resp}
\switchcolumn*

\LA
\LessonHead{Lectio II}
\switchcolumn
\EN
\LessonHead{Lesson II}
\switchcolumn*

\LA
\Cite{Jer 1:6-10}
\switchcolumn
\EN
\Cite{Jer 1:6-10}
\switchcolumn*

\LA
\noindent Et dixi, A a a, Dómine Deus: ecce néscio loqui, quia puer ego sum. Et dixit Dóminus ad me: Noli dícere, Puer sum: quóniam ad ómnia, quæ mittam te, ibis: et univérsa, quæcúmque mandávero tibi, loquéris. Ne tímeas a fácie eórum: quia tecum ego sum, ut éruam te, dicit Dóminus. Et misit Dóminus manum suam, et tétigit os meum: et dixit Dóminus ad me: Ecce dedi verba mea in ore tuo: ecce constítui te hódie super gentes, et super regna, ut evéllas, et déstruas, et dispérdas, et díssipes, et ædífices, et plantes.\par
\switchcolumn
\EN
\noindent And I said: Ah, ah, ah, Lord God: behold, I cannot speak, for I am a child. And the Lord said to me: Say not: I am a child: for thou shalt go to all that I shall send thee: and whatsoever I shall command thee, thou shalt speak. Be not afraid at their presence: for I am with thee to deliver thee, saith the Lord. And the Lord put forth his hand, and touched my mouth: and the Lord said to me: Behold I have given my words in thy mouth: Lo, I have set thee this day over the nations, and over the kingdoms, to root up, and pull down, and to waste, and to destroy, and to build, and to plant.\par
\switchcolumn*

\LA
\begin{Resp}\Rsign Elísabeth Zacharíæ magnum virum génuit, Joánnem Baptístam, præcursórem Dómini: \Star{} Qui viam Dómino præparávit in erémo.\end{Resp}
\begin{Resp}\Vsign Fuit homo missus a Deo, cui nomen erat Joánnes.\end{Resp}
\begin{Resp}\Rsign Qui viam Dómino præparávit in erémo.\end{Resp}
\switchcolumn
\EN
\begin{Resp}\Rsign Elizabeth, the wife of Zacharias, was the mother of a mighty man, even of John the Baptist, the Forerunner of the Lord, \Star{} who made straight in the desert an highway for the Lord.\end{Resp}
\begin{Resp}\Vsign There was a man sent from God, whose name was John.\end{Resp}
\begin{Resp}\Rsign Who made straight in the desert an highway for the Lord.\end{Resp}
\switchcolumn*

\LA
\LessonHead{Lectio III}
\switchcolumn
\EN
\LessonHead{Lesson III}
\switchcolumn*

\LA
\Cite{Jer 1:17-19}
\switchcolumn
\EN
\Cite{Jer 1:17-19}
\switchcolumn*

\LA
\noindent Tu ergo accínge lumbos tuos, et surge, et lóquere ad eos ómnia quæ ego præcípio tibi. Ne formídes a fácie eórum: nec enim timére te fáciam vultum eórum. Ego quippe dedi te hódie in civitátem munítam, et in colúmnam férream, et in murum ǽreum, super omnem terram, régibus Juda, princípibus ejus, et sacerdótibus, et pópulo terræ. Et bellábunt advérsum te, et non prævalébunt: quia ego tecum sum, ait Dóminus, ut líberem te.\par
\switchcolumn
\EN
\noindent Thou therefore gird up thy loins, and arise, and speak to them all that I command thee. Be not afraid at their presence: for I will make thee not to fear their countenance. For behold I have made thee this day a fortified city, and a pillar of iron, and a wall of brass, over all the land, to the kings of Juda, to the princes thereof, and to the priests, and to the people of the land. And they shall fight against thee, and shall not prevail: for I am with thee, saith the Lord, to deliver thee.\par
\switchcolumn*

\LA
\begin{Resp}\Rsign Priúsquam te formárem in útero, novi te: et ántequam exíres de ventre, sanctificávi te, \Star{} Et prophétam in géntibus dedi te.\end{Resp}
\begin{Resp}\Vsign Vir diléctus a Deo, et homínibus honorátus est.\end{Resp}
\begin{Resp}\Rsign Et prophétam in géntibus dedi te.\end{Resp}
\begin{Resp}\Vsign Glória Patri, et Fílio, \Star{} et Spirítui Sancto.\end{Resp}
\begin{Resp}\Rsign Et prophétam in géntibus dedi te.\end{Resp}
\switchcolumn
\EN
\begin{Resp}\Rsign Before I formed thee in the belly I knew thee; and before thou camest forth out of the womb I sanctified thee; \Star{} And I ordained thee a Prophet unto the nations.\end{Resp}
\begin{Resp}\Vsign A man beloved of God and men, he was had in honour.\end{Resp}
\begin{Resp}\Rsign And I ordained thee a Prophet unto the nations.\end{Resp}
\begin{Resp}\Vsign Glory be to the Father, and to the Son, \Star{} and to the Holy Ghost.\end{Resp}
\begin{Resp}\Rsign And I ordained thee a Prophet unto the nations.\end{Resp}
\switchcolumn*

\LA
\LessonHead{Lectio IV}
\switchcolumn
\EN
\LessonHead{Lesson IV}
\switchcolumn*

\LA
\LessonTitle{Sermo sancti Augustíni Epíscopi}
\switchcolumn
\EN
\LessonTitle{From the Sermons of St. Augustine, Bishop of Hippo.}
\switchcolumn*

\LA
\Source{Sermo 20 de Sanctis}
\switchcolumn
\EN
\Source{20th on the Saints.}
\switchcolumn*

\LA
\noindent Post illum sacrosánctum Dómini natális diem, nullíus hóminum nativitátem légimus celebrári, nisi solíus beáti Joánnis Baptístæ. In áliis Sanctis et eléctis Dei nóvimus illum diem coli, quo illos, post consummatiónem labórum et devíctum triumphatúmque mundum, in perpétuas æternitátes præsens hæc vita partúriit. In áliis consummáta últimi diéi mérita celebrántur: in hoc étiam prima dies, et ipsa étiam hóminis inítia consecrántur; pro hac absque dúbio causa, quia per hunc Dóminus advéntum suum, ne súbito hómines insperátum non agnóscerent, vóluit esse testátum. Joánnes autem figúra fuit véteris Testaménti, et in se formam prǽtulit legis; et ídeo Joánnes prænuntiávit Salvatórem, sicut lex grátiam præcucúrrit.\par
\switchcolumn
\EN
\noindent Outside the most holy Birthday of the Lord, we find celebrated in the Gospel the birth of only one other, namely, that of the blessed Baptist, John. As regardeth all others among God's holy and chosen ones we know that that day is observed whereon, with their work finished, and the world conquered and finally trampled down, they were born from this into a better life, even one of everlasting blessedness. In others is honoured the crowning of the struggle on their last day of dying life, but in John is honoured the first day; in him the very beginning is found hallowed. And the reason of this is, without doubt, because he was sent from God to bear witness to the coming of the Light, lest when It came It might take the darkness by surprise, and the darkness might not comprehend It. Now, John was a figure of the Old Testament, and showed in his own person a typical embodiment of the Law; and therefore John heralded beforehand the coming of the Saviour, even as the Law was our schoolmaster to bring us to the grace of Christ. (Gal. iii. 24.)\par
\switchcolumn*

\LA
\begin{Resp}\Rsign Descéndit Angelus Dómini ad Zacharíam dicens: Accipe púerum in senectúte tua: \Star{} Et habébit nomen Joánnes Baptísta.\end{Resp}
\begin{Resp}\Vsign Iste puer magnus coram Dómino: nam et manus ejus cum ipso est.\end{Resp}
\begin{Resp}\Rsign Et habébit nomen Joánnes Baptísta.\end{Resp}
\switchcolumn
\EN
\begin{Resp}\Rsign The Angel of the Lord was sent down unto Zacharias, saying: Thou shalt beget a son in thine old age; \Star{} and his name shall be called John the Baptist.\end{Resp}
\begin{Resp}\Vsign The child shall be great in the sight of the Lord; for His hand is with him.\end{Resp}
\begin{Resp}\Rsign And his name shall be called John the Baptist.\end{Resp}
\switchcolumn*

\LA
\LessonHead{Lectio V}
\switchcolumn
\EN
\LessonHead{Lesson V}
\switchcolumn*

\LA
\noindent Quod autem nondum natus de secréto matérni úteri prophetávit, et expers lucis jam testis est veritátis; hoc est intellegéndum, quod latens sub velámine et carne lítteræ, et Redemptórem mundo spíritu prædicávit, et nobis Dóminum nostrum de quodam legis útero proclamávit. Ergo quia Judǽi erravérunt a ventre, id est, a lege quæ a Christo grávida erat, erravérunt a ventre, locúti sunt falsa; ídeo hic venit in testimónium, ut testimónium perhibéret de lúmine.\par
\switchcolumn
\EN
\noindent But as touching this, that he prophesied while yet in the hidden depths of his mother's womb, and while himself lightless bore testimony to the truth, we are to understand it as a figure how that while himself wrapped round with the veil and carnal ordinances of the letter, he by the spirit preached unto the world a Redeemer, and testified that Jesus is our Lord even while for himself, working under the law, the birth of the new dispensation was still in the womb of the future, and not come to day. The Jews were estranged from the womb, that is from the Law, that womb heavy with the Christ That was to be; they went astray from the belly, speaking lies, (Ps. lvii. 4;) and therefore John came for a witness, to bear witness of the Light, that all men through him might believe. (John i. 7.)\par
\switchcolumn*

\LA
\begin{Resp}\Rsign Hic est præcúrsor diléctus, et lucérna lucens ante Dóminum: \Star{} Ipse est enim Joánnes, qui viam Dómino præparávit in erémo; sed et Agnum Dei demonstrávit, et illuminávit mentes hóminum.\end{Resp}
\begin{Resp}\Vsign Ipse præíbit ante illum in spíritu et virtúte Elíæ.\end{Resp}
\begin{Resp}\Rsign Ipse est enim Joánnes, qui viam Dómino præparávit in erémo; sed et Agnum Dei demonstrávit, et illuminávit mentes hóminum.\end{Resp}
\switchcolumn
\EN
\begin{Resp}\Rsign This is the well-beloved Fore-runner, a burning and a shining light before the Lord. \Star{} For even this is John, who made straight in the desert an highway for our God. This also is he which bare witness, saying: Behold the Lamb of God! And his light lightened the minds of men.\end{Resp}
\begin{Resp}\Vsign He shall go before Him in the spirit and power of Elias.\end{Resp}
\begin{Resp}\Rsign For even this is John, who made straight in the desert an highway for our God. This also is he which bare witness, saying Behold the Lamb of God! And his light lightened the minds of men.\end{Resp}
\switchcolumn*

\LA
\LessonHead{Lectio VI}
\switchcolumn
\EN
\LessonHead{Lesson VI}
\switchcolumn*

\LA
\noindent Quod autem Joánnes in cárcere constitútus ad Christum discípulos suos órdinat; lex ad Evangélium transmíttit. Quæ lex juxta typum Joánnis, quasi ignorántiæ clausa cárcere, in obscúro et in occúlto jacébat, et Judáica cæcitáte sensus intra lítteram tenebátur inclúsus. De hoc beátus Evangelísta prolóquitur: Ille erat lucérna ardens, id est, Spíritus Sancti igne succénsus, ut mundo ignorántiæ nocte possésso lumen salútis osténderet, et quasi inter densíssimas delictórum ténebras splendidíssimus justítiæ solem lucis suæ rádio demonstráret, et de seípso dicens: Ego vox clamántis in desérto.\par
\switchcolumn
\EN
\noindent But as for this, that when John had heard in the prison the works of Christ, he sent two of his disciples (Matth. xi. 2); this is the Law sending to the Gospel. For John here was a figure of the Law, imprisoned in ignorance, lying in the dark, and in a hidden place, and he was fettered through Jewish misunderstanding within the bonds of the letter. But of him was it said, as is written in the Blessed Evangelist He was a burning and a shining light that is to say, that, when the whole world was wrapt in the night of ignorance, this Saint was kindled by the fire of the Holy Ghost, to show before men the light of salvation, and at the hour of the thickest darkness of sin, appeared like a bright morning star to herald the rising of that Sun so right gloriously radiant, the Son of righteousness, Christ our Lord. And this is why John said of himself: I am the voice of one crying in the wilderness, Make straight the way of the Lord. (John i. 15.)\par
\switchcolumn*

\LA
\begin{Resp}\Rsign Innuébant patri ejus quem vellet vocári eum: et póstulans pugillárem, scripsit dicens: \Star{} Joánnes est nomen ejus.\end{Resp}
\begin{Resp}\Vsign Apértum est os Zacharíæ, et prophetávit dicens.\end{Resp}
\begin{Resp}\Rsign Joánnes est nomen ejus.\end{Resp}
\begin{Resp}\Vsign Glória Patri, et Fílio, \Star{} et Spirítui Sancto.\end{Resp}
\begin{Resp}\Rsign Joánnes est nomen ejus.\end{Resp}
\switchcolumn
\EN
\begin{Resp}\Rsign They made signs to his father, how he would have him called; and he asked for a writing-table, and wrote, saying: \Star{} His name is John.\end{Resp}
\begin{Resp}\Vsign The mouth of Zacharias was opened, and he prophesied, saying:\end{Resp}
\begin{Resp}\Rsign His name is John.\end{Resp}
\begin{Resp}\Vsign Glory be to the Father, and to the Son, \Star{} and to the Holy Ghost.\end{Resp}
\begin{Resp}\Rsign His name is John.\end{Resp}
\switchcolumn*

\LA
\LessonHead{Lectio VII}
\switchcolumn
\EN
\LessonHead{Lesson VII}
\switchcolumn*

\LA
\LessonTitle{Léctio sancti Evangélii secúndum Lucam}
\switchcolumn
\EN
\LessonTitle{From the Holy Gospel according to Luke}
\switchcolumn*

\LA
\Cite{Luc 1:57-69}
\switchcolumn
\EN
\Cite{Luke 1:57-69}
\switchcolumn*

\LA
\noindent Elísabeth implétum est tempus pariéndi, et péperit fílium. Et audiérunt vicíni et cognáti ejus quia magnificávit Dóminus misericórdiam suam cum illa, et congratulabántur ei. Et réliqua.\par
\switchcolumn
\EN
\noindent Elisabeth's fullness of time came that she should deliver, and she brought forth a son. And her neighbours and her cousins heard how the Lord had showed great mercy upon her, and they rejoiced with her. And so on.\par
\switchcolumn*

\LA
\LessonTitle{Homilía sancti Ambrósii Epíscopi}
\switchcolumn
\EN
\LessonTitle{Homily by St. Ambrose, Bishop of Milan.}
\switchcolumn*

\LA
\Source{Liber 2 Comm. in Lucæ cap. 1, ante finem}
\switchcolumn
\EN
\Source{Bk. ii. Comm. on Luke i.}
\switchcolumn*

\LA
\noindent Péperit fílium Elísabeth, et congratulabántur vicíni. Habet Sanctórum edítio lætítiam plurimórum, quia commúne est bonum; justítia enim commúnis est virtus. Et ídeo in ortu justi futúræ vitæ insígne præmíttitur, et grátia secutúræ virtútis exsultatióne vicinórum præfiguránte signátur. Pulchre autem tempus, quo fuit in útero Prophéta, descríbitur, ne Maríæ præséntia taceátur; sed tempus silétur infántiæ, eo quia infántiæ impediménta nescívit. Et ídeo in Evangélio nihil super eo légimus, nisi ortum ejus, et oráculum, exsultatiónem in útero, vocem in desérto.\par
\switchcolumn
\EN
\noindent Elizabeth's full time came that she should be delivered, and she brought forth a son. And her neighbours rejoiced with her. The birth of a Saint is a joy for many, for it is a good to all. Righteousness is an help to all, and therefore when a righteous man is born it is an heralding of his life, which is still to come, that the helpful excellency of his future should be hailed by the, as it were, prophetic joy of the neighbours. It is well that we should be told concerning the prophet, while he was yet in the womb, that we may know how that Mary was there; but we hear nothing of his childhood, because, we know that it was safe and strong through the nearness of the Lord, Himself then in that womb which was free from the sorrows of pregnancy. And therefore we read in the Gospel nothing touching him save his coming, the annunciation thereof to his father, the leap which he gave in the womb, and his crying in the wilderness.\par
\switchcolumn*

\LA
\begin{Resp}\Rsign Præcúrsor Dómini venit, de quo ipse testátur: \Star{} Nullus major inter natos mulíerum Joánne Baptísta.\end{Resp}
\begin{Resp}\Vsign Hic est enim prophéta, et plus quam prophéta, de quo Salvátor ait.\end{Resp}
\begin{Resp}\Rsign Nullus major inter natos mulíerum Joánne Baptísta.\end{Resp}
\switchcolumn
\EN
\begin{Resp}\Rsign The Fore-runner of the Lord cometh, to whom He Himself bare witness, saying: \Star{} Among them that are born of women there hath not risen a greater than John the Baptist.\end{Resp}
\begin{Resp}\Vsign A Prophet? Yea, and much more than a Prophet. This is he of whom the Saviour saith:\end{Resp}
\begin{Resp}\Rsign Among them that are born of women there hath not risen a greater than John the Baptist.\end{Resp}
\switchcolumn*

\LA
\LessonHead{Lectio VIII}
\switchcolumn
\EN
\LessonHead{Lesson VIII}
\switchcolumn*

\LA
\noindent Neque enim ullam infántiæ sensit ætátem, qui supra natúram, supra ætátem in útero pósitus matris, a mensúra cœpit ætátis plenitúdinis Christi. Mire sanctus Evangelísta præmitténdum putávit, quod plúrimi infántem patris nómine Zacharíam appellándum putáverint; ut advértas matri non nomen alicújus displicuísse degéneris, sed id Sancto infúsum Spíritu, quod ab Angelo ante Zacharíæ fúerat prænuntiátum. Et quidem ille mutus intimáre vocábulum fílii nequívit uxóri; sed per prophetíam Elísabeth dídicit, quod non didícerat a maríto.\par
\switchcolumn
\EN
\noindent It was not for him to feel childishness, who beyond all use of nature or of his age, when as yet he lay in his mother's womb, leapt at once unto the measure of the stature of the fullness of Christ. (Eph. iv. 13.) It is strange how that the Holy Evangelist hath judged meet to tell us that they thought to call the child Zacharias, after the name of his father, that thou mayest notice that the mother would have none of the names whereby their kindred were called, but only that name which the Holy Ghost had dictated, and which the Angel had told before unto Zacharias. The dumb man had certainly not been able to tell his wife by what name to call the child, and Elizabeth must needs have learnt by revelation what she could not have heard from her husband.\par
\switchcolumn*

\LA
\begin{Resp}\Rsign Gábriel Angelus appáruit Zacharíæ dicens: Nascétur tibi fílius, nomen ejus Joánnes vocábitur: \Star{} Et multi in nativitáte ejus gaudébunt.\end{Resp}
\begin{Resp}\Vsign Erit enim magnus coram Dómino, vinum et síceram non bibet.\end{Resp}
\begin{Resp}\Rsign Et multi in nativitáte ejus gaudébunt.\end{Resp}
\begin{Resp}\Vsign Glória Patri, et Fílio, \Star{} et Spirítui Sancto.\end{Resp}
\begin{Resp}\Rsign Et multi in nativitáte ejus gaudébunt.\end{Resp}
\switchcolumn
\EN
\begin{Resp}\Rsign The Angel Gabriel appeared unto Zacharias, and said thy wife Elizabeth shall bear thee a son, and thou shalt call his name John; \Star{} And many shall rejoice at his birth.\end{Resp}
\begin{Resp}\Vsign For he shall be great in the sight of the Lord, and shall drink neither wine nor strong drink.\end{Resp}
\begin{Resp}\Rsign And many shall rejoice at his birth.\end{Resp}
\begin{Resp}\Vsign Glory be to the Father, and to the Son, \Star{} and to the Holy Ghost.\end{Resp}
\begin{Resp}\Rsign And many shall rejoice at his birth.\end{Resp}
\switchcolumn*

\LA
\LessonHead{Lectio IX}
\switchcolumn
\EN
\LessonHead{Lesson IX}
\switchcolumn*

\LA
\noindent Joánnes est, inquit, nomen ejus; hoc est, non nos ei nomen impónimus, qui jam a Deo nomen accépit. Habet vocábulum suum, quod agnóvimus, non quod elégimus. Habent hoc mérita Sanctórum, ut a Deo nomen accípiant. Sic Jacob Israël dícitur, quia Deum vidit. Sic Dóminus noster Jesus nominátus est, ántequam natus; cui non Angelus, sed Pater nomen impósuit. Vides Angelos quæ audíerint, non quæ usurpáverint, nuntiáre. Nec miréris, si nomen múlier, quod non audívit, asséruit; quando Spíritus ei Sanctus, qui Angelo mandáverat, revelávit.\par
\switchcolumn
\EN
\noindent His name is John; that is, it is not for us to choose a name now for him to whom God hath given a name already. He hath a name, which we know, but it is not one of our choosing. To receive a name from God is one of the honours of the Saints. Thus was it that Jacob's name was no more called Jacob but Israel, because he saw God face to face. (Gen. xxxii. 28.) Thus was it that our Lord Jesus was named before He was born, with a name not given by an Angel, but by the Father. Thou seest that Angels tell that which they have been bidden to tell, not matters of their own choosing. Nor oughtest thou to wonder that Elizabeth named a name which she had not heard, since it had been revealed to her by the same Holy Ghost Who had commanded the Angel to tell it.\par
\switchcolumn*

\LA
\TeDeum{Te Deum}
\switchcolumn
\EN
\TeDeum{Te Deum}
\switchcolumn*

\end{paracol}

\par\needspace{10\baselineskip}\addvspace{1.2em}{\centering\small\textsc{Die 25 Junii} · \textsc{25 June}\par}
\Office{S. Gulielmi Abbatis}{St. William, Abbot}{Duplex}
\addcontentsline{toc}{subsection}{Die 25 Junii · S. Gulielmi Abbatis \textit{— St. William, Abbot}}
\begin{paracol}{2}
\LA
\RefLine{Lectiones I–III: de Scriptura occurrente.}
\switchcolumn
\EN
\RefLine{Lessons I–III: from the Scripture of the day.}
\switchcolumn*

\LA
\LessonHead{Lectio IV}
\switchcolumn
\EN
\LessonHead{Lesson IV}
\switchcolumn*

\LA
\noindent Guliélmus nobílibus paréntibus Vercéllis in Insúbria natus, vix quartum décimum ætátis annum expléverat, cum miro quodam pietátis ardóre flagrans, Compostellánam peregrinatiónem ad celebérrimum sancti Jacóbi templum aggréssus est, Quod iter una amíctus túnica, ac dúplici férreo círculo præcínctus, nudísque pédibus prosecútus, aspérrima frígoris et æstus, famis et sitis summo cum vitæ discrímine perpéssus est incómmoda. Revérsus in Itáliam, novam ad sanctum Dómini sepúlcrum peregrinatiónem molítur: sed quo minus propósitum exsequátur, vária atque gravíssima intercédunt impediménta, divíno númine ad altióra, et sanctióra religiósam júvenis índolem retrahénte. Porro in solículo monte biénnium inter assíduas preces, vigílias, chaméunias, et jejúnia commorátus, divína subníxus ope, cæco lumen restítuit. Cujus miráculi fama percrebrescénte, jam Guliélmus latére non póterat: quare íterum Jerosólymam cógitat, et álacris se itíneri commíttit.\par
\switchcolumn
\EN
\noindent This William was born of noble parents at Vercelli in Lombardy. He was but little over fourteen years of age, when, impelled by a strange earnestness for holiness, he undertook a pilgrimage to Compostella, to the far-famed Church of St. James. This journey he made clad in a single garment, wearing an iron girdle wound two-fold round his body, and with bare feet. He accomplished his object under the severest hardships of cold and heat, hunger and thirst, and at the great danger of his life. After his return to Italy, he undertook a new pilgrimage, this time to the Holy Sepulchre of the Lord. But in the way of fulfilling this, there arose diverse and most grievous obstacles, whereby the hand of God drew the lad to the higher and holier life of a monk. He dwelt in the town of Monte Solicolo for two years, which he passed in constant prayer, watching, sleeping upon the ground, and fasting. At the end of this time, the power of God made him the mean to restore a blind man to sight. The fame of this miracle became so noised abroad, that William could no longer remain unknown. His thoughts turned again towards Jerusalem, and he again entered cheerfully on the journey.\par
\switchcolumn*

\LA
\begin{Resp}\Rsign Honéstum fecit illum Dóminus, et custodívit eum ab inimícis, et a seductóribus tutávit illum: \Star{} Et dedit illi claritátem ætérnam.\end{Resp}
\begin{Resp}\Vsign Justum dedúxit Dóminus per vias rectas, et osténdit illi regnum Dei.\end{Resp}
\begin{Resp}\Rsign Et dedit illi claritátem ætérnam.\end{Resp}
\switchcolumn
\EN
\begin{Resp}\Rsign The Lord made him honourable, and defended him from his enemies, and kept him safe from those that lay in wait for him; \Star{} And gave him perpetual glory.\end{Resp}
\begin{Resp}\Vsign He went down with him into the pit, and left him not in bonds.\end{Resp}
\begin{Resp}\Rsign And gave him perpetual glory.\end{Resp}
\switchcolumn*

\LA
\LessonHead{Lectio V}
\switchcolumn
\EN
\LessonHead{Lesson V}
\switchcolumn*

\LA
\noindent Dei autem mónitu, qui eídem appáruit, a propósito revocátur, utílior ac fructuósior tam apud Italos, quam apud éxteras natiónes futúrus. Tum monastérium in Virgiliáni montis cacúmine, quod deínde Vírginis est appellátum, loco áspero, et inaccésso miránda exædíficat celeritáte. Sócios deínde viros, et religiósos ascíscit, eósque ad vivéndi normam evangélicis præcéptis et consíliis summópere accommodátam tum certis légibus ex beáti Benedícti institútis magna ex parte desúmptis, tum verbo et sanctíssimæ vitæ exémplis infórmat.\par
\switchcolumn
\EN
\noindent He was again hindered by a vision from God, and remained among the Italians to be more useful, and to bring forth more fruit than he would have done among strangers. With extraordinary speed, he built a monastery upon the summit of Monte Vergiliano, ever since named Monte Vergine, (between Nola and Benevento.) Thither he called around him, as his comrades, devout men, and schooled them into a way of life most closely following the commands and counsels of the Gospel, in great part by a rule taken from the constitutions of Blessed Benedict, and supplemented by his own words, and the example of his own holy life.\par
\switchcolumn*

\LA
\begin{Resp}\Rsign Amávit eum Dóminus, et ornávit eum: stolam glóriæ índuit eum, \Star{} Et ad portas paradísi coronávit eum.\end{Resp}
\begin{Resp}\Vsign Índuit eum Dóminus lorícam fídei, et ornávit eum.\end{Resp}
\begin{Resp}\Rsign Et ad portas paradísi coronávit eum.\end{Resp}
\switchcolumn
\EN
\begin{Resp}\Rsign The Lord loved him and beautified him; He clothed him with a robe of glory, \Star{} And crowned him at the gates of Paradise.\end{Resp}
\begin{Resp}\Vsign The Lord hath put on him the breast-plate of faith, and hath adorned him.\end{Resp}
\begin{Resp}\Rsign And crowned him at the gates of Paradise.\end{Resp}
\switchcolumn*

\LA
\LessonHead{Lectio VI}
\switchcolumn
\EN
\LessonHead{Lesson VI}
\switchcolumn*

\LA
\noindent Aliis deínde monastériis eréctis, clárior in dies Guliélmi facta sánctitas multos ad eum úndique viros perdúcit, sanctitátis odóre, et miraculórum fama alléctos. Nam muti loquélam, surdi audítum, áridi vigórem, varióque et immedicábili morbo laborántes, sanitátem ipsíus intercessióne recepérunt. Aquam in vinum convértit, aliáque complúra mirabília patrávit: inter quæ illud non siléndum, quod, muliércula ad ejus castitátem tentándam missa, in ardéntibus prunis humi stratis illǽsum se volutávit. De qua re cértior factus Rogérius Neápolis rex, in summam viri Dei veneratiónem addúcitur. Demum témpore sui óbitus regi, aliísque prænuntiáto, innúmeris virtútibus et miráculis clarus obdormívit in Dómino, anno salútis millésimo centésimo quadragésimo secúndo.\par
\switchcolumn
\EN
\noindent As other monasteries were raised, the holy life of William became more known day by day, and brought men to him from all quarters, drawn by the sweet savour of his godliness, and the fame of his miracles. At his prayers the dumb spake, the deaf heard, the withered were strengthened, and they that suffered under diverse and incurable diseases received health. He turned water into wine, and openly worked many other miracles. Among all these things it must be told that a wretched woman sought him to lure him to impurity, but he raked hot embers out upon the floor and cast himself down upon them, and wallowed among them, and escaped unhurt. When this thing came to the knowledge of Roger I., King of Naples, it roused in him the highest reverence for the man of God. At the last, after foretelling his own death to the King and to others, and full of good works and miracles, he fell asleep in the Lord, in the year of salvation 1142.\par
\switchcolumn*

\LA
\begin{Resp}\Rsign Iste homo perfécit ómnia quæ locútus est ei Deus, et dixit ad eum: Ingrédere in réquiem meam: \Star{} Quia te vidi justum coram me ex ómnibus géntibus.\end{Resp}
\begin{Resp}\Vsign Iste est, qui contémpsit vitam mundi, et pervénit ad cæléstia regna.\end{Resp}
\begin{Resp}\Rsign Quia te vidi justum coram me ex ómnibus géntibus.\end{Resp}
\begin{Resp}\Vsign Glória Patri, et Fílio, \Star{} et Spirítui Sancto.\end{Resp}
\begin{Resp}\Rsign Quia te vidi justum coram me ex ómnibus géntibus.\end{Resp}
\switchcolumn
\EN
\begin{Resp}\Rsign This is he which did according unto all that God commanded him; and God said unto him: Enter thou into My rest, \Star{} For thee have I seen righteous before Me among all people.\end{Resp}
\begin{Resp}\Vsign This is he which loved not his life in this world, and is come unto an everlasting kingdom.\end{Resp}
\begin{Resp}\Rsign For thee have I seen righteous before Me among all people.\end{Resp}
\begin{Resp}\Vsign Glory be to the Father, and to the Son, \Star{} and to the Holy Ghost.\end{Resp}
\begin{Resp}\Rsign For thee have I seen righteous before Me among all people.\end{Resp}
\switchcolumn*

\LA
\LessonHead{Lectio VII}
\switchcolumn
\EN
\LessonHead{Lesson VII}
\switchcolumn*

\LA
\LessonTitle{Léctio sancti Evangélii secúndum Matthǽum}
\switchcolumn
\EN
\LessonTitle{The Lesson is taken from the Holy Gospel according to Matthew}
\switchcolumn*

\LA
\Cite{Matt 19:27-29}
\switchcolumn
\EN
\Cite{Matt 19:27-29}
\switchcolumn*

\LA
\noindent In illo tempore: Dixit Petrus ad Jesum: Ecce nos reliquimus omnia, et secuti sumus te: quid ergo erit nobis? Et reliqua.\par
\switchcolumn
\EN
\noindent At that time: Peter said unto Jesus: Behold, we have forsaken all, and followed Thee: what shall we have therefore? And so on.\par
\switchcolumn*

\LA
\LessonTitle{De Homilía venerábilis Bedæ Presbýteri}
\switchcolumn
\EN
\LessonTitle{Homily by the Venerable Bede, Priest (at Jarrow and Doctor of the Church).}
\switchcolumn*

\LA
\Source{In Natali S. Benedicti}
\switchcolumn
\EN
\Source{For St Benedict’s Birthday.}
\switchcolumn*

\LA
\noindent Duo sunt órdines electórum in judício futúri: unus judicántium cum Dómino, de quibus hoc loco mémorat, qui reliquérunt ómnia, et secúti sunt illum: alius judicandórum a Dómino, qui non quidem ómnia sua páriter reliquérunt, sed de his tamen quæ habébant, quotidiánas dare eleemósynas Christi paupéribus curábant: unde et auditúri sunt in judício: Veníte, benedícti Patris mei, possidíte præparátum vobis regnum a constitutióne mundi. Esurívi enim, et dedístis mihi manducáre: sitívi, et dedístis mihi bíbere.\par
\switchcolumn
\EN
\noindent In the judgment to come, the elect will be in two classes. One class are they who have forsaken all, and followed the Lord: and these shall judge along with Him. The other class are they who have not equally forsaken all that they had, but who have been careful daily to give alms of their goods to the poor of Christ: these shall be the subjects of judgment, and these are they who shall then hear these words: “Come, ye blessed of My Father, inherit the kingdom prepared for you from the foundation of the world: for I was an hungered, and ye gave Me meat: I was thirsty, and ye gave Me drink.” (Matth. xxv. 34, 35.)\par
\switchcolumn*

\LA
\begin{Resp}\Rsign Iste est qui ante Deum magnas virtútes operátus est, et de omni corde suo laudávit Dóminum: \Star{} Ipse intercédat pro peccátis ómnium populórum.\end{Resp}
\begin{Resp}\Vsign Ecce homo sine queréla, verus Dei cultor, ábstinens se ab omni ópere malo, et pérmanens in innocéntia sua.\end{Resp}
\begin{Resp}\Rsign Ipse intercédat pro peccátis ómnium populórum.\end{Resp}
\switchcolumn
\EN
\begin{Resp}\Rsign This is he which wrought great wonders before God, and praised the Lord with all his heart. \Star{} May he pray for all people, that their sins may be forgiven unto them.\end{Resp}
\begin{Resp}\Vsign Behold a man without blame, a worshipper of God in truth, keeping himself clean from every evil work, and abiding still in his innocency.\end{Resp}
\begin{Resp}\Rsign May he pray for all people, that their sins may be forgiven unto them.\end{Resp}
\switchcolumn*

\LA
\LessonHead{Lectio VIII}
\switchcolumn
\EN
\LessonHead{Lesson VIII}
\switchcolumn*

\LA
\noindent Sed et reprobórum duos ibi futúros órdines Dómino narránte comperímus: unum eórum, qui fídei christiánæ mystériis initiáti, ópera fídei exercére contémnunt, quibus in judício testátur: Discédite a me maledícti in ignem aetérnum, qui præparátus est diábolo, et ángelis ejus: esurívi enim, et non dedístis mihi manducáre. Alterum eórum, qui fidem et mystéria Christi vel numquam suscepére, vel suscéptam per apostasíam deseruére: de quibus dicit: Qui autem non credit, jam judicatus est: quia non credit in nómine unigéniti Fílii Dei.\par
\switchcolumn
\EN
\noindent Of the reprobate also we gather, from the words of the Lord, that there will be two classes. One class are they who, being made partakers in the mystery of Christian faith, have neglected to show their faith by their works: these are they to whom it will be said at the judgment: “Depart from Me, ye cursed, into everlasting fire, prepared for the devil and his angels: for I was an-hungered, and ye gave Me no meat.” The other class are they who either have never received the faith and mysteries of Christ, or who, having received, have apostatised, and abandoned it: and touching these it is said: “But he that believeth not is condemned already, because he hath not believed in the name of the only-begotten Son of God.” (John iii. 18.)\par
\switchcolumn*

\LA
\begin{Resp}\Rsign Sint lumbi vestri præcíncti, et lucérnæ ardéntes in mánibus vestris: \Star{} Et vos símiles homínibus exspectántibus dóminum suum, quando revertátur a núptiis.\end{Resp}
\begin{Resp}\Vsign Vigiláte ergo, quia nescítis qua hora Dóminus vester ventúrus sit.\end{Resp}
\begin{Resp}\Rsign Et vos símiles homínibus exspectántibus dóminum suum, quando revertátur a núptiis.\end{Resp}
\begin{Resp}\Vsign Glória Patri, et Fílio, \Star{} et Spirítui Sancto.\end{Resp}
\begin{Resp}\Rsign Et vos símiles homínibus exspectántibus dóminum suum, quando revertátur a núptiis.\end{Resp}
\switchcolumn
\EN
\begin{Resp}\Rsign Let your loins be girded about, and your lights burning; \Star{} And ye yourselves like unto men that wait for their lord, when he will return from the wedding.\end{Resp}
\begin{Resp}\Vsign Watch therefore, for ye know not what hour your Lord doth come.\end{Resp}
\begin{Resp}\Rsign And ye yourselves like unto men that wait for their lord, when he will return from the wedding.\end{Resp}
\begin{Resp}\Vsign Glory be to the Father, and to the Son, \Star{} and to the Holy Ghost.\end{Resp}
\begin{Resp}\Rsign And ye yourselves like unto men that wait for their lord, when he will return from the wedding.\end{Resp}
\switchcolumn*

\LA
\LessonHead{Lectio IX}
\switchcolumn
\EN
\LessonHead{Lesson IX}
\switchcolumn*

\LA
\noindent Verum his cum timóre et pavóre débito paulísper commemorátis, ad lætíssima pótius Dómini et Salvatóris nostri promíssa convertámus audítum. Videámus quæ tantæ grátia pietátis: non ætérnae tantúmmodo vitæ præmia suis sequácibus, sed et præséntis múnera pollicétur exímia. Et omnis, inquit, qui relíquerit domum, vel fratres, aut soróres, aut patrem, aut matrem, aut uxórem, aut fílios, aut agros, propter nomen meum, céntuplum accípiet, et vitam ætérnam possidébit. Qui enim terrénis afféctibus sive possessiónibus pro Christi discipulátu renuntiáverit, quo plus in ejus amórem profécerit, eo plures invéniet, qui se intérno suscípere afféctu, et suis gáudeant sustentáre substántiis.\par
\switchcolumn
\EN
\noindent And now that we have touched for a moment, with fear and just dread, upon these things, let us rather turn our hearing to the right joyful promises of our Lord and Saviour. Let us look what His so great, beautiful, and fatherly love will give to such as follow Him; not the reward of life everlasting only, but gifts exceeding precious in this life also. “Every one,” saith He, that hath forsaken houses, or brethren, or sisters, or father, or mother, or wife, or children, or lands, for My Name’s sake, shall receive an hundredfold, and shall inherit everlasting life.” For every one that shall forsake earthly affections and goods, to go and be Christ’s disciple, the further he goeth on in Christ’s love, the more shall he find who will rejoice to give him a place in their hearts, and to minister to him of their substance.\par
\switchcolumn*

\LA
\TeDeum{Te Deum}
\switchcolumn
\EN
\TeDeum{Te Deum}
\switchcolumn*

\LA
\LessonHead{Lectio IX (pro commemoratione)}
\switchcolumn
\EN
\LessonHead{Lesson IX (when commemorated)}
\switchcolumn*

\LA
\LessonTitle{Commemoratio: S. Gulielmi Abbatis}
\switchcolumn
\EN
\LessonTitle{Commemoration of: St. William, Abbot}
\switchcolumn*

\LA
\noindent Guliélmus, nobílibus paréntibus Vercéllis natus, vix quartum décimum ætátis annum expléverat, cum Compostellánam peregrinatiónem miro spíritu pœniténtiæ ac pietátis ardóre perégit. Dein, nova peregrinatióne ad Christi Dómini sepúlcrum frustra tentáta, in solitário monte inter assíduas preces, vigílias et jejúnia per biénnium delítuit. Cum cæco lumen restituísset, hóminum existimatiónem fúgiens monastérium in Virgiliáni montis, quod deínde Vírginis est appellátum, loco áspero et inaccésso exædíficat; ibi sócios ascíscit, eósque certis légibus, ex beáti Benedícti institútis magna ex parte desúmptis, verbo et sanctíssimæ vitæ exémplis infórmat. Aliis deínde monastériis eréctis, clárior in dies Guliélmi facta sánctitas multos ad eum úndique viros perdúcit, frequéntium étiam miraculórum fama alléctos. Demum, témpore sui óbitus prænuntiáto, obdormívit in Dómino anno salútis millésimo centésimo quadragésimo secúndo.\par
\switchcolumn
\EN
\noindent William, born of noble parents at Vercelli, had scarcely finished his fourteenth year when he made a pilgrimage to Compostella in a wonderful spirit of penitence and devotional zeal. Then, having vainly attempted another, pilgrimage, to the tomb of Christ the Lord, he spent two years on a solitary mountain in constant prayer, in vigils and in fasting. Fleeing human renown after he had restored sight to a blind man, he built a monastery on Monte Virgiliáno, which was thereafter called Monte Vergine, in a wild and inaccessible spot. There companions joined him, and he formed them by fixed regulations taken largely from those of St. Benedict, by word, and by the example of a most holy life. Then he built other monasteries, and daily his fame as a holy man grew; so that many came to him from all parts, drawn by the report of his frequent miracles. Finally, having foretold the day of his death, he fell asleep in the Lord, in the year of salvation 1142.\par
\switchcolumn*

\LA
\TeDeum{Te Deum}
\switchcolumn
\EN
\TeDeum{Te Deum}
\switchcolumn*

\end{paracol}

\par\needspace{10\baselineskip}\addvspace{1.2em}{\centering\small\textsc{Die 25 Junii} · \textsc{25 June}\par}
\Office{Die secunda infra Octavam Nativitatis S. Joannis Baptistæ}{Second Day of the Octave of the Nativity of St. John the Baptist}{Semiduplex}
\addcontentsline{toc}{subsection}{Die 25 Junii · Die secunda infra Octavam Nativitatis S. Joannis Baptistæ \textit{— Second Day of the Octave of the Nativity of St. John the Baptist}}
\begin{paracol}{2}
\LA
\LessonHead{Lectio IV}
\switchcolumn
\EN
\LessonHead{Lesson IV}
\switchcolumn*

\LA

\switchcolumn
\EN
\Note{No English translation in the source files; the Latin is given.}
\switchcolumn*

\LA
\LessonTitle{Commemoratio: Die secunda infra Octavam Nativitatis S. Joannis Baptistæ}
\switchcolumn
\EN
\LessonTitle{Commemoration of: Second Day of the Octave of the Nativity of St. John the Baptist}
\switchcolumn*

\LA
\Source{Serm. 21 de Sanctis.}
\switchcolumn
\EN
\Source{Serm. 21 de Sanctis.}
\switchcolumn*

\LA
\noindent Sermo sanctii Augustíni Epíscopi\par
\switchcolumn
\EN
\noindent Sermo sanctii Augustíni Epíscopi\par
\switchcolumn*

\LA
Natálem sancti Joánnis, fratres caríssimi, hódie celebrámus: quod nulli umqua Sanctórum légimus fuísse concéssum. Solíus enim Dómini, et beáti Joánnis dies nativitátis in univérso mundo celebrátur et cólitur. Hunc enim stérilis péperit, illum virgo concépit: in Elísabeth sterílitas víncitur, in beáta María conceptiónis consuetúdo mutátur. Elísabeth virum cognoscéndo, fílium génuit: María Angelo crédidit, et concépit. Hóminem concépit Elísabeth, hóminem et María: sed Elísabeth solum hóminem, María Deum et hóminem.\par
\switchcolumn
\EN
Natálem sancti Joánnis, fratres caríssimi, hódie celebrámus: quod nulli umqua Sanctórum légimus fuísse concéssum. Solíus enim Dómini, et beáti Joánnis dies nativitátis in univérso mundo celebrátur et cólitur. Hunc enim stérilis péperit, illum virgo concépit: in Elísabeth sterílitas víncitur, in beáta María conceptiónis consuetúdo mutátur. Elísabeth virum cognoscéndo, fílium génuit: María Angelo crédidit, et concépit. Hóminem concépit Elísabeth, hóminem et María: sed Elísabeth solum hóminem, María Deum et hóminem.\par
\switchcolumn*

\LA
Magnus ígitur Joánnes, cujus magnitúdini étiam Salvátor testimónium perhibet, dicens: Non surréxit inter natos mulíerum major Joánne Baptísta. Præcéllit céteros, éminet univérsis: antecéllit Prophétas, supergréditur Patriárchas: et quisquis de mulíere natus est, inférior est Joánne. Dicet fortásse áliquis: Si inter natos mulíerum Joánnes major est, major est Salvatóre? Absit. Joánnes enim natus mulíeris, Christus autem vírginis natus est. Ille corruptíbilis úteri sínibus effúsus est, iste impollútæ vírginis flore progénitus.\par
\switchcolumn
\EN
Magnus ígitur Joánnes, cujus magnitúdini étiam Salvátor testimónium perhibet, dicens: Non surréxit inter natos mulíerum major Joánne Baptísta. Præcéllit céteros, éminet univérsis: antecéllit Prophétas, supergréditur Patriárchas: et quisquis de mulíere natus est, inférior est Joánne. Dicet fortásse áliquis: Si inter natos mulíerum Joánnes major est, major est Salvatóre? Absit. Joánnes enim natus mulíeris, Christus autem vírginis natus est. Ille corruptíbilis úteri sínibus effúsus est, iste impollútæ vírginis flore progénitus.\par
\switchcolumn*

\LA
Si ergo intellígimus mystérium, fratres mei, Joánnes homo est, Christus Deus est. Humiliétur homo, ut exaltétur Deus, secúndum illud quod de Dómino ipse Joánnes dixit: Illum opórtet créscere, me autem mínui. Ut humiliarétur homo, hódie natus est Joánnes, quo incípiunt decréscere dies: ut exaltétur Deus, eо die natus est Christus, quo incípiunt créscere dies. Magnum sacraméntum, fratres caríssimi. Ideo celebrámus natálem Joánnis sicut et Christi, quia et ipsa natívitas plena est mystério. Quo mystério, nisi humilitátis nostræ, sicut natívitas Christi plena est mystério altitúdinis nostræ?\par
\switchcolumn
\EN
Si ergo intellígimus mystérium, fratres mei, Joánnes homo est, Christus Deus est. Humiliétur homo, ut exaltétur Deus, secúndum illud quod de Dómino ipse Joánnes dixit: Illum opórtet créscere, me autem mínui. Ut humiliarétur homo, hódie natus est Joánnes, quo incípiunt decréscere dies: ut exaltétur Deus, eо die natus est Christus, quo incípiunt créscere dies. Magnum sacraméntum, fratres caríssimi. Ideo celebrámus natálem Joánnis sicut et Christi, quia et ipsa natívitas plena est mystério. Quo mystério, nisi humilitátis nostræ, sicut natívitas Christi plena est mystério altitúdinis nostræ?\par
\switchcolumn*

\LA
\TeDeum{Te Deum}
\switchcolumn
\EN
\TeDeum{Te Deum}
\switchcolumn*

\end{paracol}

\par\needspace{10\baselineskip}\addvspace{1.2em}{\centering\small\textsc{Die 26 Junii} · \textsc{26 June}\par}
\Office{Ss. Joannis et Pauli Martyrum}{Sts. John and Paul, Martyrs}{Duplex}
\addcontentsline{toc}{subsection}{Die 26 Junii · Ss. Joannis et Pauli Martyrum \textit{— Sts. John and Paul, Martyrs}}
\begin{paracol}{2}
\LA
\RefLine{Lectiones I–III: de Scriptura occurrente.}
\switchcolumn
\EN
\RefLine{Lessons I–III: from the Scripture of the day.}
\switchcolumn*

\LA
\LessonHead{Lectio IV}
\switchcolumn
\EN
\LessonHead{Lesson IV}
\switchcolumn*

\LA
\noindent Joánnes et Paulus fratres Románi, cum facultátibus a Constántia Constantíni fília, cui pie fidelitérque servíerant, sibi relíctis, Christi páuperes álerent; a Juliáno Apóstata in númerum familiárium suórum invitáti, líbere negavérunt se apud eum esse velle, qui a Jesu Christo defecísset. Quibus ille ad deliberándum decem dies præfínit ut nisi ad eam diem ei adhærére et Jovi sacrificáre constitúerint, sibi moriéndum esse certo sciant.\par
\switchcolumn
\EN
\noindent John and Paul were two Roman brethren, the godly and trustworthy servants of Constantia, daughter of Constantine. At her death they spent in feeding Christ's poor the property which she left them. Julian the Apostate asked them to enter his household, but they bravely answered that they would not be servants to one who had run away from the service of Jesus Christ. Julian gave them ten days to consider on their choice, whether, at the end of that time, they would cleave to him, and sacrifice to Jupiter, or most surely die.\par
\switchcolumn*

\LA
\begin{Resp}\Rsign Isti sunt duo viri misericórdiæ, qui assístunt ante Dóminum, \Star{} Dominatórem univérsæ terræ.\end{Resp}
\begin{Resp}\Vsign Isti sunt duæ olívæ, et duo candelábra lucéntia ante Dóminum.\end{Resp}
\begin{Resp}\Rsign Dominatórem univérsæ terræ.\end{Resp}
\switchcolumn
\EN
\begin{Resp}\Rsign These are two men full of mercy, who stand in the presence of the Lord, \Star{} The Lord of the whole earth.\end{Resp}
\begin{Resp}\Vsign These are two olive trees, and two candlesticks, giving light in the presence of the Lord.\end{Resp}
\begin{Resp}\Rsign The Lord of the whole earth.\end{Resp}
\switchcolumn*

\LA
\LessonHead{Lectio V}
\switchcolumn
\EN
\LessonHead{Lesson V}
\switchcolumn*

\LA
\noindent Illi intra id tempus réliqua sua bona distribuérunt paupéribus, quo expeditióres ad Dóminum migráre possent, et plures juvárent, a quibus in ætérna tabernácula reciperéntur. Die décima Terentiánus prætóriæ cohórtis præféctus, ad eos missus cum alláta Jovis effígie ut eam veneraréntur, imperatóris mandátum eis expónit, ut nisi Jovi cultum adhíbeant, moriántur. Qui, ut erant orántes respondérunt, se pro Christi fide, quem Deum mente et ore venerabántur, non dubitánter mortem subitúros.\par
\switchcolumn
\EN
\noindent This interval they spent in distributing to poor creatures all that remained of their goods, that they might be quite free to depart hence to the Lord, and so succoured many by whom they have long since been received into everlasting habitations. On the tenth day Terentian, Prefect of the Praetorian Cohort, was sent to them, bringing with him the image of Jupiter. He explained to them the command of the Emperor, that they should worship the said image or die. They were engaged in prayer, but answered him that for their loyalty to Christ, Whom their understanding acknowledged and their mouths confessed to be God, they felt no hesitation in choosing to suffer death.\par
\switchcolumn*

\LA
\begin{Resp}\Rsign Vidi conjúnctos viros, habéntes spléndidas vestes, et Angelus Dómini locútus est ad me, dicens: \Star{} Isti sunt viri sancti facti amíci Dei.\end{Resp}
\begin{Resp}\Vsign Vidi Angelum Dei fortem, volántem per médium cælum, voce magna clamántem et dicéntem.\end{Resp}
\begin{Resp}\Rsign Isti sunt viri sancti facti amíci Dei.\end{Resp}
\switchcolumn
\EN
\begin{Resp}\Rsign I saw men standing together clad in shining raiment; and the Angel of the Lord spake unto me, saying: \Star{} These men are holy, for they are the friends of God.\end{Resp}
\begin{Resp}\Vsign And I beheld a mighty Angel of God, flying through the midst of heaven, crying with a loud voice, and saying:\end{Resp}
\begin{Resp}\Rsign These men are holy, for they are the friends of God.\end{Resp}
\switchcolumn*

\LA
\LessonHead{Lectio VI}
\switchcolumn
\EN
\LessonHead{Lesson VI}
\switchcolumn*

\LA
\noindent At Terentiánus véritus, ne, si públice interficeréntur, pópulus commoverétur; domi, ubi tunc erant, abscíssis eórum capítibus, sexto Kaléndas Júlii secréto eos sepeliéndos curávit; rumorémque sparsit, Joánnem et Paulum in exsílium ejéctos esse. Verum eórum mors a spirítibus immúndis, qui multórum córpora vexábant, pervulgáta est; in quibus Terentiáni fílius, et ipse oppréssus a dǽmone, ad sepúlcrum Mártyrum perdúctus liberátus est. Quo miráculo et is in Christum crédidit, et ejus pater Terentiánus; a quo étiam horum beatórum Mártyrum vita scripta esse dícitur.\par
\switchcolumn
\EN
\noindent Terentian, to avoid the uproar, which might have been caused by their public execution, caused their heads to be cut off at home where they then were. They lifted up their last earthly testimony upon the 26th day of June, (in the year of our Lord 362.) They were privately buried, and a story set about that they had been sent into exile. The fact of their death was made generally known by the unclean spirits by whom the bodies of many were tormented, and among others that of Terentian's own son, who was possessed with a devil, and delivered by being brought to the grave of the Martyrs. By this miracle he was led to believe in Christ, and so likewise was his father Terentian, who is said to have been the writer of the life of these blessed Martyrs.\par
\switchcolumn*

\LA
\begin{Resp}\Rsign Tamquam aurum in fornáce probávit eléctos Dóminus, et quasi holocáusti hóstiam accépit illos: et in témpore erit respéctus illórum: \Star{} Quóniam donum et pax est eléctis Dei.\end{Resp}
\begin{Resp}\Vsign Qui confídunt in illum, intélligent veritátem, et fidéles in dilectióne acquiéscent illi.\end{Resp}
\begin{Resp}\Rsign Quóniam donum et pax est eléctis Dei.\end{Resp}
\begin{Resp}\Vsign Glória Patri, et Fílio, \Star{} et Spirítui Sancto.\end{Resp}
\begin{Resp}\Rsign Quóniam donum et pax est eléctis Dei.\end{Resp}
\switchcolumn
\EN
\begin{Resp}\Rsign As gold in the furnace hath the Lord tried His chosen ones, and received them as a burnt-offering, and yet a while, and they shall be regarded; \Star{} For the grace of God, and His peace, are with His chosen.\end{Resp}
\begin{Resp}\Vsign They that put their trust in Him shall understand the truth and such as be faithful in love shall abide with Him.\end{Resp}
\begin{Resp}\Rsign For the grace of God, and His peace, are with His chosen.\end{Resp}
\begin{Resp}\Vsign Glory be to the Father, and to the Son, \Star{} and to the Holy Ghost.\end{Resp}
\begin{Resp}\Rsign For the grace of God, and His peace, are with His chosen.\end{Resp}
\switchcolumn*

\LA
\LessonHead{Lectio VII}
\switchcolumn
\EN
\LessonHead{Lesson VII}
\switchcolumn*

\LA
\LessonTitle{Léctio sancti Evangélii secúndum Lucam}
\switchcolumn
\EN
\LessonTitle{From the Holy Gospel according to Luke}
\switchcolumn*

\LA
\Cite{Luc 12:1-8}
\switchcolumn
\EN
\Cite{Luke 12:1-8}
\switchcolumn*

\LA
\noindent In illo témpore: Dixit Jesus discípulis suis: Atténdite a fermento pharisæórum, quod est hypócrisis. Et réliqua.\par
\switchcolumn
\EN
\noindent At that time: Jesus said unto His disciples Beware of the leaven of the Pharisees, which is hypocrisy. And so on.\par
\switchcolumn*

\LA
\LessonTitle{Homilía sancti Bedæ Venerábilis Presbýteri}
\switchcolumn
\EN
\LessonTitle{Homily by the Venerable Bede, Priest at Jarrow and Doctor of the Church.}
\switchcolumn*

\LA
\Source{Lib. 4 in Lucam, cap. 12}
\switchcolumn
\EN
\Source{Bk. iv. on Luke, Cap. Hi.}
\switchcolumn*

\LA
\noindent De hoc ferménto Apóstolus præcépit: Itáque epulémur, non in ferménto véteri, neque in ferménto malítiæ et nequítiæ, sed in ázymis sinceritátis et veritátis. Nam, sicut módicum ferméntum totam farínæ, cui injícitur, massam corrúmpit, universámque mox conspersiónem suo sapóre commáculat; sic nimírum simulátio, cujus semel ánimum imbúerit, tota virtútum sinceritáte et veritáte fraudábit. Est ergo sensus: Atténdite ne æmulémini simulatóres; quia véniet profécto tempus, in quo et vestra virtus ómnibus, et eórum revelétur hypócrisis.\par
\switchcolumn
\EN
\noindent Touching this leaven the Apostle warneth us Therefore let us keep the feast, not with old leaven, neither with the leaven of malice and wickedness, but with the unleavened bread of sincerity and truth. (i Cor. v. 8.) For even as a little leaven doth infect the whole lump wherein it is put, and the savour thereof doth spread all abroad therein, so doth hypocrisy, when once it hath tainted the soul, drive out from it all sincerity and truth. The meaning, therefore, of this passage is this Beware, lest ye be as the hypocrites, for yet a little while, and all men shall see that ye are good, and they are evil.\par
\switchcolumn*

\LA
\begin{Resp}\Rsign Propter testaméntum Dómini, et leges patérnas, Sancti Dei perstitérunt in amóre fraternitátis: \Star{} Quia unus fuit semper spíritus in eis, et una fides.\end{Resp}
\begin{Resp}\Vsign Ecce quam bonum et quam jucúndum habitáre fratres in unum.\end{Resp}
\begin{Resp}\Rsign Quia unus fuit semper spíritus in eis, et una fides.\end{Resp}
\switchcolumn
\EN
\begin{Resp}\Rsign Because of the covenant of the Lord, and the laws of their fathers, the Saints of God abode in brotherly love, \Star{} For one spirit and one faith was ever in them.\end{Resp}
\begin{Resp}\Vsign Behold how good and how pleasant it is for brethren to dwell together in unity.\end{Resp}
\begin{Resp}\Rsign For one spirit and one faith was ever in them.\end{Resp}
\switchcolumn*

\LA
\LessonHead{Lectio VIII}
\switchcolumn
\EN
\LessonHead{Lesson VIII}
\switchcolumn*

\LA
\noindent Verum quod séquitur: Quóniam quæ in ténebris dixístis, in lúmine dicéntur; non solum in futúro, quando cuncta córdium abscóndita proferéntur ad lucem, sed et in præsénti témpore potest congruénter áccipi. Quóniam quæ inter ténebras quondam pressurárum carcerúmque umbras vel locúti vel passi sunt Apóstoli; nunc, clarificáta per orbem Ecclésia, lectis eórum áctibus, públice prædicántur. Ne terreámini ab his qui occídunt corpus. Si persecutóres Sanctórum, occísis corpóribus, non habent ámplius quid contra illos agant: ergo supervácua furunt insánia, qui mórtua Mártyrum membra, feris avibúsque discerpénda projíciunt, cum nequáquam omnipoténtiæ Dei, quin ea resuscitándo vivíficet, obsístere possint.\par
\switchcolumn
\EN
\noindent As touching what followeth, for there is nothing covered that shall not be revealed, neither hid, that shall not be known. Therefore, whatsoever ye have spoken in darkness shall be heard in the light. These words are true, not only as concerning the world which is to come, wherein the secrets of all hearts shall be made manifest, but even as concerning this present world, since now that which the Apostles spake and suffered in the darkness of persecution, and the gloom of dungeons, is, since that the Church is glorified, told of them for a memorial of them, wherever their acts are read throughout the whole world. Be not afraid of them that kill the body, for they that persecute the righteous, when they have killed the body, after that, have no more that they can do. Truly, it is a childish folly which maketh such men to cast the dead limbs of the martyrs to birds and beasts, while yet they have no strength to withstand the Almight of God, whereby He will surely quicken the same limbs and raise them up again.\par
\switchcolumn*

\LA
\begin{Resp}\Rsign Hæc est vera fratérnitas, quæ numquam pótuit violári certámine: qui effúso sánguine secúti sunt Dóminum: \Star{} Contemnéntes aulam régiam, pervenérunt ad regna cæléstia.\end{Resp}
\begin{Resp}\Vsign Ecce quam bonum et quam jucúndum habitáre fratres in unum.\end{Resp}
\begin{Resp}\Rsign Contemnéntes aulam régiam, pervenérunt ad regna cæléstia.\end{Resp}
\begin{Resp}\Vsign Glória Patri, et Fílio, \Star{} et Spirítui Sancto.\end{Resp}
\begin{Resp}\Rsign Contemnéntes aulam régiam, pervenérunt ad regna cæléstia.\end{Resp}
\switchcolumn
\EN
\begin{Resp}\Rsign Theirs is a brotherhood indeed, whose tie no storms availed to sever together they followed the Lord in the shedding of their blood. \Star{} Together they set at nought the Royal Palace; together they attained unto the kingdom of heaven.\end{Resp}
\begin{Resp}\Vsign Behold how good and how pleasant it is for brethren to dwell together in unity.\end{Resp}
\begin{Resp}\Rsign Together they set at nought the Royal Palace; together they attained unto the kingdom of heaven.\end{Resp}
\begin{Resp}\Vsign Glory be to the Father, and to the Son, \Star{} and to the Holy Ghost.\end{Resp}
\begin{Resp}\Rsign Together they set at nought the Royal Palace; together they attained unto the kingdom of heaven.\end{Resp}
\switchcolumn*

\LA
\LessonHead{Lectio IX}
\switchcolumn
\EN
\LessonHead{Lesson IX}
\switchcolumn*

\LA
\noindent Duo autem sunt génera persecutórum: unum palam sæviéntium, álterum ficte fraudulentérque blandiéntium. Contra utrúmque nos muníre atque institúere volens Salvátor, et supra ab hypócrisi pharisæórum atténdere, et hic a carníficum cæde prǽcipit non timére; quia vidélicet, post mortem nec horum crudélitas nec illórum váleat simulátio duráre. Nonne quinque pásseres véneunt dipóndio? Si minutíssima, inquit, animália, et quæ quólibet per aëra ferúntur volatília, Deus oblivísci non potest; vos, qui ad imáginem facti estis Creatóris, non debétis terréri ab his qui occídunt corpus; quia, qui irrationabília animália gubérnat, rationabília curáre non désinit.\par
\switchcolumn
\EN
\noindent If persecutors there are two kinds first, of such as do openly rage in cruelty against us and, secondly, of such as do seek, by cunning wiliness and lying, to beguile us. Against both these the Saviour willeth to guard and strengthen us, in one place warning us to be not afraid of them that kill the body, and, in another place, to beware of the leaven of the Pharisees since, when we are dead, neither the cruelty of the one class, nor the falsehood of the other, will be able any more to touch us. Are not five sparrows sold for two farthings? If God, saith the Lord, if God cannot forget the least of the works of His hands that hath life, the little birds that fly hither and thither in the air, if He cannot forget them, wherefore should ye, who are made in the image and likeness of your Maker, wherefore should ye be afraid of them that kill the body? He that is the careful Lord of the beasts, which think not, how much more shall He be careful of man which hath a reasonable soul?\par
\switchcolumn*

\LA
\TeDeum{Te Deum}
\switchcolumn
\EN
\TeDeum{Te Deum}
\switchcolumn*

\LA
\LessonHead{Lectio IX (pro commemoratione)}
\switchcolumn
\EN
\LessonHead{Lesson IX (when commemorated)}
\switchcolumn*

\LA
\LessonTitle{Commemoratio: Ss. Joannis et Pauli Martyrum}
\switchcolumn
\EN
\LessonTitle{Commemoration of: Sts. John and Paul, Martyrs}
\switchcolumn*

\LA
\noindent Joánnes et Paulus, fratres Románi, cum facultátes a Constántia, Constantíni fília, cui pie fidelitérque servíerant, sibi relíctas in Christi páuperes distribúerent, a Juliáno Apóstata in númerum familiárium suórum invitáti, líbere negavérunt se apud eum esse velle, qui a Jesu Christo defecísset. Quare præfinítum est eis spátium decem diérum, ut Jóvi sacrificáre induceréntur; quod scelus patráre cum constantíssime recusássent. Terentiáno júdice, abscíssis domi capítibus martýrii palmam meruérunt. Eórum gloriósus éxitus, a spirítibus immúndis est pervulgátus, qui multórum vexábant córpora; in quibus Terentiáni fílius, et ipse oppréssus a dǽmone, ad sepúlcrum Mártyrum liberátus est. Quo miráculo et is in Christum crédidit, et ejus pater Terentiánus; a quo étiam horum beatórum Mártyrum vita scripta esse dícitur.\par
\switchcolumn
\EN
\noindent The Roman brothers John and Paul distributed to the poor the wealth they had been left by Constantia, the daughter of Constantine, whom they had served justly and faithfully. Invited by Julian the Apostate to join the members of his household, they boldly declared that they did not wish to live in the house of a man who had abandoned Jesus Christ. They were therefore given ten days in which to be persuaded to sacrifice to Jupiter. As they steadfastly refused to commit this sin, they were beheaded in their home, at the command of Terentian the judge, thus meriting the palm of martyrdom. The news of their glorious death was spread abroad by unclean spirits, who began tormenting the bodies of many persons, among them the son of Terentian. He was freed of his diabolical tormentor at the tomb of the Martyrs. This miracle led both him and his father, Terentian, to believe in Christ; and the latter is said to have written the life of the holy Martyrs.\par
\switchcolumn*

\LA
\TeDeum{Te Deum}
\switchcolumn
\EN
\TeDeum{Te Deum}
\switchcolumn*

\end{paracol}

\par\needspace{10\baselineskip}\addvspace{1.2em}{\centering\small\textsc{Die 27 Junii} · \textsc{27 June}\par}
\Office{Quarta die infra Octavam Nativitatis S. Joannis Baptistæ}{Fourth Day of the Octave of the Nativity of St. John the Baptist}{Semiduplex}
\addcontentsline{toc}{subsection}{Die 27 Junii · Quarta die infra Octavam Nativitatis S. Joannis Baptistæ \textit{— Fourth Day of the Octave of the Nativity of St. John the Baptist}}
\begin{paracol}{2}
\LA
\RefLine{Lectiones I–III: de Scriptura occurrente.}
\switchcolumn
\EN
\RefLine{Lessons I–III: from the Scripture of the day.}
\switchcolumn*

\LA
\LessonHead{Lectio IV}
\switchcolumn
\EN
\LessonHead{Lesson IV}
\switchcolumn*

\LA
\LessonTitle{Sermo sancti Basilíi Magni}
\switchcolumn
\EN
\LessonTitle{From the Sermons of St. Basil the Great, Archbishop of Caesarea-in-Pontus.}
\switchcolumn*

\LA
\Source{Homilia 2 in Psalm. 28}
\switchcolumn
\EN
\Source{Hom ii. Ps. xxviii.)}
\switchcolumn*

\LA
\noindent Vox Dómini super aquas. Qualis vox? super quas aquas? Velut prophetíam accipiámus quod dictum est. Memíneris Joánnis, qui interrogátus a Judǽis: Tu quis es? quod respónsum dábimus iis, qui misérunt nos? respóndit: Ego vox clamántis in desérto. Igitur vox Dómini est Joánnes, Angelus a Deo missus ante fáciem Dómini, ut paráret Dómino plebem perféctam. Hæc ígitur vox super aquas, erat super Jordánem, in quo baptizábat prǽdicans pœniténtiæ baptísmum, et non solum in Jordáne, sed étiam in Ænon prope Salim, quia aquæ multæ erant illic.\par
\switchcolumn
\EN
\noindent The voice of the Lord is upon the waters. What voice is this What are these waters Let us take that which is here said as a prophecy. Thou rememberest how that this is the record of John, when the Jews (sent Priests and Levites from Jerusalem to ask him, Who art thou And he confessed and denied not, but confessed, I am not the Christ. And they asked him, What then Art thou Elias And he saith, I am not. Art thou that Prophet? And he answered, No. Then said they unto him, Who art thou that we may give an answer to them that sent us. \textasciitilde{}(What sayest thou of thyself?) He said, I am the voice of one crying in the wilderness. (John i. 19-23; Isa. xl. 3, 4.) Therefore, John is the voice of the Lord. This is he, of whom it is written, Behold, I send My messenger before thy face (Luke vii. 27) to make ready a people prepared for the Lord (i. 17.) This voice upon the waters was heard upon those of Jordan, wherein he baptized, preaching the baptism of repentance for the remission of sins (iii. 3.) Nor was this voice heard upon the waters of Jordan only, but also, in Enon near to Salim, because there was much water there. (John iii. 23.)\par
\switchcolumn*

\LA
\begin{Resp}\Rsign Descéndit Angelus Dómini ad Zacharíam dicens: Accipe púerum in senectúte tua: \Star{} Et habébit nomen Joánnes Baptísta.\end{Resp}
\begin{Resp}\Vsign Iste puer magnus coram Dómino: nam et manus ejus cum ipso est.\end{Resp}
\begin{Resp}\Rsign Et habébit nomen Joánnes Baptísta.\end{Resp}
\switchcolumn
\EN
\begin{Resp}\Rsign The Angel of the Lord was sent down unto Zacharias, saying: Thou shalt beget a son in thine old age; \Star{} and his name shall be called John the Baptist.\end{Resp}
\begin{Resp}\Vsign The child shall be great in the sight of the Lord; for His hand is with him.\end{Resp}
\begin{Resp}\Rsign And his name shall be called John the Baptist.\end{Resp}
\switchcolumn*

\LA
\LessonHead{Lectio V}
\switchcolumn
\EN
\LessonHead{Lesson V}
\switchcolumn*

\LA
\noindent Igitur vox Dómini super aquas, Joánnes est super baptísmum. Illic et Deus majestátis intónuit: venit enim vox de cælo, dicens: Hic est Fílius meus diléctus, in quo mihi complácui. Tunc étiam Dóminus super aquas multas dignátus est descéndere in baptísma Joánnis, ut compléret omnem justítiam quæ in lege est. Vox Dómini in virtúte. Auferet enim debilitátes pópuli per pœniténtiæ baptísmum, per ipsum baptízans in aqua ad pœniténtiam. In virtúte est vox, dicens: Pœniténtiam ágite, appropinquávit enim regnum cælórum: et, Fácite fructus dignos pœniténtiæ.\par
\switchcolumn
\EN
\noindent The voice of the Lord upon the waters, then, is John upon baptism. Then also the God of glory thundereth. For (when Jesus came from Nazareth of Galilee, and was baptized of John in Jordan..) there came a voice from heaven, saying, This is My Beloved Son, in Whom I am well pleased. (Mark i. 9, II.) Then also was it true that The Lord is upon many waters when He was pleased to come from Galilee to Jordan unto John, to be baptized of him... to fulfil all righteousness \textasciitilde{}(Matth. iii. 13, 15,) which is of the law. (Rom. x. 5.) The voice of the Lord is powerful powerful to heal the weaknesses of the people through the baptism of repentance, baptizing through John with water unto repentance. (Matth. iii. 11.) The voice of the Lord is powerful, which saith Repent ye, for the kingdom of heaven is at hand, and Bring forth fruits meet for repentance.\par
\switchcolumn*

\LA
\begin{Resp}\Rsign Hic est præcúrsor diléctus, et lucérna lucens ante Dóminum: \Star{} Ipse est enim Joánnes, qui viam Dómino præparávit in erémo; sed et Agnum Dei demonstrávit, et illuminávit mentes hóminum.\end{Resp}
\begin{Resp}\Vsign Ipse præíbit ante illum in spíritu et virtúte Elíæ.\end{Resp}
\begin{Resp}\Rsign Ipse est enim Joánnes, qui viam Dómino præparávit in erémo; sed et Agnum Dei demonstrávit, et illuminávit mentes hóminum.\end{Resp}
\switchcolumn
\EN
\begin{Resp}\Rsign This is the well-beloved Fore-runner, a burning and a shining light before the Lord. \Star{} For even this is John, who made straight in the desert an highway for our God. This also is he which bare witness, saying: Behold the Lamb of God! And his light lightened the minds of men.\end{Resp}
\begin{Resp}\Vsign He shall go before Him in the spirit and power of Elias.\end{Resp}
\begin{Resp}\Rsign For even this is John, who made straight in the desert an highway for our God. This also is he which bare witness, saying Behold the Lamb of God! And his light lightened the minds of men.\end{Resp}
\switchcolumn*

\LA
\LessonHead{Lectio VI}
\switchcolumn
\EN
\LessonHead{Lesson VI}
\switchcolumn*

\LA
\noindent Vox Dómini confringéntis cedros. Potest dici, quod parans Dómino pópulum perféctum, elátas impietátes et contra cognitiónem Dei exaltátas confríngens, ac cónterens, oblíqua faciébat recta. Qui enim omnem collem ac montem humíliat, hic erat qui confringébat cedros, et Dómino viam adæquábat, per hoc quod ad pœniténtiam inducébat altum, et elátum, et supérbum cor. Unde ejus præparatiónem suscípiens Dóminus, suo advéntu confrégit oppósitas poténtias, cedros Líbani figuráte dictas. Opórtet enim Dóminum regnáre, donec ponat inimícos sub pedes suos, et cedros istas commínuat.\par
\switchcolumn
\EN
\noindent The voice of the Lord breaketh the cedars. This may well be said of him who, being sent as a messenger before the face of the Lord, to make ready for Him a prepared people, made the crooked places straight, by breaking down and treading flat the haughty growths of ungodliness that had lifted themselves up to block out the acknowledgment of God. He by whom every valley was exalted, and every mountain and hill was made low, the same was he who broke the cedars, and made straight in the desert an highway for our God, by laying low in repentance the hearts that were haughty, and proud, and lifted up. The Lord availing Himself of that preparation, smote down at His coming every power that withstood Him, which are spoken of under a similitude as the cedars of Lebanon. For the Lord must reign till He hath put all enemies under His feet, (i Cor. xv. 25,) and trodden down these cedars.\par
\switchcolumn*

\LA
\begin{Resp}\Rsign Innuébant patri ejus quem vellet vocári eum: et póstulans pugillárem, scripsit dicens: \Star{} Joánnes est nomen ejus.\end{Resp}
\begin{Resp}\Vsign Apértum est os Zacharíæ, et prophetávit dicens.\end{Resp}
\begin{Resp}\Rsign Joánnes est nomen ejus.\end{Resp}
\begin{Resp}\Vsign Glória Patri, et Fílio, \Star{} et Spirítui Sancto.\end{Resp}
\begin{Resp}\Rsign Joánnes est nomen ejus.\end{Resp}
\switchcolumn
\EN
\begin{Resp}\Rsign They made signs to his father, how he would have him called; and he asked for a writing-table, and wrote, saying: \Star{} His name is John.\end{Resp}
\begin{Resp}\Vsign The mouth of Zacharias was opened, and he prophesied, saying:\end{Resp}
\begin{Resp}\Rsign His name is John.\end{Resp}
\begin{Resp}\Vsign Glory be to the Father, and to the Son, \Star{} and to the Holy Ghost.\end{Resp}
\begin{Resp}\Rsign His name is John.\end{Resp}
\switchcolumn*

\LA
\LessonHead{Lectio VII}
\switchcolumn
\EN
\LessonHead{Lesson VII}
\switchcolumn*

\LA
\LessonTitle{Léctio sancti Evangélii secúndum Lucam}
\switchcolumn
\EN
\LessonTitle{From the Holy Gospel according to Luke}
\switchcolumn*

\LA
\Cite{Luc 1:57-68}
\switchcolumn
\EN
\Cite{Luke 1:57-68}
\switchcolumn*

\LA
\noindent Elísabeth implétum est tempus pariéndi, et péperit fílium. Et audiérunt vicíni et cognáti ejus quia magnificávit Dóminus misericórdiam suam cum illa, et congratulabántur ei. Et réliqua.\par
\switchcolumn
\EN
\noindent Elizabeth's full time came that she should be delivered and she brought forth a son. And her neighbours and her cousins heard how the Lord had showed great mercy upon her and they rejoiced with her. And so on\par
\switchcolumn*

\LA
\LessonTitle{De Homilía sancti Ambrósii Epíscopi}
\switchcolumn
\EN
\LessonTitle{From the Homily by St. Ambrose, Bishop of Milan.}
\switchcolumn*

\LA
\Source{Lib. 2 in Luc. cap. 1 in fine}
\switchcolumn
\EN
\Source{Bk. ii. on Luke's}
\switchcolumn*

\LA
\noindent Et Zacharías pater ejus implétus est Spíritu Sancto, et prophetábat. Vide quam bonus Deus, et fácilis indulgére peccátis; non solum abláta restítuit, sed étiam insperáta concédit. Ille dudum mutus prophétat; hæc enim grátia Dei máxima, quod eum, qui negáverat, confitétur. Nemo ergo diffídat, nemo véterum cónscius delictórum prǽmia divína despéret. Novit Dóminus mutáre senténtiam, si tu nóveris emendáre delíctum.\par
\switchcolumn
\EN
\noindent And his father Zacharias was filled with the Holy Ghost, and prophesied. Behold how good is God, and how ready to forgive sinners Not only doth He give back that which He hath taken away, but He addeth moreover such and such things, more than either we ask or think. He that had hitherto been dumb, now speaketh prophecy. But the greatest grace of God is in this, that he which had denied, now maketh profession. Therefore let no man despair, nay, though he well knoweth what have been his past sins, let him never give up hope of a reward from God. If thou knowest how to amend thy crooked ways, God knoweth how to turn away the judgments of His anger.\par
\switchcolumn*

\LA
\begin{Resp}\Rsign Præcúrsor Dómini venit, de quo ipse testátur: \Star{} Nullus major inter natos mulíerum Joánne Baptísta.\end{Resp}
\begin{Resp}\Vsign Hic est enim prophéta, et plus quam prophéta, de quo Salvátor ait.\end{Resp}
\begin{Resp}\Rsign Nullus major inter natos mulíerum Joánne Baptísta.\end{Resp}
\switchcolumn
\EN
\begin{Resp}\Rsign The Fore-runner of the Lord cometh, to whom He Himself bare witness, saying: \Star{} Among them that are born of women there hath not risen a greater than John the Baptist.\end{Resp}
\begin{Resp}\Vsign A Prophet? Yea, and much more than a Prophet. This is he of whom the Saviour saith:\end{Resp}
\begin{Resp}\Rsign Among them that are born of women there hath not risen a greater than John the Baptist.\end{Resp}
\switchcolumn*

\LA
\LessonHead{Lectio VIII}
\switchcolumn
\EN
\LessonHead{Lesson VIII}
\switchcolumn*

\LA
\noindent Et tu puer, prophéta Altíssimi vocáberis. Pulchre, cum de Dómino prophetáret, ad prophétam sua verba convértit, ut hoc quoque benefícium esse Dómini designáret: ne cum públice enumeráret sua, quasi ingrátus tacuísse viderétur quæ agnoscébat in fílio. Sed fortásse áliqui quasi irrationábilem mentis excéssum putent, quod octo diérum allóquitur infántem. Verum, si teneámus, intelligémus profécto quod pótuit vocem patris natus audíre, qui Maríæ salutatiónem ántequam nascerétur audívit.\par
\switchcolumn
\EN
\noindent And thou, child, shalt be called the Prophet of the Highest. How gracefully, while as he prophesieth of the Lord, he turneth his address to the Prophet, making mention of this great mercy of the Lord along with the others, lest, while he openly gave thanks for his own benefits, he should seem to keep the silence of unthankfulness regarding those which he knew had been given to his boy. But some will perhaps deem it his folly that he addressed his discourse to a babe of eight days old. Verily, if we call to mind that John heard the voice of Mary's salutation when he was in his mother's womb, we shall understand how much rather he could hear the voice of his father now when he was born.\par
\switchcolumn*

\LA
\begin{Resp}\Rsign Gábriel Angelus appáruit Zacharíæ dicens: Nascétur tibi fílius, nomen ejus Joánnes vocábitur: \Star{} Et multi in nativitáte ejus gaudébunt.\end{Resp}
\begin{Resp}\Vsign Erit enim magnus coram Dómino, vinum et síceram non bibet.\end{Resp}
\begin{Resp}\Rsign Et multi in nativitáte ejus gaudébunt.\end{Resp}
\begin{Resp}\Vsign Glória Patri, et Fílio, \Star{} et Spirítui Sancto.\end{Resp}
\begin{Resp}\Rsign Et multi in nativitáte ejus gaudébunt.\end{Resp}
\switchcolumn
\EN
\begin{Resp}\Rsign The Angel Gabriel appeared unto Zacharias, and said thy wife Elizabeth shall bear thee a son, and thou shalt call his name John; \Star{} And many shall rejoice at his birth.\end{Resp}
\begin{Resp}\Vsign For he shall be great in the sight of the Lord, and shall drink neither wine nor strong drink.\end{Resp}
\begin{Resp}\Rsign And many shall rejoice at his birth.\end{Resp}
\begin{Resp}\Vsign Glory be to the Father, and to the Son, \Star{} and to the Holy Ghost.\end{Resp}
\begin{Resp}\Rsign And many shall rejoice at his birth.\end{Resp}
\switchcolumn*

\LA
\LessonHead{Lectio IX}
\switchcolumn
\EN
\LessonHead{Lesson IX}
\switchcolumn*

\LA
\noindent Sciébat profécto álias esse aures prophétæ, quæ Spíritu Dei, non córporis ætáte reserántur. Habébat intelligéndi sensum, qui exsultándi habébat afféctum. Simul illud advérte, quam paucis Elísabeth, quam multis Zacharías prophétet: et utérque Sancto replétus Spíritu loquebátur: sed disciplína servátur, ut múlier díscere magis quæ divína sunt, stúdeat, quam docére.\par
\switchcolumn
\EN
\noindent Zacharias knew well that a Prophet hath ears which open under the influence of the Spirit of God, instead of that of advancing age. He that had had sense to leap in the womb for joy, wanted not understanding. At the same time remark unto how many Zacharias prophesied, and though both he and his wife were filled with the Holy Ghost, yet all things are done in due order, and the woman studieth rather to learn the things of God than to teach them.\par
\switchcolumn*

\LA
\TeDeum{Te Deum}
\switchcolumn
\EN
\TeDeum{Te Deum}
\switchcolumn*

\end{paracol}

\par\needspace{10\baselineskip}\addvspace{1.2em}{\centering\small\textsc{Die 28 Junii} · \textsc{28 June}\par}
\Office{S. Irenæi Episcopi et Martyris}{St. Irenaeus, Bishop and Martyr}{Duplex}
\addcontentsline{toc}{subsection}{Die 28 Junii · S. Irenæi Episcopi et Martyris \textit{— St. Irenaeus, Bishop and Martyr}}
\begin{paracol}{2}
\LA
\RefLine{Lectiones I–III: de Scriptura occurrente.}
\switchcolumn
\EN
\RefLine{Lessons I–III: from the Scripture of the day.}
\switchcolumn*

\LA
\RefLine{Lectio IX: de In Vigilia Ss. Petri et Pauli Apostolorum (06-28r).}
\switchcolumn
\EN
\RefLine{Lesson IX: of Vigil of Sts. Peter and Paul, Apostles (06-28r).}
\switchcolumn*

\LA
\LessonHead{Lectio IV}
\switchcolumn
\EN
\LessonHead{Lesson IV}
\switchcolumn*

\LA
\noindent Irenæus, non longe ab urbe Smyrna in Asia proconsulari natus, jam inde a púero sese Polycarpo, Joánnis Evangelistæ discipulo eidemque episcopo Smyrnæórum, tradiderat in disciplínam. Hoc tam excellénti magistro, progréssus in doctrina præceptisque christianæ religiónis insignes fecit. Polycarpo in cælum martyrii glória sublato, etsi erat Irenæus in sacris litteris egregie versatus, quod tamen incredibili studio flagraret discéndi quæ dógmata depósiti loco custodienda ceteri accepissent, quos Apóstoli institúerant; horum quam pótuit plures convénit, quæque ab iisdem audívit, mémori mente ténuit, ea deinceps opportune advérsus hæreses allatúrus, quas cum vidéret ingénti pópuli christiáni damno latius in dies manare, diligenter copioséque reféllere cogitarat. In Gállias inde profectus, a Pothino episcopo presbyter est constitútus Ecclésiæ Lugdunénsis. Quod munus sic laborando in verbo et doctrina gessit, ut (testibus sanctis Martyribus, qui, Marco Aurelio imperatóre, strenue pro vera pietáte certarant) æmulatórem sese præstíterit testaménti Christi.\par
\switchcolumn
\EN
\noindent Irenaeus was born in proconsular Asia, not far from the city of Smyrna. There he had already as a boy entrusted himself to the teaching of Polycarp, disciple of John the Evangelist, and bishop of Smyrna. Under such an excellent master, he made remarkable progress in learning and in the precepts of the Christian religion. When Polycarp was taken up to heaven by a glorious martyrdom, although Irenaeus was eminently versed in sacred letters, nevertheless, he burned with an incredible zeal to learn what articles of belief the others who were instructed by the Apostles had received, to be preserved in the deposit of faith. For this reason he brought together as many of those men as he could, and whatever things he heard from them, he carefully retained in his mind. Thus he could advantageously bring them to bear in the future against those heresies, which he saw were being diffused more widely day by day to the great detriment of the Christian people, and he diligently planned thoroughly to confute them. Then, having set out for Gaul, he was appointed as a priest of the church of Lyons by Pothinus the bishop. And this office he discharged in such a manner, labouring both by word and by teaching, that (according to the testimony of the holy Martyrs who, when Marcus Aurelius was emperor, were engaged in a vigorous combat for the true religion) he distinguished himself as an imitator of the testament of Christ.\par
\switchcolumn*

\LA
\begin{Resp}\Rsign Honéstum fecit illum Dóminus, et custodívit eum ab inimícis, et a seductóribus tutávit illum: \Star{} Et dedit illi claritátem ætérnam.\end{Resp}
\begin{Resp}\Vsign Descendítque cum illo in fóveam, et in vínculis non derelíquit eum.\end{Resp}
\begin{Resp}\Rsign Et dedit illi claritátem ætérnam.\end{Resp}
\switchcolumn
\EN
\begin{Resp}\Rsign The Lord made him honourable, and defended him from his enemies, and kept him safe from those that lay in wait for him, \Star{} And gave him perpetual glory.\end{Resp}
\begin{Resp}\Vsign He went down with him into the pit, and left him not in bonds.\end{Resp}
\begin{Resp}\Rsign And gave him perpetual glory.\end{Resp}
\switchcolumn*

\LA
\LessonHead{Lectio V}
\switchcolumn
\EN
\LessonHead{Lesson V}
\switchcolumn*

\LA
\noindent Cum Mártyres ipsi clerusque Lugdunénsis de pace Ecclesiárum Asiæ, quam Montanistárum factio turbarat, solliciti cum primis essent, Irenæum, cujus esse potíssimum habéndam ratiónem prædicábant, unum ómnium maxime delegérunt, quem Romam ad Eleuthérium Pontificem mitterent rogátum, ut novis sectariis auctoritate Sedis Apostolicæ reprobátis, discordiárum causa tollerétur. Jam Pothinus epíscopus, martyr decesserat: huic Irenæus cum successísset, tam feliciter munus obiit episcopatus, ut sapiéntia, oratióne, exemploque suo non modo brevi cives Lugdunenses omnes, sed multos étiam aliárum Galliæ urbium íncolas superstitiónem atque errórem abjecisse, dedísseque christianæ milítiæ nómina víderit. Interea, cum de die celebrándi Paschátis orta esset contentio, ac Victor, Romanus Pontifex, Asiános episcopos ab collégis réliquis fere ómnibus dissidéntes aut prohibuísset communióne sacrórum, aut prohibére minátus esset, eum Irenæus, sequester pacis, decenter monuit, exemplisque usus Pontificum superiórum indúxit, ut ne tot Ecclésias ob ritum quem a majóribus accepisse se dicerent, avelli ab unitate catholica paterétur.\par
\switchcolumn
\EN
\noindent These very Martyrs, together with the clergy of Lyons, began to be anxious concerning the peace of the churches of Asia, which the faction of the Montanists had disturbed. And so they selected Irenaeus, whose person they considered of the greatest importance, as the one before all others whom they should send to Rome to Pope Eleutherius to ask, that, with the condemnation of the new dissidents by the authority of the Apostolic See, the cause of the dissensions might be removed. Already the bishop Pothinus had died a martyr and Irenaeus succeeded him. He applied himself so well to the duties of a bishop, that in a short time he saw not only all the citizens of Lyons, but also many of the inhabitants of other cities in Gaul cast aside their superstitions and errors, and enroll themselves in the Christian army. Meanwhile, a dispute had arisen concerning the date of the celebration of Easter. As the bishops of Asia were disagreeing with nearly all their fellow-bishops, the Roman Pontiff Victor had cut them off from the communion of the faithful. Irenaeus, however, who was zealous for peace, admonished him in a becoming manner, and urged, by examples of the practice of previous Pontiffs, that he should not suffer so many Churches to be cut off from Catholic unity, on account of a rite which they said they had received from their ancestors.\par
\switchcolumn*

\LA
\begin{Resp}\Rsign Desidérium ánimæ ejus tribuísti ei Dómine, \Star{} Et voluntáte labiórum ejus non fraudásti eum.\end{Resp}
\begin{Resp}\Vsign Quóniam prævenísti eum in benedictiónibus dulcédinis: posuísti in cápite ejus corónam de lápide pretióso.\end{Resp}
\begin{Resp}\Rsign Et voluntáte labiórum ejus non fraudásti eum.\end{Resp}
\switchcolumn
\EN
\begin{Resp}\Rsign Lord, Thou hast given him his heart's desire. \Star{} And hast not withholden the request of his lips.\end{Resp}
\begin{Resp}\Vsign For Thou hast prevented him with the blessings of sweetness: Thou hast set a crown of precious stones upon his head.\end{Resp}
\begin{Resp}\Rsign And hast not withholden the request of his lips.\end{Resp}
\switchcolumn*

\LA
\LessonHead{Lectio VI}
\switchcolumn
\EN
\LessonHead{Lesson VI}
\switchcolumn*

\LA
\noindent Multa scripsit, quæ Eusebius Cæsariénsis et sanctus Hieronymus mémorant, quorúmque pars magna intércidit injuria témporum. Exstant ejus advérsus hæreses libri quinque, anno circiter centésimo octogesimo perscripti, dum adhuc Eleutherius rem christianam publicam géreret. In tertio libro vir Dei, ab iis edoctus quos auditores constat fuísse Apostolórum, grave in primis atque præclárum de Romana Ecclésia, deque illíus episcopórum successióne, divinæ traditiónis fideli, perpetua, certíssima custode, testimónium dixit. Atque ad hanc, dixit, Ecclésiam propter potiórem principalitátem necesse est omnem conveníre Ecclésiam, hoc est eos qui sunt undique fidéles. Postremo una cum aliis prope innumerabílibus, quos ipse ad veram fidem frugemque perduxerat, martyrio coronátus migrávit in cælum, anno salútis ducentésimo secundo, quo témpore Septimius Sevérus Augustus eos omnes, qui constanter in colenda christiana religióne perstare voluíssent, in summum cruciátum dari atque intérfici imperaverat. Sancti Irenæi festum Benedíctus décimus quintus Pontifex maximus ad universam Ecclésiam extendit.\par
\switchcolumn
\EN
\noindent He wrote many works, which are mentioned by Eusebius of Caesarea and by St. Jerome, a great part of which have perished through the ravages of time. There are extant five books of his against heresies, written down about the year 180, while Eleutherius was still ruling the Christian commonwealth. In the third book, the man of God, instructed by those who, it is certain, had been hearers of the Apostles, gives to the Roman Church and to the succession of her bishops a testimony surpassing all others in weight and brilliancy, when he calleth her the faithful, perpetual, and most assured guardian of divine tradition. For he said, that with this Church it is necessary that the whole Church (that is, those in all places who are of the faithful) should agree, because of its more powerful preeminence. At length with almost countless others, whom he had himself brought over to the true faith and its practice, being crowned with martyrdom he passed to heaven in the year of salvation 202. At that time Septimius Severus Augustus had commanded that all those who wished to remain constantly steadfast in the practice of the Christian religion should be condemned to the most cruel torments and to death. The supreme Pontiff Benedict XV extended the feast of St. Irenaeus to the universal Church.\par
\switchcolumn*

\LA
\begin{Resp}\Rsign Stola jucunditátis índuit eum Dóminus: \Star{} Et corónam pulchritúdinis pósuit super caput ejus.\end{Resp}
\begin{Resp}\Vsign Cibávit illum Dóminus pane vitæ et intelléctus: et aqua sapiéntiæ salutáris potávit illum.\end{Resp}
\begin{Resp}\Rsign Et corónam pulchritúdinis pósuit super caput ejus.\end{Resp}
\begin{Resp}\Vsign Glória Patri, et Fílio, \Star{} et Spirítui Sancto.\end{Resp}
\begin{Resp}\Rsign Et corónam pulchritúdinis pósuit super caput ejus.\end{Resp}
\switchcolumn
\EN
\begin{Resp}\Rsign The Lord hath put on him a robe of honour, \Star{} And put about his head a crown of joy.\end{Resp}
\begin{Resp}\Vsign With the bread of life and understanding hath the Lord fed him, and given him the water of health and wisdom to drink.\end{Resp}
\begin{Resp}\Rsign And put about his head a crown of joy.\end{Resp}
\begin{Resp}\Vsign Glory be to the Father, and to the Son, \Star{} and to the Holy Ghost.\end{Resp}
\begin{Resp}\Rsign And put about his head a crown of joy.\end{Resp}
\switchcolumn*

\LA
\LessonHead{Lectio VII}
\switchcolumn
\EN
\LessonHead{Lesson VII}
\switchcolumn*

\LA
\LessonTitle{Léctio sancti Evangélii secúndum Matthǽum}
\switchcolumn
\EN
\LessonTitle{From the Holy Gospel according to Matthew}
\switchcolumn*

\LA
\Cite{Matt 10:28-33}
\switchcolumn
\EN
\Cite{Matt 10:28-33}
\switchcolumn*

\LA
\noindent In illo témpore: Dixit Jesus discípulis suis: Nolíte timere eos qui occídunt corpus, ánimam autem non possunt occídere; sed potius timéte eum qui potest et ánimam et corpus perdere in gehennam. Et réliqua.\par
\switchcolumn
\EN
\noindent At that time: Jesus said unto his disciples: Fear not them which kill the body, but are not able to kill the soul: but rather fear him which is able to destroy both soul and body in hell. And so forth.\par
\switchcolumn*

\LA
\LessonTitle{Homilía sancti Irenæi Epíscopi et Martyris}
\switchcolumn
\EN
\LessonTitle{A Homily by St. Irenaeus the Bishop and Martyr}
\switchcolumn*

\LA
\Source{Lib. 3 advérsus Hæreses cap. 18, alias 20, num. 5-6}
\switchcolumn
\EN
\Source{Lib. 3 adversus Haereses cap. 18, alias 20, num. 5-6}
\switchcolumn*

\LA
\noindent Sciebat Dóminus eos qui persecutiónem passuri essent; sciebat et eos qui flagéllari et occídi haberent propter eum. Erat ergo sermo ejus, adhortántis étiam illos: Nolíte timere eos qui occídunt corpus, ánimam autem non possunt occídere. Timéte autem magis eum qui habet potestátem et corpus et ánimam mittere in gehennam, et servare eas quæ essent ad eum confessiónes. Etenim ipse confessurum se promittebat coram Patre suo eos, qui confiteréntur nomen suum coram homínibus; negatúrum autem eos qui negarent eum, et confusurum qui confunderéntur contra confessiónem ejus. Et, cum hæc ita se habeant, ad tantam temeritátem progressi sunt quidam, ut étiam Mártyres spernant, et vitúperent eos qui propter Dómini confessiónem occidúntur, et sústinent ómnia a Dómino prædicta, et secúndum hoc conántur vestígia ássequi passiónis Dómini, passibilis Mártyres facti; quos et concédimus ipsis Martyribus. Cum enim inquirétur sanguis eórum, et glóriam consequéntur, tunc a Christo confundéntur omnes qui inhonoravérunt eórum martyrium.\par
\switchcolumn
\EN
\noindent The Lord knew both those who would suffer persecution; and he knew those who would be scourged and slain for his sake. His words are indeed those of exhortation: Fear not them which kill the body, but are not able to kill the soul. Rather fear him which hath power to cast both body and soul into hell, and to preserve whomsoever should confess him. For indeed he hath promised to confess before his Father those who should confess his name before men; and to deny those who should deny him, and to be ashamed of those who should be ashamed to confess him. Yet for all this, some have gone so far in audacity, as to scorn even the Martyrs, and to revile those who have been put to death for confessing the Lord, and who endure all that the Lord foretold, and so far endeavour to follow in the footsteps of the Lord's Passion, being made Martyrs of one who was capable of suffering; but these we leave to the Martyrs themselves. For when their blood shall be required, and they shall obtain glory, then all those who have dishonoured their martyrdom will be put to confusion by Christ.\par
\switchcolumn*

\LA
\begin{Resp}\Rsign Coróna áurea super caput ejus, \Star{} Expréssa signo sanctitátis, glória honóris, et opus fortitúdinis.\end{Resp}
\begin{Resp}\Vsign Quóniam prævenísti eum in benedictiónibus dulcédinis, posuísti in cápite ejus corónam de lápide pretióso.\end{Resp}
\begin{Resp}\Rsign Expréssa signo sanctitátis, glória honóris, et opus fortitúdinis.\end{Resp}
\switchcolumn
\EN
\begin{Resp}\Rsign A crown of gold upon his head, \Star{} Wherein is engraved Holiness, an ornament of honour, a costly work.\end{Resp}
\begin{Resp}\Vsign For Thou hast prevented him with the blessings of sweetness, Thou hast set a crown of precious stones upon his head.\end{Resp}
\begin{Resp}\Rsign Wherein is engraved Holiness, an ornament of honour, a costly work.\end{Resp}
\switchcolumn*

\LA
\LessonHead{Lectio VIII}
\switchcolumn
\EN
\LessonHead{Lesson VIII}
\switchcolumn*

\LA
\noindent Hoc autem idem et illis occurrit, qui dicunt eum putative passum. Si enim non vere passus est, nulla grátia ei, cum nulla fúerit passio; et nos, cum incipiémus vere pati, seducens vidébitur, adhortans nos vapulare, et álteram præbére maxíllam, si ipse illud non prior in veritáte passus est. Et, quemádmodum illos sedúxit, ut viderétur eis ipse hoc quod non erat; et nos sedúcit, adhortans perferre ea quæ ipse non pértulit. Erimus autem et super Magistrum, dum pátimur et sustinémus quæ neque passus est, neque sustinuit Magister. Sed, quóniam solus vere magister Dóminus noster, et bonus vere Fílius Dei; et pátiens Verbum Dei Patris, fílius hóminis factus. Luctátus est enim et vicit; erat enim homo pro pátribus certans, et per obediéntiam inobediéntiam persolvens. Alligávit enim fortem, et solvit infirmos, et salútem donávit plasmáti suo, destruens peccátum. Igitur qui dicunt eum putative manifestátum, neque in carne natum, neque vere hóminem factum, adhuc sub veteri sunt damnatióne.\par
\switchcolumn
\EN
\noindent This also meeteth the objection of those who say that Christ only appeared to be suffering. For if he did not really suffer, no gratitude is due unto him, when there was no suffering; and when we begin really to suffer, he would seem to be deceiving us, when exhorting us to receive blows and to turn the other cheek, if he himself had not first suffered in reality. And as he deceived them, so as to appear to them what he was not; so also he misleadeth us, exhorting us to endure things which he himself did not endure. And so we should even be higher than the Master, when we suffer and endure things, which the Master neither suffered nor endured. But, whereas our Lord alone is truly the master, he is also truly the Son of God, good and kind, the Word of God the Father made Son of Man. For he strove, and he conquered; for he was a man contending for his fathers and paying the debt of disobedience by his obedience. For he hath bound the strong man, and, he unbindeth the weak, and he hath given salvation to his creature by destroying sin. Therefore, those who say he is manifested only in appearance, and not born in the flesh, nor really made man, are still under the old condemnation.\par
\switchcolumn*

\LA
\begin{Resp}\Rsign Hic est vere Martyr, qui pro Christi nómine sánguinem suum fudit: \Star{} Qui minas júdicum non tímuit, nec terrénæ dignitátis glóriam quæsívit, sed ad cæléstia regna pervénit.\end{Resp}
\begin{Resp}\Vsign Justum dedúxit Dóminus per vias rectas, et osténdit illi regnum Dei.\end{Resp}
\begin{Resp}\Rsign Qui minas júdicum non tímuit, nec terrénæ dignitátis glóriam quæsívit, sed ad cæléstia regna pervénit.\end{Resp}
\begin{Resp}\Vsign Glória Patri, et Fílio, \Star{} et Spirítui Sancto.\end{Resp}
\begin{Resp}\Rsign Qui minas júdicum non tímuit, nec terrénæ dignitátis glóriam quæsívit, sed ad cæléstia regna pervénit.\end{Resp}
\switchcolumn
\EN
\begin{Resp}\Rsign This is a Martyr indeed, who shed his blood for Christ's Name's sake \Star{} Who feared not for the threats of judges, nor sought to be great with the glory of this world, but pressed on unto the kingdom of heaven.\end{Resp}
\begin{Resp}\Vsign The Lord guided the righteous in right paths, and showed him the kingdom of God.\end{Resp}
\begin{Resp}\Rsign Who feared not for the threats of judges, nor sought to be great with the glory of this world, but pressed on unto the kingdom of heaven.\end{Resp}
\begin{Resp}\Vsign Glory be to the Father, and to the Son, \Star{} and to the Holy Ghost.\end{Resp}
\begin{Resp}\Rsign Who feared not for the threats of judges, nor sought to be great with the glory of this world, but pressed on unto the kingdom of heaven.\end{Resp}
\switchcolumn*

\LA
\LessonHead{Lectio IX (pro commemoratione)}
\switchcolumn
\EN
\LessonHead{Lesson IX (when commemorated)}
\switchcolumn*

\LA
\LessonTitle{Commemoratio: S. Irenæi Episcopi et Martyris}
\switchcolumn
\EN
\LessonTitle{Commemoration of: St. Irenaeus, Bishop and Martyr}
\switchcolumn*

\LA
\noindent Irenǽus, non longe ab urbe Smyrna natus, jam inde a púero sese Polycárpo, Joánnis Evangelístæ discípulo eidémque epíscopo Smyrnæórum, tradíderat in disciplínam. Polycárpo in cælum martýrii glória subláto, cum incredíbili stúdio flagráret discéndi quæ dógmata depósiti loco custodiénda céteri accepíssent, quos Apóstoli institúerant; horum quam pótuit plures convénit, quæque ab iísdem audívit, mémori mente ténuit, ea deínceps opportúne advérsus hǽreses allatúrus. In Gálliam proféctus, Ecclésiæ Lugdunénsis présbyter a Photíno epíscopo est constitútus; cui cum successísset, tam felíciter munus óbiit episcopátus, ut sapiéntia, oratióne exemplóque suo non modo brevi cives lugdunénses omnes, sed multos étiam aliárum Galliárum úrbium íncolas superstitiónem atque errórem abiecísse, dedisséque christiánæ milítiæ nómina víderit. Multa scripsit, quorum magna pars intércidit injúria témporum. Exstant ejus advérsus hǽreses libri quinque, in quorum tértio libro grave imprímis atque præclárum de Romána Ecclésia, deque illíus episcopórum successióne, divínæ traditióni fidéli, perpétua, certíssima custóde, testimónium dixit. Atque ad hanc, dixit, Ecclésiam propter potiórem principalitátem necésse est omnem conveníre Ecclésiam, hoc est eos qui sunt úndique fidéles. Martýrio coronátus, migrávit in cælum anno salútis ducentésimo secúndo.\par
\switchcolumn
\EN
\noindent Irenaeus, born not far from the city of Smyrna, from his boyhood was the disciple of Polycarp, himself the disciple of John the Evangelist and the bishop of Smyrna. When Polycarp was taken up to heaven with the glory of martyrdom, Irenaeus strove with incredible zeal to learn what articles of belief the others who were instructed by the Apostles had received to be preserved in the deposit of faith. For this reason he brought together as many of these men as he could, and whatever he heard from them he carefully retained in his memory to bring out later at an opportune time against the heretics. He went to Gaul and was made priest at the church of Lyon by Photinus the bishop. When he succeeded Photinus he carried out the work of his bishopric most successfully: by his wisdom, prayer and example, in a short time, he had rid not only the citizens of Lyon, but also the inhabitants of many other cities of Gaul, from superstition and error, and had enrolled them in the army of Christ. He wrote many works, a great part of which had perished through the ravages of time. Five of his books against heretics are extant, in the third of which he gives to the Roman Church and to the succession of her bishops a testimony surpassing all others in weight and brilliancy, calling her the faithful, perpetual and sure guardian of divine tradition. For he said that it is necessary to the whole Church (that is, those who are of the faithful in all places) should agree with the Roman Church, because of its eminent primacy. Crowned with martyrdom, he went to heaven in the year of salvation 202.\par
\switchcolumn*

\LA
\TeDeum{Te Deum}
\switchcolumn
\EN
\TeDeum{Te Deum}
\switchcolumn*

\end{paracol}

\par\needspace{10\baselineskip}\addvspace{1.2em}{\centering\small\textsc{Die 28 Junii} · \textsc{28 June}\par}
\Office{In Vigilia Ss. Petri et Pauli Apostolorum}{Vigil of Sts. Peter and Paul, Apostles}{Vigilia}
\addcontentsline{toc}{subsection}{Die 28 Junii · In Vigilia Ss. Petri et Pauli Apostolorum \textit{— Vigil of Sts. Peter and Paul, Apostles}}
\begin{paracol}{2}
\LA
\LessonHead{Lectio I}
\switchcolumn
\EN
\LessonHead{Lesson I}
\switchcolumn*

\LA
\LessonTitle{Commemoratio: In Vigilia Ss. Petri et Pauli Apostolorum Léctio sancti Evangélii secúndum Joánnem}
\switchcolumn
\EN
\LessonTitle{Commemoration of: Vigil of Sts. Peter and Paul, Apostles Continuation of the Holy Gospel according to John}
\switchcolumn*

\LA
\Cite{Joannes 21:15-17}
\switchcolumn
\EN
\Cite{John 21:15-17}
\switchcolumn*

\LA
\noindent In illo témpore: Dixit Jesus Simóni Petro: Simon Joánnis, díligis me plus his? Et réliqua.\par
\switchcolumn
\EN
\noindent At that time: Jesus saith to Simon Peter, Simon, son of Jonas, lovest thou me more than these? And so on, and that which followeth.\par
\switchcolumn*

\LA
\LessonTitle{Homilía sancti Augustíni Epíscopi}
\switchcolumn
\EN
\LessonTitle{A Homily by St. Augustine the Bishop}
\switchcolumn*

\LA
\Source{Tractatus 123 in Joannem, in medio}
\switchcolumn
\EN
\Source{Tractatus 123 in Joannem. in medio}
\switchcolumn*

\LA
\noindent Rédditur negatióni trinæ trina conféssio, ne minus amóri lingua sérviat, quam timóri: et plus vocis elicuísse videátur mors ímminens, quam vita præsens. Sit amóris offícium páscere Domínicum gregem, si fuit timóris indícium negáre pastórem. Qui hoc ánimo pascunt oves Christi, ut suas velint esse, non Christi, se convincúntur amáre, non Christum: vel gloriándi vel dominándi, vel acquiréndi cupiditáte; non obediéndi, et subveniéndi, et Deo placéndi caritáte.\par
\switchcolumn
\EN
\noindent To the threefold denial there is now appended a threefold confession, that his tongue may not yield a feebler service to love than to fear, and imminent death may not appear to have elicited more from his lips than present life. Let it be the office of love to feed the Lord's flock, if it was the signal of fear to deny the Shepherd. Those who have this purpose in feeding the flock of Christ, that they may have them as their own, and not as Christ's, are convicted of loving themselves, and not Christ, from the desire either of boasting, or wielding power, or acquiring gain, and not from the love of obeying, serving and pleasing God.\par
\switchcolumn*

\LA
\TeDeum{Te Deum}
\switchcolumn
\EN
\TeDeum{Te Deum}
\switchcolumn*

\end{paracol}

\par\needspace{10\baselineskip}\addvspace{1.2em}{\centering\small\textsc{Die 29 Junii} · \textsc{29 June}\par}
\Office{SS. Apostolorum Petri et Pauli}{Sts. Peter and Paul, Apostles}{Duplex I classis cum Octava communi}
\addcontentsline{toc}{subsection}{Die 29 Junii · SS. Apostolorum Petri et Pauli \textit{— Sts. Peter and Paul, Apostles}}
\begin{paracol}{2}
\LA
\LessonHead{Lectio I}
\switchcolumn
\EN
\LessonHead{Lesson I}
\switchcolumn*

\LA
\LessonTitle{De Actibus Apostolórum}
\switchcolumn
\EN
\LessonTitle{Lesson from the Acts of the Apostles}
\switchcolumn*

\LA
\Cite{Act 3:1-5}
\switchcolumn
\EN
\Cite{Acts 3:1-5}
\switchcolumn*

\LA
\noindent Petrus autem et Joánnes ascendébant in templum ad horam oratiónis nonam. Et quidam vir, qui erat claudus ex útero matris suæ, bajulabátur; quem ponébant cotídie ad portam templi, quæ dícitur Speciósa, ut péteret eleemósynam ab introëúntibus in templum. Is, cum vidísset Petrum et Joánnem incipiéntes introíre in templum, rogábat ut eleemósynam accíperet. Intuens autem in eum Petrus cum Joánne, dixit: Réspice in nos. At ille intendébat in eos, sperans se áliquid acceptúrum ab eis.\par
\switchcolumn
\EN
\noindent Now Peter and John went up into the temple at the ninth hour of prayer. And a certain man who was lame from his mother's womb, was carried: whom they laid every day at the gate of the temple, which is called Beautiful, that he might ask alms of them that went into the temple. He, when he had seen Peter and John about to go into the temple, asked to receive an alms. But Peter with John fastening his eyes upon him, said: Look upon us. But he looked earnestly upon them, hoping that he should receive something of them.\par
\switchcolumn*

\LA
\begin{Resp}\Rsign Simon Petre, ántequam de navi vocárem te, novi te et super plebem meam príncipem te constítui: \Star{} Et claves regni cælórum trádidi tibi.\end{Resp}
\begin{Resp}\Vsign Quodcúmque ligáveris super terram erit ligátum et in cælis; et quodcúmque sólveris super terram, erit solútum et in cælis.\end{Resp}
\begin{Resp}\Rsign Et claves regni cælórum trádidi tibi.\end{Resp}
\switchcolumn
\EN
\begin{Resp}\Rsign Simon Peter, before I called thee out of the ship, I knew thee, and appointed thee for a ruler over My people. \Star{} And I have given unto thee the keys of the kingdom of heaven.\end{Resp}
\begin{Resp}\Vsign Whatsoever, thou shalt bind on earth, shall be bound in heaven; and whatsoever thou shalt loose on earth, shall be loosed in heaven.\end{Resp}
\begin{Resp}\Rsign And I have given unto thee the keys of the kingdom of heaven.\end{Resp}
\switchcolumn*

\LA
\LessonHead{Lectio II}
\switchcolumn
\EN
\LessonHead{Lesson II}
\switchcolumn*

\LA
\Cite{Act 3:6-10}
\switchcolumn
\EN
\Cite{Acts 3:6-10}
\switchcolumn*

\LA
\noindent Petrus autem dixit: Argéntum et aurum non est mihi: quod autem hábeo, hoc tibi do: In nómine Jesu Christi Nazaréni surge, et ámbula. Et, apprehénsa manu ejus déxtera, allevávit eum, et prótinus consolidátæ sunt bases ejus, et plantæ. Et exsíliens stetit, et ambulábat; et intrávit cum illis in templum ámbulans, et exsíliens, et laudans Deum. Et vidit omnis pópulus eum ambulántem, et laudántem Deum. Cognoscébant autem illum, quod ipse erat, qui ad eleemósynam sedébat ad Speciósam portam templi; et impléti sunt stupóre et éxtasi in eo, quod contígerat illi.\par
\switchcolumn
\EN
\noindent But Peter said: Silver and gold I have none; but what I have, I give thee: In the name of Jesus Christ of Nazareth, arise, and walk. And taking him by the right hand, he lifted him up, and forthwith his feet and soles received strength. And he leaping up, stood, and walked, and went in with them into the temple, walking, and leaping, and praising God. And all the people saw him walking and praising God. And they knew him, that it was he who sat begging alms at the Beautiful gate of the temple: and they were filled with wonder and amazement at that which had happened to him.\par
\switchcolumn*

\LA
\begin{Resp}\Rsign Si díligis me, Simon Petre, pasce oves meas. Dómine, tu nosti quia amo te, \Star{} Et ánimam meam pono pro te.\end{Resp}
\begin{Resp}\Vsign Si oportúerit me mori tecum, non te negábo.\end{Resp}
\begin{Resp}\Rsign Et ánimam meam pono pro te.\end{Resp}
\switchcolumn
\EN
\begin{Resp}\Rsign Simon Peter, if thou lovest Me, feed My sheep. Lord, Thou knowest that I love thee; \Star{} I will lay down my life for thy sake.\end{Resp}
\begin{Resp}\Vsign If I should die with thee, I will not deny thee.\end{Resp}
\begin{Resp}\Rsign I will lay down my life for thy sake.\end{Resp}
\switchcolumn*

\LA
\LessonHead{Lectio III}
\switchcolumn
\EN
\LessonHead{Lesson III}
\switchcolumn*

\LA
\Cite{Act 3:11-16}
\switchcolumn
\EN
\Cite{Acts 3:11-16}
\switchcolumn*

\LA
\noindent Cum tenéret autem Petrum et Joánnem, cucúrrit omnis pópulus ad eos ad pórticum, quæ appellátur Salomónis, stupéntes. Videns autem Petrus, respóndit ad pópulum: Viri Israëlítæ, quid mirámini in hoc, aut nos quid intuémini, quasi nostra virtúte aut potestáte fecérimus hunc ambuláre? Deus Abraham, et Deus Isaac, et Deus Jacob, Deus patrum nostrórum glorificávit Fílium suum Jesum, quem vos quidem tradidístis et negástis ante fáciem Piláti, judicánte illo dimítti. Vos autem Sanctum et Justum negástis, et petístis virum homicídam donári vobis: Auctórem vero vitæ interfecístis, quem Deus suscitávit a mórtuis, cujus nos testes sumus. Et in fide nóminis ejus, hunc, quem vos vidístis et nostis, confirmávit nomen ejus: et fides, quæ per eum est, dedit íntegram sanitátem istam in conspéctu ómnium vestrum.\par
\switchcolumn
\EN
\noindent And as he held Peter and John, all the people ran to them to the porch which is called Solomon's, greatly wondering. But Peter seeing, made answer to the people: Ye men of Israel, why wonder you at this? or why look you upon us, as if by our strength or power we had made this man to walk? The God of Abraham, and the God of Isaac, and the God of Jacob, the God of our fathers, hath glorified his Son Jesus, whom you indeed delivered up and denied before the face of Pilate, when he judged he should be released. But you denied the Holy One and the Just, and desired a murderer to be granted unto you. But the author of life you killed, whom God hath raised from the dead, of which we are witnesses. And in the faith of his name, this man, whom you have seen and known, hath his name strengthened; and the faith which is by him, hath given this perfect soundness in the sight of you all.\par
\switchcolumn*

\LA
\begin{Resp}\Rsign Tu es Petrus, et super hanc petram ædificábo Ecclésiam meam, et portæ ínferi non prævalébunt advérsus eam: \Star{} Et tibi dabo claves regni cælórum.\end{Resp}
\begin{Resp}\Vsign Quodcúmque ligáveris super terram, erit ligátum et in cælis; et quodcúmque sólveris super terram, erit solútum et in cælis.\end{Resp}
\begin{Resp}\Rsign Et tibi dabo claves regni cælórum.\end{Resp}
\begin{Resp}\Vsign Glória Patri, et Fílio, \Star{} et Spirítui Sancto.\end{Resp}
\begin{Resp}\Rsign Et tibi dabo claves regni cælórum.\end{Resp}
\switchcolumn
\EN
\begin{Resp}\Rsign Thou art Peter, and upon this rock I will build My Church, and the gates of hell shall not prevail against it. \Star{} And I will give unto thee the keys of the kingdom of heaven.\end{Resp}
\begin{Resp}\Vsign Whatsoever thou shalt bind on earth, shall be bound in heaven; and whatsoever thou shalt loose on earth, shall be loosed in heaven.\end{Resp}
\begin{Resp}\Rsign And I will give unto thee the keys of the kingdom of heaven.\end{Resp}
\begin{Resp}\Vsign Glory be to the Father, and to the Son, \Star{} and to the Holy Ghost.\end{Resp}
\begin{Resp}\Rsign And I will give unto thee the keys of the kingdom of heaven.\end{Resp}
\switchcolumn*

\LA
\LessonHead{Lectio IV}
\switchcolumn
\EN
\LessonHead{Lesson IV}
\switchcolumn*

\LA
\LessonTitle{Sermo sancti Leónis Papæ}
\switchcolumn
\EN
\LessonTitle{From the Sermons of Pope St. Leo the Great}
\switchcolumn*

\LA
\Source{Sermo 1 in natali App. Petri et Pauli}
\switchcolumn
\EN
\Source{First for the Birthday of the Holy Apostles Peter and Paul}
\switchcolumn*

\LA
\noindent Omnium quidem sanctárum solemnitátum, dilectíssimi, totus mundus est párticeps, et uníus fídei píetas éxigit, ut quidquid pro salúte universórum gestum recólitur, commúnibus ubíque gáudiis celebrétur. Verúmtamen hodiérna festívitas, præter illam reveréntiam quam toto terrárum orbe proméruit, speciáli et própria nostræ Urbis exsultatióne veneránda est; ut, ubi præcipuórum Apostolórum glorificátus est éxitus, ibi in die martýrii eórum sit lætítiæ principátus. Isti enim sunt viri, per quos tibi Evangélium Christi, Roma, resplénduit; et, quæ eras magístra erróris, facta es discípula veritátis.\par
\switchcolumn
\EN
\noindent Dearly beloved brethren, in the joy of all the holy Feast-days the whole world is partaker. There is but one love of God, and whatsoever is solemnly called to memory, if it hath been done for the salvation of all, must needs be worth the honour of a joyful memorial at the hands of all. Nevertheless, this feast which we are keeping today, besides that world-wide worship which it doth of right get throughout all the earth, doth deserve from this city of ours an outburst of gladness altogether special and our own. In this place it was that the two chiefest of the Apostles did so right gloriously finish their race. And upon this day whereon they lifted up that their last testimony, let it be in this place that the memory thereof receiveth the chiefest of jubilant celebrations. O Rome, these twain are the men who brought the light of the Gospel of Christ to shine upon thee! These are they by whom thou, from being the teacher of lies, wast turned into a learner of the truth.\par
\switchcolumn*

\LA
\begin{Resp}\Rsign Dómine, si tu es, jube me veníre ad te super aquas. \Star{} Et exténdens manum apprehéndit eum, et dixit Jesus: Módicæ fídei, quare dubitásti?\end{Resp}
\begin{Resp}\Vsign Cumque vidísset ventum válidum veniéntem, tímuit; et, cum cœpísset mergi, clamávit dicens: Dómine, salvum me fac.\end{Resp}
\begin{Resp}\Rsign Et exténdens manum apprehéndit eum, et dixit Jesus: Módicæ fídei, quare dubitásti?\end{Resp}
\switchcolumn
\EN
\begin{Resp}\Rsign Lord, if it be Thou, bid me come unto thee on the water; \Star{} And Jesus stretched forth His Hand, and caught him, and said unto him: O thou of little faith, wherefore didst thou doubt?\end{Resp}
\begin{Resp}\Vsign But when he saw the wind boisterous, he was afraid and, beginning to sink, he cried, saying: Lord, save me!\end{Resp}
\begin{Resp}\Rsign And Jesus stretched forth His Hand, and caught him, and said unto him: O thou of little faith, wherefore didst thou doubt?\end{Resp}
\switchcolumn*

\LA
\LessonHead{Lectio V}
\switchcolumn
\EN
\LessonHead{Lesson V}
\switchcolumn*

\LA
\noindent Isti sunt patres tui veríque pastóres, qui te regnis cæléstibus inseréndam multo mélius multóque felícius condidérunt, quam illi quorum stúdio prima mœ́nium tuórum fundaménta locáta sunt; ex quibus is qui tibi nomen dedit, fratérna te cæde fœdávit. Isti sunt, qui te ad hanc glóriam provexérunt, ut gens sancta, pópulus eléctus, cívitas sacerdotális et régia, per sacram beáti Petri Sedem caput orbis effécta, látius præsidéres religióne divína quam dominatióne terréna. Quamvis enim, multis aucta victóriis, jus impérii tui terra maríque protúleris; minus tamen est quod tibi béllicus labor súbdidit, quam quod pax christiána subjécit.\par
\switchcolumn
\EN
\noindent These twain be thy fathers, these be in good sooth thy shepherds, these twain be they who laid for thee, as touching the kingdom of heaven, better and happier foundations, than did they that first planned thine earthly ramparts, wherefrom he that gave thee thy name took occasion to pollute thee with a brother's blood. These are they who have set on thine head this thy glorious crown, that thou art become an holy nation, a chosen people, a city both Priestly and Kingly, whom the Sacred Throne of blessed Peter hath exalted till thou art become the Lady of the world, unto whom the world-wide love for God hath conceded a broader lordship than is the possession of any mere earthly empire. Thou wast once waxen great by victories, until thy power was spread haughtily over land and sea, but thy power was narrower then which the toils of war had won for thee, than that thou now hast which hath been laid at thy feet by the peace of Christ.\par
\switchcolumn*

\LA
\begin{Resp}\Rsign Surge, Petre, et índue te vestiméntis tuis, áccipe fortitúdinem ad salvándas gentes: \Star{} Quia cecidérunt caténæ de mánibus tuis.\end{Resp}
\begin{Resp}\Vsign Angelus Dómini ástitit, et lumen refúlsit in habitáculo cárceris, percussóque látere Petri, excitávit eum, dicens: Surge, velóciter.\end{Resp}
\begin{Resp}\Rsign Quia cecidérunt caténæ de mánibus tuis.\end{Resp}
\switchcolumn
\EN
\begin{Resp}\Rsign Arise, O Peter! Cast thy garment about thee, gird thee with strength for the saving of the nations. \Star{} The chains are fallen off from thine hands.\end{Resp}
\begin{Resp}\Vsign The Angel of the Lord came upon him, and a light shined in the prison and he smote Peter on the side and raised him up, saying: Arise up quickly.\end{Resp}
\begin{Resp}\Rsign The chains are fallen off from thine hands.\end{Resp}
\switchcolumn*

\LA
\LessonHead{Lectio VI}
\switchcolumn
\EN
\LessonHead{Lesson VI}
\switchcolumn*

\LA
\noindent Disposítio namque divínitus óperi máxime congruébat, ut multa regna uno confœderaréntur império, et cito pérvios habéret pópulos prædicátio generális, quos uníus tenéret régimen civitátis. Hæc autem cívitas ignórans suæ provectiónis auctórem, cum pene ómnibus dominarétur géntibus, ómnium géntium serviébat erróribus; et magnam sibi videbátur assumpsísse religiónem, quia nullam respúerat falsitátem. Unde, quantum erat per diábolum tenácius illigáta, tantum per Christum est mirabílius absolúta.\par
\switchcolumn
\EN
\noindent It well suited for the doing of the work which God had decreed that the multitude of kingdoms should be bound together under one rule, and that so the universal preaching of the Gospel should find easier entry into all peoples, since all were governed by the empire of one city. But this city, knowing not Him, Who had been pleased to make her great, used her lordship over almost all nations to make herself the minister of all their falsehoods and seemed to herself exceeding godly because there was no false god whom she rejected. But the tighter that Satan had bound her, the more wondrous was the work of Christ in setting her free.\par
\switchcolumn*

\LA
\begin{Resp}\Rsign Tu es pastor óvium, princeps Apostolórum: tibi trádidit Deus ómnia regna mundi: \Star{} Et ídeo tibi tráditæ sunt claves regni cælórum.\end{Resp}
\begin{Resp}\Vsign Quodcúmque ligáveris super terram, erit ligátum et in cælis; et quodcúmque sólveris super terram, erit solútum et in cælis.\end{Resp}
\begin{Resp}\Rsign Et ídeo tibi tráditæ sunt claves regni cælórum.\end{Resp}
\begin{Resp}\Vsign Glória Patri, et Fílio, \Star{} et Spirítui Sancto.\end{Resp}
\begin{Resp}\Rsign Et ídeo tibi tráditæ sunt claves regni cælórum.\end{Resp}
\switchcolumn
\EN
\begin{Resp}\Rsign Thou art the Shepherd of the sheep, and the Prince of the Apostles, and unto thee hath God given all the kingdoms of the world. \Star{} Therefore unto thee hath He given the keys of the kingdom of heaven.\end{Resp}
\begin{Resp}\Vsign Whatsoever thou shalt bind on earth shall be bound in heaven; and whatsoever thou shalt loose on earth shall be loosed in heaven.\end{Resp}
\begin{Resp}\Rsign Therefore unto thee hath He given the keys of the kingdom of heaven.\end{Resp}
\begin{Resp}\Vsign Glory be to the Father, and to the Son, \Star{} and to the Holy Ghost.\end{Resp}
\begin{Resp}\Rsign Therefore unto thee hath He given the keys of the kingdom of heaven.\end{Resp}
\switchcolumn*

\LA
\LessonHead{Lectio VII}
\switchcolumn
\EN
\LessonHead{Lesson VII}
\switchcolumn*

\LA
\LessonTitle{Léctio sancti Evangélii secúndum Matthǽum}
\switchcolumn
\EN
\LessonTitle{From the Holy Gospel according to Matthew}
\switchcolumn*

\LA
\Cite{Matt 16:13-19}
\switchcolumn
\EN
\Cite{Matt 16:13-19}
\switchcolumn*

\LA
\noindent In illo témpore: Venit Jesus in partes Cæsaréæ Philíppi, et interrogábat discípulos suos, dicens: Quem dicunt hómines esse Fílium hóminis? Et réliqua.\par
\switchcolumn
\EN
\noindent At that time: Jesus came into the coasts of Caesarea Philippi, and He asked His disciples, saying: Who do men say that I, the Son of man, am? And so on.\par
\switchcolumn*

\LA
\LessonTitle{Homilía sancti Hierónymi Presbýteri}
\switchcolumn
\EN
\LessonTitle{Homily by St. Jerome, Priest (at Bethlehem.)}
\switchcolumn*

\LA
\Source{Lib. 3 Comment. in Matth. cap. 16}
\switchcolumn
\EN
\Source{Bk. iii Comment. on Matth. xvi}
\switchcolumn*

\LA
\noindent Pulchre intérrogat: Quem dicunt hómines esse Fílium hóminis? quia qui de fílio hóminis loquúntur, hómines sunt; qui vero divinitátem ejus intéllegunt, non hómines, sed dii appellántur. At illi dixérunt: Alii Joánnem Baptístam, álii autem Elíam. Miror quosdam intérpretes causas errórum inquírere singulórum, et disputatiónem longíssimam téxere, quare Dóminum nostrum Jesum Christum álii Joánnem putáverint, álii Elíam, álii Jeremíam aut unum ex prophétis; cum sic potúerint erráre in Elía et Jeremía, quo modo Heródes errávit in Joánne, dicens: Quem ego decollávi Joánnem, ipse surréxit a mórtuis, et virtútes operántur in eo.\par
\switchcolumn
\EN
\noindent Who do men say that I, the Son of man, am? This question is well put, for they who speak of Him as the Son of man are men, while they that know of Him that He is God are called not men but gods. And they said: Some say that Thou art John the Baptist, some, Elias. I marvel that some commentators have thought it worth their while to search into the origin of each of these blunders, and to engage in a discussion of weary length as to why some thought that our Lord Jesus Christ was John the Baptist, some, Elias and others, Jeremias, or one of the Prophets. Their blunders concerning Elias and Jeremias were but of a piece with Herod's concerning John the Baptist: It is John, whom I beheaded; he is risen from the dead and therefore mighty works do show forth themselves in him. (Mark vi. 16, 14.)\par
\switchcolumn*

\LA
\begin{Resp}\Rsign Ego pro te rogávi, Petre, ut non defíciat fides tua: \Star{} Et tu aliquándo convérsus confírma fratres tuos.\end{Resp}
\begin{Resp}\Vsign Caro et sanguis non revelávit tibi, sed Pater meus, qui est in cælis.\end{Resp}
\begin{Resp}\Rsign Et tu aliquándo convérsus confírma fratres tuos.\end{Resp}
\switchcolumn
\EN
\begin{Resp}\Rsign Peter, I have prayed for thee, that thy faith fail not; \Star{} And when thou art converted, strengthen thy brethren.\end{Resp}
\begin{Resp}\Vsign Flesh and blood hath not revealed it unto thee, but My Father Which is in heaven.\end{Resp}
\begin{Resp}\Rsign And when thou art converted, strengthen thy brethren.\end{Resp}
\switchcolumn*

\LA
\LessonHead{Lectio VIII}
\switchcolumn
\EN
\LessonHead{Lesson VIII}
\switchcolumn*

\LA
\noindent Vos autem quem me esse dícitis? Prudens lector, atténde quod, ex consequéntibus textúque sermónis, Apóstoli nequáquam hómines, sed dii appellántur. Cum enim dixísset: Quem dicunt hómines esse Fílium hóminis? subjécit: Vos autem quem me esse dícitis? Illis, quia hómines sunt, humána opinántibus, vos qui estis dii, quem me esse existimátis? Petrus ex persóna ómnium Apostolórum profitétur: Tu es Christus Fílius Dei vivi. Deum vivum appéllat, ad distinctiónem eórum deórum, qui putántur dii, sed mórtui sunt.\par
\switchcolumn
\EN
\noindent But who say ye that I am? Mark, discreet reader, from the context, that a distinction is here drawn between the Apostles and mere men. The Apostles are called gods. Who, asketh the Lord, do men say that I am, but, on the other hand, who say ye that I am? They, being but men, deal in human speculations, but ye that are gods, who be ye persuaded that I am. And then Peter, as the representative of all the Apostles, uttered the testimony: Thou art the Christ, the Son of the living God. He calleth God living, to mark the difference between Him and all other that be called gods, and who are indeed dead.\par
\switchcolumn*

\LA
\begin{Resp}\Rsign Quem dicunt hómines esse Filium hóminis? dixit Jesus discípulis suis. Respóndens Petrus dixit: Tu es Christus Fílius Dei vivi. \Star{} Et ego dico tibi, quia tu es Petrus, et super hanc petram ædificábo Ecclésiam meam.\end{Resp}
\begin{Resp}\Vsign Beátus es, Simon Bar-Jona, quia caro et sanguis non revelávit tibi, sed Pater meus qui est in cælis.\end{Resp}
\begin{Resp}\Rsign Et ego dico tibi, quia tu es Petrus, et super hanc petram ædificábo Ecclésiam meam.\end{Resp}
\begin{Resp}\Vsign Glória Patri, et Fílio, \Star{} et Spirítui Sancto.\end{Resp}
\begin{Resp}\Rsign Et ego dico tibi, quia tu es Petrus, et super hanc petram ædificábo Ecclésiam meam.\end{Resp}
\switchcolumn
\EN
\begin{Resp}\Rsign Jesus asked His disciples, saying: Who do men say that I, the Son of Man, am? Peter answered, and said: Thou art the Christ, the Son of the living God. \Star{} And I say unto thee, that thou art Peter, and upon this rock I will build My Church.\end{Resp}
\begin{Resp}\Vsign Blessed art thou, Simon Bar-Jona, for flesh and blood hath not revealed it unto thee, but My Father Which is in heaven.\end{Resp}
\begin{Resp}\Rsign And I say unto thee, that thou art Peter, and upon this rock I will build My Church.\end{Resp}
\begin{Resp}\Vsign Glory be to the Father, and to the Son, \Star{} and to the Holy Ghost.\end{Resp}
\begin{Resp}\Rsign And I say unto thee, that thou art Peter, and upon this rock I will build My Church.\end{Resp}
\switchcolumn*

\LA
\LessonHead{Lectio IX}
\switchcolumn
\EN
\LessonHead{Lesson IX}
\switchcolumn*

\LA
\noindent Respóndens autem Jesus, dixit ei: Beátus es, Simon Bar-Jona. Testimónio de se Apóstoli reddit vicem. Petrus díxerat: Tu es Christus Fílius Dei vivi; mercédem recépit vera conféssio: Beátus es, Simon Bar-Jona. Quare? Quia non revelávit tibi caro et sanguis, sed revelávit Pater. Quod caro et sanguis reveláre non pótuit, Spíritus Sancti grátia revelátum est. Ergo ex confessióne sortítur vocábulum, quod revelatiónem ex Spíritu Sancto hábeat, cujus et fílius appellándus sit. Síquidem Bar-Jona in nostra lingua sonat Fílius colúmbæ.\par
\switchcolumn
\EN
\noindent And Jesus answered and said unto him: Blessed art thou, Simon Bar-Jona. The Apostle having testified of the Lord, the Lord in turn testifieth of the Apostle. Peter had said: Thou art the Christ, the Son of the living God, and he received, in return for that his testimony to the truth, the words: Blessed art thou, Simon Bar-Jona. Why blessed? For flesh and blood hath not revealed it unto thee, but My Father. What flesh and blood could not reveal, the grace of the Holy Ghost had revealed. Meet for him therefore, because of his confession, is his name, as the name of one who hath revelation from the Holy Ghost, and therefore is called His son. Bar-Jona is, being interpreted, The-son-of-the-Dove.\par
\switchcolumn*

\LA
\TeDeum{Te Deum}
\switchcolumn
\EN
\TeDeum{Te Deum}
\switchcolumn*

\end{paracol}

\par\needspace{10\baselineskip}\addvspace{1.2em}{\centering\small\textsc{Die 30 Junii} · \textsc{30 June}\par}
\Office{In Commemoratione S. Pauli Apostoli}{Commemoration of St. Paul the Apostle}{Duplex majus}
\addcontentsline{toc}{subsection}{Die 30 Junii · In Commemoratione S. Pauli Apostoli \textit{— Commemoration of St. Paul the Apostle}}
\begin{paracol}{2}
\LA
\LessonHead{Lectio I}
\switchcolumn
\EN
\LessonHead{Lesson I}
\switchcolumn*

\LA
\LessonTitle{De Actibus Apostolórum}
\switchcolumn
\EN
\LessonTitle{Lesson from the Acts of the Apostles}
\switchcolumn*

\LA
\Cite{Act 13:1-4}
\switchcolumn
\EN
\Cite{Acts 13:1-4}
\switchcolumn*

\LA
\noindent Erant autem in Ecclésia quæ erat Antiochíæ, prophétæ et doctóres, in quibus Bárnabas et Simon, qui vocátur Niger, et Lucius Cyrenénsis et Mánahen, qui erat Heródis tetrárchæ collactáneus, et Saulus. Ministrántibus autem illis Dómino et jejunántibus, dixit illis Spíritus Sanctus: Segregáte mihi Saulum et Bárnabam, in opus ad quod assúmpsi eos. Tunc jejunántes et orántes imponentésque eis manus dimisérunt illos. Et ipsi quidem missi a Spíritu Sancto abiérunt Seleucíam et inde navigavérunt Cyprum.\par
\switchcolumn
\EN
\noindent Now there were in the church which was at Antioch, prophets and doctors, among whom was Barnabas, and Simon who was called Niger, and Lucius of Cyrene, and Manahen, who was the foster brother of Herod the tetrarch, and Saul. And as they were ministering to the Lord, and fasting, the Holy Ghost said to them: Separate me Saul and Barnabas, for the work whereunto I have taken them. Then they, fasting and praying, and imposing their hands upon them, sent them away. So they being sent by the Holy Ghost, went to Seleucia: and from thence they sailed to Cyprus.\par
\switchcolumn*

\LA
\begin{Resp}\Rsign Qui operátus est Petro in apostolátum, operátus est et mihi inter gentes: \Star{} Et cognovérunt grátiam Dei, quæ data est mihi.\end{Resp}
\begin{Resp}\Vsign Grátia Dei in me vácua non fuit, sed grátia ejus semper in me manet.\end{Resp}
\begin{Resp}\Rsign Et cognovérunt grátiam Dei, quæ data est mihi.\end{Resp}
\switchcolumn
\EN
\begin{Resp}\Rsign He That wrought effectually in Peter to the Apostleship, the Same was mighty in me toward the Gentiles, \Star{} And they perceived the grace that was given unto me of the Lord Christ.\end{Resp}
\begin{Resp}\Vsign The grace of God which was bestowed upon me was not in vain, but His grace abideth ever in me.\end{Resp}
\begin{Resp}\Rsign And they perceived the grace that was given unto me of the Lord Christ.\end{Resp}
\switchcolumn*

\LA
\LessonHead{Lectio II}
\switchcolumn
\EN
\LessonHead{Lesson II}
\switchcolumn*

\LA
\Cite{Act 13:5-8}
\switchcolumn
\EN
\Cite{Acts 13:5-8}
\switchcolumn*

\LA
\noindent Et, cum veníssent Salamínam, prædicábant verbum Dei in synagógis Judæórum; habébant autem et Joánnem in ministério. Et, cum perambulássent univérsam ínsulam usque Paphum, invenérunt quemdam virum magum pseudoprophétam Judǽum, cui nomen erat Barjésu, qui erat cum procónsule Sérgio Paulo, viro prudénte. Hic, accersítis Bárnaba et Saulo, desiderábat audíre verbum Dei. Resistébat autem illis Elymas magus (sic enim interpretátur nomen ejus) quærens avértere procónsulem a fide.\par
\switchcolumn
\EN
\noindent And when they were come to Salamina, they preached the word of God in the synagogues of the Jews. And they had John also in the ministry. And when they had gone through the whole island, as far as Paphos, they found a certain man, a magician, a false prophet, a Jew, whose name was Bar-jesu: Who was with the proconsul Sergius Paulus, a prudent man. He sending for Barnabas and Saul, desired to hear the word of God. But Elymas the magician (for so his name is interpreted) withstood them, seeking to turn away the proconsul from the faith.\par
\switchcolumn*

\LA
\begin{Resp}\Rsign Bonum certámen certávi, cursum consummávi, fidem servávi: \Star{} Ideóque repósita est mihi coróna justítiæ.\end{Resp}
\begin{Resp}\Vsign Scio cui crédidi, et certus sum quia potens est depósitum meum serváre in illum diem.\end{Resp}
\begin{Resp}\Rsign Ideóque repósita est mihi coróna justítiæ.\end{Resp}
\switchcolumn
\EN
\begin{Resp}\Rsign I have fought a good fight, I have finished my course, I have kept the faith; \Star{} Henceforth there is laid up for me a crown of righteousness.\end{Resp}
\begin{Resp}\Vsign I know Whom I have believed, and am persuaded that He is able to keep that which I have committed unto Him against that day.\end{Resp}
\begin{Resp}\Rsign Henceforth there is laid up for me a crown of righteousness.\end{Resp}
\switchcolumn*

\LA
\LessonHead{Lectio III}
\switchcolumn
\EN
\LessonHead{Lesson III}
\switchcolumn*

\LA
\Cite{Act 13:9-13}
\switchcolumn
\EN
\Cite{Acts 13:9-13}
\switchcolumn*

\LA
\noindent Saulus autem, qui et Paulus, replétus Spíritu Sancto, íntuens in eum, dixit: O plene omni dolo et omni fallácia, fili diáboli, inimíce omnis justítiæ, non désinis subvértere vias Dómini rectas? Et nunc ecce manus Dómini super te, et eris cæcus, non videns solem usque ad tempus. Et conféstim cécidit in eum calígo et ténebræ, et circúiens quærébat qui ei manum daret. Tunc procónsul, cum vidísset factum, crédidit admírans super doctrína Dómini. Et, cum a Papho navigássent, Paulus et qui cum eo erant venérunt Pergen Pamphýliæ; Joánnes autem discédens ab eis revérsus est Jerosólymam.\par
\switchcolumn
\EN
\noindent Then Saul, otherwise Paul, filled with the Holy Ghost, looking upon him, Said: O full of all guile, and of all deceit, child of the devil, enemy of all justice, thou ceasest not to pervert the right ways of the Lord. And now behold, the hand of the Lord is upon thee, and thou shalt be blind, not seeing the sun for a time. And immediately there fell a mist and darkness upon him, and going about, he sought some one to lead him by the hand. Then the proconsul, when he had seen what was done, believed, admiring at the doctrine of the Lord. Now when Paul and they that were with him had sailed from Paphos, they came to Perge in Pamphylia. And John departing from them, returned to Jerusalem.\par
\switchcolumn*

\LA
\begin{Resp}\Rsign Repósita est mihi coróna justítiæ, \Star{} Quam reddet mihi Dóminus in illum diem justus judex.\end{Resp}
\begin{Resp}\Vsign Scio cui crédidi, et certus sum quia potens est depósitum meum serváre in illum diem.\end{Resp}
\begin{Resp}\Rsign Quam reddet mihi Dóminus in illum diem justus judex.\end{Resp}
\begin{Resp}\Vsign Glória Patri, et Fílio, \Star{} et Spirítui Sancto.\end{Resp}
\begin{Resp}\Rsign Quam reddet mihi Dóminus in illum diem justus judex.\end{Resp}
\switchcolumn
\EN
\begin{Resp}\Rsign There is laid up for me a crown of righteousness: \Star{} Which the Lord, the righteous Judge, shall give me at that day.\end{Resp}
\begin{Resp}\Vsign I know Whom I have believed, and am persuaded that He is able to keep that which I have committed unto Him against that day.\end{Resp}
\begin{Resp}\Rsign Which the Lord, the righteous Judge, shall give me at that day.\end{Resp}
\begin{Resp}\Vsign Glory be to the Father, and to the Son, \Star{} and to the Holy Ghost.\end{Resp}
\begin{Resp}\Rsign Which the Lord, the righteous Judge, shall give me at that day.\end{Resp}
\switchcolumn*

\LA
\LessonHead{Lectio IV}
\switchcolumn
\EN
\LessonHead{Lesson IV}
\switchcolumn*

\LA
\LessonTitle{Ex libro sancti Augustíni Epíscopi de grátia et líbero arbítrio}
\switchcolumn
\EN
\LessonTitle{From the Book of St. Augustine, Bishop of Hippo, touching Grace and Free-will}
\switchcolumn*

\LA
\Source{Cap. 6-7}
\switchcolumn
\EN
\Source{Chapter 6}
\switchcolumn*

\LA
\noindent Apóstolus Paulus, quem certe invénimus sine ullis méritis bonis, immo cum multis méritis malis, Dei grátiam consecútum, reddéntis bona pro malis, videámus quid dicat, sua jam propinquánte passióne, scribens ad Timótheum: Ego enim jam ímmolor, inquit, et tempus resolutiónis meæ instat. Bonum certámen certávi, cursum consummávi, fidem servávi. Ista útique jam mérita sua bona commémorat; ut post bona mérita consequátur corónam, qui post mérita mala consecútus est grátiam. Dénique atténdite quid sequátur: Súperest, inquit, mihi coróna justítiæ, quam reddet mihi Dóminus in illa die justus judex. Cui rédderet corónam justus judex, si non donásset grátiam miséricors Pater? Et quómodo esset ista coróna justítiæ, nisi præcessísset grátia, quæ justíficat ímpium? Quómodo ista débita redderétur, nisi prius illa gratúita donarétur?\par
\switchcolumn
\EN
\noindent The Apostle Paul was a man, who, when we first hear of him, had not only no merits, but a great many demerits. That man received the grace of God, Who returneth good for evil, and let us see in what sort of language, when the hour of his last sufferings was at hand, he wrote to Timothy. He saith: I am now ready to be offered, and the time of my departure is at hand. I have fought a good fight, I have finished my course, I have kept the faith. (2 Tim. iii. 6.) Here he counteth his merits, whereon a crown was immediately to follow, just as grace had followed immediately on his demerits. Listen to what cometh next Henceforth there is laid up for me a crown of righteousness, which the Lord, the righteous Judge, shall give me at that day. Unto whom would the Lord give a crown as a righteous Judge, if He had not first given grace, as a merciful Father And how would that crown be a crown of righteousness, if there had not first come grace which justifieth the ungodly How could a reward have been earned unless the power to earn had first been given unearned?\par
\switchcolumn*

\LA
\begin{Resp}\Rsign Tu es vas electiónis, sancte Paule Apóstole, prædicator veritátis in univérso mundo, \Star{} Per quem omnes gentes cognovérunt grátiam Dei.\end{Resp}
\begin{Resp}\Vsign Intercéde pro nobis ad Deum, qui te elégit.\end{Resp}
\begin{Resp}\Rsign Per quem omnes gentes cognovérunt grátiam Dei.\end{Resp}
\switchcolumn
\EN
\begin{Resp}\Rsign O Holy Apostle Paul, thou art a chosen vessel unto God, to preach the Gospel throughout the whole world \Star{} Through whom all nations have known the grace of God.\end{Resp}
\begin{Resp}\Vsign Pray for us to God Who hath chosen thee.\end{Resp}
\begin{Resp}\Rsign Through whom all nations have known the grace of God.\end{Resp}
\switchcolumn*

\LA
\LessonHead{Lectio V}
\switchcolumn
\EN
\LessonHead{Lesson V}
\switchcolumn*

\LA
\noindent Proínde considerémus ipsa mérita Apóstoli Pauli, quibus dixit corónam redditúrum júdicem justum, et videámus utrum mérita ipsíus tamquam ipsíus, id est, ex ipso ei comparáta, an dona sint Dei. Bonum, inquit, certámen certávi, cursum consummávi, fidem servávi. Primo ista bona ópera, si non ea præcessíssent cogitatiónes bonæ, nulla essent. Atténdite ítaque quid de ipsis cogitatiónibus dicat; ait enim scribens ad Corínthios: Non quia idónei sumus cogitáre áliquid a nobis, tamquam a nobismetípsis; sed sufficiéntia nostra ex Deo est. Deínde síngula inspiciámus.\par
\switchcolumn
\EN
\noindent Let us now consider what were the merits of the Apostle Paul, which entitled him to look for a crown of righteousness from the Lord, the righteous Judge, and let us see whether these merits sprang from himself or were God's gifts to him. He saith I have fought a good fight, I have finished my course, I have kept the faith. To begin with, these good works would have been nothing worth, unless they had come from good thoughts. Listen therefore to what this same Paul saith concerning good thoughts. He writeth to the Corinthians: Not that we are sufficient of ourselves to think anything, as of ourselves but our sufficiency\par
is of God. (2 Cor. iii. 5.) Now let us take point by point.\par
\switchcolumn*

\LA
\begin{Resp}\Rsign Grátia Dei sum id quod sum: \Star{} Et grátia ejus in me vacua non fuit, sed semper in me manet.\end{Resp}
\begin{Resp}\Vsign Qui operátus est Petro in apostolátum, operátus est et mihi inter gentes.\end{Resp}
\begin{Resp}\Rsign Et grátia ejus in me vacua non fuit, sed semper in me manet.\end{Resp}
\switchcolumn
\EN
\begin{Resp}\Rsign By the grace of God I am what I am. \Star{} And His grace which was bestowed upon me was not in vain, but abideth ever in me.\end{Resp}
\begin{Resp}\Vsign He that wrought effectually in Peter to the Apostleship, the Same was mighty in me toward the Gentiles.\end{Resp}
\begin{Resp}\Rsign And His grace which was bestowed upon me was not in vain, but abideth ever in me.\end{Resp}
\switchcolumn*

\LA
\LessonHead{Lectio VI}
\switchcolumn
\EN
\LessonHead{Lesson VI}
\switchcolumn*

\LA
\noindent Bonum, inquit, certámen certávi. Quæro qua virtúte certáverit, utrum quæ illi ex semetípso fúerit, an quæ désuper data sit? Sed absit ut tantus doctor ignoráverit legem Dei, cujus vox est in Deuteronómio: Ne dicas in corde tuo: Fortitúdo mea et poténtia manus meæ fecit mihi virtútem magnam hanc; sed memoráberis Dómini Dei tui, quia ipse tibi dat fortitúdinem fácere virtútem. Quid autem prodest bonum certámen, nisi sequátur victória? Et quis dat victóriam, nisi ille, de quo dicit ipse: Grátias Deo, qui dat nobis victóriam per Dóminum nostrum Jesum Christum?\par
\switchcolumn
\EN
\noindent He saith: I have fought a good fight. I should like to know in whose strength he fought; in his own strength or in strength given him from above? God forbid that this great teacher should be supposed not to have known the Law of his God, Who saith in the Book of Deuteronomy: Say not in thine heart: My power and the might of mine hand hath gotten me this wealth. But thou shalt remember the Lord thy God, for it is He That giveth thee power to get wealth. (viii. 17, 18.) And, again, what is the good of a fight except it end in victory? And who is he that giveth victory save He of Whom this very same Paul saith: Thanks be to God, Which giveth us the victory, through our Lord Jesus Christ! (1 Cor. xv. 57.)\par
\switchcolumn*

\LA
\begin{Resp}\Rsign Saulus, qui et Paulus, magnus prædicátor, \Star{} A Deo confortátus convalescébat, et confundébat Judæos.\end{Resp}
\begin{Resp}\Vsign Osténdens quia hic est Christus, Fílius Dei.\end{Resp}
\begin{Resp}\Rsign A Deo confortátus convalescébat, et confundébat Judæos.\end{Resp}
\begin{Resp}\Vsign Glória Patri, et Fílio, \Star{} et Spirítui Sancto.\end{Resp}
\begin{Resp}\Rsign A Deo confortátus convalescébat, et confundébat Judæos.\end{Resp}
\switchcolumn
\EN
\begin{Resp}\Rsign Saul, who also is called Paul, was made a great preacher \Star{} And being of God increased the more in strength, he confounded the Jews.\end{Resp}
\begin{Resp}\Vsign Proving that This is very Christ, the Son of God.\end{Resp}
\begin{Resp}\Rsign And being of God increased the more in strength, he confounded the Jews.\end{Resp}
\begin{Resp}\Vsign Glory be to the Father, and to the Son, \Star{} and to the Holy Ghost.\end{Resp}
\begin{Resp}\Rsign And being of God increased the more in strength, he confounded the Jews.\end{Resp}
\switchcolumn*

\LA
\LessonHead{Lectio VII}
\switchcolumn
\EN
\LessonHead{Lesson VII}
\switchcolumn*

\LA
\LessonTitle{Léctio sancti Evangélii secúndum Matthǽum}
\switchcolumn
\EN
\LessonTitle{From the Holy Gospel according to Matthew}
\switchcolumn*

\LA
\Cite{Matt 10:16-22}
\switchcolumn
\EN
\Cite{Matt 10:16-22}
\switchcolumn*

\LA
\noindent In illo témpore: Dixit Jesus discípulis suis: Ecce ego mitto vos sicut oves in médio lupórum. Et réliqua.\par
\switchcolumn
\EN
\noindent At that time, Jesus said unto His disciples: Behold, I send you forth as sheep in the midst of wolves. And so on.\par
\switchcolumn*

\LA
\LessonTitle{Homilía sancti Joánnis Chrysóstomi}
\switchcolumn
\EN
\LessonTitle{Homily by St. John Chrysostom, Patriarch of Constantinople}
\switchcolumn*

\LA
\Source{Homilia 34 in Matth., longe post initium}
\switchcolumn
\EN
\Source{34 on Matthew}
\switchcolumn*

\LA
\noindent Sed dícere vidétur: Nolíte turbári, si, cum vos inter lupos mitto, tamquam oves et colúmbas esse júbeo. Nam, etsi possum contrárium quoque præstáre, et non permíttere ut grave áliquid patiámini, nec lupis tamquam oves subjécti sitis, sed effícere ut leónibus terribilióres evadátis; tamen sic éxpedit fíeri: hoc et vos quoque illustrióres fáciet, et meam declarábit virtútem. Sic enim póstea dixit Paulo: Súfficit tibi grátia mea, nam virtus mea in infirmitáte perfícitur. Ipse ígitur vos ut tales essétis feci.\par
\switchcolumn
\EN
\noindent It is as though He said Let not your heart be troubled although when I send you forth as sheep in the midst of wolves, I bid you be harmless as doves for, albeit, if I would, I could now make things otherwise, and suffer not that ye should have to bear anything grievous, neither be at the mercy of the wolves as are other sheep, but on the contrary, could make you more dreadful to the lions than the lions to you nevertheless, thus must it needs be, and yourselves it will make more glorious, and My power it will wholly show forth. For thus was it that afterwards the Same Lord said unto Paul: My grace is sufficient for thee, for My strength is made perfect in weakness. (2 Cor. xii. 9.) It is I that have made you to be what ye are.\par
\switchcolumn*

\LA
\begin{Resp}\Rsign Sancte Paule Apóstole, prædicátor veritátis et doctor géntium, \Star{} Intercéde pro nobis ad Deum, qui te elégit, ut digni efficiámur grátia Dei.\end{Resp}
\begin{Resp}\Vsign Tu es vas electiónis, sancte Paule Apostole, prædicátor veritátis.\end{Resp}
\begin{Resp}\Rsign Intercéde pro nobis ad Deum, qui te elégit, ut digni efficiámur grátia Dei.\end{Resp}
\switchcolumn
\EN
\begin{Resp}\Rsign O Holy Apostle Paul, Preacher of the truth, and teacher of the Gentiles, \Star{} Pray for us to God, Who hath chosen thee, that we may be made worthy of the grace of God.\end{Resp}
\begin{Resp}\Vsign O Holy Apostle Paul, thou art a chosen vessel unto God, and a Preacher of the truth.\end{Resp}
\begin{Resp}\Rsign Pray for us to God Who hath chosen thee, that we may be made worthy of the grace of God.\end{Resp}
\switchcolumn*

\LA
\LessonHead{Lectio VIII}
\switchcolumn
\EN
\LessonHead{Lesson VIII}
\switchcolumn*

\LA
\noindent Sed inspiciámus quam prudéntiam éxigit: serpéntis certe. Nam, quemádmodum serpens totum seípsum tradit, nec mínimum curat si ipsum quoque corpus incídi necésse sit, dúmmodo caput suum íntegrum servet; eódem tu quoque modo, præter fidem cétera pérdere non cures, sive pecúnias, sive corpus, sive étiam vitam ipsam profúndere necésse sit. Fides enim caput est et radix; qua serváta, etiámsi ómnia perdas, ómnia tamen rursus majóre cum glória recuperábis. Idcírco nec símplices solum jussit esse, nec prudéntes solum; sed ambo hæc in unum míscuit, ut ea in virtútem convertántur.\par
\switchcolumn
\EN
\noindent But let us look what wisdom it is which the Lord requireth. It is the wisdom of the serpent. The serpent draweth all the rest of his body after his head, and it is no matter to him if his body be cut through, so long as he keepeth his head unharmed. Thus, O Christian, is it with thee. It is no matter to thee that for thy faith's sake thou shouldest lose all things else money, or body, or, if need be, life itself. thy faith is thy head, and the root of thy being hold fast to that, and, as long as thou hast that, although thou shouldest lose all things else, it will only be to receive them back again with interest an hundredfold. And thus it is that the Lord biddeth us, not to be single-hearted only, nor wise only, but both together, that therefrom we may be strong.\par
\switchcolumn*

\LA
\begin{Resp}\Rsign Damasci, præpósitus gentis Aretæ regis vóluit me comprehéndere: \Star{} Et a frátribus per murum demíssus sum in sporta, * Et sic evasi manus ejus in nómine Dómini.\end{Resp}
\begin{Resp}\Vsign Deus et Pater Dómini nostri Jesu Christi scit quia non mentior.\end{Resp}
\begin{Resp}\Rsign Et a frátribus per murum demissus sum in sporta.\end{Resp}
\begin{Resp}\Vsign Glória Patri, et Fílio, \Star{} et Spirítui Sancto.\end{Resp}
\begin{Resp}\Rsign Et sic evasi manus ejus in nómine Dómini.\end{Resp}
\switchcolumn
\EN
\begin{Resp}\Rsign In Damascus the governor under Aretas the king was desirous to apprehend me \Star{} And by the brethren in a basket was I let down by the wall, and so escaped I his hands, in the name of the Lord.\end{Resp}
\begin{Resp}\Vsign The God and Father of our Lord Jesus Christ knoweth that I lie not.\end{Resp}
\begin{Resp}\Rsign And by the brethren in a basket was I let down by the wall.\end{Resp}
\begin{Resp}\Vsign Glory be to the Father, and to the Son, \Star{} and to the Holy Ghost.\end{Resp}
\begin{Resp}\Rsign And by the brethren in a basket was I let down by the wall.\end{Resp}
\switchcolumn*

\LA
\LessonHead{Lectio IX}
\switchcolumn
\EN
\LessonHead{Lesson IX}
\switchcolumn*

\LA
\noindent Quod si rebus ipsis id ita fíeri vidére desíderas, lege Actuum Apostolórum librum; perspícies profécto, cum sæpe Judæórum pópulus in Apóstolos insurréxerit ac dentes exacúerit, illos, colúmbæ simplicitátem imitándo et cum decénti modéstia respondéndo, iram ipsórum superásse, furórem exstinxísse, ímpetum retardásse. Nam, cum illi dícerent: Nonne præcipiéndo præcépimus vobis, ne docerétis in nómine isto? quamvis innúmera possent édere mirácula, nihil tamen ásperum neque dixérunt neque fecérunt; sed summa cum mansuetúdine respondéntes dicébant: Si justum est vos audíre magis quam Deum, judicáte. Perspexísti simplicitátem colúmbæ, vide nunc serpéntis prudéntiam: Non enim póssumus, ínquiunt, nos quæ vídimus et audívimus, non loqui.\par
\switchcolumn
\EN
\noindent If thou wilt see how these words were brought to the proof in very deed, read the Book of the Acts of the Apostles. There thou wilt see how that oftentimes the Jewish people rose against the Apostles and gnashed on them with their teeth, but they, with dove-like guilelessness, gave them smooth answers, and turned away their wrath, and quenched their fury, and stopped their onset. When the Jews said Did not we straitly command you, that ye should not teach in this Name? although the Apostles could have worked any miracles they chose, yet they neither said nor did anything sharp, but answered them with all meekness Whether it be right in the sight of God to hearken unto you more than unto God, judge ye. (Acts v. 28 iv. 19.) Here thou hast the harmlessness of doves listen now to the wisdom of serpents We cannot but speak the things which we have seen and heard. (20.)\par
\switchcolumn*

\LA
\TeDeum{Te Deum}
\switchcolumn
\EN
\TeDeum{Te Deum}
\switchcolumn*

\end{paracol}

\SeasonTitle{Mensis Julius}{July}

\addcontentsline{toc}{section}{Mensis Julius \textit{— July}}

\par\needspace{10\baselineskip}\addvspace{1.2em}{\centering\small\textsc{Die 1 Julii} · \textsc{1 July}\par}
\Office{Pretiosissimi Sanguinis Domini Nostri Jesu Christi}{Most Precious Blood of Our Lord Jesus Christ}{Duplex I classis}
\addcontentsline{toc}{subsection}{Die 1 Julii · Pretiosissimi Sanguinis Domini Nostri Jesu Christi \textit{— Most Precious Blood of Our Lord Jesus Christ}}
\begin{paracol}{2}
\LA
\RefLine{Lectio IX: de In Octava S. Joannis Baptistæ (07-01t).}
\switchcolumn
\EN
\RefLine{Lesson IX: of Octave of the Nativity of St. John the Baptist (07-01t).}
\switchcolumn*

\LA
\LessonHead{Lectio I}
\switchcolumn
\EN
\LessonHead{Lesson I}
\switchcolumn*

\LA
\LessonTitle{De Epístola beáti Pauli Apóstoli ad Hebrǽos}
\switchcolumn
\EN
\LessonTitle{Lesson from the letter of St. Paul the Apostle to the Hebrews}
\switchcolumn*

\LA
\Cite{Heb 9:11-15}
\switchcolumn
\EN
\Cite{Heb 9:11-15}
\switchcolumn*

\LA
\noindent Christus assístens póntifex futurórum bonórum, per ámplius et perféctius tabernáculum, non manufáctum, id est, non hujus creatiónis: neque per sánguinem hircórum aut vitulórum, sed per próprium sánguinem introívit semel in Sancta, ætérna redemptióne invénta. Si enim sanguis hircórum et taurórum, et cinis vítulæ aspérsus inquinátos sanctíficat ad emundatiónem carnis: quanto magis sanguis Christi, qui per Spíritum Sanctum semetípsum óbtulit immaculátum Deo, emundábit consciéntiam nostram ab opéribus mórtuis, ad serviéndum Deo vivénti? Et ídeo novi testaménti mediátor est: ut morte intercedénte, in redemptiónem eárum prævaricatiónum, quæ erant sub prióri testaménto, repromissiónem accípiant qui vocáti sunt ætérnæ hereditátis.\par
\switchcolumn
\EN
\noindent But Christ, being come an high priest of the good things to come, by a greater and more perfect tabernacle not made with hand, that is, not of this creation: Neither by the blood of goats, or of calves, but by his own blood, entered once into the holies, having obtained eternal redemption. For if the blood of goats and of oxen, and the ashes of an heifer being sprinkled, sanctify such as are defiled, to the cleansing of the flesh: How much more shall the blood of Christ, who by the Holy Ghost offered himself unspotted unto God, cleanse our conscience from dead works, to serve the living God? And therefore he is the mediator of the new testament: that by means of his death, for the redemption of those transgressions, which were under the former testament, they that are called may receive the promise of eternal inheritance.\par
\switchcolumn*

\LA
\begin{Resp}\Rsign Jesus, ut sanctificáret per suum sánguinem pópulum, extra portam passus est: \Star{} Exeámus ígitur ad eum extra castra, impropérium ejus portántes.\end{Resp}
\begin{Resp}\Vsign Nondum enim usque ad sánguinem restitístis advérsus peccátum repugnántes.\end{Resp}
\begin{Resp}\Rsign Exeámus ígitur ad eum extra castra, impropérium ejus portántes.\end{Resp}
\switchcolumn
\EN
\begin{Resp}\Rsign Jesus also, that He might sanctify the people with His own Blood, suffered without the gate. \Star{} Let us go forth therefore unto Him without the camp, bearing His reproach.\end{Resp}
\begin{Resp}\Vsign Ye have not yet resisted unto blood, striving against sin.\end{Resp}
\begin{Resp}\Rsign Let us go forth therefore unto Him without the gate, bearing His reproach.\end{Resp}
\switchcolumn*

\LA
\LessonHead{Lectio II}
\switchcolumn
\EN
\LessonHead{Lesson II}
\switchcolumn*

\LA
\Cite{Heb 9:16-22}
\switchcolumn
\EN
\Cite{Heb 9:16-22}
\switchcolumn*

\LA
\noindent Ubi enim testaméntum est, mors necésse est intercédat testatóris. Testaméntum enim in mórtuis confirmátum est: alióquin nondum valet, dum vivit qui testátus est. Unde nec primum quidem sine sánguine dedicátum est. Lecto enim omni mandáto legis a Móyse univérso pópulo, accípiens sánguinem vitulórum et hircórum cum aqua, et lana coccínea, et hyssópo, ipsum quoque librum, et omnem pópulum aspérsit, dicens: Hic sanguis testaménti, quod mandávit ad vos Deus. Etiam tabernáculum et ómnia vasa ministérii sánguine simíliter aspérsit. Et ómnia pene in sánguine secúndum legem mundántur: et sine sánguinis effusióne non fit remíssio.\par
\switchcolumn
\EN
\noindent For where there is a testament, the death of the testator must of necessity come in. For a testament is of force, after men are dead: otherwise it is as yet of no strength, whilst the testator liveth. Whereupon neither was the first indeed dedicated without blood. For when every commandment of the law had been read by Moses to all the people, he took the blood of calves and goats, with water, and scarlet wool and hyssop, and sprinkled both the book itself and all the people, Saying: This is the blood of the testament, which God hath enjoined unto you. The tabernacle also and all the vessels of the ministry, in like manner, he sprinkled with blood. And almost all things, according to the law, are cleansed with blood: and without shedding of blood there is no remission.\par
\switchcolumn*

\LA
\begin{Resp}\Rsign Móyses sumptum sánguinem respérsit in pópulum: \Star{} Et ait: Hic est sanguis fœ́deris, quod pépigit Dóminus vobíscum.\end{Resp}
\begin{Resp}\Vsign Fide celebrávit Pascha et sánguinis effusióne, ne qui vastábat primitíva, tángeret eos.\end{Resp}
\begin{Resp}\Rsign Et ait: Hic est sanguis fœ́deris, quod pépigit Dóminus vobíscum.\end{Resp}
\switchcolumn
\EN
\begin{Resp}\Rsign Moses took the blood, and sprinkled all the people; \Star{} Saying: This is the blood of the Testament which God hath enjoined unto you.\end{Resp}
\begin{Resp}\Vsign Through faith he kept the Passover, and the sprinkling of blood, lest he that destroyed the first-born should touch them.\end{Resp}
\begin{Resp}\Rsign Saying: This is the blood of the Testament which God hath enjoined unto you.\end{Resp}
\switchcolumn*

\LA
\LessonHead{Lectio III}
\switchcolumn
\EN
\LessonHead{Lesson III}
\switchcolumn*

\LA
\Cite{Heb 10:19-24}
\switchcolumn
\EN
\Cite{Heb 10:19-24}
\switchcolumn*

\LA
\noindent Habéntes ítaque, fratres, fidúciam in intróitu sanctórum in sánguine Christi, quam initiávit nobis viam novam, et vivéntem per velámen, id est, carnem suam, et sacerdótem magnum super domum Dei: accedámus cum vero corde in plenitúdine fídei, aspérsi corda a consciéntia mala, et ablúti corpus aqua munda, teneámus spei nostræ confessiónem indeclinábilem (fidélis enim est qui repromísit), et considerémus ínvicem in provocatiónem caritátis, et bonórum óperum.\par
\switchcolumn
\EN
\noindent Having therefore, brethren, a confidence in the entering into the holies by the blood of Christ; A new and living way which he hath dedicated for us through the veil, that is to say, his flesh, And a high priest over the house of God: Let us draw near with a true heart in fulness of faith, having our hearts sprinkled from an evil conscience, and our bodies washed with clean water. Let us hold fast the confession of our hope without wavering (for he is faithful that hath promised), And let us consider one another, to provoke unto charity and to good works:\par
\switchcolumn*

\LA
\begin{Resp}\Rsign Vos, qui aliquándo erátis longe, facti estis prope in sánguine Christi: \Star{} Ipse enim est pax nostra, qui fecit útraque unum.\end{Resp}
\begin{Resp}\Vsign Complácuit per eum reconciliáre ómnia in ipsum, pacíficans per sánguinem crucis ejus, sive quæ in cælis sunt.\end{Resp}
\begin{Resp}\Rsign Ipse enim est pax nostra, qui fecit útraque unum.\end{Resp}
\begin{Resp}\Vsign Glória Patri, et Fílio, \Star{} et Spirítui Sancto.\end{Resp}
\begin{Resp}\Rsign Ipse enim est pax nostra, qui fecit útraque unum.\end{Resp}
\switchcolumn
\EN
\begin{Resp}\Rsign Ye, who sometimes were far off, are made nigh by the Blood of Christ. \Star{} For He is our Peace, Who hath made both one.\end{Resp}
\begin{Resp}\Vsign It pleased [the Father that in Him should all fulness dwell, and,] having made peace through the Blood of His Cross, by Him to reconcile all things unto Himself, [by Him, I say,] whether they be things in earth or things in heaven.\end{Resp}
\begin{Resp}\Rsign For He is our Peace, Who hath made both one.\end{Resp}
\begin{Resp}\Vsign Glory be to the Father, and to the Son, \Star{} and to the Holy Ghost.\end{Resp}
\begin{Resp}\Rsign For He is our Peace, Who hath made both one.\end{Resp}
\switchcolumn*

\LA
\LessonHead{Lectio IV}
\switchcolumn
\EN
\LessonHead{Lesson IV}
\switchcolumn*

\LA
\LessonTitle{Sermo sancti Joánnis Chrysóstomi}
\switchcolumn
\EN
\LessonTitle{From the Sermons of St. John Chrysostom, Archbishop of Constantinople}
\switchcolumn*

\LA
\Source{Homilia ad Neophytos}
\switchcolumn
\EN
\Source{Homily 84 in John, cap. 19}
\switchcolumn*

\LA
\noindent Vis sánguinis Christi audíre virtútem? Redeámus ad ejus exémplum, et priórem typum recordémur, et prístinam Scriptúram narrémus. In Ægýpto, nocte média, Ægýptiis Deus plagam décimam minabátur, ut eórum primogénita deperírent, quia primogénitum ejus pópulum detinébant. Sed, ne amáta plebs Judæórum una cum illis periclitarétur, quia unus locus continébat univérsos, remédium discretiónis invéntum est. Proínde exémplum mirábile, ut discas in veritáte virtútem. Ira divínæ indignatiónis operabátur, et domos síngulas mórtifer circuíbat. Quid ígitur Móyses? Occídite, inquit, agnum annículum, et sánguine ejus liníte jánuas. Quid ais, Móyses? Sanguis ovis rationálem hóminem liberáre consuévit? Valde, inquit; non eo quod sanguis est, sed quia Domínici sánguinis per eum demonstrátur exémplum.\par
\switchcolumn
\EN
\noindent Wouldest thou hear the power of the Blood of Christ? Then let us look at the figure thereof, let us call to mind the old type, and tell the story written in the ancient Scriptures. The Egyptians would not let God take away Israel His firstborn, And Moses said: Thus saith the Lord: About midnight will I go out into the midst of Egypt, and all the first-born in the land of Egypt shall die, from the first-born of Pharaoh that sitteth upon his throne unto the first-born of the maid-servant that is behind the mill, and all the first-born of beasts. And there shall be a great cry throughout all the land of Egypt, such as there was none like it, nor shall be like it any more. But against any of the children of Israel shall not a dog move his tongue, against man or beast that ye may know how that the Lord hath put a difference between the Egyptians and Israel. (Ex. xi. 4-7.) Then Moses called for all the elders of Israel, and said unto them: Draw out and take you a lamb according to your families and kill the Passover. And ye shall take a bunch of hyssop and dip it in the blood that is in the basin, and strike the lintel and the two side-posts with the blood. and when He seeth the blood upon the lintel and on the two side-posts, the Lord will pass over the door, and will not suffer the destroyer to come in unto your houses to smite you. (xii. 21-23.) And could the blood of a sheep save a man Yea, in good sooth not because it was blood, but because it represented in a figure the Blood of the Lord.\par
\switchcolumn*

\LA
\begin{Resp}\Rsign In timóre incolátus vestri témpore conversámini \Star{} Sciéntes quod non corruptibílibus auro vel argénto redémpti estis.\end{Resp}
\begin{Resp}\Vsign Sed pretióso sánguine quasi agni immaculáti Christi.\end{Resp}
\begin{Resp}\Rsign Sciéntes quod non corruptibílibus auro vel argénto redémpti estis.\end{Resp}
\switchcolumn
\EN
\begin{Resp}\Rsign Pass the time of your sojourning here in fear, forasmuch as: \Star{} Ye know that ye were not redeemed with corruptible things, as silver and gold\end{Resp}
\begin{Resp}\Vsign But with the Precious Blood of Christ, as of a lamb without spot.\end{Resp}
\begin{Resp}\Rsign Ye know that ye were not redeemed with corruptible things, as silver and gold.\end{Resp}
\switchcolumn*

\LA
\LessonHead{Lectio V}
\switchcolumn
\EN
\LessonHead{Lesson V}
\switchcolumn*

\LA
\noindent Nam, sicut regnántium státuæ, quæ sine causa sunt et sermóne, nonnúmquam ad se confugiéntibus homínibus, ánima et ratióne decorátis, subveníre consuevérunt, non quia sunt ære conféctæ, sed quia rétinent imáginem principálem; ita et sanguis ille, qui irrationális fuit, ánimas habéntes hómines liberávit, non quia sanguis fuit, sed quia hujus sánguinis ostendébat advéntum. Et tunc Angelus ille vastátor, cum linítos postes atque áditus pervidéret, transjécit gressus et non est ausus intráre. Nunc ergo si víderit inimícus, non póstibus impósitum sánguinem typi, sed fidélium ore lucéntem sánguinem veritátis Christi, templi póstibus dedicátum, multo magis se súbtrahet. Si enim Angelus cessit exémplo, quanto magis terrébitur inimícus, si ipsam perspéxerit veritátem? Vis et áliam hujus sánguinis scrutári virtútem? Volo. Unde primum occúrrit, inspícias, et de quo fonte manávit. De ipsa primum cruce procéssit; latus illud Domínicum inítium fuit. Latus miles apéruit, et templi sancti paríetem patefécit; et ego thesáurum præclárum invéni, et fulgéntes divítias me grátulor reperíre.\par
\switchcolumn
\EN
\noindent The statues of monarchs, mindless and speechless images though they be, have sometimes been an helpful refuge to men endowed with soul and reason, not because they are works of the brazier's skill, but because the likeness they bear is a King's. And just so did this unconscious blood deliver the lives of men, not because it was blood, but because it foreshadowed the shedding of the Blood of Jesus. On that night in Egypt, when the destroying Angel saw the blood upon the lintel and on the two side-posts, he passed over the door, and came not in unto the house. Even so now much more will the destroyer of souls flee away when he seeth, not the lintel and the two side-posts sprinkled with the blood of a lamb, but the mouth of the faithful Christian, the living dwelling of the Holy Ghost, shining with the blood of the True Messiah. If the Angel let the type be, how shall not the enemy quail before the Reality? Wouldest thou hear more of the power of that Blood I am willing. Consider from what source it welleth, from what fountain it springeth. Its fountain is the Heart of the Lord, pierced for us upon the Cross. Then came the soldiers, and brake the legs of the first, and of the other which was crucified with Him but when they came to Jesus, and saw that He was dead already, they brake not His Legs, but one of the soldiers with a spear pierced His Side, and forthwith came thereout Blood and Water, (John xix, 32-34,) whereof One is a figure of Baptism, and the other of the Sacrament (of the Altar.) One of the soldiers with a spear pierced His Side and the veil of the Temple of His Body was rent in twain. (John ii. 19-21, Matth. xxvii. 51.) O how glorious is the treasure that is laid open to me therein! How noble the riches that it is my joy there to have found!\par
\switchcolumn*

\LA
\begin{Resp}\Rsign Empti estis prétio magno: \Star{} Glorificáte et portáte Deum in córpore vestro.\end{Resp}
\begin{Resp}\Vsign Prétio empti estis: nolíte fíeri servi hóminum.\end{Resp}
\begin{Resp}\Rsign Glorificáte et portáte Deum in córpore vestro.\end{Resp}
\switchcolumn
\EN
\begin{Resp}\Rsign Ye are bought with a great price, therefore: \Star{} Glorify God and bear Him in your body.\end{Resp}
\begin{Resp}\Vsign Ye are bought with a price; be not ye the servants of men.\end{Resp}
\begin{Resp}\Rsign Glorify God and bear Him in your body.\end{Resp}
\switchcolumn*

\LA
\LessonHead{Lectio VI}
\switchcolumn
\EN
\LessonHead{Lesson VI}
\switchcolumn*

\LA
\noindent Sic et de illo agno factum est: Judǽi ovem occidérunt, et ego fructum de sacraménto cognóvi. De látere sanguis et aqua. Nolo tam fácile, audítor, tránseas tanti secréta mystérii; restat enim mihi mýstica atque secretális orátio. Dixi baptísmatis sýmbolum et mysteriórum, aquam illam et sánguinem demonstráre. Ex his enim sancta fundáta est Ecclésia per lavácri regeneratiónem, et renovatiónem Spíritus Sancti. Per baptísma, inquam, et mystéria, quæ ex látere vidéntur esse proláta. Ex látere ígitur suo Christus ædificávit Ecclésiam, sicut de látere Adam ejus conjux Heva proláta est. Nam hac de causa Paulus quoque testátur dicens: De córpore ejus et de óssibus ejus sumus; latus vidélicet illud signíficans. Nam, sicut de illo látere Deus fecit féminam procreári, sic et de suo látere Christus aquam nobis et sánguinem dedit, unde repararétur Ecclésia.\par
Recurrénte autem anno ab humáni géneris redemptióne undevícies centenário, quem ad tam ineffábile benefícium recoléndum solémni præ ómnibus sacro Jubilǽo summus Póntifex Pius undécimus celebrátum vóluit; ut pretiósi sánguinis quo redémpti sumus, Agni Immaculáti Christi uberióres dimanárent in hómines fructus, ejúsque memória fidélibus, vivídius commendarétur, idem summus Póntifex Pretiosíssimi Sánguinis Dómini nostri Jesu Christi festum, quotánnis ab univérsa Ecclésia peragéndum, ad ritum dúplicem primæ classis evéxit.\par
\switchcolumn
\EN
\noindent And so was it done concerning that Lamb: the Jews killed a sheep, and I have learned the value of the sacrament. From the Side flowed forth Blood and Water. I would not, O my hearer, that thou shouldest pass by the depths of such a mystery as this without pausing; for I have yet a mystic and mysterious discourse to deliver. I have said that the Water and Blood shewed forth symbolically baptism and the sacraments. For from these, holy Church was founded by the laver of regeneration, and the renovation of the Holy Ghost. Through baptism, I say, and through the sacraments, which seem to have issued from his Side. It was therefore out of the Side of Christ that the Church was created, just as it was out of the side of Adam that Eve was raised up to be his bride. This is the reason why Paul saith, no doubt in allusion to his Side: We are members of his Body, and of his bones. For even as God made the woman Eve out of the rib which he had taken out of the side of Adam, so hath Christ made the Church out of the Blood and Water which he made to flow for us out of his own Side. -On the occasion of the nineteenth centenary of the accomplishment of the redemption of mankind, as a fitting celebration of this ineffable blessing, Pope Pius XI decreed an extraordinary Jubilee. During that year the Supreme Pontiff, wishing that the fruits of the Precious Blood of Christ, the Lamb without spot, might redound more abundantly upon mankind and that the minds of the faithful be impressed with more vivid recollections of this same Blood as the price of their redemption, elevated the Feast of the Most Precious Blood of our Lord Jesus Christ to the rank of a double of the first class, to be celebrated as such every year by the universal Church.\par
\switchcolumn*

\LA
\begin{Resp}\Rsign Comméndat caritátem suam Deus in nobis: \Star{} Quóniam cum adhuc peccatóres essémus, secúndum tempus Christus pro nobis mórtuus est.\end{Resp}
\begin{Resp}\Vsign Multo ígitur magis nunc justificáti in sánguine ipsíus, salvi érimus ab ira per ipsum.\end{Resp}
\begin{Resp}\Rsign Quóniam cum adhuc peccatóres essémus, secúndum tempus Christus pro nobis mórtuus est.\end{Resp}
\begin{Resp}\Vsign Glória Patri, et Fílio, \Star{} et Spirítui Sancto.\end{Resp}
\begin{Resp}\Rsign Quóniam cum adhuc peccatóres essémus, secúndum tempus Christus pro nobis mórtuus est.\end{Resp}
\switchcolumn
\EN
\begin{Resp}\Rsign God commendeth His love toward us: \Star{} In that, while we were yet sinners, in due time Christ died for us.\end{Resp}
\begin{Resp}\Vsign Much more then being now justified by His Blood, we shall be saved from wrath through Him.\end{Resp}
\begin{Resp}\Rsign In that, while we were yet sinners, in due time Christ died for us.\end{Resp}
\begin{Resp}\Vsign Glory be to the Father, and to the Son, \Star{} and to the Holy Ghost.\end{Resp}
\begin{Resp}\Rsign In that, while we were yet sinners, in due time Christ died for us.\end{Resp}
\switchcolumn*

\LA
\LessonHead{Lectio VII}
\switchcolumn
\EN
\LessonHead{Lesson VII}
\switchcolumn*

\LA
\LessonTitle{Léctio sancti Evangélii secúndum Joánnem}
\switchcolumn
\EN
\LessonTitle{From the Holy Gospel according to John}
\switchcolumn*

\LA
\Cite{Joannes 19:30-35}
\switchcolumn
\EN
\Cite{John 19:30-35}
\switchcolumn*

\LA
\noindent In illo témpore: Cum accepísset Jesus acétum, dixit: Consummátum est. Et, inclináto cápite, trádidit spíritum. Et réliqua.\par
\switchcolumn
\EN
\noindent At that time: When Jesus had received the vinegar, He said: It is finished. And He bowed His Head, and gave up the ghost. And so on.\par
\switchcolumn*

\LA
\LessonTitle{Homilía sancti Augustíni Epíscopi}
\switchcolumn
\EN
\LessonTitle{Homily by St. Augustine, Bishop of Hippo}
\switchcolumn*

\LA
\Source{Tractatus 120 in Joannem}
\switchcolumn
\EN
\Source{120th Tract on John}
\switchcolumn*

\LA
\noindent Vigilánti verbo Evangelísta usus est, ut non díceret: Latus ejus percússit, aut vulnerávit, aut quid áliud, sed Apéruit; ut illic quodámmodo vitæ óstium panderétur, unde sacraménta Ecclésiæ manavérunt, sine quibus ad vitam, quæ vera vita est, non intrátur. Ille sanguis qui fusus est, in remissiónem fusus est peccatórum. Aqua illa salutáre témperat póculum; hæc et lavácrum præstat et potum. Hoc prænuntiábat quod Noë in látere arcæ óstium fácere jussus est, quo intrárent animália quæ non erant dilúvio peritúra, quibus præfigurabátur Ecclésia. Propter hoc prima múlier facta est de látere viri dormiéntis, et appelláta est vita matérque vivórum. Magnum quippe significávit bonum, ante magnum prævaricatiónis malum. Hic secúndus Adam, inclináto cápite, in cruce dormívit, ut inde formarétur ei conjux, quæ de látere dormiéntis efflúxit. O mors, unde mórtui revivíscunt! Quid isto sánguine múndius? Quid vúlnere isto salúbrius?\par
\switchcolumn
\EN
\noindent One of the soldiers with a spear pierced His Side, and forthwith came thereout Blood and Water. The Evangelist speaketh carefully. He saith not that he smote the Side, nor yet that he wounded It, nor yet anything else, but pierced It, to fling wide the entrance unto life, whence flow the Sacraments of the Church, those Sacraments without which there is no entrance unto the life which is life indeed. That Blood which was shed there was shed for the remission of sins, that Water is the water that mantleth in the cup of salvation. Therein are we washed, and thereof do we drink. Of this was it a type when it was said unto Noah: The door of the ark shalt thou set in the side thereof and of every living thing of all flesh shalt thou bring into the ark to keep them alive. (Gen. vi. 16, 19.) A figure this of the Church. Thus was it that the first woman was made from the side of her husband while he slept, and she was called (Eve, which is, being interpreted,) Life, because she was the mother of all living. (Gen. iii. 20.) This name set forth a great good, before it became associated with the bitter fruit of a great evil. And here we have the Second Adam bowing His Head, and the deep sleep of death falling upon Him upon the Cross, and He sleepeth, that the Lord God may take a thing out of His side, and may make thereof a wife for Him. O what a death was His, which quickeneth the dead! What is cleaner than His Blood? What more health-giving than His wounding?\par
\switchcolumn*

\LA
\begin{Resp}\Rsign Hic est, qui venit per aquam et sánguinem, Jesus Christus: \Star{} Non in aqua solum, sed in aqua et sánguine.\end{Resp}
\begin{Resp}\Vsign In die illa erit fons patens dómui David et habitántibus Jerúsalem in ablutiónem peccatóris.\end{Resp}
\begin{Resp}\Rsign Non in aqua solum, sed in aqua et sánguine.\end{Resp}
\switchcolumn
\EN
\begin{Resp}\Rsign This is He That came by Water and Blood, even Jesus Christ; \Star{} Not by Water only, but by Water and Blood.\end{Resp}
\begin{Resp}\Vsign In that day there shall be a fountain opened to the house of David and to the inhabitants of Jerusalem, for sin and for uncleanness.\end{Resp}
\begin{Resp}\Rsign Not by Water only, but by Water and Blood.\end{Resp}
\switchcolumn*

\LA
\LessonHead{Lectio VIII}
\switchcolumn
\EN
\LessonHead{Lesson VIII}
\switchcolumn*

\LA
\Source{Enarrat. in Psalm. 95, n. 5}
\switchcolumn
\EN
\Source{Enarr. in Psalm 95, n. 5}
\switchcolumn*

\LA
\noindent Tenebántur hómines captívi sub diábolo, et dæmónibus serviébant; sed redémpti sunt a captivitáte. Véndere enim se potuérunt, sed redímere non potuérunt. Venit Redémptor, et dedit prétium; fudit sánguinem suum, et emit orbem terrárum. Quǽritis quid émerit? Vidéte quid déderit, et inveniétis quid émerit. Sanguis Christi prétium est. Tanti quid valet? quid, nisi totus orbis? quid, nisi omnes gentes? Valde ingráti sunt prétio suo, aut multum supérbi sunt, qui dicunt, aut illud tam parum esse ut solos Afros émerit, aut se tam magnos esse pro quibus solis illud sit datum. Non ergo exsúltent, non supérbiant. Pro toto dedit, quantum dedit.\par
\switchcolumn
\EN
\noindent Then were being held bondsmen to the devil, slaves to evil spirits. But they have been redeemed from that bondage. They had been able to sell themselves, but they were not able to redeem themselves. A Redeemer came and paid the price for them. He shed His Blood, and at that cost bought the world. Ye ask what He bought? Look what He paid, and ye shall see what He bought. Christ's Blood was the price. What is His Blood worth? What, but the whole world? What, but all men? They are very unthankful for His redemption, or very proud, who say that It is only precious enough to buy the Africans, or that they themselves are so precious that It was shed only for them. Let there be an end to such conceit, an end to such vainglory. What He paid, He paid for all.\par
\switchcolumn*

\LA
\begin{Resp}\Rsign Prædestinávit nos Deus in adoptiónem filiórum per Jesum Christum, \Star{} In quo habémus redemptiónem per sánguinem ejus.\end{Resp}
\begin{Resp}\Vsign Remissiónem peccatórum secúndum divítias grátiæ ejus, quæ superabundávit in nobis.\end{Resp}
\begin{Resp}\Rsign In quo habémus redemptiónem per sánguinem ejus.\end{Resp}
\begin{Resp}\Vsign Glória Patri, et Fílio, \Star{} et Spirítui Sancto.\end{Resp}
\begin{Resp}\Rsign In quo habémus redemptiónem per sánguinem ejus.\end{Resp}
\switchcolumn
\EN
\begin{Resp}\Rsign God hath predestinated us unto the adoption of children by Jesus Christ; \Star{} In Whom we have redemption through His Blood.\end{Resp}
\begin{Resp}\Vsign The forgiveness of sins, according to the riches of His grace, wherein He hath abounded toward us.\end{Resp}
\begin{Resp}\Rsign In Whom we have redemption through His Blood.\end{Resp}
\begin{Resp}\Vsign Glory be to the Father, and to the Son, \Star{} and to the Holy Ghost.\end{Resp}
\begin{Resp}\Rsign In Whom we have redemption through His Blood.\end{Resp}
\switchcolumn*

\LA
\LessonHead{Lectio IX}
\switchcolumn
\EN
\LessonHead{Lesson IX}
\switchcolumn*

\LA
\Source{Sermo 31, alias 344}
\switchcolumn
\EN
\Source{Sermon 31, alias 344}
\switchcolumn*

\LA
\noindent Hábuit ille sánguinem, unde nos redímeret; et ad hoc accépit sánguinem, ut esset quem pro nobis rediméndis effúnderet. Sanguis Dómini tui, si vis, datus est pro te; si nolúeris esse, non est datus pro te. Forte enim dicis: Hábuit sánguinem Deus meus, quo me redímeret, sed jam, cum passus est, totum dedit; quid illi remánsit, quod det et pro me? Hoc est magnum, quia semel dedit, et pro ómnibus dedit. Sanguis Christi volénti est salus, nolénti supplícium. Quid ergo dúbitas, qui mori non vis, a secúnda pótius morte liberári? Qua liberáris, si vis tóllere crucem tuam, et sequi Dóminum; quia ille tulit suam, et quæsívit servum.\par
\switchcolumn
\EN
\noindent That Blood was his own, and thereby he redeemed us. Yea, it was to this end that he took Flesh and Blood, namely that he might shed his Blood in order to redeem us. If thou wilt accept it, the Blood of thy Lord was given for thee. If thou wilt not accept it, it was not given for thee. For perchance thou sayest: My God had Blood, with which he redeemed me, but now since he hath suffered, he hath given it all; what hath remained to him, that he may also give any of it for me? This is a great thing, because he gave once, and he gave for all. The Blood of Christ is salvation to him that doth accept it, punishment to him that doth not accept it. Why therefore dost thou hesitate to be set free from the second death, thou who dost not wish to die? By this thou art set free, if thou art willing to take up thy Cross, and follow the Lord; for he took up his Cross and sought his servant.\par
\switchcolumn*

\LA
\TeDeum{Te Deum}
\switchcolumn
\EN
\TeDeum{Te Deum}
\switchcolumn*

\end{paracol}

\par\needspace{10\baselineskip}\addvspace{1.2em}{\centering\small\textsc{Die 1 Julii} · \textsc{1 July}\par}
\Office{In Octava S. Joannis Baptistæ}{Octave of the Nativity of St. John the Baptist}{Duplex}
\addcontentsline{toc}{subsection}{Die 1 Julii · In Octava S. Joannis Baptistæ \textit{— Octave of the Nativity of St. John the Baptist}}
\begin{paracol}{2}
\LA
\LessonHead{Lectio IV}
\switchcolumn
\EN
\LessonHead{Lesson IV}
\switchcolumn*

\LA
\LessonTitle{Commemoratio: In Octava S. Joannis Baptistæ}
\switchcolumn
\EN
\LessonTitle{Commemoration of: Octave of the Nativity of St. John the Baptist}
\switchcolumn*

\LA
\Source{Homilia 3. in Nativ. S. Joannis Bapt.}
\switchcolumn
\EN
\Source{3rd for the Birth-day of St. John the Baptist}
\switchcolumn*

\LA
\noindent Sermo sancti Maximi Epíscopi.\par
\switchcolumn
\EN
\noindent From the Sermons of St. Maximus, Bishop of Turin\par
\switchcolumn*

\LA
Festivitatem præsentis diei, fratres carissimi, venerandi Joannis Baptistæ genuina nativitas consecravit: qui idcirco in hoc sæculum superna dispensatione directus est, ut non solum ipse prophetali sublimaretur gloria, sed ut et omnium étiam per eum prophetarum præconia firmarentur. Nec immerito illum præcipuo honore nunc veneramur, qui speciali quadam gratia, ob hoc Redemptorem mundi novissimus prophetavit, ut ostenderet eum primus. Hic enim solus est prophetarum, qui Dominum nostrum Jesum Christum, quem alii in longa tempora futurum præscierunt, propriis oculis videre meruit, et annuntiare præsentem.\par
\switchcolumn
\EN
Dearly beloved brethren, the feast which we this day keep hath been hallowed by the true Birth of the worshipful Baptist, John. He was a man whom God was pleased to send into this world, not only that he himself might fill the glorious office of a Prophet, but also that by him the words of all the other Prophets might be established. That we should give him, as we now do, the chiefest honour among them all, is right for he had poured upon him this special blessing, that, being the last who prophesied in the world the coming of Him Who came to redeem the world, he was the first who, when He came, pointed to Him as already come. The other Prophets, at the distance of ages, knew that our Lord Jesus Christ was to come the Prophet John alone was worthy to see Him with his very eyes, and to proclaim Him actually present.\par
\switchcolumn*

\LA
Hic est ille, quem, inspirante Deo, præscius annuntiat Isaias, dicens: Vox clamantis in deserto, Parate viam Domini. Quam congrue fratres carissimi beatus Joannes prædicatus est vox, qui cælestis Verbi et præco mittebatur, et testis. Hic est ille, cujus per Angelum Gabrielem prænuntiatur nativitas, et nomen, et meritum. Hic est, qui judicio cælestis sententiæ cunctis mortalibus antefertur, dicente Domino: Non surrexit inter natos mulierum major Joanne Baptista. Quam pulchre dictum est, inter natos mulierum non esse majorem: quia ille omnimodis major erat Joanne, qui de Virgine nascebatur.\par
\switchcolumn
\EN
This is he concerning whom the seer Isaiah, by the inspiration of God, said: The voice of him that crieth in the wilderness Prepare ye the way of the Lord. \textasciitilde{}(xl. 3.) How meetly, dearly beloved brethren, is this blessed John called a voice he who was sent to be at once the herald and the witness of the Word from heaven This is he whose birth, whose name, and whose work were foretold by the angel Gabriel. This is he who in the sight of heavenly knowledge is more honourable than any other man, as hath been declared by the Lord Among them that are born of women there hath not risen a greater than John the Baptist. \textasciitilde{}(Matth. xi. 11.) How beautifully is it said, born of women since He Who was infinitely greater than John was born of a Virgin\par
\switchcolumn*

\LA
Quibus perspectis, advertite dilectissimi, quantam huic reverentiam, quantum devotionis debeamus impendere: qui ut honorabilis sit, a sancto est Spiritu prophetatus, promissus ab Angelo, laudatus a Domino, et perpetua sanctæ mortis gloria consecratus. Decebat namque, ut mysticam nativitatem ejus mirabilis vita sequeretur, et sanctam perfectamque vitam mors Deo devota concluderet. Propter quod, fratres, rectissime toto hodie orbe, Christi Ecclesia primordia nativitatis ejus lætissima festivitate concelebrat, qui æterna mortalibus adesse gaudia, stupenti mundo testis fidelissimus revelavit.\par
\switchcolumn
\EN
In sight of these things, consider, dearly beloved brethren, in how great worship, in how great love ye are behoven to hold him, who is thus honourable, that he hath been foretold by the Holy Ghost, promised by an Angel, extolled by the Lord Himself, and hallowed by the everlasting glory of an holy death. His birth, surrounded with wonders, is the meet beginning of a wondrous life, that holy, sinless life the meet fore-runner of a death which was a gift to God. It is but just then, my brethren, that the Church of Christ should on this day with keen joy keep throughout the whole world the Feast of his Birth, even his, who arose a right faithful witness to tell a wondering world that a life of eternal gladness was lying open to the choice of dying creatures.\par
\switchcolumn*

\LA
\TeDeum{Te Deum}
\switchcolumn
\EN
\TeDeum{Te Deum}
\switchcolumn*

\end{paracol}

\par\needspace{10\baselineskip}\addvspace{1.2em}{\centering\small\textsc{Die 2 Julii} · \textsc{2 July}\par}
\Office{In Visitatione Beatæ Mariæ Virginis}{In Visitatione Beatæ Mariæ Virginis}{Duplex II classis}
\addcontentsline{toc}{subsection}{Die 2 Julii · In Visitatione Beatæ Mariæ Virginis \textit{— In Visitatione Beatæ Mariæ Virginis}}
\begin{paracol}{2}
\LA
\RefLine{Lectio IX: de Ss. Processi et Martiniani Martyrum (07-02cc).}
\switchcolumn
\EN
\RefLine{Lesson IX: of Ss. Processus et Martinian, Martyrs (07-02cc).}
\switchcolumn*

\LA
\LessonHead{Lectio I}
\switchcolumn
\EN
\LessonHead{Lesson I}
\switchcolumn*

\LA
\LessonTitle{De Cánticis canticórum}
\switchcolumn
\EN
\LessonTitle{From the Song of Songs}
\switchcolumn*

\LA
\Cite{Cant 2:1-7}
\switchcolumn
\EN
\Cite{Song 2:1-7}
\switchcolumn*

\LA
\noindent Ego flos campi et lílium convállium. Sicut lílium inter spinas, sic amíca mea inter fílias. Sicut malus inter ligna silvárum, sic diléctus meus inter fílios. Sub umbra illíus quem desideráveram sedi, et fructus ejus dulcis gútturi meo. Introdúxit me in cellam vináriam, ordinávit in me caritátem. Fulcíte me flóribus, stipáte me malis, quia amóre lángueo. Læva ejus sub cápite meo, et déxtera illíus amplexábitur me. Adjúro vos, fíliæ Jerúsalem, per cápreas cervósque campórum, ne suscitétis neque evigiláre faciátis diléctam, quoadúsque ipsa velit.\par
\switchcolumn
\EN
\noindent I am the flower of the field, and the lily of the valleys. As the lily among thorns, so is my love among the daughters. As the apple tree among the trees of the woods, so is my beloved among the sons. I sat down under his shadow, whom I desired: and his fruit was sweet to my palate. He brought me into the cellar of wine, he set in order charity in me. Stay me up with flowers, compass me about with apples: because I languish with love. His left hand is under my head, and his right hand shall embrace me. I adjure you, O ye daughters of Jerusalem, by the roes, and the harts of the, fields, that you stir not up, nor make the beloved to awake, till she please\par
\switchcolumn*

\LA
\begin{Resp}\Rsign Surge, própera, amíca mea, formósa mea, et veni: jam enim hiems tránsiit, imber ábiit et recéssit: \Star{} Vox túrturis audíta est in terra nostra.\end{Resp}
\begin{Resp}\Vsign Intrávit María in domum Zacharíæ et salutávit Elísabeth.\end{Resp}
\begin{Resp}\Rsign Vox túrturis audíta est in terra nostra.\end{Resp}
\switchcolumn
\EN
\begin{Resp}\Rsign Rise up, make haste, my love my fair one and come away; For, lo, the winter is past, the rain is over and gone. \Star{} The voice of the turtle-dove is heard in our land.\end{Resp}
\begin{Resp}\Vsign Mary entered into the house of Zacharias and saluted Elizabeth.\end{Resp}
\begin{Resp}\Rsign The voice of the turtle-dove is heard in our land.\end{Resp}
\switchcolumn*

\LA
\LessonHead{Lectio II}
\switchcolumn
\EN
\LessonHead{Lesson II}
\switchcolumn*

\LA
\Cite{Cant 2:8-13}
\switchcolumn
\EN
\Cite{Song 2:8-13}
\switchcolumn*

\LA
\noindent Vox dilécti mei: ecce iste venit sáliens in móntibus, transíliens colles. Símilis est diléctus meus cápreæ hinnulóque cervórum. En ipse stat post paríetem nostrum respíciens per fenéstras, prospíciens per cancéllos. En diléctus meus lóquitur mihi: Surge, própera, amíca mea, colúmba mea, formósa mea, et veni. Jam enim hiems tránsiit, imber ábiit et recéssit, flores apparuérunt in terra nostra, tempus putatiónis advénit, vox túrturis audíta est in terra nostra, ficus prótulit grossos suos, víneæ floréntes dedérunt odórem suum.\par
\switchcolumn
\EN
\noindent The voice of my beloved, behold he cometh leaping upon the mountains, skipping over the hills. My beloved is like a roe, or a young hart. Behold he standeth behind our wall, looking through the windows, looking through the lattices. Behold my beloved speaketh to me: Arise, make haste, my love, my dove, my beautiful one, and come. For winter is now past, the rain is over and gone. The flowers have appeared in our land, the time of pruning is come: the voice of the turtle is heard in our land: The fig tree hath put forth her green figs: the vines in flower yield their sweet smell.\par
\switchcolumn*

\LA
\begin{Resp}\Rsign Quæ est ista quæ procéssit sicut sol, et formósa tamquam Jerúsalem? \Star{} Vidérunt eam fíliæ Sion, et beátam dixérunt, et regínæ laudavérunt eam.\end{Resp}
\begin{Resp}\Vsign Et sicut dies verni circúmdabant eam flores rosárum et lília convállium.\end{Resp}
\begin{Resp}\Rsign Vidérunt eam fíliæ Sion, et beátam dixérunt, et regínæ laudavérunt eam.\end{Resp}
\switchcolumn
\EN
\begin{Resp}\Rsign Who is this that cometh up like the sun? This, comely as Jerusalem? \Star{} The daughters of Zion saw her, and called her blessed: the queens also, and they praised her.\end{Resp}
\begin{Resp}\Vsign And about her it was as the flower of roses in the spring of the year, and lilies of the valleys.\end{Resp}
\begin{Resp}\Rsign The daughters of Zion saw her, and called her blessed the queens also, and they praised her.\end{Resp}
\switchcolumn*

\LA
\LessonHead{Lectio III}
\switchcolumn
\EN
\LessonHead{Lesson III}
\switchcolumn*

\LA
\Cite{Cant 2:13-17}
\switchcolumn
\EN
\Cite{Song 2:13-17}
\switchcolumn*

\LA
\noindent Surge, amíca mea, speciósa mea, et veni, colúmba mea, in foramínibus petræ, in cavérna macériæ, osténde mihi fáciem tuam, sonet vox tua in áuribus meis: vox enim tua dulcis, et fácies tua decóra. Cápite nobis vulpes párvulas, quæ demoliúntur víneas; nam vínea nostra flóruit. Diléctus meus mihi, et ego illi, páscitur inter lília, donec aspíret dies, et inclinéntur umbræ. Revértere; símilis esto, dilécte mi, cápreæ hinnulóque cervórum super montes Bether.\par
\switchcolumn
\EN
\noindent Arise, my love, my beautiful one, and come: My dove in the clefts of the rock, in the hollow places of the wall, show me thy face, let thy voice sound in my ears: for thy voice is sweet, and thy face comely. Catch us the little foxes that destroy the vines: for our vineyard hath flourished. My beloved to me, and I to him who feedeth among the lilies, Till the day break, and the shadows retire. Return: be like, my beloved, to a roe, or to a young hart upon the mountains of Bether.\par
\switchcolumn*

\LA
\begin{Resp}\Rsign Repléta est Spíritu Sancto Elísabeth et exclamávit: Benedícta tu inter mulíeres, et benedíctus fructus ventris tui: \Star{} Et unde hoc mihi, ut véniat mater Dómini mei ad me?\end{Resp}
\begin{Resp}\Vsign Ecce enim, ut facta est vox salutatiónis tuæ in áuribus meis, exsultávit in gáudio infans in útero meo.\end{Resp}
\begin{Resp}\Rsign Et unde hoc mihi, ut véniat mater Dómini mei ad me?\end{Resp}
\begin{Resp}\Vsign Glória Patri, et Fílio, \Star{} et Spirítui Sancto.\end{Resp}
\begin{Resp}\Rsign Et unde hoc mihi, ut véniat mater Dómini mei ad me?\end{Resp}
\switchcolumn
\EN
\begin{Resp}\Rsign Elizabeth was filled with the Holy Ghost and she spake out with a loud voice, and said: Blessed art thou among women, and blessed is the Fruit of thy womb; \Star{} And whence is this to me, that the mother of my Lord should come to me?\end{Resp}
\begin{Resp}\Vsign For, lo, as soon as the voice of thy salutation sounded in mine ears, the babe leaped in my womb for joy.\end{Resp}
\begin{Resp}\Rsign And whence is this to me, that the mother of my Lord should come to me?\end{Resp}
\begin{Resp}\Vsign Glory be to the Father, and to the Son, \Star{} and to the Holy Ghost.\end{Resp}
\begin{Resp}\Rsign And whence is this to me, that the mother of my Lord should come to me?\end{Resp}
\switchcolumn*

\LA
\LessonHead{Lectio IV}
\switchcolumn
\EN
\LessonHead{Lesson IV}
\switchcolumn*

\LA
\LessonTitle{Sermo sancti Joánnis Chrysóstomi}
\switchcolumn
\EN
\LessonTitle{From the Sermons of St. John Chrysostom, Patriarch of Constantinople}
\switchcolumn*

\LA
\Source{Apud Metaphr. mense Julio}
\switchcolumn
\EN
\Source{In Metaphrast. July}
\switchcolumn*

\LA
\noindent Cum ad nos advenísset Redémptor nostri géneris, venit prótinus ad suum amícum Joánnem, dum adhuc esset in ventre matris. Quem cum ex útero in útero aspexísset Joánnes, términos natúræ concútiens, exclámat: Vídeo Dóminum, qui natúræ impósuit términos, et non exspécto tempus nascéndi. Novem ménsium tempus mihi non est hic necessárium; in me est enim, qui est ætérnus. Egrédiar ex hoc tenebróso tabernáculo, rerum admirabílium compendiósam prædicábo cognitiónem. Sum signum: significábo Christi advéntum. Sum tuba: próferam Fílii Dei in carne dispensatiónem. Tuba canam; eo ipso patérnæ linguæ benedícam, et eam traham, ut loquátur. Tuba canam, et úterum matérnum vivificábo.\par
\switchcolumn
\EN
\noindent As soon as our Redeemer was come among us, He went with haste, while as yet He was in His mother's womb, to visit His friend John. And John, in the one womb, becoming conscious of the Presence of Jesus in the other womb, dashed himself impatiently against the narrow walls of his natural prison, as though crying out: I see the very Lord who hath given nature her bounds, and I wait not for the due season of my birth. There is no need for me to linger here till nine months are ended, for He That is Eternal is with me! I will break out of my dark cell; I will proclaim my full knowledge of many wonders. I am the sign. I will show that the Christ is here. I am the trumpet; let me peal forth the news that the Son of God is come in the flesh. Let me give my trumpet-note, let me bless my father's tongue, and make it to speak again. Let me give my trumpet-note, let me quicken my mother's womb.\par
\switchcolumn*

\LA
\begin{Resp}\Rsign Ecce iste venit sáliens in móntibus, transíliens colles: \Star{} Símilis est diléctus meus cápreæ hinnulóque cervórum.\end{Resp}
\begin{Resp}\Vsign Exsultávit ut gigas ad curréndam viam, a summo cælo egréssio ejus.\end{Resp}
\begin{Resp}\Rsign Símilis est diléctus meus cápreæ hinnulóque cervórum.\end{Resp}
\switchcolumn
\EN
\begin{Resp}\Rsign Behold, he cometh leaping upon the mountains, skipping upon the hills; \Star{} My beloved is like a gazelle or a young roe-buck!\end{Resp}
\begin{Resp}\Vsign He rejoiceth as a giant to run his course, his going-forth is from the end of the heavens.\end{Resp}
\begin{Resp}\Rsign My beloved is like a gazelle or a young roe-buck!\end{Resp}
\switchcolumn*

\LA
\LessonHead{Lectio V}
\switchcolumn
\EN
\LessonHead{Lesson V}
\switchcolumn*

\LA
\noindent Vides, o dilécte, quam sit novum et admirábile mystérium. Nondum náscitur, et sáltibus lóquitur; nondum appáret, et minas inténtat; nondum ei permíttitur clamáre, et per facta audítur; nondum ducit vitam, et Deum prǽdicat; nondum áspicit lucem, et solem índicat; nondum páritur, et próperat præcúrrere. Non fert enim, præsénte Dómino, continéri; non sústinet natúræ exspectáre términos; sed conténdit rúmpere cárcerem ventris, et studet præsignificáre veniéntem Salvatórem. Accéssit, inquit, qui solvit víncula; et quid ego sédeo vinctus, et retíneor ut máneam? Venit Verbum, ut ómnia constítuat; et ego adhuc máneo deténtus? Exíbo, præcúrram, et prædicábo ómnibus: Ecce Agnus Dei, qui tollit peccátum mundi.\par
\switchcolumn
\EN
\noindent Thou seest, O beloved, how new and how strange a mystery is here! John is not born, but by leaping he speaketh; he is yet unseen, and he giveth warning; he is not yet able to cry, but by his acts he is heard. He draweth not yet the breath of life, but he preacheth God. he seeth not yet the light, but he maketh known the Sun! He is not come out of the womb, but he hasteth to play the Fore-runner in the Presence of the Lord. He cannot restrain himself; he rebelleth against the bounds set by nature, and struggleth to break out of the prison of the belly. His longing is to herald the coming Saviour. He saith, as it were: Behold, the Deliverer cometh and am I to remain still bound to abide here? The Word cometh, that He may set right all things and am I still to tarry in prison? I will go forth. I will run before Him, and cry aloud to all men: “Behold the Lamb of God Which taketh away the sin of the world.”\par
\switchcolumn*

\LA
\begin{Resp}\Rsign Congratulámini mihi, omnes qui dilígitis Dóminum: quia cum essem párvula, plácui Altíssimo, \Star{} Et de meis viscéribus génui Deum et hóminem.\end{Resp}
\begin{Resp}\Vsign Beátam me dicent omnes generatiónes, quia ancíllam húmilem respéxit Deus.\end{Resp}
\begin{Resp}\Rsign Et de meis viscéribus génui Deum et hóminem.\end{Resp}
\switchcolumn
\EN
\begin{Resp}\Rsign Rejoice with me, rejoice with me, all ye that love the Lord, for, while I was yet a little one, I pleased the Most High, \Star{} And I have brought forth from my bowels God and man.\end{Resp}
\begin{Resp}\Vsign All generations shall call me blessed, for God hath regarded the low estate of His hand-maiden.\end{Resp}
\begin{Resp}\Rsign And I have brought forth from my bowels God and man.\end{Resp}
\switchcolumn*

\LA
\LessonHead{Lectio VI}
\switchcolumn
\EN
\LessonHead{Lesson VI}
\switchcolumn*

\LA
\noindent Sed dic nobis, Joánnes, cum adhuc in tenebróso matris útero contineáris, quómodo intuéris et audis? quómodo res divínas contempláris? quómodo éxsilis et exsúltas? Magnum est, inquit, quod perágitur mystérium, et actus ab humána remótus comprehensióne. Mérito ínnovo natúram propter eum, qui est innovatúrus ea quæ sunt supra natúram. Vídeo, etsi adhuc in útero sim; quóniam in útero gestári vídeo solem justítiæ. Auribus percípio, quóniam nascor vox magni Verbi. Exclámo, quóniam Fílium Patris unigénitum consídero carne indútum. Exsúlto, quóniam univérsi Effectórem vídeo formam, hóminis suscípere. Exsílio, quóniam mundi Redemptórem cógito incorporátum. Præcúrro advéntum ejus, et quodámmodo vobis prǽeo confessióne.\par
\switchcolumn
\EN
\noindent But do thou tell us, O John, how it came to pass that while thou wast still in the darkness of thy mother's womb, thou didst see and hear? How didst thou behold the things of God? How didst thou leap and bound for joy? Great, saith John, is the mystery of that which taketh place here, far from the understanding of men are these doings. It is meet that I should do a new thing in nature for the sake of Him Who is making new things which are beyond nature. I see in the womb, because I see the Sun of righteousness in a womb. I hear, because I am coming as the herald of the Great Word. I cry out, because I espy the Only-begotten Son of the Father clad in Flesh. I bound for joy, because I see that He by Whom all things were made, hath taken upon Him the form of a servant. I leap, because I think of the Redeemer of the world being made Flesh. I run before His coming, and herald His approach unto you with this, as it were, my confession.\par
\switchcolumn*

\LA
\begin{Resp}\Rsign Beáta quæ credidísti, quóniam perficiéntur in te quæ dicta sunt tibi a Dómino. Et ait María: \Star{} Magníficat ánima mea Dóminum.\end{Resp}
\begin{Resp}\Vsign Veníte, et audíte, et narrábo quanta fecit Deus ánimæ meæ.\end{Resp}
\begin{Resp}\Rsign Magníficat ánima mea Dóminum.\end{Resp}
\begin{Resp}\Vsign Glória Patri, et Fílio, \Star{} et Spirítui Sancto.\end{Resp}
\begin{Resp}\Rsign Magníficat ánima mea Dóminum.\end{Resp}
\switchcolumn
\EN
\begin{Resp}\Rsign Blessed art thou that hast believed, for there shall be a performance of those things which were told thee from the Lord. And Mary said: \Star{} My soul doth magnify the Lord.\end{Resp}
\begin{Resp}\Vsign Come, hear, and I will declare what God hath done for my soul.\end{Resp}
\begin{Resp}\Rsign My soul doth magnify the Lord.\end{Resp}
\begin{Resp}\Vsign Glory be to the Father, and to the Son, \Star{} and to the Holy Ghost.\end{Resp}
\begin{Resp}\Rsign My soul doth magnify the Lord.\end{Resp}
\switchcolumn*

\LA
\LessonHead{Lectio VII}
\switchcolumn
\EN
\LessonHead{Lesson VII}
\switchcolumn*

\LA
\LessonTitle{Léctio sancti Evangélii secúndum Lucam}
\switchcolumn
\EN
\LessonTitle{From the Holy Gospel according to Luke}
\switchcolumn*

\LA
\Cite{Luc 1:39-47}
\switchcolumn
\EN
\Cite{Luke 1:39-47}
\switchcolumn*

\LA
\noindent In illo témpore: Exsúrgens María ábiit in montána cum festinatióne in civitáte Juda: et intrávit in domum Zacharíæ, et salutávit Elísabeth. Et réliqua.\par
\switchcolumn
\EN
\noindent And Mary arose in those days and went into the hill-country with haste, into a city of Judah. And entered into the house of Zacharias, and saluted Elizabeth. And so on.\par
\switchcolumn*

\LA
\LessonTitle{Homilía sancti Ambrósii Epíscopi}
\switchcolumn
\EN
\LessonTitle{Homily by St. Ambrose, Bishop of Milan}
\switchcolumn*

\LA
\Source{Liber 2 Comment. in Lucæ. cap. 1, post initium}
\switchcolumn
\EN
\Source{Bk. ii Comm. on Luke i}
\switchcolumn*

\LA
\noindent Contuéndum est, quia supérior venit ad inferiórem, ut inférior adjuvétur: María ad Elísabeth, Christus ad Joánnem. Dénique étiam póstea, ut sanctificáret baptísmum Joánnis, Dóminus venit ad baptísmum. Cito quoque advéntus Maríæ et præséntiæ divínæ benefícia declarántur. Vide distinctiónem singulorúmque verbórum proprietátem. Vocem prior Elísabeth audívit, sed Joánnes prior grátiam sensit. Illa natúræ órdine audívit, iste exsultávit ratióne mystérii. Illa Maríæ, iste Dómini sensit advéntum. Istæ grátiam loquúntur, illi intus operántur, pietatísque mystérium matérnis adoriúntur proféctibus; duplicíque miráculo prophétant matres spíritu parvulórum. Exsultávit infans, repléta est mater. Non prius mater repléta, quam fílius; sed, cum fílius esset replétus Spíritu Sancto, replévit et matrem.\par
\switchcolumn
\EN
\noindent We must here consider that the greater cometh unto the lesser, Mary unto Elizabeth, Christ unto John. And again afterwards, to hallow the baptism of John, the Lord came unto him to be baptized. It was soon that the blessings of the coming of Mary and of the Presence of God were made manifest. Have regard here to the distinction made, and to the special weight of every word. Elizabeth was the first to hear the voice of Mary's salutation, but John was the first to receive grace. She heard naturally, but he leaped mystically. She hailed the coming of Mary, he that of the Lord, Mary and Elizabeth spake words full of grace, but Jesus and John worked, and commenced their mystery of godliness from their mothers' beginnings, and so by twin miracles the mothers prophesied from the spirit of their unborn offspring. The babe leaped, and the mother was filled with the Holy Ghost. The mother was not filled before the son, but when the son was filled with the Holy Ghost, he filled his mother also.\par
\switchcolumn*

\LA
\begin{Resp}\Rsign Beátam me dicent omnes generatiónes, \Star{} Quia fecit mihi Dóminus magna qui potens est, et sanctum nomen ejus.\end{Resp}
\begin{Resp}\Vsign Et misericórdia ejus a progénie in progénies timéntibus eum.\end{Resp}
\begin{Resp}\Rsign Quia fecit mihi Dóminus magna qui potens est, et sanctum nomen ejus.\end{Resp}
\switchcolumn
\EN
\begin{Resp}\Rsign All generations shall call me blessed. \Star{} For He that is mighty, even the Lord, hath done to me great things; and Holy is His Name.\end{Resp}
\begin{Resp}\Vsign And his mercy is on them that fear Him, from generation to generation.\end{Resp}
\begin{Resp}\Rsign He that is mighty, even the Lord, hath done to me great things; and Holy is His Name.\end{Resp}
\switchcolumn*

\LA
\LessonHead{Lectio VIII}
\switchcolumn
\EN
\LessonHead{Lesson VIII}
\switchcolumn*

\LA
\noindent Et unde hoc mihi, ut véniat Mater Dómini mei ad me? Hoc est, Quo tantum bonum mihi áccidit, ut mater Dómini mei véniat ad me? Miráculum séntio, agnósco mystérium: Mater Dómini Verbo fœta, Deo plena est. Mansit autem María cum illa ménsibus tribus, et revérsa est in domum suam. Bene indúcitur sancta María et exhibuísse offícium, et mýsticum númerum custodísse.\par
\switchcolumn
\EN
\noindent And whence is this to me, that the Mother of my Lord should come to me? That is to say, How cometh it to pass that so great a good should befall me, as that the Mother of my Lord should come to me I feel the miracle, I acknowledge the mystery: the Mother of my Lord, pregnant with the Word, is full of God. And Mary abode with her about three months, and returned to her own house. It is meet to record how Mary showed this kindness, and abode this mystic number of months.\par
\switchcolumn*

\LA
\begin{Resp}\Rsign Felix namque es, sacra Virgo María, et omni laude digníssima: \Star{} Quia ex te ortus est sol justítiæ, * Christus, Deus noster.\end{Resp}
\begin{Resp}\Vsign Ora pro pópulo, intérveni pro clero, intercéde pro devóto femíneo sexu: séntiant omnes tuum juvámen, quicúmque célebrant tuam sanctam Visitatiónem.\end{Resp}
\begin{Resp}\Rsign Quia ex te ortus est sol justítiæ, Christus, Deus noster.\end{Resp}
\begin{Resp}\Vsign Glória Patri, et Fílio, \Star{} et Spirítui Sancto.\end{Resp}
\begin{Resp}\Rsign Christus, Deus noster.\end{Resp}
\switchcolumn
\EN
\begin{Resp}\Rsign O Holy Virgin Mary, happy indeed art thou, and right worthy of all praise \Star{} For out of thee rose the Sun of righteousness, * Even Christ our God.\end{Resp}
\begin{Resp}\Vsign Pray for the people, plead for the clergy, make intercession for all women vowed to God. May all that are keeping this thy Visitation feel the might of thine assistance.\end{Resp}
\begin{Resp}\Rsign For out of thee rose the Sun of righteousness, even Christ our God.\end{Resp}
\begin{Resp}\Vsign Glory be to the Father, and to the Son, \Star{} and to the Holy Ghost.\end{Resp}
\begin{Resp}\Rsign Even Christ our God.\end{Resp}
\switchcolumn*

\end{paracol}

\par\needspace{10\baselineskip}\addvspace{1.2em}{\centering\small\textsc{Die 2 Julii} · \textsc{2 July}\par}
\Office{Ss. Processi et Martiniani Martyrum}{Ss. Processus et Martinian, Martyrs}{Simplex}
\addcontentsline{toc}{subsection}{Die 2 Julii · Ss. Processi et Martiniani Martyrum \textit{— Ss. Processus et Martinian, Martyrs}}
\begin{paracol}{2}
\LA
\LessonHead{Lectio IX (pro commemoratione)}
\switchcolumn
\EN
\LessonHead{Lesson IX (when commemorated)}
\switchcolumn*

\LA
\LessonTitle{Commemoratio: Ss. Processi et Martiniani Martyrum}
\switchcolumn
\EN
\LessonTitle{Commemoration of: Ss. Processus et Martinian, Martyrs}
\switchcolumn*

\LA
\noindent Quo témpore Petrus et Paulus tenebántur sub custódia Mamertíni in monte Tarpéjo, duo custódes, Procéssus et Martiniánus, cum áliis quadragínta, Apostolórum prædicatióne miraculísque commóti, se ad Jesu Christi fidem convertérunt; et, cum repénte fons e saxo ortus esset, baptizáti sunt. Qui permisérunt Apóstolis, ut, si vellent, abírent. Sed Paulínus mílitum præféctus, re cógnita, Procéssum et Martiniánum a suscépto consílio revocáre conátur. Qui, cum frustra tempus contéreret, ipsórum ora saxo contúndi dentésque commínui jubet. Mox ad Jovis státuam addúctos, cum eádem constántia veneratúros se idóla negárent, ímperat equúleo torquéri, candéntibus láminis ad eórum corpus admótis, ac cædi fústibus; quibus in cruciátibus una hæc illórum vox audiebátur: Sit nomen Dómini benedíctum. Dénique conjécti in cárcerem, paulo post extra Urbem via Aurélia secúri feriúntur. Quorum córpora Lucína in prǽdio suo sepelívit, sexto Nonas Júlii; quæ póstea, in Urbem transláta, in basílica Príncipis Apostolórum cóndita sunt.\par
\switchcolumn
\EN
\noindent At what time Peter and Paul were kept in the Mamertine Prison, at the foot of the Tarpeian Rock, two of the gaolers, named Processus and Martinian, along with other forty, were moved by the preaching and miracles of the Apostles to believe in Christ, and were baptized in a spring which suddenly brake forth out of the rock. These men let the Apostles depart if they willed it. But Paulinus, Prefect of the soldiers, when he heard what was come to pass, strove to turn away Processus and Martinian from their purpose. And when he found that he but wasted time, he ordered their faces to be bruised and their teeth to be broken, with stones. Moreover, when he had had them led to the image of Jupiter, and they still boldly answered that they would not worship the gods, he caused them to be tormented on the rack, and white-hot plates of metal to be put to their flesh, and that they should be beaten with sticks. While they were suffering all these things they were heard to say only this one word Blessed be the Name of the Lord. They were afterwards cast into prison, and in a little while they were taken outside the city, and slain with the axe, upon the Aurelian Way. The Lady Lucina buried their bodies upon her own farm, upon the 2nd day of July, but they were afterwards brought into the City, and are buried in the Church of the Prince of the Apostles.\par
\switchcolumn*

\LA
\TeDeum{Te Deum}
\switchcolumn
\EN
\TeDeum{Te Deum}
\switchcolumn*

\end{paracol}

\par\needspace{10\baselineskip}\addvspace{1.2em}{\centering\small\textsc{Die 3 Julii} · \textsc{3 July}\par}
\Office{S. Leonis Papæ et Confessoris}{St. Leo II, Pope and Confessor}{Semiduplex}
\addcontentsline{toc}{subsection}{Die 3 Julii · S. Leonis Papæ et Confessoris \textit{— St. Leo II, Pope and Confessor}}
\begin{paracol}{2}
\LA
\RefLine{Lectiones I–III: de Scriptura occurrente.}
\switchcolumn
\EN
\RefLine{Lessons I–III: from the Scripture of the day.}
\switchcolumn*

\LA
\RefLine{Lectiones VII–IX: ex Commúni Summorum Pontificum Confessorum (C4b).}
\switchcolumn
\EN
\RefLine{Lessons VII–IX: from the Common of Popes (C4b).}
\switchcolumn*

\LA
\LessonHead{Lectio IV}
\switchcolumn
\EN
\LessonHead{Lesson IV}
\switchcolumn*

\LA
\noindent Leo secúndus, Póntifex máximus, Sículus, humánis et divínis lítteris Græce et Latíne doctus, músicis étiam erudítus fuit; ipse enim sacros Hymnos et Psalmos in Ecclésia ad concéntum meliórem redúxit. Probávit acta sextæ Synodi, quæ Constantinópoli celebráta est, præsidéntibus legátis apostólicæ Sedis, præsénte quoque Constantíno imperatóre, et duobus patriárchis, Constantinopolitáno et Antiochéno, ac centum septuagínta epíscopis; quam et in Latínum tránstulit.\par
\switchcolumn
\EN
\noindent Leo II was a Sicilian. He was learned in sacred and worldly letters in the Greek and Latin tongues, and was moreover an excellent musician. He re-arranged and improved the music of the sacred hymns and Psalms used in the Church. He approved the acts of, the sixth general Council, which was held at Constantinople, under the Presidency of the Legates of the Apostolic See, in the presence of the Emperor Constantine (the Bearded,) the Patriarchs of Constantinople and Antioch, and one hundred and seventy Bishops. Leo also translated the decrees of the Council from Greek into Latin.\par
\switchcolumn*

\LA
\begin{Resp}\Rsign Invéni David servum meum, óleo sancto meo unxi eum: \Star{} Manus enim mea auxiliábitur ei.\end{Resp}
\begin{Resp}\Vsign Nihil profíciet inimícus in eo, et fílius iniquitátis non nocébit ei.\end{Resp}
\begin{Resp}\Rsign Manus enim mea auxiliábitur ei.\end{Resp}
\switchcolumn
\EN
\begin{Resp}\Rsign I have found David My servant, with My holy oil have I anointed him \Star{} For My hand shall help him.\end{Resp}
\begin{Resp}\Vsign The enemy shall prevail nothing against him, nor the son of wickedness afflict him.\end{Resp}
\begin{Resp}\Rsign For My hand shall help him.\end{Resp}
\switchcolumn*

\LA
\LessonHead{Lectio V}
\switchcolumn
\EN
\LessonHead{Lesson V}
\switchcolumn*

\LA
\noindent In eo concílio Cyrus, Sergius et Pyrrhus condemnáti sunt, unam tantúmmodo voluntátem et operatiónem in Christo prædicántes. Hic fregit supérbiam antístitum Ravennátum, qui, exarchórum freti poténtia, Sedi apostólicæ non obtemperábant. Quam ob rem decrévit ut eléctio cleri Ravennátis írrita esset, nisi Románi Pontíficis auctoritáte comprobarétur.\par
\switchcolumn
\EN
\noindent It was in this Council that Cyrus, Sergius, and Pyrrhus were condemned for teaching that there is in Christ only one Will and one Working. Leo broke the pride of the Archbishops of Ravenna, who had puffed themselves up under the power of the Exarchs to set at naught the power of the Apostolic See. Wherefore he decreed that the elections of the clergy of Ravenna should be nothing worth, until they had been confirmed by the authority of the Bishop of Rome.\par
\switchcolumn*

\LA
\begin{Resp}\Rsign Pósui adjutórium super poténtem, et exaltávi eléctum de plebe mea: \Star{} Manus enim mea auxiliábitur ei.\end{Resp}
\begin{Resp}\Vsign Invéni David servum meum, óleo sancto meo unxi eum.\end{Resp}
\begin{Resp}\Rsign Manus enim mea auxiliábitur ei.\end{Resp}
\switchcolumn
\EN
\begin{Resp}\Rsign I have laid help upon one that is mighty, and have exalted one chosen out of My people \Star{} For My hand shall help him.\end{Resp}
\begin{Resp}\Vsign I have found David My servant, with My holy oil have I anointed him.\end{Resp}
\begin{Resp}\Rsign For My hand shall help him.\end{Resp}
\switchcolumn*

\LA
\LessonHead{Lectio VI}
\switchcolumn
\EN
\LessonHead{Lesson VI}
\switchcolumn*

\LA
\noindent Vere pater páuperum fuit; non enim pecúnia solum, sed opera, labóre et consíliis, egéntium viduárum et pupillórum inópiam ac solitúdinem sublevábat. Qui, dum síngulos non magis prædicatióne quam vita ad pie sanctéque vivéndum adhortarétur, obdormívit in Dómino mense sui pontificátus undécimo, quinto Nonas Julii, anno sexcentésimo octogésimo tertio, sepultúsque est in basilica sancti Petri. Ordinatióne una, mense Junio, creávit presbýteros novem, diáconos tres, epíscopos divérsis in locis vigínti tres.\par
\switchcolumn
\EN
\noindent He was a very father to the poor. Not by money only, but by his work, his labours, and his advice he relieved the poverty and loneliness of widows and orphans. He was leading all to live holy and godly lives, not by mere preaching, but by his own life, when he fell asleep in the Lord, having sat as Pope nine months and twenty-seven days. He was buried in the Church of St. Peter upon the 28th day of June. In the month of June he held one ordination whereat he ordained nine Priests, three Deacons, and twenty-three Bishops for diverse places.\par
\switchcolumn*

\LA
\begin{Resp}\Rsign Iste est, qui ante Deum magnas virtútes operátus est, et omnis terra doctrína ejus repléta est: \Star{} Ipse intercédat pro peccátis ómnium populórum.\end{Resp}
\begin{Resp}\Vsign Iste est, qui contémpsit vitam mundi, et pervénit ad cæléstia regna.\end{Resp}
\begin{Resp}\Rsign Ipse intercédat pro peccátis ómnium populórum.\end{Resp}
\begin{Resp}\Vsign Glória Patri, et Fílio, \Star{} et Spirítui Sancto.\end{Resp}
\begin{Resp}\Rsign Ipse intercédat pro peccátis ómnium populórum.\end{Resp}
\switchcolumn
\EN
\begin{Resp}\Rsign This is he which wrought great wonders before God, and the whole earth is full of his teaching \Star{} May he pray for all people, that their sins may be forgiven unto them\end{Resp}
\begin{Resp}\Vsign This is he which loved not his life in this world, and hath attained unto the kingdom of heaven.\end{Resp}
\begin{Resp}\Rsign May he pray for all people, that their sins may be forgiven unto them\end{Resp}
\begin{Resp}\Vsign Glory be to the Father, and to the Son, \Star{} and to the Holy Ghost.\end{Resp}
\begin{Resp}\Rsign May he pray for all people, that their sins may be forgiven unto them\end{Resp}
\switchcolumn*

\LA
\LessonHead{Lectio IX (pro commemoratione)}
\switchcolumn
\EN
\LessonHead{Lesson IX (when commemorated)}
\switchcolumn*

\LA
\LessonTitle{Commemoratio: S. Leonis Papæ et Confessoris}
\switchcolumn
\EN
\LessonTitle{Commemoration of: St. Leo II, Pope and Confessor}
\switchcolumn*

\LA
\noindent Leo secúndus, Sículus, humánis et divínis lítteris Græce et Latíne doctus, músicis étiam erudítus fuit; ipse enim sacros Hymnos et Psalmos in Ecclésia ad concéntum meliórem redúxit. Probávit acta sextæ Synodi, quæ, præsidéntibus legátis apostólicæ Sedis, Constantinópoli celebráta est; quam et in Latínum tránstulit. In eo Concílio Cyrus, Sérgius et Pyrrhus condemnáti sunt, unam tantúmmodo voluntátem et operatiónem in Christo prædicántes. Vere pater páuperum fuit; non enim pecúnia solum, sed opera, labóre et consíliis egéntium viduárum et pupillórum inópiam et solitúdinem sublevábat. Obdormívit in Dómino mense sui pontificátus undécimo, quinto Nonas Júlii, anno sexcentésimo octogésimo tértio, sepultúsque est in basílica sancti Petri.\par
\switchcolumn
\EN
\noindent Leo II, Supreme Pontiff, a Sicilian, was learned in sacred and profane letters in Greek and Latin, and was moreover an excellent musician; for he reduced to better harmony the sacred hymns and psalms used in the Church. He approved the acts of the sixth Council, which was held at Constantinople, and translated them into Latin. It was at this Council that Cyrus, Sergius, and Pyrrhus were condemned, for teaching that there is in Christ only one will and one operation. He was indeed a father to the poor; for not by money alone, but by his deeds, his labours, and his advice, he relieved the poverty and loneliness of needy widows and orphans. He fell asleep in the Lord on the 3rd day of July in the year 683, in the eleventh month of his pontificate, and was buried in the Basilica of St. Peter.\par
\switchcolumn*

\LA
\TeDeum{Te Deum}
\switchcolumn
\EN
\TeDeum{Te Deum}
\switchcolumn*

\end{paracol}

\par\needspace{10\baselineskip}\addvspace{1.2em}{\centering\small\textsc{Die 3 Julii} · \textsc{3 July}\par}
\Office{Quinta die infra Octavam Ss. Apostolorum Petri et Pauli}{Fifth Day Within the Octave of Sts. Peter and Paul, Apostles}{Semiduplex}
\addcontentsline{toc}{subsection}{Die 3 Julii · Quinta die infra Octavam Ss. Apostolorum Petri et Pauli \textit{— Fifth Day Within the Octave of Sts. Peter and Paul, Apostles}}
\begin{paracol}{2}
\LA
\LessonHead{Lectio IV}
\switchcolumn
\EN
\LessonHead{Lesson IV}
\switchcolumn*

\LA
\LessonTitle{Commemoratio: Quinta die infra Octavam Ss. Apostolorum Petri et Pauli}
\switchcolumn
\EN
\LessonTitle{Commemoration of: Fifth Day Within the Octave of Sts. Peter and Paul, Apostles}
\switchcolumn*

\LA
\Source{Ex Serm. 82. in natali Ss. Apost. Petri et Pauli, in fine.}
\switchcolumn
\EN
\Source{1st for the Birth-day of SS. Peter and Paul}
\switchcolumn*

\LA
\noindent De Sermone sancti Leonis Papæ.\par
\switchcolumn
\EN
\noindent From the Sermons of Pope St. Leo the Great\par
\switchcolumn*

\LA
Pretiosa est in conspectu Domini mors sanctorum ejus: nec ullo crudelitatis genere destrui potest sacramento crucis Christi fundata religio. Non minuitur persecutionibus Ecclesia, sed augetur: et semper Dominicus ager segete ditiori vestitur, dum grana, quæ singula cadunt, multiplicata nascuntur. Unde duo ista præclara divini seminis germina in quantam sobolem germinarint, beatorum millia Mártyrum protestantur, qui apostolicorum æmuli triumphorum, Urbem nostram purpuratis et longe lateque rutilantibus populis ambierunt, et quasi ex multarum honore gemmarum conserto uno diademate coronarunt.\par
\switchcolumn
\EN
Precious in the sight of the Lord is the death of His Saints and no cruelty can trample out that Religion which was founded in the mysterious Sacrifice of the Cross of Christ. Under persecution the Church waneth not, but waxeth. And the Lord's field is ever clothed with a nobler harvest, when the grains of His wheat abide not alone, but fall into the ground, and die, and bring forth fruit manifold, (John xii. 24.) What an increase hath sprung up from these two glorious grains of God's seed, is witnessed by the thousands of blessed martyrs who have rivalled the triumph of the Apostles, who gird this our city with their red and world famous armies, and have set upon her head a crown of many kingly jewels.\par
\switchcolumn*

\LA
De quorum præsidio, dilectissimi, divinitus nobis ad exemplum patientiæ, et confirmationem fidei præparato, universaliter quidem in omnium Sanctorum commemoratione lætandum est: sed in horum excellentia Patrum merito est excellentius gloriandum, quos gratia Dei in tantum apicem inter omnia Ecclesiæ membra provexit, ut eos in corpore, cui caput est Christus, quasi geminum constitueret lumen oculorum. De quorum meritis atque virtutibus, quæ omnem loquendi superant facultatem, nihil diversum, nihil debemus sentire discretum: quia illos et electio pares, et labor similes, et finis fecit æquales.\par
\switchcolumn
\EN
Up their united succour, dearly beloved brethren, which hath been given unto us by God as an example of patience and a confirmation of faith, will be the cause of our joy on the days of the commemoration of all the Saints, but amid all their glory, doth excel in glory the glory of their spiritual fathers, whom the grace of God exalted to such a pitch of eminence among all the members of the Church, that they seem as it were the two eyes of that body whereof Christ is the Head. The worthy deed and loyal courage of either of them are such as to baffle all power of man's speech worthily to express the same, nor need we draw any distinction between them. In their election they were peers, in their work fellows, and in their end equals.\par
\switchcolumn*

\LA
Sicut autem et nos experti sumus, et nostri probavere majores, credimus atque confidimus, inter omnes labores istius vitæ ad obtinendam misericordiam Dei semper nos specialium patronorum orationibus adjuvandos: ut quantum propriis peccatis deprimimur, tantum apostolicis meritis erigamur: Per Dominum nostrum Jesum Christum, cui est cum Patre et sancto Spiritu eadem potestas, una divinitas, in sæcula sæculorum. Amen.\par
\switchcolumn
\EN
But, as our own experience hath taught us, and that of our forefathers hath proved, we believe and trust that, amid all the toils whereby in this life we strive towards the obtaining of God's mercy, we shall ever be holpen by the prayers of these our special Patrons, and that when our own sins thrust us down, the good deeds of the Apostles may pull us up. Through our Lord Jesus Christ, Who, with the Father and the Holy Ghost, hath but one power, as He hath one Divine nature, for ever and ever. Amen.\par
\switchcolumn*

\LA
\TeDeum{Te Deum}
\switchcolumn
\EN
\TeDeum{Te Deum}
\switchcolumn*

\end{paracol}

\par\needspace{10\baselineskip}\addvspace{1.2em}{\centering\small\textsc{Die 4 Julii} · \textsc{4 July}\par}
\Office{Sexta die infra Octavam Ss. Apostolorum Petri et Pauli}{Sixth Day Within the Octave of the Apostles Sts. Peter and Paul}{Semiduplex}
\addcontentsline{toc}{subsection}{Die 4 Julii · Sexta die infra Octavam Ss. Apostolorum Petri et Pauli \textit{— Sixth Day Within the Octave of the Apostles Sts. Peter and Paul}}
\begin{paracol}{2}
\LA
\RefLine{Lectiones I–III: de Scriptura occurrente.}
\switchcolumn
\EN
\RefLine{Lessons I–III: from the Scripture of the day.}
\switchcolumn*

\LA
\LessonHead{Lectio IV}
\switchcolumn
\EN
\LessonHead{Lesson IV}
\switchcolumn*

\LA
\LessonTitle{De Expositióne sancti Joánnis Chrysóstomi in Epístolam ad Romános}
\switchcolumn
\EN
\LessonTitle{From the Exposition of the Epistle of St. Paul to the Romans, written by St. John Chrysostom, Patriarch of Constantinople}
\switchcolumn*

\LA
\Source{Sermo 32 in morali exhortat}
\switchcolumn
\EN
\Source{Serm. xxxii. in the Moral Exhortations}
\switchcolumn*

\LA
\noindent Cum Paulus Apóstolus grátiam Dómini nostri Jesu Christi, matrem ómnium bonórum, nobis precétur; réliquum est, ut nos tali patrocinio dignos exhibeamus, ut non hic solum vocem Pauli audiámus, sed et postquam illuc migravérimus, athletam Christi vidére mereamur. Immo, si hic audivérimus, et illic ipsum omnino vidébimus, licet non e propinquo stantes; vidébimus tamen prope regalem thronum splendentem, ubi Chérubim Deum glorificant, ubi Séraphim volant. Illic Paulum vidébimus cum Petro Sanctórum chori príncipem ac ducem, et ejus germana caritate fruemur.\par
\switchcolumn
\EN
\noindent Apostle Paul wisheth unto us the grace of our Lord Jesus Christ, as the mother of all good, and it remaineth for us to show ourselves worthy of the care of such a Protector, that we may not only listen to Paul's voice, while we are here, but when we pass away to the hereafter, may earn a sight of that great soldier of Christ. Yea if we listen to him here, we shall see him there. Not nigh, but from afar off, shall we see him see him standing near the glory of that Kingly throne where the Cherubim glorify God, where the Seraphim are flying, there shall we see Paul along with Peter, a prince and a leader of the army of the saints, and we shall rejoice in his brotherly love.\par
\switchcolumn*

\LA
\begin{Resp}\Rsign Vidi conjúnctos viros, habéntes spléndidas vestes, et Ángelus Dómini locútus est ad me, dicens: \Star{} Isti sunt viri sancti facti amíci Dei.\end{Resp}
\begin{Resp}\Vsign Vidi Ángelum Dei fortem, volántem per médium cælum, voce magna clamántem et dicéntem.\end{Resp}
\begin{Resp}\Rsign Isti sunt viri sancti facti amíci Dei.\end{Resp}
\switchcolumn
\EN
\begin{Resp}\Rsign I saw men standing together, clad in shining raiment, and the Angel of the Lord spoke unto me, saying: \Star{} These men are holy, for they are the friends of God.\end{Resp}
\begin{Resp}\Vsign I saw a strong Angel of God fly into the midst of heaven, saying with a loud voice:\end{Resp}
\begin{Resp}\Rsign These men are holy, for they are the friends of God.\end{Resp}
\switchcolumn*

\LA
\LessonHead{Lectio V}
\switchcolumn
\EN
\LessonHead{Lesson V}
\switchcolumn*

\LA
\noindent Si enim, cum hic esset, usque adeo diligebat hómines, ut, cum dissolvi et cum Christo esse cúperet, elegerit tamen hic esse; multo magis illic ferventiórem caritátem osténdet. Ego et Romam proptérea diligo, tametsi aliunde illam laudare queam, nempe a magnitúdine, ab antiquitáte, a pulchritúdine, a multitúdine, ab imperio, a divítiis, et a rebus in bello fortiter gestis. Sed, his ómnibus omissis, ob id illam beátam prædico, quod erga illos Paulus, dum viveret, adeo fuit benévolus, adeo illos amávit, et coram disseruit, et postremo vitam apud eos finívit. Cujus sanctum corpus ipsi póssident. Et proptérea cívitas illa hinc facta est insígnis magis quam ab aliis rebus ómnibus; et tamquam corpus magnum ac validum duos habet óculos fulgéntes, Sanctórum videlicet horum corpora.\par
\switchcolumn
\EN
\noindent For if, while he was yet here, he so loved men, that, although he would fain have been dissolved and been with Christ, yet he was willing still to tarry for man's sake, much greater is the tender love which he now showeth. This is why I love Rome, although if I would, there are many other things for which I might praise her her greatness, her antiquity, her beauty, her population, her empire, her wealth, or her victories. But all these I pass by, and I call Rome blessed for this cause, that Paul in his lifetime loved her children so well, was so kindly toward them, taught openly there, and at length laid down his life among them. They have there his holy body, and this alone maketh that city illustrious more than doth aught else. And just as a great and strong body hath two bright eyes, so are the bodies of these two Holy Apostles in the city of Rome.\par
\switchcolumn*

\LA
\begin{Resp}\Rsign Beáti estis, cum maledíxerint vobis hómines, et persecúti vos fúerint, et díxerint omne malum advérsum vos, mentiéntes, propter me: \Star{} Gaudéte et exsultáte, quóniam merces vestra copiósa est in cælis.\end{Resp}
\begin{Resp}\Vsign Cum vos óderint hómines, et cum separáverint vos, et exprobráverint, et ejécerint nomen vestrum tamquam malum propter Fílium hóminis.\end{Resp}
\begin{Resp}\Rsign Gaudéte et exsultáte, quóniam merces vestra copiósa est in cælis.\end{Resp}
\switchcolumn
\EN
\begin{Resp}\Rsign Blessed are ye when men shall revile you, and persecute you, and shall say all manner of evil against you falsely, for My sake; \Star{} Rejoice, and be exceeding glad, for great is your reward in heaven.\end{Resp}
\begin{Resp}\Vsign When men shall hate you, and when they shall separate you from their company, and shall reproach you, and cast out your name as evil, for the Son of Man's sake.\end{Resp}
\begin{Resp}\Rsign Rejoice, and be exceeding glad, for great is your reward in heaven.\end{Resp}
\switchcolumn*

\LA
\LessonHead{Lectio VI}
\switchcolumn
\EN
\LessonHead{Lesson VI}
\switchcolumn*

\LA
\noindent Non ita cælum splendéscit quando radios sol demittit, quemádmodum Romanórum urbs duos istos fulgóres ubíque terrárum emíttens. Hinc rapiétur Paulus, hinc Petrus. Considerate et horrete, quale spectaculum visura sit Roma; Paulum videlicet, repénte ex theca illa cum Petro resurgéntem, in occúrsum Dómini sursum ferri. Qualem rosam Christo mittet Roma! quálibus corónis duabus ornátur urbs ista! quálibus catenis áureis cincta est! quales habet fontes! Proptérea admiror hanc urbem, non propter colúmnas, neque propter aliam quamcúmque rerum spéciem, sed propter colúmnas illas Ecclésiæ. Quis mihi nunc dabit circumvolvi corpori Pauli, affigi sepúlcro, vidére púlverem corporis illíus, quæ adhuc in Christo déerant adimplentis, stígmata illíus gestántis, prædicatiónem Evangélii ubíque seminantis?\par
\switchcolumn
\EN
\noindent Not brighter is the sky when the sun doth make it all light with his beams, than is the city of Rome darting forth these twin rays of light to the uttermost bounds of the earth. There it is that Paul, there it is that Peter, will rise, and be caught up to meet the Lord in the air. (Thess. iv. 16.) Think, and thrill at the thought, of what Rome will see then, when she beholdeth Paul and Peter rising suddenly out of that coffin, to be caught up to meet the Lord. What rose will Rome offer to Christ? What twin crowns are they wherewith that city is adorned withal? What fresh springs hath she in her? Therefore it is that I marvel at that city, not because of the abundance of her gold, not because of her pillars, not because of any other loveliness that she hath, but because of these two pillars of the Church. Would that I could even now embrace the corpse of Paul that I could cling to his grave that I could see the dust of that body, which filled up those things that were behind of the sufferings of Christ (Col. i. 24,) which bore about in it the marks of the Lord Jesus (Gal. vi. 17,) and which went everywhere carrying the seed of the Gospel.\par
\switchcolumn*

\LA
\begin{Resp}\Rsign Isti sunt triumphatóres et amíci Dei, qui contemnéntes jussa príncipum, meruérunt prǽmia ætérna: \Star{} Modo coronántur, et accípiunt palmam.\end{Resp}
\begin{Resp}\Vsign Isti sunt qui venérunt ex magna tribulatióne, et lavérunt stolas suas in sánguine Agni.\end{Resp}
\begin{Resp}\Rsign Modo coronántur, et accípiunt palmam.\end{Resp}
\begin{Resp}\Vsign Glória Patri, et Fílio, \Star{} et Spirítui Sancto.\end{Resp}
\begin{Resp}\Rsign Modo coronántur, et accípiunt palmam.\end{Resp}
\switchcolumn
\EN
\begin{Resp}\Rsign These are they which have conquered, and are become the friends of God, who recked not of the commandments of princes, and earned the everlasting reward. \Star{} And now have they crowns on their heads, and palms in their hands.\end{Resp}
\begin{Resp}\Vsign These are they which came out of great tribulation, and have washed their robes in the blood of the Lamb.\end{Resp}
\begin{Resp}\Rsign And now have they crowns on their heads, and palms in their hands.\end{Resp}
\begin{Resp}\Vsign Glory be to the Father, and to the Son, \Star{} and to the Holy Ghost.\end{Resp}
\begin{Resp}\Rsign And now have they crowns on their heads, and palms in their hands.\end{Resp}
\switchcolumn*

\LA
\LessonHead{Lectio VII}
\switchcolumn
\EN
\LessonHead{Lesson VII}
\switchcolumn*

\LA
\LessonTitle{Léctio sancti Evangélii secúndum Matthǽum}
\switchcolumn
\EN
\LessonTitle{The Lesson is taken from the Holy Gospel according to Matthew}
\switchcolumn*

\LA
\Cite{Matt 19:27-29}
\switchcolumn
\EN
\Cite{Matt 19:27-29}
\switchcolumn*

\LA
\noindent In illo tempore: Dixit Petrus ad Jesum: Ecce nos reliquimus omnia, et secuti sumus te: quid ergo erit nobis? Et reliqua.\par
\switchcolumn
\EN
\noindent At that time: Peter said unto Jesus: Behold, we have forsaken all, and followed Thee: what shall we have therefore? And so on.\par
\switchcolumn*

\LA
\LessonTitle{De Homilía venerábilis Bedæ Presbýteri}
\switchcolumn
\EN
\LessonTitle{Homily by the Venerable Bede, Priest (at Jarrow and Doctor of the Church).}
\switchcolumn*

\LA
\Source{In Natali S. Benedicti}
\switchcolumn
\EN
\Source{For St Benedict’s Birthday.}
\switchcolumn*

\LA
\noindent Duo sunt órdines electórum in judício futúri: unus judicántium cum Dómino, de quibus hoc loco mémorat, qui reliquérunt ómnia, et secúti sunt illum: alius judicandórum a Dómino, qui non quidem ómnia sua páriter reliquérunt, sed de his tamen quæ habébant, quotidiánas dare eleemósynas Christi paupéribus curábant: unde et auditúri sunt in judício: Veníte, benedícti Patris mei, possidíte præparátum vobis regnum a constitutióne mundi. Esurívi enim, et dedístis mihi manducáre: sitívi, et dedístis mihi bíbere.\par
\switchcolumn
\EN
\noindent In the judgment to come, the elect will be in two classes. One class are they who have forsaken all, and followed the Lord: and these shall judge along with Him. The other class are they who have not equally forsaken all that they had, but who have been careful daily to give alms of their goods to the poor of Christ: these shall be the subjects of judgment, and these are they who shall then hear these words: “Come, ye blessed of My Father, inherit the kingdom prepared for you from the foundation of the world: for I was an hungered, and ye gave Me meat: I was thirsty, and ye gave Me drink.” (Matth. xxv. 34, 35.)\par
\switchcolumn*

\LA
\begin{Resp}\Rsign Isti sunt qui vivéntes in carne, plantavérunt Ecclésiam sánguine suo: \Star{} Cálicem Dómini bibérunt, et amíci Dei facti sunt.\end{Resp}
\begin{Resp}\Vsign In omnem terram exívit sonus eórum, et in fines orbis terræ verba eórum.\end{Resp}
\begin{Resp}\Rsign Cálicem Dómini bibérunt, et amíci Dei facti sunt.\end{Resp}
\switchcolumn
\EN
\begin{Resp}\Rsign These are they who while yet they lived in the flesh, planted the Church in their own blood; \Star{} They drank of the Lord's cup, and became the friends of God.\end{Resp}
\begin{Resp}\Vsign Their sound is gone out through all the earth, and their words to the ends of the world.\end{Resp}
\begin{Resp}\Rsign They drank of the Lord's cup, and became the friends of God.\end{Resp}
\switchcolumn*

\LA
\LessonHead{Lectio VIII}
\switchcolumn
\EN
\LessonHead{Lesson VIII}
\switchcolumn*

\LA
\noindent Sed et reprobórum duos ibi futúros órdines Dómino narránte comperímus: unum eórum, qui fídei christiánæ mystériis initiáti, ópera fídei exercére contémnunt, quibus in judício testátur: Discédite a me maledícti in ignem aetérnum, qui præparátus est diábolo, et ángelis ejus: esurívi enim, et non dedístis mihi manducáre. Alterum eórum, qui fidem et mystéria Christi vel numquam suscepére, vel suscéptam per apostasíam deseruére: de quibus dicit: Qui autem non credit, jam judicatus est: quia non credit in nómine unigéniti Fílii Dei.\par
\switchcolumn
\EN
\noindent Of the reprobate also we gather, from the words of the Lord, that there will be two classes. One class are they who, being made partakers in the mystery of Christian faith, have neglected to show their faith by their works: these are they to whom it will be said at the judgment: “Depart from Me, ye cursed, into everlasting fire, prepared for the devil and his angels: for I was an-hungered, and ye gave Me no meat.” The other class are they who either have never received the faith and mysteries of Christ, or who, having received, have apostatised, and abandoned it: and touching these it is said: “But he that believeth not is condemned already, because he hath not believed in the name of the only-begotten Son of God.” (John iii. 18.)\par
\switchcolumn*

\LA
\begin{Resp}\Rsign Isti sunt viri sancti, quos elégit Dóminus in caritáte non ficta, et dedit illis glóriam sempitérnam: \Star{} Quorum doctrína fulget Ecclésia, ut sole luna.\end{Resp}
\begin{Resp}\Vsign Sancti per fidem vicérunt regna: operáti sunt justítiam.\end{Resp}
\begin{Resp}\Rsign Quorum doctrína fulget Ecclésia, ut sole luna.\end{Resp}
\begin{Resp}\Vsign Glória Patri, et Fílio, \Star{} et Spirítui Sancto.\end{Resp}
\begin{Resp}\Rsign Quorum doctrína fulget Ecclésia, ut sole luna.\end{Resp}
\switchcolumn
\EN
\begin{Resp}\Rsign These men are saints, whom the Lord hath chosen in love unfeigned, and hath given them glory everlasting. These are they \Star{} By the light of whose teaching the Church is glorified, even as the moon is glorified by the light of the sun.\end{Resp}
\begin{Resp}\Vsign The saints through faith subdued kingdoms, wrought righteousness.\end{Resp}
\begin{Resp}\Rsign By the light of whose teaching the Church is glorified, even as the moon is glorified by the light of the sun.\end{Resp}
\begin{Resp}\Vsign Glory be to the Father, and to the Son, \Star{} and to the Holy Ghost.\end{Resp}
\begin{Resp}\Rsign By the light of whose teaching the Church is glorified, even as the moon is glorified by the light of the sun.\end{Resp}
\switchcolumn*

\LA
\LessonHead{Lectio IX}
\switchcolumn
\EN
\LessonHead{Lesson IX}
\switchcolumn*

\LA
\noindent Verum his cum timóre et pavóre débito paulísper commemorátis, ad lætíssima pótius Dómini et Salvatóris nostri promíssa convertámus audítum. Videámus quæ tantæ grátia pietátis: non ætérnae tantúmmodo vitæ præmia suis sequácibus, sed et præséntis múnera pollicétur exímia. Et omnis, inquit, qui relíquerit domum, vel fratres, aut soróres, aut patrem, aut matrem, aut uxórem, aut fílios, aut agros, propter nomen meum, céntuplum accípiet, et vitam ætérnam possidébit. Qui enim terrénis afféctibus sive possessiónibus pro Christi discipulátu renuntiáverit, quo plus in ejus amórem profécerit, eo plures invéniet, qui se intérno suscípere afféctu, et suis gáudeant sustentáre substántiis.\par
\switchcolumn
\EN
\noindent And now that we have touched for a moment, with fear and just dread, upon these things, let us rather turn our hearing to the right joyful promises of our Lord and Saviour. Let us look what His so great, beautiful, and fatherly love will give to such as follow Him; not the reward of life everlasting only, but gifts exceeding precious in this life also. “Every one,” saith He, that hath forsaken houses, or brethren, or sisters, or father, or mother, or wife, or children, or lands, for My Name’s sake, shall receive an hundredfold, and shall inherit everlasting life.” For every one that shall forsake earthly affections and goods, to go and be Christ’s disciple, the further he goeth on in Christ’s love, the more shall he find who will rejoice to give him a place in their hearts, and to minister to him of their substance.\par
\switchcolumn*

\LA
\TeDeum{Te Deum}
\switchcolumn
\EN
\TeDeum{Te Deum}
\switchcolumn*

\end{paracol}

\par\needspace{10\baselineskip}\addvspace{1.2em}{\centering\small\textsc{Die 5 Julii} · \textsc{5 July}\par}
\Office{S. Antonii Mariæ Zaccaria Confessoris}{St. Anthony Mary Zaccaria, Confessor}{Duplex}
\addcontentsline{toc}{subsection}{Die 5 Julii · S. Antonii Mariæ Zaccaria Confessoris \textit{— St. Anthony Mary Zaccaria, Confessor}}
\begin{paracol}{2}
\LA
\RefLine{Lectiones I–III: de Scriptura occurrente.}
\switchcolumn
\EN
\RefLine{Lessons I–III: from the Scripture of the day.}
\switchcolumn*

\LA
\LessonHead{Lectio IV}
\switchcolumn
\EN
\LessonHead{Lesson IV}
\switchcolumn*

\LA
\noindent Antonius Maria Zaccaría, Cremonæ in Insubria nobili genere natus, jam a púero qua futurus esset sanctitáte porténdere visus est. Eximiárum enim in eo virtútum significatiónes mature eluxérunt, pietátis in Deum ac beátam Vírginem; insígnis præsertim in páuperes misericórdiæ; quorum inópiæ sublevandæ, vel pretiósa veste sibi detracta, haud semel præsto fuit. Humanióribus litteris in pátria excultus, Ticíni philosophíæ, Patavii medicinæ addiscendæ operam dedit; utque ómnibus vitæ integritate, ita et æquálibus acúmine ingenii facile antecelluit. Lauream adeptus ac domum reversus, ubi intelléxit se Dei mónitu ad animórum magis quam córporum morbis medéndum vocari, in sacras disciplínas percipiendas sédulo incúbuit. Interea ægrotos vísere, púeros christiana doctrina informare, júvenum cœtus pietáte excólere, ætate étiam provectos ad mores emendándos frequenter hortari non déstitit. Sacris initiatus, cum primo litáret, cælésti oborto lúmine, Angelórum coróna circúmdatus stupénti pópulo apparuísse tráditur. Exinde animárum salúti impensius consúlere, depravátis móribus summa ope obsístere curæ fuit. Ad hæc ádvenas, egénos, afflictos paterno complexus afféctu, piis alloquiis atque subsidiis recreatos ita solari, ut ejus domus miserórum perfúgium haberétur, ípseque pater pátriæ atque angelus merúerit a suis civibus appellari.\par
\switchcolumn
\EN
\noindent Anthony Mary Zaccaria was born (in 1502) of a noble family, at Cremona, on the Pau. Even in his childhood marks of his future holiness became manifest. There shone brightly in him, signs of excellent graces of childlike love toward God and the Blessed Virgin, and more especially of tenderness toward the poor, for the relief of whose needs he was ready more than once to strip off his own costly dress. He studied arts at his own home, philosophy at Ticino, and medicine at Padua, and as he excelled all others in goodness, so did he surpass all his companions in intellectual power. After taking his degree he returned home, and there understood from God that his call was to the healing of souls, rather than to that of bodies. He therefore began earnestly to study theology while he continued in the meantime to visit the sick, to teach Christian doctrine to children, to excite godliness among the young, and oftentimes even to exhort the aged to amend their ways. It is said that when he first said Mass after his ordination a light broke from heaven, and he seemed to the astonished bystanders to be surrounded by a circle of angels: from that time forth he laboured more earnestly for the salvation of souls, and the struggle against evil living. His fatherly love for strangers, for the needy, and for the afflicted, and the godly exhortations and alms wherewith he entertained them, made his house to become a refuge for the wretched, and earned for himself from his fellow -citizens the title of father of the fatherland and of angels.\par
\switchcolumn*

\LA
\begin{Resp}\Rsign Honéstum fecit illum Dóminus, et custodívit eum ab inimícis, et a seductóribus tutávit illum: \Star{} Et dedit illi claritátem ætérnam.\end{Resp}
\begin{Resp}\Vsign Justum dedúxit Dóminus per vias rectas, et osténdit illi regnum Dei.\end{Resp}
\begin{Resp}\Rsign Et dedit illi claritátem ætérnam.\end{Resp}
\switchcolumn
\EN
\begin{Resp}\Rsign The Lord made him honourable, and defended him from his enemies, and kept him safe from those that lay in wait for him; \Star{} And gave him perpetual glory.\end{Resp}
\begin{Resp}\Vsign He went down with him into the pit, and left him not in bonds.\end{Resp}
\begin{Resp}\Rsign And gave him perpetual glory.\end{Resp}
\switchcolumn*

\LA
\LessonHead{Lectio V}
\switchcolumn
\EN
\LessonHead{Lesson V}
\switchcolumn*

\LA
\noindent Mediolani, cum secum agitaret uberióres in rem christianam manare posse fructus, si in vinea Dómini sibi labórum socios adscísceret, re communicáta cum Bartholomæo Ferrario et Jacobo Morígia, nobilíssimis et sanctíssimis viris, sodalitátis Clericórum regulárium fundaménta jecit; quam, ob suum in Géntium Apostolum amórem, a sancto Paulo nuncupávit. Quæ, Cleménte septimo Pontifice máximo approbante et Paulo tertio confirmánte, brevi per complures regiónes propagáta est. Sanctimoniálium quoque Angelicárum socíetas ipsum Antónium Mariam parentem et auctórem hábuit. Qui tamen adeo de se submisse sentiebat, ut nullo pacto præésse suo ordini umquam volúerit. Tanta vero fuit patiéntia, ut formidolosíssimas tempestates in suos commótas constanti animo perferret; tanta caritate, ut piis adhortatiónibus religiósos viros ad Dei amórem inflammare, sacerdótes ad apostolicam vivéndi normam revocare, patrumque familias sodalítia ad bonam frugem institúere numquam intermíserit. Immo, intérdum præláta cruce per cómpita plateasque cum suis progréssus, férvida ac veheménti oratióne aberrántes improbosque hómines ad salútem redúceret.\par
\switchcolumn
\EN
\noindent While he was at Milan he bethought him that greater Christian good might be done if he gathered round him some fellow -labourers in the Lord's vineyard, and when he had conferred thereon with those noble and holy men Bartholomew Ferrari and James Morigia, he founded the brotherhood of Clerks Regulars, to whom on account of his own great love for the Apostle of the Gentiles he gave the name of Clerks of St. Paul. Under the approbation of the Supreme Pontiff Clement VII. and the confirmation of Paul III. this brotherhood was in a short time widely spread abroad. The Congregation of nuns who are called Angelicals also regard Anthony Mary as their Father and Founder. His own thought of himself was so lowly that he never would be at the head of his own Order. In great long-suffering he bore with patience the violent storms which were raised against his Institute. In the greatness of his charity he never ceased to enkindle the members of religious orders to love toward God, to exhort priests to live Apostolic lives, and to found guilds of married men, to the bringing forth of much fruit. Somewhiles he and his disciples would walk through the streets and squares with a Cross carried before them, and there by burning and vehement harangues call to salvation the wandering and the wicked.\par
\switchcolumn*

\LA
\begin{Resp}\Rsign Amávit eum Dóminus, et ornávit eum: stolam glóriæ índuit eum, \Star{} Et ad portas paradísi coronávit eum.\end{Resp}
\begin{Resp}\Vsign Índuit eum Dóminus lorícam fídei, et ornávit eum.\end{Resp}
\begin{Resp}\Rsign Et ad portas paradísi coronávit eum.\end{Resp}
\switchcolumn
\EN
\begin{Resp}\Rsign The Lord loved him and beautified him; He clothed him with a robe of glory, \Star{} And crowned him at the gates of Paradise.\end{Resp}
\begin{Resp}\Vsign The Lord hath put on him the breast-plate of faith, and hath adorned him.\end{Resp}
\begin{Resp}\Rsign And crowned him at the gates of Paradise.\end{Resp}
\switchcolumn*

\LA
\LessonHead{Lectio VI}
\switchcolumn
\EN
\LessonHead{Lesson VI}
\switchcolumn*

\LA
\noindent Illud étiam memorándum, quod, in Jesum crucifixum amóre flagrans, crucis mystérium ab ómnibus, ad statum æris campani indícium, sexta quaque feria sub vésperas, recoléndum curávit. Sanctíssimum Christi nomen in suis scriptis passim usurpábat et in ore semper habebat; ejusdémque cruciatus, vere Pauli discipulus, in corpore suo præ se ferebat. In sacram Eucharistiam singulari caritate ferebátur; cujus et frequenter percipiendæ consuetúdinem instaurávit, et morem e sublimi throno publice in triduum adorandæ invexísse perhibétur. Pudicítiam adeo coluit, ut étiam in exsangui corpore, revivíscere visus, ejus amórem testarétur. Accéssere cæléstia dona éxtasis, lacrimárum, futurórum eventuum cognitiónis, scrutatiónis cordium, virtútis in humani generis hostem. Tandem, magnis labóribus ubíque exantlátis, Guastallæ quo pacis sequester accítus fuerat, gravi morbo correptus est. Cremonam adductus, inter suórum fletus et complexus piíssimæ matris, quam proxime obituram prædixit, superna Apostolórum visióne recreatus, sodalitátis suæ increménta prænuntians, tertio Nonas Julii anno millesimo quingentésimo trigesimo nono, sanctíssime obiit, annos natus sex supra trigínta. Cultum tanto viro, ob exímiam ejus sanctitátem et signórum copiam a christiáno pópulo statim exhíbitum, Leo décimus tertius Pontifex maximus ratum hábuit et confirmávit; eumdémque anno millesimo octingentésimo nonagesimo septimo, in festo Ascensiónis Dominicæ, solemni ritu Sanctórum fastis adscripsit.\par
\switchcolumn
\EN
\noindent It is to be remembered that in his burning love for Jesus Crucified he reminded all men of the Mystery of the Cross by the sound of a bell every Friday evening, and himself as a true disciple of Paul always bore about in his body the dying of the Lord Jesus, (2 Cor. iv. 10.) The holy Name of Christ is found everywhere in his writings and was ever in his mouth. He was moved by a singular love toward the Holy Eu charist. He established a custom of receiving it often, and is said to have brought in the practice of exposing the same upon a lofty throne for three days' adoration. Of his earnest modesty the appearance of life which was seen even in his dead body seemed a witness. Together with all these things he possessed the gifts of trance, of tears, of knowledge of things to come, of reading the thoughts of the heart, and of power against the enemy of mankind. He was worn out with toil when he was seized with his last illness at Guastalla, whither he had been called as a peacemaker. He was carried to Cremona amid the tears of his brethren and the embraces of his devoted mother, whose imminent death he foretold. He was comforted by a vision of the Apostles above, and predicted the increase of his Brotherhood. On the th day of July, in the year 1539, he died a holy death at the age of thirty-six. Christians forthwith began to honour him for his eminent sanctity and the number of his signs and wonders, which honour the Supreme Pontiff Leo XIII. approved and confirmed, and on the Feast of the Lord's Ascension in the year 1897 solemnly enrolled his name among those of the Saints.\par
\switchcolumn*

\LA
\begin{Resp}\Rsign Iste homo perfécit ómnia quæ locútus est ei Deus, et dixit ad eum: Ingrédere in réquiem meam: \Star{} Quia te vidi justum coram me ex ómnibus géntibus.\end{Resp}
\begin{Resp}\Vsign Iste est, qui contémpsit vitam mundi, et pervénit ad cæléstia regna.\end{Resp}
\begin{Resp}\Rsign Quia te vidi justum coram me ex ómnibus géntibus.\end{Resp}
\begin{Resp}\Vsign Glória Patri, et Fílio, \Star{} et Spirítui Sancto.\end{Resp}
\begin{Resp}\Rsign Quia te vidi justum coram me ex ómnibus géntibus.\end{Resp}
\switchcolumn
\EN
\begin{Resp}\Rsign This is he which did according unto all that God commanded him; and God said unto him: Enter thou into My rest, \Star{} For thee have I seen righteous before Me among all people.\end{Resp}
\begin{Resp}\Vsign This is he which loved not his life in this world, and is come unto an everlasting kingdom.\end{Resp}
\begin{Resp}\Rsign For thee have I seen righteous before Me among all people.\end{Resp}
\begin{Resp}\Vsign Glory be to the Father, and to the Son, \Star{} and to the Holy Ghost.\end{Resp}
\begin{Resp}\Rsign For thee have I seen righteous before Me among all people.\end{Resp}
\switchcolumn*

\LA
\LessonHead{Lectio VII}
\switchcolumn
\EN
\LessonHead{Lesson VII}
\switchcolumn*

\LA
\LessonTitle{Léctio sancti Evangélii secúndum Marcum}
\switchcolumn
\EN
\LessonTitle{From the Holy Gospel according to Mark}
\switchcolumn*

\LA
\Cite{Marc 10:15-21}
\switchcolumn
\EN
\Cite{Mark 10:15-21}
\switchcolumn*

\LA
\noindent In illo témpore: Dixit Jesus discípulis suis: Quisquis non recéperit regnum Dei velut párvulus, non intrábit in illud. Et réliqua.\par
\switchcolumn
\EN
\noindent At that time: Jesus said unto His Disciples: Whosoever shall not receive the Kingdom of God as a little child: he shall not enter therein. And He took them up in His arms, put His hands upon them, and blessed them. And when He was gone forth into the way, there came one running, and kneeled to Him, and asked Him. And so on.\par
\switchcolumn*

\LA
\LessonTitle{Homilía sancti Augustíni Epíscopi}
\switchcolumn
\EN
\LessonTitle{Homily by St. Augustine, Bishop (of Hippo)}
\switchcolumn*

\LA
\Source{Sermo 47 de diversis}
\switchcolumn
\EN
\Source{Sermon 47: various}
\switchcolumn*

\LA
\noindent Durum vidétur et grave quod Dóminus imperávit, ut, si quis eum vult sequi ábneget seípsum; sed non est durum nec grave quod ille ímperat, qui ádjuvat ut fiat quod ímperat. Nam illud verum est quod ei dícitur in Psalmo: Propter verba labiórum tuórum ego custodívi vias duras. Quidquid enim durum est in præceptis, ut sis lene, caritas facit. Nóvimus quanta ipse amor facit. Quid autem est, Neget se? Non præsumat de se, séntiat se hóminem, et respiciat dictum propheticum: Maledictus omnis qui spem suam ponit in hómine: subdúcat se sibi sed non deorsum versus: subdúcat se sibi, ut hæreat Deo.\par
\switchcolumn
\EN
\noindent That which the Lord commandeth hard and heavy seemeth: (Come, take up the Cross, and follow Me; and as it is written elsewhere,) If any man will come after Me, let him deny himself, and take up his cross, and follow Me. (Matth. xvi, 24) But neither hard nor heavy is that which He commandeth when He that commandeth giveth help to fulfil. That is true which is said in Ps. xvi. 5, By the words of Thy lips I have kept me to strait paths. That which is hard in the commandment love maketh easy. How great is the power of love, we know. And what signifieth this, let him deny? Let him put no trust in himself, let him feel that he is man, and let him have regard unto that which was spoken of the prophet, (Jer. xvii, 5) saying, Cursed be the man that trusteth in men. Let him mistrust himself, but not to sink; let him mistrust himself, that he may cleave unto God.\par
\switchcolumn*

\LA
\begin{Resp}\Rsign Iste est qui ante Deum magnas virtútes operátus est, et de omni corde suo laudávit Dóminum: \Star{} Ipse intercédat pro peccátis ómnium populórum.\end{Resp}
\begin{Resp}\Vsign Ecce homo sine queréla, verus Dei cultor, ábstinens se ab omni ópere malo, et pérmanens in innocéntia sua.\end{Resp}
\begin{Resp}\Rsign Ipse intercédat pro peccátis ómnium populórum.\end{Resp}
\switchcolumn
\EN
\begin{Resp}\Rsign This is he which wrought great wonders before God, and praised the Lord with all his heart. \Star{} May he pray for all people, that their sins may be forgiven unto them.\end{Resp}
\begin{Resp}\Vsign Behold a man without blame, a worshipper of God in truth, keeping himself clean from every evil work, and abiding still in his innocency.\end{Resp}
\begin{Resp}\Rsign May he pray for all people, that their sins may be forgiven unto them.\end{Resp}
\switchcolumn*

\LA
\LessonHead{Lectio VIII}
\switchcolumn
\EN
\LessonHead{Lesson VIII}
\switchcolumn*

\LA
\noindent Quo sequendus est Dóminus? Quo iit, novimus; resurréxit enim et ascéndit in cælum: illo sequendus est. Plane desperándum non est, quia ipse promísit, non quia homo aliquid potest. Jam quare desperémus, si membra illíus cápitis sumus? Bonum est illo eum sequi; sed vidéndum est, qua. Etenim verba ista Dóminus Jesus non tunc dicebat quando a mórtuis jam resurrexerat; nondum erat passus, venturus erat ad crucem, venturus ad exhonoratiónem, ad contumelias, ad flagélla, ad spinas, ad vulnera, ad insultatiónes, oppróbria, mortem. Quasi exasperáta est via: pigrum te facis, non vis sequi: Séquere. Nam quis non velit ire ad exaltatiónem? omnes delectat celsitudo, sed humilitas gradus est.\par
\switchcolumn
\EN
\noindent Whither are we to follow the Lord? Whither He is gone, we know. He is risen from the dead, and is gone up into heaven. Thither we must follow Him. And we must not despair of so doing, not because man is able to do anything, but because He is faithful That promised. (Heb. xi, n.) Why, then, should we despair, since we are members of His Body, of His Flesh, and of His Bones, Who is the Head of the Church, and He is the Saviour of the body, (Eph. xxx, 23) Good is it to follow Him, but whither we are to follow Him we must see. When the Lord Jesus uttered those words bidding us to follow Him, He had not Himself as yet arisen from the dead, He had not as yet suffered, there lay still before Him the cross, shame, mockery, scourging, thorns, wounds, outrages, insults, death. After He uttered those words, His way became very rough. Art thou slothful? willest thou not to follow Him, but follow Him all the same, for who would not follow unto glory?. All men love exaltation, but lowliness is the step to rise withal.\par
\switchcolumn*

\LA
\begin{Resp}\Rsign Sint lumbi vestri præcíncti, et lucérnæ ardéntes in mánibus vestris: \Star{} Et vos símiles homínibus exspectántibus dóminum suum, quando revertátur a núptiis.\end{Resp}
\begin{Resp}\Vsign Vigiláte ergo, quia nescítis qua hora Dóminus vester ventúrus sit.\end{Resp}
\begin{Resp}\Rsign Et vos símiles homínibus exspectántibus dóminum suum, quando revertátur a núptiis.\end{Resp}
\begin{Resp}\Vsign Glória Patri, et Fílio, \Star{} et Spirítui Sancto.\end{Resp}
\begin{Resp}\Rsign Et vos símiles homínibus exspectántibus dóminum suum, quando revertátur a núptiis.\end{Resp}
\switchcolumn
\EN
\begin{Resp}\Rsign Let your loins be girded about, and your lights burning; \Star{} And ye yourselves like unto men that wait for their lord, when he will return from the wedding.\end{Resp}
\begin{Resp}\Vsign Watch therefore, for ye know not what hour your Lord doth come.\end{Resp}
\begin{Resp}\Rsign And ye yourselves like unto men that wait for their lord, when he will return from the wedding.\end{Resp}
\begin{Resp}\Vsign Glory be to the Father, and to the Son, \Star{} and to the Holy Ghost.\end{Resp}
\begin{Resp}\Rsign And ye yourselves like unto men that wait for their lord, when he will return from the wedding.\end{Resp}
\switchcolumn*

\LA
\LessonHead{Lectio IX}
\switchcolumn
\EN
\LessonHead{Lesson IX}
\switchcolumn*

\LA
\noindent Tolle crucem tuam, et séquere Dóminum. Crux enim nostra, quam Dóminus portari a nobis jubet, ut eum expeditíssimi sequamur, quid aliud quam mortalitátem carnis hujus significat? Ipsa enim nos cruciat, donec absorbeátur mors in victoriam. Crux ergo hæc ipsa crucifigenda est, et transfigenda est clavis timoris Dei, ne solutis et liberis membris reluctántem portare non possis; sequi enim Dóminum, nisi eam portans, omnino non vales. Nam quómodo eum sequéris, si non es ejus? Qui autem Jesu Christi sunt, ait Apóstolus, carnem suam crucifixérunt cum passiónibus et desidériis.\par
\switchcolumn
\EN
\noindent Take up thy cross, and follow the Lord, and the cross that the Lord commandeth us to carry after Him, that we may follow Him most speedfully, what is it but the death of this flesh? For it is this flesh that crucifieth us until death is swallowed up in victory, (1 Cor. xv, 55) Therefore must this our own cross itself be crucified and pierced with the nails of the fear of God, lest if it be free it hamper thee, in the carrying of it, and thou canst nowise follow the Lord save in carrying it. For how canst thou follow Him if thou be none of His. And they that are Christ's, saith the Apostle, have crucified the flesh with the affections and lusts, (Gal. v, 24)\par
\switchcolumn*

\LA
\TeDeum{Te Deum}
\switchcolumn
\EN
\TeDeum{Te Deum}
\switchcolumn*

\LA
\LessonHead{Lectio IX (pro commemoratione)}
\switchcolumn
\EN
\LessonHead{Lesson IX (when commemorated)}
\switchcolumn*

\LA
\LessonTitle{Commemoratio: S. Antonii Mariæ Zaccaria Confessoris}
\switchcolumn
\EN
\LessonTitle{Commemoration of: St. Anthony Mary Zaccaria, Confessor}
\switchcolumn*

\LA
\noindent Antónius María Zaccaría, Cremónæ nóbili génere natus, jam a púero morum pudicítia et misericórdia in páuperes elúxit. Humánis lítteris, philosophíæ ac medicínæ vacans, integritáte vitæ et ingénii acúmine æquálibus antecélluit. Dei mónitu, disciplínas sédulo excóluit; mox, sacerdótio auctus, talem se prǽbuit, ut pater pátriæ atque ángelus merúerit a suis cívibus appellári. Medioláni, cum Bartholomǽo Ferrário et Jacóbo Morígia, sanctíssimis viris, sodalitátem Clericórum regulárium, a sancto Paulo nuncupátam, et Sanctimoniálium Angelicárum societátem instítuit. Sacræ Eucharístiæ cultor assíduus, públicam sanctíssimi Sacraménti expositiónem mirífice promóvit. Cæléstibus donis a Deo ditátus, magnísque labóribus oppréssus, gravem morbum cum nactus esset, Cremónæ, tértio nonas júlii, anno millésimo quingentésimo trigésimo nono, sanctíssime óbiit. Leo Papa décimus tértius cultum ei exhíbitum, ratum hábuit et confirmávit, eúmque Sanctórum catálogo adscrípsit.\par
\switchcolumn
\EN
\noindent Born at Cremona of a noble family, Anthony Mary Zacharias even from his boyhood shone by his virtuous character and his mercy to the poor. During his education in the humanities, philosophy and medicine, he excelled his companions both in holiness of life and in keenness of mind. At a sign from God, he zealously cultivated the sacred sciences. After his ordination, the zeal of his priestly life soon earned for him the titles of Father and Angel of his country, bestowed on him by his fellow-citizens. At Milan, with the holy men Bartholomew Ferrári and James Morigia, he founded the Society of Clerks Regular named after St. Paul, and the society of nuns called the Angelicals. He was zealous in adoration of the Holy Eucharist and strongly promoted the public exposition of the most holy Sacrament. Enriched by God with heavenly gifts and worn out by his great labours, he contracted a serious illness, and he died a most holy death on the 5th day of July, 1539, at Cremona. Leo XIII approved and confirmed the cult already paid to him and enrolled him in the list of the Saints.\par
\switchcolumn*

\LA
\TeDeum{Te Deum}
\switchcolumn
\EN
\TeDeum{Te Deum}
\switchcolumn*

\end{paracol}

\par\needspace{10\baselineskip}\addvspace{1.2em}{\centering\small\textsc{Die 5 Julii} · \textsc{5 July}\par}
\Office{Septima die infra Octavam Ss. Apostolorum Petri et Pauli}{Seventh Day Within the Octave of the Apostles Sts. Peter and Paul}{Semiduplex}
\addcontentsline{toc}{subsection}{Die 5 Julii · Septima die infra Octavam Ss. Apostolorum Petri et Pauli \textit{— Seventh Day Within the Octave of the Apostles Sts. Peter and Paul}}
\begin{paracol}{2}
\LA
\LessonHead{Lectio IV}
\switchcolumn
\EN
\LessonHead{Lesson IV}
\switchcolumn*

\LA

\switchcolumn
\EN
\Note{No English translation in the source files; the Latin is given.}
\switchcolumn*

\LA
\LessonTitle{Commemoratio: Septima die infra Octavam Ss. Apostolorum Petri et Pauli}
\switchcolumn
\EN
\LessonTitle{Commemoration of: Seventh Day Within the Octave of the Apostles Sts. Peter and Paul}
\switchcolumn*

\LA
\noindent Sermo S. Maximi Epíscopi\par
\switchcolumn
\EN
\noindent Sermo S. Maximi Epíscopi\par
\switchcolumn*

\LA
Non sine causa factum putemus, fratres carissimi, quod gloriosissimi Christianæ fidei principes Petrus et Paulus, una die, uno loco, unius tyranni toleravere sententiam. Una die passi sunt, ut ad Christum pariter pervenirent: uno in loco, ne lateri Roma deesset: sub uno persecutore, ut æqualis crudelitas utrumque constringeret. Dies ergo, puto, pro merito, locus pro gloria, persecutor decretus est pro virtute. Et in quo tandem loco martyrium pertulerunt? In urbe Roma, quæ principatum et caput obtinet nationum; scilicet ut ubi caput superstitionis erat, illic caput quiesceret sanctitatis; et ubi gentilium principes habitabant, illic ecclesiarum principes morarentur.\par
\switchcolumn
\EN
Non sine causa factum putemus, fratres carissimi, quod gloriosissimi Christianæ fidei principes Petrus et Paulus, una die, uno loco, unius tyranni toleravere sententiam. Una die passi sunt, ut ad Christum pariter pervenirent: uno in loco, ne lateri Roma deesset: sub uno persecutore, ut æqualis crudelitas utrumque constringeret. Dies ergo, puto, pro merito, locus pro gloria, persecutor decretus est pro virtute. Et in quo tandem loco martyrium pertulerunt? In urbe Roma, quæ principatum et caput obtinet nationum; scilicet ut ubi caput superstitionis erat, illic caput quiesceret sanctitatis; et ubi gentilium principes habitabant, illic ecclesiarum principes morarentur.\par
\switchcolumn*

\LA
Cujus autem meriti sint beatissimi Petrus et Paulus, hinc possumus intelligere, quod cum Dominus Orientis regionem propria illustraverit passione; Occidentis plagam, ne quid minus esset, vice sui apostolorum sanguine illuminare dignatus est. Et licet illius passio nobis sufficiat ad salutem, tamen étiam horum martyrium nobis contulit ad exemplum. Hodierna igitur die beati apostoli sanguinem profuderunt: sed videamus causam quare isti perpessi sunt, scilicet quod inter cetera mirabilia étiam magum illum Simonem orationibus suis de aëris vacuo præcipiti ruina prostraverunt.\par
\switchcolumn
\EN
Cujus autem meriti sint beatissimi Petrus et Paulus, hinc possumus intelligere, quod cum Dominus Orientis regionem propria illustraverit passione; Occidentis plagam, ne quid minus esset, vice sui apostolorum sanguine illuminare dignatus est. Et licet illius passio nobis sufficiat ad salutem, tamen étiam horum martyrium nobis contulit ad exemplum. Hodierna igitur die beati apostoli sanguinem profuderunt: sed videamus causam quare isti perpessi sunt, scilicet quod inter cetera mirabilia étiam magum illum Simonem orationibus suis de aëris vacuo præcipiti ruina prostraverunt.\par
\switchcolumn*

\LA
Cum enim idem Simon se Christum diceret, et tamquam filium ad patrem asseret volando se posse conscendere, atque elatus subito magicis artibus volare cæpisset: tunc Petrus fixis genibus precatus est Dominum, et precatione sancta vicit magicam levitatem. Prior enim ascendit ad Dominum oratio, quam volatus, et ante pervenit justa petitio, quam iniqua præsumptio; ante Petrus in terris positus obtinuit quod petebat quam Simon perveniret in cælestibus, quo tendebat. Tunc igitur Petrus velut vinctum illum de sublimi ære deposuit, et quodam præcipitio in saxo elidens ejus crura confregit: et hoc in opprobrio facti illius, ut qui paulo ante volare tentaverat, subito ambulare non posset, et qui pennas assumpserat, plantas amitteret.\par
\switchcolumn
\EN
Cum enim idem Simon se Christum diceret, et tamquam filium ad patrem asseret volando se posse conscendere, atque elatus subito magicis artibus volare cæpisset: tunc Petrus fixis genibus precatus est Dominum, et precatione sancta vicit magicam levitatem. Prior enim ascendit ad Dominum oratio, quam volatus, et ante pervenit justa petitio, quam iniqua præsumptio; ante Petrus in terris positus obtinuit quod petebat quam Simon perveniret in cælestibus, quo tendebat. Tunc igitur Petrus velut vinctum illum de sublimi ære deposuit, et quodam præcipitio in saxo elidens ejus crura confregit: et hoc in opprobrio facti illius, ut qui paulo ante volare tentaverat, subito ambulare non posset, et qui pennas assumpserat, plantas amitteret.\par
\switchcolumn*

\LA
\TeDeum{Te Deum}
\switchcolumn
\EN
\TeDeum{Te Deum}
\switchcolumn*

\end{paracol}

\par\needspace{10\baselineskip}\addvspace{1.2em}{\centering\small\textsc{Die 6 Julii} · \textsc{6 July}\par}
\Office{In Octava Ss. Apostolorum Petri et Pauli}{Octave of Sts. Peter and Paul}{Duplex majus}
\addcontentsline{toc}{subsection}{Die 6 Julii · In Octava Ss. Apostolorum Petri et Pauli \textit{— Octave of Sts. Peter and Paul}}
\begin{paracol}{2}
\LA
\RefLine{Lectiones I–III: de Scriptura occurrente.}
\switchcolumn
\EN
\RefLine{Lessons I–III: from the Scripture of the day.}
\switchcolumn*

\LA
\LessonHead{Lectio IV}
\switchcolumn
\EN
\LessonHead{Lesson IV}
\switchcolumn*

\LA
\LessonTitle{Sermo sancti Joánnis Chrysóstomi}
\switchcolumn
\EN
\LessonTitle{From the Sermons of St. John Chrysostom, Patriarch of Constantinople}
\switchcolumn*

\LA
\Source{Apud Metaphrasten}
\switchcolumn
\EN
\Source{In the Metaphrastes}
\switchcolumn*

\LA
\noindent Quasnam vobis, o beáti Apóstoli, referémus grátias, qui tantum pro nobis laborastis? Mémini tui, Petre, et obstupesco: recordor tui, Paule, et excédens mente ópprimor lácrimis. Quid enim dicam, aut quid loquar, vestras contemplans afflictiónes, nescio. Quot carceres sanctificastis! quot catenas decorastis! quot torménta sustinuistis! quot maledicta tolerástis! quómodo Christum portastis! quómodo prædicatióne ecclésias lætificastis! Sunt benedicta vestræ linguæ instruménta: sánguine conspersa sunt membra vestra propter Ecclésiam. Vos Christum imitáti estis in omnibus. In omnem terram éxiit vester sonus, et verba vestra in fines orbis terræ.\par
\switchcolumn
\EN
\noindent The blessed Apostles, who have toiled so much for us, what thanks shall we give you When I remember thee, O Peter, I am lost in amazement Paul, when I think of thee, my heart overwhelmeth me, and I weep. When I look at your sufferings I know not what to say or what to speak. How many prisons have ye made holy How many fetters have ye made honourable? How many torments have ye endured How many reproaches have ye borne How have ye carried Christ How have ye made the Churches glad by your preaching? Verily, your tongues were blessed instruments it was for the Church's sake that your limbs were bloody. Ye have been made in all things followers of Christ. Your sound is gone out through all the earth, and your words to the ends of the world. (Ps. xviii. 4.)\par
\switchcolumn*

\LA
\begin{Resp}\Rsign Vidi conjúnctos viros, habéntes spléndidas vestes, et Ángelus Dómini locútus est ad me, dicens: \Star{} Isti sunt viri sancti facti amíci Dei.\end{Resp}
\begin{Resp}\Vsign Vidi Ángelum Dei fortem, volántem per médium cælum, voce magna clamántem et dicéntem.\end{Resp}
\begin{Resp}\Rsign Isti sunt viri sancti facti amíci Dei.\end{Resp}
\switchcolumn
\EN
\begin{Resp}\Rsign I saw men standing together, clad in shining raiment, and the Angel of the Lord spoke unto me, saying: \Star{} These men are holy, for they are the friends of God.\end{Resp}
\begin{Resp}\Vsign I saw a strong Angel of God fly into the midst of heaven, saying with a loud voice:\end{Resp}
\begin{Resp}\Rsign These men are holy, for they are the friends of God.\end{Resp}
\switchcolumn*

\LA
\LessonHead{Lectio V}
\switchcolumn
\EN
\LessonHead{Lesson V}
\switchcolumn*

\LA
\noindent Gaudeas, Petre, cui datum est ut ligno crucis Christi fruereris. Et ad Magistri quidem similitúdinem voluísti crucifigi, non recta quidem figura, ut Christus Dóminus; sed cápite in terram verso, tamquam qui a terra in cælum iter fáceres. Beáti clavi, qui sancta illa membra penetrarunt. Tu cum omni fiducia in manus Dómini ánimam tradidísti, qui assidue ei et ejus sponsæ Ecclésiæ servísti, qui fervénti spíritu Dóminum dilexísti, ómnium Apostolórum fidelíssimus.\par
\switchcolumn
\EN
\noindent Rejoice, O Peter, who hast been gladdened by the wood of the Cross of. Christ., It was. a showing forth of thy Teacher that thou didst will to be crucified, not like the Lord Christ, standing upright, but with thine head toward the earth, as one that made a way from earth to heaven. Blessed are the nails which pierced thine holy, limbs. With sure and certain hope didst thou commend thy spirit into the hands of the Lord, thou who hadst been a faithful servant to Him and to His Bride the Church, thou who in thy warm heart hadst loved the Lord more loyally than all the Apostles.\par
\switchcolumn*

\LA
\begin{Resp}\Rsign Beáti estis, cum maledíxerint vobis hómines, et persecúti vos fúerint, et díxerint omne malum advérsum vos, mentiéntes, propter me: \Star{} Gaudéte et exsultáte, quóniam merces vestra copiósa est in cælis.\end{Resp}
\begin{Resp}\Vsign Cum vos óderint hómines, et cum separáverint vos, et exprobráverint, et ejécerint nomen vestrum tamquam malum propter Fílium hóminis.\end{Resp}
\begin{Resp}\Rsign Gaudéte et exsultáte, quóniam merces vestra copiósa est in cælis.\end{Resp}
\switchcolumn
\EN
\begin{Resp}\Rsign Blessed are ye when men shall revile you, and persecute you, and shall say all manner of evil against you falsely, for My sake; \Star{} Rejoice, and be exceeding glad, for great is your reward in heaven.\end{Resp}
\begin{Resp}\Vsign When men shall hate you, and when they shall separate you from their company, and shall reproach you, and cast out your name as evil, for the Son of Man's sake.\end{Resp}
\begin{Resp}\Rsign Rejoice, and be exceeding glad, for great is your reward in heaven.\end{Resp}
\switchcolumn*

\LA
\LessonHead{Lectio VI}
\switchcolumn
\EN
\LessonHead{Lesson VI}
\switchcolumn*

\LA
\noindent Gaudeas et tu, beate Paule, cui caput fuit gládio amputátum, cujus virtútes nullis verbis explicari possunt. Quisnam gládius sanctum guttur tuum pervasit, Dominicum, inquam, instruméntum, quod a cælo habétur in admiratióne, et quod terra reverétur? Quisnam locus tuum sánguinem excepit, lactis spécie in ejus qui te percússit, túnica apparentem? qui ánimam illíus bárbari supra modum dulciórem reddens, fidélem effécit cum sociis. Sit mihi gládius ille pro coróna, et clavi Petri pro gemmis infixis in diadémate.\par
\switchcolumn
\EN
\noindent Rejoice thou also, O blessed Paul, whose head was cut off by the sword, thou whose fearless devotion no words can express. What sword was that which divided thine holy neck, that instrument of the Lord's work, worthy that heaven should wonder at it, and earth worship it What place was that which drank in thy blood, that appeared like drops of milk upon the raiment of him who smote thee, and made the savage and his comrades to become strangely gentle and faithful Would that I could have that sword for a crown, and the nails of Peter set therein as the jewels of the diadem.\par
\switchcolumn*

\LA
\begin{Resp}\Rsign Isti sunt triumphatóres et amíci Dei, qui contemnéntes jussa príncipum, meruérunt prǽmia ætérna: \Star{} Modo coronántur, et accípiunt palmam.\end{Resp}
\begin{Resp}\Vsign Isti sunt qui venérunt ex magna tribulatióne, et lavérunt stolas suas in sánguine Agni.\end{Resp}
\begin{Resp}\Rsign Modo coronántur, et accípiunt palmam.\end{Resp}
\begin{Resp}\Vsign Glória Patri, et Fílio, \Star{} et Spirítui Sancto.\end{Resp}
\begin{Resp}\Rsign Modo coronántur, et accípiunt palmam.\end{Resp}
\switchcolumn
\EN
\begin{Resp}\Rsign These are they which have conquered, and are become the friends of God, who recked not of the commandments of princes, and earned the everlasting reward. \Star{} And now have they crowns on their heads, and palms in their hands.\end{Resp}
\begin{Resp}\Vsign These are they which came out of great tribulation, and have washed their robes in the blood of the Lamb.\end{Resp}
\begin{Resp}\Rsign And now have they crowns on their heads, and palms in their hands.\end{Resp}
\begin{Resp}\Vsign Glory be to the Father, and to the Son, \Star{} and to the Holy Ghost.\end{Resp}
\begin{Resp}\Rsign And now have they crowns on their heads, and palms in their hands.\end{Resp}
\switchcolumn*

\LA
\LessonHead{Lectio VII}
\switchcolumn
\EN
\LessonHead{Lesson VII}
\switchcolumn*

\LA
\LessonTitle{Léctio sancti Evangélii secúndum Matthǽum}
\switchcolumn
\EN
\LessonTitle{From the Holy Gospel according to Matthew}
\switchcolumn*

\LA
\Cite{Matt 14:22-33}
\switchcolumn
\EN
\Cite{Matt 14:22-23}
\switchcolumn*

\LA
\noindent In illo témpore: Cómpulit Jesus discipulos ascendere in naviculam, et præcedere eum trans fretum, donec dimítteret turbas. Et réliqua.\par
\switchcolumn
\EN
\noindent At that time Jesus constrained His disciples to get into a ship, and to go before Him unto the other side, while He sent the multitudes away. And so on.\par
\switchcolumn*

\LA
\LessonTitle{Homilía sancti Hierónymi Presbýteri}
\switchcolumn
\EN
\LessonTitle{Homily by St. Jerome, Priest at Bethlehem}
\switchcolumn*

\LA
\Source{Liber 2 Comment. in cap. 14 Matth.}
\switchcolumn
\EN
\Source{Bk. ii, Comm. on Matth. xiv}
\switchcolumn*

\LA
\noindent Discípulis Dóminus præcepit transfretare, et cómpulit ut ascenderent naviculam. Quo sermóne osténditur, invitos eos a Dómino recessisse, dum amóre Præceptoris ne punctum quidem témporis ab eo volunt separari. Et, dimissa turba, ascéndit in montem solus orare. Si fuíssent cum eo discípuli Petrus et Jacobus et Joánnes, qui víderant glóriam transformáti, fórsitan ascendíssent in montem cum eo; sed turba ad sublimia sequi non potest, nisi docúerit eam juxta mare in littore et alúerit in desérto.\par
\switchcolumn
\EN
\noindent The Lord commanded His disciples to cross over to the other side, and constrained them to get into a ship. By these expressions we perceive that they were unwilling to leave the Lord, the love of their Teacher making them desire not to lose a moment of His company. And when He had sent the multitudes away, He went up into a mountain apart to pray. Perchance, if Peter, and James, and John, who had seen Him in the glory of the Transfiguration, had been with Him, they would have gone up into the mountain with Him, but the common herd could not follow Him, save when He taught them on the sea shore, or fed them in the wilderness.\par
\switchcolumn*

\LA
\begin{Resp}\Rsign Isti sunt qui vivéntes in carne, plantavérunt Ecclésiam sánguine suo: \Star{} Cálicem Dómini bibérunt, et amíci Dei facti sunt.\end{Resp}
\begin{Resp}\Vsign In omnem terram exívit sonus eórum, et in fines orbis terræ verba eórum.\end{Resp}
\begin{Resp}\Rsign Cálicem Dómini bibérunt, et amíci Dei facti sunt.\end{Resp}
\switchcolumn
\EN
\begin{Resp}\Rsign These are they who while yet they lived in the flesh, planted the Church in their own blood; \Star{} They drank of the Lord's cup, and became the friends of God.\end{Resp}
\begin{Resp}\Vsign Their sound is gone out through all the earth, and their words to the ends of the world.\end{Resp}
\begin{Resp}\Rsign They drank of the Lord's cup, and became the friends of God.\end{Resp}
\switchcolumn*

\LA
\LessonHead{Lectio VIII}
\switchcolumn
\EN
\LessonHead{Lesson VIII}
\switchcolumn*

\LA
\noindent Quod autem ascéndit solus orare, non ad eum réferas, qui de quinque pánibus quinque míllia saturávit hóminum, exceptis párvulis et muliéribus; sed ad eum, qui, audíta morte Joánnis, secessit in solitúdinem: non quod personam Dómini separemus, sed quod ópera ejus inter Deum hominémque divisa sint. Navícula autem in médio mari jactabátur fluctibus. Recte quasi invíti et retractántes Apóstoli a Dómino recesserant, ne, illo absente, naufragia sustinérent.\par
\switchcolumn
\EN
\noindent He went up into a mountain apart to pray, not as He Who, with five loaves and two fishes, had satisfied about five thousand men, besides women and children, but as He, Who when He heard of the death of John, departed into a desert place apart (vv. 12-21.) Not that we make two Persons in the Lord but some of His works He did as God, and some as man. But the ship was now in the midst of the sea, tossed with waves. The Apostles were right to be slow and unwilling to leave the Lord, for, when He was not with them, they were in peril of shipwreck.\par
\switchcolumn*

\LA
\begin{Resp}\Rsign Isti sunt viri sancti, quos elégit Dóminus in caritáte non ficta, et dedit illis glóriam sempitérnam: \Star{} Quorum doctrína fulget Ecclésia, ut sole luna.\end{Resp}
\begin{Resp}\Vsign Sancti per fidem vicérunt regna: operáti sunt justítiam.\end{Resp}
\begin{Resp}\Rsign Quorum doctrína fulget Ecclésia, ut sole luna.\end{Resp}
\begin{Resp}\Vsign Glória Patri, et Fílio, \Star{} et Spirítui Sancto.\end{Resp}
\begin{Resp}\Rsign Quorum doctrína fulget Ecclésia, ut sole luna.\end{Resp}
\switchcolumn
\EN
\begin{Resp}\Rsign These men are saints, whom the Lord hath chosen in love unfeigned, and hath given them glory everlasting. These are they \Star{} By the light of whose teaching the Church is glorified, even as the moon is glorified by the light of the sun.\end{Resp}
\begin{Resp}\Vsign The saints through faith subdued kingdoms, wrought righteousness.\end{Resp}
\begin{Resp}\Rsign By the light of whose teaching the Church is glorified, even as the moon is glorified by the light of the sun.\end{Resp}
\begin{Resp}\Vsign Glory be to the Father, and to the Son, \Star{} and to the Holy Ghost.\end{Resp}
\begin{Resp}\Rsign By the light of whose teaching the Church is glorified, even as the moon is glorified by the light of the sun.\end{Resp}
\switchcolumn*

\LA
\LessonHead{Lectio IX}
\switchcolumn
\EN
\LessonHead{Lesson IX}
\switchcolumn*

\LA
\noindent Denique, Dómino in montis cacúmine commorante, statim ventus contrarius óritur, et turbat mare, et periclitántur Apóstoli; et tamdiu imminens naufragium persevérat, quámdiu Jesus véniat. Quarta autem vigília noctis venit ad eos ámbulans supra mare. Statiónes et vigíliæ militares in terna horárum spatia dividúntur. Quando ergo dicit, quarta vigília noctis venisse ad eos Dóminum, osténdit tota nocte periclitatos; et extremo noctis, atque in consummatióne mundi, eis auxílium præbiturum.\par
\switchcolumn
\EN
\noindent Thilst the Lord abode alone upon the top of the mountain, a contrary wind arose, and the sea raged, and the Apostles were endangered and yet the threatening shipwreck held off until Jesus came. And in the fourth watch of the night, Jesus went unto them, walking on the sea. The watches of soldiers are divided into three. When therefore it is said that the Lord came unto them in the fourth watch, it appeareth that they had been in peril all night, and that it was at the end of the night, as it will again be at the end of the world, that He came to the rescue of His disciples.\par
\switchcolumn*

\LA
\TeDeum{Te Deum}
\switchcolumn
\EN
\TeDeum{Te Deum}
\switchcolumn*

\end{paracol}

\par\needspace{10\baselineskip}\addvspace{1.2em}{\centering\small\textsc{Die 7 Julii} · \textsc{7 July}\par}
\Office{Ss. Cyrilli et Methodii Pont. et Conf.}{Sts. Cyril and Methodius, Bishops and Confessors}{Duplex}
\addcontentsline{toc}{subsection}{Die 7 Julii · Ss. Cyrilli et Methodii Pont. et Conf. \textit{— Sts. Cyril and Methodius, Bishops and Confessors}}
\begin{paracol}{2}
\LA
\RefLine{Lectiones I–III: de Scriptura occurrente.}
\switchcolumn
\EN
\RefLine{Lessons I–III: from the Scripture of the day.}
\switchcolumn*

\LA
\LessonHead{Lectio IV}
\switchcolumn
\EN
\LessonHead{Lesson IV}
\switchcolumn*

\LA
\Source{Litteræ Encyclicæ Leónis Papæ XIII}
\switchcolumn
\EN
\Source{From the Encyclical Letter of Pope Leo XIII}
\switchcolumn*

\LA
\noindent Cyrillus et Methodius fratres germani, Thessalonícæ amplíssimo loco nati, Constantinopolim mature concessérunt, ut in ipsa urbe Oriéntis príncipe humanitátis artes addíscerent. Uterque plurimum brevi profecérunt; sed maxime Cyrillus, qui tantam scientiárum laudem adeptus est, ut singuláris honoris causa Philósophus appellarétur. Deínde mónachum ágere Methodius cœpit; Cyrillus autem dignus est hábitus, cui Theodora imperátrix, auctore Ignátio patriarcha, negotium daret erudiéndi ad fidem christianam Cházaros, trans Chersonésum incoléntes, quos, præceptis suis edoctos et Dei numine instinctos, multiplici superstitióne deleta, ad Jesum Christum adjunxit. Recénti Christianórum communitáte óptime constituta, Constantinopolim rediit álacer, atque in monastérium Polychrónis, quo se jam Methodius receperat, Cyrillus ipse secessit. Interim cum res trans Chersonésum próspere gestas ad Rastilláum Moraviæ príncipem fama detulísset, is de aliquot operariis evangelicis Constantinopoli arcesséndis cum imperatóre Michaéle tertio egit. Igitur Cyrillus et Methodius illi expeditióni destináti, et in Moraviam celebri lætítia excepti, animos christiánis institutiónibus tanta vi tamque operosa industria excolendos aggrediúntur, ut non longo intervallo ea gens nomen Jesu Christo libentíssime dederit. Ad eam rem non parum sciéntia váluit dictiónis Slavónicæ, quam Cyrillus ante percéperat, multumque potuérunt sacræ utriusque Testaménti litteræ, quas proprio pópuli sermóne reddiderat; nam Cyrillus et Methodius príncipes inveniéndi fuérunt ipsas litteras, quibus est sermo ipsórum Slavórum signátus et expréssus eáque de causa ejusdem sermónis auctores non immérito habentur.\par
\switchcolumn
\EN
\noindent The brethren Cyril and Methodius were born in an honourable position at Thessalonica. As they advanced in years they went to Constantinople to study letters in the capital of the Eastern world. Both made quick progress, but most chiefly Cyril, who gained such learning that he was called for excellency the Philosopher. Methodius became a monk, but the Empress Theodora, on the recommendation of the Patriarch Ignatius, deemed Cyril worthy of receiving the task of teaching Christianity to the Khazar who dwelt beyond the Crimea. By the grace of God he so taught them that they laid aside their many superstitions and were joined to Jesus Christ. After properly establishing the new community of Christians Cyril hastened back to Constantinople, where he entered the monastery of Polychron, whither Methodius had already withdrawn himself. Rastilaw, Prince of Moravia, having heard tell of the good deeds beyond the Crimea, sent to Constantinople to the Emperor Michael III to obtain some Gospel labourers. Cyril and Methodius were sent to him, and gladly received in Moravia, and applied themselves with such power and industry to the work of Christianising souls that it was not long before that nation also joyfully submitted to Jesus Christ. To this end Cyril found of great use the knowledge of the Slavonic language, which he had already acquired, and much effect was produced by the translation of holy Scripture which he made into the language of the people. Cyril and Methodius were the inventors of the alphabet in which the language of the Slavs is characteristically expressed, and for this reason they have been not unjustly termed the fathers of Slavonic literature.\par
\switchcolumn*

\LA
\begin{Resp}\Rsign Invéni David servum meum, óleo sancto meo unxi eum: \Star{} Manus enim mea auxiliábitur ei.\end{Resp}
\begin{Resp}\Vsign Nihil profíciet inimícus in eo, et fílius iniquitátis non nocébit ei.\end{Resp}
\begin{Resp}\Rsign Manus enim mea auxiliábitur ei.\end{Resp}
\switchcolumn
\EN
\begin{Resp}\Rsign I have found David My servant, with My holy oil have I anointed him \Star{} For My hand shall help him.\end{Resp}
\begin{Resp}\Vsign The enemy shall prevail nothing against him, nor the son of wickedness afflict him.\end{Resp}
\begin{Resp}\Rsign For My hand shall help him.\end{Resp}
\switchcolumn*

\LA
\LessonHead{Lectio V}
\switchcolumn
\EN
\LessonHead{Lesson V}
\switchcolumn*

\LA
\noindent Cum rerum gestárum glóriam secundus rumor Romam nuntiásset, sanctus Nicoláus primus Pontifex maximus fratres optimos Romam conténdere jussit. Illi Romanum iter ingressi, relíquias sancti Clementis primi Pontificis maximi, quas Cyrillus Chersónæ repérerat, secum ádvehunt. Quo nuntio, Hadrianus secundus, qui Nicoláo demortuo fuerat suffectus, clero populóque comitante, obviam eis magna cum honóre significatióne progréditur. Deínde Cyrillus et Methodius de munere apostolico in quo essent sancte laborioséque versáti ad Pontificem maximum, assidénte clero, réferunt. Cum autem eo nómine ab ínvidis accusaréntur, quod sermónem Slavónicum in perfunctióne munerum sacrórum usurpavissent, causam dixére ratiónibus tam certis tamque illustribus, ut Pontifex et clerus et laudarint hómines et probarint. Tum ambo, juráti se in fide beáti Petri et Pontificum Romanórum permansuros, epíscopi ab Hadriáno consecráti sunt. Sed erat provisum divinitus, ut Cyrillus vitæ cursum Romæ cónderet, virtúte magis quam ætate matúrus. Itaque defuncti corpus, elátum fúnere publico, in ipso sepúlcro quod sibi Hadrianus exstruxerat compósitum fuit; tum ad sancti Clementis deductum, et hujus prope cíneres cónditum. Cumque veherétur per urbem inter festos Psalmórum cantus, non tam funeris quam triumphi pompa, visus est pópulus Romanus libaménta honórum cæléstium viro sanctíssimo detulisse. Methodius vero in Moraviam regréssus, ibique factus forma gregis ex animo, rei catholicæ inservire majore in dies studio ínstitit. Quin étiam Pannonios, Búlgaros, Dálmatas in fide christiáni nóminis confirmávit; in Carinthiis autem ad uníus veri Dei cultum traducéndis plurimum elaborávit.\par
\switchcolumn
\EN
\noindent Then the happy tidings of what they had done reached Rome, the Supreme Pontiff the holy Nicholas I. commanded these excellent brethren to come to Rome. When they started for Rome they brought with them the relics of the supreme Pontiff the holy Clement I which Cyril had discovered at Cherson. On hearing of their approach Adrian II, who had succeeded to the Papacy upon the death of Nicholas, went forth to meet them accompanied by the clergy and people with every sign of honour. Then Cyril and Methodius gave to the Supreme Pontiff in the presence of the clergy an account of the Apostolic office which they had discharged in so holy and toilsome a manner. When it was made blame to them by some enviers that they had used the Slavonic language for the purposes of public worship, they stated their reasons with such clearness and force that the Pontiff and clergy praised and approved them. When they had both taken an oath that they would remain in the faith of blessed Peter and of the Roman Pontiffs, they were consecrated bishops by Adrian, but it was the Will of God that Cyril, old in grace rather than in years, should close his life at Rome. His dead body received a public funeral, and was laid in the tomb which Adrian had built for himself, but it was afterwards brought to St. Clement's and buried hard by the ashes of that martyr. As it was carried through the city with joyful psalm-singing, it seemed as though the procession were rather that of a triumph than that of a funeral, and that the Roman people were offering heavenly honour to some eminent saint. Methodius went back to Moravia, and there became from his whole soul a pattern to his flock, and from day to day more zealous in the service of Catholicism. He confirmed the Pannonians, the Bulgarians, and the Dalmatians in the Christian religion, and laboured much to bring the Corinthians to the worship of the one true God.\par
\switchcolumn*

\LA
\begin{Resp}\Rsign Pósui adjutórium super poténtem, et exaltávi eléctum de plebe mea: \Star{} Manus enim mea auxiliábitur ei.\end{Resp}
\begin{Resp}\Vsign Invéni David servum meum, óleo sancto meo unxi eum.\end{Resp}
\begin{Resp}\Rsign Manus enim mea auxiliábitur ei.\end{Resp}
\switchcolumn
\EN
\begin{Resp}\Rsign I have laid help upon one that is mighty, and have exalted one chosen out of My people \Star{} For My hand shall help him.\end{Resp}
\begin{Resp}\Vsign I have found David My servant, with My holy oil have I anointed him.\end{Resp}
\begin{Resp}\Rsign For My hand shall help him.\end{Resp}
\switchcolumn*

\LA
\LessonHead{Lectio VI}
\switchcolumn
\EN
\LessonHead{Lesson VI}
\switchcolumn*

\LA
\noindent Apud Joánnem octavum, qui Hadriáno successerat, íterum de suspecta fide vialatoque more majórum accusatus, ac Romam veníre jussus, coram Joánne et epíscopis aliquot cleroque urbano, facile vicit catholicam prorsus fidem et se retinuisse constanter, et ceteros diligenter edocuísse: quod vero ad linguam Slavonicam in sacris peragéndis usurpatam, se certis de causis ex venia Hadriáni Pontificis, nec sacris Litteris repugnántibus, jure fecísse. Quapropter in re præsénti complexus Methódium Pontifex, potestátem ejus archiepiscopálem expeditionémque Slavónicam, datis étiam litteris, ratam esse jussit. Quare Methodius in Moraviam reversus assignátum sibi munus explere vigilantius perseverávit, pro quo et exsílium libenter passus est. Bohemórum príncipem ejusque uxórem ad fidem perdúxit, et in ea gente christianum nomen longe lateque vulgávit. Evangélii lumen in Polóniam invéxit, et, ut nonnulli scriptores tradunt, sede episcopali Leópoli fundata, in Moscóviam proprii nóminis digréssus, thronum pontificalem Kiowensem constítuit. Demum in Moraviam reversus est ad suos; jamque sese ábripi ad humánum éxitum sentiens, ípsemet sibi successórem designávit, clerumque et pópulum supremis præceptis ad virtútem cohortatus, ea vita, quæ sibi via in cælum fuit, placidíssime defunctus est. Uti Cyrillum Roma, sic Methódium Moravia decedentem summo honóre prosecuta est. Illórum vero festum, quod apud Slavoniæ pópulos jamdiu celebrari consueverat, Leo décimus tertius Pontifex maximus cum Officio ac Missa propria in univérsa Ecclésia quotannis agi præcépit.\par
\switchcolumn
\EN
\noindent Methodius was again accused before John VIII, the successor of Adrian, of unsoundness in faith, and transgression of the traditions of the elders he was summoned to Rome, and there easily proved, in the presence of John and of some Bishops and clergy of the city, that he had himself always firmly held the Catholic faith, and had carefully taught it to others, and that as regarded the use of the Slavonic language for public worship, he had acted lawfully from certain reasons, and the permission of Pope Adrian, and in nowise contrary to holy writ. The Pontiff therefore in this matter concurred with Methodius, and confirmed even in writing his archepiscopal authority and his mission among the Slavs. Methodius therefore went back to Moravia and resumed more earnestly than before the task committed to him, for the which also he cheerfully suffered exile. He converted the Prince of the Bohemians and his wife, and spread the Christian name far and wide among that people. He carried the light of the Gospel into Poland, and according to some writers, after establishing the see of Leopolis, went into Muscovy properly so called and established the see of Kieff. At the last he returned into Moravia, and when he felt that he was about to go the way of all flesh he named his own successor, exhorted the clergy and people for the last time to good living, and then calmly departed that life which had been to him a path to heaven. As Rome had honoured Cyril in his death, so did Moravia honour Methodius. The festival of these Saints, which had long been observed among the Slav nations, the Supreme Pontiff Leo XIII. ordered to be kept throughout the Universal Church with a special office and Mass.\par
\switchcolumn*

\LA
\begin{Resp}\Rsign Iste est, qui ante Deum magnas virtútes operátus est, et omnis terra doctrína ejus repléta est: \Star{} Ipse intercédat pro peccátis ómnium populórum.\end{Resp}
\begin{Resp}\Vsign Iste est, qui contémpsit vitam mundi, et pervénit ad cæléstia regna.\end{Resp}
\begin{Resp}\Rsign Ipse intercédat pro peccátis ómnium populórum.\end{Resp}
\begin{Resp}\Vsign Glória Patri, et Fílio, \Star{} et Spirítui Sancto.\end{Resp}
\begin{Resp}\Rsign Ipse intercédat pro peccátis ómnium populórum.\end{Resp}
\switchcolumn
\EN
\begin{Resp}\Rsign This is he which wrought great wonders before God, and the whole earth is full of his teaching \Star{} May he pray for all people, that their sins may be forgiven unto them\end{Resp}
\begin{Resp}\Vsign This is he which loved not his life in this world, and hath attained unto the kingdom of heaven.\end{Resp}
\begin{Resp}\Rsign May he pray for all people, that their sins may be forgiven unto them\end{Resp}
\begin{Resp}\Vsign Glory be to the Father, and to the Son, \Star{} and to the Holy Ghost.\end{Resp}
\begin{Resp}\Rsign May he pray for all people, that their sins may be forgiven unto them\end{Resp}
\switchcolumn*

\LA
\LessonHead{Lectio VII}
\switchcolumn
\EN
\LessonHead{Lesson VII}
\switchcolumn*

\LA
\LessonTitle{Léctio sancti Evangélii secúndum Lucam}
\switchcolumn
\EN
\LessonTitle{From the Holy Gospel according to Luke}
\switchcolumn*

\LA
\Cite{Luc 10:1-9}
\switchcolumn
\EN
\Cite{Luke 10:1-9}
\switchcolumn*

\LA
\noindent In illo témpore: Designávit Dóminus et álios septuagínta duos: et misit illos binos ante fáciem suam, in omnem civitátem et locum, quo erat ipse ventúrus. Et réliqua.\par
\switchcolumn
\EN
\noindent At that time, the Lord appointed other seventy-two also, and sent them two and two before His face into every city and place, whither He Himself would come. And so on.\par
\switchcolumn*

\LA
\LessonTitle{Homilía sancti Gregórii Papæ}
\switchcolumn
\EN
\LessonTitle{Homily by Pope St. Gregory (the Great.)}
\switchcolumn*

\LA
\Source{Homilia 17 in Evangelia}
\switchcolumn
\EN
\Source{17th Homily on the Gospels}
\switchcolumn*

\LA
\noindent Dóminus et Salvátor noster, fratres caríssimi, aliquándo nos sermónibus, aliquándo vero opéribus ádmonet. Ipsa étenim facta ejus præcépta sunt: quia dum áliquid tácitus facit, quid ágere debeámus innotéscit. Ecce enim binos in prædicatiónem discípulos mittit: quia duo sunt præcépta caritátis, Dei vidélicet amor, et próximi: et minus quam inter duos cáritas habéri non potest. Nemo enim próprie ad semetípsum habére caritátem dícitur: sed diléctio in álterum tendit, ut cáritas esse possit.\par
\switchcolumn
\EN
\noindent Dearly beloved brethren, our Lord and Saviour doth sometimes admonish us by words, and sometimes by works. Yea, His very works do themselves teach us for that which He doth silently His example still moveth us to copy. Behold how He sendeth forth His disciples to preach by two and two since there are two commandments to love, that is, a commandment to love God, and a commandment to love our neighbour and where there are not two, the one, being alone, hath not whereon to do the Lord's commandment. And no man can properly be said to love himself: for love tendeth outward toward our neighbour, if it be the love whereto the Gospel doth oblige us.\par
\switchcolumn*

\LA
\begin{Resp}\Rsign Amávit eum Dóminus, et ornávit eum: stolam glóriæ índuit eum, \Star{} Et ad portas paradísi coronávit eum.\end{Resp}
\begin{Resp}\Vsign Induit eum Dóminus lorícam fídei, et ornávit eum.\end{Resp}
\begin{Resp}\Rsign Et ad portas paradísi coronávit eum.\end{Resp}
\switchcolumn
\EN
\begin{Resp}\Rsign The Lord loved him and beautified him He clothed him with a robe of glory \Star{} And crowned him at the gates of Paradise.\end{Resp}
\begin{Resp}\Vsign The Lord hath put on him the breast-plate of faith, and hath adorned him.\end{Resp}
\begin{Resp}\Rsign And crowned him at the gates of Paradise.\end{Resp}
\switchcolumn*

\LA
\LessonHead{Lectio VIII}
\switchcolumn
\EN
\LessonHead{Lesson VIII}
\switchcolumn*

\LA
\noindent Ecce enim binos ad prædicándum discípulos Dóminus mittit: quátenus hoc nobis tácitus ínnuat, quia qui caritátem erga álterum non habet, prædicatiónis offícium suscípere nullátenus debet. Bene autem dícitur, quia misit eos ante fáciem suam in omnem civitátem et locum, quo erat ipse ventúrus. Prædicatóres enim suos Dóminus séquitur: quia prædicátio prǽvenit, et tunc ad mentis nostræ habitáculum Dóminus venit, quando verba exhortatiónis præcúrrunt: atque per hoc véritas in mente suscípitur.\par
\switchcolumn
\EN
\noindent Behold, the Lord sendeth forth His disciples to preach by two and two and thus doing, He doth silently teach us that whosoever loveth not his neighbour, such a one it behoveth not to take upon him the office of a preacher. Well also is it said that He sent them before His face into every city and place whither He Himself would come. The Lord followeth His preachers first cometh preaching, and then the Lord Himself cometh to the house of our mind, whither the word of exhortation hath come before and so cometh the truth into our mind.\par
\switchcolumn*

\LA
\begin{Resp}\Rsign Sint lumbi vestri præcíncti, et lucérnæ ardéntes in mánibus vestris: \Star{} Et vos símiles homínibus exspectántibus dóminum suum, quando revertátur a núptiis.\end{Resp}
\begin{Resp}\Vsign Vigiláte ergo, quia nescítis qua hora Dóminus vester ventúrus sit.\end{Resp}
\begin{Resp}\Rsign Et vos símiles homínibus exspectántibus dóminum suum, quando revertátur a núptiis.\end{Resp}
\begin{Resp}\Vsign Glória Patri, et Fílio, \Star{} et Spirítui Sancto.\end{Resp}
\begin{Resp}\Rsign Et vos símiles homínibus exspectántibus dóminum suum, quando revertátur a núptiis.\end{Resp}
\switchcolumn
\EN
\begin{Resp}\Rsign Let your loins be girded about, and your lights burning; \Star{} And ye yourselves like unto men that wait for their lord, when he will return from the wedding.\end{Resp}
\begin{Resp}\Vsign Watch therefore, for ye know not what hour your Lord doth come.\end{Resp}
\begin{Resp}\Rsign And ye yourselves like unto men that wait for their lord, when he will return from the wedding.\end{Resp}
\begin{Resp}\Vsign Glory be to the Father, and to the Son, \Star{} and to the Holy Ghost.\end{Resp}
\begin{Resp}\Rsign And ye yourselves like unto men that wait for their lord, when he will return from the wedding.\end{Resp}
\switchcolumn*

\LA
\LessonHead{Lectio IX}
\switchcolumn
\EN
\LessonHead{Lesson IX}
\switchcolumn*

\LA
\noindent Hinc namque eísdem prædicatóribus Isaías dicit: Paráte viam Dómini, rectas fácite sémitas Dei nostri. Hinc fíliis Psalmísta ait: Iter fácite ei, qui ascéndit super occásum. Super occásum namque Dóminus ascéndit: quia unde in passióne occúbuit, inde majórem suam glóriam resurgéndo manifestávit. Super occásum vidélicet ascéndit; quia mortem quam pértulit, resurgéndo calcávit. Ei ergo qui ascéndit super occásum, iter fácimus, cum nos ejus glóriam vestris méntibus prædicámus, ut eas et ipse post véniens, per amóris sui præséntiam illústret.\par
\switchcolumn
\EN
\noindent Therefore, to preachers saith Isaiah: “Prepare ye the way of the Lord, make straight an highway for our God.” (xl. 3.) And again the Psalmist saith: “Spread a path before him that rideth upon the West.” (lxvii. 4.) The Lord rideth upon the West above that from which in death He veiled His glory, hath He royally exalted that glory that excelleth, even the glory of His rising again. He rideth upon the West, Who, being risen again from the dead, is throned high above the death to which He bowed. Before Him, therefore, That rideth upon the West, we spread a path, when we set forth His glory before the eyes of your mind, to the end that He Himself may come after, and Himself enlighten the same your minds by His presence and His love.\par
\switchcolumn*

\LA
\TeDeum{Te Deum}
\switchcolumn
\EN
\TeDeum{Te Deum}
\switchcolumn*

\LA
\LessonHead{Lectio IX (pro commemoratione)}
\switchcolumn
\EN
\LessonHead{Lesson IX (when commemorated)}
\switchcolumn*

\LA
\LessonTitle{Commemoratio: Ss. Cyrilli et Methodii Pont. et Conf.}
\switchcolumn
\EN
\LessonTitle{Commemoration of: Sts. Cyril and Methodius, Bishops and Confessors}
\switchcolumn*

\LA
\noindent Cyríllus et Methódius, fratres germáni, Thessalonícæ amplíssimo loco nati, ab imperatóre Michaéle tértio destináti in Moráviam, brevi eam gentem ad fidem Christi adduxérunt. Cum rerum gestárum glóriam secúndus rumor Romam nuntiásset, sanctus Nicoláus primus Póntifex máximus fratres óptimos, Romam conténdere jussit; ubi ab Hadriáno, Nicolái successóre, epíscopi sunt consecráti. Sed, cum brevi Cyríllus Romæ piíssime obiísset, Methódius in Moráviam regréssus, rei cathólicæ inservíre impénsius ínstitit. Quin étiam Bohémos, Pannónios, Búlgaros, Dálmatas in fide christiána confirmávit; in Carínthiis autem ad uníus veri Dei cultum traducéndis valde elaborávit. Item Evangélii lumen in Polóniam invéxit, et, ut nonnúlli scriptóres tradunt, sede episcopáli Leópoli fundáta, in Moscóviam próprii nóminis digréssus, thronum pontificálem Kiowénsem constítuit. Demum in Moráviam revérsus, clerum et pópulum suprémis præcéptis ad virtútem cohortátus, placidíssime defúnctus est. Cyrílli et Methódii festum, apud Slavóniæ pópulos jámdiu celebrátum, Leo décimus tértius ad univérsam Ecclésiam exténdit.\par
\switchcolumn
\EN
\noindent Cyril and Methodius were brothers, born of a distinguished family in Thessalonica. The Emperor Michael III sent them into Moravia, where in a short time they brought the nation to the faith of Christ. When a favourable report of what they had done was brought to Rome, Pope St. Nicholas I ordered the brothers to come there. At Rome they were consecrated bishops by Adrian, Nicholas's successor. A short time later, however, Cyril died a most holy death in Rome, and Methodius went back to Moravia and increased his efforts on behalf of Catholicism. Moreover, he confirmed the Bohemians, the Pannonians, the Bulgarians and the Dalmatians in the Christian faith, and worked hard to bring the Corinthians to the worship of the one true God. He also brought the light of the Gospel to Poland and, as some writers say, founded the bishopric of Lemberg. Then he went to Muscovy properly so called and established the pontifical see of Kiev. At length he came back to Moravia, exhorted the clergy and people to virtue with his last words, and died peacefully. The feast day of Cyril and Methodius, already celebrated by the Slavic peoples, was extended to the universal Church by Leo XIII.\par
\switchcolumn*

\LA
\TeDeum{Te Deum}
\switchcolumn
\EN
\TeDeum{Te Deum}
\switchcolumn*

\end{paracol}

\par\needspace{10\baselineskip}\addvspace{1.2em}{\centering\small\textsc{Die 8 Julii} · \textsc{8 July}\par}
\Office{S. Elisabeth Reg. Portugaliæ Viduæ}{St. Elizabeth of Portugal, Queen and Widow}{Semiduplex}
\addcontentsline{toc}{subsection}{Die 8 Julii · S. Elisabeth Reg. Portugaliæ Viduæ \textit{— St. Elizabeth of Portugal, Queen and Widow}}
\begin{paracol}{2}
\LA
\RefLine{Lectiones I–III: de Scriptura occurrente.}
\switchcolumn
\EN
\RefLine{Lessons I–III: from the Scripture of the day.}
\switchcolumn*

\LA
\RefLine{Lectiones VII–IX: ex Commúni non Virginum non Martyrum (C7a).}
\switchcolumn
\EN
\RefLine{Lessons VII–IX: from the Common of Widows (C7a).}
\switchcolumn*

\LA
\LessonHead{Lectio IV}
\switchcolumn
\EN
\LessonHead{Lesson IV}
\switchcolumn*

\LA
\noindent Elísabeth Aragoniæ régibus ortam, Christi anno millesimo ducentésimo septuagesimo primo, in præsagium futuræ sanctimoniæ, paréntes, præter morem relícto matris aviæque nómine, a magna ejus matértera, Thuringiæ domina, sancta Elísabeth, in baptismo nominátam voluére. Ubi nata est, statim pátuit quam felix regum regnorúmque esset futura pacatrix; natalícia enim ejus lætítia perniciosas ávi patrisque dissensiónes in concordiam convértit. Pater vero, crescentis póstea fíliæ admirátus índolem, affirmábat fore, ut una Elísabeth réliquas Aragoniórum regum sánguine creatas feminas virtúte longe superáret. Sic cælestem ipsíus vitam in contemnéndo corporis ornatu, in fugiéndis voluptátibus, in jejuniis frequentandis, in divinis precibus assidue recitandis, in caritátis opéribus exercendis, veneratus, rerum suárum regnique felicitátem uníus fíliæ meritis referebat acceptam. Tandem, ubíque nota et a multis princípibus exoptata, Dionysio Lusitaniæ regi christiánis cæremoniis rite est in matrimónium collocata.\par
\switchcolumn
\EN
\noindent Elizabeth, daughter of Peter III., King of Aragon, was born in the year of Christ 1271, and it was an omen of her saintly life; that her father and mother, contrary to the usual custom, caused her to be baptized, not by the name of her mother or grandmother, but by that of her mother's aunt, the holy Lady Elizabeth of Thuringia. As soon as ever she was born, her destiny of being a peacemaker between kings and kingdoms began to appear, for the joy of her birth put an end to the ruinous quarrels of her father and grandfather. As she grew up, her father, delighted with her disposition, was used to foretell that his Elizabeth would in herself excel all the daughters of the kingly house of Aragon, and that the happiness of his own home and kingdom was all owing to this one damsel, whose heavenly life he venerated for her indifference to bodily finery, her abstinence from pleasures, her many fasts, her instancy in prayer to God, and her activity in doing works of charity. This illustrious maiden was sought in marriage by many princes, and at twelve years of age was wedded with Christian rites to Denis, King of Portugal.\par
\switchcolumn*

\LA
\begin{Resp}\Rsign Propter veritátem, et mansuetúdinem, et justítiam: \Star{} Et dedúcet te mirabíliter déxtera tua.\end{Resp}
\begin{Resp}\Vsign Spécie tua et pulchritúdine tua inténde, próspere procéde, et regna.\end{Resp}
\begin{Resp}\Rsign Et dedúcet te mirabíliter déxtera tua.\end{Resp}
\switchcolumn
\EN
\begin{Resp}\Rsign Because of truth, and meekness, and righteousness; \Star{} And thy right hand shall lead thee wonderfully.\end{Resp}
\begin{Resp}\Vsign In thy comeliness, and thy beauty, go forward, fare prosperously, and reign.\end{Resp}
\begin{Resp}\Rsign And thy right hand shall lead thee wonderfully.\end{Resp}
\switchcolumn*

\LA
\LessonHead{Lectio V}
\switchcolumn
\EN
\LessonHead{Lesson V}
\switchcolumn*

\LA
\noindent Juncta conjugio, non minórem excoléndis virtútibus quam liberis educandis operam dabat, viro placere studens, sed magis Deo. Mediam fere anni partem solo pane tolerábat et aqua; quæ in quodam ipsíus morbo divinitus versa est in vinum, cum id a médicis præscriptum bíbere recusasset. Páuperis feminæ ulcus horréndum exosculata, derepénte sanávit. Pecunias paupéribus distribuendas, ut regem latérent, hiberno témpore in rosas convértit. Vírginem cæcam a nativitáte illuminávit; multos alios solo crucis signo a gravíssimis morbis liberávit; plurima id genus miracula patrávit. Monasteria, collegia et templa non modo exstruxit, sed étiam magnifice dotávit. In regum discordiis componéndis admirábilis fuit; in privátis publicisque mortalium sublevandis calamitátibus indefessa.\par
\switchcolumn
\EN
\noindent As a wife, she gave herself up as much to the education of her children, as to her own improvement, striving in all ways, next to God, to please her husband. For nearly half the year, she was used to live on bread and water, and once, when she was ill, God changed the water into wine, which the physicians had ordered her to drink, but which she was unwilling to take. Once when she kissed a disgusting ulcer in a poor woman, it was immediately healed. One winter-time when she was giving some money to the poor, and was fain her husband should not see her alms, the coins changed into roses. She gave sight to a maiden who had been born blind, and healed many other persons of grievous sicknesses by the Sign of the Cross. The miracles of this kind, which she worked, were many. She not only built, but richly endowed convents, schools, and churches. She had a wonderful skill in making peace between kings, and toiled unweariedly to lighten all suffering, whether public or private.\par
\switchcolumn*

\LA
\begin{Resp}\Rsign Dilexísti justítiam, et odísti iniquitátem: \Star{} Proptérea unxit te Deus, Deus tuus, óleo lætítiæ.\end{Resp}
\begin{Resp}\Vsign Propter veritátem, et mansuetúdinem, et justítiam.\end{Resp}
\begin{Resp}\Rsign Proptérea unxit te Deus, Deus tuus, óleo lætítiæ.\end{Resp}
\switchcolumn
\EN
\begin{Resp}\Rsign Thou hast loved righteousness, and hated iniquity; \Star{} Therefore God, thy God, hath anointed thee with the oil of gladness.\end{Resp}
\begin{Resp}\Vsign Because of truth, and meekness, and righteousness.\end{Resp}
\begin{Resp}\Rsign Therefore God, thy God, hath anointed thee with the oil of gladness.\end{Resp}
\switchcolumn*

\LA
\LessonHead{Lectio VI}
\switchcolumn
\EN
\LessonHead{Lesson VI}
\switchcolumn*

\LA
\noindent Defuncto rege Dionysio, sicut virgínibus in prima ætate, in matrimonio conjugibus, ita viduis in solitúdine fuit ómnium virtútem exemplar. Illico enim religiosis sanctæ Claræ vestibus induta, regio fúneri constanter intérfuit, ac paulo post, Compostellam proficiscens, multa ex holosérico, argénto, auro, gemmisque donaria pro regis ánima óbtulit. Inde reversa domum, quidquid sibi carum aut pretiósum supérerat, in sacros ac pios usus convértit; absolvendóque suo vere regio Conimbricénsi vírginum cœnobio, et aléndis paupéribus, et protegéndis viduis, defendéndis pupillis, miseris omnibus juvandis intenta, non sibi, sed Deo, et mortalium ómnium commodis vivebat. Reges duos, fílium et generum, pacificatúra Stremotium nóbile oppidum veniens, morbo ex itinere contracto, ibidem a Vírgine Deípara, visitata, sanctíssime obiit, anno millesimo trecentésimo trigesimo sexto die quarta Julii. Post mortem multis miraculis cláruit, præsertim suavíssimo corporis jam per annos fere trecentos incorrúpti odore; semper étiam reginæ sanctæ cognomento celebris. Tandem anno jubilæi, et nostræ salútis millesimo sexcentésimo vigesimo quinto, totius christiáni orbis concursu et applausu, ab Urbano octavo rite inter Sanctos adscripta est.\par
\switchcolumn
\EN
\noindent King Denis died (on the 6th day of January, 1325,) and Elizabeth, who in her maidenhood had been a pattern to virgins, and in her married life to wives, now, in her loneliness, was an example to widows. Clad in the raiment of the nuns of St. Clare, she faithfully attended at the King's funeral, and soon after went to Compostella, where she offered many precious gifts, of silk, and gold, and silver, and precious stones, for the benefit of his soul. Thence she returned home, and spent in holy and godly uses everything that remained to her that was dear and costly, eager to relieve every kind of suffering. She lived, not for herself, but for God, and to be useful to mankind. She finished the convent for nuns, right worthy of a Queen, which she had founded at Coimbra. She fed the poor, defended widows, protected orphans. A war being lighted up between her son Alphonsus IV, King of Portugal, and her grandson Alphonsus V, King of Castile, she resolved to set out to reconcile them, and went to the famous city of Estremoz, upon the borders of the two kingdoms. On the journey, she caught a violent fever, of which, after a vision of the Virgin Mother of God, she died a saintly death on the 4th day of July, in the year 1336. She became illustrious for miracles after her death, especially for the sweetness of the savour of her body, which hath now remained uncorrupt for well-nigh three hundred years, and she hath always been spoken of as the Holy Queen Elizabeth. At length, in the year of our salvation 1625, which was that of the Jubilee, Urban VIII, all Christendom gathered together and approving, formally enrolled her name among those of the Saints.\par
\switchcolumn*

\LA
\begin{Resp}\Rsign Fallax grátia, et vana est pulchritúdo: \Star{} Múlier timens Dóminum, ipsa laudábitur.\end{Resp}
\begin{Resp}\Vsign Date ei de fructu mánuum suárum, et laudent eam in portis ópera ejus.\end{Resp}
\begin{Resp}\Rsign Múlier timens Dóminum, ipsa laudábitur.\end{Resp}
\begin{Resp}\Vsign Glória Patri, et Fílio, \Star{} et Spirítui Sancto.\end{Resp}
\begin{Resp}\Rsign Múlier timens Dóminum, ipsa laudábitur.\end{Resp}
\switchcolumn
\EN
\begin{Resp}\Rsign Favour is deceitful, and beauty is vain \Star{} A woman that feareth God she shall be praised.\end{Resp}
\begin{Resp}\Vsign Give her of the fruit of her hands, and let her own works praise her in the gates.\end{Resp}
\begin{Resp}\Rsign A woman that feareth God she shall be praised.\end{Resp}
\begin{Resp}\Vsign Glory be to the Father, and to the Son, \Star{} and to the Holy Ghost.\end{Resp}
\begin{Resp}\Rsign A woman that feareth God she shall be praised.\end{Resp}
\switchcolumn*

\LA
\LessonHead{Lectio IX (pro commemoratione)}
\switchcolumn
\EN
\LessonHead{Lesson IX (when commemorated)}
\switchcolumn*

\LA
\LessonTitle{Commemoratio: S. Elisabeth Reg. Portugaliæ Viduæ}
\switchcolumn
\EN
\LessonTitle{Commemoration of: St. Elizabeth of Portugal, Queen and Widow}
\switchcolumn*

\LA
\noindent Elísabeth Aragóniæ régibus orta est anno Christi millésimo ducentésimo septuagésimo primo. Natális ejus lætítia perniciósas avi patrísque dissensiónes in concórdiam convértit, ex quo statim pátuit, quam felix regum regnorúmque esset futúra pacátrix. In castigándo córpore, in précibus assídue recitándis, in caritátis opéribus exercéndis admirábilis fuit. Dionýsio Lusitániæ regi in matrimónium trádita, non minórem excoléndis virtútibus quam líberis educándis óperam dabat, viro placére studens, sed magis Deo. Monastéria, collégia et templa non modo exstrúxit, sed étiam magnífice dotávit. In regum discórdiis componéndis admirábilis fuit, in privátis publicísque mortálium sublevándis calamitátibus indeféssa et miráculis clara. Defúncto rege Dionýsio, cum hábitum Seráphici órdinis induísset, quidquid sibi carum aut pretiósum supérerat, pro regis ánima templo Compostelláno óbtulit, et in sacros ac pios usus convértit. Dénique reges duos, fílium et génerum, pacificatúra, morbo ex itínere contrácto, a Vírgine Deípara visitáta, sanctíssime óbiit. Eam, miráculis claram, Urbánus octávus inter Sanctos adscrípsit.\par
\switchcolumn
\EN
\noindent Elizabeth was born of the royal family of Aragon in the year of our Lord 1271. The joy of her birth put an end to the unhappy quarrels between her grandfather and her father, thus making it clear from the outset that she would be a blessed peacemaker between kings and kingdoms. She was remarkable for the way in which she chastised her body, for her constancy in prayer, and for her exercise of the works of charity. When she was married to Denis, King of Portugal, she devoted herself no less to the work of cultivating virtue than to that of educating her children, striving to please her husband, but still more to please God. She not only had monasteries, colleges and churches built, but gave them magnificent endowments. She was wonderful in settling the disputes of kings, unwearied in relieving the private and public calamities of her fellow-men, and famous for her miracles. When King Denis had died, she put on the habit of the Seraphic Order, and whatever she had that was dear and precious to her she offered at the church of Compostella for the soul of the king and used for works of devotion and mercy. Finally, having fallen ill as a result of a journey she made to establish peace between two kings, her son and her grandson, she died a most holy death, after receiving a visit from the Virgin Mother of God. Famous for miracles, she was enrolled among the Saints by Urban VIII.\par
\switchcolumn*

\LA
\TeDeum{Te Deum}
\switchcolumn
\EN
\TeDeum{Te Deum}
\switchcolumn*

\end{paracol}

\par\needspace{10\baselineskip}\addvspace{1.2em}{\centering\small\textsc{Die 10 Julii} · \textsc{10 July}\par}
\Office{Ss. Septem Fratrum Martyrum, ac Rufinæ et Secundæ Virginum et Martyrum}{Seven Brothers, Martyrs, and Sts. Rufina and Secunda, Virgins and Martyrs}{Semiduplex}
\addcontentsline{toc}{subsection}{Die 10 Julii · Ss. Septem Fratrum Martyrum, ac Rufinæ et Secundæ Virginum et Martyrum \textit{— Seven Brothers, Martyrs, and Sts. Rufina and Secunda, Virgins and Martyrs}}
\begin{paracol}{2}
\LA
\RefLine{Lectiones I–III: de Scriptura occurrente.}
\switchcolumn
\EN
\RefLine{Lessons I–III: from the Scripture of the day.}
\switchcolumn*

\LA
\LessonHead{Lectio IV}
\switchcolumn
\EN
\LessonHead{Lesson IV}
\switchcolumn*

\LA
\noindent Septem fratres, fílii sanctæ Felicitátis, Romæ in persecutióne Marci Aurélii Antonini a Publio præfécto primum blandítiis, deínde terróribus tentati, ut Christo renuntiántes, deos veneraréntur; et, sua virtúte, et matre hortante, in fidei confessióne perseverántes, varie necáti sunt. Januarius plumbátis cæsus; Felix et Philippus fustibus contusi; Silvanus ex altíssimo loco præceps dejéctus est; Alexander, Vitalis et Martialis cápite plectúntur. Mater eórum quarto post mense eamdem martyrii palmam consecuta est: illi sexto Idus Julii spíritum Dómino reddidérunt.\par
\switchcolumn
\EN
\noindent In the persecution at Rome under Marcus Aurelius Antoninus, there were seven brethren, sons of the holy woman Felicity, whom the Prefect Publius first essayed to cajole by kindness, and then to shake by fear, to deny Christ and worship the gods but, by their own bravery and the exhortation of their mother, they remained firm in their confession, and were all put to death in diverse ways. Januarius was lashed to death with whips loaded with lead Felix and Philip were beaten to death with cudgels Silvanus was thrown over a precipice Alexander, Vitalis, and Martial were beheaded. Their mother gained the same palm of martyrdom four months afterwards. The seven Brethren gave up their souls to God upon the 10th day of July.\par
\switchcolumn*

\LA
\begin{Resp}\Rsign Sancti tui, Dómine, mirábile consecúti sunt iter, serviéntes præcéptis tuis, ut inveniréntur illǽsi in aquis válidis: \Star{} Terra appáruit árida: et in Mari Rubro via sine impediménto.\end{Resp}
\begin{Resp}\Vsign Quóniam percússit petram, et fluxérunt aquæ, et torréntes inundavérunt.\end{Resp}
\begin{Resp}\Rsign Terra appáruit árida: et in Mari Rubro via sine impediménto.\end{Resp}
\switchcolumn
\EN
\begin{Resp}\Rsign thy Saints, O Lord, have passed a wonderful way, serving thy commandments, that they might be found without hurt in the midst of the mighty waters. \Star{} Dry land appeared, and, out of the Red Sea, a way without impediment.\end{Resp}
\begin{Resp}\Vsign He smote the rock, and the waters gushed out, and the streams overflowed.\end{Resp}
\begin{Resp}\Rsign Dry land appeared, and, out of the Red Sea, a way without impediment.\end{Resp}
\switchcolumn*

\LA
\LessonHead{Lectio V}
\switchcolumn
\EN
\LessonHead{Lesson V}
\switchcolumn*

\LA
\noindent Rufina et Secunda sorores, vírgines Romanæ, rejecto connubio Armentárii et Verini, quibus a paréntibus desponsæ fuerant, quod Jesu Christo virginitátem vovissent, Valeriáno et Gallieno imperatóribus, comprehendúntur. Quas cum nec promissis nec terróre Juníus præfectus a proposito posset abducere, Rufinam primum virgis cædi jubet; in cujus verbéribus Secunda judicem sic interpellat: Quid est, quod sorórem meam honore, me áfficis ignomínia? Jube ambas simul cædi, quæ simul Christum Deum confitemur. Quibus verbis incénsus judex imperat utramque detrudi in tenebricósum et fœtidum carcerem. Quo loco statim claríssima luce et suavíssimo odore completo, in ardénte balnei solio includúntur. Et cum inde étiam integræ evasissent, mox saxo ad collum alligato in Tiberim projectæ sunt; unde ab Angelo liberátæ, extra Urbem via Aurelia milliario décimo, cápite plectúntur. Quarum corpora, a Plautilla matrona in ejus prædio sepúlta, ac póstea in Urbem translata, in basilica Constantiniana prope baptistérium cóndita sunt.\par
\switchcolumn
\EN
\noindent The virgin sisters, Rufina and Secunda, were Romans. Their parents had betrothed them to Armentarius and Verinus, but they both consecrated their virginity by vow to Christ, and refused marriage. They were arrested in the reign of the Emperors Valerian and Gallienus. The Praefect Junius failed to change their minds either by promises or threats, and then ordered Rufina to be scourged. While the lashing was going on, Secunda said to the Judge Why dost thou judge my sister to honour and me to dishonour Be pleased to whip us both together, for we both together declare that Christ is God. The Judge was angered at these words, and ordered them both to a dark and stinking dungeon but it was presently filled with a bright light and a sweet savour. They were then shut up in a hot flue of a bath, but they came forth from it unharmed. Stones were next tied to their necks and they were cast into the river Tiber, but an Angel delivered them therefrom. In the end they were beheaded on the Aurelian Way, at the tenth mile-stone from the city. The Lady Plautilla buried their bodies upon her own farm, but they were afterwards brought into the city, and laid in the Cathedral Church of the Most Holy Saviour, hard by the Baptistery.\par
\switchcolumn*

\LA
\begin{Resp}\Rsign Vérbera carníficum non timuérunt Sancti Dei, moriéntes pro Christi nómine: \Star{} Ut herédes fíerent in domo Dómini.\end{Resp}
\begin{Resp}\Vsign Tradidérunt córpora sua propter Deum ad supplícia.\end{Resp}
\begin{Resp}\Rsign Ut herédes fíerent in domo Dómini.\end{Resp}
\switchcolumn
\EN
\begin{Resp}\Rsign The Saints of God shrank not from the stripes of the executioners, but died for Christ's Name's sake \Star{} That they might be made joint-heirs in the house of the Lord.\end{Resp}
\begin{Resp}\Vsign They gave their bodies for God's sake to death.\end{Resp}
\begin{Resp}\Rsign That they might be made joint-heirs in the house of the Lord.\end{Resp}
\switchcolumn*

\LA
\LessonHead{Lectio VI}
\switchcolumn
\EN
\LessonHead{Lesson VI}
\switchcolumn*

\LA
\LessonTitle{Sermo sancti Augustíni Epíscopi}
\switchcolumn
\EN
\LessonTitle{The Lesson is taken from the Sermons of St. Augustine, Bishop (of Hippo.)}
\switchcolumn*

\LA
\Source{Sermo 110 de diversis}
\switchcolumn
\EN
\Source{11th on various subjects.}
\switchcolumn*

\LA
\noindent Magnum spectaculum, fratres, positum est ante óculos fidei nostræ. Aure audívimus, corde vídimus optántem matrem ante se finire istam vitam fílios suos, longe contrariis votis consuetudini humanæ. Omnes enim hómines fílios suos, ex hac vita migrando, præcedere volunt, non sequi; illa autem optávit posterior mori. Non enim ammitebat fílios, sed præmittebat; nec intuebátur quam vitam finirent, sed quam inchoarent. Desinébant enim vivere, ubi quandoque fuerant morituri; et incipiébant vivere, sine fine victuri. Parum est fuísse spectatricem; miráti sumus potius hortatricem. Fœcundior virtútibus quam fœtibus: videns certántes, in quibus omnibus ipsa certabat, et in omnibus vincéntibus ipsa vincebat.\par
\switchcolumn
\EN
\noindent Wonderful is the sight, my brethren, which is set before the eyes of our faith. We have heard with our ears, and we see in our thoughts, a mother, with superhuman love, watching her sons leaving this life before her. All men would fain depart hence before their children, but she was ready to die last. Departing from her, they were not lost, but gone before. And she looked, not to the life they were ending, but to the life they were beginning. They laid aside a life which must needs end in death, and began that life wherein they are alive for ever. The least of her work was that she was an on-looker more amazing is it, when we remember that she was their exhortress. Her courage was more fruitful than her womb, and when she saw them contending and conquering, her heart contended and conquered in each.\par
\switchcolumn*

\LA
\begin{Resp}\Rsign Tamquam aurum in fornáce probávit eléctos Dóminus, et quasi holocáusti hóstiam accépit illos: et in témpore erit respéctus illórum: \Star{} Quóniam donum et pax est eléctis Dei.\end{Resp}
\begin{Resp}\Vsign Qui confídunt in illum, intélligent veritátem, et fidéles in dilectióne acquiéscent illi.\end{Resp}
\begin{Resp}\Rsign Quóniam donum et pax est eléctis Dei.\end{Resp}
\begin{Resp}\Vsign Glória Patri, et Fílio, \Star{} et Spirítui Sancto.\end{Resp}
\begin{Resp}\Rsign Quóniam donum et pax est eléctis Dei.\end{Resp}
\switchcolumn
\EN
\begin{Resp}\Rsign As gold in the furnace hath the Lord tried His chosen ones, and received them as a burnt-offering, and yet a while, and they shall be regarded; \Star{} For the grace of God, and His peace, are with His chosen.\end{Resp}
\begin{Resp}\Vsign They that put their trust in Him shall understand the truth and such as be faithful in love shall abide with Him.\end{Resp}
\begin{Resp}\Rsign For the grace of God, and His peace, are with His chosen.\end{Resp}
\begin{Resp}\Vsign Glory be to the Father, and to the Son, \Star{} and to the Holy Ghost.\end{Resp}
\begin{Resp}\Rsign For the grace of God, and His peace, are with His chosen.\end{Resp}
\switchcolumn*

\LA
\LessonHead{Lectio VII}
\switchcolumn
\EN
\LessonHead{Lesson VII}
\switchcolumn*

\LA
\LessonTitle{Léctio sancti Evangélii secúndum Matthǽum}
\switchcolumn
\EN
\LessonTitle{From the Holy Gospel according to Matthew}
\switchcolumn*

\LA
\Cite{Matt 12:46-50}
\switchcolumn
\EN
\Cite{Matt 12:46-50}
\switchcolumn*

\LA
\noindent In illo témpore: Loquénte Jesu ad turbas, ecce Mater ejus et fratres stabant foris quæréntes loqui ei. Et réliqua.\par
\switchcolumn
\EN
\noindent At that time: While Jesus yet talked to the people, behold, His mother and His brethren stood without, desiring to speak with Him. And so on.\par
\switchcolumn*

\LA
\LessonTitle{Homilía sancti Gregórii Papæ}
\switchcolumn
\EN
\LessonTitle{Homily by Pope St. Gregory the Great}
\switchcolumn*

\LA
\Source{Homilia 3 in Evangelia}
\switchcolumn
\EN
\Source{3rd on the Gospels}
\switchcolumn*

\LA
\noindent Sancti Evangélii, fratres caríssimi, brevis est léctio recitata, sed magnis mysteriórum pondéribus grávida. Jesus étenim, cónditor et redémptor noster, Matrem se nosse dissimulat, et quæ ei mater sit et qui propinqui, non per cognatiónem carnis, sed per conjunctiónem spíritus designat, dicens: Quæ est mater mea, et qui sunt fratres mei? Quicúmque enim fécerit voluntátem Patris mei, qui in cælis est, ipse meus frater, et soror, et mater est. Quibus nobis verbis quid aliud innuit, nisi quod obsequéntes jussiónibus suis multos ex Gentilitate colligit; et Judǽam, ex cujus carne est génitus, non agnoscit?\par
\switchcolumn
\EN
\noindent Dearly beloved brethren The gospel which is read this day is but very short, but it is heavy with great mysteries. Here Jesus, our Maker and Redeemer, feigneth Himself as though He knew not His Own Mother, and telleth who they be who are His mother and brethren, not by fleshly kinship, but by kinship of mind. Who is My Mother And who are My brethren Whosoever shall do the will of My Father, Which is in heaven, the same is My brother, and sister, and mother. By the which words, what are we to understand, save that He gathereth together from out of Heathendom many that are willing to obey His commandments, and that He knoweth not Jewry, whereof, according to the flesh, He is a son.\par
\switchcolumn*

\LA
\begin{Resp}\Rsign Propter testaméntum Dómini, et leges patérnas, Sancti Dei perstitérunt in amóre fraternitátis: \Star{} Quia unus fuit semper spíritus in eis, et una fides.\end{Resp}
\begin{Resp}\Vsign Ecce quam bonum et quam jucúndum habitáre fratres in unum.\end{Resp}
\begin{Resp}\Rsign Quia unus fuit semper spíritus in eis, et una fides.\end{Resp}
\switchcolumn
\EN
\begin{Resp}\Rsign Because of the covenant of the Lord, and the laws of their fathers, the Saints of God abode in brotherly love, \Star{} For one spirit and one faith was ever in them.\end{Resp}
\begin{Resp}\Vsign Behold how good and how pleasant it is for brethren to dwell together in unity.\end{Resp}
\begin{Resp}\Rsign For one spirit and one faith was ever in them.\end{Resp}
\switchcolumn*

\LA
\LessonHead{Lectio VIII}
\switchcolumn
\EN
\LessonHead{Lesson VIII}
\switchcolumn*

\LA
\noindent Sed cum is, qui voluntátem Patris fécerit, soror et frater Dómini dícitur, propter utrumque sexum qui ad fidem collígitur, mirum non est; mirándum vero valde est, quómodo étiam mater dicátur. Fideles enim discipulos fratres vocare dignátus est, dicens: Ite, nuntiáte frátribus meis. Qui ergo frater Dómini fíeri ad fidem veniéndo potúerit, quæréndum est quómodo étiam et mater esse possit.\par
\switchcolumn
\EN
\noindent Seeing that both men and women are called to the faith, we marvel not that He saith that whosoever shall do the will of His Father, the same is His brother, and sister. But it is startling to hear that the same is also His mother. His faithful disciples He is pleased to call His brethren, where He saith Go, tell My brethren, (Matth. xxviii. 10.) If then it is by joining His religion that one can become the brother of the Lord, let us see how one can become His mother.\par
\switchcolumn*

\LA
\begin{Resp}\Rsign Sancti mei, qui in carne pósiti, certámen habuístis: \Star{} Mercédem labóris ego reddam vobis.\end{Resp}
\begin{Resp}\Vsign Veníte benedícti Patris mei, percípite regnum.\end{Resp}
\begin{Resp}\Rsign Mercédem labóris ego reddam vobis.\end{Resp}
\begin{Resp}\Vsign Glória Patri, et Fílio, \Star{} et Spirítui Sancto.\end{Resp}
\begin{Resp}\Rsign Mercédem labóris ego reddam vobis.\end{Resp}
\switchcolumn
\EN
\begin{Resp}\Rsign O ye My Saints, who, being in the flesh, didst have striving \Star{} I will render unto you a reward of your labours.\end{Resp}
\begin{Resp}\Vsign Come, ye blessed of My Father, inherit the kingdom\end{Resp}
\begin{Resp}\Rsign I will render unto you a reward of your labours.\end{Resp}
\begin{Resp}\Vsign Glory be to the Father, and to the Son, \Star{} and to the Holy Ghost.\end{Resp}
\begin{Resp}\Rsign I will render unto you a reward of your labours.\end{Resp}
\switchcolumn*

\LA
\LessonHead{Lectio IX}
\switchcolumn
\EN
\LessonHead{Lesson IX}
\switchcolumn*

\LA
\noindent Sed sciéndum nobis est, quia qui Christi soror et frater est credendo, mater effícitur prædicando. Quasi enim parit Dóminum, quem cordi audiéntis infuderit; et mater ejus prædicando effícitur, si per ejus vocem amor Dómini in próximi mente generátur. Ad quam rem nobis idónee confírmandam adest beáta Felícitas, cujus hódie natalícia celebramus, quæ credéndo exstitit ancílla Christi, et prædicando facta est mater Christi. Septem quippe fílios, sicut in gestis ejus emendatióribus legitur, sic post se tímuit vivos in carne relínquere, sicut carnales paréntes solent metúere ne mórtuos præmittant.\par
\switchcolumn
\EN
\noindent But we must know that even as one becometh His brother or His sister by believing in Him, so one becometh His mother by preaching Him. Such an one, as it were, giveth birth to the Lord by causing Him to be in the hearer's heart, and by words giving His love existence in their neighbour's mind. For an example in point, behold blessed Felicity, whose Birth-day we keep today. She was one whose faith made her Christ's hand-maiden, and whose preaching made her Christ's mother. We read in the corrected edition of her Acts that she dreaded as much to leave her seven sons behind her alive in the flesh, as do worldly mothers to send theirs dead before them.\par
\switchcolumn*

\LA
\TeDeum{Te Deum}
\switchcolumn
\EN
\TeDeum{Te Deum}
\switchcolumn*

\end{paracol}

\par\needspace{10\baselineskip}\addvspace{1.2em}{\centering\small\textsc{Die 11 Julii} · \textsc{11 July}\par}
\Office{S. Pii I Papæ et Martyris}{St. Pius I, Pope and Martyr}{Simplex}
\addcontentsline{toc}{subsection}{Die 11 Julii · S. Pii I Papæ et Martyris \textit{— St. Pius I, Pope and Martyr}}
\begin{paracol}{2}
\LA
\RefLine{Lectiones I–II: de Scriptura occurrente.}
\switchcolumn
\EN
\RefLine{Lessons I–II: from the Scripture of the day.}
\switchcolumn*

\LA
\LessonHead{Lectio III (pro commemoratione)}
\switchcolumn
\EN
\LessonHead{Lesson III (when commemorated)}
\switchcolumn*

\LA
\noindent Pius, hujus nóminis primus, Aquilejénsis, Ruffini fílius, ex presbytero sanctæ Romanæ Ecclésiæ Summus Pontifex creátus est, Antonino Pio et Marco Aurelio imperatóribus augustis. Quinque ordinatiónibus, mense Decembri, episcopos duódecim, octódecim presbyteros creávit. Exstant nonnulla ab eo præcláre instituta, præsertim ut Resurréctio Dómini nónnisi die Dominico celebrarétur. Pudentis domum in ecclésiam mutávit, eamque ob præstantiam supra ceteros titulos, útpote Romani Pontificis mansiónem, título Pastoris dicávit; et in qua sæpe rem sacram fecit, et multos ad fidem conversos baptizávit ac in fidelium númerum adscripsit. Dum vero boni pastoris munus obiret, fuso pro suis ovibus et summo pastore Christo sánguine, martyrio coronátus est quinto Idus Julii, ac sepúltus in Vaticano.\par
\switchcolumn
\EN
\noindent Pius I., the son of Rufinus, was from Aquilia, and was a Priest of the holy Roman Church when he was made Supreme Pontiff, he lived under the Emperors Antoninus Pius and Marcus Aurelius, he held five ordinations in the month of December, wherein he ordained twelve Bishops and eighteen Priests. There remain several eminent ordinances of his, notably that which ruleth that the Resurrection of the Lord be not observed upon any day of the week save the Lord's Day. He turned the house of Pudens into a church, and on account of its eminence above the other churches, as being that where the Bishop of Rome dwelt, he dedicated it under the name of the Shepherd. Here he often celebrated, and baptized and numbered among the faithful many converts to the faith. While he strove to do the work of a good shepherd he shed his blood for his sheep, and for the chief Shepherd Christ. He was crowned with martyrdom upon the 11th day of July, and buried upon the Vatican Hill.\par
\switchcolumn*

\LA
\TeDeum{Te Deum}
\switchcolumn
\EN
\TeDeum{Te Deum}
\switchcolumn*

\end{paracol}

\par\needspace{10\baselineskip}\addvspace{1.2em}{\centering\small\textsc{Die 12 Julii} · \textsc{12 July}\par}
\Office{S. Joannis Gualberti Abbatis}{St. John Gualbert, Abbot}{Duplex}
\addcontentsline{toc}{subsection}{Die 12 Julii · S. Joannis Gualberti Abbatis \textit{— St. John Gualbert, Abbot}}
\begin{paracol}{2}
\LA
\RefLine{Lectiones I–III: de Scriptura occurrente.}
\switchcolumn
\EN
\RefLine{Lessons I–III: from the Scripture of the day.}
\switchcolumn*

\LA
\LessonHead{Lectio IV}
\switchcolumn
\EN
\LessonHead{Lesson IV}
\switchcolumn*

\LA
\noindent Joánnes Gualbertus, Floréntiæ nobili genere ortus, dum patri óbsequens rem militarem sequitur, Ugo, únicus ejus frater, occíditur a consanguíneo. Quem cum solum et inérmem sancto Parasceves die Joánnes, armis ac milítibus stipatus, óbvium haberet, ubi néuter álterum poterat declináre, ob sanctæ crucis reveréntiam, quam homicida supplex, mortem jamjam subiturus, brácchiis signabat, vitam ei clementer indulget. Hoste in fratrem recepto, próximum sancti Miniátis templum oratúrus ingréditur, ubi adoratam Crucifixi imáginem caput sibi fléctere cónspicit. Quo mirábili facto permotus Joánnes, Deo exínde, étiam invíto patre, militare decernit, atque ibidem propriis sibi mánibus comam totondit, ac monasticum hábitum índuit : adeoque piis ac religiosis virtútibus brevi coruscat, ut multis se perfectiónis specimen ac normam præbéret ; ita ut, ejusdem loci abbate defuncto, communi ómnium voto in superiórem eligerétur. At Dei famulus, cupiens subesse potius quam præesse, ad majora divina voluntáte servatus, ad Camaldulénsis eremi íncolam Romualdum proficíscitur, a quo cælicum sui institúti vaticinium áccipit ; tum suum ordinem sub regula sancti Benedicti apud Umbrosam vallem instituit.\par
\switchcolumn
\EN
\noindent John Gualberto was the son of a noble family at Florence. In accordance with the wishes of his father, he became a soldier. While he was in that profession, his only brother, Hew, was slain by a cousin. On a certain Good Friday, John, armed and accompanied by soldiers, met the murderer, alone and defenceless, in a narrow way, where neither could turn aside. As he was at the point to kill him, the wretch fell on his knees, and stretched out his arms in the form of the Cross, adjuring him, for the sake of that sign, to forgive him and out of reverence for the Cross he had mercy on him and spared his life. After pardoning his enemy, he went into the Church of St. Minias, which was hard by, to pray. And there he saw the image of Jesus crucified, which had that day received the worship of the faithful, bow its head to him. By this miracle John was so moved, that he laid aside soldiering, even against his father's wishes cut off his hair with his own hands, at the Convent of St. Minias, and clad himself in the garb of a monk. In a short while he so shone with all godly and monkish graces, that he became a pattern of excellency to many. When the Abbat of that house died, the monks all chose John to succeed him. But the servant of God desired to obey, more than to command, and, being kept by God for greater things, he betook himself to one Romuald, a dweller in the hermitage at Camaldoli. Through Romuald he received a revelation from heaven, and forthwith founded an Order of his own under the Rule of St. Benedict, in the valley called Vallombrosa.\par
\switchcolumn*

\LA
\begin{Resp}\Rsign Honéstum fecit illum Dóminus, et custodívit eum ab inimícis, et a seductóribus tutávit illum: \Star{} Et dedit illi claritátem ætérnam.\end{Resp}
\begin{Resp}\Vsign Justum dedúxit Dóminus per vias rectas, et osténdit illi regnum Dei.\end{Resp}
\begin{Resp}\Rsign Et dedit illi claritátem ætérnam.\end{Resp}
\switchcolumn
\EN
\begin{Resp}\Rsign The Lord made him honourable, and defended him from his enemies, and kept him safe from those that lay in wait for him; \Star{} And gave him perpetual glory.\end{Resp}
\begin{Resp}\Vsign He went down with him into the pit, and left him not in bonds.\end{Resp}
\begin{Resp}\Rsign And gave him perpetual glory.\end{Resp}
\switchcolumn*

\LA
\LessonHead{Lectio V}
\switchcolumn
\EN
\LessonHead{Lesson V}
\switchcolumn*

\LA
\noindent Deínde, plurimis ad eum ob ejus sanctitátis famam undique convolántibus, una cum iis in socios adscitis, ad hæreticam et simoníacam pravitátem exstirpandam et apostolicam fidem propagandam sédulo incumbit, innumera proptérea in se et suis incommoda expertus. Nam, ut eum ejusque socios adversarii perdant, noctu sancti Salvii cœnóbium repénte aggrediúntur, templum incendunt, ædes demoliúntur, et mónachos omnes lethali vulnere sauciant ; quos vir Dei único crucis signo incólumes prótinus reddit, et, Petro ejus mónacho per immensum ardentissimúmque ignem illæso mirabíliter transeúnte, optatam sibi et suis tranquillitátem óbtinet. Inde simoníacam labem ab Etruria éxpulit, ac in tota Italia fidem prístinæ integritáti restituit.\par
\switchcolumn
\EN
\noindent Many gathered themselves to him, drawn by the fame of his holy life. Them he took for his comrades, and laboured earnestly among them to cleanse the Church in those parts from the pollution of heresy and simony, and spread abroad the Apostolic Faith. He and his had to fight with almost countless hardships. Certain enemies broke by night into the monastery of San Salvi, to destroy John and his monks, set the church on fire, pulled down the huts, and mortally wounded all the monks but the man of God perfectly healed them all by the sign of the Cross. One of his monks named Peter also passed unhurt through a vast and raging fire. At length John and his disciples got the peace which they longed for. He purged Tuscany of the pollution of simony, and restored the faith throughout all Italy to its first purity.\par
\switchcolumn*

\LA
\begin{Resp}\Rsign Amávit eum Dóminus, et ornávit eum: stolam glóriæ índuit eum, \Star{} Et ad portas paradísi coronávit eum.\end{Resp}
\begin{Resp}\Vsign Índuit eum Dóminus lorícam fídei, et ornávit eum.\end{Resp}
\begin{Resp}\Rsign Et ad portas paradísi coronávit eum.\end{Resp}
\switchcolumn
\EN
\begin{Resp}\Rsign The Lord loved him and beautified him; He clothed him with a robe of glory, \Star{} And crowned him at the gates of Paradise.\end{Resp}
\begin{Resp}\Vsign The Lord hath put on him the breast-plate of faith, and hath adorned him.\end{Resp}
\begin{Resp}\Rsign And crowned him at the gates of Paradise.\end{Resp}
\switchcolumn*

\LA
\LessonHead{Lectio VI}
\switchcolumn
\EN
\LessonHead{Lesson VI}
\switchcolumn*

\LA
\noindent Multa fúnditus eréxit monasteria, eademque et alia ædificiis ac regulari observántia instaurata, sanctis légibus communívit. Ad egénos alendos sacram supelléctilem véndidit : ad ímprobos coërcendos eleménta sibi famulari conspéxit : ad dæmones comprimendos crucem quasi ensem adhibuit. Demum abstinentiis, vigiliis, jejuniis, oratiónibus, carnis maceratiónibus ac senio confectus, dum infirma valetúdine gravarétur. Davídica illa verba persæpe repetebat : Sitívit ánima mea ad Deum fortem vivum : quando véniam et apparébo ante fáciem Dei? Jamque morti próximus, convocatos discipulos ad fraternam concordiam cohortátur, et in brevículo, cui consepeliri vóluit, jussit hæc scribi : Ego Joánnes credo et confiteor fidem, quam sancti Apóstoli prædicavérunt, et sancti Patres in quatuor conciliis confirmavérunt. Tandem, triduano Angelórum obsequio dignátus, septuagesimum octavum annum agens, apud Passinianum, ubi summa veneratióne cólitur, migrávit ad Dóminum, anno salútis millesimo septuagesimo tertio, quarto Idus Julii. Quem Cælestinus tertius, innumeris miraculis clarum, in Sanctórum númerum rétulit.\par
\switchcolumn
\EN
\noindent He entirely built several monasteries, and furnished them and others with buildings. He restored in them the strict observance of the Rule, and gave them holy laws. He sold the furniture of the Church to feed the poor, and found the very elements subject to him to bend stubborn hearts withal. He used the Cross like a sword to drive out devils. In his old age, worn out by abstinence, watching, fasts, prayers, and punishing of the flesh, his strength utterly gave way, and he often repeated the words of David My soul thirsteth for God, for the mighty God, for the living God when shall I come and appear before God \textasciitilde{}(Ps. xli. 2.) When he was at the point of death, he gathered his disciples together and exhorted them to love one another, and, after a little while, ordered the following words to be written down, which he wished should be buried with him I, John, do believe and confess that Faith which the Holy Apostles preached, and which the Holy Fathers have ratified in the four Councils. At length, at Passignano, where he is held in the highest reverence, after a vision of angels which lasted three days, he passed away to be with the Lord, upon the 12th day of July, in the 78th year of his own age, and in that of salvation 1073. He is illustrious for countless miracles, and Celestine III. enrolled his name among those of the Saints.\par
\switchcolumn*

\LA
\begin{Resp}\Rsign Iste homo perfécit ómnia quæ locútus est ei Deus, et dixit ad eum: Ingrédere in réquiem meam: \Star{} Quia te vidi justum coram me ex ómnibus géntibus.\end{Resp}
\begin{Resp}\Vsign Iste est, qui contémpsit vitam mundi, et pervénit ad cæléstia regna.\end{Resp}
\begin{Resp}\Rsign Quia te vidi justum coram me ex ómnibus géntibus.\end{Resp}
\begin{Resp}\Vsign Glória Patri, et Fílio, \Star{} et Spirítui Sancto.\end{Resp}
\begin{Resp}\Rsign Quia te vidi justum coram me ex ómnibus géntibus.\end{Resp}
\switchcolumn
\EN
\begin{Resp}\Rsign This is he which did according unto all that God commanded him; and God said unto him: Enter thou into My rest, \Star{} For thee have I seen righteous before Me among all people.\end{Resp}
\begin{Resp}\Vsign This is he which loved not his life in this world, and is come unto an everlasting kingdom.\end{Resp}
\begin{Resp}\Rsign For thee have I seen righteous before Me among all people.\end{Resp}
\begin{Resp}\Vsign Glory be to the Father, and to the Son, \Star{} and to the Holy Ghost.\end{Resp}
\begin{Resp}\Rsign For thee have I seen righteous before Me among all people.\end{Resp}
\switchcolumn*

\LA
\LessonHead{Lectio VII}
\switchcolumn
\EN
\LessonHead{Lesson VII}
\switchcolumn*

\LA
\LessonTitle{Léctio sancti Evangélii secúndum Matthǽum}
\switchcolumn
\EN
\LessonTitle{From the Holy Gospel according to Matthew}
\switchcolumn*

\LA
\Cite{Matt 5:43-48}
\switchcolumn
\EN
\Cite{Matt 5:43-48}
\switchcolumn*

\LA
\noindent In illo témpore: Dixit Jesus discípulis suis: Audístis quia dictum est: Díliges próximum tuum, et ódio habebis inimícum tuum. Et réliqua.\par
\switchcolumn
\EN
\noindent At that time, Jesus said unto His disciples: Ye have heard that it hath been said, Thou shalt love thy neighbour and hate thine enemy. And so on.\par
\switchcolumn*

\LA
\LessonTitle{Homilía sancti Hierónymi Presbýteri}
\switchcolumn
\EN
\LessonTitle{Homily by St. Jerome, Priest at Bethlehem.}
\switchcolumn*

\LA
\Source{Liber 1 Comment. in cap. 5 Matth.}
\switchcolumn
\EN
\Source{Bk. i., Cowan. on Matth. v. and vi.}
\switchcolumn*

\LA
\noindent Ego autem dico vobis: Dilígite inimícos vestros; benefácite his qui odérunt vos. Multi præcepta Dei, imbecillitate sua, non Sanctórum viribus æstimántes, putant esse impossibília quæ præcepta sunt: et dicunt sufficere virtútibus, non odisse inimícos: ceterum diligere, plus præcipi, quam humana natúra patiátur. Sciéndum est ergo, Christum non impossibília præcipere, sed perfécta. Quæ fecit David in Saul, et in Absalom: Stephanus quoque Martyr pro inimícis lapidántibus deprecátus est: et Paulus anáthema cupit esse pro persecutóribus suis. Hæc autem Jesus et docuit et fecit, dicens: Pater, ignosce illis: quod enim faciunt, nesciunt.\par
\switchcolumn
\EN
\noindent But I say unto you Love your enemies, do good to them that hate you. There are many who judge of the commandments of the Lord by their own weakness, and not by the strength of His Saints and so deem Him to have commanded things impossible. These are they who think that not to hate their enemies is all that they are able to do and that to command us to love them, is to command more than man's nature can bear. It behoveth then to know, that this which Christ commandeth is not impossible, albeit perfect. This is what David did in respect of Saul and Absalom the martyr Stephen also prayed for his enemies, even while they were stoning him and Paul could wish that himself were accursed from Christ for his persecutors. (Rom. ix. well as taught, when He said Father, forgive them for they know not what they do. (Luke xxiii. 34.)\par
\switchcolumn*

\LA
\begin{Resp}\Rsign Iste est qui ante Deum magnas virtútes operátus est, et de omni corde suo laudávit Dóminum: \Star{} Ipse intercédat pro peccátis ómnium populórum.\end{Resp}
\begin{Resp}\Vsign Ecce homo sine queréla, verus Dei cultor, ábstinens se ab omni ópere malo, et pérmanens in innocéntia sua.\end{Resp}
\begin{Resp}\Rsign Ipse intercédat pro peccátis ómnium populórum.\end{Resp}
\switchcolumn
\EN
\begin{Resp}\Rsign This is he which wrought great wonders before God, and praised the Lord with all his heart. \Star{} May he pray for all people, that their sins may be forgiven unto them.\end{Resp}
\begin{Resp}\Vsign Behold a man without blame, a worshipper of God in truth, keeping himself clean from every evil work, and abiding still in his innocency.\end{Resp}
\begin{Resp}\Rsign May he pray for all people, that their sins may be forgiven unto them.\end{Resp}
\switchcolumn*

\LA
\LessonHead{Lectio VIII}
\switchcolumn
\EN
\LessonHead{Lesson VIII}
\switchcolumn*

\LA
\noindent In réliquis opéribus bonis intérdum potest áliquis qualemcúmque excusatiónem præténdere; ad habéndam vero dilectiónem nullus se póterit excusare. Potest mihi áliquis dicere, Non possum jejunare; numquid potest dicere, Non possum amare? Potest áliquis dicere, Virginitátem non possum servare, non possum res totas véndere et paupéribus erogare; numquid potest dicere, Non possum diligere inimícos?\par
\switchcolumn
\EN
\noindent For the leaving undone other good works, some excuse can sometimes be given but no man can give an excuse for being loveless. Such and such an one may say to me, I am not able to fast but can he say, I am not able to love Such and such an one may say, I am not able to remain a virgin, I am not able to sell all that I have and give to the poor but can he say, I am not able to love my enemies\par
\switchcolumn*

\LA
\begin{Resp}\Rsign Sint lumbi vestri præcíncti, et lucérnæ ardéntes in mánibus vestris: \Star{} Et vos símiles homínibus exspectántibus dóminum suum, quando revertátur a núptiis.\end{Resp}
\begin{Resp}\Vsign Vigiláte ergo, quia nescítis qua hora Dóminus vester ventúrus sit.\end{Resp}
\begin{Resp}\Rsign Et vos símiles homínibus exspectántibus dóminum suum, quando revertátur a núptiis.\end{Resp}
\begin{Resp}\Vsign Glória Patri, et Fílio, \Star{} et Spirítui Sancto.\end{Resp}
\begin{Resp}\Rsign Et vos símiles homínibus exspectántibus dóminum suum, quando revertátur a núptiis.\end{Resp}
\switchcolumn
\EN
\begin{Resp}\Rsign Let your loins be girded about, and your lights burning; \Star{} And ye yourselves like unto men that wait for their lord, when he will return from the wedding.\end{Resp}
\begin{Resp}\Vsign Watch therefore, for ye know not what hour your Lord doth come.\end{Resp}
\begin{Resp}\Rsign And ye yourselves like unto men that wait for their lord, when he will return from the wedding.\end{Resp}
\begin{Resp}\Vsign Glory be to the Father, and to the Son, \Star{} and to the Holy Ghost.\end{Resp}
\begin{Resp}\Rsign And ye yourselves like unto men that wait for their lord, when he will return from the wedding.\end{Resp}
\switchcolumn*

\LA
\LessonHead{Lectio IX}
\switchcolumn
\EN
\LessonHead{Lesson IX}
\switchcolumn*

\LA
\noindent Non enim ibi aut pedes laborant, currendo, aures audiendo, aut manus operando laxántur, ut nos per ipsam excusatiónem liberare conémur. Non nobis dícitur: ite ad Oriéntem, et quærite caritátem; navigate ad Occidéntem, et inveniétis dilectiónem. Intus in nostro corde est, ubi redire jubemur, dicénte prophéta: Redíte, prævaricatores, ad cor. Non enim in longinquis regiónibus invenítur quod a nobis pétitur.\par
\switchcolumn
\EN
\noindent This is a work wherein the feet are not wearied with running, nor the ears with hearing, neither do the hands fail from labour, that we should set up thereby an excuse to rid us of the duty. It is not said unto us Go to the East, and search for charity sail to the West, and ye shall find love. It is into our own inner hearts that we are to go, as saith the Prophet Bring it again to mind, O ye transgressors. (Isa. xlvi. 8.) What is asked of us is not to be found afar off.\par
\switchcolumn*

\LA
\TeDeum{Te Deum}
\switchcolumn
\EN
\TeDeum{Te Deum}
\switchcolumn*

\end{paracol}

\par\needspace{10\baselineskip}\addvspace{1.2em}{\centering\small\textsc{Die 13 Julii} · \textsc{13 July}\par}
\Office{S. Anacleti Papæ et Martyris}{St. Anacletus, Pope and Martyr}{Semiduplex}
\addcontentsline{toc}{subsection}{Die 13 Julii · S. Anacleti Papæ et Martyris \textit{— St. Anacletus, Pope and Martyr}}
\begin{paracol}{2}
\LA
\RefLine{Lectiones I–III: de Scriptura occurrente.}
\switchcolumn
\EN
\RefLine{Lessons I–III: from the Scripture of the day.}
\switchcolumn*

\LA
\RefLine{Lectiones VII–IX: ex Commúni unius Summi Pontificis et Martyris (C2b).}
\switchcolumn
\EN
\RefLine{Lessons VII–IX: from the Common of a Single Martyr Pope (C2b).}
\switchcolumn*

\LA
\LessonHead{Lectio IV}
\switchcolumn
\EN
\LessonHead{Lesson IV}
\switchcolumn*

\LA
\noindent Anacletus Atheniensis, Trajano imperatore rexit Ecclesiam. Decrevit, ut episcopus a tribus episcopis, neque a paucioribus consecraretur, et clerici sacris ordinibus publice a proprio episcopo initiarentur: et ut in Missa, peracta consecratione omnes communicarent. Beati Petri sepulcrum ornavit, Pontificumque sepulturæ locum attribuit. Fecit ordinationes duas mense Decembri, quibus creavit presbyteros quinque, diaconos tres, episcopos sex. Sedit annos novem, menses tres, dies decem. Martyrio coronatus, sepultus est in Vaticano.\par
\switchcolumn
\EN
\noindent Anacletus was an Athenian who governed the Church in the time of the Emperor Trajan. He ordained that a Bishop should be consecrated by three Bishops and no less, that clerks should be publicly ordained to Holy Orders by their own Bishop, and that in the Mass, after the Consecration, all should afterwards Communicate. He adorned the grave of Blessed Peter, and ordered a place for burying the Popes in. He held two ordinations in the month of December, wherein he ordained five Priests, three Deacons, and six Bishops. He sat as Pope nine years, three months, and ten days. He received the crown of his testimony, and was buried on the Vatican Hill.\par
\switchcolumn*

\LA
\begin{Resp}\Rsign Honéstum fecit illum Dóminus, et custodívit eum ab inimícis, et a seductóribus tutávit illum: \Star{} Et dedit illi claritátem ætérnam.\end{Resp}
\begin{Resp}\Vsign Descendítque cum illo in fóveam, et in vínculis non derelíquit eum.\end{Resp}
\begin{Resp}\Rsign Et dedit illi claritátem ætérnam.\end{Resp}
\switchcolumn
\EN
\begin{Resp}\Rsign The Lord made him honourable, and defended him from his enemies, and kept him safe from those that lay in wait for him, \Star{} And gave him perpetual glory.\end{Resp}
\begin{Resp}\Vsign He went down with him into the pit, and left him not in bonds.\end{Resp}
\begin{Resp}\Rsign And gave him perpetual glory.\end{Resp}
\switchcolumn*

\LA
\LessonHead{Lectio V}
\switchcolumn
\EN
\LessonHead{Lesson V}
\switchcolumn*

\LA
\LessonTitle{De Expositione sancti Ambrósii Epíscopi in Psalmum centesimum decimum octavum.}
\switchcolumn
\EN
\LessonTitle{From the Exposition of the hundred-and-eighteenth Psalm by St. Ambrose, Bishop of Milan.}
\switchcolumn*

\LA
\Source{Lectio v. Serm. 21.}
\switchcolumn
\EN
\Source{21st Sermon.}
\switchcolumn*

\LA
\noindent Principes persecuti sunt me gratis: et a verbis tuis trepidavit cor meum. Bene hoc Martyr dicit, quod injuste persecutionum tormenta sustineat: qui nihil rapuerit, nullum violentus oppresserit, nullius sanguinem fuderit, nullius torum putaverit esse violandum: qui nihil legibus debeat, et graviora latronum sustinere cogatur supplicia: qui loquatur juste, et non audiatur: qui loquatur plena salutis, et impugnetur: ut possit dicere, Cum loquebar illis, impugnabant me gratis. Gratis igitur persecutionem patitur, qui impugnatur sine crimine impugnatur ut noxius, cum sit in tali confessione laudabilis: impugnatur quasi veneficus, qui in nomine Domini gloriatur, cum pietas virtutum omnium fundamentum sit.\par
\switchcolumn
\EN
\noindent Princes have persecuted me I without a cause; but my heart standeth in awe of thy word. These are rightly the words of a martyr, who beareth unjustly the torments of the persecutors, who hath robbed no man, who hath violently oppressed no man, who hath shed the blood of no man, who hath imagined to defile the bed of no man, who is debtor to the laws in nothing, and who is punished more grievously than if he were a robber: who speaketh righteousness, and there is none that will hear: who speaketh salvation, and all men fight against him: who is able to say: “When I spoke unto them, they fought against me without a cause.” (Ps. cxix. 7.) They fight against him without a cause, who can lay no sin to his charge; they fight against him as an evildoer, who is by their own acknowledgment righteous: they fight against him as a warlock, who glorieth in the name of the Lord, and who doeth all things well because he doeth all things for God's sake.\par
\switchcolumn*

\LA
\begin{Resp}\Rsign Desidérium ánimæ ejus tribuísti ei Dómine, \Star{} Et voluntáte labiórum ejus non fraudásti eum.\end{Resp}
\begin{Resp}\Vsign Quóniam prævenísti eum in benedictiónibus dulcédinis: posuísti in cápite ejus corónam de lápide pretióso.\end{Resp}
\begin{Resp}\Rsign Et voluntáte labiórum ejus non fraudásti eum.\end{Resp}
\switchcolumn
\EN
\begin{Resp}\Rsign Lord, Thou hast given him his heart's desire. \Star{} And hast not withholden the request of his lips.\end{Resp}
\begin{Resp}\Vsign For Thou hast prevented him with the blessings of sweetness: Thou hast set a crown of precious stones upon his head.\end{Resp}
\begin{Resp}\Rsign And hast not withholden the request of his lips.\end{Resp}
\switchcolumn*

\LA
\LessonHead{Lectio VI}
\switchcolumn
\EN
\LessonHead{Lesson VI}
\switchcolumn*

\LA
\noindent Vere frustra impugnatur, qui apud impios et infidos impietatis arcessitur, cum fidei sit magister. Verum qui gratis impugnatur, fortis debet esse et constans, quomodo ergo subtexuit: Et a verbis tuis trepidavit cor meum? Trepidare infirmitatis est, timoris atque formidinis. Sed est étiam infirmitas ad salutem, est étiam timor sanctorum. Timete Dominum omnes sancti ejus: et, Beatus vir, qui timet Dominum. Qua ratione beatus? Quia in mandatis ejus cupit nimis.\par
\switchcolumn
\EN
\noindent They fight against him in vain who is accused of ungodliness among the ungodly and the unfaithful, because he teacheth Faith. Verily, him that is fought against without a cause it behoveth to be strong and patient. Wherefore then saith he: My heart standeth in awe of thy word? Awe is the mark of the weak, the timid, and the fearful But there is also a weakness unto salvation, there is a fear which is an holy fear. O fear the Lord, all ye His Saints. (Ps. xxxiii. 10.) And again: Blessed is the man that feareth the Lord. (Ps. cxi. i.) And wherefore is he blessed? because he delighteth greatly in His commandments.\par
\switchcolumn*

\LA
\begin{Resp}\Rsign Stola jucunditátis índuit eum Dóminus: \Star{} Et corónam pulchritúdinis pósuit super caput ejus.\end{Resp}
\begin{Resp}\Vsign Cibávit illum Dóminus pane vitæ et intelléctus: et aqua sapiéntiæ salutáris potávit illum.\end{Resp}
\begin{Resp}\Rsign Et corónam pulchritúdinis pósuit super caput ejus.\end{Resp}
\begin{Resp}\Vsign Glória Patri, et Fílio, \Star{} et Spirítui Sancto.\end{Resp}
\begin{Resp}\Rsign Et corónam pulchritúdinis pósuit super caput ejus.\end{Resp}
\switchcolumn
\EN
\begin{Resp}\Rsign The Lord hath put on him a robe of honour, \Star{} And put about his head a crown of joy.\end{Resp}
\begin{Resp}\Vsign With the bread of life and understanding hath the Lord fed him, and given him the water of health and wisdom to drink.\end{Resp}
\begin{Resp}\Rsign And put about his head a crown of joy.\end{Resp}
\begin{Resp}\Vsign Glory be to the Father, and to the Son, \Star{} and to the Holy Ghost.\end{Resp}
\begin{Resp}\Rsign And put about his head a crown of joy.\end{Resp}
\switchcolumn*

\end{paracol}

\par\needspace{10\baselineskip}\addvspace{1.2em}{\centering\small\textsc{Die 14 Julii} · \textsc{14 July}\par}
\Office{S. Bonaventuræ Episcopi Confessoris et Ecclesiæ Doctoris}{St. Bonaventure, Bishop, Confessor and Doctor of the Church}{Duplex}
\addcontentsline{toc}{subsection}{Die 14 Julii · S. Bonaventuræ Episcopi Confessoris et Ecclesiæ Doctoris \textit{— St. Bonaventure, Bishop, Confessor and Doctor of the Church}}
\begin{paracol}{2}
\LA
\RefLine{Lectiones I–III: de Scriptura occurrente.}
\switchcolumn
\EN
\RefLine{Lessons I–III: from the Scripture of the day.}
\switchcolumn*

\LA
\LessonHead{Lectio IV}
\switchcolumn
\EN
\LessonHead{Lesson IV}
\switchcolumn*

\LA
\noindent Bonaventúra, Balneorégii in Etrúria natus, a letháli morbo adhuc puer, beáti Francísci précibus, cujus religióni, si convaluísset, voto matris dicátus fúerat, evásit incólumis. Itaque adoléscens, fratrum Minórum institútum amplécti vóluit, in quo ad eam doctrinæ præstántium Alexándro de Ales magístro pervénit, ut séptimo post anno Parísiis magistérii láuream adéptus, libros Sententiárum públice summa cum laude sit interpretátus, quos étiam præcláris póstea commentáriis illustrávit. Nec sciéntiæ solum eruditióne, sed et morum integritáte vitæque innocéntia, humilitáte, mansuetúdine, terrenárum rerum contémptu et cæléstium desidério mirífice excélluit; dignus plane, qui tamquam perfectiónis exémplar haberétur, et a beáto Thoma Aquináte, cui summa caritáte conjúnctus erat, sanctus appellarétur. Is enim, cum sancti Francísci vitam illum scribéntem comperísset: Sinámus, ait, Sanctum pro Sancto laboráre.\par
\switchcolumn
\EN
\noindent Bonaventure was born at Bagnarea in Tuscany, in the year of our Lord 1221. In his infancy he was dangerously ill, and his mother made a vow that, if he recovered, she would dedicate him to the Order of Blessed Francis. While he was still a young man he entered the Order by his own wish. Under the teaching of Alexander of Hales he advanced so quickly in learning, that in seven years he lectured publicly at Paris on the Books of the Sentences, with great applause. He afterwards explained the same Books by a brilliant Commentary. After six years he was made General Master of his Order at Rome, in which office he did his duty with such wisdom and holiness as caused all men to talk of him and marvel at him.\par
\switchcolumn*

\LA
\begin{Resp}\Rsign Invéni David servum meum, óleo sancto meo unxi eum: \Star{} Manus enim mea auxiliábitur ei.\end{Resp}
\begin{Resp}\Vsign Nihil profíciet inimícus in eo, et fílius iniquitátis non nocébit ei.\end{Resp}
\begin{Resp}\Rsign Manus enim mea auxiliábitur ei.\end{Resp}
\switchcolumn
\EN
\begin{Resp}\Rsign I have found David My servant, with My holy oil have I anointed him \Star{} For My hand shall help him.\end{Resp}
\begin{Resp}\Vsign The enemy shall prevail nothing against him, nor the son of wickedness afflict him.\end{Resp}
\begin{Resp}\Rsign For My hand shall help him.\end{Resp}
\switchcolumn*

\LA
\LessonHead{Lectio V}
\switchcolumn
\EN
\LessonHead{Lesson V}
\switchcolumn*

\LA
\noindent Divíni amóris flamma succénsus, erga Christi Dómini passiónem, quam júgiter meditabátur, ac Deíparam Vírginem, cui se totum devóverat, singulári ferebátur pietátis afféctu; quem in áliis étiam verbo et exémplo excitáre, scriptísque opúsculis augére summópere stúduit. Hinc illa morum suávitas, grátia sermónis et cáritas in omnes effúsa, qua singulórum ánimos sibi arctíssime devinciébat. Quam ob rem, vix quinque et trigínta annos natus, Romæ summo ómnium consénsu generális órdinis miníster eléctus est; susceptúmque munus per duodevigínti annos admirábili prudéntia gessit ac laude sanctitátis. Plura constítuit regulári disciplínæ et amplificándo órdini utília; quem una cum áliis ordínibus mendicántibus advérsus obtrectatórum calúmnias felíciter propugnávit.\par
\switchcolumn
\EN
\noindent He was the author of many books, in which the depth of his learning and the earnestness of his godliness affect the reader and teach him at the same time. Gregory X., moved by his reputation for wisdom and holiness, created him a Cardinal and Bishop of Albano. Blessed Thomas of Aquino gave him the title of Saint even during his life -time. It fell on this wise. Thomas found him writing the Life of St. Francis, and said, Let us leave one Saint to work for the other. He departed this life at the Council of Lyons, upon the 14th day of July 1274, being of the age of 53 years, and having worked many miracles. Pope Sixtus IV. numbered him among the Saints.\par
\switchcolumn*

\LA
\begin{Resp}\Rsign Pósui adjutórium super poténtem, et exaltávi eléctum de plebe mea: \Star{} Manus enim mea auxiliábitur ei.\end{Resp}
\begin{Resp}\Vsign Invéni David servum meum, óleo sancto meo unxi eum.\end{Resp}
\begin{Resp}\Rsign Manus enim mea auxiliábitur ei.\end{Resp}
\switchcolumn
\EN
\begin{Resp}\Rsign I have laid help upon one that is mighty, and have exalted one chosen out of My people \Star{} For My hand shall help him.\end{Resp}
\begin{Resp}\Vsign I have found David My servant, with My holy oil have I anointed him.\end{Resp}
\begin{Resp}\Rsign For My hand shall help him.\end{Resp}
\switchcolumn*

\LA
\LessonHead{Lectio VI}
\switchcolumn
\EN
\LessonHead{Lesson VI}
\switchcolumn*

\LA

\switchcolumn
\EN
\Note{No English translation in the source files; the Latin is given.}
\switchcolumn*

\LA
\noindent Ad Lugdunénse concílium a beáto Gregório décimo accersítus et cardinális epíscopus Albanénsis creátus, árduis concílii rebus egrégiam navávit óperam; qua et schísmatis dissídia, compósita sunt, et ecclesiástica dógmata vindicáta. Quibus in labóribus, anno ætátis suæ quinquagésimo tértio, salútis vero millésimo ducentésimo septuagésimo quarto, summo ómnium mæróre decéssit, ab univérso concílio, ipso præsénte Románo Pontífice, fúnere honestátus. Eum Xystus quartus, plúrimis maximísque clarum miráculis, in Sanctórum númerum rétulit. Multa scripsit, in quibus summam eruditiónem cum pietátis ardóre conjúngens, lectórem docéndo movet: quare a Xysto quinto Doctóris Seráphici nómine mérito est insignítus.\par
\switchcolumn
\EN
\noindent Ad Lugdunénse concílium a beáto Gregório décimo accersítus et cardinális epíscopus Albanénsis creátus, árduis concílii rebus egrégiam navávit óperam; qua et schísmatis dissídia, compósita sunt, et ecclesiástica dógmata vindicáta. Quibus in labóribus, anno ætátis suæ quinquagésimo tértio, salútis vero millésimo ducentésimo septuagésimo quarto, summo ómnium mæróre decéssit, ab univérso concílio, ipso præsénte Románo Pontífice, fúnere honestátus. Eum Xystus quartus, plúrimis maximísque clarum miráculis, in Sanctórum númerum rétulit. Multa scripsit, in quibus summam eruditiónem cum pietátis ardóre conjúngens, lectórem docéndo movet: quare a Xysto quinto Doctóris Seráphici nómine mérito est insignítus.\par
\switchcolumn*

\LA
\begin{Resp}\Rsign Iste est, qui ante Deum magnas virtútes operátus est, et omnis terra doctrína ejus repléta est: \Star{} Ipse intercédat pro peccátis ómnium populórum.\end{Resp}
\begin{Resp}\Vsign Iste est, qui contémpsit vitam mundi, et pervénit ad cæléstia regna.\end{Resp}
\begin{Resp}\Rsign Ipse intercédat pro peccátis ómnium populórum.\end{Resp}
\begin{Resp}\Vsign Glória Patri, et Fílio, \Star{} et Spirítui Sancto.\end{Resp}
\begin{Resp}\Rsign Ipse intercédat pro peccátis ómnium populórum.\end{Resp}
\switchcolumn
\EN
\begin{Resp}\Rsign This is he which wrought great wonders before God, and the whole earth is full of his teaching \Star{} May he pray for all people, that their sins may be forgiven unto them\end{Resp}
\begin{Resp}\Vsign This is he which loved not his life in this world, and hath attained unto the kingdom of heaven.\end{Resp}
\begin{Resp}\Rsign May he pray for all people, that their sins may be forgiven unto them\end{Resp}
\begin{Resp}\Vsign Glory be to the Father, and to the Son, \Star{} and to the Holy Ghost.\end{Resp}
\begin{Resp}\Rsign May he pray for all people, that their sins may be forgiven unto them\end{Resp}
\switchcolumn*

\LA
\LessonHead{Lectio VII}
\switchcolumn
\EN
\LessonHead{Lesson VII}
\switchcolumn*

\LA
\LessonTitle{Léctio sancti Evangélii secúndum Matthǽum}
\switchcolumn
\EN
\LessonTitle{From the Holy Gospel according to Matthew}
\switchcolumn*

\LA
\Cite{Matt 5:13-19}
\switchcolumn
\EN
\Cite{Matt 5:13-19}
\switchcolumn*

\LA
\noindent In illo témpore: Dixit Jesus discípulis suis: Vos estis sal terræ. Quod si sal evanúerit, in quo saliétur? Et réliqua.\par
\switchcolumn
\EN
\noindent At that time: Jesus said unto His disciples: Ye are the salt of the earth. But if the salt have lost his savour, wherewith shall it be salted? And so on.\par
\switchcolumn*

\LA
\LessonTitle{Homilía sancti Joánnis Chrysóstomi}
\switchcolumn
\EN
\LessonTitle{Homily by St John Chrysostom, Patriarch (of Constantinople.)}
\switchcolumn*

\LA
\Source{Homilia 15 in Matthæum, sub medium}
\switchcolumn
\EN
\Source{15/A on Matth.}
\switchcolumn*

\LA
\noindent Atténdite quid díxerit, Vos estis sal terræ: per quod osténdit quam necessário ista præcípiat. Non enim de vestra, inquit, tantúmmodo vita, sed de univérso orbe vobis ratio reddénda est. Non ad duas quippe urbes, aut decem, aut vigínti, neque ad unam gentem vos mitto, sicut mittébam prophétas: sed ad omnem terram prorsus ac mare, totúmque mundum, et hunc váriis crimínibus oppréssum.\par
\switchcolumn
\EN
\noindent Consider how that the Lord saith: “Ye are the salt of the earth,” by the which figure He showeth what a necessary of life is the Gospel. By this figure, He hath us to know that they unto whom He spake have an account to render, not of their own life only, but for the whole world. Not unto two cities, saith the Lord, nor unto ten, nor unto twenty, nor unto one people, as I sent the Prophets, send I you. But I send you unto every land and sea, even unto the whole world, lying groaning, as it is, under the burden of diverse sins.\par
\switchcolumn*

\LA
\begin{Resp}\Rsign Amávit eum Dóminus, et ornávit eum: stolam glóriæ índuit eum, \Star{} Et ad portas paradísi coronávit eum.\end{Resp}
\begin{Resp}\Vsign Induit eum Dóminus lorícam fídei, et ornávit eum.\end{Resp}
\begin{Resp}\Rsign Et ad portas paradísi coronávit eum.\end{Resp}
\switchcolumn
\EN
\begin{Resp}\Rsign The Lord loved him and beautified him He clothed him with a robe of glory \Star{} And crowned him at the gates of Paradise.\end{Resp}
\begin{Resp}\Vsign The Lord hath put on him the breast-plate of faith, and hath adorned him.\end{Resp}
\begin{Resp}\Rsign And crowned him at the gates of Paradise.\end{Resp}
\switchcolumn*

\LA
\LessonHead{Lectio VIII}
\switchcolumn
\EN
\LessonHead{Lesson VIII}
\switchcolumn*

\LA
\noindent Dicéndo enim, Vos estis sal terræ, osténdit univérsam hóminum infatuátam esse natúram, et peccatórum vi corrúptam: et idcírco illas ab eis virtútes requírit, quæ máxime ad multórum salútem procurándam necessáriæ sunt atque útiles. Nam qui mansuétus est ac modéstus, et miséricors et justus, non intra se tantúmmodo hæc recte facta conclúdit, verum in aliórum quoque utilitátem præcláros hos fáciet efflúere fontes. Igitur qui corde mundo est atque pacíficus, et persecutiónem pro veritáte pátitur, nihilóminus in commúne cómmodum vitam instítuit.\par
\switchcolumn
\EN
\noindent These words, “Ye are the salt of the earth,” show unto us the whole nature of man as savourless and stinking with the strong corruption of sin. And therefore demandeth He of His Apostles such qualities as are most needful and useful to the furthering the salvation of many. He that is gentle and lowly, tender and just, shutteth not up all these good things in his own heart, but openeth these bright fountains that they may gush forth for the use of his neighbour. He whose heart is pure, and who seeketh peace, suffering persecution for the truth's sake, doth still lead a life for the good of the commonwealth.\par
\switchcolumn*

\LA
\begin{Resp}\Rsign In médio Ecclésiæ apéruit os ejus, \Star{} Et implévit eum Dóminus spíritu sapiéntiæ et intelléctus.\end{Resp}
\begin{Resp}\Vsign Jucunditátem et exsultatiónem thesaurizávit super eum.\end{Resp}
\begin{Resp}\Rsign Et implévit eum Dóminus spíritu sapiéntiæ et intelléctus.\end{Resp}
\begin{Resp}\Vsign Glória Patri, et Fílio, \Star{} et Spirítui Sancto.\end{Resp}
\begin{Resp}\Rsign Et implévit eum Dóminus spíritu sapiéntiæ et intelléctus.\end{Resp}
\switchcolumn
\EN
\begin{Resp}\Rsign In the midst of the congregation did the Lord open his mouth. \Star{} And filled him with the spirit of wisdom and understanding.\end{Resp}
\begin{Resp}\Vsign He made him rich with joy and gladness.\end{Resp}
\begin{Resp}\Rsign And filled him with the spirit of wisdom and understanding.\end{Resp}
\begin{Resp}\Vsign Glory be to the Father, and to the Son, \Star{} and to the Holy Ghost.\end{Resp}
\begin{Resp}\Rsign And filled him with the spirit of wisdom and understanding.\end{Resp}
\switchcolumn*

\LA
\LessonHead{Lectio IX}
\switchcolumn
\EN
\LessonHead{Lesson IX}
\switchcolumn*

\LA
\noindent Ne ígitur putétis, inquit, ad lévia vos ducéndos esse certámina, neque exiguárum rerum vobis ineúndam esse ratiónem. Vos estis sal terræ. Quid ígitur? Ipsine putrefácta medicáti sunt? Nequáquam: neque enim fíeri potest, ut ea quæ jam corrúpta sunt, salis perfricatióne reparéntur. Non ergo hoc fecérunt, sed ante renováta, sibíque trádita, atque ab illa jam putrédine liberáta, aspergébant sale, et in ea novitáte conservábant, quam a Dómino suscéperant. Liberáre quippe a putrédine peccatórum, Christi virtútis est: ut autem ad illa íterum non revertántur, Apostolórum curæ est ac labóris.\par
\switchcolumn
\EN
\noindent Think not, saith the Lord, that the struggle is easy whereunto ye shall be led, neither shall your reckoning be of light matters. Ye are the salt of the earth. Have ye then salted that which is corrupted? Nay, for it is impossible that that which is once corrupted can be made sound again by the rubbing it with salt. This it is not asked of them to do. But their work is to sprinkle with salt, and to keep fresh thereafter, such things as the Lord hath given over into their charge, and which He Himself hath made new, and freed from all taint, before giving them. To make sound after the corruption of sin, is the work of Christ's power alone; to preserve from falling away again, is the duty and the toil commanded to the Apostles.\par
\switchcolumn*

\LA
\TeDeum{Te Deum}
\switchcolumn
\EN
\TeDeum{Te Deum}
\switchcolumn*

\end{paracol}

\par\needspace{10\baselineskip}\addvspace{1.2em}{\centering\small\textsc{Die 15 Julii} · \textsc{15 July}\par}
\Office{S. Henrici Imperatoris Confessoris}{St. Henry the Emperor, Confessor}{Semiduplex}
\addcontentsline{toc}{subsection}{Die 15 Julii · S. Henrici Imperatoris Confessoris \textit{— St. Henry the Emperor, Confessor}}
\begin{paracol}{2}
\LA
\RefLine{Lectiones I–III: de Scriptura occurrente.}
\switchcolumn
\EN
\RefLine{Lessons I–III: from the Scripture of the day.}
\switchcolumn*

\LA
\RefLine{Lectiones VII–IX: ex Commúni Confessoris non pontificis (C5).}
\switchcolumn
\EN
\RefLine{Lessons VII–IX: from the Common of a Confessor not a Bishop (C5).}
\switchcolumn*

\LA
\LessonHead{Lectio IV}
\switchcolumn
\EN
\LessonHead{Lesson IV}
\switchcolumn*

\LA
\noindent Henrícus, cognoménto Pius, e duce Baváriæ rex Germániæ ac póstmodum Romanórum imperátor, temporális regni non conténtus angústiis, pro adipiscénda immortalitátis coróna sédulam ætérno Regi exhíbuit servitútem. Adépto enim império, religióni amplificándæ studióse incúmbens, ecclésias ab infidélibus destrúctas magnificéntius reparávit, plurimísque largitiónibus et prǽdiis locupletávit. Monastéria aliáque loca pia vel ipse ædificávit, vel assignátis redítibus auxit. Episcopátum Bambergénsem, hereditáriis ópibus fundátum, beáto Petro Romanóque Pontífici vectigálem fecit. Benedíctum octávum, a quo impérii corónam accéperat, prófugum excépit, suǽque Sedi restítuit.\par
\switchcolumn
\EN
\noindent Henry II., surnamed the Pious, became successively Duke of Bavaria, (in the year 995,) King of Germany, (in 1002,) and the Emperor of the Romans, (in 1014.) His hope soared beyond the short enjoyment of a fleeting kingdom, and he aimed at the possession of an unfading crown by living as the loyal servant of the Eternal King. After he became Emperor, he earnestly set himself to the furtherance of the cause of godliness. He restored with new splendour the Churches which had been ruined by the unbelievers, and enriched them with many offerings and possessions. Monasteries and other godly places he either built himself, or endowed them with allowances. He founded out of his own family inheritance the Bishopric of Bamberg, and made it tributary to Blessed Peter and to the Bishop of Rome. When Benedict VIII, who had set on his head the Imperial crown, was an exile, he hospitably received him, and afterwards restored him to his See.\par
\switchcolumn*

\LA
\begin{Resp}\Rsign Honéstum fecit illum Dóminus, et custodívit eum ab inimícis, et a seductóribus tutávit illum: \Star{} Et dedit illi claritátem ætérnam.\end{Resp}
\begin{Resp}\Vsign Justum dedúxit Dóminus per vias rectas, et osténdit illi regnum Dei.\end{Resp}
\begin{Resp}\Rsign Et dedit illi claritátem ætérnam.\end{Resp}
\switchcolumn
\EN
\begin{Resp}\Rsign The Lord made him honourable, and defended him from his enemies, and kept him safe from those that lay in wait for him; \Star{} And gave him perpetual glory.\end{Resp}
\begin{Resp}\Vsign He went down with him into the pit, and left him not in bonds.\end{Resp}
\begin{Resp}\Rsign And gave him perpetual glory.\end{Resp}
\switchcolumn*

\LA
\LessonHead{Lectio V}
\switchcolumn
\EN
\LessonHead{Lesson V}
\switchcolumn*

\LA
\noindent In Cassinénsi monastério, gravi deténtus infirmitáte, a sancto Benedícto, insígni miráculo, sanátus est. Románam Ecclésiam amplíssimo diplómate munerátus, eídem tuéndæ bellum advérsus Græcos suscépit, et Apúliam, diu ab illis posséssam, recuperávit. Nihil sine précibus ággredi sólitus, Angelum Dómini sanctósque Mártyres tuteláres pro se pugnántes ante áciem intérdum vidit. Divína autem protéctus ope, bárbaras natiónes précibus magis quam armis expugnávit. Pannóniam adhuc infidélem, trádita Stéphano regi soróre sua in uxórem, eóque baptizáto, ad Christi fidem perdúxit. Virginitátem raro exémplo matrimónio junxit, sanctámque Cunegúndam, cónjugem suam, propínquis ejus, morti próximus, illibátam restítuit.\par
\switchcolumn
\EN
\noindent Then he was struck down with a grievous sickness in the Monastery of Monte Cassino, he was healed by an evident miracle through the intercession of St Benedict. He was a princely benefactor to the Church of Rome, for the defence of which he entered into a war against the Greeks, and took again from them the province of Apulia, which they had long possessed. He never undertook anything until he had made it a subject of prayer. And in battle he once saw the Angel of the Lord and the Holy Martyrs (Laurence, George, and Adrian) his patrons under whose protection he had placed his army, fighting for him in front of his line. With the help of God, he prevailed against the tribes of savages more by prayer than by arms. He gave his sister in marriage to King Stephen of Hungary, whom he induced to be baptized, and so brought all that country to believe in Christ. His marriage with the holy maiden Cunegunda is one of the rare instances of the union of two virgins. When he drew near to death, he gave her back inviolate to her kinsfolk.\par
\switchcolumn*

\LA
\begin{Resp}\Rsign Amávit eum Dóminus, et ornávit eum: stolam glóriæ índuit eum, \Star{} Et ad portas paradísi coronávit eum.\end{Resp}
\begin{Resp}\Vsign Índuit eum Dóminus lorícam fídei, et ornávit eum.\end{Resp}
\begin{Resp}\Rsign Et ad portas paradísi coronávit eum.\end{Resp}
\switchcolumn
\EN
\begin{Resp}\Rsign The Lord loved him and beautified him; He clothed him with a robe of glory, \Star{} And crowned him at the gates of Paradise.\end{Resp}
\begin{Resp}\Vsign The Lord hath put on him the breast-plate of faith, and hath adorned him.\end{Resp}
\begin{Resp}\Rsign And crowned him at the gates of Paradise.\end{Resp}
\switchcolumn*

\LA
\LessonHead{Lectio VI}
\switchcolumn
\EN
\LessonHead{Lesson VI}
\switchcolumn*

\LA
\noindent Dénique, rebus ómnibus, quæ ad impérii honórem et utilitátem pertinébant, summa prudéntia dispósitis, et illústribus per Gálliam, Itáliam et Germániam religiósæ munificéntiæ vestígiis passim relíctis; postquam heróicæ virtútis suavíssimum odórem longe latéque diffúderat, sanctitáte quam sceptro clárior, ad regni cæléstis prǽmia, consummátis vitæ labóribus, a Dómino vocátus est anno salútis millésimo vigésimo quarto. Cujus corpus in ecclésia beatórum Apostolórum Petri et Pauli Bambergæ cónditum fuit; statímque ad ejus túmulum multa mirácula, Deo ipsum glorificánte, patráta sunt. Quibus póstea rite probátis, Eugénius tértius Sanctórum número illum adscrípsit.\par
\switchcolumn
\EN
\noindent He managed with great wisdom whatever could tend to the honour and usefulness of the Empire. He left in France? Italy, and Germany, splendid monuments of his godly munificence. The perfume of his saintly life spread its sweetness far and wide, and the glory of his holiness outshone the splendour of his crown. When the work of his life was done, he was called by the Lord to the possession of an eternal kingdom (on the 14th day of July,) in the year of salvation 1024. His body was buried in the Church of the Blessed Apostles Peter and Paul at Bamberg, and God glorified him by the miracles which began forthwith to take place at his grave. The same being duly proved, Eugene III numbered him among the Saints.\par
\switchcolumn*

\LA
\begin{Resp}\Rsign Iste homo perfécit ómnia quæ locútus est ei Deus, et dixit ad eum: Ingrédere in réquiem meam: \Star{} Quia te vidi justum coram me ex ómnibus géntibus.\end{Resp}
\begin{Resp}\Vsign Iste est, qui contémpsit vitam mundi, et pervénit ad cæléstia regna.\end{Resp}
\begin{Resp}\Rsign Quia te vidi justum coram me ex ómnibus géntibus.\end{Resp}
\begin{Resp}\Vsign Glória Patri, et Fílio, \Star{} et Spirítui Sancto.\end{Resp}
\begin{Resp}\Rsign Quia te vidi justum coram me ex ómnibus géntibus.\end{Resp}
\switchcolumn
\EN
\begin{Resp}\Rsign This is he which did according unto all that God commanded him; and God said unto him: Enter thou into My rest, \Star{} For thee have I seen righteous before Me among all people.\end{Resp}
\begin{Resp}\Vsign This is he which loved not his life in this world, and is come unto an everlasting kingdom.\end{Resp}
\begin{Resp}\Rsign For thee have I seen righteous before Me among all people.\end{Resp}
\begin{Resp}\Vsign Glory be to the Father, and to the Son, \Star{} and to the Holy Ghost.\end{Resp}
\begin{Resp}\Rsign For thee have I seen righteous before Me among all people.\end{Resp}
\switchcolumn*

\end{paracol}

\par\needspace{10\baselineskip}\addvspace{1.2em}{\centering\small\textsc{Die 16 Julii} · \textsc{16 July}\par}
\Office{In Commemoratione Beatæ Mariæ Virginis de Monte Carmelo}{In Commemoratione Beatæ Mariæ Virginis de Monte Carmelo}{Duplex majus}
\addcontentsline{toc}{subsection}{Die 16 Julii · In Commemoratione Beatæ Mariæ Virginis de Monte Carmelo \textit{— In Commemoratione Beatæ Mariæ Virginis de Monte Carmelo}}
\begin{paracol}{2}
\LA
\RefLine{Lectiones I–III: ex In Festis Beatae Mariae Virginis (C11).}
\switchcolumn
\EN
\RefLine{Lessons I–III: from the In Festis Beatae Mariae Virginis (C11).}
\switchcolumn*

\LA
\RefLine{Lectiones VII–IX: ex In Festis Beatae Mariae Virginis (C11).}
\switchcolumn
\EN
\RefLine{Lessons VII–IX: from the In Festis Beatae Mariae Virginis (C11).}
\switchcolumn*

\LA
\LessonHead{Lectio IV}
\switchcolumn
\EN
\LessonHead{Lesson IV}
\switchcolumn*

\LA
\noindent Cum sacra Pentecostes die Apóstoli, cælitus afflati, variis linguis loqueréntur, et, invocato augustíssimo Jesu nómine, mira multa patrarent; viri plurimi (ut fertur), qui vestigiis sanctórum prophetárum Elíæ ac Eliséi instíterant, et Joánnis Baptistæ præconio ad Christi advéntum comparáti fuerant, rerum veritáte perspecta atque probata, evangelicam fidem conféstim amplexáti sunt, ac peculiari quodam afféctu beatíssimam Vírginem (cujus colloquiis ac familiaritate feliciter frui potuére) adeo venerari cœpérunt, ut primi ómnium in eo montis Carmeli loco, ubi Elías olim ascendentem nébulam, Vírginis typo insignem, conspexerat, eidem puríssimæ Virgini sacellum construxerint.\par
\switchcolumn
\EN
\noindent Here is a story to the effect that many men who had kept a tradition of the holy Prophets Elijah and Elisha, were made ready by the preaching of John the Baptist to hail the coming of the Messiah, and that, when the Apostles having been filled with the Spirit upon the holy day of Pentecost, spake with diverse tongues and worked miracles by calling upon the Name of Jesus which is above every other name, these men, seeing and being assured of the truth, straightway embraced the faith of the Gospel, and that on account of their singular love toward the Most Blessed Virgin, whose conversation and friendship they were able to enjoy, they paid her the respect of building her a little Chapel, the first which was ever raised in honour of this same most pure Maiden, and which stood upon that part of Mount Carmel whence the servant of Elijah had in old days espied that manifest type of the Virgin, the little cloud like a man's hand, arising out of the sea. (3 Kings xviii. 44.)\par
\switchcolumn*

\LA
\begin{Resp}\Rsign Sicut cedrus exaltáta sum in Líbano, et sicut cypréssus in monte Sion: quasi myrrha elécta, \Star{} Dedi suavitátem odóris.\end{Resp}
\begin{Resp}\Vsign Et sicut cinnamómum et bálsamum aromatízans.\end{Resp}
\begin{Resp}\Rsign Dedi suavitátem odóris.\end{Resp}
\switchcolumn
\EN
\begin{Resp}\Rsign I was exalted like a cedar in Lebanon, and as a cypress-tree upon Mount Zion. Like the best myrrh \Star{} I yielded a pleasant odour.\end{Resp}
\begin{Resp}\Vsign Like cinnamon and sweet balsam.\end{Resp}
\begin{Resp}\Rsign I yielded a pleasant odour.\end{Resp}
\switchcolumn*

\LA
\LessonHead{Lectio V}
\switchcolumn
\EN
\LessonHead{Lesson V}
\switchcolumn*

\LA
\noindent Ad novum ergo sacellum sæpe quotídie conveniéntes, rítibus piis, precatiónibus ac laudibus beatíssimam Vírginem, velut singularem ordinis tutelam, colebant. Quam ob rem fratres beátæ Maríæ de Monte Carmelo passim ab ómnibus appellari cœpérunt; eumque titulum summi Pontifices non modo confírmarunt, sed et indulgéntias peculiares iis, qui eo titulo vel ordinem vel fratres síngulos nuncuparent, concessére. Nec vero nomenclatúram tantum magnificentíssima Virgo tríbuit et tutelam, verum et insigne sacri scapularis; quod beato Simoni Anglico præbuit, ut cælésti hac veste ordo ille sacer dignoscerétur, et a malis ingruéntibus protegerétur. Ac demum, cum olim in Europa ordo esset ignotus, et ob id apud Honórium tertium non pauci pro illíus exstinctióne instarent, ástitit Honorio noctu piíssima Virgo Maria, planeque jussit ut institutum et hómines benígne complecterétur.\par
\switchcolumn
\EN
\noindent Now this new Chapel they repaired oftentimes day by day, and in their sacred ceremonies, prayers, and praises, honoured the Most Blessed Virgin as the particular Guardian of their Congregation. For this reason they came to be everywhere called the Brethren of Blessed Mary, of Mount Carmel, and the Supreme Pontiffs have not only confirmed to them the right to use this name, but have granted particular indulgences to all those who so call either the Order itself, or any particular member thereof. Her name and protection are not the only gifts which the most bountiful Virgin hath given them yea, she hath given them the badge of the Holy Scapular, which she delivered to the Blessed Englishman Simon Stock, even an heavenly garment whereby this Holy Order is marked, and harnessed against all assaults. Moreover, in old times, when this Order was unknown in Europe, and not a few were instant with Honorius III. to put an end to it, the most gracious Virgin Mary appeared by night to the said Honorius, and flatly commanded him to show kindness to the Order and to the men belonging thereto.\par
\switchcolumn*

\LA
\begin{Resp}\Rsign Quæ est ista quæ procéssit sicut sol, et formósa tamquam Jerúsalem? \Star{} Vidérunt eam fíliæ Sion, et beátam dixérunt, et regínæ laudavérunt eam.\end{Resp}
\begin{Resp}\Vsign Et sicut dies verni circúmdabant eam flores rosárum et lília convállium.\end{Resp}
\begin{Resp}\Rsign Vidérunt eam fíliæ Sion, et beátam dixérunt, et regínæ laudavérunt eam.\end{Resp}
\switchcolumn
\EN
\begin{Resp}\Rsign Who is this that cometh up like the sun? This, comely as Jerusalem? \Star{} The daughters of Zion saw her, and called her blessed: the queens also, and they praised her.\end{Resp}
\begin{Resp}\Vsign And about her it was as the flower of roses in the spring of the year, and lilies of the valleys.\end{Resp}
\begin{Resp}\Rsign The daughters of Zion saw her, and called her blessed the queens also, and they praised her.\end{Resp}
\switchcolumn*

\LA
\LessonHead{Lectio VI}
\switchcolumn
\EN
\LessonHead{Lesson VI}
\switchcolumn*

\LA
\noindent Non in hac tantum sæculo ordinem sibi tam acceptum multis prærogativis beatíssima Virgo insignívit; verum et in alio (cum ubíque et poténtia et misericórdia plurimum valeat) fílios in scapularis societátem relatos, qui abstinéntiam modicam precesque paucas eis præscriptas frequentarunt, ac pro sui status ratióne castitátem coluérunt, materno plane afféctu, dum igne purgatorii expiántur, solari ac in cælestem pátriam obtentu suo quantócius pie creditur efferre. Tot ergo tantisque beneficiis ordo cumulátus, solemnem beatíssimæ Vírginis Commemoratiónem, ritu perpetuo ad ejusdem Vírginis glóriam quotannis celebrandam, instituit.\par
\switchcolumn
\EN
\noindent Many godly persons believe that it is not in this world only that the most Blessed Virgin hath marked with her favour this Order which pleaseth her so well, but that in the next world, where her power and mercy have a freer scope than here, they who belong to the Guild of the Scapular, who have practised an easy abstinence, have been regular in reciting a few prayers enjoined to them, and have kept chastity according to their state of life, are comforted by her motherly love while they are being cleansed in the fire of Purgatory, and by her help are borne forward towards their home in heaven more quickly than others. The Order loaded with so many and so great gifts, hath instituted a solemn Commemoration of the Most Blessed Virgin, to be made year after year, in perpetual observance, for the glory of the same Virgin.\par
\switchcolumn*

\LA
\begin{Resp}\Rsign Ornátam monílibus fíliam Jerúsalem Dóminus concupívit: \Star{} Et vidéntes eam fíliæ Sion, beatíssimam prædicavérunt, dicéntes: * Unguéntum effúsum nomen tuum.\end{Resp}
\begin{Resp}\Vsign Astitit regína a dextris tuis in vestítu deauráto, circúmdata varietáte.\end{Resp}
\begin{Resp}\Rsign Et vidéntes eam fíliæ Sion, beatíssimam prædicavérunt, dicéntes:\end{Resp}
\begin{Resp}\Vsign Glória Patri, et Fílio, \Star{} et Spirítui Sancto.\end{Resp}
\begin{Resp}\Rsign Unguéntum effúsum nomen tuum.\end{Resp}
\switchcolumn
\EN
\begin{Resp}\Rsign When the Lord beheld the daughter of Jerusalem adorned with her jewels, He greatly desired her beauty \Star{} And when the daughters of Zion saw her, they cried out that she was most blessed, saying thy name is as ointment poured forth.\end{Resp}
\begin{Resp}\Vsign Upon thy right hand did stand the Queen in a vesture of gold wrought about with diverse colours.\end{Resp}
\begin{Resp}\Rsign And when the daughters of Zion saw her, they cried out that she was most blessed, saying:\end{Resp}
\begin{Resp}\Vsign Glory be to the Father, and to the Son, \Star{} and to the Holy Ghost.\end{Resp}
\begin{Resp}\Rsign Thy name is as ointment poured forth.\end{Resp}
\switchcolumn*

\end{paracol}

\par\needspace{10\baselineskip}\addvspace{1.2em}{\centering\small\textsc{Die 17 Julii} · \textsc{17 July}\par}
\Office{S. Alexii Confessoris}{St. Alexis, Confessor}{Semiduplex}
\addcontentsline{toc}{subsection}{Die 17 Julii · S. Alexii Confessoris \textit{— St. Alexis, Confessor}}
\begin{paracol}{2}
\LA
\RefLine{Lectiones I–III: de Scriptura occurrente.}
\switchcolumn
\EN
\RefLine{Lessons I–III: from the Scripture of the day.}
\switchcolumn*

\LA
\LessonHead{Lectio IV}
\switchcolumn
\EN
\LessonHead{Lesson IV}
\switchcolumn*

\LA
\noindent Aléxius, Romanórum nobilíssimus, propter exímium Jesu Christi amórem prima nocte nuptiárum peculiari Dei mónitu relínquens intactam sponsam, illustrium orbis terræ ecclesiárum peregrinatiónem suscépit. Quibus in itinéribus cum ignotus septémdecim annos fuísset, aliquándo apud Edessam, Syriæ urbem, per imáginem sanctíssimæ Maríæ Vírginis, ejus nómine divulgato, inde navi discessit. Ad portum Romanum appulsus, a patre suo tamquam alienus pauper hospítio accipitur; apud quem, ómnibus incógnitus, cum decem et septem annos vixísset, relícto scripto sui nóminis, sánguinis, ac totius vitæ cursus, migrávit in cælum, Innocentio primo summo Pontifice.\par
\switchcolumn
\EN
\noindent Alexis was a member of one of the noblest Roman families. Through his exceeding great love for Jesus Christ, he received a particular command from God to leave his bride untouched upon his wedding night, and to undertake a pilgrimage to the most famous Churches of the world. For seventeen years he remained occupied in these journeys and utterly unknown. At the end of that time, his name was spoken from an image of the most holy Virgin Mary in the city of Edessa, in Syria, and when he found himself recognised he took ship from thence. He landed at Porto near Rome, and fared to the house of his own father, who gave him shelter as a strange beggar. He lived there unrecognised by any for seventeen years more, and then passed away to heaven, in the time of Pope Innocent I. He left behind him a writing giving his name, family, and the story of his life.\par
\switchcolumn*

\LA
\begin{Resp}\Rsign Honéstum fecit illum Dóminus, et custodívit eum ab inimícis, et a seductóribus tutávit illum: \Star{} Et dedit illi claritátem ætérnam.\end{Resp}
\begin{Resp}\Vsign Justum dedúxit Dóminus per vias rectas, et osténdit illi regnum Dei.\end{Resp}
\begin{Resp}\Rsign Et dedit illi claritátem ætérnam.\end{Resp}
\switchcolumn
\EN
\begin{Resp}\Rsign The Lord made him honourable, and defended him from his enemies, and kept him safe from those that lay in wait for him; \Star{} And gave him perpetual glory.\end{Resp}
\begin{Resp}\Vsign He went down with him into the pit, and left him not in bonds.\end{Resp}
\begin{Resp}\Rsign And gave him perpetual glory.\end{Resp}
\switchcolumn*

\LA
\LessonHead{Lectio V}
\switchcolumn
\EN
\LessonHead{Lesson V}
\switchcolumn*

\LA
\LessonTitle{Ex libro Moralium sancti Gregorii Papæ.}
\switchcolumn
\EN
\LessonTitle{From the Book of Moral Reflections upon Job, written by Pope St. Gregory the Great}
\switchcolumn*

\LA
\Source{Lib. 10. cap. 16. in cap. 12. Job.}
\switchcolumn
\EN
\Source{\{Bk. x, Chap. xvi, on Job xii.)}
\switchcolumn*

\LA
\noindent Deridétur justi simplícitas. Hujus mundi sapiéntia est, cor machinatiónibus tégere, sensum verbis veláre: quæ falsa sunt, vera osténdere; quæ vera sunt, falsa demonstráre. Hæc nimírum prudéntia usu a juvénibus scitur hæc a púeris prétio díscitur; hanc qui sciunt, céteros despiciéndo supérbiunt: hanc qui nésciunt, subjécti et tímidi in áliis mirántur: quia ab eis hæc éadem duplicitátis iníquitas nómine palliáta dilígitur, dum mentis pervérsitas urbánitas vocátur. Hæc sibi obsequéntibus præcipit honórum cúlmina quærere, adépta temporális glóriæ vanitáte gaudére, irrogáta ab áliis mala multiplícius réddere: cum vires súppetunt, nullis resisténtibus cédere: cum virtútis possibílitas deest, quidquid explére per malítiam non valent, hoc in pacífica bonitáte simuláre.\par
\switchcolumn
\EN
\noindent The simplicity of the righteous is made a subject of derision. The wisdom of this world hideth our true feelings by artifice, and useth language to conceal our thoughts this is the wisdom which demonstrated the truth of falsehood, and showeth the falsehood of the truth. This kind of shrewdness the young acquire by practice, and children pay for the learning it. Those who are good at this look down upon their neighbours those who are bad at it are humble and timid, and wonder at it in others they regard this astuteness too, wrong though it be, with wistful admiration, under softened epithets. Unstraightforwardness is called good breeding. The principles of the world teach those who entertain them, to try and rise to distinction, and when they have attained the bubble of glory which is so soon to pass away, to feel it sweet to have at their feet them on whom they may wreak rich revenge. These principles teach a man, as long as he is strong enough, to give way to nobody else, and, if he hath no chance by force, to try and attain his object by diplomacy.\par
\switchcolumn*

\LA
\begin{Resp}\Rsign Amávit eum Dóminus, et ornávit eum: stolam glóriæ índuit eum, \Star{} Et ad portas paradísi coronávit eum.\end{Resp}
\begin{Resp}\Vsign Índuit eum Dóminus lorícam fídei, et ornávit eum.\end{Resp}
\begin{Resp}\Rsign Et ad portas paradísi coronávit eum.\end{Resp}
\switchcolumn
\EN
\begin{Resp}\Rsign The Lord loved him and beautified him; He clothed him with a robe of glory, \Star{} And crowned him at the gates of Paradise.\end{Resp}
\begin{Resp}\Vsign The Lord hath put on him the breast-plate of faith, and hath adorned him.\end{Resp}
\begin{Resp}\Rsign And crowned him at the gates of Paradise.\end{Resp}
\switchcolumn*

\LA
\LessonHead{Lectio VI}
\switchcolumn
\EN
\LessonHead{Lesson VI}
\switchcolumn*

\LA
\noindent At contra, sapiéntia justórum est: nil per ostensiónem fíngere, sensum verbis aperíre, vera ut sunt dilígere, falsa devitáre; bona gratis exhibére, mala libéntius toleráre quam fácere; nullam injúriæ ultiónem quærere, pro veritáte contuméliam lucrum putáre. Sed hæc justórum simplícitas deridétur; quia ab hujus mundi sapiéntibus puritátis virtus fatúitas créditur. Omne enim quod innocénter ágitur, ab eis proculdúbio stultum putátur; et quidquid in ópere véritas ápprobat, carnáli sapiéntiæ fátuum sonat. Quid namque stúltius vidétur mundo, quam mentem verbis osténdere, nil cállida machinatióne simuláre, nullas injúriis contumélias réddere, pro maledicéntibus oráre, paupertátem quærere, posséssa relínquere, rapiénti non resístere, percutiénti álteram maxíllam præbére?\par
\switchcolumn
\EN
\noindent The wisdom of the righteous is the contrary of all this. They seek to avoid deception, to give their thoughts a clear expression in their words, to love the truth because it is the truth, to avoid falsehood, and rather to suffer than to inflict evil. Such are they who seek not to avenge themselves for wrong, and deem it gain to be despised for the truth's sake. This their simplicity is made a subject of derision, for such as are wise in this world believe the purity of their virtue to be simple foolery. Whatsoever is done innocently, they consider without doubt stupid. Such works as the truth approveth are idiotic, when tried by carnal standards of wisdom. After all, what stupider thing is there in this world than to express our real thoughts in our words, to keep nothing quiet by skilful tact, to repay no injuries, to pray for them which curse us, to seek poverty, to give up property, to strive not with such as take from us, to turn the other cheek to the smiter?\par
\switchcolumn*

\LA
\begin{Resp}\Rsign Iste homo perfécit ómnia quæ locútus est ei Deus, et dixit ad eum: Ingrédere in réquiem meam: \Star{} Quia te vidi justum coram me ex ómnibus géntibus.\end{Resp}
\begin{Resp}\Vsign Iste est, qui contémpsit vitam mundi, et pervénit ad cæléstia regna.\end{Resp}
\begin{Resp}\Rsign Quia te vidi justum coram me ex ómnibus géntibus.\end{Resp}
\begin{Resp}\Vsign Glória Patri, et Fílio, \Star{} et Spirítui Sancto.\end{Resp}
\begin{Resp}\Rsign Quia te vidi justum coram me ex ómnibus géntibus.\end{Resp}
\switchcolumn
\EN
\begin{Resp}\Rsign This is he which did according unto all that God commanded him; and God said unto him: Enter thou into My rest, \Star{} For thee have I seen righteous before Me among all people.\end{Resp}
\begin{Resp}\Vsign This is he which loved not his life in this world, and is come unto an everlasting kingdom.\end{Resp}
\begin{Resp}\Rsign For thee have I seen righteous before Me among all people.\end{Resp}
\begin{Resp}\Vsign Glory be to the Father, and to the Son, \Star{} and to the Holy Ghost.\end{Resp}
\begin{Resp}\Rsign For thee have I seen righteous before Me among all people.\end{Resp}
\switchcolumn*

\LA
\LessonHead{Lectio VII}
\switchcolumn
\EN
\LessonHead{Lesson VII}
\switchcolumn*

\LA
\LessonTitle{Léctio sancti Evangélii secúndum Matthǽum.}
\switchcolumn
\EN
\LessonTitle{From the Holy Gospel according to Matthew}
\switchcolumn*

\LA
\Cite{Matt 19:27-29}
\switchcolumn
\EN
\Cite{Matt 19:27-29}
\switchcolumn*

\LA
\noindent In illo témpore: Dixit Petrus ad Jesum: Ecce nos relíquimus ómnia, et secúti sumus te: quid ergo erit nobis? Et réliqua.\par
\switchcolumn
\EN
\noindent At that time, Peter said unto Jesus: Behold, we have forsaken all, and followed thee what shall we have therefore? And so on.\par
\switchcolumn*

\LA
\LessonTitle{Homilía sancti Hierónymi Presbýteri.}
\switchcolumn
\EN
\LessonTitle{Homily by St. Jerome, Priest (at Bethlehem.)}
\switchcolumn*

\LA
\Source{Lib. 3. in Matth. cap. 19.}
\switchcolumn
\EN
\Source{Bk. iii. on Matth xix.}
\switchcolumn*

\LA
\noindent Grandis fidúcia! Petrus piscátor erat, dives non fúerat, cibos manu et arte quærébat: et tamen lóquitur confidénter: Relíquimus ómnia. Et quia non súfficit tantum relínquere, jungit quod perféctum est: Et secúti sumus te. Fécimus quod jussísti: quid ígitur nobis dabis prǽmii? Jesus autem dixit illis: Amen dico vobis, quod vos qui secúti estis me, in regeneratióne, cum séderit Fílius hóminis in sede majestátis suæ, sedébitis et vos super sedes duódecim, judicántes duódecim tribus Ísraël. Non dixit: Qui reliquístis ómnia; hoc enim et Crates fecit philósophus, et multi álii divítias contempsérunt: sed, Qui secúti estis me: quod próprie Apostolórum est, atque credéntium.\par
\switchcolumn
\EN
\noindent Peter was a fisherman, he was not rich, he earned his bread by his hand and skill, and nevertheless he is thus bold, and saith confidently: We have forsaken all. And because it sufficeth not to forsake only, he addeth that which to do is to be perfect: and followed thee. We have done that which Thou hast commanded us, what reward therefore wilt Thou give us? And Jesus said unto them Amen I say unto you, that ye which have followed Me, in the regeneration, when the Son of Man shall sit in the throne of His glory, ye also shall sit upon twelve thrones, judging the twelve tribes of Israel. He said not, Ye which have forsaken all, for this did even Crates the philosopher, and they which have set nothing by riches are many, but, Ye which have followed Me. This did the Apostles, and this do believers do.\par
\switchcolumn*

\LA
\begin{Resp}\Rsign Iste est qui ante Deum magnas virtútes operátus est, et de omni corde suo laudávit Dóminum: \Star{} Ipse intercédat pro peccátis ómnium populórum.\end{Resp}
\begin{Resp}\Vsign Ecce homo sine queréla, verus Dei cultor, ábstinens se ab omni ópere malo, et pérmanens in innocéntia sua.\end{Resp}
\begin{Resp}\Rsign Ipse intercédat pro peccátis ómnium populórum.\end{Resp}
\switchcolumn
\EN
\begin{Resp}\Rsign This is he which wrought great wonders before God, and praised the Lord with all his heart. \Star{} May he pray for all people, that their sins may be forgiven unto them.\end{Resp}
\begin{Resp}\Vsign Behold a man without blame, a worshipper of God in truth, keeping himself clean from every evil work, and abiding still in his innocency.\end{Resp}
\begin{Resp}\Rsign May he pray for all people, that their sins may be forgiven unto them.\end{Resp}
\switchcolumn*

\LA
\LessonHead{Lectio VIII}
\switchcolumn
\EN
\LessonHead{Lesson VIII}
\switchcolumn*

\LA
\noindent In regeneratióne, cum séderit Fílius hóminis in sede majestátis suæ (quando et mórtui de corruptióne resúrgent incorrúpti) sedébitis et vos in sóliis judicántium, condemnántes duódecim tribus Ísraël: quia vobis credéntibus, illi crédere noluérunt. Et omnis qui relíquerit domum, vel fratres, aut soróres, aut patrem, aut matrem, aut uxórem, aut fílios, aut agros propter nomen meum, céntuplum accípiet, et vitam ætérnam possidébit. Locus iste cum illa senténtia cóngruit, in qua Salvátor lóquitur: Non veni pacem míttere, sed gládium. Veni enim separáre hóminem a patre suo, et matrem a fília, et nurum a socru: et inimíci hóminis doméstici ejus. Qui ergo propter fidem Christi, et prædicatiónem Evangélii, omnes afféctus contémpserint, atque divítias et sǽculi voluptátes: isti céntuplum recípient, et vitam ætérnam possidébunt.\par
\switchcolumn
\EN
\noindent In the regeneration, when the Son of Man shall sit in the throne of His glory, and when the dead shall rise again from corruption incorruptible, (i Cor. xv. 53,) ye also shall sit upon twelve thrones of judgment, condemning the twelve tribes of Israel, because, when ye believed in Me, they would not. (John iii. 18.) And every one that hath forsaken houses, or brethren, or sisters, or father, or mother, or wife, or children, or lands, for My Name's sake, shall receive an hundredfold, and shall inherit everlasting life. This place agreeth well with that other where the Saviour saith I came not to send peace, but a sword. For I am come to set a man at variance against his father, and the daughter against her mother, and the daughter-in-law against her mother-in-law; and a man's foes shall be they of his own household. (Matth. x. 34.) Every one, therefore, that hath set no store by affection, and riches, and the pleasures of the world, for Christ's faith's sake, and the preaching of the Gospel, shall receive an hundred-fold, and shall inherit everlasting life.\par
\switchcolumn*

\LA
\begin{Resp}\Rsign Sint lumbi vestri præcíncti, et lucérnæ ardéntes in mánibus vestris: \Star{} Et vos símiles homínibus exspectántibus dóminum suum, quando revertátur a núptiis.\end{Resp}
\begin{Resp}\Vsign Vigiláte ergo, quia nescítis qua hora Dóminus vester ventúrus sit.\end{Resp}
\begin{Resp}\Rsign Et vos símiles homínibus exspectántibus dóminum suum, quando revertátur a núptiis.\end{Resp}
\begin{Resp}\Vsign Glória Patri, et Fílio, \Star{} et Spirítui Sancto.\end{Resp}
\begin{Resp}\Rsign Et vos símiles homínibus exspectántibus dóminum suum, quando revertátur a núptiis.\end{Resp}
\switchcolumn
\EN
\begin{Resp}\Rsign Let your loins be girded about, and your lights burning; \Star{} And ye yourselves like unto men that wait for their lord, when he will return from the wedding.\end{Resp}
\begin{Resp}\Vsign Watch therefore, for ye know not what hour your Lord doth come.\end{Resp}
\begin{Resp}\Rsign And ye yourselves like unto men that wait for their lord, when he will return from the wedding.\end{Resp}
\begin{Resp}\Vsign Glory be to the Father, and to the Son, \Star{} and to the Holy Ghost.\end{Resp}
\begin{Resp}\Rsign And ye yourselves like unto men that wait for their lord, when he will return from the wedding.\end{Resp}
\switchcolumn*

\LA
\LessonHead{Lectio IX}
\switchcolumn
\EN
\LessonHead{Lesson IX}
\switchcolumn*

\LA
\noindent Ex occasióne hujus senténtiæ quidam introdúcunt mille annos post resurrectiónem, dicéntes, tunc nobis céntuplum ómnium rerum quas dimísimus, et vitam ætérnam esse reddéndam: non intelligéntes, quod, si in céteris digna sit repromíssio, in uxóribus appáreat turpitúdo: ut qui unam pro Dómino dimíserit centum recípiat in futúro. Sensus ergo iste est: Qui carnália pro Salvatóre dimíserit, spirituália recípiet: quæ comparatióne et mérito sui ita erunt, quasi si parvo número centenárius númerus comparétur.\par
\switchcolumn
\EN
\noindent By reason of these words, an hundredfold, some will have that there shall be a thousand years after the resurrection, wherein they that have forsaken all things shall receive an hundredfold of those things which they have forsaken, and shall inherit everlasting life. Such men consider not that though in other things this were worthy, as touching wives it is unseemly for it becometh us not to think that he that hath forsaken one wife in this world, shall receive an hundred wives in that which is to come. But the meaning is this, that every one that for the Saviour's sake hath forsaken earthly things, shall receive spiritual things which things, being rightly weighed against earthly things, are as though an hundredfold were weighed against one.\par
\switchcolumn*

\LA
\TeDeum{Te Deum}
\switchcolumn
\EN
\TeDeum{Te Deum}
\switchcolumn*

\end{paracol}

\par\needspace{10\baselineskip}\addvspace{1.2em}{\centering\small\textsc{Die 18 Julii} · \textsc{18 July}\par}
\Office{S. Camilli de Lellis Confessoris}{St. Camillus de Lellis, Confessor}{Duplex}
\addcontentsline{toc}{subsection}{Die 18 Julii · S. Camilli de Lellis Confessoris \textit{— St. Camillus de Lellis, Confessor}}
\begin{paracol}{2}
\LA
\RefLine{Lectiones I–III: de Scriptura occurrente.}
\switchcolumn
\EN
\RefLine{Lessons I–III: from the Scripture of the day.}
\switchcolumn*

\LA
\RefLine{Lectio IX: de Ss. Symphorosæ et Septem Filiorum Martyrum (07-18o).}
\switchcolumn
\EN
\RefLine{Lesson IX: of S. Symphorosa with her Seven Sons, Martyrs (07-18o).}
\switchcolumn*

\LA
\LessonHead{Lectio IV}
\switchcolumn
\EN
\LessonHead{Lesson IV}
\switchcolumn*

\LA
\noindent Camillus Bucclánici, Theatinæ diœcesis oppido, ex nobili Lelliórum família natus est matre sexagenaria, cui gravidæ visum est per quietem, puerulum crucis signo in péctore munítum, et ágmini puerórum idem signum gestántium præeuntem, se peperisse. Adoléscens rem militarem secutus, sæculi vitiis aliquámdiu indulsit; donec, vigesimum quintum agens ætátis annum, tanto supernæ gratiæ lúmine divinæque offensæ dolóre correptus fuit, ut uberrimo lacrimárum imbre íllico perfusus, anteáctæ vitæ sordes indesinenter abstérgere, novumque indúere hóminem firmiter decreverit. Quare ipso, quo id cóntigit, Purificatiónis beatíssimæ Vírginis festo die, ad fratres Minores, quos Capuccinos vocant, cónvolans, ut eórum número adscriberétur summis precibus exorávit. Voti compos semel atque íterum factus est; sed, fœdo úlcere, quo aliquándo laboráverat, in ejus tibia iterato recrudescente, divinæ providéntiæ majora de eo disponentis consílio humíliter se subjécit; suique victor, illíus religiónis bis expetítum et susceptum habitum bis dimísit.\par
\switchcolumn
\EN
\noindent Camillus was a son of the noble family of the Lelli, and was born at Bacchianico, a town in the Diocese of Chieti, (in the Abruzzi, in the year of our Lord 1550.) His mother was sixty years of age at the time of his birth. While she was great with child, she dreamed that she brought forth a babe bearing the mark of a Cross upon his breast, and going before a troop of other babes marked likewise. When Camillus was a young man he served as a soldier, and yielded himself for a while to the sins of the world. In the twenty-fifth year of his age light from God broke upon him and in a violent fit of tears he determined to wipe away the evil relics of his past life, and to put on the new man. That very day, being the holiday of the Purification of the Most Blessed Virgin, he ran to the Friars Minors, who are commonly called Capuchins, and implored them to enrol him among them. They granted his wishes, but God was keeping him for greater things, and on this as well as on another occasion when he made the same attempt he was forced to abandon it by the increasing virulence of a loathsome running sore in the leg, with which he was afflicted. He meekly bowed himself to the will of Providence, and conquering his own wishes twice, stripped himself of the habit of the Order, which he had sought and received.\par
\switchcolumn*

\LA
\begin{Resp}\Rsign Honéstum fecit illum Dóminus, et custodívit eum ab inimícis, et a seductóribus tutávit illum: \Star{} Et dedit illi claritátem ætérnam.\end{Resp}
\begin{Resp}\Vsign Justum dedúxit Dóminus per vias rectas, et osténdit illi regnum Dei.\end{Resp}
\begin{Resp}\Rsign Et dedit illi claritátem ætérnam.\end{Resp}
\switchcolumn
\EN
\begin{Resp}\Rsign The Lord made him honourable, and defended him from his enemies, and kept him safe from those that lay in wait for him; \Star{} And gave him perpetual glory.\end{Resp}
\begin{Resp}\Vsign He went down with him into the pit, and left him not in bonds.\end{Resp}
\begin{Resp}\Rsign And gave him perpetual glory.\end{Resp}
\switchcolumn*

\LA
\LessonHead{Lectio V}
\switchcolumn
\EN
\LessonHead{Lesson V}
\switchcolumn*

\LA
\noindent Romam profectus, in nosocomíum, quod Insanabílium dícitur, recéptus est; cujus étiam administratiónem, ob perspectas ejus virtútes sibi demandatam, summa integritate, ac sollicitúdine vere paterna peregit. Omnium ægrotum servum se réputans, eórum stérnere lectulos, sordes térgere, ulcéribus medéri, agoníque extremo piis precibus et cohortatiónibus opem ferre solemne hábuit; quibus in munéribus præclára præbuit admirábilis patiéntiæ, invictæ fortitúdinis et heróicæ caritátis exempla. Verum, cum animárum in extremis periclitántium, quod únice intendébat, levámini subsidium litterárum plurimum conférre intellígeret, trigínta duos annos natus in primis grammaticæ elementis tirocinium inter púeros íterum subire non erubuit. Sacerdotio póstea rite initiatus, nonnullis sibi adjunctis sociis, prima jecit congregatiónis Clericórum regulárium infirmis ministrántium fundamenta, írrito conatu obniténte humani generis hoste. Nam Camillus, cælésti voce e Christi crucifixi, manus étiam de ligno avulsas admirando prodigio protendentis, simulacro emissa mirabíliter confirmátus, ordinem suum a Sede apostolica approbari obtinuit; sodálibus quarto obstrictis maxime arduo voto, infirmis, quos étiam pestis infécerit, ministrandi. Quod institutum, quam foret Deo acceptum et animárum salúti proficuum, sanctus Philippus Nerius, qui Camillo a sacris confessiónibus erat, comprobávit, dum ejus alumnis decedéntium agoni opem feréntibus Angelos suggeréntes verba sæpius se vidisse testátus est.\par
\switchcolumn
\EN
\noindent He went to Rome and was received as an inmate in the Hospital for Incurables. In consequence of his eminent good qualities the administration of the Hospital was committed to his charge, and he discharged this office with the most thorough trustworthiness and with a tenderness like a father's. He counted himself the slave of all the patients, and made it a religious duty to make their beds, clean them, dress their sores, and help by godly prayers and exhortations such as were in their last agony. In doing these things he showed himself a bright example of wonderful patience, indomitable firmness, and heroic charity. He became persuaded that a knowledge of letters would make him much more useful as a comforter to the dying, who were his peculiar care, and therefore, at the age of thirty-two years, he humbly went to school again, among little boys learning the first rudiments. After a time he took Priests' orders, and, in company with some companions who joined him, he laid the first foundations of the Congregation of Regular Clerks for ministering to the sick a scheme against which the enemy of man made an unsuccessful struggle. Camillus heard a voice from heaven issue from an image of Christ Crucified, strengthening him, and saw the nailed hands stretched out from the Cross to protect him He obtained from the Apostolic see an approval of his Institute, the members of which, (besides the three vows of Poverty, Chastity, and Obedience,) take a fourth and very stern one, by which they bind themselves to serve all sick persons, even those stricken with the plague. Holy Philip Neri, who was Confessor to Camillus, testified that he had often seen Angels prompting the members of this Congregation what they should speak when they were assisting the dying, a proof how well-pleasing in the sight of God, and how useful for the salvation of souls, is this Institution.\par
\switchcolumn*

\LA
\begin{Resp}\Rsign Amávit eum Dóminus, et ornávit eum: stolam glóriæ índuit eum, \Star{} Et ad portas paradísi coronávit eum.\end{Resp}
\begin{Resp}\Vsign Índuit eum Dóminus lorícam fídei, et ornávit eum.\end{Resp}
\begin{Resp}\Rsign Et ad portas paradísi coronávit eum.\end{Resp}
\switchcolumn
\EN
\begin{Resp}\Rsign The Lord loved him and beautified him; He clothed him with a robe of glory, \Star{} And crowned him at the gates of Paradise.\end{Resp}
\begin{Resp}\Vsign The Lord hath put on him the breast-plate of faith, and hath adorned him.\end{Resp}
\begin{Resp}\Rsign And crowned him at the gates of Paradise.\end{Resp}
\switchcolumn*

\LA
\LessonHead{Lectio VI}
\switchcolumn
\EN
\LessonHead{Lesson VI}
\switchcolumn*

\LA
\noindent Arctióribus hisce vinculis ægrotántium ministerio mancipatus, mirum est qua alacritate, nullis fractus labóribus, nullis deterritus vitæ periculis, diu noctuque ad supremum usque spíritum, eórum cómmodis vigilaverit. Omnibus ómnia factus, vilíssima quæque offícia demissíssimo obsequio flexísque plerumque genibus, véluti Christum ipsum cérneret in infirmis, hílari promptoque animo arripiébat; utque ómnium indigéntiis præsto esset, generalem ordinis præfecturam, cælique delicias quibus in contemplatióne defixus affluebat, sponte dimísit. Paternus vero illíus erga míseros amor tum maxime effulsit, dum et Urbs contagióso morbo primum, deínde extrema annonæ laboráret inópia, et Nolæ in Campánia dira pestis grassarétur. Tanta denique in Deum et próximum caritate exársit, ut angelus nuncupari, et Angelórum opem in vario itinerum discrímine experíri promererétur. Prophetíæ dono et grátia sanitátum præditus, arcana quoque cordium inspéxit; ejusque precibus nunc cibária multiplicáti sunt, nunc aqua in vinum conversa. Tandem vigiliis, jejuniis et assiduis attritus labóribus, cum pelle tantum et óssibus constare viderétur, quinque molestis æque ac diútinis morbis, quos misericórdias Dómini appellabat, fortiter tolerátis, sacramentis munítus, Romæ, inter suavíssima Jesu et Maríæ nómina, ad ea verba: Mitis atque festivus Christi Jesu tibi aspectus appareat; qua prædixerat hora, obdormívit in Dómino, pridie Idus Julii, anno salútis millesimo sexcentésimo décimo quarto, ætátis suæ sexagesimo quinto. Quem, plúribus illustrem miraculis, Benedíctus décimus quartus solemni ritu Sanctórum fastis adscripsit; et Leo décimus tértius, ex Sacrórum cathólici orbis antístitum voto ac Rítuum Congregatiónis consúlto, cæléstem ómnium hospitálium et infirmórum ubíque degéntium patrónum declarávit, ipsiúsque nomen in agonizántium litaníis invocári præcépit.\par
\switchcolumn
\EN
\noindent Then he had thus given himself entirely over by these strict ties to the service of the sick, it was wonderful to see with what earnestness Camillus, broken by no weariness, and scared by no danger to himself, watched over their comfort by day and by night as long as life lasted. Becoming all things to all men, he took with cheerful readiness the most repulsive duties, discharging them with the most humble attention, and oftentimes on his knees, as though he saw Christ Himself in His suffering members. That he might be the readier to serve every one's need, he resigned the general government of his own Institute, and denied himself the indulgence in the heavenly refreshment which abundantly poured upon him, when he fixed his mind solely upon God. His tender, fatherly love toward the wretched had its brightest manifestations when Rome was stricken first by a contagious sickness, and then by famine, and Nola in Campagna suffered from a frightful plague. His love to God and to his neighbour was so glorious that he earned the nickname of Angel, and found Angels helping him in the difficulties of his diverse journeyings. He had the gifts of prophecy and healing, and could read the secret thoughts of men's hearts. At his prayer, food was multiplied, and water turned into wine. His want of sleep, fasting, and unceasing work wore him down till he seemed nothing but skin and bones. He suffered from a complication of five different painful and incurable diseases, which he was accustomed to call the Lord's mercies to him, and which he bore bravely. He died at Rome on the day which he had himself foretold, the fourteenth of July, in the year of salvation 1614, and of his own age the 65th. He had received the Sacraments, and fell asleep in the Lord in an attempt to utter the sweet names of Jesus and Mary, while the Priest was reciting the words of the Ritual: “Gentle and joyous may the Countenance of Christ Jesus appear to thee.” He was famous for many miracles, and Benedict XIV solemnly enrolled him in the Kalendar of the Saints.\par
\switchcolumn*

\LA
\begin{Resp}\Rsign Iste homo perfécit ómnia quæ locútus est ei Deus, et dixit ad eum: Ingrédere in réquiem meam: \Star{} Quia te vidi justum coram me ex ómnibus géntibus.\end{Resp}
\begin{Resp}\Vsign Iste est, qui contémpsit vitam mundi, et pervénit ad cæléstia regna.\end{Resp}
\begin{Resp}\Rsign Quia te vidi justum coram me ex ómnibus géntibus.\end{Resp}
\begin{Resp}\Vsign Glória Patri, et Fílio, \Star{} et Spirítui Sancto.\end{Resp}
\begin{Resp}\Rsign Quia te vidi justum coram me ex ómnibus géntibus.\end{Resp}
\switchcolumn
\EN
\begin{Resp}\Rsign This is he which did according unto all that God commanded him; and God said unto him: Enter thou into My rest, \Star{} For thee have I seen righteous before Me among all people.\end{Resp}
\begin{Resp}\Vsign This is he which loved not his life in this world, and is come unto an everlasting kingdom.\end{Resp}
\begin{Resp}\Rsign For thee have I seen righteous before Me among all people.\end{Resp}
\begin{Resp}\Vsign Glory be to the Father, and to the Son, \Star{} and to the Holy Ghost.\end{Resp}
\begin{Resp}\Rsign For thee have I seen righteous before Me among all people.\end{Resp}
\switchcolumn*

\LA
\LessonHead{Lectio VII}
\switchcolumn
\EN
\LessonHead{Lesson VII}
\switchcolumn*

\LA
\LessonTitle{Léctio sancti Evangélii secúndum Joánnem}
\switchcolumn
\EN
\LessonTitle{From the Holy Gospel according to John}
\switchcolumn*

\LA
\Cite{Joannes 15:12-16}
\switchcolumn
\EN
\Cite{John 15:12-16}
\switchcolumn*

\LA
\noindent In illo témpore: Dixit Jesus discípulis suis: Hoc est præcéptum meum, ut diligátis invicem, sicut diléxi vos. Et réliqua.\par
\switchcolumn
\EN
\noindent At that time, Jesus said unto His disciples: This is My commandment, That ye love one another, as I have loved you. And so on.\par
\switchcolumn*

\LA
\LessonTitle{Homilía sancti Augustíni Epíscopi}
\switchcolumn
\EN
\LessonTitle{Homily by St. Augustine, Bishop of Hippo}
\switchcolumn*

\LA
\Source{Tractátus 83 in Joánnem}
\switchcolumn
\EN
\Source{83rd Tract on John}
\switchcolumn*

\LA
\noindent Quid putamus, fratres mei? numquidnam solum ejus de ista dilectióne mandátum est, qua dilígimus invicem? Nonne est et aliud majus, ut diligámus Deum? Aut vero de sola Deus nobis dilectióne mandávit, ut alia non requiramus? Tria certe commendat Apóstolus, dicens: Manent autem fides, spes, caritas, tria hæc, major autem horum caritas. Et, si in caritate, hoc est, in dilectióne, concludúntur duo illa præcepta, major tamen dicta est esse, non sola. De fide ígitur nobis quam multa mandáta sunt, quam multa de spe! Quis potest cuncta colligere, quis enumerando sufficere? Sed intueamur, quod ait idem Apóstolus: Plenitúdo legis caritas.\par
\switchcolumn
\EN
\noindent What think ye, my brethren? Is this His only commandment, this, That we love one another? Is there not another and a greater, the commandment to love God? Or hath God commanded us only to love, so that we need seek to do no more? Surely the Apostle commendeth three things: And now abideth Faith, Hope, Charity, these three but the greatest of these is Charity. (1 Cor. xiii. 13.) And although in charity, that is, in love, he included the two first and great commandments, and charity he called the greatest, yet is charity not said to be alone. Concerning Faith, concerning Hope, how much is commanded us? Who can gather them all together? Who can reckon them all? And yet let us consider how the same Apostle saith: Love is the fulfilling of the Law. (Rom. xiii. 10.)\par
\switchcolumn*

\LA
\begin{Resp}\Rsign Iste est qui ante Deum magnas virtútes operátus est, et de omni corde suo laudávit Dóminum: \Star{} Ipse intercédat pro peccátis ómnium populórum.\end{Resp}
\begin{Resp}\Vsign Ecce homo sine queréla, verus Dei cultor, ábstinens se ab omni ópere malo, et pérmanens in innocéntia sua.\end{Resp}
\begin{Resp}\Rsign Ipse intercédat pro peccátis ómnium populórum.\end{Resp}
\switchcolumn
\EN
\begin{Resp}\Rsign This is he which wrought great wonders before God, and praised the Lord with all his heart. \Star{} May he pray for all people, that their sins may be forgiven unto them.\end{Resp}
\begin{Resp}\Vsign Behold a man without blame, a worshipper of God in truth, keeping himself clean from every evil work, and abiding still in his innocency.\end{Resp}
\begin{Resp}\Rsign May he pray for all people, that their sins may be forgiven unto them.\end{Resp}
\switchcolumn*

\LA
\LessonHead{Lectio VIII}
\switchcolumn
\EN
\LessonHead{Lesson VIII}
\switchcolumn*

\LA
\noindent Ubi ergo caritas est, quid est quod possit deésse? ubi autem non est, quid est quod possit prodesse? Dæmon credit, nec díligit: nemo díligit, qui non credit. Frustra quidem, sed tamen potest speráre véniam qui non díligit; nemo autem potest desperáre qui díligit. Itaque ubi diléctio est, ibi necessario fides et spes; et, ubi diléctio próximi, ibi necessario étiam diléctio Dei. Qui enim non díligit Deum, quómodo díligit próximum tamquam se ipsum? Quandóquidem non díligit et se ipsum. Est quippe ímpius et iníquus; qui autem díligit iniquitátem, non plane díligit, sed odit ánimam suam.\par
\switchcolumn
\EN
\noindent Where therefore, Charity is, what can be lacking or where Charity is not, what can there be availing? The devil believeth and loveth not, but there is no one that loveth and believeth not. Useless though it be, it is still possible for one that loveth not, to hope to be forgiven but for one that loveth it is impossible to give up hope. Therefore, where love is, there also must faith and hope needs be, and where there is love toward our neighbour there also must there needs be love toward God. For one that loveth not God, how can he love his neighbour as himself, seeing he hateth himself, for he is a blasphemous, wicked wretch, and the lover of wickedness is not the lover, but the deadly enemy of his own self.\par
\switchcolumn*

\LA
\begin{Resp}\Rsign Sint lumbi vestri præcíncti, et lucérnæ ardéntes in mánibus vestris: \Star{} Et vos símiles homínibus exspectántibus dóminum suum, quando revertátur a núptiis.\end{Resp}
\begin{Resp}\Vsign Vigiláte ergo, quia nescítis qua hora Dóminus vester ventúrus sit.\end{Resp}
\begin{Resp}\Rsign Et vos símiles homínibus exspectántibus dóminum suum, quando revertátur a núptiis.\end{Resp}
\begin{Resp}\Vsign Glória Patri, et Fílio, \Star{} et Spirítui Sancto.\end{Resp}
\begin{Resp}\Rsign Et vos símiles homínibus exspectántibus dóminum suum, quando revertátur a núptiis.\end{Resp}
\switchcolumn
\EN
\begin{Resp}\Rsign Let your loins be girded about, and your lights burning; \Star{} And ye yourselves like unto men that wait for their lord, when he will return from the wedding.\end{Resp}
\begin{Resp}\Vsign Watch therefore, for ye know not what hour your Lord doth come.\end{Resp}
\begin{Resp}\Rsign And ye yourselves like unto men that wait for their lord, when he will return from the wedding.\end{Resp}
\begin{Resp}\Vsign Glory be to the Father, and to the Son, \Star{} and to the Holy Ghost.\end{Resp}
\begin{Resp}\Rsign And ye yourselves like unto men that wait for their lord, when he will return from the wedding.\end{Resp}
\switchcolumn*

\end{paracol}

\par\needspace{10\baselineskip}\addvspace{1.2em}{\centering\small\textsc{Die 18 Julii} · \textsc{18 July}\par}
\Office{Ss. Symphorosæ et Septem Filiorum Martyrum}{S. Symphorosa with her Seven Sons, Martyrs}{Simplex}
\addcontentsline{toc}{subsection}{Die 18 Julii · Ss. Symphorosæ et Septem Filiorum Martyrum \textit{— S. Symphorosa with her Seven Sons, Martyrs}}
\begin{paracol}{2}
\LA
\LessonHead{Lectio IX (pro commemoratione)}
\switchcolumn
\EN
\LessonHead{Lesson IX (when commemorated)}
\switchcolumn*

\LA
\LessonTitle{Commemoratio: Ss. Symphorosæ et Septem Filiorum Martyrum}
\switchcolumn
\EN
\LessonTitle{Commemoration of: S. Symphorosa with her Seven Sons, Martyrs}
\switchcolumn*

\LA
\noindent Symphorosa Tiburtina, Getulii Mártyris uxor, ex eo septem fílios péperit, Crescéntium, Julianum, Nemesium, Primitivum, Justinum, Stacteum, et Eugenium; qui omnes propter christianæ fidei professiónem una cum matre, Hadriáno imperatóre, comprehénsi sunt. Quorum pietas, multis variisque tentáta suppliciis, cum stábilis permanéret, mater, quæ fíliis fidei magistra fúerat, dux eisdem ad martyrium éxstitit. Nam, saxo ad collum alligato, in profluentem dejícitur: cujus corpus, conquisítum a fratre ejus Eugenio, sepelítur. Postridie ejus diéi, qui fuit décimo quinto Kalendas Augusti, septem fratres singuli ad palum alligati, varie sunt interfécti: Crescentio guttur ferro transfígitur; Juliáno pectus confóditur; Nemesio cor transverberátur; Primitivo trajícitur umbilicus; Justinus membratim secátur; Stacteus telis confígitur; Eugenius a péctore in duas partes dividitur. Ita octo hostiæ Deo gratíssimæ sunt immolátæ. Corpora in altíssimam fóveam projecta sunt via Tiburtina, nono ab Urbe lápide; quæ, póstea Romam translata, cóndita sunt in ecclésia sancti Angeli in piscina.\par
\switchcolumn
\EN
\noindent Symphorosa was a woman of Tivoli, the wife of the martyr Getulius, unto whom she bore seven sons, named respectively, Crescentius, Julian, Nemesius, Primitivus, Justin, Stacteus, and Eugene, all of whom were arrested along with their mother, in the reign of the Emperor Hadrian, for professing the Christian faith. Their love was tried by many and diverse torments, and their mother who had taught them their religion, was their leader to martyrdom. A stone was tied round her neck and she was thrown into the river. Her body was found and buried by her brother Eugene. The next day, being the 18th of July, the seven brethren were tied each to a stake, and all put to death in diverse ways. Crescentius was stabbed in the throat, Julian in the breast, Nemesius in the heart, and Primitivus in the navel. Justin was hacked limb from limb. Stacteus was shot to death with darts. Eugene was cut into two parts across his breast, (from the head downwards.) Thus were these eight sacrifices of sweet savour offered up to God. Their bodies were thrown into a deep pit, on the road between Rome and Tivoli, at the ninth milestone from Rome, but were afterwards brought to Rome and buried in the Church of the Holy Angel-in-the-Fish-market.\par
\switchcolumn*

\LA
\TeDeum{Te Deum}
\switchcolumn
\EN
\TeDeum{Te Deum}
\switchcolumn*

\end{paracol}

\par\needspace{10\baselineskip}\addvspace{1.2em}{\centering\small\textsc{Die 19 Julii} · \textsc{19 July}\par}
\Office{S. Vincentii a Paulo Confessoris}{St. Vincent de Paul, Confessor}{Duplex}
\addcontentsline{toc}{subsection}{Die 19 Julii · S. Vincentii a Paulo Confessoris \textit{— St. Vincent de Paul, Confessor}}
\begin{paracol}{2}
\LA
\RefLine{Lectiones I–III: de Scriptura occurrente.}
\switchcolumn
\EN
\RefLine{Lessons I–III: from the Scripture of the day.}
\switchcolumn*

\LA
\LessonHead{Lectio IV}
\switchcolumn
\EN
\LessonHead{Lesson IV}
\switchcolumn*

\LA
\noindent Vincentius a Paulo, natióne Gallus, Podii non procul ab Aquis Tarbellis in Aquitánia natus, jam tum a púero exímiam in páuperes caritátem præ se tulit. A custódia paterni gregis ad litteras evocatus, humanas Aquis, divinas cum Tolosæ, tum Cæsaraugustæ didicit. Sacerdotio initiátus ac theologíæ laurea insignítus, in Turcas íncidit, qui captivum in Africam abduxérunt. Sed in captivitáte positus herum ipsum Christo rursus lucrifecit. Cum eo ígitur ex bárbaris oris, opitulante Deípara, sese prorípiens, ad apostolica limina iter instituit. Unde in Gálliam reversus, Clippíaci primum, mox Castelliónis parœcias sanctíssime rexit. Renuntiátus a rege primarius Sacrórum miníster in Galliæ trirémibus, mirum quo zelo et ducum et rémigum salúti operam posúerit. Moniálibus Visitatiónis a sancto Francisco Salesio præpositus, tanta prudéntia per annos circiter quadragínta eam curam sustinuit, ut maxime comprobaverit judícium sanctíssimi præsulis, qui sacerdotem Vincentio digniórem nullum se nosse fatebátur.\par
\switchcolumn
\EN
\noindent Vincent de Paul was a Frenchman by nation, and was born at (Ranquines, in the parish of) Puy, not far from Dax in Gascony, (upon the 24th day of April, in the year of salvation 1576.) From a little child he showed remarkable charity towards the poor. His father removed him from keeping his cattle, in order to give him a school education, and he learnt earthly things at Dax, and theology both at Toulouse and at Saragossa. He took Priest's orders, and a degree in Divinity. (In 1605,) he was taken prisoner by Mahommedan pirates, who carried him off, and sold him for a slave in Africa. In his slavery he converted his owner, who was an apostate, back to Christ. Under the protection of the Mother of God, Vincent escaped from Barbary. He first visited the thresholds of the Apostles, and afterwards returned to France. He was the saintly Rector first of the Parish of Clichy, and afterwards of that of Chatillon. He was appointed by the King, ChaplainGeneral for the galleys of France, and worked with extraordinary zeal for the health of the souls both of those who commanded and of the convicts who rowed. He was made Superior of the Nuns of the Visitation by St Francis de Sales, and discharged this duty for about forty years, with a wisdom which so approved itself to the judgment of their holy Founder, that he was used to say he knew no worthier Priest than Vincent.\par
\switchcolumn*

\LA
\begin{Resp}\Rsign Honéstum fecit illum Dóminus, et custodívit eum ab inimícis, et a seductóribus tutávit illum: \Star{} Et dedit illi claritátem ætérnam.\end{Resp}
\begin{Resp}\Vsign Justum dedúxit Dóminus per vias rectas, et osténdit illi regnum Dei.\end{Resp}
\begin{Resp}\Rsign Et dedit illi claritátem ætérnam.\end{Resp}
\switchcolumn
\EN
\begin{Resp}\Rsign The Lord made him honourable, and defended him from his enemies, and kept him safe from those that lay in wait for him; \Star{} And gave him perpetual glory.\end{Resp}
\begin{Resp}\Vsign He went down with him into the pit, and left him not in bonds.\end{Resp}
\begin{Resp}\Rsign And gave him perpetual glory.\end{Resp}
\switchcolumn*

\LA
\LessonHead{Lectio V}
\switchcolumn
\EN
\LessonHead{Lesson V}
\switchcolumn*

\LA
\noindent Evangelizandis paupéribus, præsertim rurícolis, ad decrépitam usque ætátem indeféssus incúbuit, eique apostolico operi tum se, tum alumnos congregatiónis, quam sub nómine Presbyterórum sæculárium Missiónis instituit, perpetuo voto a sancta Sede confírmato, speciatim obstrinxit. Quantum autem augendæ cleri disciplinæ allaboraverit, testántur erecta majórum clericórum seminaria, collatiónum de divinis inter sacerdótes frequéntia, et sacræ ordinatióni præmittenda exercítia; ad quæ, sicut et ad pios laicórum secéssus, institúti sui domicília libenter patére vóluit. Insuper, ad amplificandam fidem et pietátem, evangelicos misit operarios, non in solas Galliæ provincias, sed et in Italiam, Polóniam, Scotiam, Hiberniam, atque ad Bárbaros et Indos. Ipse vero, vita functo Ludovíco décimo tertio, cui moriénti hortator ástitit, a regína Anna Austríaca, matre Ludovíci décimi quarti, in sanctius concílium accítus, studiosíssime egit, ut nónnisi digniores ecclésiis ac monastériis præficeréntur; civiles discordiæ, singularia certamina, serpéntes erróres, quos simul sensit et exhorruit, amputaréntur; debítaque judíciis apostolicis obediéntia præstarétur ab ómnibus.\par
\switchcolumn
\EN
\noindent The preaching of the Gospel to the poor, especially peasants, was the work at which he toiled unweariedly, till he was disabled by age. To this special work he bound himself and the members of the Congregation which he founded under the missionary Congregation of Secular Priests, by a perpetual vow approved by the Holy See. How great were his labours for bettering the discipline of the clergy, is attested by the building of Seminaries for the final education of young clerks, the number of meetings of Priests to discuss holy things, and the religious exercises preparatory to Ordination, for which, as well as for godly retreats by laymen, he wished that the houses belonging to his Institute should be always freely open. To spread wider the growth of faith and godliness, he sent his Gospel laborers not only into the several provinces of France, but also into Italy, Poland, Scotland, and Ireland, and also to Barbary and India. He assisted Lewis XIII. on his death-bed, and the Queen Anne of Austria, mother of Lewis XIV., put him upon the young King's Council of Conscience during the Regency, in which position it was his unceasing effort that none but the most worthy should be named to churches and monasteries, that civil contests, duels, and creeping false doctrines, from which himself shrank as soon as he met them, should be put down, and that all men should yield the obedience which was due to the decisions of the Apostolic See.\par
\switchcolumn*

\LA
\begin{Resp}\Rsign Amávit eum Dóminus, et ornávit eum: stolam glóriæ índuit eum, \Star{} Et ad portas paradísi coronávit eum.\end{Resp}
\begin{Resp}\Vsign Índuit eum Dóminus lorícam fídei, et ornávit eum.\end{Resp}
\begin{Resp}\Rsign Et ad portas paradísi coronávit eum.\end{Resp}
\switchcolumn
\EN
\begin{Resp}\Rsign The Lord loved him and beautified him; He clothed him with a robe of glory, \Star{} And crowned him at the gates of Paradise.\end{Resp}
\begin{Resp}\Vsign The Lord hath put on him the breast-plate of faith, and hath adorned him.\end{Resp}
\begin{Resp}\Rsign And crowned him at the gates of Paradise.\end{Resp}
\switchcolumn*

\LA
\LessonHead{Lectio VI}
\switchcolumn
\EN
\LessonHead{Lesson VI}
\switchcolumn*

\LA
\noindent Nullum fuit calamitátis genus, cui paterne non occúrrerit. Fideles sub Turcárum jugo geméntes, infántes expósitos, júvenes díscolos, vírgines periclitántes, moniáles dispérsas, mulieres lapsas, ad trirémes damnatos, peregrinos infirmos, artífices inválidos, ipsosque mente captos, ac innúmeros mendícos subsidiis et hospitiis etiamnum superstítibus excepit ac pie fovit. Lotharingiam, Campaniam, Picardiam, aliasque regiónes peste, fame, belloque vastatas, prolixe refécit. Plurima ad perquirendos et sublevándos míseros sodalítia fundávit, inter quæ celebris matronárum cœtus, et late diffúsa sub nómine Caritátis puellárum socíetas. Puellas quoque tum de Cruce, tum de Providéntia, ac sanctæ Genovéfæ ad sequioris sexus educatiónem erigendas curávit. Hæc inter et alia gravíssima negotia Deo jugiter intentus, cunctis affábilis, ac sibi semper constans, simplex, rectus, húmilis, ab honóribus, divítiis ac deliciis semper abhorruit; audítus dicere, rem nullam sibi placére, prætérquam in Christo Jesu, quem in ómnibus studebat imitari. Corporis demum afflictatióne, labóribus, senioque attritus, die vigesima septima Septembris, anno salútis supra millesimum sexcentésimo sexagesimo, ætátis suæ octogesimo quinto, Parísiis, in domo sancti Lázari, quæ caput est congregatiónis Missiónis, plácide obdormívit; quem, virtútibus, meritis ac miraculis clarum, Clemens duodecimus inter Sanctos rétulit, ipsíus celebritáti die décima nona mensis Julii quotannis assignata. Hunc autem divínæ caritátis exímium heróem, de unoquóque hóminum génere óptime méritum, Leo tértius décimus, instántibus plúribus Sacrórum antistítibus, ómnium societátum caritátis, in toto cathólico orbe exsisténtium et ab eo quomodocúmque promanántium, peculiárem apud Deum patrónum declarávit et constítuit.\par
\switchcolumn
\EN
\noindent There was no kind of misery which he did not strive with fatherly tenderness to relieve. Christians groaning in Mahommedan slavery, foundlings, deformed children, young maidens exposed to danger, houseless nuns, fallen women, convicts sent to the galleys, sick foreigners, disabled workmen, lunatics, and beggars without number, all these he relieved, and devoutly housed in diverse charitable institutions which remain to this day. When Lorraine, Champagne, Picardy, and other districts were desolated by plague, famine, and war, he made immense efforts for their relief. He founded many charitable societies, to find out and succour the unfortunate. Among these are remarkable that of Matrons, and that of Sisters of Charity which hath been so widely spread. By those of the Cross, of Providence, and of St. Guinevere he aimed at bringing up young girls as school - mistresses. Amid all these and other most anxious business-matters, he remained always looking simply to God, kind to all, true to himself, plain, upright, and lowly. From all honors, riches, and pleasures, he ever shrank, and was heard to say, that nothing gave him any pleasure, except in Christ Jesus, Whom it was his wish in all things to follow. With a body worn out with hardships, work, and old age, he gently fell asleep in the house of St. Lazarus at Paris, the chief house of the Congregation of the Missions, upon the 27th day of September, in the year of salvation 1660, and of his own age the 85th. He was famous on account of his life, his works, and his miracles, and Clement XII inscribed his name among those of the saints, appointing for his feast-day the 19th day of the month of July. Finally, at the earnest prayer of many prelates, Leo XIII. proclaimed and established this hero of charity, illustrious for his services to all classes of men, as the patron before God in heaven of all charitable societies throughout the whole Catholic world which derive their origin in any way from his institution.\par
\switchcolumn*

\LA
\begin{Resp}\Rsign Iste homo perfécit ómnia quæ locútus est ei Deus, et dixit ad eum: Ingrédere in réquiem meam: \Star{} Quia te vidi justum coram me ex ómnibus géntibus.\end{Resp}
\begin{Resp}\Vsign Iste est, qui contémpsit vitam mundi, et pervénit ad cæléstia regna.\end{Resp}
\begin{Resp}\Rsign Quia te vidi justum coram me ex ómnibus géntibus.\end{Resp}
\begin{Resp}\Vsign Glória Patri, et Fílio, \Star{} et Spirítui Sancto.\end{Resp}
\begin{Resp}\Rsign Quia te vidi justum coram me ex ómnibus géntibus.\end{Resp}
\switchcolumn
\EN
\begin{Resp}\Rsign This is he which did according unto all that God commanded him; and God said unto him: Enter thou into My rest, \Star{} For thee have I seen righteous before Me among all people.\end{Resp}
\begin{Resp}\Vsign This is he which loved not his life in this world, and is come unto an everlasting kingdom.\end{Resp}
\begin{Resp}\Rsign For thee have I seen righteous before Me among all people.\end{Resp}
\begin{Resp}\Vsign Glory be to the Father, and to the Son, \Star{} and to the Holy Ghost.\end{Resp}
\begin{Resp}\Rsign For thee have I seen righteous before Me among all people.\end{Resp}
\switchcolumn*

\LA
\LessonHead{Lectio VII}
\switchcolumn
\EN
\LessonHead{Lesson VII}
\switchcolumn*

\LA
\LessonTitle{Léctio sancti Evangélii secúndum Lucam}
\switchcolumn
\EN
\LessonTitle{From the Holy Gospel according to Luke}
\switchcolumn*

\LA
\Cite{Luc 10:1-9}
\switchcolumn
\EN
\Cite{Luke 10:1-9}
\switchcolumn*

\LA
\noindent In illo témpore: Designávit Dóminus et álios septuagínta duos: et misit illos binos ante fáciem suam, in omnem civitátem et locum, quo erat ipse ventúrus. Et réliqua.\par
\switchcolumn
\EN
\noindent At that time, the Lord appointed other seventy-two also, and sent them two and two before His face into every city and place, whither He Himself would come. And so on.\par
\switchcolumn*

\LA
\LessonTitle{Homilía sancti Gregórii Papæ}
\switchcolumn
\EN
\LessonTitle{Homily by Pope St. Gregory (the Great.)}
\switchcolumn*

\LA
\Source{Homilia 17 in Evangelia}
\switchcolumn
\EN
\Source{17th Homily on the Gospels}
\switchcolumn*

\LA
\noindent Dóminus et Salvátor noster, fratres caríssimi, aliquándo nos sermónibus, aliquándo vero opéribus ádmonet. Ipsa étenim facta ejus præcépta sunt: quia dum áliquid tácitus facit, quid ágere debeámus innotéscit. Ecce enim binos in prædicatiónem discípulos mittit: quia duo sunt præcépta caritátis, Dei vidélicet amor, et próximi: et minus quam inter duos cáritas habéri non potest. Nemo enim próprie ad semetípsum habére caritátem dícitur: sed diléctio in álterum tendit, ut cáritas esse possit.\par
\switchcolumn
\EN
\noindent Dearly beloved brethren, our Lord and Saviour doth sometimes admonish us by words, and sometimes by works. Yea, His very works do themselves teach us for that which He doth silently His example still moveth us to copy. Behold how He sendeth forth His disciples to preach by two and two since there are two commandments to love, that is, a commandment to love God, and a commandment to love our neighbour and where there are not two, the one, being alone, hath not whereon to do the Lord's commandment. And no man can properly be said to love himself: for love tendeth outward toward our neighbour, if it be the love whereto the Gospel doth oblige us.\par
\switchcolumn*

\LA
\begin{Resp}\Rsign Iste est qui ante Deum magnas virtútes operátus est, et de omni corde suo laudávit Dóminum: \Star{} Ipse intercédat pro peccátis ómnium populórum.\end{Resp}
\begin{Resp}\Vsign Ecce homo sine queréla, verus Dei cultor, ábstinens se ab omni ópere malo, et pérmanens in innocéntia sua.\end{Resp}
\begin{Resp}\Rsign Ipse intercédat pro peccátis ómnium populórum.\end{Resp}
\switchcolumn
\EN
\begin{Resp}\Rsign This is he which wrought great wonders before God, and praised the Lord with all his heart. \Star{} May he pray for all people, that their sins may be forgiven unto them.\end{Resp}
\begin{Resp}\Vsign Behold a man without blame, a worshipper of God in truth, keeping himself clean from every evil work, and abiding still in his innocency.\end{Resp}
\begin{Resp}\Rsign May he pray for all people, that their sins may be forgiven unto them.\end{Resp}
\switchcolumn*

\LA
\LessonHead{Lectio VIII}
\switchcolumn
\EN
\LessonHead{Lesson VIII}
\switchcolumn*

\LA
\noindent Ecce enim binos ad prædicándum discípulos Dóminus mittit: quátenus hoc nobis tácitus ínnuat, quia qui caritátem erga álterum non habet, prædicatiónis offícium suscípere nullátenus debet. Bene autem dícitur, quia misit eos ante fáciem suam in omnem civitátem et locum, quo erat ipse ventúrus. Prædicatóres enim suos Dóminus séquitur: quia prædicátio prǽvenit, et tunc ad mentis nostræ habitáculum Dóminus venit, quando verba exhortatiónis præcúrrunt: atque per hoc véritas in mente suscípitur.\par
\switchcolumn
\EN
\noindent Behold, the Lord sendeth forth His disciples to preach by two and two and thus doing, He doth silently teach us that whosoever loveth not his neighbour, such a one it behoveth not to take upon him the office of a preacher. Well also is it said that He sent them before His face into every city and place whither He Himself would come. The Lord followeth His preachers first cometh preaching, and then the Lord Himself cometh to the house of our mind, whither the word of exhortation hath come before and so cometh the truth into our mind.\par
\switchcolumn*

\LA
\begin{Resp}\Rsign Sint lumbi vestri præcíncti, et lucérnæ ardéntes in mánibus vestris: \Star{} Et vos símiles homínibus exspectántibus dóminum suum, quando revertátur a núptiis.\end{Resp}
\begin{Resp}\Vsign Vigiláte ergo, quia nescítis qua hora Dóminus vester ventúrus sit.\end{Resp}
\begin{Resp}\Rsign Et vos símiles homínibus exspectántibus dóminum suum, quando revertátur a núptiis.\end{Resp}
\begin{Resp}\Vsign Glória Patri, et Fílio, \Star{} et Spirítui Sancto.\end{Resp}
\begin{Resp}\Rsign Et vos símiles homínibus exspectántibus dóminum suum, quando revertátur a núptiis.\end{Resp}
\switchcolumn
\EN
\begin{Resp}\Rsign Let your loins be girded about, and your lights burning; \Star{} And ye yourselves like unto men that wait for their lord, when he will return from the wedding.\end{Resp}
\begin{Resp}\Vsign Watch therefore, for ye know not what hour your Lord doth come.\end{Resp}
\begin{Resp}\Rsign And ye yourselves like unto men that wait for their lord, when he will return from the wedding.\end{Resp}
\begin{Resp}\Vsign Glory be to the Father, and to the Son, \Star{} and to the Holy Ghost.\end{Resp}
\begin{Resp}\Rsign And ye yourselves like unto men that wait for their lord, when he will return from the wedding.\end{Resp}
\switchcolumn*

\LA
\LessonHead{Lectio IX}
\switchcolumn
\EN
\LessonHead{Lesson IX}
\switchcolumn*

\LA
\noindent Hinc namque eísdem prædicatóribus Isaías dicit: Paráte viam Dómini, rectas fácite sémitas Dei nostri. Hinc fíliis Psalmísta ait: Iter fácite ei, qui ascéndit super occásum. Super occásum namque Dóminus ascéndit: quia unde in passióne occúbuit, inde majórem suam glóriam resurgéndo manifestávit. Super occásum vidélicet ascéndit; quia mortem quam pértulit, resurgéndo calcávit. Ei ergo qui ascéndit super occásum, iter fácimus, cum nos ejus glóriam vestris méntibus prædicámus, ut eas et ipse post véniens, per amóris sui præséntiam illústret.\par
\switchcolumn
\EN
\noindent Therefore, to preachers saith Isaiah: “Prepare ye the way of the Lord, make straight an highway for our God.” (xl. 3.) And again the Psalmist saith: “Spread a path before him that rideth upon the West.” (lxvii. 4.) The Lord rideth upon the West above that from which in death He veiled His glory, hath He royally exalted that glory that excelleth, even the glory of His rising again. He rideth upon the West, Who, being risen again from the dead, is throned high above the death to which He bowed. Before Him, therefore, That rideth upon the West, we spread a path, when we set forth His glory before the eyes of your mind, to the end that He Himself may come after, and Himself enlighten the same your minds by His presence and His love.\par
\switchcolumn*

\LA
\TeDeum{Te Deum}
\switchcolumn
\EN
\TeDeum{Te Deum}
\switchcolumn*

\end{paracol}

\par\needspace{10\baselineskip}\addvspace{1.2em}{\centering\small\textsc{Die 20 Julii} · \textsc{20 July}\par}
\Office{S. Hieronymi Æmiliani Confessoris}{St. Jerome Emiliani, Confessor}{Duplex}
\addcontentsline{toc}{subsection}{Die 20 Julii · S. Hieronymi Æmiliani Confessoris \textit{— St. Jerome Emiliani, Confessor}}
\begin{paracol}{2}
\LA
\RefLine{Lectiones I–III: de Scriptura occurrente.}
\switchcolumn
\EN
\RefLine{Lessons I–III: from the Scripture of the day.}
\switchcolumn*

\LA
\LessonHead{Lectio IV}
\switchcolumn
\EN
\LessonHead{Lesson IV}
\switchcolumn*

\LA
\noindent Hieronymus, e gente patrícia Æmiliáni Venetiis ortus, a prima adolescéntia milítiæ addictus, difficillimis reipublicæ tempóribus Casto novo ad Quarum in móntibus Tarvisinis præfícitur. Arce ab hostibus capta, ipse in teterrimum carcerem detrúditur, mánibus ac pédibus vinctus; cui omni humana ope destituto beatíssima Virgo, ejus precibus exorata, clemens adest, víncula solvit, et per médios hostes, qui vias omnes obséderant, in Tarvisii conspéctum incólumum ducit. Urbem ingréssus, ad Deíparæ aram, cui se vóverat, mánicas, compedes, catenas, quas secum detulerat, in accepti beneficii testimónium suspendit. Reversus Venetias, cœpit pietátis stúdia impensius cólere, in páuperes mire effúsus, sed puerórum præsertim misértus, qui paréntibus orbati, egeni et sórdidi per urbem vagabántur; quos in ædes a se conductas recepit de suo alendos, et christiánis móribus imbuéndos.\par
\switchcolumn
\EN
\noindent This Jerome was born at Venice, of the Patrician family of the Miani, (in the year of our Lord 1481.) He was trained up to be a soldier, and (in 1508,) in the most troublesome times of the Commonwealth, he commanded the fortress of Castel-nuovo at Quero, in the mountains near Treviso. After the fall of the place, he was chained hand and foot, and cast into a filthy prison. When all hope of help from man had forsaken him, the Most Blessed Virgin, in answer to his prayers, mercifully came to him, loosed his fetters, and brought him unhurt within sight of Treviso, through the midst of the enemy, who kept all the roads. As soon as he entered the city of Treviso, as an acknowledgment of the favour he had received, he hung up his chains, which he had brought away with him, at an Altar of the Mother of God, to whom he had vowed himself. After his return to Venice he gave himself up to godly works. Amid his great tenderness to all the poor, his compassion was chiefly roused by the fatherless little boys who wandered through the city starving and filthy; them he took into an house conducted by himself, where at his own cost he provided them with board, lodging, clothing, and a Christian education.\par
\switchcolumn*

\LA
\begin{Resp}\Rsign Honéstum fecit illum Dóminus, et custodívit eum ab inimícis, et a seductóribus tutávit illum: \Star{} Et dedit illi claritátem ætérnam.\end{Resp}
\begin{Resp}\Vsign Justum dedúxit Dóminus per vias rectas, et osténdit illi regnum Dei.\end{Resp}
\begin{Resp}\Rsign Et dedit illi claritátem ætérnam.\end{Resp}
\switchcolumn
\EN
\begin{Resp}\Rsign The Lord made him honourable, and defended him from his enemies, and kept him safe from those that lay in wait for him; \Star{} And gave him perpetual glory.\end{Resp}
\begin{Resp}\Vsign He went down with him into the pit, and left him not in bonds.\end{Resp}
\begin{Resp}\Rsign And gave him perpetual glory.\end{Resp}
\switchcolumn*

\LA
\LessonHead{Lectio V}
\switchcolumn
\EN
\LessonHead{Lesson V}
\switchcolumn*

\LA
\noindent Per eos dies Venetias appúlerant beátus Cajetanus et Petrus Caráfa (postmodum Paulus quartus), qui, Hierónymi spíritu novoque instituto colligéndi órphanos probato, illum in Incurabílium hospitale adduxérunt, in quo órphanos simul educaret, atque ægrotis pari caritate inservíret. Mox eorúmdem hortatu in próximam continentem profectus, Brixiæ primum, deínde Bérgomi atque Novocómi orphanotróphia eréxit: Bérgomi præsertim, ubi præter duo, pro púeris unum et pro puellis álterum, domum excipiendis, novo in illis regiónibus exemplo, muliéribus a turpi vita ad pœniténtiam conversis, apéruit. Somaschæ, in humili pago agri Bergoménsis ad Vénetæ ditiónis fines, sibi ac suis ibi sedem constítuit, formamque indúxit, congregatiónis, cui proptérea a Somascha nomen factum; quam subinde auctam et propagatam, nedum orphanórum regímini et ecclesiárum cultui, sed ad majórem christianæ reipublicæ utilitátem, adolescéntium in litteris et bonis móribus institutióni in collegiis, academíis et seminariis addictam, sanctus Pius quintus inter religiósos ordines adscripsit, ceteríque Pontifices privilegiis ornarunt.\par
\switchcolumn
\EN
\noindent In those days there came to Venice blessed Cajetan (of Tiene,) and Peter Carafa, who was afterwards Paul IV. They were pleased with the spirit of Jerome, and with his new Asylum for Orphans, and took him to the Hospital for Incurables, as well to bring up Orphans, as to extend his charity equally to the sick. Soon after, by the advice of the same, he went to the mainland, and built orphanages first at Brescia, then at Bergamo and Como. His chief foundations were at Bergamo, where besides an orphanage for little boys, and another for little girls, he opened an house of Refuge for repentant harlots, being the first institution of that kind in that part of the world. In the end he went to dwell at Somascha, a hamlet in the district of Bergamo, close to the frontiers of the Venetian territory, and there made an house for himself and his disciples, and gave shape to a congregation, which is generally called the Congregation of Somascha. This congregation grew and spread, and found its work not only in the education of orphans and the service of Churches, but also in a wider usefulness to the Christian Commonwealth, by training up lads in letters and good manners. Holy Pius V. enrolled it among the religious Orders, and other Popes have given it diverse privileges.\par
\switchcolumn*

\LA
\begin{Resp}\Rsign Amávit eum Dóminus, et ornávit eum: stolam glóriæ índuit eum, \Star{} Et ad portas paradísi coronávit eum.\end{Resp}
\begin{Resp}\Vsign Índuit eum Dóminus lorícam fídei, et ornávit eum.\end{Resp}
\begin{Resp}\Rsign Et ad portas paradísi coronávit eum.\end{Resp}
\switchcolumn
\EN
\begin{Resp}\Rsign The Lord loved him and beautified him; He clothed him with a robe of glory, \Star{} And crowned him at the gates of Paradise.\end{Resp}
\begin{Resp}\Vsign The Lord hath put on him the breast-plate of faith, and hath adorned him.\end{Resp}
\begin{Resp}\Rsign And crowned him at the gates of Paradise.\end{Resp}
\switchcolumn*

\LA
\LessonHead{Lectio VI}
\switchcolumn
\EN
\LessonHead{Lesson VI}
\switchcolumn*

\LA
\noindent Orphanis colligéndis intentus, Mediolanum proficíscitur atque Ticinum; et utrobíque collectis agmínibus puerórum, tectum, victum, vestem, magistros, nobílibus viris favéntibus, próvide constítuit. Inde Somascham redux, ómnibus ómnia factus, a nullo abhorrebat opere, quod in próximi bonum cédere prævidéret. Agricolis immixtus per agros sparsis, dum se illis adjutórem in meténdis frugibus præbet, mysteria fidei explicábat; puerórum cápita porígine fœda abstergens, et patienter tractans curábat; pútridis rusticórum vulnéribus medebátur eo successu, ut grátia curatiónum donatus censerétur. In monte, qui Somaschæ imminet, repérta specu, in illam se ábdidit; ubi, se flagellis cædens, dies integros jejunus tránsigens, oratióne in plurimum noctem protracta, super nudo saxo brevem somnum carpens, sui aliorúmque noxárum pœnas luebat. In hujus specus interiori recessu ex árido sílice exstillat aqua, precibus servi Dei, ut constant tradítio est, impetrata, quæ usque in hodiérnam diem jugiter manans, et in varias regiónes delata, ægris sanitátem plerumque conciliat. Tandem, ex contagióne quæ per omnem vallem serpebat, dum ægrotántibus inservit, et vita functos propriis humeris ad sepultúram defert, contracto morbo, annos natus sex et quinquagínta, quam paulo ante prædixerat, pretiósam mortem obiit anno millesimo quingentésimo trigesimo septimo. Quem plúribus in vita et post mortem miraculis illustrem, Benedíctus décimus quartus Beatórum, Clemens vero décimus tertius Sanctórum fastis solemniter adscripsit.\par
\switchcolumn
\EN
\noindent Jerome went to Milan and to Ticino to gather orphans together, and in both places he gathered a multitude of little boys for whom the charity of noblemen enabled him to provide board, lodging, clothes, and schooling. He returned to Somascha, and, still making himself all things to all men, refused no toil by which he saw that he could be of any use to his neighbour. He was used to go about in the fields, helping the reapers in their work, and meanwhile teaching them in the mysteries of the faith. He was very patient in cleansing and healing the heads of little boys foul with lice, and proved so successful a physician to the stinking sores of the poor country-people, that he got a reputation for having the gift of healing. He found a cave in the mountain which hangs over Somascha, and there he would hide himself, passing whole days without meat or drink, and oftentimes scourging himself, continuing in prayer long into the night, and taking his short sleep upon the bare rock, in expiation of his own sins and the sins of others. In the far end of this cave, there drippeth out of the dry stone some water, which is said by an unwavering tradition to have come there at the prayers of the man of God. It droppeth freely even to this day, and is taken to diverse places at a distance, where it often hath an healing effect upon the sick. At length an infectious disorder broke out in all the valley, and Jerome, who nursed the sick and carried the dead to burial on his own shoulders, caught it, and died a precious death, as he had himself foretold, (upon the 8th day of February,) in the fifty-seventh year of his age, and that of salvation 1537. He was famous for many miracles, both during his life and after his death Benedict X. solemnly enrolled his name among those of the Blessed, and Clement XIII inserted it in the Kalendar of the Saints.\par
\switchcolumn*

\LA
\begin{Resp}\Rsign Iste homo perfécit ómnia quæ locútus est ei Deus, et dixit ad eum: Ingrédere in réquiem meam: \Star{} Quia te vidi justum coram me ex ómnibus géntibus.\end{Resp}
\begin{Resp}\Vsign Iste est, qui contémpsit vitam mundi, et pervénit ad cæléstia regna.\end{Resp}
\begin{Resp}\Rsign Quia te vidi justum coram me ex ómnibus géntibus.\end{Resp}
\begin{Resp}\Vsign Glória Patri, et Fílio, \Star{} et Spirítui Sancto.\end{Resp}
\begin{Resp}\Rsign Quia te vidi justum coram me ex ómnibus géntibus.\end{Resp}
\switchcolumn
\EN
\begin{Resp}\Rsign This is he which did according unto all that God commanded him; and God said unto him: Enter thou into My rest, \Star{} For thee have I seen righteous before Me among all people.\end{Resp}
\begin{Resp}\Vsign This is he which loved not his life in this world, and is come unto an everlasting kingdom.\end{Resp}
\begin{Resp}\Rsign For thee have I seen righteous before Me among all people.\end{Resp}
\begin{Resp}\Vsign Glory be to the Father, and to the Son, \Star{} and to the Holy Ghost.\end{Resp}
\begin{Resp}\Rsign For thee have I seen righteous before Me among all people.\end{Resp}
\switchcolumn*

\LA
\LessonHead{Lectio VII}
\switchcolumn
\EN
\LessonHead{Lesson VII}
\switchcolumn*

\LA
\LessonTitle{Léctio sancti Evangélii secúndum Matthǽum}
\switchcolumn
\EN
\LessonTitle{From the Holy Gospel according to Matthew}
\switchcolumn*

\LA
\Cite{Matt 19:13-21}
\switchcolumn
\EN
\Cite{Matt 19:13-21}
\switchcolumn*

\LA
\noindent In illo témpore: Obláti sunt Jesu párvuli, ut manus eis impóneret, et oraret. Et réliqua.\par
\switchcolumn
\EN
\noindent At that time, there were brought unto Jesus little children that He should put His Hands on them, and pray. And so on.\par
\switchcolumn*

\LA
\LessonTitle{Homilía sancti Joánnis Chrysostomi}
\switchcolumn
\EN
\LessonTitle{Homily by St. John Chrysostom, Patriarch of Constantinople.}
\switchcolumn*

\LA
\Source{Homilia 62 in Matthæum}
\switchcolumn
\EN
\Source{Homily on Matthew}
\switchcolumn*

\LA
\noindent Cur discípuli púeros abigebant? Dignitátis causa. Quid ergo ille? Ut doceat illos modeste sapere, fastumque mundanum conculcare, et suscipit, et ulnis complectitur, talibúsque regnum cælórum pollicétur; id quod étiam dixit superius. Et nos ígitur, si volumus heredes esse cælórum, hanc virtútem cum diligéntia magna sectemur. Hoc est enim philosophíæ culmen, simplicem esse cum prudéntia; hæc vita est angelica. Anima enim puéruli ómnibus animi morbis vacua est; non memóriam retinet injuriárum, sed eas inferéntes adit ut amicos, ac si nihil factum esset. Et quamvis a matre verbéribus cædátur, eam semper quærit, et ómnibus antepónit.\par
\switchcolumn
\EN
\noindent Wherefore did the disciples rebuke them that brought them? From an idea of His dignity. What therefore did He To teach them to be lowly, and to be above the niceness of the world, He took the little children, and embraced them in His Arms, and declared that of such is the kingdom of heaven as also He had said above, (xviii. 3, 4.) And we also, if we would fain be heirs of the kingdom of heaven, let us seek with great earnestness this virtue. For this is the highest peak of philosophy, to be simple and wise this is the life of an Angel. The mind of a little child is free from all the diseases of the mind a little child keepeth no remembrance of injuries, but goeth unto such as have inflicted them, as if unto friends, and as if nothing had happened. Although his mother give him stripes, yet a little child ever seeketh her, and putteth her before all.\par
\switchcolumn*

\LA
\begin{Resp}\Rsign Iste est qui ante Deum magnas virtútes operátus est, et de omni corde suo laudávit Dóminum: \Star{} Ipse intercédat pro peccátis ómnium populórum.\end{Resp}
\begin{Resp}\Vsign Ecce homo sine queréla, verus Dei cultor, ábstinens se ab omni ópere malo, et pérmanens in innocéntia sua.\end{Resp}
\begin{Resp}\Rsign Ipse intercédat pro peccátis ómnium populórum.\end{Resp}
\switchcolumn
\EN
\begin{Resp}\Rsign This is he which wrought great wonders before God, and praised the Lord with all his heart. \Star{} May he pray for all people, that their sins may be forgiven unto them.\end{Resp}
\begin{Resp}\Vsign Behold a man without blame, a worshipper of God in truth, keeping himself clean from every evil work, and abiding still in his innocency.\end{Resp}
\begin{Resp}\Rsign May he pray for all people, that their sins may be forgiven unto them.\end{Resp}
\switchcolumn*

\LA
\LessonHead{Lectio VIII}
\switchcolumn
\EN
\LessonHead{Lesson VIII}
\switchcolumn*

\LA
\noindent Si regínam ipsi ostendas diadémate ornatam, non præfert eam matri pannis detritis vestitæ, malletque illam incultam vidére, quam regínam mirifice amíctam. Nam, quod suum, quod aliénum est, non ex paupertáte vel divítiis, sed ex amóre existimare solet; et nihil plus requirit quam necessaria, atque ut lacte repletus est, statim a mamma abscedit. Non eisdem, quibus nos, ærumnis prémitur, nec pecuniárum jactura, rebusque simílibus; nec iisdem, quibus nos, fluxis rebus lætátur, neque córporum pulchritúdinem mirátur. Ideo dicebat: Talium est enim regnum cælórum; ut ex proposito voluntátis illa operemur, quæ natúra sua púeri faciunt.\par
\switchcolumn
\EN
\noindent If thou wert to show him a Queen adorned with her crown, he would not prefer her before his own mother, in raiment how faded soever, and he would rather see her, albeit unkempt, than the Queen in all her glorious apparel. For his use is to account of things whether they be his own, or of others, not by the standard of poverty and riches, but by that of love only. He seeketh no more than he needeth. When he is satisfied with milk, he leaveth the pap. The things that press upon us, such as the loss of money, and the like, do not press upon him, nor do the same transitory things that please us, please him, neither doth he gaze with admiration at loveliness of shape. Therefore Christ said: Of such is the kingdom of heaven, to make us do by force of will what little children do by nature.\par
\switchcolumn*

\LA
\begin{Resp}\Rsign Sint lumbi vestri præcíncti, et lucérnæ ardéntes in mánibus vestris: \Star{} Et vos símiles homínibus exspectántibus dóminum suum, quando revertátur a núptiis.\end{Resp}
\begin{Resp}\Vsign Vigiláte ergo, quia nescítis qua hora Dóminus vester ventúrus sit.\end{Resp}
\begin{Resp}\Rsign Et vos símiles homínibus exspectántibus dóminum suum, quando revertátur a núptiis.\end{Resp}
\begin{Resp}\Vsign Glória Patri, et Fílio, \Star{} et Spirítui Sancto.\end{Resp}
\begin{Resp}\Rsign Et vos símiles homínibus exspectántibus dóminum suum, quando revertátur a núptiis.\end{Resp}
\switchcolumn
\EN
\begin{Resp}\Rsign Let your loins be girded about, and your lights burning; \Star{} And ye yourselves like unto men that wait for their lord, when he will return from the wedding.\end{Resp}
\begin{Resp}\Vsign Watch therefore, for ye know not what hour your Lord doth come.\end{Resp}
\begin{Resp}\Rsign And ye yourselves like unto men that wait for their lord, when he will return from the wedding.\end{Resp}
\begin{Resp}\Vsign Glory be to the Father, and to the Son, \Star{} and to the Holy Ghost.\end{Resp}
\begin{Resp}\Rsign And ye yourselves like unto men that wait for their lord, when he will return from the wedding.\end{Resp}
\switchcolumn*

\LA
\LessonHead{Lectio IX}
\switchcolumn
\EN
\LessonHead{Lesson IX}
\switchcolumn*

\LA
\noindent Quia enim pharisæi non aliunde quam a nequítia et arrogántia ad agéndum ferebántur, ideo ubíque discipulos suos simplices esse jubet, illosque subíndicat, dum hos instituit. Nihil enim ita supérbiam parit, ut principátus et primi conséssus. Quóniam ígitur discípuli per totum terrárum orbem multum honoris consecuturi erant, ipsórum animos prævenit, nec sinit eos humánum quid pati, nec honórem a vulgo expétere, vel ante alios sese efferre. Nam, etiamsi hæc parva vidéntur esse, at malis ingéntibus causam præbent. Sic enim institúti pharisæi in malórum culmen ascendérunt, salutatiónes, primos conséssus et médios requiréntes; hinc in ardentem glóriæ cupiditátem, inde vero in impietátem lapsi sunt.\par
\switchcolumn
\EN
\noindent The Pharisees' usual springs of action were spite and vanity; therefore doth the Lord everywhere command His disciples to be simple, and in teaching the one, pointeth silently at the other class. Nothing breedeth pride so much as princedom and precedence. Since, then, His disciples were to receive much honour throughout all the world, He warneth their minds beforehand, and letteth them not stumble into the snare of men, nor go seeking for honours from the mob, nor put themselves forward before others. It is true, these may seem little things, but they give occasion for very great evils. It was when they were placed in these positions, that the Pharisees fell into their direst misfortunes from looking for salutations, and foremost or good places, they got into a keen desire of distinction, and from that into ungodliness.\par
\switchcolumn*

\LA
\TeDeum{Te Deum}
\switchcolumn
\EN
\TeDeum{Te Deum}
\switchcolumn*

\end{paracol}

\par\needspace{10\baselineskip}\addvspace{1.2em}{\centering\small\textsc{Die 21 Julii} · \textsc{21 July}\par}
\Office{S. Praxedis Virginis}{St. Praxedes, Virgin}{Simplex}
\addcontentsline{toc}{subsection}{Die 21 Julii · S. Praxedis Virginis \textit{— St. Praxedes, Virgin}}
\begin{paracol}{2}
\LA
\RefLine{Lectiones I–II: de Scriptura occurrente.}
\switchcolumn
\EN
\RefLine{Lessons I–II: from the Scripture of the day.}
\switchcolumn*

\LA
\LessonHead{Lectio III (pro commemoratione)}
\switchcolumn
\EN
\LessonHead{Lesson III (when commemorated)}
\switchcolumn*

\LA
\noindent Praxédes, virgo Romana, Pudentiánæ vírginis soror, Marco Antonino imperatóre Christiános persequente, eos facultátibus, opera, consolatióne, et omni caritátis officio prosequebátur. Nam alios domi occultábat; alios ad fidei constantiam hortabátur: aliórum córpora sepeliebat; iis, qui in carcere inclusi erant, qui in ergástulis exercebántur, nulla re déerat. Quæ, cum tantam Christianórum stragem jam ferre non posset, Deum precáta est, ut, si mori expediret, se e tantis malis eríperet. Itaque duodecimo Kalendas Augusti ad pietátis præmia vocátur in cælum. Cujus corpus a Pastore presbytero in patris et sororis Pudentiánæ sepúlcrum illátum est, quod erat in cœmeterio Priscillæ, via Salaria.\par
\switchcolumn
\EN
\noindent This Praxedes was a maiden of Rome, and the sister of the maiden Pudentiana. When the Emperor Marcus Antoninus was hunting down the Christians, she followed them constantly with money, labour, comfort, and every helpful office of Christian charity. Some she hid in her house, some she exhorted to firmness in professing the faith, of some she buried the bodies. For them that were in prison, and them that were toiling in slavery, she supplied every need. At last the sight of such butchery of Christians was more than she could bear, and she implored God that if it were expedient for her to die, He would release her from such suffering. And so upon the 21st day of July she was called away to receive the reward of her godly conversation in heaven. Pastor the Priest laid her body in the grave of her father (Pudens) and her sister Pudentiana, which was in the cemetery of Priscilla, upon the Salarian Way.\par
\switchcolumn*

\LA
\TeDeum{Te Deum}
\switchcolumn
\EN
\TeDeum{Te Deum}
\switchcolumn*

\end{paracol}

\par\needspace{10\baselineskip}\addvspace{1.2em}{\centering\small\textsc{Die 22 Julii} · \textsc{22 July}\par}
\Office{S. Mariæ Magdalenæ Pœnitentis}{St. Mary Magdalene, Penitent}{Duplex}
\addcontentsline{toc}{subsection}{Die 22 Julii · S. Mariæ Magdalenæ Pœnitentis \textit{— St. Mary Magdalene, Penitent}}
\begin{paracol}{2}
\LA
\LessonHead{Lectio I}
\switchcolumn
\EN
\LessonHead{Lesson I}
\switchcolumn*

\LA
\LessonTitle{De Canticis canticórum}
\switchcolumn
\EN
\LessonTitle{De Canticis Canticorum}
\switchcolumn*

\LA
\Cite{Cant 3:1-4}
\switchcolumn
\EN
\Cite{Song 3:1-4}
\switchcolumn*

\LA
\noindent In léctulo meo, per noctes, quæsívi quem díligit ánima mea; quæsívi illum, et non invéni. Surgam, et circuíbo civitátem: per vicos et platéas, quæram quem díligit ánima mea: quæsívi illum, et non invéni. Invenérunt me vígiles qui custódiunt civitátem. Num quem díligit ánima mea vidístis? Páululum cum pertransíssem eos, invéni quem díligit ánima mea: ténui eum, nec dimíttam, donec introdúcam illum in domum matris meæ, et in cubículum genetrícis meæ.\par
\switchcolumn
\EN
\noindent In my bed by night I sought him whom my soul loveth: I sought him, and found him not. I will rise, and will go about the city: in the streets and the broad ways I will seek him whom my soul loveth: I sought him, and I found him not. The watchmen who keep the city, found me: Have you seen him, whom my soul loveth? When I had a little passed by them, I found him whom my soul loveth: I held him: and I will not let him go, till I bring him into my mother's house, and into the chamber of her that bore me.\par
\switchcolumn*

\LA
\begin{Resp}\Rsign María Magdaléne, et áltera María ibant dilúculo ad monuméntum: \Star{} Jesum quem quǽritis, non est hic, surréxit sicut locútus est, præcédet vos in Galilǽam, ibi eum vidébitis.\end{Resp}
\begin{Resp}\Vsign Et valde mane una sabbatórum véniunt ad monuméntum, orto jam sole: et introëúntes vidérunt júvenem sedéntem in dextris, qui dixit illis.\end{Resp}
\begin{Resp}\Rsign Jesum quem quǽritis, non est hic, surréxit sicut locútus est, præcédet vos in Galilǽam, ibi eum vidébitis.\end{Resp}
\switchcolumn
\EN
\begin{Resp}\Rsign Mary Magdalene and the other Mary went very early to the sepulchre. \Star{} That Jesus whom ye seek, is not here: for He is risen, as He said: He goeth before you into Galilee; there shall ye see Him.\end{Resp}
\begin{Resp}\Vsign And very early in the morning, the first day of the week, they came unto the sepulchre, at the rising of the sun; and, entering into the sepulchre, they saw a young man sitting upon the right side, who saith unto them:\end{Resp}
\begin{Resp}\Rsign That Jesus Whom ye seek is not here: for He is risen, as He said: He goeth before you into Galilee: there shall ye see Him.\end{Resp}
\switchcolumn*

\LA
\LessonHead{Lectio II}
\switchcolumn
\EN
\LessonHead{Lesson II}
\switchcolumn*

\LA
\Cite{Cant 8:1-4}
\switchcolumn
\EN
\Cite{Song 8:1-4}
\switchcolumn*

\LA
\noindent Quis mihi det te fratrem meum, sugentem úbera matris meæ, ut invéniam te foris, et deósculer te, et jam me nemo despíciat? Apprehéndam te, et ducam in domum matris meæ; ibi me docébis, et dabo tibi póculum ex vino condíto, et mustum malórum granatórum meórum. Læva ejus sub cápite meo, et déxtera illíus amplexábitur me. Adjúro vos, fíliæ Jerúsalem, ne suscitétis, neque evigiláre faciátis diléctam, donec ipsa velit.\par
\switchcolumn
\EN
\noindent Who shall give thee to me for my brother, sucking the breasts of my mother, that I may find thee without, and kiss thee, and now no man may despise me? I will take hold of thee, and bring thee Into my mother's house: there thou shalt teach me, and I will give thee a cup of spiced wine and new wine of my pomegranates. His left hand under my head, and his right hand shall embrace me. I adjure you, O daughters of Jerusalem, that you stir not up, nor awake my love till she please.\par
\switchcolumn*

\LA
\begin{Resp}\Rsign Congratulámini mihi, omnes qui dilígitis Dóminum, quia quem quærébam, appáruit mihi: \Star{} Et dum flerem ad monuméntum, vidi Dóminum meum, allelúja.\end{Resp}
\begin{Resp}\Vsign Recedéntibus discípulis, non recedébam, et amóris ejus igne succénsa, ardébam desidério.\end{Resp}
\begin{Resp}\Rsign Et dum flerem ad monuméntum, vidi Dóminum meum, allelúja.\end{Resp}
\switchcolumn
\EN
\begin{Resp}\Rsign Rejoice with me, all ye that love the Lord, for I sought Him and He hath appeared unto me \Star{} And while as I was weeping at the Sepulchre, I saw the Lord.\end{Resp}
\begin{Resp}\Vsign When His disciples were gone away, I tarried still; and the fire of love in my heart yearned after Him.\end{Resp}
\begin{Resp}\Rsign And while as I was weeping at the Sepulchre, I saw the Lord.\end{Resp}
\switchcolumn*

\LA
\LessonHead{Lectio III}
\switchcolumn
\EN
\LessonHead{Lesson III}
\switchcolumn*

\LA
\Cite{Cant 8:5-7}
\switchcolumn
\EN
\Cite{Song 8:5-7}
\switchcolumn*

\LA
\noindent Quæ est ista, quæ ascéndit de desérto, delíciis áffluens, inníxa super diléctum suum? Sub árbore malo suscitávi te, ibi corrúpta est mater tua, ibi violáta est génitrix tua. Pone me ut signáculum super cor tuum, ut signáculum super brácchium tuum, quia fortis est ut mors diléctio, dura sicut inférnus æmulátio; lámpades ejus lámpades ignis atque flammárum. Aquæ multæ non potuérunt exstínguere caritátem, nec flúmina óbruent illam.\par
\switchcolumn
\EN
\noindent Who is this that cometh up from the desert, flowing with delights, leaning upon her beloved? Under the apple tree I raised thee up: there thy mother was corrupted, there she was defloured that bore thee. Put me as a seal upon thy heart, as a seal upon thy arm, for love is strong as death, jealousy as hard as hell, the lamps thereof are fire and flames. Many waters cannot quench charity, neither can the floods drown it: if a man should give all the substance of his house for love, he shall despise it as nothing.\par
\switchcolumn*

\LA
\begin{Resp}\Rsign Tulérunt Dóminum meum, et néscio, ubi posuérunt eum. Dicunt ei Angeli: Múlier, quid ploras? surréxit sicut dixit: \Star{} Præcédet vos in Galilǽam: ibi eum vidébitis.\end{Resp}
\begin{Resp}\Vsign Cum ergo fleret, inclinávit se, et prospéxit in monuméntum: et vidit duos Angelos in albis, sedéntes, qui dicunt ei.\end{Resp}
\begin{Resp}\Rsign Præcédet vos in Galilǽam: ibi eum vidébitis.\end{Resp}
\begin{Resp}\Vsign Glória Patri, et Fílio, \Star{} et Spirítui Sancto.\end{Resp}
\begin{Resp}\Rsign Præcédet vos in Galilǽam: ibi eum vidébitis.\end{Resp}
\switchcolumn
\EN
\begin{Resp}\Rsign They have taken away my Lord, and I know not where they have laid Him. The Angels say unto her: Woman, why weepest thou? He is risen, as He said. \Star{} He goeth before you into Galilee; there shall ye see Him.\end{Resp}
\begin{Resp}\Vsign And as she wept, she stooped down and looked into the Sepulchre, and saw two Angels in white, sitting; and they say unto her:\end{Resp}
\begin{Resp}\Rsign He goeth before you into Galilee; there shall ye see Him.\end{Resp}
\begin{Resp}\Vsign Glory be to the Father, and to the Son, \Star{} and to the Holy Ghost.\end{Resp}
\begin{Resp}\Rsign He goeth before you into Galilee; there shall ye see Him.\end{Resp}
\switchcolumn*

\LA
\LessonHead{Lectio IV}
\switchcolumn
\EN
\LessonHead{Lesson IV}
\switchcolumn*

\LA
\LessonTitle{Sermo sancti Gregórii Papæ}
\switchcolumn
\EN
\LessonTitle{From the Sermons of Pope St. Gregory (the Great.)}
\switchcolumn*

\LA
\Source{Homilia 25 in Evangelia}
\switchcolumn
\EN
\Source{25th on the Gospels.}
\switchcolumn*

\LA
\noindent María Magdaléne, quæ fúerat in civitáte peccátrix, amándo veritátem, lavit lácrimis máculas críminis: et vox Veritátis implétur, qua dícitur: Dimíssa sunt ei peccáta multa, quia diléxit multum. Quæ enim prius frígida peccándo remánserat, póstmodum amándo fórtiter ardébat. Quæ a monuménto Dómini, étiam discípulis recedéntibus, non recedébat. Exquirébat quem non invénerat; flebat inquiréndo, et amóris sui igne succénsa, ejus, quem ablátum crédidit, ardébat desidério. Unde cóntigit, ut eum sola tunc vidéret, quæ remánserat ut quǽreret: quia nimírum virtus boni óperis perseverántia est.\par
\switchcolumn
\EN
\noindent Mary Magdalen, a woman in the city, which was a sinner, through love of the truth, washed away in her tears the defilement of her sins, and the words of the Truth are fulfilled which He spake: Her sins, which are many, are forgiven; for she loved much. She who had remained chilly in sin, became fiery through love. When even His disciples went away again unto their own home, Mary still stood without at the sepulchre of Christ, weeping. She sought Him Whom her soul loved, but she found Him not. She searched for Him with tears; she yearned with strong desire for Him Who, she believed, had been taken away. And thus it befell her, that being the only one who had remained to seek Him, she was the only one that saw Him. It is the truth that the backbone of a good work is perseverance.\par
\switchcolumn*

\LA
\begin{Resp}\Rsign Propter veritátem, et mansuetúdinem, et justítiam: \Star{} Et dedúcet te mirabíliter déxtera tua.\end{Resp}
\begin{Resp}\Vsign Spécie tua et pulchritúdine tua inténde, próspere procéde, et regna.\end{Resp}
\begin{Resp}\Rsign Et dedúcet te mirabíliter déxtera tua.\end{Resp}
\switchcolumn
\EN
\begin{Resp}\Rsign Because of truth, and meekness, and righteousness; \Star{} And thy right hand shall lead thee wonderfully.\end{Resp}
\begin{Resp}\Vsign In thy comeliness, and thy beauty, go forward, fare prosperously, and reign.\end{Resp}
\begin{Resp}\Rsign And thy right hand shall lead thee wonderfully.\end{Resp}
\switchcolumn*

\LA
\LessonHead{Lectio V}
\switchcolumn
\EN
\LessonHead{Lesson V}
\switchcolumn*

\LA
\noindent Quæsívit ergo prius, et mínime invénit: perseverávit ut quæreret, unde et cóntigit ut inveníret: actúmque est, ut desidéria diláta créscerent, et crescéntia cáperent quod inveníssent. Hinc est enim quod de eódem sponsa Ecclésia in Cánticis canticórum dicit: In léctulo meo per noctes quæsívi quem díligit ánima mea. Diléctum namque in léctulo quærimus, quando, in præséntis vitæ aliquántula réquie. Redemptóris nostri desidério suspirámus. Per noctem quærimus: quia, etsi jam in illo mens vígilat, tamen adhuc óculus calígat.\par
\switchcolumn
\EN
\noindent At first when she sought Him, she found Him not; she went on searching, and so it came to pass that she found Him; and this was so, to the end that her longing might grow in earnestness, and so in its earnestness might find what it sought. Hence is it that the Bride in the Song of Songs saith as representing the Church: By night on my bed I sought him whom my soul loveth. We seek on our bed for Him Whom our soul loveth, when, having got some little rest in this world, we still sigh for the Presence of our Redeemer but it is by night that we so seek Him, for though our mind may be on the alert for Him, yet still He is hidden from our eyes by the darkness that now is.\par
\switchcolumn*

\LA
\begin{Resp}\Rsign Dilexísti justítiam, et odísti iniquitátem: \Star{} Proptérea unxit te Deus, Deus tuus, óleo lætítiæ.\end{Resp}
\begin{Resp}\Vsign Propter veritátem, et mansuetúdinem, et justítiam.\end{Resp}
\begin{Resp}\Rsign Proptérea unxit te Deus, Deus tuus, óleo lætítiæ.\end{Resp}
\switchcolumn
\EN
\begin{Resp}\Rsign Thou hast loved righteousness, and hated iniquity; \Star{} Therefore God, thy God, hath anointed thee with the oil of gladness.\end{Resp}
\begin{Resp}\Vsign Because of truth, and meekness, and righteousness.\end{Resp}
\begin{Resp}\Rsign Therefore God, thy God, hath anointed thee with the oil of gladness.\end{Resp}
\switchcolumn*

\LA
\LessonHead{Lectio VI}
\switchcolumn
\EN
\LessonHead{Lesson VI}
\switchcolumn*

\LA
\noindent Sed, qui diléctum suum non invénit, restat ut surgat, civitátem circúmeat, id est, sanctam electórum Ecclésiam mente et inquisitióne percúrrat; per vicos eum et platéas quærat, id est, per angústa et lata gradiéntes aspíciat, ut, si qua inveníre in eis váleat, ejus vestígia exquírat: quia sunt nonnúlli étiam vitæ sæculáris, qui imitándum áliquid hábeant de actióne virtútis. Quæréntes autem nos vígiles invéniunt, qui custódiunt civitátem: quia sancti Patres, qui Ecclésiæ statum custódiunt, bonis nostris stúdiis occúrrunt, ut suo vel verbo vel scripto nos dóceant. Quos cum páululum pertransímus, invenímus quem dilígimus: quia Redémptor noster, etsi humilitáte homo inter hómines, divinitáte tamen super hómines fuit.\par
\switchcolumn
\EN
\noindent But if we find not Him Whom our soul loveth, it remaineth that we should rise and go about the city, that is, by thought and questioning, go through the holy Church of the elect seek Him in the streets, and in the broad ways, that is, walk anxiously looking about us both in the narrow and the broad places, that if we can, we may find His footsteps there for there are some even of those who live for the world, from whom something may be learnt to be imitated by a godly man. As we thus go wakefully about, the watchmen, that keep the city, find us; the holy Fathers, who are the watchmen of the bulwarks of the Church, come to meet our good endeavours, and to teach us either by their words or by their writings. And it needeth but a little to pass from them, but we find Him Whom our soul loveth (a little we must pass,) for albeit our Redeemer in lowliness became a man among men, yet by right of His Divine Nature He is still above men.\par
\switchcolumn*

\LA
\begin{Resp}\Rsign Fallax grátia, et vana est pulchritúdo: \Star{} Múlier timens Dóminum, ipsa laudábitur.\end{Resp}
\begin{Resp}\Vsign Date ei de fructu mánuum suárum, et laudent eam in portis ópera ejus.\end{Resp}
\begin{Resp}\Rsign Múlier timens Dóminum, ipsa laudábitur.\end{Resp}
\begin{Resp}\Vsign Glória Patri, et Fílio, \Star{} et Spirítui Sancto.\end{Resp}
\begin{Resp}\Rsign Múlier timens Dóminum, ipsa laudábitur.\end{Resp}
\switchcolumn
\EN
\begin{Resp}\Rsign Favour is deceitful, and beauty is vain \Star{} A woman that feareth God she shall be praised.\end{Resp}
\begin{Resp}\Vsign Give her of the fruit of her hands, and let her own works praise her in the gates.\end{Resp}
\begin{Resp}\Rsign A woman that feareth God she shall be praised.\end{Resp}
\begin{Resp}\Vsign Glory be to the Father, and to the Son, \Star{} and to the Holy Ghost.\end{Resp}
\begin{Resp}\Rsign A woman that feareth God she shall be praised.\end{Resp}
\switchcolumn*

\LA
\LessonHead{Lectio VII}
\switchcolumn
\EN
\LessonHead{Lesson VII}
\switchcolumn*

\LA
\LessonTitle{Léctio sancti Evangélii secúndum Lucam}
\switchcolumn
\EN
\LessonTitle{From the Holy Gospel according to Luke}
\switchcolumn*

\LA
\Cite{Luc 7:36-50}
\switchcolumn
\EN
\Cite{Luke 7:36-50}
\switchcolumn*

\LA
\noindent In illo témpore: Rogábat Jesum quidam de pharisǽis, ut manducáret cum illo. Et ingréssus domum pharisǽi discúbuit. Et réliqua.\par
\switchcolumn
\EN
\noindent At that time One of the Pharisees desired Jesus that He would eat with him. And He went into the Pharisee's house, and sat down to meat. And so on.\par
\switchcolumn*

\LA
\LessonTitle{Homilía sancti Augustíni Epíscopi}
\switchcolumn
\EN
\LessonTitle{Homily by St. Augustine, Bishop (of Hippo.)}
\switchcolumn*

\LA
\Source{Liber 50 Homilía 23, tom. 10}
\switchcolumn
\EN
\Source{Bk. 1. Hom. 23, tom. x.}
\switchcolumn*

\LA
\noindent Evangélium cum legerétur, attentíssime audístis; et res gesta narráta atque versáta est ante óculos cordis vestri. Vidístis enim, non carne, sed mente, Dóminum Jesum Christum in domo pharisǽi recumbéntem, et ab illo invitátum non fastidiéntem. Vidístis étiam in civitáte mulíerem famósam, mala útique fama, quæ erat peccátrix, non invitátam irruísse convívio ubi suus médicus recumbébat, et quæsísse pia impudéntia sanitátem; írruens, quasi importúna convívio, opportúna benefício. Nóverat enim quanto morbo laboráret; et illum ad sanándum esse idóneum, ad quem vénerat, sciébat.\par
\switchcolumn
\EN
\noindent Ye have listened carefully to the Gospel while it was being read, so that the thing told hath, as it were, passed before the eyes of your heart. Ye have seen in your mind's eye, albeit not with bodily sight, the Lord Jesus Christ sitting down to meat in the Pharisee's house, and not refusing when He is bidden of him. Ye have seen also an infamous woman of the city, one of utterly bad character, a sinner, thrusting herself in an uninvited guest, to the banquet where her Healer was sitting, and seeking health at His hands with godly shamelessness; thrusting herself in eager for mercy, as though eager for the feast. She knew under what a disease she laboured, and she knew that He unto Whom she came was mighty to cure it.\par
\switchcolumn*

\LA
\begin{Resp}\Rsign Os suum apéruit sapiéntiæ, et lex cleméntiæ in lingua ejus: considerávit sémitas domus suæ, \Star{} Et panem otiósa non comédit.\end{Resp}
\begin{Resp}\Vsign Gustávit et vidit quia bona est negotiátio ejus: non exstinguétur in nocte lucérna ejus.\end{Resp}
\begin{Resp}\Rsign Et panem otiósa non comédit.\end{Resp}
\switchcolumn
\EN
\begin{Resp}\Rsign She openeth her mouth with wisdom, and in her tongue is the law of kindness. She looketh well to the ways of her household \Star{} And eateth not the bread of idleness.\end{Resp}
\begin{Resp}\Vsign She tasteth and perceiveth that her merchandise is good. Her candle goeth not out by night.\end{Resp}
\begin{Resp}\Rsign And she eateth not the bread of idleness.\end{Resp}
\switchcolumn*

\LA
\LessonHead{Lectio VIII}
\switchcolumn
\EN
\LessonHead{Lesson VIII}
\switchcolumn*

\LA
\noindent Accéssit ergo non ad caput Dómini, sed ad pedes. Et quæ diu male ambuláverat, vestígia recta quærébat. Prius fudit lácrimas cordis, et lavit Dómini pedes obséquio confessiónis, capíllis suis tersit, osculáta est, unxit; tácita loquebátur, non sermónem promébat, sed devotiónem ostendébat. Quia ergo tétigit Dóminum rigándo, osculándo, tergéndo, unguéndo pedes ejus; pharisæus, qui invitáverat Dóminum Jesum Christum, quia ex illo génere erat hóminum superbórum, de quibus Isaías prophéta dicit: Qui dicunt, Recéde longe a me, noli me tángere, quóniam mundus sum; putávit Dóminum nescísse mulíerem.\par
\switchcolumn
\EN
\noindent She drew near therefore, not unto the Lord's Head, but unto His Feet. She that had so long walked the paths of sin betook her unto the Feet that went about doing good. She first poured forth heart-felt tears, and washed the Lord's Feet with the humble service of her acknowledgment, wiped them with her hair, kissed them, and anointed them. Her silence cried aloud, next in words but in manifested love. The Pharisee, who had desired the Lord Jesus Christ that He would eat with him, belonged to that class of proud men concerning whom the Prophet Isaiah saith people which say, Stand by thyself, come not near to me; for I am holier than thou. (lxv. 5.) When therefore he saw how this woman touched the Lord's Feet with her tears, her kisses, her hair, and her ointment, he spake within himself, saying This Man, if He were a Prophet, would have known who and what manner of woman this is, that toucheth Him; for she is a sinner.\par
\switchcolumn*

\LA
\begin{Resp}\Rsign Regnum mundi et omnem ornátum sǽculi contémpsi, propter amórem Dómini mei Jesu Christi: \Star{} Quem vidi, quem amávi, in quem crédidi, quem diléxi.\end{Resp}
\begin{Resp}\Vsign Eructávit cor meum verbum bonum: dico ego ópera mea Regi.\end{Resp}
\begin{Resp}\Rsign Quem vidi, quem amávi, in quem crédidi, quem diléxi.\end{Resp}
\begin{Resp}\Vsign Glória Patri, et Fílio, \Star{} et Spirítui Sancto.\end{Resp}
\begin{Resp}\Rsign Quem vidi, quem amávi, in quem crédidi, quem diléxi.\end{Resp}
\switchcolumn
\EN
\begin{Resp}\Rsign The kingdom of this world and all the beauty of life I have esteemed as nothing, for the excellency of the love of Jesus Christ my Lord \Star{} Whom, having seen, I loved Whom, having believed, I longed after.\end{Resp}
\begin{Resp}\Vsign My heart is overflowing with a good matter I speak of my works unto the King.\end{Resp}
\begin{Resp}\Rsign Whom, having seen, I loved Whom, having believed, I longed after.\end{Resp}
\begin{Resp}\Vsign Glory be to the Father, and to the Son, \Star{} and to the Holy Ghost.\end{Resp}
\begin{Resp}\Rsign Whom, having seen, I loved Whom, having believed, I longed after.\end{Resp}
\switchcolumn*

\LA
\LessonHead{Lectio IX}
\switchcolumn
\EN
\LessonHead{Lesson IX}
\switchcolumn*

\LA
\noindent O pharisǽe invitátor et irrísor Dómini, Dóminum pascis, et a quo pascéndus sis, non intélligis? Unde scis Dóminum nescísse, quæ fúerit illa múlier, nisi quia permíssa est accédere, nisi quia, illo patiénte, osculáta est pedes ejus, nisi quia tersit, nisi quia unxit? Hæc enim non débuit permítti fácere in pédibus mundis múlier immúnda? Ad illíus ergo pharisæi pedes si talis múlier accessísset, dictúrus erat, quod Isaías de tálibus dicit: Recéde a me, noli me tángere, quóniam mundus sum. Accéssit autem ad Dóminum immúnda, ut redíret munda; accéssit ægra, ut redíret sana; accéssit conféssa, ut redíret proféssa.\par
\switchcolumn
\EN
\noindent Pharisee, that biddest and scornest the Lord! Thou invitest the Lord to meat, and thou knowest not Him That should have given thee to eat! Whence knowest thou that the Lord knoweth not who and what manner of woman this is, save from this, that she is allowed to draw near unto Him, and that He suffereth her to kiss His Feet, to wipe them, and to anoint them? Ought not an unclean woman to have been permitted to do these things to clean feet? If such a woman had drawn near to the feet of this Pharisee, he would have said to her what Isaiah putteth into the mouth of such Stand by thyself, come not near to me, for I am holier than thou. But she came unto the Lord unclean that she might go away cleansed, sick, that she might go away healed, with confession, that she might go away with thanksgiving.\par
\switchcolumn*

\LA
\TeDeum{Te Deum}
\switchcolumn
\EN
\TeDeum{Te Deum}
\switchcolumn*

\end{paracol}

\par\needspace{10\baselineskip}\addvspace{1.2em}{\centering\small\textsc{Die 23 Julii} · \textsc{23 July}\par}
\Office{S. Apollinaris Episcopi et Martyris}{St. Apollinaris, Bishop and Martyr}{Duplex}
\addcontentsline{toc}{subsection}{Die 23 Julii · S. Apollinaris Episcopi et Martyris \textit{— St. Apollinaris, Bishop and Martyr}}
\begin{paracol}{2}
\LA
\RefLine{Lectiones I–III: de Scriptura occurrente.}
\switchcolumn
\EN
\RefLine{Lessons I–III: from the Scripture of the day.}
\switchcolumn*

\LA
\LessonHead{Lectio IV}
\switchcolumn
\EN
\LessonHead{Lesson IV}
\switchcolumn*

\LA
\noindent Apollináris cum Príncipe Apostolórum Antiochía Romam venit; a quo, ordinátus epíscopus, Ravennam ad Christi Dómini Evangélium prædicándum mittitur; ubi, cum ad Christi fidem plurimos convérteret, captus ab idolórum sacerdótibus gráviter cæsus est. Cumque ipso orante Bonifatius nobilis vir, qui diu mutus fuerat, loquerétur, ejusque fília immundo spíritu liberáta esset; íterum est in illum commóta seditio. Itaque virgis cæsus, ardéntes carbónes nudis pédibus prémere cogitur; quem cum subjéctus ignis nihil læderet, ejícitur extra urbem.\par
\switchcolumn
\EN
\noindent Apollinaris came from Antioch to Rome with the Prince of the Apostles, and was by him ordained a Bishop, and sent to Ravenna to preach the Gospel of the Lord Christ. He had already converted a great number of persons to the Christian Faith, when the idolatrous priests caught him and gave him a sharp flogging. A second riot was got up against him on account of one Boniface, a nobleman who had long been dumb, speaking, and his daughter being delivered from an unclean spirit. On this occasion Apollinaris was flogged again, and made to walk barefoot over hot embers. The fire did him no harm, and he was expelled from the city.\par
\switchcolumn*

\LA
\begin{Resp}\Rsign Honéstum fecit illum Dóminus, et custodívit eum ab inimícis, et a seductóribus tutávit illum: \Star{} Et dedit illi claritátem ætérnam.\end{Resp}
\begin{Resp}\Vsign Descendítque cum illo in fóveam, et in vínculis non derelíquit eum.\end{Resp}
\begin{Resp}\Rsign Et dedit illi claritátem ætérnam.\end{Resp}
\switchcolumn
\EN
\begin{Resp}\Rsign The Lord made him honourable, and defended him from his enemies, and kept him safe from those that lay in wait for him, \Star{} And gave him perpetual glory.\end{Resp}
\begin{Resp}\Vsign He went down with him into the pit, and left him not in bonds.\end{Resp}
\begin{Resp}\Rsign And gave him perpetual glory.\end{Resp}
\switchcolumn*

\LA
\LessonHead{Lectio V}
\switchcolumn
\EN
\LessonHead{Lesson V}
\switchcolumn*

\LA
\noindent Is vero latens aliquámdiu cum quibusdam Christiánis, inde profectus est in Æmiliam, ubi Rufini patricii filiam mortuam ad vitam revocávit; ut proptérea tota Rufini família in Jesum Christum crederet. Quare veheménter incénsus præfectus accersit Apollinarem, et cum eo gravius agit, ut finem fáciat disseminándi in urbe Christi fidem. Cujus cum Apollináris jussa neglígeret, equuleo cruciátur; in cujus plagas aqua fervens infúnditur, saxoque os túnditur: mox férreis vinculis constrictus includitur in carcere. Quarto die impositus in navem, mittitur in exsílium; ac facto naufragio venit in Mysiam, inde ad ripam Danubii, póstea in Thraciam.\par
\switchcolumn
\EN
\noindent Apollinaris lay hid for a while with certain Christians. Thence he went to Emilia, where he restored to life the dead daughter of the Patrician Rufinus, so that the whole household of Rufinus might believe in Jesus Christ. This affair greatly incensed the Praefect, who sent for Apollinaris, and earnestly dealt with him to induce him to cease spreading the Christian Faith in that city. As Apollinaris paid no heed to the Prefect's orders, he was tortured on the rack, boiling water poured on his wounds, and his mouth bruised with a stone, after which he was ironed and cast into prison. On the fourth day he was put on board a ship and sent into banishment. The ship was wrecked, and he so came to Mysia, thence to the shores of the Danube, and afterwards into Thrace.\par
\switchcolumn*

\LA
\begin{Resp}\Rsign Desidérium ánimæ ejus tribuísti ei Dómine, \Star{} Et voluntáte labiórum ejus non fraudásti eum.\end{Resp}
\begin{Resp}\Vsign Quóniam prævenísti eum in benedictiónibus dulcédinis: posuísti in cápite ejus corónam de lápide pretióso.\end{Resp}
\begin{Resp}\Rsign Et voluntáte labiórum ejus non fraudásti eum.\end{Resp}
\switchcolumn
\EN
\begin{Resp}\Rsign Lord, Thou hast given him his heart's desire. \Star{} And hast not withholden the request of his lips.\end{Resp}
\begin{Resp}\Vsign For Thou hast prevented him with the blessings of sweetness: Thou hast set a crown of precious stones upon his head.\end{Resp}
\begin{Resp}\Rsign And hast not withholden the request of his lips.\end{Resp}
\switchcolumn*

\LA
\LessonHead{Lectio VI}
\switchcolumn
\EN
\LessonHead{Lesson VI}
\switchcolumn*

\LA
\noindent Cum autem in Serápidis templo dæmon se responsa datúrum negaret, dum ibidem Petri Apóstoli discipulus morarétur; diu conquisítus, invéntus est Apollináris: qui íterum jubétur navigare. Ita reversus Ravennam, ab iisdem illis idolórum sacerdótibus accusatus, centurióni custodiendus tráditur; qui, cum occulte Christum cóleret, noctu Apollinarem dimísit. Re cógnita, satéllites eum persequúntur, et plagis in itinere confectum, quod mortuum crederent, relinquunt. Quem cum inde Christiáni sustulissent, septimo die exhortans illos ad fidei constantiam, martyrii glória clarus migrávit e vita. Cujus corpus prope murum urbis sepúltum est.\par
\switchcolumn
\EN
\noindent However, the devil in the temple of Serapis declared that he could not give oracles, while the disciple of the Apostle Peter abode in these parts, and after a long search Apollinaris was found and commanded again to take ship. Thus he went back to Ravenna, where he was denounced by the same idolatrous priests as before, and given into the keeping of a centurion. This centurion was a secret worshipper of Christ, and in the night he let Apollinaris go. When it became known some of the officers of justice followed after him, caught him on the road, beat him till they thought he was dead, and left him. Some Christians took him up, but on the seventh day, still exhorting them to stand firm in the Faith, he departed this life with the glorious splendour of martyrdom. His body was buried hard by the wall of the city.\par
\switchcolumn*

\LA
\begin{Resp}\Rsign Stola jucunditátis índuit eum Dóminus: \Star{} Et corónam pulchritúdinis pósuit super caput ejus.\end{Resp}
\begin{Resp}\Vsign Cibávit illum Dóminus pane vitæ et intelléctus: et aqua sapiéntiæ salutáris potávit illum.\end{Resp}
\begin{Resp}\Rsign Et corónam pulchritúdinis pósuit super caput ejus.\end{Resp}
\begin{Resp}\Vsign Glória Patri, et Fílio, \Star{} et Spirítui Sancto.\end{Resp}
\begin{Resp}\Rsign Et corónam pulchritúdinis pósuit super caput ejus.\end{Resp}
\switchcolumn
\EN
\begin{Resp}\Rsign The Lord hath put on him a robe of honour, \Star{} And put about his head a crown of joy.\end{Resp}
\begin{Resp}\Vsign With the bread of life and understanding hath the Lord fed him, and given him the water of health and wisdom to drink.\end{Resp}
\begin{Resp}\Rsign And put about his head a crown of joy.\end{Resp}
\begin{Resp}\Vsign Glory be to the Father, and to the Son, \Star{} and to the Holy Ghost.\end{Resp}
\begin{Resp}\Rsign And put about his head a crown of joy.\end{Resp}
\switchcolumn*

\LA
\LessonHead{Lectio VII}
\switchcolumn
\EN
\LessonHead{Lesson VII}
\switchcolumn*

\LA
\LessonTitle{Léctio sancti Evangélii secúndum Lucam.}
\switchcolumn
\EN
\LessonTitle{From the Holy Gospel according to Luke}
\switchcolumn*

\LA
\Cite{Luc 22:24-30}
\switchcolumn
\EN
\Cite{Luke 22:24-30}
\switchcolumn*

\LA
\noindent In illo témpore: Facta est contentio inter discipulos, quis eórum viderétur esse major. Et réliqua.\par
\switchcolumn
\EN
\noindent At that time There was a strife among the disciples, which of them should be accounted the greatest. And so on.\par
\switchcolumn*

\LA
\LessonTitle{Homilía sancti Ambrósii Epíscopi.}
\switchcolumn
\EN
\LessonTitle{Homily by St. Ambrose, Bishop (of Milan.)}
\switchcolumn*

\LA
\Source{Liber 10 in Lucæ, cap. 22, post initium.}
\switchcolumn
\EN
\Source{Bk. x. on Luke xxii.}
\switchcolumn*

\LA
\noindent Regnum Dei non est de hoc mundo. Non ergo æqualitátis hómini ad Deum, sed similitúdinis æmulátio est. Solus enim Christus est plena imago Dei, propter expressam in se Paternæ claritúdinis unitátem. Justus autem homo ad imáginem Dei est, si propter imitandam divinæ conversatiónis similitúdinem mundum hunc Dei cognitióne contemnat, voluptatesque terrenas verbi perceptióne despíciat, quo álimur in vitam: unde et corpus Christi édimus, ut vitæ ætérnæ possímus esse participes.\par
\switchcolumn
\EN
\noindent My kingdom, our God hath said, is not of this world. (John xviii. 36.) Man, then, must strive, not to be equal with God, but to be like unto God. Christ alone is the full image of God, as being the brightness of His Father's glory, and the express image of His Person. (Heb. i. 3.) But the righteous man is made after the image of God, in so far as, for the sake of following the example of his God, he, through knowledge of God, setteth no esteem upon this world, and looketh down upon all carnal motions of the will, through having taken in that Word whereby we are fed unto life everlasting and this is the end whereto we eat the Body of Christ, namely, that we may have eternal life. (John vi. 50.)\par
\switchcolumn*

\LA
\begin{Resp}\Rsign Coróna áurea super caput ejus, \Star{} Expréssa signo sanctitátis, glória honóris, et opus fortitúdinis.\end{Resp}
\begin{Resp}\Vsign Quóniam prævenísti eum in benedictiónibus dulcédinis, posuísti in cápite ejus corónam de lápide pretióso.\end{Resp}
\begin{Resp}\Rsign Expréssa signo sanctitátis, glória honóris, et opus fortitúdinis.\end{Resp}
\switchcolumn
\EN
\begin{Resp}\Rsign A crown of gold upon his head, \Star{} Wherein is engraved Holiness, an ornament of honour, a costly work.\end{Resp}
\begin{Resp}\Vsign For Thou hast prevented him with the blessings of sweetness, Thou hast set a crown of precious stones upon his head.\end{Resp}
\begin{Resp}\Rsign Wherein is engraved Holiness, an ornament of honour, a costly work.\end{Resp}
\switchcolumn*

\LA
\LessonHead{Lectio VIII}
\switchcolumn
\EN
\LessonHead{Lesson VIII}
\switchcolumn*

\LA
\noindent Non enim victus et potus nobis præmii loco spondétur et honoris, sed communicátio gratiæ cæléstis et vitæ. Neque duodecim throni tamquam áliqua corporalis sunt receptácula sessiónis. Sed, quia sicut secúndum divinam similitúdinem júdicat Christus cognitióne cordium, non interrogatióne factórum, virtútem remúnerans, impietatémque condemnans; ita et Apóstoli in judícium spiritale formántur remuneratióne fidei et exsecratióne perfidiæ, virtúte errórem redarguéntes, sacrílegos ódio persequéntes.\par
\switchcolumn
\EN
\noindent The reward which is promised unto us is not meat and drink, but a part in that grace and that life which come down from heaven. Neither are the twelve thrones \textasciitilde{}(Matth. xix. 28) to be understood as meaning chairs for bodies to sit in, but as meaning a dignity like that of God Himself, wherein they that have left all and followed Christ shall judge as Christ doth, not by dint of cross-examinations, but by simple knowledge of the heart, rewarding the good, and condemning the evil. Thus are the Apostles erected into a ghostly tribunal, for the rewarding of faith and the cursing of unbelief, by their might breaking back errors, and laying upon blasphemers the just punishment of hatred.\par
\switchcolumn*

\LA
\begin{Resp}\Rsign Hic est vere Martyr, qui pro Christi nómine sánguinem suum fudit: \Star{} Qui minas júdicum non tímuit, nec terrénæ dignitátis glóriam quæsívit, sed ad cæléstia regna pervénit.\end{Resp}
\begin{Resp}\Vsign Justum dedúxit Dóminus per vias rectas, et osténdit illi regnum Dei.\end{Resp}
\begin{Resp}\Rsign Qui minas júdicum non tímuit, nec terrénæ dignitátis glóriam quæsívit, sed ad cæléstia regna pervénit.\end{Resp}
\begin{Resp}\Vsign Glória Patri, et Fílio, \Star{} et Spirítui Sancto.\end{Resp}
\begin{Resp}\Rsign Qui minas júdicum non tímuit, nec terrénæ dignitátis glóriam quæsívit, sed ad cæléstia regna pervénit.\end{Resp}
\switchcolumn
\EN
\begin{Resp}\Rsign This is a Martyr indeed, who shed his blood for Christ's Name's sake \Star{} Who feared not for the threats of judges, nor sought to be great with the glory of this world, but pressed on unto the kingdom of heaven.\end{Resp}
\begin{Resp}\Vsign The Lord guided the righteous in right paths, and showed him the kingdom of God.\end{Resp}
\begin{Resp}\Rsign Who feared not for the threats of judges, nor sought to be great with the glory of this world, but pressed on unto the kingdom of heaven.\end{Resp}
\begin{Resp}\Vsign Glory be to the Father, and to the Son, \Star{} and to the Holy Ghost.\end{Resp}
\begin{Resp}\Rsign Who feared not for the threats of judges, nor sought to be great with the glory of this world, but pressed on unto the kingdom of heaven.\end{Resp}
\switchcolumn*

\LA
\LessonHead{Lectio IX}
\switchcolumn
\EN
\LessonHead{Lesson IX}
\switchcolumn*

\LA
\noindent Convertámur ígitur, et caveámus, ne in perditiónem aliqua inter nos de prælatióne possit esse contentio. Si enim contendébant Apóstoli, non excusatióni obténditur, sed cautióni propónitur. Si Petrus aliquándo convértitur, qui ad primam Dómini secutus est vocem; quis potest dicere, cito se esse conversum? Cave ergo jactantiam, cave sæculum. Ille enim confírmare jubétur fratres suos, qui dixit: Omnia dimísimus, et secúti sumus te.\par
\switchcolumn
\EN
\noindent Let us, then, turn about, and look well to it that there be no strife among us, which of us shall be accounted the greatest. That this strife arose among the Apostles, is not an excuse but a caution for us. If it was only after a while that Peter was converted Peter, who had started up at the first command of the Lord who can promise himself to be converted forthwith? Have a care, then, of boasting; have a care of the world. He who was commanded to strengthen his brethren (Luke xxii. 32) was he who was able to say, Behold, we have forsaken all, and have followed thee. (Matth. xix. 27.)\par
\switchcolumn*

\LA
\TeDeum{Te Deum}
\switchcolumn
\EN
\TeDeum{Te Deum}
\switchcolumn*

\end{paracol}

\par\needspace{10\baselineskip}\addvspace{1.2em}{\centering\small\textsc{Die 24 Julii} · \textsc{24 July}\par}
\Office{In Vigilia S. Jacobi Ap.}{Vigil of St. James the Apostle}{Vigilia}
\addcontentsline{toc}{subsection}{Die 24 Julii · In Vigilia S. Jacobi Ap. \textit{— Vigil of St. James the Apostle}}
\begin{paracol}{2}
\LA
\RefLine{Lectiones I–III: de Scriptura occurrente.}
\switchcolumn
\EN
\RefLine{Lessons I–III: from the Scripture of the day.}
\switchcolumn*

\end{paracol}

\par\needspace{10\baselineskip}\addvspace{1.2em}{\centering\small\textsc{Die 25 Julii} · \textsc{25 July}\par}
\Office{S. Jacobi Apostoli}{St. James the Apostle}{Duplex II classis}
\addcontentsline{toc}{subsection}{Die 25 Julii · S. Jacobi Apostoli \textit{— St. James the Apostle}}
\begin{paracol}{2}
\LA
\RefLine{Lectiones I–III: ex Commúni Apostolorum (C1).}
\switchcolumn
\EN
\RefLine{Lessons I–III: from the Commune Apostolorum (C1).}
\switchcolumn*

\LA
\LessonHead{Lectio IV}
\switchcolumn
\EN
\LessonHead{Lesson IV}
\switchcolumn*

\LA
\noindent Jacóbus, Zebedǽi fílius, Joánnis Apóstoli germánus frater, Galilǽus, inter primos Apóstolos vocátus cum fratre, relíctis patre ac rétibus, secútus est Dóminum, et ambo ab ipso Jesu Boanérges, id est, tonítrui fílii sunt appelláti. Is unus fuit ex tribus Apóstolis, quos Salvátor máxime diléxit, et testes esse vóluit suæ transfiguratiónis, et interésse miráculo cum archisynagógi fíliam a mórtuis excitávit, et adésse cum secéssit in montem Olivéti, Patrem oratúrus, ántequam a Judǽis comprehenderétur.\par
\switchcolumn
\EN
\noindent James, the Son of Zebedee and brother of the Apostle John, was a Galilean, and with his brother one of the first of His Apostles whom the Lord called, whileas they were in a ship with Zebedee their father, mending their nets and they immediately left the ship, and their father, and followed Him. (Matth. iv. 21, 22.) And He surnamed them Boanerges, which is, The sons of thunder. (Mark iii. 17.) Peter, and James, and John, were the three Apostles whom the Saviour loved best; them He took and brought up into an high mountain apart, and was transfigured before them, (Matth. xvii. 1,2;) when He went to the house of the ruler of the synagogue to raise his daughter from the dead, He suffered no man to follow Him save Peter, and James, and John, (Mark v. 37;) and, at the last, just before the Jews took Him, when He cometh unto a place called Gethsemane, and saith unto the disciples: Sit ye here, while I go and pray yonder. He took with Him Peter and the two sons of Zebedee. (Matth. xxvi. 36, 37.)\par
\switchcolumn*

\LA
\begin{Resp}\Rsign Vidi conjúnctos viros, habéntes spléndidas vestes, et Ángelus Dómini locútus est ad me, dicens: \Star{} Isti sunt viri sancti facti amíci Dei.\end{Resp}
\begin{Resp}\Vsign Vidi Ángelum Dei fortem, volántem per médium cælum, voce magna clamántem et dicéntem.\end{Resp}
\begin{Resp}\Rsign Isti sunt viri sancti facti amíci Dei.\end{Resp}
\switchcolumn
\EN
\begin{Resp}\Rsign I saw men standing together, clad in shining raiment, and the Angel of the Lord spoke unto me, saying: \Star{} These men are holy, for they are the friends of God.\end{Resp}
\begin{Resp}\Vsign I saw a strong Angel of God fly into the midst of heaven, saying with a loud voice:\end{Resp}
\begin{Resp}\Rsign These men are holy, for they are the friends of God.\end{Resp}
\switchcolumn*

\LA
\LessonHead{Lectio V}
\switchcolumn
\EN
\LessonHead{Lesson V}
\switchcolumn*

\LA
\noindent Post Jesu Christi ascénsum in cælum, in Judǽa et Samaría ejus divinitátem prǽdicans, plúrimos ad christiánam fidem perdúxit. Mox in Hispániam proféctus, ibi áliquos ad Christum convértit; ex quorum número septem póstea epíscopi a beáto Petro ordináti in Hispániam primi dirécti sunt. Deínde Jerosólymam revérsus, cum inter álios Hermógenem magum fídei veritáte imbuísset, Heródes Agríppa Cláudio imperatóre ad regnum elátus ut a Judǽis grátiam iníret, Jacóbum líbere Jesum Christum Deum confiténtem cápitis condemnávit. Quem cum is, qui eum dúxerat ad tribúnal, fórtiter martýrium subeúntem vidísset, statim se et ipse Christiánum esse proféssus est.\par
\switchcolumn
\EN
\noindent After that Jesus Christ was ascended into heaven, James preached how that He was God, and led many in Judaea and Samaria to the Christian Faith. A while afterward, he went to Spain, and there he brought some to Christ, of whom seven were afterwards ordained Bishops by Blessed Peter, and were the first such sent into that country. From Spain James went back to Jerusalem, where he taught the Faith to diverse persons, and, among others, to the Magician Hermogenes. Thereupon Herod Agrippa, who had been raised to the kingdom under the Emperor Claudius, to curry favour with the Jews, condemned James to death for his firm confession that Jesus Christ is God. The officer who led James to the judgment-seat, at sight of the courage wherewith he was ready to offer up his testimony, declared himself also to be a Christian.\par
\switchcolumn*

\LA
\begin{Resp}\Rsign Beáti estis, cum maledíxerint vobis hómines, et persecúti vos fúerint, et díxerint omne malum advérsum vos, mentiéntes, propter me: \Star{} Gaudéte et exsultáte, quóniam merces vestra copiósa est in cælis.\end{Resp}
\begin{Resp}\Vsign Cum vos óderint hómines, et cum separáverint vos, et exprobráverint, et ejécerint nomen vestrum tamquam malum propter Fílium hóminis.\end{Resp}
\begin{Resp}\Rsign Gaudéte et exsultáte, quóniam merces vestra copiósa est in cælis.\end{Resp}
\switchcolumn
\EN
\begin{Resp}\Rsign Blessed are ye when men shall revile you, and persecute you, and shall say all manner of evil against you falsely, for My sake; \Star{} Rejoice, and be exceeding glad, for great is your reward in heaven.\end{Resp}
\begin{Resp}\Vsign When men shall hate you, and when they shall separate you from their company, and shall reproach you, and cast out your name as evil, for the Son of Man's sake.\end{Resp}
\begin{Resp}\Rsign Rejoice, and be exceeding glad, for great is your reward in heaven.\end{Resp}
\switchcolumn*

\LA
\LessonHead{Lectio VI}
\switchcolumn
\EN
\LessonHead{Lesson VI}
\switchcolumn*

\LA
\noindent Ad supplícium cum raperéntur, pétiit ille a Jacóbo véniam; quem Jacóbus osculátus, Pax, inquit, tibi sit. Itaque utérque est secúri percússus, cum paulo ante Jacóbus paralýticum sanásset. Corpus ejus póstea Compostéllam translátum est, ubi summa celebritáte cólitur, conveniéntibus eo religiónis et voti causa ex toto terrárum orbe peregrínis. Memória ipsíus natális hodiérno die, qui translatiónis dies est, ab Ecclésia celebrátur, cum ipse circa festum Paschæ, primus Apostolórum Jerosólymis profúso sánguine testimónium Jesu Christo déderit.\par
\switchcolumn
\EN
\noindent As they were being hurried to execution, this man asked pardon of James, and the Apostle kissed him, saying, Peace be unto thee. James healed a paralytic, and immediately afterwards both the prisoners were beheaded. The body of the Apostle was afterwards taken to Compostella, (in the province of Galicia, in Spain,) where his grave is very famous. Multitudes of pilgrims from all parts of the earth betake themselves thither to pray, out of sheer piety or in fulfilment of vows. The Birthday of James is kept by the Church upon this day, which is that of the bringing of his body to Compostella. It was about Easter-time (Acts xii. 2-4) that he bore witness to Jesus Christ with his blood, at Jerusalem, being the first of the Apostles to do so.\par
\switchcolumn*

\LA
\begin{Resp}\Rsign Isti sunt triumphatóres et amíci Dei, qui contemnéntes jussa príncipum, meruérunt prǽmia ætérna: \Star{} Modo coronántur, et accípiunt palmam.\end{Resp}
\begin{Resp}\Vsign Isti sunt qui venérunt ex magna tribulatióne, et lavérunt stolas suas in sánguine Agni.\end{Resp}
\begin{Resp}\Rsign Modo coronántur, et accípiunt palmam.\end{Resp}
\begin{Resp}\Vsign Glória Patri, et Fílio, \Star{} et Spirítui Sancto.\end{Resp}
\begin{Resp}\Rsign Modo coronántur, et accípiunt palmam.\end{Resp}
\switchcolumn
\EN
\begin{Resp}\Rsign These are they which have conquered, and are become the friends of God, who recked not of the commandments of princes, and earned the everlasting reward. \Star{} And now have they crowns on their heads, and palms in their hands.\end{Resp}
\begin{Resp}\Vsign These are they which came out of great tribulation, and have washed their robes in the blood of the Lamb.\end{Resp}
\begin{Resp}\Rsign And now have they crowns on their heads, and palms in their hands.\end{Resp}
\begin{Resp}\Vsign Glory be to the Father, and to the Son, \Star{} and to the Holy Ghost.\end{Resp}
\begin{Resp}\Rsign And now have they crowns on their heads, and palms in their hands.\end{Resp}
\switchcolumn*

\LA
\LessonHead{Lectio VII}
\switchcolumn
\EN
\LessonHead{Lesson VII}
\switchcolumn*

\LA
\LessonTitle{Léctio sancti Evangélii secúndum Matthǽum}
\switchcolumn
\EN
\LessonTitle{From the Holy Gospel according to Matthew}
\switchcolumn*

\LA
\Cite{Matt 20:20-23}
\switchcolumn
\EN
\Cite{Matt 20:20-23}
\switchcolumn*

\LA
\noindent In illo témpore: Accéssit ad Jesum mater filiórum Zebedǽi cum fíliis suis, adórans et petens áliquid ab eo. Et réliqua.\par
\switchcolumn
\EN
\noindent At that time, came to Jesus the mother of Zebedee's children, with her sons, worshipping Him, and desiring a certain thing of Him. And so on.\par
\switchcolumn*

\LA
\LessonTitle{Homilía sancti Joánnis Chrysóstomi}
\switchcolumn
\EN
\LessonTitle{Homily by St. John Chrysostom}
\switchcolumn*

\LA
\Source{Homilia 66 in Matthæum}
\switchcolumn
\EN
\Source{66th on Matth.}
\switchcolumn*

\LA
\noindent Non turbétur quisquam, si ádeo imperféctos dícimus Apóstolos fuísse; nondum enim mystérium crucis erat consummátum, nondum grátia Spíritus in corda ipsórum infúsa. Quod si virtútem ipsórum díscere cupis, quales post datam grátiam Spíritus fúerint, consídera, et vidébis omnem ab illis pervérsam affectiónem fuísse superátam. Hac enim de causa eórum modo imperféctio revelátur, ut apérte percípere possis quales súbito per grátiam effécti fuérunt. Quod ígitur nihil spiritále petébant nec de cælésti regno quidquam cogitábant, perspícuum est. Sed tamen inspiciámus étiam quómodo accédant et quid dicant. Vólumus, ínquiunt, ut, quodcúmque petiérimus, fácias nobis. Ad quod Christus, Quid vultis? respóndit: non ignórans certe, sed ut eos respondére cogat et ulcus détegat, et ita medicaméntum appónat.\par
\switchcolumn
\EN
\noindent Let no man be troubled if we say that the Apostles were still imperfect, for the mystery of the Cross was not yet finished, the grace of the Spirit had not yet been shed abroad in their hearts. If thou wilt behold them in their strength, consider them such as they became after the grace of the Spirit was given them, and thou wilt perceive that they had trodden under foot every vain desire. This is the cause wherefore their present imperfection is made known unto us, that is, that thou mayest see how great a change could be forthwith wrought by grace. But nevertheless let us now look how they came unto Christ, and what they said. Master, they said, we would that Thou shouldest do for us whatsoever we shall desire. (Mark x. 35.) And He said unto them: What would ye that I should do for you? (Mark X. 36.) not, surely, that He knew not what their wish was, but that He would make them answer, and so uncover the wound, to lay a plaster upon it.\par
\switchcolumn*

\LA
\begin{Resp}\Rsign Isti sunt qui vivéntes in carne, plantavérunt Ecclésiam sánguine suo: \Star{} Cálicem Dómini bibérunt, et amíci Dei facti sunt.\end{Resp}
\begin{Resp}\Vsign In omnem terram exívit sonus eórum, et in fines orbis terræ verba eórum.\end{Resp}
\begin{Resp}\Rsign Cálicem Dómini bibérunt, et amíci Dei facti sunt.\end{Resp}
\switchcolumn
\EN
\begin{Resp}\Rsign These are they who while yet they lived in the flesh, planted the Church in their own blood; \Star{} They drank of the Lord's cup, and became the friends of God.\end{Resp}
\begin{Resp}\Vsign Their sound is gone out through all the earth, and their words to the ends of the world.\end{Resp}
\begin{Resp}\Rsign They drank of the Lord's cup, and became the friends of God.\end{Resp}
\switchcolumn*

\LA
\LessonHead{Lectio VIII}
\switchcolumn
\EN
\LessonHead{Lesson VIII}
\switchcolumn*

\LA
\noindent Illi vero, cum erubéscerent et verecúndia prohiberéntur, quóniam humáno afféctu eo devénerant, seórsum ab áliis discípulis Christum accipiéntes, interrogavérunt. Progréssi sunt enim, inquit, ne illis manifésti fíerent; et ita demum ea quæ volébant, dixérunt. Volébant autem, ut ego conício, quóniam super duódecim sedes sessúros discípulos audiérunt, primátum hujus conséssus impetráre: et præpóni quidem se céteris sciébant; Petrum vero sibi præférri formidántes, dícere ausi sunt: Dic ut unus a dextris, alter a sinístris sédeat. Et urgent dicéntes: Dic. Quid ígitur ipse? Ut significáret eos nihil pétere spiritále, sed nec scire quidem quid póstulent, non enim pétere audérent, si scirent. Nescítis, ait, quid petátis: nescítis quam magnum hoc sit, quam mirábile ac ipsas superióres excédens virtútes.\par
\switchcolumn
\EN
\noindent Their wish proceeded from earthly motives, and they were shy and ashamed to express it, and therefore they took Christ apart, and so asked Him. The Evangelist saith: For they were gone apart, that they might not be discovered of them and then they told Him what they sought. To me it seemeth most likely that they had heard how that the disciples should sit upon twelve thrones; they were fain to obtain for themselves the chiefest places at this enthronement; they knew that the Lord loved them better than the most of the others; but they feared that Peter would still be preferred before them; and therefore they made bold to say: Grant unto us that we may sit, one at thy right Hand, and the other at thy left Hand, in thy glory. (Mark x. 37.) They were even instant with Him, saying: Say that we may. And what answered He? To show that they were asking no spiritual gift, nor even knew for themselves what they were asking, nor would have asked it if they had known what it was, Jesus said unto them: Ye know not what ye ask, ye know not how great a thing, how wonderful a thing this is, a thing which even is not Mine to give.\par
\switchcolumn*

\LA
\begin{Resp}\Rsign Isti sunt viri sancti, quos elégit Dóminus in caritáte non ficta, et dedit illis glóriam sempitérnam: \Star{} Quorum doctrína fulget Ecclésia, ut sole luna.\end{Resp}
\begin{Resp}\Vsign Sancti per fidem vicérunt regna: operáti sunt justítiam.\end{Resp}
\begin{Resp}\Rsign Quorum doctrína fulget Ecclésia, ut sole luna.\end{Resp}
\begin{Resp}\Vsign Glória Patri, et Fílio, \Star{} et Spirítui Sancto.\end{Resp}
\begin{Resp}\Rsign Quorum doctrína fulget Ecclésia, ut sole luna.\end{Resp}
\switchcolumn
\EN
\begin{Resp}\Rsign These men are saints, whom the Lord hath chosen in love unfeigned, and hath given them glory everlasting. These are they \Star{} By the light of whose teaching the Church is glorified, even as the moon is glorified by the light of the sun.\end{Resp}
\begin{Resp}\Vsign The saints through faith subdued kingdoms, wrought righteousness.\end{Resp}
\begin{Resp}\Rsign By the light of whose teaching the Church is glorified, even as the moon is glorified by the light of the sun.\end{Resp}
\begin{Resp}\Vsign Glory be to the Father, and to the Son, \Star{} and to the Holy Ghost.\end{Resp}
\begin{Resp}\Rsign By the light of whose teaching the Church is glorified, even as the moon is glorified by the light of the sun.\end{Resp}
\switchcolumn*

\LA
\LessonHead{Lectio IX}
\switchcolumn
\EN
\LessonHead{Lesson IX}
\switchcolumn*

\LA
\noindent Et adjécit: Potéstis bíbere cálicem, quem ego bibitúrus sum; et baptísmo, quo ego baptízor, baptizári? Perpéndis quómodo statim ab hac opinióne ipsos remóvit, contrária eis dísserens. Nam vos, inquit, de honóribus et de corónis mecum ágitis; ego vero de luctámine atque sudóre díssero. Non præmiórum hoc tempus est, nec illa glória mea modo apparébit; sed cædis ac periculórum tempus præsens est. Pérspice autem quáliter ipso interrogatiónis modo et hortátur et állicit. Non enim dixit: Potestísne cædem subíre, potestísne vestrum effúndere sánguinem? Sed, Quonam pacto potéstis bíbere cálicem? Deínde allíciens, inquit, Quem ego bibitúrus sum; ut ipsa cum eo communicatióne labórum promptióres redderéntur.\par
\switchcolumn
\EN
\noindent And He said moreover: Are ye able to drink of the cup that I shall drink of, and to be baptized with the baptism that I am baptized with? Behold how He turneth their thoughts at once another way, speaking to them of things altogether different, as though He said, Ye come unto Me treating of honours and crowns, but I speak unto you of the wrestling and the sweat. This is not yet the time of reward, neither is My glory immediately to be revealed; but now death and danger are present with you. But consider how, by the manner of His questioning, He doth both exhort and invite them. He saith not Are ye able to bear death? Are ye able to shed your blood? but: How are ye able to drink the cup whereto He presently inviteth them, saying: the cup that I shall drink of; that He may make them readier for the strife by knowing that it is a strife which they are to share with Him.\par
\switchcolumn*

\LA
\TeDeum{Te Deum}
\switchcolumn
\EN
\TeDeum{Te Deum}
\switchcolumn*

\end{paracol}

\par\needspace{10\baselineskip}\addvspace{1.2em}{\centering\small\textsc{Die 26 Julii} · \textsc{26 July}\par}
\Office{S. Annæ Matris B.M.V.}{St. Anne, Mother of the Blessed Virgin Mary}{Duplex II classis}
\addcontentsline{toc}{subsection}{Die 26 Julii · S. Annæ Matris B.M.V. \textit{— St. Anne, Mother of the Blessed Virgin Mary}}
\begin{paracol}{2}
\LA
\RefLine{Lectiones I–III: ex Commúni non Virginum non Martyrum (C7a).}
\switchcolumn
\EN
\RefLine{Lessons I–III: from the Common of Widows (C7a).}
\switchcolumn*

\LA
\RefLine{Lectiones VII–IX: ex Commúni non Virginum non Martyrum (C7a).}
\switchcolumn
\EN
\RefLine{Lessons VII–IX: from the Common of Widows (C7a).}
\switchcolumn*

\LA
\LessonHead{Lectio IV}
\switchcolumn
\EN
\LessonHead{Lesson IV}
\switchcolumn*

\LA
\LessonTitle{Sermo sancti Joánnis Damascéni}
\switchcolumn
\EN
\LessonTitle{From the Sermons of St. John of Damascus}
\switchcolumn*

\LA
\Source{Oratio 2 de Nativ. B. Mariæ, prope finem}
\switchcolumn
\EN
\Source{2nd on the Birth of the Blessed Virgin}
\switchcolumn*

\LA
\noindent Propónitur nobis Annæ thálamus, conjugális vitæ simul et virginitátis formam réferens, illam matris, hanc fíliæ: quarum áltera recens sterilitáte liberáta est, áltera autem aliquánto post, Christi partum, ad nostræ natúræ conditiónem divíno opifício formátum, supra natúram est editúra. Mérito ígitur Anna, divíno Spíritu plena, læto hilaríque ánimo pérsonat: Congaudéte mecum, quæ promissiónis germen ex stérili ventre péperi, ac benedictiónis fructum ubéribus meis, ut optáveram, nútrio. Sterilitátis mæstítiam, éxui, ac lætum fecunditátis vestem índui. Congáudeat mecum hódie Anna illa, Phenénnæ adversária, et novum hoc atque inopinátum miráculum, quod in me gestum est, suo exémplo concélebret.\par
\switchcolumn
\EN
\noindent The home of Anne is set before us, wherein to see an example both of married and of maiden life, the one in the person of the mother, the other in that of the daughter, whereof the one hath but now ceased to be barren, and the other is in a little while destined, beyond the course of nature, to become the Mother of the Messiah by a singular birth, specially designed by God to build up anew our nature. It is with reason then that Anne, filled with the Holy Ghost, with joyful and jubilant spirit singeth aloud: Rejoice with me, for out of my barren womb I have borne the bud of promise, and, as I have longed, I nourish at my breasts the fruit of benediction. I have laid aside the mournful garments of barrenness, and put on the joyful raiment of fruitfulness. Let Hannah the adversary of Peninnah make merry with me, and join with me for fellow-feeling, in singing of this new and unhoped-for wonder that is wrought in me.\par
\switchcolumn*

\LA
\begin{Resp}\Rsign Propter veritátem, et mansuetúdinem, et justítiam: \Star{} Et dedúcet te mirabíliter déxtera tua.\end{Resp}
\begin{Resp}\Vsign Spécie tua et pulchritúdine tua inténde, próspere procéde, et regna.\end{Resp}
\begin{Resp}\Rsign Et dedúcet te mirabíliter déxtera tua.\end{Resp}
\switchcolumn
\EN
\begin{Resp}\Rsign Because of truth, and meekness, and righteousness; \Star{} And thy right hand shall lead thee wonderfully.\end{Resp}
\begin{Resp}\Vsign In thy comeliness, and thy beauty, go forward, fare prosperously, and reign.\end{Resp}
\begin{Resp}\Rsign And thy right hand shall lead thee wonderfully.\end{Resp}
\switchcolumn*

\LA
\LessonHead{Lectio V}
\switchcolumn
\EN
\LessonHead{Lesson V}
\switchcolumn*

\LA
\noindent Exsúltet Sara, seníli gáudio géstiens, meúmque ab sterilitáte concéptum præfigúrans. Cóncinant simul stériles et infecúndæ visitatiónem meam, admirábili modo cǽlitus factam. Dicant item matres omnes hac fecunditáte prǽditæ: Benedíctus, qui orántibus id quod optábant, largítus est, et fecunditátem stérili dedit, ac felicíssimum illud germen Vírginis concéssit, quæ Mater Dei secúndum carnem fuit, cujus venter cælum est, in qua habitávit is qui nullo loco capi potest. Cónsonam his nos quoque ipsi, quæ vocabátur sterilis nunc autem virgínei thálami mater éxstitit, laudem offerámus. Dicámus ad eam cum Scriptúra: Quam beáta domus David, ex qua prodiísti, et venter, in quo Deus sanctificatiónis arcam, hoc est, eam a qua ipse sine semine concéptus est, fabricávit.\par
\switchcolumn
\EN
\noindent Let Sarah be glad that was joyfully pregnant in her old age, and was a shadow cast before of my conception that hitherto have been barren. Let all the barren and fruitless break forth into singing, when they behold in what wondrous wise I have been visited from heaven. Let all mothers likewise, that like Anne are gifted with fruitfulness, say: Blessed be He That gave their desire unto them that besought Him, That gave fruitfulness unto her that was barren, and That granted unto her that from her should bud forth the joy-bringing Virgin, who, according to the flesh, was Mother of God, and whose womb was a heaven wherein He dwelt Whom no place can contain. Let us also with them offer our praises to her that was called barren, but now is become the mother of a maidchild; let us say unto her in the words of the Scripture: O how blessed is the house of David from whence thou art sprung, and that womb wherein God hath fashioned the ark of His holiness, that is, her, by whom He was Himself conceived without man's seed.\par
\switchcolumn*

\LA
\begin{Resp}\Rsign Dilexísti justítiam, et odísti iniquitátem: \Star{} Proptérea unxit te Deus, Deus tuus, óleo lætítiæ.\end{Resp}
\begin{Resp}\Vsign Propter veritátem, et mansuetúdinem, et justítiam.\end{Resp}
\begin{Resp}\Rsign Proptérea unxit te Deus, Deus tuus, óleo lætítiæ.\end{Resp}
\switchcolumn
\EN
\begin{Resp}\Rsign Thou hast loved righteousness, and hated iniquity; \Star{} Therefore God, thy God, hath anointed thee with the oil of gladness.\end{Resp}
\begin{Resp}\Vsign Because of truth, and meekness, and righteousness.\end{Resp}
\begin{Resp}\Rsign Therefore God, thy God, hath anointed thee with the oil of gladness.\end{Resp}
\switchcolumn*

\LA
\LessonHead{Lectio VI}
\switchcolumn
\EN
\LessonHead{Lesson VI}
\switchcolumn*

\LA
\noindent Vere beáta es, ac ter beáta, quæ beatitúdine donátam a Deo infántem, hoc est, Maríam, nómine quoque ipso magnópere venerándam, peperísti; ex qua Christus vitæ flos éxstitit, cujus Vírginis et gloriósus fuit ortus, et partus mundo sublímior. Nos quoque, o beatíssima fémina, tibi gratulámur; étenim nostrum ómnium spem divínitus concéssam, hoc est, promissiónis fœtum peperísti. Beáta revéra es, et beátus fructus ventris tui. Piórum autem lingua germen tuum magníficat, ac sermo omnis lætus partum tuum prǽdicat. Dignum sane quidem, ac máxime dignum est eam laudáre, quæ divína benignitáte oráculum accépit, ac talem et tantum nobis fructum édidit, ex quo dulcis Jesus pródiit.\par
\switchcolumn
\EN
\noindent Blessed art thou, and thrice blessed, whom God hath so blessed as to make thee to bring forth, as His own gift, the babe Mary, whose very name is highly honourable, out of whom Christ, the Flower of life, blossomed a maiden whose rising is glorious, and whose delivery is worth more than the world. We also, O woman most blessed, do wish thee joy. In sooth thou hast brought forth what we have all hoped for, and God hath given us, namely, the babe of promise. Blessed indeed art thou, and blessed is the fruit of thy womb. The tongues of all the godly do magnify thine offspring, and every glad word is spoken concerning her of whom thou art delivered. Meet in truth is it, and most meet to praise her who received a revelation from the goodness of God, and bore for us such and so great a fruit, from whom sweet Jesus sprang.\par
\switchcolumn*

\LA
\begin{Resp}\Rsign Fallax grátia, et vana est pulchritúdo: \Star{} Múlier timens Dóminum, ipsa laudábitur.\end{Resp}
\begin{Resp}\Vsign Date ei de fructu mánuum suárum, et laudent eam in portis ópera ejus.\end{Resp}
\begin{Resp}\Rsign Múlier timens Dóminum, ipsa laudábitur.\end{Resp}
\begin{Resp}\Vsign Glória Patri, et Fílio, \Star{} et Spirítui Sancto.\end{Resp}
\begin{Resp}\Rsign Múlier timens Dóminum, ipsa laudábitur.\end{Resp}
\switchcolumn
\EN
\begin{Resp}\Rsign Favour is deceitful, and beauty is vain \Star{} A woman that feareth God she shall be praised.\end{Resp}
\begin{Resp}\Vsign Give her of the fruit of her hands, and let her own works praise her in the gates.\end{Resp}
\begin{Resp}\Rsign A woman that feareth God she shall be praised.\end{Resp}
\begin{Resp}\Vsign Glory be to the Father, and to the Son, \Star{} and to the Holy Ghost.\end{Resp}
\begin{Resp}\Rsign A woman that feareth God she shall be praised.\end{Resp}
\switchcolumn*

\end{paracol}

\par\needspace{10\baselineskip}\addvspace{1.2em}{\centering\small\textsc{Die 27 Julii} · \textsc{27 July}\par}
\Office{S. Pantaleonis Martyris}{S. Pantaleon Martyr}{Simplex}
\addcontentsline{toc}{subsection}{Die 27 Julii · S. Pantaleonis Martyris \textit{— S. Pantaleon Martyr}}
\begin{paracol}{2}
\LA
\RefLine{Lectiones I–II: de Scriptura occurrente.}
\switchcolumn
\EN
\RefLine{Lessons I–II: from the Scripture of the day.}
\switchcolumn*

\LA
\LessonHead{Lectio III (pro commemoratione)}
\switchcolumn
\EN
\LessonHead{Lesson III (when commemorated)}
\switchcolumn*

\LA
\noindent Pantáleon Nicomediénsis, nobilis médicus, ab Hermoláo presbytero in Jesu Christi fide eruditus, baptizátus est. Qui mox patri Eustorgio persuasit, ut Christianus fieret. Quare, cum Nicomedíæ póstea Christi Dómini fidem líbere prædicaret, et ad ejus doctrinam omnes cohortarétur, Diocletiáno imperatóre, equuleo tortus, et, admotis ad ejus corpus láminis candéntibus, cruciátus est. Quam tormentórum vim æquo et forti animo ferens, ad extremum gládio percússus, martyrii corónam adeptus est.\par
\switchcolumn
\EN
\noindent This Pantaleon was of a noble family of Nicomedia, and was a physician by trade. Hermolaus the Priest taught him well the faith of Jesus Christ, and he was baptized. A little while after, he persuaded his father Eustorgius to become a Christian. He afterwards duly preached the Faith of the Lord Christ at Nicomedia, and exhorted all to embrace His doctrine. For this he was tortured under the Emperor Diocletian, first on the rack, and then by putting red-hot metal to his body. He bore all the bitterness of his torments with a quiet and brave heart, and at last received the stroke of the sword, and grasped the crown of martyrdom.\par
\switchcolumn*

\LA
\TeDeum{Te Deum}
\switchcolumn
\EN
\TeDeum{Te Deum}
\switchcolumn*

\end{paracol}

\par\needspace{10\baselineskip}\addvspace{1.2em}{\centering\small\textsc{Die 28 Julii} · \textsc{28 July}\par}
\Office{Ss. Nazarii et Celsi Martyrum, Victoris I Papæ et Martyris ac Innocentii I Papæ et Confessoris}{Sts. Nazarius and Celsus, Martyr, Pope Victor I, Martyr, and Pope Innocent I, Confessor}{Semiduplex}
\addcontentsline{toc}{subsection}{Die 28 Julii · Ss. Nazarii et Celsi Martyrum, Victoris I Papæ et Martyris ac Innocentii I Papæ et Confessoris \textit{— Sts. Nazarius and Celsus, Martyr, Pope Victor I, Martyr, and Pope Innocent I, Confessor}}
\begin{paracol}{2}
\LA
\RefLine{Lectiones I–III: de Scriptura occurrente.}
\switchcolumn
\EN
\RefLine{Lessons I–III: from the Scripture of the day.}
\switchcolumn*

\LA
\RefLine{Lectiones VII–IX: ex Commúni Plurimorum Martyrum Pontificum (C3).}
\switchcolumn
\EN
\RefLine{Lessons VII–IX: from the Commune Plurimorum Martyrum Pontificum (C3).}
\switchcolumn*

\LA
\LessonHead{Lectio IV}
\switchcolumn
\EN
\LessonHead{Lesson IV}
\switchcolumn*

\LA
\noindent Nazarius, a beato Lino papa baptizatus, cum in Gálliam profectus esset, ibi Celsum púerum, a se christiánis præceptis prius instructum, baptizávit. Qui una Trévirim euntes, Nerónis persecutióne, in mare uterque dejícitur, unde mirabíliter evasérunt. Postea Mediolanum veniéntes, cum ibi Christi fidem disseminárent, ab Anolino præfecto, constantíssime Christum Deum confiténtes, cápite plectúntur: quorum córpora extra portam Romanam sepúlta sunt. Quæ cum diu latuíssent, Dei mónitu a beato Ambrosio conspersa recénti sánguine sunt invénta, tamquam si paulo ante martyrium passi essent; unde, in urbem translata, honorifico sepúlcro contecta sunt.\par
\switchcolumn
\EN
\noindent Nazarius was baptized by the blessed Pope Linus, and afterwards went to Gaul. There he met with the boy Celsus, whom he instructed in the Christian law, and baptized. They went together to Trier, and, in the persecution under Nero, were both thrown into the sea, from which they had a marvellous escape. Later on, they came to Milan, where they spread the Faith of Christ, and as they remained firm in declaring that He is God, the Prsefect Anolinus had them beheaded. Their bodies were buried outside the Roman gate, and lay long unknown, till, by a revelation from God, Blessed Ambrose found them, smeared with fresh blood, as though they had only a short time undergone martyrdom. They were taken up from thence, carried into the city, and laid in an honourable sepulchre.\par
\switchcolumn*

\LA
\begin{Resp}\Rsign Sancti tui, Dómine, mirábile consecúti sunt iter, serviéntes præcéptis tuis, ut inveniréntur illǽsi in aquis válidis: \Star{} Terra appáruit árida: et in Mari Rubro via sine impediménto.\end{Resp}
\begin{Resp}\Vsign Quóniam percússit petram, et fluxérunt aquæ, et torréntes inundavérunt.\end{Resp}
\begin{Resp}\Rsign Terra appáruit árida: et in Mari Rubro via sine impediménto.\end{Resp}
\switchcolumn
\EN
\begin{Resp}\Rsign thy Saints, O Lord, have passed a wonderful way, serving thy commandments, that they might be found without hurt in the midst of the mighty waters. \Star{} Dry land appeared, and, out of the Red Sea, a way without impediment.\end{Resp}
\begin{Resp}\Vsign He smote the rock, and the waters gushed out, and the streams overflowed.\end{Resp}
\begin{Resp}\Rsign Dry land appeared, and, out of the Red Sea, a way without impediment.\end{Resp}
\switchcolumn*

\LA
\LessonHead{Lectio V}
\switchcolumn
\EN
\LessonHead{Lesson V}
\switchcolumn*

\LA
\noindent Victor in Africa natus, Sevéro imperatóre rexit Ecclésiam. Confirmávit decretum Pii primi, ut sacrum Pascha die Dominico celebrarétur: qui ritus ut póstea in mores inducerétur, habita sunt multis in locis concília; et in Nicæna dénique prima synodo sancítum est, ut Paschæ dies festus post quartam décimam lunam agerétur, ne Christiáni Judæos imitari videréntur. Státuit ut quavis aqua, modo naturali, si necessitas cógeret, quicúmque baptizari posset. Theódotum Coriárium Byzantinum, docéntem Christum tantummodo hóminem fuísse, ejécit ex Ecclésia. Scripsit de quæstióne Paschæ et alia quædam opúscula. Creávit duabus ordinatiónibus, mense Decembri, presbyteros quatuor, diaconos septem, episcopos per diversa loca duodecim. Martyrio coronátus, sepelítur in Vaticano, quinto Kalendas Augusti. Sedit annos novem, mensem unum, dies viginti octo.\par
\switchcolumn
\EN
\noindent Victor was by birth an African, and governed the Church in the time of the Emperor Severus. He confirmed the decree of Pius I that the Holy Passover should be kept upon the Lord's Day. To bring this rule into use, Councils were held in many places, and in the First Synod of Nice it was decided that the Holy Passover Day should be kept after the 14th day of the month (Nisan), lest the Christians should seem to be copying the Jews. Victor decided that if need be, baptism can be administered with any water, as long as it be natural. He cast out of the Church Theodotus the tanner, of Constantinople, who taught that Christ was nothing but a man. He wrote upon the subject of the Passover, and some other small works. He held two Ordinations in the month of December, wherein he ordained four Priests, seven Deacons, and twelve Bishops for diverse places. He received the crown of his testimony, and was buried at the Vatican on the 28th day of July, (in the year of our Lord 197.) He sat in the throne of Peter nine years, one month, and twenty-eight days.\par
\switchcolumn*

\LA
\begin{Resp}\Rsign Vérbera carníficum non timuérunt Sancti Dei, moriéntes pro Christi nómine: \Star{} Ut herédes fíerent in domo Dómini.\end{Resp}
\begin{Resp}\Vsign Tradidérunt córpora sua propter Deum ad supplícia.\end{Resp}
\begin{Resp}\Rsign Ut herédes fíerent in domo Dómini.\end{Resp}
\switchcolumn
\EN
\begin{Resp}\Rsign The Saints of God shrank not from the stripes of the executioners, but died for Christ's Name's sake \Star{} That they might be made joint-heirs in the house of the Lord.\end{Resp}
\begin{Resp}\Vsign They gave their bodies for God's sake to death.\end{Resp}
\begin{Resp}\Rsign That they might be made joint-heirs in the house of the Lord.\end{Resp}
\switchcolumn*

\LA
\LessonHead{Lectio VI}
\switchcolumn
\EN
\LessonHead{Lesson VI}
\switchcolumn*

\LA
\noindent Innocentius Albanénsis, sancti Hierónymi et Augustíni ætate flóruit. De quo ille ad Demetríadem vírginem: Sancti Innocentii, qui apostolicæ cathedræ et beátæ memóriæ Anastasii successor et fílius est, teneas fidem nec peregrinam, quamvis tibi prudens callídaque videaris, doctrinam recipias. Eum, tamquam justum Lot subtractum Dei providéntia, ad Ravennam servátum fuísse, scribit Orósius, ne Romani pópuli vidéret excidium. Is, Pelagio et Cælestio damnatis, contra eórum hæresim decretum fecit, ut párvuli, ex christiana étiam mulíere nati, per baptismum renasci debérent; ut in eis regeneratióne mundétur, quod generatióne contraxérunt. Probávit etiam, ut Sábbato, ob memóriam Christi Dómini sepultúræ jejunium servarétur. Sedit annos quíndecim, mensem unum, dies decem. Quatuor ordinatiónibus, mense Decembri, creávit presbyteros trigínta, diaconos quíndecim, episcopos per diversa loca quinquagínta quatuor. Sepultus est in cœmeterio ad Ursum pileátum.\par
\switchcolumn
\EN
\noindent Innocent of Albano flourished in the time of St. Jerome and St. Augustine. Concerning him, St. Jerome saith, writing to Demetrias: Keep firm hold on the faith of holy Innocent, who is the heir and the child of the Apostolic See, and of Anastasius of blessed memory, and receive not any strange doctrine, how wise and acute soever thou mayest count thyself. Orosius writeth that God kept Innocent at Ravenna that he might not see the destruction of the Roman people, even as righteous Lot was withdrawn by God's Providence (from being at the burning of Sodom.) He condemned Pelagius and Caelestius, and made a decree against their heresy, that little children, even those whose mother was a Christian, must be born again in baptism, that the new birth may wash away in them the stain which they have contracted in their conception. He approved also that a Fast should be kept upon Saturday in memory of the Lord Christ's lying in the grave on that day. He sat in the throne of Peter fifteen years, one month, and ten days. He held four Ordinations in the month of December, and made therein thirty Priests, fifteen Deacons, and forty-four Bishops for diverse places. He was buried in the Bearand-Cap Cemetery, (upon the 28th of July, in the year 417.)\par
\switchcolumn*

\LA
\begin{Resp}\Rsign Tamquam aurum in fornáce probávit eléctos Dóminus, et quasi holocáusti hóstiam accépit illos: et in témpore erit respéctus illórum: \Star{} Quóniam donum et pax est eléctis Dei.\end{Resp}
\begin{Resp}\Vsign Qui confídunt in illum, intélligent veritátem, et fidéles in dilectióne acquiéscent illi.\end{Resp}
\begin{Resp}\Rsign Quóniam donum et pax est eléctis Dei.\end{Resp}
\begin{Resp}\Vsign Glória Patri, et Fílio, \Star{} et Spirítui Sancto.\end{Resp}
\begin{Resp}\Rsign Quóniam donum et pax est eléctis Dei.\end{Resp}
\switchcolumn
\EN
\begin{Resp}\Rsign As gold in the furnace hath the Lord tried His chosen ones, and received them as a burnt-offering, and yet a while, and they shall be regarded; \Star{} For the grace of God, and His peace, are with His chosen.\end{Resp}
\begin{Resp}\Vsign They that put their trust in Him shall understand the truth and such as be faithful in love shall abide with Him.\end{Resp}
\begin{Resp}\Rsign For the grace of God, and His peace, are with His chosen.\end{Resp}
\begin{Resp}\Vsign Glory be to the Father, and to the Son, \Star{} and to the Holy Ghost.\end{Resp}
\begin{Resp}\Rsign For the grace of God, and His peace, are with His chosen.\end{Resp}
\switchcolumn*

\end{paracol}

\par\needspace{10\baselineskip}\addvspace{1.2em}{\centering\small\textsc{Die 29 Julii} · \textsc{29 July}\par}
\Office{S. Marthæ Virginis}{St. Martha, Virgin}{Semiduplex}
\addcontentsline{toc}{subsection}{Die 29 Julii · S. Marthæ Virginis \textit{— St. Martha, Virgin}}
\begin{paracol}{2}
\LA
\RefLine{Lectiones I–III: de Scriptura occurrente.}
\switchcolumn
\EN
\RefLine{Lessons I–III: from the Scripture of the day.}
\switchcolumn*

\LA
\LessonHead{Lectio IV}
\switchcolumn
\EN
\LessonHead{Lesson IV}
\switchcolumn*

\LA
\noindent Martha, nobílibus et copiosis paréntibus nata, sed Christi Dómini hospítio clarior, post ejus ascénsum in cælum, cum fratre, sorore, et Marcella pedíssequa, ac Maximino, uno ex septuagínta duobus discípulis Christi Dómini, qui totam illam domum baptizaverat, multisque aliis Christiánis, comprehensa a Judæis, in navem sine velo ac remigio impónitur, vastissimóque mari ad certum naufragium committitur. Sed navis, Deo gubernante, salvis ómnibus, Massiliam appulsa est.\par
\switchcolumn
\EN
\noindent Martha was the daughter of noble and wealthy parents, but is best known as having been the hostess of the Lord Christ. After that He was ascended into heaven, Martha, along with her brother (Lazarus,) her sister (Mary Magdalene,) her waiting-woman Marcella, Maximin, who was one of the seventy-two disciples of the Lord Christ, and who had baptized the whole of the family, and many other Christians, was taken by the Jews, and turned adrift upon the open sea in a ship without sail or oars, to meet with certain wreck, but by the governance of God the ship came to land at Marseilles with all safe.\par
\switchcolumn*

\LA
\begin{Resp}\Rsign Propter veritátem, et mansuetúdinem, et justítiam: \Star{} Et dedúcet te mirabíliter déxtera tua.\end{Resp}
\begin{Resp}\Vsign Spécie tua et pulchritúdine tua inténde, próspere procéde, et regna.\end{Resp}
\begin{Resp}\Rsign Et dedúcet te mirabíliter déxtera tua.\end{Resp}
\switchcolumn
\EN
\begin{Resp}\Rsign Because of truth, and meekness, and righteousness; \Star{} And thy right hand shall lead thee wonderfully.\end{Resp}
\begin{Resp}\Vsign In thy comeliness, and thy beauty, go forward, fare prosperously, and reign.\end{Resp}
\begin{Resp}\Rsign And thy right hand shall lead thee wonderfully.\end{Resp}
\switchcolumn*

\LA
\LessonHead{Lectio V}
\switchcolumn
\EN
\LessonHead{Lesson V}
\switchcolumn*

\LA
\noindent Eo miraculo et horum prædicatióne, primum Massilienses, mox Aquenses ac finítimæ gentes in Christum credidérunt; Lazarusque Massiliensium et Maximinus Aquensium epíscopus creatur. Magdalena vero, assueta oratióni et pédibus Dómini, ut optima parte contemplandæ cæléstis beatitúdinis, quam elégerat, fruerétur, in vastam altíssimi montis spelúncam se cóntulit; ubi trigínta annos vixit, ab omni hóminum consuetúdine disjuncta, quotidiéque per id tempus ad audiendas cæléstium laudes in altum ab Angelis elata.\par
\switchcolumn
\EN
\noindent Through this miracle and the preaching of the Saints, the people of Marseilles first, and then those of Aix, and of the uttermost tribes, believed in Christ, and Lazarus was made Bishop of Marseilles, and Maximin Bishop of Aix. Mary Magdalene sat still at Jesus' Feet, being altogether given to prayer and the contemplation of heavenly blessedness, that that good part which she had chosen might not be taken away from her, withdrew herself to a great cave in an exceeding high mountain, where she lived for thirty years, utterly cut off from all conversation with men, and every day during that time carried up by Angels into the air, to listen to them that dwell in heaven praising God.\par
\switchcolumn*

\LA
\begin{Resp}\Rsign Dilexísti justítiam, et odísti iniquitátem: \Star{} Proptérea unxit te Deus, Deus tuus, óleo lætítiæ.\end{Resp}
\begin{Resp}\Vsign Propter veritátem, et mansuetúdinem, et justítiam.\end{Resp}
\begin{Resp}\Rsign Proptérea unxit te Deus, Deus tuus, óleo lætítiæ.\end{Resp}
\switchcolumn
\EN
\begin{Resp}\Rsign Thou hast loved righteousness, and hated iniquity; \Star{} Therefore God, thy God, hath anointed thee with the oil of gladness.\end{Resp}
\begin{Resp}\Vsign Because of truth, and meekness, and righteousness.\end{Resp}
\begin{Resp}\Rsign Therefore God, thy God, hath anointed thee with the oil of gladness.\end{Resp}
\switchcolumn*

\LA
\LessonHead{Lectio VI}
\switchcolumn
\EN
\LessonHead{Lesson VI}
\switchcolumn*

\LA
\noindent Martha autem, mirábili vitæ sanctitáte et caritate, ómnium Massiliensium ánimis in sui amórem et admiratiónem adductis, in locum a viris remotum cum aliquot honestíssimis féminis se recepit; ubi summa cum laude pietátis et prudéntiæ diu vixit, ac demum, morte sua multo ante prædicta, miraculis clara migrávit ad Dóminum, quarto Kalendas Augusti. Cujus corpus apud Tarascum magnam habet veneratiónem.\par
\switchcolumn
\EN
\noindent Martha, by the wondrous holiness and charity of her life, drew upon herself the love and wonder of all the inhabitants of Marseilles. She withdrew herself in company with some other honourable women into a place out of the way of men, where she lived long, with great praise for godliness and discretion. She foretold her own death long before, and at last, illustrious for miracles, passed away to be ever with the Lord, upon the 29th day of July. Her body is held in great worship at Tarascon.\par
\switchcolumn*

\LA
\begin{Resp}\Rsign Afferéntur Regi vírgines post eam, próximæ ejus; \Star{} Afferéntur tibi in lætítia et exsultatióne.\end{Resp}
\begin{Resp}\Vsign Spécie tua et pulchritúdine tua; inténde, próspere procéde, et regna.\end{Resp}
\begin{Resp}\Rsign Afferéntur tibi in lætítia et exsultatióne.\end{Resp}
\begin{Resp}\Vsign Glória Patri, et Fílio, \Star{} et Spirítui Sancto.\end{Resp}
\begin{Resp}\Rsign Afferéntur tibi in lætítia et exsultatióne.\end{Resp}
\switchcolumn
\EN
\begin{Resp}\Rsign After her shall virgins be brought unto the King, her fellows \Star{} They shall be brought unto thee with gladness and rejoicing.\end{Resp}
\begin{Resp}\Vsign In thy comeliness and thy beauty, go forward, fare prosperously, and reign.\end{Resp}
\begin{Resp}\Rsign They shall be brought unto thee with gladness and rejoicing.\end{Resp}
\begin{Resp}\Vsign Glory be to the Father, and to the Son, \Star{} and to the Holy Ghost.\end{Resp}
\begin{Resp}\Rsign They shall be brought unto thee with gladness and rejoicing.\end{Resp}
\switchcolumn*

\LA
\LessonHead{Lectio VII}
\switchcolumn
\EN
\LessonHead{Lesson VII}
\switchcolumn*

\LA
\LessonTitle{Léctio sancti Evangélii secúndum Lucam}
\switchcolumn
\EN
\LessonTitle{From the Holy Gospel according to Luke}
\switchcolumn*

\LA
\Cite{Luc 10:38-42}
\switchcolumn
\EN
\Cite{Luke 10:38-42}
\switchcolumn*

\LA
\noindent In illo témpore: Intrávit Jesus in quoddam castéllum; et múlier quædam, Martha nómine, excépit illum in domum suam. Et réliqua.\par
\switchcolumn
\EN
\noindent At that time: Jesus entered into a certain village, and a certain woman, named Martha, received Him into her house. And so on.\par
\switchcolumn*

\LA
\LessonTitle{Homilía sancti Augustíni Epíscopi}
\switchcolumn
\EN
\LessonTitle{Homily by St. Augustine, Bishop (of Hippo)}
\switchcolumn*

\LA
\Source{Sermo 20 de Verbis Domini}
\switchcolumn
\EN
\Source{26th upon the Words of the Lord}
\switchcolumn*

\LA
\noindent Verba Dómini nostri Jesu Christi, quæ modo ex Evangélio recitáta sunt, ádmonent nos, esse unum aliquid quo tendámus, quando in hujus sæculi multitúdine laboramus. Téndimus autem adhuc peregrinántes, nondum manéntes; adhuc in via, nondum in pátria; adhuc desiderándo, nondum fruéndo. Tamen tendámus, et sine pigrítia et sine intermissióne tendámus, ut aliquándo perveníre valeamus. Martha et Maria duæ sorores erant, ambæ non solum carne, sed étiam religióne germanæ; ambæ Dómino adhæsérunt, ambæ Dómino in carne præsénti concorditer serviérunt.\par
\switchcolumn
\EN
\noindent The words of our Lord Jesus Christ which have just been read from the Gospel, give us to wit that there is one thing toward the which we are making our way, all the while that we are striving amid the diverse cares of this world. Thitherward we make our way, while we are still strangers and pilgrims, unpossessed as yet of any abiding city, still on the journey, not yet come home, still hoping, not yet enjoying. Still thitherward let us make our way, not slothfully nor by fits and starts, but so that some day we may arrive thither. Martha and Mary were sisters, not in the flesh only, but also in godliness; together, they clave unto the Lord; together, with one heart they served the Lord present in the Flesh.\par
\switchcolumn*

\LA
\begin{Resp}\Rsign Hæc est Virgo sápiens, quam Dóminus vigilántem invénit, quæ accéptis lampádibus sumpsit secum óleum: \Star{} Et veniénte Dómino, introívit cum eo ad núptias.\end{Resp}
\begin{Resp}\Vsign Média nocte clamor factus est: Ecce sponsus venit, exíte óbviam ei.\end{Resp}
\begin{Resp}\Rsign Et veniénte Dómino, introívit cum eo ad núptias.\end{Resp}
\switchcolumn
\EN
\begin{Resp}\Rsign This is one of those wise virgins, whom the Lord found watching, for when she took her lamp, she took oil with her. \Star{} And when the Lord came, she went in with him to the marriage.\end{Resp}
\begin{Resp}\Vsign At midnight there was a cry made: Behold, the Bridegroom cometh, go ye out to meet him.\end{Resp}
\begin{Resp}\Rsign And when the Lord came, she went in with him to the marriage.\end{Resp}
\switchcolumn*

\LA
\LessonHead{Lectio VIII}
\switchcolumn
\EN
\LessonHead{Lesson VIII}
\switchcolumn*

\LA
\noindent Suscépit eum Martha, sicut solent súscipi peregrini; sed tamen suscépit famula Dóminum, ægra Salvatórem, creatura Creatórem. Suscépit autem spíritu pascenda, in carne pascéndum. Voluit enim Dóminus formam servi accípere, et, accepta forma servi, in illa pasci a servis, dignatióne, non conditione. Nam et ista dignátio fuit, se præbére pascéndum. Habebat carnem, in qua esuriret quidem, et sitiret; sed nescítis, quia in erémo esuriénti Angeli ministrábant? Ergo quod pasci vóluit, pascénti præstitit. Quid autem mirum, si et de sancto Elía præstitit víduæ, quem prius, corvo ministrante, pascebat? Numquid pascéndo defécerat, quando ad víduam mittebat? Nequaquam, sed religiosam víduam, per obsequium exhibitum servo suo, benedicere disponebat.\par
\switchcolumn
\EN
\noindent Martha received Him into her house. It was just as strangers are received, but it was the handmaiden receiving her Lord, the sick receiving her Saviour, the creature receiving her Creator. She received Him, to give bodily meat unto Him by Whom she herself was to be fed unto eternal life. It had been the Lord's will to take upon Him the form of a servant, to be fed by servants, (still out of His good pleasure, not of necessity,) and in that form of a servant which He had taken upon Him. This was His good pleasure, to offer Himself as a subject for hospitality. He had Flesh, wherein He was sometimes hungered and thirsty, but know ye not how that, when He was in the desert and was an-hungered, angels came and ministered unto Him. Himself it was therefore, That gave unto them of whom He was fain to be fed, the wherewithal. And what wonder is this if we consider how that holy Elijah, coming from being fed by the ministry of ravens, asked bread of the widow of Zarephath, and himself gave her the wherewithal to feed him? Had God failed to feed Elijah when He sent him unto the widow? God forbid. He did so that He might bless that godly widow for a service rendered unto His servant.\par
\switchcolumn*

\LA
\begin{Resp}\Rsign Média nocte clamor factus est: \Star{} Ecce sponsus venit, exíte óbviam ei.\end{Resp}
\begin{Resp}\Vsign Prudéntes vírgines, aptáte vestras lámpades.\end{Resp}
\begin{Resp}\Rsign Ecce sponsus venit, exíte óbviam ei.\end{Resp}
\begin{Resp}\Vsign Glória Patri, et Fílio, \Star{} et Spirítui Sancto.\end{Resp}
\begin{Resp}\Rsign Ecce sponsus venit, exíte óbviam ei.\end{Resp}
\switchcolumn
\EN
\begin{Resp}\Rsign At midnight there was a cry made: \Star{} Behold, the Bridegroom cometh, go ye out to meet him.\end{Resp}
\begin{Resp}\Vsign Trim your lamps, O ye wise virgins.\end{Resp}
\begin{Resp}\Rsign Behold, the Bridegroom cometh, go ye out to meet him.\end{Resp}
\begin{Resp}\Vsign Glory be to the Father, and to the Son, \Star{} and to the Holy Ghost.\end{Resp}
\begin{Resp}\Rsign Behold, the Bridegroom cometh, go ye out to meet him.\end{Resp}
\switchcolumn*

\LA
\LessonHead{Lectio IX}
\switchcolumn
\EN
\LessonHead{Lesson IX}
\switchcolumn*

\LA
\noindent Sic ergo susceptus est Dóminus, tamquam hospes, qui in sua propria venit, et sui eum non recepérunt; sed quotquot recepérunt eum, dedit eis potestátem fílios Dei fíeri, adoptans servos et liberos fáciens, rédimens captivos et fáciens coheredes. Ne quis tamen vestrum fórsitan dicat: O beáti, qui Christum suscipere in domum propriam meruérunt! Noli dolére, noli murmurare quia tempóribus natus es, quando jam Dóminus non vides in carne. Non tibi ábstulit istam dignatiónem. Cum uni, inquit, ex minimis meis fecistis, mihi fecistis. Hæc de Dómino pascéndo in carne, sed pascénte in spiritu, pauca pro témpore dixérimus.\par
\switchcolumn
\EN
\noindent Thus was that same Lord received as a guest, Who came unto His own, and His own received Him not, but as many as received Him, to them gave He power to become the sons of God, adopting servants and making them children, redeeming prisoners and appointing them coheirs. Perchance some of you will say: O how blessed were they who were worthy to receive Christ as a guest into their own home! but mourn not, neither murmur, for that thou hast been born in an age wherein thou canst no more see Christ in the flesh. He hath not put the honour of receiving Him beyond thy reach. Inasmuch, saith He, as ye have done it unto one of the least of these My brethren, ye have done it unto Me. (Matth. xxv. 40.) The above remarks have occurred to me regarding the Lord considered as fed in the flesh, and I shall now touch briefly, as time permits, upon the Same, considered as the Feeder of the soul.\par
\switchcolumn*

\LA
\TeDeum{Te Deum}
\switchcolumn
\EN
\TeDeum{Te Deum}
\switchcolumn*

\end{paracol}

\par\needspace{10\baselineskip}\addvspace{1.2em}{\centering\small\textsc{Die 30 Julii} · \textsc{30 July}\par}
\Office{S. Abdon et Sennen Martyrum}{St. Abdon and Sennen, Martyr}{Simplex}
\addcontentsline{toc}{subsection}{Die 30 Julii · S. Abdon et Sennen Martyrum \textit{— St. Abdon and Sennen, Martyr}}
\begin{paracol}{2}
\LA
\RefLine{Lectiones I–II: de Scriptura occurrente.}
\switchcolumn
\EN
\RefLine{Lessons I–II: from the Scripture of the day.}
\switchcolumn*

\LA
\LessonHead{Lectio III (pro commemoratione)}
\switchcolumn
\EN
\LessonHead{Lesson III (when commemorated)}
\switchcolumn*

\LA
\noindent Abdon et Sennen Persæ, Decio imperatóre, accusáti quod córpora Christianórum, quæ inhumáta projiciebántur, in suo prædio sepelissent, jussu imperatóris comprehendúntur, et diis jubéntur sacrificare. Quod cum fácere neglígerent, et Jesum Christum Deum constantíssime prædicarent, tráditos in arctam custódiam, Romam póstea rédiens Decius vinctos duxit in triumpho. Qui cum in Urbe ad simulacra attracti essent, ea detestáti conspuérunt. Quam ob rem ursis ac leónibus objecti sunt; quos feræ non audébant attingere. Demum, gládiis trucidati, colligátis pédibus tracti sunt ante solis simulacrum. Quorum corpora, clam inde asportata, Quirínus diaconus sepelívit in suis ædibus.\par
\switchcolumn
\EN
\noindent Abdon and Sennen were Persians. In the reign of the Emperor Decius they were accused of interring, on their own farm, the bodies of Christians, which had been thrown out unburied. The Emperor commanded them to be arrested and ordered to sacrifice to the gods. This they refused to do, and persistently preached that Jesus Christ is God, whereupon they were put into strict confinement. When Decius afterwards returned to Rome, he had them led in chains in his triumph. Being thus dragged into the city and up to the idols, they abhorred and spat upon them, for which they were cast to bears and lions; the beasts were afraid to touch them. They were butchered with the sword, and the corpses, with their feet bound together, were dragged before the image of the sun. Thence they were stolen away, and the Deacon Quirinus buried them in his own house.\par
\switchcolumn*

\LA
\TeDeum{Te Deum}
\switchcolumn
\EN
\TeDeum{Te Deum}
\switchcolumn*

\end{paracol}

\par\needspace{10\baselineskip}\addvspace{1.2em}{\centering\small\textsc{Die 31 Julii} · \textsc{31 July}\par}
\Office{S. Ignatii Confessoris}{St. Ignatius of Loyola, Confessor}{Duplex majus}
\addcontentsline{toc}{subsection}{Die 31 Julii · S. Ignatii Confessoris \textit{— St. Ignatius of Loyola, Confessor}}
\begin{paracol}{2}
\LA
\RefLine{Lectiones I–III: de Scriptura occurrente.}
\switchcolumn
\EN
\RefLine{Lessons I–III: from the Scripture of the day.}
\switchcolumn*

\LA
\LessonHead{Lectio IV}
\switchcolumn
\EN
\LessonHead{Lesson IV}
\switchcolumn*

\LA
\noindent Ignatius, natióne Hispanus, nobili genere Loyolæ in Cantabria natus, primo catholici regis aulam, deínde milítiam secutus est. In propugnatióne Pampelonénsi accepto vulnere gráviter decumbens, ex fortuita piórum librórum lectióne, ad Christi Sanctorúmque sectanda vestígia mirabíliter exársit. Ad montem Serrátum profectus, ante aram beátæ Vírginis suspénsis armis, noctem éxcubans, sacræ milítiæ tirocinium pósuit. Inde, ut erat indútus sacco, traditis antea mendíco pretiósis vestibus, Manresam secessit; ubi, emendicato pane et aqua víctitans, exceptisque diébus Dominicis, jejunans, aspera catena cilicióque carnem domans, humi cubans, et férreis se flagellis cruentans, per annum commorátus est, claris adeo illustratiónibus a Deo recreatus, ut póstea dicere sólitus sit: Si sacræ Litteræ non exstarent, se tamen pro fide mori parátum ex iis solum, quæ sibi Manresæ patefécerat Dóminus. Quo témpore, homo litterárum plane rudis, admirábilem illum compósuit exercitiórum librum, Sedis apostolicæ judício et ómnium utilitate comprobátum.\par
\switchcolumn
\EN
\noindent Ignatius was a Spaniard by nation, and was born of the noble Biscayan family of Loyola, (in the year of our Lord 1491.) He followed first the Court and then the army of the Most Catholic King. At the siege of Pampeluna (in the year 1521) he received a severe wound which laid him up with a long and dangerous illness. During this time he chanced to read some godly books, and conceived from them a burning desire to follow in the footsteps of Christ and His saints. He betook himself to Monserrat, and there entered himself for the heavenly warfare, by hanging up his weapons, and watching them for a night before the Altar of the Blessed Virgin. Thence he withdrew to Manresa, clad in sackcloth, for he had before given his costly raiment to a beggar. At Manresa he lived upon bread and water, begging the bread, and fasting every day except the Lord's Day. He mastered his flesh by the use of a sharp chain and hair-cloth, slept upon the ground, and lashed himself to bloodshedding with iron scourges. Thus he dwelt for a year, feasted by God with such clear lights, that he was used afterwards to say that even if the Holy Bible had not existed, he would have been ready to die for the faith only on the evidence of those things which the Lord had shown unto him at Manresa. It was at this time that, albeit a man of little education, he put together that wonderful book entitled Spiritual Exercises, whose worth hath been attested by the judgment of the Apostolic See, and by universal usefulness.\par
\switchcolumn*

\LA
\begin{Resp}\Rsign Honéstum fecit illum Dóminus, et custodívit eum ab inimícis, et a seductóribus tutávit illum: \Star{} Et dedit illi claritátem ætérnam.\end{Resp}
\begin{Resp}\Vsign Justum dedúxit Dóminus per vias rectas, et osténdit illi regnum Dei.\end{Resp}
\begin{Resp}\Rsign Et dedit illi claritátem ætérnam.\end{Resp}
\switchcolumn
\EN
\begin{Resp}\Rsign The Lord made him honourable, and defended him from his enemies, and kept him safe from those that lay in wait for him; \Star{} And gave him perpetual glory.\end{Resp}
\begin{Resp}\Vsign He went down with him into the pit, and left him not in bonds.\end{Resp}
\begin{Resp}\Rsign And gave him perpetual glory.\end{Resp}
\switchcolumn*

\LA
\LessonHead{Lectio V}
\switchcolumn
\EN
\LessonHead{Lesson V}
\switchcolumn*

\LA
\noindent Ut vero se ad animárum lucra rite formaret, subsidium litterárum, a grammatica inter púeros exorsus, adhibére státuit. Cumque nihil interim omitteret de studio alienæ salútis, mirum est quas ubíque locórum ærúmnas ac ludibria devoráverit, asperrima quæque, et víncula et verbera pene ad mortem usque perpéssus: quibus tamen longe plura pro Dómini sui glória semper expetebat. Lutetiæ Parisiórum adjunctis sibi ex illa academía variárum natiónum sociis novem, qui omnes artium magistériis et theologíæ gradibus insignes erant, ibidem in monte Mártyrum prima ordinis fundaménta jecit, quem póstea Romæ instítuens, ad tria consueta quarto addito de Missiónibus voto, Sedi apostolicæ arctius adstrinxit; et Paulus tertius primo recepit confirmavítque, mox alii Pontifices ac Tridentina synodus probavére. Ipse autem, misso ad prædicándum Indis Evangélium sancto Francisco Xaverio, aliisque in alias mundi plagas ad religiónem propagandam disseminatis, éthnicæ superstitióni hæresique bellum indixit; eo successu continuátum, ut constans fúerit ómnium sensus, étiam pontificio confirmátus oraculo, Deum, sicut alios aliis tempóribus sanctos viros, ita Luthéro ejusdemque témporis hæreticis Ignatium et institutam ab eo societátem objecisse.\par
\switchcolumn
\EN
\noindent To make himself of greater use for the profit of souls, he determined to improve himself by education, beginning by going through the rudiments among little boys. He left nothing untried that could help towards the salvation of others, and it was marvellous what pain and mockery he cheerfully accepted on all hands, suffering ill-usage also, imprisonment and stripes almost unto death; but he was willing to suffer them all much more for the greater glory of his Master. At Paris he took to him seven comrades from the members of that University, men of different nations, but who had all taken the Degree of Master of Arts and in Divinity. With these seven he laid the first foundations of the Society of Jesus in the crypt at Montmartre, (upon the 15th day of August, in the year of Christ 1534.) When he afterwards organised the same Society at Rome he bound it by the closest bonds to the Apostolic See, adding to the three accustomed vows of Poverty, Chastity, and Obedience, a fourth, concerning Missions. Paul III. was the first Pope to receive and confirm the Institute, but it has since been approved by other Popes and by the Council of Trent. Ignatius, to spread the Faith, sent holy Francis Xavier to preach the Gospel in the Indies, and others in other parts of the world, and the war, which he thus proclaimed against paganism and heresy, was waged with such success, that it was the general belief, confirmed by the utterance of the Pope, that even as God had in other times raised up holy men specially to meet the needs of their day, so He had raised up against Luther and the heretics of that age, Ignatius and the Society which he had founded.\par
\switchcolumn*

\LA
\begin{Resp}\Rsign Amávit eum Dóminus, et ornávit eum: stolam glóriæ índuit eum, \Star{} Et ad portas paradísi coronávit eum.\end{Resp}
\begin{Resp}\Vsign Índuit eum Dóminus lorícam fídei, et ornávit eum.\end{Resp}
\begin{Resp}\Rsign Et ad portas paradísi coronávit eum.\end{Resp}
\switchcolumn
\EN
\begin{Resp}\Rsign The Lord loved him and beautified him; He clothed him with a robe of glory, \Star{} And crowned him at the gates of Paradise.\end{Resp}
\begin{Resp}\Vsign The Lord hath put on him the breast-plate of faith, and hath adorned him.\end{Resp}
\begin{Resp}\Rsign And crowned him at the gates of Paradise.\end{Resp}
\switchcolumn*

\LA
\LessonHead{Lectio VI}
\switchcolumn
\EN
\LessonHead{Lesson VI}
\switchcolumn*

\LA
\noindent Sed in primis inter Catholicos instaurare pietátem curæ fuit. Templórum nitor, catechismi traditio, conciónum ac sacramentórum frequéntia ab ipso increméntum accepére. Ipse, apertis ubíque locórum ad juventútem erudiéndam in litteris ac pietáte gymnasiis, eréctis Romæ Germanórum collegio, male nuptárum et periclitántium puellárum cœnobiis, utriusque sexus tam orphanórum quam catechumenórum dómibus, aliisque pietátis opéribus, indeféssus lucrandis Deo ánimis instábat; audítus aliquándo dicere, si optio darétur, malle se beatitúdinis incertum vivere, et interim Deo inservire et proximórum saluti, quam certum ejusdem glóriæ statim mori. In dæmones mirum exercuit impérium. Vultum ejus cælésti luce radiántem sanctus Philippus Nerius aliique conspexere. Denique, ætátis anno sexagesimo quinto, ad Dómini sui amplexum, cujus majórem glóriam in ore semper habúerat, semper in ómnibus quæsíerat, emigrávit. Quem Gregórius décimus quintus, magnis in Ecclésiam meritis et miraculis illustrem, Sanctórum fastis adscripsit, et Pius undecimus, sacrórum Antístitum votis obsecundans, ómnium Exercitiórum Spiritualium Patronum cælestem constítuit ac declarávit.\par
\switchcolumn
\EN
\noindent But the first care of Ignatius was to set forward godliness among Catholics. He was a great promoter of seemliness in the Churches, instruction in the Catechism, and often hearing Sermons and using the Sacraments. He opened schools everywhere to train up boys in godliness and good learning. At Rome he founded the German College, a home for fallen and another for imperilled girls, an orphanage for boys and another for girls, houses for converts under instruction, and other godly institutions. He never wearied in his work of gaining souls for God, and was sometimes heard to say that if he had the choice he would rather live without knowing whether he was to be among the blessed, and meanwhile work for God and the salvation of his neighbours, than know he was going to glory and die forthwith. He exercised an extraordinary power over devils. Holy Philip Neri and others saw heavenly light shining from his face. At last, (on the 31st day of July,) in the year of our Redemption 1556 and of his own age the sixty fifth, he passed away to the embrace of that Lord Whose greater glory had been the constant theme of his words and aim of all his works. He is very illustrious in the Church on account of his great deeds and miracles, and Gregory XV. enrolled him in the Kalendar of the Saints.\par
\switchcolumn*

\LA
\begin{Resp}\Rsign Iste homo perfécit ómnia quæ locútus est ei Deus, et dixit ad eum: Ingrédere in réquiem meam: \Star{} Quia te vidi justum coram me ex ómnibus géntibus.\end{Resp}
\begin{Resp}\Vsign Iste est, qui contémpsit vitam mundi, et pervénit ad cæléstia regna.\end{Resp}
\begin{Resp}\Rsign Quia te vidi justum coram me ex ómnibus géntibus.\end{Resp}
\begin{Resp}\Vsign Glória Patri, et Fílio, \Star{} et Spirítui Sancto.\end{Resp}
\begin{Resp}\Rsign Quia te vidi justum coram me ex ómnibus géntibus.\end{Resp}
\switchcolumn
\EN
\begin{Resp}\Rsign This is he which did according unto all that God commanded him; and God said unto him: Enter thou into My rest, \Star{} For thee have I seen righteous before Me among all people.\end{Resp}
\begin{Resp}\Vsign This is he which loved not his life in this world, and is come unto an everlasting kingdom.\end{Resp}
\begin{Resp}\Rsign For thee have I seen righteous before Me among all people.\end{Resp}
\begin{Resp}\Vsign Glory be to the Father, and to the Son, \Star{} and to the Holy Ghost.\end{Resp}
\begin{Resp}\Rsign For thee have I seen righteous before Me among all people.\end{Resp}
\switchcolumn*

\LA
\LessonHead{Lectio VII}
\switchcolumn
\EN
\LessonHead{Lesson VII}
\switchcolumn*

\LA
\LessonTitle{Léctio sancti Evangélii secúndum Lucam}
\switchcolumn
\EN
\LessonTitle{From the Holy Gospel according to Luke}
\switchcolumn*

\LA
\Cite{Luc 10:1-9}
\switchcolumn
\EN
\Cite{Luke 10:1-9}
\switchcolumn*

\LA
\noindent In illo témpore: Designávit Dóminus et álios septuagínta duos: et misit illos binos ante fáciem suam, in omnem civitátem et locum, quo erat ipse ventúrus. Et réliqua.\par
\switchcolumn
\EN
\noindent At that time, the Lord appointed other seventy-two also, and sent them two and two before His face into every city and place, whither He Himself would come. And so on.\par
\switchcolumn*

\LA
\LessonTitle{Homilía sancti Gregórii Papæ}
\switchcolumn
\EN
\LessonTitle{Homily by Pope St. Gregory (the Great.)}
\switchcolumn*

\LA
\Source{Homilia 17 in Evangelia}
\switchcolumn
\EN
\Source{17th Homily on the Gospels}
\switchcolumn*

\LA
\noindent Dóminus et Salvátor noster, fratres caríssimi, aliquándo nos sermónibus, aliquándo vero opéribus ádmonet. Ipsa étenim facta ejus præcépta sunt: quia dum áliquid tácitus facit, quid ágere debeámus innotéscit. Ecce enim binos in prædicatiónem discípulos mittit: quia duo sunt præcépta caritátis, Dei vidélicet amor, et próximi: et minus quam inter duos cáritas habéri non potest. Nemo enim próprie ad semetípsum habére caritátem dícitur: sed diléctio in álterum tendit, ut cáritas esse possit.\par
\switchcolumn
\EN
\noindent Dearly beloved brethren, our Lord and Saviour doth sometimes admonish us by words, and sometimes by works. Yea, His very works do themselves teach us for that which He doth silently His example still moveth us to copy. Behold how He sendeth forth His disciples to preach by two and two since there are two commandments to love, that is, a commandment to love God, and a commandment to love our neighbour and where there are not two, the one, being alone, hath not whereon to do the Lord's commandment. And no man can properly be said to love himself: for love tendeth outward toward our neighbour, if it be the love whereto the Gospel doth oblige us.\par
\switchcolumn*

\LA
\begin{Resp}\Rsign Iste est qui ante Deum magnas virtútes operátus est, et de omni corde suo laudávit Dóminum: \Star{} Ipse intercédat pro peccátis ómnium populórum.\end{Resp}
\begin{Resp}\Vsign Ecce homo sine queréla, verus Dei cultor, ábstinens se ab omni ópere malo, et pérmanens in innocéntia sua.\end{Resp}
\begin{Resp}\Rsign Ipse intercédat pro peccátis ómnium populórum.\end{Resp}
\switchcolumn
\EN
\begin{Resp}\Rsign This is he which wrought great wonders before God, and praised the Lord with all his heart. \Star{} May he pray for all people, that their sins may be forgiven unto them.\end{Resp}
\begin{Resp}\Vsign Behold a man without blame, a worshipper of God in truth, keeping himself clean from every evil work, and abiding still in his innocency.\end{Resp}
\begin{Resp}\Rsign May he pray for all people, that their sins may be forgiven unto them.\end{Resp}
\switchcolumn*

\LA
\LessonHead{Lectio VIII}
\switchcolumn
\EN
\LessonHead{Lesson VIII}
\switchcolumn*

\LA
\noindent Ecce enim binos ad prædicándum discípulos Dóminus mittit: quátenus hoc nobis tácitus ínnuat, quia qui caritátem erga álterum non habet, prædicatiónis offícium suscípere nullátenus debet. Bene autem dícitur, quia misit eos ante fáciem suam in omnem civitátem et locum, quo erat ipse ventúrus. Prædicatóres enim suos Dóminus séquitur: quia prædicátio prǽvenit, et tunc ad mentis nostræ habitáculum Dóminus venit, quando verba exhortatiónis præcúrrunt: atque per hoc véritas in mente suscípitur.\par
\switchcolumn
\EN
\noindent Behold, the Lord sendeth forth His disciples to preach by two and two and thus doing, He doth silently teach us that whosoever loveth not his neighbour, such a one it behoveth not to take upon him the office of a preacher. Well also is it said that He sent them before His face into every city and place whither He Himself would come. The Lord followeth His preachers first cometh preaching, and then the Lord Himself cometh to the house of our mind, whither the word of exhortation hath come before and so cometh the truth into our mind.\par
\switchcolumn*

\LA
\begin{Resp}\Rsign Sint lumbi vestri præcíncti, et lucérnæ ardéntes in mánibus vestris: \Star{} Et vos símiles homínibus exspectántibus dóminum suum, quando revertátur a núptiis.\end{Resp}
\begin{Resp}\Vsign Vigiláte ergo, quia nescítis qua hora Dóminus vester ventúrus sit.\end{Resp}
\begin{Resp}\Rsign Et vos símiles homínibus exspectántibus dóminum suum, quando revertátur a núptiis.\end{Resp}
\begin{Resp}\Vsign Glória Patri, et Fílio, \Star{} et Spirítui Sancto.\end{Resp}
\begin{Resp}\Rsign Et vos símiles homínibus exspectántibus dóminum suum, quando revertátur a núptiis.\end{Resp}
\switchcolumn
\EN
\begin{Resp}\Rsign Let your loins be girded about, and your lights burning; \Star{} And ye yourselves like unto men that wait for their lord, when he will return from the wedding.\end{Resp}
\begin{Resp}\Vsign Watch therefore, for ye know not what hour your Lord doth come.\end{Resp}
\begin{Resp}\Rsign And ye yourselves like unto men that wait for their lord, when he will return from the wedding.\end{Resp}
\begin{Resp}\Vsign Glory be to the Father, and to the Son, \Star{} and to the Holy Ghost.\end{Resp}
\begin{Resp}\Rsign And ye yourselves like unto men that wait for their lord, when he will return from the wedding.\end{Resp}
\switchcolumn*

\LA
\LessonHead{Lectio IX}
\switchcolumn
\EN
\LessonHead{Lesson IX}
\switchcolumn*

\LA
\noindent Hinc namque eísdem prædicatóribus Isaías dicit: Paráte viam Dómini, rectas fácite sémitas Dei nostri. Hinc fíliis Psalmísta ait: Iter fácite ei, qui ascéndit super occásum. Super occásum namque Dóminus ascéndit: quia unde in passióne occúbuit, inde majórem suam glóriam resurgéndo manifestávit. Super occásum vidélicet ascéndit; quia mortem quam pértulit, resurgéndo calcávit. Ei ergo qui ascéndit super occásum, iter fácimus, cum nos ejus glóriam vestris méntibus prædicámus, ut eas et ipse post véniens, per amóris sui præséntiam illústret.\par
\switchcolumn
\EN
\noindent Therefore, to preachers saith Isaiah: “Prepare ye the way of the Lord, make straight an highway for our God.” (xl. 3.) And again the Psalmist saith: “Spread a path before him that rideth upon the West.” (lxvii. 4.) The Lord rideth upon the West above that from which in death He veiled His glory, hath He royally exalted that glory that excelleth, even the glory of His rising again. He rideth upon the West, Who, being risen again from the dead, is throned high above the death to which He bowed. Before Him, therefore, That rideth upon the West, we spread a path, when we set forth His glory before the eyes of your mind, to the end that He Himself may come after, and Himself enlighten the same your minds by His presence and His love.\par
\switchcolumn*

\LA
\TeDeum{Te Deum}
\switchcolumn
\EN
\TeDeum{Te Deum}
\switchcolumn*

\end{paracol}

\SeasonTitle{Mensis Augustus}{August}

\addcontentsline{toc}{section}{Mensis Augustus \textit{— August}}

\par\needspace{10\baselineskip}\addvspace{1.2em}{\centering\small\textsc{Die 1 Augusti} · \textsc{1 August}\par}
\Office{S. Petri ad Vincula}{St. Peter in Chains}{Duplex majus}
\addcontentsline{toc}{subsection}{Die 1 Augusti · S. Petri ad Vincula \textit{— St. Peter in Chains}}
\begin{paracol}{2}
\LA
\RefLine{Lectio IX: de Ss. Martyrum Machabæorum (08-01r).}
\switchcolumn
\EN
\RefLine{Lesson IX: of The Holy Maccabean Martyrs (08-01r).}
\switchcolumn*

\LA
\LessonHead{Lectio I}
\switchcolumn
\EN
\LessonHead{Lesson I}
\switchcolumn*

\LA
\LessonTitle{De Actibus Apostolórum}
\switchcolumn
\EN
\LessonTitle{From the Acts of the Apostles}
\switchcolumn*

\LA
\Cite{Act 12:1-5}
\switchcolumn
\EN
\Cite{Acts 12:1-5}
\switchcolumn*

\LA
\noindent Misit Heródes rex manus ut afflígeret quosdam de Ecclésia. Occídit autem Jacóbum, fratrem Joánnis, gládio. Videns autem quia placéret Judǽis, appósuit ut apprehénderet et Petrum. Erant autem dies azymórum. Quem, cum apprehendísset, misit in cárcerem, tradens quátuor quaterniónibus mílitum custodiéndum, volens post Pascha prodúcere eum pópulo. Et Petrus quidem servabátur in cárcere; orátio autem fiébat sine intermissióne ab Ecclésia ad Deum pro eo.\par
\switchcolumn
\EN
\noindent And at the same time, Herod the king stretched forth his hands, to afflict some of the church. And he killed James, the brother of John, with the sword. And seeing that it pleased the Jews, he proceeded to take up Peter also. Now it was in the days of the Azymes. And when he had apprehended him, he cast him into prison, delivering him to four files of soldiers to be kept, intending, after the pasch, to bring him forth to the people. Peter therefore was kept in prison. But prayer was made without ceasing by the church unto God for him.\par
\switchcolumn*

\LA
\begin{Resp}\Rsign Simon Petre, ántequam de navi vocárem te, novi te et super plebem meam príncipem te constítui: \Star{} Et claves regni cælórum trádidi tibi.\end{Resp}
\begin{Resp}\Vsign Quodcúmque ligáveris super terram erit ligátum et in cælis; et quodcúmque sólveris super terram, erit solútum et in cælis.\end{Resp}
\begin{Resp}\Rsign Et claves regni cælórum trádidi tibi.\end{Resp}
\switchcolumn
\EN
\begin{Resp}\Rsign Simon Peter, before I called thee out of the ship, I knew thee, and appointed thee for a ruler over My people. \Star{} And I have given unto thee the keys of the kingdom of heaven.\end{Resp}
\begin{Resp}\Vsign Whatsoever, thou shalt bind on earth, shall be bound in heaven; and whatsoever thou shalt loose on earth, shall be loosed in heaven.\end{Resp}
\begin{Resp}\Rsign And I have given unto thee the keys of the kingdom of heaven.\end{Resp}
\switchcolumn*

\LA
\LessonHead{Lectio II}
\switchcolumn
\EN
\LessonHead{Lesson II}
\switchcolumn*

\LA
\Cite{Act 12:6-8}
\switchcolumn
\EN
\Cite{Acts 12:6-8}
\switchcolumn*

\LA
\noindent Cum autem productúrus eum esset Heródes, in ipsa nocte erat Petrus dórmiens inter duos mílites vinctus caténis duábus, et custódes ante óstium custodiébant cárcerem. Et ecce Angelus Dómini ástitit, et lumen refúlsit in habitáculo, percussóque látere Petri, excitávit eum, dicens: Surge velóciter. Et cecidérunt caténæ de mánibus ejus. Dixit autem Angelus ad eum: Præcíngere, et cálcea te cáligas tuas. Et fecit sic. Et dixit illi: Circúmda tibi vestiméntum tuum, et séquere me.\par
\switchcolumn
\EN
\noindent And when Herod would have brought him forth, the same night Peter was sleeping between two soldiers, bound with two chains: and the keepers before the door kept the prison. And behold an angel of the Lord stood by him: and a light shined in the room: and he striking Peter on the side, raised him up, saying: Arise quickly. And the chains fell off from his hands. And the angel said to him: Gird thyself, and put on thy sandals. And he did so. And he said to him: Cast thy garment about thee, and follow me.\par
\switchcolumn*

\LA
\begin{Resp}\Rsign Si díligis me, Simon Petre, pasce oves meas. Dómine, tu nosti quia amo te, \Star{} Et ánimam meam pono pro te.\end{Resp}
\begin{Resp}\Vsign Si oportúerit me mori tecum, non te negábo.\end{Resp}
\begin{Resp}\Rsign Et ánimam meam pono pro te.\end{Resp}
\switchcolumn
\EN
\begin{Resp}\Rsign Simon Peter, if thou lovest Me, feed My sheep. Lord, Thou knowest that I love thee; \Star{} I will lay down my life for thy sake.\end{Resp}
\begin{Resp}\Vsign If I should die with thee, I will not deny thee.\end{Resp}
\begin{Resp}\Rsign I will lay down my life for thy sake.\end{Resp}
\switchcolumn*

\LA
\LessonHead{Lectio III}
\switchcolumn
\EN
\LessonHead{Lesson III}
\switchcolumn*

\LA
\Cite{Act 12:9-11}
\switchcolumn
\EN
\Cite{Acts 12:9-11}
\switchcolumn*

\LA
\noindent Et éxiens sequebátur eum, et nesciébat quia verum est quod fiébat per Angelum; existimábat autem se visum vidére. Transeúntes autem primam et secúndam custódiam venérunt ad portam férream, quæ ducit ad civitátem, quæ ultro apérta est eis; et exeúntes processérunt vicum unum, et contínuo discéssit Angelus ab eo. Et Petrus ad se revérsus, dixit: Nunc scio vere quia misit Dóminus Angelum suum, et erípuit me de manu Heródis et de omni exspectatióne plebis Judæórum.\par
\switchcolumn
\EN
\noindent And going out, he followed him, and he knew not that it was true which was done by the angel: but thought he saw a vision. And passing through the first and the second ward, they came to the iron gate that leadeth to the city, which of itself opened to them. And going out, they passed on through one street: and immediately the angel departed from him. And Peter coming to himself, said: Now I know in very deed, that the Lord hath sent his angel, and hath delivered me out of the hand of Herod, and from all the expectation of the people of the Jews.\par
\switchcolumn*

\LA
\begin{Resp}\Rsign Tu es Petrus, et super hanc petram ædificábo Ecclésiam meam, et portæ ínferi non prævalébunt advérsus eam: \Star{} Et tibi dabo claves regni cælórum.\end{Resp}
\begin{Resp}\Vsign Quodcúmque ligáveris super terram, erit ligátum et in cælis; et quodcúmque sólveris super terram, erit solútum et in cælis.\end{Resp}
\begin{Resp}\Rsign Et tibi dabo claves regni cælórum.\end{Resp}
\begin{Resp}\Vsign Glória Patri, et Fílio, \Star{} et Spirítui Sancto.\end{Resp}
\begin{Resp}\Rsign Et tibi dabo claves regni cælórum.\end{Resp}
\switchcolumn
\EN
\begin{Resp}\Rsign Thou art Peter, and upon this rock I will build My Church, and the gates of hell shall not prevail against it. \Star{} And I will give unto thee the keys of the kingdom of heaven.\end{Resp}
\begin{Resp}\Vsign Whatsoever thou shalt bind on earth, shall be bound in heaven; and whatsoever thou shalt loose on earth, shall be loosed in heaven.\end{Resp}
\begin{Resp}\Rsign And I will give unto thee the keys of the kingdom of heaven.\end{Resp}
\begin{Resp}\Vsign Glory be to the Father, and to the Son, \Star{} and to the Holy Ghost.\end{Resp}
\begin{Resp}\Rsign And I will give unto thee the keys of the kingdom of heaven.\end{Resp}
\switchcolumn*

\LA
\LessonHead{Lectio IV}
\switchcolumn
\EN
\LessonHead{Lesson IV}
\switchcolumn*

\LA
\noindent Theodósio junióre imperánte, cum Eudócia ejus uxor Jerosólymam solvéndi voti causa venísset, ibi multis est affécta munéribus. Præ céteris insígne donum accépit férreæ caténæ, auro gemmísque ornátæ, quam illam esse affirmábant, qua Petrus Apóstolus ab Heróde vinctus fúerat. Eudócia, caténam pie veneráta, eam póstea Romam ad fíliam Eudóxiam misit; quæ illam Pontífici máximo détulit. Isque vicíssim illi monstrávit álteram caténam, qua, Neróne imperatóre, idem Apóstolus constríctus fúerat.\par
\switchcolumn
\EN
\noindent In (the year of our Lord 439, in) the reign of the Emperor Theodosius the younger, his wife went to Jerusalem in fulfilment of a vow, and there was gifted with many presents. Among other things, they gave her in especial an iron chain, adorned with gold and precious stones, which they affirmed to be the same wherewith the Apostle Peter had been bound by King Herod. Eudocia, with godly reverence, afterwards sent this chain to Rome, to her daughter Eudoxia, who brought it to the Pope, and the Pope in return showed to her another chain wherewith the same Apostle had been shackled under the Emperor Nero.\par
\switchcolumn*

\LA
\begin{Resp}\Rsign Dómine, si tu es, jube me veníre ad te super aquas. \Star{} Et exténdens manum apprehéndit eum, et dixit Jesus: Módicæ fídei, quare dubitásti?\end{Resp}
\begin{Resp}\Vsign Cumque vidísset ventum válidum veniéntem, tímuit; et, cum cœpísset mergi, clamávit dicens: Dómine, salvum me fac.\end{Resp}
\begin{Resp}\Rsign Et exténdens manum apprehéndit eum, et dixit Jesus: Módicæ fídei, quare dubitásti?\end{Resp}
\switchcolumn
\EN
\begin{Resp}\Rsign Lord, if it be Thou, bid me come unto thee on the water; \Star{} And Jesus stretched forth His Hand, and caught him, and said unto him: O thou of little faith, wherefore didst thou doubt?\end{Resp}
\begin{Resp}\Vsign But when he saw the wind boisterous, he was afraid and, beginning to sink, he cried, saying: Lord, save me!\end{Resp}
\begin{Resp}\Rsign And Jesus stretched forth His Hand, and caught him, and said unto him: O thou of little faith, wherefore didst thou doubt?\end{Resp}
\switchcolumn*

\LA
\LessonHead{Lectio V}
\switchcolumn
\EN
\LessonHead{Lesson V}
\switchcolumn*

\LA
\noindent Cum ígitur Póntifex Románam caténam cum ea, quæ Jerosólymis alláta fúerat, contulísset, factum est ut illæ inter se sic connecteréntur, ut non duæ, sed una caténa ab eódem artífice confécta esse viderétur. Quo miráculo tantus honor sacris illis vínculis habéri cœpit, ut proptérea hoc nómine sancti Petri ad víncula ecclésia, título Eudóxiæ, dedicáta sit in Exquíliis, ejúsque memóriæ dies festus institútus Kaléndis Augústi.\par
\switchcolumn
\EN
\noindent When then the Pope put together the Roman chain and that which had been brought from Jerusalem, it came to pass that they got so entangled the one with the other that they seemed no longer two but one chain. From this wonder these holy fetters began to receive such honour, that Eudoxia's Church of St. Peter on the Esquiline Mount was dedicated under the name of St. Peter in Chains, and a Feast - Day instituted upon the first day of August in memory of it.\par
\switchcolumn*

\LA
\begin{Resp}\Rsign Surge, Petre, et índue te vestiméntis tuis, áccipe fortitúdinem ad salvándas gentes: \Star{} Quia cecidérunt caténæ de mánibus tuis.\end{Resp}
\begin{Resp}\Vsign Angelus Dómini ástitit, et lumen refúlsit in habitáculo cárceris, percussóque látere Petri, excitávit eum, dicens: Surge, velóciter.\end{Resp}
\begin{Resp}\Rsign Quia cecidérunt caténæ de mánibus tuis.\end{Resp}
\switchcolumn
\EN
\begin{Resp}\Rsign Arise, O Peter! Cast thy garment about thee, gird thee with strength for the saving of the nations. \Star{} The chains are fallen off from thine hands.\end{Resp}
\begin{Resp}\Vsign The Angel of the Lord came upon him, and a light shined in the prison and he smote Peter on the side and raised him up, saying: Arise up quickly.\end{Resp}
\begin{Resp}\Rsign The chains are fallen off from thine hands.\end{Resp}
\switchcolumn*

\LA
\LessonHead{Lectio VI}
\switchcolumn
\EN
\LessonHead{Lesson VI}
\switchcolumn*

\LA
\noindent Quo ex témpore honos, qui eo die profánis Gentílium celebritátibus tríbui sólitus erat, Petri vínculis habéri cœpit, quæ tacta ægros sanábant et dǽmones ejiciébant. Quo in génere, anno salútis humánæ nongentésimo sexagésimo nono áccidit, ut quidam comes, Ottónis imperatóris familiáris, occupátus ab immúndo spíritu, seípsum déntibus dilaniáret. Quare is jussu imperatóris ad Joánnem Pontíficem dúcitur, qui ut sacra caténa cómitis collum áttigit, erúmpens nefárius spíritus hóminem líberum relíquit; ac deínceps in Urbe sanctórum vinculórum relígio propagáta est.\par
\switchcolumn
\EN
\noindent From that time forth the honour which before had used to be paid to the profane festivity of the Gentiles, (held in memory of the dedication of the temple of Mars, and of the birth of Claudius,) began to be turned to the Chains of Peter, whose very touch healed the sick, and drove out devils. Among other such cases there befell in the year of man's Redemption 969, that of a certain Count, a servant of the Emperor Otho, who was possessed by an unclean spirit, and tore himself with his own teeth. This man the Emperor ordered to be taken to Pope John, and as soon as he had touched the Count's neck with the hallowed chains, the foul spirit came out of him, and left him free. And thenceforward the reverence for these holy chains greatly increased in the City.\par
\switchcolumn*

\LA
\begin{Resp}\Rsign Tu es pastor óvium, princeps Apostolórum: tibi trádidit Deus ómnia regna mundi: \Star{} Et ídeo tibi tráditæ sunt claves regni cælórum.\end{Resp}
\begin{Resp}\Vsign Quodcúmque ligáveris super terram, erit ligátum et in cælis; et quodcúmque sólveris super terram, erit solútum et in cælis.\end{Resp}
\begin{Resp}\Rsign Et ídeo tibi tráditæ sunt claves regni cælórum.\end{Resp}
\begin{Resp}\Vsign Glória Patri, et Fílio, \Star{} et Spirítui Sancto.\end{Resp}
\begin{Resp}\Rsign Et ídeo tibi tráditæ sunt claves regni cælórum.\end{Resp}
\switchcolumn
\EN
\begin{Resp}\Rsign Thou art the Shepherd of the sheep, and the Prince of the Apostles, and unto thee hath God given all the kingdoms of the world. \Star{} Therefore unto thee hath He given the keys of the kingdom of heaven.\end{Resp}
\begin{Resp}\Vsign Whatsoever thou shalt bind on earth shall be bound in heaven; and whatsoever thou shalt loose on earth shall be loosed in heaven.\end{Resp}
\begin{Resp}\Rsign Therefore unto thee hath He given the keys of the kingdom of heaven.\end{Resp}
\begin{Resp}\Vsign Glory be to the Father, and to the Son, \Star{} and to the Holy Ghost.\end{Resp}
\begin{Resp}\Rsign Therefore unto thee hath He given the keys of the kingdom of heaven.\end{Resp}
\switchcolumn*

\LA
\LessonHead{Lectio VII}
\switchcolumn
\EN
\LessonHead{Lesson VII}
\switchcolumn*

\LA
\LessonTitle{Léctio sancti Evangélii secúndum Matthǽum}
\switchcolumn
\EN
\LessonTitle{From the Holy Gospel according to Matthew}
\switchcolumn*

\LA
\Cite{Matt 16:13-19}
\switchcolumn
\EN
\Cite{Matt 16:13-19}
\switchcolumn*

\LA
\noindent In illo témpore: Venit Jesus in partes Cæsaréæ Philíppi, et interrogábat discípulos suos, dicens: Quem dicunt hómines esse Fílium hóminis? Et réliqua.\par
\switchcolumn
\EN
\noindent At that time Jesus came into the coasts of Caesarea Philippi, and He asked His disciples, saying Who do men say that I, the Son of Man, am And so on.\par
\switchcolumn*

\LA
\LessonTitle{Homilía sancti Augustíni Epíscopi}
\switchcolumn
\EN
\LessonTitle{Homily by St. Augustine, Bishop (of Hippo.)}
\switchcolumn*

\LA
\Source{Ex serm. 29 de Sanctis, in medio}
\switchcolumn
\EN
\Source{29th on the Saints}
\switchcolumn*

\LA
\noindent Solus Petrus inter Apóstolos méruit audíre: Amen dico tibi, quia tu es Petrus, et super hanc petram ædificábo Ecclésiam meam; dignus certe, qui ædificándis in domo Dei pópulis lapis esset ad fundaméntum, colúmna ad sustentáculum, clavis ad regnum. De hoc ait sermo divínus: Et ponébant, inquit, infírmos suos, ut umbra saltem transeúntis Petri obumbráret eos. Si tunc opem ferre póterat umbra córporis, quanto magis nunc plenitúdo virtútis? si tunc supplicántibus próderat aura quædam pertranseúntis, quanto magis grátia nunc permanéntis? Mérito per omnes Christi ecclésias auro pretiósius habétur ferrum illud pœnálium vinculórum.\par
\switchcolumn
\EN
\noindent Peter was the only one of the Apostles who was worthy to hear the words Amen, I say unto thee that thou art Peter, and upon this rock I will build My Church. Worthy indeed must he be, who, when the nations are to be built up into a Temple of God, is chosen as the ground-stone whereon the building is to stand the pillar whereby it is to be held up, and the key wherethrough entrance is to be made into the kingdom. Concerning him the Word of God saith That they brought forth the sick into the streets, and laid them on beds and couches, that at the least the shadow of Peter passing by might overshadow some of them. (Acts v. 15.) If the shadow of his body then could give help, how much more shall the fulness of his strength give help now. If the very air, as he passed by, was then profitable to such as besought him, how much more shall his favour profit where now he abideth. It is with reason that, throughout all the Churches of Christ, the iron chains wherewith he was afflicted are reckoned more precious than gold.\par
\switchcolumn*

\LA
\begin{Resp}\Rsign Ego pro te rogávi, Petre, ut non defíciat fides tua: \Star{} Et tu aliquándo convérsus confírma fratres tuos.\end{Resp}
\begin{Resp}\Vsign Caro et sanguis non revelávit tibi, sed Pater meus, qui est in cælis.\end{Resp}
\begin{Resp}\Rsign Et tu aliquándo convérsus confírma fratres tuos.\end{Resp}
\switchcolumn
\EN
\begin{Resp}\Rsign Peter, I have prayed for thee, that thy faith fail not \Star{} And when thou art converted, strengthen thy brethren.\end{Resp}
\begin{Resp}\Vsign Flesh and blood hath not revealed it unto thee, but My Father Which is in heaven.\end{Resp}
\begin{Resp}\Rsign And when thou art converted, strengthen thy brethren.\end{Resp}
\switchcolumn*

\LA
\LessonHead{Lectio VIII}
\switchcolumn
\EN
\LessonHead{Lesson VIII}
\switchcolumn*

\LA
\noindent Si tam medicábilis fuit obumbrátio visitántis, quanto magis caténa vinciéntis? Si inánis quædam spécies vácuæ imáginis habére pótuit in se vim salútis, quanto plus de córpore meruérunt attráhere salubritátis, férreo póndere sacris impréssa membris, víncula passiónis? Si ad præsídia supplicántium tam potens fuit ante martýrium, quanto magis éfficax post triúmphum? Felíces illi nexus, qui, de mánicis et compédibus in corónam mutándi, Apóstolum contingéntes, Mártyrem reddidérunt! Felícia víncula, quæ reum suum usque ad Christi crucem, non tam condemnatúra quam consecratúra, misérunt!\par
\switchcolumn
\EN
\noindent If his shadow as a visitor was so healthful, what is his chain now that he bindeth and looseth. If his empty image in the air had healing power, how much power must have been contracted from his body by those chains, whose iron weight sank into his holy limbs during his suffering. If, before he testified, he was so mighty to aid them that called upon him, how much mightier is he now since his victory. Blessed were the links, doomed to be changed from fetters and shackles, into a crown, which by touching the Apostle, made him a Martyr. Blessed were the chains, whose prisoner left them for the Cross of Christ, and which brought him thither, not as the instruments of condemnation, but of sanctification.\par
\switchcolumn*

\LA
\begin{Resp}\Rsign Quem dicunt hómines esse Filium hóminis? dixit Jesus discípulis suis. Respóndens Petrus dixit: Tu es Christus Fílius Dei vivi. \Star{} Et ego dico tibi, quia tu es Petrus, et super hanc petram ædificábo Ecclésiam meam.\end{Resp}
\begin{Resp}\Vsign Beátus es, Simon Bar-Jona, quia caro et sanguis non revelávit tibi, sed Pater meus qui est in cælis.\end{Resp}
\begin{Resp}\Rsign Et ego dico tibi, quia tu es Petrus, et super hanc petram ædificábo Ecclésiam meam.\end{Resp}
\begin{Resp}\Vsign Glória Patri, et Fílio, \Star{} et Spirítui Sancto.\end{Resp}
\begin{Resp}\Rsign Et ego dico tibi, quia tu es Petrus, et super hanc petram ædificábo Ecclésiam meam.\end{Resp}
\switchcolumn
\EN
\begin{Resp}\Rsign Jesus asked His disciples, saying: Who do men say that I, the Son of Man, am? Peter answered, and said: Thou art the Christ, the Son of the living God. \Star{} And I say unto thee, that thou art Peter, and upon this rock I will build My Church.\end{Resp}
\begin{Resp}\Vsign Blessed art thou, Simon Bar-Jona, for flesh and blood hath not revealed it unto thee, but My Father Which is in heaven.\end{Resp}
\begin{Resp}\Rsign And I say unto thee, that thou art Peter, and upon this rock I will build My Church.\end{Resp}
\begin{Resp}\Vsign Glory be to the Father, and to the Son, \Star{} and to the Holy Ghost.\end{Resp}
\begin{Resp}\Rsign And I say unto thee, that thou art Peter, and upon this rock I will build My Church.\end{Resp}
\switchcolumn*

\end{paracol}

\par\needspace{10\baselineskip}\addvspace{1.2em}{\centering\small\textsc{Die 1 Augusti} · \textsc{1 August}\par}
\Office{Ss. Martyrum Machabæorum}{The Holy Maccabean Martyrs}{Simplex}
\addcontentsline{toc}{subsection}{Die 1 Augusti · Ss. Martyrum Machabæorum \textit{— The Holy Maccabean Martyrs}}
\begin{paracol}{2}
\LA
\LessonHead{Lectio IX (pro commemoratione)}
\switchcolumn
\EN
\LessonHead{Lesson IX (when commemorated)}
\switchcolumn*

\LA
\LessonTitle{Commemoratio: Ss. Martyrum Machabæorum}
\switchcolumn
\EN
\LessonTitle{Commemoration of: The Holy Maccabean Martyrs}
\switchcolumn*

\LA
\Source{Orátio 20 in Machabæos}
\switchcolumn
\EN
\Source{Sermon 20th on the Machabees}
\switchcolumn*

\LA
\noindent Sermo sancti Gregórii Nazianzéni\par
\switchcolumn
\EN
\noindent From the Sermons of St. Gregory of Nazianzus (Patriarch of Constantinople)\par
\switchcolumn*

\LA
Quidnam Machabǽi? Horum enim nómine dies festus præsénti frequéntia celebrátur. Qui, etsi apud multos non sunt in honóre, quod illud certámen post Christum non suscepérunt; tamen digni sunt qui ab ómnibus honoréntur, quia pro pátriis légibus et institútis fortes constantésque se præbuérunt. Nam qui martýrium subiére ante Christi passiónem, quid factúri erant, si post Christum persecutiónem passi fuíssent, ejúsque mortem, nostræ salútis causa suscéptam, ad imitándum propósitam habuíssent? Quorum enim, nullo propósito exémplo, talis tántaque virtus fuit; an non ii, exémplum illud intuéntes, fórtius in certámen descendíssent? Quin étiam mýstica quædam et occúlta rátio, mihi quidem atque ómnibus Dei amatóribus valde probábilis est, néminem eórum qui ante Christi advéntum martýrio consummáti sunt, id sine fide in Christum cónsequi potuísse.\par
\switchcolumn
\EN
What were the Machabees For it is under their name that the Festival at which you are this day assembled is kept. It is true that many persons do not hold them in honour, because they fought before the coming of Christ nevertheless they deserve to be venerated by all men, for they bore themselves bravely and faithfully in defence of the laws and ordinances of their people. They that underwent martyrdom before Christ came, what would not have been their deeds, if they had suffered persecution after He came, and had had before them for a pattern the Death which He embraced for the sake of man's salvation With no example to lead them, their bravery was what it was had they had the example before their eyes, would they not have gone down with double nerve to the battle There is a mystic and subtle idea, which seemeth very likely to me and to all lovers of God, that none of those who were crowned with martyrdom before Christ came, could have been so, unless they had had faith in Christ.\par
\switchcolumn*

\LA
\TeDeum{Te Deum}
\switchcolumn
\EN
\TeDeum{Te Deum}
\switchcolumn*

\end{paracol}

\par\needspace{10\baselineskip}\addvspace{1.2em}{\centering\small\textsc{Die 2 Augusti} · \textsc{2 August}\par}
\Office{S. Alfonsi Mariæ de Ligorio Episcopi Confessoris et Ecclesiæ Doctoris}{St. Alphonsus Liguori, Bishop, Confessor and Doctor of the Church}{Duplex}
\addcontentsline{toc}{subsection}{Die 2 Augusti · S. Alfonsi Mariæ de Ligorio Episcopi Confessoris et Ecclesiæ Doctoris \textit{— St. Alphonsus Liguori, Bishop, Confessor and Doctor of the Church}}
\begin{paracol}{2}
\LA
\RefLine{Lectiones I–III: de Scriptura occurrente.}
\switchcolumn
\EN
\RefLine{Lessons I–III: from the Scripture of the day.}
\switchcolumn*

\LA
\RefLine{Lectio IX: de S. Stephani Papæ Martyris (08-02o).}
\switchcolumn
\EN
\RefLine{Lesson IX: of S. Stephen Pope and Martyr (08-02o).}
\switchcolumn*

\LA
\LessonHead{Lectio IV}
\switchcolumn
\EN
\LessonHead{Lesson IV}
\switchcolumn*

\LA
\noindent Alfonsus Maria de Ligorio, Neapoli nobílibus paréntibus natus, ab ineunte ætate non obscura præbuit sanctitátis indicia. Eum adhuc infántem cum paréntes obtulissent sancto Francisco de Hieronymo e societáte Jesu, is bene precátus edíxit eumdem ad nonagesimum usque annum perventurum, ad episcopalem dignitátem evectum iri, maximóque Ecclésiæ bono futurum. Jam tum a puerítia, a ludis abhorrens, nóbiles ephébos ad christianam modestiam verbo et exemplo componebat. Adoléscens, dato piis sodalitátibus nómine, in publicis nosocomíis ægrotis inservire, jugi in templis oratióni vacare, ac sacra mysteria frequenter obire in deliciis habebat. Pietátem litterárum stúdiis adeo conjunxit, ut séxdecim vix annos natus utriusque juris lauream in pátria universitáte fúerit assecutus. Patri obtemperans, causárum patrocinia suscépit; in quo munere obeundo, etsi magnam sibi laudem comparasset, fori tamen pericula expertus, ejusmodi vitæ institutum ultro dimísit. Spreto ígitur præcláro conjugio sibi a patre propósito, avíta primogenitúra abdicata, et ad aram Vírginis de Mercede ense suspenso, divinis ministériis se mancipávit. Sacerdos factus, tanto zelo irruit in vitia, ut apostolico munere fungens, huc illuc pérvolans, ingéntes perditórum hóminum conversiónes perágeret. Páuperum præsertim et ruricolárum miseratus, congregatiónem Presbyterórum instituit sanctíssimi Redemptoris, qui, ipsum Redemptórem secuti, per agros, pagos, et castella paupéribus evangelizarent.\par
\switchcolumn
\EN
\noindent Alphonsus Mary de' Liguori was born of a noble family, at (Marianella, near) Naples, (on the 26th day of September, in the year of salvation 1696.) From his earliest days he gave no dark signs of holiness. When he was but a babe, his parents carried him to holy Francis de Hieronymo, of the Society of Jesus, and holy Francis, after long prayer, said that the child would live to ninety years of age, that he would become a Bishop, and that he would be a great blessing to the Church. From his childhood, he had a strong distaste to games, and by his entreaty and example, induced the noble pages (of the Court, among whom he served,) to conduct themselves with Christian decency. As a young man, he became a member of diverse godly guilds, and made it among his delights to nurse the sick in the hospitals, to spend much time in prayer in the Churches, and often to receive the Holy Sacraments. With his godliness he so joined zeal for learning, that when he was scarcely sixteen years of age he took degrees in Canon and Civil law in the University of Naples. In obedience to the wish of his father, he adopted the profession of an advocate, in which he gained great credit, but, finding dangers in the practice of the law, he entirely gave it up. He declined a very brilliant marriage which was proposed to him by his father, resigned his family inheritance as an eldest son, hung up his sword at the Altar of the Blessed Virgin Mary, styled of Ransom, and surrendered himself altogether to the service of God. He became a Priest, (in 1726,) and made so zealous an onslaught on sin, running hither and thither in the office of an Apostle, that he accomplished the conversion of multitudes of lost creatures. The poor and the country-people most chiefly roused his compassion, and (in 1742) he founded the Congregation of Priests called that of the Most Holy Redeemer, to follow the Redeemer's footsteps by preaching the Gospel to the poor throughout the fields, villages, and hamlets.\par
\switchcolumn*

\LA
\begin{Resp}\Rsign Invéni David servum meum, óleo sancto meo unxi eum: \Star{} Manus enim mea auxiliábitur ei.\end{Resp}
\begin{Resp}\Vsign Nihil profíciet inimícus in eo, et fílius iniquitátis non nocébit ei.\end{Resp}
\begin{Resp}\Rsign Manus enim mea auxiliábitur ei.\end{Resp}
\switchcolumn
\EN
\begin{Resp}\Rsign I have found David My servant, with My holy oil have I anointed him \Star{} For My hand shall help him.\end{Resp}
\begin{Resp}\Vsign The enemy shall prevail nothing against him, nor the son of wickedness afflict him.\end{Resp}
\begin{Resp}\Rsign For My hand shall help him.\end{Resp}
\switchcolumn*

\LA
\LessonHead{Lectio V}
\switchcolumn
\EN
\LessonHead{Lesson V}
\switchcolumn*

\LA
\noindent Ne autem a propósito umquam diverteret, perpetuo se voto obstrinxit nullam témporis jacturam faciéndi. Hinc animárum zelo succénsus, tum divini verbi prædicatióne, tum scriptis sacra eruditióne et pietáte refertis, ánimas Christo lucrifácere et ad perfectiórem vitam adducere studuit. Mirum sane quot odia exstinxerit, quot devios ad rectum salútis iter revocaverit. Dei Genitrícis cultor exímius, de illíus laudibus librum edidit, ac de iis dum ferventius concionando disserit, a Vírginis imagine, in eum immisso miro splendore, totus fácie coruscare et in éxtasim rapi coram universo pópulo non semel visus est. Dominicæ passiónis et sacræ Eucharistiæ contemplator assiduus, ejus cultum mirifice propagávit. Dum vero ad ejus aram oraret vel Sacrum fáceret, quod numquam omisit, præ amoris veheméntia, vel seraphicis liquescebat ardóribus, vel insólitis quatiebátur mótibus, vel abstrahebátur a sensibus. Miram vitæ innocéntiam, quam nulla umquam lethali labe fœdávit, pari cum pœniténtia socians, corpus suum inedia, férreis caténulis, ciliciis cruentáque flagellatióne castigábat. Inter hæc prophetíæ, scrutatiónis cordium, bilocatiónis et miraculórum donis incláruit.\par
\switchcolumn
\EN
\noindent That he might not turn aside from his work, he bound him- self by a vow never to lose any time. Inflamed with the love of souls, he toiled to gain them to Christ and to amend their lives, not only by preaching of the word of God, but also by writings full of holy learning and godliness. It is a marvel how many hatreds he stilled, and how many backsliders he led again into the paths of salvation. He was eminently devoted to the Mother of God, published a book on her glories, and when he was earnestly speaking thereof in his sermons, it happened more than once that all the people openly saw a strange brightness fall upon him from her image, till all his countenance shone, and he was rapt in an ecstasy. The sufferings of the Lord and the Holy Eucharist were ever before his eyes, and to them he spread abroad a wonderful love. When he was praying before the Altar of the Blessed Sacrament, or celebrating the Holy Liturgy, which he never failed to do every day, through the seraphic violence of his love, he wept burning tears, or shook with strange movements, or became altogether beside himself. He joined a wonderful innocence and purity, which he never polluted by the stain of deadly sin, to a wonderful depth of repentance, and chastised his body with hunger, iron chains, hair-cloth, and scourgings even to blood-shedding. Among all these things he was remarkable for the gift of prophecy, the power of seeing into the hearts of men, the ability to be in more places than one at the same time, and other miracles.\par
\switchcolumn*

\LA
\begin{Resp}\Rsign Pósui adjutórium super poténtem, et exaltávi eléctum de plebe mea: \Star{} Manus enim mea auxiliábitur ei.\end{Resp}
\begin{Resp}\Vsign Invéni David servum meum, óleo sancto meo unxi eum.\end{Resp}
\begin{Resp}\Rsign Manus enim mea auxiliábitur ei.\end{Resp}
\switchcolumn
\EN
\begin{Resp}\Rsign I have laid help upon one that is mighty, and have exalted one chosen out of My people \Star{} For My hand shall help him.\end{Resp}
\begin{Resp}\Vsign I have found David My servant, with My holy oil have I anointed him.\end{Resp}
\begin{Resp}\Rsign For My hand shall help him.\end{Resp}
\switchcolumn*

\LA
\LessonHead{Lectio VI}
\switchcolumn
\EN
\LessonHead{Lesson VI}
\switchcolumn*

\LA
\noindent Ab ecclesiásticis dignitátibus sibi oblátis constantíssime abhorruit; at, Clementis décimi tertii Pontificis auctoritate coactus, sanctæ Agathæ Gothórum ecclésiam gubernándum suscépit. Epíscopus, externum dumtaxat habitum, non autem sevéram vivéndi ratiónem immutávit. Eadem frugalitas, summus christianæ disciplinæ zelus, impensum in vitiis coërcéndis arcendisque erróribus, et in réliquis pastorálibus munéribus obeundis studium. Liberalis in páuperes, omnes ecclésiæ proventus iisdem distribuebat, ac, urgénte annonæ caritate, ipsam domesticam supellectilem in aléndis famélicis erogávit. Omnibus ómnia factus, sanctimoniáles ad perfectiórem vivéndi formam redégit, suæque congregatiónis moniálium monastérium constituéndum curávit. Episcopatu ob graves habitualesque morbos dimisso, ad alumnos suos, a quibus pauper discesserat, revértitur pauper. Demum, quamvis senio laboribúsque, diuturna arthrítide aliisque gravíssimis morbis fractus corpore, spíritu tamen alacrior, de cæléstibus rebus disseréndi aut scribéndi finem numquam adhibuit, donec nonagenarius, Kaléndis Augusti, anno millesimo septingentésimo octogesimo septimo, Nucériæ Paganórum, inter suórum alumnórum lácrimas, placidíssime exspirávit. Eum, inde virtútibus et miraculis clarum, Pius septimus Pontifex maximus anno millesimo octingentésimo décimo sexto Beatórum fastis; novisque fulgentem signis, Gregórius décimus sextus in festo sanctíssimæ Trinitátis, anno millesimo octingentésimo trigesimo nono, solemni ritu, Sanctórum catálogo accensuit. Tandem Pius nonus Pontifex maximus, ex sacrórum Rituum Congregatiónis consulto, universalis Ecclésiæ Doctórem declarávit.\par
\switchcolumn
\EN
\noindent But firmly and perseveringly refused all high places in the Church which were offered him, but (in 1762) Pope Clement XIII. absolutely commanded him to take the Bishopric of the Church of Santa Agata de' Goti. On becoming a Bishop, the only change which he made in the hardness of his life was that of his outer raiment. There remained, too, the same simplicity of meats, the same strong zeal for Christian discipline, the same determined will to put down sin and keep out false doctrines, and the same earnestness in all the duties of a shepherd of souls. In his tenderness to the poor, he spent among them all the revenues of his Church, and in a year of famine sold the furniture of his own house to feed his starving people. He was all things to all men and brought nuns to lead a more perfect life, while he saw to it that a monastery was opened for nuns attached to his own Congregation. On account of grievous and continual sickness, he resigned his Bishopric, and poor as when he had left them, poor he returned among his disciples. On the 1st day of August, in the year 1787, he peacefully died at Noceradei - Pagani, amid the tears of his followers. He was then ninety years of age his body was worn out with old age and hard work, and with chronic gout, and other painful maladies, but the freshness of his mind never failed to the last, in talking and writing on heavenly things. In the year 1816 Pope Pius VII., finding him famous on account of his good works and miracles, enrolled his name among those of the Blessed. God still glorified him by new signs and wonders, and on the Feast of the Most Holy Trinity, in the year 1839, Gregory XVI., with solemn pomp, numbered him among the Saints of the Church. Lastly, Pope Pius IX., in accordance with a Resolution of the Congregation of Sacred Rites, gave him the title of Doctor of the Universal Church.\par
\switchcolumn*

\LA
\begin{Resp}\Rsign Iste est, qui ante Deum magnas virtútes operátus est, et omnis terra doctrína ejus repléta est: \Star{} Ipse intercédat pro peccátis ómnium populórum.\end{Resp}
\begin{Resp}\Vsign Iste est, qui contémpsit vitam mundi, et pervénit ad cæléstia regna.\end{Resp}
\begin{Resp}\Rsign Ipse intercédat pro peccátis ómnium populórum.\end{Resp}
\begin{Resp}\Vsign Glória Patri, et Fílio, \Star{} et Spirítui Sancto.\end{Resp}
\begin{Resp}\Rsign Ipse intercédat pro peccátis ómnium populórum.\end{Resp}
\switchcolumn
\EN
\begin{Resp}\Rsign This is he which wrought great wonders before God, and the whole earth is full of his teaching \Star{} May he pray for all people, that their sins may be forgiven unto them\end{Resp}
\begin{Resp}\Vsign This is he which loved not his life in this world, and hath attained unto the kingdom of heaven.\end{Resp}
\begin{Resp}\Rsign May he pray for all people, that their sins may be forgiven unto them\end{Resp}
\begin{Resp}\Vsign Glory be to the Father, and to the Son, \Star{} and to the Holy Ghost.\end{Resp}
\begin{Resp}\Rsign May he pray for all people, that their sins may be forgiven unto them\end{Resp}
\switchcolumn*

\LA
\LessonHead{Lectio VII}
\switchcolumn
\EN
\LessonHead{Lesson VII}
\switchcolumn*

\LA
\LessonTitle{Léctio sancti Evangélii secúndum Lucam}
\switchcolumn
\EN
\LessonTitle{From the Holy Gospel according to Luke}
\switchcolumn*

\LA
\Cite{Luc 10:1-9}
\switchcolumn
\EN
\Cite{Luke 10:1-9}
\switchcolumn*

\LA
\noindent In illo témpore: Designávit Dóminus et álios septuagínta duos: et misit illos binos ante fáciem suam, in omnem civitátem et locum, quo erat ipse ventúrus. Et réliqua.\par
\switchcolumn
\EN
\noindent At that time, the Lord appointed other seventy-two also, and sent them two and two before His face into every city and place, whither He Himself would come. And so on.\par
\switchcolumn*

\LA
\LessonTitle{Homilía sancti Gregórii Papæ}
\switchcolumn
\EN
\LessonTitle{Homily by Pope St. Gregory (the Great.)}
\switchcolumn*

\LA
\Source{Homilia 17 in Evangelia}
\switchcolumn
\EN
\Source{17th Homily on the Gospels}
\switchcolumn*

\LA
\noindent Dóminus et Salvátor noster, fratres caríssimi, aliquándo nos sermónibus, aliquándo vero opéribus ádmonet. Ipsa étenim facta ejus præcépta sunt: quia dum áliquid tácitus facit, quid ágere debeámus innotéscit. Ecce enim binos in prædicatiónem discípulos mittit: quia duo sunt præcépta caritátis, Dei vidélicet amor, et próximi: et minus quam inter duos cáritas habéri non potest. Nemo enim próprie ad semetípsum habére caritátem dícitur: sed diléctio in álterum tendit, ut cáritas esse possit.\par
\switchcolumn
\EN
\noindent Dearly beloved brethren, our Lord and Saviour doth sometimes admonish us by words, and sometimes by works. Yea, His very works do themselves teach us for that which He doth silently His example still moveth us to copy. Behold how He sendeth forth His disciples to preach by two and two since there are two commandments to love, that is, a commandment to love God, and a commandment to love our neighbour and where there are not two, the one, being alone, hath not whereon to do the Lord's commandment. And no man can properly be said to love himself: for love tendeth outward toward our neighbour, if it be the love whereto the Gospel doth oblige us.\par
\switchcolumn*

\LA
\begin{Resp}\Rsign Amávit eum Dóminus, et ornávit eum: stolam glóriæ índuit eum, \Star{} Et ad portas paradísi coronávit eum.\end{Resp}
\begin{Resp}\Vsign Induit eum Dóminus lorícam fídei, et ornávit eum.\end{Resp}
\begin{Resp}\Rsign Et ad portas paradísi coronávit eum.\end{Resp}
\switchcolumn
\EN
\begin{Resp}\Rsign The Lord loved him and beautified him He clothed him with a robe of glory \Star{} And crowned him at the gates of Paradise.\end{Resp}
\begin{Resp}\Vsign The Lord hath put on him the breast-plate of faith, and hath adorned him.\end{Resp}
\begin{Resp}\Rsign And crowned him at the gates of Paradise.\end{Resp}
\switchcolumn*

\LA
\LessonHead{Lectio VIII}
\switchcolumn
\EN
\LessonHead{Lesson VIII}
\switchcolumn*

\LA
\noindent Ecce enim binos ad prædicándum discípulos Dóminus mittit: quátenus hoc nobis tácitus ínnuat, quia qui caritátem erga álterum non habet, prædicatiónis offícium suscípere nullátenus debet. Bene autem dícitur, quia misit eos ante fáciem suam in omnem civitátem et locum, quo erat ipse ventúrus. Prædicatóres enim suos Dóminus séquitur: quia prædicátio prǽvenit, et tunc ad mentis nostræ habitáculum Dóminus venit, quando verba exhortatiónis præcúrrunt: atque per hoc véritas in mente suscípitur.\par
\switchcolumn
\EN
\noindent Behold, the Lord sendeth forth His disciples to preach by two and two and thus doing, He doth silently teach us that whosoever loveth not his neighbour, such a one it behoveth not to take upon him the office of a preacher. Well also is it said that He sent them before His face into every city and place whither He Himself would come. The Lord followeth His preachers first cometh preaching, and then the Lord Himself cometh to the house of our mind, whither the word of exhortation hath come before and so cometh the truth into our mind.\par
\switchcolumn*

\LA
\begin{Resp}\Rsign In médio Ecclésiæ apéruit os ejus, \Star{} Et implévit eum Dóminus spíritu sapiéntiæ et intelléctus.\end{Resp}
\begin{Resp}\Vsign Jucunditátem et exsultatiónem thesaurizávit super eum.\end{Resp}
\begin{Resp}\Rsign Et implévit eum Dóminus spíritu sapiéntiæ et intelléctus.\end{Resp}
\begin{Resp}\Vsign Glória Patri, et Fílio, \Star{} et Spirítui Sancto.\end{Resp}
\begin{Resp}\Rsign Et implévit eum Dóminus spíritu sapiéntiæ et intelléctus.\end{Resp}
\switchcolumn
\EN
\begin{Resp}\Rsign In the midst of the congregation did the Lord open his mouth. \Star{} And filled him with the spirit of wisdom and understanding.\end{Resp}
\begin{Resp}\Vsign He made him rich with joy and gladness.\end{Resp}
\begin{Resp}\Rsign And filled him with the spirit of wisdom and understanding.\end{Resp}
\begin{Resp}\Vsign Glory be to the Father, and to the Son, \Star{} and to the Holy Ghost.\end{Resp}
\begin{Resp}\Rsign And filled him with the spirit of wisdom and understanding.\end{Resp}
\switchcolumn*

\end{paracol}

\par\needspace{10\baselineskip}\addvspace{1.2em}{\centering\small\textsc{Die 2 Augusti} · \textsc{2 August}\par}
\Office{S. Stephani Papæ Martyris}{S. Stephen Pope and Martyr}{Simplex}
\addcontentsline{toc}{subsection}{Die 2 Augusti · S. Stephani Papæ Martyris \textit{— S. Stephen Pope and Martyr}}
\begin{paracol}{2}
\LA
\LessonHead{Lectio IX (pro commemoratione)}
\switchcolumn
\EN
\LessonHead{Lesson IX (when commemorated)}
\switchcolumn*

\LA
\LessonTitle{Commemoratio: S. Stephani Papæ Martyris}
\switchcolumn
\EN
\LessonTitle{Commemoration of: S. Stephen Pope and Martyr}
\switchcolumn*

\LA
\noindent Stephanus Romanus pontificátum gessit Valeriáno et Gallieno imperatóribus. Instituit ut sacerdótes et diáconi nusquam sacris vestibus, nisi in ecclésia, uteréntur. Baptizatos ab hæreticis, íterum baptizari vetuit, rescribens sancto Cypriáno verbis illis: Nihil innovétur, nisi quod traditum est. Multos étiam ad Christum convértit, in iis Olympium tribunum cum uxore Exsuperia et fílio Theodúlo, item Nemesium tribunum cum tota família, óculis Lucillæ ejus fíliæ restitútis: qui omnes Jesu Christi Mártyres fuérunt. Sed, vehementius jam ingravescénte persecutióne imperatórum, Stephanus, convocato clero, ad martyrium suos hortabátur, in cryptisque Mártyrum assidue Missas et concília celebrábat. Qui, cum aliquándo raptus esset ab infidelibus ad templum Martis ut ei sacrificaret, líbere negávit se dæmónibus eum honórem velle tribúere, qui uni Deo deberétur. Quibus in verbis Martis statua terræmotu cóncidit templumque contrémuit. Quare ómnibus aufugiéntibus qui Stephanum tenébant, Pontifex ad suos revértitur in cœmetérium Lucinæ; quos, divinis præceptis ínstruens, sacraménto corporis Christi communicávit; ibique, dum Missárum solemnia pérficit; adveniéntibus íterum imperatórum satellítibus, ei in sua sede caput abscínditur. Corpus, cum eadem sede conspersa Mártyris sánguine, a clericis sepúltum est in cœmeterio Callísti, quarto Nonas Augusti. Vixit in pontificatu annos tres, menses tres, dies viginti duos. Habuit ordinatiónes duas mense Decembri, quibus creávit presbyteros sex, diaconos quinque, episcopos tres.\par
\switchcolumn
\EN
\noindent This Stephen was a Roman, and exercised the Popedom in the reign of the Emperors Valerian and Gallienus. It was his ordinance which forbade Priests and Deacons ever to use their hallowed garments except in the Church. He forbade a re-baptism of such as had been baptized by heretics, writing to St. Cyprian in these words: Let us have no innovations, but only what hath been handed down unto us. He turned many to Christ, and, among them, the Tribune Olympius, with his wife Exuperia, and his son Theodulus, and the Tribune Nemesius, to whose blind daughter he had given sight, along with all his household. All these were martyrs for Jesus Christ. When the persecution of the Emperors was waxing dreader and more dread, Stephen gathered together the clergy, and exhorted them to be brave in lifting up their testimony, and himself celebrated Masses and Councils in the Catacombs. He was caught by some unbelievers, and haled to the temple of Mars, to do sacrifice to that idol, but he boldly said he would never pay to devils an honour which it behoved to give to God only. As he spake these words an earthquake made the image of Mars to fall down, and all the temple to tremble. All they that held Stephen fled, and the Pope went back to his own people in the cemetery of Lucina. He there delivered to them a discourse full of the Word of God, and gave them the Communion of the Sacrament of the Body of Christ. While he was finishing the Mass, the soldiers of the Emperor again brake in upon them, and his head was cut off as he sat in his chair. The reliques of the Martyr, along with the chair stained with his blood, were buried by the clergy in the cemetery of Callistus, upon the 2nd day of August, (in the year of our Lord 257.) He lived as Pope three years, three months, and twenty two days. He held two ordinations in the month of December, and in them ordained six Priests, five Deacons, and three Bishops.\par
\switchcolumn*

\LA
\TeDeum{Te Deum}
\switchcolumn
\EN
\TeDeum{Te Deum}
\switchcolumn*

\end{paracol}

\par\needspace{10\baselineskip}\addvspace{1.2em}{\centering\small\textsc{Die 3 Augusti} · \textsc{3 August}\par}
\Office{De Inventione S. Stephani Protomartyris}{Finding of St. Stephen the First Martyr}{Semiduplex}
\addcontentsline{toc}{subsection}{Die 3 Augusti · De Inventione S. Stephani Protomartyris \textit{— Finding of St. Stephen the First Martyr}}
\begin{paracol}{2}
\LA
\LessonHead{Lectio I}
\switchcolumn
\EN
\LessonHead{Lesson I}
\switchcolumn*

\LA
\LessonTitle{De Actibus Apostolorum}
\switchcolumn
\EN
\LessonTitle{From the Acts of the Apostles}
\switchcolumn*

\LA
\Cite{Acts 7:51-54}
\switchcolumn
\EN
\Cite{Acts 7:51-54}
\switchcolumn*

\LA
\noindent Dura cervice et incircumcisis córdibus et áuribus, vos semper Spiritui Sancto resístitis: sicut patres vestri, ita et vos. Quem prophetárum non sunt persecúti patres vestri? Et occidérunt eos, qui prænuntiábunt de adventu Justi, cujus vos nunc proditores et homicidæ fuistis, qui accepístis legem in dispositióne Angelórum et non custodistis. Audiéntes autem hæc, disseccabántur córdibus suis, et stridébant déntibus in eum.\par
\switchcolumn
\EN
\noindent You stiffnecked and uncircumcised in heart and ears, you always resist the Holy Ghost: as your fathers did, so do you also. Which of the prophets have not your fathers persecuted? And they have slain them who foretold of the coming of the Just One; of whom you have been now the betrayers and murderers: Who have received the law by the disposition of angels, and have not kept it. Now hearing these things, they were cut to the heart, and they gnashed with their teeth at him.\par
\switchcolumn*

\LA
\begin{Resp}\Rsign Stéphanus autem plenus grátia et fortitúdine, \Star{} Faciébat prodígia et signa magna in pópulo.\end{Resp}
\begin{Resp}\Vsign Surrexérunt quidam de synagóga disputántes cum Stéphano: et non póterant resístere sapiéntiæ, et Spirítui qui loquebátur.\end{Resp}
\begin{Resp}\Rsign Faciébat prodígia et signa magna in pópulo.\end{Resp}
\switchcolumn
\EN
\begin{Resp}\Rsign And Stephen, full of grace and power \Star{} Did great wonders and miracles among the people.\end{Resp}
\begin{Resp}\Vsign There arose certain of the synagogue, disputing with Stephen; and they were not able to resist the wisdom, and the Spirit which spoke.\end{Resp}
\begin{Resp}\Rsign Did great wonders and miracles among the people.\end{Resp}
\switchcolumn*

\LA
\LessonHead{Lectio II}
\switchcolumn
\EN
\LessonHead{Lesson II}
\switchcolumn*

\LA
\Cite{Acts 7:55-58}
\switchcolumn
\EN
\Cite{Acts 7:55-58}
\switchcolumn*

\LA
\noindent Cum autem esset plenus Spíritu Sancto, inténdens in cælum, vidit glóriam Dei, et Jesum stantem a dextris Dei. Et ait: Ecce video cælos apertos, et Fílium hóminis stantem a dextris Dei. Exclamántes autem voce magna, continuérunt aures suas, et ímpetum fecérunt unanímiter in eum. Et ejiciéntes eum extra civitátem lapidábant: et testes deposuérunt vestiménta sua secus pedes adolescentis, qui vocabátur Saulus. Et lapidábant Stephanum, invocántem et dicéntem: Dómine Jesu, suscipe spíritum meum.\par
\switchcolumn
\EN
\noindent But he, being full of the Holy Ghost, looking up steadfastly to heaven, saw the glory of God, and Jesus standing on the right hand of God. And he said: Behold, I see the heavens opened, and the Son of man standing on the right hand of God. And they crying out with a loud voice, stopped their ears, and with one accord ran violently upon him. And casting him forth without the city, they stoned him; and the witnesses laid down their garments at the feet of a young man, whose name was Saul. And they stoned Stephen, invoking, and saying: Lord Jesus, receive my spirit.\par
\switchcolumn*

\LA
\begin{Resp}\Rsign Vidébant omnes Stéphanum, qui erant in concílio: \Star{} Et intuebántur vultum ejus tamquam vultum Angeli stantis inter illos.\end{Resp}
\begin{Resp}\Vsign Plenus grátia et fortitúdine, faciébat prodígia et signa magna in pópulo.\end{Resp}
\begin{Resp}\Rsign Et intuebántur vultum ejus tamquam vultum Angeli stantis inter illos.\end{Resp}
\switchcolumn
\EN
\begin{Resp}\Rsign All that sat in the council, looking steadfastly on Stephen, saw \Star{} His face as it had been the face of an angel standing among them.\end{Resp}
\begin{Resp}\Vsign Full of grace and power, he did great wonders and miracles among the people.\end{Resp}
\begin{Resp}\Rsign His face as it had been the face of an angel standing among them.\end{Resp}
\switchcolumn*

\LA
\LessonHead{Lectio III}
\switchcolumn
\EN
\LessonHead{Lesson III}
\switchcolumn*

\LA
\Cite{Acts 7:59; 8:1-2}
\switchcolumn
\EN
\Cite{Acts 7:59; 8:1-2}
\switchcolumn*

\LA
\noindent Positis autem genibus, clamávit voce magna: Dómine, ne statuas illis hoc peccátum. Et cum hoc dixísset, obdormívit in Dómino. Saulus autem erat consentiens neci ejus. Facta est autem in illa die persecútio magna in ecclésia quæ erat Jerosolymis; et omnes dispersi sunt per regiónes Judææ et Samaríæ, præter Apóstolos. Curavérunt autem Stephanum viri timorati, et fecérunt planctum magnum super eum.\par
\switchcolumn
\EN
\noindent And falling on his knees, he cried with a loud voice, saying: Lord, lay not this sin to their charge. And when he had said this, he fell asleep in the Lord. And Saul was consenting to his death. And at that time there was raised a great persecution against the church which was at Jerusalem; and they were all dispersed through the countries of Judea, and Samaria, except the apostles. And devout men took order for Stephen's funeral, and made great mourning over him.\par
\switchcolumn*

\LA
\begin{Resp}\Rsign Intuens in cælum beátus Stéphanus, vidit glóriam Dei, et ait: \Star{} Ecce vídeo cælos apértos, et Fílium hóminis stantem a dextris virtútis Dei.\end{Resp}
\begin{Resp}\Vsign Cum autem esset Stéphanus plenus Spíritu Sancto, inténdens in cælum, vidit glóriam Dei, et ait.\end{Resp}
\begin{Resp}\Rsign Ecce vídeo cælos apértos, et Fílium hóminis stantem a dextris virtútis Dei.\end{Resp}
\begin{Resp}\Vsign Glória Patri, et Fílio, \Star{} et Spirítui Sancto.\end{Resp}
\begin{Resp}\Rsign Ecce vídeo cælos apértos, et Fílium hóminis stantem a dextris virtútis Dei.\end{Resp}
\switchcolumn
\EN
\begin{Resp}\Rsign The blessed Stephen looked up steadfastly into heaven, and saw the glory of God, and said \Star{} Behold, I see the heavens opened, and the Son of Man standing at the right hand of the power of God.\end{Resp}
\begin{Resp}\Vsign But Stephen, being full of the Holy Ghost, looked up steadfastly into heaven, and saw the glory of God, and said\end{Resp}
\begin{Resp}\Rsign Behold, I see the heavens opened, and the Son of man standing at the right hand of the power of God.\end{Resp}
\begin{Resp}\Vsign Glory be to the Father, and to the Son, \Star{} and to the Holy Ghost.\end{Resp}
\begin{Resp}\Rsign Behold, I see the heavens opened, and the Son of Man standing at the right hand of the power of God.\end{Resp}
\switchcolumn*

\LA
\LessonHead{Lectio IV}
\switchcolumn
\EN
\LessonHead{Lesson IV}
\switchcolumn*

\LA
\noindent Sanctórum córpora Stephani Protomártyris, Gamaliélis, Nicodémi et Abibónis, quæ diu in obscúro ac sórdido loco jacúerant, Honorio imperatóre, Luciáno presbytero divínitus admónito, invénta sunt prope Jerosolymam. Cui Gamaliel cum in somnis apparuísset, gravi quadam et præclára senis spécie, locum jacéntium córporum commonstrávit, imperans ut Joánnem Jerosolymitanum antístitem adíret, agerétque cum eo ut honestius illa córpora sepelirentur.\par
\switchcolumn
\EN
\noindent In (the year of our Lord 415,) in the reign of the Emperor Honorius, a Priest named Lucian, (dwelling at Caphargamala, about twenty miles from Jerusalem,) received a message from God, in consequence of which discovery was made of the bodies of the Saints Stephen the First Martyr, Gamaliel, Nicodemus, and Abibon, (the son of Gamaliel,) which had long been lying unknown and unheeded. Lucian was asleep when Gamaliel appeared to him in a dream as a tall comely old man of worshipful presence, told him where the bodies were lying, and bade him go to John, Patriarch of Jerusalem, and deal with him that they might have more honourable burial.\par
\switchcolumn*

\LA
\begin{Resp}\Rsign Lapidábant Stéphanum invocántem et dicéntem: \Star{} Dómine Jesu Christe, áccipe spíritum meum: et ne státuas illis hoc peccátum.\end{Resp}
\begin{Resp}\Vsign Pósitis autem génibus, clamávit voce magna, dicens.\end{Resp}
\begin{Resp}\Rsign Dómine Jesu Christe, áccipe spíritum meum: et ne státuas illis hoc peccátum.\end{Resp}
\switchcolumn
\EN
\begin{Resp}\Rsign They stoned Stephen, calling upon God and saying: \Star{} Lord Jesus Christ, receive my spirit; and lay not this sin to their charge.\end{Resp}
\begin{Resp}\Vsign And he kneeled down, and cried with a loud voice, saying\end{Resp}
\begin{Resp}\Rsign Lord Jesus Christ, receive my spirit; and lay not this sin to their charge.\end{Resp}
\switchcolumn*

\LA
\LessonHead{Lectio V}
\switchcolumn
\EN
\LessonHead{Lesson V}
\switchcolumn*

\LA
\noindent Quibus auditis Jerosolymórum antístes, finitimárum urbium epíscopis presbyterisque convocatis, ad locum pergit; defossos lóculos ínvenit, unde suavíssimus odor efflabátur. Cujus rei fama commóta, magna hóminum multitúdo eo convénit, multique ex variis morbis ægroti ac débiles sani et íntegri domum rediérunt. Sacrum autem sancti Stephani corpus, quod summa tunc celebritate in sanctam ecclésiam Sion illátum est, sub Theodosio junióre Constantinópolim, inde Romam, Pelagio primo summo Pontifice, translátum, in agro Verano in sepúlcro sancti Lauréntii Mártyris collocátum est.\par
\switchcolumn
\EN
\noindent When the Patriarch of Jerusalem heard it, he called together Bishops and Priests from the neighbouring cities, and betook himself to the place, where he found the tombs hewn in the rock, and a right sweet savour flowing forth from them. The thing being noised abroad, a great multitude of people came together, and many that were sick and weak of diverse diseases returned home whole. The sacred body of holy Stephen was then carried with great pomp to the holy Church of Zion. Under the Emperor Theodosius the Younger it was taken to Constantinople and during the Popedom of Pelagius I. it was brought to Rome, where it has been laid in the sepulchre of the holy Martyr Lawrence in the Veranian Field.\par
\switchcolumn*

\LA
\begin{Resp}\Rsign Impetum fecérunt unanímiter in eum, et ejecérunt eum extra civitátem, invocántem et dicéntem: \Star{} Dómine Jesu, áccipe spíritum meum.\end{Resp}
\begin{Resp}\Vsign Et testes deposuérunt vestiménta sua secus pedes adolescéntis, qui vocabátur Saulus: et lapidábant Stéphanum invocántem et dicéntem.\end{Resp}
\begin{Resp}\Rsign Dómine Jesu, áccipe spíritum meum.\end{Resp}
\switchcolumn
\EN
\begin{Resp}\Rsign They ran upon him with one accord, and cast him out of the city, calling upon God, and saying: \Star{} Lord Jesus, receive my spirit.\end{Resp}
\begin{Resp}\Vsign And the witnesses laid down their clothes at a young man's feet, whose name was Saul; and they stoned Stephen, calling upon God, and saying\end{Resp}
\begin{Resp}\Rsign Lord Jesus, receive my spirit.\end{Resp}
\switchcolumn*

\LA
\LessonHead{Lectio VI}
\switchcolumn
\EN
\LessonHead{Lesson VI}
\switchcolumn*

\LA
\LessonTitle{Ex libro sancti Augustíni Epíscopi de Civitáte Dei}
\switchcolumn
\EN
\LessonTitle{From the Book upon The City of God, written by St. Augustine, Bishop (of Hippo)}
\switchcolumn*

\LA
\Source{Liber 22, cap. 8, circa medium}
\switchcolumn
\EN
\Source{xxii, 8}
\switchcolumn*

\LA
\noindent Ad aquas Tibilitanas, episcopo afferénte Projecto relíquias Mártyris gloriosíssimi Stephani, ad ejus memóriam veniebat magnæ multitúdinis concursus et occúrsus. Ibi cæca mulier, ut ad episcopum portántem pígnora sacra ducerétur, orávit: flores, quos ferebat, dedit; recepit, óculis admóvit, prótinus vidit. Stupéntibus qui áderant, præibat exsultans, viam carpens et viæ ducem ulterius non requírens. Memoráti memóriam Mártyris, quæ pósita est in castello Synicénsi, quod Hipponénsi coloniæ vicinum est, ejusdem loci Lucillus epíscopus, pópulo præcedénte atque sequente, portábat; fístula, cujus molestia jamdiu laboraverat, et familiaríssimi sui médici, qui eam secaret, opperiebátur manus, illíus piæ sarcinæ vectatióne repénte sanáta est.\par
\switchcolumn
\EN
\noindent When the Bishop Projectus brought some reliques of that most glorious Martyr Stephen to Tibilis, a great multitude came together and went out to meet the shrine. A blind woman prayed to be led to the Bishop who was bearing the hallowed deposit. She laid on the reliques the flowers which she was carrying, took them up again, touched her eyes with them and forthwith saw. She went forward rejoicing, at the head of the amazed procession, choosing her own path, and needing no more that any should lead her. I remember also the shrine of this same Martyr which hath been placed in the town of Synica, hard by this city of Hippo. Lucillus, Bishop of that place, was carrying it, with a multitude going before and following after when, all of a sudden, by bearing this hallowed burden, he was healed of the emerods, from which he was even then suffering, and which were being treated by a physician, an intimate friend of his, who was about to cut them.\par
\switchcolumn*

\LA
\begin{Resp}\Rsign Impii super justum jactúram fecérunt, ut eum morti tráderent: \Star{} At ille gaudens suscépit lápides, ut mererétur accípere corónam glóriæ.\end{Resp}
\begin{Resp}\Vsign Continuérunt aures suas, et ímpetum fecérunt unanímiter in eum.\end{Resp}
\begin{Resp}\Rsign At ille gaudens suscépit lápides, ut mererétur accípere corónam glóriæ.\end{Resp}
\begin{Resp}\Vsign Glória Patri, et Fílio, \Star{} et Spirítui Sancto.\end{Resp}
\begin{Resp}\Rsign At ille gaudens suscépit lápides, ut mererétur accípere corónam glóriæ.\end{Resp}
\switchcolumn
\EN
\begin{Resp}\Rsign The ungodly fell upon the righteous, to put him to death. \Star{} But he received the stones with joy, that he might earn a crown of glory.\end{Resp}
\begin{Resp}\Vsign They stopped their ears, and ran upon him with one accord.\end{Resp}
\begin{Resp}\Rsign But he received the stones with joy, that he might earn a crown of glory.\end{Resp}
\begin{Resp}\Vsign Glory be to the Father, and to the Son, \Star{} and to the Holy Ghost.\end{Resp}
\begin{Resp}\Rsign But he received the stones with joy, that he might earn a crown of glory.\end{Resp}
\switchcolumn*

\LA
\LessonHead{Lectio VII}
\switchcolumn
\EN
\LessonHead{Lesson VII}
\switchcolumn*

\LA
\LessonTitle{Léctio sancti Evangélii secúndum Matthǽum}
\switchcolumn
\EN
\LessonTitle{From the Holy Gospel according to Matthew}
\switchcolumn*

\LA
\Cite{Matt 23:34-39}
\switchcolumn
\EN
\Cite{Matt 23:34-39}
\switchcolumn*

\LA
\noindent In illo témpore: Dicébat Jesus scribis et pharisǽis: Ecce ego mitto ad vos prophétas, et sapiéntes, et scribas: et ex illis occidétis, et crucifigétis. Et réliqua.\par
\switchcolumn
\EN
\noindent At that time, Jesus said to the scribes and pharisee: I send to you prophets, and wise men, and scribes: and some of them you will put to death and crucify. And so on,\par
\switchcolumn*

\LA
\LessonTitle{Homilía sancti Hierónymi Presbýteri}
\switchcolumn
\EN
\LessonTitle{Homily by St. Jerome, Priest (at Bethlehem.)}
\switchcolumn*

\LA
\Source{Liber 4 Comment. in cap. 23 Matthæi}
\switchcolumn
\EN
\Source{Book IV, Commentary on Matth. XXIII}
\switchcolumn*

\LA
\noindent Hoc quod ántea dixerámus, Impléte mensúram patrum vestrórum; ad persónam Dómini pertinére, eo quod occidéndus esset ab eis, potest et ad discípulos ejus reférri, de quibus nunc dicit: Ecce ego mitto ad vos prophétas, et sapiéntes, et scribas. Simúlque obsérva, juxta Apóstolum scribéntem ad Corínthios, vária esse dona discipulórum Christi; álios prophétas, qui ventúra prædícent; álios sapiéntes, qui novérunt quando débeant proférre sermónem; álios scribas in lege doctíssimos, ex quibus lapidátus est Stéphanus, Paulus occísus, crucifíxus Petrus, flagelláti in Actibus Apostolórum discípuli.\par
\switchcolumn
\EN
\noindent We have already remarked that the Lord's words, Fill ye up the measure of your fathers, (32,) refer in the first place to Himself, Whom the Jews afterwards put to death. In a secondary sense it may likewise be applied to His disciples, of whom He saith, Behold, I send unto you Prophets, and wise men, and Scribes. Here observe that, according to the Apostle writing to the Corinthians, (1 Cor. xii. 4,) there are diversities of gifts among Christ's followers. Some are Prophets of that which is to come; some are wise men, who know the due season for rebuke and exhortation; some are Scribes learned in the law. And of these they stoned Stephen, slew Paul with the sword, crucified Peter, and scourged the Disciples mentioned in the Acts of the Apostles. (v, 40; xvi, 23.)\par
\switchcolumn*

\LA
\begin{Resp}\Rsign Stéphanus servus Dei, quem lapidábant Judǽi, vidit cælos apértos: vidit, et introívit: \Star{} Beátus homo, cui cæli patébant.\end{Resp}
\begin{Resp}\Vsign Cum ígitur saxórum crepitántium túrbine quaterétur, inter æthéreos aulæ cæléstis sinus divína ei cláritas fulsit.\end{Resp}
\begin{Resp}\Rsign Beátus homo, cui cæli patébant.\end{Resp}
\switchcolumn
\EN
\begin{Resp}\Rsign Stephen, the servant of God, who was stoned by the Jews, saw the heavens opened he saw and entered in. \Star{} Blessed is he, unto whom the heavens were opened.\end{Resp}
\begin{Resp}\Vsign While his poor body was crushed by the hurtling shower of stones, God's brightness broke upon him out of the heavenly palaces.\end{Resp}
\begin{Resp}\Rsign Blessed is he unto whom the heavens were opened.\end{Resp}
\switchcolumn*

\LA
\LessonHead{Lectio VIII}
\switchcolumn
\EN
\LessonHead{Lesson VIII}
\switchcolumn*

\LA
\noindent Quǽrimus, quis iste sit Zacharías fílius Barachíæ: quia multos légimus Zacharías. Et ne líbera nobis tribuerétur erróris facúltas, ádditum est: Quem occidístis inter templum et altáre. In divérsis divérsa legi: et débeo singulórum opiniónes pónere. Alii Zacharíam fílium Barachíæ dicunt, qui in duódecim prophétis undécimus est, patrísque in eo nomen conséntit: sed ubi occísus sit inter templum et altáre, Scriptúra non lóquitur; máxime cum tempóribus ejus vix ruínæ templi fúerint. Alii Zacharíam patrem Joánnis intellégi volunt, ex quibúsdam apocryphórum sómniis approbántes, quod proptérea occísus sit, quia Salvatóris prædicáverit advéntum.\par
\switchcolumn
\EN
\noindent It is a subject of dispute among commentators who is meant by Zacharias the son of Barachias. We read of several persons of the name of Zacharias. But here, in order to prevent any mistake, it is particularly said, Whom ye slew between the temple and the altar. I have read various opinions in various places upon this question, and I will give each. First, some hold that Zacharias the son of Barachias is the eleventh of the twelve Minor Prophets; and this opinion is supported by the father's name. But the Bible nowhere telleth us that this Prophet was slain between the temple and the altar; and it is hardly possible that he can have been, for in his time it could scarcely be said that even the ruins of the temple were in existence. Secondly, others maintain that this Zacharias was Zacharias, the father of John the Baptist. This interpretation is derived from the dreams of the Apocryphal Gospels, wherein it is asserted that he was martyred for preaching Christ's coming.\par
\switchcolumn*

\LA
\begin{Resp}\Rsign Patefáctæ sunt jánuæ cæli Christi Mártyri beáto Stéphano, qui in número Mártyrum invéntus est primus: \Star{} Et ídeo triúmphat in cælis coronátus.\end{Resp}
\begin{Resp}\Vsign Mortem enim, quam Salvátor noster dignátus est pro nobis pati, hanc ille primus réddidit Salvatóri.\end{Resp}
\begin{Resp}\Rsign Et ídeo triúmphat in cælis coronátus.\end{Resp}
\begin{Resp}\Vsign Glória Patri, et Fílio, \Star{} et Spirítui Sancto.\end{Resp}
\begin{Resp}\Rsign Et ídeo triúmphat in cælis coronátus.\end{Resp}
\switchcolumn
\EN
\begin{Resp}\Rsign The gates of heaven were opened to Christ's blessed martyr Stephen, and he is the first of all the martyrs. \Star{} Wherefore he reigneth crowned in heaven.\end{Resp}
\begin{Resp}\Vsign For he was the first to make an offering of his death to that Saviour Who vouchsafed to suffer death for us.\end{Resp}
\begin{Resp}\Rsign Wherefore he reigneth crowned in heaven.\end{Resp}
\begin{Resp}\Vsign Glory be to the Father, and to the Son, \Star{} and to the Holy Ghost.\end{Resp}
\begin{Resp}\Rsign Wherefore he reigneth crowned in heaven.\end{Resp}
\switchcolumn*

\LA
\LessonHead{Lectio IX}
\switchcolumn
\EN
\LessonHead{Lesson IX}
\switchcolumn*

\LA
\noindent Alii istum volunt esse Zacharíam, qui occísus est a Joas rege Judæ inter templum et altáre, sicut Regum narrat história. Sed observándum, quod ille Zacharías non sit fílius Barachíæ, sed fílius Jójadæ sacerdótis: unde et Scriptúra refert: Non fuit recordátus Joas patris ejus Jójadæ, quia sibi fecísset bona. Cum ergo et Zacharíam teneámus et occisiónis conséntiat locus: quǽrimus, quare Barachíæ dicátur fílius, et non Jójadæ? Barachía lingua nostra Benedíctus Dómini dícitur, et sacerdótis Jójadæ justítia Hebrǽo nómine demonstrátur. In Evangélio, quo utúntur Nazaréni, pro fílio Barachíæ, fílium Jójadæ reperímus scriptum.\par
\switchcolumn
\EN
\noindent A third school will have it that this Zacharias, the son of Barachias, was that Zacharias of whom we read, in Chron. xxiv. 22, that he was slain by Joash, king of Judah, between the temple and the altar. Against this it is to be remarked, that that Zacharias was not the son of Barachias, but of Jehoiada the priest; whence it is written, Joash remembered not the kindness which Jehoiada his father had done to him. The question therefore ariseth, if this opinion be true, why, the name and manner of death both agreeing with this explanation, Zacharias is called the son, not of Jehoiada, but of Barachias. In Hebrew, Barachias signifieth the Blessed of the Lord, and Jehoiada proves his Righteousness. In the Gospel used by the Nazarenes the name of Jehoiada is used instead of Barachias.\par
\switchcolumn*

\LA
\TeDeum{Te Deum}
\switchcolumn
\EN
\TeDeum{Te Deum}
\switchcolumn*

\end{paracol}

\par\needspace{10\baselineskip}\addvspace{1.2em}{\centering\small\textsc{Die 4 Augusti} · \textsc{4 August}\par}
\Office{S. Dominici Confessoris}{St. Dominic the Confessor}{Duplex majus}
\addcontentsline{toc}{subsection}{Die 4 Augusti · S. Dominici Confessoris \textit{— St. Dominic the Confessor}}
\begin{paracol}{2}
\LA
\RefLine{Lectiones I–III: de Scriptura occurrente.}
\switchcolumn
\EN
\RefLine{Lessons I–III: from the Scripture of the day.}
\switchcolumn*

\LA
\RefLine{Lectiones VII–IX: ex Commúni Confessoris non pontificis (C5).}
\switchcolumn
\EN
\RefLine{Lessons VII–IX: from the Common of a Confessor not a Bishop (C5).}
\switchcolumn*

\LA
\LessonHead{Lectio IV}
\switchcolumn
\EN
\LessonHead{Lesson IV}
\switchcolumn*

\LA
\noindent Dominicus, Calarógæ in Hispánia ex nobili Gusmanórum família natus, Paléntiæ liberálibus disciplinis et theologíæ operam dedit; quo in studio cum plurimum profecísset, prius Oxoménsis ecclésiæ canonicus regularis, deínde ordinis fratrum Prædicatórum auctor fuit. Hujus mater grávida sibi visa est, in quiete, continére in alvo cátulum ore præferentem facem, qua, éditus in lucem, orbem terrárum incenderet. Quo somnio significabátur fore ut, splendore sanctitátis ac doctrinæ, gentes ad christianam pietátem inflammaréntur. Veritátem éxitus comprobávit; id enim et præstitit per se, et per sui ordinis socios deinceps est consecutus.\par
\switchcolumn
\EN
\noindent Dominic was a Spaniard, a son of the noble family of Guzman. He was born at Calaruega, (in Old Castile,) in the year of our Lord 1170, and was sent to be taught worldly and sacred learning at Palencia. He attained great success in his studies, and became, first, a Regular Canon of the Cathedral Church of Osma, and afterwards Founder of the Order of Friars Preachers. While his mother was great with child, she dreamt that what was in her womb was a little whelp dog with a lighted torch in his mouth, and that when she was delivered of him, he set all the world on fire. By this dream was figured that burning and shining light of holiness of life and power of doctrine, whereby he should enkindle godliness throughout whole nations. The end proved the truth of the image, for this is in good sooth what his work was, a work which he hath continued to do through the Brethren of his Order.\par
\switchcolumn*

\LA
\begin{Resp}\Rsign Honéstum fecit illum Dóminus, et custodívit eum ab inimícis, et a seductóribus tutávit illum: \Star{} Et dedit illi claritátem ætérnam.\end{Resp}
\begin{Resp}\Vsign Justum dedúxit Dóminus per vias rectas, et osténdit illi regnum Dei.\end{Resp}
\begin{Resp}\Rsign Et dedit illi claritátem ætérnam.\end{Resp}
\switchcolumn
\EN
\begin{Resp}\Rsign The Lord made him honourable, and defended him from his enemies, and kept him safe from those that lay in wait for him; \Star{} And gave him perpetual glory.\end{Resp}
\begin{Resp}\Vsign He went down with him into the pit, and left him not in bonds.\end{Resp}
\begin{Resp}\Rsign And gave him perpetual glory.\end{Resp}
\switchcolumn*

\LA
\LessonHead{Lectio V}
\switchcolumn
\EN
\LessonHead{Lesson V}
\switchcolumn*

\LA
\noindent Hujus autem ingenium ac virtus maxime enituit in everténdis hæreticis, qui perniciosis erróribus Tolosátes pervertere conabántur; quo in negótio septem consumpsit annos. Postea Romam venit ad Lateranense concílium cum episcopo Tolosano, ut ordo, quem institúerat, ab Innocentio tertio confirmarétur. Quæ res dum in deliberatióne versatur. Dominicus hortatu Pontificis ad suos revértitur, ut sibi regulam delígeret. Romam rédiens, ab Honorio tertio, qui próximus Innocentio successerat, confirmatiónem ordinis Prædicatórum ímpetrat. Romæ autem duo instituit monasteria, álterum virórum, mulíerum álterum. Tres étiam mórtuos ad vitam revocávit, múltaque alia edidit miracula, quibus ordo Prædicatórum mirifice propagari cœpit.\par
\switchcolumn
\EN
\noindent But his wisdom and bravery were most chiefly shown in confounding the \textasciitilde{}(Albigensian) heretics, who were fain, with their pestilential falsehoods, to corrupt the people of the county of Toulouse. In this business he spent seven years. Afterwards he came to Rome to the Fourth Lateran Council, along with the Bishop of Toulouse, that the Order which he had founded might be confirmed by Innocent III. While the matter was under consideration, Dominick, by the advice of the Pope, went home, that he might put his Rule into shape. He returned again to Rome (in 1216,) and obtained from Honorius III., the immediate successor of Innocent, the confirmation of the Order of Preachers. At Rome Dominic founded two Convents, one for men, and one for women. He raised three dead men to life, and worked many other miracles, where through the Order of Preachers became exceedingly spread abroad.\par
\switchcolumn*

\LA
\begin{Resp}\Rsign Amávit eum Dóminus, et ornávit eum: stolam glóriæ índuit eum, \Star{} Et ad portas paradísi coronávit eum.\end{Resp}
\begin{Resp}\Vsign Índuit eum Dóminus lorícam fídei, et ornávit eum.\end{Resp}
\begin{Resp}\Rsign Et ad portas paradísi coronávit eum.\end{Resp}
\switchcolumn
\EN
\begin{Resp}\Rsign The Lord loved him and beautified him; He clothed him with a robe of glory, \Star{} And crowned him at the gates of Paradise.\end{Resp}
\begin{Resp}\Vsign The Lord hath put on him the breast-plate of faith, and hath adorned him.\end{Resp}
\begin{Resp}\Rsign And crowned him at the gates of Paradise.\end{Resp}
\switchcolumn*

\LA
\LessonHead{Lectio VI}
\switchcolumn
\EN
\LessonHead{Lesson VI}
\switchcolumn*

\LA
\noindent Verum, cum ejus ópera ubíque terrárum monasteria jam ædificaréntur, innumerabilésque hómines religiosam ac piam vitam institúerent, Bononiæ anno Christi ducentésimo vigesimo primo supra millesimum in febrem incidit. Ex qua cum se moriturum intellígeret, convocátis frátribus et alumnis suæ disciplinæ, eos ad innocéntiam et integritátem cohortátus est. Postremo caritátem, humilitátem, paupertátem, tamquam certum patrimónium eis testaménto relíquit; fratribúsque orántibus, in illis verbis, Subveníte, Sancti Dei, occúrrite, Angeli, obdormívit in Dómino, octavo Idus Augusti. Quem póstea Gregórius nonus Pontifex rétulit in Sanctórum númerum.\par
\switchcolumn
\EN
\noindent By his work convents were already during his life-time established in all countries, and almost countless persons embraced the life of prayer and godliness. In the year of Christ 1221 he fell sick of a fever at Bologna. When he understood that he was about to die, he called together his brethren and disciples, and exhorted them to innocence and uprightness. Lastly, he left to them as his legacy, love, lowliness, and poverty, and as the brethren were all praying, and the words were being said, “Come ye to help him, ye Saints of God run ye to meet him, ye Angels”, he fell asleep in the Lord. It was the 6th day of August. Pope Gregory IX. afterwards placed his name among those of the Saints.\par
\switchcolumn*

\LA
\begin{Resp}\Rsign Iste homo perfécit ómnia quæ locútus est ei Deus, et dixit ad eum: Ingrédere in réquiem meam: \Star{} Quia te vidi justum coram me ex ómnibus géntibus.\end{Resp}
\begin{Resp}\Vsign Iste est, qui contémpsit vitam mundi, et pervénit ad cæléstia regna.\end{Resp}
\begin{Resp}\Rsign Quia te vidi justum coram me ex ómnibus géntibus.\end{Resp}
\begin{Resp}\Vsign Glória Patri, et Fílio, \Star{} et Spirítui Sancto.\end{Resp}
\begin{Resp}\Rsign Quia te vidi justum coram me ex ómnibus géntibus.\end{Resp}
\switchcolumn
\EN
\begin{Resp}\Rsign This is he which did according unto all that God commanded him; and God said unto him: Enter thou into My rest, \Star{} For thee have I seen righteous before Me among all people.\end{Resp}
\begin{Resp}\Vsign This is he which loved not his life in this world, and is come unto an everlasting kingdom.\end{Resp}
\begin{Resp}\Rsign For thee have I seen righteous before Me among all people.\end{Resp}
\begin{Resp}\Vsign Glory be to the Father, and to the Son, \Star{} and to the Holy Ghost.\end{Resp}
\begin{Resp}\Rsign For thee have I seen righteous before Me among all people.\end{Resp}
\switchcolumn*

\end{paracol}

\par\needspace{10\baselineskip}\addvspace{1.2em}{\centering\small\textsc{Die 5 Augusti} · \textsc{5 August}\par}
\Office{Sanctæ Mariæ Virginis ad Nives}{Sanctæ Mariæ Virginis ad Nives}{Duplex majus}
\addcontentsline{toc}{subsection}{Die 5 Augusti · Sanctæ Mariæ Virginis ad Nives \textit{— Sanctæ Mariæ Virginis ad Nives}}
\begin{paracol}{2}
\LA
\RefLine{Lectiones I–III: ex In Festis Beatae Mariae Virginis (C11).}
\switchcolumn
\EN
\RefLine{Lessons I–III: from the In Festis Beatae Mariae Virginis (C11).}
\switchcolumn*

\LA
\RefLine{Lectiones VII–IX: ex In Festis Beatae Mariae Virginis (C11).}
\switchcolumn
\EN
\RefLine{Lessons VII–IX: from the In Festis Beatae Mariae Virginis (C11).}
\switchcolumn*

\LA
\LessonHead{Lectio IV}
\switchcolumn
\EN
\LessonHead{Lesson IV}
\switchcolumn*

\LA
\noindent Liberio summo Pontifice, Joánnes patricius Romanus et uxor pari nobilitate, cum liberos non suscepíssent, quos bonórum herédes relínquerent, suam hereditátem sanctíssimæ Virgini Dei Matri vovérunt, ab ea summis precibus assidue peténtes, ut in quod pium opus eam pecúniam potíssimum erogari vellet aliquo modo significaret. Quorum preces et vota, ex animo facta, beáta Virgo Maria, benígne áudiens, miraculo comprobávit.\par
\switchcolumn
\EN
\noindent In the time of Pope Liberius, there lived at Rome a certain nobleman named John and a noble lady his wife, who had no children to whom to leave their substance. Then they vowed that they would make the most holy Virgin Mother of God their heiress, and earnestly besought her in some way to make known to them upon what godly work she would that the money should be spent. The Blessed Virgin Mary graciously listened to their prayers and heart-felt earnestness, and by a miracle assured them of her will.\par
\switchcolumn*

\LA
\begin{Resp}\Rsign Sicut cedrus exaltáta sum in Líbano, et sicut cypréssus in monte Sion: quasi myrrha elécta, \Star{} Dedi suavitátem odóris.\end{Resp}
\begin{Resp}\Vsign Et sicut cinnamómum et bálsamum aromatízans.\end{Resp}
\begin{Resp}\Rsign Dedi suavitátem odóris.\end{Resp}
\switchcolumn
\EN
\begin{Resp}\Rsign I was exalted like a cedar in Lebanon, and as a cypress-tree upon Mount Zion. Like the best myrrh \Star{} I yielded a pleasant odour.\end{Resp}
\begin{Resp}\Vsign Like cinnamon and sweet balsam.\end{Resp}
\begin{Resp}\Rsign I yielded a pleasant odour.\end{Resp}
\switchcolumn*

\LA
\LessonHead{Lectio V}
\switchcolumn
\EN
\LessonHead{Lesson V}
\switchcolumn*

\LA
\noindent Nonis ígitur Augusti, quo témpore in Urbi maximi calóres esse solent, noctu nix partem collis Exquilini contéxit. Qua nocte Dei Mater separatim Joánnem et conjugem in somnis admonuit ut quem locum nive conspersum vidérent, in eo ecclésiam ædificarent, quæ Maríæ vírginis nómine dedicarétur; se enim ita velle ab ipsis herédem institui. Quod Joánnes ad Libérium Pontificem détulit, qui idem per sómnium sibi contigisse affirmávit.\par
\switchcolumn
\EN
\noindent On the 5th day of August, which is that time when the heat of summer waxeth greatest in Rome, a part of the Esquiline Hill was covered by night with snow. And on this same night the Mother of God appeared in a dream to John and his wife separately, and told them that on that spot, which in the morning they should see clad with snow, they should build a Church, to be dedicated in the name of the Virgin Mary, for that this was the way in which she chose that they should make her their heiress. John went and told it to Pope Liberius, who declared that he also had been visited by a like dream.\par
\switchcolumn*

\LA
\begin{Resp}\Rsign Quæ est ista quæ procéssit sicut sol, et formósa tamquam Jerúsalem? \Star{} Vidérunt eam fíliæ Sion, et beátam dixérunt, et regínæ laudavérunt eam.\end{Resp}
\begin{Resp}\Vsign Et sicut dies verni circúmdabant eam flores rosárum et lília convállium.\end{Resp}
\begin{Resp}\Rsign Vidérunt eam fíliæ Sion, et beátam dixérunt, et regínæ laudavérunt eam.\end{Resp}
\switchcolumn
\EN
\begin{Resp}\Rsign Who is this that cometh up like the sun? This, comely as Jerusalem? \Star{} The daughters of Zion saw her, and called her blessed: the queens also, and they praised her.\end{Resp}
\begin{Resp}\Vsign And about her it was as the flower of roses in the spring of the year, and lilies of the valleys.\end{Resp}
\begin{Resp}\Rsign The daughters of Zion saw her, and called her blessed the queens also, and they praised her.\end{Resp}
\switchcolumn*

\LA
\LessonHead{Lectio VI}
\switchcolumn
\EN
\LessonHead{Lesson VI}
\switchcolumn*

\LA
\noindent Quare solemni sacerdotum et pópuli supplicatióne ad collem venit nive coopertum, et in eo locum ecclésiæ designávit, quæ Joánnis et uxoris pecúnia exstructa est, póstea a Xysto tertio restituta. Variis nomínibus primum est appellata, basilica Liberii, sancta Maria ad Præsepe; sed, cum multæ jam essent in Urbe ecclésiæ sub nómine sanctæ Maríæ Vírginis, ut quæ basilica novitate miraculi ac dignitate ceteris ejusdem nóminis basilicis præstaret, vocabuli étiam excelléntia significarétur, ecclésia sanctæ Maríæ majoris dicta est. Cujus dedicatiónis memória ex nive, quæ hac die mirabíliter cécidit, anniversaria celebritate cólitur.\par
\switchcolumn
\EN
\noindent Therefore he came in a solemn procession of Priests and people to the snow-clad hill, and traced upon that spot the plan of a Church which was built with the money of John and his wife. It was afterwards rebuilt by Sixtus III. At the beginning it was called by diverse names, sometimes the Liberian Basilica, sometimes the Church of St. Mary-at-the-Manger. Howbeit, since there are in Rome many Churches called after the Holy Virgin Mary, and this Church doth excel them all, both in honour, and because of the strange sign wherewith it was dedicated, it hath come to be called the Church of St. Mary, the Greater. The memory of the dedication thereof is kept every year by a Feast-day that taketh name from the wonderful fall of snow which on this day took place.\par
\switchcolumn*

\LA
\begin{Resp}\Rsign Ornátam monílibus fíliam Jerúsalem Dóminus concupívit: \Star{} Et vidéntes eam fíliæ Sion, beatíssimam prædicavérunt, dicéntes: * Unguéntum effúsum nomen tuum.\end{Resp}
\begin{Resp}\Vsign Astitit regína a dextris tuis in vestítu deauráto, circúmdata varietáte.\end{Resp}
\begin{Resp}\Rsign Et vidéntes eam fíliæ Sion, beatíssimam prædicavérunt, dicéntes:\end{Resp}
\begin{Resp}\Vsign Glória Patri, et Fílio, \Star{} et Spirítui Sancto.\end{Resp}
\begin{Resp}\Rsign Unguéntum effúsum nomen tuum.\end{Resp}
\switchcolumn
\EN
\begin{Resp}\Rsign When the Lord beheld the daughter of Jerusalem adorned with her jewels, He greatly desired her beauty \Star{} And when the daughters of Zion saw her, they cried out that she was most blessed, saying thy name is as ointment poured forth.\end{Resp}
\begin{Resp}\Vsign Upon thy right hand did stand the Queen in a vesture of gold wrought about with diverse colours.\end{Resp}
\begin{Resp}\Rsign And when the daughters of Zion saw her, they cried out that she was most blessed, saying:\end{Resp}
\begin{Resp}\Vsign Glory be to the Father, and to the Son, \Star{} and to the Holy Ghost.\end{Resp}
\begin{Resp}\Rsign Thy name is as ointment poured forth.\end{Resp}
\switchcolumn*

\end{paracol}

\par\needspace{10\baselineskip}\addvspace{1.2em}{\centering\small\textsc{Die 6 Augusti} · \textsc{6 August}\par}
\Office{In Transfiguratione Domini Nostri Jesu Christi}{In Transfiguratione Domini Nostri Jesu Christi}{Duplex II classis}
\addcontentsline{toc}{subsection}{Die 6 Augusti · In Transfiguratione Domini Nostri Jesu Christi \textit{— In Transfiguratione Domini Nostri Jesu Christi}}
\begin{paracol}{2}
\LA
\RefLine{Lectio IX: de Ss. Xysti II Papæ, Felicissimi et Agapiti Martyrum (08-06cc).}
\switchcolumn
\EN
\RefLine{Lesson IX: of S. Xystus II Pope, Felicissimus and Agapitus, Martyrs (08-06cc).}
\switchcolumn*

\LA
\LessonHead{Lectio I}
\switchcolumn
\EN
\LessonHead{Lesson I}
\switchcolumn*

\LA
\LessonTitle{De Epístola secúnda beáti Petri Apóstoli}
\switchcolumn
\EN
\LessonTitle{Lesson from the second letter of St. Peter the Apostle}
\switchcolumn*

\LA
\Cite{2 Pet 1:10-14}
\switchcolumn
\EN
\Cite{2 Pet 1:10-14}
\switchcolumn*

\LA
\noindent Fratres: Magis satágite, ut per bona ópera certam vestram vocatiónem et electiónem faciátis: hæc enim faciéntes, non peccábitis aliquándo. Sic enim abundánter ministrábitur vobis intróitus in ætérnum regnum Dómini nostri, et Salvatóris Jesu Christi. Propter quod incípiam vos semper commonére de his: et quidem sciéntes et confirmátos vos in præsénti veritáte. Justum autem árbitror, quámdiu sum in hoc tabernáculo, suscitáre vos in commonitióne: certus quod velox est deposítio tabernáculi mei secúndum quod et Dóminus noster Jesus Christus significávit mihi.\par
\switchcolumn
\EN
\noindent Wherefore, brethren, labour the more, that by good works you may make sure your calling and election. For doing these things, you shall not sin at any time. For so an entrance shall be ministered to you abundantly into the everlasting kingdom of our Lord and Saviour Jesus Christ. For which cause I will begin to put you always in remembrance of these things: though indeed you know them, and are confirmed in the present truth. But I think it meet as long as I am in this tabernacle, to stir you up by putting you in remembrance. Being assured that the laying away of this my tabernacle is at hand, according as our Lord Jesus Christ also hath signified to me.\par
\switchcolumn*

\LA
\begin{Resp}\Rsign Surge, illumináre, Jerúsalem, quia venit lumen tuum, \Star{} Et glória Dómini super te orta est.\end{Resp}
\begin{Resp}\Vsign Et ambulábunt gentes in lúmine tuo, et reges in splendóre ortus tui.\end{Resp}
\begin{Resp}\Rsign Et glória Dómini super te orta est.\end{Resp}
\switchcolumn
\EN
\begin{Resp}\Rsign Arise, shine, O Jerusalem, for thy light is come: \Star{} And the glory of the Lord is risen upon thee.\end{Resp}
\begin{Resp}\Vsign And the Gentiles shall come to thy light, and kings to the brightness of thy rising.\end{Resp}
\begin{Resp}\Rsign And the glory of the Lord is risen upon thee.\end{Resp}
\switchcolumn*

\LA
\LessonHead{Lectio II}
\switchcolumn
\EN
\LessonHead{Lesson II}
\switchcolumn*

\LA
\Cite{2 Pet 1:15-17}
\switchcolumn
\EN
\Cite{2 Pet 1:15-17}
\switchcolumn*

\LA
\noindent Dabo autem óperam et frequénter habére vos post óbitum meum, ut horum memóriam faciátis. Non enim doctas fábulas secúti notam fécimus vobis Dómini nostri Jesu Christi virtútem et præséntiam, sed speculatóres facti illíus magnitúdinis. Accípiens enim a Deo Patre honórem et glóriam, voce delápsa ad eum hujuscémodi a magnífica glória: Hic est Fílius meus diléctus, in quo mihi complácui; ipsum audíte.\par
\switchcolumn
\EN
\noindent And I will endeavour, that you frequently have after my decease, whereby you may keep a memory of these things. For we have not by following artificial fables, made known to you the power, and presence of our Lord Jesus Christ; but we were eyewitnesses of his greatness. For he received from God the Father, honour and glory: this voice coming down to him from the excellent glory: This is my beloved Son, in whom I am well pleased; hear ye him.\par
\switchcolumn*

\LA
\begin{Resp}\Rsign In splendénte nube Spíritus Sanctus visus est, Patérna vox audíta est: \Star{} Hic est Fílius meus diléctus, in quo mihi bene complácui; ipsum audíte.\end{Resp}
\begin{Resp}\Vsign Appáruit nubes obúmbrans, et vox Patris intónuit.\end{Resp}
\begin{Resp}\Rsign Hic est Fílius meus diléctus, in quo mihi bene complácui; ipsum audíte.\end{Resp}
\switchcolumn
\EN
\begin{Resp}\Rsign The Holy Ghost was manifested in the bright cloud and the Father was heard: \Star{} This is My beloved Son, in Whom I am well pleased; hear ye Him.\end{Resp}
\begin{Resp}\Vsign There was a cloud that overshadowed them, and the voice of the Father came out thereof in thunder.\end{Resp}
\begin{Resp}\Rsign This is My beloved Son, in Whom I am well pleased; hear ye Him.\end{Resp}
\switchcolumn*

\LA
\LessonHead{Lectio III}
\switchcolumn
\EN
\LessonHead{Lesson III}
\switchcolumn*

\LA
\Cite{2 Pet 1:18-21}
\switchcolumn
\EN
\Cite{2 Pet 1:18-21}
\switchcolumn*

\LA
\noindent Et hanc vocem nos audívimus de cælo allátam, cum essémus cum ipso in monte sancto. Et habémus firmiórem prophéticum sermónem, cui bene fácitis attendéntes quasi lucérnæ lucénti in caliginóso loco, donec dies elucéscat et lúcifer oriátur in córdibus vestris, hoc primum intellegéntes quod omnis prophetía Scriptúræ própria interpretatióne non fit. Non enim voluntáte humána alláta est aliquándo prophetía; sed Spíritu Sancto inspiráti locúti sunt sancti Dei hómines.\par
\switchcolumn
\EN
\noindent And this voice we heard brought from heaven, when we were with him in the holy mount. And we have the more firm prophetical word: whereunto you do well to attend, as to a light that shineth in a dark place, until the day dawn, and the day star arise in your hearts: Understanding this first, that no prophecy of scripture is made by private interpretation. For prophecy came not by the will of man at any time: but the holy men of God spoke, inspired by the Holy Ghost.\par
\switchcolumn*

\LA
\begin{Resp}\Rsign Vidéte qualem caritátem dedit nobis Deus Pater, \Star{} Ut fílii Dei nominémur et simus.\end{Resp}
\begin{Resp}\Vsign Scimus quóniam, cum apparúerit, símiles ei érimus, quóniam vidébimus eum sícuti est.\end{Resp}
\begin{Resp}\Rsign Ut fílii Dei nominémur et simus.\end{Resp}
\begin{Resp}\Vsign Glória Patri, et Fílio, \Star{} et Spirítui Sancto.\end{Resp}
\begin{Resp}\Rsign Ut fílii Dei nominémur et simus.\end{Resp}
\switchcolumn
\EN
\begin{Resp}\Rsign Behold what manner of love God the Father hath bestowed upon us, \Star{} That we should be called, and should be, the sons of God.\end{Resp}
\begin{Resp}\Vsign We know that when He shall appear, we shall be like Him, for we shall see Him as He is.\end{Resp}
\begin{Resp}\Rsign That we should be called, and should be, the Sons of God.\end{Resp}
\begin{Resp}\Vsign Glory be to the Father, and to the Son, \Star{} and to the Holy Ghost.\end{Resp}
\begin{Resp}\Rsign That we should be called, and should be, the Sons of God.\end{Resp}
\switchcolumn*

\LA
\LessonHead{Lectio IV}
\switchcolumn
\EN
\LessonHead{Lesson IV}
\switchcolumn*

\LA
\LessonTitle{Sermo sancti Leónis Papæ}
\switchcolumn
\EN
\LessonTitle{From the Sermons of Pope St. Leo (the Great)}
\switchcolumn*

\LA
\Source{Sermo de Transfiguratione, ante medium}
\switchcolumn
\EN
\Source{On the Transfiguration}
\switchcolumn*

\LA
\noindent Aperit Dóminus coram eléctis téstibus glóriam suam, et commúnem illam cum céteris córporis formam tanto splendóre claríficat, ut et fácies ejus solis fulgóri símilis, et vestítus candóri nívium esset æquális. In qua transfiguratióne illud quidem principáliter agebátur, ut de córdibus discipulórum crucis scándalum tollerétur; nec conturbáret eórum fidem voluntáriæ humílitas passiónis, quibus reveláta esset abscónditæ excelléntia dignitátis. Sed non minóre providéntia spes sanctæ Ecclésiæ fundabátur, ut totum Christi corpus agnósceret quali esset commutatióne donándum, ut ejus sibi honóris consórtium membra promítterent, qui in cápite præfulsísset.\par
\switchcolumn
\EN
\noindent The Lord taketh chosen witnesses, and in their presence, revealeth His glory. That form of body which He had in common with other men, He so transfigured with light, that His Face did shine as the sun, and His raiment became exceeding white as snow. Of this metamorphosis the chief work was to remove from the hearts of the disciples the stumbling at the Cross. Before their eyes was unveiled the splendour of His hidden majesty, that the lowliness of His freely-chosen suffering might not confound their faith. But none the less was there here laid by the Providence of God a solid foundation for the hope of the Holy Church, whereby the whole body of Christ should know with what a change it is yet to be honoured. The members of that body whose Head hath already been transfigured in light may promise themselves a share in His glory.\par
\switchcolumn*

\LA
\begin{Resp}\Rsign Inebriáti sunt ab ubertáte domus tuæ: \Star{} Et torrénte voluptátis tuæ potásti eos.\end{Resp}
\begin{Resp}\Vsign Quóniam apud te est fons vitæ, et in lúmine tuo vidébimus lumen.\end{Resp}
\begin{Resp}\Rsign Et torrénte voluptátis tuæ potásti eos.\end{Resp}
\switchcolumn
\EN
\begin{Resp}\Rsign They were abundantly satisfied with the fatness of thy house, and \Star{} Thou madest them drink of the river of thy pleasures.\end{Resp}
\begin{Resp}\Vsign For with thee is the fountain of life, and in thy light shall we see light.\end{Resp}
\begin{Resp}\Rsign And Thou madest them drink of the river of thy pleasures.\end{Resp}
\switchcolumn*

\LA
\LessonHead{Lectio V}
\switchcolumn
\EN
\LessonHead{Lesson V}
\switchcolumn*

\LA
\noindent Confirmándis vero Apóstolis et ad omnem sciéntiam provehéndis, ália quoque in illo miráculo accéssit instrúctio. Móyses enim et Elías, lex scílicet et prophétæ apparuérunt cum Dómino loquéntes; ut veríssime in illa quinque virórum præséntia complerétur quod dictum est: in duóbus vel tribus téstibus stat omne verbum. Quid hoc stabílius, quid fírmius verbo, in cujus prædicatióne véteris et novi Testaménti cóncinit tuba, et cum evangélica doctrína, antiquárum protestatiónum instruménta concúrrunt? Astipulántur enim sibi ínvicem utriúsque Fœ́deris páginæ; et, quem sub velámine mysteriórum præcedéntia signa promíserant, maniféstum atque perspícuum præséntis glóriæ splendor osténdit.\par
\switchcolumn
\EN
\noindent For the strengthening the Apostles and bringing them forward into all knowledge, there appeared unto them Moses and Elias that is, the Law and the Prophets talking with Him. Before five witnesses did His glorification take place, as though to fulfill that which is written: At the mouth of two witnesses, or at the mouth of three witnesses, shall the matter be established. (Deut. xix. 15.) What can be more certain, what can be better attested than this matter, which is proclaimed by the trumpets of both the Old and the New Testaments, and concerning which the witness of ancient testimony uniteth with the teaching of the Gospel? The pages of either Covenant strengthen one another, and the brightness of open glory maketh manifest and distinct Him Whom the former prophecies had promised under the veil of mysteries.\par
\switchcolumn*

\LA
\begin{Resp}\Rsign Præcéptor, bonum est nos hic esse: \Star{} Faciámus hic tria tabernácula; tibi unum, Móysi unum et Elíæ unum.\end{Resp}
\begin{Resp}\Vsign Non enim sciébat quid díceret.\end{Resp}
\begin{Resp}\Rsign Faciámus hic tria tabernácula; tibi unum, Móysi unum et Elíæ unum.\end{Resp}
\switchcolumn
\EN
\begin{Resp}\Rsign Master, it is good for us to be here. \Star{} Let us make here three tabernacles, one for thee, and one for Moses, and one for Elias.\end{Resp}
\begin{Resp}\Vsign For he wist not what to say.\end{Resp}
\begin{Resp}\Rsign Let us make here three tabernacles, one for thee, and one for Moses, and one for Elias.\end{Resp}
\switchcolumn*

\LA
\LessonHead{Lectio VI}
\switchcolumn
\EN
\LessonHead{Lesson VI}
\switchcolumn*

\LA
\noindent His ergo sacramentórum revelatiónibus Petrus Apóstolus incitátus, mundána spernens et terréna fastídiens, in æternórum desidérium quodam mentis rapiebátur excéssu; et, gáudio totíus visiónis implétus, ibi cum Jesu optábat habitáre, ubi manifésta ejus glória lætabátur. Unde et ait: Dómine, bonum est nos hic esse: si vis, faciámus hic tria tabernácula; tibi unum, Móysi unum et Elíæ unum. Sed huic suggestióni Dóminus non respóndit, signíficans, non quidem ímprobum, sed inordinátum esse quod cúperet; cum salvári mundus, nisi Christi morte, non posset, et exémplo Dómini in hoc vocarétur credéntium fides, ut licet non oportéret de beatitúdinis promissiónibus dubitári, intellegerémus tamen inter tentatiónes hujus vitæ prius nobis tolerántiam postulándam esse, quam glóriam.\par
\switchcolumn
\EN
\noindent The unveiling of such mysteries roused the mind of the Apostle Peter to an outburst of longing for the things eternal, which despised and disdained the things worldly and earthly overflowing with gladness at the vision, he yearned to dwell with Jesus there, where the revelation of His glory had rejoiced him. And so he said Master, it is good for us to be here if Thou wilt, let us make here three tabernacles, one for thee, and one for Moses, and one for Elias. To this proposal the Lord answered nothing, this signifying, that what Peter wished was not wrong, but out of place, since the world could not be saved but by the death of Christ. And the Lord's example was to call the faith of believers to this, that albeit we are behoven to have no doubts concerning the promise of eternal blessedness, yet we are to understand that, amid the trials of this life, we are to seek for endurance before glory.\par
\switchcolumn*

\LA
\begin{Resp}\Rsign Si ministrátio mortis, lítteris deformáta in lapídibus, fuit in glória, ita ut non possent inténdere fílii Israël in fáciem Móysi propter glóriam vultus ejus, quæ evacuátur: \Star{} Multo magis ministrátio spíritus, quæ manet, erit in glória.\end{Resp}
\begin{Resp}\Vsign Amplióris enim glóriæ Christus præ Móyse dignus est hábitus, quanto ampliórem honórem habet domo, qui fabricávit eam.\end{Resp}
\begin{Resp}\Rsign Multo magis ministrátio spíritus, quæ manet, erit in glória.\end{Resp}
\begin{Resp}\Vsign Glória Patri, et Fílio, \Star{} et Spirítui Sancto.\end{Resp}
\begin{Resp}\Rsign Multo magis ministrátio spíritus, quæ manet, erit in glória.\end{Resp}
\switchcolumn
\EN
\begin{Resp}\Rsign If the ministration of death, written and engraven on stones, was glorious, so that the children of Israel could not steadfastly behold the face of Moses, for the glory of his countenance which is done away: \Star{} How shall not the ministration of the Spirit, which abideth, be rather glorious?\end{Resp}
\begin{Resp}\Vsign For Christ was counted worthy of more glory than Moses, inasmuch as he who hath builded the house hath more honour than the house.\end{Resp}
\begin{Resp}\Rsign How shall not the ministration of the Spirit, which abideth, be rather glorious?\end{Resp}
\begin{Resp}\Vsign Glory be to the Father, and to the Son, \Star{} and to the Holy Ghost.\end{Resp}
\begin{Resp}\Rsign How shall not the ministration of the Spirit, which abideth, be rather glorious?\end{Resp}
\switchcolumn*

\LA
\LessonHead{Lectio VII}
\switchcolumn
\EN
\LessonHead{Lesson VII}
\switchcolumn*

\LA
\LessonTitle{Léctio sancti Evangélii secúndum Matthǽum}
\switchcolumn
\EN
\LessonTitle{From the Holy Gospel according to Matthew}
\switchcolumn*

\LA
\Cite{Matt 17:1-9}
\switchcolumn
\EN
\Cite{Matt 17:1-9}
\switchcolumn*

\LA
\noindent In illo témpore: Assúmpsit Jesus Petrum, et Jacóbum, et Joánnem fratrem ejus, et duxit illos in montem excélsum seórsum: et transfigurátus est ante eos. Et réliqua.\par
\switchcolumn
\EN
\noindent At that time: Jesus took Peter, and James, and John his brother, and brought them up into an high mountain apart, and was transfigured before them. And so on.\par
\switchcolumn*

\LA
\LessonTitle{Homilía sancti Joánnis Chrysóstomi}
\switchcolumn
\EN
\LessonTitle{Homily by St. John Chrysostom (Patriarch of Constantinople)}
\switchcolumn*

\LA
\Source{Homilia 57 in Matth., in init.}
\switchcolumn
\EN
\Source{57th on Matthew}
\switchcolumn*

\LA
\noindent Quóniam multa de perículis, multa de passióne sua, multa de morte et de cæde discipulórum locútus est Dóminus, et áspera complúra atque árdua eis injúnxit; et illa quidem in præsénti vita, et jam imminébant; bona vero in spe et exspectatióne erant: ut puta, quia servárent ánimam suam, si pérderent eam; quia in glória Patris sui ventúrus sit, et prǽmia redditúrus: ut visu étiam certióres fáceret, et osténderet quidnam sit illa glória, cum qua ventúrus est, quantum cápere póterant in hac præsénti vita, illis osténdit eámque détegit, ne aut sua aut Dómini morte dóleant, et máxime Petrus.\par
\switchcolumn
\EN
\noindent Since the Lord had spoken much concerning dangers, much concerning His Own sufferings, much concerning death, and the killing of His disciples, and had laid upon them many hard and grievous things, and since all these were in this present life, and already hanging over them, whereas the good things were matter for hope and waiting as, for example, that whosoever should lose his life for His sake should find it, for that the Son of Man should come in the glory of His Father, and reward every man according to his works. (Matth. xvi. 25, 27.) Therefore, to assure them by their own eyes, and show them what the glory is wherein He will come, He manifested and unveiled it to them, as far as in this life they were able to grasp it, lest they and especially Peter should grieve over their own deaths, or the death of their Lord.\par
\switchcolumn*

\LA
\begin{Resp}\Rsign Vocávit nos Deus vocatióne sua sancta, secúndum grátiam suam, quæ manifestáta est nunc; \Star{} Per illuminatiónem Salvatóris nostri Jesu Christi.\end{Resp}
\begin{Resp}\Vsign Qui destrúxit mortem, illuminávit autem vitam in incorruptiónem.\end{Resp}
\begin{Resp}\Rsign Per illuminatiónem Salvatóris nostri Jesu Christi.\end{Resp}
\switchcolumn
\EN
\begin{Resp}\Rsign God hath called us with an holy calling, according to His own grace, which is now made manifest; \Star{} By the appearing of our Saviour Jesus Christ.\end{Resp}
\begin{Resp}\Vsign Who hath abolished death, and brought life and immortality to light.\end{Resp}
\begin{Resp}\Rsign By the appearing of our Saviour Jesus Christ.\end{Resp}
\switchcolumn*

\LA
\LessonHead{Lectio VIII}
\switchcolumn
\EN
\LessonHead{Lesson VIII}
\switchcolumn*

\LA
\noindent Et vide quid agit, cum de regno et de gehénna disserúerit. Nam in eo quod dixit: Qui ínvenit ánimam suam, perdet eam; et quicúmque perdet eam grátia mei, invéniet ipsam; et in eo quod ait: Reddet unicuíque secúndum ópera sua; et regnum et gehénnam designávit. Cum ígitur de utrísque disserúerit, regnum quidem ut óculis cernátur, concédit, gehénnam autem mínime; quóniam rudióribus atque ineptióribus illud necessárium fuísset, sed, cum illi probi essent ac perspicáces, satis fuit eos a melióribus confirmári. Hoc étiam multo magis ipsum decébat. Non tamen omníno illud prætermísit, sed et aliquándo atrocitátem gehénnæ quasi ante óculos propónit, véluti cum Lázari imáginem descrípsit, et ejus méminit qui denários centum repétiit.\par
\switchcolumn
\EN
\noindent Behold what He doth, when He treateth of heaven and hell. Where He saith Whosoever will save his life shall lose it, and whosoever will lose his life for My sake shall find it. And again He shall reward every man according to his works in these words He pointeth at heaven and hell. But although He speaketh concerning both, He giveth a glimpse of heaven only and not of hell. To see hell would have profited the brutish and stupid, but His disciples were upright and clear-sighted, and therefore for them it was enough to be strengthened by the better things. This was what suited Him the best. Yet He left not the other altogether undone. Sometimes He set the horrors of hell, as it were, before the eyes, as for instance in the parable of the rich man and Lazarus, and that of him who was fain to wring the hundred pence from his fellow-servant.\par
\switchcolumn*

\LA
\begin{Resp}\Rsign Deus, qui fecit de ténebris lumen splendéscere, illúxit in córdibus nostris \Star{} Ad illuminatiónem sciéntiæ claritátis Dei, in fácie Jesu Christi.\end{Resp}
\begin{Resp}\Vsign Exórtum est in ténebris lumen rectis corde, miséricors et miserátor, et justus.\end{Resp}
\begin{Resp}\Rsign Ad illuminatiónem sciéntiæ claritátis Dei, in fácie Jesu Christi.\end{Resp}
\begin{Resp}\Vsign Glória Patri, et Fílio, \Star{} et Spirítui Sancto.\end{Resp}
\begin{Resp}\Rsign Ad illuminatiónem sciéntiæ claritátis Dei, in fácie Jesu Christi.\end{Resp}
\switchcolumn
\EN
\begin{Resp}\Rsign God, Who commanded the light to shine out of darkness, hath shined in our hearts, \Star{} To give the light of the knowledge of the glory of God, in the Face of Jesus Christ.\end{Resp}
\begin{Resp}\Vsign Unto the upright there ariseth light in the darkness he is gracious and full of compassion, and righteous.\end{Resp}
\begin{Resp}\Rsign To give the light of the knowledge of the glory of God, in the Face of Jesus Christ.\end{Resp}
\begin{Resp}\Vsign Glory be to the Father, and to the Son, \Star{} and to the Holy Ghost.\end{Resp}
\begin{Resp}\Rsign To give the light of the knowledge of the glory of God, in the Face of Jesus Christ.\end{Resp}
\switchcolumn*

\end{paracol}

\par\needspace{10\baselineskip}\addvspace{1.2em}{\centering\small\textsc{Die 6 Augusti} · \textsc{6 August}\par}
\Office{Ss. Xysti II Papæ, Felicissimi et Agapiti Martyrum}{S. Xystus II Pope, Felicissimus and Agapitus, Martyrs}{Simplex}
\addcontentsline{toc}{subsection}{Die 6 Augusti · Ss. Xysti II Papæ, Felicissimi et Agapiti Martyrum \textit{— S. Xystus II Pope, Felicissimus and Agapitus, Martyrs}}
\begin{paracol}{2}
\LA
\LessonHead{Lectio IX (pro commemoratione)}
\switchcolumn
\EN
\LessonHead{Lesson IX (when commemorated)}
\switchcolumn*

\LA
\LessonTitle{Commemoratio: Ss. Xysti II Papæ, Felicissimi et Agapiti Martyrum}
\switchcolumn
\EN
\LessonTitle{Commemoration of: S. Xystus II Pope, Felicissimus and Agapitus, Martyrs}
\switchcolumn*

\LA
\noindent Xystus secúndus, Atheniénsis, ex philósopho Christi discípulus, in persecutióne Valeriáni accusátus quod públice Christum prædicáret, comprehénsus tráhitur in templum Martis, propósita ei capitáli pœna, nisi illi simulácro sacrificáret. Qua impietáte constantíssime recusáta, cum ad martýrium ducerétur, occurrénti sancto Lauréntio et dolénter in hunc modum interrogánti, Quo progréderis sine fílio, pater? quo, sacérdos sancte, sine minístro próperas? respóndit: Non ego te désero, fili, majóra te manent pro Christi fide certámina: post tríduum me sequéris, sacerdótem levíta; intérea, si quid in thesáuris habes, paupéribus distríbue. Eódem ígitur die interféctus est una cum Felicíssimo et Agapíto diáconis, Januário, Magno, Vincéntio et Stéphano subdiáconis; et in cœmetério Callísti sepúltus octávo Idus Augústi, céteri vero in cœmetério Prætextáti. Sedit menses úndecim, dies duódecim. Quo témpore hábuit ordinatiónem mense Decémbri, creátis presbýteris quátuor, diáconis septem, epíscopis duóbus.\par
\switchcolumn
\EN
\noindent Xystus II. was an Athenian, who, from a philosopher, became a disciple of Christ. In the persecution under Valerian he was accused of openly preaching Christ, and was seized and haled to the temple of Mars, where he was given the choice of death or offering sacrifice to the idol. He firmly refused to commit that wickedness, and as he was being led away to seal his testimony, holy Lawrence ran up to him and in his grief said to him Father, whither goest thou without thy son Holy Priest, dost thou fare hence without a Deacon? Xystus answered him: I am not leaving thee, my son there awaiteth thee for Christ's truth a sterner wrestling than mine yet three days, and thou shalt follow me, the Deacon behind the Priest and in the meanwhile, if thou hast anything in the treasury, give it to the poor. Xystus was accordingly slain upon that day, and with him the Deacons Felicissimus and Agapitus, and the Subdeacons Januarius, Magnus, Vincent, and Stephen. He was buried in the cemetery of Callistus, and they in the cemetery of Praetextatus upon the 6th day of August, (in the year of our Lord 258.) He sat (in the throne of Peter) eleven months and twelve days. During that time he held one ordination in the month of December, wherein he made four Priests, seven Deacons, and two Bishops.\par
\switchcolumn*

\LA
\TeDeum{Te Deum}
\switchcolumn
\EN
\TeDeum{Te Deum}
\switchcolumn*

\end{paracol}

\par\needspace{10\baselineskip}\addvspace{1.2em}{\centering\small\textsc{Die 7 Augusti} · \textsc{7 August}\par}
\Office{S. Cajetani Confessoris}{St. Cajetan, Confessor}{Duplex}
\addcontentsline{toc}{subsection}{Die 7 Augusti · S. Cajetani Confessoris \textit{— St. Cajetan, Confessor}}
\begin{paracol}{2}
\LA
\RefLine{Lectiones I–III: de Scriptura occurrente.}
\switchcolumn
\EN
\RefLine{Lessons I–III: from the Scripture of the day.}
\switchcolumn*

\LA
\RefLine{Lectio IX: de S. Donati Episcopi et Martyris (08-07o).}
\switchcolumn
\EN
\RefLine{Lesson IX: of St. Donatus, Bishop and Martyr (08-07o).}
\switchcolumn*

\LA
\LessonHead{Lectio IV}
\switchcolumn
\EN
\LessonHead{Lesson IV}
\switchcolumn*

\LA
\noindent Cajetanus, nobili Thienæa gente Vicéntiæ ortus, statim a matre Deiparæ Virgini oblátus est. Mira a téneris annis morum innocéntia in eo eluxit, adeo ut Sanctus ab ómnibus nuncuparétur. Juris utriusque lauream Patavii adeptus, Romam profectus est; ubi, inter prælatos a Julio secundo collocatus, et sacerdotio initiatus, tanto divini amoris æstu succénsus est, ut relícta aula se totum Deo mancipaverit. Nosocomiis proprio ære fundatis, étiam morbo pestilénti laborántibus, suis ipse mánibus inserviébat. Proximórum salúti assidua cura incumbebat, dictus proptérea Venator animarum.\par
\switchcolumn
\EN
\noindent Cajetan was born at Vicenza, (in the year 1480,) of the noble family of the lords of Tiene, and was forthwith dedicated by his mother to the Virgin Mother of God. From his childhood such wonderful innocence shone in him that all called him the Saint. He took the degree of Doctor in Civil and Canon Law at Padua, and afterwards went to Rome, where Julius II. gave him a place among the Prelates. Having taken Priest's Orders, he became so full of the fire of the love of God, that he left the Court, that he might be free to work entirely for God. He founded hospitals at his own expense, and nursed the sick, even such as were suffering from the plague, with his own hands. He laboured with such constant earnestness for the salvation of his neighbours that he got the name of the Hunter of souls.\par
\switchcolumn*

\LA
\begin{Resp}\Rsign Honéstum fecit illum Dóminus, et custodívit eum ab inimícis, et a seductóribus tutávit illum: \Star{} Et dedit illi claritátem ætérnam.\end{Resp}
\begin{Resp}\Vsign Justum dedúxit Dóminus per vias rectas, et osténdit illi regnum Dei.\end{Resp}
\begin{Resp}\Rsign Et dedit illi claritátem ætérnam.\end{Resp}
\switchcolumn
\EN
\begin{Resp}\Rsign The Lord made him honourable, and defended him from his enemies, and kept him safe from those that lay in wait for him; \Star{} And gave him perpetual glory.\end{Resp}
\begin{Resp}\Vsign He went down with him into the pit, and left him not in bonds.\end{Resp}
\begin{Resp}\Rsign And gave him perpetual glory.\end{Resp}
\switchcolumn*

\LA
\LessonHead{Lectio V}
\switchcolumn
\EN
\LessonHead{Lesson V}
\switchcolumn*

\LA
\noindent Collapsam ecclesiasticórum disciplínam ad formam apostolicæ vitæ instaurare desiderans, ordinem Clericórum regulárium instituit, qui, abdicáta rerum ómnium terrenárum sollicitúdine, nec réditus possiderent nec vitæ subsidia a fidelibus péterent, sed solis eleemosynis sponte oblátis víverent. Itaque, approbante Cleménte septimo, ad aram maximam basilicæ Vaticanæ una cum Joánne Petro Caráfa, episcopo Theatino, qui póstea Paulus quartus Pontifex maximus fuit, et aliis duobus exímiæ pietátis viris, vota solemnia emisit. In Urbis direptióne a milítibus crudelíssime vexátus ut pecúniam próderet, quam dudum in cælestes thesáuros manus páuperum deportaverant, verbera, torménta et carceres invicta patiéntia sustinuit. In suscepto vitæ instituto constantíssime perseverávit, soli divinæ providéntiæ inhærens, quam sibi numquam defuísse aliquándo miracula comprobarunt.\par
\switchcolumn
\EN
\noindent From a desire to restore the corrupted discipline of the clergy to the mould of the Apostolic life, he founded, (in 1524,) a Congregation of Clerks Regular, who should give up all care of earthly things, neither keeping any income, nor begging the needful things of life from the faithful, but living only on such alms as might be given them unasked. For this end, and with the approval of Clement VII., Cajetan himself, together with John Peter Carafa, Archbishop of Chieti, (afterwards Pope Paul IV.), and two other men of eminent godliness, took solemn vows at the High Altar of St. Peter's Church on the Vatican. When the city of Rome was sacked by the troops (of the Constable de Bourbon, in 1527,) Cajetan was most cruelly illused to make him reveal his wealth, which had long before been laid up for him in heaven by the hands of the poor, and he endured with unconquered patience stripes, torture, and imprisonment. He held on bravely in the way of life he had taken up, trusting altogether to the Providence of God, Whose unfailing care of him was sometimes attested by miracles.\par
\switchcolumn*

\LA
\begin{Resp}\Rsign Amávit eum Dóminus, et ornávit eum: stolam glóriæ índuit eum, \Star{} Et ad portas paradísi coronávit eum.\end{Resp}
\begin{Resp}\Vsign Índuit eum Dóminus lorícam fídei, et ornávit eum.\end{Resp}
\begin{Resp}\Rsign Et ad portas paradísi coronávit eum.\end{Resp}
\switchcolumn
\EN
\begin{Resp}\Rsign The Lord loved him and beautified him; He clothed him with a robe of glory, \Star{} And crowned him at the gates of Paradise.\end{Resp}
\begin{Resp}\Vsign The Lord hath put on him the breast-plate of faith, and hath adorned him.\end{Resp}
\begin{Resp}\Rsign And crowned him at the gates of Paradise.\end{Resp}
\switchcolumn*

\LA
\LessonHead{Lectio VI}
\switchcolumn
\EN
\LessonHead{Lesson VI}
\switchcolumn*

\LA
\noindent Divini cultus studium, nitórem domus Dei, sacrórum rituum observantiam et sanctíssimæ Eucharistiæ frequentiórem usum maxime promovit. Hæresum monstra et látebras non semel detéxit ac profligávit. Oratiónem ad octo passim horas jugibus lácrimis protrahebat: sæpe in éxtasim raptus ac prophetíæ dono illustris Romæ, nocte natalícia, ad præsepe Dómini infántem Jesum accípere meruit a Deípara in ulnas suas. Corpus íntegras noctes intérdum verberatiónibus affligebat; nec umquam adduci pótuit ut vitæ asperitátem emolliret, testátus in cínere velle se mori. Dénique, ex animi dolóre concepto morbo, quod offéndi plebis seditióne Deum vidéret, cælésti visióne recreatus, Neapoli migrávit in cælum; ibique corpus ejus in ecclésia sancti Pauli magna religióne cólitur. Quem, multis miraculis in vita et post mortem gloriosum, Clemens décimus, Pontifex maximus, Sanctórum número adscripsit.\par
\switchcolumn
\EN
\noindent He was a great advancer of care in the worship of God, of splendour in the house of God, of exactness in the holy ceremonies, and of the often receiving of the most holy Eucharist. The hideous forms and dark convolutions of heresy he more than once unmasked and abolished. He would remain in prayer with abundance of tears as much as eight hours at a time. He was often thrown into trances, and was celebrated for the gift of prophecy. One Christmas night at Rome, when he was praying before the Lord's manger, he was deemed worthy that the Mother of God should lay the Child Jesus in his arms. He sometimes spent the whole night in whipping himself, nor could he ever be persuaded to soften the hardness of his life, but witnessed that he was fain to die in sackcloth and ashes. In the end he fell ill with grief at the offence against God, which the people of Naples committed by rebelling (against the establishment of the Inquisition.) Refreshed by a vision from heaven, he departed thither (on the 7th day of August, 1547.) His body lieth at Naples in the Church of St. Paul, where it is held in great reverence. Pope Clement X., finding him to have been illustrious for miracles, both during his life and after his death, enrolled his name among those of the Saints.\par
\switchcolumn*

\LA
\begin{Resp}\Rsign Iste homo perfécit ómnia quæ locútus est ei Deus, et dixit ad eum: Ingrédere in réquiem meam: \Star{} Quia te vidi justum coram me ex ómnibus géntibus.\end{Resp}
\begin{Resp}\Vsign Iste est, qui contémpsit vitam mundi, et pervénit ad cæléstia regna.\end{Resp}
\begin{Resp}\Rsign Quia te vidi justum coram me ex ómnibus géntibus.\end{Resp}
\begin{Resp}\Vsign Glória Patri, et Fílio, \Star{} et Spirítui Sancto.\end{Resp}
\begin{Resp}\Rsign Quia te vidi justum coram me ex ómnibus géntibus.\end{Resp}
\switchcolumn
\EN
\begin{Resp}\Rsign This is he which did according unto all that God commanded him; and God said unto him: Enter thou into My rest, \Star{} For thee have I seen righteous before Me among all people.\end{Resp}
\begin{Resp}\Vsign This is he which loved not his life in this world, and is come unto an everlasting kingdom.\end{Resp}
\begin{Resp}\Rsign For thee have I seen righteous before Me among all people.\end{Resp}
\begin{Resp}\Vsign Glory be to the Father, and to the Son, \Star{} and to the Holy Ghost.\end{Resp}
\begin{Resp}\Rsign For thee have I seen righteous before Me among all people.\end{Resp}
\switchcolumn*

\LA
\LessonHead{Lectio VII}
\switchcolumn
\EN
\LessonHead{Lesson VII}
\switchcolumn*

\LA
\LessonTitle{Léctio sancti Evangélii secúndum Matthǽum}
\switchcolumn
\EN
\LessonTitle{From the Holy Gospel according to Matthew}
\switchcolumn*

\LA
\Cite{Matt 6:24-33}
\switchcolumn
\EN
\Cite{Matt 6:24-33}
\switchcolumn*

\LA
\noindent In illo témpore: Dixit Jesus discípulis suis: Nemo potest duobus dóminis servire. Et réliqua.\par
\switchcolumn
\EN
\noindent At that time: Jesus said unto His disciples: No man can serve two masters. And so on.\par
\switchcolumn*

\LA
\LessonTitle{Homilía sancti Augustíni Epíscopi}
\switchcolumn
\EN
\LessonTitle{Homily by St. Augustine, Bishop of Hippo}
\switchcolumn*

\LA
\Source{Liber 2 de Sermóne Dómini in monte, cap. 14}
\switchcolumn
\EN
\Source{Bk. ii. on the Lord's Sermon on the Mount, ch. xiv}
\switchcolumn*

\LA
\noindent Nemo potest duobus dóminis servire. Ad hanc ipsam intentiónem referéndum est quod consequenter exponit, dicens: Aut enim unum ódio habebit, et álterum diliget; aut álterum patiétur, et álterum contemnet. Quæ verba diligenter consideranda sunt: nam, qui sint duo domini, deinceps osténdit, cum dicit: Non potestis Deo servire, et mammonæ. Mammona apud Hebræos divítiæ appellari dicúntur. Congruit et Púnicum nomen; nam lucrum Púnice mammon dícitur. Sed qui servit mammonæ, illi útique servit, qui rebus istis terrenis mérito suæ perversitátis præpositus, magistrátus hujus sæculi a Dómino dícitur. Aut enim unum ódio habébit homo, et álterum diliget, id est, Deum; aut álterum patiétur, et álterum contemnet. Patiétur enim durum et perniciósum dóminum, quisquis servit mammonæ; sua enim cupiditate implicatus, súbditur diabolo, et non eum díligit. Quis enim est qui diligat diabolum? sed tamen patitur.\par
\switchcolumn
\EN
\noindent No man can serve two masters, and this is further explained for either he will hate the one, and love the other; or else he will hold to the one, and despise the other. These words we ought carefully to weigh, for the Lord showeth straightway who be the two masters whom we have choice of: Ye cannot serve God and Mammon. Mammon is a term which the Hebrews are said to use for riches. It is also a Carthaginian word for the Punic for gain is mammon. He which serveth mammon, serveth that evil one who hath perversely chosen to be lord of these earthly things, and is called by the Lord the prince of this world. (John xiv. 30.) Of these two masters, either a man will hate the one and love the other, that is God or he will hold to the one and despise the other. He which serveth mammon holdeth to an hard and destroying master, for he is led captive by his lust, and sold a slave to the devil, and him loveth no man is there any man that loveth the devil And yet there be that hold to him.\par
\switchcolumn*

\LA
\begin{Resp}\Rsign Iste est qui ante Deum magnas virtútes operátus est, et de omni corde suo laudávit Dóminum: \Star{} Ipse intercédat pro peccátis ómnium populórum.\end{Resp}
\begin{Resp}\Vsign Ecce homo sine queréla, verus Dei cultor, ábstinens se ab omni ópere malo, et pérmanens in innocéntia sua.\end{Resp}
\begin{Resp}\Rsign Ipse intercédat pro peccátis ómnium populórum.\end{Resp}
\switchcolumn
\EN
\begin{Resp}\Rsign This is he which wrought great wonders before God, and praised the Lord with all his heart. \Star{} May he pray for all people, that their sins may be forgiven unto them.\end{Resp}
\begin{Resp}\Vsign Behold a man without blame, a worshipper of God in truth, keeping himself clean from every evil work, and abiding still in his innocency.\end{Resp}
\begin{Resp}\Rsign May he pray for all people, that their sins may be forgiven unto them.\end{Resp}
\switchcolumn*

\LA
\LessonHead{Lectio VIII}
\switchcolumn
\EN
\LessonHead{Lesson VIII}
\switchcolumn*

\LA
\noindent Ideo, inquit, dico vobis, non habére sollicitúdinem ánimæ vestræ quid edatis, neque corpori vestro quid induátis; ne forte, quamvis jam supérflua non quærántur, propter ipsa necessaria cor duplicétur, et ad ipsa conquirenda, nostra detorqueátur intentio, cum aliquid quasi misericorditer operámur: id est, ut cum consúlere alícui videri vólumus, nostrum emoluméntum ibi potius, quam illíus utilitátem attendámus; et ideo nobis non videámur peccáre, quia non supérflua, sed necessaria sunt, quæ cónsequi vólumus.\par
\switchcolumn
\EN
\noindent Therefore, I say unto you, Take no thought for your life, what ye shall eat, or what ye shall drink; nor yet for your body, what ye shall put on lest, albeit such things are not idle, but needful to be sought after, yet the seeking for things even needful should divide the heart and our intention should be corrupted when we do something as it were mercifully that is, lest, when we would seem to be seeking another's good, it should be profit to ourselves, rather than benefit to him, that we seek and therefore we seem not to ourselves to sin, because we would seek things not idle, but needful.\par
\switchcolumn*

\LA
\begin{Resp}\Rsign Sint lumbi vestri præcíncti, et lucérnæ ardéntes in mánibus vestris: \Star{} Et vos símiles homínibus exspectántibus dóminum suum, quando revertátur a núptiis.\end{Resp}
\begin{Resp}\Vsign Vigiláte ergo, quia nescítis qua hora Dóminus vester ventúrus sit.\end{Resp}
\begin{Resp}\Rsign Et vos símiles homínibus exspectántibus dóminum suum, quando revertátur a núptiis.\end{Resp}
\begin{Resp}\Vsign Glória Patri, et Fílio, \Star{} et Spirítui Sancto.\end{Resp}
\begin{Resp}\Rsign Et vos símiles homínibus exspectántibus dóminum suum, quando revertátur a núptiis.\end{Resp}
\switchcolumn
\EN
\begin{Resp}\Rsign Let your loins be girded about, and your lights burning; \Star{} And ye yourselves like unto men that wait for their lord, when he will return from the wedding.\end{Resp}
\begin{Resp}\Vsign Watch therefore, for ye know not what hour your Lord doth come.\end{Resp}
\begin{Resp}\Rsign And ye yourselves like unto men that wait for their lord, when he will return from the wedding.\end{Resp}
\begin{Resp}\Vsign Glory be to the Father, and to the Son, \Star{} and to the Holy Ghost.\end{Resp}
\begin{Resp}\Rsign And ye yourselves like unto men that wait for their lord, when he will return from the wedding.\end{Resp}
\switchcolumn*

\end{paracol}

\par\needspace{10\baselineskip}\addvspace{1.2em}{\centering\small\textsc{Die 7 Augusti} · \textsc{7 August}\par}
\Office{S. Donati Episcopi et Martyris}{St. Donatus, Bishop and Martyr}{Simplex}
\addcontentsline{toc}{subsection}{Die 7 Augusti · S. Donati Episcopi et Martyris \textit{— St. Donatus, Bishop and Martyr}}
\begin{paracol}{2}
\LA
\LessonHead{Lectio IX (pro commemoratione)}
\switchcolumn
\EN
\LessonHead{Lesson IX (when commemorated)}
\switchcolumn*

\LA
\LessonTitle{Commemoratio: S. Donati Episcopi et Martyris}
\switchcolumn
\EN
\LessonTitle{Commemoration of: St. Donatus, Bishop and Martyr}
\switchcolumn*

\LA
\noindent Donatus, paréntibus propter Jesu Christi fidem interfectis, Arrétium in Etrúriam cum Hilarino monacho ex fuga se cóntulit: cujus étiam urbis epíscopus creátus est. Ubi, cum a Quadratiáno præfecto in Juliáni persecutióne uterque idola venerari juberétur, negántibus illis se tam nefárium scelus commissuros, Hilarinus in óculis Quadratiáni tamdiu fustibus cæsus est, dum efflaret ánimam; Donatus, veheménter et ipse cruciatus, gládio percússus est. Quorum córpora a Christiánis prope eam urbem honorifice sepúlta sunt.\par
\switchcolumn
\EN
\noindent Donatus was the child of a father and mother who had both been slain for Jesus Christ's sake. He fled with the monk Hilarinus to Arezzo in Tuscany, of which city he afterwards became Bishop. There the Praefect Quadratian, during the persecution under Julian, commanded both Hilarinus and Donatus to worship idols, and when they both refused to commit such abominable iniquity, Hilarinus was beaten with clubs before the eyes of Quadratian, until he gave up the ghost. Donatus also was savagely tortured, and slain with the sword. The Christians buried their bodies honourably hard by the city.\par
\switchcolumn*

\LA
\TeDeum{Te Deum}
\switchcolumn
\EN
\TeDeum{Te Deum}
\switchcolumn*

\end{paracol}

\par\needspace{10\baselineskip}\addvspace{1.2em}{\centering\small\textsc{Die 8 Augusti} · \textsc{8 August}\par}
\Office{Ss. Cyriaci, Largi et Smaragdi Martyrum}{Sts. Cyriacus, Largus and Smaragdus, Martyrs}{Semiduplex}
\addcontentsline{toc}{subsection}{Die 8 Augusti · Ss. Cyriaci, Largi et Smaragdi Martyrum \textit{— Sts. Cyriacus, Largus and Smaragdus, Martyrs}}
\begin{paracol}{2}
\LA
\RefLine{Lectiones I–III: de Scriptura occurrente.}
\switchcolumn
\EN
\RefLine{Lessons I–III: from the Scripture of the day.}
\switchcolumn*

\LA
\LessonHead{Lectio IV}
\switchcolumn
\EN
\LessonHead{Lesson IV}
\switchcolumn*

\LA
\noindent Cyríacus diaconus, cum Sisinio, Largo et Smaragdo diutius inclusus in carcere, multa édidit miracula, in quibus Arthémiam Diocletiáni filiam precibus a dæmone liberávit; missusque ad Saporem Persárum regem, Jobíam étiam ejus filiam a nefario spíritu erípuit. Rege vero ejus patre cum quadringentis ac trigínta aliis baptizatis, Romam rediit; ubi, Maximiáni imperatóris jussu comprehénsus, catenis vinctus ante rhedam suam tráhitur; et post dies quatuor e carcere edúctus, pice liquáta perfusus et in catasta exténsus, demum cum Largo et Smaragdo aliisque viginti secúri percússus est via Salaria, ad hortos Sallustiános. Quorum córpora in eádem via décimo septimo Kalendas Aprilis sepúlta a Joánne presbytero, póstea sexto Idus Augusti a Marcello Pontifice et Lucina nobili femina lineis velis involuta et pretiósis unguentis condíta, in ipsíus Lucinæ prædium via Ostiénsi, septimo ab Urbe lápide transláta sunt.\par
\switchcolumn
\EN
\noindent Cyriacus the Deacon was long kept in prison with Sisinius, Largus, and Smaragdus, and wrought many wonderful works. Among other things he by his prayers freed from a devil Arthemia, a daughter of Diocletian, and, being sent to Sapor, King of the Persians, also delivered his daughter Jobia from a foul spirit. He baptized the king, her father, and four hundred and thirty others, and afterwards returned to Rome. He was arrested by command of the Emperor Maximian, and dragged in chains before his chariot. Then after four days he was brought forth from prison, had boiling pitch poured upon him, was stretched on a block, \textasciitilde{}(and belaboured with a cudgel,) and at last was slain with the axe, along with Largus, Smaragdus, and twenty others, at the gardens of Sallust, on the Salarian Way. On this Way were their bodies buried by John the Priest, on the 16th day of March, (in the year of our Lord 303,) but afterwards, on the 8th of August, Pope Marcellus and the noble lady Lucina took them and wrapped them in linen, and embalmed them with costly ointments, and carried them to the farm belonging to the said lady Lucina, at the seventh milestone from Rome on the road to Ostia.\par
\switchcolumn*

\LA
\begin{Resp}\Rsign Sancti tui, Dómine, mirábile consecúti sunt iter, serviéntes præcéptis tuis, ut inveniréntur illǽsi in aquis válidis: \Star{} Terra appáruit árida: et in Mari Rubro via sine impediménto.\end{Resp}
\begin{Resp}\Vsign Quóniam percússit petram, et fluxérunt aquæ, et torréntes inundavérunt.\end{Resp}
\begin{Resp}\Rsign Terra appáruit árida: et in Mari Rubro via sine impediménto.\end{Resp}
\switchcolumn
\EN
\begin{Resp}\Rsign thy Saints, O Lord, have passed a wonderful way, serving thy commandments, that they might be found without hurt in the midst of the mighty waters. \Star{} Dry land appeared, and, out of the Red Sea, a way without impediment.\end{Resp}
\begin{Resp}\Vsign He smote the rock, and the waters gushed out, and the streams overflowed.\end{Resp}
\begin{Resp}\Rsign Dry land appeared, and, out of the Red Sea, a way without impediment.\end{Resp}
\switchcolumn*

\LA
\LessonHead{Lectio V}
\switchcolumn
\EN
\LessonHead{Lesson V}
\switchcolumn*

\LA
\LessonTitle{Sermo sancti Joánnis Chrysóstomi}
\switchcolumn
\EN
\LessonTitle{From the Sermon of S. John Chrysostom}
\switchcolumn*

\LA
\Source{Sermo 1 de Martyribus, tom. 3.}
\switchcolumn
\EN
\Source{Serm. 1. de Mart, tom. 3}
\switchcolumn*

\LA
\noindent Nemo est, qui nesciat, Mártyrum glorias ad hoc divino consílio a Dei pópulis frequentári, ut et illis débitus honor dicétur, et nobis virtútis exémpla, favénte Christo, monstréntur; ut, dum hæc ita celebrári perspícimus, cognoscámus quanta eos glória máneat in cælis, quorum natalítia táliter celebrántur in terris; quo possímus, étiam ipsi, tálibus provocári exémplis, virtúte pari, devotióne consímili ac fide; ut, Christo præstánte, dimicáre et víncere hostem possímus, ut, parta victória, cum iísdem Sanctis in regnis cæléstibus triumphémus.\par
\switchcolumn
\EN
\noindent Every man knoweth how, by the good Providence of God, the diverse glories of His Martyrs are held in such esteem by His people, that the same His Saints in all places receive worthy honour, and before us is set, by the favour of Christ, the noble example of their courage: thus are we stirred up to consider, on the occasion of these Holidays, how great glory doth abide them in heaven, whose birthdays are thus kept upon earth: thereby, also, we are roused to strive to be like them, brave, godly, and true: so that, in the strength of Christ, we, like them, may wrestle with, and conquer our enemy, and, when we have gained the same victory that they gained, may with them at last be glorified in the kingdom of heaven.\par
\switchcolumn*

\LA
\begin{Resp}\Rsign Vérbera carníficum non timuérunt Sancti Dei, moriéntes pro Christi nómine: \Star{} Ut herédes fíerent in domo Dómini.\end{Resp}
\begin{Resp}\Vsign Tradidérunt córpora sua propter Deum ad supplícia.\end{Resp}
\begin{Resp}\Rsign Ut herédes fíerent in domo Dómini.\end{Resp}
\switchcolumn
\EN
\begin{Resp}\Rsign The Saints of God shrank not from the stripes of the executioners, but died for Christ's Name's sake \Star{} That they might be made joint-heirs in the house of the Lord.\end{Resp}
\begin{Resp}\Vsign They gave their bodies for God's sake to death.\end{Resp}
\begin{Resp}\Rsign That they might be made joint-heirs in the house of the Lord.\end{Resp}
\switchcolumn*

\LA
\LessonHead{Lectio VI}
\switchcolumn
\EN
\LessonHead{Lesson VI}
\switchcolumn*

\LA
\noindent Quis est enim, qui eórum volens mérito copulári, nisi prius constántiam eórum téneat, sectétur fidem, imitétur virtútem passiónis; eórum glóriam páribus vitæ lineaméntis aut invéniat aut exquírat? Qui, etsi martýrio par esse non possit, tamen múneris tanti dignitáte se quisque bonis áctibus dignum præbeat. Adest enim clementíssimus Deus, qui desiderántibus suis aut martýrium præbeat, aut, sine martýrio, cum Sanctis præmia divína retríbuat.\par
\switchcolumn
\EN
\noindent For what man is there willing to share their reward, that if he do not first lay hold on their steadfastness, follow after the example of their faith, and imitate their brave patience, can either seek or find their glory by likeness to their lives? But whosoever doth so follow them, let him not doubt but that, though in very deed he gain not the crown of martyrdom, he is yet able by good works to make himself meet therefor. For we have a most merciful God, Which either giveth Martyrdom unto such as be willing, or, without Martyrdom, doth make them joint heirs with the Saints in the kingdom of God.\par
\switchcolumn*

\LA
\begin{Resp}\Rsign Tamquam aurum in fornáce probávit eléctos Dóminus, et quasi holocáusti hóstiam accépit illos: et in témpore erit respéctus illórum: \Star{} Quóniam donum et pax est eléctis Dei.\end{Resp}
\begin{Resp}\Vsign Qui confídunt in illum, intélligent veritátem, et fidéles in dilectióne acquiéscent illi.\end{Resp}
\begin{Resp}\Rsign Quóniam donum et pax est eléctis Dei.\end{Resp}
\begin{Resp}\Vsign Glória Patri, et Fílio, \Star{} et Spirítui Sancto.\end{Resp}
\begin{Resp}\Rsign Quóniam donum et pax est eléctis Dei.\end{Resp}
\switchcolumn
\EN
\begin{Resp}\Rsign As gold in the furnace hath the Lord tried His chosen ones, and received them as a burnt-offering, and yet a while, and they shall be regarded; \Star{} For the grace of God, and His peace, are with His chosen.\end{Resp}
\begin{Resp}\Vsign They that put their trust in Him shall understand the truth and such as be faithful in love shall abide with Him.\end{Resp}
\begin{Resp}\Rsign For the grace of God, and His peace, are with His chosen.\end{Resp}
\begin{Resp}\Vsign Glory be to the Father, and to the Son, \Star{} and to the Holy Ghost.\end{Resp}
\begin{Resp}\Rsign For the grace of God, and His peace, are with His chosen.\end{Resp}
\switchcolumn*

\LA
\LessonHead{Lectio VII}
\switchcolumn
\EN
\LessonHead{Lesson VII}
\switchcolumn*

\LA
\LessonTitle{Léctio sancti Evangélii secúndum Marcum}
\switchcolumn
\EN
\LessonTitle{From the holy Gospel according to Mark}
\switchcolumn*

\LA
\Cite{Marc 16:15-18}
\switchcolumn
\EN
\Cite{Mark 16:15-18}
\switchcolumn*

\LA
\noindent In illo témpore: Dixit Jesus discípulis suis: Euntes in mundum univérsum prædicate Evangélium omni creaturæ. Et réliqua.\par
\switchcolumn
\EN
\noindent At that time, Jesus said unto His disciples: Go ye into all the world, and preach the Gospel to every creature. And so on.\par
\switchcolumn*

\LA
\LessonTitle{Homilía sancti Gregórii Papæ}
\switchcolumn
\EN
\LessonTitle{Homily by Pope St. Gregory (the Great)}
\switchcolumn*

\LA
\Source{Homilia 29 in Evangelia, post initium}
\switchcolumn
\EN
\Source{29th on the Gospels}
\switchcolumn*

\LA
\noindent Potest omnis creaturæ nómine omnis natio Géntium designari. Ante enim dictum fuerat: In viam Géntium ne abieritis; nunc autem dícitur: Prædicate omni creaturæ: ut scilicet prius a Judǽa Apostolórum repulsa prædicátio tunc nobis in adjutórium fieret, cum hanc illa ad damnatiónis suæ testimónium superba repulísset. Sed cum discipulos ad prædicándum Véritas mittit, quid aliud in mundo facit, nisi grana séminis spargit? Et pauca grana mittit in semine, ut multárum messium fruges recipiat ex nostra fide.\par
\switchcolumn
\EN
\noindent By the words “every creature” we may understand every tribe of the Gentiles. Of aforetime it had been said, Go not into the way of the Gentiles, (Matth. x. 5,) but now, Preach the Gospel to every creature, that, since the Jews had proudly rejected the preaching of the Apostles, that might become our gain which was the seal of their condemnation. But when the Eternal Truth sendeth forth His disciples to preach, what doth He but scatter seed over the field of the world He scattered abroad a few grains for seed, that He may afterward reap an abundant harvest in our faith.\par
\switchcolumn*

\LA
\begin{Resp}\Rsign Propter testaméntum Dómini, et leges patérnas, Sancti Dei perstitérunt in amóre fraternitátis: \Star{} Quia unus fuit semper spíritus in eis, et una fides.\end{Resp}
\begin{Resp}\Vsign Ecce quam bonum et quam jucúndum habitáre fratres in unum.\end{Resp}
\begin{Resp}\Rsign Quia unus fuit semper spíritus in eis, et una fides.\end{Resp}
\switchcolumn
\EN
\begin{Resp}\Rsign Because of the covenant of the Lord, and the laws of their fathers, the Saints of God abode in brotherly love, \Star{} For one spirit and one faith was ever in them.\end{Resp}
\begin{Resp}\Vsign Behold how good and how pleasant it is for brethren to dwell together in unity.\end{Resp}
\begin{Resp}\Rsign For one spirit and one faith was ever in them.\end{Resp}
\switchcolumn*

\LA
\LessonHead{Lectio VIII}
\switchcolumn
\EN
\LessonHead{Lesson VIII}
\switchcolumn*

\LA
\noindent Neque étenim in universo mundo tanta fidelium messis exsúrgeret, si de manu Dómini super rationalem terram illa electa grana prædicántium non venissent. Sequitur: Qui crediderit et baptizátus fúerit, salvus erit: qui vero non credíderit, condemnábitur. Fortasse unusquísque apud semetípsum dicat: Ego jam credidi, salvus ero. Verum dicit, si fidem opéribus tenet. Vera étenim fides est, quæ in hoc, quod verbis dicit, móribus non contradicit. Hinc est enim quod de quibusdam falsis fidelibus Paulus dicit: Qui confiténtur se nosse Deum, factis autem negant.\par
\switchcolumn
\EN
\noindent The great harvest of faithful souls throughout the whole world would never have sprung up, if the hand of the Lord had not first scattered those chosen grains of preachers over the reasonable soil of men's minds. Then is written, He that believeth and is baptized shall be saved but he that believeth not shall be damned. Now, perchance, thou sayest in thine heart I believe, and therefore I shall be saved. True, if to thy faith thou dost add works. He only hath a living faith whose life doth not give the lie to his profession. It is of this that Paul speaketh, where he saith of certain vain believers, They profess that they know God but in works they deny Him. (Tit. i. 16.)\par
\switchcolumn*

\LA
\begin{Resp}\Rsign Sancti mei, qui in carne pósiti, certámen habuístis: \Star{} Mercédem labóris ego reddam vobis.\end{Resp}
\begin{Resp}\Vsign Veníte benedícti Patris mei, percípite regnum.\end{Resp}
\begin{Resp}\Rsign Mercédem labóris ego reddam vobis.\end{Resp}
\begin{Resp}\Vsign Glória Patri, et Fílio, \Star{} et Spirítui Sancto.\end{Resp}
\begin{Resp}\Rsign Mercédem labóris ego reddam vobis.\end{Resp}
\switchcolumn
\EN
\begin{Resp}\Rsign O ye My Saints, who, being in the flesh, didst have striving \Star{} I will render unto you a reward of your labours.\end{Resp}
\begin{Resp}\Vsign Come, ye blessed of My Father, inherit the kingdom\end{Resp}
\begin{Resp}\Rsign I will render unto you a reward of your labours.\end{Resp}
\begin{Resp}\Vsign Glory be to the Father, and to the Son, \Star{} and to the Holy Ghost.\end{Resp}
\begin{Resp}\Rsign I will render unto you a reward of your labours.\end{Resp}
\switchcolumn*

\LA
\LessonHead{Lectio IX}
\switchcolumn
\EN
\LessonHead{Lesson IX}
\switchcolumn*

\LA
\noindent Signa autem eos qui credituri sunt, hæc sequéntur: In nómine meo dæmónia ejícient, linguis loquéntur novis, serpéntes tollent: et si mortíferum quid bíberint, non eis nocébit: super ægros manus impónent, et bene habébunt. Numquidnam, fratres mei, quia ista signa non fácitis, minime creditis? Sed hæc necessaria in exordio Ecclésiæ fuérunt. Ut enim ad fidem crésceret multitúdo credéntium, miraculis fuerat nutriénda: quia et nos, cum arbústa plantamus, tamdiu eis aquam infúndimus, quoúsque ea in terra jam coaluisse videámus: et si semel radícem fíxerint, irrigátio cessabit. Hinc est enim quod Paulus dicit: Linguæ in signum sunt non fidelibus, sed infidelibus.\par
\switchcolumn
\EN
\noindent And these signs shall follow them that believe In My name they shall cast out devils, they shall speak with new tongues, they shall take up serpents and if they drink any deadly thing, it shall not hurt them they shall lay hands on the sick, and they shall recover. My brethren, these signs do not follow us. Do we, then, not believe Nay. The truth is, these things were needful when the Church was young. That she might grow by the increase of the faithful, she needed to be nourished with miracles. So we, when we plant a young tree, continually water and tend it, till we see that it hath taken firm root in the earth but when once it hath taken firm root, it can grow of itself. Hence Paul saith of tongues: Tongues are for a sign not to them that believe, but to them that believe not. (Cor. xiv. 22.)\par
\switchcolumn*

\LA
\TeDeum{Te Deum}
\switchcolumn
\EN
\TeDeum{Te Deum}
\switchcolumn*

\end{paracol}

\par\needspace{10\baselineskip}\addvspace{1.2em}{\centering\small\textsc{Die 9 Augusti} · \textsc{9 August}\par}
\Office{S. Joannis Mariæ Vianney Confessoris}{St. Jean-Marie Vianney, Confessor}{Duplex}
\addcontentsline{toc}{subsection}{Die 9 Augusti · S. Joannis Mariæ Vianney Confessoris \textit{— St. Jean-Marie Vianney, Confessor}}
\begin{paracol}{2}
\LA
\RefLine{Lectiones I–III: de Scriptura occurrente.}
\switchcolumn
\EN
\RefLine{Lessons I–III: from the Scripture of the day.}
\switchcolumn*

\LA
\RefLine{Lectiones VII–VIII: ex Commúni Confessoris non pontificis (C5).}
\switchcolumn
\EN
\RefLine{Lessons VII–VIII: from the Common of a Confessor not a Bishop (C5).}
\switchcolumn*

\LA
\RefLine{Lectio IX: de In Vigilia S. Laurentii Mart. (08-09t).}
\switchcolumn
\EN
\RefLine{Lesson IX: of Vigil of St. Lawrence, Martyr (08-09t).}
\switchcolumn*

\LA
\LessonHead{Lectio IV}
\switchcolumn
\EN
\LessonHead{Lesson IV}
\switchcolumn*

\LA
\noindent Joánnes Maria Vianney, in pago Dardilly, Diœcesis Lugdunénsis piis ruriculis ortus, ab infántia plura dedit sanctitátis indicia. Cum octennis oves custodiret, modo puerulos, ad imáginem Deíparæ genuflexos, Rosárium verbo et exemplo edocére, modo sorori vel alteri commisso grege, secretiórem locum pétere solebat, quo expeditior ante simulacrum Vírginis, oratióni vacaret. Páuperum amantíssimus, eos turmatim in patris domum dedúcere et omnímodo adjuvare in deliciis habebat. Litteris imbuendus, párocho Vici Ecully tráditus est; sed ut erat tardioris ingenii, in stúdiis fere insuperábiles expertus est difficultates. Jejunio et oratióne divinam opem implorávit, et facilitátem discéndi rogatúrus, tumulum sancti Francisci Regis, stipem quæritans, adívit, Theologíæ curriculo operose confecto, satis idoneus invéntus est, qui sacris initiarétur.\par
\switchcolumn
\EN
\noindent John Mary Vianney, born of pious peasants in the village of Dardilly in the diocese of Lyons, gave many signs of holiness from his infancy. When, at the age of eight, he was taking care of the sheep, he would sometimes by word and example instruct little boys, kneeling before a statue of the Mother of God, in the use of the Rosary; and at other times, entrusting the flock to his sister or to another child, he was wont to seek out a more retired spot, that he might more readily devote himself to prayer before an image of the Virgin. Having a very great love for the poor, he would lead them in crowds to his father's house, and he took a delight in aiding them in every way. That he might be initiated into letters, he was sent to the parish priest of the village of Ecully; but as he was very slow to understand, he encountered almost unsurmountable difficulties in his studies. Fasting and praying, he entreated the divine assistance, and, with a view to begging for a facility in learning, he approached the tomb of St. Francis Regis, earnestly beseeching him for that gift. Having most laboriously passed through the course of theology, he was found to be sufficiently suitable to receive holy orders.\par
\switchcolumn*

\LA
\begin{Resp}\Rsign Honéstum fecit illum Dóminus, et custodívit eum ab inimícis, et a seductóribus tutávit illum: \Star{} Et dedit illi claritátem ætérnam.\end{Resp}
\begin{Resp}\Vsign Justum dedúxit Dóminus per vias rectas, et osténdit illi regnum Dei.\end{Resp}
\begin{Resp}\Rsign Et dedit illi claritátem ætérnam.\end{Resp}
\switchcolumn
\EN
\begin{Resp}\Rsign The Lord made him honourable, and defended him from his enemies, and kept him safe from those that lay in wait for him; \Star{} And gave him perpetual glory.\end{Resp}
\begin{Resp}\Vsign He went down with him into the pit, and left him not in bonds.\end{Resp}
\begin{Resp}\Rsign And gave him perpetual glory.\end{Resp}
\switchcolumn*

\LA
\LessonHead{Lectio V}
\switchcolumn
\EN
\LessonHead{Lesson V}
\switchcolumn*

\LA
\noindent In pago Ecully, præeunte párocho, cujus vicarius renunciátus fuerat, potiores pastoralis perfectiónis gradus totis viribus attingere contendit. Elapso triennio, in viculum Ars qui non ita multo post diœcesi Bellicénsi adscriptus est, quasi Angelus de cælo fuit missus et omnino squalentis ac desertæ suæ parœciæ fáciem florentíssime renovávit. In consciéntiis judicándis ac moderandis ad plurimas horas quotídie assiduus, frequentem Eucharistiæ usum invéxit, pias sodalitates instaurávit: mirum autem in modum téneram in Immaculátam Vírginem ánimis pietátem índidit. Ratus vero pastoris esse, flagítia concréditæ plebis expiare, nec oratiónibus, nec vigiliis, maceratiónibus et continuis jejuniis parcebat. Tantam viri Dei virtútem cum Satan ferre non posset, eum vexatiónibus primum, dein aperto certamine adortus est; sed atrocíssimas afflictiónes patienter tolerábat Joánnes Maria.\par
\switchcolumn
\EN
\noindent In the village of Ecully, under the guidance of the parish priest, whose assistant he had been appointed, he strove with all his strength to attain to the higher degrees of pastoral perfection. After three years had gone by, he was sent, like an Angel from heaven, to the small village of Ars, which not so long after was included in the diocese of Belley, and in a most brilliant manner he entirely renewed the condition of his neglected and forsaken parish. Continually engaged for many hours daily in hearing confessions and in giving spiritual direction, he introduced the frequent reception of the Eucharist, and organized pious sodalities: and in a remarkable manner he inspired into souls a tender devotion to the Immaculate Virgin. And, deeming that it is the duty of the pastor to expiate the sins of the flock accredited to him, he spared neither prayers, nor vigils, nor mortifications and continual fastings. Since Satan could not endure such great virtues in this man of God, he assailed him, first with mere annoyances, and afterwards in open combat; but John Mary patiently endured the most malevolent injuries.\par
\switchcolumn*

\LA
\begin{Resp}\Rsign Amávit eum Dóminus, et ornávit eum: stolam glóriæ índuit eum, \Star{} Et ad portas paradísi coronávit eum.\end{Resp}
\begin{Resp}\Vsign Índuit eum Dóminus lorícam fídei, et ornávit eum.\end{Resp}
\begin{Resp}\Rsign Et ad portas paradísi coronávit eum.\end{Resp}
\switchcolumn
\EN
\begin{Resp}\Rsign The Lord loved him and beautified him; He clothed him with a robe of glory, \Star{} And crowned him at the gates of Paradise.\end{Resp}
\begin{Resp}\Vsign The Lord hath put on him the breast-plate of faith, and hath adorned him.\end{Resp}
\begin{Resp}\Rsign And crowned him at the gates of Paradise.\end{Resp}
\switchcolumn*

\LA
\LessonHead{Lectio VI}
\switchcolumn
\EN
\LessonHead{Lesson VI}
\switchcolumn*

\LA
\noindent Invitabátur sæpius a curiónibus vicínis ut, Missionariórum more, animárum saluti, qua concionando, qua confessiónes excipiéndo consúleret et singulis præsto semper erat. Studio glóriæ Dei incénsus, effécit, ut pia Missiónum exercítia in ámplius centum parœciis, constituto perpetuo censu, institueréntur. Inter hæc, Deo servum suum miraculis et charismátibus illustrante, orta est celebris illa peregrinátio, qua, per vicennium centum fere míllia hóminum cujusque ordinis et ætátis, quotannis Ars conflúxerint non solum e Galliæ et Europæ, sed étiam Américæ díssitis provinciis. Labóribus potius quam senio consumptus, prænuntiáto suæ mortis die, in osculo Dómini quiévit, die quarta Augusti, anno millesimo octingentésimo quinquagesimo nono, annos natus tres ac septuagínta. Quem multis clarum miraculis, Pius décimus inter Beatos, Pius vero undecimus inter Sanctos Cælites anno sacro adscripsit, ejusque festum ad universam Ecclésiam extendit, anno vero quinquagesimo sacerdotii sui ómnium parochórum cælestem patronum constítuit.\par
\switchcolumn
\EN
\noindent He was very often asked by the neighbouring priests to labour for the salvation of souls after the manner of the Missionaries, either by preaching sermons, or by hearing confessions, and he was always at hand in every case. Burning with zeal for the glory of God, he brought it about, that the pious exercises of Missions were established in more than an hundred parishes arranged in a continuous and permanent series. Meanwhile, as God was rendering his servant famous by miracles and by graces, there began that celebrated pilgrimage, in which, throughout a period of twenty years, nearly one hundred thousand persons of every class flocked to Ars, not only from France and from Europe, but even from the distant regions of America. Worn out by labours rather than by old age, having foretold the day of his death, he went to rest in the embrace of the Lord, on the 4th day of August, in the year 1859, and of his age the seventy-third. After he became illustrious for many miracles, Pius X added him to the number of the Blessed , and Pius XI, in the holy year numbered him with the Saints in heaven and extended his feast to the universal Church, and on the fiftieth anniversary of his own priesthood, appointed him the heavenly patron of all parish priests.\par
\switchcolumn*

\LA
\begin{Resp}\Rsign Iste homo perfécit ómnia quæ locútus est ei Deus, et dixit ad eum: Ingrédere in réquiem meam: \Star{} Quia te vidi justum coram me ex ómnibus géntibus.\end{Resp}
\begin{Resp}\Vsign Iste est, qui contémpsit vitam mundi, et pervénit ad cæléstia regna.\end{Resp}
\begin{Resp}\Rsign Quia te vidi justum coram me ex ómnibus géntibus.\end{Resp}
\begin{Resp}\Vsign Glória Patri, et Fílio, \Star{} et Spirítui Sancto.\end{Resp}
\begin{Resp}\Rsign Quia te vidi justum coram me ex ómnibus géntibus.\end{Resp}
\switchcolumn
\EN
\begin{Resp}\Rsign This is he which did according unto all that God commanded him; and God said unto him: Enter thou into My rest, \Star{} For thee have I seen righteous before Me among all people.\end{Resp}
\begin{Resp}\Vsign This is he which loved not his life in this world, and is come unto an everlasting kingdom.\end{Resp}
\begin{Resp}\Rsign For thee have I seen righteous before Me among all people.\end{Resp}
\begin{Resp}\Vsign Glory be to the Father, and to the Son, \Star{} and to the Holy Ghost.\end{Resp}
\begin{Resp}\Rsign For thee have I seen righteous before Me among all people.\end{Resp}
\switchcolumn*

\end{paracol}

\par\needspace{10\baselineskip}\addvspace{1.2em}{\centering\small\textsc{Die 9 Augusti} · \textsc{9 August}\par}
\Office{In Vigilia S. Laurentii Mart.}{Vigil of St. Lawrence, Martyr}{Vigilia Simplex}
\addcontentsline{toc}{subsection}{Die 9 Augusti · In Vigilia S. Laurentii Mart. \textit{— Vigil of St. Lawrence, Martyr}}
\begin{paracol}{2}
\LA
\LessonHead{Lectio I}
\switchcolumn
\EN
\LessonHead{Lesson I}
\switchcolumn*

\LA
\LessonTitle{Commemoratio: In Vigilia S. Laurentii Mart. Léctio sancti Evangélii secúndum Matthǽum}
\switchcolumn
\EN
\LessonTitle{Commemoration of: Vigil of St. Lawrence, Martyr Continuation of the Holy Gospel according to Matthew}
\switchcolumn*

\LA
\Cite{Matt 16:24-27}
\switchcolumn
\EN
\Cite{Matt 16:24-27}
\switchcolumn*

\LA
\noindent In illo témpore: Dixit Jesus discípulis suis: Si quis vult post me veníre, ábneget semetípsum, et tollat crucem suam, et sequátur me. Et réliqua.\par
\switchcolumn
\EN
\noindent At that time, Jesus said unto His disciples: If any man will come after Me, let him deny himself, and take up his cross, and follow Me. And so on.\par
\switchcolumn*

\LA
\LessonTitle{Homilía sancti Gregórii Papæ}
\switchcolumn
\EN
\LessonTitle{Homily by Pope St. Gregory (the Great)}
\switchcolumn*

\LA
\Source{Homilia 32 in Evangelia}
\switchcolumn
\EN
\Source{32nd on the Gospels}
\switchcolumn*

\LA
\noindent Quia Dóminus ac Redémptor noster novus homo veni in mundum, nova præcépta dedit mundo. Vitæ étenim nostræ véteri, in vítiis enutrítæ, contrarietátem oppósuit novitátis suæ. Quid enim vetus, quid carnális homo nóverat, nisi sua retinére, aliéna rápere, si posset; concupíscere, si non posset? Sed cæléstis médicus síngulis quibúsque vítiis obviántia ádhibet medicaménta. Nam, sicut arte medicínæ cálida frígidis, frígida cálidis curántur; ita Dóminus noster contrária oppósuit medicaménta peccátis, ut lúbricis continéntiam, tenácibus largitátem, iracúndis mansuetúdinem, elátis præcíperet humilitátem.\par
\switchcolumn
\EN
\noindent Our Lord and Redeemer came into the world a new Man, and gave the world new commandments. For against the ways of our old life, brought and bred up in sin, He set the contrast of His new life. It was the old way, according to the knowledge of the carnal man, for every man to keep his own goods, and, if he were able to do it, to take his neighbour's goods also, and, if he were not able to take them, at least to lust after them. But the Heavenly Physician hath medicines wherewith to meet all the diseases of sin. For, even, as by the art of the physician, things hot are healed by things cold, and things cold by things hot, so doth our Lord set against sin holiness, ordaining for the lecherous purity, for the miserly munificence, for the hot-tempered meekness, and for the proud lowliness.\par
\switchcolumn*

\LA
\TeDeum{Te Deum}
\switchcolumn
\EN
\TeDeum{Te Deum}
\switchcolumn*

\end{paracol}

\par\needspace{10\baselineskip}\addvspace{1.2em}{\centering\small\textsc{Die 10 Augusti} · \textsc{10 August}\par}
\Office{S. Laurentii Martyris}{St. Lawrence, Martyr}{Duplex II classis cum Octava simplici}
\addcontentsline{toc}{subsection}{Die 10 Augusti · S. Laurentii Martyris \textit{— St. Lawrence, Martyr}}
\begin{paracol}{2}
\LA
\LessonHead{Lectio I}
\switchcolumn
\EN
\LessonHead{Lesson I}
\switchcolumn*

\LA
\LessonTitle{De libro Ecclesiástici}
\switchcolumn
\EN
\LessonTitle{Lesson from the book of Ecclesiasticus}
\switchcolumn*

\LA
\Cite{Sir 51:1-7}
\switchcolumn
\EN
\Cite{Sir 51:1-7}
\switchcolumn*

\LA
\noindent Confitébor tibi, Dómine, Rex, et collaudábo te Deum Salvatórem meum. Confitébor nómini tuo: quóniam adjútor et protéctor factus es mihi, et liberásti corpus meum a perditióne, a láqueo linguæ iníquæ et a lábiis operántium mendácium, et in conspéctu astántium factus es mihi adjútor. Et liberásti me secúndum multitúdinem misericórdiæ nóminis tui a rugiéntibus, præparátis ad escam: de mánibus quæréntium ánimam meam, et de portis tribulatiónum, quæ circumdedérunt me; a pressúra flammæ, quæ circúmdedit me, et in médio ignis non sum æstuáta; de altitúdine ventris ínferi, et a lingua coinquináta, et a verbo mendácii, a rege iníquo, et a lingua injústa.\par
\switchcolumn
\EN
\noindent A prayer of Jesus the son of Sirach. I will give glory to thee, O Lord, O King, and I will praise thee, O God my Saviour. I will give glory to thy name: for thou hast been a helper and protector to me. And hast preserved my body from destruction, from the snare of an unjust tongue, and from the lips of them that forge lies, and in the sight of them that stood by, thou hast been my helper. And thou hast delivered me, according to the multitude of the mercy of thy name, from them that did roar, prepared to devour. Out of the hands of them that sought my life, and from the gates of afflictions, which compassed me about: From the oppression of the flame which surrounded me, and in the midst of the fire I was not burnt. From the depth of the belly of hell, and from an unclean tongue, and from lying words, from an unjust king, and from a slanderous tongue:\par
\switchcolumn*

\LA
\begin{Resp}\Rsign Levíta Lauréntius bonum opus operátus est, qui per signum crucis cæcos illuminávit, \Star{} Et thesáuros Ecclésiæ dedit paupéribus.\end{Resp}
\begin{Resp}\Vsign Dispérsit, dedit paupéribus: justítia ejus manet in sǽculum sǽculi.\end{Resp}
\begin{Resp}\Rsign Et thesáuros Ecclésiæ dedit paupéribus.\end{Resp}
\switchcolumn
\EN
\begin{Resp}\Rsign The Levite Lawrence wrought a good work, in that with the sign of the Cross he gave sight to the blind, \Star{} And dispersed among the needy the riches of the Church.\end{Resp}
\begin{Resp}\Vsign He hath dispersed, he hath given to the poor his righteousness endureth forever.\end{Resp}
\begin{Resp}\Rsign And dispersed among the needy the riches of the Church.\end{Resp}
\switchcolumn*

\LA
\LessonHead{Lectio II}
\switchcolumn
\EN
\LessonHead{Lesson II}
\switchcolumn*

\LA
\Cite{Sir 51:8-13}
\switchcolumn
\EN
\Cite{Sir 51:8-13}
\switchcolumn*

\LA
\noindent Laudábit usque ad mortem ánima mea Dóminum, et vita mea appropínquans erat in inférno deórsum. Circumdedérunt me úndique, et non erat qui adjuváret. Respíciens eram ad adjutórium hóminum, et non erat. Memorátus sum misericórdiæ tuæ, Dómine, et operatiónis tuæ, quæ a sǽculo sunt: quóniam éruis sustinéntes te, Dómine, et líberas eos de mánibus géntium. Exaltásti super terram habitatiónem meam, et pro morte defluénte deprecátus sum.\par
\switchcolumn
\EN
\noindent My soul shall praise the Lord even to death. And my life was drawing near to hell beneath. They compassed me on every side, and there was no one that would help me. I looked for the succour of men, and there was none. I remembered thy mercy, O Lord, and thy works, which are from the beginning of the world. How thou deliverest them that wait for thee, O Lord, and savest them out of the hands of the nations. Thou hast exalted my dwelling place upon the earth and I have prayed for death to pass away.\par
\switchcolumn*

\LA
\begin{Resp}\Rsign Puer meus, noli timére, quia ego tecum sum, dicit Dóminus: \Star{} Si transíeris per ignem, flamma non nocébit tibi, et odor ignis non erit in te.\end{Resp}
\begin{Resp}\Vsign Liberábo te de manu pessimórum et éruam te de manu fórtium.\end{Resp}
\begin{Resp}\Rsign Si transíeris per ignem, flamma non nocébit tibi, et odor ignis non erit in te.\end{Resp}
\switchcolumn
\EN
\begin{Resp}\Rsign O my child, be not afraid, for I am with thee, saith the Lord. \Star{} When thou passest through the fire, the flame shall not harm thee, neither shall the smell of fire pass upon thee.\end{Resp}
\begin{Resp}\Vsign And I will deliver thee out of the hand of the wicked, and I will redeem thee out of the hand of the terrible.\end{Resp}
\begin{Resp}\Rsign When thou passest through the fire the flame shall not harm thee, neither shall the smell of fire pass upon thee.\end{Resp}
\switchcolumn*

\LA
\LessonHead{Lectio III}
\switchcolumn
\EN
\LessonHead{Lesson III}
\switchcolumn*

\LA
\Cite{Sir 51:14-17}
\switchcolumn
\EN
\Cite{Sir 51:14-17}
\switchcolumn*

\LA
\noindent Invocávi Dóminum, Patrem Dómini mei, ut non derelínquat me in die tribulatiónis meæ, et in témpore superbórum, sine adjutório. Laudábo nomen tuum assídue, et collaudábo illud in confessióne, et exaudíta est orátio mea. Et liberásti me de perditióne, et eripuísti me de témpore iníquo. Proptérea confitébor, et laudem dicam tibi, et benedícam nómini Dómini.\par
\switchcolumn
\EN
\noindent I called upon the Lord, the father of my Lord, that he would not leave me in the day of my trouble, and in the time of the proud without help. I will praise thy name continually, and will praise it with thanksgiving, and my prayer was heard. And thou hast saved me from destruction, and hast delivered me from the evil time. Therefore I will give thanks, and praise thee, and bless the name of the Lord.\par
\switchcolumn*

\LA
\begin{Resp}\Rsign Strinxérunt córporis membra pósita super cratículam: ministrántibus prunas insúltat Levíta Christi: \Star{} Beáte Laurénti, Martyr Christi, intercéde pro nobis.\end{Resp}
\begin{Resp}\Vsign Mea nox obscúrum non habet, sed ómnia in luce claréscunt.\end{Resp}
\begin{Resp}\Rsign Beáte Laurénti, Martyr Christi, intercéde pro nobis.\end{Resp}
\begin{Resp}\Vsign Glória Patri, et Fílio, \Star{} et Spirítui Sancto.\end{Resp}
\begin{Resp}\Rsign Beáte Laurénti, Martyr Christi, intercéde pro nobis.\end{Resp}
\switchcolumn
\EN
\begin{Resp}\Rsign They laid him upon the grating, and stretched out his limbs. Christ's Levite held in derision them that brought the live coals. \Star{} O Lawrence, blessed witness for Christ, plead for us!\end{Resp}
\begin{Resp}\Vsign The darkness is no darkness to me, but the night is all as clear as the morning, that shineth more and more unto the perfect day.\end{Resp}
\begin{Resp}\Rsign O Lawrence, blessed witness for Christ, plead for us!\end{Resp}
\begin{Resp}\Vsign Glory be to the Father, and to the Son, \Star{} and to the Holy Ghost.\end{Resp}
\begin{Resp}\Rsign O Lawrence, blessed witness for Christ, plead for us!\end{Resp}
\switchcolumn*

\LA
\LessonHead{Lectio IV}
\switchcolumn
\EN
\LessonHead{Lesson IV}
\switchcolumn*

\LA
\LessonTitle{Sermo sancti Leónis Papæ}
\switchcolumn
\EN
\LessonTitle{The Lesson is taken From the Sermons of Pope St. Leo (the Great.)}
\switchcolumn*

\LA
\Source{In Natali S. Laurentii, post initium}
\switchcolumn
\EN
\Source{For St. Lawrence's Birthday}
\switchcolumn*

\LA
\noindent Cum furor gentílium potestátum in electíssima quæque Christi membra sævíret, ac præcípue eos qui órdinis erant sacerdotális impéteret, in levítam Lauréntium, qui non solum ministério sacramentórum sed étiam dispensatióne ecclesiásticæ substántiæ præeminébat, ímpius persecútor efférbuit, dúplicem sibi prædam de uníus viri comprehensióne promíttens; quem, si fecísset sacræ pecúniæ traditórem, fáceret étiam veræ religiónis exsórtem. Armátur ítaque gémina face homo pecúniæ cúpidus et veritátis inimícus: avarítia, ut rápiat aurum; impietáte, ut áuferat Christum. Póstulat sibi ab immaculáto sacrárii prǽsule opes ecclesiásticas, quibus avidíssimus inhiábat, inférri. Cui levíta castíssimus, ubi eas repósitas habéret osténdens, numerosíssimos sanctórum páuperum óbtulit greges, in quorum victu atque vestítu inamissíbiles condíderat facultátes, quæ tanto intégrius erant salvæ, quanto sánctius probabántur expénsæ.\par
\switchcolumn
\EN
\noindent Then the fury of the heathen power was raging against Christ's choicest members, in aiming in especial at such as were of the Priestly Order, the wicked persecutor turned fiercely on the Levite Lawrence, who was remarkable, not only as a minister of the Sacraments, but also as distributor of the property of the Church, promising himself a double prey by the taking of this one man, namely, to make him betray the consecrated treasure, and apostatise from the true faith. The wretch was thus doubly fired by his greed for money and his hatred of the truth, his greed urging him to seize the gold, and his wickedness to rob a believer of Christ. He demanded of the upright keeper of the sacred treasury, to bring him the wealth of the Church, for which his rapacity longed. But the pureminded Levite showed him where these riches were stored, by bringing before him a great multitude of holy poor, by the feeding and clothing of whom he had laid up all. that he had, in such wise, that it could be lost no more, and was now all the safer, as the way of spending it had been the holier.\par
\switchcolumn*

\LA
\begin{Resp}\Rsign Quo progréderis sine fílio, pater? quo, sacérdos sancte, sine diácono próperas? \Star{} Tu nunquam sine minístro sacrifícium offérre consuéveras.\end{Resp}
\begin{Resp}\Vsign Quid ergo in me displícuit paternitáti tuæ? numquid degénerem me probásti? Experíre utrum idóneum minístrum elégeris, cui commisísti Domínici sánguinis dispensatiónem.\end{Resp}
\begin{Resp}\Rsign Tu nunquam sine minístro sacrifícium offérre consuéveras.\end{Resp}
\switchcolumn
\EN
\begin{Resp}\Rsign Father, whither goest thou without thy son? Holy Priest, dost thou fare hence without a Deacon? \Star{} It hath never been thy use to offer sacrifice without a minister.\end{Resp}
\begin{Resp}\Vsign What therefore in me hath displeased thee, my Father? Hast thou tried me and found me unworthy to be called thy son? Make trial if I am indeed a useless servant, even I, whom thou didst choose, to commit unto me the administration of the Cup of the Blood of the Lord.\end{Resp}
\begin{Resp}\Rsign It hath never been thy use to offer sacrifice without a minister.\end{Resp}
\switchcolumn*

\LA
\LessonHead{Lectio V}
\switchcolumn
\EN
\LessonHead{Lesson V}
\switchcolumn*

\LA
\noindent Fremit ergo prædo frustrátus, et in ódium religiónis, quæ talem divitiárum usum instituísset, ardéscens, direptiónem thesáuri potióris aggréditur, apud quem nullam denariórum substántiam reperísset; ut illud depósitum, quo sacrátius erat dives, auférret. Renuntiáre Christo Lauréntium jubet, et solidíssimam illam levítici ánimi fortitúdinem diris parat urgére supplíciis; quorum ubi prima nil óbtinent, vehementióra succédunt. Láceros artus et multa vérberum sectióne conscíssos subjécto prǽcipit igne torréri; ut per cratem férream, quæ jam de fervóre contínuo vim in se habéret uréndi, conversórum altérna mutatióne membrórum, fíeret cruciátus veheméntior et pœna prodúctior.\par
\switchcolumn
\EN
\noindent The baffled thief chafed, and his hatred for the godliness which had appointed such an use of riches, flaming forth, he attempted the robbery of a dearer treasure from him in whose hands he had found no coin, even to take from him that possession wherein he had holier wealth. He commanded Lawrence to deny Christ, and made ready to assail the immovable firmness of the Levite's soul with appalling tortures, of which the failure of the first was followed by the application of others more fearful still. When his limbs had been mangled and cut by many stripes, his tormentor ordered them to be roasted over a fire. He was laid on an iron grating, the bars of which by the continual fire below, became themselves burning hot, so that by turning and rearranging his limbs upon them, his agony was kept keener, and his suffering made to last longer.\par
\switchcolumn*

\LA
\begin{Resp}\Rsign Noli me derelínquere, pater sancte, quia thesáuros tuos jam expéndi. \Star{} Non ego te désero, fili, neque derelínquo; sed majóra tibi debéntur pro fide Christi certámina.\end{Resp}
\begin{Resp}\Vsign Nos quasi senes levióris pugnæ cursum recípimus, te autem quasi júvenem manet gloriósior de tyránno triúmphus: post tríduum me sequéris sacerdótem levíta.\end{Resp}
\begin{Resp}\Rsign Non ego te désero, fili, neque derelínquo; sed majóra tibi debéntur pro fide Christi certámina.\end{Resp}
\switchcolumn
\EN
\begin{Resp}\Rsign Leave me not, holy Father, for I have spent already thy treasures. \Star{} I leave thee not, my son, neither forsake thee but there awaiteth thee for Christ's truth a wrestling sterner than mine.\end{Resp}
\begin{Resp}\Vsign We, as old men, have an easier race to run; for thee in thy youth, there is kept a more glorious victory over the enemy; yet three days, and thou shalt follow me, the Deacon behind the Priest.\end{Resp}
\begin{Resp}\Rsign I leave thee not, my son, neither forsake thee but there awaiteth thee for Christ's truth a wrestling sterner than mine.\end{Resp}
\switchcolumn*

\LA
\LessonHead{Lectio VI}
\switchcolumn
\EN
\LessonHead{Lesson VI}
\switchcolumn*

\LA
\noindent Nihil óbtines, nihil próficis, sæva crudélitas. Subtráhitur torméntis tuis matéria mortális, et, Lauréntio in cælos abeúnte, tu déficis flammis tuis. Superári cáritas Christi flamma non pótuit; et ségnior fuit ignis, qui foris ussit, quam qui intus accéndit. Sævísti persecútor in Mártyrem: sævísti, et auxísti palmam, dum ággeras pœnam. Nam quid non ad victóris glóriam ingénium tuum réperit, quando in honórem transiérunt triúmphi étiam instruménta supplícii? Gaudeámus ígitur, dilectíssimi, gáudio spiritáli, et de felicíssimo íncliti viri fine gloriémur in Dómino, qui est mirábilis in Sanctis suis, in quibus nobis et præsídium constítuit et exémplum; atque ita per univérsum mundum claríficat glóriam suam, ut a solis ortu usque ad occásum, leviticórum lúminum coruscánte fulgóre, quam clarificáta est Jerosólyma Stéphano, tam illústris fíeret Roma Lauréntio.\par
\switchcolumn
\EN
\noindent Cruel savage! thou gainest nothing, and advancest nothing. That which can die passeth by degrees beyond the reach of thy tortures, and when Lawrence departeth to heaven, thou and thy fires are conquered. The love of Christ could not be overcome by the flames, and the glow that scorched the outward man was colder than that that burnt in the inner. Thou didst rage, O thou persecutor thou didst rage against the Martyr, but by making keener his agony, thou hast but made nobler his palm. What did thine imagination fail to discover that could minister to the glory of him who conquered thee, since even the means of his execution have turned to the honour of his triumph? Wherefore, dearly beloved brethren, let us rejoice with spiritual joy, and make our boast in the Lord, Who is wonderful in His Saints, (Ps. lxvii. 36,) and hath given unto us, in them an help and an example. Let us, I say, make our boast of the extraordinary happiness of the illustrious Lawrence's end, in that same Lord Who hath so glorified the name of His servant throughout the whole world, that from the rising of the sun unto the going down thereof, wheresoever the constellation of the Levitical lights shineth, even as Jerusalem is made glorious by Stephen, so Rome is made famous by Lawrence.\par
\switchcolumn*

\LA
\begin{Resp}\Rsign Beátus Lauréntius clamávit et dixit: Deum meum colo, illi soli sérvio; \Star{} Et ídeo non tímeo torménta tua.\end{Resp}
\begin{Resp}\Vsign Mea nox obscúrum non habet, sed ómnia in luce claréscunt.\end{Resp}
\begin{Resp}\Rsign Et ídeo non tímeo torménta tua.\end{Resp}
\begin{Resp}\Vsign Glória Patri, et Fílio, \Star{} et Spirítui Sancto.\end{Resp}
\begin{Resp}\Rsign Et ídeo non tímeo torménta tua.\end{Resp}
\switchcolumn
\EN
\begin{Resp}\Rsign The blessed Lawrence cried and said: My God do I worship, and Him only do I serve \Star{} And therefore I am not afraid of thy tortures.\end{Resp}
\begin{Resp}\Vsign The darkness is no darkness to me, but the night is all as clear as the morning, that shineth more and more unto the perfect day.\end{Resp}
\begin{Resp}\Rsign And therefore I am not afraid of thy tortures.\end{Resp}
\begin{Resp}\Vsign Glory be to the Father, and to the Son, \Star{} and to the Holy Ghost.\end{Resp}
\begin{Resp}\Rsign And therefore I am not afraid of thy tortures.\end{Resp}
\switchcolumn*

\LA
\LessonHead{Lectio VII}
\switchcolumn
\EN
\LessonHead{Lesson VII}
\switchcolumn*

\LA
\LessonTitle{Léctio sancti Evangélii secúndum Joánnem}
\switchcolumn
\EN
\LessonTitle{From the Holy Gospel according to John}
\switchcolumn*

\LA
\Cite{Joannes 12:24-26}
\switchcolumn
\EN
\Cite{John 12:24-26}
\switchcolumn*

\LA
\noindent In illo témpore: Dixit Jesus discípulis suis: Amen, amen dico vobis, nisi granum fruménti cadens in terram, mórtuum fúerit, ipsum solum manet. Et réliqua.\par
\switchcolumn
\EN
\noindent At that time, Jesus said unto His disciples: Amen, Amen, I say unto you, except a corn of wheat fall into the ground and die, it abideth alone. And so on.\par
\switchcolumn*

\LA
\LessonTitle{Homilía sancti Augustíni Epíscopi}
\switchcolumn
\EN
\LessonTitle{Homily by St. Augustine, Bishop (of Hippo.)}
\switchcolumn*

\LA
\Source{Tract. 51 in Joannem, sub medium}
\switchcolumn
\EN
\Source{From Tract 51 on John}
\switchcolumn*

\LA
\noindent Ipse Dóminus Jesus erat granum mortificándum et multiplicándum; mortificándum infidelitáte Judæórum, multiplicándum fide populórum. Jam vero exhórtans ad passiónis suæ sectánda vestígia, Qui amat, inquit, ánimam suam, perdet eam. Quod duóbus modis intéllegi potest. Qui amat, perdet; id est, si amas, perdes. Si cupis vitam tenére in Christo, noli mortem timére pro Christo. Item álio modo: Qui amat ánimam suam, perdet eam; noli amáre, ne perdas; noli amáre in hac vita, ne perdas in ætérna vita.\par
\switchcolumn
\EN
\noindent The Lord Jesus was Himself a corn of wheat that was to die and bring forth much fruit to die by the unbelief of the Jews, and to bring forth much fruit in the faith of the Gentiles. He, exhorting men to follow His steps, saith: He that loveth his life shall lose it. Now, these words may be understood in two ways. First he that loveth his life shall lose it, that is, If thou love life, thou wilt lose it; if thou wilt live for ever in Christ, refuse not to die for Christ. Or secondly, he that loveth his life shall lose it; love not then that which thou shalt lose; love not this present life, so that thou be thereby in jeopardy of losing life eternal.\par
\switchcolumn*

\LA
\begin{Resp}\Rsign In cratícula te Deum non negávi, \Star{} Et ad ignem applicátus te Dóminum Jesum Christum conféssus sum.\end{Resp}
\begin{Resp}\Vsign Probásti, Dómine, cor meum, et visitásti nocte.\end{Resp}
\begin{Resp}\Rsign Et ad ignem applicátus te Dóminum Jesum Christum conféssus sum.\end{Resp}
\switchcolumn
\EN
\begin{Resp}\Rsign Upon the bars I denied thee not, O God. \Star{} And when they put me to the fire, I acknowledged thee to be the Lord, O Jesus Christ.\end{Resp}
\begin{Resp}\Vsign Lord, Thou hast proved mine heart and visited it by night.\end{Resp}
\begin{Resp}\Rsign And when they put me to the fire, I acknowledged thee to be the Lord, O Jesus Christ.\end{Resp}
\switchcolumn*

\LA
\LessonHead{Lectio VIII}
\switchcolumn
\EN
\LessonHead{Lesson VIII}
\switchcolumn*

\LA
\noindent Hoc autem, quod postérius dixi, magis habére vidétur evangélicus sensus; séquitur enim: Et qui odit ánimam suam in hoc mundo, in vitam ætérnam custódit eam. Ergo, quod supra dictum est, Qui amat, subintellégitur in hoc mundo, in vitam ætérnam ipse custódit eam. Magna et mira senténtia, quemádmodum sit hóminis in ánimam suam amor ut péreat, ódium ne péreat. Si male amáveris, tunc odísti; si bene óderis, tunc amásti. Felíces, qui odérunt custodiéndo, ne perdant amándo.\par
\switchcolumn
\EN
\noindent What this second interpretation is the meaning of the Gospel, appeareth most probable from the words which follow: And he that hateth his life in this world, shall keep it unto life eternal. From which we may suppose the sense of the first words to be: He that loveth his life in this world shall lose it unto life eternal. This is a great and marvellous saying, showing how a man may so love life as to lose life, and so hate life as to keep life. If thou love it too well, then dost thou hate it; if thou hate it with an holy hatred, then dost thou love it. Blessed are they that, lest they should so love it as to lose it, so hate it as to keep it.\par
\switchcolumn*

\LA
\begin{Resp}\Rsign O Hippólyte, si credíderis in Dóminum Jesum Christum, \Star{} Et thesáuros tibi osténdo, et vitam ætérnam promítto.\end{Resp}
\begin{Resp}\Vsign Beátus Lauréntius Hippólyto dixit: Si credis in Dóminum Jesum Christum.\end{Resp}
\begin{Resp}\Rsign Et thesáuros tibi osténdo, et vitam ætérnam promítto.\end{Resp}
\begin{Resp}\Vsign Glória Patri, et Fílio, \Star{} et Spirítui Sancto.\end{Resp}
\begin{Resp}\Rsign Et thesáuros tibi osténdo, et vitam ætérnam promítto.\end{Resp}
\switchcolumn
\EN
\begin{Resp}\Rsign O Hippolytus, if thou believest in the Lord Jesus Christ, \Star{} I lay before thee treasure an hundredfold, and promise thee life everlasting.\end{Resp}
\begin{Resp}\Vsign The blessed Lawrence said unto Hippolytus: If thou believest in the Lord Jesus Christ.\end{Resp}
\begin{Resp}\Rsign I lay before thee treasure an hundred-fold, and promise thee life everlasting.\end{Resp}
\begin{Resp}\Vsign Glory be to the Father, and to the Son, \Star{} and to the Holy Ghost.\end{Resp}
\begin{Resp}\Rsign I lay before thee treasure an hundred-fold, and promise thee life everlasting.\end{Resp}
\switchcolumn*

\LA
\LessonHead{Lectio IX}
\switchcolumn
\EN
\LessonHead{Lesson IX}
\switchcolumn*

\LA
\noindent Sed vide ne tibi subrépat, ut teípsum velis interímere, sic intellegéndo, quod debes odísse in hoc mundo ánimam tuam. Hinc enim quidam malígni atque pervérsi hómines, et in seípsis crudelióres et sceleratióres homicídæ, flammis se donant, aquis se præfócant, præcipítio se collídunt, et péreunt. Hoc Christus non dócuit; immo étiam diábolo præcipítium suggerénti respóndit: Redi retro, sátana: scriptum est, Non tentábis Dóminum, Deum tuum. Petro autem dixit, signíficans qua morte clarificatúrus erat Deum: Cum esses júnior, cingébas te, et ambulábas quo volébas; cum autem senúeris, alter te cinget et feret quo tu non vis. Ubi satis expréssit, non a seípso, sed ab álio debére occídi, qui vestígia séquitur Christi.\par
\switchcolumn
\EN
\noindent Beware, lest thou take these words: He that hateth his life in this world shall keep it unto life eternal as some do, for an approval of suicide. Some evil and perverse men, bloody and guilty murderers of themselves, do indeed throw themselves into the fire, drown themselves in water, and cast themselves down precipices, and so perish. This is not the teaching of Christ, Who, when the devil would have Him cast Himself down from an high place, answered: Get thee behind Me, Satan. It is written, Thou shalt not tempt the Lord thy God. (Matth. iv. 5-7.) Who also said to Peter, signifying by what death he should glorify God: When thou wast young thou girdedst thyself and walkedst whither thou wouldest; but when thou shalt be old, another shall gird thee, and carry thee whither thou wouldest not. (John xxi. 18.) From which it is evident that he that would follow Christ's footsteps, must be slain, not by himself, but by another.\par
\switchcolumn*

\LA
\TeDeum{Te Deum}
\switchcolumn
\EN
\TeDeum{Te Deum}
\switchcolumn*

\end{paracol}

\par\needspace{10\baselineskip}\addvspace{1.2em}{\centering\small\textsc{Die 11 Augusti} · \textsc{11 August}\par}
\Office{Ss. Tiburtii et Susannæ Virginis, Martyrum}{Sts. Tiburtius and Susanna, Martyrs}{Simplex}
\addcontentsline{toc}{subsection}{Die 11 Augusti · Ss. Tiburtii et Susannæ Virginis, Martyrum \textit{— Sts. Tiburtius and Susanna, Martyrs}}
\begin{paracol}{2}
\LA
\RefLine{Lectiones I–II: de Scriptura occurrente.}
\switchcolumn
\EN
\RefLine{Lessons I–II: from the Scripture of the day.}
\switchcolumn*

\LA
\LessonHead{Lectio III (pro commemoratione)}
\switchcolumn
\EN
\LessonHead{Lesson III (when commemorated)}
\switchcolumn*

\LA
\noindent Tiburtius, Chromátii præfecti Urbis fílius, sancti Sebastiáni ópera Christianus, cum ob eam causam ad Fabianum judicem adductus esset múltaque apud illum de Christi fide prædicáret, excandéscens judex paviméntum candéntibus carbónibus sterni jubet. Mox, Tiburti, inquit, vel diis nostris sacrífices oportet, vel per istos carbónes nudis pédibus tibi incedéndum est. At ille, crucis signo se muniens fidentérque ámbulans in pruna, Disce, inquit, ex hoc solum esse Deum, quem Christiáni colunt; prunæ enim mihi flores vidéntur. Quod cum magicis artibus tribuerétur, extra Urbem ductus Tiburtius, ac via Lavicana tertio ab Urbe lápide gládio percússus, ibi a Christiánis sepelítur. Quo die Susanna virgo nobilíssima, quod Galerii Maximiáni, fílii Diocletiáni imperatóris, conjúgium recusaret, ut quæ virginitátem Deo vóverat; post multa tormentórum genera, quibus varie tentátum est sanctum Vírginis propositum, domi suæ, jussu imperatóris, gládio percússa, ad duplex virginitátis et martyrii præmium migrávit in cælum.\par
\switchcolumn
\EN
\noindent Tibutrius was son to Chromatius, Praefect of the city of Rome, and was converted to Christianity by holy Sebastian. On this account he was brought before Fabian the judge, and spake boldly in his presence many things concerning belief in Christ. Then Fabian broke out in anger and caused the pavement to be spread with live coals. Now, Tiburtius, said he, thou must either sacrifice to our gods, or walk barefoot on these coals. Tiburtius armed himself with the sign of the Cross and walked boldly on the coals. Learn from this, said he, that there is no God but He whom the Christians worship; for the coals are to me like flowers. For this (in the year 286) he was credited with art magic, led forth without the city and smitten with the sword at the third milestone on the Lavican Road, where he was buried by the Christians. On the same day, (about the year 295,) the noble maiden Susanna, having refused the offer of marriage of Galerius Maximianus, son to the Emperor Diocletian, because she had made a vow of her virginity to God, after diverse torments wherewith her holy resolution was tried, was smitten with the sword, in her own house, by order of the Emperor, and passed to heaven to receive the double reward of virginity and martyrdom.\par
\switchcolumn*

\LA
\TeDeum{Te Deum}
\switchcolumn
\EN
\TeDeum{Te Deum}
\switchcolumn*

\end{paracol}

\par\needspace{10\baselineskip}\addvspace{1.2em}{\centering\small\textsc{Die 12 Augusti} · \textsc{12 August}\par}
\Office{S. Claræ Virginis}{St. Clare, Virgin}{Duplex}
\addcontentsline{toc}{subsection}{Die 12 Augusti · S. Claræ Virginis \textit{— St. Clare, Virgin}}
\begin{paracol}{2}
\LA
\RefLine{Lectiones I–III: de Scriptura occurrente.}
\switchcolumn
\EN
\RefLine{Lessons I–III: from the Scripture of the day.}
\switchcolumn*

\LA
\RefLine{Lectiones VII–IX: ex Commúni Virginum tantum (C6a).}
\switchcolumn
\EN
\RefLine{Lessons VII–IX: from the Common of a Virgin not a Martyr (C6a).}
\switchcolumn*

\LA
\LessonHead{Lectio IV}
\switchcolumn
\EN
\LessonHead{Lesson IV}
\switchcolumn*

\LA
\noindent Clara nóbilis virgo, Assísii nata in Umbria, sanctum Francíscum concívem suum imitáta, cuncta sua bona in eleemósynas et páuperum subsídia distríbuit et convértit. De sǽculi strépitu fúgiens, in campéstrem declinávit ecclésiam, ibíque ab eódem beáto Francísco recépta tonsúra, consanguíneis ipsam redúcere conántibus fórtiter réstitit. Et dénique ad ecclésiam sancti Damiáni fuit per eúmdem addúcta, ubi ei Dóminus plures sócias aggregávit; et sic ipsa sacrárum sorórum collégium instítuit, quarum régimen, nímia sancti Francísci devícta importunitáte, recépit. Suum monastérium sollícite ac prudénter, in timóre Dómini ac plena órdinis observántia, annis quadragínta duóbus mirabíliter gubernávit; ejus enim vita erat áliis erudítio et doctrína, unde céteræ vivéndi régulam didicérunt.\par
\switchcolumn
\EN
\noindent The noble maiden Clara was born at Assisi in Umbria, (in the year 1193.) In imitation of her holy fellow-citizen Francis, she distributed all her goods among the poor and needy. She fled from the din of the world, and (on the 18th day of March, 1212,) betook herself to the Church (of St. Mary of the Angels) in the fields, where blessed Francis cut her hair. She stoutly resisted the efforts of her family to make her come back, and after a while Francis took her to the Church of St. Damian, where the Lord gathered around her several companions. Thus she founded an holy Sisterhood, which, at the earnest entreaty of holy Francis, she governed. For two-and-forty years she directed her monastery with wonderful care and wisdom in the fear of the Lord and the full keeping of the Rule. Her own life was an instruction and teaching for the rest, whence others learnt to order their own.\par
\switchcolumn*

\LA
\begin{Resp}\Rsign Propter veritátem, et mansuetúdinem, et justítiam: \Star{} Et dedúcet te mirabíliter déxtera tua.\end{Resp}
\begin{Resp}\Vsign Spécie tua et pulchritúdine tua inténde, próspere procéde, et regna.\end{Resp}
\begin{Resp}\Rsign Et dedúcet te mirabíliter déxtera tua.\end{Resp}
\switchcolumn
\EN
\begin{Resp}\Rsign Because of truth, and meekness, and righteousness; \Star{} And thy right hand shall lead thee wonderfully.\end{Resp}
\begin{Resp}\Vsign In thy comeliness, and thy beauty, go forward, fare prosperously, and reign.\end{Resp}
\begin{Resp}\Rsign And thy right hand shall lead thee wonderfully.\end{Resp}
\switchcolumn*

\LA
\LessonHead{Lectio V}
\switchcolumn
\EN
\LessonHead{Lesson V}
\switchcolumn*

\LA
\noindent Ut, carne depréssa, spíritu convalésceret, nudam humum et intérdum sarménta pro lecto habébat, et pro pulvinári sub cápite durum lignum. Una túnica cum mantéllo de vili et híspido panno utebátur, áspero cilício nonnúmquam adhíbito juxta carnem. Tanta se frenábat abstinéntia, ut longo témpore tribus in hebdómada diébus nihil pénitus pro sui córporis aliménto gustáverit; réliquis autem diébus tali se cibórum parvitáte restríngens, ut áliæ, quómodo subsístere póterat, miraréntur. Binas quotánnis (ántequam ægrotáret) quadragésimas, solo pane et aqua refécta, jejunábat. Vigíliis ínsuper et oratiónibus assídue dédita, in his præcípue dies noctésque expendébat. Diútinus perpléxa languóribus, cum ad exercítium corporále non posset súrgere per se ipsam, sorórum suffrágio levabátur; et, fulciméntis ad tergum appósitis, laborábat própriis mánibus, ne in suis étiam esset infirmitátibus otiósa. Amátrix præcípua paupertátis, ab ea pro nulla umquam necessitáte discéssit; et possessiónes pro sorórum sustentatióne a Gregório nono oblátas constantíssime recusávit.\par
\switchcolumn
\EN
\noindent What she might wax stronger in spirit by keeping the body down, she made her bed on the bare ground, sometimes with little twigs, and with hard wood for a pillow. Her dress was a gown and cloak of mean and rough cloth, and she sometimes wore hair-cloth next the skin. She bridled herself with such abstinence, that for a long time she took no bodily nourishment whatever upon three days in the week. Upon the remaining days she ate so little that the others wondered how she lived. As long as her health allowed it, she kept two Lents every year, during which she fasted upon bread and water. Moreover, she was instant in watching and prayer, wherein she chiefly spent both her days and nights. She suffered from constant illnesses, and when she could not herself rise to bodily work, she sat up with the help of the sisters, and with her back propped, worked with her hands, that she might not be idle even in the midst of her weaknesses. She was an eminent lover of poverty, from which no need ever made her swerve, and she persistently refused the possessions which were offered to the sisters by Gregory IX. for their support.\par
\switchcolumn*

\LA
\begin{Resp}\Rsign Dilexísti justítiam, et odísti iniquitátem: \Star{} Proptérea unxit te Deus, Deus tuus, óleo lætítiæ.\end{Resp}
\begin{Resp}\Vsign Propter veritátem, et mansuetúdinem, et justítiam.\end{Resp}
\begin{Resp}\Rsign Proptérea unxit te Deus, Deus tuus, óleo lætítiæ.\end{Resp}
\switchcolumn
\EN
\begin{Resp}\Rsign Thou hast loved righteousness, and hated iniquity; \Star{} Therefore God, thy God, hath anointed thee with the oil of gladness.\end{Resp}
\begin{Resp}\Vsign Because of truth, and meekness, and righteousness.\end{Resp}
\begin{Resp}\Rsign Therefore God, thy God, hath anointed thee with the oil of gladness.\end{Resp}
\switchcolumn*

\LA
\LessonHead{Lectio VI}
\switchcolumn
\EN
\LessonHead{Lesson VI}
\switchcolumn*

\LA
\noindent Multis et váriis miráculis virtus suæ sanctitátis effúlsit. Cuídam de soróribus sui monastérii loquélam restítuit expedítam; álteri aurem surdam apéruit; laborántem febre, tuméntem hydrópisi, plagátam fístula, aliásque áliis oppréssas languóribus liberávit. Fratrem de órdine Minórum ab insániæ passióne sanávit. Cum óleum in monastério totáliter defecísset, Clara accépit úrceum atque lavit, et invéntus est óleo, benefício divínæ largitátis, implétus. Uníus panis medietátem ádeo multiplicávit, ut soróribus quinquagínta suffécerit. Saracénis Assísium obsidéntibus et Claræ monastérium invádere conántibus, ægra se ad portam afférri vóluit, unáque vas in quo sanctíssimum Eucharístiæ sacraméntum erat inclúsum, ibíque orávit: Ne tradas, Dómine, béstiis ánimas confiténtes tibi, et custódi fámulas tuas, quas pretióso sánguine redemísti. In cujus oratióne ea vox audíta est: Ego vos semper custódiam. Saracéni autem partim se fugæ mandárunt, partim, qui murum ascénderant, capti óculis, præcípites cecidérunt. Ipsa dénique Virgo, cum in extrémis ágeret, a cándido beatárum Vírginum cœtu (inter quas una eminéntior ac fulgídior apparébat) visitáta, ac sacra Eucharístia sumpta, et peccatórum indulgéntia ab Innocéntio quarto ditáta, tértio Idus Augústi ánimam Deo réddidit. Post óbitum vero quam plúrimis miráculis resplendéntem, Alexánder quartus inter sanctas Vírgines rétulit.\par
\switchcolumn
\EN
\noindent The power of her holy life shone forth in many and diverse miracles. She restored the use of speech to one of the sisters in her convent for another she opened a deaf ear she healed one sick with fever, one swollen with dropsy, one troubled with an hollow oozing ulcer, and others afflicted with diverse ailments. She cured a brother of the Order of Friars Minors of raging madness. Once when all the oil in the house was spent, she took the vessel and washed it, and it was found filled with oil by the goodness of God. She multiplied half a loaf till it was enough to satisfy fifty sisters. When the Saracens (attached to the army of Frederick II.) attacked Assisi, (in the year 1239,) and were fain to break into Clare's monastery, she being sick, caused herself to be carried to the door, and likewise the vessel in which was held the Most Holy Sacrament of the Eucharist, and there she prayed, saying O Lord, deliver not unto beasts the souls of them that praise thee, but preserve thine handmaids whom Thou hast redeemed with thy Precious Blood. Whereupon a voice was heard which said I will always preserve you. Some of the Saracens took to flight, and others who had mounted the wall became blind, and fell down headlong. When Clare herself was at the point of death she beheld a white multitude of blessed Virgins, with one among them nobler and brighter than the rest. Having received the Holy Eucharist, and a Plenary Indulgence from Innocent IV., she resigned her soul to God upon the\par
12th day of August (1253.) After her death she became illustrious for very many miracles, and Alexander IV. enrolled her name among those of the Holy Virgins.\par
\switchcolumn*

\LA
\begin{Resp}\Rsign Afferéntur Regi vírgines post eam, próximæ ejus; \Star{} Afferéntur tibi in lætítia et exsultatióne.\end{Resp}
\begin{Resp}\Vsign Spécie tua et pulchritúdine tua; inténde, próspere procéde, et regna.\end{Resp}
\begin{Resp}\Rsign Afferéntur tibi in lætítia et exsultatióne.\end{Resp}
\begin{Resp}\Vsign Glória Patri, et Fílio, \Star{} et Spirítui Sancto.\end{Resp}
\begin{Resp}\Rsign Afferéntur tibi in lætítia et exsultatióne.\end{Resp}
\switchcolumn
\EN
\begin{Resp}\Rsign After her shall virgins be brought unto the King, her fellows \Star{} They shall be brought unto thee with gladness and rejoicing.\end{Resp}
\begin{Resp}\Vsign In thy comeliness and thy beauty, go forward, fare prosperously, and reign.\end{Resp}
\begin{Resp}\Rsign They shall be brought unto thee with gladness and rejoicing.\end{Resp}
\begin{Resp}\Vsign Glory be to the Father, and to the Son, \Star{} and to the Holy Ghost.\end{Resp}
\begin{Resp}\Rsign They shall be brought unto thee with gladness and rejoicing.\end{Resp}
\switchcolumn*

\end{paracol}

\par\needspace{10\baselineskip}\addvspace{1.2em}{\centering\small\textsc{Die 13 Augusti} · \textsc{13 August}\par}
\Office{Ss. Hippolyti et Cassiani Martyrum}{Sts. Hippolytus and Cassian, Martyrs}{Simplex}
\addcontentsline{toc}{subsection}{Die 13 Augusti · Ss. Hippolyti et Cassiani Martyrum \textit{— Sts. Hippolytus and Cassian, Martyrs}}
\begin{paracol}{2}
\LA
\RefLine{Lectiones I–II: de Scriptura occurrente.}
\switchcolumn
\EN
\RefLine{Lessons I–II: from the Scripture of the day.}
\switchcolumn*

\LA
\LessonHead{Lectio III (pro commemoratione)}
\switchcolumn
\EN
\LessonHead{Lesson III (when commemorated)}
\switchcolumn*

\LA
\noindent Hippolytus, a sancto Lauréntio baptizatus, domi suæ, dum Eucharistiam sumeret, comprehénsus, et ad Valerianum imperatórem adductus, ab eo de suæ religiónis professióne interrogátus, líbere se Christianum proféssus est. Quam ob rem fustibus cæditur; quibus in verbéribus cum ejus fides constantior invenirétur, munéribus et honórum promissis tentátur. Quæ cum ómnia frustra diceréntur, præfecto occidendus traditur. Qui, domum Hippolyti veniens, ut ejus facultates publicaret, totam familiam christianam esse cognoscit; atque, iis a christiana fide frustra deterritis, primum plumbátis cæsa Concordia, Hippolyti nutríce, quæ ceteros confirmábat, réliquos extra portam Tiburtinam occídi jubet. Hippolytus, indómitis equis raptátus per loca tríbulis et cárduis cónsita, lacerato corpore, spíritum Deo réddidit, unáque cum réliquis a Justino presbytero ad agrum Veranum sepúltus est. Eadem die, ad Forum Syllæ, crudelíssimo supplicio afféctus est Cassianus Martyr; qui, vinctis post terga mánibus, puerórum, quos erudiebat, férreis stylis configendus excarnificandúsque traditur. Quorum quanto erat infirmior vis, tanto ejus pœna martyrii gravior ac diuturnior, pálmaque illustrior.\par
\switchcolumn
\EN
\noindent Hippolytus was one of those baptized by St. Lawrence. He was arrested in his own house while he was taking the Holy Communion. He was brought before the Emperor Valerian, and, when he was asked by him touching his religious profession, he freely confessed that he was a Christian. Wherefore he was beaten with clubs, but when his faith was found only the bolder under the blows, he was tempted with promises of gifts and honours. Then when words were found only to be thrown away upon him, he was given over to the Prefect to be put to death. The Praefect went to the house of Hippolytus to take possession of his goods, and there found that all the household were Christians. He strove in vain to awe them into the denial of their faith, and then ordered Concordia, the nurse of Hippolytus, who was encouraging the rest, to be beaten to death with whips loaded with lead, and afterwards the others to be slain outside the gate that leadeth toward Tivoli. Hippolytus was tied to wild horses which dragged him through rough places full of briars and thistles, until with a mangled body he resigned his soul to God. Justin the Priest buried him along with the others. On the same day, at Imola, the martyr Cassian was put to a most cruel death. He was a schoolmaster, and was given up to his scholars, with his hands bound behind his back, to be stabbed and torn to death with steel pens. Owing to the weakness of the means, the suffering of his martyrdom was very grievous and long, and his palm all the more glorious.\par
\switchcolumn*

\LA
\TeDeum{Te Deum}
\switchcolumn
\EN
\TeDeum{Te Deum}
\switchcolumn*

\end{paracol}

\par\needspace{10\baselineskip}\addvspace{1.2em}{\centering\small\textsc{Die 14 Augusti} · \textsc{14 August}\par}
\Office{In Vigilia Assumptionis B.M.V.}{In Vigilia Assumptionis B.M.V.}{Simplex}
\addcontentsline{toc}{subsection}{Die 14 Augusti · In Vigilia Assumptionis B.M.V. \textit{— In Vigilia Assumptionis B.M.V.}}
\begin{paracol}{2}
\LA
\LessonHead{Lectio I}
\switchcolumn
\EN
\LessonHead{Lesson I}
\switchcolumn*

\LA
\LessonTitle{Léctio sancti Evangélii secúndum Lucam}
\switchcolumn
\EN
\LessonTitle{The Lesson is taken from the Holy Gospel according to Luke}
\switchcolumn*

\LA
\Cite{Luc 11:27-28}
\switchcolumn
\EN
\Cite{Luke 11:27-28}
\switchcolumn*

\LA
\noindent In illo témpore: Loquénte Jesu ad turbas, extóllens vocem quædam múlier de turba dixit illi: Beátus venter qui te portávit. Et réliqua.\par
\switchcolumn
\EN
\noindent At that time: It came to pass as Jesus spoke, a certain woman of the company lifted up her voice, and said unto him: Blessed is the womb that bore thee, and the paps which thou hast sucked. And so on, and that which followeth.\par
\switchcolumn*

\LA
\LessonTitle{Homilía sancti Joánnis Chrysóstomi}
\switchcolumn
\EN
\LessonTitle{Homily by St. John Chrysostom}
\switchcolumn*

\LA
\Source{In Joann. c. 2 Homilia 20, circa finem}
\switchcolumn
\EN
\Source{Homily 20 on John ch. 2, toward the end}
\switchcolumn*

\LA
\noindent Cum audíeris mulíerem illam dicéntem: Beátus venter qui te portávit, et úbera quæ tu suxísti; deínde Dóminum respondéntem: Quinímmo beáti qui áudiunt verbum Dei et custódiunt illud; ea senténtia dictum exístima, non quod Matrem neglégeret, sed nihil ei utilitátis matris nomen allatúrum osténderet, nisi bonitáte et fide præstáret. Quod si Maríæ nihil, sine virtúte, matérna cáritas erat profutúra; longe minus nobis patris, fratris, matris, fílii bónitas, nisi áliquid nostrum afferámus.\par
\switchcolumn
\EN
\noindent When you give ear to the saying of that woman: Blessed is the womb that bare thee, and the paps which thou hast sucked: and to the Lord's reply: Yea, rather, blessed are they that hear the word of God and keep it: think not that he made this observation as slighting his Mother, but as wishing to shew that it would profit her nothing to be called his Mother, unless she excelled in goodness and faith. Now, if a mother's love would avail Mary nothing without virtue, much less will it avail us to be a good father, brother, mother, or son, unless we are good in ourselves.\par
\switchcolumn*

\LA
\begin{Resp}\Rsign Dómine, Pater et Deus vitæ meæ, ne derelínquas me in cogitátu malígno: extolléntiam oculórum meórum ne déderis mihi, et desidérium malígnum avérte a me, Dómine; aufer a me concupiscéntiam, \Star{} Et ánimo irreverénti et infruníto ne tradas me, Dómine.\end{Resp}
\begin{Resp}\Vsign Ne derelínquas me, Dómine, ne accréscant ignorántiæ meæ, nec multiplicéntur delícta mea.\end{Resp}
\begin{Resp}\Rsign Et ánimo irreverénti et infruníto ne tradas me, Dómine.\end{Resp}
\switchcolumn
\EN
\begin{Resp}\Rsign O Lord, Father and God of my life, leave me not to evil counsels; give me not a proud look, but turn away from me an haughty mind, O Lord turn away from me concupiscence, \Star{} And give me not over unto an impudent and froward mind, O Lord!\end{Resp}
\begin{Resp}\Vsign Leave me not, O Lord, lest mine ignorance increase, and my sins abound.\end{Resp}
\begin{Resp}\Rsign And give me not over unto an impudent and froward mind, O Lord.\end{Resp}
\switchcolumn*

\LA
\LessonHead{Lectio II}
\switchcolumn
\EN
\LessonHead{Lesson II}
\switchcolumn*

\LA
\noindent In nullo namque álio, nisi in solis sui ipsíus virtútibus, post divínam grátiam, de salúte cuípiam sperándum est. Nam, si id profutúrum erat per se Maríæ, profuísset étiam Judǽis, quorum consanguíneus Christus secúndum carnem; profuísset civitáti, in qua natus est; profuísset frátribus. Atqui, dum fratres rerum suárum curam habuérunt, nihil eis propinquitátis nomen prófuit, sed cum réliquo mundo damnáti erant.\par
\switchcolumn
\EN
\noindent For indeed, salvation for anyone, apart from the divine grace, is to be hoped for in nothing else but his own virtues. For if her kinship in itself could have profited Mary, it would also have profited the Jews, for Christ was their kinsman according to the flesh; it would have profited the city in which he was born; it would have profited his brethren. Yet as long as his brethren cared only for their own interests, their relationship to Christ profited them nothing, but they were condemned with the rest of the world.\par
\switchcolumn*

\LA
\begin{Resp}\Rsign Magna enim sunt judícia tua, Dómine, et inenarrabília verba tua: \Star{} Magnificásti pópulum tuum et honorásti.\end{Resp}
\begin{Resp}\Vsign Transtulísti illos per Mare Rubrum et transvexísti eos per aquam nímiam.\end{Resp}
\begin{Resp}\Rsign Magnificásti pópulum tuum et honorásti.\end{Resp}
\switchcolumn
\EN
\begin{Resp}\Rsign Great are thy judgments, O Lord, and thy words cannot be expressed. \Star{} Thou didst make thy people mighty and honourable.\end{Resp}
\begin{Resp}\Vsign Thou broughtest them through the Red Sea, and leddest them through much water.\end{Resp}
\begin{Resp}\Rsign Thou didst make thy people mighty and honourable.\end{Resp}
\switchcolumn*

\LA
\LessonHead{Lectio III}
\switchcolumn
\EN
\LessonHead{Lesson III}
\switchcolumn*

\LA
\noindent Tunc autem admiratióni esse cœpérunt, quando própria claruérunt virtúte. Pátria vero nihil inde consecúta cécidit, et incéndio absúmpta est; concíves mísere interémpti periére; consanguínei secúndum carnem nil ad salútem lucráti sunt, deficiénte virtútis patrocínio. Verum Apóstoli ante omnes claríssimi evasérunt, cum se ad veram et expeténdam ejus familiaritátem consuetudinémque per obediéntiam contulérunt. Hinc intellégimus, fide semper nobis opus esse, et vita quæ virtútibus lúceat: hæc dumtáxat salvos nos fácere póterit.\par
\switchcolumn
\EN
\noindent Then only did they begin to be worthy of admiration, when they shone by their own virtue. His native land, indeed, having gained nothing from its own connection with him, fell and was burnt by fire; his fellow-citizens were put to death and perished miserably; his kindred according to the flesh gained nothing towards their salvation; insofar as all these lacked the protection of virtue. But of them all, the Apostles became the most renowned, since by obedience they joined themselves to him in a right and desirable friendship and companionship. From this we learn, that we always have need of faith, and a life shining with virtues; since this alone will have power to save us.\par
\switchcolumn*

\LA
\begin{Resp}\Rsign Quæ sunt in corde hóminum, óculi tui vident, Dómine, et in libro tuo ómnia scribéntur: \Star{} Homo videt in fácie, Deus autem in corde.\end{Resp}
\begin{Resp}\Vsign Omnia enim corda scrutátur, et univérsas méntium cogitatiónes intéllegit.\end{Resp}
\begin{Resp}\Rsign Homo videt in fácie, Deus autem in corde.\end{Resp}
\begin{Resp}\Vsign Glória Patri, et Fílio, \Star{} et Spirítui Sancto.\end{Resp}
\begin{Resp}\Rsign Homo videt in fácie, Deus autem in corde.\end{Resp}
\switchcolumn
\EN
\begin{Resp}\Rsign Lord, thine eyes behold all that is in the heart of man, and in thy book are they all written. \Star{} Man looketh on the outward appearance, but God looketh on the heart.\end{Resp}
\begin{Resp}\Vsign For He searcheth all hearts, and understandeth all the imaginations of the thoughts.\end{Resp}
\begin{Resp}\Rsign Man looketh on the outward appearance, but God looketh on the heart.\end{Resp}
\begin{Resp}\Vsign Glory be to the Father, and to the Son, \Star{} and to the Holy Ghost.\end{Resp}
\begin{Resp}\Rsign Man looketh on the outward appearance, but God looketh on the heart.\end{Resp}
\switchcolumn*

\end{paracol}

\par\needspace{10\baselineskip}\addvspace{1.2em}{\centering\small\textsc{Die 15 Augusti} · \textsc{15 August}\par}
\Office{In Assumptione Beatæ Mariæ Virginis}{In Assumptione Beatæ Mariæ Virginis}{Duplex I classis cum Octava communi}
\addcontentsline{toc}{subsection}{Die 15 Augusti · In Assumptione Beatæ Mariæ Virginis \textit{— In Assumptione Beatæ Mariæ Virginis}}
\begin{paracol}{2}
\LA
\LessonHead{Lectio I}
\switchcolumn
\EN
\LessonHead{Lesson I}
\switchcolumn*

\LA
\LessonTitle{De libro Génesis}
\switchcolumn
\EN
\LessonTitle{Lesson from the book of Genesis}
\switchcolumn*

\LA
\Cite{Gen 3:9-15}
\switchcolumn
\EN
\Cite{Gen 3:9-15}
\switchcolumn*

\LA
\noindent Vocavítque Dóminus Deus Adam, et dixit ei: Ubi es? Qui ait: Vocem tuam audívi in paradiso; et tímui eo quod nudus essem, et abscóndi me. Cui dixit: Quis enim indicávit tibi quod nudus esses, nisi quod ex ligno de quo præcéperam tibi ne coméderes, comedísti? Dixítque Adam: Múlier quam dedísti mihi sóciam, dedit mihi de ligno, et comédi. Et dixit Dóminus Deus ad mulíerem: Quare hoc fecísti? Quæ respóndit: Serpens decépit me, et comédi. Et ait Dóminus Deus ad serpéntem: Quia fecísti hoc, maledíctus es inter ómnia animántia et béstias terræ; super pectus tuum gradiéris, et terram cómedes cunctis diébus vitæ tuæ. Inimicítias ponam inter te et mulíerem, et semen tuum et semen illíus; ipsa cónteret caput tuum, et tu insidiáberis calcáneo ejus.\par
\switchcolumn
\EN
\noindent And the Lord God called Adam, and said to him: Where art thou? And he said: I heard thy voice in paradise; and I was afraid, because I was naked, and I hid myself. And he said to him: And who hath told thee that thou wast naked, but that thou hast eaten of the tree whereof I commanded thee that thou shouldst not eat? And Adam said: The woman, whom thou gavest me to be my companion, gave me of the tree, and I did eat. And the Lord God said to the woman: Why hast thou done this? And she answered: The serpent deceived me, and I did eat. And the Lord God said to the serpent: Because thou hast done this thing, thou art cursed among all cattle, and beasts of the earth: upon thy breast shalt thou go, and earth shalt thou eat all the days of thy life. I will put enmities between thee and the woman, and thy seed and her seed: she shall crush thy head, and thou shalt lie in wait for her heel.\par
\switchcolumn*

\LA
\begin{Resp}\Rsign Vidi speciósam sicut colúmbam, ascendéntem désuper rivos aquárum, cujus inæstimábilis odor erat nimis in vestiméntis ejus; \Star{} Et sicut dies verni circúmdabant eam flores rosárum et lília convállium.\end{Resp}
\begin{Resp}\Vsign Quæ est ista quæ ascéndit per desértum sicut vírgula fumi ex aromátibus myrrhæ et thuris?\end{Resp}
\begin{Resp}\Rsign Et sicut dies verni circúmdabant eam flores rosárum et lília convállium.\end{Resp}
\switchcolumn
\EN
\begin{Resp}\Rsign I saw her, when, fair like a dove, she winged her flight above the rivers of waters. The priceless savour of her perfumes hung heavy in her garments. \Star{} And about her it was as the flower of roses in the spring of the year, and lilies of the valleys.\end{Resp}
\begin{Resp}\Vsign Who is this that cometh out of the wilderness like a pillar of smoke, perfumed with myrrh and frankincense?\end{Resp}
\begin{Resp}\Rsign And about her it was as the flower of roses in the spring of the year, and lilies of the valleys.\end{Resp}
\switchcolumn*

\LA
\LessonHead{Lectio II}
\switchcolumn
\EN
\LessonHead{Lesson II}
\switchcolumn*

\LA
\LessonTitle{De Epístola prima beáti Pauli Apóstoli ad Corínthios}
\switchcolumn
\EN
\LessonTitle{From the First Epistle of Blessed Paul the Apostle to the Corinthians}
\switchcolumn*

\LA
\Cite{1 Cor 15:20-26}
\switchcolumn
\EN
\Cite{1 Cor 15:20-26}
\switchcolumn*

\LA
\noindent Christus resurréxit a mórtuis primítiæ dormiéntium, quóniam quidem per hóminem mors, et per hóminem resurréctio mortuórum. Et sicut in Adam omnes moriúntur, ita et in Christo omnes vivificabúntur; unusquísque autem in suo órdine: primítiæ Christus; deínde ii qui sunt Christi, qui in advéntu ejus credidérunt. Deínde finis, cum tradíderit regnum Deo et Patri, cum evacuáverit omnem principátum, et potestátem, et virtútem. Opórtet autem illum regnáre, donec ponat omnes inimícos sub pédibus ejus. Novíssima autem inimíca destruétur mors.\par
\switchcolumn
\EN
\noindent But now Christ is risen from the dead, the firstfruits of them that sleep: For by a man came death, and by a man the resurrection of the dead. And as in Adam all die, so also in Christ all shall be made alive. But every one in his own order: the firstfruits Christ, then they that are of Christ, who have believed in his coming. Afterwards the end, when he shall have delivered up the kingdom to God and the Father, when he shall have brought to nought all principality, and power, and virtue. For he must reign, until he hath put all his enemies under his feet. And the enemy death shall be destroyed last.\par
\switchcolumn*

\LA
\begin{Resp}\Rsign Sicut cedrus exaltáta sum in Líbano, et sicut cypréssus in monte Sion: quasi myrrha elécta, \Star{} Dedi suavitátem odóris.\end{Resp}
\begin{Resp}\Vsign Et sicut cinnamómum et bálsamum aromatízans.\end{Resp}
\begin{Resp}\Rsign Dedi suavitátem odóris.\end{Resp}
\switchcolumn
\EN
\begin{Resp}\Rsign I was exalted like a cedar in Lebanon, and as a cypress-tree upon Mount Zion. Like the best myrrh \Star{} I yielded a pleasant odour.\end{Resp}
\begin{Resp}\Vsign Like cinnamon and sweet balsam.\end{Resp}
\begin{Resp}\Rsign I yielded a pleasant odour.\end{Resp}
\switchcolumn*

\LA
\LessonHead{Lectio III}
\switchcolumn
\EN
\LessonHead{Lesson III}
\switchcolumn*

\LA
\Cite{1 Cor 15:53-57}
\switchcolumn
\EN
\Cite{1 Cor 15:53-57}
\switchcolumn*

\LA
\noindent Opórtet enim corruptíbile hoc indúere incorruptiónem, et mortále hoc indúere immortalitátem. Cum autem mortále hoc indúerit immortalitátem, tunc fiet sermo qui scriptus est: Absórpta est mors in victória. Ubi est, mors, victória tua? ubi est, mors, stímulus tuus? Stímulus autem mortis peccátum est, virtus vero peccáti lex. Deo autem grátias, qui dedit nobis victóriam per Dóminum Jesum Christum.\par
\switchcolumn
\EN
\noindent For this corruptible must put on incorruption; and this mortal must put on immortality. And when this mortal hath put on immortality, then shall come to pass the saying that is written: Death is swallowed up in victory. O death, where is thy victory? O death, where is thy sting? Now the sting of death is sin: and the power of sin is the law. But thanks be to God, who hath given us the victory through our Lord Jesus Christ.\par
\switchcolumn*

\LA
\begin{Resp}\Rsign Quæ est ista quæ procéssit sicut sol, et formósa tamquam Jerúsalem? \Star{} Vidérunt eam fíliæ Sion, et beátam dixérunt, et regínæ laudavérunt eam.\end{Resp}
\begin{Resp}\Vsign Et sicut dies verni circúmdabant eam flores rosárum et lília convállium.\end{Resp}
\begin{Resp}\Rsign Vidérunt eam fíliæ Sion, et beátam dixérunt, et regínæ laudavérunt eam.\end{Resp}
\begin{Resp}\Vsign Glória Patri, et Fílio, \Star{} et Spirítui Sancto.\end{Resp}
\begin{Resp}\Rsign Vidérunt eam fíliæ Sion, et beátam dixérunt, et regínæ laudavérunt eam.\end{Resp}
\switchcolumn
\EN
\begin{Resp}\Rsign Who is this that cometh up like the sun? This, comely as Jerusalem? \Star{} The daughters of Zion saw her, and called her blessed: the queens also, and they praised her.\end{Resp}
\begin{Resp}\Vsign And about her it was as the flower of roses in the spring of the year, and lilies of the valleys.\end{Resp}
\begin{Resp}\Rsign The daughters of Zion saw her, and called her blessed the queens also, and they praised her.\end{Resp}
\begin{Resp}\Vsign Glory be to the Father, and to the Son, \Star{} and to the Holy Ghost.\end{Resp}
\begin{Resp}\Rsign The daughters of Zion saw her, and called her blessed the queens also, and they praised her.\end{Resp}
\switchcolumn*

\LA
\LessonHead{Lectio IV}
\switchcolumn
\EN
\LessonHead{Lesson IV}
\switchcolumn*

\LA
\LessonTitle{Sermo sancti Joánnis Damascéni}
\switchcolumn
\EN
\LessonTitle{From the Sermons of St. John of Damascus.}
\switchcolumn*

\LA
\Source{Oratio 2 de Dormit. B. M. V., post init.}
\switchcolumn
\EN
\Source{2nd on the Falling-asleep of Blessed Mary.}
\switchcolumn*

\LA
\noindent Hódie sacra et animáta arca Dei vivéntis, quæ suum in útero concépit Creatórem, requiéscit in templo Dómini, quod nullis est exstrúctum mánibus. Et David exsúltat ejus parens, et cum eo choros ducunt Angeli, célebrant Archángeli, Virtútes gloríficant, Principátus exsúltant, Potestátes collætántur, gaudent Dominatiónes, Throni festum diem agunt, laudant Chérubim, glóriam ejus prǽdicant Séraphim. Hódie Eden novi Adam paradísum súscipit animátum, in quo solúta est condemnátio, in quo plantátum est lignum vitæ, in quo opérta fuit nostra núditas.\par
\switchcolumn
\EN
\noindent This day the holy and animated Ark of the living God, which had held within it its own Maker, is borne to rest in that Temple of the Lord, which is not made with hands. David, whence it sprang, leapeth before it, and in company with him the Angels dance, the Archangels sing aloud, the Virtues ascribe glory, the Princedoms shout for joy, the Powers make merry, the Lordships rejoice, the Thrones keep holiday, the Cherubim utter praise, and the Seraphim proclaim its glory. This day the Eden of the new Adam receiveth the living garden of delight, wherein the condemnation was annulled, wherein the Tree of Life was planted, wherein our nakedness was covered.\par
\switchcolumn*

\LA
\begin{Resp}\Rsign Ornátam monílibus fíliam Jerúsalem Dóminus concupívit: \Star{} Et vidéntes eam fíliæ Sion, beatíssimam prædicavérunt, dicéntes: Unguéntum effúsum nomen tuum.\end{Resp}
\begin{Resp}\Vsign Astitit regína a dextris tuis in vestítu deauráto, circúmdata varietáte.\end{Resp}
\begin{Resp}\Rsign Et vidéntes eam fíliæ Sion, beatíssimam prædicavérunt, dicéntes: Unguéntum effúsum nomen tuum.\end{Resp}
\switchcolumn
\EN
\begin{Resp}\Rsign When the Lord beheld the daughter of Jerusalem adorned with her jewels, He greatly desired her beauty \Star{} And when the daughters of Zion saw her, they cried out that she was most blessed, saying thy name is as ointment poured forth.\end{Resp}
\begin{Resp}\Vsign Upon thy right hand did stand the Queen in a vesture of gold wrought about with diverse colours.\end{Resp}
\begin{Resp}\Rsign And when the daughters of Zion saw her, they cried out that she was most blessed, saying thy name is as ointment poured forth.\end{Resp}
\switchcolumn*

\LA
\LessonHead{Lectio V}
\switchcolumn
\EN
\LessonHead{Lesson V}
\switchcolumn*

\LA
\noindent Hódie Virgo immaculáta, quæ nullis terrénis inquináta est afféctibus, sed cæléstibus educáta cogitatiónibus, non in terram revérsa est; sed, cum esset animátum cælum, in cæléstibus tabernáculis collocátur. Ex qua enim ómnibus vera vita manávit, quómodo illa mortem gustáret? Sed cedit legi latæ ab eo quem génuit; et, ut fília véteris Adam, véterem senténtiam súbiit (nam et ejus Fílius, qui est vita ipsa, eam non recusávit); ut autem Dei vivéntis Mater, ad illum ipsum digne assúmitur.\par
\switchcolumn
\EN
\noindent This day the stainless maiden, who had been defiled by no earthly lust, but ennobled by heavenly desires, returned not to dust, but, being herself a living heaven, took her place among the heavenly mansions. From her true life had flowed for all men, and how should she taste of death? But she yielded obedience to the law established by Him to Whom she had given birth, and, as the daughter of the old Adam, underwent the old sentence, which even her Son, Who is the very Life Itself, had not refused; but, as the Mother of the living God, she was worthily taken by Him unto Himself.\par
\switchcolumn*

\LA
\begin{Resp}\Rsign Beátam me dicent omnes generatiónes, \Star{} Quia fecit mihi Dóminus magna qui potens est, et sanctum nomen ejus.\end{Resp}
\begin{Resp}\Vsign Et misericórdia ejus a progénie in progénies timéntibus eum.\end{Resp}
\begin{Resp}\Rsign Quia fecit mihi Dóminus magna qui potens est, et sanctum nomen ejus.\end{Resp}
\switchcolumn
\EN
\begin{Resp}\Rsign All generations shall call me blessed. \Star{} For He that is mighty, even the Lord, hath done to me great things; and Holy is His Name.\end{Resp}
\begin{Resp}\Vsign And his mercy is on them that fear Him, from generation to generation.\end{Resp}
\begin{Resp}\Rsign He that is mighty, even the Lord, hath done to me great things; and Holy is His Name.\end{Resp}
\switchcolumn*

\LA
\LessonHead{Lectio VI}
\switchcolumn
\EN
\LessonHead{Lesson VI}
\switchcolumn*

\LA
\noindent Ex Actis Pii Papæ duodécimi\par
\switchcolumn
\EN
\noindent From the Acts of Pope Pius XII\par
\switchcolumn*

\LA
\LessonTitle{Quóniam vero univérsa Ecclésia fidem in corpóream beátæ Maríæ Vírginis Assumptiónem per sæculórum decúrsum manifestávit, et totíus orbis Epíscopi prope unánimi consensióne petiérunt ut hæc véritas, quæ Sacris Lítteris innítitur, Christifidélium ánimis pénitus est ínsita, ceterísque revelátis veritátibus plane cónsona, tamquam divínæ et cathólicæ fídei dogma definirétur, Pius duodécimus Póntifex Máximus, totíus Ecclésiæ votis ánnuens, státuit hoc Beátæ Maríæ Vírginis privilégium solémniter renuntiáre. Itaque die prima Novémbris anni máximi Jubilǽi millésimi nongentésimi quinquagésimi, Romæ, in foro ad sancti Petri Basílicam paténte, plurimórum Sanctæ Románæ Ecclésiæ Cardinálium atque Episcopórum ex díssitis étiam regiónibus astánte cœtu, coram ingénti Christifidélium multitúdine, univérso cathólico orbe plaudénte, corpóream Beátæ Maríæ Vírginis Assumptiónem in cælum infallíbili oráculo in hæc verba proclamávit: Postquam súpplices étiam atque étiam ad Deum admóvimus preces, ac Veritátis Spíritus lumen invocávimus, ad Omnipoténtis Dei glóriam, qui peculiárem benevoléntiam suam Maríæ Vírgini dilargítus est, ad sui Fílii honórem, immortális sæculórum Regis ac peccáti mortísque victóris, ad ejúsdem augústæ Matris augéndam glóriam et ad totíus Ecclésiæ gáudium exsultationémque, auctoritáte Dómini Nostri Jesu Christi, Beatórum Apostolórum Petri et Pauli ac Nostra pronuntiámus, declarámus et definímus revelátum dogma esse: Immaculátam Deíparam semper Vírginem Maríam, expléto terréstris vitæ cursu, fuísse córpore et ánima ad cæléstem glóriam assúmptam.}
\switchcolumn
\EN
\LessonTitle{Since indeed the universal Church hath at all times and throughout the ages manifested faith in the bodily Assumption of the Blessed Virgin Mary, and since the Bishops of the whole world by an almost unanimous agreement have petitioned that this truth, which is enshrined in Sacred Scripture and deeply rooted in the souls of Christ's faithful, and is also truly in accord with other revealed truths, should be defined as a dogma of the divine and Catholic Faith, Pope Pius XII, acceding to the requests of the whole Church, decreed that this privilege of the Blessed Virgin Mary be solemnly proclaimed, and thus, on the first day of November of the year of the Great Jubilee, nineteen hundred and fifty, at Rome, in the open square before the Basilica of St. Peter, surrounded by a throng of many Cardinals and Bishops of the Holy Roman Church who had come from distant parts of the earth, and before a great multitude of the faithful, with the whole Catholic world rejoicing, proclaimed in these words and with infallible statement the bodily Assumption of the Blessed Virgin Mary into heaven: Wherefore, having offered to God continual prayers of supplication, and having invoked the light of the Spirit of Truth, to the glory of Almighty God who hath enriched the Virgin Mary with his special favour; in honour of his Son, the immortal King of ages and victor over sin and death; for the increase of the glory of the same august Mother, and for the joy and exultation of the whole Church, by the authority of Our Lord Jesus Christ, of the holy Apostles Peter and Paul, and by our own authority, we pronounce, declare and define it to be a divinely revealed dogma that: The Immaculate Mother of God, Mary ever Virgin, was, at the end of her earthly life, assumed body and soul into heavenly glory.}
\switchcolumn*

\LA
\begin{Resp}\Rsign Beáta es, Virgo María, quæ Dóminum portásti, Creatórem mundi: \Star{} Genuísti qui te fecit, et in ætérnum pérmanes Virgo.\end{Resp}
\begin{Resp}\Vsign Ave, María, grátia plena; Dóminus tecum.\end{Resp}
\begin{Resp}\Rsign Genuísti qui te fecit, et in ætérnum pérmanes Virgo.\end{Resp}
\begin{Resp}\Vsign Glória Patri, et Fílio, \Star{} et Spirítui Sancto.\end{Resp}
\begin{Resp}\Rsign Genuísti qui te fecit, et in ætérnum pérmanes Virgo.\end{Resp}
\switchcolumn
\EN
\begin{Resp}\Rsign Blessed art thou, O Virgin Mary, who hast carried the Lord, the Maker of the world. \Star{} Thou hast borne Him Who created thee, and thou abidest a virgin for ever.\end{Resp}
\begin{Resp}\Vsign Hail, Mary, full of grace. the Lord is with thee.\end{Resp}
\begin{Resp}\Rsign Thou hast borne Him Who created thee, and thou abidest a virgin for ever.\end{Resp}
\begin{Resp}\Vsign Glory be to the Father, and to the Son, \Star{} and to the Holy Ghost.\end{Resp}
\begin{Resp}\Rsign Thou hast borne Him Who created thee, and thou abidest a virgin for ever.\end{Resp}
\switchcolumn*

\LA
\LessonHead{Lectio VII}
\switchcolumn
\EN
\LessonHead{Lesson VII}
\switchcolumn*

\LA
\LessonTitle{Léctio sancti Evangélii secúndum Lucam}
\switchcolumn
\EN
\LessonTitle{From the Holy Gospel according to Luke}
\switchcolumn*

\LA
\Cite{Luc 1:41-50}
\switchcolumn
\EN
\Cite{Luke 1:41-50}
\switchcolumn*

\LA
\noindent In illo témpore: Repléta est Spíritu Sancto Elísabeth et exclamávit voce magna, et dixit: Benedícta tu inter mulíeres. Et réliqua.\par
\switchcolumn
\EN
\noindent At that time: Elisabeth was filled with the Holy Ghost: and she spoke out with a loud voice, and said, Blessed art thou among women. And so on.\par
\switchcolumn*

\LA
\LessonTitle{Homilía sancti Petri Canísii Presbýteri}
\switchcolumn
\EN
\LessonTitle{A Homily by St. Peter Canisius, the Priest}
\switchcolumn*

\LA
\Source{De Maria Deipara Virgine, lib. 5, c. 6}
\switchcolumn
\EN
\Source{De Maria Deipara Virgine, lib. 5, c. 6}
\switchcolumn*

\LA
\noindent Férias Deíparæ sacras, ut frequénter, ita reverénter perágit Ecclésia, certo intéllegens, Deo gratum, et christifidélibus dignum, hoc esse offícium, quo inter beátos et beátas omnes beatíssima Dómini et Dei nostri mater, certis anni tempóribus ac públicis cæremóniis sæpenúmero celebrátur. Inter has autem férias, quæ tot sǽculis religióse háctenus observántur, festum Assumptiónis máximum reputátur, locúmque príncipem tenet. Nec dies quidem ullus Maríæ lǽtior feliciórque cóntigit, si novam et ánimæ et córporis felicitátem eo die illi collátam rite contemplémur: ideóque si ántea unquam, tunc vel máxime spíritus et ánima et corpus ejus in Deum vivum mirabíliter exsultávit, ípsaque summo jure dícere pótuit: Respéxit humilitátem ancíllæ suæ: Ecce enim ex hoc beátam me dicent omnes generatiónes: quia fecit mihi magna qui potens est.\par
\switchcolumn
\EN
\noindent The Church frequently and reverently keepeth feastdays dedicated to the Mother of God, realizing that it is a work pleasing to God and worthy of the faithful if many feasts, with fixed dates and public ceremonies, are celebrated in honour of the most blessed of all the blessed in heaven, the Mother of our Lord and God. Among all these feasts which have been celebrated so devotedly for so many years, even unto the present day, the Feast of the Assumption is considered the greatest and holdeth chief place. Indeed there was no happier or more joyful day for Mary, if we duly consider the happiness of both body and soul granted to her on that day. Then especially, as never before, her spirit, soul and body rejoiced wondrously in the living God and she could rightfully say: He hath regarded the lowliness of his handmaiden; for behold, all generations shall call me blessed; for he that is mighty hath magnified me.\par
\switchcolumn*

\LA
\begin{Resp}\Rsign Diffúsa est grátia in lábiis tuis: \Star{} Proptérea benedíxit te Deus in ætérnum.\end{Resp}
\begin{Resp}\Vsign Myrrha, et gutta, et cásia a vestiméntis tuis, a dómibus ebúrneis, ex quibus delectavérunt te fíliæ regum in honóre tuo.\end{Resp}
\begin{Resp}\Rsign Proptérea benedíxit te Deus in ætérnum.\end{Resp}
\switchcolumn
\EN
\begin{Resp}\Rsign Grace is poured into thy lips \Star{} Therefore God hath blessed thee for ever.\end{Resp}
\begin{Resp}\Vsign thy garments smell of myrrh, and aloes, and cassia, out of the ivory palaces whereby kings' daughters among thy honourable women have made thee glad.\end{Resp}
\begin{Resp}\Rsign Therefore God hath blessed thee for ever.\end{Resp}
\switchcolumn*

\LA
\LessonHead{Lectio VIII}
\switchcolumn
\EN
\LessonHead{Lesson VIII}
\switchcolumn*

\LA
\noindent Quare nos étiam, qui te Filiúmque tuum amámus, o ter beáta et vere augústa Mater, non póssumus huic præcláræ et incomparábili felicitáti tuæ non ex ánimo gratulári, quandóquidem quæcúmque a Dómino tibi et de te dicta sunt, tam pulchro éxitu concludúntur, ac omni ex parte tandem perficiúntur. Beáta quæ non modo credidísti, sed fídei étiam omnísque virtútis fructum et finem hódie consecúta es, et ejus quem tantópere amásti atque desiderásti, jucundíssimo conspéctu nunc demum pérfrui meruísti. Emmanuélem in castéllum mundi hujus intrántem, tamquam hóspitem, hóspita excepísti: hódie ab illo vicíssim in regále suum palátium excíperis, et tamquam mater tali Salomóne digna magnífice primum honoráris.\par
\switchcolumn
\EN
\noindent O thrice blessed and truly august Mother, it is for this reason that we who love thee and thy Son cannot refrain from congratulating thee with all sincerity upon thine admirable and incomparable happiness, especially since everything that hath been said to thee and about thee by the Lord, is brought to a conclusion by thy beautiful passing away from this life, and in every wise hath been perfectly fulfilled. Blessed art thou who hast not only believed but hast this day attained unto the end of faith and the fruit of all virtue, and now at last hast merited to enjoy the most pleasing sight of him whom thou didst love and desire so greatly. Thyself a guest, thou didst receive Emmanuel who as a guest did enter into thee, as into a mighty fortress in this world; and today, thou in turn art received by him into his royal mansion, and magnificently welcomed with the highest honour, as befitteth one found worthy to be the Mother of such a Solomon.\par
\switchcolumn*

\LA
\begin{Resp}\Rsign Beáta es Virgo María, Dei Génetrix, quæ credidísti Dómino: perfécta sunt in te quæ dicta sunt tibi: ecce exaltáta es super choros Angelórum: \Star{} Intercéde pro nobis ad Dóminum, Deum nostrum.\end{Resp}
\begin{Resp}\Vsign Ave, María, grátia plena, Dóminus tecum.\end{Resp}
\begin{Resp}\Rsign Intercéde pro nobis ad Dóminum, Deum nostrum.\end{Resp}
\begin{Resp}\Vsign Glória Patri, et Fílio, \Star{} et Spirítui Sancto.\end{Resp}
\begin{Resp}\Rsign Intercéde pro nobis ad Dóminum, Deum nostrum.\end{Resp}
\switchcolumn
\EN
\begin{Resp}\Rsign O Virgin Mary, Mother of God, blessed art thou that didst believe the Lord, for there hath been a performance of those things which were told thee from the Lord. Behold, thou art exalted over choirs of Angels. \Star{} Plead for us with the Lord our God.\end{Resp}
\begin{Resp}\Vsign Hail, Mary, full of grace, the Lord is with thee.\end{Resp}
\begin{Resp}\Rsign Plead for us with the Lord our God.\end{Resp}
\begin{Resp}\Vsign Glory be to the Father, and to the Son, \Star{} and to the Holy Ghost.\end{Resp}
\begin{Resp}\Rsign Plead for us with the Lord our God.\end{Resp}
\switchcolumn*

\LA
\LessonHead{Lectio IX}
\switchcolumn
\EN
\LessonHead{Lesson IX}
\switchcolumn*

\LA
\noindent Felix dies, qui tam pretiósum munus de hujus mundi desérto transmísit, et ad sanctam æternámque civitátem ita perdúxit, ut commúne et incredíbile gáudium, nec minor admirátio cælítibus ómnibus crearétur. Felix dies, qui sponsæ languéntis tam ardéntibus quam diutúrnis desidériis satisfécit, ut quod quæsíverat inveníret, quod petíverat accíperet, quod exspectárat secúra possidéret, nimírum in illa summi æterníque boni perfécta visióne atque fruitióne pénitus conquiéscens. Felix dies, qui ancíllam Dómini humíllimam eo usque provéxit ac éxtulit, ut præstantíssima cæli Regína, mundíque Dómina efficerétur, et ne áltius quidem posset assúrgere, quóniam in Regni sólio sublimáta, post Christum gloriósa resédit. Felix hic planéque venerándus dies est, qui Regínam et Matrem nobis poténtem simul et cleméntem in regno Dei constítuit ac confirmávit, ut eámdem, quæ Júdicis mater assídue manet, matrem misericórdiæ nobis favéntem, apud Christum patrocinántem, nostrǽque salútis negótia fidéliter procurántem habeámus.\par
\switchcolumn
\EN
\noindent O blessed day which sent so precious a gift from the desert of this world, and carried it to the holy and eternal city, so that universal and unheard of joy no less than admiration welled up in all the blessed in heaven. O blessed day, that fulfilled the long and ardent yearning of the gentle spouse, so that she might find what she had sought, that she might receive what she asked; that what she awaited she might possess securely, resting safely at last in that perfect vision and inward joy of the eternal and all-great Goodness. O blessed day which raised up and so highly exalted this most humble handmaiden of the Lord that she might become the most glorious Queen of Heaven and the mistress of the world. Indeed she could not have risen to more sublime heights since she had been elevated to the very Throne of the heavenly kingdom, and thus was established in glory next after Christ. O blessed and truly honourable is this day which constituted and confirmed for us a Queen and Mother who is at once powerful and merciful in the kingdom of God, that we might have her, who ever remaineth the Mother of the Judge, as a Mother of mercy protecting us and interceding for us with Christ, unceasingly watching over the work of our salvation.\par
\switchcolumn*

\LA
\TeDeum{Te Deum}
\switchcolumn
\EN
\TeDeum{Te Deum}
\switchcolumn*

\end{paracol}

\par\needspace{10\baselineskip}\addvspace{1.2em}{\centering\small\textsc{Die 16 Augusti} · \textsc{16 August}\par}
\Office{S. Joachim Confessoris, Patris B. M. V.}{St. Joachim, Confessor, Father of the Blessed Virgin Mary}{Duplex II classis}
\addcontentsline{toc}{subsection}{Die 16 Augusti · S. Joachim Confessoris, Patris B. M. V. \textit{— St. Joachim, Confessor, Father of the Blessed Virgin Mary}}
\begin{paracol}{2}
\LA
\RefLine{Lectiones I–III: ex Commúni Confessoris non pontificis (C5).}
\switchcolumn
\EN
\RefLine{Lessons I–III: from the Common of a Confessor not a Bishop (C5).}
\switchcolumn*

\LA
\LessonHead{Lectio IV}
\switchcolumn
\EN
\LessonHead{Lesson IV}
\switchcolumn*

\LA
\LessonTitle{Sermo sancti Epiphánii Epíscopi}
\switchcolumn
\EN
\LessonTitle{From the Sermons of St. Epiphanius, Bishop (of Pavia)}
\switchcolumn*

\LA
\Source{Orat. de Laud. Virg. sub init.}
\switchcolumn
\EN
\Source{Discourse on the Praises of the Virgin}
\switchcolumn*

\LA
\noindent De radíce Jesse ortus est rex David, et de tribu regis David sancta Virgo: sancta, inquam, et sanctórum virórum fília, cujus paréntes fuérunt Jóachim et Anna; qui quidem in vita sua Deo placuérunt, atque étiam fructum ejúsmodi germinavérunt, sanctam Vírginem Maríam, templum simul et matrem Dei. Jóachim porro, Anna et María, hi tres Trinitáti palam sacrifícium laudis offerébant. Jóachim enim interpretátur Præparátio Dómini, eo quod ex illo præparátum sit templum Dómini, nempe Virgo. Anna rursus simíliter Grátia interpretátur, proptérea quod Jóachim et Anna grátiam accepérunt, ut accedéntibus précibus, talem fructum germinárent, sanctam Vírginem adépti; Jóachim síquidem precabátur in monte, et Anna in horto suo.\par
\switchcolumn
\EN
\noindent From the Root of Jesse sprang King David, and from the stock of King David, the Holy Virgin. Holy I call her, and the daughter of holy men. Her father and mother were Joachim and Anne, who pleased God in their lives, and brought forth an offspring well pleasing to Him, even the Holy Virgin Mary, at once the Temple and the Mother of God. These three, Joachim, Anne, and Mary, clearly offered up unto the Trinity a sacrifice of praise. For the name Joachim being interpreted, signifieth “the preparation of the Lord”, and out of him was prepared the Temple of the Lord, namely, the Virgin. The name Anne signifieth grace, and she and Joachim did indeed receive a grace when, in answer to their prayers, they generated such an offspring, compassing the Holy Virgin. Joachim prayed upon the mountain and Anne in her garden.\par
\switchcolumn*

\LA
\begin{Resp}\Rsign Honéstum fecit illum Dóminus, et custodívit eum ab inimícis, et a seductóribus tutávit illum: \Star{} Et dedit illi claritátem ætérnam.\end{Resp}
\begin{Resp}\Vsign Justum dedúxit Dóminus per vias rectas, et osténdit illi regnum Dei.\end{Resp}
\begin{Resp}\Rsign Et dedit illi claritátem ætérnam.\end{Resp}
\switchcolumn
\EN
\begin{Resp}\Rsign The Lord made him honourable, and defended him from his enemies, and kept him safe from those that lay in wait for him; \Star{} And gave him perpetual glory.\end{Resp}
\begin{Resp}\Vsign He went down with him into the pit, and left him not in bonds.\end{Resp}
\begin{Resp}\Rsign And gave him perpetual glory.\end{Resp}
\switchcolumn*

\LA
\LessonHead{Lectio V}
\switchcolumn
\EN
\LessonHead{Lesson V}
\switchcolumn*

\LA
\LessonTitle{Sermo sancti Joánnis Damascéni}
\switchcolumn
\EN
\LessonTitle{From the Sermons of St. John of Damascus}
\switchcolumn*

\LA
\Source{Oratio 1 de Virg. Mariæ Nativit., circa principium}
\switchcolumn
\EN
\Source{1st on the Birth of the Blessed Virgin Mary}
\switchcolumn*

\LA
\noindent Quóniam futúrum erat ut Dei Génetrix et Virgo ex Anna orirétur, natúra grátiæ fœtum antevértere mínime ausa est; verum tantísper exspectávit, dum grátia fructum suum produxísset. Síquidem oportébat eam primogénitam in lucem edi, quæ rerum ómnium conditárum Primogénitum, in quo ómnia coagmentáta sunt, paritúra erat. O par beátum Jóachim et Anna! vobis omnis creatúra obstrícta est. Per vos enim donum, ómnium donórum præstantíssimum, Creatóri óbtulit, nempe castam matrem, quæ sola Creatóre digna erat.\par
\switchcolumn
\EN
\noindent Since it was to be that the Virgin Mother of God should be born of Anne, nature dared not to produce any other child before this child of grace, but humbly waited until grace should have produced her's. It behoved that she should come into the world as a first-born, who was to bear the First-born of every creature, even Him by Whom all things were made. O blessed couple, Joachim and Anne unto you is all creation laid under debt, since through you creation hath offered to the Creator this noblest of gifts, namely, that chaste mother, who alone was worthy of the Creator.\par
\switchcolumn*

\LA
\begin{Resp}\Rsign Amávit eum Dóminus, et ornávit eum: stolam glóriæ índuit eum, \Star{} Et ad portas paradísi coronávit eum.\end{Resp}
\begin{Resp}\Vsign Índuit eum Dóminus lorícam fídei, et ornávit eum.\end{Resp}
\begin{Resp}\Rsign Et ad portas paradísi coronávit eum.\end{Resp}
\switchcolumn
\EN
\begin{Resp}\Rsign The Lord loved him and beautified him; He clothed him with a robe of glory, \Star{} And crowned him at the gates of Paradise.\end{Resp}
\begin{Resp}\Vsign The Lord hath put on him the breast-plate of faith, and hath adorned him.\end{Resp}
\begin{Resp}\Rsign And crowned him at the gates of Paradise.\end{Resp}
\switchcolumn*

\LA
\LessonHead{Lectio VI}
\switchcolumn
\EN
\LessonHead{Lesson VI}
\switchcolumn*

\LA
\noindent Exsúlta, Jóachim, quóniam ex fília tua Fílius natus est nobis; et vocátur nomen ejus magni consílii, hoc est, salútis totíus mundi Angelus. Pudóre afficiátur Nestórius, ac manum ori impónat. Puer hic Deus est. Quonam ígitur modo ea Dei Génetrix non sit, quæ péperit? Si quis sanctam Dei Genetrícem non confitétur, a Deitáte remótus est. Mea non est hæc orátio, quamquam alióqui mea; hanc enim diviníssimam hereditátem a theólogo patre Gregório accépi. O beátum par Jóachim et Anna! Ac profécto ex ventris vestri fructu immaculáti agnoscímini, quemádmodum Christus quodam loco dixit: Ex frúctibus eórum cognoscétis eos. Ut Deo gratum erat ac dignum ea quæ a vobis orta est, vitæ vestræ ratiónes instituístis. Caste enim ac sancte múnere vestro functi, virginitátis thesáurum produxístis.\par
\switchcolumn
\EN
\noindent Rejoice, O Joachim, from whose daughter a Child hath unto us been born, and His name is called The Angel of the Great Counsel, that is, of the salvation of the whole world. Let Nestorius be ashamed, and put his hand upon his mouth. Her Child is God. How then can His Mother be other than Mother of God? Whosoever acknowledgeth not the Holy Mother of God is far from God. This saying is not mine, although it is mine in all other senses than that of authority. I have received it as a most godly legacy from Father Gregory the Divine. O blessed couple, Joachim and Anne. Christ saith in a certain place: By their fruits ye shall know them (Matth. iii. 20) and ye are known by the fruit of your chaste loins. That that which should be born of you might be worthy and well-pleasing in the sight of God, ye ordered your own lives by rule. In the chaste and holy exercise of your natural gift, ye produced the treasure of virginity.\par
\switchcolumn*

\LA
\begin{Resp}\Rsign Iste homo perfécit ómnia quæ locútus est ei Deus, et dixit ad eum: Ingrédere in réquiem meam: \Star{} Quia te vidi justum coram me ex ómnibus géntibus.\end{Resp}
\begin{Resp}\Vsign Iste est, qui contémpsit vitam mundi, et pervénit ad cæléstia regna.\end{Resp}
\begin{Resp}\Rsign Quia te vidi justum coram me ex ómnibus géntibus.\end{Resp}
\begin{Resp}\Vsign Glória Patri, et Fílio, \Star{} et Spirítui Sancto.\end{Resp}
\begin{Resp}\Rsign Quia te vidi justum coram me ex ómnibus géntibus.\end{Resp}
\switchcolumn
\EN
\begin{Resp}\Rsign This is he which did according unto all that God commanded him; and God said unto him: Enter thou into My rest, \Star{} For thee have I seen righteous before Me among all people.\end{Resp}
\begin{Resp}\Vsign This is he which loved not his life in this world, and is come unto an everlasting kingdom.\end{Resp}
\begin{Resp}\Rsign For thee have I seen righteous before Me among all people.\end{Resp}
\begin{Resp}\Vsign Glory be to the Father, and to the Son, \Star{} and to the Holy Ghost.\end{Resp}
\begin{Resp}\Rsign For thee have I seen righteous before Me among all people.\end{Resp}
\switchcolumn*

\LA
\LessonHead{Lectio VII}
\switchcolumn
\EN
\LessonHead{Lesson VII}
\switchcolumn*

\LA
\LessonTitle{Léctio sancti Evangélii secúndum Matthǽum}
\switchcolumn
\EN
\LessonTitle{From the Holy Gospel according to Matthew}
\switchcolumn*

\LA
\Cite{Matt 1:1-16}
\switchcolumn
\EN
\Cite{Matt 1:1-16}
\switchcolumn*

\LA
\noindent Liber generatiónis Jesu Christi, fílii David, fílii Abraham. Abraham génuit Isaac, Isaac autem génuit Jacob. Et réliqua.\par
\switchcolumn
\EN
\noindent The book of the generation of Jesus Christ, the son of David, the son of Abraham. Abraham begat Isaac, and Isaac begat Jacob. And so on.\par
\switchcolumn*

\LA
\LessonTitle{Homilía sancti Joánnis Damascéni}
\switchcolumn
\EN
\LessonTitle{Homily by St. John of Damascus}
\switchcolumn*

\LA
\Source{Lib. 4 de Fide orthodoxa, cap. 15 de Domini genealogia et sanctae Dei Genetricis}
\switchcolumn
\EN
\Source{Bk. iv on the Orthodox Faith}
\switchcolumn*

\LA
\noindent Quod Joseph ex davídica tribu oríginem dúxerit, sanctíssimi Evangelístæ Matthǽus et Lucas líquido demonstrárunt. Verum, hoc inter eos discríminis est, quod Matthǽus ex Davíde per Salomónem Joséphum dedúcit; Lucas autem per Nathan. At vero sanctæ Vírginis ortum utérque siléntio prætériit. Quocírca scire óperæ prétium est, nec apud Hebrǽos nec apud Scriptúram sacram hoc in more pósitum fuísse, ut mulíerum genus recenserétur. Verum, hoc demum lege cautum erat ne tribus ulla uxóres ex áltera tribu accérseret. Ac proínde Joseph, qui ex tribu davídica ortum trahébat justitiámque colébat (hanc enim laudem ei tríbuit divínum Evangélium), sanctam Vírginem haudquáquam præter legis præscríptum despondísset, nisi ex eódem sceptro genus duxísset. Ob idque satis hábuit Evangelísta demonstrásse, unde Joseph ortum tráxerit.\par
\switchcolumn
\EN
\noindent That Joseph sprang from the lineage of David, the most holy Evangelists Matthew and Luke have clearly shown. There is this difference between them, that Matthew traceth the pedigree from David through Solomon, and Luke through Nathan. But both of them pass in silence over the descent of the Holy Virgin. In explanation of this we shall find on investigation that among the Jews, and in the Holy Scriptures, it hath never been in use to chronicle the pedigrees of women. But the Law containeth a warning against the tribes intermarrying one with another. Joseph was of the same tribe as David, (namely, Judah,) and since he was a just man (this is the praise which the Gospel of God giveth him,) he would not have espoused the Holy Virgin, unless she had been of the same race, as such an union would not have been in accordance with the commandment of the law. For this reason the Evangelist held it enough to have shown the descent of Joseph.\par
\switchcolumn*

\LA
\begin{Resp}\Rsign Iste est qui ante Deum magnas virtútes operátus est, et de omni corde suo laudávit Dóminum: \Star{} Ipse intercédat pro peccátis ómnium populórum.\end{Resp}
\begin{Resp}\Vsign Ecce homo sine queréla, verus Dei cultor, ábstinens se ab omni ópere malo, et pérmanens in innocéntia sua.\end{Resp}
\begin{Resp}\Rsign Ipse intercédat pro peccátis ómnium populórum.\end{Resp}
\switchcolumn
\EN
\begin{Resp}\Rsign This is he which wrought great wonders before God, and praised the Lord with all his heart. \Star{} May he pray for all people, that their sins may be forgiven unto them.\end{Resp}
\begin{Resp}\Vsign Behold a man without blame, a worshipper of God in truth, keeping himself clean from every evil work, and abiding still in his innocency.\end{Resp}
\begin{Resp}\Rsign May he pray for all people, that their sins may be forgiven unto them.\end{Resp}
\switchcolumn*

\LA
\LessonHead{Lectio VIII}
\switchcolumn
\EN
\LessonHead{Lesson VIII}
\switchcolumn*

\LA
\noindent Igitur ex stirpe Nathan, fílii David, Levi génuit Melchi et Panthérem. Panther autem génuit Barpanthérem (namque ita vocabátur). Barpánther rursus génuit Jóachim. Jóachim dénique génuit sanctam Dei Genetrícem. Rursus, ex stirpe Salomónis, fílii David, Mathan ex uxóre sua génuit Jacob. Mórtuo autem Mathan, Melchi ex tribu Nathan, fílius Levi ac frater Panthéris, uxórem ipsíus Mathan, quæ étiam Jacóbi mater erat, matrimónio sibi copulávit, atque ex ea génuit Heli. Ita uteríni fratres erant Jacob et Heli; ille nimírum ex tribu Salomónis, hic ex tribu Nathan oriúndus.\par
\switchcolumn
\EN
\noindent Therefore, from the stock of Nathan the son of David, Levi begat Melchi and Panther and Panther begat Bar-Panther (for so was he called) and Bar-Panther begat Joachim and Joachim begat the Holy Mother of God. Again, from the stock of Solomon the son of David, Mathan of his wife begat Jacob. And after Mathan was dead, Melchi, of the family of Nathan, son to Levi and brother to Panther, took to him in marriage her that had been the wife of Mathan, and was mother of Jacob, and begat of her Heli. So that Jacob and Heli were half-brothers, sons of the same mother, but one of the stock of Solomon and the other of the stock of Nathan.\par
\switchcolumn*

\LA
\begin{Resp}\Rsign Sint lumbi vestri præcíncti, et lucérnæ ardéntes in mánibus vestris: \Star{} Et vos símiles homínibus exspectántibus dóminum suum, quando revertátur a núptiis.\end{Resp}
\begin{Resp}\Vsign Vigiláte ergo, quia nescítis qua hora Dóminus vester ventúrus sit.\end{Resp}
\begin{Resp}\Rsign Et vos símiles homínibus exspectántibus dóminum suum, quando revertátur a núptiis.\end{Resp}
\begin{Resp}\Vsign Glória Patri, et Fílio, \Star{} et Spirítui Sancto.\end{Resp}
\begin{Resp}\Rsign Et vos símiles homínibus exspectántibus dóminum suum, quando revertátur a núptiis.\end{Resp}
\switchcolumn
\EN
\begin{Resp}\Rsign Let your loins be girded about, and your lights burning; \Star{} And ye yourselves like unto men that wait for their lord, when he will return from the wedding.\end{Resp}
\begin{Resp}\Vsign Watch therefore, for ye know not what hour your Lord doth come.\end{Resp}
\begin{Resp}\Rsign And ye yourselves like unto men that wait for their lord, when he will return from the wedding.\end{Resp}
\begin{Resp}\Vsign Glory be to the Father, and to the Son, \Star{} and to the Holy Ghost.\end{Resp}
\begin{Resp}\Rsign And ye yourselves like unto men that wait for their lord, when he will return from the wedding.\end{Resp}
\switchcolumn*

\LA
\LessonHead{Lectio IX}
\switchcolumn
\EN
\LessonHead{Lesson IX}
\switchcolumn*

\LA
\noindent Porro Heli, qui ex tribu Nathan erat, nullis líberis suscéptis, e vita migrávit; eáque de causa Jacob ipsíus frater, qui ex tribu Salomónis flúxerat, ipsíus uxórem accépit, fratríque suo semen excitávit ac Joséphum progénuit. Joseph ítaque natúra quidem fílius Jacob erat, a Salomóne oriúndus; legis autem ratióne patrem habébat Heli, ex Nathan oriúndum. Quæ cum ita sint, Jóachim lectíssimam illam ac summis láudibus dignam mulíerem Annam matrimónio sibi copulávit. Verum, quemádmodum prisca illa Anna, cum sterilitátis morbo laboráret, per oratiónem ac promissiónem, Samuélem procreávit; eódem modo hæc étiam, per obsecratiónem et promissiónem, Dei Genetrícem a Deo accépit, ut ne hic quoque cuíquam ex illústribus matrónis céderet. Itaque grátia (nam hoc sonat Annæ vocábulum) Dóminam parit (id enim Maríæ nómine significátur). Vere étenim rerum ómnium conditárum Dómina facta est, cum Creatóris Mater éxstitit.\par
\switchcolumn
\EN
\noindent Then Heli, who was the descendant of Nathan, died, having had no children. Whereupon, Jacob, who was the descendant of Solomon, took to himself the widow of his half-brother Heli, to raise up seed unto his brother, and begat of her Joseph. So Joseph was by nature the son of Jacob the descendant of Solomon, but in the eye of the Law he had for his father Heli the descendant of Nathan. Things being thus, Joachim took to wife that most eminent and praiseworthy woman, Anne. And even as the patient Hannah, being stricken with barrenness, by prayer and promise became the mother of Samuel, so likewise this woman also through prayer and promise received from God the Mother of God, that in fruitfulness she might not be behind any of the famous matrons. And thus grace (for such is the signification of the name of Anne) is mother of the Lady (for such is the signification of the name of Mary.) And indeed she became the Lady of every creature, since she hath been mother of the Creator.\par
\switchcolumn*

\LA
\TeDeum{Te Deum}
\switchcolumn
\EN
\TeDeum{Te Deum}
\switchcolumn*

\end{paracol}

\par\needspace{10\baselineskip}\addvspace{1.2em}{\centering\small\textsc{Die 17 Augusti} · \textsc{17 August}\par}
\Office{S. Hyacinthi Confessoris}{St. Hyacinth, Confessor}{Duplex}
\addcontentsline{toc}{subsection}{Die 17 Augusti · S. Hyacinthi Confessoris \textit{— St. Hyacinth, Confessor}}
\begin{paracol}{2}
\LA
\RefLine{Lectiones I–III: de Scriptura occurrente.}
\switchcolumn
\EN
\RefLine{Lessons I–III: from the Scripture of the day.}
\switchcolumn*

\LA
\RefLine{Lectiones VII–VIII: ex Commúni Confessoris non pontificis (C5).}
\switchcolumn
\EN
\RefLine{Lessons VII–VIII: from the Common of a Confessor not a Bishop (C5).}
\switchcolumn*

\LA
\RefLine{Lectio IX: de In Octava S. Laurentii Martyris (08-17t).}
\switchcolumn
\EN
\RefLine{Lesson IX: of Octave of St. Lawrence, Martyr (08-17t).}
\switchcolumn*

\LA
\LessonHead{Lectio IV}
\switchcolumn
\EN
\LessonHead{Lesson IV}
\switchcolumn*

\LA
\noindent Hyacínthus Polónus, nobílibus et christiánis paréntibus, in Camiénsi villa episcopátus Vratislaviénsis natus est. A puerítia litteris instructus, post datam jurisprudéntiæ et sacris litteris operam, inter canonicos Cracovienses adscitus, insígni morum pietáte et summa eruditióne ceteros antecelluit. Romæ in Prædicatórum ordinem ab ipso institutore sancto Dominico adscriptus, perféctam vivéndi ratiónem, quam ab ipso didicerat, usque ad finem vitæ sanctíssime retinuit. Virginitátem perpetuo coluit; modestiam, patiéntiam, humilitátem, abstinéntiam ceterasque virtútes, ut certum reliogiosæ vitæ patrimónium, adamávit.\par
\switchcolumn
\EN
\noindent Hyacinth was a Pole, and was born (in the year 1185) of the noble and Christian family (of the Counts of Odrowatz,) in the town of Camien, in the diocese of Wratislaw. He was trained up in learning from his youth, and after studying law and theology, became a Canon of Crakow, where he was eminent above his fellows by the singular godliness of his life and the depth of his learning. Being at Rome, (in 1218,) he was received into the Order of Friars Preachers by the Founder, St. Dominic, himself, and kept in holiness to the end of his life the rule of perfect living which he had learnt from him. He remained always a virgin, and loved modesty, long-suffering, lowliness, self-restraint, and all other good graces as his heritage in the life of a Friar.\par
\switchcolumn*

\LA
\begin{Resp}\Rsign Honéstum fecit illum Dóminus, et custodívit eum ab inimícis, et a seductóribus tutávit illum: \Star{} Et dedit illi claritátem ætérnam.\end{Resp}
\begin{Resp}\Vsign Justum dedúxit Dóminus per vias rectas, et osténdit illi regnum Dei.\end{Resp}
\begin{Resp}\Rsign Et dedit illi claritátem ætérnam.\end{Resp}
\switchcolumn
\EN
\begin{Resp}\Rsign The Lord made him honourable, and defended him from his enemies, and kept him safe from those that lay in wait for him; \Star{} And gave him perpetual glory.\end{Resp}
\begin{Resp}\Vsign He went down with him into the pit, and left him not in bonds.\end{Resp}
\begin{Resp}\Rsign And gave him perpetual glory.\end{Resp}
\switchcolumn*

\LA
\LessonHead{Lectio V}
\switchcolumn
\EN
\LessonHead{Lesson V}
\switchcolumn*

\LA
\noindent Caritáte in Deum fervens, integras sæpe noctes fundéndis précibus castigandóque córpori insúmens, nullum eidem levaméntum, nisi lapidi innixus sive humi cubans, adhibebat. Remissus in pátriam, Frisaci primum in itinere amplíssimum sui ordinis monastérium, mox Cracoviæ álterum eréxit. Inde per alias Poloniæ regni provincias aliis quatuor exædificatis, incredibile dictu est quantum verbi Dei prædicatióne et vitæ innocéntia apud omnes profécerit. Nullum diem prætermisit, quo non præclára aliqua fidei, pietátis atque innocéntiæ arguménta præstiterit.\par
\switchcolumn
\EN
\noindent In the heat of his love for God, he sometimes passed whole nights in pouring forth prayers and chastising his body, to which he never gave rest but in leaning against a stone or lying upon the ground. He was sent back to his own country, and, on the way, founded a very large house of his Order at Friesach and soon afterwards another at Crakow. In other provinces of the kingdom of Poland he founded four others, and it passeth belief what success he had with all kinds of men, by his preaching of the Word of God, and the innocency of his life. Not a day passed wherein he did not display some bright gift of faith, godliness, or innocency.\par
\switchcolumn*

\LA
\begin{Resp}\Rsign Amávit eum Dóminus, et ornávit eum: stolam glóriæ índuit eum, \Star{} Et ad portas paradísi coronávit eum.\end{Resp}
\begin{Resp}\Vsign Índuit eum Dóminus lorícam fídei, et ornávit eum.\end{Resp}
\begin{Resp}\Rsign Et ad portas paradísi coronávit eum.\end{Resp}
\switchcolumn
\EN
\begin{Resp}\Rsign The Lord loved him and beautified him; He clothed him with a robe of glory, \Star{} And crowned him at the gates of Paradise.\end{Resp}
\begin{Resp}\Vsign The Lord hath put on him the breast-plate of faith, and hath adorned him.\end{Resp}
\begin{Resp}\Rsign And crowned him at the gates of Paradise.\end{Resp}
\switchcolumn*

\LA
\LessonHead{Lectio VI}
\switchcolumn
\EN
\LessonHead{Lesson VI}
\switchcolumn*

\LA
\noindent Sanctíssime viri studium erga proximórum salútem maximis Deus miraculis illustrávit. Inter quæ illud insigne, quod Vándalum fluvium prope Visogradum aquis redundántem, nullo navigio usus, trajecit, sociis quoque expanso super undas pallio traductis. Admirábili vitæ genere ad quadragínta prope annos post professiónem perducto, mortis die suis frátribus prænuntiáto, ipso assumptæ Vírginis festo, Horis canonicis persolutis, sacramentis ecclesiásticis summa cum veneratióne perceptis, iis verbis: In manus tuas, Dómine; spíritum Deo réddidit, anno salútis millesimo ducentésimo quinquagesimo septimo. Quem miraculis, étiam post óbitum, illustrem, Clemens Papa octavus in Sanctórum númerum rétulit.\par
\switchcolumn
\EN
\noindent The zeal of this most holy man for the salvation of his neighbors was that which God marked by His greatest miracles. Among these is famous the time when coming to the River Vistula near Wisgrade, and finding it in flood, he crossed it without a boat, drawing over also his three companions standing upon the waves upon his outspread mantle. He led a wonderful life for nearly forty years after his profession, and then foretold to his brethren the day of his death. Upon the very day of the Feast of the Assumption of the Virgin, he finished the recitation of the Office of the Church, received the Sacraments with the utmost reverence, and then with the words, Into thy hands, O Lord, gave up his soul to God in the year of salvation 1257. He was illustrious for miracles even after his death, and Pope Clement VIII. numbered him among the saints.\par
\switchcolumn*

\LA
\begin{Resp}\Rsign Iste homo perfécit ómnia quæ locútus est ei Deus, et dixit ad eum: Ingrédere in réquiem meam: \Star{} Quia te vidi justum coram me ex ómnibus géntibus.\end{Resp}
\begin{Resp}\Vsign Iste est, qui contémpsit vitam mundi, et pervénit ad cæléstia regna.\end{Resp}
\begin{Resp}\Rsign Quia te vidi justum coram me ex ómnibus géntibus.\end{Resp}
\begin{Resp}\Vsign Glória Patri, et Fílio, \Star{} et Spirítui Sancto.\end{Resp}
\begin{Resp}\Rsign Quia te vidi justum coram me ex ómnibus géntibus.\end{Resp}
\switchcolumn
\EN
\begin{Resp}\Rsign This is he which did according unto all that God commanded him; and God said unto him: Enter thou into My rest, \Star{} For thee have I seen righteous before Me among all people.\end{Resp}
\begin{Resp}\Vsign This is he which loved not his life in this world, and is come unto an everlasting kingdom.\end{Resp}
\begin{Resp}\Rsign For thee have I seen righteous before Me among all people.\end{Resp}
\begin{Resp}\Vsign Glory be to the Father, and to the Son, \Star{} and to the Holy Ghost.\end{Resp}
\begin{Resp}\Rsign For thee have I seen righteous before Me among all people.\end{Resp}
\switchcolumn*

\end{paracol}

\par\needspace{10\baselineskip}\addvspace{1.2em}{\centering\small\textsc{Die 17 Augusti} · \textsc{17 August}\par}
\Office{In Octava S. Laurentii Martyris}{Octave of St. Lawrence, Martyr}{Simplex}
\addcontentsline{toc}{subsection}{Die 17 Augusti · In Octava S. Laurentii Martyris \textit{— Octave of St. Lawrence, Martyr}}
\begin{paracol}{2}
\LA
\LessonHead{Lectio IV}
\switchcolumn
\EN
\LessonHead{Lesson IV}
\switchcolumn*

\LA
\LessonTitle{Commemoratio: In Octava S. Laurentii Martyris}
\switchcolumn
\EN
\LessonTitle{Commemoration of: Octave of St. Lawrence, Martyr}
\switchcolumn*

\LA
\Source{Tractatus 27. in fine.}
\switchcolumn
\EN
\Source{(21th.)}
\switchcolumn*

\LA
\noindent Ex Tractatu sancti Augustíni Epíscopi in Joánnem.\par
\switchcolumn
\EN
\noindent From one of the Tracts upon John, written by St. Augustine, Bishop (of Hippo.)\par
\switchcolumn*

\LA
Certi sumus, fratres, quia omnes, qui sumus in corpore Domini, et manemus in illo, ut et ipse maneat in nobis, in hoc sæculo necesse habemus usque in finem inter malos vivere; non inter illos dico malos, qui blasphemant Christum; raro enim jam inveniuntur, qui lingua blasphemant, sed multi, qui vita. Necesse est ergo, ut inter illos usque in finem vivamus. Sed quid est quod ait, Qui manet in me, et ego in illo? Quid, nisi quod Martyres audiebant: Qui perseveraverit usque in finem, hic salvus erit?\par
\switchcolumn
\EN
Brethren, we know that all we who are in the body of the Lord, and abide in Him that He also may abide in us, must needs in this world live even unto the end among the wicked; by the wicked I mean not such as blaspheme Christ, for it is now but few who are found who blaspheme Him with their tongue, but many there be who do it by their lives. Among such we must needs live, even unto the end. But what is it that He said, He that abideth in Me and I in him? (John xv. 5.) Is it not that whereunto the Martyrs gave ear, He that endureth to the end shall be saved? (Matth. x. 22.)\par
\switchcolumn*

\LA
Quomodo mansit in illo sanctus Laurentius, cujus hodie festum diem celebramus? Mansit usque ad tentationem, mansit usque ad tyrannicam interrogationem, mansit usque ad acerrimam comminationem, mansit usque ad peremptionem. Parum est: usque ad immanem excruciationem mansit. Non enim occisus est cito, sed cruciatus est in igne. Diu vivere permissus est: immo non diu vivere permissus est, sed tarde mori compulsus est. In illa ergo longa morte, in illis tormentis, quia bene manducaverat, et bene biberat, tamquam illa esca saginatus, et illo calice ebrius, tormenta non sensit.\par
\switchcolumn
\EN
In what wise abode in Him the holy Lawrence whose Feast-day we keep to-day? He abode in Him until trial came; he abode in Him until He was browbeaten by the tyrant; he abode in Him until his fearful sentence came; he abode in Him even unto death. That is saying but a little of him; he abode in Him even in savage torture. He was not slain quickly, but tortured in the fire. He was allowed to remain long alive, and yet not allowed to remain long alive, but slowly put to death. During that long death, in all that agony, since he had well eaten and well drunk, filled with that Food, and drunken with that Cup, he felt not the torture.\par
\switchcolumn*

\LA
Dixerat autem illi Xystus Martyr sanctus, cujus diem quinto abhinc retro die celebravimus: Noli merere fili: (episcopus enim erat ille, iste diaconus) noli merere, inquit, sequeris me post triduum. Triduum autem dixit medium inter diem passionis sancti Xysti, et diem hodiernam passionis sancti Laurentii. Triduum est medium. O consolatio! Non ait, Noli merere fili, desinet persecutio, et securus eris: sed, Noli merere: quo ego præcedo, tu sequeris; nec consecutio tua differtur. Triduum medium erit, et mecum eris. Accepit oraculum, vicit diabolum, pervenit ad triumphum.\par
\switchcolumn
\EN
Nor the holy Martyr Xystus, whose day we kept five days ago, had said unto him: Be not sorrowful, my son (for he was the Bishop and Lawrence the Deacon,) be not sorrowful, said he, yet three days and thou shalt follow me. These three days of which he spake were the three clear days which elapse between the day of the passion of holy Xystus, and this day, the day of the passion of holy Lawrence. The three days are those between. O what a strange comfort! He saith not: Be not sorrowful, my son; the persecution will come to an end, and thou wilt be safe, but Be not sorrowful, whither I go, thither shalt thou follow me, and the time will not be long. Three clear days shall pass, and then thou shalt be with me. Lawrence received the prophecy, conquered the devil, and entered into his triumph.\par
\switchcolumn*

\LA
\TeDeum{Te Deum}
\switchcolumn
\EN
\TeDeum{Te Deum}
\switchcolumn*

\end{paracol}

\par\needspace{10\baselineskip}\addvspace{1.2em}{\centering\small\textsc{Die 18 Augusti} · \textsc{18 August}\par}
\Office{Quarta die infra Octavam S. Assumptionis Beatæ Mariæ Virginis}{Quarta die infra Octavam S. Assumptionis Beatæ Mariæ Virginis}{Semiduplex}
\addcontentsline{toc}{subsection}{Die 18 Augusti · Quarta die infra Octavam S. Assumptionis Beatæ Mariæ Virginis \textit{— Quarta die infra Octavam S. Assumptionis Beatæ Mariæ Virginis}}
\begin{paracol}{2}
\LA
\RefLine{Lectio IX: de S. Agapiti Martyris (08-18r).}
\switchcolumn
\EN
\RefLine{Lesson IX: of St. Agapitus, Martyr (08-18r).}
\switchcolumn*

\LA
\LessonHead{Lectio I}
\switchcolumn
\EN
\LessonHead{Lesson I}
\switchcolumn*

\LA
\LessonTitle{De Canticis canticorum}
\switchcolumn
\EN
\LessonTitle{From the Song of Songs}
\switchcolumn*

\LA
\Cite{Cant 4:1-4}
\switchcolumn
\EN
\Cite{Song 4:1-4}
\switchcolumn*

\LA
\noindent (Sponsus) Quam pulchra es, amica mea ! quam pulchra es ! Oculi tui columbarum, absque eo quod intrinsecus latet. Capilli tui sicut greges caprarum quæ ascenderunt de monte Galaad. Dentes tui sicut greges tonsarum quæ ascenderunt de lavacro; omnes gemellis fœtibus, et sterilis non est inter eas. Sicut vitta coccinea labia tua, et eloquium tuum dulce. Sicut fragmen mali punici, ita genæ tuæ, absque eo quod intrinsecus latet. Sicut turris David collum tuum, quæ ædificata est cum propugnaculis; mille clypei pendent ex ea, omnis armatura fortium.\par
\switchcolumn
\EN
\noindent thy teeth as flocks of sheep, that are shorn which come up from the washing, all with twins, and there is none barren among them. thy lips are as a scarlet lace: and thy speech sweet. thy cheeks are as a piece of a pomegranate, besides that which lieth hid within. thy neck, is as the tower of David, which is built with bulwarks: a thousand bucklers hang upon it, all the armour of valiant men.\par
How beautiful art thou, my love, how beautiful art thou! thy eyes are doves' eyes, besides what is hid within. thy hair is as flocks of goats, which Come up from mount Galaad.\par
\switchcolumn*

\LA
\begin{Resp}\Rsign Vidi speciósam sicut colúmbam, ascendéntem désuper rivos aquárum, cujus inæstimábilis odor erat nimis in vestiméntis ejus; \Star{} Et sicut dies verni circúmdabant eam flores rosárum et lília convállium.\end{Resp}
\begin{Resp}\Vsign Quæ est ista quæ ascéndit per desértum sicut vírgula fumi ex aromátibus myrrhæ et thuris?\end{Resp}
\begin{Resp}\Rsign Et sicut dies verni circúmdabant eam flores rosárum et lília convállium.\end{Resp}
\switchcolumn
\EN
\begin{Resp}\Rsign I saw her, when, fair like a dove, she winged her flight above the rivers of waters. The priceless savour of her perfumes hung heavy in her garments. \Star{} And about her it was as the flower of roses in the spring of the year, and lilies of the valleys.\end{Resp}
\begin{Resp}\Vsign Who is this that cometh out of the wilderness like a pillar of smoke, perfumed with myrrh and frankincense?\end{Resp}
\begin{Resp}\Rsign And about her it was as the flower of roses in the spring of the year, and lilies of the valleys.\end{Resp}
\switchcolumn*

\LA
\LessonHead{Lectio II}
\switchcolumn
\EN
\LessonHead{Lesson II}
\switchcolumn*

\LA
\Cite{Cant 4:7-10}
\switchcolumn
\EN
\Cite{Song 4:7-10}
\switchcolumn*

\LA
\noindent Tota pulchra es, amíca mea, et mácula non est in te. Veni de Líbano, sponsa mea: veni de Líbano, veni, coronáberis: de cápite Amana, de vértice Sanir et Hermon, de cubílibus leónum, de móntibus pardórum. Vulnerásti cor meum, soror mea, sponsa; vulnerásti cor meum in uno oculórum tuórum, et in uno crine colli tui. Quam pulchræ sunt mammæ tuæ, soror mea sponsa! pulchrióra sunt úbera tua vino, et odor unguentórum tuórum super ómnia arómata.\par
\switchcolumn
\EN
\noindent Thou art all fair, O my love, and there is not a spot in thee. Come from Libanus, my spouse, come from Libanus, come: thou shalt be crowned from the top of Amana, from the top of Sanir and Hermon, from the dens of the lions, from the mountains of the leopards. Thou hast wounded my heart, my sister, my spouse, thou hast wounded my heart with one of thy eyes, and with one hair of thy neck. How beautiful are thy breasts, my sister, my spouse! thy breasts are more beautiful than wine, and the sweet smell of thy ointments above all aromatical spices.\par
\switchcolumn*

\LA
\begin{Resp}\Rsign Sicut cedrus exaltáta sum in Líbano, et sicut cypréssus in monte Sion: quasi myrrha elécta, \Star{} Dedi suavitátem odóris.\end{Resp}
\begin{Resp}\Vsign Et sicut cinnamómum et bálsamum aromatízans.\end{Resp}
\begin{Resp}\Rsign Dedi suavitátem odóris.\end{Resp}
\switchcolumn
\EN
\begin{Resp}\Rsign I was exalted like a cedar in Lebanon, and as a cypress-tree upon Mount Zion. Like the best myrrh \Star{} I yielded a pleasant odour.\end{Resp}
\begin{Resp}\Vsign Like cinnamon and sweet balsam.\end{Resp}
\begin{Resp}\Rsign I yielded a pleasant odour.\end{Resp}
\switchcolumn*

\LA
\LessonHead{Lectio III}
\switchcolumn
\EN
\LessonHead{Lesson III}
\switchcolumn*

\LA
\Cite{Cant 4:11-15}
\switchcolumn
\EN
\Cite{Song 4:11-15}
\switchcolumn*

\LA
\noindent Favus distíllans lábia tua, sponsa; mel et lac sub lingua tua: et odor vestimentórum tuórum sicut odor thuris. Hortus conclúsus soror mea, sponsa, hortus conclúsus, fons signátus. Emissiónes tuæ paradísus malórum punicórum, cum pomórum frúctibus, cypri cum nardo. Nardus et crocus, fístula et cinnamómum, cum univérsis lignis Líbani; myrrha et áloe, cum ómnibus primis unguéntis. Fons hortórum, púteus aquárum vivéntium, quæ fluunt ímpetu de Líbano.\par
\switchcolumn
\EN
\noindent thy lips, my spouse, are as a dropping honeycomb, honey and milk are under thy tongue; and the smell of thy garments, as the smell of frankincense. My sister, my spouse, is a garden enclosed, a garden enclosed, a fountain sealed up. thy plants are a paradise of pomegranates with the fruits of the orchard. Cypress with spikenard. Spikenard and saffron, sweet cane and cinnamon, with all the trees of Libanus, myrrh and aloes with all the chief perfumes. The fountain of gardens: the well of living waters, which run with a strong stream from Libanus.\par
\switchcolumn*

\LA
\begin{Resp}\Rsign Quæ est ista quæ procéssit sicut sol, et formósa tamquam Jerúsalem? \Star{} Vidérunt eam fíliæ Sion, et beátam dixérunt, et regínæ laudavérunt eam.\end{Resp}
\begin{Resp}\Vsign Et sicut dies verni circúmdabant eam flores rosárum et lília convállium.\end{Resp}
\begin{Resp}\Rsign Vidérunt eam fíliæ Sion, et beátam dixérunt, et regínæ laudavérunt eam.\end{Resp}
\begin{Resp}\Vsign Glória Patri, et Fílio, \Star{} et Spirítui Sancto.\end{Resp}
\begin{Resp}\Rsign Vidérunt eam fíliæ Sion, et beátam dixérunt, et regínæ laudavérunt eam.\end{Resp}
\switchcolumn
\EN
\begin{Resp}\Rsign Who is this that cometh up like the sun? This, comely as Jerusalem? \Star{} The daughters of Zion saw her, and called her blessed: the queens also, and they praised her.\end{Resp}
\begin{Resp}\Vsign And about her it was as the flower of roses in the spring of the year, and lilies of the valleys.\end{Resp}
\begin{Resp}\Rsign The daughters of Zion saw her, and called her blessed the queens also, and they praised her.\end{Resp}
\begin{Resp}\Vsign Glory be to the Father, and to the Son, \Star{} and to the Holy Ghost.\end{Resp}
\begin{Resp}\Rsign The daughters of Zion saw her, and called her blessed the queens also, and they praised her.\end{Resp}
\switchcolumn*

\LA
\LessonHead{Lectio IV}
\switchcolumn
\EN
\LessonHead{Lesson IV}
\switchcolumn*

\LA
\noindent Ex Constitutione Apostolica Pii Papæ duodecimi\par
\switchcolumn
\EN
\noindent The Lesson is taken from the Apostolic Constitution of Pope Pius XII\par
\switchcolumn*

\LA
\LessonTitle{Sanctorum Patrum ac theologorum argumenta considerationesque Sacris Litteris, tamquam ultimo fundamento, nituntur; quæ quidem almam Dei Matrem nobis véluti ante oculos proponunt divino Filio suo conjunctissimam, ejusque semper participantem sortem. Quamóbrem quasi impossibile videtur eam cérnere, quæ Christum concepit, péperit, suo lacte aluit, eumque inter ulnas habuit pectorique obstrinxit suo, ab eodem post terrestrem hanc vitam, etsi non anima, corpore tamen separatam. Cum Redemptor noster Mariæ Filius sit, haud poterat profecto, útpote divinæ legis observator perfectissimus, præter Æternum Patrem, Matrem quoque suam dilectissimam non honorare. Atqui, cum eam posset tam magno honore exornare, ut eam a sepulcri corruptione servaret incólumem, id reapse fecisse credendum est.}
\switchcolumn
\EN
\LessonTitle{All these proofs and considerations of the holy Fathers and Theologians are based upon the Sacred Scriptures as the final foundation; they establish the blessed Mother of God before our eyes, as it were, as most closely united to her divine Son and always sharing his lot. Therefore, it seemeth impossible to think of her who conceived Christ, gave birth to him, gave him milk, held him in her arms, and clasped him to her heart, as being, after her earthly life, separated from him in body if not in soul. Since our Redeemer is the Son of Mary, he could surely not do otherwise, as the most perfect observer of the divine law, than to honour his most beloved Mother in addition to honouring his eternal Father. And, since it was possible for him to give her this great honour, that she might be preserved from the corruption of the grave, we must believe that he really did so.}
\switchcolumn*

\LA
\begin{Resp}\Rsign Ornátam monílibus fíliam Jerúsalem Dóminus concupívit: \Star{} Et vidéntes eam fíliæ Sion, beatíssimam prædicavérunt, dicéntes: Unguéntum effúsum nomen tuum.\end{Resp}
\begin{Resp}\Vsign Astitit regína a dextris tuis in vestítu deauráto, circúmdata varietáte.\end{Resp}
\begin{Resp}\Rsign Et vidéntes eam fíliæ Sion, beatíssimam prædicavérunt, dicéntes: Unguéntum effúsum nomen tuum.\end{Resp}
\switchcolumn
\EN
\begin{Resp}\Rsign When the Lord beheld the daughter of Jerusalem adorned with her jewels, He greatly desired her beauty \Star{} And when the daughters of Zion saw her, they cried out that she was most blessed, saying thy name is as ointment poured forth.\end{Resp}
\begin{Resp}\Vsign Upon thy right hand did stand the Queen in a vesture of gold wrought about with diverse colours.\end{Resp}
\begin{Resp}\Rsign And when the daughters of Zion saw her, they cried out that she was most blessed, saying thy name is as ointment poured forth.\end{Resp}
\switchcolumn*

\LA
\LessonHead{Lectio V}
\switchcolumn
\EN
\LessonHead{Lesson V}
\switchcolumn*

\LA
\noindent Maxime autem illud memorandum est, inde a sæculo secundo, Mariam Virginem a sanctis Patribus véluti novam Hevam proponi novo Adæ, etsi subjectam, arctissime conjunctam in certamine illo adversus inferorum hostem, quod quemadmodum in protoevangelio præsignificatur, ad plenissimam deventurum erat victoriam de peccato ac de morte, quæ semper in gentium Apostoli scriptis inter se copulantur. Quamóbrem, sicut gloriosa Christi anástasis essentialis pars fuit ac postremum hujus victoriæ tropæum, ita Beatæ Virginis commune cum Filio suo certamen virginei corporis « glorificatione » concludendum erat; ut enim idem Apostolus ait, « cum ... mortale hoc induerit immortalitatem, tunc fiet sermo, qui scriptus est: absorpta est mors in victoria. »\par
\switchcolumn
\EN
\noindent And this indeed should be borne in mind, that as far back as the second century the holy Fathers represented Mary as the new Eve of a new Adam, and closely joined to him (although dependently upon him) in the fight against the hellish enemy which, as already foretold in the protogospel, would end in complete victory over sin and death, which two phases are always joined in the writings of the Apostle of the Gentiles. Wherefore, just as the essential glory of Christ's resurrection was a part and final trophy of this victory, so also the association of Mary with her Son in this common struggle was to end with the glorification of her virginal body; for as the same Apostle saith : When this mortal shall have put on immortality, then shall be brought to pass the saying that is written, Death is swallowed up in victory.\par
\switchcolumn*

\LA
\begin{Resp}\Rsign Beátam me dicent omnes generatiónes, \Star{} Quia fecit mihi Dóminus magna qui potens est, et sanctum nomen ejus.\end{Resp}
\begin{Resp}\Vsign Et misericórdia ejus a progénie in progénies timéntibus eum.\end{Resp}
\begin{Resp}\Rsign Quia fecit mihi Dóminus magna qui potens est, et sanctum nomen ejus.\end{Resp}
\switchcolumn
\EN
\begin{Resp}\Rsign All generations shall call me blessed. \Star{} For He that is mighty, even the Lord, hath done to me great things; and Holy is His Name.\end{Resp}
\begin{Resp}\Vsign And his mercy is on them that fear Him, from generation to generation.\end{Resp}
\begin{Resp}\Rsign He that is mighty, even the Lord, hath done to me great things; and Holy is His Name.\end{Resp}
\switchcolumn*

\LA
\LessonHead{Lectio VI}
\switchcolumn
\EN
\LessonHead{Lesson VI}
\switchcolumn*

\LA
\noindent Idcirco augusta Dei Mater, Jesu Christo, inde ab omni æternitate, « uno eodemque decreto » prædestinationis, arcano modo conjuncta, immaculata in suo conceptu, in divina maternitate sua integerrima virgo, generosa Divini Redemptoris socia, qui plenum de peccato ejusque consectariis deportavit triumphum, id tandem assecuta est, quasi supremam suorum privilegiorum coronam, ut a sepulcri corruptione servaretur immunis, utque, quemadmodum jam Filius suus, devicta morte, corpore et anima ad supernam cæli gloriam eveheretur, ubi Regina refulgéret ad ejusdem sui Filii dexteram, immortalis sæculorum Regis.\par
\switchcolumn
\EN
\noindent Therefore the august Mother of God, from all eternity united in a hidden way with Jesus Christ by one and the same decree of predestination, immaculate in her conception, an inviolate virgin in her divine motherhood, the gracious cooperator of the Divine Redeemer who triumphed completely over sin and its consequences, was finally granted as the supreme crown of her privileges that she should be preserved free from the corruption of the grave, and that, like her own Son, having conquered death, she might be taken up, body and soul, to the exalted glory of heaven, where she sitteth in splendour at the right hand of her very Son, the immortal King of ages.\par
\switchcolumn*

\LA
\begin{Resp}\Rsign Beáta es, Virgo María, quæ Dóminum portásti, Creatórem mundi: \Star{} Genuísti qui te fecit, et in ætérnum pérmanes Virgo.\end{Resp}
\begin{Resp}\Vsign Ave, María, grátia plena; Dóminus tecum.\end{Resp}
\begin{Resp}\Rsign Genuísti qui te fecit, et in ætérnum pérmanes Virgo.\end{Resp}
\begin{Resp}\Vsign Glória Patri, et Fílio, \Star{} et Spirítui Sancto.\end{Resp}
\begin{Resp}\Rsign Genuísti qui te fecit, et in ætérnum pérmanes Virgo.\end{Resp}
\switchcolumn
\EN
\begin{Resp}\Rsign Blessed art thou, O Virgin Mary, who hast carried the Lord, the Maker of the world. \Star{} Thou hast borne Him Who created thee, and thou abidest a virgin for ever.\end{Resp}
\begin{Resp}\Vsign Hail, Mary, full of grace. the Lord is with thee.\end{Resp}
\begin{Resp}\Rsign Thou hast borne Him Who created thee, and thou abidest a virgin for ever.\end{Resp}
\begin{Resp}\Vsign Glory be to the Father, and to the Son, \Star{} and to the Holy Ghost.\end{Resp}
\begin{Resp}\Rsign Thou hast borne Him Who created thee, and thou abidest a virgin for ever.\end{Resp}
\switchcolumn*

\LA
\LessonHead{Lectio VII}
\switchcolumn
\EN
\LessonHead{Lesson VII}
\switchcolumn*

\LA
\LessonTitle{Lectio sancti Evangelii secundum Lucam}
\switchcolumn
\EN
\LessonTitle{From the Holy Gospel according to Luke}
\switchcolumn*

\LA
\Cite{Luc 1:41-50}
\switchcolumn
\EN
\Cite{Luke 1:41-50}
\switchcolumn*

\LA
\noindent In illo tempore: Repleta est Spiritu Sancto Elisabeth et exclamavit voce magna, et dixit: Benedicta tu inter mulíeres. Et réliqua.\par
\switchcolumn
\EN
\noindent At that time : Elisabeth was filled with the Holy Ghost: and she spake out with a loud voice, and said, Blessed art thou among women. And so on, and that which followeth.\par
\switchcolumn*

\LA
\LessonTitle{Homilia sancti Petri Damiani Episcopi}
\switchcolumn
\EN
\LessonTitle{A Homily by St. Peter Damian the Bishop}
\switchcolumn*

\LA
\Source{In Nativitate B. M. V.}
\switchcolumn
\EN
\Source{In Nativitate B. M. V.}
\switchcolumn*

\LA
\noindent Virgo Dei Genitrix, cujus pulchritudinem sol et luna mirantur, súbveni, dómina, clamantibus ad te jugiter: Revértere, revértere Sunamítis, revértere, revértere, ut intueamur te. Tu benedicta, et super benedicta, revértere per potentiam. Fecit in te magna qui potens est, et data est tibi omnis potestas in cælo et in terra. Nihil tibi impossibile, cui possibile est desperatos in spem beatitudinis relevare. Quomodo enim illa potestas tuæ potentiæ poterit obviare, quæ de carne tua carnis suscepit originem? Accedis enim ante illud aureum humanæ, reconciliationis altare, non solum rogans, sed imperans, dómina, non ancilla. Moveat te natura, potentia moveat, quia quanto potentior, tanto misericordior esse debebis. Potestati enim cedit ad gloriam, injurias ulcisci nolle cum possit. Revértere per amórem. Scio, dómina, quia benignissima es, et amas nos amóre invincibili, quos in te et per te Filius tuus et Deus tuus summa dilectione dilexit. Quis scit quoties refrigeras iram Judicis, cum justitiæ virtus a præsentia deitatis egréditur?\par
\switchcolumn
\EN
\noindent O Virgin Mother of God, at whose beauty the sun and the moon stand in awe, come to the aid of them that unceasingly call upon thee, O Lady. Turn thou unto us, o turn thou unto us, O Sulamitess, turn thou unto us, o turn thou unto us that we may behold thee, O blessed, yea more than blessed Maid, turn thou unto us in thy power. He that is mighty hath magnified thee, and hath given thee all power in heaven and in earth. Nothing is impossible to thee, to whom it is possible to shew the most desolate souls the hope of happiness. For how can that Power ignore thy power, when it received its fleshly origin from thy flesh? Thou standest before that golden altar of reconciliation, not only asking, but commanding, as mistress rather than handmaid. Let thy nature move thee, because the more power thou art, the more merciful must thou needs be. For it doth surely redound to the glory of power, to be unwilling to exact vengeance for wrong. Look upon us through love; I know, O Lady, that thou art the most blessed and that thou lovest with an unconquerable love them that in thee and through thee, thy Son and thy God hath loved with a great love. Who knoweth how many times thou didst turn away the anger of the Judge when the virtue of justice went forth from the presence of God.\par
\switchcolumn*

\LA
\begin{Resp}\Rsign Diffúsa est grátia in lábiis tuis: \Star{} Proptérea benedíxit te Deus in ætérnum.\end{Resp}
\begin{Resp}\Vsign Myrrha, et gutta, et cásia a vestiméntis tuis, a dómibus ebúrneis, ex quibus delectavérunt te fíliæ regum in honóre tuo.\end{Resp}
\begin{Resp}\Rsign Proptérea benedíxit te Deus in ætérnum.\end{Resp}
\switchcolumn
\EN
\begin{Resp}\Rsign Grace is poured into thy lips \Star{} Therefore God hath blessed thee for ever.\end{Resp}
\begin{Resp}\Vsign thy garments smell of myrrh, and aloes, and cassia, out of the ivory palaces whereby kings' daughters among thy honourable women have made thee glad.\end{Resp}
\begin{Resp}\Rsign Therefore God hath blessed thee for ever.\end{Resp}
\switchcolumn*

\LA
\LessonHead{Lectio VIII}
\switchcolumn
\EN
\LessonHead{Lesson VIII}
\switchcolumn*

\LA
\noindent Revértere per singularitatem. In manibus tuis sunt thesauri miserationum Domini, et sola electa es, cui gratia tanta conceditur. Absit ut cesset manus tua, cum occasionem quæras salvandi miseros et misericordiam effundendi! neque enim tua gloria minúitur, sed augetur, cum pœnitentes ad veniam, justificati ad gloriam assumuntur. Revértere ergo, Sunamítis, id est despecta, cujus animam pertransivit gladius quæ fabri uxor appellata fuisti. Ad quid? Ut intueamur te. Summa gloria est post Deum te vidére, adhærére tibi, et in tuæ protectionis munimine demorari. Audi nos. Nam et Filius nihil negans honorat te, qui est Deus benedictus in sæcula sæculorum. Amen.\par
\switchcolumn
\EN
\noindent Turn thou unto us through thy very uniqueness. In thy hands are the treasures of the Lord's mercy, and thou alone art the chosen one to whom such great grace is given. Heaven forbid thy hand should ever weaken, for thou wilt never have to seek any occasion of saving the unfortunate or of pouring you thy mercy. Nor doth thy glory ever grow less, but it ever increaseth, for penitents are taken back into favour, and just souls are raised into glory. Turn thou unto us, O Sulamitess, that is, one who was despised, whose soul the sword did pierce and who hath been called the spouse of a carpenter. And why do we ask thee to turn unto us? So that we may behold thee. The greatest glory after seeing God is to see thee, to cling to thee, and to dwell in the fastness of thy protection. Hear thou us, for thy Son doth honour thee and denieth thee nothing, and he is God who is blessed forever and ever. Amen.\par
\switchcolumn*

\LA
\begin{Resp}\Rsign Beáta es Virgo María, Dei Génetrix, quæ credidísti Dómino: perfécta sunt in te quæ dicta sunt tibi: ecce exaltáta es super choros Angelórum: \Star{} Intercéde pro nobis ad Dóminum, Deum nostrum.\end{Resp}
\begin{Resp}\Vsign Ave, María, grátia plena, Dóminus tecum.\end{Resp}
\begin{Resp}\Rsign Intercéde pro nobis ad Dóminum, Deum nostrum.\end{Resp}
\begin{Resp}\Vsign Glória Patri, et Fílio, \Star{} et Spirítui Sancto.\end{Resp}
\begin{Resp}\Rsign Intercéde pro nobis ad Dóminum, Deum nostrum.\end{Resp}
\switchcolumn
\EN
\begin{Resp}\Rsign O Virgin Mary, Mother of God, blessed art thou that didst believe the Lord, for there hath been a performance of those things which were told thee from the Lord. Behold, thou art exalted over choirs of Angels. \Star{} Plead for us with the Lord our God.\end{Resp}
\begin{Resp}\Vsign Hail, Mary, full of grace, the Lord is with thee.\end{Resp}
\begin{Resp}\Rsign Plead for us with the Lord our God.\end{Resp}
\begin{Resp}\Vsign Glory be to the Father, and to the Son, \Star{} and to the Holy Ghost.\end{Resp}
\begin{Resp}\Rsign Plead for us with the Lord our God.\end{Resp}
\switchcolumn*

\end{paracol}

\par\needspace{10\baselineskip}\addvspace{1.2em}{\centering\small\textsc{Die 18 Augusti} · \textsc{18 August}\par}
\Office{S. Agapiti Martyris}{St. Agapitus, Martyr}{Simplex}
\addcontentsline{toc}{subsection}{Die 18 Augusti · S. Agapiti Martyris \textit{— St. Agapitus, Martyr}}
\begin{paracol}{2}
\LA
\LessonHead{Lectio IX (pro commemoratione)}
\switchcolumn
\EN
\LessonHead{Lesson IX (when commemorated)}
\switchcolumn*

\LA
\LessonTitle{Commemoratio: S. Agapiti Martyris}
\switchcolumn
\EN
\LessonTitle{Commemoration of S. Agapitus, Martyr}
\switchcolumn*

\LA
\noindent Agapitus Prænestínus, quindecim annos natus, Aureliáno imperatóre, martyrii cupidíssimus, cum propter constantiam religiónis, imperatóris jussu, primum nervis diutíssime cæsus, deínde in tetrum carcerem conjéctus esset, ut nihil omnino per quatuor dies gustaret, e custódia est edúctus. Et, ardéntibus carbónibus capiti ejus impositis, Deo agens grátias, íterum afféctus verbéribus, nudus ita pédibus suspénsus est, ut ingens fumus e subjecto igne os ejus obrúeret. Tum fervens aqua in ejus ventrem effúsa est, maxillǽque confractæ. Quo témpore judex, e tribunali lapsus, paulo post mórtuus est. Ea re incénso imperatóre, sanctum juvenem jubénte feris óbjici, cum illæ non audérent attingere, Præneste gládio percússus est.\par
\switchcolumn
\EN
\noindent Agapitus was a young man of Palestrina, who eagerly accepted martyrdom at the age of fifteen years, under the Emperor Aurelian. On account of his firmness in his religion, the Emperor ordered him first to receive an exceedingly long hiding with cat-gut, and then to be thrown into a foul dungeon, where he remained for four days without food. Being brought out of prison, live embers were put upon his head, but, whereas he still gave God thanks, he was whipped again, and hung up naked by the feet, in such wise that a thick smoke from a fire kindled under his face might pour into his mouth. Afterwards, boiling water was poured upon his belly, and his jaws were broken. Then presently the judge fell from his judgement seat, and shortly after died. Whereupon the Emperor was enraged, and commanded the holy youth to be thrown to wild beasts, but these dared not to touch him, and he was stricken by the sword at Palestrina.\par
\switchcolumn*

\LA
\TeDeum{Te Deum}
\switchcolumn
\EN
\TeDeum{Te Deum}
\switchcolumn*

\end{paracol}

\par\needspace{10\baselineskip}\addvspace{1.2em}{\centering\small\textsc{Die 19 Augusti} · \textsc{19 August}\par}
\Office{S. Joannis Eudes Confessoris}{St. John Eudes, Confessor}{Duplex}
\addcontentsline{toc}{subsection}{Die 19 Augusti · S. Joannis Eudes Confessoris \textit{— St. John Eudes, Confessor}}
\begin{paracol}{2}
\LA
\RefLine{Lectiones I–III: de Scriptura occurrente.}
\switchcolumn
\EN
\RefLine{Lessons I–III: from the Scripture of the day.}
\switchcolumn*

\LA
\RefLine{Lectiones VII–IX: ex Commúni Confessoris non pontificis (C5).}
\switchcolumn
\EN
\RefLine{Lessons VII–IX: from the Common of a Confessor not a Bishop (C5).}
\switchcolumn*

\LA
\LessonHead{Lectio IV}
\switchcolumn
\EN
\LessonHead{Lesson IV}
\switchcolumn*

\LA
\noindent Joannes, anno millesimo sexcentesimo primo, in pago vulgo Ri, Sagiensis diœcesis, e piis honestisque parentibus ortus est. Adhuc puer, Angelorum pane refectus, perpetuam castitatem alacriter vovit. In Cadomense collegium a patribus Societatis Jesu directum exceptus, singulari pietate emicuit; et Mariæ Virginis tutelæ sese committens, speciale fœdus initum cum ea vix adolescens suo sanguine signavit. Litterarum ac philosophiæ curriculo multa cum laude peracto, spretisque sibi oblatis nuptiis, Berulliani Oratorii congregationi nomen dedit, et sacerdotio Parisiis auctus est. Mira caritate erga proximum flagravit: nam apud plura loca, corporibus, asiatica lue perculsis, animisque curandis sedulam operam dedit. Domus Oratorianæ Cadomi rector factus, quum diu cogitaret ad Ecclesiæ ministerium juvenes idoneos instituere, a sodalibus, quibuscum viginti annos vixerat, divina ope implorata, licet ægre, forti animo discessit.\par
\switchcolumn
\EN
\noindent John was born in the year 1601, of pious and respectable parents, at a village commonly known as Ri, in the diocese of Seez. While still a boy, when he was fed with the bread of Angels, he cheerfully made a vow of perpetual chastity. Having been received at the College of Caen, directed by the Fathers of the Society of Jesus, he was conspicuous for a remarkable piety; and, committing himself to the protection of the Virgin Mary, when still a youth he signed with his own blood, the special covenant he had entered into with her. Having completed his courses of letters and of philosophy with great distinction, and having spurned opportunities of marriage which had been arranged for him, he enrolled himself with the Congregation of the Oratory de Bérulle, and was ordained priest at Paris. He was on fire with a marvellous love towards his neighbour: for he took the most constant pains in caring for both the souls and bodies of those smitten with the Asiatic plague, in many different places. He was made Rector of the Oratorian house at Caen, but since he had been thinking for a long time of educating suitable young men for the ministry of the Church, earnestly asking for the divine assistance, with a brave spirit he most regretfully departed from the associates with whom he had lived for twenty years.\par
\switchcolumn*

\LA
\begin{Resp}\Rsign Honéstum fecit illum Dóminus, et custodívit eum ab inimícis, et a seductóribus tutávit illum: \Star{} Et dedit illi claritátem ætérnam.\end{Resp}
\begin{Resp}\Vsign Justum dedúxit Dóminus per vias rectas, et osténdit illi regnum Dei.\end{Resp}
\begin{Resp}\Rsign Et dedit illi claritátem ætérnam.\end{Resp}
\switchcolumn
\EN
\begin{Resp}\Rsign The Lord made him honourable, and defended him from his enemies, and kept him safe from those that lay in wait for him; \Star{} And gave him perpetual glory.\end{Resp}
\begin{Resp}\Vsign He went down with him into the pit, and left him not in bonds.\end{Resp}
\begin{Resp}\Rsign And gave him perpetual glory.\end{Resp}
\switchcolumn*

\LA
\LessonHead{Lectio V}
\switchcolumn
\EN
\LessonHead{Lesson V}
\switchcolumn*

\LA
\noindent Quinque igitur sacerdotes sibi consocians, anno millesimo sexcentesimo quadragesimo tertio, die festo Annuntiationis beatæ Mariæ Virginis, congregationem presbyterorum instituit, cui sanctissima Jesu et Mariæ nomina dedit, et Cadomi primum seminarium aperuit, quod in Normannia et Britannia minori complura, eodem auctore, subsecuta sunt. Pro peccatricibus feminis, ad christianam vitam revocandis, ordinem Dominæ nostræ a Caritate fundavit; cujus nobilissimæ arboris ramus est congregatio Andegavensis a Bono Pastore. Insuper societatem a Matris Dei Corde admirabili, et alia caritatis opera condidit. Scripta plura præclare edidit, et missionarius Apostolicus tot pagos, oppida et urbes ipsamque regiam aulam usque ad extremam ætatem, evangelizavit.\par
\switchcolumn
\EN
\noindent Accordingly, associating five priests with himself, in the year 1643, on the feastday of the Annunciation of the Blessed Virgin Mary, he founded a Congregation of Priests, to whom he gave the most holy names of Jesus and Mary, and opened the first seminary at Caen; and a great many others followed immediately in Normandy and Brittany, also founded by him. For the recalling of sinful women to a Christian life, he founded the Order of Our Lady of Charity; of which most noble tree, the Congregation of the Good Shepherd of Angers is a branch. Furthermore, he founded the Society of the Admirable Heart of the Mother of God, and other charitable institutions. He was the author of many excellent treatises, and laboured as an Apostolic Missionary to the very end of his life, preaching the Gospel in very many villages, towns, and cities, and even in the royal court.\par
\switchcolumn*

\LA
\begin{Resp}\Rsign Amávit eum Dóminus, et ornávit eum: stolam glóriæ índuit eum, \Star{} Et ad portas paradísi coronávit eum.\end{Resp}
\begin{Resp}\Vsign Índuit eum Dóminus lorícam fídei, et ornávit eum.\end{Resp}
\begin{Resp}\Rsign Et ad portas paradísi coronávit eum.\end{Resp}
\switchcolumn
\EN
\begin{Resp}\Rsign The Lord loved him and beautified him; He clothed him with a robe of glory, \Star{} And crowned him at the gates of Paradise.\end{Resp}
\begin{Resp}\Vsign The Lord hath put on him the breast-plate of faith, and hath adorned him.\end{Resp}
\begin{Resp}\Rsign And crowned him at the gates of Paradise.\end{Resp}
\switchcolumn*

\LA
\LessonHead{Lectio VI}
\switchcolumn
\EN
\LessonHead{Lesson VI}
\switchcolumn*

\LA
\noindent Singulare ejus studium emicuit in salutari devotione promovenda erga sacratissima Corda Jesu et Mariæ, quorum liturgicum cultum eisdem præstandum non sine aliquo divino afflatu, primus omnium excogitavit ideoque ejusdem cultus pater, doctor et apostolus habitus est. Jansenistarum doctrinis fortiter resistens, immutatum erga Petri cathedram obsequium servavit, et pro suis inimicis, tamquam pro fratribus, assidue Deum precatus est. Tot laboribus, potius quam annis fractus, cupiens dissolvi et esse cum Christo, die decima nona Augusti, anno millesimo sexcentesimo octogesimo, suavia Jesu et Mariæ nomina sæpius repetens, placide exspiravit. Quem pluribus miraculis clarum Pius Papa decimus Beatorum albo adscripsit eumque novis signis fulgentem Pius Papa undecimus, anno sacro, in die dominica Pentecostes inter Sanctos retulit ejusque Officium ac Missam ad universam Ecclesiam exendit.\par
\switchcolumn
\EN
\noindent His matchless zeal was very conspicuous in promoting the salutary devotion towards the most sacred Hearts of Jesus and Mary, whose liturgical worship he was the first of all to devise, although not without some divine inspiration. He is therefore held to be the father, the teacher, and the apostle of that worship. Courageously withstanding the doctrines of the Jansenists, he preserved unalterable obedience towards the Chair of Peter, and he constantly prayed to God, both for his enemies as well as for his brethren. Broken by so many labours, rather than by years, desiring to be freed and to be with Christ, on the 19th day of August, 1680, frequently repeating the sweet names of Jesus and Mary, he died in peace. As he became illustrious by many miracles, Pope Pius X added him to the list of the Blessed, and as he still shone forth with new signs and wonders, Pope Pius XI, in the holy year and on the day of Pentecost, placed him among the Saints, and extended his Office and Mass to the universal Church.\par
\switchcolumn*

\LA
\begin{Resp}\Rsign Iste homo perfécit ómnia quæ locútus est ei Deus, et dixit ad eum: Ingrédere in réquiem meam: \Star{} Quia te vidi justum coram me ex ómnibus géntibus.\end{Resp}
\begin{Resp}\Vsign Iste est, qui contémpsit vitam mundi, et pervénit ad cæléstia regna.\end{Resp}
\begin{Resp}\Rsign Quia te vidi justum coram me ex ómnibus géntibus.\end{Resp}
\begin{Resp}\Vsign Glória Patri, et Fílio, \Star{} et Spirítui Sancto.\end{Resp}
\begin{Resp}\Rsign Quia te vidi justum coram me ex ómnibus géntibus.\end{Resp}
\switchcolumn
\EN
\begin{Resp}\Rsign This is he which did according unto all that God commanded him; and God said unto him: Enter thou into My rest, \Star{} For thee have I seen righteous before Me among all people.\end{Resp}
\begin{Resp}\Vsign This is he which loved not his life in this world, and is come unto an everlasting kingdom.\end{Resp}
\begin{Resp}\Rsign For thee have I seen righteous before Me among all people.\end{Resp}
\begin{Resp}\Vsign Glory be to the Father, and to the Son, \Star{} and to the Holy Ghost.\end{Resp}
\begin{Resp}\Rsign For thee have I seen righteous before Me among all people.\end{Resp}
\switchcolumn*

\end{paracol}

\par\needspace{10\baselineskip}\addvspace{1.2em}{\centering\small\textsc{Die 20 Augusti} · \textsc{20 August}\par}
\Office{S. Bernardi Abbatis et Ecclesiæ Doctoris}{St. Bernard of Clairvaux, Abbot and Doctor of the Church}{Duplex}
\addcontentsline{toc}{subsection}{Die 20 Augusti · S. Bernardi Abbatis et Ecclesiæ Doctoris \textit{— St. Bernard of Clairvaux, Abbot and Doctor of the Church}}
\begin{paracol}{2}
\LA
\RefLine{Lectiones I–III: de Scriptura occurrente.}
\switchcolumn
\EN
\RefLine{Lessons I–III: from the Scripture of the day.}
\switchcolumn*

\LA
\LessonHead{Lectio IV}
\switchcolumn
\EN
\LessonHead{Lesson IV}
\switchcolumn*

\LA
\noindent Bernárdus, Fóntanis in Burgúndia honésto loco natus, adolescéns propter egrégiam formam veheménter sollicitátus a muliéribus, numquam de senténtia coléndæ castitátis dimovéri pótuit. Quas diáboli tentatiónes ut effúgeret, duos et vigínti annos natus, monastérium Cisterciénse, unde hic ordo incépit, et quod tum sanctitáte florébat, íngredi constítuit. Quo Bernárdi consílio cógnito, fratres summópere conáti sunt eum a propósito deterrére: in quo ipse eloquéntior ac felícior fuit. Nam sic eos aliósque multos in suam perdúxit senténtiam, ut cum eo trigínta júvenes eámdem religiónem suscéperint. Mónachus jejúnio ita déditus erat, ut quóties suméndus esset cibus, tóties torméntum subíre viderétur. In vigíliis étiam et oratiónibus mirífice se exercébat: et christiánam paupertátem colens, quasi cæléstem vitam agébat in terris, ab omni caducárum rerum cura et cupiditáte aliénam.\par
\switchcolumn
\EN
\noindent Bernard was born (in the year of salvation 1091) at a decent place in Burgundy called Fontaines. On account of extraordinary good looks, he was as a boy very much sought after by women, but he could never be turned aside from his resolution to keep chaste. To fly from these temptations of the devil, he determined at two-and-twenty years of age to enter the Monastery of Citeaux, whence the Cistercian Order took its rise. When this resolution of Bernard's became known, his brothers did all their diligence to change his purpose, but he only became the more eloquent and happy about it. Them and others he so brought over to his mind, that thirty young men entered the same Order along with him. As a monk he was so given to fasting, that as often as he had to eat, so often he seemed to be in pain. He exercised himself wonderfully in watching and prayer, and was a great lover of Christian poverty. Thus he led on earth an heavenly life, purged of all care and desire for transitory things.\par
\switchcolumn*

\LA
\begin{Resp}\Rsign Honéstum fecit illum Dóminus, et custodívit eum ab inimícis, et a seductóribus tutávit illum: \Star{} Et dedit illi claritátem ætérnam.\end{Resp}
\begin{Resp}\Vsign Justum dedúxit Dóminus per vias rectas, et osténdit illi regnum Dei.\end{Resp}
\begin{Resp}\Rsign Et dedit illi claritátem ætérnam.\end{Resp}
\switchcolumn
\EN
\begin{Resp}\Rsign The Lord made him honourable, and defended him from his enemies, and kept him safe from those that lay in wait for him; \Star{} And gave him perpetual glory.\end{Resp}
\begin{Resp}\Vsign He went down with him into the pit, and left him not in bonds.\end{Resp}
\begin{Resp}\Rsign And gave him perpetual glory.\end{Resp}
\switchcolumn*

\LA
\LessonHead{Lectio V}
\switchcolumn
\EN
\LessonHead{Lesson V}
\switchcolumn*

\LA
\noindent Elucébat in eo humílitas, misericórdia, benígnitas: contemplatióni autem sic addíctus erat, ut vix sénsibus, nisi ad offícia pietátis, uterétur: in quibus tamen prudéntiæ laude excellébat. Quo in stúdio occupátus, Genuénsem, ac Mediolanénsem, aliósque episcopátus oblátos recusávit, proféssus se tanti offícii múnere indígnum esse. Abbas factus Claravallénsis, multis in locis ædificávit monastéria, in quibus præclára Bernárdi institútio ac disciplína diu viguit. Romæ, sanctórum Vincéntii et Anastásii monastério ab Innocéntio secúndo Papa restitúto præfécit abbátem illum, qui póstea Eugénius tértius Summus Póntifex fuit, ad quem étiam librum misit de Consideratióne.\par
\switchcolumn
\EN
\noindent He was a burning and shining light of lowliness, mercifulness, and kindness. His concentration of thought was such, that he hardly used his senses except to do good works, in which latter he acted with admirable wisdom. Thus occupied, he refused the Bishoprics of Genoa, Milan, and others, which were offered to him, declaring that he was unworthy of so high a sphere of duty. Being made Abbot of Clairvaux (in 1115,) he built monasteries in many places, wherein the excellent rules and discipline of Bernard long flourished. When Pope Innocent II., (in 1138,) restored the monastery of St. Vincent and St. Anastasius at Rome, Bernard set over it the Abbot who was afterwards the Supreme Pontiff Eugene III., and who is also the same to whom he addressed his book upon Consideration.\par
\switchcolumn*

\LA
\begin{Resp}\Rsign Amávit eum Dóminus, et ornávit eum: stolam glóriæ índuit eum, \Star{} Et ad portas paradísi coronávit eum.\end{Resp}
\begin{Resp}\Vsign Índuit eum Dóminus lorícam fídei, et ornávit eum.\end{Resp}
\begin{Resp}\Rsign Et ad portas paradísi coronávit eum.\end{Resp}
\switchcolumn
\EN
\begin{Resp}\Rsign The Lord loved him and beautified him; He clothed him with a robe of glory, \Star{} And crowned him at the gates of Paradise.\end{Resp}
\begin{Resp}\Vsign The Lord hath put on him the breast-plate of faith, and hath adorned him.\end{Resp}
\begin{Resp}\Rsign And crowned him at the gates of Paradise.\end{Resp}
\switchcolumn*

\LA
\LessonHead{Lectio VI}
\switchcolumn
\EN
\LessonHead{Lesson VI}
\switchcolumn*

\LA
\noindent Multa prætérea scripsit, in quibus appáret, eum doctrína pótius divínitus trádita, quam labóre comparáta, instrúctum fuísse. In summa virtútum laude exorátus a máximis princípibus de eórum componéndis controvérsiis, et de ecclesiásticis rebus constituéndis, sǽpius in Itáliam venit. Innocéntium item secúndum Pontíficem Máximum in confutándo schísmate Petri Leónis, cum apud imperatórem et Henrícum Angliæ regem, tum in concílio Pisis coácto, egrégie adjúvit. Dénique tres et sexagínta annos natus, obdormívit in Dómino, ac miráculis illústris, ab Alexándro tértio Papa inter sanctos relátus est. Pius vero octávus Póntifex Máximus ex sacrórum Rítuum Congregatiónis consílio sanctum Bernárdum universális Ecclésiæ Doctórem declarávit et confirmávit, nec non Missam et Offícium de Doctóribus ab ómnibus recitári jussit, atque indulgéntias plenárias quotánnis in perpétuum órdinis Cisterciénsium ecclésias visitántibus die hujus sancti festo concéssit.\par
\switchcolumn
\EN
\noindent He was the author of many writings, in which it is manifest that his teaching was rather given him of God, than gained by hard work. In consequence of his high reputation for excellence, he was called by the most exalted Princes to act as arbiter of their disputes, and for this end, and to settle affairs of the Church, he often went to Italy. He was an eminent helper to Pope Innocent II., in putting down the schism of Peter Leoni, and worked to this end, both at the Courts of the Emperor and of Henry King of England, and in the Council of Pisa. He fell asleep in the Lord, (at Clairvaux, on the 20th day of August,) in the (year 1153, the) sixty-third year of his age. He was famous for miracles, and Pope Alexander III. numbered him among the Saints. Pope Pius VIII., acting on the advice of the Congregation of Sacred Rites, declared and confirmed St. Bernard a Doctor of the Universal Church. He also commanded that all should use the Mass and Office for him as for a Doctor, and granted perpetual yearly plenary indulgences to all who should visit Churches of the Cistercian Order upon the Feast-day of this Saint.\par
\switchcolumn*

\LA
\begin{Resp}\Rsign Iste homo perfécit ómnia quæ locútus est ei Deus, et dixit ad eum: Ingrédere in réquiem meam: \Star{} Quia te vidi justum coram me ex ómnibus géntibus.\end{Resp}
\begin{Resp}\Vsign Iste est, qui contémpsit vitam mundi, et pervénit ad cæléstia regna.\end{Resp}
\begin{Resp}\Rsign Quia te vidi justum coram me ex ómnibus géntibus.\end{Resp}
\begin{Resp}\Vsign Glória Patri, et Fílio, \Star{} et Spirítui Sancto.\end{Resp}
\begin{Resp}\Rsign Quia te vidi justum coram me ex ómnibus géntibus.\end{Resp}
\switchcolumn
\EN
\begin{Resp}\Rsign This is he which did according unto all that God commanded him; and God said unto him: Enter thou into My rest, \Star{} For thee have I seen righteous before Me among all people.\end{Resp}
\begin{Resp}\Vsign This is he which loved not his life in this world, and is come unto an everlasting kingdom.\end{Resp}
\begin{Resp}\Rsign For thee have I seen righteous before Me among all people.\end{Resp}
\begin{Resp}\Vsign Glory be to the Father, and to the Son, \Star{} and to the Holy Ghost.\end{Resp}
\begin{Resp}\Rsign For thee have I seen righteous before Me among all people.\end{Resp}
\switchcolumn*

\LA
\LessonHead{Lectio VII}
\switchcolumn
\EN
\LessonHead{Lesson VII}
\switchcolumn*

\LA
\LessonTitle{Léctio sancti Evangélii secúndum Lucam}
\switchcolumn
\EN
\LessonTitle{From the Holy Gospel according to Luke}
\switchcolumn*

\LA
\Cite{Luc 10:1-9}
\switchcolumn
\EN
\Cite{Luke 10:1-9}
\switchcolumn*

\LA
\noindent In illo témpore: Designávit Dóminus et álios septuagínta duos: et misit illos binos ante fáciem suam, in omnem civitátem et locum, quo erat ipse ventúrus. Et réliqua.\par
\switchcolumn
\EN
\noindent At that time, the Lord appointed other seventy-two also, and sent them two and two before His face into every city and place, whither He Himself would come. And so on.\par
\switchcolumn*

\LA
\LessonTitle{Homilía sancti Gregórii Papæ}
\switchcolumn
\EN
\LessonTitle{Homily by Pope St. Gregory (the Great.)}
\switchcolumn*

\LA
\Source{Homilia 17 in Evangelia}
\switchcolumn
\EN
\Source{17th Homily on the Gospels}
\switchcolumn*

\LA
\noindent Dóminus et Salvátor noster, fratres caríssimi, aliquándo nos sermónibus, aliquándo vero opéribus ádmonet. Ipsa étenim facta ejus præcépta sunt: quia dum áliquid tácitus facit, quid ágere debeámus innotéscit. Ecce enim binos in prædicatiónem discípulos mittit: quia duo sunt præcépta caritátis, Dei vidélicet amor, et próximi: et minus quam inter duos cáritas habéri non potest. Nemo enim próprie ad semetípsum habére caritátem dícitur: sed diléctio in álterum tendit, ut cáritas esse possit.\par
\switchcolumn
\EN
\noindent Dearly beloved brethren, our Lord and Saviour doth sometimes admonish us by words, and sometimes by works. Yea, His very works do themselves teach us for that which He doth silently His example still moveth us to copy. Behold how He sendeth forth His disciples to preach by two and two since there are two commandments to love, that is, a commandment to love God, and a commandment to love our neighbour and where there are not two, the one, being alone, hath not whereon to do the Lord's commandment. And no man can properly be said to love himself: for love tendeth outward toward our neighbour, if it be the love whereto the Gospel doth oblige us.\par
\switchcolumn*

\LA
\begin{Resp}\Rsign Iste est qui ante Deum magnas virtútes operátus est, et de omni corde suo laudávit Dóminum: \Star{} Ipse intercédat pro peccátis ómnium populórum.\end{Resp}
\begin{Resp}\Vsign Ecce homo sine queréla, verus Dei cultor, ábstinens se ab omni ópere malo, et pérmanens in innocéntia sua.\end{Resp}
\begin{Resp}\Rsign Ipse intercédat pro peccátis ómnium populórum.\end{Resp}
\switchcolumn
\EN
\begin{Resp}\Rsign This is he which wrought great wonders before God, and praised the Lord with all his heart. \Star{} May he pray for all people, that their sins may be forgiven unto them.\end{Resp}
\begin{Resp}\Vsign Behold a man without blame, a worshipper of God in truth, keeping himself clean from every evil work, and abiding still in his innocency.\end{Resp}
\begin{Resp}\Rsign May he pray for all people, that their sins may be forgiven unto them.\end{Resp}
\switchcolumn*

\LA
\LessonHead{Lectio VIII}
\switchcolumn
\EN
\LessonHead{Lesson VIII}
\switchcolumn*

\LA
\noindent Ecce enim binos ad prædicándum discípulos Dóminus mittit: quátenus hoc nobis tácitus ínnuat, quia qui caritátem erga álterum non habet, prædicatiónis offícium suscípere nullátenus debet. Bene autem dícitur, quia misit eos ante fáciem suam in omnem civitátem et locum, quo erat ipse ventúrus. Prædicatóres enim suos Dóminus séquitur: quia prædicátio prǽvenit, et tunc ad mentis nostræ habitáculum Dóminus venit, quando verba exhortatiónis præcúrrunt: atque per hoc véritas in mente suscípitur.\par
\switchcolumn
\EN
\noindent Behold, the Lord sendeth forth His disciples to preach by two and two and thus doing, He doth silently teach us that whosoever loveth not his neighbour, such a one it behoveth not to take upon him the office of a preacher. Well also is it said that He sent them before His face into every city and place whither He Himself would come. The Lord followeth His preachers first cometh preaching, and then the Lord Himself cometh to the house of our mind, whither the word of exhortation hath come before and so cometh the truth into our mind.\par
\switchcolumn*

\LA
\begin{Resp}\Rsign In médio Ecclésiæ apéruit os ejus, \Star{} Et implévit eum Dóminus spíritu sapiéntiæ et intelléctus.\end{Resp}
\begin{Resp}\Vsign Jucunditátem et exsultatiónem thesaurizávit super eum.\end{Resp}
\begin{Resp}\Rsign Et implévit eum Dóminus spíritu sapiéntiæ et intelléctus.\end{Resp}
\begin{Resp}\Vsign Glória Patri, et Fílio, \Star{} et Spirítui Sancto.\end{Resp}
\begin{Resp}\Rsign Et implévit eum Dóminus spíritu sapiéntiæ et intelléctus.\end{Resp}
\switchcolumn
\EN
\begin{Resp}\Rsign In the midst of the congregation did the Lord open his mouth. \Star{} And filled him with the spirit of wisdom and understanding.\end{Resp}
\begin{Resp}\Vsign He made him rich with joy and gladness.\end{Resp}
\begin{Resp}\Rsign And filled him with the spirit of wisdom and understanding.\end{Resp}
\begin{Resp}\Vsign Glory be to the Father, and to the Son, \Star{} and to the Holy Ghost.\end{Resp}
\begin{Resp}\Rsign And filled him with the spirit of wisdom and understanding.\end{Resp}
\switchcolumn*

\LA
\LessonHead{Lectio IX}
\switchcolumn
\EN
\LessonHead{Lesson IX}
\switchcolumn*

\LA
\noindent Hinc namque eísdem prædicatóribus Isaías dicit: Paráte viam Dómini, rectas fácite sémitas Dei nostri. Hinc fíliis Psalmísta ait: Iter fácite ei, qui ascéndit super occásum. Super occásum namque Dóminus ascéndit: quia unde in passióne occúbuit, inde majórem suam glóriam resurgéndo manifestávit. Super occásum vidélicet ascéndit; quia mortem quam pértulit, resurgéndo calcávit. Ei ergo qui ascéndit super occásum, iter fácimus, cum nos ejus glóriam vestris méntibus prædicámus, ut eas et ipse post véniens, per amóris sui præséntiam illústret.\par
\switchcolumn
\EN
\noindent Therefore, to preachers saith Isaiah: “Prepare ye the way of the Lord, make straight an highway for our God.” (xl. 3.) And again the Psalmist saith: “Spread a path before him that rideth upon the West.” (lxvii. 4.) The Lord rideth upon the West above that from which in death He veiled His glory, hath He royally exalted that glory that excelleth, even the glory of His rising again. He rideth upon the West, Who, being risen again from the dead, is throned high above the death to which He bowed. Before Him, therefore, That rideth upon the West, we spread a path, when we set forth His glory before the eyes of your mind, to the end that He Himself may come after, and Himself enlighten the same your minds by His presence and His love.\par
\switchcolumn*

\LA
\TeDeum{Te Deum}
\switchcolumn
\EN
\TeDeum{Te Deum}
\switchcolumn*

\end{paracol}

\par\needspace{10\baselineskip}\addvspace{1.2em}{\centering\small\textsc{Die 21 Augusti} · \textsc{21 August}\par}
\Office{S. Joannæ Franciscæ Frémiot de Chantal Viduæ}{St. Jeanne-Françoise Frémiot de Chantal, Widow;Duplex}{Duplex}
\addcontentsline{toc}{subsection}{Die 21 Augusti · S. Joannæ Franciscæ Frémiot de Chantal Viduæ \textit{— St. Jeanne-Françoise Frémiot de Chantal, Widow;Duplex}}
\begin{paracol}{2}
\LA
\RefLine{Lectiones I–III: de Scriptura occurrente.}
\switchcolumn
\EN
\RefLine{Lessons I–III: from the Scripture of the day.}
\switchcolumn*

\LA
\RefLine{Lectiones VII–IX: ex Commúni non Virginum non Martyrum (C7a).}
\switchcolumn
\EN
\RefLine{Lessons VII–IX: from the Common of Widows (C7a).}
\switchcolumn*

\LA
\LessonHead{Lectio IV}
\switchcolumn
\EN
\LessonHead{Lesson IV}
\switchcolumn*

\LA
\noindent Joanna Francisca Frémiot de Chantal, Divione in Burgundia clarissimis orta natalibus, ab ineunte ætate eximiæ sanctitatis non obscuras edidit significationes. Eam enim vix quinquennem nobilem quemdam Calvinistam solida supra ætatem argumentatione perstrinxisse ferunt, collatumque ab eo munusculum flammis illico tradidisse, in hæc verba: En quomodo hæretici apud inferos comburentur, qui loquenti Christo fidem detrectant. Matre orbata, Deiparæ Virginis tutelæ se commendavit, et famulam, quæ ad mundi amórem eam alliciebat, ab se rejecit. Nihil puerile in moribus exprimens, a sæculi deliciis abhorrens, martyriumque anhelans, religioni ac pietati impense studebat. Baroni de Chantal nuptui a patre tradita, virtutibus omnibus excolendis operam dedit, liberos, famulos, aliosque sibi subjectos in fidei doctrina, bonisque moribus imbuere satagens. Profusa liberalitate pauperum inopiam sublevabat, annona divinitus non raro multiplicata: quo factum est, ut nemini se umquam Christi nomine roganti stipem abnegaturam spoponderit.\par
\switchcolumn
\EN
\noindent Jane Frances Fremiot de Chantal was born of parents of the highest rank, at Dijon in Burgundy, (on the 23rd day of January, in the year 1575.) From her earliest childhood she gave no dark promise of a life of eminent holiness. It is said that when she was scarcely fifteen years of age she confuted with precocious acuteness a Presbyterian nobleman, and when he gave her a little present she put it in the fire, saying That is how heretics will burn in hell for not believing Christ when He speaketh. On the death of her mother, she placed herself under the keeping of the Virgin Mother of God, and discharged a maid who strove to beguile her into loving the world. She had nothing youthful about her ways. She shrank from the pleasures of life. She had a strong wish that she might die a martyr. She devoted herself unweariedly to religion and godliness. Her father gave her in marriage, (at twenty years of age,) to the Baron de Chantal, and she strove to excel in all the duties and graces of a wife. She made it her work to see that her children, her servants, and all others under her authority were taught the doctrines of the faith and the practice of good living. She relieved the sufferings of the poor by plentiful almsgiving, for which purposes God not unfrequently miraculously multiplied her money. And so it came to pass that no one ever asked her for food in Christ's name and was refused it.\par
\switchcolumn*

\LA
\begin{Resp}\Rsign Propter veritátem, et mansuetúdinem, et justítiam: \Star{} Et dedúcet te mirabíliter déxtera tua.\end{Resp}
\begin{Resp}\Vsign Spécie tua et pulchritúdine tua inténde, próspere procéde, et regna.\end{Resp}
\begin{Resp}\Rsign Et dedúcet te mirabíliter déxtera tua.\end{Resp}
\switchcolumn
\EN
\begin{Resp}\Rsign Because of truth, and meekness, and righteousness; \Star{} And thy right hand shall lead thee wonderfully.\end{Resp}
\begin{Resp}\Vsign In thy comeliness, and thy beauty, go forward, fare prosperously, and reign.\end{Resp}
\begin{Resp}\Rsign And thy right hand shall lead thee wonderfully.\end{Resp}
\switchcolumn*

\LA
\LessonHead{Lectio V}
\switchcolumn
\EN
\LessonHead{Lesson V}
\switchcolumn*

\LA
\noindent Viro in venatione interempto, perfectioris vitæ consilium iniens, continentiæ voto se obstrinxit. Viri necem non solum æquo animo tulit, sed, in publicum indultæ veniæ testimonium, occisoris filium e sacro fonte suscipere sui victrix elegit. Modica familia, tenui victu atque vestitu contenta, pretiosas vestes in pios usus convertit. Quidquid a domesticis curis supererat temporis, precibus, piis lectionibus, laborique impendebat. Numquam adduci potuit ut alteras nuptias, quamvis utiles et honorificas, iniret. Ne autem a proposito castimoniæ observandæ in posterum dimoveretur, illius voto innovato, sanctissimum Jesu Christi nomen candenti ferro pectori insculpsit. Ardentius in dies caritate fervescens, pauperes, derelictos, ægros, teterrimisque morbis infectos ad se adducendos curabat; eosque non hospitio tantum excipiebat, solabatur, fovebat, verum étiam sordidas eorumdem vestes depurgabat, laceras reficiebat, et manantibus fœtido pure ulceribus labia admovere non exhorrebat.\par
\switchcolumn
\EN
\noindent When her husband was accidentally killed out shooting, and in her widowhood she determined to embrace the more excellent way, and took a vow not to marry again. She bore her bereavement with resignation to the Will of God, and so far overcame her horror of the gentleman who had fired the shot, that, to show she attributed no blame to him, she stood god-mother to his little boy. She was quite content with few servants and plain cookery and dress, and sold her rich wardrobe for the benefit of charities. She received offers of second marriage which would have been both politic and honourable, but never was induced to accept one of them, and to harden herself in her intention of remaining in her widowhood, she renewed her vow to that effect, and branded on her chest with a hot iron the most holy name of Jesus Christ. Her love grew tenderer every day, and she had brought to her the starving, the abandoned, the diseased, and those who were afflicted with the most sickening disorders. Them she not only sheltered, comforted, and nursed, but washed and mended their filthy and ragged garments, and shrank not from putting her mouth to their sores oozing with disgusting matter. Whatever of her time was not taken up by her household duties she spent in prayer, reading godly books, and working.\par
\switchcolumn*

\LA
\begin{Resp}\Rsign Dilexísti justítiam, et odísti iniquitátem: \Star{} Proptérea unxit te Deus, Deus tuus, óleo lætítiæ.\end{Resp}
\begin{Resp}\Vsign Propter veritátem, et mansuetúdinem, et justítiam.\end{Resp}
\begin{Resp}\Rsign Proptérea unxit te Deus, Deus tuus, óleo lætítiæ.\end{Resp}
\switchcolumn
\EN
\begin{Resp}\Rsign Thou hast loved righteousness, and hated iniquity; \Star{} Therefore God, thy God, hath anointed thee with the oil of gladness.\end{Resp}
\begin{Resp}\Vsign Because of truth, and meekness, and righteousness.\end{Resp}
\begin{Resp}\Rsign Therefore God, thy God, hath anointed thee with the oil of gladness.\end{Resp}
\switchcolumn*

\LA
\LessonHead{Lectio VI}
\switchcolumn
\EN
\LessonHead{Lesson VI}
\switchcolumn*

\LA
\noindent A sancto Francisco Salesio, quo spiritus moderatore usa fuit, divinam voluntatem edocta, proprium parentem, socerum, filium denique ipsum, quem étiam vocationi obsistentem, sua e domo egrediens, pedibus calcare non dubitavit, invicta constantia deseruit, et sacri instituti Visitationis sanctæ Mariæ fundamenta jecit. Ejus instituti leges integerrime custodivit, et adeo paupertatis fuit amans, ut vel necessaria sibi deesse gauderet. Christianæ vero animi demissionis et obedientiæ, virtutum denique omnium perfectissimum exemplar se præbuit Altiores in corde suo ascensiones disponens, arduissimo efficiendi semper id quod perfectius esse intelligeret, voto se obstrinxit. Denique, sacro Visitationis instituto ejus potissimum opera longe lateque diffuso, verbo, exemplo, et scriptis étiam divina sapientia refertis, ad pietatem et caritatem sororibus excitatis, meritis referta, et sacramentis rite susceptis, Molinis, anno millesimo sexcentesimo quadragesimo primo, die decima tertia Decembris, migravit ad Dominum, ejusque animam, occurrente sancto Francisco Salesio, in cælos deferri sanctus Vincentius a Paulo procul distans aspexit Ejus corpus postea Annecium translatum fuit: eamque miraculis ante et post obitum claram Benedictus decimus quartus beatorum, Clemens vero decimus tertius Pontifex Maximus albo sanctorum adjecit. Festum autem ejusdem die duodecimo Kalendas Septembris ab universa Ecclesia Clemens decimus quartus Pontifex Maximus celebrari præcepit.\par
\switchcolumn
\EN
\noindent She used the services of St. Francis de Sales as her spiritual adviser, and when she learnt from him what was the will of God, she scrupled not to disregard the wishes of her own father, brother-in-law, and even of her son, whom she left with calm determination, went forth from her home, and founded the holy Institution of the Sisters of the Visitation of St. Mary, (at Annecy, upon Trinity Sunday 1610.) She most rigidly kept the rules of this Institute, and loved so well to be poor, that it made her glad to lack even the necessaries of life. She showed herself a model of Christian lowliness, obedience, and all graces. Having settled in her heart still to go up higher and higher towards the Temple of the Lord, she bound herself by a most difficult vow always to do that which she should understand to be best. It was chiefly through her labour that the holy Institute of the Visitation became spread far and wide, and she stirred up the sisters to godliness and love by her words, by her example, and by writings full of Divine wisdom. She duly received the Sacraments before her death, and then, at Moulins, on the 13th day of December, in the year 1641, departed hence, to be for ever with the Lord. St. Vincent of Paul, who was far distant, in a vision beheld her soul borne to heaven, and St. Francis de Sales coming to meet it. Her body was afterwards taken to Annecy. She was famous for miracles both before and after her death, and Pope Benedict XIV. enrolled her among the Blessed, and Pope Clement XIII. among the Saints. Pope Clement XIV. ordered her Feast-day to be kept by the whole Church upon the twenty-first day of August.\par
\switchcolumn*

\LA
\begin{Resp}\Rsign Fallax grátia, et vana est pulchritúdo: \Star{} Múlier timens Dóminum, ipsa laudábitur.\end{Resp}
\begin{Resp}\Vsign Date ei de fructu mánuum suárum, et laudent eam in portis ópera ejus.\end{Resp}
\begin{Resp}\Rsign Múlier timens Dóminum, ipsa laudábitur.\end{Resp}
\begin{Resp}\Vsign Glória Patri, et Fílio, \Star{} et Spirítui Sancto.\end{Resp}
\begin{Resp}\Rsign Múlier timens Dóminum, ipsa laudábitur.\end{Resp}
\switchcolumn
\EN
\begin{Resp}\Rsign Favour is deceitful, and beauty is vain \Star{} A woman that feareth God she shall be praised.\end{Resp}
\begin{Resp}\Vsign Give her of the fruit of her hands, and let her own works praise her in the gates.\end{Resp}
\begin{Resp}\Rsign A woman that feareth God she shall be praised.\end{Resp}
\begin{Resp}\Vsign Glory be to the Father, and to the Son, \Star{} and to the Holy Ghost.\end{Resp}
\begin{Resp}\Rsign A woman that feareth God she shall be praised.\end{Resp}
\switchcolumn*

\end{paracol}

\par\needspace{10\baselineskip}\addvspace{1.2em}{\centering\small\textsc{Die 22 Augusti} · \textsc{22 August}\par}
\Office{Immaculati Cordis Beatæ Mariæ Virginis}{Immaculati Cordis Beatæ Mariæ Virginis}{Duplex II classis}
\addcontentsline{toc}{subsection}{Die 22 Augusti · Immaculati Cordis Beatæ Mariæ Virginis \textit{— Immaculati Cordis Beatæ Mariæ Virginis}}
\begin{paracol}{2}
\LA
\RefLine{Lectiones I–III: ex In Festis Beatae Mariae Virginis (C11).}
\switchcolumn
\EN
\RefLine{Lessons I–III: from the In Festis Beatae Mariae Virginis (C11).}
\switchcolumn*

\LA
\RefLine{Lectio IX: de Ss. Timothei, Hippolyti et Symphoriani Martyres (08-22cc).}
\switchcolumn
\EN
\RefLine{Lesson IX: of Sts. Timothy, Hippolytus and Symphorianus, Martyrs (08-22cc).}
\switchcolumn*

\LA
\LessonHead{Lectio IV}
\switchcolumn
\EN
\LessonHead{Lesson IV}
\switchcolumn*

\LA
\LessonTitle{Sermo sancti Bernardíni Senénsis}
\switchcolumn
\EN
\LessonTitle{From a Sermon by St. Bernardine of Siena}
\switchcolumn*

\LA
\Source{Ex sermone 9 de visitatione}
\switchcolumn
\EN
\Source{From sermon 9 on the Visitation}
\switchcolumn*

\LA
\noindent Quis mortálium, nisi divíno tutus oráculo, de vera Dei et hóminis Genetríce quidquam módicum, sive grande præsúmat incircumcísis, immo pollútis lábiis nomináre, quam Pater ante sǽcula Deus perpétuam prædestinávit in Vírginem, digníssimam Fílius elégit in Matrem, Spíritus Sanctus omnis grátiæ domicílium præparávit? Quibus verbis ego homúnculus sensus altíssimos virgínei Cordis, sanctíssimo ore prolátos, éfferam, quibus non súfficit lingua ómnium Angelórum? Dóminus enim ait: Bonus homo de bono thesáuro cordis profert bona; quod verbum potest étiam esse thesáurus. Quis inter puros hómines mélior homo potest excogitári, quam illa, quæ méruit éffici Mater Dei, quæ novem ménsibus in corde et in útero suo ipsum Deum hospitáta est? Quis thesáurus mélior, quam ipse divínus amor quo fornáceum cor Vírginis ardens erat?\par
\switchcolumn
\EN
\noindent What man, unless secure in a divine oracle, may presume to speak with impure, indeed with polluted lips, anything little or great about the true Parent of God and of man, whom the Father before all ages predestined a perpetual Virgin, whom the Son chose as his most worthy Mother, whom the Holy Ghost prepared as the dwelling place of every grace? With what words shall I, a lowly man, give expression to the highest sentiments of the virginal Heart uttered by the holiest mouth, for which the tongues of all the Angels do not suffice? For the Lord saith: A good man bringeth forth good things from the good treasure of his heart; and this word can also be a treasure. Among pure mortals who can be conceived of as better than she who was worthy to be the Mother of God, who for nine months had as a guest in her heart and in her womb God himself? What better treasure than the divine love itself, which was burning in the Heart of the Virgin as in a furnace?\par
\switchcolumn*

\LA
\begin{Resp}\Rsign Sicut cedrus exaltáta sum in Líbano, et sicut cypréssus in monte Sion: quasi myrrha elécta, \Star{} Dedi suavitátem odóris.\end{Resp}
\begin{Resp}\Vsign Et sicut cinnamómum et bálsamum aromatízans.\end{Resp}
\begin{Resp}\Rsign Dedi suavitátem odóris.\end{Resp}
\switchcolumn
\EN
\begin{Resp}\Rsign I was exalted like a cedar in Lebanon, and as a cypress-tree upon Mount Zion. Like the best myrrh \Star{} I yielded a pleasant odour.\end{Resp}
\begin{Resp}\Vsign Like cinnamon and sweet balsam.\end{Resp}
\begin{Resp}\Rsign I yielded a pleasant odour.\end{Resp}
\switchcolumn*

\LA
\LessonHead{Lectio V}
\switchcolumn
\EN
\LessonHead{Lesson V}
\switchcolumn*

\LA
\noindent De hoc ígitur Corde quasi de fornáce divíni ardóris Virgo beáta prótulit verba bona, id est, verba ardentíssimæ caritátis. Sicut enim a vase summo et óptimo vino pleno, non potest exíre nisi óptimum vinum; aut sicut a fornáce summi ardóris non egréditur nisi incéndium fervens; sic quippe a Christi Matre exíre non pótuit verbum, nisi summi summéque divíni amóris atque ardóris. Sapiéntis quoque dóminæ et matrónæ est pauca verba, sólida tamen atque sententiósa habére; proínde septem vícibus quasi septem verba tantum miræ senténtiæ et virtútis a Christi benedictíssima Matre legúntur dicta, ut mýstice ostendátur ipsam fuísse plenam grátia septifórmi. Cum Angelo bis tantúmmodo est locúta. Cum Elísabeth bis étiam. Cum Fílio étiam bis, semel in templo, secúndo in núptiis. Cum minístris semel. Et in his ómnibus semper valde parum locúta est; excépto quod in laude Dei et gratiárum actióne se ámplius dilatávit, scílicet, quum ait: Magníficat ánima mea Dóminum. Ubi non cum hómine, sed cum Deo locúta fuit. Hæc septem verba secúndum septem amóris procéssus et actus sub miro gradu et órdine sunt proláta; quasi sint septem flammæ fornácei Cordis ejus.\par
\switchcolumn
\EN
\noindent And so, from this Heart as from a furnace of divine ardor the Blessed Virgin brought forth good words, that is, words of the most ardent charity. For as from a vessel full of the richest and best wine only good wine can be poured; or as from a furnace of intense heat only a burning fire is emítted; so indeed from the Mother of Christ no word can go forth except of the greatest and most intense divine love and ardor. It is also the mark of a wise woman and matron to speak few words, but words that are effective and full of meaning; and so seven times, as it were, seven words of such wonderful meaning and virtue are read as having been uttered by the Most Blessed Mother of Christ, that mystically it may be shown she was full of the sevenfold grace. To the Angel twice only did she speak; to Elizabeth also twice; with her Son likewise twice, once in the temple, and a second time at the marriage feast; and once to the attendants. And on all those occasions she always said very little; with this one exception that she spake at length in the praise of God and in thanksgiving, namely, when she said: My soul doth magnify the Lord. But here she did not speak with man, but with God. Those seven words were spoken in a wonderful degree and order according to the seven courses and acts of love; as if they were seven flames from the furnace of her Heart.\par
\switchcolumn*

\LA
\begin{Resp}\Rsign Quæ est ista quæ procéssit sicut sol, et formósa tamquam Jerúsalem? \Star{} Vidérunt eam fíliæ Sion, et beátam dixérunt, et regínæ laudavérunt eam.\end{Resp}
\begin{Resp}\Vsign Et sicut dies verni circúmdabant eam flores rosárum et lília convállium.\end{Resp}
\begin{Resp}\Rsign Vidérunt eam fíliæ Sion, et beátam dixérunt, et regínæ laudavérunt eam.\end{Resp}
\switchcolumn
\EN
\begin{Resp}\Rsign Who is this that cometh up like the sun? This, comely as Jerusalem? \Star{} The daughters of Zion saw her, and called her blessed: the queens also, and they praised her.\end{Resp}
\begin{Resp}\Vsign And about her it was as the flower of roses in the spring of the year, and lilies of the valleys.\end{Resp}
\begin{Resp}\Rsign The daughters of Zion saw her, and called her blessed the queens also, and they praised her.\end{Resp}
\switchcolumn*

\LA
\LessonHead{Lectio VI}
\switchcolumn
\EN
\LessonHead{Lesson VI}
\switchcolumn*

\LA
\Source{Ex ecclesiasticis documentis}
\switchcolumn
\EN
\Source{From ecclesiastical documents}
\switchcolumn*

\LA
\noindent Cultum litúrgicum, quo Cordi Immaculáto Vírginis Maríæ débitus tribúitur honor, cuíque plures viri sancti ac mulíeres viam parárunt, ipsa Apostólica Sedes primum approbávit ineúnte sǽculo undevicésimo, cum Pius Papa séptimus festum Puríssimi Cordis Maríæ Vírginis instítuit, ab ómnibus diœcésibus et religiósis famíliis, quæ id petiíssent, pie sanctéque agéndum: quod póstmodum Pius Papa nonus Offício ac Missa própria auxit. Ardens autem stúdium atque optátum, jam sǽculo décimo séptimo exórtum et in dies invaléscens, ut nempe ejúsmodi festum, majóri solemnitáte donátum, totíus Ecclésiæ commúne efficerétur, Summus Póntifex Pius duodécimus benígne excípiens anno millésimo nongentésimo quadragésimo secúndo, bello atrocíssimo per orbem fere totum ingravescénte, infínitas populórum ærúmnas míserans, pro sua in Matrem cæléstem pietáte ac fidúcia genus hóminum univérsum illíus Cordi benigníssimo obsecratióne solémni eníxe commendávit, atque in honórem ejúsdem Immaculáti Cordis festum cum Offício et Missa própriis in perpétuum ubíque celebrándum indíxit.\par
\switchcolumn
\EN
\noindent The liturgical worship, through which due honor is given to the Immaculate Heart of the Virgin Mary, and for which many holy men and women have prepared the way, the Apostolic See itself first approved in the beginning of the nineteenth century, when Pope Pius VII instituted the feast of the Most Pure Heart of the Virgin Mary, to be piously and reverently celebrated by all the dioceses and religious families who had asked for it. Afterwards Pope Pius IX added an Office and a proper Mass to it. But an ardent desire and longing, which had arisen in the seventeenth century, grew day by day, that namely, the same Feast, given greater solemnity, might be spread to the entire Church. In 1942, Pope Pius XII, graciously acceding to this wish, and during the terrible war then ravaging almost the entire world, pitying the infinite hardships of men, and because of his devotion and confidence in our heavenly Mother, in solemn supplication earnestly entrusted the entire human race to her most generous Heart, and in honor of the same Immaculate Heart he ordered a Feast to be kept forever with its proper Office and Mass.\par
\switchcolumn*

\LA
\begin{Resp}\Rsign Ornátam monílibus fíliam Jerúsalem Dóminus concupívit: \Star{} Et vidéntes eam fíliæ Sion, beatíssimam prædicavérunt, dicéntes: * Unguéntum effúsum nomen tuum.\end{Resp}
\begin{Resp}\Vsign Astitit regína a dextris tuis in vestítu deauráto, circúmdata varietáte.\end{Resp}
\begin{Resp}\Rsign Et vidéntes eam fíliæ Sion, beatíssimam prædicavérunt, dicéntes:\end{Resp}
\begin{Resp}\Vsign Glória Patri, et Fílio, \Star{} et Spirítui Sancto.\end{Resp}
\begin{Resp}\Rsign Unguéntum effúsum nomen tuum.\end{Resp}
\switchcolumn
\EN
\begin{Resp}\Rsign When the Lord beheld the daughter of Jerusalem adorned with her jewels, He greatly desired her beauty \Star{} And when the daughters of Zion saw her, they cried out that she was most blessed, saying thy name is as ointment poured forth.\end{Resp}
\begin{Resp}\Vsign Upon thy right hand did stand the Queen in a vesture of gold wrought about with diverse colours.\end{Resp}
\begin{Resp}\Rsign And when the daughters of Zion saw her, they cried out that she was most blessed, saying:\end{Resp}
\begin{Resp}\Vsign Glory be to the Father, and to the Son, \Star{} and to the Holy Ghost.\end{Resp}
\begin{Resp}\Rsign Thy name is as ointment poured forth.\end{Resp}
\switchcolumn*

\LA
\LessonHead{Lectio VII}
\switchcolumn
\EN
\LessonHead{Lesson VII}
\switchcolumn*

\LA
\LessonTitle{Léctio sancti Evangélii secúndum Joánnem}
\switchcolumn
\EN
\LessonTitle{From the Holy Gospel according to John}
\switchcolumn*

\LA
\Cite{Joannes 19:25-27}
\switchcolumn
\EN
\Cite{John 19:25-27}
\switchcolumn*

\LA
\noindent In illo témpore: Stabant juxta crucem Jesu Mater ejus, et soror Matris ejus María Cléophæ, et María Magdaléne. Et réliqua.\par
\switchcolumn
\EN
\noindent In that time: There stood by the cross of Jesus, his mother, and his mother's sister, Mary of Cleophas, and Mary Magdalen. And so on.\par
\switchcolumn*

\LA
\LessonTitle{Homilía sancti Robérti Bellarmíno Epíscopi}
\switchcolumn
\EN
\LessonTitle{A Homily by St. Robert Bellarmine, Bishop}
\switchcolumn*

\LA
\Source{De septem verbis Christi in Cruce, cap. 12}
\switchcolumn
\EN
\Source{On the Seven Words of Christ on the Cross cap. 12}
\switchcolumn*

\LA
\noindent Onus et jugum impósitum a Dómino sancto Joánni, ut Vírginis Matris curam géreret, vere fuit jugum suáve et onus leve. Quis enim non libentíssime cohabitáret Matri illi, quæ Verbum incarnátum in útero novem ménsibus portávit, et illi totos trigínta annos devotíssime dulcissiméque cohabitávit? Quis non invídeat dilécto Dómini, qui in abséntia Fílii Dei præséntiam obtínuit Matris Dei? Sed, nisi fallor, póssumus et nos a benignitáte Verbi nostri causa incarnáti et ex dilectióne nímia nostri causa crucifíxi, précibus impetráre, ut dicat et nobis: Ecce Mater tua; et Matri suæ de nobis dicat: Ecce fílius tuus.\par
\switchcolumn
\EN
\noindent The burden and yoke which our Lord imposed on St. John, that he take care of his Virgin Mother, was indeed a sweet yoke and a light burden. Who indeed would not esteem it a happiness to dwell under the same roof with her, who for nine months had borne in her womb the Incarnate Word, and for thirty years had enjoyed the sweetest and happiest communication of sentiments with him? Who doth not envy the chosen disciple of our Lord, who in the absence of the Son of God, was given the presence of the Mother of God? yet, if I am not mistaken, we can obtain by our prayers that our most kind Lord, who became man for our sakes and was crucified for love of us, should say to us: Behold thy Mother; and should say to his Mother for each one of us: Behold thy son.\par
\switchcolumn*

\LA
\begin{Resp}\Rsign Felix namque es, sacra Virgo María, et omni laude digníssima: \Star{} Quia ex te ortus est sol justítiæ, Christus, Deus noster.\end{Resp}
\begin{Resp}\Vsign Ora pro pópulo, intérveni pro clero, intercéde pro devóto femíneo sexu: séntiant omnes tuum juvámen, quicúmque célebrant tuam sanctam festivitátem.\end{Resp}
\begin{Resp}\Rsign Quia ex te ortus est sol justítiæ, Christus, Deus noster.\end{Resp}
\switchcolumn
\EN
\begin{Resp}\Rsign O Holy Virgin Mary, happy indeed art thou, and right worthy of all praise \Star{} For out of thee rose the Sun of righteousness, even Christ our God.\end{Resp}
\begin{Resp}\Vsign Pray for the people, plead for the clergy, make intercession for all women vowed to God. May all that are keeping this thine holy Feast day feel the might of thine assistance.\end{Resp}
\begin{Resp}\Rsign For out of thee rose the Sun of righteousness, even Christ our God.\end{Resp}
\switchcolumn*

\LA
\LessonHead{Lectio VIII}
\switchcolumn
\EN
\LessonHead{Lesson VIII}
\switchcolumn*

\LA
\noindent Non est avárus pius Dóminus gratiárum, dúmmodo ad thronum grátiæ ejus cum fide et fidúcia et non ficto corde, sed vero et sincéro accedámus. Qui nos coherédes esse vóluit regni Patris sui, non dedignábitur certe nos coherédes habére amóris Matris suæ. Sed nec ipsa Virgo benigníssima gravábitur multitúdine filiórum, cum sinum amplíssimum hábeat et valde cúpiat, nullum períre ex his, quos Fílius suus tam pretióso sánguine et tam pretiósa morte redémit. Adeámus ergo cum fidúcia ad thronum grátiæ Christi, et supplíciter nec sine lácrimis ab eo petámus, ut de unoquóque nostrum Matri suæ dicat: Ecce fílius tuus; et unicuíque nostrum de Matre sua dicat: Ecce Mater tua.\par
\switchcolumn
\EN
\noindent Our good Lord is not avaricious of his graces, if only we approach the throne of grace with faith and confidence, with a true and sincere but not a false heart. He who desireth to have us co-heirs in the kingdom of his Father, will certainly not disdain to have us co-heirs in the love of his Mother. Nor will the most benign Virgin herself take it amiss to have a countless number of children, since she hath an heart capable of embracing all of us, and ardently desireth that not even one of those souls should perish whom her divine Son redeemed with his precious Blood, and his still more precious death. And so let us approach with confidence the throne of the grace of Christ, and with tears let us humbly beg of him to say to his Mother for each of us: Behold thy son; and to each one of us concerning his Mother: Behold thy Mother.\par
\switchcolumn*

\LA
\begin{Resp}\Rsign Beátam me dicent omnes generatiónes, \Star{} Quia fecit mihi Dóminus magna qui potens est, et sanctum nomen ejus.\end{Resp}
\begin{Resp}\Vsign Et misericórdia ejus a progénie in progénies timéntibus eum.\end{Resp}
\begin{Resp}\Rsign Quia fecit mihi Dóminus magna qui potens est, et sanctum nomen ejus.\end{Resp}
\begin{Resp}\Vsign Glória Patri, et Fílio, \Star{} et Spirítui Sancto.\end{Resp}
\begin{Resp}\Rsign Quia fecit mihi Dóminus magna qui potens est, et sanctum nomen ejus.\end{Resp}
\switchcolumn
\EN
\begin{Resp}\Rsign All generations shall call me blessed. \Star{} For He that is mighty, even the Lord, hath done to me great things; and Holy is His Name.\end{Resp}
\begin{Resp}\Vsign And his mercy is on them that fear Him, from generation to generation.\end{Resp}
\begin{Resp}\Rsign He that is mighty, even the Lord, hath done to me great things; and Holy is His Name.\end{Resp}
\begin{Resp}\Vsign Glory be to the Father, and to the Son, \Star{} and to the Holy Ghost.\end{Resp}
\begin{Resp}\Rsign He that is mighty, even the Lord, hath done to me great things; and Holy is His Name.\end{Resp}
\switchcolumn*

\end{paracol}

\par\needspace{10\baselineskip}\addvspace{1.2em}{\centering\small\textsc{Die 22 Augusti} · \textsc{22 August}\par}
\Office{Ss. Timothei, Hippolyti et Symphoriani Martyres}{Sts. Timothy, Hippolytus and Symphorianus, Martyrs}{Simplex}
\addcontentsline{toc}{subsection}{Die 22 Augusti · Ss. Timothei, Hippolyti et Symphoriani Martyres \textit{— Sts. Timothy, Hippolytus and Symphorianus, Martyrs}}
\begin{paracol}{2}
\LA
\LessonHead{Lectio IX (pro commemoratione)}
\switchcolumn
\EN
\LessonHead{Lesson IX (when commemorated)}
\switchcolumn*

\LA
\LessonTitle{Commemoratio: Ss. Timothei, Hippolyti et Symphoriani Martyres}
\switchcolumn
\EN
\LessonTitle{Commemoration of: Sts. Timothy, Hippolytus and Symphorianus, Martyrs}
\switchcolumn*

\LA
\noindent Timotheus Antiochenus, Romam veniens, Melchiade Summo Pontifice, cum per annum ibi Christi fidem prædicasset, a Tarquinio Urbis præfecto conjicitur in vincula: et post diuturnas carceris ærumnas ad idola perducitur, ut eis sacrificet. Quam impietatem summa libertate detestatus, acerbissime cæditur, et excarnificatum corpus viva calce perfunditur. In quibus aliusque suppliciis constans Martyr capite plectitur. Corpus via Ostiensi prope sepulchrum beati Pauli Apostoli sepelitur. Quo die etiam, Alexandro imperatore, apud Ostia iberina Hippolytus episcopus Portuensis ob præclaram fidei confessionem manibus pedibusque ligatis, in altam foveam aquis plenam præcipitatus, martyrio coronatus est, et ibidem a Christianis sepultus. Quo item die, Aureliano imperatore, propter eamdem fidem Augustodúni Symphorianus adolescens varie tortus est. Qui dum ad ultimum supplicium duceretur, matrem ita clamantem audiens: Nate, nate, memento æternæ vitæ, cælum suspice, et ibi regnantem intuere: tibi enim vita non eripitur, sed mutatur in melius: fortiter Jesu Christi causa carnifici collum præbuit.\par
\switchcolumn
\EN
\noindent Timothy came from Antioch to Rome in the time of Pope Melchiades. He had preached the faith of Christ there for a year, when he was thrown into irons by Tarquinius, Praefect of the city. After suffering a long imprisonment he was brought to the idols to offer them sacrifice. He refused right boldly to commit this great sin, and was thereupon savagely scourged, and his raw body covered with quick-lime. He steadily persisted in his testimony under these and other tortures, and at last was beheaded, (in the year 311.) His body is buried upon the road to Ostia, hard by the sepulchre of the blessed Apostle Paul. On the same day, under the Emperor Alexander, and at Ostia, Hippolytus, Bishop of Porto, on account of his illustrious confession of the faith, had his hands and feet bound, and was thrown into a deep pit full of water, and so received the crown of his testimony. The Christians buried him there. Also on the same day, \textasciitilde{}(in the year 180,) under the Emperor Aurelian, and at Autun, the young lad Symphorian was tortured in diverse ways for professing the same faith. As he was being led to die, he heard his mother crying out to him My child, my child think of life eternal Look to heaven and to Him That reigneth there! thy life is not being taken away, but changed for a better. And so, for Jesus Christ's sake, he bravely offered his neck to the executioner.\par
\switchcolumn*

\LA
\TeDeum{Te Deum}
\switchcolumn
\EN
\TeDeum{Te Deum}
\switchcolumn*

\end{paracol}

\par\needspace{10\baselineskip}\addvspace{1.2em}{\centering\small\textsc{Die 23 Augusti} · \textsc{23 August}\par}
\Office{S. Philippi Benitii Confessoris}{St. Philip Benizi, Confessor}{Duplex}
\addcontentsline{toc}{subsection}{Die 23 Augusti · S. Philippi Benitii Confessoris \textit{— St. Philip Benizi, Confessor}}
\begin{paracol}{2}
\LA
\RefLine{Lectiones I–III: de Scriptura occurrente.}
\switchcolumn
\EN
\RefLine{Lessons I–III: from the Scripture of the day.}
\switchcolumn*

\LA
\RefLine{Lectio IX: de In Vigilia S.Bartholomæi Apostoli (08-23o).}
\switchcolumn
\EN
\RefLine{Lesson IX: of Vigil of St. Bartholomew, Apostle (08-23o).}
\switchcolumn*

\LA
\LessonHead{Lectio IV}
\switchcolumn
\EN
\LessonHead{Lesson IV}
\switchcolumn*

\LA
\noindent Philippus ex nobili Benitiorum familia Florentiæ natus, futuræ sanctitatis jam inde ab incunabulis indicium præbuit. Vix enim quintum ætatis mensem ingressus linguam in voces mirifice solvit, hortatusque fuit matrem, ut Deiparæ Servis eleemosynam impertiret. Adolescens, dum Parisiis litterarum studia cum pietatis ardore conjungeret, plurimos ad cælestis patriæ desiderium inflammavit. Reversus in patriam, et singulari visione a beatissima Virgine in Servorum suorum familiam nuper institutam vocatus, in Senarii montis antrum concessit, ubi asperam quidem jugi corporis castigatione, sed Christi Domini cruciatum meditatione suavem vitam duxit: indeque per universam pene Europam, magnamque Asiæ partem, quam evangelicis prædicationibus obivit, sodalitia Septem Dolorum Dei Matris instituit, suumque ordinem eximio virtutum exemplo propagavit.\par
\switchcolumn
\EN
\noindent This Philip was a scion of the noble Florentine family of the Benizi from his very cradle he showed signs of holiness. When he had scarcely entered the fifth month of his life, his cries marvellously assumed the form of words, entreating his mother to give some alms to the servants of the Mother of God. While he was a young man at Paris studying letters, but ever of a fervent piety, he stirred up in many the love of our Fatherland which is in heaven. After his return to his own country, the most Blessed Virgin appeared to him in a vision, and specially called on him to enter the Order of her Servants, which had then been newly founded. He withdrew himself to a cave on Monte Senario, where he led a life hard as touching the chastisement of the flesh, but sweet with thoughts of the agonies of Christ. Thence he came forth and went through nearly all Europe and great part of Asia, preaching the Gospel, founding Guilds everywhere in honour of the Seven Sorrows of the Mother of God, and extending his Order by the wonderful example of his own holy life.\par
\switchcolumn*

\LA
\begin{Resp}\Rsign Honéstum fecit illum Dóminus, et custodívit eum ab inimícis, et a seductóribus tutávit illum: \Star{} Et dedit illi claritátem ætérnam.\end{Resp}
\begin{Resp}\Vsign Justum dedúxit Dóminus per vias rectas, et osténdit illi regnum Dei.\end{Resp}
\begin{Resp}\Rsign Et dedit illi claritátem ætérnam.\end{Resp}
\switchcolumn
\EN
\begin{Resp}\Rsign The Lord made him honourable, and defended him from his enemies, and kept him safe from those that lay in wait for him; \Star{} And gave him perpetual glory.\end{Resp}
\begin{Resp}\Vsign He went down with him into the pit, and left him not in bonds.\end{Resp}
\begin{Resp}\Rsign And gave him perpetual glory.\end{Resp}
\switchcolumn*

\LA
\LessonHead{Lectio V}
\switchcolumn
\EN
\LessonHead{Lesson V}
\switchcolumn*

\LA
\noindent Divinæ caritatis et catholicæ fidei dilatandæ ardore vehementer accensus, sui ordinis Generalis reluctans atque invitus renuntiatus, fratres ad prædicandum Christi Evangelium in Scythiam misit: ipse vero plurimas Italiæ urbes concursans, gliscentes in eis civium discordias composuit, multasque ad Romani Pontificis obedientiam revocavit; nihilque de studio alienæ salutis omittens, perditissimos homines e vitiorum cœno ad pœnitentiam ac Jesu Christi amórem perduxit. Orationi summopere addictus, sæpe in extasim rapi visus est. Virginitatem vero adeo coluit, ut ad extremum usque spiritum voluntariis ac durissimis suppliciis illibatam custodierit.\par
\switchcolumn
\EN
\noindent He was forced against his own wishes to undertake the duties of General of his Order, and, in his love of God and of the spreading of the Catholic Faith, sent forth brethren to preach the Gospel of Christ in Russia. He himself went through many cities of Italy, stilled the raging quarrels of the inhabitants, and recalled many of them to their obedience to the Bishop of Rome. He left nothing undone to forward the salvation of his neighbour, and brought the most depraved wretches to leave the slough of their sins, to do penance, and to love Jesus Christ. He was most earnest in prayer, and was often seen to fall into trances while engaged in it. Virginity he so prized that to his very last breath he kept it unsullied by dint of self-imposed and stern penances.\par
\switchcolumn*

\LA
\begin{Resp}\Rsign Amávit eum Dóminus, et ornávit eum: stolam glóriæ índuit eum, \Star{} Et ad portas paradísi coronávit eum.\end{Resp}
\begin{Resp}\Vsign Índuit eum Dóminus lorícam fídei, et ornávit eum.\end{Resp}
\begin{Resp}\Rsign Et ad portas paradísi coronávit eum.\end{Resp}
\switchcolumn
\EN
\begin{Resp}\Rsign The Lord loved him and beautified him; He clothed him with a robe of glory, \Star{} And crowned him at the gates of Paradise.\end{Resp}
\begin{Resp}\Vsign The Lord hath put on him the breast-plate of faith, and hath adorned him.\end{Resp}
\begin{Resp}\Rsign And crowned him at the gates of Paradise.\end{Resp}
\switchcolumn*

\LA
\LessonHead{Lectio VI}
\switchcolumn
\EN
\LessonHead{Lesson VI}
\switchcolumn*

\LA
\noindent Effloruit in eo jugiter singularis erga pauperes misericordia, sed præcipue cum apud Camilianum agri Senensis vicum leproso nudo eleemosynam petenti propriam, qua indutus erat, vestem fuit elargitus: qua ille contectus, statim a lepra mundatus est. Cujus miraculi cum longe lateque fama manasset, nonnulli ex cardinalibus, qui Viterbium, Clemente quarto vita functo, pro successore deligendo convenerant, in Philippum, cujus cælestem étiam prudentiam perspectam habebant, intenderunt. Quo comperto vir Dei, ne forte pastoralis regiminis onus subire cogeretur, apud Tuniatum montem tamdiu delituit, donec Gregorius decimus Pontifex Maximus fuerit renuntiatus: ubi balneis, quæ étiam hodie sancti Philippi vocantur, virtutem sanandi morbos suis precibus impetravit. Denique Tuderti, anno millesimo ducentesimo octogesimo quinto in Christi Domini e cruce pendentis amplexu, quem suum appellabat librum, sanctissime ex hac vita migravit. Ad ejus tumulum cæci visum, claudi gressum, mortui vitam receperunt. Quibus aliisque plurimis fulgentem signis Clemens decimus Pontifex Maximus sanctorum numero adscripsit.\par
\switchcolumn
\EN
\noindent Everywhere appeared in him an extraordinary pity towards the poor, whereof it is a famous instance that at the village of Camiliano in the territory of Sienna he gave his own garment to a naked leper who asked him for an alms, and as soon as the said leper had cast it about him he was straightway cleansed of his leprosy. The fame of this miracle spread far and wide, and some of the Cardinals who had assembled at Viterbo after the death of Clement IV., to elect a successor to him, cast their eyes upon Philip, with whose heavenly wisdom they were also acquainted. When the man of God found how things stood, lest he should be constrained to take upon him the burden of the Pastoral Office, he went and hid himself on Montagnate, until Gregory X. had been proclaimed Pope. By his prayers he obtained medicinal powers for the waters in these mountains, which are still called St. Philip's Baths. At length, (on the 22nd of August,) in the year 1285, he departed this life in a most holy manner at Todi, while embracing the image of Christ hanging upon the Cross, which he called his book. At his grave the blind received their sight, the lame walked, and the dead were raised. Pope Clement X., finding him famous for these and many other great signs and wonders, enrolled his name among those of the Saints.\par
\switchcolumn*

\LA
\begin{Resp}\Rsign Iste homo perfécit ómnia quæ locútus est ei Deus, et dixit ad eum: Ingrédere in réquiem meam: \Star{} Quia te vidi justum coram me ex ómnibus géntibus.\end{Resp}
\begin{Resp}\Vsign Iste est, qui contémpsit vitam mundi, et pervénit ad cæléstia regna.\end{Resp}
\begin{Resp}\Rsign Quia te vidi justum coram me ex ómnibus géntibus.\end{Resp}
\begin{Resp}\Vsign Glória Patri, et Fílio, \Star{} et Spirítui Sancto.\end{Resp}
\begin{Resp}\Rsign Quia te vidi justum coram me ex ómnibus géntibus.\end{Resp}
\switchcolumn
\EN
\begin{Resp}\Rsign This is he which did according unto all that God commanded him; and God said unto him: Enter thou into My rest, \Star{} For thee have I seen righteous before Me among all people.\end{Resp}
\begin{Resp}\Vsign This is he which loved not his life in this world, and is come unto an everlasting kingdom.\end{Resp}
\begin{Resp}\Rsign For thee have I seen righteous before Me among all people.\end{Resp}
\begin{Resp}\Vsign Glory be to the Father, and to the Son, \Star{} and to the Holy Ghost.\end{Resp}
\begin{Resp}\Rsign For thee have I seen righteous before Me among all people.\end{Resp}
\switchcolumn*

\LA
\LessonHead{Lectio VII}
\switchcolumn
\EN
\LessonHead{Lesson VII}
\switchcolumn*

\LA
\LessonTitle{Léctio sancti Evangélii secúndum Lucam}
\switchcolumn
\EN
\LessonTitle{From the Holy Gospel according to Luke}
\switchcolumn*

\LA
\Cite{Luc 12:32-34}
\switchcolumn
\EN
\Cite{Luke 12:32-35}
\switchcolumn*

\LA
\noindent In illo témpore: Dixit Jesus discípulis suis: Nolíte timére pusíllus grex, quia complácuit Patri vestro dare vobis regnum. Et réliqua.\par
\switchcolumn
\EN
\noindent At that time, Jesus said unto His disciples: Fear not, little flock, for it is your Father's good pleasure to give you the kingdom. And so on.\par
\switchcolumn*

\LA
\LessonTitle{Homilía sancti Bedæ Venerábilis Presbýteri}
\switchcolumn
\EN
\LessonTitle{Homily by the Venerable Bede, Priest at Jarrow, and Doctor of the Church.}
\switchcolumn*

\LA
\Source{Lib. 4, Cap. 54, in Luc. 12}
\switchcolumn
\EN
\Source{Bk. iv. Ch. 54 on Luke xii.}
\switchcolumn*

\LA
\noindent Pusíllum gregem electórum, vel ob comparatiónem majóris númeri reprobórum, vel pótius ob humilitátis devotiónem nóminat: quia vidélicet Ecclésiam suam quantálibet numerositáte jam dilatátam, tamen usque ad finem mundi humilitáte vult créscere, et ad promíssum regnum humilitáte perveníre. Ideóque ejus labóres blande consolátus, quam regnum Dei tantum quærére prǽcipit, eídem regnum a Patre dandum complácita benignitáte promíttit.\par
\switchcolumn
\EN
\noindent The elect are called a little flock, perchance because the reprobate are far more in number than they, but, more probably, because they love to be lowly, since it is God's will that however much His Church should grow in numbers, she should grow with lowliness even unto the end of the world, and should enter lowly into that kingdom which is hers by His promise. That kingdom He promiseth to her here, when He biddeth her to seek only the kingdom of God, and, to comfort her in her travail, He doth so sweetly and so graciously say that her Father will give it to her.\par
\switchcolumn*

\LA
\begin{Resp}\Rsign Iste est qui ante Deum magnas virtútes operátus est, et de omni corde suo laudávit Dóminum: \Star{} Ipse intercédat pro peccátis ómnium populórum.\end{Resp}
\begin{Resp}\Vsign Ecce homo sine queréla, verus Dei cultor, ábstinens se ab omni ópere malo, et pérmanens in innocéntia sua.\end{Resp}
\begin{Resp}\Rsign Ipse intercédat pro peccátis ómnium populórum.\end{Resp}
\switchcolumn
\EN
\begin{Resp}\Rsign This is he which wrought great wonders before God, and praised the Lord with all his heart. \Star{} May he pray for all people, that their sins may be forgiven unto them.\end{Resp}
\begin{Resp}\Vsign Behold a man without blame, a worshipper of God in truth, keeping himself clean from every evil work, and abiding still in his innocency.\end{Resp}
\begin{Resp}\Rsign May he pray for all people, that their sins may be forgiven unto them.\end{Resp}
\switchcolumn*

\LA
\LessonHead{Lectio VIII}
\switchcolumn
\EN
\LessonHead{Lesson VIII}
\switchcolumn*

\LA
\noindent Véndite quæ possidétis, et date eleemósynam. Nolíte, inquit, timére, ne propter regnum Dei militántibus, hujus vitæ necessária desint: quin étiam posséssa propter eleemósynam véndite. Quod tunc digne fit, quando quis semel pro Dómino suis ómnibus spretis, nihilóminus post hæc labóre mánuum, unde et victum transígere, et eleemósynam dare queat, operátur. Unde gloriátur Apóstolus, dicens: Argéntum, et aurum, aut vestem nullíus concupívi: ipsi scitis, quóniam ad ea quæ mihi opus erant, et his qui mecum sunt, ministravérunt manus istæ. Ómnia osténdi vobis, quóniam sic laborántes opórtet suscípere infírmos.\par
\switchcolumn
\EN
\noindent Sell what ye have and give alms. Fear not, He saith, lest, while ye fight for the kingdom of God, ye should lack such things as are needful for this life, nay rather, sell even that which ye have, and give alms. This doth, whosoever for the Lord's sake leaveth all that he hath, and then worketh with his hands, that so he may have to eat, and withal to give alms. In this doth the Apostle boast himself, saying I have coveted no man's silver, or gold, or apparel, as ye yourselves know for these hands have ministered unto my necessities, and to them that were with me. I have showed you all things, how that so labouring ye ought to support the weak. (Acts xx. 33, 34, 35.)\par
\switchcolumn*

\LA
\begin{Resp}\Rsign Sint lumbi vestri præcíncti, et lucérnæ ardéntes in mánibus vestris: \Star{} Et vos símiles homínibus exspectántibus dóminum suum, quando revertátur a núptiis.\end{Resp}
\begin{Resp}\Vsign Vigiláte ergo, quia nescítis qua hora Dóminus vester ventúrus sit.\end{Resp}
\begin{Resp}\Rsign Et vos símiles homínibus exspectántibus dóminum suum, quando revertátur a núptiis.\end{Resp}
\begin{Resp}\Vsign Glória Patri, et Fílio, \Star{} et Spirítui Sancto.\end{Resp}
\begin{Resp}\Rsign Et vos símiles homínibus exspectántibus dóminum suum, quando revertátur a núptiis.\end{Resp}
\switchcolumn
\EN
\begin{Resp}\Rsign Let your loins be girded about, and your lights burning; \Star{} And ye yourselves like unto men that wait for their lord, when he will return from the wedding.\end{Resp}
\begin{Resp}\Vsign Watch therefore, for ye know not what hour your Lord doth come.\end{Resp}
\begin{Resp}\Rsign And ye yourselves like unto men that wait for their lord, when he will return from the wedding.\end{Resp}
\begin{Resp}\Vsign Glory be to the Father, and to the Son, \Star{} and to the Holy Ghost.\end{Resp}
\begin{Resp}\Rsign And ye yourselves like unto men that wait for their lord, when he will return from the wedding.\end{Resp}
\switchcolumn*

\end{paracol}

\par\needspace{10\baselineskip}\addvspace{1.2em}{\centering\small\textsc{Die 23 Augusti} · \textsc{23 August}\par}
\Office{In Vigilia S.Bartholomæi Apostoli}{Vigil of St. Bartholomew, Apostle}{Simplex}
\addcontentsline{toc}{subsection}{Die 23 Augusti · In Vigilia S.Bartholomæi Apostoli \textit{— Vigil of St. Bartholomew, Apostle}}
\begin{paracol}{2}
\LA
\LessonHead{Lectio I}
\switchcolumn
\EN
\LessonHead{Lesson I}
\switchcolumn*

\LA
\LessonTitle{Commemoratio: In Vigilia S.Bartholomæi Apostoli Léctio sancti Evangélii secúndum Joánnem.}
\switchcolumn
\EN
\LessonTitle{Commemoration of: Vigil of St. Bartholomew, Apostle Continuation of the Holy Gospel according to John}
\switchcolumn*

\LA
\Cite{Joannes 15:12-16}
\switchcolumn
\EN
\Cite{John 15:12-16}
\switchcolumn*

\LA
\noindent In illo témpore: Dixit Jesus discípulis suis: Hoc est præcéptum meum, ut diligátis ínvicem, sicut diléxi vos. Et réliqua.\par
\switchcolumn
\EN
\noindent In that time Jesus said to his disciples: This is my commandment, that you love one another, as I have loved you. And so on.\par
\switchcolumn*

\LA
\LessonTitle{Homilía S. Gregórii Papæ.}
\switchcolumn
\EN
\LessonTitle{Homily by Pope St Gregory (the Great.)}
\switchcolumn*

\LA
\Source{Homilia 27. in Evangelia.}
\switchcolumn
\EN
\Source{21th on the Gospels}
\switchcolumn*

\LA
\noindent Cum cuncta sacra elóquia Domínicis plena sint præcéptis quid est quod de dilectióne, quasi de singulári mandáto, Dóminus dicit, Hoc est præcéptum meum, ut diligátis ínvicem: nisi quia omne mandátum de sola dilectióne est: et omnia unum præcéptum sunt? Quia quidquid præcípitur, in sola caritáte solidátur. Ut enim multi árboris rami ex una radíce pródeunt: sic multæ virtútes ex una caritáte generántur. Nec habet áliquid viriditátis ramus boni óperis, si non manet in radíce caritátis.\par
\switchcolumn
\EN
\noindent All the holy words of the Lord are full of His commandments. Why, then, speaketh the Lord of the commandment to love one another as if He gave no other commandment? This, saith He, is My commandment, That ye love one another. Is it not because love is the one object of all His commandments, and all His commandments are one? For, even as a tree, having but one root, bringeth forth many branches, so, if the root be love, many virtues do spring therefrom. Neither is the branch of good works green, if it abide not in the root of love.\par
\switchcolumn*

\LA
\TeDeum{Te Deum}
\switchcolumn
\EN
\TeDeum{Te Deum}
\switchcolumn*

\end{paracol}

\par\needspace{10\baselineskip}\addvspace{1.2em}{\centering\small\textsc{Die 24 Augusti} · \textsc{24 August}\par}
\Office{S. Bartholomæi Apostoli}{St. Bartholomew the Apostle}{Duplex II classis}
\addcontentsline{toc}{subsection}{Die 24 Augusti · S. Bartholomæi Apostoli \textit{— St. Bartholomew the Apostle}}
\begin{paracol}{2}
\LA
\RefLine{Lectiones I–III: ex Commúni Apostolorum (C1).}
\switchcolumn
\EN
\RefLine{Lessons I–III: from the Commune Apostolorum (C1).}
\switchcolumn*

\LA
\LessonHead{Lectio IV}
\switchcolumn
\EN
\LessonHead{Lesson IV}
\switchcolumn*

\LA
\noindent Bartholomǽus Apóstolus, Galilǽus, cum in Indiam citeriórem, quæ ei in orbis terrárum sortitióne ad prædicándum Jesu Christi Evangélium obvénerat, progréssus esset, advéntum Dómini Jesu juxta sancti Matthǽi Evangélium illis géntibus prædicávit. Sed, cum in ea província plúrimos ad Jesum Christum convertísset, multos labóres calamitatésque perpéssus, venit in majórem Arméniam.\par
\switchcolumn
\EN
\noindent The Apostle Bartholomew was a Galilean. In the division of the world among the Apostles it fell to his lot to preach the Gospel of Jesus Christ in hither India. He went thither and preached to those nations the coming of the Lord Jesus, according to the Gospel of St. Matthew. When he had turned many in that province to Jesus Christ, and had endured many toils and woes, he came into the Greater Armenia.\par
\switchcolumn*

\LA
\begin{Resp}\Rsign Vidi conjúnctos viros, habéntes spléndidas vestes, et Ángelus Dómini locútus est ad me, dicens: \Star{} Isti sunt viri sancti facti amíci Dei.\end{Resp}
\begin{Resp}\Vsign Vidi Ángelum Dei fortem, volántem per médium cælum, voce magna clamántem et dicéntem.\end{Resp}
\begin{Resp}\Rsign Isti sunt viri sancti facti amíci Dei.\end{Resp}
\switchcolumn
\EN
\begin{Resp}\Rsign I saw men standing together, clad in shining raiment, and the Angel of the Lord spoke unto me, saying: \Star{} These men are holy, for they are the friends of God.\end{Resp}
\begin{Resp}\Vsign I saw a strong Angel of God fly into the midst of heaven, saying with a loud voice:\end{Resp}
\begin{Resp}\Rsign These men are holy, for they are the friends of God.\end{Resp}
\switchcolumn*

\LA
\LessonHead{Lectio V}
\switchcolumn
\EN
\LessonHead{Lesson V}
\switchcolumn*

\LA
\noindent Ibi Polýmium regem et cónjugem ejus ac prætérea duódecim civitátes ad christiánam fidem perdúxit. Quæ res in eum magnam invídiam concitávit illíus gentis sacerdótum. Nam usque ádeo Astýagem, Polýmii regis fratrem, in Apóstolum incendérunt, ut is vivo Bartholomǽo pellem crudéliter détrahi jússerit ac caput abscíndi; quo in martýrio ánimam Deo réddidit.\par
\switchcolumn
\EN
\noindent There he brought to the Christian faith Polymius the King, and his wife, and likewise the inhabitants of twelve cities. This stirred up a great hatred against him among the priests of that nation. They so inflamed against the Apostle Astyages the brother of King Polymius, that he savagely ordered Bartholomew to be flayed alive and beheaded; under the which martyrdom he gave up his soul to God.\par
\switchcolumn*

\LA
\begin{Resp}\Rsign Beáti estis, cum maledíxerint vobis hómines, et persecúti vos fúerint, et díxerint omne malum advérsum vos, mentiéntes, propter me: \Star{} Gaudéte et exsultáte, quóniam merces vestra copiósa est in cælis.\end{Resp}
\begin{Resp}\Vsign Cum vos óderint hómines, et cum separáverint vos, et exprobráverint, et ejécerint nomen vestrum tamquam malum propter Fílium hóminis.\end{Resp}
\begin{Resp}\Rsign Gaudéte et exsultáte, quóniam merces vestra copiósa est in cælis.\end{Resp}
\switchcolumn
\EN
\begin{Resp}\Rsign Blessed are ye when men shall revile you, and persecute you, and shall say all manner of evil against you falsely, for My sake; \Star{} Rejoice, and be exceeding glad, for great is your reward in heaven.\end{Resp}
\begin{Resp}\Vsign When men shall hate you, and when they shall separate you from their company, and shall reproach you, and cast out your name as evil, for the Son of Man's sake.\end{Resp}
\begin{Resp}\Rsign Rejoice, and be exceeding glad, for great is your reward in heaven.\end{Resp}
\switchcolumn*

\LA
\LessonHead{Lectio VI}
\switchcolumn
\EN
\LessonHead{Lesson VI}
\switchcolumn*

\LA
\noindent Ejus corpus Albáni, quæ est urbs majóris Arméniæ, ubi is passus fúerat, sepúltus est. Quod póstea ad Líparam ínsulam delátum, inde Benevéntum translátum est. Postrémo Romam ab Ottóne tértio imperatóre portátum, in Tíberis ínsula, in ecclésia ejus nómine Deo dicáta, collocátum fuit.\par
\switchcolumn
\EN
\noindent His body was buried at the town of Albanopolis in the Greater Armenia, where he had suffered. It was afterwards taken to the Island of Lipari, and thence carried to Benevento. Lastly, the Emperor Otto III brought it to Rome, where it was laid in the Church dedicated to God in his name on the Island in the Tiber.\par
\switchcolumn*

\LA
\begin{Resp}\Rsign Isti sunt triumphatóres et amíci Dei, qui contemnéntes jussa príncipum, meruérunt prǽmia ætérna: \Star{} Modo coronántur, et accípiunt palmam.\end{Resp}
\begin{Resp}\Vsign Isti sunt qui venérunt ex magna tribulatióne, et lavérunt stolas suas in sánguine Agni.\end{Resp}
\begin{Resp}\Rsign Modo coronántur, et accípiunt palmam.\end{Resp}
\begin{Resp}\Vsign Glória Patri, et Fílio, \Star{} et Spirítui Sancto.\end{Resp}
\begin{Resp}\Rsign Modo coronántur, et accípiunt palmam.\end{Resp}
\switchcolumn
\EN
\begin{Resp}\Rsign These are they which have conquered, and are become the friends of God, who recked not of the commandments of princes, and earned the everlasting reward. \Star{} And now have they crowns on their heads, and palms in their hands.\end{Resp}
\begin{Resp}\Vsign These are they which came out of great tribulation, and have washed their robes in the blood of the Lamb.\end{Resp}
\begin{Resp}\Rsign And now have they crowns on their heads, and palms in their hands.\end{Resp}
\begin{Resp}\Vsign Glory be to the Father, and to the Son, \Star{} and to the Holy Ghost.\end{Resp}
\begin{Resp}\Rsign And now have they crowns on their heads, and palms in their hands.\end{Resp}
\switchcolumn*

\LA
\LessonHead{Lectio VII}
\switchcolumn
\EN
\LessonHead{Lesson VII}
\switchcolumn*

\LA
\LessonTitle{Léctio sancti Evangélii secúndum Lucam}
\switchcolumn
\EN
\LessonTitle{From the Holy Gospel according to Luke}
\switchcolumn*

\LA
\Cite{Luc 6:12-19}
\switchcolumn
\EN
\Cite{Luke 6:12-19}
\switchcolumn*

\LA
\noindent In illo témpore: Exiit Jesus in montem oráre et erat pernóctans in oratióne Dei. Et, cum dies factus esset, vocávit discípulos suos. Et réliqua.\par
\switchcolumn
\EN
\noindent At that time Jesus went out into a mountain to pray, and continued all night in prayer to God. And when it was day, He called unto Him His disciples. And so on.\par
\switchcolumn*

\LA
\LessonTitle{Homilía sancti Ambrósii Epíscopi}
\switchcolumn
\EN
\LessonTitle{Homily by St. Ambrose, Bishop (of Milan.)}
\switchcolumn*

\LA
\Source{Lib. 5 Comment. in Luc. cap. 6, post initium}
\switchcolumn
\EN
\Source{Bk. v. Comment. on Luke 6}
\switchcolumn*

\LA
\noindent Omnes magni, omnes sublímes montem ascéndunt. Non enim cuicúmque prophéta dicit: Ascénde in montem excélsum, qui evangelízas Sion: exálta in virtúte vocem tuam, qui evangelízas Jerúsalem. Non vestígiis corporálibus, sed factis sublimióribus in hunc montem ascénde, et séquere Christum, ut ipse esse mons possis: Montes enim in circúitu ejus. Quare in Evangélio invénies solos cum Dómino montem ascendísse discípulos. Orat ergo Dóminus, non ut pro se óbsecret, sed ut pro me ímpetret. Nam, etsi ómnia posúerit Pater in potestáte Fílii; Fílius tamen, ut hóminis formam impléret, obsecrándum Patrem putat esse pro nobis, quia advocátus est noster.\par
\switchcolumn
\EN
\noindent All they who go up into the mountain are the great and the aspiring. It is not to every man that the Prophet saith: O thou that tellest good tidings to Zion, get thee up into the high mountain, Thou that tellest good tidings to Jerusalem, lift up thy voice with strength. (Isa. xl. 9.) Not with bodily feet, but by high deeds get thee up into this mountain, and follow Christ, that thou mayest be a mountain thyself. Therefore it is that thou findest in the Gospel that none but His disciples went up into the mountain with the Lord. The Lord therefore prayeth, not to entreat anything for Himself, but to obtain somewhat for me. For, albeit the Father had given the Son power over all flesh, that He might give eternal life to as many as He had given Him, (John xvii. 2,) the Son Himself, being found in fashion as a man, (Phil. ii. 8,) thinketh well to pray the Father (John xiv. 16) on our behalf, inasmuch as He is our Advocate with the Father, (1 John ii. 1.)\par
\switchcolumn*

\LA
\begin{Resp}\Rsign Isti sunt qui vivéntes in carne, plantavérunt Ecclésiam sánguine suo: \Star{} Cálicem Dómini bibérunt, et amíci Dei facti sunt.\end{Resp}
\begin{Resp}\Vsign In omnem terram exívit sonus eórum, et in fines orbis terræ verba eórum.\end{Resp}
\begin{Resp}\Rsign Cálicem Dómini bibérunt, et amíci Dei facti sunt.\end{Resp}
\switchcolumn
\EN
\begin{Resp}\Rsign These are they who while yet they lived in the flesh, planted the Church in their own blood; \Star{} They drank of the Lord's cup, and became the friends of God.\end{Resp}
\begin{Resp}\Vsign Their sound is gone out through all the earth, and their words to the ends of the world.\end{Resp}
\begin{Resp}\Rsign They drank of the Lord's cup, and became the friends of God.\end{Resp}
\switchcolumn*

\LA
\LessonHead{Lectio VIII}
\switchcolumn
\EN
\LessonHead{Lesson VIII}
\switchcolumn*

\LA
\noindent Et erat, inquit, pernóctans in oratióne Dei. Spécies tibi, Christiáne, datur, forma præscríbitur, quam débeas æmulári. Quid enim te pro salúte tua fácere opórtet, quando pro te Christus in oratióne pernóctat? Quid te fácere cónvenit, cum vis áliquod pietátis offícium adoríri, quando Christus, missúrus Apóstolos, orávit prius, et solus orávit? Nec usquam álibi, si non fallor, orásse cum Apóstolis reperítur; ubíque solus óbsecrat. Dei enim consílium humána vota non cápiunt, nec quisquam interiórum potest esse párticeps Christi.\par
\switchcolumn
\EN
\noindent And continued all night in prayer to God. Herein, O Christian, a pattern is set before thee, an example is given thee, after the which thou oughtest to aspire. What doth it not behove thee to do for thy salvation, when Christ spent an whole night in prayer for the same What doth it become thee to do, when thou willest some good work, when Christ prayed before He sent forth His Apostles He prayed first, and He prayed alone. Neither, unless I am mistaken, do we anywhere find that He ever joined in prayer with His disciples. He ever prayed alone. Human desires cannot grasp the counsel of God, nor can any man, however spiritually minded, share the thoughts of God.\par
\switchcolumn*

\LA
\begin{Resp}\Rsign Isti sunt viri sancti, quos elégit Dóminus in caritáte non ficta, et dedit illis glóriam sempitérnam: \Star{} Quorum doctrína fulget Ecclésia, ut sole luna.\end{Resp}
\begin{Resp}\Vsign Sancti per fidem vicérunt regna: operáti sunt justítiam.\end{Resp}
\begin{Resp}\Rsign Quorum doctrína fulget Ecclésia, ut sole luna.\end{Resp}
\begin{Resp}\Vsign Glória Patri, et Fílio, \Star{} et Spirítui Sancto.\end{Resp}
\begin{Resp}\Rsign Quorum doctrína fulget Ecclésia, ut sole luna.\end{Resp}
\switchcolumn
\EN
\begin{Resp}\Rsign These men are saints, whom the Lord hath chosen in love unfeigned, and hath given them glory everlasting. These are they \Star{} By the light of whose teaching the Church is glorified, even as the moon is glorified by the light of the sun.\end{Resp}
\begin{Resp}\Vsign The saints through faith subdued kingdoms, wrought righteousness.\end{Resp}
\begin{Resp}\Rsign By the light of whose teaching the Church is glorified, even as the moon is glorified by the light of the sun.\end{Resp}
\begin{Resp}\Vsign Glory be to the Father, and to the Son, \Star{} and to the Holy Ghost.\end{Resp}
\begin{Resp}\Rsign By the light of whose teaching the Church is glorified, even as the moon is glorified by the light of the sun.\end{Resp}
\switchcolumn*

\LA
\LessonHead{Lectio IX}
\switchcolumn
\EN
\LessonHead{Lesson IX}
\switchcolumn*

\LA
\noindent Vocávit, inquit, discípulos suos, et elégit duódecim ex ipsis; quos ad propagándum auxílium salútis humánæ per terrárum orbem satóres fídei destináret. Simul advérte cæléste consílium: non sapiéntes áliquos, non dívites, non nóbiles, sed piscatóres et publicános, quos dirígeret, elégit; ne traduxísse prudéntia, ne redemísse divítiis, ne poténtiæ nobilitatísque auctoritáte traxísse áliquos ad suam grátiam viderétur; ut veritátis rátio, non disputatiónis grátia prævaléret.\par
\switchcolumn
\EN
\noindent The Evangelist continueth thus: And when it was day, He called unto Him His disciples and of them He chose twelve whom He sent forth to help the salvation of men by sowing the seed of the faith throughout the whole world. Consider here the counsel of heaven. He chose out for His mission men, not wise, nor rich, nor noble, but fishermen and publicans, lest He should seem to have converted any to His grace by skill, or bought them with money, or drawn them by the power and authority of greatness, and the simple force of the truth, not the charms of argument, might have the victory.\par
\switchcolumn*

\LA
\TeDeum{Te Deum}
\switchcolumn
\EN
\TeDeum{Te Deum}
\switchcolumn*

\end{paracol}

\par\needspace{10\baselineskip}\addvspace{1.2em}{\centering\small\textsc{Die 25 Augusti} · \textsc{25 August}\par}
\Office{S. Ludovici Regis Franciæ Confessoris}{St. Louis, Confessor}{Semiduplex}
\addcontentsline{toc}{subsection}{Die 25 Augusti · S. Ludovici Regis Franciæ Confessoris \textit{— St. Louis, Confessor}}
\begin{paracol}{2}
\LA
\RefLine{Lectiones I–III: de Scriptura occurrente.}
\switchcolumn
\EN
\RefLine{Lessons I–III: from the Scripture of the day.}
\switchcolumn*

\LA
\LessonHead{Lectio IV}
\switchcolumn
\EN
\LessonHead{Lesson IV}
\switchcolumn*

\LA
\noindent Ludovicus nonus Galliæ rex, duodecim annos natus, patre amisso, et in Blanchæ matris sanctissima disciplina educatus, cum jam vigesimum annum in regno ageret, in morbum incidit: quo tempore cogitavit de recuperanda possessione Jerosolymorum. Quamobrem ubi convaluisset, vexillum ab episcopo Parisiensi accepit: deínde mare cum ingenti exercitu trajiciens, primo prælio Saracenos fugavit. Sed cum ex pestilentia magna militum multitudo periisset, victus ipse captusque est.\par
\switchcolumn
\EN
\noindent Louis IX, King of France, (was born on the 25 th day of April, in the year of our Lord 1215.) At the age of twelve years he lost his father. He was brought up under the godly care of his mother, Blanche of Castile. In the twentieth year of his reign he fell grievously sick, and the thought then occurred to him of delivering Jerusalem out of the hands of the Moslems. On his health being restored, he received a banner from the Bishop of Paris, and crossed the sea (to Egypt) with a very great army. In his first battle he put the Saracens to flight, but, a great number of the soldiers perishing by disease, he was himself conquered and taken prisoner.\par
\switchcolumn*

\LA
\begin{Resp}\Rsign Honéstum fecit illum Dóminus, et custodívit eum ab inimícis, et a seductóribus tutávit illum: \Star{} Et dedit illi claritátem ætérnam.\end{Resp}
\begin{Resp}\Vsign Justum dedúxit Dóminus per vias rectas, et osténdit illi regnum Dei.\end{Resp}
\begin{Resp}\Rsign Et dedit illi claritátem ætérnam.\end{Resp}
\switchcolumn
\EN
\begin{Resp}\Rsign The Lord made him honourable, and defended him from his enemies, and kept him safe from those that lay in wait for him; \Star{} And gave him perpetual glory.\end{Resp}
\begin{Resp}\Vsign He went down with him into the pit, and left him not in bonds.\end{Resp}
\begin{Resp}\Rsign And gave him perpetual glory.\end{Resp}
\switchcolumn*

\LA
\LessonHead{Lectio V}
\switchcolumn
\EN
\LessonHead{Lesson V}
\switchcolumn*

\LA
\noindent Rebus postea cum Saracenis compositis, liber rex exercitusque dimittitur. Quinque annis in Oriente commoratus, plurimos Christianos a barbarorum servitute redemit, multos étiam infideles ad Christi fidem convertit: præterea aliquot Christianorum urbes refecit suis sumptibus. Interim mater ejus migrat e vita: quare domum redire cogitur, ubi totum se dedit pietatis officiis.\par
\switchcolumn
\EN
\noindent The King afterwards entered into treaty with the Saracens, and he and his army departed in peace. He remained five years in the East, during which he redeemed great numbers of Christians from slavery among the unbelievers, and also brought many of the unbelievers themselves to believe in Christ. Moreover he rebuilt several cities of the Christians at his own cost. Meanwhile, his mother departed this life, whereby he was constrained to return home, where he gave himself up entirely to works of godliness.\par
\switchcolumn*

\LA
\begin{Resp}\Rsign Amávit eum Dóminus, et ornávit eum: stolam glóriæ índuit eum, \Star{} Et ad portas paradísi coronávit eum.\end{Resp}
\begin{Resp}\Vsign Índuit eum Dóminus lorícam fídei, et ornávit eum.\end{Resp}
\begin{Resp}\Rsign Et ad portas paradísi coronávit eum.\end{Resp}
\switchcolumn
\EN
\begin{Resp}\Rsign The Lord loved him and beautified him; He clothed him with a robe of glory, \Star{} And crowned him at the gates of Paradise.\end{Resp}
\begin{Resp}\Vsign The Lord hath put on him the breast-plate of faith, and hath adorned him.\end{Resp}
\begin{Resp}\Rsign And crowned him at the gates of Paradise.\end{Resp}
\switchcolumn*

\LA
\LessonHead{Lectio VI}
\switchcolumn
\EN
\LessonHead{Lesson VI}
\switchcolumn*

\LA
\noindent Multa ædificavit monasteria, et pauperum hospitia: beneficentia egentes sublevabat: frequens visebat ægrotos, quibus ipse non solum suis sumptibus omnia suppeditabat, sed etiam, quæ opus erant, manibus ministrabat. Vestitu vulgari utebatur, cilicio ac jejunio corpus assidue affligebat. Sed cum iterum transmisisset, bellum Saracenis illaturus, jamque castra in eorum conspectu posuisset, pestilentia decessit in illa oratione: Introibo in domum tuam; adorabo ad templum sanctum tuum, et confitebor nomini tuo. Ejus corpus postea Lutétiam Parisiorum translatum est, quod in celebri sancti Dionysii templo asservatur et colitur; caput vero in sacra æde sanctæ Capellæ. Ipse clarus miraculis a Bonifacio Papa octavo in sanctorum numerum est relatus.\par
\switchcolumn
\EN
\noindent He built many monasteries, and charitable institutions for the poor. By his alms he relieved the needy, and often visited the sick, for whom he not only provided at his own cost, but waited on them with his own hands with such things as they wanted. He wore a plain dress and constantly chastised his body with hair-cloth and fasting. (In the year 1270) he crossed the sea (to Tunis) to make war again upon the Saracens. His camp was pitched in sight of the enemy, but he was seized with pestilence, and died uttering the words I will come into thy house I will worship toward thy holy temple, and praise thy Name. (Pss. v, 8; cxxxvii, 2.) His body was afterwards carried to Paris, and it is kept and honoured in the famous Abbey Church of St. Denis, but his head in the oratory called La Sainte Chapelle. He was renowned for miracles, and Pope Boniface VIII. enrolled his name among those of the Saints.\par
\switchcolumn*

\LA
\begin{Resp}\Rsign Iste homo perfécit ómnia quæ locútus est ei Deus, et dixit ad eum: Ingrédere in réquiem meam: \Star{} Quia te vidi justum coram me ex ómnibus géntibus.\end{Resp}
\begin{Resp}\Vsign Iste est, qui contémpsit vitam mundi, et pervénit ad cæléstia regna.\end{Resp}
\begin{Resp}\Rsign Quia te vidi justum coram me ex ómnibus géntibus.\end{Resp}
\begin{Resp}\Vsign Glória Patri, et Fílio, \Star{} et Spirítui Sancto.\end{Resp}
\begin{Resp}\Rsign Quia te vidi justum coram me ex ómnibus géntibus.\end{Resp}
\switchcolumn
\EN
\begin{Resp}\Rsign This is he which did according unto all that God commanded him; and God said unto him: Enter thou into My rest, \Star{} For thee have I seen righteous before Me among all people.\end{Resp}
\begin{Resp}\Vsign This is he which loved not his life in this world, and is come unto an everlasting kingdom.\end{Resp}
\begin{Resp}\Rsign For thee have I seen righteous before Me among all people.\end{Resp}
\begin{Resp}\Vsign Glory be to the Father, and to the Son, \Star{} and to the Holy Ghost.\end{Resp}
\begin{Resp}\Rsign For thee have I seen righteous before Me among all people.\end{Resp}
\switchcolumn*

\LA
\LessonHead{Lectio VII}
\switchcolumn
\EN
\LessonHead{Lesson VII}
\switchcolumn*

\LA
\LessonTitle{Léctio sancti Evangélii secúndum Lucam}
\switchcolumn
\EN
\LessonTitle{From the Holy Gospel according to Luke}
\switchcolumn*

\LA
\Cite{Luc 19:12-26}
\switchcolumn
\EN
\Cite{Luke 19:12-25}
\switchcolumn*

\LA
\noindent In illo témpore: Dixit Jesus discípulis suis parábolam hanc: Homo quidam nóbilis ábiit in regiónem longinquam, accípere sibi regnum, et revérti. Et réliqua.\par
\switchcolumn
\EN
\noindent At that time, Jesus spoke this parable unto His disciples: A certain nobleman went into a far country to receive for himself a kingdom and to return. And so on.\par
\switchcolumn*

\LA
\LessonTitle{Homilía sancti Ambrósii Epíscopi}
\switchcolumn
\EN
\LessonTitle{Homily by St. Ambrose, Bishop (of Milan.)}
\switchcolumn*

\LA
\Source{Lib. 8. in Lucam.}
\switchcolumn
\EN
\Source{Book viii. on Luke.}
\switchcolumn*

\LA
\noindent Bonus ordo, ut vocatúrus gentes, et Judǽos jussúrus intérfici, qui noluérunt regnáre supra se Christum, hanc præmítteret comparatiónem, ne dicerétur: Nihil déderat pópulo Judæórum, unde póterat mélior fíeri? ut quid ab eo qui nihil recépit, exígitur? Non medíocris ista est mna, quam supra múlier evangélica quia non invénit, lucérnam accéndit, lúmine quærit admóto, gratulátur invéntam.\par
\switchcolumn
\EN
\noindent It is well ordered that, being about to call the Gentiles, and to command the destruction of those Jews, who would not have Christ to reign over them, He should put forth first this parable; lest it should be said He had given the Jews no means of becoming better. How can they be asked to repay who have received nothing? That is not a piece of silver of little worth, which, when the woman before mentioned in this Gospel (xv. 8) hath lost, she lighteth a candle, and sweepeth the house, and searcheth diligently until she findeth it.\par
\switchcolumn*

\LA
\begin{Resp}\Rsign Iste est qui ante Deum magnas virtútes operátus est, et de omni corde suo laudávit Dóminum: \Star{} Ipse intercédat pro peccátis ómnium populórum.\end{Resp}
\begin{Resp}\Vsign Ecce homo sine queréla, verus Dei cultor, ábstinens se ab omni ópere malo, et pérmanens in innocéntia sua.\end{Resp}
\begin{Resp}\Rsign Ipse intercédat pro peccátis ómnium populórum.\end{Resp}
\switchcolumn
\EN
\begin{Resp}\Rsign This is he which wrought great wonders before God, and praised the Lord with all his heart. \Star{} May he pray for all people, that their sins may be forgiven unto them.\end{Resp}
\begin{Resp}\Vsign Behold a man without blame, a worshipper of God in truth, keeping himself clean from every evil work, and abiding still in his innocency.\end{Resp}
\begin{Resp}\Rsign May he pray for all people, that their sins may be forgiven unto them.\end{Resp}
\switchcolumn*

\LA
\LessonHead{Lectio VIII}
\switchcolumn
\EN
\LessonHead{Lesson VIII}
\switchcolumn*

\LA
\noindent Dénique ex una decem mnas álius fecit, álius quinque. Fortásse iste morália habet, quia quinque sunt córporis sensus: ille duplícia, id est, mýstica legis, et morália probitátis. Unde et Matthǽus quinque talénta, et duo talénta pósuit: in quinque taléntis ut sint morália, in duóbus utrúmque, mýsticum atque morále. Ita quod número inférius, re ubérius.\par
\switchcolumn
\EN
\noindent A single pound one gained ten and another five pounds. Perchance by him which had the five pounds is signified he which practiseth well, since the body hath five senses, and by him which had the ten, that is, double the other, he which is learned and orthodox in the deep things of doctrine, as well as upright in his practical life. Hence also in Matthew we have five talents and two talents the five talents signifying good practice, and the two talents precept and practice together. So that that which counteth as the greater number is but a fraction of the lesser number.\par
\switchcolumn*

\LA
\begin{Resp}\Rsign Sint lumbi vestri præcíncti, et lucérnæ ardéntes in mánibus vestris: \Star{} Et vos símiles homínibus exspectántibus dóminum suum, quando revertátur a núptiis.\end{Resp}
\begin{Resp}\Vsign Vigiláte ergo, quia nescítis qua hora Dóminus vester ventúrus sit.\end{Resp}
\begin{Resp}\Rsign Et vos símiles homínibus exspectántibus dóminum suum, quando revertátur a núptiis.\end{Resp}
\begin{Resp}\Vsign Glória Patri, et Fílio, \Star{} et Spirítui Sancto.\end{Resp}
\begin{Resp}\Rsign Et vos símiles homínibus exspectántibus dóminum suum, quando revertátur a núptiis.\end{Resp}
\switchcolumn
\EN
\begin{Resp}\Rsign Let your loins be girded about, and your lights burning; \Star{} And ye yourselves like unto men that wait for their lord, when he will return from the wedding.\end{Resp}
\begin{Resp}\Vsign Watch therefore, for ye know not what hour your Lord doth come.\end{Resp}
\begin{Resp}\Rsign And ye yourselves like unto men that wait for their lord, when he will return from the wedding.\end{Resp}
\begin{Resp}\Vsign Glory be to the Father, and to the Son, \Star{} and to the Holy Ghost.\end{Resp}
\begin{Resp}\Rsign And ye yourselves like unto men that wait for their lord, when he will return from the wedding.\end{Resp}
\switchcolumn*

\LA
\LessonHead{Lectio IX}
\switchcolumn
\EN
\LessonHead{Lesson IX}
\switchcolumn*

\LA
\noindent Et hic póssumus decem mnas decem verba intellígere, id est, legis doctrínam; quinque autem mnas, magistéria disciplínæ. Sed legisperítum in ómnibus volo esse perféctum; non enim in sermóne, sed in virtúte est regnum Dei. Bene autem, quia de Judǽis dicit, duo soli multiplicátam pecúniam déferunt, non útique æris, sed dispensatiónis usúris. Ália est enim pecúniæ fœnébris, ália doctrínæ cæléstis usúra.\par
\switchcolumn
\EN
\noindent And here we may also understand by the ten pounds the ten words, that is, the Commandments, and by the five pounds, the enforcement of their teaching. But I would that a lawyer should be in all things perfect. For the kingdom of God is not in word but in power. (1 Cor. iv. 20.) Meet also is it, that, in speaking of Jews, Christ should represent only two as bringing in increased capital, for these talents are talents not of money but of grace, and to increase money by usury is a very different thing from improving heavenly revelation by the like means.\par
\switchcolumn*

\LA
\TeDeum{Te Deum}
\switchcolumn
\EN
\TeDeum{Te Deum}
\switchcolumn*

\end{paracol}

\par\needspace{10\baselineskip}\addvspace{1.2em}{\centering\small\textsc{Die 26 Augusti} · \textsc{26 August}\par}
\Office{S. Zephyrini Papæ et Martyris}{St. Zephyrinus, Pope and Martyr}{Simplex}
\addcontentsline{toc}{subsection}{Die 26 Augusti · S. Zephyrini Papæ et Martyris \textit{— St. Zephyrinus, Pope and Martyr}}
\begin{paracol}{2}
\LA
\RefLine{Lectiones I–II: de Scriptura occurrente.}
\switchcolumn
\EN
\RefLine{Lessons I–II: from the Scripture of the day.}
\switchcolumn*

\LA
\LessonHead{Lectio III (pro commemoratione)}
\switchcolumn
\EN
\LessonHead{Lesson III (when commemorated)}
\switchcolumn*

\LA
\noindent Zephyrinus Romanus Severo imperatore ad regendam Ecclesiam assumptus, sancivit, ut qui ordinandi essent, opportuno tempore et multis præsentibus clericis et laicis, de more sacris initiarentur; doctique ac spectatæ vitæ homines ad id officii munus deligerentur. Decrevit præterea, ut rem divinam facienti episcopo sacerdotes omnes astarent. Idem instituit ut patriarcha, primas, metropolitanus adversus episcopum non ferant sententiam, nisi apostolica auctoritate fulti. Vixit in pontificatu annos decem et octo, dies decem et octo. Habuit ordinationes quatuor mense Decembri, quibus creavit presbyteros tredecim, diaconos septem, episcopos per diversa loca tredecim. Antonino imperatore martyrio coronatus est, et sepultus via Appia prope cemeterium Callisti, septimo Kalendas Septembris.\par
\switchcolumn
\EN
\noindent Pope Zephyrinus was a Roman, who was called to govern the Church, (in the year 202,) during the reign of the Emperor Severus. It was he who decreed that they who are to be ordained should be ordained only at a fit time, and in the presence of many clerks and laymen, as was indeed already the custom, and that none but learned men and well known and spoken of should be set apart to that office. He decreed also that when the Bishop celebrated the Holy Liturgy, all the Priests should be present around him. Also he decreed that no Patriarch, Primate, or Metropolitan should pronounce sentence on a Bishop, unless they were charged with the authority of the Apostolic See. He lived as Pope eighteen years. He held four December ordinations, wherein he made thirteen Priests, seven Deacons, and thirteen Bishops for diverse places. He received the crown of his testimony under the Emperor Antonine, and was buried on the Appian Way, near the cemetery of Callistus, upon the 26th day of August, (in the year 219.)\par
\switchcolumn*

\LA
\TeDeum{Te Deum}
\switchcolumn
\EN
\TeDeum{Te Deum}
\switchcolumn*

\end{paracol}

\par\needspace{10\baselineskip}\addvspace{1.2em}{\centering\small\textsc{Die 27 Augusti} · \textsc{27 August}\par}
\Office{S. Josephi Calasanctii Confessoris}{St. Joseph Calasanz, Confessor}{Duplex}
\addcontentsline{toc}{subsection}{Die 27 Augusti · S. Josephi Calasanctii Confessoris \textit{— St. Joseph Calasanz, Confessor}}
\begin{paracol}{2}
\LA
\RefLine{Lectiones I–III: de Scriptura occurrente.}
\switchcolumn
\EN
\RefLine{Lessons I–III: from the Scripture of the day.}
\switchcolumn*

\LA
\LessonHead{Lectio IV}
\switchcolumn
\EN
\LessonHead{Lesson IV}
\switchcolumn*

\LA
\noindent Josephus Calasanctius a Matre Dei, Petraltæ in Aragonia nobili genere natus, a teneris annis futuræ in pueros caritatis et eorum institutionis indicia præbuit. Nam adhuc parvulus eos ad se convocatos in mysteriis fidei et sacris precibus erudiebat. Humanis divinisque litteris egregie doctus cum studiis theologicis Valentiæ operam daret, nobilis potentisque feminæ illecebris fortiter superatis, virginitatem, quam Deo voverat, inoffensam insigni victoria servavit. Sacerdos ex voto factus, a compluribus episcopis in Castellæ Novæ, Aragoniæ et Catalauniæ regnis in partem laboris ascitus, exspectationem omnium vicit, pravis ubique moribus emendatis, ecclesiastica disciplina restituta, inimicitiis cruentisque factionibus mirifice exstinctis. At cælesti visione et Dei voce frequenter admonitus, Romam profectus est.\par
\switchcolumn
\EN
\noindent Joseph Calasanz, called of the Mother of God, was born of a noble family at Petralta in Aragon, (on the 15th day of September, in the year of Christ 1556.) From his tender years he began to show that fondness for children, and that gift of instructing them for which he was afterwards distinguished. He called them around him when he was still but a child himself, and taught them the mysteries of the faith and godly prayers. He was deeply learned in profane and sacred letters, and it was while he was studying theology at Valencia that he bravely overcame the wiles of a noble and powerful lady and, by a brilliant victory, kept untarnished that virginity which he had vowed to God. He became a Priest in consequence of a vow, and was summoned by many Bishops in the kingdoms of New Castile, Aragon, and Catalonia to help them in their work, wherein he surpassed the hopes of all, correcting depraved manners, restoring the discipline of the Church, and marvellously putting an end to hatreds and bloody feuds. But in obedience to a vision from heaven and many warnings from the voice of God, he left Spain and went to Rome.\par
\switchcolumn*

\LA
\begin{Resp}\Rsign Honéstum fecit illum Dóminus, et custodívit eum ab inimícis, et a seductóribus tutávit illum: \Star{} Et dedit illi claritátem ætérnam.\end{Resp}
\begin{Resp}\Vsign Justum dedúxit Dóminus per vias rectas, et osténdit illi regnum Dei.\end{Resp}
\begin{Resp}\Rsign Et dedit illi claritátem ætérnam.\end{Resp}
\switchcolumn
\EN
\begin{Resp}\Rsign The Lord made him honourable, and defended him from his enemies, and kept him safe from those that lay in wait for him; \Star{} And gave him perpetual glory.\end{Resp}
\begin{Resp}\Vsign He went down with him into the pit, and left him not in bonds.\end{Resp}
\begin{Resp}\Rsign And gave him perpetual glory.\end{Resp}
\switchcolumn*

\LA
\LessonHead{Lectio V}
\switchcolumn
\EN
\LessonHead{Lesson V}
\switchcolumn*

\LA
\noindent In Urbe summa vitæ asperitate, vigiliis et jejuniis corpus affligens, in orationibus et cælestium rerum contemplatione dies noctesque versabatur, septem ejusdem Urbis ecclesias singulis fere noctibus obire solitus: quem inde morem complures annos servavit. Dato piis sodalitatibus nomine, mirum quanto ardore pauperes, infirmos potissimum, aut carceribus detentos eleemosynis omnique pietatis officio sublevaret. Lue Urbem depopulante, una cum sancto Camillo, tanto fuit actus impetu caritatis, ut præter subsidia ægrotis pauperibus large collata, ipsa étiam defunctorum cadavera suis humeris tumulanda transferret. Verum cum divinitus accepisset, se ad informandos intelligentiæ ac pietatis spiritu adolescentulos præcipue pauperes, destinari, ordinem Clericorum regularium pauperum Matris Dei scholarum piarum fundavit, qui peculiarem curam circa puerorum eruditionem ex proprio instituto profiterentur: ipsumque ordinem a Clemente octavo, a Paulo quinto aliisque summis Pontificibus magnopere probatum, brevi tempore per plurimas Europæ provincias et regna mirabiliter propagavit. In hoc autem tot labores perpessus est, ac tot ærumnas invicto animo toleravit, ut omnium voce miraculum fortitudinis, et sancti Jobi exemplum diceretur.\par
\switchcolumn
\EN
\noindent In Rome he afflicted his body with extraordinary hardness of living, with watching, and with fasting, and so passed his days and nights in prayer, and in the contemplation of heavenly things. He was used to visit the Seven Churches almost every night, a custom which he kept for many years. Having joined several godly Brotherhoods, it was strange how eagerly he relieved the poor by alms and every sort of kindness, choosing especially the sick and the imprisoned. When the city was ravaged by a pestilence, such was the charitable zeal with which he joined in the labours of St. Camillus de' Lelli, that besides the great help which he brought to the sick poor, he would even carry the bodies of the dead on his own shoulders to burial. Having understood from God that his call was to bring up children in godliness and good learning, he founded the Order of the Poor Regular Clerks of the Pious Schools of the Mother of God, who profess as the special object of their Institute a singular care for the teaching of the poor. This Institute received the warm approval of Clement VIII., Paul V., and other Popes, and in a short time obtained a marvellous extension through many provinces and kingdoms of Europe. In this work Joseph Casalanz underwent so many toils, and patiently bore so many griefs, that he was proclaimed by all men a wonder of endurance and a very image of holy Job.\par
\switchcolumn*

\LA
\begin{Resp}\Rsign Amávit eum Dóminus, et ornávit eum: stolam glóriæ índuit eum, \Star{} Et ad portas paradísi coronávit eum.\end{Resp}
\begin{Resp}\Vsign Índuit eum Dóminus lorícam fídei, et ornávit eum.\end{Resp}
\begin{Resp}\Rsign Et ad portas paradísi coronávit eum.\end{Resp}
\switchcolumn
\EN
\begin{Resp}\Rsign The Lord loved him and beautified him; He clothed him with a robe of glory, \Star{} And crowned him at the gates of Paradise.\end{Resp}
\begin{Resp}\Vsign The Lord hath put on him the breast-plate of faith, and hath adorned him.\end{Resp}
\begin{Resp}\Rsign And crowned him at the gates of Paradise.\end{Resp}
\switchcolumn*

\LA
\LessonHead{Lectio VI}
\switchcolumn
\EN
\LessonHead{Lesson VI}
\switchcolumn*

\LA
\noindent Quamvis ordini universo præesset, totisque viribus ad animarum salutem incumberet, numquam tamen intermisit pueros, præsertim pauperiores, erudire, quorum scholas verrere, eosque domum comitari consuevit. In eo summæ patientiæ et humilitatis munere, valetudine étiam infirma, duos et quinquaginta annos perseveravit: dignus propterea, quem crebris Deus miraculis coram discipulis illustraret, et cui beatissima Virgo cum puero Jesu, illis orantibus benedicente, appareret. Amplissimis interim dignitatibus repudiatis prophetia, abdita cordium et absentia cognoscendi donis et miraculis clarus, Deiparæ Virginis, quam singulari pietate et ipse ab infantia coluit, et suis maxime commendavit, aliorumque cælitum frequenti apparitione dignatus, cum obitus sui diem, et ordinis tunc prope eversi restitutionem atque incrementum prænuntiasset, secundum et nonagesimum annum agens, Romæ obdormivit in Domino, octavo Kalendas Septembris, anno millesimo sexcentesimo quadragesimo octavo. Ejus cor et lingua post sæculum integra et incorrupta reperta sunt. Ipse vero multis post obitum quoque signis a Deo illustratus, primum a Benedicto decimo quarto beatorum cultu decoratus fuit, ac deínde a Clemente decimo tertio inter sanctos solemniter est relatus.\par
\switchcolumn
\EN
\noindent Then when he was at the head of his whole Order, and toiling with all his might for the salvation of souls, he never ceased to teach children, especially the poor, to sweep out the school-rooms, and to accompany the scholars home. Thus in spite of broken health he worked on for two and fifty years, with the greatest longsuffering and lowliness. He won that God should glorify him by many miracles worked in the presence of his disciples, and that the Most Blessed Virgin should appear to him, with the Child Jesus in her arms, blessing them as they prayed. He refused wealthy preferments when they were offered to him. He was eminent for the gift of prophecy, for the power of reading the secrets of the heart, of knowing distant events, and of miracles. The Virgin Mother of God, to whom from his childhood he had had an especial love, and other heavenly ones, honoured him by often allowing him to see them. He foretold the day of his own death, and the restoration and growth of his Order, which seemed at that time to be almost entirely destroyed. He fell asleep in the Lord at Rome, upon the 25 th day of August, in the year of salvation 1648, and of his own age the 92 nd. An hundred years after his death his heart and tongue were found whole and incorrupt. God glorified him by many miracles even after his death, and he was first crowned by Benedict XIV. with the honours paid to the Blessed, and then solemnly enrolled by Clement XIII. among the Saints.\par
\switchcolumn*

\LA
\begin{Resp}\Rsign Iste homo perfécit ómnia quæ locútus est ei Deus, et dixit ad eum: Ingrédere in réquiem meam: \Star{} Quia te vidi justum coram me ex ómnibus géntibus.\end{Resp}
\begin{Resp}\Vsign Iste est, qui contémpsit vitam mundi, et pervénit ad cæléstia regna.\end{Resp}
\begin{Resp}\Rsign Quia te vidi justum coram me ex ómnibus géntibus.\end{Resp}
\begin{Resp}\Vsign Glória Patri, et Fílio, \Star{} et Spirítui Sancto.\end{Resp}
\begin{Resp}\Rsign Quia te vidi justum coram me ex ómnibus géntibus.\end{Resp}
\switchcolumn
\EN
\begin{Resp}\Rsign This is he which did according unto all that God commanded him; and God said unto him: Enter thou into My rest, \Star{} For thee have I seen righteous before Me among all people.\end{Resp}
\begin{Resp}\Vsign This is he which loved not his life in this world, and is come unto an everlasting kingdom.\end{Resp}
\begin{Resp}\Rsign For thee have I seen righteous before Me among all people.\end{Resp}
\begin{Resp}\Vsign Glory be to the Father, and to the Son, \Star{} and to the Holy Ghost.\end{Resp}
\begin{Resp}\Rsign For thee have I seen righteous before Me among all people.\end{Resp}
\switchcolumn*

\LA
\LessonHead{Lectio VII}
\switchcolumn
\EN
\LessonHead{Lesson VII}
\switchcolumn*

\LA
\LessonTitle{Léctio sancti Evangélii secúndum Matthǽum.}
\switchcolumn
\EN
\LessonTitle{From the Holy Gospel according to Matthew}
\switchcolumn*

\LA
\Cite{Matt 18:1-5}
\switchcolumn
\EN
\Cite{Matt 18:1-5}
\switchcolumn*

\LA
\noindent In illo témpore: Accesserunt discipuli ad Jesum, dicentes: Quis, putas, major est in regno cælorum? Et réliqua.\par
\switchcolumn
\EN
\noindent At that time Came the disciples unto Jesus, saying Who is the greatest in the kingdom of heaven And so on.\par
\switchcolumn*

\LA
\LessonTitle{Homilia sancti Joannis Chrysostomi.}
\switchcolumn
\EN
\LessonTitle{Homily by St. John Chrysostom, Patriarch (of Constantinople.)}
\switchcolumn*

\LA
\Source{In cap. 18. Matth. Hom. 60.}
\switchcolumn
\EN
\Source{60th on Matth. xviii.}
\switchcolumn*

\LA
\noindent Videte ne aliquem istorum contempseritis parvulorum, quia eorum Angeli Patris mei faciem semper aspiciunt, et quia ego propter eos veni, et hæc Patris mei voluntas est. Ad tuendos conservandosque pusilios diligentiores nos reddit. Perspicis quam ingentia in tutelam tenujum mœnia erexerit, et quantum studium curamque habeat, ne perdantur tum quia supremas despicientibus eos pœnas statuit, tum quia summam pollicetur mercedem his, qui curam eorum suscipiunt; idque tam suo, quam Patris exemplo corroborat.\par
\switchcolumn
\EN
\noindent Take heed, saith Jesus, that ye despise not one of these little ones, for I say unto you, that in heaven their angels do always behold the face of My Father and that for their sake am I come, and this is the will of My Father. Hereby the Lord stirreth us up to guard and save these little ones. Thou seest how mighty are the walls which He raiseth to protect little children, and how great thought and care He hath lest they should be lost, threatening on the one hand the uttermost punishment against whosoever shall offend one of these little ones which believe in Him, (6) and promising on the other hand, the highest reward to whosoever shall receive one such little child in His Name, (5) and this His teaching He giveth both in His Own, and in His Father's Name.\par
\switchcolumn*

\LA
\begin{Resp}\Rsign Iste est qui ante Deum magnas virtútes operátus est, et de omni corde suo laudávit Dóminum: \Star{} Ipse intercédat pro peccátis ómnium populórum.\end{Resp}
\begin{Resp}\Vsign Ecce homo sine queréla, verus Dei cultor, ábstinens se ab omni ópere malo, et pérmanens in innocéntia sua.\end{Resp}
\begin{Resp}\Rsign Ipse intercédat pro peccátis ómnium populórum.\end{Resp}
\switchcolumn
\EN
\begin{Resp}\Rsign This is he which wrought great wonders before God, and praised the Lord with all his heart. \Star{} May he pray for all people, that their sins may be forgiven unto them.\end{Resp}
\begin{Resp}\Vsign Behold a man without blame, a worshipper of God in truth, keeping himself clean from every evil work, and abiding still in his innocency.\end{Resp}
\begin{Resp}\Rsign May he pray for all people, that their sins may be forgiven unto them.\end{Resp}
\switchcolumn*

\LA
\LessonHead{Lectio VIII}
\switchcolumn
\EN
\LessonHead{Lesson VIII}
\switchcolumn*

\LA
\noindent Dominum igitur étiam nos imitemur, et nihil pro fratribus omittamus, étiam eorum quæ humilia viliaque nimium videntur; sed si administratione nostra étiam opus fuerit, quamvis tenuis atque abjectus quidem, cui administrandum sit, fuerit, quamvis ardua nobis res atque laboris plena esse videatur, omnia hæc pro fratris salute tolerabiliora facilioraque, oro, videantur: tanto enim studio tantaque cura Deus dignam esse animam ostendit, ut neque Filio suo pepercerit.\par
\switchcolumn
\EN
\noindent Let us therefore take example by the Lord, and let us leave nothing undone for the good of any of our brethren, even for such as seem to us the least and lowliest, but if there be any need that we should serve any, low and outcast though he be, let us serve him though the thing look hard to us and calling for a great deal of work, let such things, I pray, be looked on as light and easy if they be required for our neighbour's salvation, for of such price and such care did God count his soul to be worth, that He spared not to purchase it, even His Own Son. (Rom. viii. 32.)\par
\switchcolumn*

\LA
\begin{Resp}\Rsign Sint lumbi vestri præcíncti, et lucérnæ ardéntes in mánibus vestris: \Star{} Et vos símiles homínibus exspectántibus dóminum suum, quando revertátur a núptiis.\end{Resp}
\begin{Resp}\Vsign Vigiláte ergo, quia nescítis qua hora Dóminus vester ventúrus sit.\end{Resp}
\begin{Resp}\Rsign Et vos símiles homínibus exspectántibus dóminum suum, quando revertátur a núptiis.\end{Resp}
\begin{Resp}\Vsign Glória Patri, et Fílio, \Star{} et Spirítui Sancto.\end{Resp}
\begin{Resp}\Rsign Et vos símiles homínibus exspectántibus dóminum suum, quando revertátur a núptiis.\end{Resp}
\switchcolumn
\EN
\begin{Resp}\Rsign Let your loins be girded about, and your lights burning; \Star{} And ye yourselves like unto men that wait for their lord, when he will return from the wedding.\end{Resp}
\begin{Resp}\Vsign Watch therefore, for ye know not what hour your Lord doth come.\end{Resp}
\begin{Resp}\Rsign And ye yourselves like unto men that wait for their lord, when he will return from the wedding.\end{Resp}
\begin{Resp}\Vsign Glory be to the Father, and to the Son, \Star{} and to the Holy Ghost.\end{Resp}
\begin{Resp}\Rsign And ye yourselves like unto men that wait for their lord, when he will return from the wedding.\end{Resp}
\switchcolumn*

\LA
\LessonHead{Lectio IX}
\switchcolumn
\EN
\LessonHead{Lesson IX}
\switchcolumn*

\LA
\noindent Si non est nobis satis ad salutem, quod virtuose ipsi vivamus, sed oportet aliorum salutem re ipsa desiderare; cum neque nos recte vivamus, neque alios hortemur, quid respondebimus? Quæ nobis spes salutis reliqua erit? Quid majus, quam animis moderari, quam adolescentulorum fingere mores? Omni certe pictore, omni certe statuario, ceterisque hujusmodi omnibus excellentiorem hunc duco, qui juvenum animos fingere non ignoret.\par
\switchcolumn
\EN
\noindent If it be not enough for our salvation that we should ourselves live well, but we must also seek the salvation of others, what shall we answer, if we neither live well ourselves, nor exhort others What hope that we shall be saved is then left to us What more important task is there than to train up minds, and teach to the young how to live He that is skilled to mould well the minds of children I reckon a nobler workman than any painter or sculptor, or such like artist.\par
\switchcolumn*

\LA
\TeDeum{Te Deum}
\switchcolumn
\EN
\TeDeum{Te Deum}
\switchcolumn*

\end{paracol}

\par\needspace{10\baselineskip}\addvspace{1.2em}{\centering\small\textsc{Die 28 Augusti} · \textsc{28 August}\par}
\Office{S. Augustini Episcopi et Confessoris et Ecclesiæ Doctoris}{St. Augustine of Hippo, Bishop, Confessor, and Doctor of the Church}{Duplex}
\addcontentsline{toc}{subsection}{Die 28 Augusti · S. Augustini Episcopi et Confessoris et Ecclesiæ Doctoris \textit{— St. Augustine of Hippo, Bishop, Confessor, and Doctor of the Church}}
\begin{paracol}{2}
\LA
\RefLine{Lectiones I–III: de Scriptura occurrente.}
\switchcolumn
\EN
\RefLine{Lessons I–III: from the Scripture of the day.}
\switchcolumn*

\LA
\RefLine{Lectiones VII–IX: ex Commúni Doctoris Pontificis (C4a).}
\switchcolumn
\EN
\RefLine{Lessons VII–IX: from the Common of a Doctor Bishop (C4a).}
\switchcolumn*

\LA
\LessonHead{Lectio IV}
\switchcolumn
\EN
\LessonHead{Lesson IV}
\switchcolumn*

\LA
\noindent Augustinus, Tagaste in Africa honestis parentibus natus, ac puer docilitate ingenii æquales longe superans, brevi omnibus doctrina antecelluit. Adolescens, dum esset Carthagine, in Manichæorum hæresim incidit. Postea Romam profectus, inde Mediolanum missus ut rhetoricam doceret, cum ibi frequens Ambrósii episcopi esset auditor, ejus opera incensus studio catholicæ fidei, annos natus triginta tres ab ipso baptizatur. Reversus in Africam, cum religione vitæ sanctimoniam conjungens, a Valerio notæ sanctitatis episcopo Hipponensi presbyter factus est. Quo tempore familiam instituit religiosorum, quibuscum victu communi eodemque cultu utens, eos ad apostolicæ vitæ doctrinæque disciplinam diligentissime erudiebat. Sed cum vigeret Manichæorum hæresis, vehementius in illam invehi cœpit, Fortunatumque hæresiarcham confutavit.\par
\switchcolumn
\EN
\noindent Augustine was born of honourable parents at Tagaste in Africa, (upon the 13th day of November, in the year of our Lord 354.) As a boy his great intellectual sharpness caused him to distance all his companions in learning. When he was living at Carthage as a young man, he fell into the heresy of the Manichaeans. He afterwards went to Rome, and was thence sent to Milan to teach Rhetoric. At Milan he often went to hear the sermons of Bishop Ambrose, by whose labours he was drawn to the Catholic Church, and by whom he was baptized (on Holy Saturday, 387,) at the age of thirty-three. After his return to Africa, (in 388,) Valerius, the illustrious and saintly Bishop of Hippo, finding him to unite holiness of life, with Catholic profession, made him a Priest, (about the end of 390.) At this time he founded a sort of family of godly men, who lived and worshipped in common with him, and whom he earnestly formed upon the model of the Apostolic life and teaching. The Manichaean heresy flaming forth with violence, he began strongly to attack it, and confounded the arch-heretic Fortunatus.\par
\switchcolumn*

\LA
\begin{Resp}\Rsign Invéni David servum meum, óleo sancto meo unxi eum: \Star{} Manus enim mea auxiliábitur ei.\end{Resp}
\begin{Resp}\Vsign Nihil profíciet inimícus in eo, et fílius iniquitátis non nocébit ei.\end{Resp}
\begin{Resp}\Rsign Manus enim mea auxiliábitur ei.\end{Resp}
\switchcolumn
\EN
\begin{Resp}\Rsign I have found David My servant, with My holy oil have I anointed him \Star{} For My hand shall help him.\end{Resp}
\begin{Resp}\Vsign The enemy shall prevail nothing against him, nor the son of wickedness afflict him.\end{Resp}
\begin{Resp}\Rsign For My hand shall help him.\end{Resp}
\switchcolumn*

\LA
\LessonHead{Lectio V}
\switchcolumn
\EN
\LessonHead{Lesson V}
\switchcolumn*

\LA
\noindent Hac Augustíni pietate commotus Valerius, eum adjutorem adhibuit episcopalis officii. Nihil illo fuit humilius, nihil continentius. Lectus ac vestitus moderatus, vulgaris mensa, quam semper sacra vel lectione vel disputatione condiebat. Tanta benignitate fuit in pauperes, ut, cum non esset alia facultas, sacra vasa frangeret ad eorum inopiam sustentandam. Feminarum, et in eis sororis, et fratris filiæ, contubernium familiaritatemque vitavit: quippe qui diceret, etsi propinquæ mulieres suspectæ non essent, tamen quæ ad eas ventitarent, posse suspicionem efficere. Nullum finem fecit prædicandi Dei verbum, nisi gravi morbo oppressus. Hæreticos perpetuo insectatus et coram et scriptis, ac nullo loco passus consistere, Africam a Manichæorum, Donatistarum, Pelagianorum, aliorumque præterea hæreticorum errore magna ex parte liberavit.\par
\switchcolumn
\EN
\noindent Valerius, moved by the godly zeal of Augustine, (in December 395,) joined him with himself as an assistant in his duties of Bishop (and dying in the year following, was succeeded by him.) He was lowly and pure in the highest degree. His furniture and dress were plain, and his food of the commonest sort, which he always seasoned when at table by either reading some religious book, or arguing upon some religious subject. His tenderness to the poor was such that, failing all other resources, he broke up the hallowed vessels to relieve their wants. It was his rule not to dwell or be very close friends with any woman, a rule which he did not relax even in the case of his sister and niece, for he was accustomed to say, that although no scandal could arise in the case of such near kinswomen, yet it might arise concerning the women friends who sought their company. He never ceased to preach the Word of God, until he was disabled by heavy sickness. He was always an hard follower after heretics, and by his words and his writings never suffered them to rest anywhere. In great measure, he purged Africa of the Manicheans, Donatists, Pelagians, and other heretics.\par
\switchcolumn*

\LA
\begin{Resp}\Rsign Pósui adjutórium super poténtem, et exaltávi eléctum de plebe mea: \Star{} Manus enim mea auxiliábitur ei.\end{Resp}
\begin{Resp}\Vsign Invéni David servum meum, óleo sancto meo unxi eum.\end{Resp}
\begin{Resp}\Rsign Manus enim mea auxiliábitur ei.\end{Resp}
\switchcolumn
\EN
\begin{Resp}\Rsign I have laid help upon one that is mighty, and have exalted one chosen out of My people \Star{} For My hand shall help him.\end{Resp}
\begin{Resp}\Vsign I have found David My servant, with My holy oil have I anointed him.\end{Resp}
\begin{Resp}\Rsign For My hand shall help him.\end{Resp}
\switchcolumn*

\LA
\LessonHead{Lectio VI}
\switchcolumn
\EN
\LessonHead{Lesson VI}
\switchcolumn*

\LA
\noindent Tam multa pie, subtiliter et copiose scripsit, ut christianam doctrinam maxime illustrarit. Quem in primis secuti sunt, qui postea theologicam disciplinam via et ratione tradiderunt. Wandalis Africam bello vastantibus, et Hipponem tertium jam mensem obsidentibus, in febrim incidit. Itaque cum discessum e vita sibi instare intelligeret, Psalmos David qui ad pœnitentiam pertinent, in conspectu positos profusis lacrimis legebat. Solebat autem dicere neminem, etsi nullius sceleris sibi conscius esset, committere debere, ut sine pænitentia migraret e vita. Ergo sensibus integris, in oratione defixus, astantibus fratribus, quos ad caritatem, pietatem, virtutesque omnes erat adhortatus, migravit in cælum. Vixit annos septuaginta sex, in episcopatu ad triginta sex. Cujus corpus primum in Sardiniam delatum, deínde a Luitprando, Longobardorum rege, magno pretio redemptum, Ticinum translatum est, ibique honorifice conditum.\par
\switchcolumn
\EN
\noindent He wrote so much, and that with such godliness and understanding, that he is to be held among the very chiefest of them by whom the teachings of Christianity have been shown forth. He is one of the first of those whom later theologians have followed, in method, and in argument. He fell sick of a fever what time the Vandals were laying Africa waste, and when they were busy in the third month of besieging Hippo. When he understood that his departure from this present life was at hand, he caused the Psalms of David which most speak the language of repentance to be placed before him, and read them with tears, for he was wont to say that even if a man's conscience were to accuse him of no sin, he should not dare to leave this world except as a penitent. His senses remained vigorous to the last, and it was while rapt in prayer, in the presence of the brethren whom he had exhorted to love, godliness, and all goodness, that he departed for heaven, (upon the 28th of August, 430.) He lived 76 years, whereof he had been a Bishop nearly thirty six. His body was first carried to Sardinia, but Luitprand, King of the Lombards, afterwards bought it for a great price, and took it to Ticino, where it is honourably buried.\par
\switchcolumn*

\LA
\begin{Resp}\Rsign Iste est, qui ante Deum magnas virtútes operátus est, et omnis terra doctrína ejus repléta est: \Star{} Ipse intercédat pro peccátis ómnium populórum.\end{Resp}
\begin{Resp}\Vsign Iste est, qui contémpsit vitam mundi, et pervénit ad cæléstia regna.\end{Resp}
\begin{Resp}\Rsign Ipse intercédat pro peccátis ómnium populórum.\end{Resp}
\begin{Resp}\Vsign Glória Patri, et Fílio, \Star{} et Spirítui Sancto.\end{Resp}
\begin{Resp}\Rsign Ipse intercédat pro peccátis ómnium populórum.\end{Resp}
\switchcolumn
\EN
\begin{Resp}\Rsign This is he which wrought great wonders before God, and the whole earth is full of his teaching \Star{} May he pray for all people, that their sins may be forgiven unto them\end{Resp}
\begin{Resp}\Vsign This is he which loved not his life in this world, and hath attained unto the kingdom of heaven.\end{Resp}
\begin{Resp}\Rsign May he pray for all people, that their sins may be forgiven unto them\end{Resp}
\begin{Resp}\Vsign Glory be to the Father, and to the Son, \Star{} and to the Holy Ghost.\end{Resp}
\begin{Resp}\Rsign May he pray for all people, that their sins may be forgiven unto them\end{Resp}
\switchcolumn*

\end{paracol}

\par\needspace{10\baselineskip}\addvspace{1.2em}{\centering\small\textsc{Die 29 Augusti} · \textsc{29 August}\par}
\Office{In Decollatione S. Joannis Baptistæ}{Beheading of St. John the Baptist}{Duplex majus}
\addcontentsline{toc}{subsection}{Die 29 Augusti · In Decollatione S. Joannis Baptistæ \textit{— Beheading of St. John the Baptist}}
\begin{paracol}{2}
\LA
\RefLine{Lectio IX: de S. Sabinæ Martyris (08-29cc).}
\switchcolumn
\EN
\RefLine{Lesson IX: of St. Sabina, Martyr (08-29cc).}
\switchcolumn*

\LA
\LessonHead{Lectio I}
\switchcolumn
\EN
\LessonHead{Lesson I}
\switchcolumn*

\LA
\LessonTitle{Incipit liber Jeremíæ Prophétæ}
\switchcolumn
\EN
\LessonTitle{Lesson from the book of Jeremias}
\switchcolumn*

\LA
\Cite{Jer 1:1-5}
\switchcolumn
\EN
\Cite{Jer 1:1-5}
\switchcolumn*

\LA
\noindent Verba Jeremíæ fílii Helcíæ, de sacerdótibus, qui fuérunt in Anathoth, in terra Bénjamin. Quod factum est verbum Dómini ad eum in diébus Josíæ fílii Amon regis Juda, in tertiodécimo anno regni ejus. Et factum est in diébus Jóakim fílii Josíæ regis Juda, usque ad consummatiónem undécimi anni Sedecíæ fílii Josíæ regis Juda, usque ad transmigratiónem Jerúsalem, in mense quinto. Et factum est verbum Dómini ad me, dicens: Priúsquam te formárem in útero, novi te: et ántequam exíres de vulva, sanctificávi te, et prophétam in géntibus dedi te.\par
\switchcolumn
\EN
\noindent The words of Jeremias the son of Helcias, of the priests that were in Anathoth, in the land of Benjamin. The word of the Lord which came to him in the days of Josias the son of Amon king of Juda, in the thirteenth year of his reign. And which came to him in the days of Joakim the son of Josias king of Juda, unto the end of the eleventh year of Sedecias the son of Josias king of Juda, even unto the carrying away of Jerusalem captive, in the fifth month. And the word of the Lord came to me, saying: Before I formed thee in the bowels of thy mother, I knew thee: and before thou camest forth out of the womb, I sanctified thee, and made thee a prophet unto the nations.\par
\switchcolumn*

\LA
\begin{Resp}\Rsign Misit Heródes rex manus, ac ténuit Joánnem et vinxit eum in cárcere, quia metuébat eum propter Herodíadem, \Star{} Quam túlerat fratri suo Philíppo uxórem.\end{Resp}
\begin{Resp}\Vsign Arguébat Heródem Joánnes propter Herodíadem.\end{Resp}
\begin{Resp}\Rsign Quam túlerat fratri suo Philíppo uxórem.\end{Resp}
\switchcolumn
\EN
\begin{Resp}\Rsign Herod the King sent forth, and laid hold upon John, and bound him in prison, for he feared him, for Herodias' sake, \Star{} His brother Philip's wife, for he had married her.\end{Resp}
\begin{Resp}\Vsign For John had rebuked Herod, for Herodias' sake.\end{Resp}
\begin{Resp}\Rsign His brother Philip's wife, for he had married her.\end{Resp}
\switchcolumn*

\LA
\LessonHead{Lectio II}
\switchcolumn
\EN
\LessonHead{Lesson II}
\switchcolumn*

\LA
\Cite{Jer 1:6-10}
\switchcolumn
\EN
\Cite{Jer 1:6-10}
\switchcolumn*

\LA
\noindent Et dixi, A a a, Dómine Deus: ecce néscio loqui, quia puer ego sum. Et dixit Dóminus ad me: Noli dícere, Puer sum: quóniam ad ómnia, quæ mittam te, ibis: et univérsa, quæcúmque mandávero tibi, loquéris. Ne tímeas a fácie eórum: quia tecum ego sum, ut éruam te, dicit Dóminus. Et misit Dóminus manum suam, et tétigit os meum: et dixit Dóminus ad me: Ecce dedi verba mea in ore tuo: ecce constítui te hódie super gentes, et super regna, ut evéllas, et déstruas, et dispérdas, et díssipes, et ædífices, et plantes.\par
\switchcolumn
\EN
\noindent And I said: Ah, ah, ah, Lord God: behold, I cannot speak, for I am a child. And the Lord said to me: Say not: I am a child: for thou shalt go to all that I shall send thee: and whatsoever I shall command thee, thou shalt speak. Be not afraid at their presence: for I am with thee to deliver thee, saith the Lord. And the Lord put forth his hand, and touched my mouth: and the Lord said to me: Behold I have given my words in thy mouth: Lo, I have set thee this day over the nations, and over the kingdoms, to root up, and pull down, and to waste, and to destroy, and to build, and to plant.\par
\switchcolumn*

\LA
\begin{Resp}\Rsign Joánnes Baptísta arguébat Heródem \Star{} Propter Herodíadem, quam túlerat fratri suo vivénti uxórem.\end{Resp}
\begin{Resp}\Vsign Misso Heródes spiculatóre, præcépit amputári caput Joánnis in cárcere.\end{Resp}
\begin{Resp}\Rsign Propter Herodíadem, quam tulerat fratri suo vivénti uxórem.\end{Resp}
\switchcolumn
\EN
\begin{Resp}\Rsign John the Baptist had rebuked Herod, \Star{} For Herodias' sake, his brother's wife, whom he had married while his brother was yet alive.\end{Resp}
\begin{Resp}\Vsign Herod sent an executioner, and commanded to behead John in the prison.\end{Resp}
\begin{Resp}\Rsign For Herodias' sake, his brother's wife, whom he had married while his brother was yet alive.\end{Resp}
\switchcolumn*

\LA
\LessonHead{Lectio III}
\switchcolumn
\EN
\LessonHead{Lesson III}
\switchcolumn*

\LA
\Cite{Jer 1:17-19}
\switchcolumn
\EN
\Cite{Jer 1:17-19}
\switchcolumn*

\LA
\noindent Tu ergo accínge lumbos tuos, et surge, et lóquere ad eos ómnia quæ ego præcípio tibi. Ne formídes a fácie eórum: nec enim timére te fáciam vultum eórum. Ego quippe dedi te hódie in civitátem munítam, et in colúmnam férream, et in murum ǽreum, super omnem terram, régibus Juda, princípibus ejus, et sacerdótibus, et pópulo terræ. Et bellábunt advérsum te, et non prævalébunt: quia ego tecum sum, ait Dóminus, ut líberem te.\par
\switchcolumn
\EN
\noindent Thou therefore gird up thy loins, and arise, and speak to them all that I command thee. Be not afraid at their presence: for I will make thee not to fear their countenance. For behold I have made thee this day a fortified city, and a pillar of iron, and a wall of brass, over all the land, to the kings of Juda, to the princes thereof, and to the priests, and to the people of the land. And they shall fight against thee, and shall not prevail: for I am with thee, saith the Lord, to deliver thee.\par
\switchcolumn*

\LA
\begin{Resp}\Rsign Puéllæ saltánti imperávit mater: \Star{} Nihil áliud petas, nisi caput Joánnis. * Et contristátus est rex propter jusjurándum et propter simul discumbéntes.\end{Resp}
\begin{Resp}\Vsign Ait puélla matri suæ: Quid petam? At illa ait.\end{Resp}
\begin{Resp}\Rsign Nihil áliud petas, nisi caput Joánnis.\end{Resp}
\begin{Resp}\Vsign Glória Patri, et Fílio, \Star{} et Spirítui Sancto.\end{Resp}
\begin{Resp}\Rsign Et contristátus est rex propter jusjurándum et propter simul discumbéntes.\end{Resp}
\switchcolumn
\EN
\begin{Resp}\Rsign The damsel danced, and her mother charged her, saying: \Star{} See thou ask nothing but only the head of John. And the king was sorry, for his oath's sake, and for their sakes which sat with him.\end{Resp}
\begin{Resp}\Vsign The damsel said unto her mother What shall I ask? And she said:\end{Resp}
\begin{Resp}\Rsign See thou ask nothing, but only the head of John. And the king was sorry, for his oath's sake, and for their sakes which sat with him.\end{Resp}
\begin{Resp}\Vsign Glory be to the Father, and to the Son, \Star{} and to the Holy Ghost.\end{Resp}
\begin{Resp}\Rsign And the king was sorry, for his oath's sake, and for their sakes which sat with him.\end{Resp}
\switchcolumn*

\LA
\LessonHead{Lectio IV}
\switchcolumn
\EN
\LessonHead{Lesson IV}
\switchcolumn*

\LA
\LessonTitle{Ex libro sancti Ambrósii Epíscopi de Virgínibus}
\switchcolumn
\EN
\LessonTitle{From the Book upon Virgins, written by St. Ambrose, Bishop (of Milan.)}
\switchcolumn*

\LA
\Source{Lib. 3 post initium}
\switchcolumn
\EN
\Source{ii. 6}
\switchcolumn*

\LA
\noindent Quóniam beáti Joánnis Baptístæ non strictim prætereúnda est recordátio, ínterest ut quis et a quibus et quam ob causam, quo modo et quo témpore sit occísus, advértere debeámus. Ab adúlteris justus occíditur, et a reis in júdicem capitális scéleris pœna convértitur. Deínde prǽmium saltatrícis, mors est Prophétæ. Postrémo (quod étiam omnes bárbari horrére consuevérunt) inter épulas atque convívia consummándæ crudelitátis profértur edíctum; et a convívio ad cárcerem, de cárcere ad convívium ferális flagítii circumfértur obséquium. Quanta in uno facínore sunt crímina!\par
\switchcolumn
\EN
\noindent We must not hurry by the record of the Blessed Baptist John. We must ask what he was, and by whom, and why, and how, and when he was slain. He was a righteous man murdered by adulterers. The guilty passed upon their judge the sentence of death. Moreover, the death of the Prophet was the fee of a dancing-girl. And lastly, there was a feature about it from which even savages shrink; the order for completing the atrocity was given amid the merriment of a dinner-party. From banquet to prison, from prison to banquet, that was the course run by the servants of the murderer. How many horrors does this simple crime embrace within its details?\par
\switchcolumn*

\LA
\begin{Resp}\Rsign Justus germinábit sicut lílium: \Star{} Et florébit in ætérnum ante Dóminum.\end{Resp}
\begin{Resp}\Vsign Plantátus in domo Dómini, in átriis domus Dei nostri.\end{Resp}
\begin{Resp}\Rsign Et florébit in ætérnum ante Dóminum.\end{Resp}
\switchcolumn
\EN
\begin{Resp}\Rsign The righteous shall grow as the lily \Star{} Yea, he shall flourish in the presence of the Lord for ever.\end{Resp}
\begin{Resp}\Vsign Those that be planted in the house of the Lord, shall flourish in the courts of the house of our God.\end{Resp}
\begin{Resp}\Rsign Yea, he shall flourish in the presence of the Lord for ever.\end{Resp}
\switchcolumn*

\LA
\LessonHead{Lectio V}
\switchcolumn
\EN
\LessonHead{Lesson V}
\switchcolumn*

\LA
\noindent Quis non, cum e convívio ad carcerem cursari vidéret, putaret Prophétam jussum esse dimitti? Quis, inquam, cum audísset natalem esse Herodis, solemne convivium, puellæ optiónem eligéndi quod vellet datam; missum ad Joánnem ob solutiónem non arbitrarétur? Quid crudelitáti cum deliciis? quid cum funéribus voluptati? Rápitur ad pœnam Prophéta convivali témpore, convivali præcépto, quo non cuperet vel absolvi: perímitur gládio, caput ejus affertur in disco. Hoc crudelitáti férculum debebátur, quo insatiáta épulis féritas vescerétur.\par
\switchcolumn
\EN
\noindent Who is there, that, on seeing the messenger hasten from the dinner-table to the prison, would not have forthwith concluded that he carried an order for the Prophet's release. If any one had heard that it was Herod's birth-day, and that he was giving a great feast, and that he had offered a damsel the choice of whatever she listed, and that thereupon a messenger had been sent to John's dungeon. If any one, I say, had heard this, what would he have supposed? He would have concluded that the damsel had asked and obtained John's freedom. What have executions in common with dinners, or death with gaiety? While the banquet was going on, the Prophet was hurried to death, by an order from the reveller whom he had not troubled even by a prayer for release. He was slain with the sword, and his head was served up in a plate. This was the new dish demanded by a cruelty which the Feast had been powerless to glut.\par
\switchcolumn*

\LA
\begin{Resp}\Rsign Iste cognóvit justítiam, et vidit mirabília magna, et exorávit Altíssimum: \Star{} Et invéntus est in número Sanctórum.\end{Resp}
\begin{Resp}\Vsign Iste est, qui contempsit vitam mundi, et pervenit ad cæléstia regna.\end{Resp}
\begin{Resp}\Rsign Et invéntus est in número Sanctórum.\end{Resp}
\switchcolumn
\EN
\begin{Resp}\Rsign This is he which knew righteousness, and saw great wonders, and made his prayer unto the Most High \Star{} And he is numbered among the Saints.\end{Resp}
\begin{Resp}\Vsign This is he which loved not his life in this world, and is come unto an everlasting kingdom.\end{Resp}
\begin{Resp}\Rsign And he is numbered among the Saints.\end{Resp}
\switchcolumn*

\LA
\LessonHead{Lectio VI}
\switchcolumn
\EN
\LessonHead{Lesson VI}
\switchcolumn*

\LA
\noindent Intuére, rex acerbíssime, tuo spectacula digna convivio. Pórrige déxteram, ne quid sævítiæ tuæ desit; ut inter digitos tuos rivi défluant sacri cruoris. Et, quóniam non exsaturari épulis fames, non restingui póculis pótuit inauditæ sævítiæ sitis; bibe sánguinem scaturiéntibus adhuc venis exsecti cápitis profluentem. Cerne óculos in ipsa morte sceleris tui testes, aversántes conspéctum deliciárum. Claudúntur lúmina non tam mortis necessitate quam horrore luxuriæ. Os aureum illud exsangue, cujus senténtiam ferre non poteras, conticescit, et adhuc timétur.\par
\switchcolumn
\EN
\noindent Look, savage King, look at a decoration which suiteth well with thy banquet. Put out thine hand, so as to lose no part of the luxury of cruelty, and let the streams of the sacred blood run between thy fingers. Thine hunger the dinner hath been unable to satisfy, thy cups have not been able to quench thine inhuman thirst. Suck, suck the blood which the still palpitating veins are discharging from the place where the neck has been severed. Look at the eyes. Even in death they remain the eyes of a witness of thine uncleanness, but they are closing themselves upon the spectacle of thy pleasures. Those eyes indeed are shutting but it seems not so much from the laws of natural death, as from horror at the scene of thine enjoyment. The golden mouth, whose bloodless lips are silent now, can repeat no more the denunciation which thou couldest not bear to hear, and still thou art afraid of it.\par
\switchcolumn*

\LA
\begin{Resp}\Rsign Honestum fecit illum Dóminus, et custodívit eum ab inimícis, et a seductóribus tutávit illum: \Star{} Et dedit illi claritátem ætérnam.\end{Resp}
\begin{Resp}\Vsign Descendítque cum illo in fóveam, et in vinculis non derelíquit eum.\end{Resp}
\begin{Resp}\Rsign Et dedit illi claritátem ætérnam.\end{Resp}
\begin{Resp}\Vsign Glória Patri, et Fílio, \Star{} et Spirítui Sancto.\end{Resp}
\begin{Resp}\Rsign Et dedit illi claritátem ætérnam.\end{Resp}
\switchcolumn
\EN
\begin{Resp}\Rsign The Lord made him honourable, and defended him from his enemies, and kept him safe from those that lay in wait for him, \Star{} And gave him perpetual glory.\end{Resp}
\begin{Resp}\Vsign He went down with him into the pit, and left him not in bonds.\end{Resp}
\begin{Resp}\Rsign And gave him perpetual glory.\end{Resp}
\begin{Resp}\Vsign Glory be to the Father, and to the Son, \Star{} and to the Holy Ghost.\end{Resp}
\begin{Resp}\Rsign And gave him perpetual glory.\end{Resp}
\switchcolumn*

\LA
\LessonHead{Lectio VII}
\switchcolumn
\EN
\LessonHead{Lesson VII}
\switchcolumn*

\LA
\LessonTitle{Léctio sancti Evangélii secúndum Marcum}
\switchcolumn
\EN
\LessonTitle{From the Holy Gospel according to Mark}
\switchcolumn*

\LA
\Cite{Marc 6:17-29}
\switchcolumn
\EN
\Cite{Mark 6:17-29}
\switchcolumn*

\LA
\noindent In illo témpore: Misit Heródes ac ténuit Joánnem, et vinxit eum in cárcere propter Herodíadem, uxórem Philíppi fratris sui, quia duxerat eam. Et réliqua.\par
\switchcolumn
\EN
\noindent At that time Herod had sent forth, and laid hold upon John, and bound him in prison, for Herodias' sake, his brother Philip's wife, for he had married her. And so on.\par
\switchcolumn*

\LA
\LessonTitle{Homilía sancti Augustíni Epíscopi}
\switchcolumn
\EN
\LessonTitle{Homily by St. Augustine, Bishop (of Hippo.)}
\switchcolumn*

\LA
\Source{Sermo 10 in novis. Serm.}
\switchcolumn
\EN
\Source{New Sermons.}
\switchcolumn*

\LA
\noindent Cum sanctum Evangélium legerétur, crudele spectaculum ante óculos nostros constitutum est: caput sancti Joánnis in disco, feralis missus crudelitátis propter ódium veritátis. Puella saltat, et sævit mater; et inter lascivias et delicias convivántium témere jurátur, et ímpie, quod jurátur, implétur. Factum est Joánni, quod ipse prædixerat; de Dómino enim Jesu Christo dixerat: Illum oportet crescere, me autem minui. Iste minutus est in cápite, ille crevit in cruce. Odium péperit véritas. Non pótuit æquo animo tolerári quod homo Dei sanctus monebat; qui útique salútem eórum quærebat, quos sic monebat. Respondérunt illi mala pro bonis.\par
\switchcolumn
\EN
\noindent The reading of the Holy Gospel hath set a scene of cruelty before our eyes even the head of St. John in a charger; a message of death sent forth to discharge the bloody commands of one that hateth the truth; a damsel dancing, and a mother rabid; a rash oath sworn in the midst of uncleanness and the revels of a supper, and a wicked fulfillment of the oath so sworn. It befell unto John according to his own saying. For he had said concerning the Lord Jesus Christ: He must increase, but I must decrease (John iii. 30,) so John decreased by an head, and Christ's height was made higher upon the Cross. The truth drew hatred. It could not be borne in patience that the holy man of God should utter a rebuke, albeit he sought by his rebuke nothing but the soul's health of them to whom he addressed it. They repaid him evil for good.\par
\switchcolumn*

\LA
\begin{Resp}\Rsign Desidérium ánimæ ejus tribuísti ei Dómine, \Star{} Et voluntáte labiórum ejus non fraudásti eum.\end{Resp}
\begin{Resp}\Vsign Quóniam prævenísti eum in benedictiónibus dulcédinis: posuísti in cápite ejus corónam de lápide pretióso.\end{Resp}
\begin{Resp}\Rsign Et voluntáte labiórum ejus non fraudásti eum.\end{Resp}
\switchcolumn
\EN
\begin{Resp}\Rsign Lord, Thou hast given him his heart's desire. \Star{} And hast not withholden the request of his lips.\end{Resp}
\begin{Resp}\Vsign For Thou hast prevented him with the blessings of sweetness: Thou hast set a crown of precious stones upon his head.\end{Resp}
\begin{Resp}\Rsign And hast not withholden the request of his lips.\end{Resp}
\switchcolumn*

\LA
\LessonHead{Lectio VIII}
\switchcolumn
\EN
\LessonHead{Lesson VIII}
\switchcolumn*

\LA
\noindent Quid enim ille diceret, nisi quo plenus erat? et quid illi respondérent, nisi quo pleni erant? Ille tríticum seminávit, sed spinas invénit. Dicebat regi: Non licet tibi habere uxórem fratris tui. Vincebat enim regem libído: tenebat apud se prohíbitam uxórem fratris sui. Sed eum tamen sic libebat, ut non sæviret; honorábat eum, a quo verum audiebat. Sed mulier detestábilis ódium concipiebat, quod aliquándo dato témpore páreret. Quando autem parturiebat, péperit filiam, filiam saltántem.\par
\switchcolumn
\EN
\noindent For what could he say but that whereof he was full? And what could they answer him but that whereof they were full? He sowed wheat, and found thorns. He had said unto the King: It is not lawful for thee to have thy brother's wife. Lust had got the better of the King, and he kept a woman whom it was not lawful for him to have, even his brother's wife. But she pleased him, so that his cruelty was lulled. He respected the Saint who had spoken the truth to him. But the horrible woman conceived hatred, and by-and-by brought it forth. When she brought forth, she brought forth a girl, a dancing-girl.\par
\switchcolumn*

\LA
\begin{Resp}\Rsign Stola jucunditátis índuit eum Dóminus: \Star{} Et corónam pulchritúdinis pósuit super caput ejus.\end{Resp}
\begin{Resp}\Vsign Cibávit illum Dóminus pane vitæ et intelléctus: et aqua sapiéntiæ salutáris potávit illum.\end{Resp}
\begin{Resp}\Rsign Et corónam pulchritúdinis pósuit super caput ejus.\end{Resp}
\begin{Resp}\Vsign Glória Patri, et Fílio, \Star{} et Spirítui Sancto.\end{Resp}
\begin{Resp}\Rsign Et corónam pulchritúdinis pósuit super caput ejus.\end{Resp}
\switchcolumn
\EN
\begin{Resp}\Rsign The Lord hath put on him a robe of honour, \Star{} And put about his head a crown of joy.\end{Resp}
\begin{Resp}\Vsign With the bread of life and understanding hath the Lord fed him, and given him the water of health and wisdom to drink.\end{Resp}
\begin{Resp}\Rsign And put about his head a crown of joy.\end{Resp}
\begin{Resp}\Vsign Glory be to the Father, and to the Son, \Star{} and to the Holy Ghost.\end{Resp}
\begin{Resp}\Rsign And put about his head a crown of joy.\end{Resp}
\switchcolumn*

\end{paracol}

\par\needspace{10\baselineskip}\addvspace{1.2em}{\centering\small\textsc{Die 29 Augusti} · \textsc{29 August}\par}
\Office{S. Sabinæ Martyris}{St. Sabina, Martyr}{Simplex}
\addcontentsline{toc}{subsection}{Die 29 Augusti · S. Sabinæ Martyris \textit{— St. Sabina, Martyr}}
\begin{paracol}{2}
\LA
\LessonHead{Lectio IX (pro commemoratione)}
\switchcolumn
\EN
\LessonHead{Lesson IX (when commemorated)}
\switchcolumn*

\LA
\LessonTitle{Commemoratio: S. Sabinæ Martyris}
\switchcolumn
\EN
\LessonTitle{Commemoration of: St. Sabina, Martyr}
\switchcolumn*

\LA
\noindent Sabina, mulier Romana, Valentini viri claríssimi uxor, a Seráphia vírgine christianæ fidei præceptis instituta, post sanctæ Vírginis martyrium, collectas ejus relíquias piis exsequiis sepelívit. Quæ propter eam causam paulo post, Hadriáno imperatóre, comprehensa, Elpidio júdici sistitur. Cui is: Tu ne illa Sabina et génere et matrimonio nobilíssima? At illa, Sum, inquit: sed Dómino meo Jesu Christo grátias ago, qui me, intercessióne Seráphiæ famulæ suæ, e dæmonum potestate liberávit. Quam varie tentatam ut propositum mutaret, cum a fidei constántia movére non posset, præfectus, pronuntiáta senténtia quod deos contemneret, cápitis damnávit. Ejus corpus a Christiánis in eodem sepúlcro cónditum est, in quo ipsa magistram fidei suæ Seráphiam posúerat.\par
\switchcolumn
\EN
\noindent Sabina was a Roman lady, the wife of a distinguished nobleman named Valentine. The Christian faith was taught to her by a maiden named Seraphia. After the martyrdom of this holy virgin, Sabina gathered together her reliques, and buried them with godly service. For this cause she was in a little while arrested, under the Emperor Hadrian, and brought before the Judge Elpidius. Art thou, said he, the same Sabina who is so distinguished for her blood and for her marriage She answered “I am, but I give thanks to my Lord Jesus Christ for having delivered me through the prayers of His hand-maiden Seraphia from the troubling of devils.” Diverse attempts were made to make her change her mind, but when they proved in vain the Praefect passed sentence of death upon her for despising the gods. The Christians laid her body in the same grave in which she had herself laid that of Seraphia, her teacher in the faith.\par
\switchcolumn*

\LA
\TeDeum{Te Deum}
\switchcolumn
\EN
\TeDeum{Te Deum}
\switchcolumn*

\end{paracol}

\par\needspace{10\baselineskip}\addvspace{1.2em}{\centering\small\textsc{Die 30 Augusti} · \textsc{30 August}\par}
\Office{S. Rosæ a Sancta Maria Limanæ Virginis}{St. Rose of Lima, Virgin}{Duplex}
\addcontentsline{toc}{subsection}{Die 30 Augusti · S. Rosæ a Sancta Maria Limanæ Virginis \textit{— St. Rose of Lima, Virgin}}
\begin{paracol}{2}
\LA
\RefLine{Lectiones I–III: de Scriptura occurrente.}
\switchcolumn
\EN
\RefLine{Lessons I–III: from the Scripture of the day.}
\switchcolumn*

\LA
\RefLine{Lectiones VII–VIII: ex Commúni Virginum tantum (C6a).}
\switchcolumn
\EN
\RefLine{Lessons VII–VIII: from the Common of a Virgin not a Martyr (C6a).}
\switchcolumn*

\LA
\RefLine{Lectio IX: de Ss. Felicis et Adaucti Martyrum (08-30o).}
\switchcolumn
\EN
\RefLine{Lesson IX: of St. Felix and Adauctus, Martyrs (08-30o).}
\switchcolumn*

\LA
\LessonHead{Lectio IV}
\switchcolumn
\EN
\LessonHead{Lesson IV}
\switchcolumn*

\LA
\noindent Primus Americæ Meridionalis flos sanctitatis, virgo Rosa, christianis parentibus Limæ progenita, mox ab incunabulis claruit futuræ sanctimoniæ indiciis. Nam vultus infantis mirabiliter in rosæ effugiem transfiguratus, huic nomini occasionem dedit: cui postea Virgo Deipara cognomen adjecit, jubens vocari deinceps Rosam a sancta Maria. Quinquennis votum perpetuæ virginitatis emisit. Adultior, ne a parentibus ad nuptias cogeretur, clam sibimet venustissimam capitis cæsariem præscidit. Jejuniis supra humánum modum addicta, integras Quadragesimas transegit, pane abstinens, ac dietim solis quinque granulis mali citrini victitans.\par
\switchcolumn
\EN
\noindent The first flower of holiness which came to full blossom in South America, was the maiden Rose. She was born at Lima, of a Christian father and mother, (upon the 20th of April, in the year 1586,) and was remarkable from her childhood for marks of saintliness. The occasion of her name was a strange likeness to a rose, which her face assumed when she was a babe. To this name she afterwards added that of the Virgin Mother of God, desiring to be called St. Mary's Rose. At the age of fifteen years she uttered a vow of perpetual virginity. As she grew older, lest her parents should force her to marry, she polled her head of all her hair, which was very beautiful. She fasted to a degree almost superhuman, passing whole Lents without taking bread, and eating day by day only five pips of a lime.\par
\switchcolumn*

\LA
\begin{Resp}\Rsign Propter veritátem, et mansuetúdinem, et justítiam: \Star{} Et dedúcet te mirabíliter déxtera tua.\end{Resp}
\begin{Resp}\Vsign Spécie tua et pulchritúdine tua inténde, próspere procéde, et regna.\end{Resp}
\begin{Resp}\Rsign Et dedúcet te mirabíliter déxtera tua.\end{Resp}
\switchcolumn
\EN
\begin{Resp}\Rsign Because of truth, and meekness, and righteousness; \Star{} And thy right hand shall lead thee wonderfully.\end{Resp}
\begin{Resp}\Vsign In thy comeliness, and thy beauty, go forward, fare prosperously, and reign.\end{Resp}
\begin{Resp}\Rsign And thy right hand shall lead thee wonderfully.\end{Resp}
\switchcolumn*

\LA
\LessonHead{Lectio V}
\switchcolumn
\EN
\LessonHead{Lesson V}
\switchcolumn*

\LA
\noindent Habitu tertii ordinis sancti Dominici assumpto, pristinas vitæ austeritates duplicavit: oblongo asperrimoque cilicio sparsim mimisculas acus innexuit: sub velo coronam densis aculeis introrsus obarmatam interdiu noctuque gestavit. Sanctæ Catharinæ Senensis ardua premens vestigia, catena ferrea, triplici nexu circumducta, lumbos cinxit. Lectulum sibi e truncis nodosis composuit, horumque vacuas commissuras fragminibus testarum implevit. Cellulam sibi angustissimam struxit in extremo horti angulo, ubi cælestium contemplationi dedita, crebris disciplinis, inedia, vigiliis corpusculum extenuans, at spiritu vegetata, larvas dæmonum, frequenti certamine victrix, impavide protrivit ac superavit.\par
\switchcolumn
\EN
\noindent She took the habit of the Third Order of St. Dominic, and then doubled her former severities. She wore a long and very rough hair-cloth, into which she inserted small pins. She wore day and night under her veil a crown, the inner side of which was armed with pricks. In imitation of the hard steps of St Katharine of Sienna, she girded her loins with a threefold iron chain. She made to herself a bed of knotty sticks, and filled the gaps with broken bits of potsherd. She built herself a very small hut in the farthest corner of the garden, where she gave herself up to thoughts of heavenly things, and to punishing her body with often scourging, starvation, and sleeplessness. But she waxed strong in spirit, and though she often had to fight with evil ghosts, she conquered them, fearlessly prostrated them, and triumphed over them.\par
\switchcolumn*

\LA
\begin{Resp}\Rsign Dilexísti justítiam, et odísti iniquitátem: \Star{} Proptérea unxit te Deus, Deus tuus, óleo lætítiæ.\end{Resp}
\begin{Resp}\Vsign Propter veritátem, et mansuetúdinem, et justítiam.\end{Resp}
\begin{Resp}\Rsign Proptérea unxit te Deus, Deus tuus, óleo lætítiæ.\end{Resp}
\switchcolumn
\EN
\begin{Resp}\Rsign Thou hast loved righteousness, and hated iniquity; \Star{} Therefore God, thy God, hath anointed thee with the oil of gladness.\end{Resp}
\begin{Resp}\Vsign Because of truth, and meekness, and righteousness.\end{Resp}
\begin{Resp}\Rsign Therefore God, thy God, hath anointed thee with the oil of gladness.\end{Resp}
\switchcolumn*

\LA
\LessonHead{Lectio VI}
\switchcolumn
\EN
\LessonHead{Lesson VI}
\switchcolumn*

\LA
\noindent Aegritudinum tormentis, domesticorum insultibus, linguarum morsibus dire agitata, nondum satis pro merito se affligi quærebatur. Per quindecim annos ad plusculas horas desolatione spiritus et ariditate miserrime contabescens, forti animo tulit agones omni morte amariores. Exinde cœpit supernis abundare deliciis, illustrari visionibus, colliquescere seraphicis ardoribus. Angelo tutelori, sanctæ Catharinæ Senensi, Virgini Deiparæ inter assiduas apparitiones mire familiaris, a Christo has voces audire meruit: Rosa cordis mei, tu mihi sponsa esto. Denique Sponsi hujus paradiso feliciter invectam, plurimisque ante et post obitum miraculis coruscam, Clemens decimus Pontifex Maximus sanctarum Virginum catalogo ritu solemni adscripsit.\par
\switchcolumn
\EN
\noindent She suffered greatly from painful illnesses, from the maltreatment of the servants, and from slanderous accusations, but still complained that she did not suffer as much as she deserved. For fifteen years she pined in misery from desolation and dryness of spirit, bravely enduring torments worse than any form of death. After this period she began to overflow with consolation, to be enlightened by visions, and to melt with love like a Seraph's. She attained, by the frequency of visions, to a strange personal familiarity with her Guardian Angel, with St. Katharine of Sienna, and with the Virgin Mother of God, and she earned from Christ the words, Rose of My Heart, be thou My bride. She was famous for many miracles, both before and after she departed hence, and was happily transplanted into the Bridegroom's garden, (upon the 24th of August 1617, being aged 31 years.) Pope Clement X. with solemn pomp inscribed her name in the list of holy maidens.\par
\switchcolumn*

\LA
\begin{Resp}\Rsign Afferéntur Regi vírgines post eam, próximæ ejus; \Star{} Afferéntur tibi in lætítia et exsultatióne.\end{Resp}
\begin{Resp}\Vsign Spécie tua et pulchritúdine tua; inténde, próspere procéde, et regna.\end{Resp}
\begin{Resp}\Rsign Afferéntur tibi in lætítia et exsultatióne.\end{Resp}
\begin{Resp}\Vsign Glória Patri, et Fílio, \Star{} et Spirítui Sancto.\end{Resp}
\begin{Resp}\Rsign Afferéntur tibi in lætítia et exsultatióne.\end{Resp}
\switchcolumn
\EN
\begin{Resp}\Rsign After her shall virgins be brought unto the King, her fellows \Star{} They shall be brought unto thee with gladness and rejoicing.\end{Resp}
\begin{Resp}\Vsign In thy comeliness and thy beauty, go forward, fare prosperously, and reign.\end{Resp}
\begin{Resp}\Rsign They shall be brought unto thee with gladness and rejoicing.\end{Resp}
\begin{Resp}\Vsign Glory be to the Father, and to the Son, \Star{} and to the Holy Ghost.\end{Resp}
\begin{Resp}\Rsign They shall be brought unto thee with gladness and rejoicing.\end{Resp}
\switchcolumn*

\end{paracol}

\par\needspace{10\baselineskip}\addvspace{1.2em}{\centering\small\textsc{Die 30 Augusti} · \textsc{30 August}\par}
\Office{Ss. Felicis et Adaucti Martyrum}{St. Felix and Adauctus, Martyrs}{Simplex}
\addcontentsline{toc}{subsection}{Die 30 Augusti · Ss. Felicis et Adaucti Martyrum \textit{— St. Felix and Adauctus, Martyrs}}
\begin{paracol}{2}
\LA
\LessonHead{Lectio IX (pro commemoratione)}
\switchcolumn
\EN
\LessonHead{Lesson IX (when commemorated)}
\switchcolumn*

\LA
\LessonTitle{Commemoratio: Ss. Felicis et Adaucti Martyrum}
\switchcolumn
\EN
\LessonTitle{Commemoration of: St. Felix and Adauctus, Martyrs}
\switchcolumn*

\LA
\noindent Felix, Diocletiano et Maximiano imperatoribus, propter susceptam Christi religionem comprehensus, in Serapidis templum adductus est. Cui sacrificare cum juberetur, os simulacri conspuit. Quo facto, statim serea statua corruit. Quod cum iterum ac tertio in sede Mercurii Dianæque factum esset, impietatis et magicæ artis accusatus, equuleo torquetur. Mox ad secundum ab Urbe lapidem via Ostiensi ducitur, ut securi feriretur. Cui inter viam oblatus quidam Christianus, cum Felicem agnoscens ad martyrium duci videret, Ego quoque, clara voce inquit, eadem, qua iste, lege vivo ego eumdem Jesum Christum colo. Itaque Felicem osculatus, cum eo securi percutitur, tertio Kalendas Septembris. Cujus nomen cum ignotum esset Christianis, is Adaucti nomine nobilitatus est, quod sancto Martyri Felici adauctus sit ad coronam\par
\switchcolumn
\EN
\noindent Felix was arrested in the reign of the Emperors Diocletian and Maximian, on the charge of having embraced the Christian Faith, and was brought to the temple of Serapis. When he was ordered to offer sacrifice, he spat in the face of the brazen idol, which thereupon fell down. When this happened a second and third time in the temples of Mercury and Diana, he was accused of impiety and magic, and tortured upon the rack. It was not long, however, before he was led out to the second mile-stone upon the road to Ostia, to be smitten with the axe. As they were on the way thither, they chanced to meet a certain Christian, who, when he knew that Felix was going to finish his testimony, said aloud, I live by the same law as he doth I worship the same Christ Jesus. And therewith he kissed Felix, and they were beheaded together, upon the 30th day of August. What the name of the second person was the Christians never knew, and he is therefore honoured under the title of Him-who-was-added x that is, added to the company of the Holy Martyr Felix in winning of the crown.\par
\switchcolumn*

\LA
\TeDeum{Te Deum}
\switchcolumn
\EN
\TeDeum{Te Deum}
\switchcolumn*

\end{paracol}

\par\needspace{10\baselineskip}\addvspace{1.2em}{\centering\small\textsc{Die 31 Augusti} · \textsc{31 August}\par}
\Office{S. Raymundi Nonnati Confessoris}{St. Raymond Nonnatus, Confessor}{Duplex}
\addcontentsline{toc}{subsection}{Die 31 Augusti · S. Raymundi Nonnati Confessoris \textit{— St. Raymond Nonnatus, Confessor}}
\begin{paracol}{2}
\LA
\RefLine{Lectiones I–III: de Scriptura occurrente.}
\switchcolumn
\EN
\RefLine{Lessons I–III: from the Scripture of the day.}
\switchcolumn*

\LA
\RefLine{Lectiones VII–IX: ex Commúni Confessoris non pontificis (C5).}
\switchcolumn
\EN
\RefLine{Lessons VII–IX: from the Common of a Confessor not a Bishop (C5).}
\switchcolumn*

\LA
\LessonHead{Lectio IV}
\switchcolumn
\EN
\LessonHead{Lesson IV}
\switchcolumn*

\LA
\noindent Raymundus, Nonnatus cognomento dictus, quia præter communem naturæ legem e mortuæ matris dissecto latere in lucem eductus fuit, Portelli in Catalaunia piis et nobilibus parentibus ortus, ab ipsa infantia futuræ sanctitatis indicia dedit. Nam puerilia oblectamenta, mundique illecebras respuens, ita pietati operam dabat, ut omnes in puero adultam virtutem admirarentur. Crescente vero ætate, litterarum studiis incubuit: sed mox jubente patre vitam ruri agens, sacellum sancti Nicolai in Portelli finibus situm crebro adibat, ut sacram Deiparæ imaginem, quæ in eo summa fidelium veneratione étiam nunc colitur, visitaret. Ibi effusus in preces, ipsam Dei parentem, ut se in filium adoptare viamque salutis ac scientiam sanctorum edocere dignaretur, enixe deprecabatur.\par
\switchcolumn
\EN
\noindent This Raymond is commonly called the Unborn, because his was one of the rare cases in which the child is not brought into the world in the course of nature, but by a surgical operation after the death of the mother. He was the son of godly and noble parents, at Portel, (in the dioecese of Urgel,) in Catalonia. The tokens of his holy after-life appeared even in his childhood. The things that delight children, and the attractions of the world, had no charm for him. He was so earnest in godliness that all men marvelled at his habits of premature old age. As he grew older, he gave himself to the study of letters, but, at the command of his father, turned to farming. He went often to the Chapel of St Nicolas, in the suburbs of Portel, to visit the sacred image of the Mother of God, which is still sought with great tenderness by the faithful. There he poured forth his soul in prayer, and earnestly entreated the Mother of God herself to be pleased to take him for her son, to show him the way wherein it should be safe for him to walk, and to teach him the science of the Saints.\par
\switchcolumn*

\LA
\begin{Resp}\Rsign Honéstum fecit illum Dóminus, et custodívit eum ab inimícis, et a seductóribus tutávit illum: \Star{} Et dedit illi claritátem ætérnam.\end{Resp}
\begin{Resp}\Vsign Justum dedúxit Dóminus per vias rectas, et osténdit illi regnum Dei.\end{Resp}
\begin{Resp}\Rsign Et dedit illi claritátem ætérnam.\end{Resp}
\switchcolumn
\EN
\begin{Resp}\Rsign The Lord made him honourable, and defended him from his enemies, and kept him safe from those that lay in wait for him; \Star{} And gave him perpetual glory.\end{Resp}
\begin{Resp}\Vsign He went down with him into the pit, and left him not in bonds.\end{Resp}
\begin{Resp}\Rsign And gave him perpetual glory.\end{Resp}
\switchcolumn*

\LA
\LessonHead{Lectio V}
\switchcolumn
\EN
\LessonHead{Lesson V}
\switchcolumn*

\LA
\noindent Nec defuit votis ejus benignissima Virgo. Ab ipsa enim intellexit gratissimum sibi fore, si religionem sub titulo de Mercede, seu de Misericordia redemptionis captivorum, ea suggerante nuper fundatam, ingrederetur. Qua monitione percepta, Barcinonem statim profectus, illud tam præcellentis erga proximum caritatis institutum amplexus est. Regulari igitur militiæ adscriptus, virginitatem, quam pridem beatæ Virgini consecraverat, perpetuo coluit, ceterisque virtutibus enituit, caritate præsertim erga Christianos, qui sub potestate paganorum miseram in captivitate vitam degebant. Hos ut redimeret, in Africam missus, cum jam multos a servitute liberasset, ne, consumpta pecunia, aliis item in proximo abnegandæ fidei discrimine constitutis deesset, se ipsum pignori dedit; sed cum ardentissimo salutis animarum desiderio succensus, plures Mahometanos suis concionibus ad Christum converteret, in arctam custodiam a barbaris conjectus, variisque suppliciis cruciatus, mox labiis perforatis et sera ferrea clausis, crudele martyrium diu sustinuit.\par
\switchcolumn
\EN
\noindent And the most gracious Maiden was not deaf to his prayers. From her he understood that it would please her right well, if he would join the Religious Order which had just been founded at her own inspiration, styled of Ransom or of Mercy, for buying up and freeing slaves. As soon as he had received this intimation from her, he went to Barcelona, and entered the Institute so nobly dedicated to love for our neighbour. Once enlisted in the Regular Army, he guarded unspotted for ever the virginity which he had already consecrated to the Blessed Virgin for ever. But he was a bright and shining light of all other good words and works, especially of tender compassion for Christians who were passing a life of grievous bondage in the possession of unbelieving masters. To free such he was sent into Africa, and delivered many. But his money ran short, and as there were still many in imminent danger of denying the faith, he pawned himself. He was enkindled with a most vehement longing for the salvation of souls, and by his exhortations brought diverse Mohammedans to Christ. The Moors therefore threw him into close prison, and put him to diverse tortures, at last making holes through his lips and locking them together with an iron padlock, which horrid cruelty he long bore.\par
\switchcolumn*

\LA
\begin{Resp}\Rsign Amávit eum Dóminus, et ornávit eum: stolam glóriæ índuit eum, \Star{} Et ad portas paradísi coronávit eum.\end{Resp}
\begin{Resp}\Vsign Índuit eum Dóminus lorícam fídei, et ornávit eum.\end{Resp}
\begin{Resp}\Rsign Et ad portas paradísi coronávit eum.\end{Resp}
\switchcolumn
\EN
\begin{Resp}\Rsign The Lord loved him and beautified him; He clothed him with a robe of glory, \Star{} And crowned him at the gates of Paradise.\end{Resp}
\begin{Resp}\Vsign The Lord hath put on him the breast-plate of faith, and hath adorned him.\end{Resp}
\begin{Resp}\Rsign And crowned him at the gates of Paradise.\end{Resp}
\switchcolumn*

\LA
\LessonHead{Lectio VI}
\switchcolumn
\EN
\LessonHead{Lesson VI}
\switchcolumn*

\LA
\noindent Ob hæc et alia fortiter gesta sanctitatis ejus fama longe lateque diffúsa est. Qua permotus Gregorius nonus, in amplissimum sanctæ Romanæ Ecclesiæ Cardinalium collegium Raymundum adscripsit: sed vir Dei in ea dignitate ab omni pompa abhorrens, religiosæ humilitatis tenacissimus semper fuit. Romam vero pergens, statim ac Cardonam pervenit, extremo morbo confectus, ecclesiasticis sacramentis muniri summis precibus postulavit. Cumque morbus ingravesceret, et sacerdos diutius tardaret, Angelorum ministerio, sub specie religiosorum sui ordinis apparentium, salutari viatico refectus fuit. Quo sumpto, et gratiis Deo peractis, migravit ad Dominum Dominica ultima Augusti, anno millesimo ducentesimo quadragesimo. Mortui corpus, cum circa locum sepulturæ contentio orta esset, arcæ inclusum, et mulæ cæcæ impositum, ad sacellum sancti Nicolai Dei nutu delatum fuit, ut ibi tumularetur, ubi prima jecerat sanctioris vitæ fundamenta. Illic constructo sui ordinis cænobio, a confluentibus voti causa ex universa Catalaunia fidelibus populis honoratur, variis miraculis et signis gloriosus.\par
\switchcolumn
\EN
\noindent The account of these, and other brave things that he did, he got the name of a Saint far and wide. Gregory IX. was moved thereby to make Raymond a Cardinal of the Holy Roman Church, but in this place of honour the man of God shrank from all outward show, and clung ever tightly to the lowliness that beseemeth a Religious man. He had started for Rome, (in obedience to the command of the Pope,) but had only got as far as Cardona, (six miles from Barcelona,) when he was seized with his last illness, and earnestly called for the strengthening Sacraments of the Church. But his position became critical, and the Priest had not arrived. Then Angels came unto him, clad in the habit of his own Order, and ministered unto him the wholesome Provision for the last journey. When he had taken It, he gave God thanks, and departed hence to be ever with the Lord. It was the last Lord's Day in August, 1240. After his death there was some dispute arose as to where his body should be buried so they shut it up in a box, and laid it upon a blind mule, and the beast was guided by God to carry it to the chapel of St. Nicolas, that he might be buried where he had laid the foundations of his nobler life. There was built there a Convent of his Order, and the faithful come together thither from all parts of Catalonia to honour him, and he is famous for diverse signs and wonders.\par
\switchcolumn*

\LA
\begin{Resp}\Rsign Iste homo perfécit ómnia quæ locútus est ei Deus, et dixit ad eum: Ingrédere in réquiem meam: \Star{} Quia te vidi justum coram me ex ómnibus géntibus.\end{Resp}
\begin{Resp}\Vsign Iste est, qui contémpsit vitam mundi, et pervénit ad cæléstia regna.\end{Resp}
\begin{Resp}\Rsign Quia te vidi justum coram me ex ómnibus géntibus.\end{Resp}
\begin{Resp}\Vsign Glória Patri, et Fílio, \Star{} et Spirítui Sancto.\end{Resp}
\begin{Resp}\Rsign Quia te vidi justum coram me ex ómnibus géntibus.\end{Resp}
\switchcolumn
\EN
\begin{Resp}\Rsign This is he which did according unto all that God commanded him; and God said unto him: Enter thou into My rest, \Star{} For thee have I seen righteous before Me among all people.\end{Resp}
\begin{Resp}\Vsign This is he which loved not his life in this world, and is come unto an everlasting kingdom.\end{Resp}
\begin{Resp}\Rsign For thee have I seen righteous before Me among all people.\end{Resp}
\begin{Resp}\Vsign Glory be to the Father, and to the Son, \Star{} and to the Holy Ghost.\end{Resp}
\begin{Resp}\Rsign For thee have I seen righteous before Me among all people.\end{Resp}
\switchcolumn*

\end{paracol}

\SeasonTitle{Mensis September}{September}

\addcontentsline{toc}{section}{Mensis September \textit{— September}}

\par\needspace{10\baselineskip}\addvspace{1.2em}{\centering\small\textsc{Die 1 Septembris} · \textsc{1 September}\par}
\Office{S. Ægidii Abbatis}{St. Giles, Abbot}{Simplex}
\addcontentsline{toc}{subsection}{Die 1 Septembris · S. Ægidii Abbatis \textit{— St. Giles, Abbot}}
\begin{paracol}{2}
\LA
\RefLine{Lectiones I–II: de Scriptura occurrente.}
\switchcolumn
\EN
\RefLine{Lessons I–II: from the Scripture of the day.}
\switchcolumn*

\LA
\LessonHead{Lectio III (pro commemoratione)}
\switchcolumn
\EN
\LessonHead{Lesson III (when commemorated)}
\switchcolumn*

\LA
\noindent Ægidius Atheniénsis, regiæ stirpis, a prima ætate divinis litteris et caritátis officiis ita déditus fuit, ut nihil præterea curare viderétur. Itaque, paréntibus mórtuis, totum patrimónium in páuperes erogávit; quin étiam túnicam exuit, ut ægrotum egentem tégeret, qua ille indútus, statim conváluit. Sed multis deinceps clarior miraculis, timens sui nóminis celebritátem, Arelátem ad beátum Cæsárium contendit. A quo post biennium discédens, secessit in eremum; ubi diutius herbárum radicibus et cervæ lacte, quæ statis ad eum horis veniebat, admirábili sanctitáte vixit. Quæ cerva, insequéntibus quodam die cánibus regiis, cum in antrum Ægidii refugísset, Galliæ regem ímpulit, ut ab eo summis precibus péteret, ut in loco speluncæ monastérium éxstrui paterétur. Cujus administratiónem, flagitante rege, invitus suscépit; eoque munere aliquot annis prudenter piéque gesto, migrávit in cælum.\par
\switchcolumn
\EN
\noindent The holy Abbot Giles was by birth an Athenian, and of Royal lineage. From his youth he showed ever such a love for sacred learning and for works of charity, that he seemed to care for nothing else. When his father and mother were dead, he bestowed his whole inheritance upon the poor. He took off even his own coat, to clothe a poor sick man withal, and the sick man was healed forthwith as soon as he put it on him. As Giles became famous for working miracles, he fled from glory among men, and betook him to Arles, (in France,) to the company of blessed Caesarius. After the space of two years he departed thence, and went into the desert, for he lived in wonderful holiness for a long while upon the roots of herbs and the milk of an hind, which came to him at regular hours. This hind was chased one day by the King's hounds, and took refuge in Giles's cave. Thereby the King of France was moved earnestly to entreat of him that he would suffer a monastery to be built in the place where this cave was. At the instant desire of the King, he took the rule of this monastery, albeit himself unwilling, and discharged this duty wisely and godly for some years, until he passed away to heaven.\par
\switchcolumn*

\LA
\TeDeum{Te Deum}
\switchcolumn
\EN
\TeDeum{Te Deum}
\switchcolumn*

\end{paracol}

\par\needspace{10\baselineskip}\addvspace{1.2em}{\centering\small\textsc{Die 2 Septembris} · \textsc{2 September}\par}
\Office{S. Stephani Hungariæ Regis Confessoris}{St. Stephen, King of Hungary, Confessor}{Semiduplex}
\addcontentsline{toc}{subsection}{Die 2 Septembris · S. Stephani Hungariæ Regis Confessoris \textit{— St. Stephen, King of Hungary, Confessor}}
\begin{paracol}{2}
\LA
\RefLine{Lectiones I–III: de Scriptura occurrente.}
\switchcolumn
\EN
\RefLine{Lessons I–III: from the Scripture of the day.}
\switchcolumn*

\LA
\LessonHead{Lectio IV}
\switchcolumn
\EN
\LessonHead{Lesson IV}
\switchcolumn*

\LA
\noindent Stephanus in Hungariam Christi fidem et regium nomen invéxit. Regia coróna a Romano Pontifice impetrata, ejusque jussu in regem inunctus, regnum Sedi apostolicæ obtulit. Varia pietátis domicília Romæ, Jerosolymis, Constantinopoli: in Hungaria archiepiscopátum Strigoniensem, episcopátus decem admirábili religióne et munificéntia fundavit. Par in páuperes amor et liberalitas, quos véluti Christum ipsum complectens, neminem a se mærentem ac vacuum umquam dismisit; quin ad eórum inópiam sublevandam, amplíssimis facultátibus erogatis, domesticam quoque supellectilem exímia benignitate frequenter distríbuit. Suis ínsuper mánibus lavare páuperum pedes, noctu solus et ignotus nosocomía frequentare, decumbéntibus inservire ac cetera caritátis offícia exhibere consuevit. Quarum virtútum merito, illíus déxtera, resoluto cetero corpore, incorrupta permansit.\par
\switchcolumn
\EN
\noindent This Stephen (was the son of Geysa, fourth Duke of the Hungarians, and was born at Gran in the year 977.) He it was who first gave to Hungary the faith of Christ and the name of a kingdom. He obtained the Kingly crown from the Bishop of Rome, and being by command of the same anointed King, he made an offering of his kingdom to the Apostolic See. With wonderful devotion and bounty he founded diverse godly houses at Rome, Jerusalem, and Constantinople, and in Hungary the Archbishopric of Gran and ten other Sees. Toward the poor he had the same love and bounty. He greeted them as though they were Christ Himself, and never sent any one away sorrowing and empty. He spent vast sums in relieving their poverty, and also often parted among them with exceeding tenderness even the furniture of his house. Moreover it was his use to wash the feet of the poor with his own hands, and to go in the night, alone and unknown, to the hospitals, and to wait on them that lay there, and show them other deeds of kindness. It was the reward of these good works that, when the rest of his body decayed, his right hand remained uncorrupt.\par
\switchcolumn*

\LA
\begin{Resp}\Rsign Honéstum fecit illum Dóminus, et custodívit eum ab inimícis, et a seductóribus tutávit illum: \Star{} Et dedit illi claritátem ætérnam.\end{Resp}
\begin{Resp}\Vsign Justum dedúxit Dóminus per vias rectas, et osténdit illi regnum Dei.\end{Resp}
\begin{Resp}\Rsign Et dedit illi claritátem ætérnam.\end{Resp}
\switchcolumn
\EN
\begin{Resp}\Rsign The Lord made him honourable, and defended him from his enemies, and kept him safe from those that lay in wait for him; \Star{} And gave him perpetual glory.\end{Resp}
\begin{Resp}\Vsign He went down with him into the pit, and left him not in bonds.\end{Resp}
\begin{Resp}\Rsign And gave him perpetual glory.\end{Resp}
\switchcolumn*

\LA
\LessonHead{Lectio V}
\switchcolumn
\EN
\LessonHead{Lesson V}
\switchcolumn*

\LA
\noindent Orándi studio noctes pene totas ducebat insomnes; atque in cæléstium rerum contemplatióne defixus, interdum extra sensus raptus, sublimis in áëra ferri visus fuit. Perduellium conspiratiónes ac validiórum hostium ímpetus, miro prorsus modo, non semel oratiónis præsidio evitavit. Susceptum ex Ghisella Bavárica, sancti Henrici imperatóris sorore, quam sibi matrimonio junxerat, Emerícum fílium tanta morum disciplína talíque pietáte enutrívit, quantum ejus póstea sanctitas declaravit. Regni vero negotia ita dispósuit, ut, accitis undique prudentíssimis et sanctíssimis viris, nihil umquam sine illórum consílio molirétur; humillimis interim precibus in cínere et cilício Deum déprecans, ut univérsum Hungariæ regnum, antequam e vita migraret, catholicum vidére mererétur; vere, propter ingens dilatandæ fidei studium, illíus gentis Apóstolus nuncupatus, facta a Romano Pontifice ipsi posterisque régibus præferendæ crucis potestate.\par
\switchcolumn
\EN
\noindent He passed almost whole nights in earnest prayer, and when totally rapt in the thought of heavenly things, he sometimes became beside himself, and was seen to rise off the ground into the air. In more than one instance he strangely escaped through the power of prayer from rebellion, treason, and the onslaughts of mighty foes. He married Gisela of Bavaria, sister to the holy Emperor Henry, and begat on her Emeric, whom he trained up in such manners and godliness, as are shown by his also becoming a Saint. To carry on the business of his kingdom, he gathered together from all quarters the most learned and godly men, and took nothing in hand without their advice. Meanwhile he entreated of God by the most lowly supplications, offered up in sack-cloth and ashes, that, before he departed this life, he might see all Hungary Catholic. On account of his excellent zeal for the spread of the Faith he is called the Apostle of that nation, and the Bishop of Rome gave to him and to his successors the right to have a Cross carried before them.\par
\switchcolumn*

\LA
\begin{Resp}\Rsign Amávit eum Dóminus, et ornávit eum: stolam glóriæ índuit eum, \Star{} Et ad portas paradísi coronávit eum.\end{Resp}
\begin{Resp}\Vsign Índuit eum Dóminus lorícam fídei, et ornávit eum.\end{Resp}
\begin{Resp}\Rsign Et ad portas paradísi coronávit eum.\end{Resp}
\switchcolumn
\EN
\begin{Resp}\Rsign The Lord loved him and beautified him; He clothed him with a robe of glory, \Star{} And crowned him at the gates of Paradise.\end{Resp}
\begin{Resp}\Vsign The Lord hath put on him the breast-plate of faith, and hath adorned him.\end{Resp}
\begin{Resp}\Rsign And crowned him at the gates of Paradise.\end{Resp}
\switchcolumn*

\LA
\LessonHead{Lectio VI}
\switchcolumn
\EN
\LessonHead{Lesson VI}
\switchcolumn*

\LA
\noindent Dei Genitricem, quam ardentíssime venerabátur, amplíssimo in ejus honórem constructo templo, Hungariæ patrónam instítuit; ab eádem vicíssim Vírgine recéptus in cælum ipso suæ Assumptiónis die, quem Húngari e sancti regis instituto Magnæ Dominæ diem appellant. Sacrum ejus corpus, suavíssimo fragrans odore, liquore cælésti scatens, inter multa et varia miracula, Romani Pontificis jussu, nobiliórem in locum translátum est atque honorificentius cónditum. Ejus autem festum Innocentius undecimus Pontifex maximus, quarto Nonas Septembris, ob insignem victoriam ab exercitu Leopoldi primi, Romanórum elécti imperatóris et Hungariæ regis, eádem die in Budæ expugnatióne, ope divina, e Turcis reportatam, celebrándum instituit.\par
\switchcolumn
\EN
\noindent He had a burning zeal to honour the Mother of God. He built a very great Church in her honour, and made her Patroness of Hungary. In return, the same Virgin received him into heaven, (in the year 1038,) upon the day of her own Assumption, which the Hungarians, by the example of the holy King, call “the Great Lady's Day.” His hallowed body yielded the sweetest savour, and reeked with an heavenly liquid, and amid many and diverse wonders it was removed by command of the Bishop of Rome into a more noble place, and more honourably buried. Pope Innocent XI. ordered his Feast to be held upon the 2nd day of September, on account of the famous victory over the Turks which was gained upon this day, (in the year 1686,) when the army of Leopold I., Emperor (elect) of the Romans, and King of Hungary, wrested from them, by the help of God, the city of Buda.\par
\switchcolumn*

\LA
\begin{Resp}\Rsign Iste homo perfécit ómnia quæ locútus est ei Deus, et dixit ad eum: Ingrédere in réquiem meam: \Star{} Quia te vidi justum coram me ex ómnibus géntibus.\end{Resp}
\begin{Resp}\Vsign Iste est, qui contémpsit vitam mundi, et pervénit ad cæléstia regna.\end{Resp}
\begin{Resp}\Rsign Quia te vidi justum coram me ex ómnibus géntibus.\end{Resp}
\begin{Resp}\Vsign Glória Patri, et Fílio, \Star{} et Spirítui Sancto.\end{Resp}
\begin{Resp}\Rsign Quia te vidi justum coram me ex ómnibus géntibus.\end{Resp}
\switchcolumn
\EN
\begin{Resp}\Rsign This is he which did according unto all that God commanded him; and God said unto him: Enter thou into My rest, \Star{} For thee have I seen righteous before Me among all people.\end{Resp}
\begin{Resp}\Vsign This is he which loved not his life in this world, and is come unto an everlasting kingdom.\end{Resp}
\begin{Resp}\Rsign For thee have I seen righteous before Me among all people.\end{Resp}
\begin{Resp}\Vsign Glory be to the Father, and to the Son, \Star{} and to the Holy Ghost.\end{Resp}
\begin{Resp}\Rsign For thee have I seen righteous before Me among all people.\end{Resp}
\switchcolumn*

\LA
\LessonHead{Lectio VII}
\switchcolumn
\EN
\LessonHead{Lesson VII}
\switchcolumn*

\LA
\LessonTitle{Léctio sancti Evangélii secúndum Lucam}
\switchcolumn
\EN
\LessonTitle{From the Holy Gospel according to Luke}
\switchcolumn*

\LA
\Cite{Luc 19:12-26}
\switchcolumn
\EN
\Cite{Luke 19:12-25}
\switchcolumn*

\LA
\noindent In illo témpore: Dixit Jesus discípulis suis parábolam hanc: Homo quidam nóbilis ábiit in regiónem longinquam, accípere sibi regnum, et revérti. Et réliqua.\par
\switchcolumn
\EN
\noindent At that time, Jesus spoke this parable unto His disciples: A certain nobleman went into a far country to receive for himself a kingdom and to return. And so on.\par
\switchcolumn*

\LA
\LessonTitle{Homilía sancti Ambrósii Epíscopi}
\switchcolumn
\EN
\LessonTitle{Homily by St. Ambrose, Bishop (of Milan.)}
\switchcolumn*

\LA
\Source{Lib. 8. in Lucam.}
\switchcolumn
\EN
\Source{Book viii. on Luke.}
\switchcolumn*

\LA
\noindent Bonus ordo, ut vocatúrus gentes, et Judǽos jussúrus intérfici, qui noluérunt regnáre supra se Christum, hanc præmítteret comparatiónem, ne dicerétur: Nihil déderat pópulo Judæórum, unde póterat mélior fíeri? ut quid ab eo qui nihil recépit, exígitur? Non medíocris ista est mna, quam supra múlier evangélica quia non invénit, lucérnam accéndit, lúmine quærit admóto, gratulátur invéntam.\par
\switchcolumn
\EN
\noindent It is well ordered that, being about to call the Gentiles, and to command the destruction of those Jews, who would not have Christ to reign over them, He should put forth first this parable; lest it should be said He had given the Jews no means of becoming better. How can they be asked to repay who have received nothing? That is not a piece of silver of little worth, which, when the woman before mentioned in this Gospel (xv. 8) hath lost, she lighteth a candle, and sweepeth the house, and searcheth diligently until she findeth it.\par
\switchcolumn*

\LA
\begin{Resp}\Rsign Iste est qui ante Deum magnas virtútes operátus est, et de omni corde suo laudávit Dóminum: \Star{} Ipse intercédat pro peccátis ómnium populórum.\end{Resp}
\begin{Resp}\Vsign Ecce homo sine queréla, verus Dei cultor, ábstinens se ab omni ópere malo, et pérmanens in innocéntia sua.\end{Resp}
\begin{Resp}\Rsign Ipse intercédat pro peccátis ómnium populórum.\end{Resp}
\switchcolumn
\EN
\begin{Resp}\Rsign This is he which wrought great wonders before God, and praised the Lord with all his heart. \Star{} May he pray for all people, that their sins may be forgiven unto them.\end{Resp}
\begin{Resp}\Vsign Behold a man without blame, a worshipper of God in truth, keeping himself clean from every evil work, and abiding still in his innocency.\end{Resp}
\begin{Resp}\Rsign May he pray for all people, that their sins may be forgiven unto them.\end{Resp}
\switchcolumn*

\LA
\LessonHead{Lectio VIII}
\switchcolumn
\EN
\LessonHead{Lesson VIII}
\switchcolumn*

\LA
\noindent Dénique ex una decem mnas álius fecit, álius quinque. Fortásse iste morália habet, quia quinque sunt córporis sensus: ille duplícia, id est, mýstica legis, et morália probitátis. Unde et Matthǽus quinque talénta, et duo talénta pósuit: in quinque taléntis ut sint morália, in duóbus utrúmque, mýsticum atque morále. Ita quod número inférius, re ubérius.\par
\switchcolumn
\EN
\noindent A single pound one gained ten and another five pounds. Perchance by him which had the five pounds is signified he which practiseth well, since the body hath five senses, and by him which had the ten, that is, double the other, he which is learned and orthodox in the deep things of doctrine, as well as upright in his practical life. Hence also in Matthew we have five talents and two talents the five talents signifying good practice, and the two talents precept and practice together. So that that which counteth as the greater number is but a fraction of the lesser number.\par
\switchcolumn*

\LA
\begin{Resp}\Rsign Sint lumbi vestri præcíncti, et lucérnæ ardéntes in mánibus vestris: \Star{} Et vos símiles homínibus exspectántibus dóminum suum, quando revertátur a núptiis.\end{Resp}
\begin{Resp}\Vsign Vigiláte ergo, quia nescítis qua hora Dóminus vester ventúrus sit.\end{Resp}
\begin{Resp}\Rsign Et vos símiles homínibus exspectántibus dóminum suum, quando revertátur a núptiis.\end{Resp}
\begin{Resp}\Vsign Glória Patri, et Fílio, \Star{} et Spirítui Sancto.\end{Resp}
\begin{Resp}\Rsign Et vos símiles homínibus exspectántibus dóminum suum, quando revertátur a núptiis.\end{Resp}
\switchcolumn
\EN
\begin{Resp}\Rsign Let your loins be girded about, and your lights burning; \Star{} And ye yourselves like unto men that wait for their lord, when he will return from the wedding.\end{Resp}
\begin{Resp}\Vsign Watch therefore, for ye know not what hour your Lord doth come.\end{Resp}
\begin{Resp}\Rsign And ye yourselves like unto men that wait for their lord, when he will return from the wedding.\end{Resp}
\begin{Resp}\Vsign Glory be to the Father, and to the Son, \Star{} and to the Holy Ghost.\end{Resp}
\begin{Resp}\Rsign And ye yourselves like unto men that wait for their lord, when he will return from the wedding.\end{Resp}
\switchcolumn*

\LA
\LessonHead{Lectio IX}
\switchcolumn
\EN
\LessonHead{Lesson IX}
\switchcolumn*

\LA
\noindent Et hic póssumus decem mnas decem verba intellígere, id est, legis doctrínam; quinque autem mnas, magistéria disciplínæ. Sed legisperítum in ómnibus volo esse perféctum; non enim in sermóne, sed in virtúte est regnum Dei. Bene autem, quia de Judǽis dicit, duo soli multiplicátam pecúniam déferunt, non útique æris, sed dispensatiónis usúris. Ália est enim pecúniæ fœnébris, ália doctrínæ cæléstis usúra.\par
\switchcolumn
\EN
\noindent And here we may also understand by the ten pounds the ten words, that is, the Commandments, and by the five pounds, the enforcement of their teaching. But I would that a lawyer should be in all things perfect. For the kingdom of God is not in word but in power. (1 Cor. iv. 20.) Meet also is it, that, in speaking of Jews, Christ should represent only two as bringing in increased capital, for these talents are talents not of money but of grace, and to increase money by usury is a very different thing from improving heavenly revelation by the like means.\par
\switchcolumn*

\LA
\TeDeum{Te Deum}
\switchcolumn
\EN
\TeDeum{Te Deum}
\switchcolumn*

\end{paracol}

\par\needspace{10\baselineskip}\addvspace{1.2em}{\centering\small\textsc{Die 3 Septembris} · \textsc{3 September}\par}
\Office{S. Pii X Papæ Confessoris}{St. Pius X, Pope and Confessor}{Duplex}
\addcontentsline{toc}{subsection}{Die 3 Septembris · S. Pii X Papæ Confessoris \textit{— St. Pius X, Pope and Confessor}}
\begin{paracol}{2}
\LA
\RefLine{Lectiones I–III: de Scriptura occurrente.}
\switchcolumn
\EN
\RefLine{Lessons I–III: from the Scripture of the day.}
\switchcolumn*

\LA
\LessonHead{Lectio IV}
\switchcolumn
\EN
\LessonHead{Lesson IV}
\switchcolumn*

\LA
\noindent Pius Papa décimus, cui nomen ántea Joséphus Sarto, in Venetórum pago natus est, quem Riése vocant, paréntibus quidem humílibus, sed probitáte ac pietáte conspícuus. Inter Seminárii Patavíni alúmnos adscríptus, ita pietáte ac doctrína profécit, ut condiscípulis exémplo, moderatóribus admiratióni esset. Sacerdótio initiátus, in óppido Tómbolo primum, qua vicárius cooperátor, dein Salatiáni qua párochus, per plures annos adlaborávit; quibus in obeúndis munéribus, tanta caritátis effusióne, tanto sacerdotáli zelo et sanctitáte vitæ excélluit, ut Epíscopus Tarvisínus inter canónicos cathedrális ecclésiæ eum cooptáret, eúmque Cúriæ episcopális cancellárium simúlque Seminárii diœcesáni spirituálem moderatórem renuntiáret. Hæc offícia tam egrégie persecútus, a Leóne tértio décimo, cui erat probatíssimus, Mantuánæ ecclésiæ Antístes fuit renuntiátus.\par
\switchcolumn
\EN
\noindent Pope Pius X, whose name previously was Joseph Sarto, was born in the village of Riese in the Venetian province, to humble parents remarkable for their godliness and piety. He enrolled among the students in the seminary of Padua, where he exhibited such piety and learning that he was both an example to his fellow students and the admiration of his teachers. Upon his ordination to the priesthood, he labored for several years first as curate in the town of Tombolo, then as pastor at Salzano. He applied himself to his duties with such a constant flow of charity and such priestly zeal, and was so distinguished by the holiness of his life, that the Bishop of Treviso appointed him as a canon of the cathedral church and made him the chancellor of the bishop's curia, as well as spiritual director of the diocesan seminary. His performance in these duties was so outstanding and so highly impressed Leo XIII, that he made him bishop of the Church of Mantua.\par
\switchcolumn*

\LA
\begin{Resp}\Rsign Invéni David servum meum, óleo sancto meo unxi eum: \Star{} Manus enim mea auxiliábitur ei.\end{Resp}
\begin{Resp}\Vsign Nihil profíciet inimícus in eo, et fílius iniquitátis non nocébit ei.\end{Resp}
\begin{Resp}\Rsign Manus enim mea auxiliábitur ei.\end{Resp}
\switchcolumn
\EN
\begin{Resp}\Rsign I have found David My servant, with My holy oil have I anointed him \Star{} For My hand shall help him.\end{Resp}
\begin{Resp}\Vsign The enemy shall prevail nothing against him, nor the son of wickedness afflict him.\end{Resp}
\begin{Resp}\Rsign For My hand shall help him.\end{Resp}
\switchcolumn*

\LA
\LessonHead{Lectio V}
\switchcolumn
\EN
\LessonHead{Lesson V}
\switchcolumn*

\LA
\noindent Boni pastóris nullam partem déserens, eo máxime conténdit, ut juvéntus in sortem Dómini vocáta rite ad sacra instituerétur, piæ consociatiónes novis augéscerent increméntis, rítibus divíni cultus plus decóris ac pietátis accéderet. Præcépta quibus cívitas christiána nítitur, áltius proclamáre non désiit, et qui vitam ínopem ipse ducébat, paupéribus numquam omísit afférre levámen. Tot ígitur suffragántibus méritis, inter purpurátos Patres adléctus et Venetiárum Patriárcha creátus est. Dénique post Leónis décimi tértii óbitum, cum Patrum Cardinálium suffrágia in eum coaléscerent, cumque ipse supplicatiónibus et lácrimis tantum munus a se avértere frustra conátus esset, suasiónibus tandem cedens, « accépto in crucem », inquit, et Summi Pontificátus ápicem ut crucem a Deo sibi oblátam, demísso sed forti ánimo suscépit.\par
\switchcolumn
\EN
\noindent Lacking in nothing that makes a good pastor, he labored particularly to teach young men called to the priesthood, as well as fostering the growth of devout associations and the beauty and dignity of divine worship. He would ever affirm and promote the laws upon which Christian civilization depend, and while leading himself a life of poverty, never missed the opportunity to alleviate the burden of poverty in others. Because of his great merits, he was made a cardinal and created Patriarch of Venice. After the death of Pope Leo XIII, when the votes of the College of Cardinals began to increase in his favor, he tried in vain with supplications and tears to be relieved of so heavy a burden. Finally he ceded to their persuasions, saying I accept the cross. Thus he accepted the crown of the supreme pontificate as a cross, offering himself to God, with a resigned but stedfast spirit.\par
\switchcolumn*

\LA
\begin{Resp}\Rsign Pósui adjutórium super poténtem, et exaltávi eléctum de plebe mea: \Star{} Manus enim mea auxiliábitur ei.\end{Resp}
\begin{Resp}\Vsign Invéni David servum meum, óleo sancto meo unxi eum.\end{Resp}
\begin{Resp}\Rsign Manus enim mea auxiliábitur ei.\end{Resp}
\switchcolumn
\EN
\begin{Resp}\Rsign I have laid help upon one that is mighty, and have exalted one chosen out of My people \Star{} For My hand shall help him.\end{Resp}
\begin{Resp}\Vsign I have found David My servant, with My holy oil have I anointed him.\end{Resp}
\begin{Resp}\Rsign For My hand shall help him.\end{Resp}
\switchcolumn*

\LA
\LessonHead{Lectio VI}
\switchcolumn
\EN
\LessonHead{Lesson VI}
\switchcolumn*

\LA
\noindent In Petri cáthedra constitútus, nihil de prístina vitæ ratióne remísit. Humilitáte præsértim, simplicitáte ac paupertáte refúlsit, ita ut in suo testaménto scríbere potúerit: « Pauper natus sum, pauper vixi, pauper mori cúpio ». Humílitas vero ánimi fortitúdinem in eo alébat, cum de Dei glória, Ecclésiæ libertáte, animarúmque salúte agerétur. Vir acérrimi ingénii et propósiti tenax, inter vicésimi ineúntis sæculi procéllas, Ecclésiam fírmiter rexit, et præclaríssimis ornávit institútis. Músicam sacram ad prístinum splendórem ac dignitátem revocávit; sacrórum Bibliórum stúdiis príncipem sedem Romæ constítuit; Románam Cúriam sapiénter reformávit; leges de fidélibus per catechísmum instituéndis restítuit; Eucharísticæ mensæ crebriórem, imo et cotidiánam consuetúdinem indúxit, ejúsque accéssum púeris quoque a primo ratiónis usu apéruit; actiónis cathólicæ increménta sédulo promóvit; sólidæ cleri institutióni provídit, ádditis quoque semináriis per regiónes dispósitis; sacerdótes omnes ad interiórem vitam coléndam alléxit; leges Ecclésiæ in unum corpus redégit; erróres perniciosíssimos, modernísmi appellatióne comprehénsos, damnávit atque evéllit; civíle vétitum, quod dicunt, in Pontíficis Máximi electióne rejécit. Tandem labóribus fractus ac mæróre conféctus ob bellum Europæum tunc exórtum, die vicésima mensis Augústi anni millésimi nongentésimi décimi quarti, ad cæléste præmium evolávit. Eum ubíque terrárum sanctitátis fama clarum miraculísque fulgéntem, Pius Papa duodécimus, cuncto plaudénte orbe, in Sanctórum númerum rétulit.\par
\switchcolumn
\EN
\noindent Placed upon the chair of Peter, he gave up nothing of his former way of life. He shone especially in humility, simplicity and poverty, so that he was able to write in his last testament: I was born in poverty, I lived in poverty, and I wish to die in poverty. His humility, however, nourished his soul with strength, when it concerned the glory of God, the liberty of Holy Church, and the salvation of souls. A man of passionate temperament and of firm purpose, he ruled the Church firmly as it entered into the twentieth century, and adorned it with brilliant teachings. He restored the sacred music to its pristine glory and dignity; he established Rome as the principal centre for the study of the Holy Bible; he ordered the reform of the Roman Curia with great wisdom; he restored the laws concerning the faithful for the instruction of the catechism; he introduced the custom of more frequent and even daily reception of the Holy Eucharist, as well as permitting its reception by children as soon as they reach the age of reason; he zealously promoted the growth of Catholic action; he provided for the sound education of clerics and increased the number of seminaries in their diverse regions; he encouraged every priest in the practice of the interior life; he brought the laws of the Church together into one body; he condemned and suppressed those most pernicious errors known collectively as Modernism; he suppressed the custom of civil veto at the election of a Supreme Pontiff. Finally worn out with his labors and overcome with grief at the European war which had just begun, he went to his heavenly reward on the twentieth day of August in the year 1914. Renowned throughout all the world for the fame of his holiness and miracles, Pope Pius XII, with the approbation of the whole world, numbered him among the Saints.\par
\switchcolumn*

\LA
\begin{Resp}\Rsign Iste est, qui ante Deum magnas virtútes operátus est, et omnis terra doctrína ejus repléta est: \Star{} Ipse intercédat pro peccátis ómnium populórum.\end{Resp}
\begin{Resp}\Vsign Iste est, qui contémpsit vitam mundi, et pervénit ad cæléstia regna.\end{Resp}
\begin{Resp}\Rsign Ipse intercédat pro peccátis ómnium populórum.\end{Resp}
\begin{Resp}\Vsign Glória Patri, et Fílio, \Star{} et Spirítui Sancto.\end{Resp}
\begin{Resp}\Rsign Ipse intercédat pro peccátis ómnium populórum.\end{Resp}
\switchcolumn
\EN
\begin{Resp}\Rsign This is he which wrought great wonders before God, and the whole earth is full of his teaching \Star{} May he pray for all people, that their sins may be forgiven unto them\end{Resp}
\begin{Resp}\Vsign This is he which loved not his life in this world, and hath attained unto the kingdom of heaven.\end{Resp}
\begin{Resp}\Rsign May he pray for all people, that their sins may be forgiven unto them\end{Resp}
\begin{Resp}\Vsign Glory be to the Father, and to the Son, \Star{} and to the Holy Ghost.\end{Resp}
\begin{Resp}\Rsign May he pray for all people, that their sins may be forgiven unto them\end{Resp}
\switchcolumn*

\LA
\LessonHead{Lectio VII}
\switchcolumn
\EN
\LessonHead{Lesson VII}
\switchcolumn*

\LA
\LessonTitle{Léctio sancti Evangélii secúndum Joánnem}
\switchcolumn
\EN
\LessonTitle{From the Holy Gospel according to John}
\switchcolumn*

\LA
\Cite{Joannes 21:15-17}
\switchcolumn
\EN
\Cite{John 21:15-17}
\switchcolumn*

\LA
\noindent In illo témpore: Dixit Jesus Simóni Petro: Simon Joánnis, díligis me plus his? Et réliqua.\par
\switchcolumn
\EN
\noindent At that time: Jesus saith to Simon Peter, Simon, son of Jonas, lovest thou me more than these? And so on, and that which followeth.\par
\switchcolumn*

\LA
\LessonTitle{Homilía sancti Augustíni Epíscopi}
\switchcolumn
\EN
\LessonTitle{A Homily by St. Augustine, the Bishop}
\switchcolumn*

\LA
\Source{Tractatus 123 in Joánnem. in medio}
\switchcolumn
\EN
\Source{Tractatus 123 in Joannem. in medio}
\switchcolumn*

\LA
\noindent Rédditur negatióni trinæ trina conféssio, ne minus amóri lingua sérviat, quam timóri: et plus vocis elicuísse videátur mors ímminens, quam vita præsens. Sit amóris offícium páscere Domínicum gregem, si fuit timóris indícium negáre pastórem. Qui hoc ánimo pascunt oves Christi, ut suas velint esse, non Christi, se convincúntur amáre, non Christum: vel gloriándi vel dominándi, vel acquiréndi cupiditáte; non obediéndi, et subveniéndi, et Deo placéndi caritáte.\par
\switchcolumn
\EN
\noindent To the threefold denial there is now appended a threefold confession, that his tongue may not yield a feebler service to love than to fear, and imminent death may not appear to have elicited more from his lips than present life. Let it be the office of love to feed the Lord's flock, if it was the signal of fear to deny the Shepherd. Those who have this purpose in feeding the flock of Christ, that they may have them as their own, and not as Christ's, are convicted of loving themselves, and not Christ, from the desire either of boasting, or wielding power, or acquiring gain, and not from the love of obeying, serving and pleasing God.\par
\switchcolumn*

\LA
\begin{Resp}\Rsign Amávit eum Dóminus, et ornávit eum: stolam glóriæ índuit eum, \Star{} Et ad portas paradísi coronávit eum.\end{Resp}
\begin{Resp}\Vsign Induit eum Dóminus lorícam fídei, et ornávit eum.\end{Resp}
\begin{Resp}\Rsign Et ad portas paradísi coronávit eum.\end{Resp}
\switchcolumn
\EN
\begin{Resp}\Rsign The Lord loved him and beautified him He clothed him with a robe of glory \Star{} And crowned him at the gates of Paradise.\end{Resp}
\begin{Resp}\Vsign The Lord hath put on him the breast-plate of faith, and hath adorned him.\end{Resp}
\begin{Resp}\Rsign And crowned him at the gates of Paradise.\end{Resp}
\switchcolumn*

\LA
\LessonHead{Lectio VIII}
\switchcolumn
\EN
\LessonHead{Lesson VIII}
\switchcolumn*

\LA
\noindent Contra hos ergo vígilat tótidem inculcáta ista vox Christi, quos Apóstolos gemit sua quærere, non quæ Jesu Christi. Nam quid est áliud, si díligis me, pasce oves meas; quam si dicerétur, si me díligis, non te páscere cógita? Sed oves meas, et sicut meas pasce, non sicut tuas; glóriam meam in eis quære, non tuam; domínium meum, non tuum; lucra mea, non tua; ne sis in eórum societáte qui pértinent ad témpora periculósa, seípsos amántes, et cétera quæ huic malórum inítio connectúntur.\par
\switchcolumn
\EN
\noindent Against such, therefore, there standeth as a wakeful sentinel this thrice inculcated utterance of Christ, of whom the Apostle complaineth that they seek their own, not the things that are of Christ Jesus. For what else signifieth the words: Lovest thou me? Feed my sheep: than if it were said: If thou lovest me, think not of feeding thyself, but feed my sheep as mine, and not as thine own; seek my glory in them, and not thine own; my dominion, and not thine; my gain, and not thine; lest thou be found in the fellowship of them that belong to the perilous times, lovers of their own selves, and all else that is joined on to this beginning of evils.\par
\switchcolumn*

\LA
\begin{Resp}\Rsign Sint lumbi vestri præcíncti, et lucérnæ ardéntes in mánibus vestris: \Star{} Et vos símiles homínibus exspectántibus dóminum suum, quando revertátur a núptiis.\end{Resp}
\begin{Resp}\Vsign Vigiláte ergo, quia nescítis qua hora Dóminus vester ventúrus sit.\end{Resp}
\begin{Resp}\Rsign Et vos símiles homínibus exspectántibus dóminum suum, quando revertátur a núptiis.\end{Resp}
\begin{Resp}\Vsign Glória Patri, et Fílio, \Star{} et Spirítui Sancto.\end{Resp}
\begin{Resp}\Rsign Et vos símiles homínibus exspectántibus dóminum suum, quando revertátur a núptiis.\end{Resp}
\switchcolumn
\EN
\begin{Resp}\Rsign Let your loins be girded about, and your lights burning; \Star{} And ye yourselves like unto men that wait for their lord, when he will return from the wedding.\end{Resp}
\begin{Resp}\Vsign Watch therefore, for ye know not what hour your Lord doth come.\end{Resp}
\begin{Resp}\Rsign And ye yourselves like unto men that wait for their lord, when he will return from the wedding.\end{Resp}
\begin{Resp}\Vsign Glory be to the Father, and to the Son, \Star{} and to the Holy Ghost.\end{Resp}
\begin{Resp}\Rsign And ye yourselves like unto men that wait for their lord, when he will return from the wedding.\end{Resp}
\switchcolumn*

\LA
\LessonHead{Lectio IX}
\switchcolumn
\EN
\LessonHead{Lesson IX}
\switchcolumn*

\LA
\noindent Mérito dícitur Petro: díligis me, et respóndet, amo te, eíque refértur, pasce agnos meos, et hoc íterum, hoc tértio. Ubi étiam demonstrátur unum atque idem esse amórem et dilectiónem; nam étiam Dóminus novíssime non ait, díligis me, sed amas me. Non ergo nos, sed ipsum amémus, et in pascéndis óvibus ejus ea quæ sunt ejus, non ea quæ sunt nostra, quærámus. Néscio enim quo inexplicábili modo, quisquis seípsum, non Deum amat, non se amat: et quisquis Deum, non seípsum, amat, ipse se amat. Qui enim non potest vívere de se, móritur útique amándo se: non ergo se amat, qui ne vivat se amat.\par
\switchcolumn
\EN
\noindent With great propriety, therefore, is Peter addressed: Lovest thou me? and found replying: I love thee; and the command applied to him: Feed my lambs, and this a second and a third time. We have it also demonstrated here that love (amor) and liking (dilectio) are one and the same thing; for the Lord also in the last question said not: Dost thou like me? (Diligis me?): but: Dost thou love me? (Amas me?) Let us, then, love not ourselves but him; and in feeding his sheep, let us be seeking the things which are his, not the things which are our own. For in some inexplicable way, I know not what, every one that loveth himself, and not God, loveth not himself; and whoever loveth God, and not himself, he it is that loveth himself. For he that cannot live by himself will certainly die by loving himself; he therefore loveth not himself that loveth himself to his own loss of life.\par
\switchcolumn*

\LA
\TeDeum{Te Deum}
\switchcolumn
\EN
\TeDeum{Te Deum}
\switchcolumn*

\end{paracol}

\PartTitle{Commune Sanctorum}{Common of the Saints}

\addcontentsline{toc}{chapter}{Commune Sanctorum \textit{— Common of the Saints}}

\Office{Commune Apostolorum}{Commune Apostolorum}{}
\addcontentsline{toc}{subsection}{Commune Apostolorum \textit{— Commune Apostolorum}}
\begin{paracol}{2}
\LA
\LessonHead{Lectio I}
\switchcolumn
\EN
\LessonHead{Lesson I}
\switchcolumn*

\LA
\LessonTitle{De Epístola prima beáti Pauli Apóstoli ad Corínthios}
\switchcolumn
\EN
\LessonTitle{Lesson from the first letter of St. Paul the Apostle to the Corinthians}
\switchcolumn*

\LA
\Cite{1 Cor 4:1-5}
\switchcolumn
\EN
\Cite{1 Cor 4:1-5}
\switchcolumn*

\LA
\noindent Sic nos exístimet homo ut minístros Christi, et dispensatóres mysteriórum Dei. Hic jam quǽritur inter dispensatóres ut fidélis quis inveniátur. Mihi autem pro mínimo est ut a vobis júdicer, aut ab humáno die: sed neque meípsum júdico. Nihil enim mihi cónscius sum, sed non in hoc justificátus sum: qui autem júdicat me, Dóminus est. Ítaque nolíte ante tempus judicáre, quoadúsque véniat Dóminus: qui et illuminábit abscóndita tenebrárum, et manifestábit consília córdium: et tunc laus erit unicuíque a Deo.\par
\switchcolumn
\EN
\noindent Let a man so account of us as of the ministers of Christ, and the dispensers of the mysteries of God. Here now it is required among the dispensers, that a man be found faithful. But to me it is a very small thing to be judged by you, or by man's day; but neither do I judge my own self. For I am not conscious to myself of any thing, yet am I not hereby justified; but he that judgeth me, is the Lord. Therefore judge not before the time; until the Lord come, who both will bring to light the hidden things of darkness, and will make manifest the counsels of the hearts; and then shall every man have praise from God.\par
\switchcolumn*

\LA
\begin{Resp}\Rsign Ecce ego mitto vos sicut oves in médio lupórum, dicit Dóminus: \Star{} Estóte ergo prudéntes sicut serpéntes, et símplices sicut colúmbæ.\end{Resp}
\begin{Resp}\Vsign Dum lucem habétis, crédite in lucem, ut fílii lucis sitis.\end{Resp}
\begin{Resp}\Rsign Estóte ergo prudéntes sicut serpéntes, et símplices sicut colúmbæ.\end{Resp}
\switchcolumn
\EN
\begin{Resp}\Rsign Behold, I send you as sheep in the midst of wolves, said the Lord: \Star{} Be ye, therefore, wise as serpents and simple as doves.\end{Resp}
\begin{Resp}\Vsign Whilst you have the light, believe in the light, that you may be the children of light.\end{Resp}
\begin{Resp}\Rsign Be ye therefore wise as serpents and simple as doves.\end{Resp}
\switchcolumn*

\LA
\LessonHead{Lectio II}
\switchcolumn
\EN
\LessonHead{Lesson II}
\switchcolumn*

\LA
\Cite{1 Cor 4:6-9}
\switchcolumn
\EN
\Cite{1 Cor 4:6-9}
\switchcolumn*

\LA
\noindent Hæc autem, fratres, transfigurávi in me et Apóllo, propter vos: ut in nobis discátis, ne supra quam scriptum est, unus advérsus álterum inflétur pro álio. Quis enim te discérnit? Quid autem habes quod non accepísti? Si autem accepísti, quid gloriáris quasi non accéperis? Jam saturáti estis, jam dívites facti estis: sine nobis regnátis: et útinam regnétis, ut et nos vobíscum regnémus. Puto enim quod Deus nos Apóstolos novíssimos osténdit, tamquam morti destinátos: quia spectáculum facti sumus mundo, et Ángelis, et homínibus.\par
\switchcolumn
\EN
\noindent But these things, brethren, I have in a figure transferred to myself and to Apollo, for your sakes; that in us you may learn, that one be not puffed up against the other for another, above that which is written. For who distinguisheth thee? Or what hast thou that thou hast not received? And if thou hast received, why dost thou glory, as if thou hadst not received it? You are now full; you are now become rich; you reign without us; and I would to God you did reign, that we also might reign with you. For I think that God hath set forth us apostles, the last, as it were men appointed to death: we are made a spectacle to the world, and to angels, and to men.\par
\switchcolumn*

\LA
\begin{Resp}\Rsign Tóllite jugum meum super vos, dicit Dóminus, et díscite a me, quia mitis sum et húmilis corde: \Star{} Jugum enim meum suáve est, et onus meum leve.\end{Resp}
\begin{Resp}\Vsign Et inveniétis réquiem animábus vestris.\end{Resp}
\begin{Resp}\Rsign Jugum enim meum suáve est, et onus meum leve.\end{Resp}
\switchcolumn
\EN
\begin{Resp}\Rsign Take up my yoke upon you, said the Lord, and learn of me, because I am meek, and humble of heart. \Star{} For my yoke is sweet and my burden light.\end{Resp}
\begin{Resp}\Vsign And you shall find rest to your souls.\end{Resp}
\begin{Resp}\Rsign For my yoke is sweet and my burden light.\end{Resp}
\switchcolumn*

\LA
\LessonHead{Lectio III}
\switchcolumn
\EN
\LessonHead{Lesson III}
\switchcolumn*

\LA
\Cite{1 Cor 4:10-15}
\switchcolumn
\EN
\Cite{1 Cor 4:10-15}
\switchcolumn*

\LA
\noindent Nos stulti propter Christum, vos autem prudéntes in Christo: nos infírmi, vos autem fortes: vos nóbiles, nos autem ignóbiles. Usque in hanc horam et esurímus, et sitímus, et nudi sumus, et cólaphis cǽdimur, et instábiles sumus, et laborámus operántes mánibus nostris: maledícimur, et benedícimus: persecutiónem pátimur, et sustinémus: blasphemámur, et obsecrámus: tamquam purgaménta hujus mundi facti sumus, ómnium peripséma usque adhuc. Non ut confúndam vos, hæc scribo, sed ut fílios meos caríssimos móneo. Nam si decem míllia pædagogórum habeátis in Christo, sed non multos patres, nam in Christo Jesu per Evangélium ego vos génui.\par
\switchcolumn
\EN
\noindent We are fools for Christ's sake, but you are wise in Christ; we are weak, but you are strong; you are honourable, but we without honour. Even unto this hour we both hunger and thirst, and are naked, and are buffeted, and have no fixed abode; And we labour, working with our own hands: we are reviled, and we bless; we are persecuted, and we suffer it. We are blasphemed, and we entreat; we are made as the refuse of this world, the offscouring of all even until now. I write not these things to confound you; but I admonish you as my dearest children. For if you have ten thousand instructors in Christ, yet not many fathers. For in Christ Jesus, by the gospel, I have begotten you.\par
\switchcolumn*

\LA
\begin{Resp}\Rsign Dum stetéritis ante reges et prǽsides, nolíte cogitáre quómodo aut quid loquámini: \Star{} Dábitur enim vobis in illa hora, quid loquámini.\end{Resp}
\begin{Resp}\Vsign Non enim vos estis qui loquímini: sed Spíritus Patris vestri, qui lóquitur in vobis.\end{Resp}
\begin{Resp}\Rsign Dábitur enim vobis in illa hora, quid loquámini.\end{Resp}
\begin{Resp}\Vsign Glória Patri, et Fílio, \Star{} et Spirítui Sancto.\end{Resp}
\begin{Resp}\Rsign Dábitur enim vobis in illa hora, quid loquámini.\end{Resp}
\switchcolumn
\EN
\begin{Resp}\Rsign But when they shall deliver you up to the judges, take no thought how or what to speak: \Star{} for it shall be given you in that hour what to speak:\end{Resp}
\begin{Resp}\Vsign For it is not you that speak, but the spirit of your Father that speaketh in you.\end{Resp}
\begin{Resp}\Rsign For it shall be given you in that hour what to speak.\end{Resp}
\begin{Resp}\Vsign Glory be to the Father, and to the Son, \Star{} and to the Holy Ghost.\end{Resp}
\begin{Resp}\Rsign For it shall be given you in that hour what to speak.\end{Resp}
\switchcolumn*

\end{paracol}

\Office{Commune unius Martyris Tempore Paschali}{Common of a Single Martyr in Paschaltide}{}
\addcontentsline{toc}{subsection}{Commune unius Martyris Tempore Paschali \textit{— Common of a Single Martyr in Paschaltide}}
\begin{paracol}{2}
\LA
\LessonHead{Lectio VII}
\switchcolumn
\EN
\LessonHead{Lesson VII}
\switchcolumn*

\LA
\LessonTitle{Léctio sancti Evangélii secúndum Joánnem}
\switchcolumn
\EN
\LessonTitle{From the holy Gospel according to John}
\switchcolumn*

\LA
\Cite{Joannes 15:1-7}
\switchcolumn
\EN
\Cite{John 15:1-7}
\switchcolumn*

\LA
\noindent In illo témpore: Dixit Jesus discípulis suis: Ego sum vitis vera: et Pater meus agrícola est. Et réliqua.\par
\switchcolumn
\EN
\noindent In that time Jesus said to his disciples: I an the true vine; and my Father is the husbandman. And so on.\par
\switchcolumn*

\LA
\LessonTitle{Homilía sancti Augustíni Epíscopi}
\switchcolumn
\EN
\LessonTitle{Homily by St. Augustine, Bishop (of Hippo.)}
\switchcolumn*

\LA
\Source{Tract. 80 in Joannem}
\switchcolumn
\EN
\Source{Tract 80 on John}
\switchcolumn*

\LA
\noindent Iste locus evangélicus, fratres, ubi se dicit Dóminus vitem, et discípulos suos pálmites, secúndum hoc dicit, quod est caput Ecclésiæ, nosque membra ejus, mediátor Dei et hóminum, homo Christus Jesus. Uníus quippe natúræ sunt vitis et pálmites. Propter quod cum esset Deus, cujus natúræ non sumus, factus est homo, ut in illo esset vitis humána natúra, cujus et nos hómines pálmites esse possémus.\par
\switchcolumn
\EN
\noindent Dearly beloved brethren, this passage of the Gospel, wherein the Lord saith that He is the vine, and that His disciples are the branches, is to be taken in that sense wherein it is also said, that He is the Head of the Church, (Eph. v. 23,) and that we are the members of Him (30) Who is the Mediator between God and men, the man Christ Jesus (1 Tim. ii. 5.) The vine and his branches are of one and the same nature. Therefore, seeing that He was God, of which nature we are not, He was made man, to the end that He might have in Himself this vine, that is, the manhood, whereof we men can be made branches.\par
\switchcolumn*

\LA
\begin{Resp}\Rsign Ego sum vitis vera, et vos pálmites: \Star{} Qui manet in me, et ego in eo, hic fert fructum multum, allelúja, allelúja.\end{Resp}
\begin{Resp}\Vsign Sicut diléxit me Pater, et ego diléxi vos.\end{Resp}
\begin{Resp}\Rsign Qui manet in me, et ego in eo, hic fert fructum multum, allelúja, allelúja.\end{Resp}
\switchcolumn
\EN
\begin{Resp}\Rsign I am the vine: you the branches. \Star{} He that abideth in me, and I in him, the same beareth much fruit, alleluia, alleluia.\end{Resp}
\begin{Resp}\Vsign As the Father hath loved me, I also have loved you.\end{Resp}
\begin{Resp}\Rsign He that abideth in me, and I in him, the same beareth much fruit, alleluia, alleluia.\end{Resp}
\begin{Resp}\Vsign Glory be to the Father, and to the Son, \Star{} and to the Holy Ghost.\end{Resp}
\begin{Resp}\Rsign He that abideth in me, and I in him, the same beareth much fruit, alleluia, alleluia.\end{Resp}
\switchcolumn*

\LA
\LessonHead{Lectio VIII}
\switchcolumn
\EN
\LessonHead{Lesson VIII}
\switchcolumn*

\LA
\noindent Quid ergo est, Ego sum vitis vera? Numquid ut ádderet, vera, hoc ad eam vitem rétulit, unde ista similitúdo transláta est? Sic enim dícitur vitis per similitúdinem, non per proprietátem: quemádmodum dícitur ovis, agnus, leo, petra, lapis anguláris, et cétera hujúsmodi, quæ magis ipsa sunt vera, ex quibus ducúntur istæ similitúdines, non proprietátes. Sed cum dicit, Ego sum vitis vera: ab illa se útique discérnit, cui dícitur: Quómodo convérsa es in amaritúdinem vitis aliéna? Nam quo pacto est vitis vera, quæ exspectáta est ut fáceret uvam, fecit autem spinas?\par
\switchcolumn
\EN
\noindent Why saith He: I am the true vine? As touching this word true, hath He not here regard to that other parable of a vine, the like figure whereto He doth here apply to Himself? (Jer. ii. 21.) Here is He called a vine, not plainly, but in parable, as also He is called elsewhere a sheep, (Isa. liii. 7, Acts viii. 32,) a lamb, (John i. 36,) a lion, (Apoc. v. 5,) a rock, (1 Cor. x. 4,) a corner-stone, \textasciitilde{}(Eph. ii. 20,) and other things of the like kind. But these things are in themselves that which they seem to be, albeit He is called by their names, not plainly, but in a parable, and herein are they different from that vine, whereof in this place He taketh on Him the name. For when He saith: I am the true vine, doth He not make distinction between Himself, and that which indeed seemed to be a vine, but to which it is said: How art thou turned into the degenerate plant of a strange vine unto Me? (Jer. ii. 21.) For by what title shall that plant be called other than a false vine, whereto they looked that she should bring forth grapes, and she brought forth thorns?\par
\switchcolumn*

\LA
\begin{Resp}\Rsign Cándidi facti sunt Nazarǽi ejus, allelúja: splendórem Deo dedérunt, allelúja: \Star{} Et sicut lac coaguláti sunt, allelúja, allelúja.\end{Resp}
\begin{Resp}\Vsign Candidióres nive, nitidióres lacte, rubicundióres ébore antíquo, sapphíro pulchrióres.\end{Resp}
\begin{Resp}\Rsign Et sicut lac coaguláti sunt, allelúja, allelúja.\end{Resp}
\begin{Resp}\Vsign Glória Patri, et Fílio, \Star{} et Spirítui Sancto.\end{Resp}
\begin{Resp}\Rsign Et sicut lac coaguláti sunt, allelúja, allelúja.\end{Resp}
\switchcolumn
\EN
\begin{Resp}\Rsign Her Nazarites are become pure, Alleluia: they reflect the glory of God, Alleluia. \Star{} They are whiter than milk. Alleluia, Alleluia.\end{Resp}
\begin{Resp}\Vsign They are purer than snow, they are whiter than milk, they are more ruddy in body than coral, their polishing is of sapphire.\end{Resp}
\begin{Resp}\Rsign They are whiter than milk. Alleluia, Alleluia.\end{Resp}
\begin{Resp}\Vsign Glory be to the Father, and to the Son, \Star{} and to the Holy Ghost.\end{Resp}
\begin{Resp}\Rsign They are whiter than milk. Alleluia, Alleluia.\end{Resp}
\switchcolumn*

\LA
\LessonHead{Lectio IX}
\switchcolumn
\EN
\LessonHead{Lesson IX}
\switchcolumn*

\LA
\noindent Ego sum, inquit, vitis vera: et Pater meus agrícola est. Numquid unum sunt agrícola et vitis? Secúndum hoc ergo vitis Christus, secúndum quod ait: Pater major me est. Secúndum autem id, quod ait: Ego, et Pater unum sumus; et ipse agrícola est: nec talis, quales sunt, qui extrínsecus operándo éxhibent ministérium: sed talis, ut det étiam intrínsecus increméntum. Nam neque qui plantat est áliquid, neque qui rigat: sed, qui increméntum dat, Deus. Sed útique Deus est Christus, quia Deus erat Verbum: unde ipse, et Pater unum sunt. Et, si Verbum caro factum est, quod non erat, manet quod erat.\par
\switchcolumn
\EN
\noindent He saith: I am the true vine, and My Father is the husbandman. Is the vine one with the husbandman? These words then are to be taken in that sense wherein He also saith: My Father is greater than I. (John xiv. 28.) In this sense is He the vine, and the Father is the husbandman. But again, in regard to those words: I and the Father are one, and again and: My Father is the husbandman, we understand that They are not the vine and the husbandman, after the manner of a vine, and the husbandman that from without doth care for and keep it, but after the manner of a vine and Him That from within doth make it to bring forth fruit. For neither is he that planteth anything, neither he that watereth, but God that giveth the increase. (1 Cor. iii. 7.) But Christ is God, for the Word was God. (John i. 1.) Therefore He and the Father are one: and, albeit the Word was made flesh, (John i. 14,) which, before, He was not, He ceased not to be still That Which He was.\par
\switchcolumn*

\LA
\TeDeum{Te Deum}
\switchcolumn
\EN
\TeDeum{Te Deum}
\switchcolumn*

\end{paracol}

\Office{Commune unius Summi Pontificis et Martyris}{Common of a Single Martyr Pope}{}
\addcontentsline{toc}{subsection}{Commune unius Summi Pontificis et Martyris \textit{— Common of a Single Martyr Pope}}
\begin{paracol}{2}
\LA
\LessonHead{Lectio VII}
\switchcolumn
\EN
\LessonHead{Lesson VII}
\switchcolumn*

\LA
\LessonTitle{Léctio sancti Evangélii secúndum Matthǽum}
\switchcolumn
\EN
\LessonTitle{From the Holy Gospel according to Matthew}
\switchcolumn*

\LA
\Cite{Matt 16:13-19}
\switchcolumn
\EN
\Cite{Matt 16:13-19}
\switchcolumn*

\LA
\noindent In illo témpore: Venit Jesus in partes Cæsaréæ Philíppi, et interrogábat discípulos suos, dicens: Quem dicunt hómines esse Fílium hóminis? Et réliqua.\par
\switchcolumn
\EN
\noindent At that time: When Jesus came into the coasts of Caesarea Philippi, he asked his disciples, saying, Whom do men say that I the Son of man am? And so on.\par
\switchcolumn*

\LA
\LessonTitle{Homilía sancti Leónis Papæ}
\switchcolumn
\EN
\LessonTitle{A Homily by St. Leo the Pope}
\switchcolumn*

\LA
\Source{Sermo 2 in anniversario assumpt. suæ, ante medium}
\switchcolumn
\EN
\Source{Sermo 2 in anniversario assumpt. suæ ante medium}
\switchcolumn*

\LA
\noindent Cum, sicut evangélica lectióne reserátum est, interrogásset Dóminus discípulos, quem ipsum (multis divérsa opinántibus) créderent; respondissétque beátus Petrus, dicens: Tu es Christus Fílius Dei vivi; Dóminus ait: Beátus es, Simon Bar-Jona, quia caro et sanguis non revelávit tibi, sed Pater meus, qui in cælis est: et ego dico tibi, quia tu es Petrus, et super hanc petram ædificábo Ecclésiam meam, et portæ ínferi non prævalébunt advérsus eam. Et tibi dabo claves regni cælórum: et quodcúmque ligáveris super terram, erit ligátum et in cælis: et quodcúmque sólveris super terram, erit solútum et in cælis. Manet ergo disposítio veritátis, et beátus Petrus, in accépta fortitúdine petræ persevérans, suscépta Ecclésiæ gubernácula non relíquit.\par
\switchcolumn
\EN
\noindent When the Lord, as we read in the Gospel, asked his disciples who did men, amid their diverse speculations, believe him the Son of Man to be, blessed Peter answered and said: Thou art the Christ, the Son of the living God. And the Lord answered and said unto him: Blessed art thou, Simon Bar-Jona: for flesh and blood hath not revealed it unto thee, but my Father, which is in heaven: and I say also unto thee: That thou art Peter, and upon this rock I will build my Church, and the gates of hell shall not prevail against it; and I will give unto thee the keys of the kingdom of heaven; and whatsoever thou shalt bind on earth shall be bound in heaven; and whatsoever thou shalt loose on earth shall be loosed in heaven. But the dispensation of truth perdures, and blessed Peter, persevering in the strength of the rock which he hath received, hath not relinquished the position he assumed at the helm of the Church.\par
\switchcolumn*

\LA
\begin{Resp}\Rsign Coróna áurea super caput ejus, \Star{} Expréssa signo sanctitátis, glória honóris, et opus fortitúdinis.\end{Resp}
\begin{Resp}\Vsign Quóniam prævenísti eum in benedictiónibus dulcédinis, posuísti in cápite ejus corónam de lápide pretióso.\end{Resp}
\begin{Resp}\Rsign Expréssa signo sanctitátis, glória honóris, et opus fortitúdinis.\end{Resp}
\switchcolumn
\EN
\begin{Resp}\Rsign A crown of gold upon his head, \Star{} Wherein is engraved Holiness, an ornament of honour, a costly work.\end{Resp}
\begin{Resp}\Vsign For Thou hast prevented him with the blessings of sweetness, Thou hast set a crown of precious stones upon his head.\end{Resp}
\begin{Resp}\Rsign Wherein is engraved Holiness, an ornament of honour, a costly work.\end{Resp}
\switchcolumn*

\LA
\LessonHead{Lectio VIII}
\switchcolumn
\EN
\LessonHead{Lesson VIII}
\switchcolumn*

\LA
\noindent In univérsa namque Ecclésia, Tu es Christus Fílius Dei vivi, quotídie Petrus dicit; et omnis lingua, quæ confitétur Dóminum, magistério hujus vocis imbúitur. Hæc fides diábolum vincit et captivórum ejus víncula dissólvit. Hæc érutos mundo, ínserit cælo, et portæ ínferi advérsus eam prævalére non possunt. Tanta enim divínitus soliditáte muníta est, ut eam neque hærética umquam corrúmpere právitas, nec pagána potúerit superáre perfídia. His ítaque modis, dilectíssimi, rationábile obséquio celebrétur hodiérna festívitas: ut in persóna humilitátis meæ ille intelligátur, ille honorétur, in quo et ómnium pastórum sollicitúdo, cum commendatárum sibi óvium custódia persevérat, et cujus étiam dígnitas in indígno heréde non déficit.\par
\switchcolumn
\EN
\noindent In the universal Church it is as if Peter were still saying every day: Thou art the Christ, the Son of the living God. For every tongue which confesseth the Lord is taught that confession by the teaching of Peter. This is the Faith that overcometh the devil and looseth the bonds of his prisoners. This is the Faith which maketh men free of the world and bringeth them to heaven, and the gates of hell are impotent to prevail against it. This is the rock which God hath fortified with such ramparts of salvation, that the contagion of heresy will never be able to infect it, nor idolatry and unbelief to overcome it. And therefore, dearly beloved, we celebrate today’s festival with reasonable obedience, that in my humble person he may be acknowledged and honoured who doth continue to care for all the shepherds as well as sheep entrusted unto him, and who doth lose none of his dignity even in an unworthy successor.\par
\switchcolumn*

\LA
\begin{Resp}\Rsign Hic est vere Martyr, qui pro Christi nómine sánguinem suum fudit: \Star{} Qui minas júdicum non tímuit, nec terrénæ dignitátis glóriam quæsívit, sed ad cæléstia regna pervénit.\end{Resp}
\begin{Resp}\Vsign Justum dedúxit Dóminus per vias rectas, et osténdit illi regnum Dei.\end{Resp}
\begin{Resp}\Rsign Qui minas júdicum non tímuit, nec terrénæ dignitátis glóriam quæsívit, sed ad cæléstia regna pervénit.\end{Resp}
\begin{Resp}\Vsign Glória Patri, et Fílio, \Star{} et Spirítui Sancto.\end{Resp}
\begin{Resp}\Rsign Qui minas júdicum non tímuit, nec terrénæ dignitátis glóriam quæsívit, sed ad cæléstia regna pervénit.\end{Resp}
\switchcolumn
\EN
\begin{Resp}\Rsign This is a Martyr indeed, who shed his blood for Christ's Name's sake \Star{} Who feared not for the threats of judges, nor sought to be great with the glory of this world, but pressed on unto the kingdom of heaven.\end{Resp}
\begin{Resp}\Vsign The Lord guided the righteous in right paths, and showed him the kingdom of God.\end{Resp}
\begin{Resp}\Rsign Who feared not for the threats of judges, nor sought to be great with the glory of this world, but pressed on unto the kingdom of heaven.\end{Resp}
\begin{Resp}\Vsign Glory be to the Father, and to the Son, \Star{} and to the Holy Ghost.\end{Resp}
\begin{Resp}\Rsign Who feared not for the threats of judges, nor sought to be great with the glory of this world, but pressed on unto the kingdom of heaven.\end{Resp}
\switchcolumn*

\LA
\LessonHead{Lectio IX}
\switchcolumn
\EN
\LessonHead{Lesson IX}
\switchcolumn*

\LA
\noindent Cum ergo cohortatiónes nostras áuribus vestræ sanctitátis adhibémus, ipsum vobis, cujus vice fúngimur, loqui crédite: quia et illíus vos afféctu monémus, et non áliud vobis, quam quod dócuit, prædicámus; obsecrántes, ut succíncti lumbos mentis vestræ, castam et sóbriam vitam in Dei timóre ducátis. Coróna mea, sicut Apóstolus ait, et gáudium vos estis, si fides vestra, quæ ab inítio Evangélii in univérso mundo prædicáta est, in dilectióne et sanctitáte permánserit. Nam licet omnem Ecclésiam, quæ in toto est orbe terrárum, cunctis opórteat florére virtútibus; vos tamen præcípue inter céteros pópulos decet méritis pietátis excéllere, quos in ipsa apostólicæ petræ arce fundátos, et Dóminus noster Jesus Christus cum ómnibus redémit, et beátus Apóstolus Petrus præ ómnibus erudívit.\par
\switchcolumn
\EN
\noindent When, therefore, we address our exhortations to your godly ears, believe ye that ye are hearing him speak whose office we are discharging. Yea, it is with his love for you that we warn you. And we preach unto you no other thing than that which he taught, entreating you as did he: Gird up the loins of your mind; be sober; be ye holy in all manner of living; pass the time of your sojourning here in the fear of God. My disciples, dearly beloved, ye are to me as the disciples of the Apostle Paul were to him, namely: My crown and joy; if so be that your faith, abide, still in all lowliness and holiness, like unto the first times of the Gospel. For although the whole Church, which is in all the world, should indeed abound in all the virtues, it becometh especially you among all others to excel in acts of piety, founded as ye be on the very citadel of the Apostolic Rock ye who have not only been redeemed with the rest of men by our Lord Jesus Christ, but who have been instructed by the blessed Apostle Peter far beyond all others.\par
\switchcolumn*

\LA
\TeDeum{Te Deum}
\switchcolumn
\EN
\TeDeum{Te Deum}
\switchcolumn*

\end{paracol}

\Office{Commune Plurimorum Martyrum Pontificum}{Commune Plurimorum Martyrum Pontificum}{}
\addcontentsline{toc}{subsection}{Commune Plurimorum Martyrum Pontificum \textit{— Commune Plurimorum Martyrum Pontificum}}
\begin{paracol}{2}
\LA
\LessonHead{Lectio VII}
\switchcolumn
\EN
\LessonHead{Lesson VII}
\switchcolumn*

\LA
\LessonTitle{Léctio sancti Evangélii secúndum Lucam}
\switchcolumn
\EN
\LessonTitle{From the Holy Gospel according to Luke}
\switchcolumn*

\LA
\Cite{Luc 21:9-19}
\switchcolumn
\EN
\Cite{Luke 21:9-19}
\switchcolumn*

\LA
\noindent In illo témpore: Dixit Jesus discípulis suis: Cum audiéritis prǽlia, et seditiónes, nolíte terréri: opórtet primum hæc fíeri, sed nondum statim finis. Et réliqua.\par
\switchcolumn
\EN
\noindent At that time, Jesus said unto His disciples: When ye shall hear of wars and commotions, be not terrified for these things must first come to pass; but the end is not by and by. And so on.\par
\switchcolumn*

\LA
\LessonTitle{Homilía sancti Gregórii Papæ}
\switchcolumn
\EN
\LessonTitle{Homily by Pope St Gregory (the Great.)}
\switchcolumn*

\LA
\Source{Homilia 35 in Evangel.}
\switchcolumn
\EN
\Source{35th on the Gospels.}
\switchcolumn*

\LA
\noindent Dóminus ac Redémptor noster peritúri mundi præcurréntia mala denúntiat, ut eo minus pertúrbent veniéntia, quo fúerint præscíta. Minus enim jácula fériunt, quæ prævidéntur: et nos tolerabílius mundi mala suscípimus, si contra hæc per præsciéntiæ clípeum munímur. Ecce enim dicit: Cum audiéritis prǽlia et seditiónes, nolíte terréri: opórtet enim primum hæc fíeri, sed nondum statim finis. Pensánda sunt verba Redemptóris nostri, per quæ nos áliud intérius, áliud extérius passúros esse denúntiat; bella quippe ad hostes pértinent, seditiónes ad cives. Ut ergo nos índicet intérius exteriúsque turbári, áliud nos fatétur ab hóstibus, áliud a frátribus pérpeti.\par
\switchcolumn
\EN
\noindent Our Lord and Redeemer willeth us to know what shall be the signs that the end of the world is at hand, to the end that ye may be the less terrified, when that cometh whereof ye have already had warning. Darts strike less which are seen coming: and the plagues of the earth will be to us more bearable, if we are harnessed against them with the shield of foreknowledge. Behold, how He saith When ye shall hear of wars and commotions be not terrified for these things must first come to pass; but the end is not by and by. It behoveth us to ponder these words of our Redeemer, wherein He warneth us of suffering, from without, and from within. Wars are the work of a foreign enemy, commotions of the citizens. Therefore, that He may let us know that we shall be troubled from within and from without, He showeth that our wrestling shall be in part against strangers, and in part against our brethren.\par
\switchcolumn*

\LA
\begin{Resp}\Rsign Propter testaméntum Dómini, et leges patérnas, Sancti Dei perstitérunt in amóre fraternitátis: \Star{} Quia unus fuit semper spíritus in eis, et una fides.\end{Resp}
\begin{Resp}\Vsign Ecce quam bonum et quam jucúndum habitáre fratres in unum.\end{Resp}
\begin{Resp}\Rsign Quia unus fuit semper spíritus in eis, et una fides.\end{Resp}
\switchcolumn
\EN
\begin{Resp}\Rsign Because of the covenant of the Lord, and the laws of their fathers, the Saints of God abode in brotherly love, \Star{} For one spirit and one faith was ever in them.\end{Resp}
\begin{Resp}\Vsign Behold how good and how pleasant it is for brethren to dwell together in unity.\end{Resp}
\begin{Resp}\Rsign For one spirit and one faith was ever in them.\end{Resp}
\switchcolumn*

\LA
\LessonHead{Lectio VIII}
\switchcolumn
\EN
\LessonHead{Lesson VIII}
\switchcolumn*

\LA
\noindent Sed his malis præveniéntibus, quia non statim finis sequátur, adjúngit: Surget gens contra gentem, et regnum advérsus regnum: et terræmótus magni erunt per loca, et pestiléntiæ, et fames, terrorésque de cælo: et signa magna erunt. Última tribulátio multis tribulatiónibus prævenítur: et per crebra mala, quæ prævéniunt, indicántur mala perpétua, quæ subsequéntur. Et ídeo post bella et seditiónes non statim finis: quia multa debent mala præcúrrere, ut malum váleant sine fine nuntiáre.\par
\switchcolumn
\EN
\noindent When these woes come, the end is not by and by. And He saith further Nation shall rise against nation, and kingdom against kingdom; and great earthquakes shall be in diverse places, and pestilences, and famines, and fearful sights and great signs shall there be from heaven. Before the last tribulation cometh, shall come many other tribulations: and, by the many woes which shall come first, shall be foreshadowed the everlasting woe which shall come in the end. And therefore, after wars and commotions, the end is not yet by and by many woes must come first, to give warning of the woe that hath no end.\par
\switchcolumn*

\LA
\begin{Resp}\Rsign Sancti mei, qui in carne pósiti, certámen habuístis: \Star{} Mercédem labóris ego reddam vobis.\end{Resp}
\begin{Resp}\Vsign Veníte benedícti Patris mei, percípite regnum.\end{Resp}
\begin{Resp}\Rsign Mercédem labóris ego reddam vobis.\end{Resp}
\begin{Resp}\Vsign Glória Patri, et Fílio, \Star{} et Spirítui Sancto.\end{Resp}
\begin{Resp}\Rsign Mercédem labóris ego reddam vobis.\end{Resp}
\switchcolumn
\EN
\begin{Resp}\Rsign O ye My Saints, who, being in the flesh, didst have striving \Star{} I will render unto you a reward of your labours.\end{Resp}
\begin{Resp}\Vsign Come, ye blessed of My Father, inherit the kingdom\end{Resp}
\begin{Resp}\Rsign I will render unto you a reward of your labours.\end{Resp}
\begin{Resp}\Vsign Glory be to the Father, and to the Son, \Star{} and to the Holy Ghost.\end{Resp}
\begin{Resp}\Rsign I will render unto you a reward of your labours.\end{Resp}
\switchcolumn*

\LA
\LessonHead{Lectio IX}
\switchcolumn
\EN
\LessonHead{Lesson IX}
\switchcolumn*

\LA
\noindent Sed cum tot signa perturbatiónis dicta sint, opórtet, ut eórum consideratiónem bréviter per síngula perstringámus: quia necésse est, ut ália e cælo, ália e terra, ália ab eleméntis, ália ab homínibus patiámur. Ait enim: Surget gens contra gentem, ecce perturbátio hóminum: erunt terræmótus magni per loca, ecce respéctus iræ désuper: erunt pestiléntiæ ecce inæquálitas córporum: erit fames, ecce sterílitas terræ: terrorésque de cælo et tempestátes, ecce inæquálitas áëris. Quia ergo ómnia consummánda sunt, ante consummatiónem ómnia perturbántur: et qui in cunctis delíquimus, in cunctis ferímur: ut impleátur quod dícitur, Et pugnábit pro eo orbis terrárum contra insensátos.\par
\switchcolumn
\EN
\noindent But, forasmuch as the signs and troubles whereof the Lord speaketh are so manifold, we must needs shortly consider each for, of necessity, we must suffer some things from heaven, some from the earth, some from the powers of nature, and some from men. For where He saith Nation shall rise against nation He speaketh concerning the troubling of men where great earthquakes shall be in diverse places concerning wrath from above: where: and pestilences concerning the frailty of the body where and famines concerning the barrenness of the earth: where fearful signs from heaven, and tempests concerning commotions of the air. As, then, all things shall have an end, so, before the end, shall all things be troubled and we who have sinned and come short in all things, shall in all things be afflicted, that it may be fulfilled that is written: and the world shall fight with Him against the unwise. (Wisd v. 21.)\par
\switchcolumn*

\LA
\TeDeum{Te Deum}
\switchcolumn
\EN
\TeDeum{Te Deum}
\switchcolumn*

\end{paracol}

\Office{Commune unius Confessoris Pontificis}{Common of a Confessor Bishop}{}
\addcontentsline{toc}{subsection}{Commune unius Confessoris Pontificis \textit{— Common of a Confessor Bishop}}
\begin{paracol}{2}
\LA
\LessonHead{Lectio VII}
\switchcolumn
\EN
\LessonHead{Lesson VII}
\switchcolumn*

\LA
\LessonTitle{Léctio sancti Evangélii secúndum Matthǽum}
\switchcolumn
\EN
\LessonTitle{From the Holy Gospel according to Matthew}
\switchcolumn*

\LA
\Cite{Matt 25:14-23}
\switchcolumn
\EN
\Cite{Matt 25:14-23}
\switchcolumn*

\LA
\noindent In illo témpore: Dixit Jesus discípulis suis parábolam hanc: Homo péregre proficíscens, vocávit servos suos, et trádidit illis bona sua. Et réliqua.\par
\switchcolumn
\EN
\noindent At that time Jesus spake unto His disciples this parable A man, travelling into a far country, called his own servants, and delivered unto them his goods. And so on.\par
\switchcolumn*

\LA
\LessonTitle{Homilía sancti Gregórii Papæ}
\switchcolumn
\EN
\LessonTitle{Homily by Pope St Gregory (the Great.)}
\switchcolumn*

\LA
\Source{Homilia 9 in Evangelia}
\switchcolumn
\EN
\Source{9th on the Gospels.}
\switchcolumn*

\LA
\noindent Léctio sancti Evangélii, fratres caríssimi, solícite consideráre nos ádmonet, ne nos, qui plus céteris in hoc mundo accepísse áliquid cérnimur, ab Auctóre mundi grávius inde judicémur. Cum enim augéntur dona, ratiónes étiam crescunt donórum. Tanto ergo esse humílior atque ad serviéndum Deo prómptior quisque debet ex múnere, quanto se obligatiórem esse cónspicit in reddénda ratióne. Ecce homo, qui péregre proficíscitur, servos suos vocat, eisque ad negótium talénta partítur. Post multum vero témporis positúrus ratiónem revértitur: bene operántes pro apportáto lucro remúnerat, servum vero a bono ópere torpéntem damnat.\par
\switchcolumn
\EN
\noindent Dearly beloved brethren, this Lesson from the Holy Gospel moveth us to take good heed lest we, who are seen in this world to have received more than others, should thereby bring ourselves into greater condemnation from the Maker of this world. To whom much is given, of the same is much required. Therefore, let him that receiveth much, strive to be all the more lowly, and all the more ready to do God service, for his very gifts’ sake, knowing that he will be obliged to give account thereof. Behold, a man, travelling into a far country, calleth his own servants, and delivereth unto them talents, to the end that they may trade therewith. After a long time, the lord of those servants cometh, and reckoneth with them, and to them that have done well He rendereth a reward of their labours, but that servant which was careless of his master’s work He condemneth.\par
\switchcolumn*

\LA
\begin{Resp}\Rsign Amávit eum Dóminus, et ornávit eum: stolam glóriæ índuit eum, \Star{} Et ad portas paradísi coronávit eum.\end{Resp}
\begin{Resp}\Vsign Induit eum Dóminus lorícam fídei, et ornávit eum.\end{Resp}
\begin{Resp}\Rsign Et ad portas paradísi coronávit eum.\end{Resp}
\switchcolumn
\EN
\begin{Resp}\Rsign The Lord loved him and beautified him He clothed him with a robe of glory \Star{} And crowned him at the gates of Paradise.\end{Resp}
\begin{Resp}\Vsign The Lord hath put on him the breast-plate of faith, and hath adorned him.\end{Resp}
\begin{Resp}\Rsign And crowned him at the gates of Paradise.\end{Resp}
\switchcolumn*

\LA
\LessonHead{Lectio VIII}
\switchcolumn
\EN
\LessonHead{Lesson VIII}
\switchcolumn*

\LA
\noindent Quis ítaque iste homo est qui péregre proficíscitur, nisi Redémptor noster, qui, in ea carne quam assúmpserat, ábiit in cælum? Carnis enim locus próprius terra est; quæ quasi ad peregrína dúcitur, dum per Redemptórem nostrum in cælo collocátur. Sed homo iste, péregre proficíscens, servis suis bona sua trádidit; quia fidélibus suis spirituália dona concéssit. Et uni quidem quinque talénta, álii duo, álii vero commísit unum. Quinque étenim sunt córporis sensus, vidélicet: visus, audítus, gustus, odorátus et tactus. Quinque ergo taléntis, donum quinque sénsuum, id est, exteriórum sciéntia, exprímitur. Duóbus vero, intelléctus et operátio designátur. Uníus autem talénti nómine, intelléctus tantúmmodo designátur.\par
\switchcolumn
\EN
\noindent What other, then, is that man travelling into a far country but our Redeemer, Who is gone up from us into heaven in that Flesh Which He had taken into Himself? For the earth is the home of the Flesh, Which travelleth into a far country when our Redeemer giveth It a place in heaven. But that man travelling into a far country delivered unto his servants his goods and so doth our Redeemer give spiritual gifts unto His faithful people. “And unto one he gave five talents, to another two, and to another one.” There are five bodily senses that is, sight, hearing, taste, smell, and touch. By the five talents therefore are signified the five senses, that is, outward knowledge. By the two, wit and work. And by the figure of the one talent, understanding, which is alone.\par
\switchcolumn*

\LA
\begin{Resp}\Rsign Sint lumbi vestri præcíncti, et lucérnæ ardéntes in mánibus vestris: \Star{} Et vos símiles homínibus exspectántibus dóminum suum, quando revertátur a núptiis.\end{Resp}
\begin{Resp}\Vsign Vigiláte ergo, quia nescítis qua hora Dóminus vester ventúrus sit.\end{Resp}
\begin{Resp}\Rsign Et vos símiles homínibus exspectántibus dóminum suum, quando revertátur a núptiis.\end{Resp}
\begin{Resp}\Vsign Glória Patri, et Fílio, \Star{} et Spirítui Sancto.\end{Resp}
\begin{Resp}\Rsign Et vos símiles homínibus exspectántibus dóminum suum, quando revertátur a núptiis.\end{Resp}
\switchcolumn
\EN
\begin{Resp}\Rsign Let your loins be girded about, and your lights burning; \Star{} And ye yourselves like unto men that wait for their lord, when he will return from the wedding.\end{Resp}
\begin{Resp}\Vsign Watch therefore, for ye know not what hour your Lord doth come.\end{Resp}
\begin{Resp}\Rsign And ye yourselves like unto men that wait for their lord, when he will return from the wedding.\end{Resp}
\begin{Resp}\Vsign Glory be to the Father, and to the Son, \Star{} and to the Holy Ghost.\end{Resp}
\begin{Resp}\Rsign And ye yourselves like unto men that wait for their lord, when he will return from the wedding.\end{Resp}
\switchcolumn*

\LA
\LessonHead{Lectio IX}
\switchcolumn
\EN
\LessonHead{Lesson IX}
\switchcolumn*

\LA
\noindent Sed is, qui quinque talénta accéperat, ália quinque lucrátus est: quia sunt nonnúlli, qui, etsi intérna ac mýstica penetráre nésciunt, pro intentióne tamen supérnæ pátriæ docent recta quos possunt; de ipsis exterióribus, quæ accepérunt, duplum taléntum portant; dumque se a carnis petulántia et a terrenárum rerum ámbitu atque a visibílium voluptáte custódiunt, ab his étiam álios admonéndo compéscunt. Et sunt nonnúlli, qui, quasi duóbus taléntis ditáti, intelléctum atque operatiónem percípiunt, subtília de intérnis intélligunt, mira in exterióribus operántur; cumque et intelligéndo et operándo áliis prǽdicant, quasi duplicátum de negótio lucrum repórtant.\par
\switchcolumn
\EN
\noindent And so he that had received five talents, gained other five talents for some there be who, while yet they are not able to go on unto things inward and mystic, do yet so desire our Fatherland which is above, that they teach well all whom they can, and of those very outward things which they have received make gain double. These are they which keep themselves clean from the unruly motions of the flesh, and from the lust of the world, and from the delight of things which are seen, and, by their preaching, keep other men also clean from all these things. And some there are who receive, as their two talents, the power to think and the power to work. These are they which inwardly understand dark things, and outwardly work wonders. And these, since they preach unto others, both through their understanding and their works, gain, as it were, double, for the talents which they have received.\par
\switchcolumn*

\LA
\TeDeum{Te Deum}
\switchcolumn
\EN
\TeDeum{Te Deum}
\switchcolumn*

\end{paracol}

\Office{Commune Doctoris Pontificis}{Common of a Doctor Bishop}{}
\addcontentsline{toc}{subsection}{Commune Doctoris Pontificis \textit{— Common of a Doctor Bishop}}
\begin{paracol}{2}
\LA
\LessonHead{Lectio VII}
\switchcolumn
\EN
\LessonHead{Lesson VII}
\switchcolumn*

\LA
\LessonTitle{Léctio sancti Evangélii secúndum Matthǽum}
\switchcolumn
\EN
\LessonTitle{From the Holy Gospel according to Matthew}
\switchcolumn*

\LA
\Cite{Matt 5:13-19}
\switchcolumn
\EN
\Cite{Matt 5:13-19}
\switchcolumn*

\LA
\noindent In illo témpore: Dixit Jesus discípulis suis: Vos estis sal terræ. Quod si sal evanúerit, in quo saliétur? Et réliqua.\par
\switchcolumn
\EN
\noindent At that time: Jesus said unto his disciples: Ye are the salt of the earth: But if the salt have lost his savour, wherewith shall it be salted? And so on, and that which followeth.\par
\switchcolumn*

\LA
\LessonTitle{Homilía sancti Augustíni Epíscopi}
\switchcolumn
\EN
\LessonTitle{A Homily by St. Augustine the Bishop}
\switchcolumn*

\LA
\Source{Lib. 1 de Sermone Domini in monte, cap. 6}
\switchcolumn
\EN
\Source{Lib. 1 de Sermóne Dómini in monte, cap. 6}
\switchcolumn*

\LA
\noindent Osténdit Dóminus fátuos esse judicándos, qui temporálium bonórum vel cópiam sectántes vel inópiam metuéntes, amíttunt ætérna, quæ nec dari possunt ab homínibus nec auférri. Itaque, si sal infatuátum fúerit, in quo saliétur? Id est, si vos, per quos condiéndi sunt quodámmodo pópuli, metu persecutiónum temporálium amiséritis regna cælórum; qui erunt hómines, per quos a vobis error auferátur, cum vos elégerit Deus, per quos errórem áuferat ceterórum?\par
\switchcolumn
\EN
\noindent The Lord would have us understand how that men do lose their power of savouring others with righteousness when they are willing to place their eternal welfare in jeopardy for the sake of any témporal advantage, like as attainment of ease or luxury, or escape from suffering or toil. For that which is eternal, unlike things of this world, can neither be bestowed by men, nor by them taken away. Hence, when he asketh: If the salt have lost his savour, wherewith shall it be salted? he would have us understand the question to be: If ye, by whom mankind is preserved from corruption, be willing to lose the kingdom of heaven so as to escape trials or persecutions in this world, who is there to preserve you from corruption, seeing ye are they that God hath chosen to preserve all others from corruption?\par
\switchcolumn*

\LA
\begin{Resp}\Rsign Amávit eum Dóminus, et ornávit eum: stolam glóriæ índuit eum, \Star{} Et ad portas paradísi coronávit eum.\end{Resp}
\begin{Resp}\Vsign Induit eum Dóminus lorícam fídei, et ornávit eum.\end{Resp}
\begin{Resp}\Rsign Et ad portas paradísi coronávit eum.\end{Resp}
\switchcolumn
\EN
\begin{Resp}\Rsign The Lord loved him and beautified him He clothed him with a robe of glory \Star{} And crowned him at the gates of Paradise.\end{Resp}
\begin{Resp}\Vsign The Lord hath put on him the breast-plate of faith, and hath adorned him.\end{Resp}
\begin{Resp}\Rsign And crowned him at the gates of Paradise.\end{Resp}
\switchcolumn*

\LA
\LessonHead{Lectio VIII}
\switchcolumn
\EN
\LessonHead{Lesson VIII}
\switchcolumn*

\LA
\noindent Ergo ad níhilum valet sal infatuátum, nisi ut mittátur foras et calcétur ab homínibus. Non ítaque calcátur ab homínibus qui pátitur persecutiónem; sed qui, persecutiónem timéndo, infatuátur. Calcári enim non potest nisi inférior; sed inférior non est, qui, quamvis córpore multa in terra sustíneat, corde tamen fixus in cælo est.\par
\switchcolumn
\EN
\noindent Those that should be the salt of the earth, but have lost their savour, are thenceforth good for nothing, but to be cast out, and to be trodden under foot of men. But no one that suffereth persecution is truly said to be trodden under foot of men. Rather, that one is truly trodden under foot of men who through fear of persecution hath lost the savour of righteousness. For no one can be trodden upon, unless he be beneath him which treadeth upon him. And certainly no one who hath his heart in heaven, no matter how grievously he doth suffer in his body on earth, is rightly said to be beneath anyone who misuseth him.\par
\switchcolumn*

\LA
\begin{Resp}\Rsign In médio Ecclésiæ apéruit os ejus, \Star{} Et implévit eum Dóminus spíritu sapiéntiæ et intelléctus.\end{Resp}
\begin{Resp}\Vsign Jucunditátem et exsultatiónem thesaurizávit super eum.\end{Resp}
\begin{Resp}\Rsign Et implévit eum Dóminus spíritu sapiéntiæ et intelléctus.\end{Resp}
\begin{Resp}\Vsign Glória Patri, et Fílio, \Star{} et Spirítui Sancto.\end{Resp}
\begin{Resp}\Rsign Et implévit eum Dóminus spíritu sapiéntiæ et intelléctus.\end{Resp}
\switchcolumn
\EN
\begin{Resp}\Rsign In the midst of the congregation did the Lord open his mouth. \Star{} And filled him with the spirit of wisdom and understanding.\end{Resp}
\begin{Resp}\Vsign He made him rich with joy and gladness.\end{Resp}
\begin{Resp}\Rsign And filled him with the spirit of wisdom and understanding.\end{Resp}
\begin{Resp}\Vsign Glory be to the Father, and to the Son, \Star{} and to the Holy Ghost.\end{Resp}
\begin{Resp}\Rsign And filled him with the spirit of wisdom and understanding.\end{Resp}
\switchcolumn*

\LA
\LessonHead{Lectio IX}
\switchcolumn
\EN
\LessonHead{Lesson IX}
\switchcolumn*

\LA
\noindent Vos estis lumen mundi. Quómodo dixit supérius sal terræ, sic nunc dicit lumen mundi. Nam, neque supérius ista terra accipiénda est, quam pédibus corpóreis calcámus; sed hómines, qui in terra hábitant, vel étiam peccatóres, quorum condiéndis et exstinguéndis putóribus apostólicum salem Dóminus misit. Et hic mundum non cælum et terram, sed hómines qui sunt in mundo vel díligunt mundum, opórtet intélligi; quibus illuminándis Apóstoli missi sunt. Non potest cívitas abscóndi super montem pósita; id est, fundáta super insígnem magnámque justítiam, quam signíficat étiam ipse mons, in quo dísputat Dóminus.\par
\switchcolumn
\EN
\noindent Ye are the light of the world. And we are to understand the word World in the same sense as the word Earth when he spoke above of the salt of the earth, that is, not that earth whereupon we walk with our bodily feet, but the men which dwell upon the earth; in other words, sinners, for the sweetening and correction of whose corruption, the Lord hath sent his Apostles, as it were, as so much salt. And so by the world we are to understand, not the heaven and the earth, but the men who are in the world and love the world, for the enlightening of whom the Apostles have been sent. A city that is set on a hill cannot be hid: that is, what is founded upon the heights of righteousness, whereof the mountain upon which the Lord gave this discourse was itself a figure, is magnificent in the eyes of all men.\par
\switchcolumn*

\LA
\TeDeum{Te Deum}
\switchcolumn
\EN
\TeDeum{Te Deum}
\switchcolumn*

\end{paracol}

\Office{Commune Summorum Pontificum Confessorum}{Common of Popes}{}
\addcontentsline{toc}{subsection}{Commune Summorum Pontificum Confessorum \textit{— Common of Popes}}
\begin{paracol}{2}
\LA
\LessonHead{Lectio VII}
\switchcolumn
\EN
\LessonHead{Lesson VII}
\switchcolumn*

\LA
\LessonTitle{Léctio sancti Evangélii secúndum Matthǽum}
\switchcolumn
\EN
\LessonTitle{From the Holy Gospel according to Matthew}
\switchcolumn*

\LA
\Cite{Matt 16:13-19}
\switchcolumn
\EN
\Cite{Matt 16:13-19}
\switchcolumn*

\LA
\noindent In illo témpore: Venit Jesus in partes Cæsaréæ Philíppi, et interrogábat discípulos suos, dicens: Quem dicunt hómines esse Fílium hóminis? Et réliqua.\par
\switchcolumn
\EN
\noindent At that time: When Jesus came into the coasts of Caesarea Philippi, he asked his disciples, saying, Whom do men say that I the Son of man am? And so on.\par
\switchcolumn*

\LA
\LessonTitle{Homilía sancti Leónis Papæ}
\switchcolumn
\EN
\LessonTitle{A Homily by St. Leo the Pope}
\switchcolumn*

\LA
\Source{Sermo 2 in anniversario assumpt. suæ, ante medium}
\switchcolumn
\EN
\Source{Sermo 2 in anniversario assumpt. suæ ante medium}
\switchcolumn*

\LA
\noindent Cum, sicut evangélica lectióne reserátum est, interrogásset Dóminus discípulos, quem ipsum (multis divérsa opinántibus) créderent; respondissétque beátus Petrus, dicens: Tu es Christus Fílius Dei vivi; Dóminus ait: Beátus es, Simon Bar-Jona, quia caro et sanguis non revelávit tibi, sed Pater meus, qui in cælis est: et ego dico tibi, quia tu es Petrus, et super hanc petram ædificábo Ecclésiam meam, et portæ ínferi non prævalébunt advérsus eam. Et tibi dabo claves regni cælórum: et quodcúmque ligáveris super terram, erit ligátum et in cælis: et quodcúmque sólveris super terram, erit solútum et in cælis. Manet ergo disposítio veritátis, et beátus Petrus, in accépta fortitúdine petræ persevérans, suscépta Ecclésiæ gubernácula non relíquit.\par
\switchcolumn
\EN
\noindent When the Lord, as we read in the Gospel, asked his disciples who did men, amid their diverse speculations, believe him the Son of Man to be, blessed Peter answered and said: Thou art the Christ, the Son of the living God. And the Lord answered and said unto him: Blessed art thou, Simon Bar-Jona: for flesh and blood hath not revealed it unto thee, but my Father, which is in heaven: and I say also unto thee: That thou art Peter, and upon this rock I will build my Church, and the gates of hell shall not prevail against it; and I will give unto thee the keys of the kingdom of heaven; and whatsoever thou shalt bind on earth shall be bound in heaven; and whatsoever thou shalt loose on earth shall be loosed in heaven. But the dispensation of truth perdures, and blessed Peter, persevering in the strength of the rock which he hath received, hath not relinquished the position he assumed at the helm of the Church.\par
\switchcolumn*

\LA
\begin{Resp}\Rsign Amávit eum Dóminus, et ornávit eum: stolam glóriæ índuit eum, \Star{} Et ad portas paradísi coronávit eum.\end{Resp}
\begin{Resp}\Vsign Induit eum Dóminus lorícam fídei, et ornávit eum.\end{Resp}
\begin{Resp}\Rsign Et ad portas paradísi coronávit eum.\end{Resp}
\switchcolumn
\EN
\begin{Resp}\Rsign The Lord loved him and beautified him He clothed him with a robe of glory \Star{} And crowned him at the gates of Paradise.\end{Resp}
\begin{Resp}\Vsign The Lord hath put on him the breast-plate of faith, and hath adorned him.\end{Resp}
\begin{Resp}\Rsign And crowned him at the gates of Paradise.\end{Resp}
\switchcolumn*

\LA
\LessonHead{Lectio VIII}
\switchcolumn
\EN
\LessonHead{Lesson VIII}
\switchcolumn*

\LA
\noindent In univérsa namque Ecclésia, Tu es Christus Fílius Dei vivi, quotídie Petrus dicit; et omnis lingua, quæ confitétur Dóminum, magistério hujus vocis imbúitur. Hæc fides diábolum vincit et captivórum ejus víncula dissólvit. Hæc érutos mundo, ínserit cælo, et portæ ínferi advérsus eam prævalére non possunt. Tanta enim divínitus soliditáte muníta est, ut eam neque hærética umquam corrúmpere právitas, nec pagána potúerit superáre perfídia. His ítaque modis, dilectíssimi, rationábile obséquio celebrétur hodiérna festívitas: ut in persóna humilitátis meæ ille intelligátur, ille honorétur, in quo et ómnium pastórum sollicitúdo, cum commendatárum sibi óvium custódia persevérat, et cujus étiam dígnitas in indígno heréde non déficit.\par
\switchcolumn
\EN
\noindent In the universal Church it is as if Peter were still saying every day: Thou art the Christ, the Son of the living God. For every tongue which confesseth the Lord is taught that confession by the teaching of Peter. This is the Faith that overcometh the devil and looseth the bonds of his prisoners. This is the Faith which maketh men free of the world and bringeth them to heaven, and the gates of hell are impotent to prevail against it. This is the rock which God hath fortified with such ramparts of salvation, that the contagion of heresy will never be able to infect it, nor idolatry and unbelief to overcome it. And therefore, dearly beloved, we celebrate today’s festival with reasonable obedience, that in my humble person he may be acknowledged and honoured who doth continue to care for all the shepherds as well as sheep entrusted unto him, and who doth lose none of his dignity even in an unworthy successor.\par
\switchcolumn*

\LA
\begin{Resp}\Rsign Sint lumbi vestri præcíncti, et lucérnæ ardéntes in mánibus vestris: \Star{} Et vos símiles homínibus exspectántibus dóminum suum, quando revertátur a núptiis.\end{Resp}
\begin{Resp}\Vsign Vigiláte ergo, quia nescítis qua hora Dóminus vester ventúrus sit.\end{Resp}
\begin{Resp}\Rsign Et vos símiles homínibus exspectántibus dóminum suum, quando revertátur a núptiis.\end{Resp}
\begin{Resp}\Vsign Glória Patri, et Fílio, \Star{} et Spirítui Sancto.\end{Resp}
\begin{Resp}\Rsign Et vos símiles homínibus exspectántibus dóminum suum, quando revertátur a núptiis.\end{Resp}
\switchcolumn
\EN
\begin{Resp}\Rsign Let your loins be girded about, and your lights burning; \Star{} And ye yourselves like unto men that wait for their lord, when he will return from the wedding.\end{Resp}
\begin{Resp}\Vsign Watch therefore, for ye know not what hour your Lord doth come.\end{Resp}
\begin{Resp}\Rsign And ye yourselves like unto men that wait for their lord, when he will return from the wedding.\end{Resp}
\begin{Resp}\Vsign Glory be to the Father, and to the Son, \Star{} and to the Holy Ghost.\end{Resp}
\begin{Resp}\Rsign And ye yourselves like unto men that wait for their lord, when he will return from the wedding.\end{Resp}
\switchcolumn*

\LA
\LessonHead{Lectio IX}
\switchcolumn
\EN
\LessonHead{Lesson IX}
\switchcolumn*

\LA
\noindent Cum ergo cohortatiónes nostras áuribus vestræ sanctitátis adhibémus, ipsum vobis, cujus vice fúngimur, loqui crédite: quia et illíus vos afféctu monémus, et non áliud vobis, quam quod dócuit, prædicámus; obsecrántes, ut succíncti lumbos mentis vestræ, castam et sóbriam vitam in Dei timóre ducátis. Coróna mea, sicut Apóstolus ait, et gáudium vos estis, si fides vestra, quæ ab inítio Evangélii in univérso mundo prædicáta est, in dilectióne et sanctitáte permánserit. Nam licet omnem Ecclésiam, quæ in toto est orbe terrárum, cunctis opórteat florére virtútibus; vos tamen præcípue inter céteros pópulos decet méritis pietátis excéllere, quos in ipsa apostólicæ petræ arce fundátos, et Dóminus noster Jesus Christus cum ómnibus redémit, et beátus Apóstolus Petrus præ ómnibus erudívit.\par
\switchcolumn
\EN
\noindent When, therefore, we address our exhortations to your godly ears, believe ye that ye are hearing him speak whose office we are discharging. Yea, it is with his love for you that we warn you. And we preach unto you no other thing than that which he taught, entreating you as did he: Gird up the loins of your mind; be sober; be ye holy in all manner of living; pass the time of your sojourning here in the fear of God. My disciples, dearly beloved, ye are to me as the disciples of the Apostle Paul were to him, namely: My crown and joy; if so be that your faith, abide, still in all lowliness and holiness, like unto the first times of the Gospel. For although the whole Church, which is in all the world, should indeed abound in all the virtues, it becometh especially you among all others to excel in acts of piety, founded as ye be on the very citadel of the Apostolic Rock ye who have not only been redeemed with the rest of men by our Lord Jesus Christ, but who have been instructed by the blessed Apostle Peter far beyond all others.\par
\switchcolumn*

\LA
\TeDeum{Te Deum}
\switchcolumn
\EN
\TeDeum{Te Deum}
\switchcolumn*

\end{paracol}

\Office{Commune Confessoris non pontificis}{Common of a Confessor not a Bishop}{}
\addcontentsline{toc}{subsection}{Commune Confessoris non pontificis \textit{— Common of a Confessor not a Bishop}}
\begin{paracol}{2}
\LA
\LessonHead{Lectio I}
\switchcolumn
\EN
\LessonHead{Lesson I}
\switchcolumn*

\LA
\LessonTitle{De libro Ecclesiástici}
\switchcolumn
\EN
\LessonTitle{Lesson from the book of Ecclesiasticus}
\switchcolumn*

\LA
\Cite{Sir 31:8-11}
\switchcolumn
\EN
\Cite{Sir 31:8-11}
\switchcolumn*

\LA
\noindent Beátus vir, qui invéntus est sine mácula, et qui post aurum non ábiit, nec sperávit in pecúnia et thesáuris. Quis est hic? et laudábimus eum: fecit enim mirabília in vita sua. Qui probátus est in illo, et perféctus est, erit illi glória ætérna: Qui pótuit tránsgredi, et non est transgréssus; fácere mala, et non fecit: ídeo stabilíta sunt bona illíus in Dómino, et eleemósynas illíus enarrábit omnis ecclésia sanctórum.\par
\switchcolumn
\EN
\noindent Blessed is the rich man that is found without blemish: and that hath not gone after gold, nor put his trust in money nor in treasures. Who is he, and we will praise him? for he hath done wonderful things in his life. Who hath been tried thereby, and made perfect, he shall have glory everlasting. He that could have transgressed, and hath not transgressed: and could do evil things, and hath not done them: Therefore are his goods established in the Lord, and all the church of the saints shall declare his alms\par
\switchcolumn*

\LA
\begin{Resp}\Rsign Euge, serve bone et fidélis, quia in pauca fuísti fidélis, supra multa te constítuam: \Star{} Intra in gáudium Dómini tui.\end{Resp}
\begin{Resp}\Vsign Dómine, quinque talénta tradidísti mihi, ecce ália quinque superlucrátus sum.\end{Resp}
\begin{Resp}\Rsign Intra in gáudium Dómini tui.\end{Resp}
\switchcolumn
\EN
\begin{Resp}\Rsign Well done, thou good and faithful servant, thou hast been faithful over a few things. I will make thee ruler over many things: \Star{} Enter thou into the joy of thy Lord.\end{Resp}
\begin{Resp}\Vsign Lord, thou deliveredst unto me five talents; behold, I have gained beside them five talents more.\end{Resp}
\begin{Resp}\Rsign Enter thou into the joy of thy Lord.\end{Resp}
\switchcolumn*

\LA
\LessonHead{Lectio II}
\switchcolumn
\EN
\LessonHead{Lesson II}
\switchcolumn*

\LA
\Cite{Sir 32:18-20; 32:28; 33:1-3}
\switchcolumn
\EN
\Cite{Sir 32:18-20; 32:28; 33:1-3}
\switchcolumn*

\LA
\noindent Qui timet Dóminum excípiet doctrínam ejus: et qui vigiláverint ad illum invénient benedictiónem. Qui quærit legem, replébitur ab ea: et qui insidióse agit, scandalizábitur in ea. Qui timent Dóminum invénient judícium justum, et justítias quasi lumen accéndent. Qui credit Deo atténdit mandátis: et qui confídit in illo non minorábitur. Timénti Dóminum non occúrrent mala: sed in tentatióne Deus illum conservábit, et liberábit a malis. Sápiens non odit mandáta et justítias, et non illidétur quasi in procélla navis. Homo sensátus credit legi Dei, et lex illi fidélis.\par
\switchcolumn
\EN
\noindent He that feareth the Lord, will receive his discipline: and they that will seek him early, shall find a blessing. He that seeketh the law, shall be filled with it: and he that dealeth deceitfully, shall meet with a stumblingblock therein. They that fear the Lord, shall find just judgment, and shall kindle justice as a light. He that believeth God, taketh heed to the commandments: and he that trusteth in him, shall fare never the worse. No evils shall happen to him that feareth the Lord, but in temptation God will keep him, and deliver him from evils. A wise man hateth not the commandments and justices, and he shall not be dashed in pieces as a ship in a storm. A man of understanding is faithful to the law of God, and the law is faithful to him.\par
\switchcolumn*

\LA
\begin{Resp}\Rsign Justus germinábit sicut lílium: \Star{} Et florébit in ætérnum ante Dóminum.\end{Resp}
\begin{Resp}\Vsign Plantátus in domo Dómini, in átriis domus Dei nostri.\end{Resp}
\begin{Resp}\Rsign Et florébit in ætérnum ante Dóminum.\end{Resp}
\switchcolumn
\EN
\begin{Resp}\Rsign The righteous shall grow as the lily; \Star{} Yea, he shall flourish in the presence of the Lord for ever.\end{Resp}
\begin{Resp}\Vsign Those that be planted in the house of the Lord, shall flourish in the courts of the house of our God.\end{Resp}
\begin{Resp}\Rsign Yea, he shall flourish in the presence of the Lord for ever.\end{Resp}
\switchcolumn*

\LA
\LessonHead{Lectio III}
\switchcolumn
\EN
\LessonHead{Lesson III}
\switchcolumn*

\LA
\Cite{Sir 34:14-20}
\switchcolumn
\EN
\Cite{Sir 34:14-20}
\switchcolumn*

\LA
\noindent Spíritus timéntium Deum quǽritur, et in respéctu illíus benedicétur. Spes enim illórum in salvántem illos, et óculi Dei in diligéntes se. Qui timet Dóminum nihil trepidábit: et non pavébit, quóniam ipse est spes ejus. Timéntis Dóminum, beáta est ánima ejus. Ad quem réspicit, et quis est fortitúdo ejus? Óculi Dómini super timéntes eum: protéctor poténtiæ, firmaméntum virtútis, tégimen ardóris, et umbráculum meridiáni: deprecátio offensiónis, et adjutórium casus: exáltans ánimam, et illúminans óculos, dans sanitátem, et vitam, et benedictiónem.\par
\switchcolumn
\EN
\noindent The spirit of those that fear God; is sought after, and by his regard shall be blessed. For their hope is on him that saveth them, and the eyes of God are upon them that love him. He that feareth the Lord shall tremble at nothing, and shall not be afraid for he is his hope. The soul of him that feareth the Lord is blessed. To whom doth he look, and who in his strength? The eyes of the Lord are upon them that fear him, he is their powerful protector, and strong stay, a defence from the heat, and a cover from the sun at noon, A preservation from stumbling, and a help from falling; he raiseth up the soul, and enlighteneth the eyes, and giveth health, and life, and blessing.\par
\switchcolumn*

\LA
\begin{Resp}\Rsign Iste cognóvit justítiam, et vidit mirabília magna, et exorávit Altíssimum: \Star{} Et invéntus est in número Sanctórum.\end{Resp}
\begin{Resp}\Vsign Iste est qui contémpsit vitam mundi, et pervénit ad cæléstia regna.\end{Resp}
\begin{Resp}\Rsign Et invéntus est in número Sanctórum.\end{Resp}
\begin{Resp}\Vsign Glória Patri, et Fílio, \Star{} et Spirítui Sancto.\end{Resp}
\begin{Resp}\Rsign Et invéntus est in número Sanctórum.\end{Resp}
\switchcolumn
\EN
\begin{Resp}\Rsign This is he which knew righteousness, and saw great wonders, and made his prayer unto the Most High; \Star{} And he is numbered among the Saints.\end{Resp}
\begin{Resp}\Vsign This is he which loved not his life in this world, and is come unto an everlasting kingdom.\end{Resp}
\begin{Resp}\Rsign And he is numbered among the Saints.\end{Resp}
\begin{Resp}\Vsign Glory be to the Father, and to the Son, \Star{} and to the Holy Ghost.\end{Resp}
\begin{Resp}\Rsign And he is numbered among the Saints.\end{Resp}
\switchcolumn*

\LA
\LessonHead{Lectio VII}
\switchcolumn
\EN
\LessonHead{Lesson VII}
\switchcolumn*

\LA
\LessonTitle{Léctio sancti Evangélii secúndum Lucam}
\switchcolumn
\EN
\LessonTitle{From the Holy Gospel according to Luke}
\switchcolumn*

\LA
\Cite{Luc 12:35-40}
\switchcolumn
\EN
\Cite{Luke 12:35-40}
\switchcolumn*

\LA
\noindent In illo témpore: Dixit Jesus discípulis suis: Sint lumbi vestri præcincti, et lucernæ ardentes in mánibus vestris. Et réliqua.\par
\switchcolumn
\EN
\noindent At that time, Jesus said unto His disciples: Let your loins be girded about, and your lights burning. And so on.\par
\switchcolumn*

\LA
\LessonTitle{Homilía sancti Gregórii Papæ}
\switchcolumn
\EN
\LessonTitle{Homily by Pope St. Gregory (the Great.)}
\switchcolumn*

\LA
\Source{Homilia 13 in Evang.}
\switchcolumn
\EN
\Source{13th on the Gospels.}
\switchcolumn*

\LA
\noindent Sancti Evangélii, fratres caríssimi, apérta vobis est léctio recitata. Sed ne alíquibus ipsa ejus planíties alta fortásse videátur, eam sub brevitáte transcúrrimus, quátenus ejus exposítio ita nesciéntibus fiat cógnita, ut tamen sciéntibus non sit onerósa. Dóminus dicit: Sint lumbi vestri præcíncti. Lumbos enim præcíngimus, cum carnis luxúriam per continéntiam coarctámus. Sed quia minus est, mala non ágere, nisi étiam quisque stúdeat, et bonis opéribus insudáre, prótinus ádditur: Et lucérnæ ardéntes in mánibus vestris. Lucérnas quippe ardéntes in mánibus tenémus, cum per bona ópera próximis nostris lucis exémpla monstrámus. De quibus profécto opéribus Dóminus dicit: Lúceat lux vestra coram homínibus, ut vídeant ópera vestra bona, et gloríficent Patrem vestrum qui in cælis est.\par
\switchcolumn
\EN
\noindent Dearly beloved brethren, the words of the Holy Gospel, which have just been read, lie open before you, and, lest their very plainness should make them seem to some to be hard, we will go through them with such shortness as that neither may they which understand not remain unenlightened, nor they which understand be wearied. The Lord saith Let your loins be girded about. Now, we gird our loins about, when by continency we master the lustful inclination of the flesh. But, forasmuch as it sufficeth not for a man to abstain from evil deeds, if he strive not to join thereto the earnest doing of good works, it is immediately added And your lights burning. Our lights burn when, by good works, we give bright example to our neighbour; concerning which works the Lord saith Let your light so shine before men, that they may see your good works, and glorify your Father Which is in heaven. (Matth. v. 16.)\par
\switchcolumn*

\LA
\begin{Resp}\Rsign Iste est qui ante Deum magnas virtútes operátus est, et de omni corde suo laudávit Dóminum: \Star{} Ipse intercédat pro peccátis ómnium populórum.\end{Resp}
\begin{Resp}\Vsign Ecce homo sine queréla, verus Dei cultor, ábstinens se ab omni ópere malo, et pérmanens in innocéntia sua.\end{Resp}
\begin{Resp}\Rsign Ipse intercédat pro peccátis ómnium populórum.\end{Resp}
\switchcolumn
\EN
\begin{Resp}\Rsign This is he which wrought great wonders before God, and praised the Lord with all his heart. \Star{} May he pray for all people, that their sins may be forgiven unto them.\end{Resp}
\begin{Resp}\Vsign Behold a man without blame, a worshipper of God in truth, keeping himself clean from every evil work, and abiding still in his innocency.\end{Resp}
\begin{Resp}\Rsign May he pray for all people, that their sins may be forgiven unto them.\end{Resp}
\switchcolumn*

\LA
\LessonHead{Lectio VIII}
\switchcolumn
\EN
\LessonHead{Lesson VIII}
\switchcolumn*

\LA
\noindent Duo autem sunt, quæ jubéntur, et lumbos restríngere, et lucérnas tenére: ut et mundítia sit castitátis in córpore, et lumen veritátis in operatióne. Redemptóri étenim nostro unum sine áltero placére nequáquam potest: si aut is qui bona agit, adhuc luxúriæ inquinaménta non déserit: aut is qui castitáte præéminet, necdum se per bona ópera exércet. Nec cástitas ergo magna est sine bono ópere, nec opus bonum est áliquod sine castitáte. Sed et si utrúmque ágitur, restat, ut quisquis ille est, spe ad supérnam pátriam tendat, et nequáquam se a vítiis pro mundi hujus honestáte contíneat.\par
\switchcolumn
\EN
\noindent Here, then, are two commandments, to gird our loins about, and to keep our lights burning the cleanness of purity in our body, and the light of the truth in our works. Whoso hath the one and not the other, pleaseth not thereby our Redeemer; that is, he pleaseth Him not which doth good works, but bridleth not himself from the pollutions of lust, neither he which is eminent in chastity, but exerciseth not himself in good works. Neither is chastity a great thing without good works, nor good works anything without chastity. And if any man do both, it remaineth that he must look by hope toward our Fatherland above, and not have for his reason wherethrough he turneth himself away from vice, the love of honour in this present world.\par
\switchcolumn*

\LA
\begin{Resp}\Rsign Sint lumbi vestri præcíncti, et lucérnæ ardéntes in mánibus vestris: \Star{} Et vos símiles homínibus exspectántibus dóminum suum, quando revertátur a núptiis.\end{Resp}
\begin{Resp}\Vsign Vigiláte ergo, quia nescítis qua hora Dóminus vester ventúrus sit.\end{Resp}
\begin{Resp}\Rsign Et vos símiles homínibus exspectántibus dóminum suum, quando revertátur a núptiis.\end{Resp}
\begin{Resp}\Vsign Glória Patri, et Fílio, \Star{} et Spirítui Sancto.\end{Resp}
\begin{Resp}\Rsign Et vos símiles homínibus exspectántibus dóminum suum, quando revertátur a núptiis.\end{Resp}
\switchcolumn
\EN
\begin{Resp}\Rsign Let your loins be girded about, and your lights burning; \Star{} And ye yourselves like unto men that wait for their lord, when he will return from the wedding.\end{Resp}
\begin{Resp}\Vsign Watch therefore, for ye know not what hour your Lord doth come.\end{Resp}
\begin{Resp}\Rsign And ye yourselves like unto men that wait for their lord, when he will return from the wedding.\end{Resp}
\begin{Resp}\Vsign Glory be to the Father, and to the Son, \Star{} and to the Holy Ghost.\end{Resp}
\begin{Resp}\Rsign And ye yourselves like unto men that wait for their lord, when he will return from the wedding.\end{Resp}
\switchcolumn*

\LA
\LessonHead{Lectio IX}
\switchcolumn
\EN
\LessonHead{Lesson IX}
\switchcolumn*

\LA
\noindent Et vos símiles homínibus exspectántibus dóminum suum, quando revertátur a núptiis: ut cum vénerit et pulsáverit, conféstim apériant ei. Venit quippe Dóminus, cum ad judícium próperat: pulsat vero, cum jam per ægritúdinis moléstias esse mortem vicínam desígnat. Cui conféstim aperímus, si hunc cum amóre suscípimus. Aperíre enim júdici pulsánti non vult, qui exíre de córpore trépidat: et vidére eum, quem contempsísse se méminit, júdicem formídat. Qui autem de sua spe et operatióne secúrus est, pulsanti conféstim áperit, quia lætus júdicem sústinet; et, cum tempus propínquæ mortis advénerit, de glória retributiónis hilaréscit.\par
\switchcolumn
\EN
\noindent And ye yourselves like unto men that wait for their lord, when he will return from the wedding that, when he cometh and knocketh, they may open unto him immediately. The Lord cometh at the hour of judgment He knocketh when, by the pains of sickness, He biddeth us know that death is nigh. To Him open we immediately, if we receive Him in love. Whoso feareth to leave this body, will not open to the Judge when He knocketh, for he dreadeth to see that Judge, Whom he knoweth that he hath despised. But whosoever knoweth that his hope and works are built upon a good foundation, when he heareth the Judge knock, openeth to Him immediately, for to such an one that coming is blessed, yea, when the hour of death is at hand, such an one haileth with gladness a glorious reward.\par
\switchcolumn*

\LA
\TeDeum{Te Deum}
\switchcolumn
\EN
\TeDeum{Te Deum}
\switchcolumn*

\end{paracol}

\Office{Commune Doctoris non Pontificis}{Common of a Doctor not a Bishop}{}
\addcontentsline{toc}{subsection}{Commune Doctoris non Pontificis \textit{— Common of a Doctor not a Bishop}}
\begin{paracol}{2}
\LA
\LessonHead{Lectio VII}
\switchcolumn
\EN
\LessonHead{Lesson VII}
\switchcolumn*

\LA
\LessonTitle{Léctio sancti Evangélii secúndum Matthǽum}
\switchcolumn
\EN
\LessonTitle{From the Holy Gospel according to Matthew}
\switchcolumn*

\LA
\Cite{Matt 5:13-19}
\switchcolumn
\EN
\Cite{Matt 5:13-19}
\switchcolumn*

\LA
\noindent In illo témpore: Dixit Jesus discípulis suis: Vos estis sal terræ. Quod si sal evanúerit, in quo saliétur? Et réliqua.\par
\switchcolumn
\EN
\noindent At that time: Jesus said unto his disciples: Ye are the salt of the earth: But if the salt have lost his savour, wherewith shall it be salted? And so on, and that which followeth.\par
\switchcolumn*

\LA
\LessonTitle{Homilía sancti Augustíni Epíscopi}
\switchcolumn
\EN
\LessonTitle{A Homily by St. Augustine the Bishop}
\switchcolumn*

\LA
\Source{Lib. 1 de Sermone Domini in monte, cap. 6}
\switchcolumn
\EN
\Source{Lib. 1 de Sermóne Dómini in monte, cap. 6}
\switchcolumn*

\LA
\noindent Osténdit Dóminus fátuos esse judicándos, qui temporálium bonórum vel cópiam sectántes vel inópiam metuéntes, amíttunt ætérna, quæ nec dari possunt ab homínibus nec auférri. Itaque, si sal infatuátum fúerit, in quo saliétur? Id est, si vos, per quos condiéndi sunt quodámmodo pópuli, metu persecutiónum temporálium amiséritis regna cælórum; qui erunt hómines, per quos a vobis error auferátur, cum vos elégerit Deus, per quos errórem áuferat ceterórum?\par
\switchcolumn
\EN
\noindent The Lord would have us understand how that men do lose their power of savouring others with righteousness when they are willing to place their eternal welfare in jeopardy for the sake of any témporal advantage, like as attainment of ease or luxury, or escape from suffering or toil. For that which is eternal, unlike things of this world, can neither be bestowed by men, nor by them taken away. Hence, when he asketh: If the salt have lost his savour, wherewith shall it be salted? he would have us understand the question to be: If ye, by whom mankind is preserved from corruption, be willing to lose the kingdom of heaven so as to escape trials or persecutions in this world, who is there to preserve you from corruption, seeing ye are they that God hath chosen to preserve all others from corruption?\par
\switchcolumn*

\LA
\begin{Resp}\Rsign Iste est qui ante Deum magnas virtútes operátus est, et de omni corde suo laudávit Dóminum: \Star{} Ipse intercédat pro peccátis ómnium populórum.\end{Resp}
\begin{Resp}\Vsign Ecce homo sine queréla, verus Dei cultor, ábstinens se ab omni ópere malo, et pérmanens in innocéntia sua.\end{Resp}
\begin{Resp}\Rsign Ipse intercédat pro peccátis ómnium populórum.\end{Resp}
\switchcolumn
\EN
\begin{Resp}\Rsign This is he which wrought great wonders before God, and praised the Lord with all his heart. \Star{} May he pray for all people, that their sins may be forgiven unto them.\end{Resp}
\begin{Resp}\Vsign Behold a man without blame, a worshipper of God in truth, keeping himself clean from every evil work, and abiding still in his innocency.\end{Resp}
\begin{Resp}\Rsign May he pray for all people, that their sins may be forgiven unto them.\end{Resp}
\switchcolumn*

\LA
\LessonHead{Lectio VIII}
\switchcolumn
\EN
\LessonHead{Lesson VIII}
\switchcolumn*

\LA
\noindent Ergo ad níhilum valet sal infatuátum, nisi ut mittátur foras et calcétur ab homínibus. Non ítaque calcátur ab homínibus qui pátitur persecutiónem; sed qui, persecutiónem timéndo, infatuátur. Calcári enim non potest nisi inférior; sed inférior non est, qui, quamvis córpore multa in terra sustíneat, corde tamen fixus in cælo est.\par
\switchcolumn
\EN
\noindent Those that should be the salt of the earth, but have lost their savour, are thenceforth good for nothing, but to be cast out, and to be trodden under foot of men. But no one that suffereth persecution is truly said to be trodden under foot of men. Rather, that one is truly trodden under foot of men who through fear of persecution hath lost the savour of righteousness. For no one can be trodden upon, unless he be beneath him which treadeth upon him. And certainly no one who hath his heart in heaven, no matter how grievously he doth suffer in his body on earth, is rightly said to be beneath anyone who misuseth him.\par
\switchcolumn*

\LA
\begin{Resp}\Rsign In médio Ecclésiæ apéruit os ejus, \Star{} Et implévit eum Dóminus spíritu sapiéntiæ et intelléctus.\end{Resp}
\begin{Resp}\Vsign Jucunditátem et exsultatiónem thesaurizávit super eum.\end{Resp}
\begin{Resp}\Rsign Et implévit eum Dóminus spíritu sapiéntiæ et intelléctus.\end{Resp}
\begin{Resp}\Vsign Glória Patri, et Fílio, \Star{} et Spirítui Sancto.\end{Resp}
\begin{Resp}\Rsign Et implévit eum Dóminus spíritu sapiéntiæ et intelléctus.\end{Resp}
\switchcolumn
\EN
\begin{Resp}\Rsign In the midst of the congregation did the Lord open his mouth. \Star{} And filled him with the spirit of wisdom and understanding.\end{Resp}
\begin{Resp}\Vsign He made him rich with joy and gladness.\end{Resp}
\begin{Resp}\Rsign And filled him with the spirit of wisdom and understanding.\end{Resp}
\begin{Resp}\Vsign Glory be to the Father, and to the Son, \Star{} and to the Holy Ghost.\end{Resp}
\begin{Resp}\Rsign And filled him with the spirit of wisdom and understanding.\end{Resp}
\switchcolumn*

\LA
\LessonHead{Lectio IX}
\switchcolumn
\EN
\LessonHead{Lesson IX}
\switchcolumn*

\LA
\noindent Vos estis lumen mundi. Quómodo dixit supérius sal terræ, sic nunc dicit lumen mundi. Nam, neque supérius ista terra accipiénda est, quam pédibus corpóreis calcámus; sed hómines, qui in terra hábitant, vel étiam peccatóres, quorum condiéndis et exstinguéndis putóribus apostólicum salem Dóminus misit. Et hic mundum non cælum et terram, sed hómines qui sunt in mundo vel díligunt mundum, opórtet intélligi; quibus illuminándis Apóstoli missi sunt. Non potest cívitas abscóndi super montem pósita; id est, fundáta super insígnem magnámque justítiam, quam signíficat étiam ipse mons, in quo dísputat Dóminus.\par
\switchcolumn
\EN
\noindent Ye are the light of the world. And we are to understand the word World in the same sense as the word Earth when he spoke above of the salt of the earth, that is, not that earth whereupon we walk with our bodily feet, but the men which dwell upon the earth; in other words, sinners, for the sweetening and correction of whose corruption, the Lord hath sent his Apostles, as it were, as so much salt. And so by the world we are to understand, not the heaven and the earth, but the men who are in the world and love the world, for the enlightening of whom the Apostles have been sent. A city that is set on a hill cannot be hid: that is, what is founded upon the heights of righteousness, whereof the mountain upon which the Lord gave this discourse was itself a figure, is magnificent in the eyes of all men.\par
\switchcolumn*

\LA
\TeDeum{Te Deum}
\switchcolumn
\EN
\TeDeum{Te Deum}
\switchcolumn*

\end{paracol}

\Office{Commune Virginum tantum}{Common of a Virgin not a Martyr}{}
\addcontentsline{toc}{subsection}{Commune Virginum tantum \textit{— Common of a Virgin not a Martyr}}
\begin{paracol}{2}
\LA
\LessonHead{Lectio VII}
\switchcolumn
\EN
\LessonHead{Lesson VII}
\switchcolumn*

\LA
\LessonTitle{Léctio sancti Evangélii secúndum Matthǽum}
\switchcolumn
\EN
\LessonTitle{From the Holy Gospel according to St. Matthew}
\switchcolumn*

\LA
\Cite{Matt 25:1-13}
\switchcolumn
\EN
\Cite{Matt 25:1}
\switchcolumn*

\LA
\noindent In illo témpore: Dixit Jesus discípulis suis parábolam hanc: Símile erit regnum cælórum decem virgínibus, quæ, accipiéntes lámpades suas, exiérunt óbviam sponso et sponsæ. Et réliqua.\par
\switchcolumn
\EN
\noindent At that time, Jesus said to His disciples: The Kingdom of heaven shall be likened unto ten virgins, which took their lamps, and went forth to meet the Bridegroom and the Bride. And so on.\par
\switchcolumn*

\LA
\LessonTitle{Homilía sancti Gregórii Papæ}
\switchcolumn
\EN
\LessonTitle{Homily by Pope St. Gregory the Great}
\switchcolumn*

\LA
\Source{Homilía 12 in Evangelia}
\switchcolumn
\EN
\Source{12th on the Gospels}
\switchcolumn*

\LA
\noindent Sæpe vos, fratres caríssimi, admóneo prava ópera fúgere, mundi hujus inquinaménta devitáre, sed hodiérna sancti Evangélii lectióne compéllor dícere, ut, et bona, quæ ágitis, cum magna cautéla teneátis; ne per hoc, quod a vobis rectum géritur, favor aut grátia humána requirátur; ne appetítus laudis subrépat, et, quod foris osténditur, intus a mercéde vacuétur. Ecce enim Redemptóris voce decem vírgines, et omnes dicúntur vírgines, et tamen intra beatitúdinis jánuam non omnes sunt recéptæ; quia eárum quædam, dum de virginitáte sua glóriam foris éxpetunt, in vasis suis óleum habére noluérunt.\par
\switchcolumn
\EN
\noindent Dearly beloved brethren oftentimes do I warn you to fly corrupt conversation, and to keep yourselves unspotted from the world. But the portion which is this day read from the Holy Gospel doth oblige me to say that even to these good things which ye do, ye must needs take all careful heed. Look ye well to it, that, when ye work righteousness, ye do it not as seeking the praise and admiration of men, for if the lust of praise do once creep in, that which seemeth so fair without, loseth its reward within. Behold how the Redeemer speaketh of these ten virgins. He calleth them all virgins, yet entered not all of them into the door of blessedness, for there were some of them who sought outwardly the honour of virginity, but would take no oil within their vessels with their lamps.\par
\switchcolumn*

\LA
\begin{Resp}\Rsign Hæc est Virgo sápiens, quam Dóminus vigilántem invénit, quæ accéptis lampádibus sumpsit secum óleum: \Star{} Et veniénte Dómino, introívit cum eo ad núptias.\end{Resp}
\begin{Resp}\Vsign Média nocte clamor factus est: Ecce sponsus venit, exíte óbviam ei.\end{Resp}
\begin{Resp}\Rsign Et veniénte Dómino, introívit cum eo ad núptias.\end{Resp}
\switchcolumn
\EN
\begin{Resp}\Rsign This is one of those wise virgins, whom the Lord found watching, for when she took her lamp, she took oil with her. \Star{} And when the Lord came, she went in with him to the marriage.\end{Resp}
\begin{Resp}\Vsign At midnight there was a cry made: Behold, the Bridegroom cometh, go ye out to meet him.\end{Resp}
\begin{Resp}\Rsign And when the Lord came, she went in with him to the marriage.\end{Resp}
\switchcolumn*

\LA
\LessonHead{Lectio VIII}
\switchcolumn
\EN
\LessonHead{Lesson VIII}
\switchcolumn*

\LA
\noindent Sed prius quæréndum nobis est quid sit regnum cælórum, aut cur decem virgínibus comparétur, quæ étiam vírgines prudéntes et fátuæ dicántur. Dum enim cælórum regnum constat quia reprobórum nullus ingréditur, étiam fátuis virgínibus cur símile esse perhibétur? Sed sciéndum nobis est quod sæpe in sacro elóquio regnum cælórum præséntis témporis Ecclésia dícitur. De quo álio in loco Dóminus dicit: Mittet Fílius hóminis Ángelos suos, et cólligent de regno ejus ómnia scándala. Neque enim in illo regno beatitúdinis, in quo pax summa est, inveníri scándala póterunt, quæ colligántur.\par
\switchcolumn
\EN
\noindent First of all, it is for us to ask What is the kingdom of Heaven? And wherefore shall the same be likened unto ten virgins, whereof, albeit five were wise, yet five were foolish For if the kingdom of heaven be such that there shall in no wise enter into it anything that defileth, neither whatsoever worketh abomination, or maketh a lie, (Apoc. xxi. 27,) how can it be like unto five virgins which were foolish? But we must know that, in the word of God, the kingdom of heaven doth oftentimes signify the Church as she now is, touching the which the Lord saith in another place “The Son of Man shall send forth His Angels, and they shall gather out of His kingdom all things that offend.” (Matth. xiii. 41.) In that kingdom of Blessedness, wherein peace shall have her perfect reign, there shall be nothing found that offendeth for the angels to gather out.\par
\switchcolumn*

\LA
\begin{Resp}\Rsign Média nocte clamor factus est: \Star{} Ecce sponsus venit, exíte óbviam ei.\end{Resp}
\begin{Resp}\Vsign Prudéntes vírgines, aptáte vestras lámpades.\end{Resp}
\begin{Resp}\Rsign Ecce sponsus venit, exíte óbviam ei.\end{Resp}
\begin{Resp}\Vsign Glória Patri, et Fílio, \Star{} et Spirítui Sancto.\end{Resp}
\begin{Resp}\Rsign Ecce sponsus venit, exíte óbviam ei.\end{Resp}
\switchcolumn
\EN
\begin{Resp}\Rsign At midnight there was a cry made: \Star{} Behold, the Bridegroom cometh, go ye out to meet him.\end{Resp}
\begin{Resp}\Vsign Trim your lamps, O ye wise virgins.\end{Resp}
\begin{Resp}\Rsign Behold, the Bridegroom cometh, go ye out to meet him.\end{Resp}
\begin{Resp}\Vsign Glory be to the Father, and to the Son, \Star{} and to the Holy Ghost.\end{Resp}
\begin{Resp}\Rsign Behold, the Bridegroom cometh, go ye out to meet him.\end{Resp}
\switchcolumn*

\LA
\LessonHead{Lectio IX}
\switchcolumn
\EN
\LessonHead{Lesson IX}
\switchcolumn*

\LA
\noindent In quinque autem córporis sénsibus unusquísque subsístit; geminátus autem quinárius denárium pérficit. Et, quia ex utróque sexu fidélium multitúdo collígitur, sancta Ecclésia decem virgínibus símilis esse denuntiátur. In qua quia mali cum bonis et réprobi cum eléctis admíxti sunt, recte símilis virgínibus prudéntibus et fátuis esse perhibétur. Sunt namque pleríque continéntes, qui ab appetítu se exterióri custódiunt et spe ad interióra rapiúntur, carnem mácerant, et toto desidério ad supérnam pátriam anhélant, ætérna prǽmia éxpetunt, pro labóribus suis recípere laudes humánas nolunt. Hi nimírum glóriam suam non in ore hóminum ponunt, sed intra consciéntiam cóntegunt. Et sunt pleríque, qui corpus per abstinéntiam afflígunt, sed de ipsa sua abstinéntia humános favóres éxpetunt.\par
\switchcolumn
\EN
\noindent The body of every man doth consist of five senses, and five being doubled, is ten. Forasmuch, therefore, as the whole body of the faithful doth consist of two sexes, the Holy Church is likened unto ten virgins. And forasmuch as in the Church the good are for the present mingled with the bad, and the reprobate with the elect, it is rightly said that, of the ten virgins, five are wise and five are foolish. There are many who have self-control, which do keep themselves from lusting after things outward, whose hope beareth them to things inward, who chastise the flesh, who long with intense home-sickness for their Fatherland which is in heaven, who seek an eternal reward, and who will not to receive for their labours the praise of men. These are they who reckon their glory, not in the mouths of men, but in the testimony of their own conscience. And many there be likewise who afflict the body by self-control, and yet who seek for their self-control applause from men.\par
\switchcolumn*

\LA
\TeDeum{Te Deum}
\switchcolumn
\EN
\TeDeum{Te Deum}
\switchcolumn*

\end{paracol}

\Office{Commune non Virginum non Martyrum}{Common of Widows}{}
\addcontentsline{toc}{subsection}{Commune non Virginum non Martyrum \textit{— Common of Widows}}
\begin{paracol}{2}
\LA
\LessonHead{Lectio I}
\switchcolumn
\EN
\LessonHead{Lesson I}
\switchcolumn*

\LA
\LessonTitle{De Parábolis Salomónis}
\switchcolumn
\EN
\LessonTitle{Lesson from the book of Proverbs}
\switchcolumn*

\LA
\Cite{Prov 31:10-17}
\switchcolumn
\EN
\Cite{Prov 31:10-17}
\switchcolumn*

\LA
\noindent Mulíerem fortem quis invéniet? procul et de últimis fínibus prétium ejus. Confídit in ea cor viri sui, et spóliis non indigébit. Reddet ei bonum, et non malum, ómnibus diébus vitæ suæ. Quæsívit lanam et linum, et operáta est consília mánuum suárum. Facta est quasi navis institóris, de longe portans panem suum. Et de nocte surréxit, dedítque prædam domésticis suis, et cibária ancíllis suis. Considerávit agrum, et emit eum; de fructu mánuum suárum plantávit víneam. Accínxit fortitúdine lumbos suos, et roborávit bráchium suum.\par
\switchcolumn
\EN
\noindent Who shall find a valiant woman? Far and from the uttermost coasts is the price of her. The heart of her husband trusteth in her, and he shall have no need of spoils. She will render him good, and not evil, all the days of her life. She hath sought wool and flax, and hath wrought by the counsel of her hands. She is like the merchant's ship, she bringeth her bread from afar. And she hath risen in the night, and given a prey to her household, and victuals to her maidens. She hath considered a field, and bought it: with the fruit of her hands she hath planted a vineyard. She hath girded her loins with strength, and hath strengthened her arm\par
\switchcolumn*

\LA
\begin{Resp}\Rsign Veni, elécta mea, et ponam in te thronum meum \Star{} Quia concupívit Rex spéciem tuam.\end{Resp}
\begin{Resp}\Vsign Spécie tua et pulchritúdine tua inténde, próspere procéde, et regna.\end{Resp}
\begin{Resp}\Rsign Quia concupívit Rex spéciem tuam.\end{Resp}
\switchcolumn
\EN
\begin{Resp}\Rsign Come, O My chosen one, and I will establish My throne in thee; \Star{} For the King hath greatly desired thy beauty.\end{Resp}
\begin{Resp}\Vsign In thy comeliness and thy beauty, go forward, fare prosperously, and reign.\end{Resp}
\begin{Resp}\Rsign For the King hath greatly desired thy beauty.\end{Resp}
\switchcolumn*

\LA
\LessonHead{Lectio II}
\switchcolumn
\EN
\LessonHead{Lesson II}
\switchcolumn*

\LA
\Cite{Prov 31:18-24}
\switchcolumn
\EN
\Cite{Prov 31:18-24}
\switchcolumn*

\LA
\noindent Gustávit, et vidit quia bona est negotiátio ejus; non extinguétur in nocte lucérna ejus. Manum suam misit ad fórtia, et dígiti ejus apprehendérunt fusum. Manum suam apéruit ínopi, et palmas suas exténdit ad páuperem. Non timébit dómui suæ a frigóribus nivis; omnes enim doméstici ejus vestíti sunt duplícibus. Stragulátam vestem fecit sibi; byssus et púrpura induméntum ejus. Nóbilis in portis vir ejus, quando séderit cum senatóribus terræ. Síndonem fecit, et véndidit, et cíngulum trádidit Chananǽo.\par
\switchcolumn
\EN
\noindent She hath tasted and seen that her traffic is good: her lamp shall not be put out in the night. She hath put out her hand to strong things, and her fingers have taken hold of the spindle. She hath opened her hand to the needy, and stretched out her hands to the poor. She shall not fear for her house in the cold of snow: for all her domestics are clothed with double garments. She hath made for herself clothing of tapestry: fine linen, and purple is her covering. Her husband is honourable in the gates, when he sitteth among the senators of the land. She made fine linen, and sold it, and delivered a girdle to the Chanaanite.\par
\switchcolumn*

\LA
\begin{Resp}\Rsign Diffúsa est grátia in lábiis tuis, \Star{} Proptérea benedíxit te Deus in ætérnum.\end{Resp}
\begin{Resp}\Vsign Spécie tua et pulchritúdine tua inténde, próspere procéde, et regna.\end{Resp}
\begin{Resp}\Rsign Proptérea benedíxit te Deus in ætérnum.\end{Resp}
\switchcolumn
\EN
\begin{Resp}\Rsign Grace is poured into thy lips, therefore; \Star{} God hath blessed thee for ever.\end{Resp}
\begin{Resp}\Vsign In thy comeliness and thy beauty, go forward, fare prosperously, and reign.\end{Resp}
\begin{Resp}\Rsign God hath blessed thee for ever.\end{Resp}
\switchcolumn*

\LA
\LessonHead{Lectio III}
\switchcolumn
\EN
\LessonHead{Lesson III}
\switchcolumn*

\LA
\Cite{Prov 31:25-31}
\switchcolumn
\EN
\Cite{Prov 31:25-31}
\switchcolumn*

\LA
\noindent Fortitúdo et decor induméntum ejus, et ridébit in die novíssimo. Os suum apéruit sapiéntiæ, et lex cleméntiæ in lingua ejus. Considerávit sémitas domus suæ, et panem otiósa non comédit. Surrexérunt fílii ejus, et beatíssimam prædicavérunt; vir ejus, et laudávit eam. Multæ fíliæ congregavérunt divítias; tu supergréssa es univérsas. Fallax grátia, et vana est pulchritúdo: múlier timens Dóminum, ipsa laudábitur. Date ei de fructu mánuum suárum, et laudent eam in portis ópera ejus.\par
\switchcolumn
\EN
\noindent Strength and beauty are her clothing, and she shall laugh in the latter day. She hath opened her mouth to wisdom, and the law of clemency is on her tongue. She hath looked well to the paths of her house, and hath not eaten her bread idle. Her children rose up, and called her blessed: her husband, and he praised her. Many daughters have gathered together riches: thou hast surpassed them all. Favour is deceitful, and beauty is vain: the woman that feareth the Lord, she shall be praised. Give her of the fruit of her hands: and let her works praise her in the gates.\par
\switchcolumn*

\LA
\begin{Resp}\Rsign Spécie tua et pulchritúdine tua \Star{} Inténde, próspere procéde, et regna.\end{Resp}
\begin{Resp}\Vsign Diffúsa est grátia in lábiis tuis, proptérea benedíxit te Deus in ætérnum.\end{Resp}
\begin{Resp}\Rsign Inténde, próspere procéde, et regna.\end{Resp}
\begin{Resp}\Vsign Glória Patri, et Fílio, \Star{} et Spirítui Sancto.\end{Resp}
\begin{Resp}\Rsign Inténde, próspere procéde, et regna.\end{Resp}
\switchcolumn
\EN
\begin{Resp}\Rsign In thy comeliness and thy beauty; \Star{} Go forward, fare prosperously, and reign.\end{Resp}
\begin{Resp}\Vsign Grace is poured into thy lips, therefore God hath blessed thee for ever.\end{Resp}
\begin{Resp}\Rsign Go forward, fare prosperously, and reign.\end{Resp}
\begin{Resp}\Vsign Glory be to the Father, and to the Son, \Star{} and to the Holy Ghost.\end{Resp}
\begin{Resp}\Rsign Go forward, fare prosperously, and reign.\end{Resp}
\switchcolumn*

\LA
\LessonHead{Lectio VII}
\switchcolumn
\EN
\LessonHead{Lesson VII}
\switchcolumn*

\LA
\LessonTitle{Léctio sancti Evangélii secúndum Matthǽum}
\switchcolumn
\EN
\LessonTitle{From the Holy Gospel according to Matthew}
\switchcolumn*

\LA
\Cite{Matt 13:44-52}
\switchcolumn
\EN
\Cite{Matt 13:44}
\switchcolumn*

\LA
\noindent In illo témpore: Dixit Jesus discípulis suis parábolam hanc: Simile est regnum cælórum thesáuro abscóndito in agro. Et réliqua.\par
\switchcolumn
\EN
\noindent At that time Jesus spake unto His disciples this parable The kingdom of heaven is like unto treasure hid in a field. And so on.\par
\switchcolumn*

\LA
\LessonTitle{Homilía sancti Gregórii Papæ}
\switchcolumn
\EN
\LessonTitle{Homily by Pope St Gregory the Great}
\switchcolumn*

\LA
\Source{Homilia 11 in Evangelia}
\switchcolumn
\EN
\Source{11th on the Gospels}
\switchcolumn*

\LA
\noindent Cælórum regnum, fratres caríssimi, idcírco terrénis rebus símile dícitur, ut, ex his quæ ánimus novit, surgat ad incógnita, quæ non novit: quátenus exémplo visibílium se ad invisibília rápiat, et, per ea quæ usu dídicit quasi confricátus incaléscat; ut per hoc, quod scit notum dilígere, discat et incógnita amáre. Ecce enim cælórum regnum thesáuro abscóndito in agro comparátur; quem, qui invénit homo, abscóndit, et præ gáudio illíus vadit, et vendit univérsa quæ habet, et emit agrum illum.\par
\switchcolumn
\EN
\noindent Dearly beloved brethren, the kingdom of heaven is likened unto the things of earth, to the end that by the mean of things which we know, our mind may rise to the contemplation of the things which we know not by the example of things which are seen, may fix her gaze on things which are not seen by the touch of things which she useth, may be warmed towards the things which she useth not; by things which she knoweth and loveth, to love also the things which she knoweth not. For, behold, “the kingdom of heaven is likened unto treasure hid in a field, the which when a man hath found, he hideth, and, for joy thereof, goeth and selleth all that he hath and buyeth that field.”\par
\switchcolumn*

\LA
\begin{Resp}\Rsign Os suum apéruit sapiéntiæ, et lex cleméntiæ in lingua ejus: considerávit sémitas domus suæ, \Star{} Et panem otiósa non comédit.\end{Resp}
\begin{Resp}\Vsign Gustávit et vidit quia bona est negotiátio ejus: non exstinguétur in nocte lucérna ejus.\end{Resp}
\begin{Resp}\Rsign Et panem otiósa non comédit.\end{Resp}
\switchcolumn
\EN
\begin{Resp}\Rsign She openeth her mouth with wisdom, and in her tongue is the law of kindness. She looketh well to the ways of her household \Star{} And eateth not the bread of idleness.\end{Resp}
\begin{Resp}\Vsign She tasteth and perceiveth that her merchandise is good. Her candle goeth not out by night.\end{Resp}
\begin{Resp}\Rsign And she eateth not the bread of idleness.\end{Resp}
\switchcolumn*

\LA
\LessonHead{Lectio VIII}
\switchcolumn
\EN
\LessonHead{Lesson VIII}
\switchcolumn*

\LA
\noindent Qua in re hoc quoque notándum est, quod invéntus thesáurus abscónditur, ut servétur: quia stúdium cæléstis desidérii a malígnis spirítibus custodíre non súfficit, qui hoc ab humánis láudibus non abscóndit. In præsénti étenim vita, quasi in via sumus, qua ad pátriam pérgimus. Malígni autem spíritus iter nostrum quasi quidam latrúnculi óbsident. Deprædári ergo desíderat, qui thesáurum públice portat in via. Hoc autem dico, non ut próximi ópera nostra bona non vídeant, cum scriptum sit: Vídeant ópera vestra bona, et gloríficent Patrem vestrum qui in cælis est; sed, ut per hoc quod ágimus, laudes extérius non quærámus. Sic autem sit opus in público, quátenus inténtio máneat in occúlto; ut, et de bono ópere próximis præbeámus exémplum, et tamen per intentiónem, qua Deo soli placére quǽrimus, semper optémus secrétum.\par
\switchcolumn
\EN
\noindent And herein we must remark that the treasure, when once it hath been found, is hidden to keep it safe. He whose intimate yearnings after God are not hidden from the praise of men, is open thereby to the attacks of evil spirits. In this life we are, as it were, journeying homewards on a road beset by evil spirits who are like highwaymen. He therefore inviteth robbery who carrieth his treasure ostentatiously. Doubtless our neighbour should be able to see our good works, as it is written: “Let your light so shine before men that they may see your good works, and glorify your Father which is in heaven.” But this is not to be understood to mean that we are to seek the praise of men by what we do. Rather, let us in such wise work in the open that the inner intention of devotion is not advertised. So we shall give an example to our neighbour, and yet keep hidden, except from the sight of God, our purpose of pleasing him.\par
\switchcolumn*

\LA
\begin{Resp}\Rsign Regnum mundi et omnem ornátum sǽculi contémpsi, propter amórem Dómini mei Jesu Christi: \Star{} Quem vidi, quem amávi, in quem crédidi, quem diléxi.\end{Resp}
\begin{Resp}\Vsign Eructávit cor meum verbum bonum: dico ego ópera mea Regi.\end{Resp}
\begin{Resp}\Rsign Quem vidi, quem amávi, in quem crédidi, quem diléxi.\end{Resp}
\begin{Resp}\Vsign Glória Patri, et Fílio, \Star{} et Spirítui Sancto.\end{Resp}
\begin{Resp}\Rsign Quem vidi, quem amávi, in quem crédidi, quem diléxi.\end{Resp}
\switchcolumn
\EN
\begin{Resp}\Rsign The kingdom of this world and all the beauty of life I have esteemed as nothing, for the excellency of the love of Jesus Christ my Lord \Star{} Whom, having seen, I loved Whom, having believed, I longed after.\end{Resp}
\begin{Resp}\Vsign My heart is overflowing with a good matter I speak of my works unto the King.\end{Resp}
\begin{Resp}\Rsign Whom, having seen, I loved Whom, having believed, I longed after.\end{Resp}
\begin{Resp}\Vsign Glory be to the Father, and to the Son, \Star{} and to the Holy Ghost.\end{Resp}
\begin{Resp}\Rsign Whom, having seen, I loved Whom, having believed, I longed after.\end{Resp}
\switchcolumn*

\LA
\LessonHead{Lectio IX}
\switchcolumn
\EN
\LessonHead{Lesson IX}
\switchcolumn*

\LA
\noindent Thesáurus autem cæléste est desidérium; ager vero, in quo thesáurus abscónditur, disciplína stúdii cæléstis. Quem profécto agrum, vénditis ómnibus, cómparat, qui, voluptátibus carnis renúntians, cuncta sua terréna desidéria per disciplínæ cæléstis custódiam calcat: ut nihil jam quod caro blandítur, líbeat; nihil, quod carnálem vitam trucídat, spíritus perhorréscat.\par
\switchcolumn
\EN
\noindent The treasure is the desire for heaven, the field wherein it is hidden is the earnest observance wherewith this desire is surrounded. Whosoever turneth his back upon the enjoyments of the flesh, and by earnest striving heavenward, putteth all earthly lusts under the feet of discipline, so that he smileth back no more when the flesh smileth at him, and shuddereth no more at anything that can only kill the body whosoever doth thus, hath sold all that he had, and bought that field.\par
\switchcolumn*

\LA
\TeDeum{Te Deum}
\switchcolumn
\EN
\TeDeum{Te Deum}
\switchcolumn*

\end{paracol}

\Office{Sanctæ Mariæ Sabbato (a Trinitate usque ad Adventum)}{Our Lady's Saturday (from Trinity to Advent)}{Simplex}
\addcontentsline{toc}{subsection}{Sanctæ Mariæ Sabbato (a Trinitate usque ad Adventum) \textit{— Our Lady's Saturday (from Trinity to Advent)}}
\begin{paracol}{2}
\LA
\RefLine{Lectiones I–II: de Scriptura occurrente.}
\switchcolumn
\EN
\RefLine{Lessons I–II: from the Scripture of the day.}
\switchcolumn*

\LA
\LessonHead{Lectio III (mense Junio)}
\switchcolumn
\EN
\LessonHead{Lesson III (in June)}
\switchcolumn*

\LA
\LessonTitle{Sermo sancti Bernárdi Abbátis}
\switchcolumn
\EN
\LessonTitle{Sermon of St. Bernard, Abbot}
\switchcolumn*

\LA
\Source{De Verbis Apocalypsis cap. 12, Signum magnum}
\switchcolumn
\EN
\Source{On the words of the Apocalypse, ch. 12, Signum magnum}
\switchcolumn*

\LA
\noindent Veheménter quidem nobis, dilectíssimi, vir unus et múlier una nocuére; sed, grátias Deo, per unum nihilóminus virum et mulíerum unam ómnia restaurántur, nec sine magno fǽnore gratiárum. Neque enim sicut delíctum, ita et donum; sed excédit damni æstimatiónem benefícii magnitúdo. Sic nimírum prudentíssimus et clementíssimus ártifex, quod quassátum fúerat, non confrégit, sed utílius omníno refécit, ut vidélicet nobis novum formáret Adam ex véteri, Hevam transfúnderet in Maríam.\par
\switchcolumn
\EN
\noindent It was indeed a serious injury that one man and one woman inflicted on us, dearly beloved; but thanks to God, it was also by one Man and one woman that all things were restored, and with a great increase of grace too. For “not like the offense is gift”, on the contrary the benefits received are greater than the loss sustained. Yes, that was how the Maker supreme in good judgment and in kindness plied His craft: what had been bruised, He did not break. Rather, He remade it completely in such a way as to be in more advantage to us: out of the old Adam He made a new Man; Eve He transformed into Mary.\par
\switchcolumn*

\LA
\TeDeum{Te Deum}
\switchcolumn
\EN
\TeDeum{Te Deum}
\switchcolumn*

\LA
\LessonHead{Lectio III (mense Julio)}
\switchcolumn
\EN
\LessonHead{Lesson III (in July)}
\switchcolumn*

\LA
\LessonTitle{Ex Epístola sancti Ambrósii Epíscopi ad Sirícium Papam}
\switchcolumn
\EN
\LessonTitle{From the letter of Blessed Ambrose, Bishop to Pope Siricius}
\switchcolumn*

\LA
\Source{Epist. 81, alias 7, circa medium}
\switchcolumn
\EN
\Source{Epist. 81, alias 7 circa medium}
\switchcolumn*

\LA
\noindent Non excédit fidem, quod homo exívit de vírgine, quando petra fontem prófluum scaturívit, ferrum super aquas natávit, ambulávit homo super aquas. Ergo si hóminem unda portávit, non pótuit hóminem virgo generáre, atque hóminem, de quo légimus: Et mittet illis Dóminus hóminem, qui salvos fáciet eos et notus erit Dóminus Ægýptiis? In véteri ítaque Testaménto virgo Hebræórum per mare duxit exércitum: in novo Testaménto Virgo, géneris aula cæléstis, elécta est ad salútem.\par
\switchcolumn
\EN
\noindent That a man was born from a virgin is not beyond belief, since a fountain of water gushed from a rock, and iron floated on water, and a man walked on the sea. Therefore if the waves bore a man, cannot a virgin give birth to a man, and to that man of whom we read: The Lord shall send them a Saviour, a man, and he shall deliver them, and the Lord shall be known to Egypt? In the Old Testament, a virgin led the Hebrew hosts through the Red Sea: in the New Testament, a Virgin was chosen as the palace of the heavenly Birth, to bring us salvation.\par
\switchcolumn*

\LA
\TeDeum{Te Deum}
\switchcolumn
\EN
\TeDeum{Te Deum}
\switchcolumn*

\LA
\LessonHead{Lectio III (mense Augusto)}
\switchcolumn
\EN
\LessonHead{Lesson III (in August)}
\switchcolumn*

\LA
\LessonTitle{De Expositióne sancti Gregórii Papæ in libros Regum}
\switchcolumn
\EN
\LessonTitle{From the Commentary of Pope St. Gregory on the books of Kings}
\switchcolumn*

\LA
\Source{In 1 Reg. 1}
\switchcolumn
\EN
\Source{On I Kings I}
\switchcolumn*

\LA
\noindent Fuit vir unus de Ramáthaim Sophim, de monte Ephraim. Potest, hujus montis nómine, beatíssima semper Virgo María, Dei Génetrix, designári. Mons quippe fuit, quæ omnem eléctæ creatúræ altitúdinem, electiónis suæ dignitáte, transcéndit. An non mons sublímis María, quæ, ut ad conceptiónem ætérni Verbi pertíngeret, meritórum vérticem, supra omnes Angelórum choros, usque ad sólium Deitátis eréxit? Hujus enim montis præcellentíssimam dignitátem Isaías vatícinans, ait: Erit in novíssimis diébus præparátus mons domus Dómini in vértice móntium. Mons quippe in vértice móntium fuit, quia altitúdo Maríæ supra omnes Sanctos refúlsit.\par
\switchcolumn
\EN
\noindent “There was a man of Ramathaim-Sophim, of mount Ephraim.” The name of this mountain can be taken to designate the Most Blessed Virgin Mary, Mother of God. She is a mountain in as much as, among chosen creatures, the dignity of her calling was above that of any others. Surely Mary was the highest of mountains, was she not, when she heaped up merits beyond all the choirs of Angels, to the very throne of God, that she might be worthy to conceive the eternal Word? It was this mountain whose supreme dignity Isaias was prophesying when he said: “In days to come, the mountain of the Lord's house shall be established as the highest mountain.” Mary was indeed the highest mountain, since her summit gleamed above those of all the Saints.\par
\switchcolumn*

\LA
\TeDeum{Te Deum}
\switchcolumn
\EN
\TeDeum{Te Deum}
\switchcolumn*

\LA
\LessonHead{Lectio III (mense Septembri)}
\switchcolumn
\EN
\LessonHead{Lesson III (in September)}
\switchcolumn*

\LA
\LessonTitle{Ex Epístola sancti Leónis Papæ ad Pulchériam Augústam}
\switchcolumn
\EN
\LessonTitle{From the epistle of Pope St. Leo to Pulcheria Augusta}
\switchcolumn*

\LA
\Source{Epist. 13, ante medium}
\switchcolumn
\EN
\Source{Epistle 13 in the first half}
\switchcolumn*

\LA
\noindent Sacraméntum reconcilatiónis nostræ, ante témpora ætérna dispósitum, nullæ implébant figúræ; quia nondum supervénerat Spíritus Sanctus in Vírginem nec virtus Altíssimi obumbráverat ei, ut, et intra intemeráta víscera, ædificánte sibi Sapiéntia domum, Verbum caro fíeret, et, forma Dei ac forma servi in unam conveniénte persónam, Creátor témporum nascerétur in témpore, et, per quem facta sunt ómnia, ipse inter ómnia gignerétur. Nisi enim novus homo, factus in similitúdinem carnis peccáti, nostram suscíperet vetustátem, et, consubstantiális Patri, consubstantiális esse dignarétur et matri, naturámque sibi nostram solus a peccáto liber uníret; sub jugo diáboli generáliter tenerétur humána captívas.\par
\switchcolumn
\EN
\noindent And the fulfillment of the mystery of our atonement, which was ordained from all eternity, was not assisted by any figures because the Holy Spirit had not yet come upon the Virgin, and the power of the Most High had not over-shadowed her: so that Wisdom building herself a house within her undefiled body, the Word became flesh; and the form of God and the form of a slave coming together into one person, the Creator of times was born in time; and He Himself through whom all things were made, was brought forth in the midst of all things. For if the New Man had not been made in the likeness of sinful flesh, and taken on Him our old nature, and being consubstantial with the Father, had deigned to be consubstantial with His mother also, and being alone free from sin, had united our nature to Him the whole human race would be held in bondage beneath the Devil's yoke.\par
\switchcolumn*

\LA
\TeDeum{Te Deum}
\switchcolumn
\EN
\TeDeum{Te Deum}
\switchcolumn*

\end{paracol}

\Office{In Festis Beatae Mariae Virginis}{In Festis Beatae Mariae Virginis}{}
\addcontentsline{toc}{subsection}{In Festis Beatae Mariae Virginis \textit{— In Festis Beatae Mariae Virginis}}
\begin{paracol}{2}
\LA
\LessonHead{Lectio I}
\switchcolumn
\EN
\LessonHead{Lesson I}
\switchcolumn*

\LA
\LessonTitle{De Parábolis Salomónis}
\switchcolumn
\EN
\LessonTitle{Lesson from the book of Proverbs}
\switchcolumn*

\LA
\Cite{Prov 8:12-17}
\switchcolumn
\EN
\Cite{Prov 8:12-17}
\switchcolumn*

\LA
\noindent Ego sapiéntia hábito in consílio et erudítis intérsum cogitatiónibus. Timor Dómini odit malum: arrogántiam, et supérbiam, et viam pravam, et os bilíngue detéstor. Meum est consílium et ǽquitas, mea est prudéntia, mea est fortitúdo. Per me reges regnant, et legum conditóres justa decérnunt; per me príncipes ímperant, et poténtes decérnunt justítiam. Ego diligéntes me díligo; et qui mane vígilant ad me, invénient me.\par
\switchcolumn
\EN
\noindent I wisdom dwell in counsel, and am present in learned thoughts. The fear of the Lord hateth evil: I hate arrogance, and pride, and every wicked way, and a mouth with a double tongue. Counsel and equity is mine, prudence is mine, strength is mine. By me kings reign, and lawgivers decree just things, By me princes rule, and the mighty decree justice. I love them that love me: and they that in the morning early watch for me, shall find me.\par
\switchcolumn*

\LA
\begin{Resp}\Rsign Sancta et immaculáta virgínitas, quibus te láudibus éfferam, néscio: \Star{} Quia quem cæli cápere non póterant, tuo grémio contulísti.\end{Resp}
\begin{Resp}\Vsign Benedícta tu in muliéribus, et benedíctus fructus ventris tui.\end{Resp}
\begin{Resp}\Rsign Quia quem cæli cápere non póterant, tuo grémio contulisti.\end{Resp}
\switchcolumn
\EN
\begin{Resp}\Rsign O how holy and how spotless is thy virginity; I am too dull to praise thee: \Star{} For thou hast borne in thy breast Him Whom the heavens cannot contain.\end{Resp}
\begin{Resp}\Vsign Blessed art thou among women, and blessed is the fruit of thy womb.\end{Resp}
\begin{Resp}\Rsign For thou hast borne in thy breast Him Whom the heavens cannot contain.\end{Resp}
\switchcolumn*

\LA
\LessonHead{Lectio II}
\switchcolumn
\EN
\LessonHead{Lesson II}
\switchcolumn*

\LA
\Cite{Prov 8:18-25}
\switchcolumn
\EN
\Cite{Prov 8:18-25}
\switchcolumn*

\LA
\noindent Mecum sunt divítiæ et glória, opes supérbæ et justítia. Mélior est enim fructus meus auro et lápide pretióso, et genímina mea argénto elécto. In viis justítiæ ámbulo, in médio semitárum judícii, ut ditem diligéntes me et thesáuros eórum répleam. Dóminus possidébit me in inítio viárum suárum, ántequam quidquam fáceret a princípio. Ab ætérno ordináta sum et ex antíquis, ántequam terra fíeret. Nondum erant abýssi, et ego jam concépta eram; necdum fontes aquárum erúperant, necdum montes gravi mole constíterant; ante colles ego parturiébar.\par
\switchcolumn
\EN
\noindent With me are riches and glory, glorious riches and justice. For my fruit is better than gold and the precious stone, and my blossoms than choice silver. I walk in the way of justice, in the midst of the paths of judgment, That I may enrich them that love me, and may fill their treasures. The Lord possessed me in the beginning of his ways, before he made any thing from the beginning. I was set up from eternity, and of old before the earth was made. The depths were not as yet, and I was already conceived. neither had the fountains of waters as yet sprung out: The mountains with their huge bulk had not as yet been established: before the hills I was brought forth:\par
\switchcolumn*

\LA
\begin{Resp}\Rsign Congratulámini mihi, omnes qui dilígitis Dóminum: quia cum essem párvula, plácui Altíssimo, \Star{} Et de meis viscéribus génui Deum et hóminem.\end{Resp}
\begin{Resp}\Vsign Beátam me dicent omnes generatiónes, quia ancíllam húmilem respéxit Deus.\end{Resp}
\begin{Resp}\Rsign Et de meis viscéribus génui Deum et hóminem.\end{Resp}
\switchcolumn
\EN
\begin{Resp}\Rsign Rejoice with me, all ye that love the Lord, for while I was yet a little one, I pleased the Most High. \Star{} And I have brought forth from my bowels God and man.\end{Resp}
\begin{Resp}\Vsign All generations shall call me blessed, since the Lord hath regarded the lowliness of his handmaiden.\end{Resp}
\begin{Resp}\Rsign And I have brought forth from my bowels God and man.\end{Resp}
\switchcolumn*

\LA
\LessonHead{Lectio III}
\switchcolumn
\EN
\LessonHead{Lesson III}
\switchcolumn*

\LA
\Cite{Prov 8:34-36; 9:1-5}
\switchcolumn
\EN
\Cite{Prov 8:34-36; 9:1-5}
\switchcolumn*

\LA
\noindent Beátus homo qui audit me, et qui vígilat ad fores meas cotídie, et obsérvat ad postes óstii mei. Qui me invénerit, invéniet vitam, et háuriet salútem a Dómino; qui autem in me peccáverit, lædet ánimam suam. Omnes, qui me odérunt, díligunt mortem. Sapiéntia ædificávit sibi domum, excídit colúmnas septem. Immolávit víctimas suas, míscuit vinum et propósuit mensam suam. Misit ancíllas suas ut vocárent ad arcem et ad mœ́nia civitátis: Si quis est párvulus, véniat ad me. Et insipiéntibus locúta est: Veníte, comédite panem meum, et bíbite vinum quod míscui vobis.\par
\switchcolumn
\EN
\noindent Blessed is the man that heareth me, and that watcheth daily at my gates, and waiteth at the posts of my doors. He that shall find me, shall find life, and shall have salvation from the Lord: But he that shall sin against me, shall hurt his own soul. All that hate me love death. Wisdom hath built herself a house, she hath hewn out her seven pillars. She hath slain her victims, mingled her wine, and set forth her table. She hath sent her maids to invite to the tower, and to the walls of the city: Whosoever is a little one, let him come to me. And to the unwise she said: Come, eat my bread, and drink the wine which I have mingled for you\par
\switchcolumn*

\LA
\begin{Resp}\Rsign Beáta es, Virgo María, quæ Dóminum portásti, Creatórem mundi: \Star{} Genuísti qui te fecit, et in ætérnum pérmanes Virgo.\end{Resp}
\begin{Resp}\Vsign Ave, María, grátia plena; Dóminus tecum.\end{Resp}
\begin{Resp}\Rsign Genuísti qui te fecit, et in ætérnum pérmanes Virgo.\end{Resp}
\begin{Resp}\Vsign Glória Patri, et Fílio, \Star{} et Spirítui Sancto.\end{Resp}
\begin{Resp}\Rsign Genuísti qui te fecit, et in ætérnum pérmanes Virgo.\end{Resp}
\switchcolumn
\EN
\begin{Resp}\Rsign Blessed art thou, O Virgin Mary, who hast carried the Lord, the Maker of the world. \Star{} Thou hast borne Him Who created thee, and thou abidest a virgin for ever.\end{Resp}
\begin{Resp}\Vsign Hail, Mary, full of grace. the Lord is with thee.\end{Resp}
\begin{Resp}\Rsign Thou hast borne Him Who created thee, and thou abidest a virgin for ever.\end{Resp}
\begin{Resp}\Vsign Glory be to the Father, and to the Son, \Star{} and to the Holy Ghost.\end{Resp}
\begin{Resp}\Rsign Thou hast borne Him Who created thee, and thou abidest a virgin for ever.\end{Resp}
\switchcolumn*

\LA
\LessonHead{Lectio VII}
\switchcolumn
\EN
\LessonHead{Lesson VII}
\switchcolumn*

\LA
\LessonTitle{Léctio sancti Evangélii secúndum Lucam}
\switchcolumn
\EN
\LessonTitle{From the Holy Gospel according to Luke}
\switchcolumn*

\LA
\Cite{Luc 11:27-28}
\switchcolumn
\EN
\Cite{Luke 11:27-28}
\switchcolumn*

\LA
\noindent In illo témpore: Loquénte Jesu ad turbas, extóllens vocem quædam múlier de turba dixit illi: Beátus venter qui te portávit. Et réliqua.\par
\switchcolumn
\EN
\noindent At that time, as Jesus spoke unto the multitudes, a certain woman of the company lifted up her voice and said unto Him: Blessed is the womb that bore thee. And so on.\par
\switchcolumn*

\LA
\LessonTitle{Homilía sancti Bedæ Venerábilis Presbýteri}
\switchcolumn
\EN
\LessonTitle{Homily by the Venerable Bede}
\switchcolumn*

\LA
\Source{Lib 4. cap. 49. in Luc. 11.}
\switchcolumn
\EN
\Source{Lib 4 Cap 49 in Luke 11}
\switchcolumn*

\LA
\noindent Magnæ devotiónis et fídei hæc múlier osténditur, quæ, scribis et pharisǽis Dóminum tentántibus simul et blasphemántibus, tanta ejus incarnatiónem præ ómnibus sinceritáte cognóscit, tanta fidúcia confitétur, ut et præséntium prócerum calúmniam, et futurórum confúndat hæreticórum perfídiam. Nam, sicut tunc Judǽi, Sancti Spíritus ópera blasphemándo, verum consubstantialémque Patri Dei Fílium negábant; sic hærétici póstea, negándo Maríam semper Vírginem, Sancti Spíritus operánte virtúte, nascitúro cum humánis membris Unigénito Dei, carnis suæ matériam ministrásse, verum consubstantialémque matri Fílium hóminis fatéri non debére dixérunt.\par
\switchcolumn
\EN
\noindent It is plain that this was a woman of great earnestness and faith. The Scribes and Pharisees were at once tempting and blaspheming the Lord, but this woman so clearly grasped His Incarnation, and so bravely confessed the same, that she confounded both the lies of the great men who were present, and the faithlessness of the heretics who were yet to come. Even as the Jews then, blaspheming the works of the Holy Ghost, denied the very Son of God Who is of one substance with the Father, so afterwards did the heretics, by denying that Mary always a Virgin did, under the operation of the Holy Ghost, supply flesh to the Only begotten One of God, when He was about being born in an human Body, even so, I say, did the heretics deny that the Son of Man should be called a true Son, Who is of one substance with His Mother.\par
\switchcolumn*

\LA
\begin{Resp}\Rsign Felix namque es, sacra Virgo María, et omni laude digníssima: \Star{} Quia ex te ortus est sol justítiæ, Christus, Deus noster.\end{Resp}
\begin{Resp}\Vsign Ora pro pópulo, intérveni pro clero, intercéde pro devóto femíneo sexu: séntiant omnes tuum juvámen, quicúmque célebrant tuam sanctam festivitátem.\end{Resp}
\begin{Resp}\Rsign Quia ex te ortus est sol justítiæ, Christus, Deus noster.\end{Resp}
\switchcolumn
\EN
\begin{Resp}\Rsign O Holy Virgin Mary, happy indeed art thou, and right worthy of all praise \Star{} For out of thee rose the Sun of righteousness, even Christ our God.\end{Resp}
\begin{Resp}\Vsign Pray for the people, plead for the clergy, make intercession for all women vowed to God. May all that are keeping this thine holy Feast day feel the might of thine assistance.\end{Resp}
\begin{Resp}\Rsign For out of thee rose the Sun of righteousness, even Christ our God.\end{Resp}
\switchcolumn*

\LA
\LessonHead{Lectio VIII}
\switchcolumn
\EN
\LessonHead{Lesson VIII}
\switchcolumn*

\LA
\noindent Sed, si caro Verbi Dei secúndum carnem nascéntis a carne Vírginis matris pronuntiátur extránea, sine causa venter qui eam portásset, úbera quæ lactássent, beatificántur. Dicit autem Apóstolus: Quia misit Deus Fílium suum, factum ex mulíere, factum sub lege. Neque audiéndi sunt, qui legéndum putant: Natum ex mulíere, factum sub lege, sed Factum ex mulíere; quia concéptus ex útero virgináli, carnem non de níhilo, non aliúnde, sed matérna traxit ex carne. Alióquin nec vere Fílius hóminis dicerétur, qui oríginem non habéret ex hómine. Et nos ígitur, his contra Eutychen dictis, extollámus vocem cum Ecclésia cathólica, cujus hæc múlier typum gessit, extollámus et mentem de médio turbárum, dicamúsque Salvatóri: Beátus venter qui te portávit, et úbera quæ suxísti. Vere enim beáta parens, quæ, sicut quidam ait, Eníxa est puérpera Regem, qui cælum terrámque tenet per sǽcula.\par
\switchcolumn
\EN
\noindent If we shall say that the Flesh, Wherewith the Son of God was born in the flesh, was something outside of the flesh of the Virgin His Mother, without reason should we bless the womb that bore Him, and the paps which He hath sucked. But the Apostle saith: “God sent forth His Son, made of a woman, made under the law”, (Gal. iv. 4,) and they are not to be listened to who read this passage: “Born of a woman, made under the law.” He was made of a woman, for He was conceived in a virgin's womb, and took His Flesh, not from nothing, not from elsewhere, but from the flesh of His Mother. Otherwise, and if He had not been sprung of a woman, He could not with truth be called the Son of man. Let us therefore, denying the doctrine of Eutyches, lift up our voice, along with the Universal Church, whereof that woman was a figure, let us lift up our heart as well as our voice from the company, and say unto the Saviour: Blessed is the womb that bare thee, and the paps which thou hast sucked. Blessed Mother, of whom one hath said thou art His Mother, Who reigns over earth and over heaven for ever.\par
\switchcolumn*

\LA
\begin{Resp}\Rsign Beátam me dicent omnes generatiónes, \Star{} Quia fecit mihi Dóminus magna qui potens est, et sanctum nomen ejus.\end{Resp}
\begin{Resp}\Vsign Et misericórdia ejus a progénie in progénies timéntibus eum.\end{Resp}
\begin{Resp}\Rsign Quia fecit mihi Dóminus magna qui potens est, et sanctum nomen ejus.\end{Resp}
\begin{Resp}\Vsign Glória Patri, et Fílio, \Star{} et Spirítui Sancto.\end{Resp}
\begin{Resp}\Rsign Quia fecit mihi Dóminus magna qui potens est, et sanctum nomen ejus.\end{Resp}
\switchcolumn
\EN
\begin{Resp}\Rsign All generations shall call me blessed. \Star{} For He that is mighty, even the Lord, hath done to me great things; and Holy is His Name.\end{Resp}
\begin{Resp}\Vsign And his mercy is on them that fear Him, from generation to generation.\end{Resp}
\begin{Resp}\Rsign He that is mighty, even the Lord, hath done to me great things; and Holy is His Name.\end{Resp}
\begin{Resp}\Vsign Glory be to the Father, and to the Son, \Star{} and to the Holy Ghost.\end{Resp}
\begin{Resp}\Rsign He that is mighty, even the Lord, hath done to me great things; and Holy is His Name.\end{Resp}
\switchcolumn*

\LA
\LessonHead{Lectio IX}
\switchcolumn
\EN
\LessonHead{Lesson IX}
\switchcolumn*

\LA
\noindent Quinímmo beáti qui áudiunt verbum Dei et custódiunt. Pulchre Salvátor attestatióni mulíeris ánnuit, non eam tantúmmodo quæ Verbum Dei corporáliter generáre merúerat, sed et omnes qui idem Verbum spiritáliter audítu fídei concípere, et boni óperis custódia vel in suo vel in proximórum corde párere et quasi álere studúerint, asséverans esse beátos; quia, et éadem Dei Génetrix, et inde quidem beáta quia Verbi incarnándi minístra facta est temporális, sed inde multo beátior quia ejúsdem semper amándi custos manébat ætérna.\par
\switchcolumn
\EN
\noindent Yea, rather, blessed are they that hear the Word of God and keep it. How nobly doth the Saviour say “Yea” to the woman's blessing, declaring also that not only is she blessed who was meet to give bodily birth to the Word of God, but that all they who spiritually conceive the same Word by the hearing of faith, and, by keeping it through good works, bring it forth and, as it were, carefully nurse it, in their own hearts, and in the hearts of their neighbours, are also blessed. Yea, and that the very Mother of God herself was blessed in being for a while the handmaid of the Word of God made Flesh, but that she was much more blessed in this, that through her love she keepeth Him for ever.\par
\switchcolumn*

\LA
\TeDeum{Te Deum}
\switchcolumn
\EN
\TeDeum{Te Deum}
\switchcolumn*

\end{paracol}
\backmatter
\IndexOfOffices
\end{document}
